
\import{naprochelibrary.tex}

\section{§12 Topological Spaces}

% from §12 Item 1: topology on a set
\begin{definition}\label{topology_on}
    $T$ is a topology on $X$ iff
    $T\subseteq\pow{X}$ and
    $\emptyset\in T$ and
    $X\in T$ and
    $T$ is closed under arbitrary unions and
    $T$ is closed under binary intersections.
\end{definition}

% from §12 Item 2: topological space with topology
\begin{definition}\label{topological_space}
    $X$ is a topological space with topology $T$ iff $T$ is a topology on $X$.
\end{definition}

% from §12 Item 3: open set
\begin{definition}\label{open_in}
    $U$ is open in $X$ with respect to $T$ iff $T$ is a topology on $X$ and $U\in T$.
\end{definition}

% from §12 Item 4: discrete topology
\begin{definition}\label{discrete_topology}
    $\discrete{X} = \pow{X}$.
\end{definition}

% from §12 Item 5: indiscrete topology
\begin{definition}\label{indiscrete_topology}
    $\indiscrete{X} = \{\emptyset, X\}$.
\end{definition}

% from §12 Item 6: finite complement topology
\begin{definition}\label{finite_complement_topology}
    $\finiteComplement{X} = \{U\in\pow{X} \mid \text{$X\setminus U$ is finite or $X\setminus U = X$}\}$.
\end{definition}

% from §12 Item 7: countable complement topology
\begin{definition}\label{countable_complement_topology}
    $\countableComplement{X} = \{U\in\pow{X} \mid \text{$X\setminus U$ is countable or $X\setminus U = X$}\}$.
\end{definition}

% from §12 Item 8: finer topology
\begin{definition}\label{finer_topology}
    $T_1$ is finer than $T$ on $X$ iff
    $T$ is a topology on $X$ and
    $T_1$ is a topology on $X$ and
    $T\subseteq T_1$.
\end{definition}

% from §12 Item 9: strictly finer topology
\begin{definition}\label{strictly_finer_topology}
    $T_1$ is strictly finer than $T$ on $X$ iff $T_1$ is finer than $T$ on $X$ and $T\neq T_1$.
\end{definition}

% from §12 Item 10: coarser topology
\begin{definition}\label{coarser_topology}
    $T$ is coarser than $T_1$ on $X$ iff $T_1$ is finer than $T$ on $X$.
\end{definition}

% from §12 Item 11: strictly coarser topology
\begin{definition}\label{strictly_coarser_topology}
    $T$ is strictly coarser than $T_1$ on $X$ iff $T$ is coarser than $T_1$ on $X$ and $T\neq T_1$.
\end{definition}

% from §12 Item 12: comparable topologies
\begin{definition}\label{comparable_topologies}
    $T$ is comparable with $T_1$ on $X$ iff
    $T$ is a topology on $X$ and
    $T_1$ is a topology on $X$ and
    either $T\subseteq T_1$ or $T_1\subseteq T$.
\end{definition}

\section{§13 Basis for a Topology}

% from §13 Item 1: basis on a set
\begin{definition}\label{basis_on}
    $B$ is a basis on $X$ iff
    $B\subseteq\pow{X}$ and
    for all $x\in X$ there exists $B_1\in B$ such that $x\in B_1$ and
    for all $B_1,B_2,x$ such that $B_1\in B$ and $B_2\in B$ and $x\in B_1\inter B_2$
    there exists $B_3\in B$ such that $x\in B_3$ and $B_3\subseteq B_1\inter B_2$.
\end{definition}

% from §13 Item 2: topology generated by a basis
\begin{definition}\label{topology_generated_by_basis}
    $\basisTopology{B}{X} = \{U\in\pow{X} \mid \text{for all $x\in U$ there exists $B_1\in B$ such that $x\in B_1$ and $B_1\subseteq U$}\}$.
\end{definition}

% from §13 Item 3: basis for a topology
\begin{definition}\label{basis_for_topology}
    $B$ is a basis for $T$ on $X$ iff $B$ is a basis on $X$ and $T = \basisTopology{B}{X}$.
\end{definition}

% from §13 Item 4: topology generated by a basis is a topology
\begin{lemma}\label{basis_topology_is_topology}
    Assume $B$ is a basis on $X$.
    Then $\basisTopology{B}{X}$ is a topology on $X$.
\end{lemma}
\begin{proof}
    $\basisTopology{B}{X}\subseteq\pow{X}$.
    For all $x\in\emptyset$ there exists $B_1\in B$ such that $x\in B_1$ and $B_1\subseteq\emptyset$.
    Hence $\emptyset\in\basisTopology{B}{X}$ by \cref{powerset_bottom,topology_generated_by_basis}.
    $X\in\basisTopology{B}{X}$ by \cref{basis_on,subseteq,pow_iff,subseteq_refl,topology_generated_by_basis}.
    For all $M$ such that $M\subseteq\basisTopology{B}{X}$ we have $\unions{M}\in\basisTopology{B}{X}$
        by \cref{powerset_closed_unions,unions_iff,subseteq,topology_generated_by_basis}.
    Show that for all $U,V$ such that $U\in\basisTopology{B}{X}$ and $V\in\basisTopology{B}{X}$ we have $U\inter V\in\basisTopology{B}{X}$.
    \begin{subproof}
        Fix $U$.
        Fix $V$.
        Assume $U\in\basisTopology{B}{X}$ and $V\in\basisTopology{B}{X}$.
        Show that for all $x$ such that $x\in U\inter V$ there exists $B_3\in B$ such that $x\in B_3$ and $B_3\subseteq U\inter V$.
        \begin{subproof}
            Fix $x$ such that $x\in U\inter V$.
            Then $x\in X$ by \cref{subseteq,pow_iff,inter_elim_left}.
            There exists $B_1\in B$ such that $x\in B_1$ and $B_1\subseteq U$ by \cref{inter_elim_left,topology_generated_by_basis}.
            There exists $B_2\in B$ such that $x\in B_2$ and $B_2\subseteq V$ by \cref{inter_elim_right,topology_generated_by_basis}.
            Then there exists $B_3\in B$ such that $x\in B_3$ and $B_3\subseteq B_1\inter B_2$ by \cref{basis_on,inter_intro}.
            Then $B_3\subseteq U\inter V$.
        \end{subproof}
        Therefore $U\inter V\in\basisTopology{B}{X}$ by \cref{subseteq,pow_iff,inter_subseteq,topology_generated_by_basis}.
    \end{subproof}
    Therefore $\basisTopology{B}{X}$ is a topology on $X$ by \cref{topology_on}.
\end{proof}

% from §13 helper lemma 1: basis elements are open in the generated topology
\begin{lemma}\label{basis_elem_open_in_generated}
    Assume $B$ is a basis on $X$.
    Suppose $B_1\in B$.
    Then $B_1\in\basisTopology{B}{X}$.
\end{lemma}
\begin{proof}
    Follows by \cref{subseteq_refl,basis_on,subseteq,pow_iff,topology_generated_by_basis}.
\end{proof}

% from §13 Item 5: unions of basis elements
\begin{lemma}\label{basis_topology_unions}
    Assume $B$ is a basis for $T$ on $X$.
    Then $T = \{U\in\pow{X} \mid \exists M\subseteq B. U = \unions{M}\}$.
\end{lemma}
\begin{proof}
    $T = \basisTopology{B}{X}$ by \cref{basis_for_topology}.
    It suffices to show that for all $U$ we have
    $U\in\basisTopology{B}{X}$ iff there exists $M\subseteq B$ such that $U = \unions{M}$.
    Fix $U$.
        Show that if $U\in\basisTopology{B}{X}$ then there exists $M\subseteq B$ such that $U = \unions{M}$.
        \begin{subproof}
            Assume $U\in\basisTopology{B}{X}$.
            Let $M = \{B_1\in B \mid B_1\subseteq U\}$.
            Hence $M\subseteq B$ and $U = \unions{M}$
                by \cref{topology_generated_by_basis,unions_iff,subseteq,subseteq_antisymmetric}.
        \end{subproof}
        If there exists $M\subseteq B$ such that $U = \unions{M}$ then $U\in\basisTopology{B}{X}$
            by \cref{basis_topology_is_topology,basis_for_topology,basis_elem_open_in_generated,subseteq,topology_on}.
\end{proof}

% from §13 Item 6: criterion for a collection of open sets to be a basis
\begin{lemma}\label{open_collection_basis}
    Assume $T$ is a topology on $X$.
    Assume for all $C_1\in C$ we have $C_1$ is open in $X$ with respect to $T$.
    Assume for all $U,x$ such that $U\in T$ and $x\in U$ there exists $C_1\in C$ such that $x\in C_1$ and $C_1\subseteq U$.
    Then $C$ is a basis for $T$ on $X$.
\end{lemma}
\begin{proof}
    For all $x$ such that $x\in X$ there exists $C_1\in C$ such that $x\in C_1$ and $C_1\subseteq X$ by \cref{topology_on}.
    For all $C_1,C_2,x$ such that $C_1\in C$ and $C_2\in C$ and $x\in C_1\inter C_2$
        there exists $C_3\in C$ such that $x\in C_3$ and $C_3\subseteq C_1\inter C_2$
        by \cref{open_in,topology_on}.
    Hence $C$ is a basis on $X$ by \cref{open_in,topology_on,subseteq,basis_on}.
    $T\subseteq\basisTopology{C}{X}$ by \cref{topology_on,subseteq,topology_generated_by_basis}.
    Therefore $\basisTopology{C}{X}\subseteq T$
        by \cref{basis_topology_unions,basis_for_topology,basis_topology_is_topology,open_in,subseteq,topology_on}.
    Hence $C$ is a basis for $T$ on $X$ by \cref{subseteq,subseteq_antisymmetric,basis_for_topology}.
\end{proof}

% from §13 Item 7: bases and finer topologies
\begin{lemma}\label{basis_finer_criterion}
    Assume $B$ is a basis for $T$ on $X$.
    Assume $B_1$ is a basis for $T_1$ on $X$.
    Then $T_1$ is finer than $T$ on $X$ iff
    for all $x,B_2$ such that $x\in X$ and $B_2\in B$ and $x\in B_2$
    there exists $B_3\in B_1$ such that $x\in B_3$ and $B_3\subseteq B_2$.
\end{lemma}
\begin{proof}
    If for all $x,B_2$ such that $x\in X$ and $B_2\in B$ and $x\in B_2$
        there exists $B_3\in B_1$ such that $x\in B_3$ and $B_3\subseteq B_2$
        then $T_1$ is finer than $T$ on $X$
        by \cref{basis_for_topology,topology_generated_by_basis,basis_on,subseteq,powerset_elim,subseteq_transitive,basis_topology_is_topology,finer_topology}.
    If $T_1$ is finer than $T$ on $X$ then
        for all $x,B_2$ such that $x\in X$ and $B_2\in B$ and $x\in B_2$
        there exists $B_3\in B_1$ such that $x\in B_3$ and $B_3\subseteq B_2$
        by \cref{basis_elem_open_in_generated,basis_for_topology,finer_topology,subseteq,topology_generated_by_basis}.
\end{proof}

% from §13 Item 8: half-open interval on R
\begin{definition}\label{interval_closed_open}
    $\intervalclosedopen{a}{b} = \{x\in\reals \mid a \leq x \land x < b\}$.
\end{definition}

% from §13 Item 9: standard basis on R
\begin{definition}\label{standard_basis_reals}
    $\standardBasisR = \{U\in\pow{\reals} \mid \exists a,b\in\reals. a < b \land U = \intervalopen{a}{b}\}$.
\end{definition}

% from §13 Item 10: standard topology on R
\begin{definition}\label{standard_topology_reals}
    $\standardTopologyR = \basisTopology{\standardBasisR}{\reals}$.
\end{definition}

% from §13 Item 11: lower limit basis on R
\begin{definition}\label{lower_limit_basis_reals}
    $\lowerLimitBasisR = \{U\in\pow{\reals} \mid \exists a,b\in\reals. a < b \land U = \intervalclosedopen{a}{b}\}$.
\end{definition}

% from §13 Item 12: lower limit topology on R
\begin{definition}\label{lower_limit_topology_reals}
    $\lowerLimitTopologyR = \basisTopology{\lowerLimitBasisR}{\reals}$.
\end{definition}

% from §13 Item 14: K set
\begin{definition}\label{k_set}
    $\kset = \{x\in\reals \mid \exists n\in\naturalsPlus. x = 1 / n\}$.
\end{definition}

% from §13 Item 15: K basis on R
\begin{definition}\label{k_basis_reals}
    $\kBasisR = \{U\in\pow{\reals} \mid \exists a,b\in\reals. a < b \land (U = \intervalopen{a}{b} \lor U = \intervalopen{a}{b}\setminus \kset)\}$.
\end{definition}

% from §13 Item 16: K-topology on R
\begin{definition}\label{k_topology_reals}
    $\kTopologyR = \basisTopology{\kBasisR}{\reals}$.
\end{definition}

% from §13 helper lemma 2: split a weak inequality
\begin{lemma}\label{real_leq_split}
    Suppose $a \leq b$.
    Then $a < b$ or $a = b$.
\end{lemma}
\begin{proof}
    Trivial.
\end{proof}

% from §13 helper lemma 3: substitute equality on the left
\begin{lemma}\label{real_lt_subst_left}
    Suppose $a = b$ and $b < c$.
    Then $a < c$.
\end{lemma}
\begin{proof}
    Trivial.
\end{proof}

% from §13 helper lemma 4: substitute equality on the right
\begin{lemma}\label{real_lt_subst_right}
    Suppose $a < b$ and $b = c$.
    Then $a < c$.
\end{lemma}
\begin{proof}
    Trivial.
\end{proof}

% from §13 helper lemma 5: negative zero
\begin{lemma}\label{real_neg_zero}
    $\neg{0} = 0$.
\end{lemma}
\begin{proof}
    Follows by \cref{real_add_ident,real_add_inverse_unique}.
\end{proof}

% from §13 helper lemma 6: refine open intervals
\begin{lemma}\label{open_interval_refinement}
    Suppose $a,b,c,d,x\in\reals$.
    Suppose $a < x$ and $x < b$ and $c < x$ and $x < d$.
    Then there exist $e,f\in\reals$ such that $e < x$ and $x < f$ and
    $\intervalopen{e}{f} \subseteq \intervalopen{a}{b}\inter \intervalopen{c}{d}$.
\end{lemma}
    \begin{proof}
        Let $A = \{a,c\}$.
        $A\subseteq\reals$ by \cref{upair_elim,subseteq}.
        Take $m\in A$ such that $m$ is a maximum of $A$
            by \cref{pair_is_finite,upair_intro_left,emptyset,finite_real_maximum}.
        Let $B = \{b,d\}$.
        $B\subseteq\reals$ by \cref{upair_elim,subseteq}.
        Take $n\in B$ such that $n$ is a minimum of $B$
            by \cref{pair_is_finite,upair_intro_left,emptyset,finite_real_minimum}.
        Then $m < x$ and $x < n$ by \cref{upair_elim}.
        Let $e = m$.
        Let $f = n$.
        Hence $e\in\reals$ and $f\in\reals$ and $e < x$ and $x < f$ and
            $\intervalopen{e}{f}\subseteq \intervalopen{a}{b}\inter \intervalopen{c}{d}$
            by \cref{intervalopen,real_maximum,upair_intro_left,upair_intro_right,real_minimum,real_leq_split,real_lt_subst_left,real_lt_subst_right,real_lt_trans,inter_intro,subseteq}.
\end{proof}

% from §13 helper lemma 7: refine half-open intervals
\begin{lemma}\label{closedopen_interval_refinement}
    Suppose $a,b,c,d,x\in\reals$.
    Suppose $a \leq x$ and $x < b$ and $c \leq x$ and $x < d$.
    Then there exist $e,f\in\reals$ such that $e \leq x$ and $x < f$ and
    $\intervalclosedopen{e}{f} \subseteq \intervalclosedopen{a}{b}\inter \intervalclosedopen{c}{d}$.
\end{lemma}
\begin{proof}
    Let $A = \{a,c\}$.
    $A\subseteq\reals$ by \cref{upair_elim,subseteq}.
    Take $m\in A$ such that $m$ is a maximum of $A$
        by \cref{pair_is_finite,upair_intro_left,emptyset,finite_real_maximum}.
    Let $B = \{b,d\}$.
    $B\subseteq\reals$ by \cref{upair_elim,subseteq}.
    Take $n\in B$ such that $n$ is a minimum of $B$
        by \cref{pair_is_finite,upair_intro_left,emptyset,finite_real_minimum}.
    Then $m \leq x$ and $x < n$ by \cref{upair_elim}.
    Let $e = m$.
    Let $f = n$.
    Hence $e\in\reals$ and $f\in\reals$ and $e \leq x$ and $x < f$ and
        $\intervalclosedopen{e}{f}\subseteq \intervalclosedopen{a}{b}\inter \intervalclosedopen{c}{d}$
        by \cref{interval_closed_open,real_maximum,upair_intro_left,upair_intro_right,real_minimum,real_leq_trans,real_leq_split,real_lt_subst_right,real_lt_trans,inter_intro,subseteq}.
\end{proof}

% from §13 helper lemma 8: open intervals are subsets of R
\begin{lemma}\label{intervalopen_subset_reals}
    Suppose $a\in\reals$ and $b\in\reals$.
    Then $\intervalopen{a}{b}\subseteq\reals$.
\end{lemma}
\begin{proof}
    Follows by \cref{intervalopen,subseteq}.
\end{proof}

% from §13 helper lemma 9: half-open intervals are subsets of R
\begin{lemma}\label{intervalclosedopen_subset_reals}
    Suppose $a\in\reals$ and $b\in\reals$.
    Then $\intervalclosedopen{a}{b}\subseteq\reals$.
\end{lemma}
\begin{proof}
    Follows by \cref{interval_closed_open,subseteq}.
\end{proof}

% from §13 helper lemma 10: K-set elements are positive
\begin{lemma}\label{k_set_positive}
    Suppose $x\in\kset$.
    Then $0 < x$.
\end{lemma}
\begin{proof}
    Take $n\in\naturalsPlus$ such that $x = 1 / n$ by \cref{k_set}.
    $1\in\reals$ and $n\in\reals$ by \cref{real_naturals_subset_reals,real_one_in_naturals,subseteq}.
    Hence $0 < n$ by \cref{real_naturals_subset_reals,real_one_in_naturals,subseteq,real_one_leq_naturalsplus,real_leq_split,real_zero_lt_one,real_add_ident,real_lt_subst_right,real_lt_trans}.
    Hence $0 < x$ by \cref{real_inv_pos,real_pos_nonzero,real_mul_inverse,real_div,real_mul_ident,real_lt_subst_right}.
\end{proof}

% from §13 helper lemma 11: rational between a negative real and zero
\begin{lemma}\label{rational_between_negative_and_zero}
    Suppose $a\in\reals$ and $a < 0$.
    Then there exists $p\in\rationals$ such that $p\in\intervalopen{a}{0}$.
\end{lemma}
\begin{proof}
    $0\in\reals$ and $\neg{a}\in\reals$ by \cref{real_add_ident,real_add_inverse}.
    Hence $0 < \neg{a}$ by \cref{real_lt_neg,real_neg_zero,real_lt_subst_left}.
    Take $N\in\naturalsPlus$ such that $0 < 1 / N$ and $1 / N < \neg{a}$ by \cref{real_small_reciprocal}.
    Hence $0 < N$ by \cref{real_naturals_subset_reals,subseteq,real_mul_ident,real_add_ident,real_zero_lt_one,real_one_leq_naturalsplus,real_lt_subst_right,real_lt_trans}.
    Therefore $1 / N\in\reals$ by \cref{real_naturals_subset_reals,subseteq,real_mul_ident,real_add_ident,real_zero_lt_one,real_one_leq_naturalsplus,real_lt_subst_right,real_lt_trans,real_pos_nonzero,real_mul_inverse,real_div,real_mul_closed,real_one_in_naturals}.
    Let $p = \neg{(1 / N)}$.
    Hence $p\in\intervalopen{a}{0}$ by \cref{real_add_inverse,real_lt_neg,real_neg_zero,real_lt_subst_right,real_lt_subst_left,real_neg_neg,intervalopen}.
    $1\in\naturalsPlus$ and $1\in\reals$ by \cref{real_one_in_naturals,real_mul_ident}.
    Then $\neg{1}\in\integers$ and $N\in\integers$ by \cref{real_add_inverse,real_add_comm,real_integers}.
    Hence $N\in\integers\setminus\{0\}$ by \cref{real_pos_nonzero,singleton_iff,setminus_intro}.
    $\inv{N}\in\reals$ by \cref{real_naturals_subset_reals,subseteq,real_pos_nonzero,real_mul_inverse}.
    Then $p = \neg{1} \rmul \inv{N}$ by \cref{real_add_inverse,real_div,real_mul_neg_right,real_mul_comm}.
    Hence $p = (\neg{1}) / N$ by \cref{real_div}.
    Therefore $p\in\rationals$ by \cref{real_rationals}.
\end{proof}

% from §13 Item 17: standard basis is a basis on R
\begin{lemma}\label{standard_basis_is_basis}
    $\standardBasisR$ is a basis on $\reals$.
\end{lemma}
\begin{proof}
    $\standardBasisR\subseteq\pow{\reals}$ by \cref{standard_basis_reals,subseteq}.
    For all $x$ such that $x\in\reals$ we have $x\in\intervalopen{x + \neg{1}}{x + 1}$ and
        $\intervalopen{x + \neg{1}}{x + 1}\in\standardBasisR$
        by \cref{real_naturals_subset_reals,real_one_in_naturals,subseteq,real_add_inverse,real_add_closed,real_zero_lt_one,real_add_ident,real_lt_neg,real_neg_zero,real_lt_subst_right,real_lt_add,real_add_comm,real_lt_trans,intervalopen,intervalopen_subset_reals,powerset_intro,standard_basis_reals}.
    For all $B_1,B_2,x$ such that $B_1\in\standardBasisR$ and $B_2\in\standardBasisR$ and $x\in B_1\inter B_2$
        there exists $B_3\in\standardBasisR$ such that $x\in B_3$ and $B_3\subseteq B_1\inter B_2$
        by \cref{standard_basis_reals,inter_elim_left,inter_elim_right,intervalopen,open_interval_refinement,intervalopen_subset_reals,powerset_intro,real_lt_trans}.
    Hence $\standardBasisR$ is a basis on $\reals$ by \cref{basis_on}.
\end{proof}

% from §13 Item 18: lower limit basis is a basis on R
\begin{lemma}\label{lower_limit_basis_is_basis}
    $\lowerLimitBasisR$ is a basis on $\reals$.
\end{lemma}
\begin{proof}
    $\lowerLimitBasisR\subseteq\pow{\reals}$ by \cref{lower_limit_basis_reals,subseteq}.
    For all $x$ such that $x\in\reals$ we have $x\in\intervalclosedopen{x}{x + 1}$ and
        $\intervalclosedopen{x}{x + 1}\in\lowerLimitBasisR$
        by \cref{real_naturals_subset_reals,real_one_in_naturals,subseteq,real_add_closed,real_zero_lt_one,real_add_ident,real_lt_add,real_add_comm,real_lt_subst_left,interval_closed_open,intervalclosedopen_subset_reals,powerset_intro,lower_limit_basis_reals}.
    For all $B_1,B_2,x$ such that $B_1\in\lowerLimitBasisR$ and $B_2\in\lowerLimitBasisR$ and $x\in B_1\inter B_2$
        there exists $B_3\in\lowerLimitBasisR$ such that $x\in B_3$ and $B_3\subseteq B_1\inter B_2$
        by \cref{lower_limit_basis_reals,inter_elim_left,inter_elim_right,interval_closed_open,closedopen_interval_refinement,intervalclosedopen_subset_reals,powerset_intro,real_leq_split,real_lt_subst_left,real_lt_trans}.
    Hence $\lowerLimitBasisR$ is a basis on $\reals$ by \cref{basis_on}.
\end{proof}

% from §13 Item 19: K basis is a basis on R
\begin{lemma}\label{k_basis_is_basis}
    $\kBasisR$ is a basis on $\reals$.
\end{lemma}
\begin{proof}
    $\kBasisR\subseteq\pow{\reals}$ by \cref{k_basis_reals,subseteq}.
    For all $x$ such that $x\in\reals$ we have $x\in\intervalopen{x + \neg{1}}{x + 1}$ and
        $\intervalopen{x + \neg{1}}{x + 1}\in\kBasisR$
        by \cref{real_naturals_subset_reals,real_one_in_naturals,subseteq,real_add_inverse,real_add_closed,real_zero_lt_one,real_add_ident,real_lt_neg,real_neg_zero,real_lt_subst_right,real_lt_add,real_add_comm,real_lt_trans,intervalopen,intervalopen_subset_reals,powerset_intro,k_basis_reals}.
    Show that for all $B_1,B_2,x$ such that $B_1\in\kBasisR$ and $B_2\in\kBasisR$ and $x\in B_1\inter B_2$
        there exists $B_3\in\kBasisR$ such that $x\in B_3$ and $B_3\subseteq B_1\inter B_2$.
    \begin{subproof}
        Fix $B_1,B_2$.
        Fix $x$ such that $B_1\in\kBasisR$ and $B_2\in\kBasisR$ and $x\in B_1\inter B_2$.
        Take $a,b\in\reals$ such that $a < b$ and
            $(B_1 = \intervalopen{a}{b} \lor B_1 = \intervalopen{a}{b}\setminus \kset)$
            by \cref{k_basis_reals}.
        Take $c,d\in\reals$ such that $c < d$ and
            $(B_2 = \intervalopen{c}{d} \lor B_2 = \intervalopen{c}{d}\setminus \kset)$
            by \cref{k_basis_reals}.
        \begin{byCase}
        \caseOf{$B_1 = \intervalopen{a}{b}$ and $B_2 = \intervalopen{c}{d}$.}
            Then $x\in\reals$ and $a < x$ and $x < b$ and $c < x$ and $x < d$
                by \cref{inter_elim_left,inter_elim_right,intervalopen}.
            Take $e,f\in\reals$ such that $e < x$ and $x < f$ and
                $\intervalopen{e}{f} \subseteq \intervalopen{a}{b}\inter \intervalopen{c}{d}$
                by \cref{open_interval_refinement}.
            Hence there exists $B_3\in\kBasisR$ such that $x\in B_3$ and $B_3\subseteq B_1\inter B_2$
                by \cref{intervalopen_subset_reals,powerset_intro,k_basis_reals,real_lt_trans,intervalopen}.
        \caseOf{$B_1 = \intervalopen{a}{b}\setminus \kset$ and $B_2 = \intervalopen{c}{d}$.}
            Then $x\in\reals$ and $a < x$ and $x < b$ and $c < x$ and $x < d$ and $x\notin\kset$
                by \cref{inter_elim_left,inter_elim_right,setminus_elim_left,setminus_elim_right,intervalopen}.
            Take $e,f\in\reals$ such that $e < x$ and $x < f$ and
                $\intervalopen{e}{f} \subseteq \intervalopen{a}{b}\inter \intervalopen{c}{d}$
                by \cref{open_interval_refinement}.
            Hence there exists $B_3\in\kBasisR$ such that $x\in B_3$ and $B_3\subseteq B_1\inter B_2$
                by \cref{setminus_elim_left,intervalopen,subseteq,powerset_intro,k_basis_reals,real_lt_trans,setminus_intro,setminus_elim_right,inter_elim_left,inter_elim_right,inter_intro}.
        \caseOf{$B_1 = \intervalopen{a}{b}$ and $B_2 = \intervalopen{c}{d}\setminus \kset$.}
            Then $x\in\reals$ and $a < x$ and $x < b$ and $c < x$ and $x < d$ and $x\notin\kset$
                by \cref{inter_elim_left,inter_elim_right,setminus_elim_left,setminus_elim_right,intervalopen}.
            Take $e,f\in\reals$ such that $e < x$ and $x < f$ and
                $\intervalopen{e}{f} \subseteq \intervalopen{a}{b}\inter \intervalopen{c}{d}$
                by \cref{open_interval_refinement}.
            Hence there exists $B_3\in\kBasisR$ such that $x\in B_3$ and $B_3\subseteq B_1\inter B_2$
                by \cref{setminus_elim_left,intervalopen,subseteq,powerset_intro,k_basis_reals,real_lt_trans,setminus_intro,setminus_elim_right,inter_elim_left,inter_elim_right,inter_intro}.
        \caseOf{$B_1 = \intervalopen{a}{b}\setminus \kset$ and $B_2 = \intervalopen{c}{d}\setminus \kset$.}
            Then $x\in\reals$ and $a < x$ and $x < b$ and $c < x$ and $x < d$ and $x\notin\kset$
                by \cref{inter_elim_left,inter_elim_right,setminus_elim_left,setminus_elim_right,intervalopen}.
            Take $e,f\in\reals$ such that $e < x$ and $x < f$ and
                $\intervalopen{e}{f} \subseteq \intervalopen{a}{b}\inter \intervalopen{c}{d}$
                by \cref{open_interval_refinement}.
            Hence there exists $B_3\in\kBasisR$ such that $x\in B_3$ and $B_3\subseteq B_1\inter B_2$
                by \cref{setminus_elim_left,intervalopen,subseteq,powerset_intro,k_basis_reals,real_lt_trans,setminus_intro,setminus_elim_right,inter_elim_left,inter_elim_right,inter_intro}.
        \end{byCase}
    \end{subproof}
    Hence $\kBasisR$ is a basis on $\reals$ by \cref{basis_on}.
\end{proof}

% from §13 Item 20: lower limit topology strictly finer than standard topology
\begin{lemma}\label{lower_limit_strictly_finer_standard}
    $\lowerLimitTopologyR$ is strictly finer than $\standardTopologyR$ on $\reals$.
\end{lemma}
\begin{proof}
    Let $T = \standardTopologyR$.
    Let $T_1 = \lowerLimitTopologyR$.
    $\standardBasisR$ is a basis for $T$ on $\reals$ and
    $\lowerLimitBasisR$ is a basis for $T_1$ on $\reals$
        by \cref{standard_basis_is_basis,lower_limit_basis_is_basis,basis_for_topology,standard_topology_reals,lower_limit_topology_reals}.
    For all $x,B_2$ such that $x\in\reals$ and $B_2\in\standardBasisR$ and $x\in B_2$
        there exists $B_3\in\lowerLimitBasisR$ such that $x\in B_3$ and $B_3\subseteq B_2$
        by \cref{standard_basis_reals,intervalopen,intervalclosedopen_subset_reals,powerset_intro,lower_limit_basis_reals,interval_closed_open,real_leq_split,real_lt_subst_right,real_lt_trans,subseteq}.
    Hence $T_1$ is finer than $T$ on $\reals$ by \cref{basis_finer_criterion}.
    Let $U = \intervalclosedopen{0}{1}$.
    Then $0\in U$ by \cref{interval_closed_open,real_add_ident,real_naturals_subset_reals,real_one_in_naturals,subseteq,real_zero_lt_one}.
    Therefore $U\in T_1$ by \cref{real_add_ident,real_naturals_subset_reals,real_one_in_naturals,subseteq,real_zero_lt_one,intervalclosedopen_subset_reals,powerset_intro,lower_limit_basis_reals,basis_elem_open_in_generated,basis_for_topology,lower_limit_basis_is_basis,lower_limit_topology_reals}.
    Show that $U\notin T$.
    \begin{subproof}
        Suppose not.
        There exists $B_2\in\standardBasisR$ such that $0\in B_2$ and $B_2\subseteq U$
            by \cref{interval_closed_open,topology_generated_by_basis,standard_topology_reals}.
        Take $a,b\in\reals$ such that $a < b$ and $B_2 = \intervalopen{a}{b}$ by \cref{standard_basis_reals}.
        Hence $a < 0$ and $0 < b$ by \cref{intervalopen}.
        Take $p$ such that $p\in\intervalopen{a}{0}$ by \cref{rational_between_negative_and_zero}.
        Therefore $p\in B_2$ by \cref{intervalopen,real_lt_trans}.
        $p\notin U$ by \cref{intervalopen,interval_closed_open,real_add_ident,real_lt_subst_left,real_lt_trans,real_lt_irreflexive}.
        Contradiction by \cref{subseteq}.
    \end{subproof}
    Hence $T_1$ is strictly finer than $T$ on $\reals$ by \cref{strictly_finer_topology}.
\end{proof}

% from §13 Item 21: K-topology strictly finer than standard topology
\begin{lemma}\label{k_strictly_finer_standard}
    $\kTopologyR$ is strictly finer than $\standardTopologyR$ on $\reals$.
\end{lemma}
\begin{proof}
    Let $T = \standardTopologyR$.
    Let $T_2 = \kTopologyR$.
    $\standardBasisR$ is a basis for $T$ on $\reals$ and
    $\kBasisR$ is a basis for $T_2$ on $\reals$
        by \cref{standard_basis_is_basis,k_basis_is_basis,basis_for_topology,standard_topology_reals,k_topology_reals}.
    For all $x,B_2$ such that $x\in\reals$ and $B_2\in\standardBasisR$ and $x\in B_2$
        there exists $B_3\in\kBasisR$ such that $x\in B_3$ and $B_3\subseteq B_2$
        by \cref{k_basis_reals,standard_basis_reals,subseteq_refl}.
    Hence $T_2$ is finer than $T$ on $\reals$ by \cref{basis_finer_criterion}.
    Let $U = \intervalopen{\neg{1}}{1}\setminus \kset$.
    $1\in\reals$ and $\neg{1}\in\reals$ by \cref{real_naturals_subset_reals,real_one_in_naturals,subseteq,real_add_inverse}.
    $0 < 1$ and $0\in\reals$ and $\neg{1} < 0$ by \cref{real_zero_lt_one,real_add_ident,real_lt_neg,real_neg_zero,real_lt_subst_right}.
    Therefore $\neg{1} < 1$ by \cref{real_lt_trans}.
    Then $0\in U$ by \cref{intervalopen,k_set_positive,real_lt_irreflexive,setminus_intro}.
    Therefore $U\in T_2$ by \cref{setminus_elim_left,intervalopen,subseteq,powerset_intro,k_basis_reals,basis_elem_open_in_generated,basis_for_topology,k_basis_is_basis,k_topology_reals}.
    Show that $U\notin T$.
    \begin{subproof}
        Suppose not.
        There exists $B_2\in\standardBasisR$ such that $0\in B_2$ and $B_2\subseteq U$
            by \cref{topology_generated_by_basis,standard_topology_reals}.
        Take $a,b\in\reals$ such that $a < b$ and $B_2 = \intervalopen{a}{b}$ by \cref{standard_basis_reals}.
        Hence $a < 0$ and $0 < b$ by \cref{intervalopen}.
        Take $N\in\naturalsPlus$ such that $0 < 1 / N$ and $1 / N < b$ by \cref{real_small_reciprocal}.
        Then $0 < N$ by \cref{real_naturals_subset_reals,subseteq,real_mul_ident,real_add_ident,real_zero_lt_one,real_one_leq_naturalsplus,real_lt_subst_right,real_lt_trans}.
        Hence $1 / N\in\reals$ by \cref{real_naturals_subset_reals,subseteq,real_pos_nonzero,real_mul_inverse,real_div,real_mul_closed}.
        Therefore $a < 1 / N$ by \cref{real_lt_trans}.
        Then $1 / N\in B_2$ by \cref{intervalopen}.
        Hence $1 / N\notin U$ by \cref{k_set,setminus_elim_right}.
        Contradiction by \cref{subseteq}.
    \end{subproof}
    Hence $T_2$ is strictly finer than $T$ on $\reals$ by \cref{strictly_finer_topology}.
\end{proof}

% from §13 Item 22: lower limit and K topologies are not comparable
\begin{lemma}\label{lower_limit_k_not_comparable}
    $\lowerLimitTopologyR$ is not comparable with $\kTopologyR$ on $\reals$.
\end{lemma}
\begin{proof}
    Let $T_1 = \lowerLimitTopologyR$.
    Let $T_2 = \kTopologyR$.
    Let $U_1 = \intervalclosedopen{0}{1}$.
    Let $U_2 = \intervalopen{\neg{1}}{1}\setminus \kset$.
    $0\in\reals$ and $1\in\reals$ and $0 < 1$
        by \cref{real_add_ident,real_naturals_subset_reals,real_one_in_naturals,subseteq,real_zero_lt_one}.
    Therefore $U_1\in T_1$ by \cref{intervalclosedopen_subset_reals,powerset_intro,lower_limit_basis_reals,basis_elem_open_in_generated,basis_for_topology,lower_limit_basis_is_basis,lower_limit_topology_reals}.
    $\neg{1}\in\reals$ by \cref{real_add_inverse}.
    Hence $\neg{1} < 1$ by \cref{real_zero_lt_one,real_add_ident,real_lt_neg,real_neg_zero,real_lt_subst_right,real_lt_trans}.
    Therefore $U_2\in T_2$ by \cref{setminus_elim_left,intervalopen,subseteq,powerset_intro,k_basis_reals,basis_elem_open_in_generated,basis_for_topology,k_basis_is_basis,k_topology_reals}.
    Show that $U_1\notin T_2$.
    \begin{subproof}
        Suppose not.
        There exists $B_2\in\kBasisR$ such that $0\in B_2$ and $B_2\subseteq U_1$
            by \cref{interval_closed_open,topology_generated_by_basis,k_topology_reals}.
        Take $a,b\in\reals$ such that $a < b$ and
            $(B_2 = \intervalopen{a}{b} \lor B_2 = \intervalopen{a}{b}\setminus \kset)$
            by \cref{k_basis_reals}.
        Hence $a < 0$ and $0 < b$ by \cref{intervalopen,setminus_elim_left}.
        Take $p$ such that $p\in\intervalopen{a}{0}$ by \cref{rational_between_negative_and_zero}.
        Then $p\in\reals$ and $p < 0$ by \cref{intervalopen}.
        Hence $p\in\intervalopen{a}{b}$ by \cref{intervalopen,real_lt_trans}.
        $p\notin\kset$ by \cref{k_set_positive,real_lt_trans,real_lt_irreflexive}.
        Therefore $p\in B_2$ by \cref{setminus_intro}.
        $p\notin U_1$ by \cref{interval_closed_open,real_add_ident,real_lt_subst_left,real_lt_trans,real_lt_irreflexive}.
        Contradiction by \cref{subseteq}.
    \end{subproof}
    Show that $U_2\notin T_1$.
    \begin{subproof}
        Suppose not.
        There exists $B_1\in\lowerLimitBasisR$ such that $0\in B_1$ and $B_1\subseteq U_2$
            by \cref{real_zero_lt_one,real_lt_neg,real_neg_zero,real_lt_subst_right,intervalopen,k_set_positive,real_lt_irreflexive,setminus_intro,topology_generated_by_basis,lower_limit_topology_reals}.
        Take $a,b\in\reals$ such that $a < b$ and $B_1 = \intervalclosedopen{a}{b}$ by \cref{lower_limit_basis_reals}.
        Hence $a \leq 0$ and $0 < b$ by \cref{interval_closed_open}.
        Take $N\in\naturalsPlus$ such that $0 < 1 / N$ and $1 / N < b$ by \cref{real_small_reciprocal}.
        Then $0 < N$ by \cref{real_naturals_subset_reals,subseteq,real_mul_ident,real_add_ident,real_zero_lt_one,real_one_leq_naturalsplus,real_lt_subst_right,real_lt_trans}.
        Therefore $1 / N\in\reals$ by \cref{real_naturals_subset_reals,subseteq,real_pos_nonzero,real_mul_inverse,real_div,real_mul_closed}.
        Hence $a < 1 / N$ by \cref{real_leq_split,real_lt_subst_left,real_lt_trans}.
        Then $1 / N\in B_1$ by \cref{interval_closed_open}.
        Therefore $1 / N\notin U_2$ by \cref{k_set,setminus_elim_right}.
        Contradiction by \cref{subseteq}.
    \end{subproof}
    $T_1$ is not comparable with $T_2$ on $\reals$ by \cref{comparable_topologies,subseteq}.
\end{proof}

% from §13 Item 23: subbasis on a set
\begin{definition}\label{subbasis_on}
    $S$ is a subbasis on $X$ iff $S\subseteq\pow{X}$ and $\unions{S} = X$.
\end{definition}

% from §13 Item 24: finite intersections of subbasis elements
\begin{definition}\label{finite_intersections}
    $\finiteIntersections{S} = \{B\in\pow{\unions{S}} \mid \text{there exists $F\subseteq S$ such that $F$ is finite and inhabited and $B = \inters{F}$}\}$.
\end{definition}

% from §13 Item 25: topology generated by a subbasis
\begin{definition}\label{topology_generated_by_subbasis}
    $\subbasisTopology{S}{X} = \basisTopology{\finiteIntersections{S}}{X}$.
\end{definition}

% from §13 Item 26: finite intersections of a subbasis form a basis
\begin{lemma}\label{subbasis_finite_intersections_basis}
    Assume $S$ is a subbasis on $X$.
    Then $\finiteIntersections{S}$ is a basis on $X$.
\end{lemma}
\begin{proof}
    For all $B_1,B_2,x$ such that $B_1\in\finiteIntersections{S}$ and $B_2\in\finiteIntersections{S}$ and $x\in B_1\inter B_2$
        there exists $B_3\in\finiteIntersections{S}$ such that $x\in B_3$ and $B_3\subseteq B_1\inter B_2$
        by \cref{finite_intersections,union_subsets_is_subset,union_of_finite_is_finite,union_intro_left,pow_iff,inter_subseteq,powerset_intro,inters_distrib_union,subseteq_refl}.
    Hence for all $x$ such that $x\in X$ there exists $B\in\finiteIntersections{S}$ such that $x\in B$ and
        for all $B_1,B_2,x$ such that $B_1\in\finiteIntersections{S}$ and $B_2\in\finiteIntersections{S}$ and $x\in B_1\inter B_2$
        there exists $B_3\in\finiteIntersections{S}$ such that $x\in B_3$ and $B_3\subseteq B_1\inter B_2$
        by \cref{subbasis_on,unions_iff,subseteq,pow_iff,powerset_intro,singleton_subset_intro,pair_is_finite,upair_intro_left,subseteq_finite_is_finite,singleton_intro,inters_singleton,finite_intersections}.
    Therefore $\finiteIntersections{S}$ is a basis on $X$ by \cref{subbasis_on,finite_intersections,pow_iff,powerset_intro,subseteq,basis_on}.
\end{proof}

\section{§14 The Order Topology}

% from §14 Item 1: more than one element
\begin{definition}\label{more_than_one_element}
    $X$ has more than one element iff there exist $x, y\in X$ such that $x\neq y$.
\end{definition}

% from §14 Item 2: simple order relation
\begin{definition}\label{simple_order_relation}
    $R$ is a simple order relation on $X$ iff
    $R$ is a binary relation on $X$ and
    $R$ is transitive and
    $R$ is asymmetric and
    $R$ is connex on $X$.
\end{definition}

% from §14 Item 3: smallest element
\begin{definition}\label{smallest_element}
    $a$ is a smallest element of $X$ with respect to $R$ iff
    $a\in X$ and for all $x\in X$ we have $a\mathrel{R}x$ or $a = x$.
\end{definition}

% from §14 Item 4: largest element
\begin{definition}\label{largest_element}
    $b$ is a largest element of $X$ with respect to $R$ iff
    $b\in X$ and for all $x\in X$ we have $x\mathrel{R}b$ or $x = b$.
\end{definition}

% from §14 Item 5: open interval in an ordered set
\begin{definition}\label{intervalopen_rel}
    $\intervalopenrel{R}{X}{a}{b} = \{x\in X \mid a\mathrel{R}x \land x\mathrel{R}b\}$.
\end{definition}

% from §14 Item 6: half-open interval in an ordered set
\begin{definition}\label{intervalopenclosed_rel}
    $\intervalopenclosedrel{R}{X}{a}{b} = \{x\in X \mid a\mathrel{R}x \land (x\mathrel{R}b \lor x = b)\}$.
\end{definition}

% from §14 Item 7: half-open interval in an ordered set
\begin{definition}\label{intervalclosedopen_rel}
    $\intervalclosedopenrel{R}{X}{a}{b} = \{x\in X \mid (a\mathrel{R}x \lor x = a) \land x\mathrel{R}b\}$.
\end{definition}

% from §14 Item 8: closed interval in an ordered set
\begin{definition}\label{intervalclosed_rel}
    $\intervalclosedrel{R}{X}{a}{b} = \{x\in X \mid (a\mathrel{R}x \lor x = a) \land (x\mathrel{R}b \lor x = b)\}$.
\end{definition}

% from §14 Item 9: order basis
\begin{definition}\label{order_basis}
    $\orderBasis{R}{X} = \{B\in\pow{X} \mid
        (\exists a, b\in X. a\mathrel{R}b \land B = \intervalopenrel{R}{X}{a}{b}) \lor
        (\exists b\in X. \exists a_0\in X. (\forall x\in X. a_0\mathrel{R}x \lor a_0 = x) \land B = \intervalclosedopenrel{R}{X}{a_0}{b}) \lor
        (\exists a\in X. \exists b_0\in X. (\forall x\in X. x\mathrel{R}b_0 \lor x = b_0) \land B = \intervalopenclosedrel{R}{X}{a}{b_0})\}$.
\end{definition}

% from §14 Item 10: order topology
\begin{definition}\label{order_topology}
    $T$ is the order topology on $X$ with respect to $R$ iff
    $R$ is a simple order relation on $X$ and
    $X$ has more than one element and
    $T = \basisTopology{\orderBasis{R}{X}}{X}$.
\end{definition}

% from §14 helper lemma 1: distinct element in a set with more than one element
\begin{lemma}\label{more_than_one_element_has_other}
    Assume $X$ has more than one element.
    Suppose $x\in X$.
    Then there exists $y\in X$ such that $y\neq x$.
\end{lemma}
\begin{proof}
    Follows by \cref{more_than_one_element}.
\end{proof}

% from §14 helper lemma 2: non-smallest has a predecessor
\begin{lemma}\label{non_smallest_has_predecessor}
    Assume $R$ is a simple order relation on $X$.
    Suppose $x\in X$.
    Suppose $x$ is not a smallest element of $X$ with respect to $R$.
    Then there exists $y\in X$ such that $y\mathrel{R}x$.
\end{lemma}
\begin{proof}
    Follows by \cref{simple_order_relation,connex,smallest_element}.
\end{proof}

% from §14 helper lemma 3: non-largest has a successor
\begin{lemma}\label{non_largest_has_successor}
    Assume $R$ is a simple order relation on $X$.
    Suppose $x\in X$.
    Suppose $x$ is not a largest element of $X$ with respect to $R$.
    Then there exists $y\in X$ such that $x\mathrel{R}y$.
\end{lemma}
\begin{proof}
    Follows by \cref{simple_order_relation,connex,largest_element}.
\end{proof}

% from §14 helper lemma 4: open interval inclusion by endpoints
\begin{lemma}\label{intervalopenrel_subset_by_endpoints}
    Assume $R$ is transitive.
    Suppose $a,a_1,b,b_1\in X$.
    Suppose $a_1\mathrel{R}a$ or $a_1 = a$.
    Suppose $b\mathrel{R}b_1$ or $b = b_1$.
    Then $\intervalopenrel{R}{X}{a}{b} \subseteq \intervalopenrel{R}{X}{a_1}{b_1}$.
\end{lemma}
\begin{proof}
    Follows by \cref{intervalopen_rel,subseteq,transitive}.
\end{proof}

% from §14 helper lemma 5: closed-open interval inclusion by right endpoint
\begin{lemma}\label{intervalclosedopenrel_subset_by_right}
    Assume $R$ is transitive.
    Suppose $a,b,b_1\in X$.
    Suppose $b\mathrel{R}b_1$ or $b = b_1$.
    Then $\intervalclosedopenrel{R}{X}{a}{b} \subseteq \intervalclosedopenrel{R}{X}{a}{b_1}$.
\end{lemma}
\begin{proof}
    Follows by \cref{intervalclosedopen_rel,subseteq,transitive}.
\end{proof}

% from §14 helper lemma 6: open-closed interval inclusion by left endpoint
\begin{lemma}\label{intervalopenclosedrel_subset_by_left}
    Assume $R$ is transitive.
    Suppose $a,a_1,b\in X$.
    Suppose $a_1\mathrel{R}a$ or $a_1 = a$.
    Then $\intervalopenclosedrel{R}{X}{a}{b} \subseteq \intervalopenclosedrel{R}{X}{a_1}{b}$.
\end{lemma}
\begin{proof}
    Follows by \cref{intervalopenclosed_rel,subseteq,transitive}.
\end{proof}
% from §14 helper lemma 7: basis element at a smallest element
\begin{lemma}\label{orderbasis_elem_at_smallest}
    Assume $R$ is a simple order relation on $X$.
    Suppose $x$ is a smallest element of $X$ with respect to $R$.
    Suppose $B\in\orderBasis{R}{X}$.
    Suppose $x\in B$.
    Then there exists $b\in X$ such that $x\mathrel{R}b$ and $B = \intervalclosedopenrel{R}{X}{x}{b}$.
\end{lemma}
\begin{proof}
    $((\exists a, b\in X. a\mathrel{R}b \land B = \intervalopenrel{R}{X}{a}{b}) \lor
        (\exists b\in X. \exists a_0\in X. (\forall x\in X. a_0\mathrel{R}x \lor a_0 = x) \land B = \intervalclosedopenrel{R}{X}{a_0}{b}) \lor
        (\exists a\in X. \exists b_0\in X. (\forall x\in X. x\mathrel{R}b_0 \lor x = b_0) \land B = \intervalopenclosedrel{R}{X}{a}{b_0}))$
        by \cref{order_basis}.
    \begin{byCase}
    \caseOf{$\exists a, b\in X. a\mathrel{R}b \land B = \intervalopenrel{R}{X}{a}{b}$.}
        Take $a, b\in X$ such that $a\mathrel{R}b$ and $B = \intervalopenrel{R}{X}{a}{b}$.
        Then $a\mathrel{R}x$ by \cref{intervalopen_rel}.
        $x\mathrel{R}a$ or $x = a$ by \cref{smallest_element}.
        Suppose not.
        Contradiction by \cref{simple_order_relation,asymmetric}.
    \caseOf{$\exists b\in X. \exists a_0\in X. (\forall x\in X. a_0\mathrel{R}x \lor a_0 = x) \land B = \intervalclosedopenrel{R}{X}{a_0}{b}$.}
        Take $a_0, b\in X$ such that
            $(\forall x\in X. a_0\mathrel{R}x \lor a_0 = x)$ and $B = \intervalclosedopenrel{R}{X}{a_0}{b}$.
        Then $(a_0\mathrel{R}x \lor x = a_0)$ and $x\mathrel{R}b$ by \cref{intervalclosedopen_rel}.
        Suppose not.
        Then $a_0\mathrel{R}x$ and $x\mathrel{R}a_0$ by \cref{intervalclosedopen_rel,smallest_element}.
        Contradiction by \cref{simple_order_relation,asymmetric}.
        Hence $a_0 = x$.
        Therefore there exists $b\in X$ such that $x\mathrel{R}b$ and $B = \intervalclosedopenrel{R}{X}{x}{b}$.
    \caseOf{$\exists a\in X. \exists b_0\in X. (\forall x\in X. x\mathrel{R}b_0 \lor x = b_0) \land B = \intervalopenclosedrel{R}{X}{a}{b_0}$.}
        Take $a, b_0\in X$ such that
            $(\forall x\in X. x\mathrel{R}b_0 \lor x = b_0)$ and $B = \intervalopenclosedrel{R}{X}{a}{b_0}$.
        Then $a\mathrel{R}x$ by \cref{intervalopenclosed_rel}.
        $x\mathrel{R}a$ or $x = a$ by \cref{smallest_element}.
        Suppose not.
        Contradiction by \cref{simple_order_relation,asymmetric}.
    \end{byCase}
\end{proof}
% from §14 helper lemma 8: basis element at a largest element
\begin{lemma}\label{orderbasis_elem_at_largest}
    Assume $R$ is a simple order relation on $X$.
    Suppose $x$ is a largest element of $X$ with respect to $R$.
    Suppose $B\in\orderBasis{R}{X}$.
    Suppose $x\in B$.
    Then there exists $a\in X$ such that $a\mathrel{R}x$ and $B = \intervalopenclosedrel{R}{X}{a}{x}$.
\end{lemma}
\begin{proof}
    $((\exists a, b\in X. a\mathrel{R}b \land B = \intervalopenrel{R}{X}{a}{b}) \lor
        (\exists b\in X. \exists a_0\in X. (\forall x\in X. a_0\mathrel{R}x \lor a_0 = x) \land B = \intervalclosedopenrel{R}{X}{a_0}{b}) \lor
        (\exists a\in X. \exists b_0\in X. (\forall x\in X. x\mathrel{R}b_0 \lor x = b_0) \land B = \intervalopenclosedrel{R}{X}{a}{b_0}))$
        by \cref{order_basis}.
    \begin{byCase}
    \caseOf{$\exists a, b\in X. a\mathrel{R}b \land B = \intervalopenrel{R}{X}{a}{b}$.}
        Take $a, b\in X$ such that $a\mathrel{R}b$ and $B = \intervalopenrel{R}{X}{a}{b}$.
        Then $x\mathrel{R}b$ by \cref{intervalopen_rel}.
        $b\mathrel{R}x$ or $b = x$ by \cref{largest_element}.
        Suppose not.
        Contradiction by \cref{simple_order_relation,asymmetric}.
    \caseOf{$\exists b\in X. \exists a_0\in X. (\forall x\in X. a_0\mathrel{R}x \lor a_0 = x) \land B = \intervalclosedopenrel{R}{X}{a_0}{b}$.}
        Take $a_0, b\in X$ such that
            $(\forall x\in X. a_0\mathrel{R}x \lor a_0 = x)$ and $B = \intervalclosedopenrel{R}{X}{a_0}{b}$.
        Then $x\mathrel{R}b$ by \cref{intervalclosedopen_rel}.
        $b\mathrel{R}x$ or $b = x$ by \cref{largest_element}.
        Suppose not.
        Contradiction by \cref{simple_order_relation,asymmetric}.
    \caseOf{$\exists a\in X. \exists b_0\in X. (\forall x\in X. x\mathrel{R}b_0 \lor x = b_0) \land B = \intervalopenclosedrel{R}{X}{a}{b_0}$.}
        Take $a, b_0\in X$ such that
            $(\forall x\in X. x\mathrel{R}b_0 \lor x = b_0)$ and $B = \intervalopenclosedrel{R}{X}{a}{b_0}$.
        Then $a\mathrel{R}x$ and $(x\mathrel{R}b_0 \lor x = b_0)$ by \cref{intervalopenclosed_rel}.
        Suppose not.
        Then $x\mathrel{R}b_0$ and $b_0\mathrel{R}x$ by \cref{intervalopenclosed_rel,largest_element}.
        Contradiction by \cref{simple_order_relation,asymmetric}.
        Hence $b_0 = x$.
        Therefore there exists $a\in X$ such that $a\mathrel{R}x$ and $B = \intervalopenclosedrel{R}{X}{a}{x}$.
    \end{byCase}
\end{proof}

% from §14 helper lemma 9: open interval inside a basis element
\begin{lemma}\label{orderbasis_elem_contains_open_interval}
    Assume $R$ is a simple order relation on $X$.
    Suppose $x\in X$.
    Suppose $x$ is not a smallest element of $X$ with respect to $R$.
    Suppose $x$ is not a largest element of $X$ with respect to $R$.
    Suppose $B\in\orderBasis{R}{X}$.
    Suppose $x\in B$.
    Then there exist $a,b\in X$ such that $a\mathrel{R}x$ and $x\mathrel{R}b$ and
        $\intervalopenrel{R}{X}{a}{b} \subseteq B$.
\end{lemma}
\begin{proof}
    $((\exists a, b\in X. a\mathrel{R}b \land B = \intervalopenrel{R}{X}{a}{b}) \lor
        (\exists b\in X. \exists a_0\in X. (\forall x\in X. a_0\mathrel{R}x \lor a_0 = x) \land B = \intervalclosedopenrel{R}{X}{a_0}{b}) \lor
        (\exists a\in X. \exists b_0\in X. (\forall x\in X. x\mathrel{R}b_0 \lor x = b_0) \land B = \intervalopenclosedrel{R}{X}{a}{b_0}))$
        by \cref{order_basis}.
    \begin{byCase}
    \caseOf{$\exists a, b\in X. a\mathrel{R}b \land B = \intervalopenrel{R}{X}{a}{b}$.}
        Take $a, b\in X$ such that $a\mathrel{R}b$ and $B = \intervalopenrel{R}{X}{a}{b}$.
        Then $a\mathrel{R}x$ and $x\mathrel{R}b$ and $\intervalopenrel{R}{X}{a}{b} \subseteq B$
            by \cref{intervalopen_rel,subseteq_refl}.
    \caseOf{$\exists b\in X. \exists a_0\in X. (\forall x\in X. a_0\mathrel{R}x \lor a_0 = x) \land B = \intervalclosedopenrel{R}{X}{a_0}{b}$.}
        Take $a_0, b\in X$ such that
            $(\forall x\in X. a_0\mathrel{R}x \lor a_0 = x)$ and $B = \intervalclosedopenrel{R}{X}{a_0}{b}$.
        Hence $a_0\mathrel{R}x$ and $x\mathrel{R}b$ and $\intervalopenrel{R}{X}{a_0}{b} \subseteq B$
            by \cref{intervalopen_rel,intervalclosedopen_rel,smallest_element,subseteq}.
    \caseOf{$\exists a\in X. \exists b_0\in X. (\forall x\in X. x\mathrel{R}b_0 \lor x = b_0) \land B = \intervalopenclosedrel{R}{X}{a}{b_0}$.}
        Take $a, b_0\in X$ such that
            $(\forall x\in X. x\mathrel{R}b_0 \lor x = b_0)$ and $B = \intervalopenclosedrel{R}{X}{a}{b_0}$.
        Therefore $a\mathrel{R}x$ and $x\mathrel{R}b_0$ and $\intervalopenrel{R}{X}{a}{b_0} \subseteq B$
            by \cref{intervalopen_rel,intervalopenclosed_rel,largest_element,subseteq}.
    \end{byCase}
\end{proof}

% from §14 helper lemma 10: maximum of two elements in a simple order
\begin{lemma}\label{max_of_two_elements}
    Assume $R$ is a simple order relation on $X$.
    Suppose $a_1, a_2\in X$.
    Then there exists $a\in X$ such that
        $(a = a_1 \lor a = a_2) \land
        (a_1\mathrel{R}a \lor a_1 = a) \land
        (a_2\mathrel{R}a \lor a_2 = a)$.
\end{lemma}
\begin{proof}
    \begin{byCase}
    \caseOf{$a_1 = a_2$.}
        Then there exists $a\in X$ such that
            $(a = a_1 \lor a = a_2) \land
            (a_1\mathrel{R}a \lor a_1 = a) \land
            (a_2\mathrel{R}a \lor a_2 = a)$.
    \caseOf{$a_1 \neq a_2$ and $a_1\mathrel{R}a_2$.}
        Hence there exists $a\in X$ such that
            $(a = a_1 \lor a = a_2) \land
            (a_1\mathrel{R}a \lor a_1 = a) \land
            (a_2\mathrel{R}a \lor a_2 = a)$.
    \caseOf{$a_1 \neq a_2$ and $a_2\mathrel{R}a_1$.}
        Therefore there exists $a\in X$ such that
            $(a = a_1 \lor a = a_2) \land
            (a_1\mathrel{R}a \lor a_1 = a) \land
            (a_2\mathrel{R}a \lor a_2 = a)$.
    \end{byCase}
\end{proof}

% from §14 helper lemma 11: minimum of two elements in a simple order
\begin{lemma}\label{min_of_two_elements}
    Assume $R$ is a simple order relation on $X$.
    Suppose $b_1, b_2\in X$.
    Then there exists $b\in X$ such that
        $(b = b_1 \lor b = b_2) \land
        (b\mathrel{R}b_1 \lor b = b_1) \land
        (b\mathrel{R}b_2 \lor b = b_2)$.
\end{lemma}
\begin{proof}
    \begin{byCase}
    \caseOf{$b_1 = b_2$.}
        Then there exists $b\in X$ such that
            $(b = b_1 \lor b = b_2) \land
            (b\mathrel{R}b_1 \lor b = b_1) \land
            (b\mathrel{R}b_2 \lor b = b_2)$.
    \caseOf{$b_1 \neq b_2$ and $b_1\mathrel{R}b_2$.}
        Hence there exists $b\in X$ such that
            $(b = b_1 \lor b = b_2) \land
            (b\mathrel{R}b_1 \lor b = b_1) \land
            (b\mathrel{R}b_2 \lor b = b_2)$.
    \caseOf{$b_1 \neq b_2$ and $b_2\mathrel{R}b_1$.}
        Therefore there exists $b\in X$ such that
            $(b = b_1 \lor b = b_2) \land
            (b\mathrel{R}b_1 \lor b = b_1) \land
            (b\mathrel{R}b_2 \lor b = b_2)$.
    \end{byCase}
\end{proof}

% from §14 helper lemma 12: left endpoint from a disjunction
\begin{lemma}\label{rel_from_or_left}
    Suppose $a = a_1$ or $a = a_2$.
    Suppose $a_1\mathrel{R}x$ and $a_2\mathrel{R}x$.
    Then $a\mathrel{R}x$.
\end{lemma}
\begin{proof}
    Trivial.
\end{proof}

% from §14 helper lemma 13: right endpoint from a disjunction
\begin{lemma}\label{rel_from_or_right}
    Suppose $b = b_1$ or $b = b_2$.
    Suppose $x\mathrel{R}b_1$ and $x\mathrel{R}b_2$.
    Then $x\mathrel{R}b$.
\end{lemma}
\begin{proof}
    Trivial.
\end{proof}

% from §14 Item 11: order basis is a basis on X
\begin{lemma}\label{order_basis_is_basis}
    Assume $R$ is a simple order relation on $X$.
    Assume $X$ has more than one element.
    Then $\orderBasis{R}{X}$ is a basis on $X$.
\end{lemma}
\begin{proof}
    $\orderBasis{R}{X}\subseteq\pow{X}$ by \cref{order_basis,subseteq}.
    Show that for all $x\in X$ there exists $B_1\in\orderBasis{R}{X}$ such that $x\in B_1$.
    \begin{subproof}
        Fix $x\in X$.
        \begin{byCase}
        \caseOf{$x$ is a smallest element of $X$ with respect to $R$.}
            Take $y\in X$ such that $y\neq x$ by \cref{more_than_one_element_has_other}.
            $x\mathrel{R}y$ by \cref{smallest_element}.
            Let $B_1 = \intervalclosedopenrel{R}{X}{x}{y}$.
            Then $B_1\in\orderBasis{R}{X}$ and $x\in B_1$ by \cref{intervalclosedopen_rel,subseteq,powerset_intro,smallest_element,order_basis}.
        \caseOf{$x$ is not a smallest element of $X$ with respect to $R$.}
            \begin{byCase}
            \caseOf{$x$ is a largest element of $X$ with respect to $R$.}
                Take $y\in X$ such that $y\neq x$ by \cref{more_than_one_element_has_other}.
                $y\mathrel{R}x$ by \cref{largest_element}.
                Let $B_1 = \intervalopenclosedrel{R}{X}{y}{x}$.
                Then $B_1\in\orderBasis{R}{X}$ and $x\in B_1$ by \cref{intervalopenclosed_rel,subseteq,powerset_intro,largest_element,order_basis}.
            \caseOf{$x$ is not a largest element of $X$ with respect to $R$.}
                Take $a\in X$ such that $a\mathrel{R}x$ by \cref{non_smallest_has_predecessor}.
                Take $b\in X$ such that $x\mathrel{R}b$ by \cref{non_largest_has_successor}.
                Let $B_1 = \intervalopenrel{R}{X}{a}{b}$.
                Then $B_1\in\orderBasis{R}{X}$ and $x\in B_1$ by \cref{simple_order_relation,transitive,intervalopen_rel,subseteq,powerset_intro,order_basis}.
            \end{byCase}
        \end{byCase}
    \end{subproof}
    Show that for all $B_1,B_2,x$ such that $B_1\in\orderBasis{R}{X}$ and $B_2\in\orderBasis{R}{X}$ and $x\in B_1\inter B_2$
        there exists $B_3\in\orderBasis{R}{X}$ such that $x\in B_3$ and $B_3\subseteq B_1\inter B_2$.
    \begin{subproof}
        Fix $B_1,B_2$.
        Fix $x$ such that $B_1\in\orderBasis{R}{X}$ and $B_2\in\orderBasis{R}{X}$ and $x\in B_1\inter B_2$.
        $x\in B_1$ and $x\in B_2$ by \cref{inter_elim_left,inter_elim_right}.
        Hence $x\in X$ by \cref{order_basis,powerset_elim,subseteq}.
        \begin{byCase}
        \caseOf{$x$ is a smallest element of $X$ with respect to $R$.}
            Take $b_1\in X$ such that $x\mathrel{R}b_1$ and $B_1 = \intervalclosedopenrel{R}{X}{x}{b_1}$ by \cref{orderbasis_elem_at_smallest}.
            Take $b_2\in X$ such that $x\mathrel{R}b_2$ and $B_2 = \intervalclosedopenrel{R}{X}{x}{b_2}$ by \cref{orderbasis_elem_at_smallest}.
            Take $b\in X$ such that
                $(b = b_1 \lor b = b_2) \land
                (b\mathrel{R}b_1 \lor b = b_1) \land
                (b\mathrel{R}b_2 \lor b = b_2)$
                by \cref{min_of_two_elements}.
            Let $B_3 = \intervalclosedopenrel{R}{X}{x}{b}$.
            Then $B_3\in\orderBasis{R}{X}$ and $x\in B_3$
                by \cref{order_basis,intervalclosedopen_rel,subseteq,powerset_intro,smallest_element}.
            Hence $B_3\subseteq B_1\inter B_2$
                by \cref{simple_order_relation,intervalclosedopenrel_subset_by_right,subseteq_inter_iff}.
        \caseOf{$x$ is not a smallest element of $X$ with respect to $R$.}
            \begin{byCase}
            \caseOf{$x$ is a largest element of $X$ with respect to $R$.}
                Take $a_1\in X$ such that $a_1\mathrel{R}x$ and $B_1 = \intervalopenclosedrel{R}{X}{a_1}{x}$ by \cref{orderbasis_elem_at_largest}.
                Take $a_2\in X$ such that $a_2\mathrel{R}x$ and $B_2 = \intervalopenclosedrel{R}{X}{a_2}{x}$ by \cref{orderbasis_elem_at_largest}.
                Take $a\in X$ such that
                    $(a = a_1 \lor a = a_2) \land
                    (a_1\mathrel{R}a \lor a_1 = a) \land
                    (a_2\mathrel{R}a \lor a_2 = a)$
                    by \cref{max_of_two_elements}.
                Let $B_3 = \intervalopenclosedrel{R}{X}{a}{x}$.
                Then $B_3\in\orderBasis{R}{X}$ and $x\in B_3$
                    by \cref{order_basis,intervalopenclosed_rel,subseteq,powerset_intro,largest_element}.
                Hence $B_3\subseteq B_1\inter B_2$
                    by \cref{simple_order_relation,intervalopenclosedrel_subset_by_left,subseteq_inter_iff}.
            \caseOf{$x$ is not a largest element of $X$ with respect to $R$.}
                Take $a_1,b_1\in X$ such that $a_1\mathrel{R}x$ and $x\mathrel{R}b_1$ and
                    $\intervalopenrel{R}{X}{a_1}{b_1} \subseteq B_1$ by \cref{orderbasis_elem_contains_open_interval}.
                Take $a_2,b_2\in X$ such that $a_2\mathrel{R}x$ and $x\mathrel{R}b_2$ and
                    $\intervalopenrel{R}{X}{a_2}{b_2} \subseteq B_2$ by \cref{orderbasis_elem_contains_open_interval}.
                Take $a\in X$ such that
                    $(a = a_1 \lor a = a_2) \land
                    (a_1\mathrel{R}a \lor a_1 = a) \land
                    (a_2\mathrel{R}a \lor a_2 = a)$
                    by \cref{max_of_two_elements}.
                Take $b\in X$ such that
                    $(b = b_1 \lor b = b_2) \land
                    (b\mathrel{R}b_1 \lor b = b_1) \land
                    (b\mathrel{R}b_2 \lor b = b_2)$
                    by \cref{min_of_two_elements}.
                Let $B_3 = \intervalopenrel{R}{X}{a}{b}$.
                $R$ is transitive by \cref{simple_order_relation}.
                Therefore $a\mathrel{R}b$ by \hyperref[transitive]{transitivity}.
                Then $B_3\in\orderBasis{R}{X}$ by \cref{intervalopen_rel,subseteq,powerset_intro,order_basis}.
                Hence $x\in B_3$ by \cref{rel_from_or_left,rel_from_or_right,intervalopen_rel}.
                $B_3\subseteq B_1$ by \cref{intervalopenrel_subset_by_endpoints,subseteq_transitive}.
                $B_3\subseteq B_2$ by \cref{intervalopenrel_subset_by_endpoints,subseteq_transitive}.
                Therefore $B_3\subseteq B_1\inter B_2$ by \cref{subseteq_inter_iff}.
            \end{byCase}
        \end{byCase}
    \end{subproof}
    Therefore $\orderBasis{R}{X}$ is a basis on $X$ by \cref{basis_on}.
\end{proof}

% from §14 Item 12: open right ray
\begin{definition}\label{openray_right}
    $\openrayright{R}{X}{a} = \{x\in X \mid a\mathrel{R}x\}$.
\end{definition}

% from §14 Item 13: open left ray
\begin{definition}\label{openray_left}
    $\openrayleft{R}{X}{a} = \{x\in X \mid x\mathrel{R}a\}$.
\end{definition}

% from §14 Item 14: closed right ray
\begin{definition}\label{closedray_right}
    $\closedrayright{R}{X}{a} = \{x\in X \mid a\mathrel{R}x \lor x = a\}$.
\end{definition}

% from §14 Item 15: closed left ray
\begin{definition}\label{closedray_left}
    $\closedrayleft{R}{X}{a} = \{x\in X \mid x\mathrel{R}a \lor x = a\}$.
\end{definition}

% from §14 Item 16: open rays
\begin{definition}\label{open_rays}
    $\openRays{R}{X} = \{U\in\pow{X} \mid \exists a\in X. (U = \openrayright{R}{X}{a} \lor U = \openrayleft{R}{X}{a})\}$.
\end{definition}

% from §14 helper lemma 14: open interval as intersection of open rays
\begin{lemma}\label{intervalopenrel_as_intersection_openrays}
    $\intervalopenrel{R}{X}{a}{b} = \openrayright{R}{X}{a} \inter \openrayleft{R}{X}{b}$.
\end{lemma}
\begin{proof}
    For all $x$ we have $x\in \intervalopenrel{R}{X}{a}{b}$ iff $x\in \openrayright{R}{X}{a} \inter \openrayleft{R}{X}{b}$
        by \cref{intervalopen_rel,openray_right,openray_left,inter}.
    Follows by set extensionality.
\end{proof}

% from §14 helper lemma 15: closed-open interval at a smallest element is a left ray
\begin{lemma}\label{intervalclosedopenrel_smallest_eq_openrayleft}
    Assume $a_0$ is a smallest element of $X$ with respect to $R$.
    Then $\intervalclosedopenrel{R}{X}{a_0}{b} = \openrayleft{R}{X}{b}$.
\end{lemma}
\begin{proof}
    For all $x$ we have $x\in \intervalclosedopenrel{R}{X}{a_0}{b}$ iff $x\in \openrayleft{R}{X}{b}$
        by \cref{intervalclosedopen_rel,openray_left,smallest_element}.
    Follows by set extensionality.
\end{proof}

% from §14 helper lemma 16: open-closed interval at a largest element is a right ray
\begin{lemma}\label{intervalopenclosedrel_largest_eq_openrayright}
    Assume $b_0$ is a largest element of $X$ with respect to $R$.
    Then $\intervalopenclosedrel{R}{X}{a}{b_0} = \openrayright{R}{X}{a}$.
\end{lemma}
\begin{proof}
    For all $x$ we have $x\in \intervalopenclosedrel{R}{X}{a}{b_0}$ iff $x\in \openrayright{R}{X}{a}$
        by \cref{intervalopenclosed_rel,openray_right,largest_element}.
    Follows by set extensionality.
\end{proof}

% from §14 helper lemma 17: open rays form a subbasis on X
\begin{lemma}\label{open_rays_subbasis_on}
    Assume $R$ is a simple order relation on $X$.
    Assume $X$ has more than one element.
    Then $\openRays{R}{X}$ is a subbasis on $X$.
\end{lemma}
\begin{proof}
    Let $S = \openRays{R}{X}$.
    Then $S\subseteq\pow{X}$ by \cref{open_rays,subseteq}.
    Hence $\unions{S}\subseteq X$ by \cref{unions_subseteq_of_powerset_is_subseteq}.
    For all $x\in X$ there exists $y\in X$ such that $y\neq x$ by \cref{more_than_one_element_has_other}.
    Therefore for all $x\in X$ we have $x\in \unions{S}$
        by \cref{simple_order_relation,connex,openray_left,subseteq,powerset_intro,open_rays,unions_iff,openray_right}.
    Hence $\unions{S} = X$ by \cref{subseteq,subseteq_antisymmetric}.
    Therefore $S$ is a subbasis on $X$ by \cref{subbasis_on}.
\end{proof}

% from §14 helper lemma 18: order basis elements are finite intersections of open rays
\begin{lemma}\label{order_basis_in_finite_intersections}
    Assume $R$ is a simple order relation on $X$.
    Assume $X$ has more than one element.
    Suppose $B\in\orderBasis{R}{X}$.
    Then $B\in\finiteIntersections{\openRays{R}{X}}$.
\end{lemma}
\begin{proof}
    Let $S = \openRays{R}{X}$.
    Then $B\in\pow{\unions{S}}$ and
        $((\exists a, b\in X. a\mathrel{R}b \land B = \intervalopenrel{R}{X}{a}{b}) \lor
        (\exists b\in X. \exists a_0\in X. (\forall x\in X. a_0\mathrel{R}x \lor a_0 = x) \land B = \intervalclosedopenrel{R}{X}{a_0}{b}) \lor
        (\exists a\in X. \exists b_0\in X. (\forall x\in X. x\mathrel{R}b_0 \lor x = b_0) \land B = \intervalopenclosedrel{R}{X}{a}{b_0}))$
        by \cref{open_rays_subbasis_on,subbasis_on,order_basis,pow_iff,powerset_intro}.
    \begin{byCase}
    \caseOf{$\exists a, b\in X. a\mathrel{R}b \land B = \intervalopenrel{R}{X}{a}{b}$.}
        Take $a, b\in X$ such that $a\mathrel{R}b$ and $B = \intervalopenrel{R}{X}{a}{b}$.
        Let $F = \{\openrayright{R}{X}{a}, \openrayleft{R}{X}{b}\}$.
        Then $F\subseteq S$ by \cref{openray_right,openray_left,subseteq,powerset_intro,open_rays,upair_iff}.
        $F$ is finite and $F$ is inhabited by \cref{pair_is_finite,upair_intro_left}.
        Therefore $B\in\finiteIntersections{S}$
            by \cref{intervalopenrel_as_intersection_openrays,inter_as_inters,finite_intersections}.
    \caseOf{$\exists b\in X. \exists a_0\in X. (\forall x\in X. a_0\mathrel{R}x \lor a_0 = x) \land B = \intervalclosedopenrel{R}{X}{a_0}{b}$.}
        Take $a_0, b\in X$ such that
            $(\forall x\in X. a_0\mathrel{R}x \lor a_0 = x)$ and $B = \intervalclosedopenrel{R}{X}{a_0}{b}$.
        Let $F = \{\openrayleft{R}{X}{b}\}$.
        Hence $F\subseteq S$ by \cref{openray_left,subseteq,powerset_intro,open_rays,singleton_subset_intro}.
        $F$ is finite and $F$ is inhabited.
        Therefore $B\in\finiteIntersections{S}$
            by \cref{smallest_element,intervalclosedopenrel_smallest_eq_openrayleft,inters_singleton,finite_intersections}.
    \caseOf{$\exists a\in X. \exists b_0\in X. (\forall x\in X. x\mathrel{R}b_0 \lor x = b_0) \land B = \intervalopenclosedrel{R}{X}{a}{b_0}$.}
        Take $a, b_0\in X$ such that
            $(\forall x\in X. x\mathrel{R}b_0 \lor x = b_0)$ and $B = \intervalopenclosedrel{R}{X}{a}{b_0}$.
        Let $F = \{\openrayright{R}{X}{a}\}$.
        Hence $F\subseteq S$ by \cref{openray_right,subseteq,powerset_intro,open_rays,singleton_subset_intro}.
        $F$ is finite and $F$ is inhabited.
        Therefore $B\in\finiteIntersections{S}$
            by \cref{largest_element,intervalopenclosedrel_largest_eq_openrayright,inters_singleton,finite_intersections}.
    \end{byCase}
\end{proof}

% from §14 helper lemma 19: finite intersections in a topology are open
\begin{lemma}\label{finite_intersections_open_in_topology}
    Assume $T$ is a topology on $X$.
    Suppose $F\subseteq T$.
    Suppose $F$ is finite and inhabited.
    Then $\inters{F}\in T$.
\end{lemma}
\begin{proof}
    Take $k\in\naturals$ such that there exists a bijection from $k$ to $F$ by \cref{finite}.
    Let $A = \{ n\in\naturals \mid \text{for all $G,g$ such that $g$ is a bijection from $n$ to $G$ and $G\subseteq T$ and $G$ is inhabited we have $\inters{G}\in T$} \}$.
    For all $G,g$ such that $g$ is a bijection from $\emptyset$ to $G$ and $G\subseteq T$ and $G$ is inhabited we have $\inters{G}\in T$.
    Then $\emptyset\in A$.
    Show that for all $n\in A$ we have $\suc{n}\in A$.
    \begin{subproof}
            Fix $n\in A$.
            Show that for all $G,g$ such that $g$ is a bijection from $\suc{n}$ to $G$ and $G\subseteq T$ and $G$ is inhabited we have $\inters{G}\in T$.
            \begin{subproof}
                Fix $G$.
                Fix $g$ such that $g$ is a bijection from $\suc{n}$ to $G$ and $G\subseteq T$ and $G$ is inhabited.
                $g$ is a function by \cref{bijection_is_function}.
                $\dom{g} = \suc{n}$ by \cref{bijections_dom}.
                Then $n\in\dom{g}$ by \cref{suc_intro_self}.
                Let $a = g(n)$.
                Let $G_1 = \img{g}{n}$.
                For all $x\in n$ we have $x\in\suc{n}$ by \cref{suc}.
                Hence $n\subseteq \suc{n}$ by \cref{subseteq}.
                Then $G_1\subseteq G$ by \cref{img_subseteq,bijections_dom,img_dom,bijections,ran_of_surj}.
                Therefore $G_1\subseteq T$ by \cref{subseteq_transitive}.
                $g\in\surj{\suc{n}}{G}$ by \cref{bijections}.
                Hence $\ran{g} = G$ by \cref{ran_of_surj}.
                $(n,a)\in g$ by \cref{function_apply_elim}.
                Hence $a\in G$ by \cref{ran_iff}.
                Then $a\in T$ by \cref{elem_subseteq}.
                Show that for all $z\in G$ we have $z\in G_1\union\{a\}$.
                    \begin{subproof}
                        Fix $z\in G$.
                        Thus there exists $x\in\suc{n}$ such that $g(x)=z$ by \cref{bijections,surj}.
                        $x\in n$ or $x = n$ by \cref{suc}.
                        \begin{byCase}
                        \caseOf{$x\in n$.}
                            Thus $z\in G_1\union\{a\}$ by \cref{function_apply_elim,img_iff,union_intro_left}.
                        \caseOf{$x = n$.}
                            Thus $z\in G_1\union\{a\}$ by \cref{function_apply_elim,union_intro_right,singleton_intro}.
                        \end{byCase}
                \end{subproof}
                Hence $G = G_1\union\{a\}$
                    by \cref{subseteq,singleton_subset_intro,union_subsets_is_subset,subseteq_antisymmetric}.
                \begin{byCase}
                \caseOf{$n = \emptyset$.}
                    Then $\inters{G}\in T$ by \cref{img_emptyset,union_emptyset,union_comm,inters_singleton}.
                \caseOf{$n\neq\emptyset$.}
                    Then $G_1$ is inhabited by \cref{nonempty_inhabited,function_apply_elim,img_iff}.
                    Let $h = \restrl{g}{n}$.
                    $h$ is a function by \cref{restrl_of_function_is_function}.
                    For all $x\in n$ we have $x\in\dom{g}$ by \cref{suc,bijections_dom}.
                    Hence $\dom{h} = n$ by \cref{subseteq,dom_of_restrl_eq_inter,inter_absorb_supseteq_left}.
                    Then $h\in\funs{n}{G_1}$ by \cref{function_apply_elim,restrl_subseteq,elem_subseteq,img_iff,funs_intro}.
                    For all $y\in G_1$ there exists $x\in n$ such that $h(x) = y$
                        by \cref{img_iff,function_apply_iff,restrl_of_function_apply}.
                    Hence $h$ is a bijection from $n$ to $G_1$ by \cref{surj,restrl_of_function_apply,bijections,injective_function}.
                    Then $\inters{G_1}\in T$.
                    Hence $G = \cons{a}{G_1}$ by \cref{union_comm,union_eq_cons}.
                    Therefore $\inters{G}\in T$ by \cref{inters_cons,topology_on}.
                \end{byCase}
            \end{subproof}
    \end{subproof}
    Hence $A$ is an inductive set.
    Then $k\in A$ by \cref{num_naturals_smallest_inductive_set,subseteq}.
    Therefore for all $G,g$ such that $g$ is a bijection from $k$ to $G$ and $G\subseteq T$ and $G$ is inhabited we have $\inters{G}\in T$.
    Hence $\inters{F}\in T$.
\end{proof}

% from §14 Item 17: open rays form a subbasis
\begin{lemma}\label{open_rays_subbasis}
    Assume $R$ is a simple order relation on $X$.
    Assume $X$ has more than one element.
    Assume $T$ is the order topology on $X$ with respect to $R$.
    Then $\subbasisTopology{\openRays{R}{X}}{X} = T$.
\end{lemma}
\begin{proof}
    Let $S = \openRays{R}{X}$.
    Let $B = \orderBasis{R}{X}$.
    Let $T_1 = \subbasisTopology{S}{X}$.
    $T = \basisTopology{B}{X}$ by \cref{order_topology}.
    $B$ is a basis for $T$ on $X$ by \cref{order_basis_is_basis,basis_for_topology,order_topology}.
    $\finiteIntersections{S}$ is a basis for $T_1$ on $X$
        by \cref{open_rays_subbasis_on,subbasis_finite_intersections_basis,basis_for_topology,topology_generated_by_subbasis}.
    For all $x,B_2$ such that $x\in X$ and $B_2\in B$ and $x\in B_2$
        there exists $B_3\in\finiteIntersections{S}$ such that $x\in B_3$ and $B_3\subseteq B_2$
        by \cref{order_basis_in_finite_intersections,subseteq_refl}.
    Hence $T\subseteq T_1$ by \cref{basis_finer_criterion,finer_topology}.
    Show that for all $U\in S$ we have $U\in T$.
    \begin{subproof}
        Fix $U\in S$.
        Take $a\in X$ such that $U = \openrayright{R}{X}{a}$ or $U = \openrayleft{R}{X}{a}$ by \cref{open_rays}.
        \begin{byCase}
        \caseOf{$U = \openrayright{R}{X}{a}$.}
            \begin{byCase}
            \caseOf{$\exists b_0\in X. (\forall x\in X. x\mathrel{R}b_0 \lor x = b_0)$.}
                Take $b_0\in X$ such that $b_0$ is a largest element of $X$ with respect to $R$.
                Then $U = \intervalopenclosedrel{R}{X}{a}{b_0}$ by \cref{intervalopenclosedrel_largest_eq_openrayright}.
                Hence $\exists a\in X. \exists b_0\in X. (\forall x\in X. x\mathrel{R}b_0 \lor x = b_0) \land U = \intervalopenclosedrel{R}{X}{a}{b_0}$.
                Therefore $U\in B$ and $U\in T$ by \cref{intervalopenclosed_rel,subseteq,powerset_intro,order_basis,basis_elem_open_in_generated,order_basis_is_basis}.
            \caseOf{it is wrong that there exists $b_0\in X$ such that for all $x\in X$ we have $x\mathrel{R}b_0$ or $x = b_0$.}
                Show that for all $x\in U$ there exists $B_1\in B$ such that $x\in B_1$ and $B_1\subseteq U$.
                \begin{subproof}
                    Fix $x\in U$.
                    Then $x\in X$ and $a\mathrel{R}x$ by \cref{openray_right}.
                    Hence $x$ is not a largest element of $X$ with respect to $R$.
                    Take $b\in X$ such that $x\mathrel{R}b$ by \cref{non_largest_has_successor}.
                    Let $B_1 = \intervalopenrel{R}{X}{a}{b}$.
                    Then $B_1\in B$ and $x\in B_1$
                        by \cref{simple_order_relation,transitive,order_basis,intervalopen_rel,subseteq,powerset_intro}.
                    Hence $B_1\subseteq U$ by \cref{intervalopen_rel,openray_right,subseteq}.
                \end{subproof}
                Therefore $U\in T$ by \cref{openray_right,subseteq,powerset_intro,topology_generated_by_basis}.
            \end{byCase}
        \caseOf{$U = \openrayleft{R}{X}{a}$.}
            \begin{byCase}
            \caseOf{$\exists a_0\in X. (\forall x\in X. a_0\mathrel{R}x \lor a_0 = x)$.}
                Take $a_0\in X$ such that $a_0$ is a smallest element of $X$ with respect to $R$.
                Then $U = \intervalclosedopenrel{R}{X}{a_0}{a}$ by \cref{intervalclosedopenrel_smallest_eq_openrayleft}.
                Hence $\exists b\in X. \exists a_0\in X. (\forall x\in X. a_0\mathrel{R}x \lor a_0 = x) \land U = \intervalclosedopenrel{R}{X}{a_0}{b}$.
                Therefore $U\in B$ and $U\in T$ by \cref{intervalclosedopen_rel,subseteq,powerset_intro,order_basis,basis_elem_open_in_generated,order_basis_is_basis}.
            \caseOf{it is wrong that there exists $a_0\in X$ such that for all $x\in X$ we have $a_0\mathrel{R}x$ or $a_0 = x$.}
                Show that for all $x\in U$ there exists $B_1\in B$ such that $x\in B_1$ and $B_1\subseteq U$.
                \begin{subproof}
                    Fix $x\in U$.
                    Then $x\in X$ and $x\mathrel{R}a$ by \cref{openray_left}.
                    Hence $x$ is not a smallest element of $X$ with respect to $R$.
                    Take $b\in X$ such that $b\mathrel{R}x$ by \cref{non_smallest_has_predecessor}.
                    Let $B_1 = \intervalopenrel{R}{X}{b}{a}$.
                    Then $B_1\in B$ and $x\in B_1$
                        by \cref{simple_order_relation,transitive,order_basis,intervalopen_rel,subseteq,powerset_intro}.
                    Hence $B_1\subseteq U$ by \cref{intervalopen_rel,openray_left,subseteq}.
                \end{subproof}
                Therefore $U\in T$ by \cref{openray_left,subseteq,powerset_intro,topology_generated_by_basis}.
            \end{byCase}
        \end{byCase}
    \end{subproof}
    For all $x,B_2$ such that $x\in X$ and $B_2\in\finiteIntersections{S}$ and $x\in B_2$
        there exists $B_3\in B$ such that $x\in B_3$ and $B_3\subseteq B_2$
        by \cref{finite_intersections,subseteq,order_basis_is_basis,basis_topology_is_topology,order_topology,finite_intersections_open_in_topology,topology_generated_by_basis}.
    Hence $T_1\subseteq T$ by \cref{basis_finer_criterion,finer_topology}.
    Therefore $T_1 = T$ by \cref{subseteq_antisymmetric}.
\end{proof}

\section{§15 The Product Topology on $X \times Y$}

% from §15 Item 1: product basis on X×Y
\begin{definition}\label{product_basis}
    $\productBasis{T}{T_1}{X}{Y} = \{B\in\pow{X\times Y} \mid \exists U\in T. \exists V\in T_1. B = U\times V\}$.
\end{definition}

% from §15 Item 2: product topology on X×Y
\begin{definition}\label{product_topology}
    $\productTopology{T}{T_1}{X}{Y} = \basisTopology{\productBasis{T}{T_1}{X}{Y}}{X\times Y}$.
\end{definition}

% from §15 Item 3: product basis is a basis
\begin{lemma}\label{product_basis_is_basis}
    Assume $T$ is a topology on $X$ and $T_1$ is a topology on $Y$.
    Then $\productBasis{T}{T_1}{X}{Y}$ is a basis on $X\times Y$.
\end{lemma}
\begin{proof}
    Let $B = \productBasis{T}{T_1}{X}{Y}$.
    Thus $B\subseteq\pow{X\times Y}$ by \cref{product_basis,subseteq}.
    For all $p\in X\times Y$ there exists $B_1\in B$ such that $p\in B_1$
        by \cref{topology_on,product_basis,subseteq_refl,times_subseteq_left,times_subseteq_right,subseteq_transitive,powerset_intro}.
    Show that for all $B_1,B_2,p$ such that $B_1\in B$ and $B_2\in B$ and $p\in B_1\inter B_2$
        there exists $B_3\in B$ such that $p\in B_3$ and $B_3\subseteq B_1\inter B_2$.
    \begin{subproof}
        Fix $B_1,B_2$.
        Fix $p$ such that $B_1\in B$ and $B_2\in B$ and $p\in B_1\inter B_2$.
        Take $U_1,V_1$ such that $U_1\in T$ and $V_1\in T_1$ and $B_1 = U_1\times V_1$ by \cref{product_basis}.
        Take $U_2,V_2$ such that $U_2\in T$ and $V_2\in T_1$ and $B_2 = U_2\times V_2$ by \cref{product_basis}.
        Let $B_3 = (U_1\inter U_2)\times (V_1\inter V_2)$.
        Thus $B_3\in B$ and $p\in B_3$ and $B_3\subseteq B_1\inter B_2$
            by \cref{topology_on,powerset_elim,subseteq,inter_subseteq,times_subseteq_left,times_subseteq_right,subseteq_transitive,powerset_intro,product_basis,inter_times_elim,inter_times_intro}.
    \end{subproof}
    Thus $B$ is a basis on $X\times Y$ by \cref{basis_on}.
\end{proof}

% from §15 Item 4: bases for X and Y give a basis for X×Y
\begin{lemma}\label{product_basis_from_bases}
    Assume $B$ is a basis for $T$ on $X$ and $C$ is a basis for $T_1$ on $Y$.
    Let $D = \{D_1\in\pow{X\times Y} \mid \exists B_1\in B. \exists C_1\in C. D_1 = B_1\times C_1\}$.
    Then $D$ is a basis for $\productTopology{T}{T_1}{X}{Y}$ on $X\times Y$.
\end{lemma}
\begin{proof}
    Let $T_2 = \productTopology{T}{T_1}{X}{Y}$.
    $B$ is a basis on $X$ and $C$ is a basis on $Y$
        and $T = \basisTopology{B}{X}$ and $T_1 = \basisTopology{C}{Y}$ by \cref{basis_for_topology}.
    Hence $T$ is a topology on $X$ and $T_1$ is a topology on $Y$
        and $\productBasis{T}{T_1}{X}{Y}$ is a basis on $X\times Y$
        and $T_2$ is a topology on $X\times Y$
        by \cref{basis_topology_is_topology,product_basis_is_basis,product_topology}.
    For all $D_1$ such that $D_1\in D$ we have $D_1$ is open in $X\times Y$ with respect to $T_2$
        by \cref{subseteq,basis_elem_open_in_generated,basis_for_topology,open_in,product_basis,product_topology}.
    Show that for all $U,p$ such that $U\in T_2$ and $p\in U$ there exists $D_1\in D$ such that $p\in D_1$ and $D_1\subseteq U$.
    \begin{subproof}
        Fix $U$.
        Fix $p$ such that $U\in T_2$ and $p\in U$.
        Take $B_0\in\productBasis{T}{T_1}{X}{Y}$ such that $p\in B_0$ and $B_0\subseteq U$
            by \cref{product_topology,topology_generated_by_basis}.
        Take $U_1,V_1$ such that $U_1\in T$ and $V_1\in T_1$ and $B_0 = U_1\times V_1$ by \cref{product_basis}.
        Take $x,y$ such that $p = (x,y)$ and $x\in U_1$ and $y\in V_1$ by \cref{times_elem_is_tuple}.
        Take $B_1\in B$ such that $x\in B_1$ and $B_1\subseteq U_1$ by \cref{basis_for_topology,topology_generated_by_basis,open_in}.
        Take $C_1\in C$ such that $y\in C_1$ and $C_1\subseteq V_1$ by \cref{basis_for_topology,topology_generated_by_basis,open_in}.
        Let $D_1 = B_1\times C_1$.
        Then $D_1\in D$
            by \cref{basis_on,subseteq,powerset_elim,times_subseteq_left,times_subseteq_right,subseteq_transitive,powerset_intro}.
        Hence $p\in D_1$ and $D_1\subseteq U$
            by \cref{times,inter_intro,times_subseteq_left,times_subseteq_right,subseteq_transitive}.
    \end{subproof}
    Therefore $D$ is a basis for $T_2$ on $X\times Y$ by \cref{open_collection_basis}.
\end{proof}

% from §15 Item 5: projection onto the first factor
\begin{definition}\label{projection_first}
    $\piOne{X}{Y} = \{((x,y),x) \mid x\in X, y\in Y\}$.
\end{definition}

% from §15 Item 6: projection onto the second factor
\begin{definition}\label{projection_second}
    $\piTwo{X}{Y} = \{((x,y),y) \mid x\in X, y\in Y\}$.
\end{definition}

% from §15 helper lemma 1: preimage of the first projection
\begin{lemma}\label{preimg_pi1}
    Suppose $U\subseteq X$.
    Then $\preimg{\piOne{X}{Y}}{U} = U\times Y$.
\end{lemma}
\begin{proof}
    For all $p$ we have $p\in\preimg{\piOne{X}{Y}}{U}$ iff $p\in U\times Y$
        by \cref{preimg_iff,projection_first,pair_eq_iff,times,times_elem_is_tuple,subseteq}.
    Follows by set extensionality.
\end{proof}

% from §15 helper lemma 2: preimage of the second projection
\begin{lemma}\label{preimg_pi2}
    Suppose $V\subseteq Y$.
    Then $\preimg{\piTwo{X}{Y}}{V} = X\times V$.
\end{lemma}
\begin{proof}
    For all $p$ we have $p\in\preimg{\piTwo{X}{Y}}{V}$ iff $p\in X\times V$
        by \cref{preimg_iff,projection_second,pair_eq_iff,times,times_elem_is_tuple,subseteq}.
    Follows by set extensionality.
\end{proof}

% from §15 Item 7: left projection subbasis collection
\begin{definition}\label{product_subbasis_left}
    $\productSubbasisLeft{T}{X}{Y} = \{A \mid \exists U\in T. A = \preimg{\piOne{X}{Y}}{U}\}$.
\end{definition}

% from §15 Item 8: right projection subbasis collection
\begin{definition}\label{product_subbasis_right}
    $\productSubbasisRight{T_1}{X}{Y} = \{A \mid \exists V\in T_1. A = \preimg{\piTwo{X}{Y}}{V}\}$.
\end{definition}

% from §15 Item 9: projection subbasis for the product topology
\begin{lemma}\label{product_topology_subbasis}
    Assume $T$ is a topology on $X$ and $T_1$ is a topology on $Y$.
    Let $S = \productSubbasisLeft{T}{X}{Y}\union \productSubbasisRight{T_1}{X}{Y}$.
    Then $\subbasisTopology{S}{X\times Y} = \productTopology{T}{T_1}{X}{Y}$.
\end{lemma}
\begin{proof}
    Let $T_2 = \productTopology{T}{T_1}{X}{Y}$.
    Let $T_3 = \subbasisTopology{S}{X\times Y}$.
    Let $S_1 = \productSubbasisLeft{T}{X}{Y}$.
    Let $S_2 = \productSubbasisRight{T_1}{X}{Y}$.
    $\productBasis{T}{T_1}{X}{Y}$ is a basis on $X\times Y$ and $T_2$ is a topology on $X\times Y$
        by \cref{product_basis_is_basis,basis_topology_is_topology,product_topology}.
    For all $A$ such that $A\in S_1$ we have $A\in\pow{X\times Y}$
        by \cref{product_subbasis_left,topology_on,subseteq,powerset_elim,preimg_pi1,times_subseteq_left,powerset_intro}.
    For all $A$ such that $A\in S_2$ we have $A\in\pow{X\times Y}$
        by \cref{product_subbasis_right,topology_on,subseteq,powerset_elim,preimg_pi2,times_subseteq_right,powerset_intro}.
    Hence $S\subseteq\pow{X\times Y}$ by \cref{union_iff,subseteq}.
    Let $A = \preimg{\piOne{X}{Y}}{X}$.
    Then $A\in S$ and $A = X\times Y$ by \cref{topology_on,product_subbasis_left,union_intro_left,preimg_pi1,subseteq_refl}.
    Hence $\unions{S} = X\times Y$
        by \cref{unions_iff,subseteq,unions_subseteq_of_powerset_is_subseteq,subseteq_antisymmetric}.
    Therefore $S$ is a subbasis on $X\times Y$ and $\finiteIntersections{S}$ is a basis on $X\times Y$
        by \cref{subbasis_on,subbasis_finite_intersections_basis}.
    Hence $\finiteIntersections{S}$ is a basis for $T_3$ on $X\times Y$ by \cref{basis_for_topology,topology_generated_by_subbasis}.
    Show that for all $U$ such that $U\in S_1$ we have $U\in T_2$.
    \begin{subproof}
        Fix $U$ such that $U\in S_1$.
        Take $V\in T$ such that $U = \preimg{\piOne{X}{Y}}{V}$ by \cref{product_subbasis_left}.
        Then $U\in T_2$ by \cref{topology_on,subseteq,powerset_elim,preimg_pi1,product_basis,subseteq_refl,times_subseteq_left,times_subseteq_right,subseteq_transitive,powerset_intro,product_basis_is_basis,basis_elem_open_in_generated,product_topology}.
    \end{subproof}
    Show that for all $U$ such that $U\in S_2$ we have $U\in T_2$.
    \begin{subproof}
        Fix $U$ such that $U\in S_2$.
        Take $V\in T_1$ such that $U = \preimg{\piTwo{X}{Y}}{V}$ by \cref{product_subbasis_right}.
        Then $U\in T_2$ by \cref{topology_on,subseteq,powerset_elim,preimg_pi2,product_basis,subseteq_refl,times_subseteq_left,times_subseteq_right,subseteq_transitive,powerset_intro,product_basis_is_basis,basis_elem_open_in_generated,product_topology}.
    \end{subproof}
    Hence for all $U$ such that $U\in S$ we have $U\in T_2$ by \cref{union_iff}.
    $\finiteIntersections{S}\subseteq T_2$ by \cref{finite_intersections,subseteq,finite_intersections_open_in_topology}.
    $T_3\subseteq T_2$ by \cref{basis_topology_unions,topology_generated_by_subbasis,subseteq,topology_on}.
    Show that for all $B$ such that $B\in\productBasis{T}{T_1}{X}{Y}$ we have $B\in T_3$.
    \begin{subproof}
        Fix $B$ such that $B\in\productBasis{T}{T_1}{X}{Y}$.
        Take $U,V$ such that $U\in T$ and $V\in T_1$ and $B = U\times V$ by \cref{product_basis}.
        Then $U\subseteq X$ and $V\subseteq Y$ by \cref{topology_on,subseteq,powerset_elim}.
        Let $A_1 = \preimg{\piOne{X}{Y}}{U}$.
        Let $A_2 = \preimg{\piTwo{X}{Y}}{V}$.
        Hence $A_1\in S$ and $A_2\in S$ by \cref{product_subbasis_left,product_subbasis_right,union_intro_left,union_intro_right}.
        Let $F = \{A_1, A_2\}$.
        $F\subseteq S$ and $F$ is finite and $F$ is inhabited
            by \cref{subseteq,upair_elim,pair_is_finite,upair_intro_left}.
        Therefore $B\in\finiteIntersections{S}$ and $B\in T_3$
            by \cref{preimg_pi1,preimg_pi2,inter_absorb_supseteq_right,inter_absorb_supseteq_left,inter_times,inter_as_inters,subbasis_on,times_subseteq_left,times_subseteq_right,subseteq_transitive,powerset_intro,finite_intersections,basis_elem_open_in_generated,topology_generated_by_subbasis}.
    \end{subproof}
    For all $U$ such that $U\in T_2$ we have $U\in T_3$
        by \cref{product_basis_is_basis,basis_for_topology,product_topology,basis_topology_unions,subseteq,basis_topology_is_topology,topology_generated_by_subbasis,topology_on}.
    Therefore $T_2\subseteq T_3$ and $T_2 = T_3$ by \cref{subseteq,subseteq_antisymmetric}.
\end{proof}

\section{§16 The Subspace Topology}

% from §16 Item 1: subspace topology
\begin{definition}\label{subspace_topology}
    $\subspaceTopology{T}{Y} = \{A \mid \exists U\in T. A = Y\inter U\}$.
\end{definition}

% from §16 Item 2: subspace of a space
\begin{definition}\label{subspace_of}
    $Y$ is a subspace of $X$ with respect to $T$ iff $T$ is a topology on $X$ and $Y\subseteq X$.
\end{definition}

% from §16 helper lemma 1: subspace topology is a topology
\begin{lemma}\label{subspace_topology_is_topology}
    Assume $T$ is a topology on $X$.
    Assume $Y\subseteq X$.
    Then $\subspaceTopology{T}{Y}$ is a topology on $Y$.
\end{lemma}
\begin{proof}
    Let $T_1 = \subspaceTopology{T}{Y}$.
    Then $T_1\subseteq\pow{Y}$ by \cref{subspace_topology,inter_lower_left,powerset_intro,subseteq}.
    Hence $\emptyset\in T_1$ and $Y\in T_1$
        by \cref{topology_on,inter_emptyset,inter_absorb_supseteq_right,subspace_topology}.
    Show that for all $M$ such that $M\subseteq T_1$ we have $\unions{M}\in T_1$.
    \begin{subproof}
        Fix $M$ such that $M\subseteq T_1$.
        Let $N = \{U\in T \mid \exists A\in M. A = Y\inter U\}$.
        $N\subseteq T$ by \cref{subseteq}.
        For all $x$ we have $x\in\unions{M}$ iff $x\in Y\inter\unions{N}$
            by \cref{unions_iff,subseteq,subspace_topology,inter,inter_intro}.
        Hence $\unions{M} = Y\inter\unions{N}$ by set extensionality.
        Therefore $\unions{M}\in T_1$ by \cref{topology_on,subspace_topology}.
    \end{subproof}
    Show that for all $A,B$ such that $A\in T_1$ and $B\in T_1$ we have $A\inter B\in T_1$.
    \begin{subproof}
        Fix $A$.
        Fix $B$ such that $A\in T_1$ and $B\in T_1$.
        Take $U,V\in T$ such that $A = Y\inter U$ and $B = Y\inter V$ by \cref{subspace_topology}.
        For all $x$ we have $x\in A\inter B$ iff $x\in Y\inter (U\inter V)$ by \cref{inter,inter_intro}.
        Hence $A\inter B = Y\inter (U\inter V)$ by set extensionality.
        Therefore $A\inter B\in T_1$ by \cref{topology_on,subspace_topology}.
    \end{subproof}
    Hence $T_1$ is closed under arbitrary unions and closed under binary intersections.
    Therefore $T_1$ is a topology on $Y$ by \cref{topology_on}.
\end{proof}
% from §16 Item 3: basis for the subspace topology
\begin{lemma}\label{subspace_basis}
    Assume $B$ is a basis for $T$ on $X$.
    Assume $Y\subseteq X$.
    Let $C = \{A\in\pow{Y} \mid \exists B_1\in B. A = B_1\inter Y\}$.
    Then $C$ is a basis for $\subspaceTopology{T}{Y}$ on $Y$.
\end{lemma}
\begin{proof}
    Let $T_1 = \subspaceTopology{T}{Y}$.
    $B$ is a basis on $X$ and $T = \basisTopology{B}{X}$ and $T_1$ is a topology on $Y$
        by \cref{basis_for_topology,basis_topology_is_topology,subspace_topology_is_topology}.
    For all $C_1$ such that $C_1\in C$ we have $C_1$ is open in $Y$ with respect to $T_1$
        by \cref{subseteq,basis_elem_open_in_generated,basis_for_topology,inter_comm,subspace_topology,open_in}.
    Show that for all $U,x$ such that $U\in T_1$ and $x\in U$
        there exists $C_1\in C$ such that $x\in C_1$ and $C_1\subseteq U$.
    \begin{subproof}
        Fix $U$.
        Fix $x$ such that $U\in T_1$ and $x\in U$.
        Take $U_1\in T$ such that $U = Y\inter U_1$ by \cref{subspace_topology}.
        Thus $x\in Y$ and $x\in U_1$ by \cref{inter}.
        Take $B_1\in B$ such that $x\in B_1$ and $B_1\subseteq U_1$
            by \cref{basis_for_topology,topology_generated_by_basis}.
        Let $C_1 = B_1\inter Y$.
        Then $C_1\in C$.
        Hence $x\in C_1$ by \cref{inter_intro}.
        Therefore $C_1\subseteq U$ by \cref{inter_lower_left,subseteq_transitive,inter_lower_right,subseteq_inter_iff}.
    \end{subproof}
    Therefore $C$ is a basis for $T_1$ on $Y$ by \cref{open_collection_basis}.
\end{proof}

% from §16 Item 4: open in a subspace implies open in the space if the subspace is open
\begin{lemma}\label{open_in_subspace_open_in_space}
    Assume $Y$ is a subspace of $X$ with respect to $T$.
    Assume $U$ is open in $Y$ with respect to $\subspaceTopology{T}{Y}$.
    Assume $Y$ is open in $X$ with respect to $T$.
    Then $U$ is open in $X$ with respect to $T$.
\end{lemma}
\begin{proof}
    Take $V\in T$ such that $U = Y\inter V$ by \cref{subspace_of,open_in,subspace_topology}.
    Therefore $U$ is open in $X$ with respect to $T$ by \cref{topology_on,open_in}.
\end{proof}

% from §16 Item 5: product topology and subspace topology
\begin{theorem}\label{subspace_product_topology}
    Assume $A$ is a subspace of $X$ with respect to $T$.
    Assume $B$ is a subspace of $Y$ with respect to $T_1$.
    Then $\productTopology{\subspaceTopology{T}{A}}{\subspaceTopology{T_1}{B}}{A}{B}
        = \subspaceTopology{\productTopology{T}{T_1}{X}{Y}}{A\times B}$.
\end{theorem}
\begin{proof}
    Let $T_A = \subspaceTopology{T}{A}$.
    Let $T_B = \subspaceTopology{T_1}{B}$.
    Let $T_2 = \productTopology{T}{T_1}{X}{Y}$.
    Let $T_3 = \subspaceTopology{T_2}{A\times B}$.
    Let $C = \productBasis{T_A}{T_B}{A}{B}$.
    $T$ is a topology on $X$ and $T_1$ is a topology on $Y$ and
        $A\subseteq X$ and $B\subseteq Y$ and
        $T_A$ is a topology on $A$ and $T_B$ is a topology on $B$
        by \cref{subspace_of,subspace_topology_is_topology}.
    $\productBasis{T}{T_1}{X}{Y}$ is a basis on $X\times Y$ and $T_2$ is a topology on $X\times Y$
        by \cref{product_basis_is_basis,basis_topology_is_topology,product_topology}.
    Hence $T_3$ is a topology on $A\times B$
        by \cref{times_subseteq_left,times_subseteq_right,subseteq_transitive,subspace_topology_is_topology}.
    Show that for all $C_1$ such that $C_1\in C$ we have $C_1$ is open in $A\times B$ with respect to $T_3$.
    \begin{subproof}
        Fix $C_1$ such that $C_1\in C$.
        Take $U_A, V_B$ such that $U_A\in T_A$ and $V_B\in T_B$ and $C_1 = U_A\times V_B$ by \cref{product_basis}.
        Take $U\in T$ such that $U_A = A\inter U$ by \cref{subspace_topology}.
        Take $V\in T_1$ such that $V_B = B\inter V$ by \cref{subspace_topology}.
        Then $U\times V\in T_2$ and $C_1\in T_3$ and $C_1$ is open in $A\times B$ with respect to $T_3$
            by \cref{subspace_topology,topology_on,subseteq,powerset_elim,times_subseteq_left,times_subseteq_right,subseteq_transitive,powerset_intro,product_basis,basis_elem_open_in_generated,product_topology,inter_times,open_in}.
    \end{subproof}
    Show that for all $U,p$ such that $U\in T_3$ and $p\in U$
        there exists $C_1\in C$ such that $p\in C_1$ and $C_1\subseteq U$.
    \begin{subproof}
        Fix $U$.
        Fix $p$ such that $U\in T_3$ and $p\in U$.
        Take $W\in T_2$ such that $U = (A\times B)\inter W$ by \cref{subspace_topology}.
        Thus $p\in A\times B$ and $p\in W$ by \cref{inter}.
        Take $W_1\in\productBasis{T}{T_1}{X}{Y}$ such that $p\in W_1$ and $W_1\subseteq W$
            by \cref{product_basis_is_basis,basis_for_topology,product_topology,topology_generated_by_basis}.
        Take $U_1, V_1$ such that $U_1\in T$ and $V_1\in T_1$ and $W_1 = U_1\times V_1$ by \cref{product_basis}.
        Let $U_A = A\inter U_1$.
        Let $V_B = B\inter V_1$.
        Then $U_A\in T_A$ and $V_B\in T_B$ by \cref{subspace_topology}.
        Let $C_1 = U_A\times V_B$.
        Hence $C_1\in C$
            by \cref{topology_on,subseteq,powerset_elim,times_subseteq_left,times_subseteq_right,subseteq_transitive,powerset_intro,product_basis}.
        Take $a,b$ such that $p = (a,b)$ and $a\in A$ and $b\in B$ by \cref{times_elem_is_tuple}.
        Take $a_1,b_1$ such that $p = (a_1,b_1)$ and $a_1\in U_1$ and $b_1\in V_1$ by \cref{times_elem_is_tuple}.
        Hence $p\in C_1$ by \cref{pair_eq_iff,inter_intro,times}.
        Therefore $C_1\subseteq U$
            by \cref{inter_times,inter_lower_left,inter_lower_right,subseteq_transitive,subseteq_inter_iff}.
    \end{subproof}
    Hence $C$ is a basis for $T_3$ on $A\times B$ and $T_3 = \basisTopology{C}{A\times B}$
        by \cref{open_collection_basis,basis_for_topology}.
    Therefore $\productTopology{T_A}{T_B}{A}{B} = T_3$ by \cref{product_topology}.
\end{proof}
% from §16 Item 6: convex subset of an ordered set
\begin{definition}\label{convex_in}
    $Y$ is convex in $X$ with respect to $R$ iff
    $Y\subseteq X$ and
    for all $a,b$ such that $a\in Y$ and $b\in Y$ and $a\mathrel{R}b$
        for all $x$ such that $x\in X$ and $a\mathrel{R}x$ and $x\mathrel{R}b$ we have $x\in Y$.
\end{definition}

% from §16 helper lemma 2: point left of one element of a convex set is left of all of it
\begin{lemma}\label{convex_outside_left}
    Assume $Y$ is convex in $X$ with respect to $R$.
    Assume $R$ is a simple order relation on $X$.
    Suppose $c\in X$ and $c\notin Y$.
    Suppose $y_0\in Y$ and $c\mathrel{R}y_0$.
    Then for all $y$ such that $y\in Y$ we have $c\mathrel{R}y$.
\end{lemma}
\begin{proof}
    Fix $y$ such that $y\in Y$.
    Suppose not.
    Then $y\in X$ and $y\neq c$ by \cref{convex_in,subseteq}.
    Thus $c\in Y$ by \cref{convex_in,subseteq,simple_order_relation,connex,transitive}.
    Contradiction.
\end{proof}

% from §16 helper lemma 3: point right of one element of a convex set is right of all of it
\begin{lemma}\label{convex_outside_right}
    Assume $Y$ is convex in $X$ with respect to $R$.
    Assume $R$ is a simple order relation on $X$.
    Suppose $c\in X$ and $c\notin Y$.
    Suppose $y_0\in Y$ and $y_0\mathrel{R}c$.
    Then for all $y$ such that $y\in Y$ we have $y\mathrel{R}c$.
\end{lemma}
\begin{proof}
    Fix $y$ such that $y\in Y$.
    Suppose not.
    Then $y\in X$ and $y\neq c$ by \cref{convex_in,subseteq}.
    Thus $c\in Y$ by \cref{convex_in,subseteq,simple_order_relation,connex,transitive}.
    Contradiction.
\end{proof}

% from §16 Item 7: convex subspace order topology equals subspace topology
\begin{theorem}\label{convex_order_subspace_topology}
    Assume $T$ is the order topology on $X$ with respect to $R$.
    Assume $Y$ is convex in $X$ with respect to $R$.
    Assume $R$ is a simple order relation on $Y$.
    Assume $Y$ has more than one element.
    Then $\basisTopology{\orderBasis{R}{Y}}{Y} = \subspaceTopology{T}{Y}$.
\end{theorem}
\begin{proof}
    Let $T_Y = \basisTopology{\orderBasis{R}{Y}}{Y}$.
    Let $T_S = \subspaceTopology{T}{Y}$.
    Let $S_Y = \openRays{R}{Y}$.
    Let $S_X = \openRays{R}{X}$.
    $R$ is a simple order relation on $X$ and $X$ has more than one element and
        $T = \basisTopology{\orderBasis{R}{X}}{X}$ by \cref{order_topology}.
    $Y\subseteq X$ by \cref{convex_in}.
    $R$ is a simple order relation on $Y$ and $Y$ has more than one element.
    Hence $\orderBasis{R}{X}$ is a basis on $X$ and $\orderBasis{R}{X}$ is a basis for $T$ on $X$
        by \cref{order_basis_is_basis,basis_for_topology,order_topology}.
    Then $T$ is a topology on $X$ and $T_S$ is a topology on $Y$
        by \cref{basis_topology_is_topology,subspace_topology_is_topology}.
    Hence $\orderBasis{R}{Y}$ is a basis on $Y$ and $T_Y$ is a topology on $Y$
        and $T_Y$ is the order topology on $Y$ with respect to $R$
        by \cref{order_basis_is_basis,basis_topology_is_topology,order_topology}.
    Therefore $\subbasisTopology{S_Y}{Y} = T_Y$ and $\subbasisTopology{S_X}{X} = T$ by \cref{open_rays_subbasis}.

    Show that for all $U$ such that $U\in S_X$ we have $U\in T$.
    \begin{subproof}
        Fix $U$ such that $U\in S_X$.
        Let $F = \{U\}$.
        $F$ is finite and $F$ is inhabited.
        Therefore $U\in\finiteIntersections{S_X}$ and $U\in T$
            by \cref{inters_singleton,singleton_subset_intro,open_rays_subbasis_on,subbasis_on,open_rays,powerset_elim,subseteq_transitive,powerset_intro,finite_intersections,subbasis_finite_intersections_basis,basis_elem_open_in_generated,topology_generated_by_subbasis}.
    \end{subproof}

    Show that for all $U$ such that $U\in S_Y$ we have $U\in T_Y$.
    \begin{subproof}
        Fix $U$ such that $U\in S_Y$.
        Let $F = \{U\}$.
        $F$ is finite and $F$ is inhabited.
        Therefore $U\in\finiteIntersections{S_Y}$ and $U\in T_Y$
            by \cref{inters_singleton,singleton_subset_intro,open_rays_subbasis_on,subbasis_on,open_rays,powerset_elim,subseteq_transitive,powerset_intro,finite_intersections,subbasis_finite_intersections_basis,basis_elem_open_in_generated,topology_generated_by_subbasis}.
    \end{subproof}

    Show that for all $U$ such that $U\in S_Y$ we have $U\in T_S$.
    \begin{subproof}
        Fix $U$ such that $U\in S_Y$.
        Take $a\in Y$ such that $U = \openrayright{R}{Y}{a}$ or $U = \openrayleft{R}{Y}{a}$ by \cref{open_rays}.
        \begin{byCase}
        \caseOf{$U = \openrayright{R}{Y}{a}$.}
            Let $U_1 = \openrayright{R}{X}{a}$.
            Then $U_1\in S_X$ and $U_1\in T$ by \cref{openray_right,subseteq,powerset_intro,open_rays}.
            Therefore $U\in T_S$ by \cref{openray_right,subseteq,inter_intro,inter,subseteq_antisymmetric,subspace_topology}.
        \caseOf{$U = \openrayleft{R}{Y}{a}$.}
            Let $U_1 = \openrayleft{R}{X}{a}$.
            Then $U_1\in S_X$ and $U_1\in T$ by \cref{openray_left,subseteq,powerset_intro,open_rays}.
            Therefore $U\in T_S$ by \cref{openray_left,subseteq,inter_intro,inter,subseteq_antisymmetric,subspace_topology}.
        \end{byCase}
    \end{subproof}

    For all $U$ such that $U\in T_Y$ there exists $M\subseteq\finiteIntersections{S_Y}$ such that $U = \unions{M}$
        by \cref{open_rays_subbasis_on,subbasis_finite_intersections_basis,basis_for_topology,topology_generated_by_subbasis,basis_topology_unions}.
    For all $B$ such that $B\in\finiteIntersections{S_Y}$ we have $B\in T_S$
        by \cref{finite_intersections,subseteq,finite_intersections_open_in_topology}.
    Hence for all $U$ such that $U\in T_Y$ we have $U\in T_S$ by \cref{subseteq,topology_on}.
    Hence $T_Y\subseteq T_S$ by \cref{subseteq}.

    Show that for all $c\in X$ we have $Y\inter\openrayright{R}{X}{c}\in T_Y$.
    \begin{subproof}
        Fix $c$ such that $c\in X$.
        \begin{byCase}
        \caseOf{$c\in Y$.}
            For all $x$ such that $x\in Y\inter\openrayright{R}{X}{c}$ we have $x\in \openrayright{R}{Y}{c}$
                by \cref{inter,openray_right}.
            For all $x$ such that $x\in \openrayright{R}{Y}{c}$ we have $x\in Y\inter\openrayright{R}{X}{c}$
                by \cref{openray_right,subseteq,inter_intro}.
            Therefore $Y\inter\openrayright{R}{X}{c}\in T_Y$
                by \cref{subseteq_antisymmetric,openray_right,subseteq,powerset_intro,open_rays}.
            \caseOf{$c\notin Y$.}
                Take $y_0, y_1\in Y$ such that $y_0\neq y_1$ by \cref{more_than_one_element}.
                Then $y_0\neq c$ and $(c\mathrel{R}y_0 \lor y_0\mathrel{R}c)$
                    by \cref{subseteq,simple_order_relation,connex}.
            \begin{byCase}
            \caseOf{$c\mathrel{R}y_0$.}
                For all $y$ such that $y\in Y$ we have $c\mathrel{R}y$
                    by \cref{convex_outside_left}.
                Therefore $Y\subseteq \openrayright{R}{X}{c}$ by \cref{subseteq,elem_subseteq,openray_right}.
                Hence $Y\inter\openrayright{R}{X}{c}\in T_Y$ by \cref{inter_absorb_supseteq_right,topology_on}.
            \caseOf{$y_0\mathrel{R}c$.}
                For all $y$ such that $y\in Y$ we have $y\mathrel{R}c$
                    by \cref{convex_outside_right}.
                Show that $Y\inter\openrayright{R}{X}{c} = \emptyset$.
                \begin{subproof}
                    Suppose not.
                    Take $x$ such that $x\in Y\inter\openrayright{R}{X}{c}$ by \cref{nonempty_inhabited}.
                    Then $c\mathrel{R}x$ and $x\mathrel{R}c$ by \cref{inter,openray_right}.
                    Contradiction by \cref{simple_order_relation,asymmetric}.
                \end{subproof}
                Therefore $Y\inter\openrayright{R}{X}{c}\in T_Y$ by \cref{topology_on}.
            \end{byCase}
        \end{byCase}
    \end{subproof}

    Show that for all $c\in X$ we have $Y\inter\openrayleft{R}{X}{c}\in T_Y$.
    \begin{subproof}
        Fix $c$ such that $c\in X$.
        \begin{byCase}
        \caseOf{$c\in Y$.}
            For all $x$ such that $x\in Y\inter\openrayleft{R}{X}{c}$ we have $x\in \openrayleft{R}{Y}{c}$
                by \cref{inter,openray_left}.
            For all $x$ such that $x\in \openrayleft{R}{Y}{c}$ we have $x\in Y\inter\openrayleft{R}{X}{c}$
                by \cref{openray_left,subseteq,inter_intro}.
            Therefore $Y\inter\openrayleft{R}{X}{c}\in T_Y$
                by \cref{subseteq_antisymmetric,openray_left,subseteq,powerset_intro,open_rays}.
        \caseOf{$c\notin Y$.}
            Take $y_0, y_1\in Y$ such that $y_0\neq y_1$ by \cref{more_than_one_element}.
            Then $y_0\neq c$ and $(c\mathrel{R}y_0 \lor y_0\mathrel{R}c)$
                by \cref{subseteq,simple_order_relation,connex}.
            \begin{byCase}
            \caseOf{$y_0\mathrel{R}c$.}
                For all $y$ such that $y\in Y$ we have $y\mathrel{R}c$
                    by \cref{convex_outside_right}.
                Therefore $Y\subseteq \openrayleft{R}{X}{c}$ by \cref{subseteq,elem_subseteq,openray_left}.
                Hence $Y\inter\openrayleft{R}{X}{c}\in T_Y$ by \cref{inter_absorb_supseteq_right,topology_on}.
            \caseOf{$c\mathrel{R}y_0$.}
                For all $y$ such that $y\in Y$ we have $c\mathrel{R}y$
                    by \cref{convex_outside_left}.
                Show that $Y\inter\openrayleft{R}{X}{c} = \emptyset$.
                \begin{subproof}
                    Suppose not.
                    Take $x$ such that $x\in Y\inter\openrayleft{R}{X}{c}$ by \cref{nonempty_inhabited}.
                    Then $x\mathrel{R}c$ and $c\mathrel{R}x$ by \cref{inter,openray_left}.
                    Contradiction by \cref{simple_order_relation,asymmetric}.
                \end{subproof}
                Therefore $Y\inter\openrayleft{R}{X}{c}\in T_Y$ by \cref{topology_on}.
            \end{byCase}
        \end{byCase}
    \end{subproof}

    Show that for all $B$ such that $B\in\orderBasis{R}{X}$ we have $Y\inter B\in T_Y$.
    \begin{subproof}
        Fix $B$ such that $B\in\orderBasis{R}{X}$.
        $((\exists a, b\in X. a\mathrel{R}b \land B = \intervalopenrel{R}{X}{a}{b}) \lor
            (\exists b\in X. \exists a_0\in X. (\forall x\in X. a_0\mathrel{R}x \lor a_0 = x) \land B = \intervalclosedopenrel{R}{X}{a_0}{b}) \lor
            (\exists a\in X. \exists b_0\in X. (\forall x\in X. x\mathrel{R}b_0 \lor x = b_0) \land B = \intervalopenclosedrel{R}{X}{a}{b_0}))$
            by \cref{order_basis}.
        \begin{byCase}
        \caseOf{$\exists a, b\in X. a\mathrel{R}b \land B = \intervalopenrel{R}{X}{a}{b}$.}
            Take $a, b\in X$ such that $a\mathrel{R}b$ and $B = \intervalopenrel{R}{X}{a}{b}$.
            Thus $B = \openrayright{R}{X}{a} \inter \openrayleft{R}{X}{b}$ by \cref{intervalopenrel_as_intersection_openrays}.
            $Y\inter\openrayright{R}{X}{a}\in T_Y$ and $Y\inter\openrayleft{R}{X}{b}\in T_Y$
                and $T_Y$ is closed under binary intersections by \cref{topology_on}.
            For all $x$ such that $x\in Y\inter B$
                we have $x\in (Y\inter\openrayright{R}{X}{a}) \inter (Y\inter\openrayleft{R}{X}{b})$
                by \cref{inter,inter_intro}.
            Hence $Y\inter B\subseteq (Y\inter\openrayright{R}{X}{a}) \inter (Y\inter\openrayleft{R}{X}{b})$ by \cref{subseteq}.
            For all $x$ such that $x\in (Y\inter\openrayright{R}{X}{a}) \inter (Y\inter\openrayleft{R}{X}{b})$
                we have $x\in Y\inter B$ by \cref{inter,inter_intro}.
            Therefore $(Y\inter\openrayright{R}{X}{a}) \inter (Y\inter\openrayleft{R}{X}{b})\subseteq Y\inter B$ by \cref{subseteq}.
            Hence $Y\inter B\in T_Y$ by \cref{subseteq_antisymmetric,topology_on}.
        \caseOf{$\exists b\in X. \exists a_0\in X. (\forall x\in X. a_0\mathrel{R}x \lor a_0 = x) \land B = \intervalclosedopenrel{R}{X}{a_0}{b}$.}
            Take $a_0, b\in X$ such that
                $(\forall x\in X. a_0\mathrel{R}x \lor a_0 = x)$ and $B = \intervalclosedopenrel{R}{X}{a_0}{b}$.
            Hence $B = \openrayleft{R}{X}{b}$ and $Y\inter B = Y\inter\openrayleft{R}{X}{b}$
                by \cref{smallest_element,intervalclosedopenrel_smallest_eq_openrayleft}.
            Therefore $Y\inter B\in T_Y$.
        \caseOf{$\exists a\in X. \exists b_0\in X. (\forall x\in X. x\mathrel{R}b_0 \lor x = b_0) \land B = \intervalopenclosedrel{R}{X}{a}{b_0}$.}
            Take $a, b_0\in X$ such that
                $(\forall x\in X. x\mathrel{R}b_0 \lor x = b_0)$ and $B = \intervalopenclosedrel{R}{X}{a}{b_0}$.
            Hence $B = \openrayright{R}{X}{a}$ and $Y\inter B = Y\inter\openrayright{R}{X}{a}$
                by \cref{largest_element,intervalopenclosedrel_largest_eq_openrayright}.
            Therefore $Y\inter B\in T_Y$.
        \end{byCase}
    \end{subproof}

    Show that for all $U$ such that $U\in T_S$ we have $U\in T_Y$.
    \begin{subproof}
            Fix $U$ such that $U\in T_S$.
            Take $V\in T$ such that $U = Y\inter V$ by \cref{subspace_topology}.
            Then there exists $M\subseteq\orderBasis{R}{X}$ such that $V = \unions{M}$ by \cref{basis_topology_unions}.
            Let $N = \{A\in\pow{Y} \mid \exists B\in M. A = Y\inter B\}$.
            Show that for all $x$ such that $x\in U$ we have $x\in\unions{N}$.
            \begin{subproof}
                Fix $x$ such that $x\in U$.
                Thus $x\in Y$ and $x\in V$ by \cref{inter}.
                Take $B\in M$ such that $x\in B$ by \cref{unions_iff}.
                Then $x\in Y\inter B$ by \cref{inter_intro}.
                Hence $Y\inter B\in N$.
                Therefore $x\in\unions{N}$ by \cref{unions_iff}.
            \end{subproof}
            Hence $U\subseteq\unions{N}$ by \cref{subseteq}.
            For all $x$ such that $x\in\unions{N}$ we have $x\in U$ by \cref{unions_iff,inter,inter_intro}.
            Therefore $\unions{N}\subseteq U$ by \cref{subseteq}.
            For all $A$ such that $A\in N$ we have $A\in T_Y$ by \cref{subseteq}.
            Therefore $U\in T_Y$ by \cref{subseteq,topology_on,subseteq_antisymmetric}.
    \end{subproof}
    Therefore $T_S\subseteq T_Y$ by \cref{subseteq}.

    Hence $T_Y = T_S$ by \cref{subseteq_antisymmetric}.
\end{proof}
\section{§17 Closed Sets and Limit Points}

% from §17 Item 1: closed set
\begin{definition}\label{closed_in}
    $A$ is closed in $X$ with respect to $T$ iff
    $A\subseteq X$ and $X\setminus A$ is open in $X$ with respect to $T$.
\end{definition}

% from §17 Item 2: empty set is closed
\begin{lemma}\label{closed_emptyset}
    Assume $T$ is a topology on $X$.
    Then $\emptyset$ is closed in $X$ with respect to $T$.
\end{lemma}
\begin{proof}
    Follows by \cref{closed_in,emptyset_subseteq,setminus_emptyset,topology_on,open_in}.
\end{proof}

% from §17 Item 3: the whole space is closed
\begin{lemma}\label{closed_whole_space}
    Assume $T$ is a topology on $X$.
    Then $X$ is closed in $X$ with respect to $T$.
\end{lemma}
\begin{proof}
    Follows by \cref{closed_in,subseteq_refl,setminus_self,topology_on,open_in}.
\end{proof}

% from §17 Item 4: intersections of closed sets are closed
\begin{lemma}\label{closed_inters}
    Assume $T$ is a topology on $X$.
    Suppose $C\subseteq\pow{X}$.
    Suppose for all $A$ such that $A\in C$ we have $A$ is closed in $X$ with respect to $T$.
    Then $\inters{C}$ is closed in $X$ with respect to $T$.
\end{lemma}
\begin{proof}
    Then $\inters{C}\subseteq X$ by \cref{unions_subseteq_of_powerset_is_subseteq,inters,elem_subseteq,subseteq}.
    \begin{byCase}
    \caseOf{$C = \emptyset$.}
        Show that $\inters{C} = \emptyset$.
        \begin{subproof}
            Suppose not.
            Take $x$ such that $x\in\inters{C}$ by \cref{nonempty_inhabited}.
            Contradiction by \cref{inters,unions_emptyset,emptyset}.
        \end{subproof}
        Therefore $\inters{C}$ is closed in $X$ with respect to $T$ by \cref{closed_emptyset}.
    \caseOf{$C \neq \emptyset$.}
        Then $C$ is inhabited by \cref{nonempty_inhabited}.
        Let $D = \{U\in T \mid \exists A\in C. U = X\setminus A\}$.
        Show that $X\setminus\inters{C} = \unions{D}$.
        \begin{subproof}
            Show that for all $x$ such that $x\in X\setminus\inters{C}$ we have $x\in\unions{D}$.
            \begin{subproof}
                Fix $x$ such that $x\in X\setminus\inters{C}$.
                Suppose not.
                Then $x\notin\unions{D}$.
                Show that for all $A$ such that $A\in C$ we have $x\in A$.
                \begin{subproof}
                    Fix $A$ such that $A\in C$.
                    Suppose not.
                    Then $x\in\unions{D}$ by \cref{setminus_elim_left,setminus_intro,closed_in,open_in,unions_iff}.
                    Contradiction.
                \end{subproof}
                Hence $x\in\inters{C}$ by \cref{inters_intro}.
                Contradiction by \cref{setminus_elim_right}.
            \end{subproof}
            Therefore $X\setminus\inters{C}\subseteq\unions{D}$ by \cref{subseteq}.
            Show that for all $x$ such that $x\in\unions{D}$ we have $x\in X\setminus\inters{C}$.
            \begin{subproof}
                Fix $x$ such that $x\in\unions{D}$.
                Take $U\in D$ such that $x\in U$ by \cref{unions_iff}.
                Take $A\in C$ such that $U = X\setminus A$.
                Then $x\in X$ and $x\notin A$ by \cref{setminus_elim_left,setminus_elim_right}.
                Then $x\notin\inters{C}$.
                Therefore $x\in X\setminus\inters{C}$ by \cref{setminus_intro}.
            \end{subproof}
            Hence $\unions{D}\subseteq X\setminus\inters{C}$ by \cref{subseteq}.
            Therefore $X\setminus\inters{C} = \unions{D}$ by \cref{subseteq_antisymmetric}.
        \end{subproof}
        Hence $X\setminus\inters{C}$ is open in $X$ with respect to $T$ by \cref{subseteq,topology_on,open_in}.
        Hence $\inters{C}$ is closed in $X$ with respect to $T$ by \cref{closed_in}.
    \end{byCase}
\end{proof}

% from §17 helper lemma 1: complements of subsets are injective
\begin{lemma}\label{complement_eq_implies_eq}
    Suppose $A\subseteq X$ and $B\subseteq X$.
    Suppose $X\setminus A = X\setminus B$.
    Then $A = B$.
\end{lemma}
\begin{proof}
    Follows by \cref{subseteq_refl,double_complement}.
\end{proof}

% from §17 Item 5: finite unions of closed sets are closed
\begin{lemma}\label{closed_finite_unions}
    Assume $T$ is a topology on $X$.
    Suppose $F\subseteq\pow{X}$.
    Suppose $F$ is finite.
    Suppose for all $A$ such that $A\in F$ we have $A$ is closed in $X$ with respect to $T$.
    Then $\unions{F}$ is closed in $X$ with respect to $T$.
\end{lemma}
\begin{proof}
    \begin{byCase}
    \caseOf{$F = \emptyset$.}
        Therefore $\unions{F}$ is closed in $X$ with respect to $T$ by \cref{unions_emptyset,closed_emptyset}.
    \caseOf{$F \neq \emptyset$.}
        Let $G = \{U\in T \mid \exists A\in F. U = X\setminus A\}$.
        $G$ is inhabited by \cref{nonempty_inhabited,closed_in,open_in}.
        Show that for all $x$ such that $x\in X\setminus\unions{F}$ we have $x\in\inters{G}$.
        \begin{subproof}
            Fix $x$ such that $x\in X\setminus\unions{F}$.
            Thus $x\in X$ and $x\notin\unions{F}$ by \cref{setminus_elim_left,setminus_elim_right}.
            Show that for all $U$ such that $U\in G$ we have $x\in U$.
            \begin{subproof}
                Fix $U$ such that $U\in G$.
                Take $A\in F$ such that $U = X\setminus A$.
                Suppose not.
                Then $x\notin U$.
                Thus $x\notin X\setminus A$.
                Show that $x\in A$.
                \begin{subproof}
                    Suppose not.
                    Then $x\notin A$.
                    Then $x\in X\setminus A$ by \cref{setminus_intro}.
                    Contradiction.
                \end{subproof}
                Hence $x\in\unions{F}$ by \cref{unions_iff}.
                Contradiction.
            \end{subproof}
            Hence $x\in\inters{G}$ by \cref{inters_intro}.
        \end{subproof}
        Therefore $X\setminus\unions{F}\subseteq\inters{G}$ by \cref{subseteq}.
        Show that for all $x$ such that $x\in\inters{G}$ we have $x\in X\setminus\unions{F}$.
        \begin{subproof}
            Fix $x$ such that $x\in\inters{G}$.
            Take $A$ such that $A\in F$.
            Hence $x\in X\setminus A$ by \cref{closed_in,open_in,inters_destr}.
            Show that $x\notin\unions{F}$.
            \begin{subproof}
                Suppose not.
                Then $x\in\unions{F}$.
                Take $A_1\in F$ such that $x\in A_1$ by \cref{unions_iff}.
                Hence $x\notin A_1$ by \cref{closed_in,open_in,inters_destr,setminus_elim_right}.
                Contradiction.
            \end{subproof}
            Therefore $x\in X\setminus\unions{F}$ by \cref{inters_destr,setminus_elim_left,setminus_intro}.
        \end{subproof}
        Hence $\inters{G}\subseteq X\setminus\unions{F}$ by \cref{subseteq}.
        Therefore $X\setminus\unions{F} = \inters{G}$ by \cref{subseteq_antisymmetric}.
        $G\subseteq T$ by \cref{subseteq}.
        Take $k$ such that $k$ is a natural number and there exists a bijection from $k$ to $F$ by \cref{finite}.
        Take $f$ such that $f$ is a bijection from $k$ to $F$.
        Let $g = \{(n, X\setminus f(n)) \mid n\in k\}$.
        Show that $g$ is a function to $G$ and $\dom{g} = k$.
        \begin{subproof}
            For all $w$ such that $w\in g$ there exist $x,y$ such that $w = (x,y)$.
            For all $n,y,y'$ such that $(n,y)\in g$ and $(n,y')\in g$ we have $y = y'$ by \cref{pair_eq_iff}.
            Hence $g$ is a function by \cref{relation,rightunique}.
            For all $n$ such that $n\in k$ we have $n\in\dom{g}$ by \cref{dom_intro}.
            Hence $k\subseteq\dom{g}$ by \cref{subseteq}.
            For all $n$ such that $n\in\dom{g}$ we have $n\in k$ by \cref{dom_iff,pair_eq_iff}.
            Hence $\dom{g}\subseteq k$ by \cref{subseteq}.
            Hence $\dom{g} = k$ by \cref{subseteq_antisymmetric}.
            For all $n$ such that $n\in\dom{g}$ we have $f(n)\in F$ and $X\setminus f(n)\in G$ and
                $(n, X\setminus f(n))\in g$ by \cref{bijections,bijections_to_funs,funs_elim,closed_in,open_in}.
            Thus for all $n$ such that $n\in\dom{g}$ we have $g(n)\in G$ by \cref{function_apply_iff}.
            Hence $g$ is a function to $G$.
        \end{subproof}
        Then $g\in\funs{k}{G}$ by \cref{funs_intro}.
        Hence $g\in\surj{k}{G}$ by \cref{bijections,surj,function_apply_iff}.
        Thus for all $n,m$ such that $n\in k$ and $m\in k$ and $g(n) = g(m)$ we have $n = m$
            by \cref{function_apply_iff,bijections,bijections_to_funs,funs_elim,subseteq,pow_iff,complement_eq_implies_eq,injective_function}.
        Therefore $g$ is injective by \cref{injective_function}.
        Hence $g$ is a bijection from $k$ to $G$ by \cref{bijections}.
        Therefore $G$ is finite by \cref{finite}.
        Therefore $X\setminus\unions{F}$ is open in $X$ with respect to $T$
            by \cref{subseteq,finite_intersections_open_in_topology,open_in}.
        Hence $\unions{F}$ is closed in $X$ with respect to $T$
            by \cref{unions_subseteq_of_powerset_is_subseteq,closed_in}.
    \end{byCase}
\end{proof}

% from §17 Item 6: closed sets in subspaces
\begin{theorem}\label{closed_in_subspace_iff}
    Assume $Y$ is a subspace of $X$ with respect to $T$.
    Then for all $A$ we have
    $A$ is closed in $Y$ with respect to $\subspaceTopology{T}{Y}$
    iff there exists $C$ such that $C$ is closed in $X$ with respect to $T$ and $A = C\inter Y$.
\end{theorem}
\begin{proof}
    $T$ is a topology on $X$ and $Y\subseteq X$ by \cref{subspace_of}.
    Let $T_Y = \subspaceTopology{T}{Y}$.
    Fix $A$.
    Show that if $A$ is closed in $Y$ with respect to $T_Y$ then
        there exists $C$ such that $C$ is closed in $X$ with respect to $T$ and $A = C\inter Y$.
    \begin{subproof}
        Assume $A$ is closed in $Y$ with respect to $T_Y$.
        Then $Y\setminus A\in T_Y$ by \cref{closed_in,open_in}.
        Take $U\in T$ such that $Y\setminus A = Y\inter U$ by \cref{subspace_topology}.
        Let $C = X\setminus U$.
        Hence $U\subseteq X$ by \cref{open_in,topology_on,subseteq,pow_iff}.
        Hence $C$ is closed in $X$ with respect to $T$ by
            \cref{setminus_subseteq,setminus_setminus,inter_comm,topology_on,open_in,subseteq,pow_iff,inter_absorb_supseteq_right,closed_in}.
        For all $x$ such that $x\in A$ we have $x\in Y$ by \cref{closed_in,elem_subseteq}.
        Show that for all $x$ such that $x\in A$ we have $x\notin U$.
        \begin{subproof}
            Fix $x$ such that $x\in A$.
            Suppose not.
            Then $x\in U$.
            Thus $x\in Y\setminus A$ by \cref{closed_in,elem_subseteq,inter_intro}.
            Contradiction by \cref{setminus_elim_right}.
        \end{subproof}
        Hence $A\subseteq Y\setminus U$ by \cref{setminus_intro,subseteq}.
        Show that for all $x$ such that $x\in Y\setminus U$ we have $x\in A$.
        \begin{subproof}
            Fix $x$ such that $x\in Y\setminus U$.
            Suppose not.
            Then $x\notin A$.
            Thus $x\notin U$ by \cref{setminus_elim_left,setminus_intro,inter,setminus_elim_right}.
            Contradiction.
        \end{subproof}
        Therefore $Y\setminus U\subseteq A$ by \cref{subseteq}.
        Hence $A = Y\inter (X\setminus U)$ by \cref{subseteq_antisymmetric,setminus_eq_inter_complement}.
        Therefore $A = C\inter Y$ by \cref{inter_comm}.
    \end{subproof}
    Show that if there exists $C$ such that $C$ is closed in $X$ with respect to $T$ and $A = C\inter Y$ then
        $A$ is closed in $Y$ with respect to $T_Y$.
    \begin{subproof}
        Assume there exists $C$ such that $C$ is closed in $X$ with respect to $T$ and $A = C\inter Y$.
        Take $C$ such that $C$ is closed in $X$ with respect to $T$ and $A = C\inter Y$.
        For all $x$ such that $x\in Y\setminus A$ we have $x\in Y\setminus C$.
        Hence $Y\setminus A\subseteq Y\setminus C$ by \cref{subseteq}.
        For all $x$ such that $x\in Y\setminus C$ we have $x\in Y\setminus A$.
        Therefore $Y\setminus C\subseteq Y\setminus A$ by \cref{subseteq}.
        Hence $Y\setminus A = Y\setminus C$ by \cref{subseteq_antisymmetric}.
        Therefore $A$ is closed in $Y$ with respect to $T_Y$ by \cref{subseteq_antisymmetric,setminus_eq_inter_complement,closed_in,open_in,subspace_topology,inter_lower_right,subspace_topology_is_topology}.
    \end{subproof}
\end{proof}

% from §17 Item 7: closed in a closed subspace
\begin{theorem}\label{closed_in_closed_subspace}
    Assume $T$ is a topology on $X$.
    Suppose $Y\subseteq X$.
    Suppose $Y$ is closed in $X$ with respect to $T$.
    Suppose $A$ is closed in $Y$ with respect to $\subspaceTopology{T}{Y}$.
    Then $A$ is closed in $X$ with respect to $T$.
\end{theorem}
\begin{proof}
    Take $C$ such that $C$ is closed in $X$ with respect to $T$ and $A = C\inter Y$
        by \cref{closed_in_subspace_iff,subspace_of}.
    Let $D = \{C, Y\}$.
    Hence $D\subseteq\pow{X}$ and for all $B$ such that $B\in D$ we have $B$ is closed in $X$ with respect to $T$
        by \cref{upair_iff,closed_in,pow_iff,subseteq}.
    Therefore $A$ is closed in $X$ with respect to $T$ by \cref{closed_inters,inter_as_inters}.
\end{proof}

% from §17 Item 8: interior of a set
\begin{definition}\label{interior}
    $\interior{A}{X}{T} = \{x\in X \mid \exists U\in T. x\in U \land U\subseteq A\}$.
\end{definition}

% from §17 Item 9: closure of a set
\begin{definition}\label{closure}
    $\closure{A}{X}{T} = \{x\in X \mid \text{for all $C$ such that $C$ is closed in $X$ with respect to $T$ and for all $y\in A$ we have $y\in C$ we have $x\in C$}\}$.
\end{definition}

% from §17 Item 10: interior is open
\begin{lemma}\label{interior_open_in}
    Assume $T$ is a topology on $X$.
    Then $\interior{A}{X}{T}$ is open in $X$ with respect to $T$.
\end{lemma}
\begin{proof}
    Let $S = \{U\in T \mid U\subseteq A\}$.
    Hence $\interior{A}{X}{T}\subseteq\unions{S}$ by \cref{interior,unions_iff,subseteq}.
    Therefore $\unions{S}\subseteq\interior{A}{X}{T}$
        by \cref{unions_iff,topology_on,subseteq,pow_iff,elem_subseteq,interior,subseteq}.
    Therefore $\interior{A}{X}{T}$ is open in $X$ with respect to $T$ by \cref{subseteq,topology_on,subseteq_antisymmetric,open_in}.
\end{proof}

% from §17 Item 11: closure is closed
\begin{lemma}\label{closure_closed_in}
    Assume $T$ is a topology on $X$.
    Suppose $A\subseteq X$.
    Then $\closure{A}{X}{T}$ is closed in $X$ with respect to $T$.
\end{lemma}
\begin{proof}
    Let $F = \{C\in\pow{X} \mid \text{$C$ is closed in $X$ with respect to $T$ and for all $x\in A$ we have $x\in C$}\}$.
    Hence $F$ is inhabited by \cref{closed_whole_space,powerset_top,subseteq}.
    Hence $\closure{A}{X}{T}\subseteq\inters{F}$ by \cref{closure,inters_iff_forall,subseteq}.
    Show that for all $x$ such that $x\in\inters{F}$ we have $x\in\closure{A}{X}{T}$.
    \begin{subproof}
        Fix $x$ such that $x\in\inters{F}$.
        Hence for all $C$ such that $C\in F$ we have $x\in C$ by \cref{inters_iff_forall}.
        Take $D$ such that $D\in F$.
        Then $x\in X$ by \cref{inters_iff_forall,pow_iff,elem_subseteq}.
        For all $C$ such that $C$ is closed in $X$ with respect to $T$ and
            for all $y\in A$ we have $y\in C$ we have $x\in C$.
        Therefore $x\in\closure{A}{X}{T}$ by \cref{closure}.
    \end{subproof}
    Hence $\inters{F}\subseteq\closure{A}{X}{T}$ by \cref{subseteq}.
    Therefore $\closure{A}{X}{T} = \inters{F}$ by \cref{subseteq_antisymmetric}.
    Hence $F\subseteq\pow{X}$ by \cref{subseteq}.
    For all $C$ such that $C\in F$ we have $C$ is closed in $X$ with respect to $T$.
    Therefore $\inters{F}$ is closed in $X$ with respect to $T$ by \cref{closed_inters}.
    Hence $\closure{A}{X}{T}$ is closed in $X$ with respect to $T$.
\end{proof}

% from §17 Item 12: interior subset
\begin{lemma}\label{interior_subseteq}
    Assume $T$ is a topology on $X$.
    Then $\interior{A}{X}{T}\subseteq A$.
\end{lemma}
\begin{proof}
    Follows by \cref{interior,elem_subseteq,subseteq}.
\end{proof}

% from §17 Item 13: subset of closure
\begin{lemma}\label{subseteq_closure}
    Assume $T$ is a topology on $X$.
    Suppose $A\subseteq X$.
    Then $A\subseteq\closure{A}{X}{T}$.
\end{lemma}
\begin{proof}
    Follows by \cref{elem_subseteq,closure,subseteq}.
\end{proof}

% from §17 helper lemma 2: closure contained in closed superset
\begin{lemma}\label{closure_subseteq_closed}
    Assume $T$ is a topology on $X$.
    Suppose $A\subseteq X$.
    Suppose $C$ is closed in $X$ with respect to $T$.
    Suppose $A\subseteq C$.
    Then $\closure{A}{X}{T}\subseteq C$.
\end{lemma}
\begin{proof}
    Follows by \cref{closure,subseteq}.
\end{proof}

% from §17 Item 14: open set equals its interior
\begin{lemma}\label{open_eq_interior}
    Assume $T$ is a topology on $X$.
    Suppose $A$ is open in $X$ with respect to $T$.
    Then $A = \interior{A}{X}{T}$.
\end{lemma}
\begin{proof}
    Follows by \cref{open_in,subseteq_refl,topology_on,subseteq,pow_iff,elem_subseteq,interior,interior_subseteq,subseteq_antisymmetric}.
\end{proof}

% from §17 Item 15: closed set equals its closure
\begin{lemma}\label{closed_eq_closure}
    Assume $T$ is a topology on $X$.
    Suppose $A$ is closed in $X$ with respect to $T$.
    Then $A = \closure{A}{X}{T}$.
\end{lemma}
\begin{proof}
    Follows by \cref{closed_in,subseteq_closure,subseteq_refl,closure_subseteq_closed,subseteq_antisymmetric}.
\end{proof}

% from §17 Item 16: closure in subspace
\begin{theorem}\label{closure_in_subspace}
    Assume $Y$ is a subspace of $X$ with respect to $T$.
    Suppose $A\subseteq Y$.
    Let $T_Y = \subspaceTopology{T}{Y}$.
    Then $\closure{A}{Y}{T_Y} = \closure{A}{X}{T}\inter Y$.
\end{theorem}
\begin{proof}
    $T$ is a topology on $X$ and $Y\subseteq X$ and $A\subseteq X$ by \cref{subspace_of,subseteq,subseteq_transitive}.
    Let $B = \closure{A}{Y}{T_Y}$.
    Let $C = \closure{A}{X}{T}$.
    Hence $C\inter Y$ is closed in $Y$ with respect to $T_Y$
        by \cref{closure_closed_in,closed_in_subspace_iff,subspace_of}.
    Therefore $A\subseteq C\inter Y$ by \cref{subseteq_closure,subseteq_inter_iff}.
    $T_Y$ is a topology on $Y$ by \cref{subspace_topology_is_topology}.
    Hence $B\subseteq C\inter Y$ by \cref{closure_subseteq_closed}.
    $B$ is closed in $Y$ with respect to $T_Y$ by \cref{closure_closed_in}.
    Take $D$ such that $D$ is closed in $X$ with respect to $T$ and $B = D\inter Y$
        by \cref{closed_in_subspace_iff,subspace_of}.
    Therefore $A\subseteq D$ by \cref{inter_lower_left,subseteq_closure,subseteq_transitive}.
    Hence $C\subseteq D$ by \cref{closure_subseteq_closed}.
    Therefore $C\inter Y\subseteq B$ by \cref{inter,elem_subseteq,inter_intro,subseteq}.
    Hence $B = C\inter Y$ by \cref{subseteq_antisymmetric}.
\end{proof}

% from §17 Item 17: intersects
\begin{definition}\label{intersects}
    $A$ intersects $B$ iff $A\inter B \neq \emptyset$.
\end{definition}
% from §17 Item 18: closure and open sets
\begin{theorem}\label{closure_iff_intersects_open}
    Assume $T$ is a topology on $X$.
    Suppose $A\subseteq X$.
    Suppose $x\in X$.
    Then $x\in\closure{A}{X}{T}$ iff for all $U$ such that
        $U$ is open in $X$ with respect to $T$ and $x\in U$ we have $U$ intersects $A$.
\end{theorem}
\begin{proof}
    Show that if $x\in\closure{A}{X}{T}$ then for all $U$ such that
        $U$ is open in $X$ with respect to $T$ and $x\in U$ we have $U$ intersects $A$.
    \begin{subproof}
        Assume $x\in\closure{A}{X}{T}$.
            Fix $U$ such that $U$ is open in $X$ with respect to $T$ and $x\in U$.
        Suppose not.
        Hence $A\subseteq X\setminus U$ by \cref{intersects,inter_comm,subseteq_setminus}.
        Therefore $X\setminus U$ is closed in $X$ with respect to $T$ by
            \cref{setminus_setminus,inter_comm,topology_on,open_in,subseteq,pow_iff,inter_absorb_supseteq_right,setminus_subseteq,closed_in}.
        Hence $x\notin U$ by \cref{closure_subseteq_closed,elem_subseteq,setminus_elim_right}.
        Contradiction.
    \end{subproof}
    Show that if for all $U$ such that $U$ is open in $X$ with respect to $T$ and $x\in U$ we have $U$ intersects $A$
        then $x\in\closure{A}{X}{T}$.
    \begin{subproof}
        Assume for all $U$ such that $U$ is open in $X$ with respect to $T$ and $x\in U$ we have $U$ intersects $A$.
        Show that for all $C$ such that $C$ is closed in $X$ with respect to $T$ and
            for all $y\in A$ we have $y\in C$ we have $x\in C$.
        \begin{subproof}
            Fix $C$ such that $C$ is closed in $X$ with respect to $T$ and
                for all $y\in A$ we have $y\in C$.
            Suppose not.
            Let $U = X\setminus C$.
            Hence $U$ is open in $X$ with respect to $T$ and $x\in U$ by \cref{closed_in,setminus_intro}.
            For all $y$ such that $y\in U\inter A$ we have $y\in C$ and $y\notin C$
                by \cref{inter,setminus_elim_right}.
            Hence $U$ does not intersect $A$ by \cref{intersects,nonempty_inhabited}.
            Contradiction.
        \end{subproof}
        Hence $x\in\closure{A}{X}{T}$ by \cref{closure}.
    \end{subproof}
\end{proof}

% from §17 Item 19: closure and basis elements
\begin{theorem}\label{closure_iff_intersects_basis}
    Assume $T$ is a topology on $X$.
    Suppose $B$ is a basis for $T$ on $X$.
    Suppose $A\subseteq X$.
    Suppose $x\in X$.
    Then $x\in\closure{A}{X}{T}$ iff for all $B_1$ such that $B_1\in B$ and $x\in B_1$ we have $B_1$ intersects $A$.
\end{theorem}
\begin{proof}
    If $x\in\closure{A}{X}{T}$ then for all $B_1$ such that $B_1\in B$ and $x\in B_1$ we have $B_1$ intersects $A$
        by \cref{basis_for_topology,basis_elem_open_in_generated,open_in,closure_iff_intersects_open}.
    Show that if for all $B_1$ such that $B_1\in B$ and $x\in B_1$ we have $B_1$ intersects $A$ then
        $x\in\closure{A}{X}{T}$.
    \begin{subproof}
        Assume for all $B_1$ such that $B_1\in B$ and $x\in B_1$ we have $B_1$ intersects $A$.
        For all $U$ such that $U$ is open in $X$ with respect to $T$ and $x\in U$ we have $U$ intersects $A$
            by \cref{open_in,basis_for_topology,topology_generated_by_basis,intersects,nonempty_inhabited,inter,elem_subseteq,inter_intro,emptyset}.
        Hence $x\in\closure{A}{X}{T}$ by \cref{closure_iff_intersects_open}.
    \end{subproof}
\end{proof}

% from §17 Item 20: neighborhood
\begin{definition}\label{neighborhood}
    $U$ is a neighborhood of $x$ in $X$ with respect to $T$ iff
        $U$ is open in $X$ with respect to $T$ and $x\in U$.
\end{definition}

% from §17 Item 21: limit point
\begin{definition}\label{limit_point}
    $x$ is a limit point of $A$ in $X$ with respect to $T$ iff
        for all $U$ such that $U$ is a neighborhood of $x$ in $X$ with respect to $T$
        we have $U$ intersects $A\setminus\{x\}$.
\end{definition}

% from §17 Item 22: set of limit points
\begin{definition}\label{limit_points}
    $\limitPoints{A}{X}{T} = \{x\in X \mid \text{$x$ is a limit point of $A$ in $X$ with respect to $T$}\}$.
\end{definition}
% from §17 Item 23: closure equals set plus limit points
\begin{theorem}\label{closure_union_limit_points}
    Assume $T$ is a topology on $X$.
    Suppose $A\subseteq X$.
    Then $\closure{A}{X}{T} = A\union\limitPoints{A}{X}{T}$.
\end{theorem}
\begin{proof}
    Show that for all $x$ such that $x\in A\union\limitPoints{A}{X}{T}$ we have $x\in\closure{A}{X}{T}$.
    \begin{subproof}
        Fix $x$ such that $x\in A\union\limitPoints{A}{X}{T}$.
        Then $x\in A$ or $x\in\limitPoints{A}{X}{T}$ by \cref{union_iff}.
        \begin{byCase}
        \caseOf{$x\in A$.}
            Hence $x\in\closure{A}{X}{T}$ by \cref{subseteq_closure,elem_subseteq}.
        \caseOf{$x\in\limitPoints{A}{X}{T}$.}
            Hence $x\in\closure{A}{X}{T}$
                by \cref{neighborhood,limit_points,limit_point,intersects,nonempty_inhabited,inter,setminus_elim_left,inter_intro,emptyset,closure_iff_intersects_open}.
        \end{byCase}
    \end{subproof}
    Therefore $A\union\limitPoints{A}{X}{T}\subseteq\closure{A}{X}{T}$ by \cref{subseteq}.
    Show that for all $x$ such that $x\in\closure{A}{X}{T}$ we have $x\in A\union\limitPoints{A}{X}{T}$.
    \begin{subproof}
        Fix $x$ such that $x\in\closure{A}{X}{T}$.
        \begin{byCase}
        \caseOf{$x\in A$.}
            Hence $x\in A\union\limitPoints{A}{X}{T}$ by \cref{union_intro_left}.
        \caseOf{$x\notin A$.}
            Show that $x\in\limitPoints{A}{X}{T}$.
            \begin{subproof}
                Show that for all $U$ such that $U$ is a neighborhood of $x$ in $X$ with respect to $T$ we have
                    $U$ intersects $A\setminus\{x\}$.
                \begin{subproof}
                Fix $U$ such that $U$ is a neighborhood of $x$ in $X$ with respect to $T$.
                    Hence $U$ intersects $A$ by \cref{neighborhood,closure,closure_iff_intersects_open}.
                    Take $y$ such that $y\in U\inter A$ by \cref{intersects,nonempty_inhabited}.
                    Thus $y\in U$ and $y\in A$ by \cref{inter}.
                    $y\neq x$.
                    Therefore $U$ intersects $A\setminus\{x\}$ by \cref{singleton_iff,setminus_intro,inter_intro,emptyset,intersects}.
                \end{subproof}
                Hence $x\in\limitPoints{A}{X}{T}$ by \cref{closure,limit_point,limit_points}.
            \end{subproof}
            Therefore $x\in A\union\limitPoints{A}{X}{T}$ by \cref{union_intro_right}.
        \end{byCase}
    \end{subproof}
    Hence $\closure{A}{X}{T}\subseteq A\union\limitPoints{A}{X}{T}$ by \cref{subseteq}.
    Therefore $\closure{A}{X}{T} = A\union\limitPoints{A}{X}{T}$ by \cref{subseteq_antisymmetric}.
\end{proof}

% from §17 Item 24: closed iff contains limit points
\begin{corollary}\label{closed_iff_contains_limit_points}
    Assume $T$ is a topology on $X$.
    Suppose $A\subseteq X$.
    Then $A$ is closed in $X$ with respect to $T$ iff $\limitPoints{A}{X}{T}\subseteq A$.
\end{corollary}
\begin{proof}
    If $A$ is closed in $X$ with respect to $T$ then $\limitPoints{A}{X}{T}\subseteq A$ by
        \cref{closed_eq_closure,closure_union_limit_points,union_comm,union_eq_self_implies_subseteq}.
    If $\limitPoints{A}{X}{T}\subseteq A$ then $A$ is closed in $X$ with respect to $T$ by
        \cref{union_absorb_subseteq_right,closure_union_limit_points,closure_closed_in}.
\end{proof}

% from §17 Item 25: Hausdorff space
\begin{definition}\label{hausdorff}
    $X$ is Hausdorff with respect to $T$ iff
        $T$ is a topology on $X$ and
        for all $x_1,x_2$ such that $x_1\in X$ and $x_2\in X$ and $x_1\neq x_2$
        there exist $U_1,U_2$ such that
            $U_1$ is a neighborhood of $x_1$ in $X$ with respect to $T$ and
            $U_2$ is a neighborhood of $x_2$ in $X$ with respect to $T$ and
            $U_1$ is disjoint from $U_2$.
\end{definition}

% from §17 helper lemma 3: finite preserved by bijection
\begin{lemma}\label{finite_bijection}
    Suppose $A$ is finite.
    Suppose $f$ is a bijection from $A$ to $B$.
    Then $B$ is finite.
\end{lemma}
\begin{proof}
    Follows by \cref{finite,bijection_circ}.
\end{proof}

% from §17 helper lemma 4: bijection to singleton family
\begin{lemma}\label{singleton_family_bijection}
    Let $S = \{\{x\} \mid x\in A\}$.
    Then there exists a bijection from $A$ to $S$.
\end{lemma}
\begin{proof}
    Let $f = \{(x,\{x\}) \mid x\in A\}$.
    For all $w$ such that $w\in f$ there exist $x,y$ such that $w = (x,y)$.
    For all $x,y,z$ such that $(x,y)\in f$ and $(x,z)\in f$ we have $y = z$ by \cref{pair_eq_iff}.
    Hence $f$ is a function by \cref{relation,rightunique}.
    Hence $f\in\funs{A}{S}$ by
        \cref{dom_iff,pair_eq_iff,subseteq,subseteq_antisymmetric,function_apply_iff,funs_intro}.
    Therefore $f\in\surj{A}{S}$ by \cref{ran_iff,pair_eq_iff,subseteq,subseteq_antisymmetric,funs_surj_iff}.
    Therefore $f\in\bijections{A}{S}$ by \cref{pair_eq_iff,singleton_elim,singleton_intro,injective,bijections}.
    Hence there exists a bijection from $A$ to $S$.
\end{proof}

% from §17 Item 26: finite point sets are closed in Hausdorff spaces
\begin{theorem}\label{finite_point_set_closed_hausdorff}
    Assume $X$ is Hausdorff with respect to $T$.
    Suppose $A\subseteq X$.
    Suppose $A$ is finite.
    Then $A$ is closed in $X$ with respect to $T$.
\end{theorem}
\begin{proof}
    $T$ is a topology on $X$ by \cref{hausdorff}.
    Show that for all $x$ such that $x\in X$ we have $\{x\}$ is closed in $X$ with respect to $T$.
    \begin{subproof}
        Fix $x$ such that $x\in X$.
        Show that for all $y$ such that $y\in\limitPoints{\{x\}}{X}{T}$ we have $y\in\{x\}$.
        \begin{subproof}
            Fix $y$ such that $y\in\limitPoints{\{x\}}{X}{T}$.
            Hence $y\in X$ and $y$ is a limit point of $\{x\}$ in $X$ with respect to $T$ by \cref{limit_points}.
            Suppose not.
            Take $U_1,U_2$ such that
                $U_1$ is a neighborhood of $y$ in $X$ with respect to $T$ and
                $U_2$ is a neighborhood of $x$ in $X$ with respect to $T$ and
                $U_1$ is disjoint from $U_2$ by \cref{singleton_iff,hausdorff}.
            For all $z$ such that $z\in U_1\inter (\{x\}\setminus\{y\})$ we have $x\in U_1$ and $x\in U_2$
                by \cref{inter,setminus_elim_left,singleton_elim,neighborhood}.
            Hence $U_1$ does not intersect $\{x\}\setminus\{y\}$ by \cref{intersects,nonempty_inhabited,inter_intro,emptyset,intersects,disjoint}.
            Hence $y$ is not a limit point of $\{x\}$ in $X$ with respect to $T$ by \cref{limit_point}.
            Contradiction.
        \end{subproof}
        Therefore $\{x\}$ is closed in $X$ with respect to $T$ by \cref{subseteq,singleton_subset_intro,closed_iff_contains_limit_points}.
    \end{subproof}
    Let $F = \{\{x\} \mid x\in A\}$.
    Hence $F$ is finite and $F\subseteq\pow{X}$ by
        \cref{singleton_family_bijection,finite_bijection,elem_subseteq,singleton_subset_intro,powerset_intro,subseteq}.
    For all $B$ such that $B\in F$ we have $B$ is closed in $X$ with respect to $T$.
    Hence $A$ is closed in $X$ with respect to $T$ by
        \cref{closed_finite_unions,singleton_intro,unions_iff,subseteq,singleton_elim,subseteq_antisymmetric}.
\end{proof}

% from §17 Item 27: T1 axiom
\begin{definition}\label{t1_axiom}
    $X$ satisfies the T-one axiom with respect to $T$ iff
        $T$ is a topology on $X$ and
        for all $A$ such that $A\subseteq X$ and $A$ is finite
        we have $A$ is closed in $X$ with respect to $T$.
\end{definition}

% from §17 Item 28: limit points in $T_1$ spaces
\begin{theorem}\label{limit_point_infinite_neighborhood}
    Assume $X$ satisfies the T-one axiom with respect to $T$.
    Suppose $A\subseteq X$.
    Suppose $x\in X$.
    Then $x$ is a limit point of $A$ in $X$ with respect to $T$ iff
        for all $U$ such that $U$ is a neighborhood of $x$ in $X$ with respect to $T$ we have $U\inter A$ is infinite.
\end{theorem}
\begin{proof}
    Show that if $x$ is a limit point of $A$ in $X$ with respect to $T$ then
        for all $U$ such that $U$ is a neighborhood of $x$ in $X$ with respect to $T$ we have $U\inter A$ is infinite.
    \begin{subproof}
        Assume $x$ is a limit point of $A$ in $X$ with respect to $T$.
        Show that for all $U$ such that $U$ is a neighborhood of $x$ in $X$ with respect to $T$ we have $U\inter A$ is infinite.
        \begin{subproof}
            Fix $U$ such that $U$ is a neighborhood of $x$ in $X$ with respect to $T$.
            Suppose not.
            Then $U\inter A$ is finite.
            Let $B = U\inter (A\setminus\{x\})$.
            Hence $B\subseteq U\inter A$ by \cref{inter,setminus_elim_left,subseteq,subseteq_inter_iff}.
            Therefore $B$ is finite by \cref{subseteq_finite_is_finite}.
            Hence $B\subseteq X$ and $T$ is a topology on $X$ and $B$ is closed in $X$ with respect to $T$ by
                \cref{inter,subseteq,subseteq_transitive,t1_axiom}.
            Therefore $U\inter (X\setminus B)$ is open in $X$ with respect to $T$ by
                \cref{neighborhood,open_in,closed_in,topology_on,t1_axiom}.
            $x\in U$ by \cref{neighborhood}.
            $x\notin B$.
            Hence $x\in X\setminus B$ by \cref{setminus_intro}.
            Therefore $x\in U\inter (X\setminus B)$ by \cref{inter_intro}.
            Hence $U\inter (X\setminus B)$ is a neighborhood of $x$ in $X$ with respect to $T$ by \cref{neighborhood}.
            For all $y$ such that $y\in (U\inter (X\setminus B)) \inter (A\setminus\{x\})$ we have
                $y\in B$ and $y\notin B$ by \cref{inter,inter_intro,setminus_elim_right}.
            Hence $U\inter (X\setminus B)$ does not intersect $A\setminus\{x\}$ by \cref{intersects,nonempty_inhabited}.
            Contradiction by \cref{limit_point}.
        \end{subproof}
    \end{subproof}
    Show that if for all $U$ such that $U$ is a neighborhood of $x$ in $X$ with respect to $T$ we have $U\inter A$ is infinite then
        $x$ is a limit point of $A$ in $X$ with respect to $T$.
    \begin{subproof}
        Assume for all $U$ such that $U$ is a neighborhood of $x$ in $X$ with respect to $T$ we have $U\inter A$ is infinite.
        Show that for all $U$ such that $U$ is a neighborhood of $x$ in $X$ with respect to $T$ we have
            $U$ intersects $A\setminus\{x\}$.
        \begin{subproof}
            Fix $U$ such that $U$ is a neighborhood of $x$ in $X$ with respect to $T$.
            Suppose not.
            Then $U\inter (A\setminus\{x\}) = \emptyset$.
            For all $y$ such that $y\in U\inter A$ we have $y\in\{x\}$.
            Therefore $U\inter A$ is finite by \cref{subseteq,upair_iff,singleton_subset_intro,pair_is_finite,subseteq_finite_is_finite}.
            Contradiction.
        \end{subproof}
        Therefore $x$ is a limit point of $A$ in $X$ with respect to $T$ by \cref{limit_point}.
    \end{subproof}
\end{proof}

% from §17 Item 29: sequence in a set
\begin{definition}\label{sequence_in}
    $s$ is a sequence in $X$ iff $s\in\funs{\naturalsPlus}{X}$.
\end{definition}

% from §17 Item 30: convergence of a sequence
\begin{definition}\label{converges_to}
    $s$ converges to $x$ in $X$ with respect to $T$ iff
        $s$ is a sequence in $X$ and
        for all $U$ such that $U$ is a neighborhood of $x$ in $X$ with respect to $T$
        there exists $N\in\naturalsPlus$ such that for all $n$ such that $n\in\naturalsPlus$ if $N\leq n$ then $s(n)\in U$.
\end{definition}

% from §17 Item 31: limits of sequences are unique in Hausdorff spaces
\begin{theorem}\label{sequence_limit_unique}
    Assume $X$ is Hausdorff with respect to $T$.
    Suppose $s$ is a sequence in $X$.
    Suppose $x\in X$.
    Suppose $y\in X$.
    Suppose $s$ converges to $x$ in $X$ with respect to $T$.
    Suppose $s$ converges to $y$ in $X$ with respect to $T$.
    Then $x = y$.
\end{theorem}
\begin{proof}
    Suppose not.
    Then $x\neq y$.
    Take $U_1,U_2$ such that
        $U_1$ is a neighborhood of $x$ in $X$ with respect to $T$ and
        $U_2$ is a neighborhood of $y$ in $X$ with respect to $T$ and
        $U_1$ is disjoint from $U_2$ by \cref{hausdorff}.
    Take $N_1\in\naturalsPlus$ such that for all $n$ such that $n\in\naturalsPlus$ if $N_1\leq n$ then $s(n)\in U_1$ by \cref{converges_to}.
    Take $N_2\in\naturalsPlus$ such that for all $n$ such that $n\in\naturalsPlus$ if $N_2\leq n$ then $s(n)\in U_2$ by \cref{converges_to}.
    Let $A = \{N_1,N_2\}$.
    Take $m\in A$ such that for all $a\in A$ we have $a\leq m$ by
        \cref{upair_iff,subseteq,pair_is_finite,emptyset,finite_naturals_maximum}.
    Hence $N_1\leq m$ and $N_2\leq m$ by \cref{upair_iff}.
    Therefore $s(m)\in U_1$ and $s(m)\in U_2$.
    Hence $U_1$ intersects $U_2$ by \cref{inter_intro,emptyset,intersects}.
    Contradiction by \cref{disjoint}.
\end{proof}

% from §17 helper lemma 5: open right ray is open in order topology
\begin{lemma}\label{openrayright_open_in_order_topology}
    Assume $R$ is a simple order relation on $X$.
    Assume $X$ has more than one element.
    Assume $T$ is the order topology on $X$ with respect to $R$.
    Suppose $a\in X$.
    Then $\openrayright{R}{X}{a}$ is open in $X$ with respect to $T$.
\end{lemma}
\begin{proof}
    Let $S = \openRays{R}{X}$.
    $S$ is a subbasis on $X$ and $\unions{S} = X$ by \cref{open_rays_subbasis_on,subbasis_on}.
    Let $U = \openrayright{R}{X}{a}$.
    Hence $U\in S$ and $U\in\pow{X}$ by \cref{openray_right,subseteq,powerset_intro,open_rays}.
    Let $F = \{U\}$.
    $F$ is finite and $F$ is inhabited.
    Therefore $U\in\finiteIntersections{S}$ and $U\in T$
        by \cref{inters_singleton,singleton_subset_intro,powerset_elim,subseteq_transitive,powerset_intro,finite_intersections,subbasis_finite_intersections_basis,basis_elem_open_in_generated,topology_generated_by_subbasis,open_rays_subbasis}.
    Hence $U$ is open in $X$ with respect to $T$ by
        \cref{order_basis_is_basis,basis_topology_is_topology,order_topology,open_in}.
\end{proof}

% from §17 helper lemma 6: open left ray is open in order topology
\begin{lemma}\label{openrayleft_open_in_order_topology}
    Assume $R$ is a simple order relation on $X$.
    Assume $X$ has more than one element.
    Assume $T$ is the order topology on $X$ with respect to $R$.
    Suppose $a\in X$.
    Then $\openrayleft{R}{X}{a}$ is open in $X$ with respect to $T$.
\end{lemma}
\begin{proof}
    Let $S = \openRays{R}{X}$.
    $S$ is a subbasis on $X$ and $\unions{S} = X$ by \cref{open_rays_subbasis_on,subbasis_on}.
    Let $U = \openrayleft{R}{X}{a}$.
    Hence $U\in S$ and $U\in\pow{X}$ by \cref{openray_left,subseteq,powerset_intro,open_rays}.
    Let $F = \{U\}$.
    $F$ is finite and $F$ is inhabited.
    Therefore $U\in\finiteIntersections{S}$ and $U\in T$
        by \cref{inters_singleton,singleton_subset_intro,powerset_elim,subseteq_transitive,powerset_intro,finite_intersections,subbasis_finite_intersections_basis,basis_elem_open_in_generated,topology_generated_by_subbasis,open_rays_subbasis}.
    Hence $U$ is open in $X$ with respect to $T$ by
        \cref{order_basis_is_basis,basis_topology_is_topology,order_topology,open_in}.
\end{proof}

% from §17 Item 32: order topology is Hausdorff
\begin{theorem}\label{order_topology_hausdorff}
    Assume $R$ is a simple order relation on $X$.
    Assume $X$ has more than one element.
    Assume $T$ is the order topology on $X$ with respect to $R$.
    Then $X$ is Hausdorff with respect to $T$.
\end{theorem}
\begin{proof}
    Show that for all $x_1,x_2$ such that $x_1\in X$ and $x_2\in X$ and $x_1\neq x_2$ there exist $U_1,U_2$ such that
        $U_1$ is a neighborhood of $x_1$ in $X$ with respect to $T$ and
        $U_2$ is a neighborhood of $x_2$ in $X$ with respect to $T$ and
        $U_1$ is disjoint from $U_2$.
    \begin{subproof}
        Fix $x_1$.
        Fix $x_2$ such that $x_1\in X$ and $x_2\in X$ and $x_1\neq x_2$.
        Then $x_1\mathrel{R}x_2$ or $x_2\mathrel{R}x_1$ by \cref{simple_order_relation,connex}.
        \begin{byCase}
        \caseOf{$x_1\mathrel{R}x_2$.}
            \begin{byCase}
            \caseOf{there exists $z\in X$ such that $x_1\mathrel{R}z$ and $z\mathrel{R}x_2$.}
                Take $z\in X$ such that $x_1\mathrel{R}z$ and $z\mathrel{R}x_2$.
                Let $U_1 = \openrayleft{R}{X}{z}$.
                Let $U_2 = \openrayright{R}{X}{z}$.
                Hence $U_1$ is a neighborhood of $x_1$ in $X$ with respect to $T$ and
                    $U_2$ is a neighborhood of $x_2$ in $X$ with respect to $T$
                    by \cref{openrayleft_open_in_order_topology,openrayright_open_in_order_topology,openray_left,openray_right,neighborhood}.
                $U_1$ is disjoint from $U_2$ by
                    \cref{disjoint,openray_left,openray_right,simple_order_relation,asymmetric}.
            \caseOf{there exists no $z\in X$ such that $x_1\mathrel{R}z$ and $z\mathrel{R}x_2$.}
                Let $U_1 = \openrayleft{R}{X}{x_2}$.
                Let $U_2 = \openrayright{R}{X}{x_1}$.
                Therefore $U_1$ is a neighborhood of $x_1$ in $X$ with respect to $T$ and
                    $U_2$ is a neighborhood of $x_2$ in $X$ with respect to $T$
                    by \cref{openrayleft_open_in_order_topology,openrayright_open_in_order_topology,openray_left,openray_right,neighborhood}.
                $U_1$ is disjoint from $U_2$ by \cref{disjoint,openray_right,openray_left}.
            \end{byCase}
        \caseOf{$x_2\mathrel{R}x_1$.}
            \begin{byCase}
            \caseOf{there exists $z\in X$ such that $x_2\mathrel{R}z$ and $z\mathrel{R}x_1$.}
                Take $z\in X$ such that $x_2\mathrel{R}z$ and $z\mathrel{R}x_1$.
                Let $U_2 = \openrayleft{R}{X}{z}$.
                Let $U_1 = \openrayright{R}{X}{z}$.
                Hence $U_1$ is a neighborhood of $x_1$ in $X$ with respect to $T$ and
                    $U_2$ is a neighborhood of $x_2$ in $X$ with respect to $T$
                    by \cref{openrayright_open_in_order_topology,openrayleft_open_in_order_topology,openray_left,openray_right,neighborhood}.
                $U_1$ is disjoint from $U_2$ by
                    \cref{disjoint,openray_right,openray_left,simple_order_relation,asymmetric}.
            \caseOf{there exists no $z\in X$ such that $x_2\mathrel{R}z$ and $z\mathrel{R}x_1$.}
                Let $U_2 = \openrayleft{R}{X}{x_1}$.
                Let $U_1 = \openrayright{R}{X}{x_2}$.
                Therefore $U_1$ is a neighborhood of $x_1$ in $X$ with respect to $T$ and
                    $U_2$ is a neighborhood of $x_2$ in $X$ with respect to $T$
                    by \cref{openrayright_open_in_order_topology,openrayleft_open_in_order_topology,openray_left,openray_right,neighborhood}.
                $U_1$ is disjoint from $U_2$ by \cref{disjoint,openray_right,openray_left}.
            \end{byCase}
        \end{byCase}
    \end{subproof}
    Hence $X$ is Hausdorff with respect to $T$ by \cref{order_basis_is_basis,basis_topology_is_topology,order_topology,hausdorff}.
\end{proof}

% from §17 Item 33: product of Hausdorff spaces is Hausdorff
\begin{theorem}\label{product_hausdorff}
    Assume $X$ is Hausdorff with respect to $T$.
    Assume $Y$ is Hausdorff with respect to $T_1$.
    Let $T_2 = \productTopology{T}{T_1}{X}{Y}$.
    Then $X\times Y$ is Hausdorff with respect to $T_2$.
\end{theorem}
\begin{proof}
    $T$ is a topology on $X$ and $T_1$ is a topology on $Y$ by \cref{hausdorff}.
    $\productBasis{T}{T_1}{X}{Y}$ is a basis on $X\times Y$ by \cref{product_basis_is_basis}.
    Hence $T_2$ is a topology on $X\times Y$ by \cref{product_basis_is_basis,basis_topology_is_topology,product_topology}.
    Show that for all $p_1,p_2$ such that $p_1\in X\times Y$ and $p_2\in X\times Y$ and $p_1\neq p_2$
        there exist $U_1,U_2$ such that
            $U_1$ is a neighborhood of $p_1$ in $X\times Y$ with respect to $T_2$ and
            $U_2$ is a neighborhood of $p_2$ in $X\times Y$ with respect to $T_2$ and
            $U_1$ is disjoint from $U_2$.
    \begin{subproof}
        Fix $p_1$.
        Fix $p_2$ such that $p_1\in X\times Y$ and $p_2\in X\times Y$ and $p_1\neq p_2$.
        Take $x_1,y_1$ such that $x_1\in X$ and $y_1\in Y$ and $p_1 = (x_1,y_1)$ by \cref{times_elem_is_tuple}.
        Take $x_2,y_2$ such that $x_2\in X$ and $y_2\in Y$ and $p_2 = (x_2,y_2)$ by \cref{times_elem_is_tuple}.
        $x_1\neq x_2$ or $y_1\neq y_2$ by \cref{pair_eq_iff}.
        \begin{byCase}
        \caseOf{$x_1\neq x_2$.}
            Take $U_1,U_2$ such that
                $U_1$ is a neighborhood of $x_1$ in $X$ with respect to $T$ and
                $U_2$ is a neighborhood of $x_2$ in $X$ with respect to $T$ and
                $U_1$ is disjoint from $U_2$ by \cref{hausdorff}.
            Let $V_1 = Y$.
            Let $V_2 = Y$.
            Hence $V_1$ is open in $Y$ with respect to $T_1$ and
                $V_2$ is open in $Y$ with respect to $T_1$ by \cref{topology_on,open_in}.
            Let $W_1 = U_1\times V_1$.
            Let $W_2 = U_2\times V_2$.
            Hence $W_1\subseteq X\times Y$ and $W_2\subseteq X\times Y$ by
                \cref{open_in,neighborhood,topology_on,subseteq,pow_iff,times_subseteq_left,times_subseteq_right,subseteq_transitive}.
            Therefore $W_1\in\productBasis{T}{T_1}{X}{Y}$ and $W_2\in\productBasis{T}{T_1}{X}{Y}$
                by \cref{open_in,neighborhood,powerset_intro,product_basis}.
            Hence $W_1\in T_2$ and $W_2\in T_2$ by \cref{basis_elem_open_in_generated,product_topology}.
            Hence $W_1$ is a neighborhood of $p_1$ in $X\times Y$ with respect to $T_2$ and
                $W_2$ is a neighborhood of $p_2$ in $X\times Y$ with respect to $T_2$ by
                \cref{open_in,neighborhood,times_tuple_intro}.
            $W_1$ is disjoint from $W_2$ by \cref{disjoint,times_elem_is_tuple,times_tuple_elim}.
        \caseOf{$y_1\neq y_2$.}
            Take $V_1,V_2$ such that
                $V_1$ is a neighborhood of $y_1$ in $Y$ with respect to $T_1$ and
                $V_2$ is a neighborhood of $y_2$ in $Y$ with respect to $T_1$ and
                $V_1$ is disjoint from $V_2$ by \cref{hausdorff}.
            Let $U_1 = X$.
            Let $U_2 = X$.
            Hence $U_1$ is open in $X$ with respect to $T$ and
                $U_2$ is open in $X$ with respect to $T$ by \cref{topology_on,open_in}.
            Let $W_1 = U_1\times V_1$.
            Let $W_2 = U_2\times V_2$.
            Hence $W_1\subseteq X\times Y$ and $W_2\subseteq X\times Y$ by
                \cref{open_in,neighborhood,topology_on,subseteq,pow_iff,times_subseteq_left,times_subseteq_right,subseteq_transitive}.
            Therefore $W_1\in\productBasis{T}{T_1}{X}{Y}$ and $W_2\in\productBasis{T}{T_1}{X}{Y}$
                by \cref{open_in,neighborhood,powerset_intro,product_basis}.
            Hence $W_1\in T_2$ and $W_2\in T_2$ by \cref{basis_elem_open_in_generated,product_topology}.
            Hence $W_1$ is a neighborhood of $p_1$ in $X\times Y$ with respect to $T_2$ and
                $W_2$ is a neighborhood of $p_2$ in $X\times Y$ with respect to $T_2$ by
                \cref{open_in,neighborhood,times_tuple_intro}.
            $W_1$ is disjoint from $W_2$ by \cref{disjoint,times_elem_is_tuple,times_tuple_elim}.
        \end{byCase}
    \end{subproof}
    Therefore $X\times Y$ is Hausdorff with respect to $T_2$ by \cref{hausdorff}.
\end{proof}

% from §17 Item 34: subspace of a Hausdorff space is Hausdorff
\begin{theorem}\label{subspace_hausdorff}
    Assume $X$ is Hausdorff with respect to $T$.
    Suppose $Y\subseteq X$.
    Suppose $T_1 = \subspaceTopology{T}{Y}$.
    Then $Y$ is Hausdorff with respect to $T_1$.
\end{theorem}
\begin{proof}
    $T_1$ is a topology on $Y$ by \cref{hausdorff,subspace_topology_is_topology}.
    Show that for all $y_1,y_2$ such that $y_1\in Y$ and $y_2\in Y$ and $y_1\neq y_2$
        there exist $V_1,V_2$ such that
            $V_1$ is a neighborhood of $y_1$ in $Y$ with respect to $T_1$ and
            $V_2$ is a neighborhood of $y_2$ in $Y$ with respect to $T_1$ and
            $V_1$ is disjoint from $V_2$.
    \begin{subproof}
        Fix $y_1$.
        Fix $y_2$ such that $y_1\in Y$ and $y_2\in Y$ and $y_1\neq y_2$.
        Take $U_1,U_2$ such that
            $U_1$ is a neighborhood of $y_1$ in $X$ with respect to $T$ and
            $U_2$ is a neighborhood of $y_2$ in $X$ with respect to $T$ and
            $U_1$ is disjoint from $U_2$ by \cref{elem_subseteq,hausdorff}.
        Let $V_1 = Y\inter U_1$.
        Let $V_2 = Y\inter U_2$.
        Hence $V_1$ is a neighborhood of $y_1$ in $Y$ with respect to $T_1$ and
            $V_2$ is a neighborhood of $y_2$ in $Y$ with respect to $T_1$ by
            \cref{neighborhood,open_in,subspace_topology,inter_intro}.
        $V_1$ is disjoint from $V_2$ by \cref{disjoint,inter}.
    \end{subproof}
    Therefore $Y$ is Hausdorff with respect to $T_1$ by \cref{hausdorff}.
\end{proof}

\section{§18 Continuous Functions}

% from §18 Item 1: continuous function
\begin{definition}\label{continuous}
    $f$ is continuous from $X$ to $Y$ with respect to $T$ and $T_1$ iff
        $f$ is a function from $X$ to $Y$ and
        $T$ is a topology on $X$ and
        $T_1$ is a topology on $Y$ and
        for all $V$ such that $V$ is open in $Y$ with respect to $T_1$
        we have $\preimg{f}{V}$ is open in $X$ with respect to $T$.
\end{definition}

% from §18 helper lemma 1: preimage of an intersection under a function
\begin{lemma}\label{preimg_inter_function}
    Suppose $f$ is a function.
    Then $\preimg{f}{A\inter B} = \preimg{f}{A}\inter \preimg{f}{B}$.
\end{lemma}
\begin{proof}
    Follows by \cref{preimg_inter,inter,preimg_iff,function_rightunique,inter_intro,subseteq,subseteq_antisymmetric}.
\end{proof}

% from §18 helper lemma 2: image of a preimage is contained in the set
\begin{lemma}\label{img_preimg_subseteq}
    Suppose $f$ is a function.
    Then $\img{f}{\preimg{f}{B}}\subseteq B$.
\end{lemma}
\begin{proof}
    Follows by \cref{img_of_function_elim,inter_elim_right,preimg_iff,function_apply_elim,inter_elim_left,function_rightunique,subseteq}.
\end{proof}

% from §18 helper lemma 3: preimage of a complement
\begin{lemma}\label{preimg_setminus_function}
    Assume $f$ is a function from $X$ to $Y$.
    Suppose $B\subseteq Y$.
    Then $\preimg{f}{Y\setminus B} = X\setminus \preimg{f}{B}$.
\end{lemma}
\begin{proof}
    Hence $\preimg{f}{Y\setminus B}\subseteq X\setminus \preimg{f}{B}$ by
        \cref{preimg_iff,funs,funs_is_function,subseteq,preimg_subseteq_dom,function_rightunique,setminus_elim_right,setminus_intro}.
    For all $x$ such that $x\in X\setminus \preimg{f}{B}$ there exists $y\in Y\setminus B$ such that
        $x\mathrel{f} y$.
    For all $x$ such that $x\in X\setminus \preimg{f}{B}$ we have
        $x\in\preimg{f}{Y\setminus B}$ by \cref{preimg_iff}.
    Therefore $X\setminus \preimg{f}{B}\subseteq \preimg{f}{Y\setminus B}$ by \cref{subseteq}.
    Hence $\preimg{f}{Y\setminus B} = X\setminus \preimg{f}{B}$ by \cref{subseteq_antisymmetric}.
\end{proof}

% from §18 Item 2: continuity implies image of closure is contained in closure of image
\begin{lemma}\label{continuous_closure_image_subset}
    Assume $f$ is continuous from $X$ to $Y$ with respect to $T$ and $T_1$.
    Suppose $A\subseteq X$.
    Then $\img{f}{\closure{A}{X}{T}}\subseteq \closure{\img{f}{A}}{Y}{T_1}$.
\end{lemma}
\begin{proof}
    Hence $f\in\funs{X}{Y}$ and $f$ is a function and
        $T$ is a topology on $X$ and $T_1$ is a topology on $Y$
        by \cref{continuous,funs_is_function}.
    Hence $f\in\rels{X}{Y}$ by \cref{funs_subseteq_rels,subseteq}.
    Therefore $\img{f}{A}\subseteq Y$ by
        \cref{funs_subseteq_rels,subseteq,rels_ran_subseteq,img_subseteq_ran,subseteq_transitive}.
    Show that for all $y$ such that $y\in\img{f}{\closure{A}{X}{T}}$ we have
        $y\in\closure{\img{f}{A}}{Y}{T_1}$.
        \begin{subproof}
            Fix $y$ such that $y\in\img{f}{\closure{A}{X}{T}}$.
            Take $x\in\dom{f}\inter \closure{A}{X}{T}$ such that $y = f(x)$ by \cref{img_of_function_elim}.
            Hence $x\in\dom{f}$ and $x\mathrel{f} y$ and $y\in Y$
                by \cref{inter_elim_left,function_apply_elim,rels_type_ran}.
            Therefore $x\in\closure{A}{X}{T}$ and $x\in X$ by \cref{inter_elim_right,closure}.
        Show that for all $V$ such that $V$ is open in $Y$ with respect to $T_1$ and $y\in V$
            we have $V$ intersects $\img{f}{A}$.
        \begin{subproof}
            Fix $V$ such that $V$ is open in $Y$ with respect to $T_1$ and $y\in V$.
            Let $U = \preimg{f}{V}$.
            Hence $U$ is open in $X$ with respect to $T$ and $x\in U$ by
                \cref{continuous,function_apply_elim,inter_elim_left,preimg_iff}.
            Take $z$ such that $z\in U\inter A$ by \cref{closure_iff_intersects_open,intersects,nonempty_inhabited}.
            Take $w\in V$ such that $z\mathrel{f} w$ by \cref{inter,preimg_iff}.
            Hence $V$ intersects $\img{f}{A}$ by \cref{inter,function_apply_iff,inter_intro,img_of_function_intro,emptyset,intersects}.
        \end{subproof}
        Therefore $y\in\closure{\img{f}{A}}{Y}{T_1}$ by \cref{closure_iff_intersects_open}.
    \end{subproof}
    Hence $\img{f}{\closure{A}{X}{T}}\subseteq \closure{\img{f}{A}}{Y}{T_1}$ by \cref{subseteq}.
\end{proof}

% from §18 Item 3: closure condition implies preimages of closed sets are closed
\begin{lemma}\label{closure_image_subset_implies_preimg_closed}
    Assume $f$ is a function from $X$ to $Y$ and
        $T$ is a topology on $X$ and
        $T_1$ is a topology on $Y$.
    Suppose for all $A$ such that $A\subseteq X$ we have
        $\img{f}{\closure{A}{X}{T}}\subseteq \closure{\img{f}{A}}{Y}{T_1}$.
    Then for all $B$ such that $B$ is closed in $Y$ with respect to $T_1$
        we have $\preimg{f}{B}$ is closed in $X$ with respect to $T$.
\end{lemma}
\begin{proof}
    Fix $B$ such that $B$ is closed in $Y$ with respect to $T_1$.
    Let $A = \preimg{f}{B}$.
    Hence $f$ is a function and $f$ is right-unique and $f$ is a relation and $\dom{f} = X$
        by \cref{funs,funs_is_function,function_rightunique,funs_is_relation}.
    Hence $\ran{f}\subseteq Y$ by \cref{funs_subseteq_rels,subseteq,rels_ran_subseteq}.
    $A\subseteq X$ by \cref{preimg_subseteq_dom}.
    For all $x$ such that $x\in\closure{A}{X}{T}$ we have $x\in A$ by
        \cref{closure,funs,inter_intro,img_of_function_intro,subseteq,img_preimg_subseteq,img_subseteq_ran,subseteq_transitive,closure_subseteq_closed,function_apply_elim,preimg_iff}.
    Hence $A$ is closed in $X$ with respect to $T$ by
        \cref{subseteq,subseteq_closure,subseteq_antisymmetric,closure_closed_in}.
\end{proof}

% from §18 Item 4: preimages of closed sets imply continuity
\begin{lemma}\label{preimg_closed_implies_continuous}
    Assume $f$ is a function from $X$ to $Y$ and
        $T$ is a topology on $X$ and
        $T_1$ is a topology on $Y$.
    Suppose for all $B$ such that $B$ is closed in $Y$ with respect to $T_1$
        we have $\preimg{f}{B}$ is closed in $X$ with respect to $T$.
    Then $f$ is continuous from $X$ to $Y$ with respect to $T$ and $T_1$.
\end{lemma}
\begin{proof}
    Hence $f\in\funs{X}{Y}$ and $\dom{f} = X$ by \cref{funs}.
    Show that for all $V$ such that $V$ is open in $Y$ with respect to $T_1$
        we have $\preimg{f}{V}$ is open in $X$ with respect to $T$.
    \begin{subproof}
        Fix $V$ such that $V$ is open in $Y$ with respect to $T_1$.
        Let $B = Y\setminus V$.
        Hence $Y\setminus B = V$ by \cref{setminus_subseteq,setminus_setminus,open_in,topology_on,subseteq,pow_iff,inter_absorb_supseteq_left}.
        Therefore $B$ is closed in $Y$ with respect to $T_1$ by \cref{setminus_subseteq,closed_in}.
        Hence $\preimg{f}{B}$ is closed in $X$ with respect to $T$ by \cref{subseteq}.
        Therefore $\preimg{f}{V}$ is open in $X$ with respect to $T$ by
            \cref{closed_in,funs_is_function,preimg_setminus_function,subseteq,open_in}.
    \end{subproof}
    Hence $f$ is continuous from $X$ to $Y$ with respect to $T$ and $T_1$ by \cref{continuous}.
\end{proof}

% from §18 Item 5: continuity at a point
\begin{definition}\label{continuous_at}
    $f$ is continuous at $x$ from $X$ to $Y$ with respect to $T$ and $T_1$ iff
        $f$ is a function from $X$ to $Y$ and
        $T$ is a topology on $X$ and
        $T_1$ is a topology on $Y$ and
        $x\in X$ and
        for all $V$ such that $V$ is a neighborhood of $f(x)$ in $Y$ with respect to $T_1$
            there exists $U$ such that
                $U$ is a neighborhood of $x$ in $X$ with respect to $T$ and
                $\img{f}{U}\subseteq V$.
\end{definition}

% from §18 Item 6: continuity implies continuity at a point
\begin{lemma}\label{continuous_implies_continuous_at}
    Assume $f$ is continuous from $X$ to $Y$ with respect to $T$ and $T_1$.
    Suppose $x\in X$.
    Then $f$ is continuous at $x$ from $X$ to $Y$ with respect to $T$ and $T_1$.
\end{lemma}
\begin{proof}
    Follows by \cref{neighborhood,continuous,funs,function_apply_elim,preimg_iff,funs_is_function,img_preimg_subseteq,continuous_at}.
\end{proof}

% from §18 Item 7: continuity at each point implies continuity
\begin{lemma}\label{continuous_at_implies_continuous}
    Assume $f$ is a function from $X$ to $Y$ and
        $T$ is a topology on $X$ and
        $T_1$ is a topology on $Y$.
    Suppose for all $x\in X$ we have $f$ is continuous at $x$ from $X$ to $Y$ with respect to $T$ and $T_1$.
    Then $f$ is continuous from $X$ to $Y$ with respect to $T$ and $T_1$.
\end{lemma}
\begin{proof}
    Hence $f\in\funs{X}{Y}$ and $\dom{f} = X$ by \cref{funs}.
    Show that for all $V$ such that $V$ is open in $Y$ with respect to $T_1$
        we have $\preimg{f}{V}$ is open in $X$ with respect to $T$.
    \begin{subproof}
        Fix $V$ such that $V$ is open in $Y$ with respect to $T_1$.
        Let $M = \{U\in T \mid U\subseteq \preimg{f}{V}\}$.
        Show that for all $x$ such that $x\in\preimg{f}{V}$ we have $x\in\unions{M}$.
        \begin{subproof}
            Fix $x$ such that $x\in\preimg{f}{V}$.
            Take $U$ such that
                $U$ is a neighborhood of $x$ in $X$ with respect to $T$ and
                $\img{f}{U}\subseteq V$ by
                \cref{preimg_subseteq_dom,subseteq,funs,funs_is_function,preimg_iff,function_apply_iff,neighborhood,continuous_at}.
            Therefore $x\in\unions{M}$ by
                \cref{preimg_subseteq_dom,subseteq,funs,funs_is_function,neighborhood,open_in,topology_on,pow_iff,inter_intro,img_of_function_intro,function_apply_elim,preimg_iff,unions_iff}.
        \end{subproof}
        Hence $\preimg{f}{V}\subseteq \unions{M}$ by \cref{subseteq}.
        Therefore $\unions{M}\in T$ and $\unions{M}\subseteq \preimg{f}{V}$
            by \cref{subseteq,pow_iff,topology_on,unions_subseteq_of_powerset_is_subseteq}.
        Hence $\preimg{f}{V}$ is open in $X$ with respect to $T$ by \cref{subseteq_antisymmetric,open_in}.
    \end{subproof}
    Therefore $f$ is continuous from $X$ to $Y$ with respect to $T$ and $T_1$ by \cref{continuous}.
\end{proof}

% from §18 Item 8: homeomorphism
\begin{definition}\label{homeomorphism}
    $f$ is a homeomorphism between $X$ and $Y$ with respect to $T$ and $T_1$ iff
        $f$ is a bijection from $X$ to $Y$ and
        $f$ is continuous from $X$ to $Y$ with respect to $T$ and $T_1$ and
        $\converse{f}$ is continuous from $Y$ to $X$ with respect to $T_1$ and $T$.
\end{definition}

% from §18 Item 9: topological property
\begin{definition}\label{topological_property}
    $P$ is a topological property iff
        for all $X,Y,T,T_1,f$ such that
            $f$ is a homeomorphism between $X$ and $Y$ with respect to $T$ and $T_1$
        we have $(X,T)\in P$ iff $(Y,T_1)\in P$.
\end{definition}

% from §18 Item 10: topological embedding
\begin{definition}\label{topological_embedding}
    $f$ is a topological embedding from $X$ to $Y$ with respect to $T$ and $T_1$ iff
        $f$ is a function from $X$ to $Y$ and
        $f$ is injective and
        $f$ is continuous from $X$ to $Y$ with respect to $T$ and $T_1$ and
        $f$ is a homeomorphism between $X$ and $\img{f}{X}$ with respect to
            $T$ and $\subspaceTopology{T_1}{\img{f}{X}}$.
\end{definition}
% from §18 Item 11: constant function is continuous
\begin{lemma}\label{constant_continuous}
    Assume $T$ is a topology on $X$ and $T_1$ is a topology on $Y$.
    Suppose $f$ is a function from $X$ to $Y$.
    Suppose there exists $y_0\in Y$ such that for all $x\in X$ we have $f(x) = y_0$.
    Then $f$ is continuous from $X$ to $Y$ with respect to $T$ and $T_1$.
\end{lemma}
\begin{proof}
    Show that for all $V$ such that $V$ is open in $Y$ with respect to $T_1$
        we have $\preimg{f}{V}$ is open in $X$ with respect to $T$.
    \begin{subproof}
        Fix $V$ such that $V$ is open in $Y$ with respect to $T_1$.
        Take $y_0\in Y$ such that for all $x\in X$ we have $f(x) = y_0$.
        \begin{byCase}
        \caseOf{$y_0\in V$.}
            Hence $\preimg{f}{V}$ is open in $X$ with respect to $T$ by
                \cref{funs,funs_is_function,funs_is_relation,function_apply_elim,preimg_iff,subseteq,preimg_subseteq_dom,subseteq_transitive,subseteq_antisymmetric,open_in,topology_on}.
        \caseOf{$y_0\notin V$.}
            Therefore $\preimg{f}{V}$ is open in $X$ with respect to $T$ by
                \cref{preimg_iff,funs,funs_is_function,function_rightunique,funs_is_relation,preimg_subseteq_dom,subseteq,function_apply_iff,emptyset_subseteq,subseteq_antisymmetric,open_in,topology_on}.
        \end{byCase}
    \end{subproof}
    Hence $f$ is continuous from $X$ to $Y$ with respect to $T$ and $T_1$ by \cref{continuous}.
\end{proof}

% from §18 Item 12: inclusion is continuous
\begin{lemma}\label{inclusion_continuous}
    Assume $A$ is a subspace of $X$ with respect to $T$.
    Let $T_A = \subspaceTopology{T}{A}$.
    Let $j = \identity{A}$.
    Then $j$ is continuous from $A$ to $X$ with respect to $T_A$ and $T$.
\end{lemma}
\begin{proof}
    $T_A$ is a topology on $A$ by \cref{subspace_topology_is_topology,subspace_of}.
    Then $T$ is a topology on $X$ and $A\subseteq X$ by \cref{subspace_of}.
    Hence $j\in\funs{A}{A}$ and $j\in\funs{A}{X}$ by \cref{id_elem_funs,funs_weaken_codom}.
    For all $U$ such that $U$ is open in $X$ with respect to $T$ we have
        $\preimg{j}{U} = A\inter U$ by
        \cref{preimg_iff,id_iff,inter_intro,inter,subseteq,subseteq_antisymmetric}.
    For all $U$ such that $U$ is open in $X$ with respect to $T$ we have
        $\preimg{j}{U}$ is open in $A$ with respect to $T_A$ by
        \cref{preimg_iff,id_iff,inter_intro,inter,subseteq,subseteq_antisymmetric,subspace_topology,open_in}.
    Hence $j$ is continuous from $A$ to $X$ with respect to $T_A$ and $T$ by \cref{continuous}.
\end{proof}

% from §18 Item 13: composites of continuous functions are continuous
\begin{lemma}\label{continuous_composite}
    Assume $f$ is continuous from $X$ to $Y$ with respect to $T$ and $T_1$.
    Assume $g$ is continuous from $Y$ to $Z$ with respect to $T_1$ and $T_2$.
    Then $g\circ f$ is continuous from $X$ to $Z$ with respect to $T$ and $T_2$.
\end{lemma}
\begin{proof}
    Follows by \cref{continuous,preimg_iff,circ_iff,subseteq,subseteq_antisymmetric,funs_circ}.
\end{proof}

% from §18 Item 14: restricting the domain
\begin{lemma}\label{continuous_restrict_domain}
    Assume $f$ is continuous from $X$ to $Y$ with respect to $T$ and $T_1$.
    Suppose $A$ is a subspace of $X$ with respect to $T$.
    Let $T_A = \subspaceTopology{T}{A}$.
    Let $g = \restrl{f}{A}$.
    Then $g$ is continuous from $A$ to $Y$ with respect to $T_A$ and $T_1$.
\end{lemma}
\begin{proof}
    $T_A$ is a topology on $A$ by \cref{subspace_topology_is_topology,subspace_of}.
    For all $V$ such that $V$ is open in $Y$ with respect to $T_1$ we have
        $\preimg{g}{V} = A\inter \preimg{f}{V}$ by
        \cref{preimg_iff,restrl_iff,inter,inter_intro,subseteq,subseteq_antisymmetric}.
    Hence for all $V$ such that $V$ is open in $Y$ with respect to $T_1$ we have
        $\preimg{g}{V}$ is open in $A$ with respect to $T_A$ by
        \cref{preimg_iff,restrl_iff,inter,inter_intro,subseteq,subseteq_antisymmetric,continuous,open_in,subspace_topology}.
    Then $f\in\funs{X}{Y}$ and $f$ is a function and $g$ is a function and $\dom{g} = A$
        by \cref{continuous,funs_is_function,restrl_of_function_is_function,dom_of_restrl_eq_inter,funs,inter_absorb_supseteq_left,subspace_of}.
    Hence $g$ is right-unique and $g$ is a relation.
    Then $A\subseteq X$ and $\dom{f} = X$ and $A\subseteq \dom{f}$
        by \cref{subspace_of,funs,subseteq_transitive}.
    $\ran{g} = \img{f}{A}$ and $\img{f}{A}\subseteq \ran{f}$ by \cref{restrl_ran,img_subseteq_ran}.
    Therefore $f\in\rels{X}{Y}$ and $\ran{f}\subseteq Y$ by \cref{funs_subseteq_rels,subseteq,rels_ran_subseteq}.
    Hence $g\in\funs{A}{Y}$ by
        \cref{function_apply_elim,ran_iff,subseteq,funs_intro,subseteq_transitive}.
    Therefore $T_1$ is a topology on $Y$ and
        $g$ is continuous from $A$ to $Y$ with respect to $T_A$ and $T_1$ by \cref{continuous}.
\end{proof}

% from §18 Item 15: restricting the range
\begin{lemma}\label{continuous_restrict_range}
    Assume $f$ is continuous from $X$ to $Y$ with respect to $T$ and $T_1$.
    Suppose $Z$ is a subspace of $Y$ with respect to $T_1$.
    Suppose $\img{f}{X}\subseteq Z$.
    Let $T_Z = \subspaceTopology{T_1}{Z}$.
    Then $f$ is continuous from $X$ to $Z$ with respect to $T$ and $T_Z$.
\end{lemma}
\begin{proof}
    $f\in\funs{X}{Y}$ and $\dom{f} = X$ by \cref{continuous,funs}.
    Hence $T_1$ is a topology on $Y$ and $T_Z$ is a topology on $Z$
        by \cref{continuous,subspace_topology_is_topology,subspace_of}.
    For all $B$ such that $B$ is open in $Z$ with respect to $T_Z$
        there exists $U\in T_1$ such that $B = Z\inter U$ by \cref{subspace_topology,open_in}.
    Thus for all $B$ such that $B$ is open in $Z$ with respect to $T_Z$
        we have $\preimg{f}{B}$ is open in $X$ with respect to $T$ by
        \cref{preimg_iff,inter_elim_right,subseteq,preimg_subseteq_dom,funs,funs_is_function,function_rightunique,funs_is_relation,function_apply_iff,img_of_function_intro,inter_intro,subseteq_antisymmetric,continuous,open_in}.
    Hence $f\in\funs{X}{Z}$ by
        \cref{function_apply_elim,ran_iff,subseteq,funs_intro,funs_is_function,img_dom,funs}.
    Therefore $f$ is continuous from $X$ to $Z$ with respect to $T$ and $T_Z$ by \cref{continuous}.
\end{proof}
% from §18 Item 16: expanding the range
\begin{lemma}\label{continuous_expand_range}
    Assume $f$ is continuous from $X$ to $Y$ with respect to $T$ and $T_1$.
    Suppose $Y$ is a subspace of $Z$ with respect to $T_2$.
    Let $T_Y = \subspaceTopology{T_2}{Y}$.
    Suppose $T_1 = T_Y$.
    Let $j = \identity{Y}$.
    Then $j\circ f$ is continuous from $X$ to $Z$ with respect to $T$ and $T_2$.
\end{lemma}
\begin{proof}
    Follows by \cref{inclusion_continuous,subseteq,continuous_composite}.
\end{proof}

% from §18 Item 17: local formulation of continuity
\begin{lemma}\label{continuous_local}
    Assume $T$ is a topology on $X$.
    Assume $T_1$ is a topology on $Y$.
    Assume $f$ is a function from $X$ to $Y$.
    Suppose $M\subseteq T$.
    Suppose $\unions{M} = X$.
    Suppose for all $U$ such that $U\in M$ we have
        $\restrl{f}{U}$ is continuous from $U$ to $Y$ with respect to
            $\subspaceTopology{T}{U}$ and $T_1$.
    Then $f$ is continuous from $X$ to $Y$ with respect to $T$ and $T_1$.
\end{lemma}
\begin{proof}
    Show that for all $V$ such that $V$ is open in $Y$ with respect to $T_1$
        we have $\preimg{f}{V}$ is open in $X$ with respect to $T$.
    \begin{subproof}
        Fix $V$ such that $V$ is open in $Y$ with respect to $T_1$.
        Let $N = \{W\in T \mid W\subseteq \preimg{f}{V}\}$.
        Show that for all $x$ such that $x\in\preimg{f}{V}$ we have $x\in\unions{N}$.
        \begin{subproof}
            Fix $x$ such that $x\in\preimg{f}{V}$.
            Take $U\in M$ such that $x\in U$ by \cref{unions_iff,subseteq,preimg_subseteq_dom,funs}.
            Let $g = \restrl{f}{U}$.
            Hence $\preimg{g}{V}\in N$ by \cref{continuous,open_in,subseteq,topology_on,pow_iff,subspace_of,open_in_subspace_open_in_space,preimg_iff,restrl_iff}.
            Therefore $x\in\unions{N}$ by \cref{restrl_iff,preimg_iff,unions_iff}.
        \end{subproof}
        Hence $\preimg{f}{V}\subseteq \unions{N}$ by \cref{subseteq}.
        Therefore $\preimg{f}{V}$ is open in $X$ with respect to $T$ by
            \cref{subseteq,pow_iff,topology_on,unions_subseteq_of_powerset_is_subseteq,subseteq_antisymmetric,open_in}.
    \end{subproof}
    Therefore $f$ is continuous from $X$ to $Y$ with respect to $T$ and $T_1$ by \cref{continuous}.
\end{proof}

% from §18 helper lemma 4: union of two compatible functions is a function
\begin{lemma}\label{union_compatible_functions}
    Suppose $f$ is a function.
    Suppose $g$ is a function.
    Suppose for all $x$ such that $x\in\dom{f}\inter\dom{g}$ we have $f(x) = g(x)$.
    Then $f\union g$ is a function.
\end{lemma}
\begin{proof}
    $f$ is a relation and $g$ is a relation and $f\union g$ is a relation
        by \cref{funs_is_relation,union_relations_is_relation}.
    Show that for all $x,y,z$ such that $(x,y)\in f\union g$ and $(x,z)\in f\union g$ we have $y = z$.
    \begin{subproof}
        Fix $x$.
        Fix $y$.
        Fix $z$ such that $(x,y)\in f\union g$ and $(x,z)\in f\union g$.
        \begin{byCase}
        \caseOf{$(x,y)\in f$ and $(x,z)\in f$.}
            Hence $y = z$ by \cref{function_rightunique,rightunique}.
        \caseOf{$(x,y)\in g$ and $(x,z)\in g$.}
            Therefore $y = z$ by \cref{function_rightunique,rightunique}.
        \caseOf{$(x,y)\in f$ and $(x,z)\in g$.}
            Then $y = z$ by \cref{dom_intro,inter_intro,subseteq,function_apply_iff}.
        \caseOf{$(x,y)\in g$ and $(x,z)\in f$.}
            Therefore $y = z$ by \cref{dom_intro,inter_intro,subseteq,function_apply_iff}.
        \end{byCase}
    \end{subproof}
    Hence $f\union g$ is a function by \cref{rightunique,relation}.
\end{proof}

% from §18 Item 18: pasting lemma
\begin{theorem}\label{pasting_lemma}
    Assume $X = A\union B$.
    Assume $A$ is closed in $X$ with respect to $T$.
    Assume $B$ is closed in $X$ with respect to $T$.
    Assume $f$ is continuous from $A$ to $Y$ with respect to $\subspaceTopology{T}{A}$ and $T_1$.
    Assume $g$ is continuous from $B$ to $Y$ with respect to $\subspaceTopology{T}{B}$ and $T_1$.
    Suppose for all $x$ such that $x\in A\inter B$ we have $f(x) = g(x)$.
    Let $h = f\union g$.
    Then $h$ is continuous from $X$ to $Y$ with respect to $T$ and $T_1$.
\end{theorem}
\begin{proof}
    $f\in\funs{A}{Y}$ and $g\in\funs{B}{Y}$ and $f$ is a function and $g$ is a function
        and $\dom{f} = A$ and $\dom{g} = B$ by \cref{continuous,funs_is_function,funs}.
    $h$ is a function by \cref{inter,funs,subseteq,union_compatible_functions}.
    Show that for all $C$ such that $C$ is closed in $Y$ with respect to $T_1$
        we have $\preimg{h}{C}$ is closed in $X$ with respect to $T$.
    \begin{subproof}
        Fix $C$ such that $C$ is closed in $Y$ with respect to $T_1$.
        Let $A_1 = \preimg{f}{C}$.
        Let $B_1 = \preimg{g}{C}$.
        Let $D = Y\setminus C$.
        $\preimg{f}{D}$ is open in $A$ with respect to $\subspaceTopology{T}{A}$ and
            $\preimg{g}{D}$ is open in $B$ with respect to $\subspaceTopology{T}{B}$ by \cref{closed_in,continuous}.
        $\preimg{f}{D} = A\setminus A_1$ and $\preimg{g}{D} = B\setminus B_1$ by
            \cref{continuous,closed_in,setminus_subseteq,preimg_setminus_function}.
        Hence $A_1\subseteq A$ and $B_1\subseteq B$ by \cref{continuous,preimg_subseteq_dom,funs}.
        Therefore $A_1 = A\setminus \preimg{f}{D}$ and $B_1 = B\setminus \preimg{g}{D}$
            by \cref{subseteq,setminus_setminus,inter_absorb_supseteq_left}.
        Then $A_1$ is closed in $A$ with respect to $\subspaceTopology{T}{A}$ and
            $B_1$ is closed in $B$ with respect to $\subspaceTopology{T}{B}$ and
            $X\setminus A$ is open in $X$ with respect to $T$ by \cref{closed_in}.
        $A_1$ is closed in $X$ with respect to $T$ and $B_1$ is closed in $X$ with respect to $T$
            by \cref{closed_in,open_in,topology_on,closed_in_closed_subspace}.
        For all $x$ such that $x\in\preimg{h}{C}$ we have $x\in A_1\union B_1$ by
            \cref{preimg_iff,union_iff}.
        Hence $\preimg{h}{C}\subseteq A_1\union B_1$ by \cref{subseteq}.
        For all $x$ such that $x\in A_1\union B_1$ we have $x\in\preimg{h}{C}$ by
            \cref{union_iff,preimg_iff,union_intro_left,union_intro_right}.
        Therefore $A_1\union B_1\subseteq \preimg{h}{C}$ by \cref{subseteq}.
        Hence $\preimg{h}{C} = A_1\union B_1$ by \cref{subseteq_antisymmetric}.
        Then $X\setminus (A_1\union B_1)$ is open in $X$ with respect to $T$
            by \cref{closed_in,open_in,topology_on,setminus_union,subseteq}.
        Therefore $A_1\union B_1$ is closed in $X$ with respect to $T$ by
            \cref{closed_in,preimg_subseteq_dom,funs,subseteq_transitive,union_subsets_is_subset}.
        Hence $\preimg{h}{C}$ is closed in $X$ with respect to $T$ by \cref{subseteq}.
    \end{subproof}
    Therefore $\dom{h} = X$ by \cref{dom_union,funs,subseteq}.
    Then $f\in\rels{A}{Y}$ and $g\in\rels{B}{Y}$ and $\ran{h} = \ran{f}\union\ran{g}$
        by \cref{funs_subseteq_rels,subseteq,ran_union}.
    Hence $\ran{h}\subseteq Y$ by \cref{rels_ran_subseteq,union_subsets_is_subset}.
    Therefore $h\in\funs{X}{Y}$ by \cref{function_apply_elim,ran_iff,subseteq,funs_intro}.
    Therefore $h$ is continuous from $X$ to $Y$ with respect to $T$ and $T_1$
        by \cref{continuous,closed_in,open_in,preimg_closed_implies_continuous}.
\end{proof}

% from §18 helper lemma 5: first projection is continuous
\begin{lemma}\label{projection_one_continuous}
    Assume $T$ is a topology on $X$ and $T_1$ is a topology on $Y$.
    Let $T_2 = \productTopology{T}{T_1}{X}{Y}$.
    Then $\piOne{X}{Y}$ is continuous from $X\times Y$ to $X$ with respect to $T_2$ and $T$.
\end{lemma}
\begin{proof}
    $T_2$ is a topology on $X\times Y$ by \cref{basis_topology_is_topology,product_basis_is_basis,product_topology}.
    For all $p$ we have for all $x$ we have for all $x'$ such that
        $p\mathrel{\piOne{X}{Y}} x$ and $p\mathrel{\piOne{X}{Y}} x'$ we have $x = x'$
        by \cref{projection_first,pair_eq_iff}.
    Therefore $\piOne{X}{Y}$ is a function
        by \cref{projection_first,pair_eq_iff,rightunique,relation}.
    For all $p\in\dom{\piOne{X}{Y}}$ we have $\apply{\piOne{X}{Y}}{p}\in X$ by
        \cref{dom_iff,projection_first,pair_eq_iff,function_apply_intro}.
    For all $p$ such that $p\in X\times Y$ we have $p\in\dom{\piOne{X}{Y}}$ by
        \cref{times_elem_is_tuple,projection_first,dom_intro}.
    Hence $X\times Y\subseteq \dom{\piOne{X}{Y}}$ by \cref{subseteq}.
    For all $p$ such that $p\in\dom{\piOne{X}{Y}}$ we have $p\in X\times Y$ by
        \cref{dom_iff,projection_first,pair_eq_iff,times}.
    Therefore $\dom{\piOne{X}{Y}}\subseteq X\times Y$ by \cref{subseteq}.
    Hence $\dom{\piOne{X}{Y}} = X\times Y$ by \cref{subseteq_antisymmetric}.
    Then $\piOne{X}{Y}\in\funs{X\times Y}{X}$ by \cref{dom_iff,projection_first,pair_eq_iff,function_apply_intro,funs_intro}.
    For all $U$ such that $U$ is open in $X$ with respect to $T$ we have
        $\preimg{\piOne{X}{Y}}{U}$ is open in $X\times Y$ with respect to $T_2$ by
        \cref{open_in,topology_on,pow_iff,subseteq,times_subseteq_left,product_basis,product_basis_is_basis,basis_elem_open_in_generated,product_topology,preimg_pi1}.
    Therefore $\piOne{X}{Y}$ is continuous from $X\times Y$ to $X$ with respect to $T_2$ and $T$ by \cref{continuous}.
\end{proof}

% from §18 helper lemma 6: second projection is continuous
\begin{lemma}\label{projection_two_continuous}
    Assume $T$ is a topology on $X$ and $T_1$ is a topology on $Y$.
    Let $T_2 = \productTopology{T}{T_1}{X}{Y}$.
    Then $\piTwo{X}{Y}$ is continuous from $X\times Y$ to $Y$ with respect to $T_2$ and $T_1$.
\end{lemma}
\begin{proof}
    $T_2$ is a topology on $X\times Y$ by \cref{basis_topology_is_topology,product_basis_is_basis,product_topology}.
    For all $p$ we have for all $y$ we have for all $y'$ such that
        $p\mathrel{\piTwo{X}{Y}} y$ and $p\mathrel{\piTwo{X}{Y}} y'$ we have $y = y'$
        by \cref{projection_second,pair_eq_iff}.
    Therefore $\piTwo{X}{Y}$ is a function
        by \cref{projection_second,pair_eq_iff,rightunique,relation}.
    For all $p\in\dom{\piTwo{X}{Y}}$ we have $\apply{\piTwo{X}{Y}}{p}\in Y$ by
        \cref{dom_iff,projection_second,pair_eq_iff,function_apply_intro}.
    For all $p$ such that $p\in X\times Y$ we have $p\in\dom{\piTwo{X}{Y}}$ by
        \cref{times_elem_is_tuple,projection_second,dom_intro}.
    Hence $X\times Y\subseteq \dom{\piTwo{X}{Y}}$ by \cref{subseteq}.
    For all $p$ such that $p\in\dom{\piTwo{X}{Y}}$ we have $p\in X\times Y$ by
        \cref{dom_iff,projection_second,pair_eq_iff,times}.
    Therefore $\dom{\piTwo{X}{Y}}\subseteq X\times Y$ by \cref{subseteq}.
    Hence $\dom{\piTwo{X}{Y}} = X\times Y$ by \cref{subseteq_antisymmetric}.
    Then $\piTwo{X}{Y}\in\funs{X\times Y}{Y}$ by \cref{dom_iff,projection_second,pair_eq_iff,function_apply_intro,funs_intro}.
    For all $V$ such that $V$ is open in $Y$ with respect to $T_1$ we have
        $\preimg{\piTwo{X}{Y}}{V}$ is open in $X\times Y$ with respect to $T_2$ by
        \cref{open_in,topology_on,pow_iff,subseteq,times_subseteq_right,product_basis,product_basis_is_basis,product_topology,basis_elem_open_in_generated,preimg_pi2}.
    Therefore $\piTwo{X}{Y}$ is continuous from $X\times Y$ to $Y$ with respect to $T_2$ and $T_1$ by \cref{continuous}.
\end{proof}

% from §18 Item 19: maps into products (forward)
\begin{lemma}\label{continuous_map_into_product_forward}
    Assume $T_1$ is a topology on $X$ and $T_2$ is a topology on $Y$.
    Assume $f$ is continuous from $A$ to $X\times Y$ with respect to $T$ and
        $\productTopology{T_1}{T_2}{X}{Y}$.
    Let $f_1 = \piOne{X}{Y}\circ f$.
    Let $f_2 = \piTwo{X}{Y}\circ f$.
    Then $f_1$ is continuous from $A$ to $X$ with respect to $T$ and $T_1$ and
        $f_2$ is continuous from $A$ to $Y$ with respect to $T$ and $T_2$.
\end{lemma}
\begin{proof}
    Follows by \cref{projection_one_continuous,projection_two_continuous,continuous_composite}.
\end{proof}

% from §18 Item 20: maps into products (backward)
\begin{lemma}\label{continuous_map_into_product_backward}
    Assume $f_1$ is continuous from $A$ to $X$ with respect to $T$ and $T_1$.
    Assume $f_2$ is continuous from $A$ to $Y$ with respect to $T$ and $T_2$.
    Let $f = \{(a,(f_1(a),f_2(a))) \mid a\in A\}$.
    Then $f$ is continuous from $A$ to $X\times Y$ with respect to $T$ and
        $\productTopology{T_1}{T_2}{X}{Y}$.
\end{lemma}
\begin{proof}
    Let $T_3 = \productTopology{T_1}{T_2}{X}{Y}$.
    Show that for all $U$ such that $U$ is open in $X\times Y$ with respect to $T_3$
        we have $\preimg{f}{U}$ is open in $A$ with respect to $T$.
    \begin{subproof}
        Fix $U$ such that $U$ is open in $X\times Y$ with respect to $T_3$.
        Take $M\subseteq \productBasis{T_1}{T_2}{X}{Y}$ such that $U = \unions{M}$ by
            \cref{open_in,continuous,basis_for_topology,product_basis_is_basis,product_topology,basis_topology_unions}.
        Let $N = \{\preimg{f}{B} \mid B\in M\}$.
        Show that for all $C$ such that $C\in N$ we have $C\in T$.
        \begin{subproof}
            Fix $C$ such that $C\in N$.
            Take $B\in M$ such that $C = \preimg{f}{B}$.
            Take $U_1,V_1$ such that $U_1\in T_1$ and $V_1\in T_2$ and $B = U_1\times V_1$
                by \cref{subseteq,product_basis}.
            For all $a$ such that $a\in\preimg{f}{B}$ we have
                $a\in\preimg{f_1}{U_1}\inter \preimg{f_2}{V_1}$ by
                \cref{preimg_iff,times_elem_is_tuple,pair_eq_iff,continuous,funs_elim,function_apply_elim,inter_intro}.
            Hence $\preimg{f}{B}\subseteq \preimg{f_1}{U_1}\inter \preimg{f_2}{V_1}$ by \cref{subseteq}.
            For all $a$ such that $a\in\preimg{f_1}{U_1}\inter \preimg{f_2}{V_1}$ we have
                $a\in\preimg{f}{B}$ by
                \cref{inter,preimg_iff,continuous,funs_elim,function_apply_intro,dom_intro,times_tuple_intro,pair_eq_iff}.
            Therefore $\preimg{f_1}{U_1}\inter \preimg{f_2}{V_1}\subseteq \preimg{f}{B}$ by \cref{subseteq}.
            Hence $C\in T$ by \cref{continuous,open_in,topology_on,subseteq_antisymmetric}.
        \end{subproof}
        $T$ is a topology on $A$ and $\unions{N}\in T$ by \cref{continuous,topology_on,subseteq}.
        Hence $\preimg{f}{U} = \unions{N}$ by
            \cref{preimg_iff,unions_iff,subseteq,subseteq_antisymmetric}.
        Hence $\preimg{f}{U}$ is open in $A$ with respect to $T$ by \cref{open_in,subseteq}.
    \end{subproof}
    For all $x,y,z$ such that $(x,y)\in f$ and $(x,z)\in f$ we have $y = z$ by \cref{pair_eq_iff}.
    Hence $f$ is a function by \cref{relation,pair_eq_iff,rightunique}.
    Hence $\dom{f} = A$ by \cref{dom_iff,pair_eq_iff,dom_intro,subseteq,subseteq_antisymmetric}.
    Then $f\in\funs{A}{X\times Y}$ by
        \cref{continuous,funs_elim,times_tuple_intro,function_apply_intro,funs_intro}.
    Therefore $f$ is continuous from $A$ to $X\times Y$ with respect to $T$ and $T_3$ by
        \cref{continuous,basis_topology_is_topology,product_basis_is_basis,product_topology}.
\end{proof}
\section{§19 The Product Topology}

% from §19 Item 1: J-tuples
\begin{definition}\label{j_tuple}
    $f$ is a $J$-tuple of elements of $X$ iff $f\in\funs{J}{X}$.
\end{definition}

% from §19 Item 2: the set of J-tuples
\begin{definition}\label{tuple_power}
    $\tuplePower{X}{J} = \funs{J}{X}$.
\end{definition}

% from §19 Item 3: product of an indexed family
\begin{definition}\label{indexed_product}
    Let $A$ be a function.
    $\productFamily{A} = \{f\in\funs{\dom{A}}{\unions{\ran{A}}} \mid \forall \alpha\in\dom{A}. \apply{f}{\alpha}\in\apply{A}{\alpha}\}$.
\end{definition}

% from §19 Item 4: box basis
\begin{definition}\label{box_basis}
    $\boxBasis{T}{X} = \{B\in\pow{\productFamily{X}} \mid \exists U.
        U\in\funs{\dom{X}}{\pow{\unions{\ran{X}}}}\land
        (\forall \alpha\in\dom{X}. \apply{U}{\alpha}\in\apply{T}{\alpha})\land
        B = \productFamily{U}\}$.
\end{definition}

% from §19 Item 5: box topology
\begin{definition}\label{box_topology}
    $\boxTopology{T}{X} = \basisTopology{\boxBasis{T}{X}}{\productFamily{X}}$.
\end{definition}

% from §19 Item 6: projection mapping
\begin{definition}\label{projection_map_indexed}
    $\piIndex{X}{\beta} = \{(x,\apply{x}{\beta}) \mid x\in\productFamily{X}\}$.
\end{definition}

% from §19 Item 7: product subbasis for a fixed index
\begin{definition}\label{product_subbasis_indexed}
    $\productSubbasisAt{T}{X}{\beta} = \{A\in\pow{\productFamily{X}} \mid
        \exists U\in\apply{T}{\beta}. A = \preimg{\piIndex{X}{\beta}}{U}\}$.
\end{definition}

% from §19 Item 8: product subbasis
\begin{definition}\label{product_subbasis_union_indexed}
    $\productSubbasis{T}{X} = \{A\in\pow{\productFamily{X}} \mid
        \exists \beta\in\dom{X}. A\in\productSubbasisAt{T}{X}{\beta}\}$.
\end{definition}

% from §19 Item 9: product topology on indexed products
\begin{definition}\label{product_topology_indexed}
    $\productTopologyIndexed{T}{X} = \subbasisTopology{\productSubbasis{T}{X}}{\productFamily{X}}$.
\end{definition}

% from §19 Item 10: box basis is a basis on the indexed product
\begin{lemma}\label{box_basis_is_basis}
    Assume $X$ is a function.
    Assume $T$ is a function.
    Assume $\dom{T} = \dom{X}$.
    Assume that for all $\alpha\in\dom{X}$ we have $\apply{T}{\alpha}$ is a topology on $\apply{X}{\alpha}$.
    Then $\boxBasis{T}{X}$ is a basis on $\productFamily{X}$.
\end{lemma}
    \begin{proof}
        Let $B = \boxBasis{T}{X}$.
        $B\subseteq\pow{\productFamily{X}}$ by \cref{box_basis,subseteq}.
        Hence $X\in\funs{\dom{X}}{\pow{\unions{\ran{X}}}}$ by
            \cref{funs_intro,function_apply_elim,ran_intro,subseteq_pow_unions,subseteq}.
        For all $\alpha\in\dom{X}$ we have $\apply{X}{\alpha}\in\apply{T}{\alpha}$ by \cref{topology_on}.
        Therefore for all $x$ such that $x\in\productFamily{X}$ there exists $B_1\in B$ such that $x\in B_1$ by
            \cref{subseteq_refl,powerset_intro,box_basis}.
        Show that for all $B_1,B_2,x$ such that $B_1\in B$ and $B_2\in B$ and $x\in B_1\inter B_2$
            there exists $B_3\in B$ such that $x\in B_3$ and $B_3\subseteq B_1\inter B_2$.
    \begin{subproof}
        Fix $B_1$.
        Fix $B_2$.
        Fix $x$ such that $B_1\in B$ and $B_2\in B$ and $x\in B_1\inter B_2$.
        Take $U_1$ such that
            $U_1\in\funs{\dom{X}}{\pow{\unions{\ran{X}}}}$ and
            $B_1 = \productFamily{U_1}$ and
            for all $\alpha\in\dom{X}$ we have $\apply{U_1}{\alpha}\in\apply{T}{\alpha}$
            by \cref{box_basis}.
        $\dom{U_1} = \dom{X}$ by \cref{funs_elim}.
        Take $U_2$ such that
            $U_2\in\funs{\dom{X}}{\pow{\unions{\ran{X}}}}$ and
            $B_2 = \productFamily{U_2}$ and
            for all $\alpha\in\dom{X}$ we have $\apply{U_2}{\alpha}\in\apply{T}{\alpha}$
            by \cref{box_basis}.
        $\dom{U_2} = \dom{X}$ by \cref{funs_elim}.
        Let $R = \{(\alpha,\apply{U_1}{\alpha}\inter\apply{U_2}{\alpha}) \mid \alpha\in\dom{X}\}$.
        Hence $R$ is a relation by \cref{times_elem_is_tuple,relation}.
        Let $U_3 = R$.
        Show that for all $a$ we have for all $y$ we have for all $z$ if $(a,y)\in U_3$ and $(a,z)\in U_3$ then $y = z$.
        \begin{subproof}
            Fix $a$.
            Fix $y$.
            Fix $z$ such that $(a,y)\in U_3$ and $(a,z)\in U_3$.
            Take $\alpha_1\in\dom{X}$ such that $(a,y) = (\alpha_1,\apply{U_1}{\alpha_1}\inter\apply{U_2}{\alpha_1})$.
            Take $\alpha_2\in\dom{X}$ such that $(a,z) = (\alpha_2,\apply{U_1}{\alpha_2}\inter\apply{U_2}{\alpha_2})$.
            Hence $a = \alpha_1$ and $y = \apply{U_1}{\alpha_1}\inter\apply{U_2}{\alpha_1}$ and
                $a = \alpha_2$ and $z = \apply{U_1}{\alpha_2}\inter\apply{U_2}{\alpha_2}$ by \cref{pair_eq_iff}.
            Therefore $y = \apply{U_1}{a}\inter\apply{U_2}{a}$.
            Then $z = \apply{U_1}{a}\inter\apply{U_2}{a}$.
            Hence $y = z$.
        \end{subproof}
        Thus $U_3$ is a function by \cref{rightunique,relation}.
        Show that $\dom{U_3} = \dom{X}$.
        \begin{subproof}
            For all $\alpha\in\dom{X}$ we have $\alpha\in\dom{U_3}$ by \cref{dom}.
            Hence $\dom{X}\subseteq\dom{U_3}$ by \cref{subseteq}.
            For all $\alpha\in\dom{U_3}$ we have $\alpha\in\dom{X}$.
            Therefore $\dom{U_3}\subseteq\dom{X}$ by \cref{subseteq}.
            Hence $\dom{U_3} = \dom{X}$ by \cref{subseteq_antisymmetric}.
        \end{subproof}
        For all $\alpha\in\dom{X}$ we have $\apply{U_3}{\alpha} = \apply{U_1}{\alpha}\inter\apply{U_2}{\alpha}$
            by \cref{function_apply_intro}.
        For all $\alpha\in\dom{X}$ we have $\apply{U_3}{\alpha}\in\pow{\unions{\ran{X}}}$ by
            \cref{funs_elim,pow_iff,inter_subseteq,powerset_intro}.
        Therefore $U_3\in\funs{\dom{X}}{\pow{\unions{\ran{X}}}}$ by
            \cref{funs_intro,funs_elim,pow_iff,inter_subseteq,powerset_intro}.
        For all $\alpha\in\dom{X}$ we have $\apply{U_3}{\alpha}\in\apply{T}{\alpha}$ by \cref{topology_on}.
        Let $B_3 = \productFamily{U_3}$.
        For all $z$ such that $z\in B_3$ we have for all $\alpha\in\dom{X}$ we have $\apply{z}{\alpha}\in\apply{X}{\alpha}$ by
            \cref{indexed_product,topology_on,subseteq,pow_iff}.
        Therefore for all $z$ such that $z\in B_3$ we have $z\in\productFamily{X}$ by
            \cref{indexed_product,topology_on,subseteq,pow_iff,funs_intro,funs_elim,function_apply_elim,ran_intro,unions_iff}.
        Hence $B_3\in\pow{\productFamily{X}}$ by \cref{subseteq,powerset_intro}.
        Therefore $B_3\in B$ by \cref{box_basis}.
        Hence $x\in\productFamily{U_1}$ and $x\in\productFamily{U_2}$ by
            \cref{inter_elim_left,inter_elim_right,subseteq_refl,elem_subseteq}.
        For all $\alpha\in\dom{X}$ we have $\apply{x}{\alpha}\in\apply{U_3}{\alpha}$
            by \cref{indexed_product,inter_intro}.
        Therefore $x\in B_3$ by
            \cref{indexed_product,inter_intro,funs_intro,funs_elim,function_apply_elim,ran_intro,unions_iff}.
        For all $z$ such that $z\in B_3$ we have $z\in B_1$ and $z\in B_2$ by
            \cref{indexed_product,inter_elim_left,funs_intro,funs_elim,function_apply_elim,ran_intro,unions_iff,inter_elim_right}.
        Therefore for all $z$ such that $z\in B_3$ we have $z\in B_1\inter B_2$ by \cref{inter_intro}.
        Therefore $B_3\subseteq B_1\inter B_2$ by \cref{subseteq}.
    \end{subproof}
    Hence $B$ is a basis on $\productFamily{X}$ by \cref{basis_on}.
\end{proof}

% from §19 Item 11: product subbasis is a subbasis on the indexed product
\begin{lemma}\label{product_subbasis_is_subbasis}
    Assume $X$ is a function.
    Assume $T$ is a function.
    Assume $\dom{T} = \dom{X}$.
    Assume $\dom{X}$ is inhabited.
    Assume that for all $\alpha\in\dom{X}$ we have $\apply{T}{\alpha}$ is a topology on $\apply{X}{\alpha}$.
    Then $\productSubbasis{T}{X}$ is a subbasis on $\productFamily{X}$.
\end{lemma}
\begin{proof}
    Let $S = \productSubbasis{T}{X}$.
    $S\subseteq\pow{\productFamily{X}}$ by \cref{product_subbasis_union_indexed,product_subbasis_indexed,subseteq}.
    Take $\beta$ such that $\beta\in\dom{X}$.
    Let $A = \preimg{\piIndex{X}{\beta}}{\apply{X}{\beta}}$.
    Hence $A\in S$ by
        \cref{topology_on,product_subbasis_indexed,product_subbasis_union_indexed,preimg_iff,projection_map_indexed,pair_eq_iff,subseteq,powerset_intro}.
    Therefore $\unions{S} = \productFamily{X}$ by
        \cref{indexed_product,projection_map_indexed,preimg_iff,unions_iff,subseteq,unions_subseteq_of_powerset_is_subseteq,subseteq_antisymmetric}.
    Hence $S$ is a subbasis on $\productFamily{X}$ by \cref{subbasis_on}.
\end{proof}

% from §19 Item 12: finite intersections of product subbasis elements
\begin{lemma}\label{product_subbasis_finite_intersections_basis}
    Assume $X$ is a function.
    Assume $T$ is a function.
    Assume $\dom{T} = \dom{X}$.
    Assume $\dom{X}$ is inhabited.
    Assume that for all $\alpha\in\dom{X}$ we have $\apply{T}{\alpha}$ is a topology on $\apply{X}{\alpha}$.
    Then $\finiteIntersections{\productSubbasis{T}{X}}$ is a basis for $\productTopologyIndexed{T}{X}$ on $\productFamily{X}$.
\end{lemma}
\begin{proof}
    Follows by \cref{product_subbasis_is_subbasis,product_topology_indexed,subbasis_finite_intersections_basis,basis_for_topology,topology_generated_by_subbasis}.
\end{proof}

% from §19 helper lemma 1: choice function on a relation
\begin{axiom}\label{choice_function}
    Suppose $R$ is a relation.
    Suppose for all $x\in A$ there exists $y$ such that $x\mathrel{R} y$.
    Then there exists $f$ such that $f$ is a function on $A$ and for all $x\in A$ we have $x\mathrel{R}\apply{f}{x}$.
\end{axiom}

% from §19 Item 13: basis for the box topology from bases
\begin{lemma}\label{box_basis_from_bases}
    Assume $X$ is a function.
    Assume $T$ is a function.
    Assume $B$ is a function.
    Assume $\dom{T} = \dom{X}$.
    Assume $\dom{B} = \dom{X}$.
    Assume that for all $\alpha\in\dom{X}$ we have $\apply{B}{\alpha}$ is a basis for $\apply{T}{\alpha}$ on $\apply{X}{\alpha}$.
    Let $C = \{A\in\pow{\productFamily{X}} \mid
        \exists U. U\in\funs{\dom{X}}{\pow{\unions{\ran{X}}}}\land
        (\forall \alpha\in\dom{X}. \apply{U}{\alpha}\in\apply{B}{\alpha})\land
        A = \productFamily{U}\}$.
    Then $C$ is a basis for $\boxTopology{T}{X}$ on $\productFamily{X}$.
\end{lemma}
\begin{proof}
    Let $T_0 = \boxTopology{T}{X}$.
    Let $B_0 = \boxBasis{T}{X}$.
    For all $\alpha\in\dom{X}$ we have $\apply{T}{\alpha}$ is a topology on $\apply{X}{\alpha}$
        by \cref{basis_for_topology,basis_topology_is_topology}.
    Hence $T_0$ is a topology on $\productFamily{X}$ by
        \cref{box_basis_is_basis,basis_topology_is_topology,box_topology}.
    Therefore for all $A$ such that $A\in C$ we have $A$ is open in $\productFamily{X}$ with respect to $T_0$
        by \cref{subseteq,basis_for_topology,basis_elem_open_in_generated,box_basis,box_basis_is_basis,box_topology,open_in}.
    Show that for all $U,x$ such that $U\in T_0$ and $x\in U$ there exists $C_1\in C$ such that $x\in C_1$ and $C_1\subseteq U$.
    \begin{subproof}
        Fix $U$.
        Fix $x$ such that $U\in T_0$ and $x\in U$.
        Take $B_1\in B_0$ such that $x\in B_1$ and $B_1\subseteq U$ by \cref{box_topology,topology_generated_by_basis}.
        Take $V$ such that
            $V\in\funs{\dom{X}}{\pow{\unions{\ran{X}}}}$ and
            $B_1 = \productFamily{V}$ and
            for all $\alpha\in\dom{X}$ we have $\apply{V}{\alpha}\in\apply{T}{\alpha}$
            by \cref{box_basis}.
        $\dom{V} = \dom{X}$ by \cref{funs_elim}.
        For all $\alpha\in\dom{X}$ there exists $W$ such that
            $W\in\apply{B}{\alpha}$ and $\apply{x}{\alpha}\in W$ and $W\subseteq\apply{V}{\alpha}$
            by \cref{indexed_product,open_in,basis_for_topology,topology_generated_by_basis}.
        Let $R = \{r\in\dom{X}\times\pow{\unions{\ran{X}}} \mid \apply{x}{\fst{r}}\in\snd{r} \land \snd{r}\subseteq\apply{V}{\fst{r}} \land \snd{r}\in\apply{B}{\fst{r}}\}$.
        Hence $R$ is a relation by \cref{times_elem_is_tuple,relation}.
        Let $A = \dom{X}$.
        Show that for all $\alpha\in A$ there exists $W$ such that $\alpha\mathrel{R} W$.
        \begin{subproof}
            Fix $\alpha$ such that $\alpha\in A$.
            Take $W$ such that $W\in\apply{B}{\alpha}$ and $\apply{x}{\alpha}\in W$ and $W\subseteq\apply{V}{\alpha}$.
            Hence $W\in\pow{\unions{\ran{X}}}$ by
                \cref{topology_on,subseteq,pow_iff,function_apply_elim,ran_intro,unions_iff,powerset_intro}.
            Therefore $(\alpha,W)\in\dom{X}\times\pow{\unions{\ran{X}}}$ by \cref{times_tuple_intro}.
            Then $\apply{x}{\fst{(\alpha,W)}}\in\snd{(\alpha,W)}$ and
                $\snd{(\alpha,W)}\subseteq\apply{V}{\fst{(\alpha,W)}}$ and
                $\snd{(\alpha,W)}\in\apply{B}{\fst{(\alpha,W)}}$ by \cref{fst_eq,snd_eq}.
            Hence $\alpha\mathrel{R} W$.
        \end{subproof}
        Take $U_1$ such that $U_1$ is a function on $A$ and for all $\alpha\in A$ we have $\alpha\mathrel{R}\apply{U_1}{\alpha}$
            by \cref{choice_function}.
        For all $\alpha\in\dom{X}$ we have $\apply{U_1}{\alpha}\in\apply{B}{\alpha}$ and
            $\apply{x}{\alpha}\in\apply{U_1}{\alpha}$ and $\apply{U_1}{\alpha}\subseteq\apply{V}{\alpha}$
            by \cref{fst_eq,snd_eq}.
        For all $\alpha\in\dom{X}$ we have $\apply{U_1}{\alpha}\in\pow{\unions{\ran{X}}}$ by
            \cref{topology_on,subseteq,pow_iff,function_apply_elim,ran_intro,unions_iff,powerset_intro}.
        Therefore $U_1\in\funs{\dom{X}}{\pow{\unions{\ran{X}}}}$ by
            \cref{funs_intro,topology_on,subseteq,pow_iff,function_apply_elim,ran_intro,unions_iff,powerset_intro}.
        Let $C_1 = \productFamily{U_1}$.
        For all $z$ such that $z\in C_1$ we have for all $\alpha\in\dom{X}$ we have $\apply{z}{\alpha}\in\apply{X}{\alpha}$ by
            \cref{indexed_product,topology_on,subseteq,pow_iff,subseteq_transitive}.
        Therefore for all $z$ such that $z\in C_1$ we have $z\in\productFamily{X}$ by
            \cref{indexed_product,topology_on,subseteq,pow_iff,funs_intro,funs_elim,function_apply_elim,ran_intro,unions_iff}.
        Hence $C_1\in\pow{\productFamily{X}}$ by \cref{subseteq,powerset_intro}.
        Therefore $C_1\in C$.
        For all $\alpha\in\dom{X}$ we have $\apply{x}{\alpha}\in\apply{U_1}{\alpha}$.
        Hence $x\in C_1$ by
            \cref{indexed_product,funs_intro,funs_elim,function_apply_elim,ran_intro,unions_iff}.
        For all $z$ such that $z\in C_1$ we have $z\in U$ by
            \cref{indexed_product,subseteq,funs_intro,funs_elim,function_apply_elim,ran_intro,unions_iff}.
        Therefore $C_1\subseteq U$ by \cref{subseteq}.
    \end{subproof}
    Hence $C$ is a basis for $T_0$ on $\productFamily{X}$ by \cref{open_collection_basis}.
\end{proof}

% from §19 Item 14: basis for the product topology from bases
\begin{lemma}\label{product_basis_from_bases_indexed}
    Assume $X$ is a function.
    Assume $T$ is a function.
    Assume $B$ is a function.
    Assume $\dom{T} = \dom{X}$.
    Assume $\dom{B} = \dom{X}$.
    Assume $\dom{X}$ is inhabited.
    Assume that for all $\alpha\in\dom{X}$ we have $\apply{B}{\alpha}$ is a basis for $\apply{T}{\alpha}$ on $\apply{X}{\alpha}$.
    Let $S$ be a set such that for all $A$ we have
        $A\in S$ iff there exists $\beta\in\dom{X}$ such that there exists $U\in\apply{B}{\beta}$ such that
            $A = \preimg{\piIndex{X}{\beta}}{U}$.
    Then $\finiteIntersections{S}$ is a basis for $\productTopologyIndexed{T}{X}$ on $\productFamily{X}$.
\end{lemma}
\begin{proof}
    Let $T_0 = \productTopologyIndexed{T}{X}$.
    Let $T_1 = \subbasisTopology{S}{\productFamily{X}}$.
    For all $A$ such that $A\in S$ we have $A\in\pow{\productFamily{X}}$ by
        \cref{preimg_iff,projection_map_indexed,pair_eq_iff,subseteq,powerset_intro}.
    Therefore $S\subseteq\pow{\productFamily{X}}$ by \cref{subseteq}.
    For all $x$ such that $x\in\productFamily{X}$ we have $x\in\unions{S}$ by
        \cref{indexed_product,basis_for_topology,basis_on,projection_map_indexed,preimg_iff,unions_iff}.
    Therefore $\unions{S} = \productFamily{X}$ by
        \cref{subseteq,indexed_product,projection_map_indexed,preimg_iff,unions_iff,unions_subseteq_of_powerset_is_subseteq,subseteq_antisymmetric}.
    Hence $S$ is a subbasis on $\productFamily{X}$ by \cref{subbasis_on}.
    Therefore $\finiteIntersections{S}$ is a basis for $T_1$ on $\productFamily{X}$
        by \cref{subbasis_finite_intersections_basis,basis_for_topology,topology_generated_by_subbasis}.
    Show that for all $C$ such that $C\in S$ we have $C\in T_1$.
        \begin{subproof}
        Fix $C$ such that $C\in S$.
        Therefore $C\in\finiteIntersections{S}$ and $C\in T_1$
            by \cref{subbasis_on,subseteq,pow_iff,singleton_subset_intro,upair_iff,pair_is_finite,subseteq_finite_is_finite,singleton_intro,inters_singleton,finite_intersections,basis_elem_open_in_generated,basis_for_topology}.
    \end{subproof}
    Therefore $S\subseteq T_1$ by \cref{subseteq}.
    For all $\alpha\in\dom{X}$ we have $\apply{T}{\alpha}$ is a topology on $\apply{X}{\alpha}$
        by \cref{basis_for_topology,basis_topology_is_topology}.
    Let $S_0 = \productSubbasis{T}{X}$.
    $T_0 = \subbasisTopology{S_0}{\productFamily{X}}$ by \cref{product_topology_indexed}.
    Hence $\finiteIntersections{S_0}$ is a basis for $T_0$ on $\productFamily{X}$ by
        \cref{product_subbasis_finite_intersections_basis,product_topology_indexed}.
    Therefore $T_0$ is a topology on $\productFamily{X}$ and
        $T_1$ is a topology on $\productFamily{X}$ by \cref{basis_topology_is_topology,basis_for_topology}.
    Show that for all $A$ such that $A\in S$ we have $A\in T_0$.
    \begin{subproof}
        Fix $A$ such that $A\in S$.
        Take $\beta$ such that $\beta\in\dom{X}$ and there exists $U\in\apply{B}{\beta}$ such that
            $A = \preimg{\piIndex{X}{\beta}}{U}$.
        Take $U\in\apply{B}{\beta}$ such that $A = \preimg{\piIndex{X}{\beta}}{U}$.
        Therefore $A\in\finiteIntersections{S_0}$ and $A\in T_0$
            by \cref{basis_for_topology,basis_elem_open_in_generated,preimg_iff,projection_map_indexed,pair_eq_iff,subseteq,powerset_intro,product_subbasis_indexed,product_subbasis_union_indexed,product_subbasis_is_subbasis,subbasis_on,pow_iff,singleton_subset_intro,upair_iff,pair_is_finite,subseteq_finite_is_finite,singleton_intro,inters_singleton,finite_intersections}.
    \end{subproof}
    For all $B$ such that $B\in\finiteIntersections{S}$ we have $B\in T_0$ by
        \cref{finite_intersections,subseteq,finite_intersections_open_in_topology}.
    Therefore $\finiteIntersections{S}\subseteq T_0$ by \cref{subseteq}.
    For all $U$ such that $U\in T_1$ we have $U\in T_0$ by \cref{basis_topology_unions,subseteq,topology_on}.
    Hence $T_1\subseteq T_0$ by \cref{subseteq}.
    Show that for all $A$ such that $A\in S_0$ we have $A\in T_1$.
    \begin{subproof}
        Fix $A$ such that $A\in S_0$.
        Take $\beta\in\dom{X}$ such that $A\in\productSubbasisAt{T}{X}{\beta}$ by \cref{product_subbasis_union_indexed}.
        Take $U\in\apply{T}{\beta}$ such that $A = \preimg{\piIndex{X}{\beta}}{U}$ by \cref{product_subbasis_indexed}.
        Take $M\subseteq\apply{B}{\beta}$ such that $U = \unions{M}$ by \cref{basis_for_topology,basis_topology_unions}.
        Let $N = \{\preimg{\piIndex{X}{\beta}}{V} \mid V\in M\}$.
        Hence $N\subseteq T_1$ by \cref{subseteq,subseteq_transitive}.
        Therefore $A\subseteq\unions{N}$ by \cref{preimg_iff,unions_iff,subseteq}.
        Therefore $\unions{N}\subseteq A$ by \cref{unions_iff,subseteq,preimg_subseteq}.
        Therefore $A\in T_1$ by \cref{subseteq_antisymmetric,topology_on,subseteq}.
    \end{subproof}
    For all $B$ such that $B\in\finiteIntersections{S_0}$ we have $B\in T_1$ by
        \cref{finite_intersections,subseteq,finite_intersections_open_in_topology}.
    Therefore $\finiteIntersections{S_0}\subseteq T_1$ by \cref{subseteq}.
    For all $U$ such that $U\in T_0$ we have $U\in T_1$ by \cref{basis_topology_unions,subseteq,topology_on}.
    Hence $T_0\subseteq T_1$ by \cref{subseteq}.
    Hence $T_0 = T_1$ by \cref{subseteq_antisymmetric}.
    Therefore $\finiteIntersections{S}$ is a basis for $T_0$ on $\productFamily{X}$ by \cref{basis_for_topology}.
\end{proof}
% from §19 helper lemma 2: replace one coordinate in a function
\begin{lemma}\label{replace_coordinate_function}
    Assume $X$ is a function.
    Assume $\alpha\in\dom{X}$.
    Assume $U\subseteq\unions{\ran{X}}$.
    Then there exists $V$ such that
        $V\in\funs{\dom{X}}{\pow{\unions{\ran{X}}}}$ and
        $\apply{V}{\alpha} = U$ and
        for all $\beta\in\dom{X}$ such that $\beta\neq\alpha$ we have $\apply{V}{\beta} = \apply{X}{\beta}$.
\end{lemma}
\begin{proof}
    Let $R_1 = \{(\alpha,U)\}$.
    Let $R_2 = \{(\beta,\apply{X}{\beta}) \mid \beta\in\dom{X}\setminus\{\alpha\}\}$.
    Let $R = R_1\union R_2$.
    Show that for all $w$ such that $w\in R$ there exist $x,y$ such that $w = (x,y)$.
    \begin{subproof}
        Fix $w$ such that $w\in R$.
        Hence $w\in R_1$ or $w\in R_2$ by \cref{union_iff}.
        \begin{byCase}
        \caseOf{$w\in R_1$.}
            Therefore $w = (\alpha,U)$ by \cref{singleton_elim}.
        \caseOf{$w\in R_2$.}
            Take $\beta\in\dom{X}\setminus\{\alpha\}$ such that $w = (\beta,\apply{X}{\beta})$.
        \end{byCase}
    \end{subproof}
    Hence $R$ is a relation by \cref{relation}.
    Show that for all $\beta\in\dom{X}$ there exists $y$ such that $\beta\mathrel{R} y$.
    \begin{subproof}
        Fix $\beta$ such that $\beta\in\dom{X}$.
        \begin{byCase}
        \caseOf{$\beta = \alpha$.}
            Hence $\beta\mathrel{R} U$ by \cref{singleton_intro,union_iff}.
        \caseOf{$\beta\neq\alpha$.}
            Therefore $\beta\mathrel{R}\apply{X}{\beta}$ by \cref{singleton_iff,setminus_intro,union_iff}.
        \end{byCase}
    \end{subproof}
    Take $V$ such that $V$ is a function on $\dom{X}$ and for all $\beta\in\dom{X}$ we have $\beta\mathrel{R}\apply{V}{\beta}$
        by \cref{choice_function}.
    Then $V$ is a function and $\dom{V} = \dom{X}$.
    Show that for all $\beta\in\dom{X}$ we have $\apply{V}{\beta}\in\pow{\unions{\ran{X}}}$.
    \begin{subproof}
        Fix $\beta$ such that $\beta\in\dom{X}$.
        Therefore $(\beta,\apply{V}{\beta})\in R_1$ or $(\beta,\apply{V}{\beta})\in R_2$ by \cref{union_iff}.
        \begin{byCase}
        \caseOf{$(\beta,\apply{V}{\beta})\in R_1$.}
            Hence $\apply{V}{\beta}\in\pow{\unions{\ran{X}}}$ by \cref{singleton_elim,pair_eq_iff,powerset_intro,subseteq}.
        \caseOf{$(\beta,\apply{V}{\beta})\in R_2$.}
            Take $\gamma\in\dom{X}\setminus\{\alpha\}$ such that $(\beta,\apply{V}{\beta}) = (\gamma,\apply{X}{\gamma})$.
            Therefore $\apply{V}{\beta}\in\pow{\unions{\ran{X}}}$ by \cref{pair_eq_iff,function_apply_elim,ran_intro,unions_iff,subseteq,powerset_intro}.
        \end{byCase}
    \end{subproof}
    Therefore $V\in\funs{\dom{X}}{\pow{\unions{\ran{X}}}}$ by \cref{funs_intro}.
    $(\alpha,\apply{V}{\alpha})\notin R_2$ by \cref{pair_eq_iff,setminus_intro,setminus_elim_right,singleton_intro}.
    Hence $\apply{V}{\alpha} = U$ by \cref{union_iff,singleton_elim,pair_eq_iff}.
    Show that for all $\beta\in\dom{X}$ such that $\beta\neq\alpha$ we have $\apply{V}{\beta} = \apply{X}{\beta}$.
    \begin{subproof}
        Fix $\beta$ such that $\beta\in\dom{X}$ and $\beta\neq\alpha$.
        $(\beta,\apply{V}{\beta})\notin R_1$ by \cref{singleton_elim,pair_eq_iff}.
        Therefore $(\beta,\apply{V}{\beta})\in R_2$ by \cref{union_iff}.
        Take $\gamma\in\dom{X}\setminus\{\alpha\}$ such that $(\beta,\apply{V}{\beta}) = (\gamma,\apply{X}{\gamma})$.
        Hence $\beta = \gamma$ and $\apply{V}{\beta} = \apply{X}{\beta}$ by \cref{pair_eq_iff}.
    \end{subproof}
\end{proof}
% from §19 Item 15: products of subspaces for the box topology
\begin{lemma}\label{product_subspace_box}
    Assume $X$ is a function.
    Assume $A$ is a function.
    Assume $T$ is a function.
    Assume $\dom{A} = \dom{X}$.
    Assume $\dom{T} = \dom{X}$.
    Assume that for all $\alpha\in\dom{X}$ we have $\apply{A}{\alpha}\subseteq\apply{X}{\alpha}$.
    Assume that for all $\alpha\in\dom{X}$ we have $\apply{T}{\alpha}$ is a topology on $\apply{X}{\alpha}$.
    Then $\productFamily{A}$ is a subspace of $\productFamily{X}$ with respect to $\boxTopology{T}{X}$.
\end{lemma}
\begin{proof}
    Therefore $\productFamily{A}$ is a subspace of $\productFamily{X}$ with respect to $\boxTopology{T}{X}$
        by \cref{box_basis_is_basis,basis_topology_is_topology,box_topology,indexed_product,subseteq,funs_elim,function_apply_elim,ran_intro,unions_iff,funs_intro,subspace_of}.
\end{proof}

% from §19 Item 16: products of subspaces for the product topology
\begin{lemma}\label{product_subspace_product}
    Assume $X$ is a function.
    Assume $A$ is a function.
    Assume $T$ is a function.
    Assume $\dom{A} = \dom{X}$.
    Assume $\dom{T} = \dom{X}$.
    Assume $\dom{X}$ is inhabited.
    Assume that for all $\alpha\in\dom{X}$ we have $\apply{A}{\alpha}\subseteq\apply{X}{\alpha}$.
    Assume that for all $\alpha\in\dom{X}$ we have $\apply{T}{\alpha}$ is a topology on $\apply{X}{\alpha}$.
    Then $\productFamily{A}$ is a subspace of $\productFamily{X}$ with respect to $\productTopologyIndexed{T}{X}$.
\end{lemma}
\begin{proof}
    Let $T_0 = \productTopologyIndexed{T}{X}$.
    Therefore $\productFamily{A}$ is a subspace of $\productFamily{X}$ with respect to $T_0$
        by \cref{product_subbasis_finite_intersections_basis,basis_topology_is_topology,basis_for_topology,indexed_product,subseteq,funs_elim,function_apply_elim,ran_intro,unions_iff,funs_intro,subspace_of}.
\end{proof}

% from §19 Item 17: product of Hausdorff spaces with the box topology
\begin{lemma}\label{product_hausdorff_box}
    Assume $X$ is a function.
    Assume $T$ is a function.
    Assume $\dom{T} = \dom{X}$.
    Assume that for all $\alpha\in\dom{X}$ we have $\apply{X}{\alpha}$ is Hausdorff with respect to $\apply{T}{\alpha}$.
    Then $\productFamily{X}$ is Hausdorff with respect to $\boxTopology{T}{X}$.
\end{lemma}
\begin{proof}
    Let $T_0 = \boxTopology{T}{X}$.
    Let $B_0 = \boxBasis{T}{X}$.
    For all $\alpha\in\dom{X}$ we have $\apply{T}{\alpha}$ is a topology on $\apply{X}{\alpha}$ by \cref{hausdorff}.
    $\boxBasis{T}{X}$ is a basis on $\productFamily{X}$ by \cref{box_basis_is_basis}.
    Hence $T_0$ is a topology on $\productFamily{X}$ by \cref{basis_topology_is_topology,box_topology}.
    Show that for all $x_1,x_2$ such that $x_1\in\productFamily{X}$ and $x_2\in\productFamily{X}$ and $x_1\neq x_2$
        there exist $U_1,U_2$ such that
            $U_1$ is a neighborhood of $x_1$ in $\productFamily{X}$ with respect to $T_0$ and
            $U_2$ is a neighborhood of $x_2$ in $\productFamily{X}$ with respect to $T_0$ and
            $U_1$ is disjoint from $U_2$.
    \begin{subproof}
        Fix $x_1$.
        Fix $x_2$ such that $x_1\in\productFamily{X}$ and $x_2\in\productFamily{X}$ and $x_1\neq x_2$.
        Take $w$ such that either $w\in x_1$ and $w\notin x_2$ or $w\notin x_1$ and $w\in x_2$
            by \cref{neq_witness}.
        \begin{byCase}
        \caseOf{$w\in x_1$ and $w\notin x_2$.}
            Hence $x_1$ is a function and $\dom{x_1} = \dom{X}$ and
                $x_2$ is a function and $\dom{x_2} = \dom{X}$ by \cref{indexed_product,funs_elim}.
            Take $\alpha\in\dom{X}$ such that $w = (\alpha,\apply{x_1}{\alpha})$ by \cref{funs_elim,function_member_elim}.
            Therefore $\apply{x_1}{\alpha}\neq\apply{x_2}{\alpha}$ by \cref{function_apply_elim,pair_eq_iff}.
            $\apply{x_1}{\alpha}\in\apply{X}{\alpha}$ and
                $\apply{x_2}{\alpha}\in\apply{X}{\alpha}$ by \cref{indexed_product}.
            Take $U_1,U_2$ such that
                $U_1$ is a neighborhood of $\apply{x_1}{\alpha}$ in $\apply{X}{\alpha}$ with respect to $\apply{T}{\alpha}$ and
                $U_2$ is a neighborhood of $\apply{x_2}{\alpha}$ in $\apply{X}{\alpha}$ with respect to $\apply{T}{\alpha}$ and
                $U_1$ is disjoint from $U_2$ by \cref{hausdorff}.
            Take $V_1$ such that
                $V_1\in\funs{\dom{X}}{\pow{\unions{\ran{X}}}}$ and
                $\apply{V_1}{\alpha} = U_1$ and
                for all $\beta\in\dom{X}$ such that $\beta\neq\alpha$ we have $\apply{V_1}{\beta} = \apply{X}{\beta}$
                by \cref{neighborhood,open_in,hausdorff,topology_on,subseteq,pow_iff,function_apply_elim,ran_intro,unions_iff,subseteq_transitive,replace_coordinate_function}.
            Take $V_2$ such that
                $V_2\in\funs{\dom{X}}{\pow{\unions{\ran{X}}}}$ and
                $\apply{V_2}{\alpha} = U_2$ and
                for all $\beta\in\dom{X}$ such that $\beta\neq\alpha$ we have $\apply{V_2}{\beta} = \apply{X}{\beta}$
                by \cref{neighborhood,open_in,hausdorff,topology_on,subseteq,pow_iff,function_apply_elim,ran_intro,unions_iff,subseteq_transitive,replace_coordinate_function}.
            Let $W_1 = \productFamily{V_1}$.
            Let $W_2 = \productFamily{V_2}$.
            Show that for all $x$ such that $x\in W_1$ we have $x\in\productFamily{X}$.
            \begin{subproof}
                Fix $x$ such that $x\in W_1$.
                Hence $x$ is a function and $\dom{x} = \dom{X}$ by \cref{indexed_product,funs_elim}.
                Show that for all $\beta\in\dom{X}$ we have $\apply{x}{\beta}\in\apply{X}{\beta}$.
                \begin{subproof}
                    Fix $\beta$ such that $\beta\in\dom{X}$.
                    Hence $\apply{x}{\beta}\in\apply{V_1}{\beta}$ by \cref{funs_elim,indexed_product}.
                    \begin{byCase}
                    \caseOf{$\beta = \alpha$.}
                        Hence $\apply{x}{\beta}\in\apply{X}{\beta}$ by \cref{neighborhood,open_in,hausdorff,topology_on,subseteq,pow_iff,elem_subseteq}.
                    \caseOf{$\beta\neq\alpha$.}
                        Therefore $\apply{x}{\beta}\in\apply{X}{\beta}$ by \cref{subseteq_refl,elem_subseteq}.
                    \end{byCase}
                \end{subproof}
                Hence $x\in\productFamily{X}$ by \cref{funs_intro,indexed_product,function_apply_elim,ran_intro,unions_iff}.
            \end{subproof}
            Show that for all $x$ such that $x\in W_2$ we have $x\in\productFamily{X}$.
            \begin{subproof}
                Fix $x$ such that $x\in W_2$.
                Hence $x$ is a function and $\dom{x} = \dom{X}$ by \cref{indexed_product,funs_elim}.
                Show that for all $\beta\in\dom{X}$ we have $\apply{x}{\beta}\in\apply{X}{\beta}$.
                \begin{subproof}
                    Fix $\beta$ such that $\beta\in\dom{X}$.
                    Hence $\apply{x}{\beta}\in\apply{V_2}{\beta}$ by \cref{funs_elim,indexed_product}.
                    \begin{byCase}
                    \caseOf{$\beta = \alpha$.}
                        Hence $\apply{x}{\beta}\in\apply{X}{\beta}$ by \cref{neighborhood,open_in,hausdorff,topology_on,subseteq,pow_iff,elem_subseteq}.
                    \caseOf{$\beta\neq\alpha$.}
                        Therefore $\apply{x}{\beta}\in\apply{X}{\beta}$ by \cref{subseteq_refl,elem_subseteq}.
                    \end{byCase}
                \end{subproof}
                Hence $x\in\productFamily{X}$ by \cref{funs_intro,indexed_product,function_apply_elim,ran_intro,unions_iff}.
            \end{subproof}
            Therefore $W_1$ is open in $\productFamily{X}$ with respect to $T_0$ and
                $W_2$ is open in $\productFamily{X}$ with respect to $T_0$
                by \cref{neighborhood,open_in,topology_on,subseteq,pow_iff,box_basis,basis_elem_open_in_generated,box_topology,open_in}.
            For all $\beta\in\dom{X}$ we have $\apply{x_1}{\beta}\in\apply{V_1}{\beta}$
                by \cref{neighborhood,indexed_product,subseteq}.
            Hence $x_1$ is a function and $\dom{x_1} = \dom{X}$ and
                $\dom{V_1} = \dom{X}$ by \cref{funs_elim,indexed_product}.
            For all $\beta\in\dom{X}$ we have $\apply{x_2}{\beta}\in\apply{V_2}{\beta}$
                by \cref{neighborhood,indexed_product,subseteq}.
            Hence $x_2$ is a function and $\dom{x_2} = \dom{X}$ and
                $\dom{V_2} = \dom{X}$ by \cref{funs_elim,indexed_product}.
            Therefore $W_1$ is a neighborhood of $x_1$ in $\productFamily{X}$ with respect to $T_0$ and
                $W_2$ is a neighborhood of $x_2$ in $\productFamily{X}$ with respect to $T_0$
                by \cref{funs_intro,indexed_product,funs_elim,function_apply_elim,ran_intro,unions_iff,neighborhood}.
            Show that $W_1$ is disjoint from $W_2$.
            \begin{subproof}
                Suppose not.
                Then there exists $p$ such that $p\in W_1$ and $p\in W_2$.
                Hence $\apply{p}{\alpha}\in U_1$ and $\apply{p}{\alpha}\in U_2$ by \cref{funs_elim,indexed_product}.
                Then $U_1$ intersects $U_2$ by \cref{inter_intro,emptyset,intersects}.
                Contradiction by \cref{disjoint}.
            \end{subproof}
        \caseOf{$w\notin x_1$ and $w\in x_2$.}
            Hence $x_1$ is a function and $\dom{x_1} = \dom{X}$ and
                $x_2$ is a function and $\dom{x_2} = \dom{X}$ by \cref{indexed_product,funs_elim}.
            Take $\alpha\in\dom{X}$ such that $w = (\alpha,\apply{x_2}{\alpha})$ by \cref{funs_elim,function_member_elim}.
            Therefore $\apply{x_1}{\alpha}\neq\apply{x_2}{\alpha}$ by \cref{function_apply_elim,pair_eq_iff}.
            $\apply{x_1}{\alpha}\in\apply{X}{\alpha}$ and
                $\apply{x_2}{\alpha}\in\apply{X}{\alpha}$ by \cref{indexed_product}.
            Take $U_1,U_2$ such that
                $U_1$ is a neighborhood of $\apply{x_1}{\alpha}$ in $\apply{X}{\alpha}$ with respect to $\apply{T}{\alpha}$ and
                $U_2$ is a neighborhood of $\apply{x_2}{\alpha}$ in $\apply{X}{\alpha}$ with respect to $\apply{T}{\alpha}$ and
                $U_1$ is disjoint from $U_2$ by \cref{hausdorff}.
            Take $V_1$ such that
                $V_1\in\funs{\dom{X}}{\pow{\unions{\ran{X}}}}$ and
                $\apply{V_1}{\alpha} = U_1$ and
                for all $\beta\in\dom{X}$ such that $\beta\neq\alpha$ we have $\apply{V_1}{\beta} = \apply{X}{\beta}$
                by \cref{neighborhood,open_in,hausdorff,topology_on,subseteq,pow_iff,function_apply_elim,ran_intro,unions_iff,subseteq_transitive,replace_coordinate_function}.
            Take $V_2$ such that
                $V_2\in\funs{\dom{X}}{\pow{\unions{\ran{X}}}}$ and
                $\apply{V_2}{\alpha} = U_2$ and
                for all $\beta\in\dom{X}$ such that $\beta\neq\alpha$ we have $\apply{V_2}{\beta} = \apply{X}{\beta}$
                by \cref{neighborhood,open_in,hausdorff,topology_on,subseteq,pow_iff,function_apply_elim,ran_intro,unions_iff,subseteq_transitive,replace_coordinate_function}.
            Let $W_1 = \productFamily{V_1}$.
            Let $W_2 = \productFamily{V_2}$.
            Show that for all $x$ such that $x\in W_1$ we have $x\in\productFamily{X}$.
            \begin{subproof}
                Fix $x$ such that $x\in W_1$.
                Hence $x$ is a function and $\dom{x} = \dom{X}$ by \cref{indexed_product,funs_elim}.
                Show that for all $\beta\in\dom{X}$ we have $\apply{x}{\beta}\in\apply{X}{\beta}$.
                \begin{subproof}
                    Fix $\beta$ such that $\beta\in\dom{X}$.
                    Hence $\apply{x}{\beta}\in\apply{V_1}{\beta}$ by \cref{funs_elim,indexed_product}.
                    \begin{byCase}
                    \caseOf{$\beta = \alpha$.}
                        Hence $\apply{x}{\beta}\in\apply{X}{\beta}$ by \cref{neighborhood,open_in,hausdorff,topology_on,subseteq,pow_iff,elem_subseteq}.
                    \caseOf{$\beta\neq\alpha$.}
                        Therefore $\apply{x}{\beta}\in\apply{X}{\beta}$ by \cref{subseteq_refl,elem_subseteq}.
                    \end{byCase}
                \end{subproof}
                Hence $x\in\productFamily{X}$ by \cref{funs_intro,indexed_product,function_apply_elim,ran_intro,unions_iff}.
            \end{subproof}
            Show that for all $x$ such that $x\in W_2$ we have $x\in\productFamily{X}$.
            \begin{subproof}
                Fix $x$ such that $x\in W_2$.
                Hence $x$ is a function and $\dom{x} = \dom{X}$ by \cref{indexed_product,funs_elim}.
                Show that for all $\beta\in\dom{X}$ we have $\apply{x}{\beta}\in\apply{X}{\beta}$.
                \begin{subproof}
                    Fix $\beta$ such that $\beta\in\dom{X}$.
                    Hence $\apply{x}{\beta}\in\apply{V_2}{\beta}$ by \cref{funs_elim,indexed_product}.
                    \begin{byCase}
                    \caseOf{$\beta = \alpha$.}
                        Hence $\apply{x}{\beta}\in\apply{X}{\beta}$ by \cref{neighborhood,open_in,hausdorff,topology_on,subseteq,pow_iff,elem_subseteq}.
                    \caseOf{$\beta\neq\alpha$.}
                        Therefore $\apply{x}{\beta}\in\apply{X}{\beta}$ by \cref{subseteq_refl,elem_subseteq}.
                    \end{byCase}
                \end{subproof}
                Hence $x\in\productFamily{X}$ by \cref{funs_intro,indexed_product,function_apply_elim,ran_intro,unions_iff}.
            \end{subproof}
            Therefore $W_1$ is open in $\productFamily{X}$ with respect to $T_0$ and
                $W_2$ is open in $\productFamily{X}$ with respect to $T_0$
                by \cref{neighborhood,open_in,topology_on,subseteq,pow_iff,box_basis,basis_elem_open_in_generated,box_topology,open_in}.
            For all $\beta\in\dom{X}$ we have $\apply{x_1}{\beta}\in\apply{V_1}{\beta}$
                by \cref{neighborhood,indexed_product,subseteq}.
            Hence $x_1$ is a function and $\dom{x_1} = \dom{X}$ and
                $\dom{V_1} = \dom{X}$ by \cref{funs_elim,indexed_product}.
            For all $\beta\in\dom{X}$ we have $\apply{x_2}{\beta}\in\apply{V_2}{\beta}$
                by \cref{neighborhood,indexed_product,subseteq}.
            Hence $x_2$ is a function and $\dom{x_2} = \dom{X}$ and
                $\dom{V_2} = \dom{X}$ by \cref{funs_elim,indexed_product}.
            Therefore $W_1$ is a neighborhood of $x_1$ in $\productFamily{X}$ with respect to $T_0$ and
                $W_2$ is a neighborhood of $x_2$ in $\productFamily{X}$ with respect to $T_0$
                by \cref{funs_intro,indexed_product,funs_elim,function_apply_elim,ran_intro,unions_iff,neighborhood}.
            Show that $W_1$ is disjoint from $W_2$.
            \begin{subproof}
                Suppose not.
                Then there exists $p$ such that $p\in W_1$ and $p\in W_2$.
                Hence $\apply{p}{\alpha}\in U_1$ and $\apply{p}{\alpha}\in U_2$ by \cref{funs_elim,indexed_product}.
                Then $U_1$ intersects $U_2$ by \cref{inter_intro,emptyset,intersects}.
                Contradiction by \cref{disjoint}.
            \end{subproof}
        \end{byCase}
    \end{subproof}
    Hence $\productFamily{X}$ is Hausdorff with respect to $T_0$ by \cref{hausdorff}.
\end{proof}

% from §19 Item 18: product of Hausdorff spaces with the product topology
\begin{lemma}\label{product_hausdorff_product}
    Assume $X$ is a function.
    Assume $T$ is a function.
    Assume $\dom{T} = \dom{X}$.
    Assume $\dom{X}$ is inhabited.
    Assume that for all $\alpha\in\dom{X}$ we have $\apply{X}{\alpha}$ is Hausdorff with respect to $\apply{T}{\alpha}$.
    Then $\productFamily{X}$ is Hausdorff with respect to $\productTopologyIndexed{T}{X}$.
\end{lemma}
\begin{proof}
    Let $T_0 = \productTopologyIndexed{T}{X}$.
    Let $S = \productSubbasis{T}{X}$.
    For all $\alpha\in\dom{X}$ we have $\apply{T}{\alpha}$ is a topology on $\apply{X}{\alpha}$ by \cref{hausdorff}.
    $S$ is a subbasis on $\productFamily{X}$ by \cref{product_subbasis_is_subbasis}.
    Hence $\finiteIntersections{S}$ is a basis for $T_0$ on $\productFamily{X}$ by \cref{product_subbasis_finite_intersections_basis}.
    Therefore $T_0$ is a topology on $\productFamily{X}$ by \cref{basis_topology_is_topology,basis_for_topology}.
    Show that for all $x_1,x_2$ such that $x_1\in\productFamily{X}$ and $x_2\in\productFamily{X}$ and $x_1\neq x_2$
        there exist $U_1,U_2$ such that
            $U_1$ is a neighborhood of $x_1$ in $\productFamily{X}$ with respect to $T_0$ and
            $U_2$ is a neighborhood of $x_2$ in $\productFamily{X}$ with respect to $T_0$ and
            $U_1$ is disjoint from $U_2$.
    \begin{subproof}
        Fix $x_1$.
        Fix $x_2$ such that $x_1\in\productFamily{X}$ and $x_2\in\productFamily{X}$ and $x_1\neq x_2$.
        Take $w$ such that either $w\in x_1$ and $w\notin x_2$ or $w\notin x_1$ and $w\in x_2$
            by \cref{neq_witness}.
        \begin{byCase}
        \caseOf{$w\in x_1$ and $w\notin x_2$.}
            Thus $x_1$ is a function and $\dom{x_1} = \dom{X}$ and
                $x_2$ is a function and $\dom{x_2} = \dom{X}$ by \cref{indexed_product,funs_elim}.
            Take $\alpha\in\dom{X}$ such that $w = (\alpha,\apply{x_1}{\alpha})$ by \cref{funs_elim,function_member_elim}.
            Therefore $\apply{x_1}{\alpha}\neq\apply{x_2}{\alpha}$ by \cref{function_apply_elim,pair_eq_iff}.
            $\apply{x_1}{\alpha}\in\apply{X}{\alpha}$ and
                $\apply{x_2}{\alpha}\in\apply{X}{\alpha}$ by \cref{indexed_product}.
            Take $U_1,U_2$ such that
                $U_1$ is a neighborhood of $\apply{x_1}{\alpha}$ in $\apply{X}{\alpha}$ with respect to $\apply{T}{\alpha}$ and
                $U_2$ is a neighborhood of $\apply{x_2}{\alpha}$ in $\apply{X}{\alpha}$ with respect to $\apply{T}{\alpha}$ and
                $U_1$ is disjoint from $U_2$ by \cref{hausdorff}.
            Let $W_1 = \preimg{\piIndex{X}{\alpha}}{U_1}$.
            Let $W_2 = \preimg{\piIndex{X}{\alpha}}{U_2}$.
            For all $x$ such that $x\in\dom{\piIndex{X}{\alpha}}$ we have $x\in\productFamily{X}$
                by \cref{dom,projection_map_indexed,pair_eq_iff}.
            Hence $\dom{\piIndex{X}{\alpha}}\subseteq\productFamily{X}$ by \cref{subseteq}.
            For all $x$ such that $x\in\productFamily{X}$ we have $x\in\dom{\piIndex{X}{\alpha}}$ by \cref{projection_map_indexed,dom}.
            $W_1\in\pow{\productFamily{X}}$ and $W_2\in\pow{\productFamily{X}}$
                by \cref{subseteq,subseteq_antisymmetric,preimg_subseteq_dom,subseteq_transitive,powerset_intro}.
            Hence $W_1\in S$ and $W_2\in S$ by \cref{neighborhood,open_in,product_subbasis_indexed,product_subbasis_union_indexed}.
            Therefore $W_1\in\finiteIntersections{S}$ and $W_2\in\finiteIntersections{S}$ by
                \cref{subbasis_on,unions_iff,subseteq,pow_iff,singleton_subset_intro,upair_iff,pair_is_finite,subseteq_finite_is_finite,singleton_intro,inters_singleton,finite_intersections}.
            $W_1$ is open in $\productFamily{X}$ with respect to $T_0$ and
                $W_2$ is open in $\productFamily{X}$ with respect to $T_0$
                by \cref{finite_intersections,basis_elem_open_in_generated,basis_for_topology,open_in}.
            Hence $W_1$ is a neighborhood of $x_1$ in $\productFamily{X}$ with respect to $T_0$ and
                $W_2$ is a neighborhood of $x_2$ in $\productFamily{X}$ with respect to $T_0$
                by \cref{neighborhood,projection_map_indexed,preimg_iff}.
            Show that $W_1$ is disjoint from $W_2$.
            \begin{subproof}
                Suppose not.
                Then there exists $p$ such that $p\in W_1$ and $p\in W_2$.
                Hence $\apply{p}{\alpha}\in U_1$ and $\apply{p}{\alpha}\in U_2$
                    by \cref{preimg_iff,projection_map_indexed,pair_eq_iff}.
                Thus $U_1$ intersects $U_2$ by \cref{inter_intro,emptyset,intersects}.
                Contradiction by \cref{disjoint}.
            \end{subproof}
        \caseOf{$w\notin x_1$ and $w\in x_2$.}
            Thus $x_1$ is a function and $\dom{x_1} = \dom{X}$ and
                $x_2$ is a function and $\dom{x_2} = \dom{X}$ by \cref{indexed_product,funs_elim}.
            Take $\alpha\in\dom{X}$ such that $w = (\alpha,\apply{x_2}{\alpha})$ by \cref{funs_elim,function_member_elim}.
            Therefore $\apply{x_1}{\alpha}\neq\apply{x_2}{\alpha}$ by \cref{function_apply_elim,pair_eq_iff}.
            $\apply{x_1}{\alpha}\in\apply{X}{\alpha}$ and
                $\apply{x_2}{\alpha}\in\apply{X}{\alpha}$ by \cref{indexed_product}.
            Take $U_1,U_2$ such that
                $U_1$ is a neighborhood of $\apply{x_1}{\alpha}$ in $\apply{X}{\alpha}$ with respect to $\apply{T}{\alpha}$ and
                $U_2$ is a neighborhood of $\apply{x_2}{\alpha}$ in $\apply{X}{\alpha}$ with respect to $\apply{T}{\alpha}$ and
                $U_1$ is disjoint from $U_2$ by \cref{hausdorff}.
            Let $W_1 = \preimg{\piIndex{X}{\alpha}}{U_1}$.
            Let $W_2 = \preimg{\piIndex{X}{\alpha}}{U_2}$.
            For all $x$ such that $x\in\dom{\piIndex{X}{\alpha}}$ we have $x\in\productFamily{X}$
                by \cref{dom,projection_map_indexed,pair_eq_iff}.
            Hence $\dom{\piIndex{X}{\alpha}}\subseteq\productFamily{X}$ by \cref{subseteq}.
            For all $x$ such that $x\in\productFamily{X}$ we have $x\in\dom{\piIndex{X}{\alpha}}$ by \cref{projection_map_indexed,dom}.
            $W_1\in\pow{\productFamily{X}}$ and $W_2\in\pow{\productFamily{X}}$
                by \cref{subseteq,subseteq_antisymmetric,preimg_subseteq_dom,subseteq_transitive,powerset_intro}.
            Hence $W_1\in S$ and $W_2\in S$ by \cref{neighborhood,open_in,product_subbasis_indexed,product_subbasis_union_indexed}.
            Therefore $W_1\in\finiteIntersections{S}$ and $W_2\in\finiteIntersections{S}$ by
                \cref{subbasis_on,unions_iff,subseteq,pow_iff,singleton_subset_intro,upair_iff,pair_is_finite,subseteq_finite_is_finite,singleton_intro,inters_singleton,finite_intersections}.
            $W_1$ is open in $\productFamily{X}$ with respect to $T_0$ and
                $W_2$ is open in $\productFamily{X}$ with respect to $T_0$
                by \cref{finite_intersections,basis_elem_open_in_generated,basis_for_topology,open_in}.
            Hence $W_1$ is a neighborhood of $x_1$ in $\productFamily{X}$ with respect to $T_0$ and
                $W_2$ is a neighborhood of $x_2$ in $\productFamily{X}$ with respect to $T_0$
                by \cref{neighborhood,projection_map_indexed,preimg_iff}.
            Show that $W_1$ is disjoint from $W_2$.
            \begin{subproof}
                Suppose not.
                Then there exists $p$ such that $p\in W_1$ and $p\in W_2$.
                Hence $\apply{p}{\alpha}\in U_1$ and $\apply{p}{\alpha}\in U_2$
                    by \cref{preimg_iff,projection_map_indexed,pair_eq_iff}.
                Thus $U_1$ intersects $U_2$ by \cref{inter_intro,emptyset,intersects}.
                Contradiction by \cref{disjoint}.
            \end{subproof}
        \end{byCase}
    \end{subproof}
    Therefore $\productFamily{X}$ is Hausdorff with respect to $T_0$ by \cref{hausdorff}.
\end{proof}

% from §19 helper lemma 5: closure family
\begin{definition}\label{closure_family}
    $\closureFamily{A}{X}{T} =
        \{(\alpha,\closure{\apply{A}{\alpha}}{\apply{X}{\alpha}}{\apply{T}{\alpha}}) \mid \alpha\in\dom{X}\}$.
\end{definition}

% from §19 Item 19: closure of products with the box topology
\begin{lemma}\label{product_closure_box}
    Assume $X$ is a function.
    Assume $A$ is a function.
    Assume $T$ is a function.
    Assume $\dom{A} = \dom{X}$.
    Assume $\dom{T} = \dom{X}$.
    Assume that for all $\alpha\in\dom{X}$ we have $\apply{A}{\alpha}\subseteq\apply{X}{\alpha}$.
    Assume that for all $\alpha\in\dom{X}$ we have $\apply{T}{\alpha}$ is a topology on $\apply{X}{\alpha}$.
    Let $C = \closureFamily{A}{X}{T}$.
    Then $\productFamily{C} = \closure{\productFamily{A}}{\productFamily{X}}{\boxTopology{T}{X}}$.
\end{lemma}
\begin{proof}
    Let $T_0 = \boxTopology{T}{X}$.
    Let $B_0 = \boxBasis{T}{X}$.
    $B_0$ is a basis on $\productFamily{X}$ by \cref{box_basis_is_basis}.
    Hence $T_0$ is a topology on $\productFamily{X}$ by \cref{basis_topology_is_topology,box_topology}.
    $\productFamily{A}$ is a subspace of $\productFamily{X}$ with respect to $T_0$
        and $\productFamily{A}\subseteq\productFamily{X}$ by \cref{product_subspace_box,subspace_of}.
    $C$ is a relation by \cref{closure_family,relation}.
    For all $a,b_1,b_2$ such that $(a,b_1)\in C$ and $(a,b_2)\in C$ we have $b_1 = b_2$
        by \cref{closure_family,pair_eq_iff}.
    Therefore $C$ is a function by \cref{rightunique}.
    Show that for all $x$ such that $x\in\productFamily{C}$ we have
        $x\in\closure{\productFamily{A}}{\productFamily{X}}{T_0}$.
    \begin{subproof}
                Fix $x$ such that $x\in\productFamily{C}$.
                $\dom{C} = \dom{X}$ by \cref{closure_family,dom,pair_eq_iff,subseteq,subseteq_antisymmetric}.
                Thus $x$ is a function and $\dom{x} = \dom{C}$ by \cref{indexed_product,funs_elim}.
            For all $\alpha\in\dom{X}$ we have $\apply{x}{\alpha}\in\apply{X}{\alpha}$
                by \cref{indexed_product,closure_family,function_apply_intro,closure}.
            Hence $x\in\productFamily{X}$ by \cref{funs_intro,indexed_product,function_apply_elim,ran_intro,unions_iff}.
            Show that for all $B_1$ such that $B_1\in B_0$ and $x\in B_1$ we have $B_1$ intersects $\productFamily{A}$.
            \begin{subproof}
                Fix $B_1$ such that $B_1\in B_0$ and $x\in B_1$.
                Take $U$ such that
                    $U\in\funs{\dom{X}}{\pow{\unions{\ran{X}}}}$ and
                    $B_1 = \productFamily{U}$ and
                    for all $\alpha\in\dom{X}$ we have $\apply{U}{\alpha}\in\apply{T}{\alpha}$
                    by \cref{box_basis}.
                $\dom{U} = \dom{X}$ by \cref{funs_elim}.
                For all $\alpha\in\dom{X}$ we have $\apply{x}{\alpha}\in\apply{U}{\alpha}$ by \cref{indexed_product}.
                For all $\alpha\in\dom{X}$ we have $\apply{U}{\alpha}$ intersects $\apply{A}{\alpha}$
                    by \cref{open_in,closure_family,dom,indexed_product,function_apply_intro,closure_iff_intersects_open}.
                Let $R = \{w\in\dom{X}\times\unions{\ran{X}} \mid
                    \exists \alpha\in\dom{X}.\exists a\in\apply{A}{\alpha}\inter\apply{U}{\alpha}. w = (\alpha,a)\}$.
                For all $w$ such that $w\in R$ there exist $y,z$ such that $w = (y,z)$.
                Thus $R$ is a relation by \cref{relation}.
                For all $\alpha\in\dom{X}$ there exists $a$ such that $\alpha\mathrel{R} a$
                    by \cref{intersects,nonempty_inhabited,inter,elem_subseteq,function_apply_elim,ran_intro,unions_iff,times_tuple_intro}.
                Take $f$ such that $f$ is a function on $\dom{X}$ and for all $\alpha\in\dom{X}$ we have
                    $\alpha\mathrel{R}\apply{f}{\alpha}$ by \cref{choice_function}.
                Hence $f$ is a function and $\dom{f} = \dom{X}$.
                For all $\alpha\in\dom{X}$ we have $\apply{f}{\alpha}\in\apply{A}{\alpha}$ and
                    $\apply{f}{\alpha}\in\apply{U}{\alpha}$ by \cref{pair_eq_iff,inter}.
                Thus $f\in\productFamily{A}$ by \cref{funs_intro,indexed_product,function_apply_elim,ran_intro,unions_iff}.
                Therefore $B_1$ intersects $\productFamily{A}$ by \cref{funs_intro,indexed_product,funs_elim,function_apply_elim,ran_intro,unions_iff,inter_intro,emptyset,intersects}.
            \end{subproof}
            Hence $x\in\closure{\productFamily{A}}{\productFamily{X}}{T_0}$
                by \cref{box_basis_is_basis,basis_for_topology,box_topology,closure_iff_intersects_basis}.
    \end{subproof}
    Therefore $\productFamily{C}\subseteq \closure{\productFamily{A}}{\productFamily{X}}{T_0}$ by \cref{subseteq}.
    Show that for all $x$ such that $x\in\closure{\productFamily{A}}{\productFamily{X}}{T_0}$ we have
        $x\in\productFamily{C}$.
    \begin{subproof}
            Fix $x$ such that $x\in\closure{\productFamily{A}}{\productFamily{X}}{T_0}$.
            Thus $x\in\productFamily{X}$ by \cref{closure}.
            $\dom{C} = \dom{X}$ by \cref{closure_family,dom,pair_eq_iff,subseteq,subseteq_antisymmetric}.
            Show that for all $\alpha\in\dom{X}$ we have
                $\apply{x}{\alpha}\in\closure{\apply{A}{\alpha}}{\apply{X}{\alpha}}{\apply{T}{\alpha}}$.
            \begin{subproof}
                Fix $\alpha$ such that $\alpha\in\dom{X}$.
                Show that for all $V$ such that
                    $V$ is open in $\apply{X}{\alpha}$ with respect to $\apply{T}{\alpha}$ and $\apply{x}{\alpha}\in V$
                    we have $V$ intersects $\apply{A}{\alpha}$.
                \begin{subproof}
                    Fix $V$ such that $V$ is open in $\apply{X}{\alpha}$ with respect to $\apply{T}{\alpha}$ and $\apply{x}{\alpha}\in V$.
                    Take $U$ such that
                        $U\in\funs{\dom{X}}{\pow{\unions{\ran{X}}}}$ and
                        $\apply{U}{\alpha} = V$ and
                        for all $\beta\in\dom{X}$ such that $\beta\neq\alpha$ we have $\apply{U}{\beta} = \apply{X}{\beta}$
                        by \cref{open_in,topology_on,subseteq,pow_iff,function_apply_elim,ran_intro,unions_iff,subseteq_transitive,replace_coordinate_function}.
                    For all $\beta\in\dom{X}$ we have $\apply{U}{\beta}\in\apply{T}{\beta}$ by \cref{open_in,topology_on}.
                    Let $B_1 = \productFamily{U}$.
                    For all $y$ such that $y\in B_1$ we have $y\in\productFamily{X}$
                        by \cref{indexed_product,funs_elim,topology_on,elem_subseteq,pow_iff,funs_intro,function_apply_elim,ran_intro,unions_iff}.
                    Hence $B_1$ is open in $\productFamily{X}$ with respect to $T_0$ by \cref{subseteq,pow_iff,box_basis,basis_elem_open_in_generated,box_topology,open_in}.
                    For all $\beta\in\dom{X}$ we have $\apply{x}{\beta}\in\apply{U}{\beta}$ by \cref{indexed_product,subseteq}.
                    Thus $x$ is a function and $\dom{x} = \dom{X}$ and
                        $\dom{U} = \dom{X}$ by \cref{indexed_product,funs_elim}.
                    Thus $\dom{x} = \dom{U}$.
                    Thus $x\in B_1$ by \cref{funs_intro,indexed_product,funs_elim,function_apply_elim,ran_intro,unions_iff}.
                    Now for all $U$ such that $U$ is open in $\productFamily{X}$ with respect to $T_0$ and $x\in U$
                        we have $U$ intersects $\productFamily{A}$ by \cref{closure_iff_intersects_open}.
                    Hence $\productFamily{A}$ intersects $B_1$.
                    Take $f$ such that $f\in\productFamily{A}\inter B_1$ by \cref{intersects,nonempty_inhabited}.
                    Thus $f\in\productFamily{A}$ and $f\in B_1$ by \cref{inter}.
                    Thus $\apply{f}{\alpha}\in\apply{A}{\alpha}$ and $\apply{f}{\alpha}\in\apply{U}{\alpha}$ by \cref{inter,indexed_product,funs_elim}.
                    Thus $V$ intersects $\apply{A}{\alpha}$ by \cref{inter_intro,emptyset,intersects}.
                \end{subproof}
                Therefore $\apply{x}{\alpha}\in\closure{\apply{A}{\alpha}}{\apply{X}{\alpha}}{\apply{T}{\alpha}}$ by \cref{indexed_product,closure_iff_intersects_open}.
            \end{subproof}
            For all $\alpha\in\dom{X}$ we have $\apply{x}{\alpha}\in\apply{C}{\alpha}$ by \cref{closure_family,function_apply_intro}.
            Thus $x$ is a function and $\dom{x} = \dom{X}$ by \cref{indexed_product,funs_elim}.
            Therefore $x\in\productFamily{C}$ by \cref{funs_intro,indexed_product,funs_elim,closure_family,dom,pair_eq_iff,subseteq,subseteq_antisymmetric,function_apply_elim,ran_intro,unions_iff}.
    \end{subproof}
    Hence $\closure{\productFamily{A}}{\productFamily{X}}{T_0}\subseteq\productFamily{C}$ by \cref{subseteq}.
    Therefore $\productFamily{C} = \closure{\productFamily{A}}{\productFamily{X}}{T_0}$ by \cref{subseteq_antisymmetric}.
\end{proof}
% from §19 Item 20: closure of products with the product topology
\begin{lemma}\label{product_closure_product}
    Assume $X$ is a function.
    Assume $A$ is a function.
    Assume $T$ is a function.
    Assume $\dom{A} = \dom{X}$.
    Assume $\dom{T} = \dom{X}$.
    Assume $\dom{X}$ is inhabited.
    Assume that for all $\alpha\in\dom{X}$ we have $\apply{A}{\alpha}\subseteq\apply{X}{\alpha}$.
    Assume that for all $\alpha\in\dom{X}$ we have $\apply{T}{\alpha}$ is a topology on $\apply{X}{\alpha}$.
    Let $C = \closureFamily{A}{X}{T}$.
    Then $\productFamily{C} = \closure{\productFamily{A}}{\productFamily{X}}{\productTopologyIndexed{T}{X}}$.
\end{lemma}
\begin{proof}
    Let $T_0 = \productTopologyIndexed{T}{X}$.
    Let $T_1 = \boxTopology{T}{X}$.
    Let $S = \productSubbasis{T}{X}$.
    $C$ is a relation by \cref{closure_family,relation}.
    For all $a,b_1,b_2$ such that $(a,b_1)\in C$ and $(a,b_2)\in C$ we have $b_1 = b_2$
        by \cref{closure_family,pair_eq_iff}.
    Therefore $C$ is a function by \cref{rightunique}.
    $\productFamily{A}$ is a subspace of $\productFamily{X}$ with respect to $T_0$
        and $\productFamily{A}\subseteq\productFamily{X}$ by \cref{product_subspace_product,subspace_of}.
    Show that for all $x$ such that $x\in\closure{\productFamily{A}}{\productFamily{X}}{T_0}$ we have
        $x\in\productFamily{C}$.
    \begin{subproof}
            Fix $x$ such that $x\in\closure{\productFamily{A}}{\productFamily{X}}{T_0}$.
            Thus $x\in\productFamily{X}$ by \cref{closure}.
            $\dom{C} = \dom{X}$ by \cref{closure_family,dom,pair_eq_iff,subseteq,subseteq_antisymmetric}.
            Show that for all $\alpha\in\dom{X}$ we have
                $\apply{x}{\alpha}\in\closure{\apply{A}{\alpha}}{\apply{X}{\alpha}}{\apply{T}{\alpha}}$.
            \begin{subproof}
                Fix $\alpha$ such that $\alpha\in\dom{X}$.
                Show that for all $V$ such that
                    $V$ is open in $\apply{X}{\alpha}$ with respect to $\apply{T}{\alpha}$ and $\apply{x}{\alpha}\in V$
                    we have $V$ intersects $\apply{A}{\alpha}$.
                \begin{subproof}
                    Fix $V$ such that $V$ is open in $\apply{X}{\alpha}$ with respect to $\apply{T}{\alpha}$ and $\apply{x}{\alpha}\in V$.
                    Let $W = \preimg{\piIndex{X}{\alpha}}{V}$.
                    Thus $W\in\productSubbasisAt{T}{X}{\alpha}$ by
                        \cref{preimg_iff,projection_map_indexed,pair_eq_iff,subseteq,pow_iff,open_in,product_subbasis_indexed}.
                    Thus $W\in S$ by \cref{product_subbasis_indexed,product_subbasis_union_indexed}.
                    Thus $W\in\finiteIntersections{S}$ by
                        \cref{unions_iff,subseteq,pow_iff,singleton_subset_intro,upair_iff,pair_is_finite,subseteq_finite_is_finite,singleton_intro,inters_singleton,finite_intersections}.
                    Therefore $W$ is open in $\productFamily{X}$ with respect to $T_0$ by
                        \cref{product_subbasis_is_subbasis,product_subbasis_finite_intersections_basis,basis_elem_open_in_generated,basis_for_topology,basis_topology_is_topology,open_in}.
                    Thus $x\in W$ by \cref{projection_map_indexed,preimg_iff}.
                    Hence $W$ intersects $\productFamily{A}$ by \cref{basis_for_topology,basis_topology_is_topology,product_subbasis_finite_intersections_basis,subspace_of,closure_iff_intersects_open}.
                    Take $f$ such that $f\in W\inter\productFamily{A}$ by \cref{intersects,nonempty_inhabited}.
                    Thus $f\in W$ and $f\in\productFamily{A}$ by \cref{inter}.
                    Hence $\apply{f}{\alpha}\in V$ and $\apply{f}{\alpha}\in\apply{A}{\alpha}$
                        by \cref{inter,preimg_iff,projection_map_indexed,pair_eq_iff,indexed_product}.
                    Therefore $V$ intersects $\apply{A}{\alpha}$ by \cref{inter_intro,emptyset,intersects}.
                \end{subproof}
                Therefore $\apply{x}{\alpha}\in\closure{\apply{A}{\alpha}}{\apply{X}{\alpha}}{\apply{T}{\alpha}}$
                    by \cref{indexed_product,closure_iff_intersects_open}.
            \end{subproof}
            For all $\alpha\in\dom{X}$ we have $\apply{x}{\alpha}\in\apply{C}{\alpha}$ by \cref{closure_family,function_apply_intro}.
            Therefore $x\in\productFamily{C}$ by \cref{funs_intro,indexed_product,funs_elim,closure_family,dom,pair_eq_iff,subseteq,subseteq_antisymmetric,function_apply_elim,ran_intro,unions_iff}.
    \end{subproof}
    Hence $\closure{\productFamily{A}}{\productFamily{X}}{T_0}\subseteq\productFamily{C}$ by \cref{subseteq}.
    $\productFamily{C} = \closure{\productFamily{A}}{\productFamily{X}}{T_1}$ by \cref{product_closure_box}.
    Show that for all $U$ such that $U\in T_0$ we have $U\in T_1$.
    \begin{subproof}
                Fix $U$ such that $U\in T_0$.
                Then $U\in\basisTopology{\finiteIntersections{S}}{\productFamily{X}}$ by
                    \cref{product_subbasis_is_subbasis,subbasis_finite_intersections_basis,product_topology_indexed,topology_generated_by_subbasis}.
                Let $M = \{B\in\finiteIntersections{S} \mid B\subseteq U\}$.
                For all $x$ such that $x\in U$ there exists $B\in\finiteIntersections{S}$ such that $x\in B$ and $B\subseteq U$
                    by \cref{topology_generated_by_basis}.
                Thus for all $x$ such that $x\in U$ we have $x\in\unions{M}$ by \cref{unions_iff}.
                Therefore $U = \unions{M}$ by \cref{unions_iff,subseteq,subseteq_antisymmetric}.
                Show that for all $B$ such that $B\in\finiteIntersections{S}$ we have $B\in T_1$.
                \begin{subproof}
                    Fix $B$ such that $B\in\finiteIntersections{S}$.
                    Take $F\subseteq S$ such that $F$ is finite and inhabited and $B = \inters{F}$ by \cref{finite_intersections}.
                    Show that for all $A$ such that $A\in F$ we have $A\in T_1$.
                    \begin{subproof}
                        Fix $A$ such that $A\in F$.
                        Take $\beta\in\dom{X}$ such that $A\in\productSubbasisAt{T}{X}{\beta}$ by \cref{subseteq,product_subbasis_union_indexed}.
                        Take $V\in\apply{T}{\beta}$ such that $A = \preimg{\piIndex{X}{\beta}}{V}$ by \cref{product_subbasis_indexed}.
                        Show that for all $x$ such that $x\in A$ there exists $B_1\in\boxBasis{T}{X}$ such that $x\in B_1$ and $B_1\subseteq A$.
                        \begin{subproof}
                            Fix $x$ such that $x\in A$.
                            Take $U$ such that
                                $U\in\funs{\dom{X}}{\pow{\unions{\ran{X}}}}$ and
                                $\apply{U}{\beta} = V$ and
                                for all $\gamma\in\dom{X}$ such that $\gamma\neq\beta$ we have $\apply{U}{\gamma} = \apply{X}{\gamma}$
                                by \cref{open_in,topology_on,subseteq,pow_iff,function_apply_elim,ran_intro,unions_iff,subseteq_transitive,replace_coordinate_function}.
                            For all $\gamma\in\dom{X}$ we have $\apply{U}{\gamma}\in\apply{T}{\gamma}$ by \cref{open_in,topology_on}.
                            Let $B_1 = \productFamily{U}$.
                            For all $z$ such that $z\in B_1$ we have $z\in\productFamily{X}$
                                by \cref{indexed_product,funs_elim,topology_on,elem_subseteq,pow_iff,funs_intro,function_apply_elim,ran_intro,unions_iff}.
                            Thus $B_1\in\pow{\productFamily{X}}$ by \cref{subseteq,pow_iff}.
                            $B_1\in\boxBasis{T}{X}$ by \cref{box_basis}.
                            Thus $x\in\productFamily{X}$ and $\apply{x}{\beta}\in V$ by \cref{preimg_iff,projection_map_indexed,pair_eq_iff}.
                            For all $\gamma\in\dom{X}$ we have $\apply{x}{\gamma}\in\apply{U}{\gamma}$ by \cref{indexed_product,subseteq}.
                            Thus $x\in\funs{\dom{X}}{\unions{\ran{X}}}$ by \cref{indexed_product}.
                            Thus $x$ is a function and $\dom{x} = \dom{X}$ and
                                $\dom{U} = \dom{X}$ by \cref{funs_elim}.
                            Thus $\dom{x} = \dom{U}$.
                            For all $\gamma\in\dom{U}$ we have $\apply{x}{\gamma}\in\unions{\ran{U}}$
                                by \cref{funs_elim,function_apply_elim,ran_intro,unions_iff}.
                            Thus $x\in\funs{\dom{U}}{\unions{\ran{U}}}$ by \cref{funs_intro}.
                            Thus $x\in B_1$ by \cref{indexed_product}.
                            Show that for all $z$ such that $z\in B_1$ we have $z\in A$.
                            \begin{subproof}
                                Fix $z$ such that $z\in B_1$.
                                Thus $z\in\productFamily{U}$.
                                Thus $\beta\in\dom{U}$ by \cref{funs_elim}.
                                Thus $\apply{z}{\beta}\in\apply{U}{\beta}$ by \cref{indexed_product}.
                                Thus $\apply{z}{\beta}\in V$.
                                Thus $z\in\productFamily{X}$ by \cref{indexed_product,funs_elim,topology_on,elem_subseteq,pow_iff,funs_intro,function_apply_elim,ran_intro,unions_iff}.
                                Thus $z\in A$ by \cref{projection_map_indexed,preimg_iff}.
                            \end{subproof}
                            Thus $B_1\subseteq A$ by \cref{subseteq}.
                        \end{subproof}
                        Therefore $A\in T_1$ by
                            \cref{preimg_iff,projection_map_indexed,pair_eq_iff,subseteq,pow_iff,topology_generated_by_basis,box_topology}.
                    \end{subproof}
                    Thus $F\subseteq T_1$ by \cref{subseteq}.
                    Hence $B\in T_1$ by
                        \cref{box_basis_is_basis,basis_topology_is_topology,box_topology,finite_intersections_open_in_topology}.
                \end{subproof}
                Therefore $U\in T_1$ by
                    \cref{subseteq,box_basis_is_basis,basis_topology_is_topology,box_topology,topology_on}.
    \end{subproof}
    Hence $T_0\subseteq T_1$ by \cref{subseteq}.
    Show that for all $x$ such that $x\in\closure{\productFamily{A}}{\productFamily{X}}{T_1}$ we have
        $x\in\closure{\productFamily{A}}{\productFamily{X}}{T_0}$.
    \begin{subproof}
            Fix $x$ such that $x\in\closure{\productFamily{A}}{\productFamily{X}}{T_1}$.
            Show that for all $U$ such that
                $U$ is open in $\productFamily{X}$ with respect to $T_0$ and $x\in U$
                we have $U$ intersects $\productFamily{A}$.
            \begin{subproof}
                Fix $U$ such that $U$ is open in $\productFamily{X}$ with respect to $T_0$ and $x\in U$.
                Hence $U$ is open in $\productFamily{X}$ with respect to $T_1$
                    by \cref{subseteq,open_in,box_basis_is_basis,basis_topology_is_topology,box_topology}.
                Therefore $U$ intersects $\productFamily{A}$ by \cref{open_in,topology_on,subseteq,pow_iff,elem_subseteq,subspace_of,closure_iff_intersects_open}.
            \end{subproof}
            Hence $\finiteIntersections{S}$ is a basis for $T_0$ on $\productFamily{X}$
                by \cref{product_subbasis_is_subbasis,product_subbasis_finite_intersections_basis}.
            Therefore $T_0$ is a topology on $\productFamily{X}$ by \cref{basis_topology_is_topology,basis_for_topology}.
            $\productFamily{A}\subseteq\productFamily{X}$ by \cref{subspace_of}.
            $\boxBasis{T}{X}$ is a basis on $\productFamily{X}$ and
                $T_1$ is a topology on $\productFamily{X}$ by \cref{box_basis_is_basis,basis_topology_is_topology,box_topology}.
            $\productFamily{X}$ is closed in $\productFamily{X}$ with respect to $T_1$ by \cref{closed_whole_space}.
            Thus $x\in\productFamily{X}$ by \cref{closure_subseteq_closed,elem_subseteq}.
            Hence $x\in\closure{\productFamily{A}}{\productFamily{X}}{T_0}$ by \cref{closure_iff_intersects_open}.
    \end{subproof}
    Hence $\closure{\productFamily{A}}{\productFamily{X}}{T_1}\subseteq\closure{\productFamily{A}}{\productFamily{X}}{T_0}$ by \cref{subseteq}.
    Therefore $\productFamily{C}\subseteq\closure{\productFamily{A}}{\productFamily{X}}{T_0}$ by \cref{subseteq}.
    Therefore $\productFamily{C} = \closure{\productFamily{A}}{\productFamily{X}}{T_0}$ by \cref{subseteq_antisymmetric}.
\end{proof}

% from §19 helper lemma 3: projection map is a function
\begin{lemma}\label{projection_map_indexed_function}
    Assume $X$ is a function.
    Assume $\beta\in\dom{X}$.
    Then $\piIndex{X}{\beta}\in\funs{\productFamily{X}}{\apply{X}{\beta}}$.
\end{lemma}
\begin{proof}
    Let $P = \piIndex{X}{\beta}$.
    $P$ is a relation by \cref{projection_map_indexed,relation}.
    For all $x,y_1,y_2$ such that $(x,y_1)\in P$ and $(x,y_2)\in P$ we have $y_1 = y_2$
        by \cref{projection_map_indexed,pair_eq_iff}.
    Hence $P\in\funs{\productFamily{X}}{\apply{X}{\beta}}$
        by \cref{rightunique,function_apply_elim,projection_map_indexed,dom,pair_eq_iff,indexed_product,subseteq,subseteq_antisymmetric,funs_intro}.
\end{proof}

% from §19 helper lemma 4: image of a finite set under a function is finite
\begin{lemma}\label{finite_image_function}
    Assume $F$ is finite.
    Assume $g$ is a function on $F$.
    Then $\img{g}{F}$ is finite.
\end{lemma}
\begin{proof}
    Let $B = \img{g}{F}$.
    Let $R = \{w\in B\times F \mid \exists y\in B.\exists x\in F. w = (y,x) \land \apply{g}{x} = y\}$.
    $R$ is a relation by \cref{relation}.
    For all $y\in B$ there exists $x$ such that $y\mathrel{R} x$.
    Take $s$ such that $s$ is a function on $B$ and for all $y\in B$ we have $y\mathrel{R}\apply{s}{y}$ by \cref{choice_function}.
    Hence $s$ is a function and $\dom{s} = B$.
    For all $y\in B$ we have $\apply{s}{y}\in F$ and $\apply{g}{\apply{s}{y}} = y$ by \cref{pair_eq_iff}.
    Let $C = \ran{s}$.
    For all $y_1,y_2$ such that $y_1\in B$ and $y_2\in B$ and $\apply{s}{y_1} = \apply{s}{y_2}$ we have $y_1 = y_2$
        by \cref{subseteq}.
    Therefore $s$ is injective.
    For all $x$ such that $x\in C$ we have $x\in F$.
    Hence $C$ is finite by \cref{subseteq,subseteq_finite_is_finite}.
    Therefore $s\in\bijections{B}{C}$ by \cref{function_apply_elim,ran_intro,funs_intro,funs_surj_iff,bijections}.
    Hence $\converse{s}$ is a bijection from $C$ to $B$ by \cref{bijection_converse_is_bijection}.
    Therefore $B$ is finite by \cref{finite_bijection}.
\end{proof}

% from §19 Item 21: continuity into an indexed product
\begin{lemma}\label{continuous_map_into_indexed_product}
    Assume $X$ is a function.
    Assume $T$ is a function.
    Assume $\dom{T} = \dom{X}$.
    Assume $\dom{X}$ is inhabited.
    Assume that for all $\alpha\in\dom{X}$ we have $\apply{T}{\alpha}$ is a topology on $\apply{X}{\alpha}$.
    Assume $T_A$ is a topology on $A$.
    Assume $f\in\funs{A}{\productFamily{X}}$.
    Then $f$ is continuous from $A$ to $\productFamily{X}$ with respect to $T_A$ and $\productTopologyIndexed{T}{X}$ iff
        for all $\alpha\in\dom{X}$ we have $\piIndex{X}{\alpha}\circ f$ is continuous from $A$ to $\apply{X}{\alpha}$ with respect to $T_A$ and $\apply{T}{\alpha}$.
\end{lemma}
\begin{proof}
    Let $T_0 = \productTopologyIndexed{T}{X}$.
    Let $S = \productSubbasis{T}{X}$.
    For all $\alpha\in\dom{X}$ we have $\apply{T}{\alpha}$ is a topology on $\apply{X}{\alpha}$ by \cref{subseteq}.
    Show that if $f$ is continuous from $A$ to $\productFamily{X}$ with respect to $T_A$ and $T_0$ then
        for all $\alpha\in\dom{X}$ we have $\piIndex{X}{\alpha}\circ f$ is continuous from $A$ to $\apply{X}{\alpha}$
        with respect to $T_A$ and $\apply{T}{\alpha}$.
    \begin{subproof}
        Assume $f$ is continuous from $A$ to $\productFamily{X}$ with respect to $T_A$ and $T_0$.
        Fix $\alpha$ such that $\alpha\in\dom{X}$.
        Let $g = \piIndex{X}{\alpha}\circ f$.
        $\piIndex{X}{\alpha}\in\funs{\productFamily{X}}{\apply{X}{\alpha}}$ and
            $g\in\funs{A}{\apply{X}{\alpha}}$ by \cref{projection_map_indexed_function,funs_circ,continuous}.
        Show that for all $V$ such that $V$ is open in $\apply{X}{\alpha}$ with respect to $\apply{T}{\alpha}$
            we have $\preimg{g}{V}$ is open in $A$ with respect to $T_A$.
        \begin{subproof}
            Fix $V$ such that $V$ is open in $\apply{X}{\alpha}$ with respect to $\apply{T}{\alpha}$.
            Let $W = \preimg{\piIndex{X}{\alpha}}{V}$.
            Thus $W\in\pow{\productFamily{X}}$
                by \cref{projection_map_indexed_function,funs_elim,preimg_subseteq_dom,subseteq_transitive,pow_iff}.
            $V\in\apply{T}{\alpha}$ and $W\in\productSubbasisAt{T}{X}{\alpha}$ and $W\in S$
                by \cref{open_in,product_subbasis_indexed,product_subbasis_union_indexed}.
            Let $F = \{W\}$.
            $F\subseteq S$ and $F$ is finite and inhabited and $\inters{F} = W$ by
                \cref{singleton_subset_intro,upair_iff,pair_is_finite,subseteq_finite_is_finite,singleton_intro,inters_singleton}.
            Therefore $W\in\finiteIntersections{S}$ by \cref{unions_iff,subseteq,pow_iff,finite_intersections}.
            $W$ is open in $\productFamily{X}$ with respect to $T_0$ by
                \cref{product_subbasis_is_subbasis,product_subbasis_finite_intersections_basis,basis_elem_open_in_generated,basis_for_topology,basis_topology_is_topology,open_in}.
            Hence $\preimg{g}{V}$ is open in $A$ with respect to $T_A$ by
                \cref{continuous,circ_iff,preimg_iff,subseteq,subseteq_antisymmetric}.
        \end{subproof}
        Therefore $g$ is continuous from $A$ to $\apply{X}{\alpha}$ with respect to $T_A$ and $\apply{T}{\alpha}$ by \cref{continuous}.
    \end{subproof}
    Show that if for all $\alpha\in\dom{X}$ we have $\piIndex{X}{\alpha}\circ f$ is continuous from $A$ to $\apply{X}{\alpha}$
        with respect to $T_A$ and $\apply{T}{\alpha}$ then $f$ is continuous from $A$ to $\productFamily{X}$ with respect to $T_A$ and $T_0$.
    \begin{subproof}
        Assume for all $\alpha\in\dom{X}$ we have $\piIndex{X}{\alpha}\circ f$ is continuous from $A$ to $\apply{X}{\alpha}$
            with respect to $T_A$ and $\apply{T}{\alpha}$.
        Show that for all $U$ such that $U$ is open in $\productFamily{X}$ with respect to $T_0$
            we have $\preimg{f}{U}$ is open in $A$ with respect to $T_A$.
        \begin{subproof}
            Fix $U$ such that $U$ is open in $\productFamily{X}$ with respect to $T_0$.
            Take $M\subseteq \finiteIntersections{S}$ such that $U = \unions{M}$ by
                \cref{open_in,product_subbasis_is_subbasis,product_subbasis_finite_intersections_basis,basis_topology_unions}.
            Let $N = \{\preimg{f}{B} \mid B\in M\}$.
            Show that for all $C$ such that $C\in N$ we have $C$ is open in $A$ with respect to $T_A$.
            \begin{subproof}
                Fix $C$ such that $C\in N$.
                Take $B\in M$ such that $C = \preimg{f}{B}$.
                Take $F\subseteq S$ such that $F$ is finite and inhabited and $B = \inters{F}$ by \cref{subseteq,finite_intersections}.
                Let $H = \{\preimg{f}{A} \mid A\in F\}$.
                Show that for all $D$ such that $D\in H$ we have $D$ is open in $A$ with respect to $T_A$.
                \begin{subproof}
                    Fix $D$ such that $D\in H$.
                    Take $A\in F$ such that $D = \preimg{f}{A}$.
                    Take $\beta\in\dom{X}$ such that $A\in\productSubbasisAt{T}{X}{\beta}$ by \cref{subseteq,product_subbasis_union_indexed}.
                    Take $V\in\apply{T}{\beta}$ such that $A = \preimg{\piIndex{X}{\beta}}{V}$ by \cref{product_subbasis_indexed}.
                    Let $g = \piIndex{X}{\beta}\circ f$.
                    Therefore $D$ is open in $A$ with respect to $T_A$ by
                        \cref{subseteq,open_in,continuous,circ_iff,preimg_iff,subseteq_antisymmetric}.
                \end{subproof}
                Show that $H$ is finite.
                \begin{subproof}
                    Let $h = \{(A,\preimg{f}{A}) \mid A\in F\}$.
                    $h$ is a relation by \cref{relation}.
                    For all $x,y_1,y_2$ such that $(x,y_1)\in h$ and $(x,y_2)\in h$ we have $y_1 = y_2$
                        by \cref{pair_eq_iff}.
                    Hence $h$ is right-unique by \cref{rightunique}.
                    For all $x$ such that $x\in\dom{h}$ we have $x\in F$.
                    For all $A$ such that $A\in F$ we have $A\in\dom{h}$ by \cref{dom_iff}.
                    Thus $\dom{h} = F$ by \cref{subseteq,subseteq_antisymmetric}.
                    Therefore $h$ is a function on $F$.
                    For all $D$ such that $D\in H$ we have $D\in\img{h}{F}$.
                    Therefore $H$ is finite by \cref{finite_image_function,subseteq,subseteq_finite_is_finite}.
                \end{subproof}
                Take $A$ such that $A\in F$.
                Hence $H$ is inhabited and for all $D$ such that $D\in H$ we have $D\in T_A$
                    by \cref{open_in}.
                Therefore $\inters{H}$ is open in $A$ with respect to $T_A$ by
                    \cref{subseteq,topology_on,finite_intersections_open_in_topology,open_in}.
                Show that $C = \inters{H}$.
                \begin{subproof}
                    Show that for all $a$ such that $a\in C$ we have $a\in\inters{H}$.
                    \begin{subproof}
                        Fix $a$ such that $a\in C$.
                        Take $z\in B$ such that $a\mathrel{f} z$ by \cref{preimg_iff}.
                        Thus for all $A\in F$ we have $z\in A$ by \cref{subseteq,inters_iff_forall}.
                        For all $D\in H$ we have $a\in D$.
                        Thus $a\in\inters{H}$ by \cref{inters_iff_forall}.
                    \end{subproof}
                    Hence $C\subseteq\inters{H}$ by \cref{subseteq}.
                    Show that for all $a$ such that $a\in\inters{H}$ we have $a\in C$.
                    \begin{subproof}
                        Fix $a$ such that $a\in\inters{H}$.
                        For all $A\in F$ we have $a\in\preimg{f}{A}$ by \cref{img_iff,inters_iff_forall}.
                        Show that $a\in\preimg{f}{B}$.
                        \begin{subproof}
                            Take $y$ such that $a\mathrel{f} y$ by \cref{preimg_iff,preimg_subseteq_dom}.
                            Show that for all $A\in F$ we have $y\in A$.
                            \begin{subproof}
                                Fix $A$ such that $A\in F$.
                                Take $y_A\in A$ such that $a\mathrel{f} y_A$ by \cref{preimg_iff}.
                                Thus $y = y_A$ by \cref{continuous,funs_is_function,function_rightunique,rightunique}.
                                Thus $y\in A$.
                            \end{subproof}
                            Thus $a\in\preimg{f}{B}$ by \cref{inters_iff_forall,subseteq,preimg_iff}.
                        \end{subproof}
                        Thus $a\in C$.
                    \end{subproof}
                    $\inters{H}\subseteq C$ by \cref{subseteq}.
                    Therefore $C = \inters{H}$ by \cref{subseteq_antisymmetric}.
                \end{subproof}
                Hence $C$ is open in $A$ with respect to $T_A$.
            \end{subproof}
            Therefore $N\subseteq T_A$ by \cref{subseteq,open_in}.
            $T_A$ is a topology on $A$ and $\unions{N}\in T_A$ by \cref{continuous,topology_on}.
            Therefore $\preimg{f}{U} = \unions{N}$ by
                \cref{preimg_iff,unions_iff,subseteq,subseteq_antisymmetric}.
            Hence $\preimg{f}{U}$ is open in $A$ with respect to $T_A$ by \cref{open_in}.
        \end{subproof}
        Therefore $f$ is continuous from $A$ to $\productFamily{X}$ with respect to $T_A$ and $T_0$ by
            \cref{product_subbasis_is_subbasis,product_subbasis_finite_intersections_basis,basis_topology_is_topology,basis_for_topology,continuous}.
    \end{subproof}
\end{proof}
\section{§20 The Metric Topology}

% from §20 helper lemma 1: metric nonnegative
\begin{definition}\label{metric_nonnegative}
    $d$ is nonnegative on $X$ iff
    for all $x,y\in X$ we have $d((x,y)) \geq 0$.
\end{definition}

% from §20 helper lemma 2: metric zero characterization
\begin{definition}\label{metric_zero_characterization}
    $d$ is zero-characterized on $X$ iff
    for all $x,y\in X$ we have $d((x,y)) = 0$ iff $x = y$.
\end{definition}

% from §20 helper lemma 3: metric symmetry
\begin{definition}\label{metric_symmetric}
    $d$ is symmetric on $X$ iff
    for all $x,y\in X$ we have $d((x,y)) = d((y,x))$.
\end{definition}

% from §20 helper lemma 4: metric triangle inequality
\begin{definition}\label{metric_triangle_inequality}
    $d$ satisfies the triangle inequality on $X$ iff
    for all $x,y,z\in X$ we have $d((x,y)) + d((y,z)) \geq d((x,z))$.
\end{definition}

% from §20 Item 1: metric on a set
\begin{definition}\label{metric_on}
    $d$ is a metric on $X$ iff
    $d\in\funs{X\times X}{\reals}$ and
    $d$ is nonnegative on $X$ and
    $d$ is zero-characterized on $X$ and
    $d$ is symmetric on $X$ and
    $d$ satisfies the triangle inequality on $X$.
\end{definition}

% from §20 Item 2: metric ball
\begin{definition}\label{metric_ball}
    $\metricBall{d}{X}{x}{\epsilon} = \{y\in X \mid d((x,y)) < \epsilon\}$.
\end{definition}

% from §20 Item 3: metric basis
\begin{definition}\label{metric_basis}
    $\metricBasis{d}{X} = \{B\in\pow{X} \mid \exists x\in X. \exists \epsilon\in\reals. 0 < \epsilon \land B = \metricBall{d}{X}{x}{\epsilon}\}$.
\end{definition}

% from §20 Item 4: metric topology induced by a metric
\begin{definition}\label{metric_topology}
    $\metricTopology{d}{X} = \basisTopology{\metricBasis{d}{X}}{X}$.
\end{definition}

% from §20 Item 5: metric basis is a basis
\begin{lemma}\label{metric_basis_is_basis}
    Assume $d$ is a metric on $X$.
    Then $\metricBasis{d}{X}$ is a basis on $X$.
\end{lemma}
\begin{proof}
    Let $B = \metricBasis{d}{X}$.
    Hence $B\subseteq\pow{X}$ by \cref{metric_basis,metric_ball,subseteq,powerset_intro}.
    For all $x$ such that $x\in X$ there exists $B_1\in B$ such that $x\in B_1$
        by \cref{real_naturals_subset_reals,real_one_in_naturals,subseteq,real_zero_lt_one,metric_on,metric_zero_characterization,real_lt_subst_left,metric_ball,powerset_intro,metric_basis}.
    Show that for all $B_1,B_2,x$ such that $B_1\in B$ and $B_2\in B$ and $x\in B_1\inter B_2$
        there exists $B_3\in B$ such that $x\in B_3$ and $B_3\subseteq B_1\inter B_2$.
    \begin{subproof}
        Fix $B_1$.
        Fix $B_2$.
        Fix $x$ such that $B_1\in B$ and $B_2\in B$ and $x\in B_1\inter B_2$.
        Take $x_1,\epsilon_1$ such that $x_1\in X$ and $\epsilon_1\in\reals$ and $0 < \epsilon_1$ and $B_1 = \metricBall{d}{X}{x_1}{\epsilon_1}$
            by \cref{metric_basis}.
        Take $x_2,\epsilon_2$ such that $x_2\in X$ and $\epsilon_2\in\reals$ and $0 < \epsilon_2$ and $B_2 = \metricBall{d}{X}{x_2}{\epsilon_2}$
            by \cref{metric_basis}.
        Thus $x\in X$ and $d((x_1,x)) < \epsilon_1$ and $d((x_2,x)) < \epsilon_2$
            by \cref{inter_elim_left,inter_elim_right,metric_ball}.
        $d\in\funs{X\times X}{\reals}$ and $d((x_1,x))\in\reals$ and $d((x_2,x))\in\reals$
            by \cref{metric_on,times_tuple_intro,funs_type_apply}.
        Let $r_1 = \epsilon_1 - d((x_1,x))$.
        Let $r_2 = \epsilon_2 - d((x_2,x))$.
        $0 < r_1$ and $0 < r_2$ by \cref{metric_on,times_tuple_intro,funs_type_apply,real_lt_imp_pos_diff}.
        Let $A = \{r_1, r_2\}$.
        $A$ is finite and $A\neq\emptyset$ by \cref{pair_is_finite,upair_intro_left,emptyset}.
        For all $y$ such that $y\in A$ we have $y\in\reals$.
        Therefore $A\subseteq\reals$ by \cref{subseteq}.
        Take $m\in A$ such that $m$ is a minimum of $A$ by \cref{finite_real_minimum}.
        Thus $0 < m$ by \cref{upair_elim}.
        Let $B_3 = \metricBall{d}{X}{x}{m}$.
        Hence $B_3\in B$ and $x\in B_3$ by
            \cref{subseteq,metric_ball,powerset_intro,metric_basis,metric_on,metric_zero_characterization,real_lt_subst_left}.
        Show that for all $z$ such that $z\in B_3$ we have $z\in B_1\inter B_2$.
        \begin{subproof}
            Fix $z$ such that $z\in B_3$.
                Thus $z\in X$ and $d((x,z)) < m$ by \cref{metric_ball}.
                Show that $z\in B_1$.
                \begin{subproof}
                    Thus $d((x,z)) < r_1$ by
                        \cref{real_minimum,upair_intro_left,times_tuple_intro,funs_type_apply,subseteq,real_sub,real_add_closed,real_leq_split,real_lt_trans,real_lt_subst_right}.
                    Thus $d((x,z)) + d((x_1,x)) < \epsilon_1$ by
                        \cref{times_tuple_intro,funs_type_apply,real_sub,real_add_closed,real_lt_add,real_add_inverse,real_add_assoc,real_add_comm,real_add_ident,real_lt_subst_right}.
                    Thus $d((x_1,x)) + d((x,z)) < \epsilon_1$ by
                        \cref{funs_type_apply,times_tuple_intro,real_add_comm,real_lt_subst_left}.
                    Thus $d((x_1,z)) \leq d((x_1,x)) + d((x,z))$ by \cref{metric_on,metric_triangle_inequality}.
                    Thus $d((x_1,z)) < \epsilon_1$ by
                        \cref{times_tuple_intro,funs_type_apply,real_add_closed,real_leq_split,real_lt_trans,real_lt_subst_left}.
                    Thus $z\in B_1$ by \cref{metric_ball}.
                \end{subproof}
                Show that $z\in B_2$.
                \begin{subproof}
                    Thus $d((x,z)) < r_2$ by
                        \cref{real_minimum,upair_intro_right,times_tuple_intro,funs_type_apply,subseteq,real_sub,real_add_inverse,real_add_closed,real_leq_split,real_lt_trans,real_lt_subst_right}.
                    Thus $d((x,z)) + d((x_2,x)) < \epsilon_2$ by
                        \cref{times_tuple_intro,funs_type_apply,real_sub,real_add_inverse,real_add_closed,real_lt_add,real_add_assoc,real_add_comm,real_add_ident,real_lt_subst_right}.
                    Thus $d((x_2,x)) + d((x,z)) < \epsilon_2$ by
                        \cref{funs_type_apply,times_tuple_intro,real_add_comm,real_lt_subst_left}.
                    Thus $d((x_2,z)) \leq d((x_2,x)) + d((x,z))$ by \cref{metric_on,metric_triangle_inequality}.
                    Thus $d((x_2,z)) < \epsilon_2$ by
                        \cref{times_tuple_intro,funs_type_apply,real_add_closed,real_leq_split,real_lt_trans,real_lt_subst_left}.
                    Thus $z\in B_2$ by \cref{metric_ball}.
                \end{subproof}
                Thus $z\in B_1\inter B_2$ by \cref{inter_intro}.
        \end{subproof}
        Hence $B_3\subseteq B_1\inter B_2$ by \cref{subseteq}.
    \end{subproof}
    Therefore $B$ is a basis on $X$ by \cref{basis_on}.
\end{proof}

% from §20 Item 6: metric open sets
\begin{lemma}\label{metric_open_iff}
    Assume $d$ is a metric on $X$.
    Let $T = \metricTopology{d}{X}$.
    Then for all $U$ such that $U\subseteq X$ we have
    $U$ is open in $X$ with respect to $T$ iff
    for all $y\in U$ there exists $\delta\in\reals$ such that $0 < \delta$ and $\metricBall{d}{X}{y}{\delta}\subseteq U$.
\end{lemma}
\begin{proof}
    Let $B = \metricBasis{d}{X}$.
    $T = \basisTopology{B}{X}$ by \cref{metric_topology}.
    Fix $U$ such that $U\subseteq X$.
    Show that if $U$ is open in $X$ with respect to $T$ then
        for all $y\in U$ there exists $\delta\in\reals$ such that $0 < \delta$ and $\metricBall{d}{X}{y}{\delta}\subseteq U$.
    \begin{subproof}
        Assume $U$ is open in $X$ with respect to $T$.
        Fix $y$ such that $y\in U$.
        There exists $B_1\in B$ such that $y\in B_1$ and $B_1\subseteq U$ by \cref{open_in,topology_generated_by_basis}.
        Take $x,\epsilon$ such that $x\in X$ and $\epsilon\in\reals$ and $0 < \epsilon$ and
            $B_1 = \metricBall{d}{X}{x}{\epsilon}$ by \cref{metric_basis}.
        Thus $y\in X$ and $d((x,y)) < \epsilon$ by \cref{metric_ball}.
        $d\in\funs{X\times X}{\reals}$ and $d((x,y))\in\reals$
            by \cref{metric_on,times_tuple_intro,funs_type_apply}.
        Let $\delta = \epsilon - d((x,y))$.
        $0 < \delta$ by \cref{metric_on,times_tuple_intro,funs_type_apply,real_lt_imp_pos_diff}.
        Show that for all $z$ such that $z\in\metricBall{d}{X}{y}{\delta}$ we have $z\in B_1$.
        \begin{subproof}
            Fix $z$ such that $z\in\metricBall{d}{X}{y}{\delta}$.
            Thus $z\in X$ and $d((y,z)) < \delta$ by \cref{metric_ball}.
            Thus $d((x,y)) + d((y,z)) < \epsilon$ by
                \cref{times_tuple_intro,funs_type_apply,real_sub,real_add_inverse,real_add_closed,real_lt_subst_right,real_lt_add,real_add_assoc,real_add_comm,real_add_ident,real_lt_subst_left}.
            Thus $d((x,z)) \leq d((x,y)) + d((y,z))$ by \cref{metric_on,metric_triangle_inequality}.
            Thus $d((x,z)) < \epsilon$ by
                \cref{times_tuple_intro,funs_type_apply,real_add_closed,real_leq_split,real_lt_trans,real_lt_subst_left}.
            Thus $z\in B_1$ by \cref{metric_ball}.
        \end{subproof}
        Therefore $\metricBall{d}{X}{y}{\delta}\subseteq U$ by \cref{subseteq,subseteq_transitive}.
    \end{subproof}
    Show that if for all $y\in U$ there exists $\delta\in\reals$ such that $0 < \delta$ and $\metricBall{d}{X}{y}{\delta}\subseteq U$
        then $U$ is open in $X$ with respect to $T$.
    \begin{subproof}
        Assume for all $y\in U$ there exists $\delta\in\reals$ such that $0 < \delta$ and $\metricBall{d}{X}{y}{\delta}\subseteq U$.
        For all $y\in U$ there exists $B_1\in B$ such that $y\in B_1$ and $B_1\subseteq U$
            by \cref{subseteq,metric_on,metric_zero_characterization,real_lt_subst_left,metric_ball,powerset_intro,metric_basis}.
        Therefore $U$ is open in $X$ with respect to $T$ by
            \cref{metric_basis_is_basis,basis_topology_is_topology,powerset_intro,topology_generated_by_basis,open_in}.
    \end{subproof}
\end{proof}

% from §20 Item 7: discrete metric
\begin{definition}\label{discrete_metric}
    $\discreteMetric{X} = \{w \mid \exists p\in X\times X. \exists z\in\reals.
        w = (p,z) \land
        ((\fst{p} = \snd{p} \land z = 0) \lor (\fst{p}\neq \snd{p} \land z = 1))\}$.
\end{definition}
% from §20 Item 8: discrete metric is a metric
\begin{lemma}\label{discrete_metric_is_metric}
    Then $\discreteMetric{X}$ is a metric on $X$.
\end{lemma}
\begin{proof}
    Let $d = \discreteMetric{X}$.
    For all $p,z_1,z_2$ such that $(p,z_1)\in d$ and $(p,z_2)\in d$ we have $z_1 = z_2$
        by \cref{discrete_metric,pair_eq_iff}.
    Hence $d$ is right-unique by \cref{rightunique}.
    Then $d$ is a relation by \cref{discrete_metric,relation}.
    $d$ is a function.
    For all $p$ such that $p\in\dom{d}$ we have $d(p)\in\reals$
        by \cref{dom_iff,discrete_metric,pair_eq_iff,function_apply_intro}.
    Hence $d$ is a function to $\reals$.
    Then $\dom{d}\subseteq X\times X$ by \cref{dom_iff,discrete_metric,pair_eq_iff,subseteq}.
    For all $p$ such that $p\in X\times X$ we have $p\in\dom{d}$
        by \cref{real_add_ident,real_naturals_subset_reals,real_one_in_naturals,subseteq,discrete_metric,dom_intro}.
    Therefore $d\in\funs{X\times X}{\reals}$ by \cref{subseteq,subseteq_antisymmetric,funs_intro}.
    Show that for all $x,y$ such that $x\in X$ and $y\in X$ we have $d((x,y)) \geq 0$.
    \begin{subproof}
        Fix $x$.
        Fix $y$ such that $x\in X$ and $y\in X$.
        \begin{byCase}
        \caseOf{$x = y$.}
            Thus $d((x,y)) \geq 0$ by \cref{times_tuple_intro,real_add_ident,fst_eq,snd_eq,discrete_metric,function_apply_intro,funs_elim}.
        \caseOf{$x \neq y$.}
            Thus $d((x,y)) \geq 0$ by
                \cref{times_tuple_intro,real_naturals_subset_reals,real_one_in_naturals,subseteq,fst_eq,snd_eq,discrete_metric,function_apply_intro,funs_elim,real_zero_lt_one,real_lt_subst_right}.
        \end{byCase}
    \end{subproof}
    Hence $d$ is nonnegative on $X$ by \cref{metric_nonnegative}.
    Show that for all $x,y$ such that $x\in X$ and $y\in X$ we have $d((x,y)) = 0$ iff $x = y$.
    \begin{subproof}
        Fix $x$.
        Fix $y$ such that $x\in X$ and $y\in X$.
        Show that if $d((x,y)) = 0$ then $x = y$.
        \begin{subproof}
            Assume $d((x,y)) = 0$.
            Thus $((x,y),0)\in d$ by \cref{funs_elim,times_tuple_intro,subseteq,function_apply_iff}.
            Take $p, z$ such that
                $p\in X\times X$ and $z\in\reals$ and
                $((x,y),0) = (p,z)$ and
                $((\fst{p} = \snd{p} \land z = 0) \lor (\fst{p}\neq \snd{p} \land z = 1))$
                by \cref{discrete_metric}.
            $p = (x,y)$ and $z = 0$ by \cref{pair_eq_iff}.
            Hence $((\fst{p} = \snd{p} \land 0 = 0) \lor (\fst{p} \neq \snd{p} \land 0 = 1))$.
            Hence $\fst{p} = \snd{p}$.
            Hence $x = y$ by \cref{fst_eq,snd_eq}.
        \end{subproof}
        If $x = y$ then $d((x,y)) = 0$
            by \cref{times_tuple_intro,real_add_ident,fst_eq,snd_eq,discrete_metric,function_apply_intro,funs_elim}.
    \end{subproof}
    Therefore $d$ is zero-characterized on $X$ by \cref{metric_zero_characterization}.
    Show that for all $x,y$ such that $x\in X$ and $y\in X$ we have $d((x,y)) = d((y,x))$.
    \begin{subproof}
        Fix $x$.
        Fix $y$ such that $x\in X$ and $y\in X$.
        \begin{byCase}
        \caseOf{$x = y$.}
            Thus $d((x,y)) = d((y,x))$ by
                \cref{times_tuple_intro,real_add_ident,fst_eq,snd_eq,discrete_metric,function_apply_intro,funs_elim}.
        \caseOf{$x \neq y$.}
            Then $y \neq x$.
            Thus $d((x,y)) = d((y,x))$ by
                \cref{times_tuple_intro,real_naturals_subset_reals,real_one_in_naturals,subseteq,fst_eq,snd_eq,discrete_metric,function_apply_intro,funs_elim}.
        \end{byCase}
    \end{subproof}
    Hence $d$ is symmetric on $X$ by \cref{metric_symmetric}.
    Show that for all $x,y,z$ such that $x\in X$ and $y\in X$ and $z\in X$ we have
        $d((x,y)) + d((y,z)) \geq d((x,z))$.
    \begin{subproof}
        Fix $x$.
        Fix $y$.
        Fix $z$ such that $x\in X$ and $y\in X$ and $z\in X$.
        \begin{byCase}
        \caseOf{$x = z$.}
            Then $d((x,z)) = 0$ by \cref{times_tuple_intro,real_add_ident,fst_eq,snd_eq,discrete_metric,function_apply_intro,funs_elim}.
            $0\in\reals$ by \cref{real_add_ident}.
            Hence $d((x,y))\in\reals$ and $d((y,z))\in\reals$ by \cref{funs_type_apply,times_tuple_intro}.
            Therefore $d((x,y)) + d((y,z))\in\reals$ by \cref{real_add_closed}.
            Now $0 \leq d((x,y))$ and $0 \leq d((y,z))$ by \cref{metric_nonnegative}.
            Thus $0 + d((y,z)) \leq d((x,y)) + d((y,z))$ by \cref{real_leq_add}.
            Hence $0 \leq d((x,y)) + d((y,z))$ by \cref{real_add_ident,real_leq_trans}.
            Then $d((x,z)) \leq 0$.
            Therefore $d((x,z)) \leq d((x,y)) + d((y,z))$ by \cref{real_leq_trans}.
            Hence $d((x,y)) + d((y,z)) \geq d((x,z))$.
        \caseOf{$x \neq z$.}
            Thus $d((x,z)) = 1$ by
                \cref{times_tuple_intro,real_naturals_subset_reals,real_one_in_naturals,subseteq,fst_eq,snd_eq,discrete_metric,function_apply_intro,funs_elim}.
            Show that $d((x,y)) + d((y,z)) \geq 1$.
            \begin{subproof}
                \begin{byCase}
                \caseOf{$x = y$.}
                    Thus $d((y,z)) = 1$ by
                        \cref{times_tuple_intro,real_naturals_subset_reals,real_one_in_naturals,subseteq,fst_eq,snd_eq,discrete_metric,function_apply_intro,funs_elim}.
                    Thus $d((x,y))\in\reals$ by \cref{funs_type_apply,times_tuple_intro}.
                    Thus $0 \leq d((x,y))$ by \cref{metric_nonnegative}.
                    Thus $d((x,y)) + d((y,z)) \geq 1$ by
                        \cref{real_add_ident,real_naturals_subset_reals,real_one_in_naturals,subseteq,real_leq_add,real_add_comm,real_leq_trans}.
                \caseOf{$x \neq y$.}
                    Thus $d((x,y)) = 1$ by
                        \cref{times_tuple_intro,real_naturals_subset_reals,real_one_in_naturals,subseteq,fst_eq,snd_eq,discrete_metric,function_apply_intro,funs_elim}.
                    Thus $d((y,z))\in\reals$ by \cref{funs_type_apply,times_tuple_intro}.
                    Thus $0 \leq d((y,z))$ by \cref{metric_nonnegative}.
                    Thus $d((x,y)) + d((y,z)) \geq 1$ by
                        \cref{real_add_ident,real_naturals_subset_reals,real_one_in_naturals,subseteq,real_leq_add,real_add_comm,real_leq_trans}.
                \end{byCase}
            \end{subproof}
            Hence $1 \leq d((x,y)) + d((y,z))$.
            Then $d((x,z)) \leq 1$.
            Therefore $d((x,z)) \leq d((x,y)) + d((y,z))$ by \cref{real_leq_trans}.
            Hence $d((x,y)) + d((y,z)) \geq d((x,z))$.
        \end{byCase}
    \end{subproof}
    Hence $d$ is a metric on $X$ by \cref{metric_triangle_inequality,metric_on}.
\end{proof}
% from §20 Item 9: discrete metric topology
\begin{lemma}\label{discrete_metric_topology}
    Then $\metricTopology{\discreteMetric{X}}{X} = \discrete{X}$.
\end{lemma}
\begin{proof}
    Let $d = \discreteMetric{X}$.
    Let $T = \metricTopology{d}{X}$.
    $d$ is a metric on $X$ by \cref{discrete_metric_is_metric}.
    Thus $\metricBasis{d}{X}$ is a basis on $X$ and $T$ is a topology on $X$ and $T\subseteq\discrete{X}$
        by \cref{metric_basis_is_basis,basis_topology_is_topology,metric_topology,topology_on,discrete_topology}.
    Show that for all $U$ such that $U\in\discrete{X}$ we have $U\in T$.
    \begin{subproof}
        Fix $U$ such that $U\in\discrete{X}$.
        Thus $U\subseteq X$ by \cref{discrete_topology,powerset_elim}.
        Show that for all $y\in U$ there exists $\delta\in\reals$ such that
            $0 < \delta$ and $\metricBall{d}{X}{y}{\delta}\subseteq U$.
        \begin{subproof}
            Fix $y$ such that $y\in U$.
            $1\in\reals$ and $0 < 1$ by \cref{real_naturals_subset_reals,real_one_in_naturals,subseteq,real_zero_lt_one}.
            Let $\delta = 1$.
            Show that for all $z$ such that $z\in\metricBall{d}{X}{y}{\delta}$ we have $z\in U$.
            \begin{subproof}
                Fix $z$ such that $z\in\metricBall{d}{X}{y}{\delta}$.
                Thus $z\in X$ and $d((y,z)) < 1$ by \cref{metric_ball}.
                Show that $z = y$.
                \begin{subproof}
                    Suppose not.
                    Thus $d((y,z)) = 1$ by
                        \cref{subseteq,times_tuple_intro,fst_eq,snd_eq,discrete_metric,function_apply_intro,funs_elim,metric_on}.
                    Contradiction by \cref{real_lt_subst_right,real_lt_irreflexive}.
                \end{subproof}
                Hence $z\in U$.
            \end{subproof}
            Therefore $\metricBall{d}{X}{y}{\delta}\subseteq U$ by \cref{subseteq}.
        \end{subproof}
        Hence $U\in T$ by \cref{discrete_topology,powerset_elim,metric_open_iff,open_in}.
    \end{subproof}
    Therefore $T = \discrete{X}$ by \cref{subseteq,subseteq_antisymmetric}.
\end{proof}

% from §20 Item 10: standard metric on $\reals$
\begin{definition}\label{standard_metric_r}
    $\standardMetricR = \{w \mid \exists p\in\reals\times\reals. \exists z\in\reals.
        w = (p,z) \land z = \abs{\fst{p} - \snd{p}}\}$.
\end{definition}

% from §20 Item 11: standard metric is a metric
\begin{lemma}\label{standard_metric_is_metric}
    Then $\standardMetricR$ is a metric on $\reals$.
\end{lemma}
\begin{proof}
    Let $d = \standardMetricR$.
    For all $p,z_1,z_2$ such that $(p,z_1)\in d$ and $(p,z_2)\in d$ we have $z_1 = z_2$
        by \cref{standard_metric_r,pair_eq_iff}.
    Hence $d$ is right-unique by \cref{rightunique}.
    Then $d$ is a relation by \cref{standard_metric_r,relation}.
    $d$ is a function.
    For all $p$ such that $p\in\dom{d}$ we have $d(p)\in\reals$
        by \cref{dom_iff,standard_metric_r,pair_eq_iff,function_apply_intro}.
    Hence $d$ is a function to $\reals$.
    Then $\dom{d}\subseteq\reals\times\reals$
        by \cref{dom_iff,standard_metric_r,pair_eq_iff,subseteq}.
    For all $p$ such that $p\in\reals\times\reals$ we have $p\in\dom{d}$
        by \cref{standard_metric_r,dom_intro,times_elem_is_tuple,fst_eq,snd_eq,real_sub,real_add_inverse,real_add_closed,real_abs_in_reals}.
    Hence $d\in\funs{\reals\times\reals}{\reals}$ by \cref{subseteq,subseteq_antisymmetric,funs_intro}.
    For all $x,y$ such that $x\in\reals$ and $y\in\reals$ we have $d((x,y)) \geq 0$ by
        \cref{times_tuple_intro,standard_metric_r,function_apply_intro,funs_elim,fst_eq,snd_eq,real_sub,real_add_inverse,real_add_closed,real_abs_in_reals,real_abs_nonnegative}.
    Hence $d$ is nonnegative on $\reals$ by \cref{metric_nonnegative}.
    Show that for all $x,y$ such that $x\in\reals$ and $y\in\reals$ we have $d((x,y)) = 0$ iff $x = y$.
    \begin{subproof}
        Fix $x$.
        Fix $y$ such that $x\in\reals$ and $y\in\reals$.
        If $d((x,y)) = 0$ then $x = y$ by
            \cref{times_tuple_intro,standard_metric_r,function_apply_intro,funs_elim,fst_eq,snd_eq,real_sub,real_add_inverse,real_add_closed,real_abs_in_reals,real_abs_eq_zero,real_add_cancel}.
        If $x = y$ then $d((x,y)) = 0$ by
            \cref{times_tuple_intro,standard_metric_r,function_apply_intro,funs_elim,fst_eq,snd_eq,real_sub,real_add_inverse,real_add_ident,real_add_closed,real_abs_in_reals,real_abs_nonneg}.
    \end{subproof}
    Therefore $d$ is zero-characterized on $\reals$ by \cref{metric_zero_characterization}.
    Show that for all $x,y$ such that $x\in\reals$ and $y\in\reals$ we have $d((x,y)) = d((y,x))$.
    \begin{subproof}
        Fix $x$.
        Fix $y$ such that $x\in\reals$ and $y\in\reals$.
        Let $p_1 = (x,y)$.
        Let $p_2 = (y,x)$.
        $p_1\in\reals\times\reals$ and $p_2\in\reals\times\reals$ by \cref{times_tuple_intro}.
        Let $z_1 = \abs{\fst{p_1} - \snd{p_1}}$.
        Let $z_2 = \abs{\fst{p_2} - \snd{p_2}}$.
        Thus $z_1\in\reals$ and $z_2\in\reals$
            by \cref{fst_eq,snd_eq,real_sub,real_add_inverse,real_add_closed,real_abs_in_reals}.
        Let $w_1 = (p_1,z_1)$.
        Let $w_2 = (p_2,z_2)$.
        $w_1\in d$ and $w_2\in d$ by \cref{standard_metric_r}.
        Thus $d((x,y)) = \abs{(x - y)}$ and $d((y,x)) = \abs{(y - x)}$
            by \cref{function_apply_intro,funs_elim,fst_eq,snd_eq}.
        Thus $d((x,y)) = d((y,x))$ by \cref{real_abs_sub_swap}.
    \end{subproof}
    Hence $d$ is symmetric on $\reals$ by \cref{metric_symmetric}.
    Show that for all $x,y,z$ such that $x\in\reals$ and $y\in\reals$ and $z\in\reals$ we have
        $d((x,y)) + d((y,z)) \geq d((x,z))$.
    \begin{subproof}
        Fix $x$.
        Fix $y$.
        Fix $z$ such that $x\in\reals$ and $y\in\reals$ and $z\in\reals$.
        $\neg{y}\in\reals$ and $\neg{z}\in\reals$ by \cref{real_add_inverse}.
        $x - y = x + \neg{y}$ by \cref{real_sub}.
        $y - z = y + \neg{z}$ by \cref{real_sub}.
        $x - z = x + \neg{z}$ by \cref{real_sub}.
        Thus $\abs{(x - y)}\in\reals$ and $\abs{(y - z)}\in\reals$ and $\abs{(x - z)}\in\reals$
            by \cref{real_sub,real_add_inverse,real_add_closed,real_abs_in_reals}.
        Let $p_1 = (x,y)$.
        Let $p_2 = (y,z)$.
        Let $p_3 = (x,z)$.
        $p_1\in\reals\times\reals$ and $p_2\in\reals\times\reals$ and $p_3\in\reals\times\reals$
            by \cref{times_tuple_intro}.
        $\fst{p_1} = x$ and $\snd{p_1} = y$ and $\fst{p_2} = y$ and $\snd{p_2} = z$ and
            $\fst{p_3} = x$ and $\snd{p_3} = z$ by \cref{fst_eq,snd_eq}.
        Thus $\fst{p_1} - \snd{p_1} = x - y$.
        Thus $\fst{p_2} - \snd{p_2} = y - z$.
        Thus $\fst{p_3} - \snd{p_3} = x - z$.
        Thus $\fst{p_1} - \snd{p_1}\in\reals$.
        Thus $\fst{p_2} - \snd{p_2}\in\reals$.
        Thus $\fst{p_3} - \snd{p_3}\in\reals$.
        Let $z_1 = \abs{\fst{p_1} - \snd{p_1}}$.
        Let $z_2 = \abs{\fst{p_2} - \snd{p_2}}$.
        Let $z_3 = \abs{\fst{p_3} - \snd{p_3}}$.
        Thus $z_1\in\reals$ and $z_2\in\reals$ and $z_3\in\reals$
            by \cref{fst_eq,snd_eq,real_sub,real_add_inverse,real_add_closed,real_abs_in_reals}.
        Let $w_1 = (p_1,z_1)$.
        Let $w_2 = (p_2,z_2)$.
        Let $w_3 = (p_3,z_3)$.
        $w_1\in d$ and $w_2\in d$ and $w_3\in d$ by \cref{standard_metric_r}.
        Thus $d((x,y)) = \abs{(x - y)}$ and $d((y,z)) = \abs{(y - z)}$ and $d((x,z)) = \abs{(x - z)}$
            by \cref{function_apply_intro,funs_elim,fst_eq,snd_eq}.
        $x - y = x + \neg{y}$ and $y - z = y + \neg{z}$ and
            $(x - y) + (y - z) = (x + \neg{y}) + (y + \neg{z})$ by \cref{real_sub}.
        $(x + \neg{y}) + (y + \neg{z}) = x + (\neg{y} + (y + \neg{z}))$ and
            $\neg{y} + (y + \neg{z}) = (\neg{y} + y) + \neg{z}$ by \cref{real_add_assoc}.
        $\neg{y} + y = y + \neg{y}$ by \cref{real_add_comm}.
        $y + \neg{y} = 0$ by \cref{real_add_inverse}.
        Thus $\neg{y} + (y + \neg{z}) = 0 + \neg{z}$.
        $0 + \neg{z} = \neg{z}$ by \cref{real_add_ident}.
        Thus $(x - y) + (y - z) = x + \neg{z}$.
        $x + \neg{z} = x - z$ by \cref{real_sub}.
        Thus $\abs{((x - y) + (y - z))} = \abs{(x - z)}$.
        Thus $d((x,z)) \leq \abs{(x - y)} + \abs{(y - z)}$ by \cref{real_triangle_inequality}.
        Thus $d((x,z)) \leq d((x,y)) + d((y,z))$.
        Thus $d((x,y)) + d((y,z)) \geq d((x,z))$.
    \end{subproof}
    Hence $d$ is a metric on $\reals$ by \cref{metric_triangle_inequality,metric_on}.
\end{proof}

% from §20 helper lemma 6: evaluation of the standard metric
\begin{lemma}\label{standard_metric_apply}
    Assume $x\in\reals$ and $y\in\reals$ and $d = \standardMetricR$.
    Then $d((x,y)) = \abs{(x - y)}$.
\end{lemma}
\begin{proof}
    Let $p = (x,y)$.
    $(p,\abs{\fst{p} - \snd{p}})\in d$ and $d((x,y)) = \abs{(x - y)}$
        by \cref{fst_eq,snd_eq,real_sub,real_add_inverse,real_add_closed,real_abs_in_reals,times_tuple_intro,standard_metric_r,standard_metric_is_metric,metric_on,function_apply_intro,funs_elim}.
\end{proof}

% from §20 Item 12: standard metric topology
\begin{lemma}\label{standard_metric_topology}
    Then $\metricTopology{\standardMetricR}{\reals} = \standardTopologyR$.
\end{lemma}
\begin{proof}
    Let $d = \standardMetricR$.
    Let $T = \metricTopology{d}{\reals}$.
    Let $T_1 = \standardTopologyR$.
    Let $B = \metricBasis{d}{\reals}$.
    Let $B_1 = \standardBasisR$.
    $d$ is a metric on $\reals$ by \cref{standard_metric_is_metric}.
    Hence $B$ is a basis for $T$ on $\reals$ by \cref{metric_basis_is_basis,basis_for_topology,metric_topology}.
    $B_1$ is a basis for $T_1$ on $\reals$ by \cref{standard_basis_is_basis,basis_for_topology,standard_topology_reals}.
    Show that for all $x,B_2$ such that $x\in\reals$ and $B_2\in B_1$ and $x\in B_2$
        there exists $B_3\in B$ such that $x\in B_3$ and $B_3\subseteq B_2$.
    \begin{subproof}
            Fix $x$.
            Fix $B_2$ such that $x\in\reals$ and $B_2\in B_1$ and $x\in B_2$.
            Take $a,b\in\reals$ such that $a < b$ and $B_2 = \intervalopen{a}{b}$ by \cref{standard_basis_reals}.
            Thus $a < x$ and $x < b$ by \cref{intervalopen}.
            Let $r_1 = x - a$.
            Let $r_2 = b - x$.
            $0 < r_1$ and $0 < r_2$ by \cref{real_lt_imp_pos_diff}.
            Let $A = \{r_1, r_2\}$.
            $A$ is finite and $A\neq\emptyset$ by \cref{pair_is_finite,upair_intro_left,emptyset}.
            For all $y$ such that $y\in A$ we have $y\in\reals$.
            Thus $A\subseteq\reals$ by \cref{subseteq}.
            Take $m\in A$ such that $m$ is a minimum of $A$ by \cref{finite_real_minimum}.
            Thus $0 < m$ by \cref{upair_elim,real_lt_subst_right}.
            Let $B_3 = \metricBall{d}{\reals}{x}{m}$.
            Hence $B_3\in B$ and $x\in B_3$ by
                \cref{metric_ball,subseteq,powerset_intro,metric_basis,metric_on,metric_zero_characterization,real_lt_subst_left}.
            Show that for all $z$ such that $z\in B_3$ we have $z\in B_2$.
            \begin{subproof}
                    Fix $z$ such that $z\in B_3$.
                    Thus $z\in\reals$ and $d((x,z)) < m$ by \cref{metric_ball}.
                    Thus $\abs{(x - z)} < m$ by \cref{standard_metric_apply,real_lt_subst_left}.
                    Thus $x - z\in\reals$ by \cref{real_sub,real_add_inverse,real_add_closed}.
                    Thus $\abs{(x - z)}\in\reals$ by \cref{real_sub,real_add_inverse,real_add_closed,real_abs_in_reals}.
                    $r_1\in A$ by \cref{upair_intro_left}.
                    Thus $m < r_1$ or $m = r_1$ by \cref{real_minimum,subseteq,real_leq_split}.
                    Thus $\abs{(x - z)} < r_1$ by \cref{real_lt_subst_right,real_lt_trans}.
                    $x - z \leq \abs{(x - z)}$ by \cref{real_leq_abs}.
                    Thus $x - z < \abs{(x - z)}$ or $x - z = \abs{(x - z)}$ by \cref{real_leq_split}.
                    Thus $x - z < r_1$ by \cref{real_lt_subst_left,real_lt_trans}.
                        Thus $x - z < x - a$ by \cref{real_lt_subst_right}.
                        $x - a\in\reals$ by \cref{real_sub,real_add_closed}.
                        $\neg{(x - a)} < \neg{(x - z)}$ by \cref{real_lt_neg}.
                        $\neg{(x - a)} = a - x$ and $\neg{(x - z)} = z - x$ by \cref{real_sub_swap_neg}.
                        Thus $a - x < z - x$ by \cref{real_lt_subst_left,real_lt_subst_right}.
                        $\neg{x}\in\reals$ by \cref{real_add_inverse}.
                        $a - x = a + \neg{x}$ and $z - x = z + \neg{x}$ by \cref{real_sub}.
                        Thus $a - x\in\reals$ by \cref{real_add_closed}.
                        Thus $z - x\in\reals$ by \cref{real_add_closed}.
                        $(a - x) + x < (z - x) + x$ by \cref{real_lt_add}.
                        $\neg{x} + x = 0$ by \cref{real_add_comm,real_add_inverse}.
                        Thus $(a - x) + x = a + 0$ and $(z - x) + x = z + 0$
                            by \cref{real_add_assoc,real_add_eq_subst}.
                        $a + 0 = a$ and $z + 0 = z$ by \cref{real_add_ident,real_add_comm}.
                        Thus $(a - x) + x = a$ and $(z - x) + x = z$.
                        Thus $a < z$.
                    $r_2\in A$ by \cref{upair_intro_right}.
                    Thus $m < r_2$ or $m = r_2$ by \cref{real_minimum,subseteq,real_leq_split}.
                    Thus $\abs{(x - z)} < r_2$ by \cref{real_lt_subst_right,real_lt_trans}.
                        $\neg{(x - z)}\in\reals$ by \cref{real_add_inverse}.
                        $\neg{(x - z)} \leq \abs{(x - z)}$ by \cref{real_neg_leq_abs}.
                        Thus $\neg{(x - z)} < \abs{(x - z)}$ or $\neg{(x - z)} = \abs{(x - z)}$ by \cref{real_leq_split}.
                        Thus $\neg{(x - z)} < r_2$ by \cref{real_lt_subst_left,real_lt_trans}.
                        $\neg{(x - z)} = z - x$ by \cref{real_sub_swap_neg}.
                        Thus $z - x < r_2$ by \cref{real_lt_subst_left}.
                        Thus $z - x < b - x$ by \cref{real_lt_subst_right}.
                        $z - x\in\reals$ and $b - x\in\reals$ by \cref{real_sub,real_add_closed}.
                        $(z - x) + x < (b - x) + x$ by \cref{real_lt_add}.
                        $z - x = z + \neg{x}$ and $b - x = b + \neg{x}$ by \cref{real_sub}.
                        $\neg{x} + x = 0$ by \cref{real_add_comm,real_add_inverse}.
                        Thus $(z - x) + x = z + 0$ and $(b - x) + x = b + 0$
                            by \cref{real_add_assoc,real_add_eq_subst}.
                        $z + 0 = z$ and $b + 0 = b$ by \cref{real_add_ident,real_add_comm}.
                        Thus $(z - x) + x = z$ and $(b - x) + x = b$.
                        Thus $z < b$.
                    Thus $z\in B_2$ by \cref{intervalopen}.
                \end{subproof}
            Therefore $B_3\subseteq B_2$ by \cref{subseteq}.
    \end{subproof}
    Hence $T$ is finer than $T_1$ on $\reals$ by \cref{basis_finer_criterion}.
    Show that for all $x,B_2$ such that $x\in\reals$ and $B_2\in B$ and $x\in B_2$
        there exists $B_3\in B_1$ such that $x\in B_3$ and $B_3\subseteq B_2$.
    \begin{subproof}
            Fix $x$.
            Fix $B_2$ such that $x\in\reals$ and $B_2\in B$ and $x\in B_2$.
            Take $x_0,\epsilon\in\reals$ such that $0 < \epsilon$ and $B_2 = \metricBall{d}{\reals}{x_0}{\epsilon}$
                by \cref{metric_basis}.
            Thus $d((x_0,x)) < \epsilon$ by \cref{metric_ball}.
            $d\in\funs{\reals\times\reals}{\reals}$ by \cref{metric_on}.
            Thus $d((x_0,x))\in\reals$ by \cref{times_tuple_intro,funs_type_apply}.
            $\neg{d((x_0,x))}\in\reals$ by \cref{real_add_inverse}.
            Let $\delta = \epsilon - d((x_0,x))$.
            Thus $\delta\in\reals$ by \cref{real_sub,real_add_inverse,real_add_closed}.
            $0 < \delta$ by \cref{real_lt_imp_pos_diff}.
            $\neg{\delta}\in\reals$ by \cref{real_add_inverse}.
            $x - \delta = x + \neg{\delta}$ by \cref{real_sub}.
            $0\in\reals$ by \cref{real_add_ident}.
            Thus $\neg{\delta} < \neg{0}$ by \cref{real_lt_neg}.
            $\neg{0} = 0$ by \cref{real_add_inverse_unique,real_add_ident}.
            Thus $\neg{\delta} < 0$ by \cref{real_lt_subst_right}.
            $\neg{\delta} + x < 0 + x$ by \cref{real_lt_add}.
            Thus $x + \neg{\delta} < 0 + x$ by \cref{real_add_comm,real_lt_subst_left}.
            Thus $x - \delta < 0 + x$ by \cref{real_lt_subst_left,real_sub}.
            $0 + x = x$ by \cref{real_add_ident}.
            Thus $x - \delta < x$ by \cref{real_lt_subst_right}.
            $0 + x < \delta + x$ by \cref{real_lt_add}.
            $\delta + x = x + \delta$ by \cref{real_add_comm}.
            Thus $x < x + \delta$ by \cref{real_lt_subst_left,real_lt_subst_right}.
            Let $B_3 = \intervalopen{x - \delta}{x + \delta}$.
            Hence $x - \delta\in\reals$ and $x + \delta\in\reals$ by \cref{real_sub,real_add_inverse,real_add_closed}.
            Therefore $x - \delta < x + \delta$ by \cref{real_lt_trans}.
            Hence $B_3\in B_1$ by \cref{intervalopen_subset_reals,powerset_intro,standard_basis_reals}.
            Then $x\in B_3$ by \cref{intervalopen}.
            Show that for all $z$ such that $z\in B_3$ we have $z\in B_2$.
                \begin{subproof}
                    Fix $z$ such that $z\in B_3$.
                    Thus $x - \delta < z$ and $z < x + \delta$ and $z\in\reals$ by \cref{intervalopen}.
                    Thus $x - z\in\reals$ and $z - x\in\reals$ by \cref{real_sub,real_add_inverse,real_add_closed}.
                    $\neg{x}\in\reals$ by \cref{real_add_inverse}.
                    $z + \neg{x} < (x + \delta) + \neg{x}$ by \cref{real_lt_add}.
                    $z + \neg{x} = z - x$ by \cref{real_sub}.
                    $(x + \delta) + \neg{x} = x + (\delta + \neg{x})$ by \cref{real_add_assoc}.
                    $\delta + \neg{x} = \neg{x} + \delta$ by \cref{real_add_comm}.
                    $x + (\neg{x} + \delta) = (x + \neg{x}) + \delta$ by \cref{real_add_assoc}.
                    $x + \neg{x} = 0$ by \cref{real_add_inverse}.
                    Thus $(x + \delta) + \neg{x} = 0 + \delta$ by \cref{real_add_assoc,real_add_comm,real_add_inverse}.
                    $0 + \delta = \delta$ by \cref{real_add_ident}.
                    Thus $z - x < \delta$ by \cref{real_lt_subst_left,real_lt_subst_right}.
                    $(x - \delta) + \neg{x} < z + \neg{x}$ by \cref{real_lt_add}.
                    $(x - \delta) + \neg{x} = \neg{\delta}$ by \cref{real_sub,real_add_assoc,real_add_inverse,real_add_comm,real_add_ident}.
                    Thus $\neg{\delta} < z - x$ by \cref{real_lt_subst_left,real_lt_subst_right}.
                    Thus $\neg{(z - x)} < \neg{(\neg{\delta})}$ by \cref{real_lt_neg}.
                    $\neg{(z - x)} = x - z$ by \cref{real_sub_swap_neg}.
                    $\delta = \neg{(\neg{\delta})}$ by \cref{real_neg_neg}.
                    Thus $x - z < \delta$ by \cref{real_lt_subst_left,real_lt_subst_right}.
                    Thus $x - z \leq \delta$ and $\neg{(x - z)} \leq \delta$.
                    Thus $\abs{(x - z)} \leq \delta$ by \cref{real_abs_leq_if_pm_leq}.
                    Thus $\abs{(x - z)}\in\reals$ by \cref{real_abs_in_reals}.
                    Thus $\abs{(x - z)} < \delta$.
                    Thus $d((x,z)) = \abs{(x - z)}$ by \cref{standard_metric_apply}.
                    Thus $d((x,z)) < \delta$.
                    Thus $d((x_0,z)) \leq d((x_0,x)) + d((x,z))$ by \cref{metric_on,metric_triangle_inequality}.
                    Thus $d((x_0,x)) + d((x,z))\in\reals$ and $d((x_0,x)) + \delta\in\reals$
                        by \cref{funs_type_apply,real_add_closed}.
                    $d((x,z)) + d((x_0,x)) < \delta + d((x_0,x))$ by \cref{real_lt_add}.
                    $d((x,z)) + d((x_0,x)) = d((x_0,x)) + d((x,z))$ and
                        $\delta + d((x_0,x)) = d((x_0,x)) + \delta$ by \cref{real_add_comm}.
                    Thus $d((x_0,x)) + d((x,z)) < d((x_0,x)) + \delta$
                        by \cref{real_lt_subst_left,real_lt_subst_right}.
                    $(x_0,z)\in\reals\times\reals$ by \cref{times_tuple_intro}.
                    Thus $d((x_0,z))\in\reals$ by \cref{funs_type_apply}.
                    Thus $d((x_0,z)) < d((x_0,x)) + d((x,z))$ or
                        $d((x_0,z)) = d((x_0,x)) + d((x,z))$
                        by \cref{times_tuple_intro,funs_type_apply,real_leq_split}.
                    Thus $d((x_0,z)) < d((x_0,x)) + \delta$ by \cref{real_lt_subst_left,real_lt_trans}.
                    $\delta = \epsilon + \neg{d((x_0,x))}$ by \cref{real_sub}.
                    $d((x_0,x)) + \delta = \delta + d((x_0,x))$ by \cref{real_add_comm}.
                    Thus $d((x_0,x)) + \delta = (\epsilon + \neg{d((x_0,x))}) + d((x_0,x))$
                        by \cref{real_add_eq_subst}.
                    $(\epsilon + \neg{d((x_0,x))}) + d((x_0,x)) =
                        \epsilon + (\neg{d((x_0,x))} + d((x_0,x)))$ by \cref{real_add_assoc}.
                    $\neg{d((x_0,x))} + d((x_0,x)) = 0$ by \cref{real_add_comm,real_add_inverse}.
                    Thus $d((x_0,x)) + \delta = \epsilon + 0$ by \cref{real_add_eq_subst}.
                    $\epsilon + 0 = \epsilon$ by \cref{real_add_ident,real_add_comm}.
                    Thus $d((x_0,x)) + \delta = \epsilon$.
                    Thus $d((x_0,z)) < \epsilon$ by \cref{real_lt_subst_right}.
                    Thus $z\in B_2$ by \cref{metric_ball}.
                \end{subproof}
            Therefore $B_3\subseteq B_2$ by \cref{subseteq}.
    \end{subproof}
    Hence $T_1$ is finer than $T$ on $\reals$ by \cref{basis_finer_criterion}.
    Therefore $T = T_1$ by \cref{finer_topology,subseteq_antisymmetric}.
\end{proof}

% from §20 Item 13: metrizable space
\begin{definition}\label{metrizable}
    $X$ is metrizable with respect to $T$ iff
    $T$ is a topology on $X$ and
    there exists $d$ such that $d$ is a metric on $X$ and $T = \metricTopology{d}{X}$.
\end{definition}

% from §20 Item 14: metric space
\begin{definition}\label{metric_space}
    $X$ is a metric space with metric $d$ and topology $T$ iff
    $d$ is a metric on $X$ and $T = \metricTopology{d}{X}$.
\end{definition}

% from §20 Item 15: bounded subset
\begin{definition}\label{metric_bounded}
    $A$ is bounded in $X$ with respect to $d$ iff
    $d$ is a metric on $X$ and
    $A\subseteq X$ and
    there exists $M\in\reals$ such that for all $a_1,a_2\in A$ we have $d((a_1,a_2)) \leq M$.
\end{definition}

% from §20 helper lemma 5: distance set of a subset
\begin{definition}\label{metric_distance_set}
    $\metricDistanceSet{A}{d} = \{d((a_1,a_2)) \mid a_1\in A, a_2\in A\}$.
\end{definition}

% from §20 Item 16: diameter of a set
\begin{definition}\label{metric_diameter}
    $r$ is a diameter of $A$ in $X$ with respect to $d$ iff
    $A$ is bounded in $X$ with respect to $d$ and
    $A\neq\emptyset$ and
    $r$ is a supremum of $\metricDistanceSet{A}{d}$.
\end{definition}

% from §20 Item 17: standard bounded metric
\begin{definition}\label{bounded_metric}
    $\boundedMetric{d} = \{w \mid \exists p\in\dom{d}. \exists z\in\reals.
        w = (p,z) \land ((z = d(p) \land d(p) < 1) \lor (z = 1 \land 1 \leq d(p)))\}$.
\end{definition}

% from §20 Item 18: standard bounded metric is a metric
\begin{theorem}\label{bounded_metric_is_metric}
    Assume $d$ is a metric on $X$.
    Then $\boundedMetric{d}$ is a metric on $X$.
\end{theorem}
\begin{proof}
    Let $e = \boundedMetric{d}$.
    Show that $e\in\funs{X\times X}{\reals}$.
    \begin{subproof}
        Show that for all $p,z_1,z_2$ such that $(p,z_1)\in e$ and $(p,z_2)\in e$ we have $z_1 = z_2$.
        \begin{subproof}
                    Fix $p$.
                    Fix $z_1$.
                    Fix $z_2$ such that $(p,z_1)\in e$ and $(p,z_2)\in e$.
                    Take $p_1, w_1$ such that
                        $p_1\in\dom{d}$ and $w_1\in\reals$ and
                        $(p,z_1) = (p_1,w_1)$ and
                        $((w_1 = d(p_1) \land d(p_1) < 1) \lor (w_1 = 1 \land 1 \leq d(p_1)))$
                        by \cref{bounded_metric}.
                    Take $p_2, w_2$ such that
                        $p_2\in\dom{d}$ and $w_2\in\reals$ and
                        $(p,z_2) = (p_2,w_2)$ and
                        $((w_2 = d(p_2) \land d(p_2) < 1) \lor (w_2 = 1 \land 1 \leq d(p_2)))$
                        by \cref{bounded_metric}.
                    Hence $((z_1 = d(p) \land d(p) < 1) \lor (z_1 = 1 \land 1 \leq d(p)))$
                        and $((z_2 = d(p) \land d(p) < 1) \lor (z_2 = 1 \land 1 \leq d(p)))$
                        by \cref{bounded_metric,pair_eq_iff}.
                    Then $d(p)\in\reals$ by \cref{pair_eq_iff,metric_on,funs_elim}.
                    $1\in\reals$ by \cref{real_naturals_subset_reals,real_one_in_naturals,subseteq}.
                    \begin{byCase}
                    \caseOf{$d(p) < 1$.}
                        Hence $z_1 = z_2$.
                    \caseOf{it is wrong that $d(p) < 1$.}
                        Hence $z_1 = z_2$.
                    \end{byCase}
                \end{subproof}
                Hence $e$ is a function by \cref{bounded_metric,relation,rightunique}.
        For all $p\in\dom{e}$ we have $e(p)\in\reals$
            by \cref{dom_iff,bounded_metric,pair_eq_iff,function_apply_intro}.
        Hence $e$ is a function to $\reals$.
        Then $\dom{e}\subseteq X\times X$
            by \cref{dom_iff,bounded_metric,pair_eq_iff,metric_on,funs_elim,subseteq}.
        Show that for all $p$ such that $p\in X\times X$ we have $p\in\dom{e}$.
        \begin{subproof}
                Fix $p$ such that $p\in X\times X$.
                Thus $p\in\dom{d}$ and $d(p)\in\reals$ by \cref{metric_on,funs_elim,subseteq,funs_type_apply}.
                $1\in\reals$ by \cref{real_naturals_subset_reals,real_one_in_naturals,subseteq}.
                \begin{byCase}
                \caseOf{$d(p) < 1$.}
                    Thus $p\in\dom{e}$ by \cref{bounded_metric,dom_intro}.
                \caseOf{it is wrong that $d(p) < 1$.}
                    Thus $p\in\dom{e}$ by \cref{real_lt_trichotomy,bounded_metric,dom_intro}.
                \end{byCase}
        \end{subproof}
        Hence $e\in\funs{X\times X}{\reals}$ by \cref{subseteq,subseteq_antisymmetric,funs_intro}.
    \end{subproof}
    Show that for all $x,y$ such that $x\in X$ and $y\in X$ we have $e((x,y)) \geq 0$.
    \begin{subproof}
        Fix $x$.
        Fix $y$ such that $x\in X$ and $y\in X$.
        Let $p = (x,y)$.
        Thus $p\in\dom{d}$ and $d(p)\in\reals$ by \cref{times_tuple_intro,metric_on,funs_elim,subseteq,funs_type_apply}.
        $1\in\reals$ by \cref{real_naturals_subset_reals,real_one_in_naturals,subseteq}.
        \begin{byCase}
        \caseOf{$d(p) < 1$.}
            Thus $e((x,y)) \geq 0$ by \cref{bounded_metric,function_apply_intro,funs_elim,metric_on,metric_nonnegative}.
        \caseOf{it is wrong that $d(p) < 1$.}
            Thus $e((x,y)) \geq 0$ by \cref{real_lt_trichotomy,bounded_metric,function_apply_intro,funs_elim,real_zero_lt_one}.
        \end{byCase}
    \end{subproof}
    Hence $e$ is nonnegative on $X$ by \cref{metric_nonnegative}.
    Show that for all $x,y$ such that $x\in X$ and $y\in X$ we have $e((x,y)) = 0$ iff $x = y$.
    \begin{subproof}
        Fix $x$.
        Fix $y$ such that $x\in X$ and $y\in X$.
        Show that if $e((x,y)) = 0$ then $x = y$.
        \begin{subproof}
            Suppose not.
            Thus $e((x,y)) = 0$ and $x \neq y$.
            Let $p = (x,y)$.
            Thus $(p,0)\in e$ by \cref{times_tuple_intro,function_apply_iff,funs_elim}.
            Take $p_0, w$ such that
                $p_0\in\dom{d}$ and $w\in\reals$ and
                $(p,0) = (p_0,w)$ and
                $((w = d(p_0) \land d(p_0) < 1) \lor (w = 1 \land 1 \leq d(p_0)))$
                by \cref{bounded_metric}.
            Hence $((0 = d(p) \land d(p) < 1) \lor (0 = 1 \land 1 \leq d(p)))$
                by \cref{bounded_metric,pair_eq_iff}.
            Then $0\in\reals$ by \cref{pair_eq_iff}.
            \begin{byCase}
            \caseOf{$0 = d(p) \land d(p) < 1$.}
                Thus $x = y$ by \cref{metric_on,metric_zero_characterization}.
            \caseOf{$0 = 1 \land 1 \leq d(p)$.}
                Contradiction by \cref{real_zero_lt_one,real_lt_subst_right,real_lt_irreflexive}.
            \end{byCase}
        \end{subproof}
        Hence if $x = y$ then $e((x,y)) = 0$
            by \cref{times_tuple_intro,metric_on,metric_zero_characterization,real_zero_lt_one,real_lt_subst_left,funs_elim,subseteq,funs_type_apply,bounded_metric,function_apply_intro}.
    \end{subproof}
    Therefore $e$ is zero-characterized on $X$ by \cref{metric_zero_characterization}.
    Show that for all $x,y$ such that $x\in X$ and $y\in X$ we have $e((x,y)) = e((y,x))$.
    \begin{subproof}
        Fix $x$.
        Fix $y$ such that $x\in X$ and $y\in X$.
        Let $p_1 = (x,y)$.
        Let $p_2 = (y,x)$.
        Thus $d((x,y)) = d((y,x))$ by \cref{times_tuple_intro,metric_on,metric_symmetric}.
        Thus $p_1\in\dom{d}$ and $p_2\in\dom{d}$ and $d(p_1)\in\reals$ and $d(p_2)\in\reals$
            by \cref{times_tuple_intro,metric_on,funs_elim,subseteq,funs_type_apply}.
        $1\in\reals$ by \cref{real_naturals_subset_reals,real_one_in_naturals,subseteq}.
        \begin{byCase}
        \caseOf{$d(p_1) < 1$.}
            Thus $d(p_2) < 1$ by \cref{real_lt_subst_left}.
            Let $w_1 = (p_1,d(p_1))$.
            Let $w_2 = (p_2,d(p_2))$.
            $w_1\in e$ and $w_2\in e$ by \cref{bounded_metric}.
            Thus $e((x,y)) = e((y,x))$ by \cref{function_apply_intro,funs_elim}.
        \caseOf{it is wrong that $d(p_1) < 1$.}
            Thus $1 \leq d(p_1)$.
            Thus $1 \leq d(p_2)$ by \cref{real_lt_subst_left}.
            Let $w_1 = (p_1,1)$.
            Let $w_2 = (p_2,1)$.
            $w_1\in e$ and $w_2\in e$ by \cref{bounded_metric}.
            Thus $e((x,y)) = e((y,x))$ by \cref{function_apply_intro,funs_elim}.
        \end{byCase}
    \end{subproof}
    Hence $e$ is symmetric on $X$ by \cref{metric_symmetric}.
    Show that for all $x,y,z$ such that $x\in X$ and $y\in X$ and $z\in X$ we have
        $e((x,y)) + e((y,z)) \geq e((x,z))$.
    \begin{subproof}
        Fix $x$.
        Fix $y$.
        Fix $z$ such that $x\in X$ and $y\in X$ and $z\in X$.
        Let $p_1 = (x,y)$.
        Let $p_2 = (y,z)$.
        Let $p_3 = (x,z)$.
        $p_1\in X\times X$ and $p_2\in X\times X$ and $p_3\in X\times X$ by \cref{times_tuple_intro}.
        Thus $p_1\in\dom{d}$ and $p_2\in\dom{d}$ and $p_3\in\dom{d}$ and
            $d(p_1)\in\reals$ and $d(p_2)\in\reals$ and $d(p_3)\in\reals$
            by \cref{metric_on,funs_elim,subseteq,funs_type_apply}.
        $1\in\reals$ by \cref{real_naturals_subset_reals,real_one_in_naturals,subseteq}.
        Thus $d(p_3) \leq d(p_1) + d(p_2)$ by \cref{metric_on,metric_triangle_inequality}.
        \begin{byCase}
        \caseOf{$d(p_1) < 1$.}
            \begin{byCase}
            \caseOf{$d(p_2) < 1$.}
                Let $w_1 = (p_1,d(p_1))$.
                Let $w_2 = (p_2,d(p_2))$.
                $w_1\in e$ and $w_2\in e$ by \cref{bounded_metric}.
                Thus $e(p_1) = d(p_1)$ and $e(p_2) = d(p_2)$ by \cref{function_apply_intro,funs_elim}.
                Show that $e(p_3) \leq d(p_3)$.
                \begin{subproof}
                    \begin{byCase}
                    \caseOf{$d(p_3) < 1$.}
                        Let $w_3 = (p_3,d(p_3))$.
                        $w_3\in e$ by \cref{bounded_metric}.
                        Thus $e(p_3) \leq d(p_3)$ by \cref{function_apply_intro,funs_elim}.
                    \caseOf{it is wrong that $d(p_3) < 1$.}
                        Thus $1 \leq d(p_3)$.
                        Let $w_3 = (p_3,1)$.
                        $w_3\in e$ by \cref{bounded_metric}.
                        Thus $e(p_3) \leq d(p_3)$ by \cref{function_apply_intro,funs_elim}.
                    \end{byCase}
                \end{subproof}
                Hence $e(p_3)\in\reals$ by \cref{funs_elim,funs_type_apply}.
                $d(p_1) + d(p_2)\in\reals$ by \cref{real_add_closed}.
                Therefore $e(p_3) \leq d(p_1) + d(p_2)$ by \cref{real_leq_trans}.
                Then $e(p_1)\in\reals$ and $e(p_2)\in\reals$ by \cref{funs_elim}.
                Hence $e(p_1) + e(p_2) = d(p_1) + e(p_2)$ and
                    $d(p_1) + e(p_2) = d(p_1) + d(p_2)$ by \cref{real_add_eq_subst}.
                Therefore $e(p_1) + e(p_2) = d(p_1) + d(p_2)$.
                Then $d(p_1) + d(p_2) \leq e(p_1) + e(p_2)$.
                Hence $e(p_3) \leq e(p_1) + e(p_2)$ by \cref{real_leq_trans}.
            \caseOf{it is wrong that $d(p_2) < 1$.}
                Thus $1 \leq d(p_2)$.
                Let $w_1 = (p_1,d(p_1))$.
                Let $w_2 = (p_2,1)$.
                $w_1\in e$ and $w_2\in e$ by \cref{bounded_metric}.
                Thus $e(p_1) = d(p_1)$ and $e(p_2) = 1$ by \cref{function_apply_intro,funs_elim}.
                Show that $e(p_3) \leq 1$.
                \begin{subproof}
                    \begin{byCase}
                    \caseOf{$d(p_3) < 1$.}
                        Let $w_3 = (p_3,d(p_3))$.
                        $w_3\in e$ by \cref{bounded_metric}.
                        Thus $e(p_3) \leq 1$ by \cref{function_apply_intro,funs_elim,real_lt_subst_left}.
                    \caseOf{it is wrong that $d(p_3) < 1$.}
                        Thus $1 \leq d(p_3)$.
                        Let $w_3 = (p_3,1)$.
                        $w_3\in e$ by \cref{bounded_metric}.
                        Thus $e(p_3) \leq 1$ by \cref{function_apply_intro,funs_elim}.
                    \end{byCase}
                \end{subproof}
                Thus $0 \leq e(p_1)$ by \cref{metric_nonnegative}.
                Thus $e(p_1)\in\reals$ and $e(p_2)\in\reals$ and $e(p_3)\in\reals$ and
                    $e(p_1) + e(p_2)\in\reals$ by \cref{funs_type_apply,real_add_closed}.
                $0\in\reals$ and $1\in\reals$ by \cref{real_add_ident,real_naturals_subset_reals,real_one_in_naturals,subseteq}.
                Thus $0 + 1 \leq e(p_1) + 1$ by \cref{real_leq_add}.
                $0 + 1 = 1$ and $e(p_1) + 1 = 1 + e(p_1)$ by \cref{real_add_ident,real_add_comm}.
                Thus $1 \leq e(p_1) + e(p_2)$ and $e(p_3) \leq e(p_1) + e(p_2)$ by \cref{real_leq_trans}.
            \end{byCase}
        \caseOf{it is wrong that $d(p_1) < 1$.}
            Thus $1 \leq d(p_1)$.
            Let $w_1 = (p_1,1)$.
            $w_1\in e$ by \cref{bounded_metric}.
            Thus $e(p_1) = 1$ by \cref{function_apply_intro,funs_elim}.
            Show that $e(p_3) \leq 1$.
            \begin{subproof}
                \begin{byCase}
                \caseOf{$d(p_3) < 1$.}
                    Let $w_3 = (p_3,d(p_3))$.
                    $w_3\in e$ by \cref{bounded_metric}.
                    Thus $e(p_3) \leq 1$ by \cref{function_apply_intro,funs_elim,real_lt_subst_left}.
                \caseOf{it is wrong that $d(p_3) < 1$.}
                    Thus $1 \leq d(p_3)$.
                    Let $w_3 = (p_3,1)$.
                    $w_3\in e$ by \cref{bounded_metric}.
                    Thus $e(p_3) \leq 1$ by \cref{function_apply_intro,funs_elim}.
                \end{byCase}
            \end{subproof}
            Thus $0 \leq e(p_2)$ by \cref{metric_nonnegative}.
            Thus $e(p_1)\in\reals$ and $e(p_2)\in\reals$ and $e(p_3)\in\reals$ and
                $e(p_1) + e(p_2)\in\reals$ by \cref{funs_type_apply,real_add_closed}.
            $0\in\reals$ and $1\in\reals$ by \cref{real_add_ident,real_naturals_subset_reals,real_one_in_naturals,subseteq}.
            Thus $0 + 1 \leq e(p_2) + 1$ by \cref{real_leq_add}.
            $0 + 1 = 1$ and $e(p_2) + 1 = 1 + e(p_2)$ by \cref{real_add_ident,real_add_comm}.
            Thus $1 \leq e(p_1) + e(p_2)$ and $e(p_3) \leq e(p_1) + e(p_2)$ by \cref{real_leq_trans}.
        \end{byCase}
    \end{subproof}
    Hence $e$ is a metric on $X$ by \cref{metric_triangle_inequality,metric_on}.
\end{proof}
% from §20 helper lemma 6: bounded metric agrees with the original metric below one
\begin{lemma}\label{bounded_metric_apply_below_one}
    Assume $d$ is a metric on $X$.
    Let $e = \boundedMetric{d}$.
    Suppose $p\in X\times X$.
    Suppose $d(p) < 1$.
    Then $e(p) = d(p)$.
\end{lemma}
\begin{proof}
    Thus $e(p) = d(p)$ by \cref{bounded_metric,bounded_metric_is_metric,metric_on,funs_elim,subseteq,funs_type_apply,function_apply_intro}.
\end{proof}

% from §20 helper lemma 7: bounded metric agrees with the original metric when its value is below one
\begin{lemma}\label{bounded_metric_apply_lt_one}
    Assume $d$ is a metric on $X$.
    Let $e = \boundedMetric{d}$.
    Suppose $p\in X\times X$.
    Suppose $e(p) < 1$.
    Then $e(p) = d(p)$.
\end{lemma}
\begin{proof}
    Hence $(e(p) = d(p) \land d(p) < 1) \lor (e(p) = 1 \land 1 \leq d(p))$
        by \cref{bounded_metric_is_metric,metric_on,funs_elim,subseteq,function_apply_iff,bounded_metric,pair_eq_iff}.
    Then $e(p) = d(p)$.
\end{proof}

% from §20 Item 19: standard bounded metric induces the same topology
\begin{theorem}\label{bounded_metric_topology}
    Assume $d$ is a metric on $X$.
    Then $\metricTopology{\boundedMetric{d}}{X} = \metricTopology{d}{X}$.
\end{theorem}
\begin{proof}
    Let $e = \boundedMetric{d}$.
    $e$ is a metric on $X$ by \cref{bounded_metric_is_metric}.
    Let $T = \metricTopology{d}{X}$.
    Let $T_1 = \metricTopology{e}{X}$.
    Show that for all $U$ such that $U\in T_1$ we have $U\in T$.
    \begin{subproof}
        Fix $U$ such that $U\in T_1$.
        Hence $U$ is open in $X$ with respect to $T_1$ and $U\subseteq X$
            by \cref{metric_basis_is_basis,metric_topology,basis_topology_is_topology,basis_for_topology,open_in,topology_on,subseteq,powerset_elim}.
        Show that for all $y\in U$ there exists $\delta\in\reals$ such that $0 < \delta$ and $\metricBall{d}{X}{y}{\delta}\subseteq U$.
        \begin{subproof}
                Fix $y$ such that $y\in U$.
                Take $\epsilon\in\reals$ such that $0 < \epsilon$ and $\metricBall{e}{X}{y}{\epsilon}\subseteq U$
                    by \cref{metric_open_iff}.
                $1\in\reals$ by \cref{real_naturals_subset_reals,real_one_in_naturals,subseteq}.
                Thus $\epsilon < 1$ or $\epsilon = 1$ or $1 < \epsilon$ by \cref{real_lt_trichotomy}.
                \begin{byCase}
                \caseOf{$\epsilon < 1$.}
                    Let $\delta = \epsilon$.
                    For all $z$ such that $z\in\metricBall{d}{X}{y}{\delta}$ we have $z\in U$
                        by \cref{metric_ball,subseteq,times_tuple_intro,metric_on,funs_type_apply,real_lt_trans,bounded_metric_apply_below_one,real_lt_subst_left}.
                    Therefore $\metricBall{d}{X}{y}{\delta}\subseteq U$ by \cref{subseteq}.
                \caseOf{$\epsilon = 1$.}
                    Let $\delta = \epsilon$.
                    For all $z$ such that $z\in\metricBall{d}{X}{y}{\delta}$ we have $z\in U$
                        by \cref{metric_ball,real_lt_subst_right,bounded_metric_apply_below_one,subseteq,times_tuple_intro,metric_on,funs_type_apply,real_lt_subst_left}.
                    Therefore $\metricBall{d}{X}{y}{\delta}\subseteq U$ by \cref{subseteq}.
                \caseOf{$1 < \epsilon$.}
                    Let $\delta = 1$.
                    For all $z$ such that $z\in\metricBall{d}{X}{y}{\delta}$ we have $z\in U$
                        by \cref{metric_ball,real_lt_subst_right,bounded_metric_apply_below_one,metric_on,funs_type_apply,subseteq,times_tuple_intro,real_lt_subst_left,real_lt_trans}.
                    Therefore $\metricBall{d}{X}{y}{\delta}\subseteq U$ by \cref{subseteq}.
                \end{byCase}
        \end{subproof}
        Hence $U\in T$ by \cref{metric_open_iff,open_in}.
    \end{subproof}
    Hence $T_1\subseteq T$ by \cref{subseteq}.
    Show that for all $U$ such that $U\in T$ we have $U\in T_1$.
    \begin{subproof}
        Fix $U$ such that $U\in T$.
        Hence $U$ is open in $X$ with respect to $T$ and $U\subseteq X$
            by \cref{metric_basis_is_basis,metric_topology,basis_topology_is_topology,basis_for_topology,open_in,topology_on,subseteq,powerset_elim}.
        Show that for all $y\in U$ there exists $\delta\in\reals$ such that $0 < \delta$ and $\metricBall{e}{X}{y}{\delta}\subseteq U$.
        \begin{subproof}
                Fix $y$ such that $y\in U$.
                Take $\epsilon\in\reals$ such that $0 < \epsilon$ and $\metricBall{d}{X}{y}{\epsilon}\subseteq U$
                    by \cref{metric_open_iff}.
                $1\in\reals$ by \cref{real_naturals_subset_reals,real_one_in_naturals,subseteq}.
                Let $A = \{\epsilon,1\}$.
                $A$ is finite by \cref{pair_is_finite}.
                $\epsilon\in A$ and $1\in A$ by \cref{upair_intro_left,upair_intro_right}.
                Hence $A\neq\emptyset$ by \cref{upair_intro_left,emptyset}.
                Therefore $A\subseteq\reals$ by \cref{upair_elim,subseteq}.
                Take $m\in A$ such that $m$ is a minimum of $A$ by \cref{finite_real_minimum}.
                Hence $0 < m$ by \cref{upair_elim,real_zero_lt_one,real_lt_subst_right}.
                For all $z$ such that $z\in\metricBall{e}{X}{y}{m}$ we have $z\in U$
                    by \cref{metric_ball,subseteq,times_tuple_intro,metric_on,funs_elim,funs_type_apply,real_lt_subst_left,real_minimum,real_leq_split,real_lt_subst_right,real_lt_trans,bounded_metric_apply_lt_one}.
                Therefore $\metricBall{e}{X}{y}{m}\subseteq U$ by \cref{subseteq}.
        \end{subproof}
        Hence $U\in T_1$ by \cref{metric_open_iff,open_in}.
    \end{subproof}
    Therefore $T_1 = T$ by \cref{subseteq,subseteq_antisymmetric}.
\end{proof}

% from §20 Item 20: $\reals^n$
\begin{definition}\label{reals_power}
    $\realsPower{J} = \tuplePower{\reals}{J}$.
\end{definition}

% from §20 Item 21: euclidean norm
\begin{definition}\label{euclidean_norm}
    $r$ is a euclidean norm of $x$ with respect to $n$ iff
    $x\in\realsPower{\initialSegment{n}}$ and
    $r\in\reals$ and
    there exists $a\in\funs{\initialSegment{n}}{\reals}$ such that
        for all $i\in\initialSegment{n}$ we have $a(i) = \exp{x(i)}{1+1}$ and
        there exists $s$ such that $s$ is the sum of $a$ and $r = \sqrt{s}$.
\end{definition}
% from §20 helper lemma 6: sum of a finite sequence predicate
\begin{definition}\label{sum_finite_sequence_pred}
    $\sumFiniteSequence{a}{s}$ iff
    there exists $n\in\naturalsPlus$ such that
    $a\in\funs{\initialSegment{n}}{\reals}$ and
    there exists $t\in\funs{\initialSegment{n}}{\reals}$ such that
    $t(1) = a(1)$ and
    $s = t(n)$ and
    for all $k\in\naturalsPlus$ such that $k + 1 \leq n$ we have
    $t(k + 1) = t(k) + a(k + 1)$.
\end{definition}

% from §20 helper lemma 7: sum of a finite sequence is unique
\begin{lemma}\label{sum_finite_sequence_unique}
    Suppose $\sumFiniteSequence{a}{s_1}$.
    Suppose $\sumFiniteSequence{a}{s_2}$.
    Then $s_1 = s_2$.
\end{lemma}
\begin{proof}
    Take $n_1$ such that $n_1\in\naturalsPlus$ and $a\in\funs{\initialSegment{n_1}}{\reals}$ and
        there exists $t_1\in\funs{\initialSegment{n_1}}{\reals}$ such that
            $t_1(1) = a(1)$ and
            $s_1 = t_1(n_1)$ and
            for all $k\in\naturalsPlus$ such that $k + 1 \leq n_1$ we have $t_1(k + 1) = t_1(k) + a(k + 1)$
        by \cref{sum_finite_sequence_pred}.
    Take $n_2$ such that $n_2\in\naturalsPlus$ and $a\in\funs{\initialSegment{n_2}}{\reals}$ and
        there exists $t_2\in\funs{\initialSegment{n_2}}{\reals}$ such that
            $t_2(1) = a(1)$ and
            $s_2 = t_2(n_2)$ and
            for all $k\in\naturalsPlus$ such that $k + 1 \leq n_2$ we have $t_2(k + 1) = t_2(k) + a(k + 1)$
        by \cref{sum_finite_sequence_pred}.
    Take $t_1$ such that $t_1\in\funs{\initialSegment{n_1}}{\reals}$ and
        $t_1(1) = a(1)$ and
        $s_1 = t_1(n_1)$ and
        for all $k\in\naturalsPlus$ such that $k + 1 \leq n_1$ we have $t_1(k + 1) = t_1(k) + a(k + 1)$.
    Take $t_2$ such that $t_2\in\funs{\initialSegment{n_2}}{\reals}$ and
        $t_2(1) = a(1)$ and
        $s_2 = t_2(n_2)$ and
        for all $k\in\naturalsPlus$ such that $k + 1 \leq n_2$ we have $t_2(k + 1) = t_2(k) + a(k + 1)$.
    Hence $n_1 = n_2$
        by \cref{funs_elim,subseteq_antisymmetric,subseteq,initial_segment_naturalsplus,real_naturals_subset_reals,real_leq_antisym}.
    Let $n = n_1$.
    Now $t_1\in\funs{\initialSegment{n}}{\reals}$ and $t_2\in\funs{\initialSegment{n}}{\reals}$.
    Let $A = \{k\in\naturalsPlus \mid \text{if $k \leq n$ then $t_1(k) = t_2(k)$}\}$.
    $A\subseteq\naturalsPlus$ by \cref{subseteq}.
    Hence $1\in A$ by \cref{real_one_in_naturals,real_one_leq_naturalsplus}.
    Show that for all $k\in A$ we have $k + 1\in A$.
    \begin{subproof}
        Fix $k$ such that $k\in A$.
        Show that if $k + 1 \leq n$ then $t_1(k + 1) = t_2(k + 1)$.
        \begin{subproof}
            Assume $k + 1 \leq n$.
            $k\in\reals$ and $n\in\reals$ and $0\in\reals$ and $1\in\reals$
                by \cref{real_naturals_subset_reals,subseteq,real_add_ident,real_one_in_naturals}.
            Hence $k < k + 1$
                by \cref{real_zero_lt_one,real_lt_add,real_add_ident,real_add_comm,real_lt_subst_left}.
            $k \leq k + 1$.
            $k + 1\in\reals$ by \cref{real_add_closed}.
            Hence $k \leq n$ by \cref{real_leq_trans}.
            Then $k\in\initialSegment{n}$ by \cref{subseteq,initial_segment_naturalsplus}.
            $t_1(k) = t_2(k)$.
            Therefore $k + 1\in\initialSegment{n}$ by \cref{real_naturals_closed_succ,initial_segment_naturalsplus}.
            Hence $a(k + 1)\in\reals$ and $t_1(k)\in\reals$ and $t_2(k)\in\reals$ by \cref{funs_type_apply}.
            Then $t_1(k + 1) = t_1(k) + a(k + 1)$ and $t_2(k + 1) = t_2(k) + a(k + 1)$.
            Therefore $t_1(k + 1) = t_2(k + 1)$ by \cref{real_add_eq_subst}.
        \end{subproof}
        Hence $k + 1\in A$.
    \end{subproof}
    Therefore for all $k\in\naturalsPlus$ we have $k\in A$ by \cref{real_naturals_induction}.
    For all $k\in\naturalsPlus$ we have if $k \leq n$ then $t_1(k) = t_2(k)$.
    Now $\dom{t_1} = \initialSegment{n}$ and $\dom{t_2} = \initialSegment{n}$ and
        $t_1$ is a function and $t_2$ is a function by \cref{funs_elim}.
    $t_1\subseteq t_2$ by \cref{subseteq,initial_segment_naturalsplus,fun_subseteq}.
    Then $t_2\subseteq t_1$ by \cref{subseteq,initial_segment_naturalsplus,fun_subseteq}.
    Hence $t_1 = t_2$ by \cref{subseteq_antisymmetric}.
    Therefore $s_1 = s_2$.
\end{proof}

% from §20 helper lemma 8: squared coordinate sequence
\begin{definition}\label{coord_square_sequence}
    $\coordSquareSequence{n}{p} = \{w \mid \exists i\in\initialSegment{n}.
        w = (i, \exp{\apply{\fst{p}}{i} - \apply{\snd{p}}{i}}{1+1})\}$.
\end{definition}

% from §20 Item 22: euclidean metric
\begin{definition}\label{euclidean_metric}
    $\euclideanMetric{n} = \{w \mid \exists p\in\realsPower{\initialSegment{n}}\times\realsPower{\initialSegment{n}}. \exists r\in\reals.
        \exists s\in\reals.
        w = (p,r) \land
        \sumFiniteSequence{\coordSquareSequence{n}{p}}{s} \land
        r = \sqrt{s}\}$.
\end{definition}

% from §20 helper lemma 9: coordinate difference set
\begin{definition}\label{coord_diff_set}
    $\coordDiffSet{n}{x}{y} = \{\abs{x(i) - y(i)} \mid i\in\initialSegment{n}\}$.
\end{definition}

% from §20 Item 23: square metric
\begin{definition}\label{square_metric}
    $\squareMetric{n} = \{w \mid \exists p\in\realsPower{\initialSegment{n}}\times\realsPower{\initialSegment{n}}.
        \exists x\in\realsPower{\initialSegment{n}}. \exists y\in\realsPower{\initialSegment{n}}. \exists r\in\reals.
        w = (p,r) \land p = (x,y) \land
        r\in\coordDiffSet{n}{x}{y} \land
        (\forall t\in\coordDiffSet{n}{x}{y}. t \leq r)\}$.
\end{definition}

% from §20 Item 24: comparison of metric topologies
\begin{lemma}\label{metric_topology_finer_iff}
    Suppose $d$ is a metric on $X$.
    Suppose $d_1$ is a metric on $X$.
    Let $T = \metricTopology{d}{X}$.
    Let $T_1 = \metricTopology{d_1}{X}$.
    Then $T_1$ is finer than $T$ on $X$ iff
    for all $x\in X$ we have
    for all $\epsilon\in\reals$ such that $0 < \epsilon$ there exists $\delta\in\reals$
        such that $0 < \delta$ and
        $\metricBall{d_1}{X}{x}{\delta} \subseteq \metricBall{d}{X}{x}{\epsilon}$.
\end{lemma}
\begin{proof}
    Show that if $T_1$ is finer than $T$ on $X$ then
        for all $x\in X$ we have
        for all $\epsilon\in\reals$ such that $0 < \epsilon$ there exists $\delta\in\reals$
            such that $0 < \delta$ and
            $\metricBall{d_1}{X}{x}{\delta} \subseteq \metricBall{d}{X}{x}{\epsilon}$.
    \begin{subproof}
        Assume $T_1$ is finer than $T$ on $X$.
        Fix $x\in X$.
        Fix $\epsilon\in\reals$.
        Assume $0 < \epsilon$.
        Let $U = \metricBall{d}{X}{x}{\epsilon}$.
        Hence $U\subseteq X$ by \cref{metric_ball,subseteq}.
        Therefore $U\in T$ by \cref{metric_ball,powerset_intro,metric_basis,metric_basis_is_basis,basis_elem_open_in_generated,basis_for_topology,metric_topology}.
        Then $U$ is open in $X$ with respect to $T_1$
            by \cref{finer_topology,subseteq,basis_topology_is_topology,metric_topology,metric_basis_is_basis,basis_for_topology,open_in}.
        Hence $x\in U$ by \cref{metric_ball,metric_zero_characterization,real_lt_subst_left,metric_on}.
        Therefore there exists $\delta\in\reals$ such that $0 < \delta$ and
            $\metricBall{d_1}{X}{x}{\delta} \subseteq \metricBall{d}{X}{x}{\epsilon}$
            by \cref{metric_open_iff}.
    \end{subproof}
    Show that if
        for all $x\in X$ we have
        for all $\epsilon\in\reals$ such that $0 < \epsilon$ there exists $\delta\in\reals$
            such that $0 < \delta$ and
            $\metricBall{d_1}{X}{x}{\delta} \subseteq \metricBall{d}{X}{x}{\epsilon}$
        then $T_1$ is finer than $T$ on $X$.
    \begin{subproof}
        Assume for all $x\in X$ we have
            for all $\epsilon\in\reals$ such that $0 < \epsilon$ there exists $\delta\in\reals$
                such that $0 < \delta$ and
                $\metricBall{d_1}{X}{x}{\delta} \subseteq \metricBall{d}{X}{x}{\epsilon}$.
        Show that for all $U$ such that $U\in T$ we have $U\in T_1$.
        \begin{subproof}
            Fix $U$ such that $U\in T$.
            Hence $U\subseteq X$ by \cref{basis_topology_is_topology,metric_topology,metric_basis_is_basis,basis_for_topology,topology_on,subseteq,powerset_elim}.
            Show that for all $x\in U$ there exists $\delta\in\reals$ such that $0 < \delta$ and
                $\metricBall{d_1}{X}{x}{\delta} \subseteq U$.
            \begin{subproof}
                Fix $x\in U$.
                Take $\epsilon\in\reals$ such that $0 < \epsilon$ and
                    $\metricBall{d}{X}{x}{\epsilon} \subseteq U$.
                $x\in X$ by \cref{subseteq}.
                Therefore there exists $\delta\in\reals$ such that $0 < \delta$ and
                    $\metricBall{d_1}{X}{x}{\delta} \subseteq U$
                    by \cref{subseteq_transitive}.
            \end{subproof}
            Hence $U\in T_1$ by \cref{metric_open_iff,open_in}.
        \end{subproof}
        Therefore $T\subseteq T_1$ by \cref{subseteq}.
        Hence $T_1$ is finer than $T$ on $X$
            by \cref{basis_topology_is_topology,metric_topology,metric_basis_is_basis,basis_for_topology,finer_topology}.
    \end{subproof}
\end{proof}

% from §20 helper lemma 13: initial segment is finite
\begin{lemma}\label{initial_segment_finite}
    Suppose $n\in\naturalsPlus$.
    Then $\initialSegment{n}$ is finite.
\end{lemma}
\begin{proof}
    Let $A = \{n\in\naturalsPlus \mid \text{$\initialSegment{n}$ is finite}\}$.
    $A\subseteq\naturalsPlus$ by \cref{subseteq}.
    Let $B = \initialSegment{1}$.
    Thus for all $k$ such that $k\in B$ we have $k = 1$.
    Hence for all $k$ such that $k\in B$ we have $k\in\{1,1\}$ by \cref{upair_intro_left}.
    Hence $B$ is finite by \cref{subseteq,pair_is_finite,subseteq_finite_is_finite}.
    Therefore $1\in A$.
    Show that for all $n\in A$ we have $n + 1\in A$.
    \begin{subproof}
        Fix $n$ such that $n\in A$.
        Thus $\initialSegment{n}$ is finite.
        $n + 1\in\naturalsPlus$ by \cref{subseteq,real_naturals_closed_succ}.
        Let $B = \initialSegment{n + 1}$.
        For all $k$ such that $k\in B$ we have $k\in \initialSegment{n}\union\{n + 1,n + 1\}$.
        Hence $B$ is finite by \cref{subseteq,pair_is_finite,union_of_finite_is_finite,subseteq_finite_is_finite}.
        Therefore $n + 1\in A$.
    \end{subproof}
    Therefore $\initialSegment{n}$ is finite by \cref{real_naturals_induction,subseteq}.
\end{proof}

% from §20 helper lemma 14: product family of constant reals
\begin{lemma}\label{product_family_reals_power}
    Suppose $J$ is inhabited.
    Let $X = \{(i,\reals) \mid i\in J\}$.
    Then $\productFamily{X} = \realsPower{J}$.
\end{lemma}
\begin{proof}
    Hence $X$ is a relation by \cref{relation}.
    For all $i,y_1,y_2$ such that $(i,y_1)\in X$ and $(i,y_2)\in X$ we have $y_1 = y_2$.
    Therefore $X$ is a function by \cref{rightunique}.
    Hence $X$ is a function on $J$
        by \cref{dom_intro,dom_iff,pair_eq_iff,subseteq,subseteq_antisymmetric}.
    For all $y$ such that $y\in\ran{X}$ we have $y = \reals$ by \cref{ran_iff,pair_eq_iff}.
    For all $z$ such that $z\in\unions{\ran{X}}$ we have $z\in\reals$ by \cref{unions_iff}.
    For all $z$ such that $z\in\reals$ we have $z\in\unions{\ran{X}}$
        by \cref{ran_intro,unions_iff}.
    Hence $\unions{\ran{X}} = \reals$ by \cref{subseteq,subseteq_antisymmetric}.
    Therefore $\productFamily{X}\subseteq\realsPower{J}$
        by \cref{indexed_product,funs_elim,function_on_weaken_codom,funs_intro,tuple_power,reals_power,subseteq}.
    For all $f$ such that $f\in\realsPower{J}$ we have $f\in\productFamily{X}$
        by \cref{reals_power,tuple_power,funs_elim,function_on_weaken_codom,funs_intro,subseteq,function_apply_intro,indexed_product}.
    Therefore $\productFamily{X} = \realsPower{J}$ by \cref{subseteq,subseteq_antisymmetric}.
\end{proof}

% from §20 helper lemma 15: coordinate difference set is finite
\begin{lemma}\label{coord_diff_set_finite}
    Suppose $n\in\naturalsPlus$.
    Let $J = \initialSegment{n}$.
    Suppose $x\in\realsPower{J}$.
    Suppose $y\in\realsPower{J}$.
    Then $\coordDiffSet{n}{x}{y}\subseteq\reals$ and
         $\coordDiffSet{n}{x}{y}$ is finite and
         $\coordDiffSet{n}{x}{y}$ is inhabited.
\end{lemma}
\begin{proof}
    $1\in\naturalsPlus$ by \cref{real_one_in_naturals}.
    Let $D = \coordDiffSet{n}{x}{y}$.
    Hence $D\subseteq\reals$
        by \cref{coord_diff_set,reals_power,tuple_power,funs_type_apply,real_add_inverse,real_add_closed,real_sub,real_abs_in_reals,subseteq}.
    Let $g = \{(i,\abs{x(i) - y(i)}) \mid i\in J\}$.
    Hence $g$ is a relation by \cref{coord_diff_set,relation}.
    For all $i,t_1,t_2$ such that $(i,t_1)\in g$ and $(i,t_2)\in g$ we have $t_1 = t_2$ by \cref{pair_eq_iff}.
    Hence $g$ is a function by \cref{rightunique}.
    Hence $g$ is a function on $J$ by \cref{dom_intro,dom_iff,pair_eq_iff,subseteq,subseteq_antisymmetric}.
    For all $t$ we have $t\in D$ iff $t\in\img{g}{J}$ by \cref{coord_diff_set,img_iff}.
    Therefore $D$ is finite by \cref{setext,initial_segment_finite,finite_image_function}.
    Therefore $D$ is inhabited by \cref{real_one_in_naturals,real_one_leq_naturalsplus,initial_segment_naturalsplus,coord_diff_set}.
\end{proof}

% from §20 helper lemma 16: square metric is a function
\begin{lemma}\label{square_metric_function}
    Suppose $n\in\naturalsPlus$.
    Then $\squareMetric{n}$ is a function.
\end{lemma}
\begin{proof}
    Let $d = \squareMetric{n}$.
    Show that for all $p,r_1,r_2$ such that $(p,r_1)\in d$ and $(p,r_2)\in d$ we have $r_1 = r_2$.
    \begin{subproof}
        Fix $p$.
        Fix $r_1$.
        Fix $r_2$ such that $(p,r_1)\in d$ and $(p,r_2)\in d$.
        Take $x_1,y_1\in\realsPower{\initialSegment{n}}$ such that
            $p = (x_1,y_1)$ and
            $r_1\in\coordDiffSet{n}{x_1}{y_1}$ and
            for all $t\in\coordDiffSet{n}{x_1}{y_1}$ we have $t \leq r_1$
            by \cref{square_metric,pair_eq_iff}.
        Take $x_2,y_2\in\realsPower{\initialSegment{n}}$ such that
            $p = (x_2,y_2)$ and
            $r_2\in\coordDiffSet{n}{x_2}{y_2}$ and
            for all $t\in\coordDiffSet{n}{x_2}{y_2}$ we have $t \leq r_2$
            by \cref{square_metric,pair_eq_iff}.
        Hence $x_1 = x_2$ and $y_1 = y_2$ by \cref{pair_eq_iff}.
        Then $r_1 = r_2$ by \cref{coord_diff_set_finite,subseteq,real_leq_antisym}.
    \end{subproof}
    Therefore $d$ is a function by \cref{square_metric,relation,rightunique}.
\end{proof}

% from §20 helper lemma 17: absolute value bound from interval
\begin{lemma}\label{real_abs_lt_from_interval}
    Suppose $a\in\reals$.
    Suppose $b\in\reals$.
    Suppose $\epsilon\in\reals$.
    Suppose $0 < \epsilon$.
    Suppose $a - \epsilon < b$ and $b < a + \epsilon$.
    Then $\abs{a - b} < \epsilon$.
\end{lemma}
\begin{proof}
    Hence $b - a\in\reals$ by \cref{real_add_inverse,real_sub,real_add_closed}.
    Hence $b + \neg{a} < (a + \epsilon) + \neg{a}$ by \cref{real_add_inverse,real_add_closed,real_lt_add}.
    Therefore $(a + \epsilon) + \neg{a} = 0 + \epsilon$ by \cref{real_add_assoc,real_add_comm,real_add_inverse}.
    Thus $b - a < \epsilon$ by \cref{real_add_ident,real_sub}.
    $\neg{\epsilon}\in\reals$ by \cref{real_add_inverse}.
    Hence $a + \neg{\epsilon} < b$ by \cref{real_sub,real_lt_subst_left}.
    Therefore $(a + \neg{\epsilon}) + \neg{a} < b + \neg{a}$ by \cref{real_add_inverse,real_add_closed,real_lt_add}.
    Then $(a - \epsilon) + \neg{a} = 0 + \neg{\epsilon}$ by \cref{real_sub,real_add_assoc,real_add_comm,real_add_inverse}.
    Hence $\neg{\epsilon} < b - a$ by \cref{real_add_ident,real_sub}.
    Show that $\abs{b - a} < \epsilon$.
    \begin{subproof}
        \begin{byCase}
        \caseOf{$b - a = 0$.}
            Hence $\abs{b - a} < \epsilon$ by \cref{real_abs_nonneg,real_lt_subst_left}.
        \caseOf{$b - a \neq 0$.}
            Then $b - a < 0$ or $0 < b - a$ by \cref{real_add_ident,real_lt_trichotomy}.
            Hence $\abs{b - a} < \epsilon$
                by \cref{real_abs_neg,real_neg_neg,real_add_inverse,real_lt_neg,real_lt_subst_right,real_abs_nonneg,real_lt_subst_left}.
        \end{byCase}
    \end{subproof}
    Therefore $\abs{a - b} < \epsilon$ by \cref{real_abs_sub_swap}.
\end{proof}

% from §20 helper lemma 18: sum of nonnegative finite sequence is nonnegative
\begin{lemma}\label{sum_finite_sequence_nonneg}
    Suppose $\sumFiniteSequence{a}{s}$.
    Suppose for all $i\in\dom{a}$ we have $0 \leq a(i)$.
    Then $0 \leq s$.
\end{lemma}
\begin{proof}
    Take $n$ such that $n\in\naturalsPlus$ and $a\in\funs{\initialSegment{n}}{\reals}$ and
        there exists $t\in\funs{\initialSegment{n}}{\reals}$ such that
            $t(1) = a(1)$ and $s = t(n)$ and
            for all $k\in\naturalsPlus$ such that $k + 1 \leq n$ we have $t(k + 1) = t(k) + a(k + 1)$
        by \cref{sum_finite_sequence_pred}.
    Take $t$ such that $t\in\funs{\initialSegment{n}}{\reals}$ and
            $t(1) = a(1)$ and $s = t(n)$ and
            for all $k\in\naturalsPlus$ such that $k + 1 \leq n$ we have $t(k + 1) = t(k) + a(k + 1)$.
    Hence $\dom{a} = \initialSegment{n}$ by \cref{funs_elim}.
    Let $A = \{k\in\naturalsPlus \mid \text{if $k \leq n$ then $0 \leq t(k)$}\}$.
    $A\subseteq\naturalsPlus$ by \cref{subseteq}.
    $1\in\naturalsPlus$ and $1 \leq n$ by \cref{real_one_in_naturals,real_one_leq_naturalsplus}.
    Hence $0 \leq t(1)$.
    Therefore $1\in A$.
    Show that for all $k\in A$ we have $k + 1\in A$.
    \begin{subproof}
        Fix $k$ such that $k\in A$.
        $k\in\naturalsPlus$ and $k + 1\in\naturalsPlus$ by \cref{subseteq,real_naturals_closed_succ}.
        Show that if $k + 1 \leq n$ then $0 \leq t(k + 1)$.
        \begin{subproof}
            Assume $k + 1 \leq n$.
            $k\in\reals$ and $1\in\reals$ and $0\in\reals$
                by \cref{real_naturals_subset_reals,real_one_in_naturals,subseteq,real_add_ident}.
            Hence $k < k + 1$
                by \cref{real_zero_lt_one,real_lt_add,real_add_ident,real_add_comm,real_lt_subst_left}.
            $k + 1 < n$ or $k + 1 = n$.
            $k + 1\in\reals$ and $n\in\reals$ by \cref{real_add_closed,real_naturals_subset_reals,subseteq}.
            Hence $k \leq n$ by \cref{real_lt_subst_right,real_lt_trans}.
            Then $k\in\initialSegment{n}$ by \cref{initial_segment_naturalsplus}.
            Then $0 \leq t(k)$.
            Therefore $k + 1\in\initialSegment{n}$ by \cref{initial_segment_naturalsplus}.
            Hence $k + 1\in\dom{a}$ by \cref{funs_elim}.
            Therefore $0 \leq a(k + 1)$.
            Hence $a(k + 1)\in\reals$ and $t(k)\in\reals$ by \cref{funs_type_apply}.
            Then $0 + a(k + 1) \leq t(k) + a(k + 1)$ by \cref{real_leq_add}.
            Thus $0 + a(k + 1)\in\reals$ and $t(k) + a(k + 1)\in\reals$ by \cref{real_add_closed}.
            $0 + a(k + 1) = a(k + 1)$ by \cref{real_add_ident}.
            Hence $0 \leq t(k) + a(k + 1)$ by \cref{real_leq_trans}.
            Then $t(k + 1) = t(k) + a(k + 1)$.
            Therefore $0 \leq t(k + 1)$.
        \end{subproof}
        Hence $k + 1\in A$.
    \end{subproof}
    Therefore for all $k\in\naturalsPlus$ we have $k\in A$ by \cref{real_naturals_induction}.
    Therefore $0 \leq s$ by \cref{subseteq}.
\end{proof}

% from §20 helper lemma 19: terms are bounded by the sum
\begin{lemma}\label{sum_finite_sequence_term_leq}
    Suppose $\sumFiniteSequence{a}{s}$.
    Suppose for all $i\in\dom{a}$ we have $0 \leq a(i)$.
    Then for all $i\in\dom{a}$ we have $a(i) \leq s$.
\end{lemma}
\begin{proof}
    Take $n$ such that $n\in\naturalsPlus$ and $a\in\funs{\initialSegment{n}}{\reals}$ and
        there exists $t\in\funs{\initialSegment{n}}{\reals}$ such that
            $t(1) = a(1)$ and $s = t(n)$ and
            for all $k\in\naturalsPlus$ such that $k + 1 \leq n$ we have $t(k + 1) = t(k) + a(k + 1)$
        by \cref{sum_finite_sequence_pred}.
    Take $t$ such that $t\in\funs{\initialSegment{n}}{\reals}$ and
            $t(1) = a(1)$ and $s = t(n)$ and
            for all $k\in\naturalsPlus$ such that $k + 1 \leq n$ we have $t(k + 1) = t(k) + a(k + 1)$.
    Hence $\dom{a} = \initialSegment{n}$ and $\dom{t} = \initialSegment{n}$ by \cref{funs_elim}.
    Let $A = \{k\in\naturalsPlus \mid \text{if $k \leq n$ then
        for all $j\in\initialSegment{k}$ we have $t(j) \leq t(k)$}\}$.
    $A\subseteq\naturalsPlus$ by \cref{subseteq}.
    For all $j\in\initialSegment{1}$ we have $t(j) \leq t(1)$.
    Therefore $1\in A$.
    Show that for all $k\in A$ we have $k + 1\in A$.
    \begin{subproof}
        Fix $k$ such that $k\in A$.
        $k\in\naturalsPlus$ and $k + 1\in\naturalsPlus$ by \cref{subseteq,real_naturals_closed_succ}.
        Show that if $k + 1 \leq n$ then
            for all $j\in\initialSegment{k + 1}$ we have $t(j) \leq t(k + 1)$.
        \begin{subproof}
            Assume $k + 1 \leq n$.
            $k\in\reals$ and $1\in\reals$ and $0\in\reals$
                by \cref{real_naturals_subset_reals,real_one_in_naturals,subseteq,real_add_ident}.
            Hence $k < k + 1$
                by \cref{real_zero_lt_one,real_lt_add,real_add_ident,real_add_comm,real_lt_subst_left}.
            $k + 1 < n$ or $k + 1 = n$.
            $k + 1\in\reals$ and $n\in\reals$ by \cref{real_add_closed,real_naturals_subset_reals,subseteq}.
            Hence $k \leq n$ by \cref{real_lt_subst_right,real_lt_trans}.
            If $k \leq n$ then for all $j\in\initialSegment{k}$ we have $t(j) \leq t(k)$.
            Hence $0 \leq a(k + 1)$ by \cref{real_naturals_closed_succ,initial_segment_naturalsplus}.
            Therefore $a(k + 1)\in\reals$ and $t(k)\in\reals$ by \cref{initial_segment_naturalsplus,funs_elim,funs_type_apply}.
            Hence $t(k) \leq t(k + 1)$ by \cref{real_leq_add,real_add_ident,real_add_comm,real_leq_trans}.
            Show that for all $j\in\initialSegment{k + 1}$ we have $t(j) \leq t(k + 1)$.
            \begin{subproof}
                Fix $j$ such that $j\in\initialSegment{k + 1}$.
                Thus $j\in\naturalsPlus$ and $j \leq k + 1$ by \cref{initial_segment_naturalsplus}.
                $j = k + 1$ or $j \neq k + 1$.
                \begin{byCase}
                \caseOf{$j = k + 1$.}
                    Then $t(j) \leq t(k + 1)$.
                \caseOf{$j \neq k + 1$.}
                    Hence $j \leq k$ by \cref{real_naturals_leq_succ_elim}.
                    Then $k \leq n$.
                    $j\in\reals$ and $k\in\reals$ and $n\in\reals$
                        by \cref{real_naturals_subset_reals,subseteq}.
                    Therefore $j \leq n$ by \cref{real_leq_trans}.
                    Hence $j\in\initialSegment{k}$ by \cref{initial_segment_naturalsplus}.
                    Then $t(j) \leq t(k)$.
                    Therefore $t(j)\in\reals$ and $t(k + 1)\in\reals$ by \cref{initial_segment_naturalsplus,funs_elim,funs_type_apply}.
                    Hence $t(j) \leq t(k + 1)$ by \cref{real_leq_trans}.
                \end{byCase}
            \end{subproof}
        \end{subproof}
        Therefore $k + 1\in A$.
    \end{subproof}
    Therefore for all $k\in\naturalsPlus$ we have $k\in A$ by \cref{real_naturals_induction}.
    Show that for all $i\in\dom{a}$ we have $a(i) \leq s$.
    \begin{subproof}
        Fix $i$ such that $i\in\dom{a}$.
        Hence $i\in\naturalsPlus$ and $i \leq n$ by \cref{initial_segment_naturalsplus}.
        Therefore $t(i) \leq t(n)$ by \cref{subseteq,initial_segment_naturalsplus}.
        Show that $a(i) \leq t(i)$.
        \begin{subproof}
            $i = 1$ or $i \neq 1$.
            \begin{byCase}
            \caseOf{$i = 1$.}
                Then $a(i) = t(i)$.
            \caseOf{$i \neq 1$.}
                Take $m\in\naturalsPlus$ such that $i = m + 1$ by \cref{real_naturals_predecessor}.
                $m\in\reals$ and $i\in\reals$ and $n\in\reals$ and $1\in\reals$ and $0\in\reals$
                    by \cref{real_naturals_subset_reals,subseteq,real_one_in_naturals,real_add_ident}.
                $0 < 1$ by \cref{real_zero_lt_one}.
                $0 + m < 1 + m$ by \cref{real_lt_add}.
                $0 + m = m$ by \cref{real_add_ident}.
                $1 + m = m + 1$ by \cref{real_add_comm}.
                Hence $m < m + 1$ by \cref{real_lt_subst_left}.
                Then $m < i$.
                Therefore $m \leq i$.
                Hence $m \leq n$ by \cref{real_leq_trans}.
                Therefore $m\in A$ by \cref{subseteq}.
                $1 \leq m$ by \cref{real_one_leq_naturalsplus}.
                Hence $1\in\initialSegment{m}$ by \cref{initial_segment_naturalsplus}.
                Then $t(1) \leq t(m)$.
                Hence $1\in\dom{a}$ by \cref{real_one_in_naturals,real_one_leq_naturalsplus,initial_segment_naturalsplus}.
                Therefore $0 \leq t(1)$.
                Now $m\in\dom{t}$ by \cref{initial_segment_naturalsplus}.
                Then $t(1)\in\reals$ and $t(m)\in\reals$ by \cref{funs_type_apply}.
                Hence $0 \leq t(m)$ by \cref{real_leq_trans}.
                Therefore $0 \leq a(i)$.
                Then $a(i)\in\reals$ and $t(m)\in\reals$ by \cref{funs_type_apply}.
                Hence $t(i) = t(m) + a(i)$.
                Therefore $a(i) \leq t(i)$ by \cref{real_add_closed,real_leq_add,real_add_ident,real_leq_trans}.
            \end{byCase}
        \end{subproof}
        Hence $a(i)\in\reals$ and $t(i)\in\reals$ and $t(n)\in\reals$
            by \cref{funs_type_apply,initial_segment_naturalsplus,funs_elim}.
        Therefore $a(i) \leq s$ by \cref{real_leq_trans}.
    \end{subproof}
\end{proof}

% from §20 helper lemma 20: square metric value exists
\begin{lemma}\label{square_metric_value_exists}
    Suppose $n\in\naturalsPlus$.
    Let $J = \initialSegment{n}$.
    Suppose $x\in\realsPower{J}$.
    Suppose $y\in\realsPower{J}$.
    Then there exists $r\in\reals$ such that $((x,y),r)\in\squareMetric{n}$ and
        for all $i\in J$ we have $\abs{x(i) - y(i)} \leq r$.
\end{lemma}
\begin{proof}
    $1\in\naturalsPlus$ by \cref{real_one_in_naturals}.
    Let $D = \coordDiffSet{n}{x}{y}$.
    $D\subseteq\reals$ and $D$ is finite and $D\neq\emptyset$ by \cref{coord_diff_set_finite,emptyset}.
    Take $r\in D$ such that $r$ is a maximum of $D$ by \cref{finite_real_maximum}.
    Hence $((x,y),r)\in\squareMetric{n}$ by \cref{real_maximum,times_tuple_intro,subseteq,square_metric}.
    For all $i\in J$ we have $\abs{x(i) - y(i)} \leq r$
        by \cref{coord_diff_set,real_maximum}.
\end{proof}
% from §20 helper lemma 21: bounded finite sum
\begin{lemma}\label{sum_finite_sequence_leq_n_times}
    Suppose $n\in\naturalsPlus$.
    Suppose $a\in\funs{\initialSegment{n}}{\reals}$.
    Suppose $\sumFiniteSequence{a}{s}$.
    Suppose $M\in\reals$.
    Suppose for all $i\in\initialSegment{n}$ we have $0 \leq a(i)$ and $a(i) \leq M$.
    Then $s \leq n \rmul M$.
\end{lemma}
\begin{proof}
    Take $n_0$ such that $n_0\in\naturalsPlus$ and $a\in\funs{\initialSegment{n_0}}{\reals}$ and
        there exists $t\in\funs{\initialSegment{n_0}}{\reals}$ such that
            $t(1) = a(1)$ and $s = t(n_0)$ and
            for all $k\in\naturalsPlus$ such that $k + 1 \leq n_0$ we have $t(k + 1) = t(k) + a(k + 1)$
        by \cref{sum_finite_sequence_pred}.
    Take $t$ such that $t\in\funs{\initialSegment{n_0}}{\reals}$ and
        $t(1) = a(1)$ and $s = t(n_0)$ and
        for all $k\in\naturalsPlus$ such that $k + 1 \leq n_0$ we have $t(k + 1) = t(k) + a(k + 1)$.
    Hence $n = n_0$
        by \cref{funs_elim,subseteq_antisymmetric,subseteq,initial_segment_naturalsplus,real_naturals_subset_reals,real_leq_antisym}.
    Therefore $s = t(n)$.
    Now $t\in\funs{\initialSegment{n}}{\reals}$ and $\dom{t} = \initialSegment{n}$ by \cref{funs_elim}.
    Let $A = \{k\in\naturalsPlus \mid \text{if $k \leq n$ then $t(k) \leq k \rmul M$}\}$.
    $A\subseteq\naturalsPlus$ by \cref{subseteq}.
    Hence $1\in\dom{a}$ by \cref{real_one_in_naturals,real_one_leq_naturalsplus,funs_elim,initial_segment_naturalsplus}.
    Therefore $0 \leq a(1)$ and $a(1) \leq M$.
    Then $t(1) \leq 1 \rmul M$
        by \cref{real_naturals_subset_reals,real_one_in_naturals,subseteq,real_mul_ident}.
    Therefore $1\in A$.
    Show that for all $k\in A$ we have $k + 1\in A$.
    \begin{subproof}
        Fix $k$ such that $k\in A$.
        $k\in\naturalsPlus$ and $k + 1\in\naturalsPlus$ by \cref{subseteq,real_naturals_closed_succ}.
        Show that if $k + 1 \leq n$ then $t(k + 1) \leq (k + 1) \rmul M$.
        \begin{subproof}
            Assume $k + 1 \leq n$.
            $k\in\reals$ and $1\in\reals$ and $0\in\reals$
                by \cref{real_naturals_subset_reals,real_one_in_naturals,subseteq,real_add_ident}.
            Hence $k < k + 1$
                by \cref{real_zero_lt_one,real_lt_add,real_add_ident,real_add_comm,real_lt_subst_left}.
            $k + 1 < n$ or $k + 1 = n$.
            $k + 1\in\reals$ and $n\in\reals$ by \cref{real_add_closed,real_naturals_subset_reals,subseteq}.
            Hence $k \leq n$ by \cref{real_lt_subst_right,real_lt_trans}.
            Hence $t(k) \leq k \rmul M$.
            Therefore $0 \leq a(k + 1)$ and $a(k + 1) \leq M$ by \cref{initial_segment_naturalsplus,funs_elim}.
            Hence $a(k + 1)\in\reals$ and $t(k)\in\reals$ by \cref{initial_segment_naturalsplus,funs_elim,funs_type_apply}.
            Then $k \rmul M\in\reals$ by \cref{real_mul_closed}.
            Therefore $t(k) + a(k + 1)\in\reals$ by \cref{real_add_closed}.
            $t(k) + a(k + 1) \leq (k \rmul M) + a(k + 1)$ and
                $a(k + 1) + (k \rmul M) \leq M + (k \rmul M)$ by \cref{real_leq_add}.
            $a(k + 1) + (k \rmul M) = (k \rmul M) + a(k + 1)$ and
                $M + (k \rmul M) = (k \rmul M) + M$ by \cref{real_add_comm}.
            $(k \rmul M) + a(k + 1)\in\reals$ and $(k \rmul M) + M\in\reals$
                by \cref{real_add_closed}.
            Hence $(k \rmul M) + a(k + 1) \leq (k \rmul M) + M$ by \cref{real_leq_trans}.
            Therefore $t(k) + a(k + 1) \leq (k \rmul M) + M$ by \cref{real_leq_trans}.
            $t(k + 1) = t(k) + a(k + 1)$.
            $(k \rmul M) + M = (k + 1) \rmul M$
                by \cref{real_mul_ident,real_add_closed,real_mul_comm,real_distrib}.
            Hence $t(k + 1) \leq (k + 1) \rmul M$ by \cref{real_leq_trans}.
        \end{subproof}
        Therefore $k + 1\in A$.
    \end{subproof}
    Therefore for all $k\in\naturalsPlus$ we have $k\in A$ by \cref{real_naturals_induction}.
    Then $t(n) \leq n \rmul M$ by \cref{subseteq}.
    Therefore $s \leq n \rmul M$.
\end{proof}

% from §20 helper lemma 22: squares preserve order on nonnegative reals
\begin{lemma}\label{real_nonneg_leq_iff_sq_leq}
    Suppose $a\in\reals$.
    Suppose $b\in\reals$.
    Suppose $0 \leq a$.
    Suppose $0 \leq b$.
    Then $a \leq b$ iff $\exp{a}{1+1} \leq \exp{b}{1+1}$.
\end{lemma}
\begin{proof}
    Show that if $a \leq b$ then $\exp{a}{1+1} \leq \exp{b}{1+1}$.
    \begin{subproof}
        Assume $a \leq b$.
        Hence $a < b$ or $a = b$ by \cref{real_leq_split}.
        \begin{byCase}
        \caseOf{$a = b$.}
            Then $\exp{a}{1+1} \leq \exp{b}{1+1}$.
        \caseOf{$a < b$.}
            $0 < a$ or $a = 0$ by \cref{real_leq_split}.
            \begin{byCase}
            \caseOf{$a = 0$.}
                Hence $0 < b$ by \cref{real_lt_subst_left}.
                Therefore $\exp{a}{1+1} \leq \exp{b}{1+1}$
                    by \cref{real_pow_one,real_one_in_naturals,real_pow_succ,real_mul_zero,real_sq_nonnegative}.
            \caseOf{$0 < a$.}
                Therefore $0 < b$ by \cref{real_lt_trans}.
                Hence $\exp{a}{1+1} < \exp{b}{1+1}$ by \cref{real_pos_lt_iff_sq_lt}.
                Therefore $\exp{a}{1+1} \leq \exp{b}{1+1}$.
            \end{byCase}
        \end{byCase}
    \end{subproof}
    Show that if $\exp{a}{1+1} \leq \exp{b}{1+1}$ then $a \leq b$.
    \begin{subproof}
        Assume $\exp{a}{1+1} \leq \exp{b}{1+1}$.
        Suppose not.
        Hence $b < a$.
        $0 < b$ or $b = 0$ by \cref{real_leq_split}.
        \begin{byCase}
        \caseOf{$b = 0$.}
            Hence $0 < a$ by \cref{real_lt_subst_left}.
            Then $a$ is positive.
            Hence $\exp{a}{1+1} = a \rmul a$ by \cref{real_one_in_naturals,real_pow_succ,real_pow_one}.
            Therefore $\exp{a}{1+1}\in\reals$ by \cref{real_mul_closed}.
            Hence $\exp{a}{1+1} > 0$ by \cref{real_mul_pos_pos_gt}.
            Therefore $\exp{b}{1+1} = 0$ by \cref{real_pow_one,real_one_in_naturals,real_pow_succ,real_mul_zero}.
            Then $\exp{b}{1+1}\in\reals$ by \cref{real_add_ident}.
            Hence $\exp{b}{1+1} < \exp{a}{1+1}$ by \cref{real_lt_subst_left}.
            Therefore $\exp{a}{1+1} < \exp{b}{1+1}$ or $\exp{a}{1+1} = \exp{b}{1+1}$ by \cref{real_leq_split}.
            \begin{byCase}
            \caseOf{$\exp{a}{1+1} < \exp{b}{1+1}$.}
                Hence $\exp{a}{1+1} < \exp{a}{1+1}$ by \cref{real_lt_trans}.
                Contradiction by \cref{real_lt_irreflexive}.
            \caseOf{$\exp{a}{1+1} = \exp{b}{1+1}$.}
                Hence $\exp{b}{1+1} < \exp{b}{1+1}$ by \cref{real_lt_subst_right}.
                Contradiction by \cref{real_lt_irreflexive}.
            \end{byCase}
        \caseOf{$0 < b$.}
            Therefore $0 < a$ by \cref{real_lt_trans}.
            Hence $\exp{b}{1+1} < \exp{a}{1+1}$ by \cref{real_pos_lt_iff_sq_lt}.
            Therefore $\exp{a}{1+1} = a \rmul a$ by \cref{real_one_in_naturals,real_pow_succ,real_pow_one}.
            Then $\exp{a}{1+1}\in\reals$ by \cref{real_mul_closed}.
            Hence $\exp{b}{1+1} = b \rmul b$ by \cref{real_one_in_naturals,real_pow_succ,real_pow_one}.
            Therefore $\exp{b}{1+1}\in\reals$ by \cref{real_mul_closed}.
            Then $\exp{a}{1+1} < \exp{b}{1+1}$ or $\exp{a}{1+1} = \exp{b}{1+1}$ by \cref{real_leq_split}.
            \begin{byCase}
            \caseOf{$\exp{a}{1+1} < \exp{b}{1+1}$.}
                Hence $\exp{a}{1+1} < \exp{a}{1+1}$ by \cref{real_lt_trans}.
                Contradiction by \cref{real_lt_irreflexive}.
            \caseOf{$\exp{a}{1+1} = \exp{b}{1+1}$.}
                Hence $\exp{b}{1+1} < \exp{b}{1+1}$ by \cref{real_lt_subst_right}.
                Contradiction by \cref{real_lt_irreflexive}.
            \end{byCase}
        \end{byCase}
    \end{subproof}
\end{proof}

% from §20 helper lemma 23: coordinate square sequence element
\begin{lemma}\label{coord_square_sequence_elem}
    Suppose $n\in\naturalsPlus$.
    Let $J = \initialSegment{n}$.
    Suppose $x\in\realsPower{J}$.
    Suppose $y\in\realsPower{J}$.
    Let $a = \coordSquareSequence{n}{(x,y)}$.
    Suppose $i\in J$.
    Then $(i,\exp{x(i) - y(i)}{1+1})\in a$.
\end{lemma}
\begin{proof}
    Follows by \cref{coord_square_sequence,fst_eq,snd_eq}.
\end{proof}

% from §20 helper lemma 24: euclidean and square metric bounds
\begin{lemma}\label{euclidean_square_metric_bounds}
    Suppose $n\in\naturalsPlus$.
    Let $J = \initialSegment{n}$.
    Suppose $x\in\realsPower{J}$.
    Suppose $y\in\realsPower{J}$.
    Suppose $((x,y),r)\in\euclideanMetric{n}$.
    Suppose $((x,y),m)\in\squareMetric{n}$.
    Then $m \leq r$ and $r \leq \sqrt{n} \rmul m$.
\end{lemma}
\begin{proof}
    $1\in\naturalsPlus$ by \cref{real_one_in_naturals}.
    Let $D = \coordDiffSet{n}{x}{y}$.
    Take $s\in\reals$ such that
        $\sumFiniteSequence{\coordSquareSequence{n}{(x,y)}}{s}$ and $r = \sqrt{s}$
        by \cref{euclidean_metric,pair_eq_iff}.
    Let $a = \coordSquareSequence{n}{(x,y)}$.
    Take $n_1$ such that $n_1\in\naturalsPlus$ and $a\in\funs{\initialSegment{n_1}}{\reals}$ and
        there exists $t\in\funs{\initialSegment{n_1}}{\reals}$ such that
            $t(1) = a(1)$ and $s = t(n_1)$ and
            for all $k\in\naturalsPlus$ such that $k + 1 \leq n_1$ we have $t(k + 1) = t(k) + a(k + 1)$
        by \cref{sum_finite_sequence_pred}.
    Thus $a$ is a function and $\dom{a} = \initialSegment{n_1}$ by \cref{funs_elim}.
    Take $x_0,y_0\in\realsPower{J}$ such that
        $(x,y) = (x_0,y_0)$ and
        $m\in\coordDiffSet{n}{x_0}{y_0}$ and
        for all $t\in\coordDiffSet{n}{x_0}{y_0}$ we have $t \leq m$
        by \cref{square_metric,pair_eq_iff}.
    Hence $x = x_0$ and $y = y_0$ by \cref{pair_eq_iff}.
    Therefore $m\in D$ and for all $t\in D$ we have $t \leq m$.
    Take $i_0\in J$ such that $m = \abs{x(i_0) - y(i_0)}$ by \cref{coord_diff_set}.
    For all $i\in J$ we have $0 \leq a(i)$
        by \cref{coord_square_sequence_elem,reals_power,tuple_power,funs_type_apply,real_add_inverse,real_add_closed,real_sub,real_sq_nonnegative,function_apply_intro}.
    Hence $\dom{a} = J$
        by \cref{coord_square_sequence_elem,dom_intro,subseteq,dom_iff,coord_square_sequence,pair_eq_iff,subseteq_antisymmetric}.
    Hence $a\in\funs{J}{\reals}$ by \cref{funs_elim,funs_intro}.
    Therefore $0 \leq s$ by \cref{sum_finite_sequence_nonneg}.
    $r\in\reals$ and $r \geq 0$ and $r \rmul r = s$ by \cref{positive_square_root}.
    Hence $a(i_0) \leq s$ by \cref{coord_square_sequence_elem,dom_intro,sum_finite_sequence_term_leq}.
    Then $x(i_0) - y(i_0)\in\reals$
        by \cref{reals_power,tuple_power,funs_type_apply,real_add_inverse,real_add_closed,real_sub}.
    Therefore $\exp{m}{1+1} \leq s$
        by \cref{real_abs_square,coord_square_sequence_elem,function_apply_intro}.
    Hence $\exp{r}{1+1} = s$ by \cref{real_pow_succ,real_pow_one}.
    Therefore $m \leq r$ by \cref{coord_diff_set_finite,subseteq,real_abs_nonnegative,real_nonneg_leq_iff_sq_leq}.
        For all $i\in\initialSegment{n}$ we have $a(i) \leq \exp{m}{1+1}$
            by \cref{coord_diff_set,coord_diff_set_finite,subseteq,reals_power,tuple_power,funs_type_apply,real_add_inverse,real_add_closed,real_sub,real_abs_in_reals,real_abs_nonnegative,real_add_ident,real_leq_trans,real_nonneg_leq_iff_sq_leq,real_abs_square,coord_square_sequence_elem,function_apply_intro}.
        For all $i\in\initialSegment{n}$ we have $0 \leq a(i)$ and $a(i) \leq \exp{m}{1+1}$.
    $m\in\reals$ and $\exp{m}{1+1} = m \rmul m$ and $\exp{m}{1+1}\in\reals$
        by \cref{coord_diff_set_finite,subseteq,real_pow_succ,real_pow_one,real_mul_closed}.
    Therefore $s \leq n \rmul \exp{m}{1+1}$ by \cref{sum_finite_sequence_leq_n_times}.
    $n\in\reals$ by \cref{real_naturals_subset_reals,subseteq}.
    $0 < 1$ and $0\in\reals$ by \cref{real_zero_lt_one,real_add_ident}.
    $1\in\reals$ by \cref{real_naturals_subset_reals,subseteq}.
    Hence $0 \leq 1$.
    $1 \leq n$ by \cref{real_one_leq_naturalsplus}.
    Therefore $0 \leq n$ by \cref{real_leq_trans}.
    Hence $0 < n$.
        $\sqrt{n}\in\reals$ and $\sqrt{n} \geq 0$ and $\sqrt{n} \rmul \sqrt{n} = n$ by \cref{positive_square_root}.
        $\sqrt{n} \rmul m\in\reals$ by \cref{real_mul_closed}.
        $\exp{\sqrt{n} \rmul m}{1+1} = (\sqrt{n} \rmul m) \rmul (\sqrt{n} \rmul m)$ by \cref{real_pow_succ,real_pow_one}.
        $(\sqrt{n} \rmul m) \rmul (\sqrt{n} \rmul m) = (\sqrt{n} \rmul \sqrt{n}) \rmul (m \rmul m)$ by \cref{real_mul_assoc,real_mul_comm}.
        Hence $\exp{\sqrt{n} \rmul m}{1+1} = n \rmul (m \rmul m)$.
        $\exp{m}{1+1} = m \rmul m$ by \cref{real_pow_succ,real_pow_one}.
        Therefore $\exp{\sqrt{n} \rmul m}{1+1} = n \rmul \exp{m}{1+1}$.
        Hence $\exp{r}{1+1} \leq \exp{\sqrt{n} \rmul m}{1+1}$ by \cref{real_leq_trans}.
        Show that $\sqrt{n} \rmul m \geq 0$.
        \begin{subproof}
            Thus $m \geq 0$.
            $0 < m$ or $m = 0$ by \cref{real_leq_split}.
            \begin{byCase}
            \caseOf{$m = 0$.}
                Hence $\sqrt{n} \rmul m \geq 0$ by \cref{real_mul_comm,real_mul_zero}.
            \caseOf{$0 < m$.}
                Hence $m$ is positive.
                Suppose not.
                Therefore $\sqrt{n} = 0$.
                Hence $n = 0$ by \cref{positive_square_root,real_mul_zero}.
                Therefore $0 < 0$ by \cref{real_lt_subst_right}.
                Contradiction by \cref{real_lt_irreflexive}.
                Hence $\sqrt{n}$ is positive.
                Therefore $\sqrt{n} \rmul m \geq 0$ by \cref{real_mul_pos_pos_gt}.
            \end{byCase}
        \end{subproof}
        Therefore $r \leq \sqrt{n} \rmul m$ by \cref{real_nonneg_leq_iff_sq_leq}.
    Therefore $m \leq r$ and $r \leq \sqrt{n} \rmul m$.
\end{proof}
% from §20 helper lemma 25: interval bound from absolute value
\begin{lemma}\label{real_interval_from_abs_lt}
    Suppose $a\in\reals$.
    Suppose $b\in\reals$.
    Suppose $\epsilon\in\reals$.
    Suppose $0 < \epsilon$.
    Suppose $\abs{a - b} < \epsilon$.
    Then $a - \epsilon < b$ and $b < a + \epsilon$.
\end{lemma}
\begin{proof}
    Hence $\abs{a - b} \leq \epsilon$ and $\epsilon$ is nonnegative.
    Hence $b - \epsilon \leq a$ by \cref{real_abs_sub_leq_imp_left_bound}.
    Therefore $a \leq b + \epsilon$ by \cref{real_abs_sub_leq_imp_right_bound}.
    $\neg{\epsilon}\in\reals$ and $b - \epsilon\in\reals$ by \cref{real_add_inverse,real_sub,real_add_closed}.
    Hence $b \leq a + \epsilon$
        by \cref{real_leq_add,real_sub,real_add_assoc,real_add_inverse,real_add_comm,real_add_ident,real_leq_trans}.
    Therefore $a - \epsilon \leq b$
        by \cref{real_add_closed,real_leq_add,real_sub,real_add_assoc,real_add_inverse,real_add_ident,real_add_comm}.
    Suppose not.
    Hence $b = a + \epsilon$.
    Therefore $\abs{a - b} = \epsilon$
        by \cref{real_sub,real_add_inverse,real_add_assoc,real_add_comm,real_add_ident,real_abs_nonneg,real_abs_sub_swap}.
    Hence $\epsilon < \epsilon$ by \cref{real_lt_subst_left}.
    Contradiction by \cref{real_lt_irreflexive}.
    Therefore $b < a + \epsilon$ by \cref{real_leq_split}.
    Suppose not.
    Hence $a - \epsilon = b$.
    $a - \epsilon = a + \neg{\epsilon}$ by \cref{real_sub}.
    Therefore $b = a + \neg{\epsilon}$.
    Hence $\neg{b} = \neg{a} + \epsilon$ by \cref{real_neg_addition,real_neg_neg}.
    $a - b = a + \neg{b}$ by \cref{real_sub}.
    $\neg{a}\in\reals$ and $\epsilon\in\reals$ by \cref{real_add_inverse}.
    Hence $a - b = \epsilon$ by \cref{real_add_assoc,real_add_inverse,real_add_ident}.
    Therefore $\abs{a - b} = \epsilon$ by \cref{real_abs_nonneg}.
    Hence $\epsilon < \epsilon$ by \cref{real_lt_subst_left}.
    Contradiction by \cref{real_lt_irreflexive}.
    Therefore $a - \epsilon < b$ by \cref{real_leq_split}.
    Therefore $a - \epsilon < b$ and $b < a + \epsilon$.
\end{proof}

% from §20 helper lemma 26: coordinate square sequence is a function
\begin{lemma}\label{coord_square_sequence_function}
    Suppose $n\in\naturalsPlus$.
    Let $J = \initialSegment{n}$.
    Suppose $x\in\realsPower{J}$.
    Suppose $y\in\realsPower{J}$.
    Then $\coordSquareSequence{n}{(x,y)}\in\funs{J}{\reals}$.
\end{lemma}
\begin{proof}
    Let $a = \coordSquareSequence{n}{(x,y)}$.
    For all $i,t_1,t_2$ such that $(i,t_1)\in a$ and $(i,t_2)\in a$ we have $t_1 = t_2$
        by \cref{coord_square_sequence,pair_eq_iff}.
    Therefore $a\in\funs{J}{\reals}$
        by \cref{coord_square_sequence,relation,rightunique,coord_square_sequence_elem,dom_intro,subseteq,dom_iff,pair_eq_iff,subseteq_antisymmetric,reals_power,tuple_power,funs_type_apply,real_add_inverse,real_add_closed,real_sub,real_one_in_naturals,real_pow_succ,real_pow_one,real_mul_closed,function_apply_intro,funs_intro}.
\end{proof}

% from §20 helper lemma 27: sum finite sequence successor closure
\begin{lemma}\label{sum_finite_sequence_exists_succ}
    Let $A = \{n\in\naturalsPlus \mid \text{for all $a\in\funs{\initialSegment{n}}{\reals}$ there exists $s\in\reals$ such that $\sumFiniteSequence{a}{s}$}\}$.
    Then for all $n\in A$ we have $n + 1\in A$.
\end{lemma}
\begin{proof}
    Fix $n$ such that $n\in A$.
    Thus for all $a\in\funs{\initialSegment{n}}{\reals}$ there exists $s\in\reals$
        such that $\sumFiniteSequence{a}{s}$.
    Thus $n + 1\in\naturalsPlus$ by \cref{subseteq,real_naturals_closed_succ}.
    Show that for all $a\in\funs{\initialSegment{n + 1}}{\reals}$ there exists $s\in\reals$ such that $\sumFiniteSequence{a}{s}$.
    \begin{subproof}
        Fix $a$ such that $a\in\funs{\initialSegment{n + 1}}{\reals}$.
        Let $a_0 = \{(i,a(i)) \mid i\in\initialSegment{n}\}$.
        For all $w$ such that $w\in a_0$ there exist $i,t$ such that $w = (i,t)$.
        For all $i,t_1,t_2$ such that $(i,t_1)\in a_0$ and $(i,t_2)\in a_0$ we have $t_1 = t_2$.
        Thus $a_0$ is a function by \cref{relation,rightunique}.
        Hence $\dom{a_0} = \initialSegment{n}$ by \cref{dom_intro,dom_iff,pair_eq_iff,subseteq,subseteq_antisymmetric}.
        For all $i$ such that $i\in\dom{a_0}$ we have $a_0(i)\in\reals$
            by \cref{subseteq,initial_segment_naturalsplus,real_naturals_subset_reals,real_one_in_naturals,real_add_ident,real_add_closed,real_zero_lt_one,real_lt_add,real_add_comm,real_lt_subst_left,real_lt_trans,funs_elim,funs_type_apply,function_apply_intro}.
                    Hence $a_0\in\funs{\initialSegment{n}}{\reals}$ by \cref{funs_intro}.
                Take $s_0\in\reals$ such that $\sumFiniteSequence{a_0}{s_0}$.
                Take $n_0$ such that $n_0\in\naturalsPlus$ and $a_0\in\funs{\initialSegment{n_0}}{\reals}$ and
                    there exists $t_0\in\funs{\initialSegment{n_0}}{\reals}$ such that
                        $t_0(1) = a_0(1)$ and $s_0 = t_0(n_0)$ and
                        for all $k\in\naturalsPlus$ such that $k + 1 \leq n_0$ we have $t_0(k + 1) = t_0(k) + a_0(k + 1)$
                    by \cref{sum_finite_sequence_pred}.
                Take $t_0$ such that $t_0\in\funs{\initialSegment{n_0}}{\reals}$ and
                    $t_0(1) = a_0(1)$ and $s_0 = t_0(n_0)$ and
                    for all $k\in\naturalsPlus$ such that $k + 1 \leq n_0$ we have $t_0(k + 1) = t_0(k) + a_0(k + 1)$.
                    Hence $\initialSegment{n_0} = \initialSegment{n}$ by \cref{funs_elim,subseteq_antisymmetric,subseteq}.
                    Therefore $\dom{t_0} = \initialSegment{n}$ by \cref{funs_elim}.
                $t_0$ is a function by \cref{funs_elim}.
                $t_0\in\funs{\initialSegment{n}}{\reals}$ by \cref{funs_intro}.
                    Hence $t_0(n)\in\reals$ by \cref{initial_segment_naturalsplus,funs_type_apply}.
                    Therefore $a(n + 1)\in\reals$ by \cref{initial_segment_naturalsplus,funs_type_apply}.
                Let $v = t_0(n) + a(n + 1)$.
                $v\in\reals$ by \cref{real_add_closed}.
                Let $g = \{(n + 1,v)\}$.
                For all $w$ such that $w\in g$ there exist $x,y$ such that $w = (x,y)$.
                For all $x,y_1,y_2$ such that $(x,y_1)\in g$ and $(x,y_2)\in g$ we have $y_1 = y_2$
                    by \cref{singleton_elim,pair_eq_iff}.
                    Therefore $g$ is a function by \cref{relation,rightunique}.
                Show that for all $x$ such that $x\in\dom{t_0}\inter\dom{g}$ we have $t_0(x) = g(x)$.
                \begin{subproof}
                    Fix $x$ such that $x\in\dom{t_0}\inter\dom{g}$.
                    Suppose not.
                        Therefore $x\in\dom{g}$.
                    Take $y$ such that $(x,y)\in g$ by \cref{dom_iff}.
                        Hence $x = n + 1$ by \cref{singleton_elim,pair_eq_iff}.
                        Then $x\in\dom{t_0}$.
                        Hence $x\in\initialSegment{n}$ by \cref{funs_elim}.
                        Therefore $x\in\naturalsPlus$ by \cref{initial_segment_naturalsplus}.
                        Then $x < n$ or $x = n$ by \cref{initial_segment_naturalsplus}.
                    $x\in\reals$ and $n\in\reals$ and $1\in\reals$ and $0\in\reals$
                        by \cref{real_naturals_subset_reals,subseteq,real_one_in_naturals,real_add_ident}.
                    Hence $n < n + 1$
                        by \cref{real_zero_lt_one,real_lt_add,real_add_ident,real_add_comm,real_lt_subst_left}.
                        Therefore $n < x$ by \cref{real_lt_subst_right}.
                    Therefore $n < n$ by \cref{real_lt_trans,real_lt_subst_right}.
                    Contradiction by \cref{real_lt_irreflexive}.
                \end{subproof}
                Let $t = t_0\union g$.
                $t$ is a function and $\dom{t} = \dom{t_0}\union\dom{g}$ by \cref{union_compatible_functions,dom_union}.
                For all $i$ such that $i\in\dom{g}$ we have $i\in\{n + 1\}$
                    by \cref{dom_iff,singleton_elim,pair_eq_iff,singleton_intro}.
                For all $i$ such that $i\in\{n + 1\}$ we have $i\in\dom{g}$
                    by \cref{singleton_elim,singleton_intro,pair_eq_iff,dom_intro}.
                    Hence $\dom{g} = \{n + 1\}$ and $\dom{t} = \initialSegment{n}\union\{n + 1\}$
                        by \cref{subseteq,subseteq_antisymmetric,funs_elim}.
                For all $i$ such that $i\in\dom{t_0}$ we have $t(i)\in\reals$
                    by \cref{funs_type_apply,function_apply_elim,union_iff,function_apply_intro}.
                For all $i$ such that $i\in\dom{g}$ we have $t(i)\in\reals$
                    by \cref{dom_iff,singleton_elim,pair_eq_iff,union_iff,function_apply_intro}.
                For all $i\in\dom{t}$ we have $t(i)\in\reals$.
                    Hence $1\in\dom{t_0}$ by \cref{real_one_leq_naturalsplus,initial_segment_naturalsplus,real_one_in_naturals,funs_elim}.
                Hence $t(1) = t_0(1)$ by \cref{function_apply_elim,union_iff,function_apply_intro}.
                $(1,a(1))\in a_0$.
                    Therefore $a_0(1) = a(1)$ by \cref{funs_elim,function_apply_intro}.
                    Hence $t(1) = a(1)$.
                Show that for all $k\in\naturalsPlus$ such that $k + 1 \leq n + 1$ we have
                    $t(k + 1) = t(k) + a(k + 1)$.
                \begin{subproof}
                    Fix $k$ such that $k\in\naturalsPlus$ and $k + 1 \leq n + 1$.
                    $k + 1\in\naturalsPlus$ by \cref{real_naturals_closed_succ}.
                    $k + 1 = n + 1$ or $k + 1 \neq n + 1$.
                    \begin{byCase}
                    \caseOf{$k + 1 = n + 1$.}
                        $k\in\reals$ and $n\in\reals$ and $1\in\reals$ by \cref{real_naturals_subset_reals,subseteq,real_one_in_naturals}.
                            Hence $k\in\dom{t_0}$ by \cref{real_add_cancel,initial_segment_naturalsplus,funs_elim}.
                        $(k,t_0(k))\in t_0$ by \cref{function_apply_elim}.
                            Therefore $(k,t_0(k))\in t$ by \cref{union_iff}.
                            Hence $t(k) = t_0(k)$ by \cref{function_apply_intro}.
                        $t_0(n)\in\reals$ by \cref{funs_elim,initial_segment_naturalsplus,funs_type_apply}.
                            Therefore $t(k) = t_0(n)$.
                        Hence $(k + 1,v)\in g$ by \cref{singleton_intro,pair_eq_iff}.
                            Hence $(k + 1,v)\in t$ by \cref{union_iff}.
                            Therefore $t(k + 1) = v$ by \cref{function_apply_intro}.
                            Hence $t(k + 1) = t(k) + a(k + 1)$.
                    \caseOf{$k + 1 \neq n + 1$.}
                            Hence $k + 1\in\initialSegment{n}$ by \cref{real_naturals_leq_succ_elim,initial_segment_naturalsplus}.
                        $k\in\reals$ and $n\in\reals$ and $1\in\reals$ and $0\in\reals$
                            by \cref{real_naturals_subset_reals,subseteq,real_one_in_naturals,real_add_ident}.
                        $k + 1\in\reals$ by \cref{real_add_closed}.
                            Hence $k < k + 1$
                                by \cref{real_zero_lt_one,real_lt_add,real_add_ident,real_add_comm,real_lt_subst_left}.
                            Then $k + 1 < n$ or $k + 1 = n$ by \cref{initial_segment_naturalsplus}.
                        Hence $k < n$ by \cref{real_lt_trans,real_lt_subst_right}.
                            Hence $k\in\initialSegment{n}$ by \cref{initial_segment_naturalsplus}.
                        $k\in\dom{t_0}$ and $k + 1\in\dom{t_0}$ by \cref{funs_elim,subseteq}.
                        Hence $t(k) = t_0(k)$ and $t(k + 1) = t_0(k + 1)$
                            by \cref{function_apply_elim,union_iff,function_apply_intro}.
                            Therefore $k + 1\in\initialSegment{n_0}$.
                            Then $k + 1 < n_0$ or $k + 1 = n_0$ by \cref{initial_segment_naturalsplus}.
                            Hence $t_0(k + 1) = t_0(k) + a_0(k + 1)$.
                            Therefore $t(k + 1) = t(k) + a(k + 1)$
                                by \cref{funs_elim,function_apply_intro}.
                    \end{byCase}
                \end{subproof}
                Let $s = t(n + 1)$.
                Hence $s\in\reals$ by \cref{singleton_intro,union_iff,dom_intro,funs_type_apply}.
                $n + 1\in\naturalsPlus$ by \cref{real_naturals_closed_succ}.
                For all $i$ such that $i\in\initialSegment{n}\union\{n + 1\}$ we have $i\in\initialSegment{n + 1}$
                    by \cref{union_iff,initial_segment_naturalsplus,real_naturals_subset_reals,subseteq,real_one_in_naturals,real_add_ident,real_add_closed,real_zero_lt_one,real_lt_add,real_add_comm,real_lt_subst_left,real_lt_trans,singleton_elim}.
                Show that for all $i$ such that $i\in\initialSegment{n + 1}$ we have $i\in\initialSegment{n}\union\{n + 1\}$.
                \begin{subproof}
                    Fix $i$ such that $i\in\initialSegment{n + 1}$.
                        Hence $i\in\naturalsPlus$ by \cref{initial_segment_naturalsplus}.
                        Then $i < n + 1$ or $i = n + 1$ by \cref{initial_segment_naturalsplus}.
                    \begin{byCase}
                    \caseOf{$i = n + 1$.}
                            Hence $i\in\initialSegment{n}\union\{n + 1\}$ by \cref{singleton_intro,union_iff}.
                    \caseOf{$i < n + 1$.}
                            Hence $i \leq n + 1$.
                            Therefore $i \neq n + 1$.
                            Hence $i\in\initialSegment{n}$ by \cref{real_naturals_leq_succ_elim,initial_segment_naturalsplus}.
                            Hence $i\in\initialSegment{n}\union\{n + 1\}$ by \cref{union_iff}.
                    \end{byCase}
                \end{subproof}
                    Therefore $\initialSegment{n}\union\{n + 1\} = \initialSegment{n + 1}$ by \cref{subseteq,subseteq_antisymmetric}.
                    Hence $\dom{t} = \initialSegment{n + 1}$.
                    Therefore $t\in\funs{\initialSegment{n + 1}}{\reals}$ by \cref{funs_intro}.
                    Hence $\sumFiniteSequence{a}{s}$ by \cref{sum_finite_sequence_pred}.
        \end{subproof}
            Therefore $n + 1\in A$.
\end{proof}

% from §20 helper lemma 28: sum of a finite sequence exists
\begin{lemma}\label{sum_finite_sequence_exists}
    Suppose $m\in\naturalsPlus$.
    Suppose $f\in\funs{\initialSegment{m}}{\reals}$.
    Then there exists $s\in\reals$ such that $\sumFiniteSequence{f}{s}$.
\end{lemma}
\begin{proof}
    Let $A = \{n\in\naturalsPlus \mid \text{for all $a\in\funs{\initialSegment{n}}{\reals}$ there exists $s\in\reals$ such that $\sumFiniteSequence{a}{s}$}\}$.
        $A\subseteq\naturalsPlus$ by \cref{subseteq}.
        For all $a\in\funs{\initialSegment{1}}{\reals}$ there exists $s\in\reals$ such that $\sumFiniteSequence{a}{s}$
            by \cref{initial_segment_naturalsplus,real_one_in_naturals,funs_type_apply,real_one_leq_naturalsplus,real_naturals_subset_reals,subseteq,real_add_closed,real_leq_add,real_add_ident,real_zero_lt_one,real_lt_add,real_lt_subst_left,real_lt_trans,real_lt_subst_right,real_lt_irreflexive,sum_finite_sequence_pred}.
        Hence $1\in A$ by \cref{real_one_in_naturals}.
        Then for all $n\in A$ we have $n + 1\in A$ by \cref{sum_finite_sequence_exists_succ}.
        Therefore for all $n\in\naturalsPlus$ we have $n\in A$ by \cref{real_naturals_induction}.
    Therefore there exists $s\in\reals$ such that $\sumFiniteSequence{f}{s}$ by \cref{subseteq}.
\end{proof}

% from §20 helper definition 1: coordinate sum sequence
\begin{definition}\label{coord_sum_sequence}
    $\coordSumSequence{n}{a}{b} = \{w \mid \exists i\in\initialSegment{n}.
        w = (i, a(i) + b(i))\}$.
\end{definition}

% from §20 helper definition 2: coordinate product sequence
\begin{definition}\label{coord_product_sequence}
    $\coordProductSequence{n}{a}{b} = \{w \mid \exists i\in\initialSegment{n}.
        w = (i, a(i) \rmul b(i))\}$.
\end{definition}

% from §20 helper definition 3: coordinate square sequence
\begin{definition}\label{coord_square_sequence_mul}
    $\coordSquareSequenceMul{n}{a} = \{w \mid \exists i\in\initialSegment{n}.
        w = (i, a(i) \rmul a(i))\}$.
\end{definition}

% from §20 helper definition 4: coordinate difference sequence
\begin{definition}\label{coord_diff_sequence}
    $\coordDiffSequence{n}{x}{y} = \{w \mid \exists i\in\initialSegment{n}.
        w = (i, x(i) - y(i))\}$.
\end{definition}

% from §20 helper definition 5: coordinate scalar sequence
\begin{definition}\label{coord_scalar_sequence}
    $\coordScalarSequence{n}{c}{a} = \{w \mid \exists i\in\initialSegment{n}.
        w = (i, c \rmul a(i))\}$.
\end{definition}

% from §20 helper lemma 29: coordinate sum sequence is a function
\begin{lemma}\label{coord_sum_sequence_function}
    Suppose $n\in\naturalsPlus$.
    Let $J = \initialSegment{n}$.
    Suppose $a\in\funs{J}{\reals}$.
    Suppose $b\in\funs{J}{\reals}$.
    Then $\coordSumSequence{n}{a}{b}\in\funs{J}{\reals}$.
\end{lemma}
\begin{proof}
    Let $c = \coordSumSequence{n}{a}{b}$.
    For all $i,t_1,t_2$ such that $(i,t_1)\in c$ and $(i,t_2)\in c$ we have $t_1 = t_2$
        by \cref{coord_sum_sequence,pair_eq_iff}.
    Therefore $c\in\funs{J}{\reals}$
        by \cref{coord_sum_sequence,relation,rightunique,dom_intro,subseteq,dom_iff,pair_eq_iff,subseteq_antisymmetric,funs_type_apply,real_add_closed,function_apply_intro,funs_intro}.
\end{proof}

% from §20 helper lemma 30: coordinate product sequence is a function
\begin{lemma}\label{coord_product_sequence_function}
    Suppose $n\in\naturalsPlus$.
    Let $J = \initialSegment{n}$.
    Suppose $a\in\funs{J}{\reals}$.
    Suppose $b\in\funs{J}{\reals}$.
    Then $\coordProductSequence{n}{a}{b}\in\funs{J}{\reals}$.
\end{lemma}
\begin{proof}
    Let $c = \coordProductSequence{n}{a}{b}$.
    For all $i,t_1,t_2$ such that $(i,t_1)\in c$ and $(i,t_2)\in c$ we have $t_1 = t_2$
        by \cref{coord_product_sequence,pair_eq_iff}.
    Therefore $c\in\funs{J}{\reals}$
        by \cref{coord_product_sequence,relation,rightunique,dom_intro,subseteq,dom_iff,pair_eq_iff,subseteq_antisymmetric,funs_type_apply,real_mul_closed,function_apply_intro,funs_intro}.
\end{proof}

% from §20 helper lemma 31: coordinate square sequence is a function
\begin{lemma}\label{coord_square_sequence_mul_function}
    Suppose $n\in\naturalsPlus$.
    Let $J = \initialSegment{n}$.
    Suppose $a\in\funs{J}{\reals}$.
    Then $\coordSquareSequenceMul{n}{a}\in\funs{J}{\reals}$.
\end{lemma}
\begin{proof}
    Let $c = \coordSquareSequenceMul{n}{a}$.
    For all $i,t_1,t_2$ such that $(i,t_1)\in c$ and $(i,t_2)\in c$ we have $t_1 = t_2$
        by \cref{coord_square_sequence_mul,pair_eq_iff}.
    Therefore $c\in\funs{J}{\reals}$
        by \cref{coord_square_sequence_mul,relation,rightunique,dom_intro,subseteq,dom_iff,pair_eq_iff,subseteq_antisymmetric,funs_type_apply,real_mul_closed,function_apply_intro,funs_intro}.
\end{proof}

% from §20 helper lemma 32: coordinate difference sequence is a function
\begin{lemma}\label{coord_diff_sequence_function}
    Suppose $n\in\naturalsPlus$.
    Let $J = \initialSegment{n}$.
    Suppose $x\in\realsPower{J}$.
    Suppose $y\in\realsPower{J}$.
    Then $\coordDiffSequence{n}{x}{y}\in\funs{J}{\reals}$.
\end{lemma}
\begin{proof}
    Let $c = \coordDiffSequence{n}{x}{y}$.
    For all $i,t_1,t_2$ such that $(i,t_1)\in c$ and $(i,t_2)\in c$ we have $t_1 = t_2$
        by \cref{coord_diff_sequence,pair_eq_iff}.
    Therefore $c\in\funs{J}{\reals}$
        by \cref{coord_diff_sequence,relation,rightunique,dom_intro,subseteq,dom_iff,pair_eq_iff,subseteq_antisymmetric,reals_power,tuple_power,funs_type_apply,real_add_inverse,real_add_closed,real_sub,function_apply_intro,funs_intro}.
\end{proof}

% from §20 helper lemma 33: coordinate scalar sequence is a function
\begin{lemma}\label{coord_scalar_sequence_function}
    Suppose $n\in\naturalsPlus$.
    Let $J = \initialSegment{n}$.
    Suppose $c\in\reals$.
    Suppose $a\in\funs{J}{\reals}$.
    Then $\coordScalarSequence{n}{c}{a}\in\funs{J}{\reals}$.
\end{lemma}
\begin{proof}
    Let $b = \coordScalarSequence{n}{c}{a}$.
    For all $i,t_1,t_2$ such that $(i,t_1)\in b$ and $(i,t_2)\in b$ we have $t_1 = t_2$
        by \cref{coord_scalar_sequence,pair_eq_iff}.
    Therefore $b\in\funs{J}{\reals}$
        by \cref{coord_scalar_sequence,relation,rightunique,dom_intro,subseteq,dom_iff,pair_eq_iff,subseteq_antisymmetric,funs_type_apply,real_mul_closed,function_apply_intro,funs_intro}.
\end{proof}

% from §20 helper lemma 34: strict inequality implies weak inequality
\begin{lemma}\label{real_lt_imp_leq}
    Suppose $a < b$.
    Then $a \leq b$.
\end{lemma}
\begin{proof}
    Trivial.
\end{proof}

% from §20 helper lemma 34a: predecessor stays in initial segment
\begin{lemma}\label{initial_segment_pred_from_succ_bound}
    Suppose $n\in\naturalsPlus$.
    Suppose $k\in\naturalsPlus$ and $k + 1 \leq n$.
    Then $k\in\initialSegment{n}$.
\end{lemma}
\begin{proof}
    Follows by \cref{real_add_ident,real_naturals_subset_reals,real_one_in_naturals,subseteq,real_zero_lt_one,real_lt_add,real_add_comm,real_lt_subst_left,real_lt_imp_leq,real_add_closed,real_leq_trans,initial_segment_naturalsplus}.
\end{proof}

% from §20 helper lemma 34b: initial segment is contained in its successor segment
\begin{lemma}\label{initial_segment_subset_succ}
    Suppose $n\in\naturalsPlus$.
    Then $\initialSegment{n}\subseteq\initialSegment{n + 1}$.
\end{lemma}
\begin{proof}
    Follows by \cref{initial_segment_naturalsplus,real_naturals_subset_reals,subseteq,real_add_ident,real_one_in_naturals,real_zero_lt_one,real_lt_add,real_add_comm,real_add_closed,real_lt_subst_left,real_lt_imp_leq,real_leq_trans}.
\end{proof}
% from §20 helper lemma 35: sum of a coordinate sum sequence
\begin{lemma}\label{sum_finite_sequence_add}
    Suppose $n\in\naturalsPlus$.
    Let $J = \initialSegment{n}$.
    Suppose $a\in\funs{J}{\reals}$.
    Suppose $b\in\funs{J}{\reals}$.
    Suppose $\sumFiniteSequence{a}{s}$.
    Suppose $\sumFiniteSequence{b}{t}$.
    Let $c = \coordSumSequence{n}{a}{b}$.
    Then $\sumFiniteSequence{c}{s + t}$.
\end{lemma}
\begin{proof}
    Take $n_1$ such that $n_1\in\naturalsPlus$ and $a\in\funs{\initialSegment{n_1}}{\reals}$ and
        there exists $u\in\funs{\initialSegment{n_1}}{\reals}$ such that
            $u(1) = a(1)$ and $s = u(n_1)$ and
            for all $k\in\naturalsPlus$ such that $k + 1 \leq n_1$ we have $u(k + 1) = u(k) + a(k + 1)$
        by \cref{sum_finite_sequence_pred}.
    Take $n_2$ such that $n_2\in\naturalsPlus$ and $b\in\funs{\initialSegment{n_2}}{\reals}$ and
        there exists $v\in\funs{\initialSegment{n_2}}{\reals}$ such that
            $v(1) = b(1)$ and $t = v(n_2)$ and
            for all $k\in\naturalsPlus$ such that $k + 1 \leq n_2$ we have $v(k + 1) = v(k) + b(k + 1)$
        by \cref{sum_finite_sequence_pred}.
    Take $u$ such that $u\in\funs{\initialSegment{n_1}}{\reals}$ and
        $u(1) = a(1)$ and $s = u(n_1)$ and
        for all $k\in\naturalsPlus$ such that $k + 1 \leq n_1$ we have $u(k + 1) = u(k) + a(k + 1)$.
    Take $v$ such that $v\in\funs{\initialSegment{n_2}}{\reals}$ and
        $v(1) = b(1)$ and $t = v(n_2)$ and
        for all $k\in\naturalsPlus$ such that $k + 1 \leq n_2$ we have $v(k + 1) = v(k) + b(k + 1)$.
    Thus $n_1 = n$ and $n_2 = n$
        by \cref{subseteq_antisymmetric,subseteq,funs_elim,initial_segment_naturalsplus,real_naturals_subset_reals,real_leq_antisym}.
    Let $w = \coordSumSequence{n}{u}{v}$.
    Then $w\in\funs{J}{\reals}$ and $w(1) = c(1)$
        by \cref{real_one_in_naturals,real_one_leq_naturalsplus,initial_segment_naturalsplus,coord_sum_sequence,funs_elim,function_apply_intro,coord_sum_sequence_function}.
    Show that for all $k\in\naturalsPlus$ such that $k + 1 \leq n$ we have $w(k + 1) = w(k) + c(k + 1)$.
    \begin{subproof}
        Fix $k$ such that $k\in\naturalsPlus$ and $k + 1 \leq n$.
        Then $k + 1\in J$ and $k\in J$
            by \cref{real_naturals_closed_succ,initial_segment_naturalsplus,initial_segment_pred_from_succ_bound}.
        Thus $w(k + 1) = u(k + 1) + v(k + 1)$ and $w(k) = u(k) + v(k)$
            by \cref{coord_sum_sequence,funs_elim,function_apply_intro}.
        Thus $a(k + 1)\in\reals$ and $b(k + 1)\in\reals$ and $u(k)\in\reals$ and $v(k)\in\reals$
            by \cref{funs_type_apply}.
        Thus $u(k) + a(k + 1)\in\reals$ and $v(k) + b(k + 1)\in\reals$ and $u(k) + v(k)\in\reals$
            and $a(k + 1) + b(k + 1)\in\reals$ and $v(k) + a(k + 1)\in\reals$ by \cref{real_add_closed}.
        Thus $u(k + 1) = u(k) + a(k + 1)$ and $v(k + 1) = v(k) + b(k + 1)$.
        Thus $u(k + 1) + v(k + 1) = (u(k) + a(k + 1)) + (v(k) + b(k + 1))$ by \cref{real_add_eq_subst}.
        Thus $(u(k) + a(k + 1)) + (v(k) + b(k + 1)) = (u(k) + v(k)) + (a(k + 1) + b(k + 1))$
            by \cref{real_add_assoc,real_add_comm}.
        Thus $w(k + 1) = w(k) + (a(k + 1) + b(k + 1))$ by \cref{real_add_eq_subst}.
        Therefore $w(k + 1) = w(k) + c(k + 1)$
            by \cref{coord_sum_sequence,coord_sum_sequence_function,funs_elim,function_apply_intro,real_add_eq_subst}.
    \end{subproof}
    Hence $s + t = w(n)$
        by \cref{initial_segment_naturalsplus,coord_sum_sequence,funs_elim,function_apply_intro}.
    Therefore $\sumFiniteSequence{c}{s + t}$ by \cref{coord_sum_sequence_function,sum_finite_sequence_pred}.
\end{proof}

% from §20 helper lemma 36: sum of a coordinate scalar sequence
\begin{lemma}\label{sum_finite_sequence_scalar}
    Suppose $n\in\naturalsPlus$.
    Let $J = \initialSegment{n}$.
    Suppose $c\in\reals$.
    Suppose $a\in\funs{J}{\reals}$.
    Suppose $\sumFiniteSequence{a}{s}$.
    Let $b = \coordScalarSequence{n}{c}{a}$.
    Then $\sumFiniteSequence{b}{c \rmul s}$.
\end{lemma}
\begin{proof}
    Take $n_1$ such that $n_1\in\naturalsPlus$ and $a\in\funs{\initialSegment{n_1}}{\reals}$ and
        there exists $u\in\funs{\initialSegment{n_1}}{\reals}$ such that
            $u(1) = a(1)$ and $s = u(n_1)$ and
            for all $k\in\naturalsPlus$ such that $k + 1 \leq n_1$ we have $u(k + 1) = u(k) + a(k + 1)$
        by \cref{sum_finite_sequence_pred}.
    Take $u$ such that $u\in\funs{\initialSegment{n_1}}{\reals}$ and
        $u(1) = a(1)$ and $s = u(n_1)$ and
        for all $k\in\naturalsPlus$ such that $k + 1 \leq n_1$ we have $u(k + 1) = u(k) + a(k + 1)$.
    Thus $n_1 = n$
        by \cref{subseteq_antisymmetric,subseteq,funs_elim,initial_segment_naturalsplus,real_naturals_subset_reals,real_leq_antisym}.
    Let $w = \coordScalarSequence{n}{c}{u}$.
    Then $w\in\funs{J}{\reals}$ and $w(1) = b(1)$
        by \cref{real_one_in_naturals,real_one_leq_naturalsplus,initial_segment_naturalsplus,coord_scalar_sequence,funs_elim,function_apply_intro,coord_scalar_sequence_function}.
    For all $k\in\naturalsPlus$ such that $k + 1 \leq n$ we have $w(k + 1) = w(k) + b(k + 1)$
        by \cref{real_naturals_closed_succ,initial_segment_naturalsplus,initial_segment_pred_from_succ_bound,coord_scalar_sequence,funs_elim,function_apply_intro,funs_type_apply,real_add_eq_subst,real_distrib,coord_scalar_sequence_function}.
    Hence $c \rmul s = w(n)$
        by \cref{initial_segment_naturalsplus,coord_scalar_sequence,funs_elim,function_apply_intro}.
    Therefore $\sumFiniteSequence{b}{c \rmul s}$ by \cref{coord_scalar_sequence_function,sum_finite_sequence_pred}.
\end{proof}

% from §20 helper lemma 37: split sum at the last term
\begin{lemma}\label{sum_finite_sequence_split_last}
    Suppose $n\in\naturalsPlus$.
    Let $J = \initialSegment{n}$.
    Let $J_1 = \initialSegment{n + 1}$.
    Suppose $a\in\funs{J_1}{\reals}$.
    Suppose $\sumFiniteSequence{a}{s}$.
    Let $a_0 = \restrl{a}{J}$.
    Then $a_0\in\funs{J}{\reals}$ and
        there exists $s_0\in\reals$ such that $\sumFiniteSequence{a_0}{s_0}$ and $s = s_0 + a(n + 1)$.
\end{lemma}
\begin{proof}
    $a$ is a function and $\dom{a} = J_1$ by \cref{funs_elim}.
    Then $a_0$ is a function and $\dom{a_0} = J_1\inter J$ by \cref{restrl_of_function_is_function,dom_of_restrl_eq_inter}.
    Then $J\subseteq J_1$ by \cref{initial_segment_subset_succ}.
    Therefore $\dom{a_0} = J$ by \cref{initial_segment_subset_succ,inter_absorb_supseteq_left,subseteq}.
    For all $i$ such that $i\in\dom{a_0}$ we have $a_0(i)\in\reals$
        by \cref{subseteq,funs_elim,funs_type_apply,restrl_of_function_apply}.
    Hence $a_0\in\funs{J}{\reals}$ by \cref{funs_intro}.
    Take $n_0$ such that $n_0\in\naturalsPlus$ and $a\in\funs{\initialSegment{n_0}}{\reals}$ and
        there exists $t\in\funs{\initialSegment{n_0}}{\reals}$ such that
            $t(1) = a(1)$ and $s = t(n_0)$ and
            for all $k\in\naturalsPlus$ such that $k + 1 \leq n_0$ we have $t(k + 1) = t(k) + a(k + 1)$
        by \cref{sum_finite_sequence_pred}.
    Take $t$ such that $t\in\funs{\initialSegment{n_0}}{\reals}$ and
        $t(1) = a(1)$ and $s = t(n_0)$ and
        for all $k\in\naturalsPlus$ such that $k + 1 \leq n_0$ we have $t(k + 1) = t(k) + a(k + 1)$.
    Then $\dom{a} = \initialSegment{n_0}$ by \cref{funs_elim}.
    Hence $\initialSegment{n_0} = J_1$ by \cref{subseteq_antisymmetric,subseteq,funs_elim}.
    Therefore $n_0 = n + 1$
        by \cref{initial_segment_naturalsplus,subseteq,real_naturals_closed_succ,real_naturals_subset_reals,real_leq_antisym}.
    Let $t_0 = \restrl{t}{J}$.
    $t$ is a function and $\dom{t} = J_1$ by \cref{funs_elim}.
    Then $t_0$ is a function and $\dom{t_0} = J_1\inter J$ by \cref{restrl_of_function_is_function,dom_of_restrl_eq_inter}.
    Then $J\subseteq J_1$ by \cref{initial_segment_subset_succ}.
    Therefore $\dom{t_0} = J$ by \cref{initial_segment_subset_succ,inter_absorb_supseteq_left,subseteq}.
    For all $i$ such that $i\in\dom{t_0}$ we have $t_0(i)\in\reals$
        by \cref{subseteq,funs_elim,funs_type_apply,restrl_of_function_apply}.
    Hence $t_0\in\funs{J}{\reals}$ by \cref{funs_intro}.
    Thus $1\in J$ by \cref{real_one_in_naturals,real_one_leq_naturalsplus,initial_segment_naturalsplus}.
    Thus $t_0(1) = a_0(1)$ by \cref{restrl_of_function_apply,subseteq}.
    For all $k\in\naturalsPlus$ such that $k + 1 \leq n$ we have
        $t_0(k + 1) = t_0(k) + a_0(k + 1)$
        by \cref{real_naturals_closed_succ,initial_segment_naturalsplus,initial_segment_pred_from_succ_bound,restrl_of_function_apply,subseteq,real_add_eq_subst}.
    Let $s_0 = t_0(n)$.
    $n\in J$ by \cref{initial_segment_naturalsplus}.
    Hence $\sumFiniteSequence{a_0}{s_0}$ by \cref{funs_type_apply,sum_finite_sequence_pred}.
    $n + 1\in J_1$ by \cref{initial_segment_naturalsplus,real_naturals_closed_succ}.
    Thus $t_0(n) + a(n + 1) = s$
        by \cref{restrl_of_function_apply,subseteq,real_add_eq_subst}.
    Therefore there exists $s_0\in\reals$ such that $\sumFiniteSequence{a_0}{s_0}$ and $s = s_0 + a(n + 1)$ by \cref{real_add_comm}.
\end{proof}

% from §20 helper lemma 38: zero sum of a nonnegative finite sequence
\begin{lemma}\label{sum_finite_sequence_zero_terms}
    Suppose $n\in\naturalsPlus$.
    Let $J = \initialSegment{n}$.
    Suppose $a\in\funs{J}{\reals}$.
    Suppose $\sumFiniteSequence{a}{s}$.
    Suppose $s = 0$.
    Suppose for all $i\in J$ we have $0 \leq a(i)$.
    Then for all $i\in J$ we have $a(i) = 0$.
\end{lemma}
\begin{proof}
    Let $A = \{m\in\naturalsPlus \mid \text{for all $b\in\funs{\initialSegment{m}}{\reals}$ for all $t\in\reals$
        if $\sumFiniteSequence{b}{t}$ and $t = 0$ and for all $i\in\initialSegment{m}$ we have $0 \leq b(i)$ then
        for all $i\in\initialSegment{m}$ we have $b(i) = 0$}\}$.
    $A\subseteq\naturalsPlus$ by \cref{subseteq}.
    Show that for all $b\in\funs{\initialSegment{1}}{\reals}$ for all $t\in\reals$
        if $\sumFiniteSequence{b}{t}$ and $t = 0$ and for all $i\in\initialSegment{1}$ we have $0 \leq b(i)$ then
        for all $i\in\initialSegment{1}$ we have $b(i) = 0$.
    \begin{subproof}
            Fix $b\in\funs{\initialSegment{1}}{\reals}$.
            Fix $t\in\reals$.
            Assume $\sumFiniteSequence{b}{t}$ and $t = 0$ and for all $i\in\initialSegment{1}$ we have $0 \leq b(i)$.
            Take $n_0$ such that $n_0\in\naturalsPlus$ and $b\in\funs{\initialSegment{n_0}}{\reals}$ and
                there exists $u\in\funs{\initialSegment{n_0}}{\reals}$ such that
                    $u(1) = b(1)$ and $t = u(n_0)$ and
                    for all $k\in\naturalsPlus$ such that $k + 1 \leq n_0$ we have $u(k + 1) = u(k) + b(k + 1)$
                by \cref{sum_finite_sequence_pred}.
            Thus $\initialSegment{n_0} = \initialSegment{1}$ by \cref{subseteq_antisymmetric,subseteq,funs_elim}.
            Thus $n_0 = 1$
                by \cref{initial_segment_naturalsplus,subseteq,real_one_in_naturals,real_naturals_subset_reals,real_leq_antisym}.
            Take $u$ such that $u\in\funs{\initialSegment{n_0}}{\reals}$ and
                $u(1) = b(1)$ and $t = u(n_0)$ and
                for all $k\in\naturalsPlus$ such that $k + 1 \leq n_0$ we have $u(k + 1) = u(k) + b(k + 1)$.
            Thus $b(1) = 0$.
            For all $i\in\initialSegment{1}$ we have $b(i) = 0$
                by \cref{initial_segment_naturalsplus,real_one_leq_naturalsplus,real_naturals_subset_reals,subseteq,real_leq_antisym}.
    \end{subproof}
    Hence $1\in A$ by \cref{real_one_in_naturals}.
    Show that for all $k\in A$ we have $k + 1\in A$.
    \begin{subproof}
        Fix $k$ such that $k\in A$.
        Show that for all $b\in\funs{\initialSegment{k + 1}}{\reals}$ for all $t\in\reals$
            if $\sumFiniteSequence{b}{t}$ and $t = 0$ and for all $i\in\initialSegment{k + 1}$ we have $0 \leq b(i)$ then
            for all $i\in\initialSegment{k + 1}$ we have $b(i) = 0$.
        \begin{subproof}
                Fix $b\in\funs{\initialSegment{k + 1}}{\reals}$.
                Fix $t\in\reals$.
                Assume $\sumFiniteSequence{b}{t}$ and $t = 0$ and for all $i\in\initialSegment{k + 1}$ we have $0 \leq b(i)$.
                Let $J = \initialSegment{k}$.
                Let $J_1 = \initialSegment{k + 1}$.
                Let $b_0 = \restrl{b}{J}$.
                $k\in\naturalsPlus$ by \cref{subseteq}.
                Then $b_0\in\funs{J}{\reals}$ and
                    there exists $s_0\in\reals$ such that $\sumFiniteSequence{b_0}{s_0}$ and $t = s_0 + b(k + 1)$
                    by \cref{sum_finite_sequence_split_last}.
                Take $s_0\in\reals$ such that $\sumFiniteSequence{b_0}{s_0}$ and $t = s_0 + b(k + 1)$.
                For all $i\in J$ we have $0 \leq b_0(i)$
                    by \cref{initial_segment_subset_succ,subseteq,funs_elim,restrl_of_function_apply}.
                Thus for all $i\in\dom{b_0}$ we have $0 \leq b_0(i)$ by \cref{funs_elim,subseteq}.
                Hence $0 \leq s_0$ by \cref{sum_finite_sequence_nonneg}.
                $k + 1\in\naturalsPlus$ by \cref{real_naturals_closed_succ}.
                Thus $k + 1\in J_1$ by \cref{initial_segment_naturalsplus}.
                Thus $0 \leq b(k + 1)$.
                Thus $s_0 + b(k + 1) = 0$.
                $s_0\in\reals$ and $b(k + 1)\in\reals$ by \cref{funs_type_apply,restrl_of_function_apply,subseteq}.
                Then $s_0 \leq s_0 + b(k + 1)$ by \cref{real_leq_add,real_add_ident,real_add_comm}.
                Hence $s_0 + b(k + 1) \leq 0$.
                Therefore $s_0 \leq 0$ by \cref{real_leq_trans}.
                Thus $s_0 = 0$ by \cref{real_leq_antisym}.
                Hence $b(k + 1) = 0$ by \cref{real_add_eq_subst,funs_type_apply,real_add_ident}.
                Thus for all $b_1\in\funs{J}{\reals}$ for all $t_1\in\reals$
                    if $\sumFiniteSequence{b_1}{t_1}$ and $t_1 = 0$ and for all $i\in J$ we have $0 \leq b_1(i)$ then
                    for all $i\in J$ we have $b_1(i) = 0$.
                Therefore for all $i\in J$ we have $b_0(i) = 0$.
                For all $i\in J$ we have $b(i) = 0$
                    by \cref{funs_elim,initial_segment_subset_succ,subseteq,restrl_of_function_apply}.
                For all $i\in J_1$ we have $b(i) = 0$
                    by \cref{initial_segment_naturalsplus,real_naturals_subset_reals,subseteq,real_leq_split,real_naturals_leq_succ_elim}.
        \end{subproof}
        Therefore $k + 1\in A$ by \cref{real_naturals_closed_succ}.
    \end{subproof}
    Hence $n\in A$ by \cref{real_naturals_induction,subseteq}.
    Then for all $b_1\in\funs{J}{\reals}$ for all $t_1\in\reals$
        if $\sumFiniteSequence{b_1}{t_1}$ and $t_1 = 0$ and for all $i\in J$ we have $0 \leq b_1(i)$ then
        for all $i\in J$ we have $b_1(i) = 0$.
    Therefore for all $i\in J$ we have $a(i) = 0$.
\end{proof}

% from §20 helper lemma 39: sum of squares of a scalar sequence
\begin{lemma}\label{sum_finite_sequence_square_scalar}
    Suppose $n\in\naturalsPlus$.
    Let $J = \initialSegment{n}$.
    Suppose $c\in\reals$.
    Suppose $a\in\funs{J}{\reals}$.
    Suppose $\sumFiniteSequence{\coordSquareSequenceMul{n}{a}}{s}$.
    Let $b = \coordScalarSequence{n}{c}{a}$.
    Then $\sumFiniteSequence{\coordSquareSequenceMul{n}{b}}{(c \rmul c) \rmul s}$.
\end{lemma}
\begin{proof}
    Let $u = \coordSquareSequenceMul{n}{a}$.
    Let $v = \coordSquareSequenceMul{n}{b}$.
    Let $w = \coordScalarSequence{n}{c \rmul c}{u}$.
    Hence $b\in\funs{J}{\reals}$ and $u\in\funs{J}{\reals}$ and $v\in\funs{J}{\reals}$
        by \cref{coord_scalar_sequence_function,coord_square_sequence_mul_function}.
    $c \rmul c\in\reals$ by \cref{real_mul_closed}.
    Therefore $w\in\funs{J}{\reals}$ by \cref{coord_scalar_sequence_function}.
    For all $i\in J$ we have $v(i) = w(i)$
        by \cref{coord_scalar_sequence,coord_square_sequence_mul,funs_elim,function_apply_intro,funs_type_apply,real_mul_closed,real_mul_assoc,real_mul_comm}.
    $v$ is a function and $w$ is a function and $\dom{v} = J$ and $\dom{w} = J$ by \cref{funs_elim}.
    Hence $v = w$ by \cref{funs_elim,subseteq,fun_subseteq,subseteq_antisymmetric}.
    Therefore $\sumFiniteSequence{v}{(c \rmul c) \rmul s}$ by \cref{sum_finite_sequence_scalar}.
\end{proof}
% from §20 helper lemma 40: sum of squares of a coordinate sum sequence
\begin{lemma}\label{sum_finite_sequence_square_sum}
    Suppose $n\in\naturalsPlus$.
    Let $J = \initialSegment{n}$.
    Suppose $a\in\funs{J}{\reals}$.
    Suppose $b\in\funs{J}{\reals}$.
    Suppose $\sumFiniteSequence{\coordSquareSequenceMul{n}{a}}{s_1}$.
    Suppose $\sumFiniteSequence{\coordSquareSequenceMul{n}{b}}{s_2}$.
    Suppose $\sumFiniteSequence{\coordProductSequence{n}{a}{b}}{s}$.
    Let $c = \coordSumSequence{n}{a}{b}$.
    Then $\sumFiniteSequence{\coordSquareSequenceMul{n}{c}}{s_1 + s_2 + (1+1) \rmul s}$.
\end{lemma}
\begin{proof}
    Let $u = \coordSquareSequenceMul{n}{a}$.
    Let $v = \coordSquareSequenceMul{n}{b}$.
    Let $p = \coordProductSequence{n}{a}{b}$.
    $1\in\naturalsPlus$ and $1\in\reals$ and $1 + 1\in\reals$
        by \cref{real_one_in_naturals,real_naturals_subset_reals,subseteq,real_add_closed}.
    Let $p_2 = \coordScalarSequence{n}{1+1}{p}$.
    Let $w = \coordSumSequence{n}{u}{v}$.
    Let $d = \coordSumSequence{n}{w}{p_2}$.
    Let $q = \coordSquareSequenceMul{n}{c}$.
    $u\in\funs{J}{\reals}$ and $v\in\funs{J}{\reals}$ by \cref{coord_square_sequence_mul_function}.
    $p\in\funs{J}{\reals}$ by \cref{coord_product_sequence_function}.
    Hence $p_2\in\funs{J}{\reals}$ and $w\in\funs{J}{\reals}$ and $d\in\funs{J}{\reals}$ and
        $c\in\funs{J}{\reals}$ and $q\in\funs{J}{\reals}$
        by \cref{coord_scalar_sequence_function,coord_sum_sequence_function,coord_square_sequence_mul_function}.
    Show that for all $i\in J$ we have $q(i) = d(i)$.
    \begin{subproof}
            Fix $i$ such that $i\in J$.
            Thus $a(i)\in\reals$ and $b(i)\in\reals$ by \cref{funs_type_apply}.
            Thus $a(i) + b(i)\in\reals$ by \cref{real_add_closed}.
            Thus $c(i) = a(i) + b(i)$
                by \cref{real_add_closed,coord_sum_sequence,funs_elim,function_apply_intro}.
            Thus $q(i) = c(i) \rmul c(i)$ by \cref{coord_square_sequence_mul,funs_elim,function_apply_intro}.
            Then $q(i) = (a(i) + b(i)) \rmul (a(i) + b(i))$.
            Hence $(a(i) + b(i)) \rmul (a(i) + b(i)) =
                (a(i) + b(i)) \rmul a(i) + (a(i) + b(i)) \rmul b(i)$ by \cref{real_distrib}.
            Thus $(a(i) + b(i)) \rmul a(i) = a(i) \rmul (a(i) + b(i))$ by \cref{real_mul_comm}.
            Thus $a(i) \rmul (a(i) + b(i)) = (a(i) \rmul a(i)) + (a(i) \rmul b(i))$ by \cref{real_distrib}.
            Thus $(a(i) + b(i)) \rmul a(i) = (a(i) \rmul a(i)) + (a(i) \rmul b(i))$.
            Thus $(a(i) + b(i)) \rmul b(i) = b(i) \rmul (a(i) + b(i))$ by \cref{real_mul_comm}.
            Thus $b(i) \rmul (a(i) + b(i)) = (b(i) \rmul a(i)) + (b(i) \rmul b(i))$ by \cref{real_distrib}.
            Thus $(a(i) + b(i)) \rmul b(i) = (b(i) \rmul a(i)) + (b(i) \rmul b(i))$.
            Thus $q(i) =
                ((a(i) \rmul a(i)) + (a(i) \rmul b(i))) + ((b(i) \rmul a(i)) + (b(i) \rmul b(i)))$ by \cref{real_add_eq_subst}.
            Thus $a(i) \rmul a(i)\in\reals$ and $a(i) \rmul b(i)\in\reals$ and $b(i) \rmul b(i)\in\reals$
                and $b(i) \rmul a(i) = a(i) \rmul b(i)$ by \cref{real_mul_closed,real_mul_comm}.
            Thus $q(i) =
                ((a(i) \rmul a(i)) + (a(i) \rmul b(i))) + ((a(i) \rmul b(i)) + (b(i) \rmul b(i)))$ by \cref{real_add_eq_subst}.
            Thus $(a(i) \rmul b(i)) + (b(i) \rmul b(i))\in\reals$ by \cref{real_add_closed}.
            Thus $q(i) =
                (a(i) \rmul a(i)) + ((a(i) \rmul b(i)) + ((a(i) \rmul b(i)) + (b(i) \rmul b(i))))$ by \cref{real_add_assoc}.
            Thus $(a(i) \rmul b(i)) + ((a(i) \rmul b(i)) + (b(i) \rmul b(i))) =
                ((a(i) \rmul b(i)) + (a(i) \rmul b(i))) + (b(i) \rmul b(i))$ by \cref{real_add_assoc}.
            Thus $q(i) =
                (a(i) \rmul a(i)) + (((a(i) \rmul b(i)) + (a(i) \rmul b(i))) + (b(i) \rmul b(i)))$ by \cref{real_add_eq_subst}.
            Thus $(a(i) \rmul b(i)) + (a(i) \rmul b(i))\in\reals$ by \cref{real_add_closed}.
            Thus $q(i) =
                ((a(i) \rmul a(i)) + ((a(i) \rmul b(i)) + (a(i) \rmul b(i)))) + (b(i) \rmul b(i))$ by \cref{real_add_assoc}.
            Thus $(1 + 1) \rmul (a(i) \rmul b(i)) = (a(i) \rmul b(i)) \rmul (1 + 1)$ by \cref{real_mul_comm}.
            Thus $(a(i) \rmul b(i)) \rmul (1 + 1) =
                (a(i) \rmul b(i)) \rmul 1 + (a(i) \rmul b(i)) \rmul 1$ by \cref{real_distrib}.
            Thus $(a(i) \rmul b(i)) \rmul 1 = a(i) \rmul b(i)$ by \cref{real_mul_comm,real_mul_ident}.
            Thus $(1 + 1) \rmul (a(i) \rmul b(i)) =
                (a(i) \rmul b(i)) + (a(i) \rmul b(i))$ by \cref{real_add_eq_subst}.
            Thus $q(i) =
                ((a(i) \rmul a(i)) + ((1 + 1) \rmul (a(i) \rmul b(i)))) + (b(i) \rmul b(i))$ by \cref{real_add_eq_subst}.
            Thus $q(i) =
                (a(i) \rmul a(i)) + ((1 + 1) \rmul (a(i) \rmul b(i))) + (b(i) \rmul b(i))$ by \cref{real_add_assoc}.
            Thus $u(i) = a(i) \rmul a(i)$ and $v(i) = b(i) \rmul b(i)$ and $p(i) = a(i) \rmul b(i)$
                by \cref{coord_square_sequence_mul,coord_product_sequence,funs_elim,function_apply_intro}.
            Thus $p_2(i) = (1 + 1) \rmul p(i)$ by \cref{coord_scalar_sequence,funs_elim,function_apply_intro}.
            Then $p_2(i) = (1 + 1) \rmul (a(i) \rmul b(i))$.
            Hence $w(i) = u(i) + v(i)$ by \cref{coord_sum_sequence,funs_elim,function_apply_intro}.
            Therefore $d(i) = w(i) + p_2(i)$ by \cref{coord_sum_sequence,funs_elim,function_apply_intro}.
            Thus $w(i) = (a(i) \rmul a(i)) + (b(i) \rmul b(i))$ by \cref{real_add_eq_subst}.
            Thus $d(i) =
                ((a(i) \rmul a(i)) + (b(i) \rmul b(i))) + ((1 + 1) \rmul (a(i) \rmul b(i)))$ by \cref{real_add_eq_subst}.
            Thus $d(i) =
                (a(i) \rmul a(i)) + ((1 + 1) \rmul (a(i) \rmul b(i))) + (b(i) \rmul b(i))$ by \cref{real_mul_closed,real_add_assoc,real_add_comm}.
            Therefore $q(i) = d(i)$.
    \end{subproof}
    $q$ is a function and $d$ is a function and $\dom{q} = J$ and $\dom{d} = J$ by \cref{funs_elim}.
    Hence $q = d$ by \cref{funs_elim,subseteq,fun_subseteq,subseteq_antisymmetric}.
    Then $\sumFiniteSequence{d}{(s_1 + s_2) + (1+1) \rmul s}$
        by \cref{sum_finite_sequence_add,sum_finite_sequence_scalar}.
    Hence $\sumFiniteSequence{q}{(s_1 + s_2) + (1+1) \rmul s}$.
    Therefore $\sumFiniteSequence{\coordSquareSequenceMul{n}{c}}{s_1 + s_2 + (1+1) \rmul s}$
        by \cref{real_add_assoc}.
\end{proof}

% from §20 helper lemma 41a: finite sum value is real
\begin{lemma}\label{sum_finite_sequence_value_in_reals}
    Suppose $n\in\naturalsPlus$.
    Suppose $a\in\funs{\initialSegment{n}}{\reals}$.
    Suppose $\sumFiniteSequence{a}{s}$.
    Then $s\in\reals$.
\end{lemma}
\begin{proof}
    Thus $s\in\reals$ by \cref{sum_finite_sequence_pred,initial_segment_naturalsplus,funs_type_apply}.
\end{proof}

% from §20 helper lemma 41b: coordinate square sequence terms are nonnegative
\begin{lemma}\label{coord_square_sequence_mul_term_nonneg}
    Suppose $n\in\naturalsPlus$.
    Let $J = \initialSegment{n}$.
    Suppose $a\in\funs{J}{\reals}$.
    Then for all $i\in J$ we have $0 \leq \apply{\coordSquareSequenceMul{n}{a}}{i}$.
\end{lemma}
\begin{proof}
    For all $i\in J$ we have $0 \leq \apply{\coordSquareSequenceMul{n}{a}}{i}$
        by \cref{coord_square_sequence_mul,coord_square_sequence_mul_function,function_apply_intro,funs_elim,funs_type_apply,real_sq_nonnegative,real_one_in_naturals,real_pow_succ,real_pow_one}.
\end{proof}

% from §20 helper lemma 41c: finite sum of coordinate squares is nonnegative
\begin{lemma}\label{sum_finite_sequence_square_nonneg}
    Suppose $n\in\naturalsPlus$.
    Let $J = \initialSegment{n}$.
    Suppose $a\in\funs{J}{\reals}$.
    Suppose $\sumFiniteSequence{\coordSquareSequenceMul{n}{a}}{s}$.
    Then $0 \leq s$.
\end{lemma}
\begin{proof}
    Thus $0 \leq s$
        by \cref{coord_square_sequence_mul_function,funs_elim,coord_square_sequence_mul_term_nonneg,sum_finite_sequence_nonneg}.
\end{proof}

% from §20 helper lemma 41: Cauchy-Schwarz for finite sequences
\begin{lemma}\label{cauchy_schwarz_finite_sequences}
    Suppose $n\in\naturalsPlus$.
    Let $J = \initialSegment{n}$.
    Suppose $a\in\funs{J}{\reals}$.
    Suppose $b\in\funs{J}{\reals}$.
    Suppose $\sumFiniteSequence{\coordProductSequence{n}{a}{b}}{s}$.
    Suppose $\sumFiniteSequence{\coordSquareSequenceMul{n}{a}}{s_1}$.
    Suppose $\sumFiniteSequence{\coordSquareSequenceMul{n}{b}}{s_2}$.
    Then $s \rmul s \leq s_1 \rmul s_2$.
\end{lemma}
\begin{proof}
    Let $a_1 = \coordScalarSequence{n}{s_2}{a}$.
    Let $b_1 = \coordScalarSequence{n}{\neg{s}}{b}$.
    Let $c = \coordSumSequence{n}{a_1}{b_1}$.
    Let $p = \coordProductSequence{n}{a}{b}$.
    Let $p_1 = \coordProductSequence{n}{a_1}{b_1}$.
    $s_2\in\reals$ and $s\in\reals$ and $s_1\in\reals$ and $\neg{s}\in\reals$
        by \cref{coord_square_sequence_mul_function,coord_product_sequence_function,sum_finite_sequence_value_in_reals,real_add_inverse}.
    $a_1\in\funs{J}{\reals}$ and $b_1\in\funs{J}{\reals}$ by \cref{coord_scalar_sequence_function}.
    $p\in\funs{J}{\reals}$ and $p_1\in\funs{J}{\reals}$ by \cref{coord_product_sequence_function}.
    $\sumFiniteSequence{\coordSquareSequenceMul{n}{a_1}}{(s_2 \rmul s_2) \rmul s_1}$ by \cref{sum_finite_sequence_square_scalar}.
    Then $\sumFiniteSequence{\coordSquareSequenceMul{n}{b_1}}{(s \rmul s) \rmul s_2}$
        by \cref{sum_finite_sequence_square_scalar,real_mul_neg_right,real_mul_comm,real_mul_closed,real_neg_neg}.
    Let $r = \coordScalarSequence{n}{s_2 \rmul (\neg{s})}{p}$.
    $s_2 \rmul (\neg{s})\in\reals$ by \cref{real_mul_closed}.
    Hence $r\in\funs{J}{\reals}$ by \cref{real_mul_closed,coord_scalar_sequence_function}.
    Show that for all $i\in J$ we have $p_1(i) = r(i)$.
    \begin{subproof}
        Fix $i$ such that $i\in J$.
        Thus $a_1(i) = s_2 \rmul a(i)$ and $b_1(i) = \neg{s} \rmul b(i)$
            by \cref{coord_scalar_sequence,funs_elim,function_apply_intro}.
        Thus $p_1(i) = a_1(i) \rmul b_1(i)$
            by \cref{coord_product_sequence,funs_elim,function_apply_intro}.
        Thus $p_1(i) = (s_2 \rmul a(i)) \rmul (\neg{s} \rmul b(i))$.
        Thus $a(i)\in\reals$ and $b(i)\in\reals$ by \cref{funs_type_apply}.
        Thus $p_1(i) = (s_2 \rmul \neg{s}) \rmul (a(i) \rmul b(i))$
            by \cref{real_mul_assoc,real_mul_comm,real_mul_closed}.
        Thus $p(i) = a(i) \rmul b(i)$ by \cref{coord_product_sequence,funs_elim,function_apply_intro}.
        Thus $r(i) = (s_2 \rmul (\neg{s})) \rmul p(i)$
            by \cref{coord_scalar_sequence,funs_elim,function_apply_intro}.
        Thus $r(i) = (s_2 \rmul (\neg{s})) \rmul (a(i) \rmul b(i))$.
        Thus $p_1(i) = r(i)$.
    \end{subproof}
    $p_1$ is a function and $r$ is a function and $\dom{p_1} = J$ and $\dom{r} = J$ by \cref{funs_elim}.
    Therefore $p_1 = r$ by \cref{funs_elim,subseteq,fun_subseteq,subseteq_antisymmetric}.
    Then $\sumFiniteSequence{r}{(s_2 \rmul (\neg{s})) \rmul s}$ by \cref{sum_finite_sequence_scalar}.
    Hence $\sumFiniteSequence{p_1}{(s_2 \rmul (\neg{s})) \rmul s}$.
    $c\in\funs{J}{\reals}$ and $\coordSquareSequenceMul{n}{c}\in\funs{J}{\reals}$
        by \cref{coord_sum_sequence_function,coord_square_sequence_mul_function}.
    Take $S\in\reals$ such that $\sumFiniteSequence{\coordSquareSequenceMul{n}{c}}{S}$ by \cref{sum_finite_sequence_exists}.
    Then $\sumFiniteSequence{\coordSquareSequenceMul{n}{c}}{(s_2 \rmul s_2) \rmul s_1 + (s \rmul s) \rmul s_2 + (1+1) \rmul ((s_2 \rmul (\neg{s})) \rmul s)}$
        by \cref{sum_finite_sequence_square_sum}.
    Hence $S = (s_2 \rmul s_2) \rmul s_1 + (s \rmul s) \rmul s_2 + (1+1) \rmul ((s_2 \rmul (\neg{s})) \rmul s)$
        by \cref{sum_finite_sequence_unique}.
    Therefore $0 \leq S$ by \cref{sum_finite_sequence_square_nonneg}.
    $(s_2 \rmul (\neg{s})) \rmul s = s_2 \rmul ((\neg{s}) \rmul s)$ by \cref{real_mul_assoc}.
    $(\neg{s}) \rmul s = s \rmul (\neg{s})$ by \cref{real_mul_comm}.
    Therefore $(s_2 \rmul (\neg{s})) \rmul s = \neg{(s_2 \rmul (s \rmul s))}$
        by \cref{real_mul_assoc,real_mul_closed,real_mul_neg_right}.
    Thus $(1+1) \rmul ((s_2 \rmul (\neg{s})) \rmul s) =
        (1+1) \rmul \neg{(s_2 \rmul (s \rmul s))}$.
    $1\in\naturalsPlus$ and $1\in\reals$ by \cref{real_one_in_naturals,real_naturals_subset_reals,subseteq}.
    $1 + 1\in\reals$ by \cref{real_add_closed}.
    Thus $s_2 \rmul (s \rmul s)\in\reals$ by \cref{real_mul_closed}.
    Then $(1+1) \rmul \neg{(s_2 \rmul (s \rmul s))} =
        \neg{((1+1) \rmul (s_2 \rmul (s \rmul s)))}$ by \cref{real_mul_neg_right}.
    Thus $S = (s_2 \rmul s_2) \rmul s_1 + (s \rmul s) \rmul s_2 +
        \neg{((1+1) \rmul (s_2 \rmul (s \rmul s)))}$ by \cref{real_add_eq_subst}.
    $s \rmul s\in\reals$ by \cref{real_mul_closed}.
    $(s \rmul s) \rmul s_2 = s_2 \rmul (s \rmul s)$ by \cref{real_mul_comm}.
    Thus $S = s_2 \rmul (s_2 \rmul s_1) + s_2 \rmul (s \rmul s) +
        \neg{((1+1) \rmul (s_2 \rmul (s \rmul s)))}$ by \cref{real_mul_assoc}.
    Thus $s_2 \rmul (s_2 \rmul s_1)\in\reals$ and $s_2 \rmul (s \rmul s)\in\reals$ by \cref{real_mul_closed}.
    Thus $\neg{((1+1) \rmul (s_2 \rmul (s \rmul s)))}\in\reals$ by \cref{real_add_inverse,real_mul_closed}.
    Thus $S = s_2 \rmul (s_2 \rmul s_1) +
        (s_2 \rmul (s \rmul s) + \neg{((1+1) \rmul (s_2 \rmul (s \rmul s)))})$ by \cref{real_add_assoc}.
    Thus $(1+1) \rmul (s_2 \rmul (s \rmul s)) =
        (s_2 \rmul (s \rmul s)) + (s_2 \rmul (s \rmul s))$
        by \cref{real_mul_comm,real_distrib,real_mul_ident,real_add_eq_subst}.
    Hence $\neg{((1+1) \rmul (s_2 \rmul (s \rmul s)))} =
        \neg{(s_2 \rmul (s \rmul s))} + \neg{(s_2 \rmul (s \rmul s))}$
        by \cref{real_add_eq_subst,real_neg_addition}.
    Thus $s_2 \rmul (s \rmul s) + \neg{((1+1) \rmul (s_2 \rmul (s \rmul s)))} =
        \neg{(s_2 \rmul (s \rmul s))}$
        by \cref{real_add_eq_subst,real_add_assoc,real_add_inverse,real_add_ident}.
    Thus $S = s_2 \rmul (s_2 \rmul s_1) + \neg{(s_2 \rmul (s \rmul s))}$ by \cref{real_add_eq_subst}.
    Thus $s_2 \rmul s_1\in\reals$ and $s \rmul s\in\reals$ by \cref{real_mul_closed}.
    Thus $\neg{(s \rmul s)}\in\reals$ by \cref{real_add_inverse}.
    Thus $s_2 \rmul (s_2 \rmul s_1 + \neg{(s \rmul s)}) =
        s_2 \rmul (s_2 \rmul s_1) + s_2 \rmul \neg{(s \rmul s)}$ by \cref{real_distrib}.
    Thus $s_2 \rmul \neg{(s \rmul s)} = \neg{(s_2 \rmul (s \rmul s))}$ by \cref{real_mul_neg_right}.
    Thus $s_2 \rmul (s_2 \rmul s_1 + \neg{(s \rmul s)}) =
        s_2 \rmul (s_2 \rmul s_1) + \neg{(s_2 \rmul (s \rmul s))}$ by \cref{real_add_eq_subst}.
    Hence $S = s_2 \rmul (s_2 \rmul s_1 + \neg{(s \rmul s)})$.
    Therefore $S = s_2 \rmul (s_2 \rmul s_1 - s \rmul s)$ by \cref{real_sub}.
    Then $0 \leq s_2$ by \cref{sum_finite_sequence_square_nonneg}.
    $0 < s_2$ or $s_2 = 0$ by \cref{real_leq_split}.
    \begin{byCase}
    \caseOf{$s_2 = 0$.}
            For all $i\in J$ we have $b(i) = 0$
                by \cref{coord_square_sequence_mul_term_nonneg,coord_square_sequence_mul_function,sum_finite_sequence_zero_terms,coord_square_sequence_mul,function_apply_intro,funs_elim,funs_type_apply,real_mul_zero_product}.
            For all $i\in J$ we have $p(i) = 0$
                by \cref{coord_product_sequence,funs_elim,function_apply_intro,funs_type_apply,real_mul_comm,real_mul_zero}.
            $0\in\reals$ by \cref{real_add_ident}.
            Thus $\sumFiniteSequence{p}{s}$.
            Thus $s \leq 0$ by \cref{sum_finite_sequence_leq_n_times,real_naturals_subset_reals,subseteq,real_mul_comm,real_mul_zero}.
            For all $i\in\dom{p}$ we have $0 \leq p(i)$ by \cref{funs_elim,subseteq}.
            Thus $0 \leq s$ by \cref{sum_finite_sequence_nonneg}.
            Thus $s = 0$ by \cref{real_leq_antisym}.
            Thus $s \rmul s = s_1 \rmul s_2$ by \cref{real_mul_comm,real_mul_zero}.
            Therefore $s \rmul s \leq s_1 \rmul s_2$.
    \caseOf{$0 < s_2$.}
            Suppose not.
            Show that $s_2 \rmul s_1 - s \rmul s < 0$.
            \begin{subproof}
                Thus $s \rmul s < s_2 \rmul s_1$ or $s \rmul s = s_2 \rmul s_1$ or $s_2 \rmul s_1 < s \rmul s$
                    by \cref{real_lt_trichotomy}.
                \begin{byCase}
                \caseOf{$s \rmul s < s_2 \rmul s_1$.}
                    Then $s \rmul s \leq s_2 \rmul s_1$ by \cref{real_lt_imp_leq}.
                    Contradiction.
                \caseOf{$s \rmul s = s_2 \rmul s_1$.}
                    Then $s \rmul s \leq s_2 \rmul s_1$.
                    Contradiction.
                \caseOf{$s_2 \rmul s_1 < s \rmul s$.}
                    Thus $0 < s \rmul s - s_2 \rmul s_1$ by \cref{real_lt_imp_pos_diff}.
                    Thus $s \rmul s - s_2 \rmul s_1\in\reals$ by \cref{real_sub,real_add_closed,real_add_inverse,real_mul_closed}.
                    Thus $s_2 \rmul s_1 - s \rmul s = \neg{(s \rmul s - s_2 \rmul s_1)}$ by \cref{real_sub_swap_neg}.
                    Thus $s_2 \rmul s_1 - s \rmul s < 0$ by \cref{real_pos_imp_neg}.
                \end{byCase}
            \end{subproof}
            Thus $0\in\reals$ by \cref{real_add_ident}.
            Thus $s_2 \rmul s_1\in\reals$ by \cref{real_mul_closed}.
            Thus $s \rmul s\in\reals$ by \cref{real_mul_closed}.
            Thus $\neg{(s \rmul s)}\in\reals$ by \cref{real_add_inverse}.
            Thus $s_2 \rmul s_1 - s \rmul s\in\reals$ by \cref{real_sub,real_add_closed}.
            Thus $(s_2 \rmul s_1 - s \rmul s) \rmul s_2 < 0 \rmul s_2$ by \cref{real_lt_mul_pos}.
            Thus $0 \rmul s_2 = 0$ by \cref{real_mul_zero}.
            Thus $(s_2 \rmul s_1 - s \rmul s) \rmul s_2 < 0$ by \cref{real_lt_subst_right}.
            Thus $(s_2 \rmul s_1 - s \rmul s) \rmul s_2 = s_2 \rmul (s_2 \rmul s_1 - s \rmul s)$ by \cref{real_mul_comm}.
            Thus $s_2 \rmul (s_2 \rmul s_1 - s \rmul s) < 0$ by \cref{real_lt_subst_left}.
            Thus $S < 0$ by \cref{real_lt_subst_left}.
            Thus $0 < 0$ by \cref{real_lt_trans,real_lt_subst_right}.
            Contradiction by \cref{real_lt_irreflexive}.
            $0\in\reals$ and $s_2 \rmul s_1\in\reals$ and $s \rmul s\in\reals$ and
                $\neg{(s \rmul s)}\in\reals$ and $s_2 \rmul s_1 - s \rmul s \in\reals$
                by \cref{real_add_ident,real_mul_closed,real_add_inverse,real_sub,real_add_closed}.
            Thus $0 + (s \rmul s) \leq (s_2 \rmul s_1 - s \rmul s) + (s \rmul s)$ and
                $0 + (s \rmul s) = s \rmul s$ by \cref{real_leq_add,real_add_ident}.
            Thus $(s_2 \rmul s_1 - s \rmul s) + (s \rmul s) = s_2 \rmul s_1$
                by \cref{real_sub,real_add_assoc,real_add_comm,real_add_inverse,real_add_ident}.
            Thus $s \rmul s \leq s_2 \rmul s_1$ by \cref{real_leq_trans}.
            Therefore $s \rmul s \leq s_1 \rmul s_2$ by \cref{real_mul_comm}.
    \end{byCase}
\end{proof}

% from §20 helper lemma 42: Minkowski for finite sequences
\begin{lemma}\label{minkowski_finite_sequences}
    Suppose $n\in\naturalsPlus$.
    Let $J = \initialSegment{n}$.
    Suppose $a\in\funs{J}{\reals}$.
    Suppose $b\in\funs{J}{\reals}$.
    Suppose $\sumFiniteSequence{\coordSquareSequenceMul{n}{a}}{s_1}$.
    Suppose $\sumFiniteSequence{\coordSquareSequenceMul{n}{b}}{s_2}$.
    Let $c = \coordSumSequence{n}{a}{b}$.
    Suppose $\sumFiniteSequence{\coordSquareSequenceMul{n}{c}}{s}$.
    Then $\sqrt{s} \leq \sqrt{s_1} + \sqrt{s_2}$.
\end{lemma}
\begin{proof}
    Let $p = \coordProductSequence{n}{a}{b}$.
    Take $t\in\reals$ such that $\sumFiniteSequence{p}{t}$
        by \cref{coord_product_sequence_function,sum_finite_sequence_exists}.
    Then $\sumFiniteSequence{\coordSquareSequenceMul{n}{c}}{s_1 + s_2 + (1+1) \rmul t}$ by \cref{sum_finite_sequence_square_sum}.
    Hence $s = s_1 + s_2 + (1+1) \rmul t$ by \cref{sum_finite_sequence_unique}.
    $s_1\in\reals$ and $s_2\in\reals$ by \cref{coord_square_sequence_mul_function,sum_finite_sequence_value_in_reals}.
    $1\in\naturalsPlus$ and $1\in\reals$ and $1 + 1\in\reals$
        by \cref{real_one_in_naturals,real_naturals_subset_reals,subseteq,real_add_closed}.
    Therefore $(1+1) \rmul t\in\reals$ and $s_1 + s_2\in\reals$ and $s\in\reals$
        by \cref{real_one_in_naturals,real_naturals_subset_reals,subseteq,real_add_closed,real_mul_closed}.
    $0 \leq s_1$ and $0 \leq s_2$ by \cref{sum_finite_sequence_square_nonneg}.
    $0 \leq s$ by \cref{coord_sum_sequence_function,sum_finite_sequence_square_nonneg}.
    Let $r = \sqrt{s}$.
    Let $r_1 = \sqrt{s_1}$.
    Let $r_2 = \sqrt{s_2}$.
    $r\in\reals$ and $r \geq 0$ and $r \rmul r = s$ by \cref{positive_square_root}.
    $r_1\in\reals$ and $r_1 \geq 0$ and $r_1 \rmul r_1 = s_1$ by \cref{positive_square_root}.
    $r_2\in\reals$ and $r_2 \geq 0$ and $r_2 \rmul r_2 = s_2$ by \cref{positive_square_root}.
    Show that $0 \leq r_1 \rmul r_2$.
    \begin{subproof}
    $0 < r_1$ or $r_1 = 0$ by \cref{real_leq_split}.
    \begin{byCase}
        \caseOf{$r_1 = 0$.}
            Then $r_1 \rmul r_2 \geq 0$ by \cref{real_mul_comm,real_mul_zero}.
        \caseOf{$0 < r_1$.}
            $0 < r_2$ or $r_2 = 0$ by \cref{real_leq_split}.
            \begin{byCase}
            \caseOf{$r_2 = 0$.}
                Then $r_1 \rmul r_2 \geq 0$ by \cref{real_mul_comm,real_mul_zero}.
            \caseOf{$0 < r_2$.}
                Then $r_1 \rmul r_2 \geq 0$ by \cref{real_mul_pos_pos_gt,real_lt_imp_leq}.
            \end{byCase}
    \end{byCase}
    \end{subproof}
    $\abs{t}\in\reals$ and $\abs{t} \geq 0$ and $r_1 \rmul r_2\in\reals$
        by \cref{real_abs_in_reals,real_abs_nonnegative,real_mul_closed}.
    Then $\abs{t} \leq r_1 \rmul r_2$ iff
        $\exp{\abs{t}}{1+1} \leq \exp{r_1 \rmul r_2}{1+1}$ by \cref{real_nonneg_leq_iff_sq_leq}.
    Hence $\exp{\abs{t}}{1+1} = t \rmul t$ by \cref{real_abs_square,real_pow_succ,real_pow_one}.
    Therefore $t \rmul t \leq s_1 \rmul s_2$ by \cref{cauchy_schwarz_finite_sequences}.
    Thus $\exp{r_1 \rmul r_2}{1+1} = s_1 \rmul s_2$
        by \cref{real_pow_succ,real_pow_one,real_mul_assoc,real_mul_comm}.
    Hence $\exp{\abs{t}}{1+1} \leq \exp{r_1 \rmul r_2}{1+1}$ by \cref{real_leq_trans}.
    Therefore $\abs{t} \leq r_1 \rmul r_2$ by \cref{real_nonneg_leq_iff_sq_leq}.
    Hence $t \leq r_1 \rmul r_2$ by \cref{real_leq_abs,real_leq_trans}.
    Show that $(1+1) \rmul t \leq (1+1) \rmul (r_1 \rmul r_2)$.
    \begin{subproof}
    Then $t < r_1 \rmul r_2$ or $t = r_1 \rmul r_2$ by \cref{real_leq_split}.
    Hence $0 < 1 + 1$
        by \cref{real_add_ident,real_naturals_subset_reals,real_one_in_naturals,subseteq,real_zero_lt_one,real_lt_add,real_lt_subst_left,real_lt_trans}.
    \begin{byCase}
        \caseOf{$t = r_1 \rmul r_2$.}
            Then $(1+1) \rmul t \leq (1+1) \rmul (r_1 \rmul r_2)$.
        \caseOf{$t < r_1 \rmul r_2$.}
            $t\in\reals$ and $r_1 \rmul r_2\in\reals$ and $1 + 1\in\reals$
                by \cref{real_mul_closed,real_add_closed}.
            Thus $t \rmul (1+1) < (r_1 \rmul r_2) \rmul (1+1)$ by \cref{real_lt_mul_pos}.
            Thus $t \rmul (1+1) = (1+1) \rmul t$ by \cref{real_mul_comm}.
            Thus $(1+1) \rmul t < (r_1 \rmul r_2) \rmul (1+1)$ by \cref{real_lt_subst_left}.
            Thus $(r_1 \rmul r_2) \rmul (1+1) = (1+1) \rmul (r_1 \rmul r_2)$ by \cref{real_mul_comm}.
            Thus $(1+1) \rmul t < (1+1) \rmul (r_1 \rmul r_2)$ by \cref{real_lt_subst_right}.
            Therefore $(1+1) \rmul t \leq (1+1) \rmul (r_1 \rmul r_2)$.
    \end{byCase}
    \end{subproof}
    Thus $(1+1) \rmul t\in\reals$ and $(1+1) \rmul (r_1 \rmul r_2)\in\reals$ by \cref{real_mul_closed}.
    Thus $(1+1) \rmul t + s_1 \leq (1+1) \rmul (r_1 \rmul r_2) + s_1$ by \cref{real_leq_add}.
    Thus $(1+1) \rmul t + s_1 = s_1 + (1+1) \rmul t$ and
        $(1+1) \rmul (r_1 \rmul r_2) + s_1 = s_1 + (1+1) \rmul (r_1 \rmul r_2)$ by \cref{real_add_comm}.
    Thus $s_1 + (1+1) \rmul t \leq s_1 + (1+1) \rmul (r_1 \rmul r_2)$ by \cref{real_leq_trans}.
    Thus $s_1 + (1+1) \rmul t\in\reals$ and $s_1 + (1+1) \rmul (r_1 \rmul r_2)\in\reals$ by \cref{real_add_closed}.
    Thus $(s_1 + (1+1) \rmul t) + s_2 \leq (s_1 + (1+1) \rmul (r_1 \rmul r_2)) + s_2$ by \cref{real_leq_add}.
    Thus $(s_1 + (1+1) \rmul t) + s_2 = s_1 + ((1+1) \rmul t + s_2)$ by \cref{real_add_assoc}.
    Thus $(1+1) \rmul t + s_2 = s_2 + (1+1) \rmul t$ by \cref{real_add_comm}.
    Thus $s_1 + ((1+1) \rmul t + s_2) = s_1 + (s_2 + (1+1) \rmul t)$ by \cref{real_add_eq_subst}.
    Thus $s_1 + (s_2 + (1+1) \rmul t) = (s_1 + s_2) + (1+1) \rmul t$ by \cref{real_add_assoc}.
    Thus $(s_1 + (1+1) \rmul t) + s_2 = (s_1 + s_2) + (1+1) \rmul t$.
    Thus $(s_1 + (1+1) \rmul (r_1 \rmul r_2)) + s_2 = s_1 + ((1+1) \rmul (r_1 \rmul r_2) + s_2)$ and
        $(1+1) \rmul (r_1 \rmul r_2) + s_2 = s_2 + (1+1) \rmul (r_1 \rmul r_2)$ by \cref{real_add_assoc,real_add_comm}.
    Thus $s_1 + ((1+1) \rmul (r_1 \rmul r_2) + s_2) = s_1 + (s_2 + (1+1) \rmul (r_1 \rmul r_2))$ by \cref{real_add_eq_subst}.
    Thus $s_1 + (s_2 + (1+1) \rmul (r_1 \rmul r_2)) = (s_1 + s_2) + (1+1) \rmul (r_1 \rmul r_2)$ by \cref{real_add_assoc}.
    Thus $(s_1 + (1+1) \rmul (r_1 \rmul r_2)) + s_2 = (s_1 + s_2) + (1+1) \rmul (r_1 \rmul r_2)$.
    Thus $(s_1 + s_2) + (1+1) \rmul t \leq (s_1 + s_2) + (1+1) \rmul (r_1 \rmul r_2)$ by \cref{real_leq_trans}.
    Therefore $s \leq (s_1 + s_2) + (1+1) \rmul (r_1 \rmul r_2)$ by \cref{real_leq_trans}.
    Then $r_1 + r_2\in\reals$ by \cref{real_add_closed}.
    Thus $(r_1 + r_2) \rmul (r_1 + r_2) =
        (r_1 + r_2) \rmul r_1 + (r_1 + r_2) \rmul r_2$ by \cref{real_distrib}.
    Thus $(r_1 + r_2) \rmul r_1 = (r_1 \rmul r_1) + (r_1 \rmul r_2)$
        by \cref{real_mul_comm,real_distrib}.
    Thus $(r_1 + r_2) \rmul r_2 = (r_2 \rmul r_1) + (r_2 \rmul r_2)$
        by \cref{real_mul_comm,real_distrib}.
    Thus $(r_1 + r_2) \rmul (r_1 + r_2) =
        ((r_1 \rmul r_1) + (r_1 \rmul r_2)) + ((r_2 \rmul r_1) + (r_2 \rmul r_2))$ by \cref{real_add_eq_subst}.
    Thus $r_2 \rmul r_1 = r_1 \rmul r_2$ by \cref{real_mul_comm}.
    Thus $(r_1 + r_2) \rmul (r_1 + r_2) =
        ((r_1 \rmul r_1) + (r_1 \rmul r_2)) + ((r_1 \rmul r_2) + (r_2 \rmul r_2))$ by \cref{real_add_eq_subst}.
    Thus $(r_1 \rmul r_2) + (r_2 \rmul r_2)\in\reals$ by \cref{real_add_closed}.
    Thus $(r_1 + r_2) \rmul (r_1 + r_2) =
        (r_1 \rmul r_1) + ((r_1 \rmul r_2) + ((r_1 \rmul r_2) + (r_2 \rmul r_2)))$ by \cref{real_add_assoc}.
    Thus $(r_1 \rmul r_2) + ((r_1 \rmul r_2) + (r_2 \rmul r_2)) =
        ((r_1 \rmul r_2) + (r_1 \rmul r_2)) + (r_2 \rmul r_2)$ by \cref{real_add_assoc}.
    Thus $(r_1 + r_2) \rmul (r_1 + r_2) =
        (r_1 \rmul r_1) + (((r_1 \rmul r_2) + (r_1 \rmul r_2)) + (r_2 \rmul r_2))$ by \cref{real_add_eq_subst}.
    Thus $(r_1 \rmul r_2) + (r_1 \rmul r_2)\in\reals$ by \cref{real_add_closed}.
    Thus $(r_1 + r_2) \rmul (r_1 + r_2) =
        ((r_1 \rmul r_1) + ((r_1 \rmul r_2) + (r_1 \rmul r_2))) + (r_2 \rmul r_2)$ by \cref{real_add_assoc}.
    Thus $(1 + 1) \rmul (r_1 \rmul r_2) =
        (r_1 \rmul r_2) + (r_1 \rmul r_2)$
        by \cref{real_mul_comm,real_distrib,real_mul_ident,real_add_eq_subst}.
    Thus $(r_1 + r_2) \rmul (r_1 + r_2) =
        ((r_1 \rmul r_1) + ((1 + 1) \rmul (r_1 \rmul r_2))) + (r_2 \rmul r_2)$ by \cref{real_add_eq_subst}.
    Thus $(r_1 + r_2) \rmul (r_1 + r_2) =
        (r_1 \rmul r_1) + ((1 + 1) \rmul (r_1 \rmul r_2)) + (r_2 \rmul r_2)$ by \cref{real_add_assoc}.
    Thus $r_1 \rmul r_1 = s_1$ and $r_2 \rmul r_2 = s_2$.
    Thus $(r_1 + r_2) \rmul (r_1 + r_2) =
        s_1 + ((1 + 1) \rmul (r_1 \rmul r_2)) + s_2$ by \cref{real_add_eq_subst}.
    Thus $(r_1 + r_2) \rmul (r_1 + r_2) =
        s_1 + s_2 + (1 + 1) \rmul (r_1 \rmul r_2)$ by \cref{real_add_assoc,real_add_comm}.
    Thus $r_1 + r_2\in\reals$ and $0\in\reals$ and $0 \leq r_1$
        by \cref{real_add_closed,real_add_ident}.
    Thus $0 + r_2 \leq r_1 + r_2$ and $0 + r_2 = r_2$ by \cref{real_leq_add,real_add_ident}.
    Thus $r_2 \leq r_1 + r_2$ and $0 \leq r_2$ and $0 \leq r_1 + r_2$ by \cref{real_leq_trans}.
    Thus $\exp{r}{1+1} = s$ by \cref{real_pow_succ,real_pow_one}.
    Thus $\exp{r_1 + r_2}{1+1} = s_1 + s_2 + (1 + 1) \rmul (r_1 \rmul r_2)$
        by \cref{real_pow_succ,real_pow_one,real_add_eq_subst}.
    Thus $\exp{r}{1+1} \leq \exp{r_1 + r_2}{1+1}$ by \cref{real_leq_trans}.
    Hence $r \leq r_1 + r_2$ by \cref{real_nonneg_leq_iff_sq_leq}.
    Therefore $\sqrt{s} \leq \sqrt{s_1} + \sqrt{s_2}$.
\end{proof}

% from §20 helper lemma 43: euclidean metric value exists
\begin{lemma}\label{euclidean_metric_value_exists}
    Suppose $n\in\naturalsPlus$.
    Let $J = \initialSegment{n}$.
    Suppose $x\in\realsPower{J}$.
    Suppose $y\in\realsPower{J}$.
    Then there exists $r\in\reals$ such that $((x,y),r)\in\euclideanMetric{n}$.
\end{lemma}
\begin{proof}
    Let $a = \coordSquareSequence{n}{(x,y)}$.
    $a\in\funs{J}{\reals}$ by \cref{coord_square_sequence_function}.
    Take $s\in\reals$ such that $\sumFiniteSequence{a}{s}$ by \cref{sum_finite_sequence_exists}.
    For all $i\in\dom{a}$ we have $0 \leq a(i)$
        by \cref{funs_elim,reals_power,tuple_power,funs_type_apply,real_add_inverse,real_add_closed,real_sub,real_sq_nonnegative,coord_square_sequence_elem,function_apply_intro}.
    Thus $0 \leq s$ by \cref{sum_finite_sequence_nonneg}.
    Let $r = \sqrt{s}$.
    Therefore there exists $r\in\reals$ such that $((x,y),r)\in\euclideanMetric{n}$
        by \cref{positive_square_root,times_tuple_intro,euclidean_metric}.
\end{proof}
% from §20 helper lemma 44: euclidean metric is a function
\begin{lemma}\label{euclidean_metric_function}
    Suppose $n\in\naturalsPlus$.
    Then $\euclideanMetric{n}$ is a function.
\end{lemma}
\begin{proof}
    Follows by \cref{euclidean_metric,relation,rightunique,pair_eq_iff,sum_finite_sequence_unique}.
\end{proof}

% from §20 helper lemma 44a: chained subtraction
\begin{lemma}\label{real_sub_add_chain}
    Suppose $x\in\reals$.
    Suppose $y\in\reals$.
    Suppose $z\in\reals$.
    Then $(x - y) + (y - z) = x - z$.
\end{lemma}
\begin{proof}
    $\neg{y}\in\reals$ and $\neg{z}\in\reals$ by \cref{real_add_inverse}.
    Thus $(x - y) + (y - z) = x + ((\neg{y} + y) + \neg{z})$
        by \cref{real_sub,real_add_assoc,real_add_closed}.
    Thus $(x - y) + (y - z) = x + (0 + \neg{z})$
        by \cref{real_add_comm,real_add_inverse,real_add_eq_subst}.
    Therefore $(x - y) + (y - z) = x - z$ by \cref{real_add_ident,real_sub}.
\end{proof}

% from §20 helper lemma 45: square metric is a metric
\begin{lemma}\label{square_metric_is_metric}
    Suppose $n\in\naturalsPlus$.
    Let $J = \initialSegment{n}$.
    Let $X = \realsPower{J}$.
    Then $\squareMetric{n}$ is a metric on $X$.
\end{lemma}
\begin{proof}
    Let $d = \squareMetric{n}$.
    $d$ is a function by \cref{square_metric_function}.
    Hence $X\times X\subseteq\dom{d}$
        by \cref{times_elem_is_tuple,square_metric_value_exists,pair_eq_iff,dom_intro,subseteq}.
    Therefore $\dom{d}\subseteq X\times X$ by \cref{dom_iff,square_metric,pair_eq_iff,times_tuple_intro,subseteq}.
    Hence $\dom{d} = X\times X$ by \cref{subseteq_antisymmetric}.
    For all $p\in\dom{d}$ we have $d(p)\in\reals$
        by \cref{dom_iff,square_metric,pair_eq_iff,coord_diff_set_finite,subseteq,function_apply_intro}.
    Thus $d\in\funs{X\times X}{\reals}$ by \cref{funs_intro}.
    For all $x,y$ such that $x\in X$ and $y\in X$ we have $d((x,y)) \geq 0$
        by \cref{square_metric_value_exists,square_metric,pair_eq_iff,coord_diff_set,reals_power,tuple_power,funs_type_apply,real_sub,real_add_closed,real_add_inverse,real_abs_nonnegative,function_apply_intro}.
    Hence $d$ is nonnegative on $X$ by \cref{metric_nonnegative}.
    For all $x,y$ such that $x\in X$ and $y\in X$ we have if $d((x,y)) = 0$ then $x = y$
        by \cref{square_metric_value_exists,function_apply_intro,reals_power,tuple_power,funs_type_apply,real_sub,real_add_closed,real_add_inverse,real_abs_in_reals,real_abs_nonnegative,real_leq_antisym,real_abs_eq_zero,real_add_cancel,funs_elim,subseteq,fun_subseteq,subseteq_antisymmetric}.
    For all $x,y$ such that $x\in X$ and $y\in X$ we have if $x = y$ then $d((x,y)) = 0$
        by \cref{square_metric_value_exists,coord_diff_set,reals_power,tuple_power,funs_type_apply,real_sub,real_add_inverse,real_add_eq_subst,real_add_ident,real_abs_nonneg,square_metric,pair_eq_iff,function_apply_intro}.
    Therefore $d$ is zero-characterized on $X$ by \cref{metric_zero_characterization}.
    For all $x,y$ such that $x\in X$ and $y\in X$ we have $d((x,y)) = d((y,x))$
        by \cref{square_metric_value_exists,coord_diff_set,reals_power,tuple_power,funs_type_apply,real_abs_sub_swap,square_metric,pair_eq_iff,real_leq_antisym,function_apply_intro}.
    Hence $d$ is symmetric on $X$ by \cref{metric_symmetric}.
    Show that for all $x,y,z$ such that $x\in X$ and $y\in X$ and $z\in X$ we have
        $d((x,y)) + d((y,z)) \geq d((x,z))$.
    \begin{subproof}
            Fix $x$.
            Fix $y$.
            Fix $z$ such that $x\in X$ and $y\in X$ and $z\in X$.
            Take $r_1\in\reals$ such that $((x,y),r_1)\in d$ and
                for all $i\in J$ we have $\abs{x(i) - y(i)} \leq r_1$
                by \cref{square_metric_value_exists}.
            Take $r_2\in\reals$ such that $((y,z),r_2)\in d$ and
                for all $i\in J$ we have $\abs{y(i) - z(i)} \leq r_2$
                by \cref{square_metric_value_exists}.
            Take $r_3\in\reals$ such that $((x,z),r_3)\in d$ and
                for all $i\in J$ we have $\abs{x(i) - z(i)} \leq r_3$
                by \cref{square_metric_value_exists}.
            Hence $r_3\in\coordDiffSet{n}{x}{z}$ by \cref{square_metric,pair_eq_iff}.
            Take $i_0\in J$ such that $r_3 = \abs{x(i_0) - z(i_0)}$ by \cref{coord_diff_set}.
            $x(i_0)\in\reals$ and $y(i_0)\in\reals$ and $z(i_0)\in\reals$ and
                $\neg{y(i_0)}\in\reals$ and $\neg{z(i_0)}\in\reals$ and
                $x(i_0) - y(i_0)\in\reals$ and $y(i_0) - z(i_0)\in\reals$
                by \cref{reals_power,tuple_power,funs_type_apply,real_add_inverse,real_sub,real_add_closed}.
            Therefore $(x(i_0) - y(i_0)) + (y(i_0) - z(i_0)) = x(i_0) - z(i_0)$
                by \cref{real_sub_add_chain}.
            Thus $\abs{(x(i_0) - y(i_0)) + (y(i_0) - z(i_0))} \leq \abs{x(i_0) - y(i_0)} + \abs{y(i_0) - z(i_0)}$ by \cref{real_triangle_inequality}.
            Hence $\abs{x(i_0) - z(i_0)} \leq \abs{x(i_0) - y(i_0)} + \abs{y(i_0) - z(i_0)}$.
            Therefore $r_3 \leq \abs{x(i_0) - y(i_0)} + \abs{y(i_0) - z(i_0)}$.
            $\abs{x(i_0) - y(i_0)} \leq r_1$.
            $\abs{x(i_0) - y(i_0)}\in\reals$ and $\abs{y(i_0) - z(i_0)}\in\reals$ by \cref{real_abs_in_reals}.
            Thus $\abs{x(i_0) - y(i_0)} + \abs{y(i_0) - z(i_0)} \leq r_1 + \abs{y(i_0) - z(i_0)}$ and
                $\abs{x(i_0) - y(i_0)} + \abs{y(i_0) - z(i_0)}\in\reals$ and
                $r_1 + \abs{y(i_0) - z(i_0)}\in\reals$ by \cref{real_leq_add,real_add_closed}.
            Hence $r_3 \leq r_1 + \abs{y(i_0) - z(i_0)}$ by \cref{real_leq_trans}.
            $\abs{y(i_0) - z(i_0)} \leq r_2$.
            Therefore $r_1 + \abs{y(i_0) - z(i_0)} \leq r_1 + r_2$
                by \cref{real_leq_add,real_add_comm}.
            $r_1 + r_2\in\reals$ by \cref{real_add_closed}.
            Hence $r_3 \leq r_1 + r_2$ by \cref{real_leq_trans}.
            Therefore $d((x,z)) \leq d((x,y)) + d((y,z))$ by \cref{function_apply_intro}.
            Hence $d((x,y)) + d((y,z)) \geq d((x,z))$.
    \end{subproof}
    Therefore $d$ satisfies the triangle inequality on $X$ by \cref{metric_triangle_inequality}.
    Hence $d$ is a metric on $X$ by \cref{metric_on}.
\end{proof}

% from §20 helper lemma 46: square sequence of a difference sequence
\begin{lemma}\label{coord_square_sequence_diff}
    Suppose $n\in\naturalsPlus$.
    Let $J = \initialSegment{n}$.
    Suppose $x\in\realsPower{J}$.
    Suppose $y\in\realsPower{J}$.
    Let $a = \coordDiffSequence{n}{x}{y}$.
    Then $\coordSquareSequenceMul{n}{a} = \coordSquareSequence{n}{(x,y)}$.
\end{lemma}
\begin{proof}
    Let $u = \coordSquareSequenceMul{n}{a}$.
    Let $v = \coordSquareSequence{n}{(x,y)}$.
    For all $w$ we have $w\in u$ iff $w\in v$
        by \cref{coord_square_sequence_mul,coord_diff_sequence_function,funs_elim,coord_diff_sequence,function_apply_intro,reals_power,tuple_power,funs_type_apply,real_add_inverse,real_add_closed,real_sub,real_one_in_naturals,real_pow_succ,real_pow_one,coord_square_sequence,fst_eq,snd_eq}.
    Hence $u = v$ by \cref{setext}.
\end{proof}

% from §20 helper lemma 47: difference sequence of a sum
\begin{lemma}\label{coord_diff_sequence_sum}
    Suppose $n\in\naturalsPlus$.
    Let $J = \initialSegment{n}$.
    Suppose $x\in\realsPower{J}$.
    Suppose $y\in\realsPower{J}$.
    Suppose $z\in\realsPower{J}$.
    Let $a = \coordDiffSequence{n}{x}{y}$.
    Let $b = \coordDiffSequence{n}{y}{z}$.
    Let $c = \coordSumSequence{n}{a}{b}$.
    Then $c = \coordDiffSequence{n}{x}{z}$.
\end{lemma}
\begin{proof}
    Let $d = \coordDiffSequence{n}{x}{z}$.
    For all $w$ we have $w\in c$ iff $w\in d$
        by \cref{coord_sum_sequence,coord_diff_sequence,coord_diff_sequence_function,funs_elim,function_apply_intro,reals_power,tuple_power,funs_type_apply,real_sub_add_chain}.
    Hence $c = d$ by \cref{setext}.
\end{proof}

% from §20 Item 25: euclidean metric is a metric
\begin{theorem}\label{euclidean_metric_is_metric}
    Suppose $n\in\naturalsPlus$.
    Let $J = \initialSegment{n}$.
    Let $X = \realsPower{J}$.
    Then $\euclideanMetric{n}$ is a metric on $X$.
\end{theorem}
\begin{proof}
    Let $d = \euclideanMetric{n}$.
    $d$ is a function by \cref{euclidean_metric_function}.
    Hence $X\times X\subseteq\dom{d}$
        by \cref{times_elem_is_tuple,euclidean_metric_value_exists,pair_eq_iff,dom_intro,subseteq}.
    Therefore $\dom{d}\subseteq X\times X$
        by \cref{dom_iff,euclidean_metric,pair_eq_iff,times_elem_is_tuple,times_tuple_intro,subseteq}.
    Hence $\dom{d} = X\times X$ by \cref{subseteq_antisymmetric}.
    For all $p\in\dom{d}$ we have $d(p)\in\reals$
        by \cref{dom_iff,euclidean_metric,pair_eq_iff,function_apply_intro}.
    Thus $d\in\funs{X\times X}{\reals}$ by \cref{funs_intro}.
    Show that for all $x,y$ such that $x\in X$ and $y\in X$ we have $d((x,y)) \geq 0$.
    \begin{subproof}
        Fix $x$.
        Fix $y$ such that $x\in X$ and $y\in X$.
        Let $a = \coordDiffSequence{n}{x}{y}$.
        Take $r\in\reals$ such that $((x,y),r)\in d$ by \cref{euclidean_metric_value_exists}.
        Take $s\in\reals$ such that $\sumFiniteSequence{\coordSquareSequence{n}{(x,y)}}{s}$ and $r = \sqrt{s}$
            by \cref{euclidean_metric,pair_eq_iff}.
        $\coordSquareSequenceMul{n}{a} = \coordSquareSequence{n}{(x,y)}$ by \cref{coord_square_sequence_diff}.
        Then $\sumFiniteSequence{\coordSquareSequenceMul{n}{a}}{s}$.
        Thus $0 \leq s$ by \cref{coord_diff_sequence_function,sum_finite_sequence_square_nonneg}.
        Hence $r \geq 0$ by \cref{positive_square_root}.
        Therefore $d((x,y)) \geq 0$ by \cref{function_apply_intro}.
    \end{subproof}
    Hence $d$ is nonnegative on $X$ by \cref{metric_nonnegative}.
    For all $x,y$ such that $x\in X$ and $y\in X$ we have $d((x,y)) = 0$ iff $x = y$
        by \cref{euclidean_metric_value_exists,function_apply_intro,square_metric_value_exists,euclidean_square_metric_bounds,square_metric_is_metric,metric_on,metric_nonnegative,funs_elim,real_leq_antisym,metric_zero_characterization,real_add_ident,real_zero_lt_one,real_lt_imp_leq,real_one_leq_naturalsplus,real_naturals_subset_reals,real_one_in_naturals,subseteq,real_leq_trans,positive_square_root,real_mul_comm,real_mul_zero}.
    Therefore $d$ is zero-characterized on $X$ by \cref{metric_zero_characterization}.
    Show that for all $x,y$ such that $x\in X$ and $y\in X$ we have $d((x,y)) = d((y,x))$.
    \begin{subproof}
            Fix $x$.
            Fix $y$ such that $x\in X$ and $y\in X$.
            Take $r\in\reals$ such that $((x,y),r)\in d$ by \cref{euclidean_metric_value_exists}.
            Take $s\in\reals$ such that $((y,x),s)\in d$ by \cref{euclidean_metric_value_exists}.
            Take $s_1\in\reals$ such that
                $\sumFiniteSequence{\coordSquareSequence{n}{(x,y)}}{s_1}$ and $r = \sqrt{s_1}$
                by \cref{euclidean_metric,pair_eq_iff}.
            Take $s_2\in\reals$ such that
                $\sumFiniteSequence{\coordSquareSequence{n}{(y,x)}}{s_2}$ and $s = \sqrt{s_2}$
                by \cref{euclidean_metric,pair_eq_iff}.
            Let $a = \coordSquareSequence{n}{(x,y)}$.
            Let $b = \coordSquareSequence{n}{(y,x)}$.
            For all $w$ we have $w\in a$ iff $w\in b$
                by \cref{coord_square_sequence,fst_eq,snd_eq,reals_power,tuple_power,funs_type_apply,real_add_inverse,real_add_closed,real_sub,real_abs_sub_swap,real_abs_square}.
            Hence $a = b$ by \cref{setext}.
            Therefore $s_1 = s_2$ by \cref{sum_finite_sequence_unique}.
            Then $r = s$.
            Hence $d((x,y)) = d((y,x))$ by \cref{function_apply_intro}.
    \end{subproof}
    Therefore $d$ is symmetric on $X$ by \cref{metric_symmetric}.
    Show that for all $x,y,z$ such that $x\in X$ and $y\in X$ and $z\in X$ we have
        $d((x,y)) + d((y,z)) \geq d((x,z))$.
    \begin{subproof}
            Fix $x$.
            Fix $y$.
            Fix $z$ such that $x\in X$ and $y\in X$ and $z\in X$.
            Take $r_1\in\reals$ such that $((x,y),r_1)\in d$ by \cref{euclidean_metric_value_exists}.
            Take $r_2\in\reals$ such that $((y,z),r_2)\in d$ by \cref{euclidean_metric_value_exists}.
            Take $r_3\in\reals$ such that $((x,z),r_3)\in d$ by \cref{euclidean_metric_value_exists}.
            Take $s_1\in\reals$ such that
                $\sumFiniteSequence{\coordSquareSequence{n}{(x,y)}}{s_1}$ and $r_1 = \sqrt{s_1}$
                by \cref{euclidean_metric,pair_eq_iff}.
            Take $s_2\in\reals$ such that
                $\sumFiniteSequence{\coordSquareSequence{n}{(y,z)}}{s_2}$ and $r_2 = \sqrt{s_2}$
                by \cref{euclidean_metric,pair_eq_iff}.
            Take $s_3\in\reals$ such that
                $\sumFiniteSequence{\coordSquareSequence{n}{(x,z)}}{s_3}$ and $r_3 = \sqrt{s_3}$
                by \cref{euclidean_metric,pair_eq_iff}.
            Let $a = \coordDiffSequence{n}{x}{y}$.
            Let $b = \coordDiffSequence{n}{y}{z}$.
            Let $c = \coordSumSequence{n}{a}{b}$.
            Hence $\sqrt{s_3} \leq \sqrt{s_1} + \sqrt{s_2}$
                by \cref{coord_diff_sequence_function,coord_sum_sequence_function,coord_square_sequence_diff,coord_diff_sequence_sum,minkowski_finite_sequences}.
            Therefore $r_3 \leq r_1 + r_2$.
            Hence $d((x,z)) \leq d((x,y)) + d((y,z))$ by \cref{function_apply_intro}.
            Therefore $d((x,y)) + d((y,z)) \geq d((x,z))$.
    \end{subproof}
    Hence $d$ satisfies the triangle inequality on $X$ by \cref{metric_triangle_inequality}.
    Therefore $d$ is a metric on $X$ by \cref{metric_on}.
\end{proof}
% from §20 Item 26: euclidean metric topology equals product topology
\begin{theorem}\label{euclidean_metric_product_topology}
    Suppose $n\in\naturalsPlus$.
    Let $J = \initialSegment{n}$.
    Let $X = \{(i,\reals) \mid i\in J\}$.
    Let $T = \{(i,\standardTopologyR) \mid i\in J\}$.
    Let $T_0 = \productTopologyIndexed{T}{X}$.
    Then $\metricTopology{\euclideanMetric{n}}{\realsPower{J}} = T_0$.
\end{theorem}
\begin{proof}
    Let $X_0 = \realsPower{J}$.
    Let $d = \euclideanMetric{n}$.
    Let $d_1 = \squareMetric{n}$.
    Let $T_E = \metricTopology{d}{X_0}$.
    Let $T_S = \metricTopology{d_1}{X_0}$.
    Hence $J$ is inhabited by \cref{real_one_in_naturals,real_one_leq_naturalsplus,initial_segment_naturalsplus}.
    For all $w$ such that $w\in X$ there exist $x,y$ such that $w = (x,y)$.
    For all $i,y_1,y_2$ such that $(i,y_1)\in X$ and $(i,y_2)\in X$ we have $y_1 = y_2$.
    Therefore $X$ is a function by \cref{relation,rightunique}.
    For all $w$ such that $w\in T$ there exist $x,y$ such that $w = (x,y)$.
    For all $i,y_1,y_2$ such that $(i,y_1)\in T$ and $(i,y_2)\in T$ we have $y_1 = y_2$.
    Therefore $T$ is a function by \cref{relation,rightunique}.
    Hence $J\subseteq\dom{X}$ by \cref{dom_intro,subseteq}.
    For all $i$ such that $i\in\dom{X}$ we have $i\in J$.
    Therefore $\dom{X}\subseteq J$ by \cref{subseteq}.
    Hence $\dom{X} = J$ by \cref{subseteq_antisymmetric}.
    Hence $J\subseteq\dom{T}$ by \cref{dom_intro,subseteq}.
    For all $i$ such that $i\in\dom{T}$ we have $i\in J$.
    Therefore $\dom{T}\subseteq J$ by \cref{subseteq}.
    Hence $\dom{T} = J$ by \cref{subseteq_antisymmetric}.
    For all $\alpha\in\dom{X}$ we have $\apply{X}{\alpha} = \reals$ and $\apply{T}{\alpha} = \standardTopologyR$
        by \cref{subseteq,function_apply_intro}.
    Thus for all $\alpha\in\dom{X}$ we have $\apply{T}{\alpha}$ is a topology on $\apply{X}{\alpha}$
        by \cref{subseteq,function_apply_intro,standard_basis_is_basis,basis_topology_is_topology,standard_topology_reals}.
    Let $S_0 = \productSubbasis{T}{X}$.
    Let $B_0 = \finiteIntersections{S_0}$.
    $B_0$ is a basis for $T_0$ on $\productFamily{X}$ by \cref{product_subbasis_finite_intersections_basis}.
    Hence $B_0$ is a basis for $T_0$ on $X_0$ by \cref{product_family_reals_power}.
    $d$ is a metric on $X_0$ and $d_1$ is a metric on $X_0$ by \cref{euclidean_metric_is_metric,square_metric_is_metric}.
    Let $B_S = \metricBasis{d_1}{X_0}$.
    Therefore $B_S$ is a basis for $T_S$ on $X_0$
        by \cref{metric_basis_is_basis,basis_for_topology,metric_topology}.
    Show that for all $x\in X_0$ we have
        for all $\epsilon\in\reals$ such that $0 < \epsilon$ there exists $\delta\in\reals$
            such that $0 < \delta$ and
            $\metricBall{d}{X_0}{x}{\delta} \subseteq \metricBall{d_1}{X_0}{x}{\epsilon}$.
    \begin{subproof}
                Fix $x$ such that $x\in X_0$.
                Fix $\epsilon$ such that $\epsilon\in\reals$ and $0 < \epsilon$.
                Let $\delta = \epsilon$.
                $0 < \delta$.
                For all $y$ such that $y\in\metricBall{d}{X_0}{x}{\delta}$ we have $y\in\metricBall{d_1}{X_0}{x}{\epsilon}$
                    by \cref{metric_ball,euclidean_metric_value_exists,square_metric_value_exists,euclidean_square_metric_bounds,euclidean_metric_function,square_metric_function,real_leq_split,real_lt_subst_left,real_lt_trans,function_apply_intro}.
                Hence $\metricBall{d}{X_0}{x}{\delta} \subseteq \metricBall{d_1}{X_0}{x}{\epsilon}$ by \cref{subseteq}.
    \end{subproof}
    Hence $T_E$ is finer than $T_S$ on $X_0$ by \cref{metric_topology_finer_iff}.
    Show that for all $x\in X_0$ we have
        for all $\epsilon\in\reals$ such that $0 < \epsilon$ there exists $\delta\in\reals$
            such that $0 < \delta$ and
            $\metricBall{d_1}{X_0}{x}{\delta} \subseteq \metricBall{d}{X_0}{x}{\epsilon}$.
    \begin{subproof}
                Fix $x$ such that $x\in X_0$.
                Fix $\epsilon$ such that $\epsilon\in\reals$ and $0 < \epsilon$.
                $0 < 1$ and $0\in\reals$ and $1\in\reals$ and $n\in\reals$
                    by \cref{real_zero_lt_one,real_add_ident,real_naturals_subset_reals,real_one_in_naturals,subseteq}.
                Hence $0 < n$ by \cref{real_one_leq_naturalsplus,real_lt_subst_right,real_lt_trans}.
                Therefore $\sqrt{n}\in\reals$ and $\sqrt{n} \geq 0$ and $\sqrt{n} \rmul \sqrt{n} = n$ by \cref{positive_square_root}.
                Hence $\sqrt{n}$ is positive.
                Let $\delta = \epsilon \rmul \inv{\sqrt{n}}$.
                $\sqrt{n} \neq 0$ by \cref{real_pos_nonzero}.
                $\inv{\sqrt{n}}\in\reals$ and $\inv{\sqrt{n}} > 0$ by \cref{real_mul_inverse,real_inv_pos}.
                $\epsilon \rmul \inv{\sqrt{n}}\in\reals$ by \cref{real_mul_closed}.
                Thus $\delta > 0$ by \cref{real_mul_pos_pos_gt}.
                Show that for all $y$ such that $y\in\metricBall{d_1}{X_0}{x}{\delta}$ we have $y\in\metricBall{d}{X_0}{x}{\epsilon}$.
                \begin{subproof}
                    Fix $y$ such that $y\in\metricBall{d_1}{X_0}{x}{\delta}$.
                    Thus $y\in X_0$ and $d_1((x,y)) < \delta$ by \cref{metric_ball}.
                    Take $m\in\reals$ such that $((x,y),m)\in d_1$ by \cref{square_metric_value_exists}.
                    Take $r\in\reals$ such that $((x,y),r)\in d$ by \cref{euclidean_metric_value_exists}.
                    Hence $m < \delta$ by \cref{square_metric_function,function_apply_intro,real_lt_subst_left}.
                    Therefore $r \leq \sqrt{n} \rmul m$ by \cref{euclidean_square_metric_bounds}.
                    Hence $\sqrt{n} \rmul m < \sqrt{n} \rmul \delta$ by \cref{real_lt_mul_pos,real_mul_comm}.
                    Therefore $\sqrt{n} \rmul m < \epsilon$ by
                        \cref{real_lt_subst_right,real_mul_assoc,real_mul_comm,real_mul_inverse,real_mul_ident}.
                    Hence $d((x,y)) < \epsilon$ by
                        \cref{euclidean_metric_function,real_mul_closed,real_leq_split,real_lt_subst_left,real_lt_trans,function_apply_intro}.
                    Therefore $y\in\metricBall{d}{X_0}{x}{\epsilon}$ by \cref{metric_ball}.
                \end{subproof}
                Therefore $\metricBall{d_1}{X_0}{x}{\delta} \subseteq \metricBall{d}{X_0}{x}{\epsilon}$ by \cref{subseteq}.
    \end{subproof}
    Therefore $T_S$ is finer than $T_E$ on $X_0$ by \cref{metric_topology_finer_iff}.
    Hence $T_S\subseteq T_E$ and $T_E\subseteq T_S$ by \cref{finer_topology}.
    Therefore $T_E = T_S$ by \cref{subseteq_antisymmetric}.
    Show that for all $x,B_2$ such that $x\in X_0$ and $B_2\in B_S$ and $x\in B_2$
        there exists $B_3\in B_0$ such that $x\in B_3$ and $B_3\subseteq B_2$.
    \begin{subproof}
                Fix $x$.
                Fix $B_2$ such that $x\in X_0$ and $B_2\in B_S$ and $x\in B_2$.
                Take $x_0,\epsilon$ such that $x_0\in X_0$ and $\epsilon\in\reals$ and $0 < \epsilon$ and
                    $B_2 = \metricBall{d_1}{X_0}{x_0}{\epsilon}$ by \cref{metric_basis}.
                Let $R_g = \{w\in J\times\pow{\productFamily{X}} \mid \exists i\in J.
                    w = (i,\preimg{\piIndex{X}{i}}{\intervalopen{x_0(i) - \epsilon}{x_0(i) + \epsilon}})\}$.
                For all $w$ such that $w\in R_g$ there exist $i,A$ such that $w = (i,A)$.
                Hence $R_g$ is a relation by \cref{relation}.
                For all $i\in J$ we have
                    $\preimg{\piIndex{X}{i}}{\intervalopen{x_0(i) - \epsilon}{x_0(i) + \epsilon}}\in\pow{\productFamily{X}}$
                    by \cref{preimg_iff,projection_map_indexed,pair_eq_iff,subseteq,powerset_intro}.
                Thus for all $i\in J$ there exists $A$ such that $i\mathrel{R_g} A$
                    by \cref{times_tuple_intro}.
                Take $g$ such that $g$ is a function on $J$ and for all $i\in J$ we have $i\mathrel{R_g}\apply{g}{i}$
                    by \cref{choice_function}.
                Let $F = \img{g}{J}$.
                Hence $F$ is finite by \cref{initial_segment_finite,finite_image_function}.
                Take $i$ such that $i\in J$.
                Therefore $\apply{g}{i}\in F$ by \cref{function_apply_elim,img_iff}.
                Hence $F\neq\emptyset$ by \cref{function_apply_elim,img_iff,emptyset}.
                Show that for all $A$ such that $A\in F$ we have $A\in S_0$.
                \begin{subproof}
                    Fix $A$ such that $A\in F$.
                    Take $i\in J$ such that $(i,A)\in g$ by \cref{img_iff}.
                    Then $i\mathrel{R_g}\apply{g}{i}$.
                    Take $i_0\in J$ such that
                        $(i,\apply{g}{i}) = (i_0,\preimg{\piIndex{X}{i_0}}{\intervalopen{x_0(i_0) - \epsilon}{x_0(i_0) + \epsilon}})$.
                    Therefore $A = \preimg{\piIndex{X}{i}}{\intervalopen{x_0(i) - \epsilon}{x_0(i) + \epsilon}}$
                        by \cref{function_apply_intro,pair_eq_iff}.
                    Thus $x_0(i)\in\reals$ by \cref{reals_power,tuple_power,funs_type_apply}.
                    Thus $x_0(i) - \epsilon\in\reals$ and $x_0(i) + \epsilon\in\reals$
                        by \cref{real_add_inverse,real_sub,real_add_closed}.
                    $0 < \epsilon$ and $0\in\reals$ by \cref{real_zero_lt_one,real_add_ident}.
                    Thus $x_0(i) - \epsilon < x_0(i)$ by \cref{real_sub_pos_lt}.
                    Thus $x_0(i) < x_0(i) + \epsilon$ by \cref{real_lt_add,real_add_comm,real_add_ident}.
                    Thus $x_0(i) - \epsilon < x_0(i) + \epsilon$ by \cref{real_lt_trans}.
                    Thus $\intervalopen{x_0(i) - \epsilon}{x_0(i) + \epsilon}\in\standardTopologyR$
                        by \cref{intervalopen_subset_reals,powerset_intro,standard_basis_reals,basis_elem_open_in_generated,standard_basis_is_basis,standard_topology_reals}.
                    Let $U = \intervalopen{x_0(i) - \epsilon}{x_0(i) + \epsilon}$.
                    Thus $U\in\apply{T}{i}$ by \cref{function_apply_intro}.
                    Hence $A\in\productSubbasisAt{T}{X}{i}$
                        by \cref{preimg_iff,projection_map_indexed,pair_eq_iff,subseteq,powerset_intro,product_subbasis_indexed}.
                    Therefore $A\in S_0$ by \cref{product_subbasis_union_indexed}.
                \end{subproof}
                Hence $F\subseteq S_0$ by \cref{subseteq}.
                Show that for all $y$ such that $y\in B_2$ we have $y\in\inters{F}$.
                \begin{subproof}
                            Fix $y$ such that $y\in B_2$.
                            Thus $y\in X_0$ and $d_1((x_0,y)) < \epsilon$ by \cref{metric_ball}.
                            Take $m\in\reals$ such that $((x_0,y),m)\in d_1$ and
                                for all $j\in J$ we have $\abs{x_0(j) - y(j)} \leq m$ by \cref{square_metric_value_exists}.
                            Thus $m < \epsilon$ by \cref{square_metric_function,function_apply_intro,real_lt_subst_left}.
                            Show that for all $A$ such that $A\in F$ we have $y\in A$.
                            \begin{subproof}
                                Fix $A$ such that $A\in F$.
                                Take $i\in J$ such that $(i,A)\in g$ by \cref{img_iff}.
                                Thus $i\mathrel{R_g}\apply{g}{i}$.
                                Take $i_0\in J$ such that
                                    $(i,\apply{g}{i}) = (i_0,\preimg{\piIndex{X}{i_0}}{\intervalopen{x_0(i_0) - \epsilon}{x_0(i_0) + \epsilon}})$.
                                Hence $A = \preimg{\piIndex{X}{i}}{\intervalopen{x_0(i) - \epsilon}{x_0(i) + \epsilon}}$
                                    by \cref{function_apply_intro,pair_eq_iff}.
                                Thus $(y,\apply{y}{i})\in\piIndex{X}{i}$ by
                                    \cref{product_family_reals_power,subseteq,projection_map_indexed}.
                                Thus $x_0(i)\in\reals$ and $\apply{y}{i}\in\reals$
                                    by \cref{reals_power,tuple_power,funs_type_apply}.
                                Thus $x_0(i) - \apply{y}{i}\in\reals$
                                    by \cref{real_add_inverse,real_sub,real_add_closed}.
                                Thus $\abs{x_0(i) - \apply{y}{i}}\in\reals$ by \cref{real_abs_in_reals}.
                                Thus $\abs{x_0(i) - \apply{y}{i}} \leq m$.
                                Therefore $\abs{x_0(i) - \apply{y}{i}} < \epsilon$ by \cref{real_leq_split,real_lt_subst_left,real_lt_trans}.
                                Hence $x_0(i) - \epsilon < \apply{y}{i}$ and $\apply{y}{i} < x_0(i) + \epsilon$
                                    by \cref{real_interval_from_abs_lt}.
                                Therefore $\apply{y}{i}\in\intervalopen{x_0(i) - \epsilon}{x_0(i) + \epsilon}$ by \cref{intervalopen}.
                                Hence $y\in A$ by \cref{preimg_iff}.
                            \end{subproof}
                            Therefore $y\in\inters{F}$ by \cref{inters_iff_forall}.
                \end{subproof}
                Hence $B_2\subseteq\inters{F}$ by \cref{subseteq}.
                Show that for all $y$ such that $y\in\inters{F}$ we have $y\in B_2$.
                \begin{subproof}
                            Fix $y$ such that $y\in\inters{F}$.
                            Take $i$ such that $i\in J$.
                            Hence $\apply{g}{i}\in F$ by \cref{function_apply_elim,img_iff}.
                            Therefore $y\in\apply{g}{i}$ by \cref{inters_iff_forall}.
                            Then $i\mathrel{R_g}\apply{g}{i}$.
                            Take $i_0\in J$ such that
                                $(i,\apply{g}{i}) = (i_0,\preimg{\piIndex{X}{i_0}}{\intervalopen{x_0(i_0) - \epsilon}{x_0(i_0) + \epsilon}})$.
                            Thus $\apply{g}{i} = \preimg{\piIndex{X}{i}}{\intervalopen{x_0(i) - \epsilon}{x_0(i) + \epsilon}}$ by \cref{pair_eq_iff}.
                            Thus $y\in\preimg{\piIndex{X}{i}}{\intervalopen{x_0(i) - \epsilon}{x_0(i) + \epsilon}}$.
                            Take $z\in\intervalopen{x_0(i) - \epsilon}{x_0(i) + \epsilon}$ such that
                                $y\mathrel{\piIndex{X}{i}} z$ by \cref{preimg_iff}.
                            Thus $(y,z)\in\piIndex{X}{i}$.
                            Take $y_0\in\productFamily{X}$ such that $(y,z) = (y_0,\apply{y_0}{i})$ by \cref{projection_map_indexed}.
                            Thus $y = y_0$ by \cref{pair_eq_iff}.
                            Thus $y\in X_0$ by \cref{product_family_reals_power,subseteq}.
                            Take $m\in\reals$ such that $((x_0,y),m)\in d_1$ and
                                for all $j\in J$ we have $\abs{x_0(j) - y(j)} \leq m$ by \cref{square_metric_value_exists}.
                            Show that for all $t\in\coordDiffSet{n}{x_0}{y}$ we have $t < \epsilon$.
                            \begin{subproof}
                                Fix $t$ such that $t\in\coordDiffSet{n}{x_0}{y}$.
                                Take $i_0\in J$ such that $t = \abs{x_0(i_0) - y(i_0)}$ by \cref{coord_diff_set}.
                                Thus $y\in\apply{g}{i_0}$ by \cref{function_apply_elim,img_iff,inters_iff_forall}.
                                Thus $y\in\preimg{\piIndex{X}{i_0}}{\intervalopen{x_0(i_0) - \epsilon}{x_0(i_0) + \epsilon}}$
                                    by \cref{function_apply_intro,pair_eq_iff}.
                                Take $z\in\intervalopen{x_0(i_0) - \epsilon}{x_0(i_0) + \epsilon}$ such that
                                    $y\mathrel{\piIndex{X}{i_0}} z$ by \cref{preimg_iff}.
                                Thus $(y,z)\in\piIndex{X}{i_0}$.
                                Take $y_0\in\productFamily{X}$ such that $(y,z) = (y_0,\apply{y_0}{i_0})$ by \cref{projection_map_indexed}.
                                Thus $y = y_0$ and $z = \apply{y}{i_0}$ by \cref{pair_eq_iff}.
                                Thus $x_0(i_0) - \epsilon < y(i_0)$ and $y(i_0) < x_0(i_0) + \epsilon$ by \cref{intervalopen}.
                                Thus $\abs{x_0(i_0) - y(i_0)} < \epsilon$
                                    by \cref{reals_power,tuple_power,funs_type_apply,real_abs_lt_from_interval}.
                                Thus $t < \epsilon$.
                            \end{subproof}
                            Therefore $d_1((x_0,y)) < \epsilon$ by
                                \cref{square_metric,pair_eq_iff,square_metric_function,function_apply_intro,real_lt_subst_left}.
                            Hence $y\in B_2$ by \cref{metric_ball}.
                \end{subproof}
                Therefore $\inters{F}\subseteq B_2$ by \cref{subseteq}.
                Hence $B_2 = \inters{F}$ by \cref{subseteq_antisymmetric}.
                Therefore $B_2\in B_0$
                    by \cref{inters_iff_forall,unions_iff,subseteq,powerset_intro,finite_intersections}.
                Hence $x\in B_2$ and $B_2\subseteq B_2$.
    \end{subproof}
    Therefore $T_0$ is finer than $T_S$ on $X_0$ by \cref{basis_finer_criterion}.
    Show that for all $x,B_2$ such that $x\in X_0$ and $B_2\in B_0$ and $x\in B_2$
        there exists $B_3\in B_S$ such that $x\in B_3$ and $B_3\subseteq B_2$.
    \begin{subproof}
                Fix $x$.
                Fix $B_2$ such that $x\in X_0$ and $B_2\in B_0$ and $x\in B_2$.
                Take $F\subseteq S_0$ such that $F$ is finite and inhabited and $B_2 = \inters{F}$ by \cref{finite_intersections}.
                Let $R = \{w\in F\times\reals \mid \exists \beta\in J.\exists U\in\standardTopologyR.
                    \fst{w} = \preimg{\piIndex{X}{\beta}}{U} \land \apply{x}{\beta}\in U \land
                    0 < \snd{w} \land \metricBall{\standardMetricR}{\reals}{\apply{x}{\beta}}{\snd{w}} \subseteq U\}$.
                For all $w$ such that $w\in R$ there exist $y,z$ such that $w = (y,z)$.
                Hence $R$ is a relation by \cref{relation}.
                Show that for all $A\in F$ there exists $\epsilon\in\reals$ such that $A\mathrel{R}\epsilon$.
                \begin{subproof}
                    Fix $A$ such that $A\in F$.
                    Hence $x\in A$ by \cref{inters_iff_forall}.
                    Take $\beta\in\dom{X}$ such that $A\in\productSubbasisAt{T}{X}{\beta}$
                        by \cref{subseteq,product_subbasis_union_indexed}.
                    Take $U\in\apply{T}{\beta}$ such that $A = \preimg{\piIndex{X}{\beta}}{U}$ by \cref{product_subbasis_indexed}.
                    Therefore $\beta\in J$ by \cref{subseteq}.
                    Hence $\apply{x}{\beta}\in U$
                        by \cref{product_family_reals_power,subseteq,preimg_iff,projection_map_indexed,pair_eq_iff}.
                    Thus $U\in\standardTopologyR$ by \cref{function_apply_intro}.
                    $\standardBasisR$ is a basis on $\reals$ by \cref{standard_basis_is_basis}.
                    Thus $U\subseteq\reals$
                        by \cref{basis_topology_is_topology,standard_topology_reals,topology_on,subseteq,pow_iff}.
                    Thus $U\in\metricTopology{\standardMetricR}{\reals}$ by \cref{standard_metric_topology}.
                    $\standardMetricR$ is a metric on $\reals$ by \cref{standard_metric_is_metric}.
                    Thus $U$ is open in $\reals$ with respect to $\metricTopology{\standardMetricR}{\reals}$ by \cref{open_in,basis_topology_is_topology,metric_basis_is_basis,metric_topology}.
                    Take $\epsilon\in\reals$ such that $0 < \epsilon$ and
                        $\metricBall{\standardMetricR}{\reals}{\apply{x}{\beta}}{\epsilon}\subseteq U$
                        by \cref{metric_open_iff}.
                    Let $w = (A,\epsilon)$.
                    Hence $A\mathrel{R}\epsilon$ by
                        \cref{times_tuple_intro,fst_eq,snd_eq,pair_eq_iff}.
                \end{subproof}
                Take $e$ such that $e$ is a function on $F$ and for all $A\in F$ we have $A\mathrel{R}\apply{e}{A}$
                    by \cref{choice_function}.
                Let $E = \img{e}{F}$.
                Hence $E$ is finite by \cref{finite_image_function}.
                Take $A$ such that $A\in F$.
                Therefore $\apply{e}{A}\in E$ by \cref{function_apply_elim,img_iff}.
                Hence $E\neq\emptyset$ by \cref{function_apply_elim,img_iff,emptyset}.
                Show that for all $t$ such that $t\in E$ we have $t\in\reals$ and $0 < t$.
                \begin{subproof}
                    Fix $t$ such that $t\in E$.
                    Take $A$ such that $A\in F$ and $(A,t)\in e$ by \cref{img_iff}.
                    Hence $A\mathrel{R}t$.
                    Therefore $t\in\reals$ and $0 < t$.
                \end{subproof}
                Hence $E\subseteq\reals$ by \cref{subseteq}.
                Take $\delta\in E$ such that $\delta$ is a minimum of $E$ by \cref{finite_real_minimum}.
                Therefore $\delta\in\reals$ and $0 < \delta$ by \cref{subseteq}.
                Let $B_3 = \metricBall{d_1}{X_0}{x}{\delta}$.
                Hence $B_3\in B_S$ by \cref{metric_ball,subseteq,powerset_intro,metric_basis}.
                Show that for all $y$ such that $y\in B_3$ we have $y\in B_2$.
                \begin{subproof}
                        Fix $y$ such that $y\in B_3$.
                        Thus $y\in X_0$ and $d_1((x,y)) < \delta$ by \cref{metric_ball}.
                        Take $m\in\reals$ such that $((x,y),m)\in d_1$ and
                            for all $j\in J$ we have $\abs{x(j) - y(j)} \leq m$ by \cref{square_metric_value_exists}.
                        Thus $m < \delta$ by \cref{square_metric_function,function_apply_intro,real_lt_subst_left}.
                        Show that for all $A\in F$ we have $y\in A$.
                        \begin{subproof}
                            Fix $A$ such that $A\in F$.
                            Hence $(A,\apply{e}{A})\in R$.
                            Let $w = (A,\apply{e}{A})$.
                            Thus $w\in R$.
                            Take $\beta,U$ such that $\beta\in J$ and $U\in\standardTopologyR$ and
                                $\fst{w} = \preimg{\piIndex{X}{\beta}}{U}$ and $\apply{x}{\beta}\in U$ and
                                $0 < \snd{w}$ and
                                $\metricBall{\standardMetricR}{\reals}{\apply{x}{\beta}}{\snd{w}} \subseteq U$.
                            $\fst{w} = A$ and $\snd{w} = \apply{e}{A}$ by \cref{fst_eq,snd_eq}.
                            Therefore $A = \preimg{\piIndex{X}{\beta}}{U}$.
                            Hence $0 < \apply{e}{A}$.
                            Therefore $\metricBall{\standardMetricR}{\reals}{\apply{x}{\beta}}{\apply{e}{A}} \subseteq U$.
                            Hence $\delta \leq \apply{e}{A}$ by \cref{function_apply_elim,img_iff,real_minimum}.
                            Therefore $m < \apply{e}{A}$ by \cref{real_leq_split,real_lt_subst_right,real_lt_trans}.
                            Thus $x(\beta)\in\reals$ and $\apply{y}{\beta}\in\reals$
                                by \cref{reals_power,tuple_power,funs_type_apply}.
                            Thus $\abs{x(\beta) - \apply{y}{\beta}}\in\reals$
                                by \cref{reals_power,tuple_power,funs_type_apply,real_sub,real_add_inverse,real_add_closed,real_abs_in_reals}.
                            Thus $\abs{x(\beta) - \apply{y}{\beta}} \leq m$.
                            Thus $\abs{x(\beta) - \apply{y}{\beta}} < \apply{e}{A}$ by \cref{real_lt_subst_left,real_lt_trans}.
                            Therefore $\apply{\standardMetricR}{(\apply{x}{\beta},\apply{y}{\beta})} < \apply{e}{A}$
                                by \cref{standard_metric_apply,real_lt_subst_left}.
                            Hence $\apply{y}{\beta}\in\metricBall{\standardMetricR}{\reals}{\apply{x}{\beta}}{\apply{e}{A}}$ by \cref{metric_ball}.
                            Therefore $\apply{y}{\beta}\in U$ by \cref{subseteq}.
                            Hence $y\in\preimg{\piIndex{X}{\beta}}{U}$ by
                                \cref{metric_ball,product_family_reals_power,subseteq,projection_map_indexed,preimg_iff}.
                            Therefore $y\in A$.
                        \end{subproof}
                        Hence $y\in\inters{F}$ by \cref{inters_iff_forall}.
                        Therefore $y\in B_2$ by \cref{subseteq_antisymmetric}.
                \end{subproof}
                Hence $B_3\subseteq B_2$ by \cref{subseteq}.
                Therefore $x\in B_3$ and $B_3\subseteq B_2$ by \cref{metric_ball,metric_zero_characterization,metric_on}.
    \end{subproof}
    Hence $T_S$ is finer than $T_0$ on $X_0$ by \cref{basis_finer_criterion}.
    Therefore $T_0\subseteq T_S$ and $T_S\subseteq T_0$ by \cref{finer_topology}.
    Hence $T_S = T_0$ by \cref{subseteq_antisymmetric}.
    Therefore $T_E = T_0$.
\end{proof}

% from §20 helper lemma 10: uniform coordinate metric set
\begin{definition}\label{uniform_metric_set}
    $\uniformMetricSet{J}{x}{y} = \{\apply{\boundedMetric{\standardMetricR}}{(x(\alpha),y(\alpha))} \mid \alpha\in J\}$.
\end{definition}

% from §20 helper lemma 11: supremum predicate
\begin{definition}\label{real_supremum_pred}
    $\realSupremum{p}{A}$ iff $p$ is a supremum of $A$.
\end{definition}

% from §20 Item 27: uniform metric
\begin{definition}\label{uniform_metric}
    $\uniformMetric{J} = \{w \mid \exists p\in\realsPower{J}\times\realsPower{J}.
        \exists x\in\realsPower{J}. \exists y\in\realsPower{J}. \exists r\in\reals.
        w = (p,r) \land p = (x,y) \land
        \realSupremum{r}{\uniformMetricSet{J}{x}{y}}\}$.
\end{definition}

% from §20 Item 28: uniform topology
\begin{definition}\label{uniform_topology}
    $\uniformTopology{J} = \metricTopology{\uniformMetric{J}}{\realsPower{J}}$.
\end{definition}

% from §20 Item 29: uniform topology is finer than product topology
\begin{theorem}\label{uniform_finer_product}
    Let $X = \{(\alpha,\reals) \mid \alpha\in J\}$.
    Let $T = \{(\alpha,\standardTopologyR) \mid \alpha\in J\}$.
    Let $T_0 = \productTopologyIndexed{T}{X}$.
    Then $\uniformTopology{J}$ is finer than $T_0$ on $\realsPower{J}$.
\end{theorem}
\begin{proof}
    Omitted. % do not touch
\end{proof}


% from §20 Item 30: uniform topology is coarser than box topology
\begin{theorem}\label{uniform_coarser_box}
    Let $X = \{(\alpha,\reals) \mid \alpha\in J\}$.
    Let $T = \{(\alpha,\standardTopologyR) \mid \alpha\in J\}$.
    Let $T_0 = \boxTopology{T}{X}$.
    Then $\uniformTopology{J}$ is coarser than $T_0$ on $\realsPower{J}$.
\end{theorem}
\begin{proof}
    Omitted. % do not touch
\end{proof}

% from §20 helper lemma 12: omega coordinate metric set
\begin{definition}\label{omega_metric_set}
    $\omegaMetricSet{x}{y} = \{\apply{\boundedMetric{\standardMetricR}}{(x(i),y(i))} / i \mid i\in\naturalsPlus\}$.
\end{definition}

% from §20 Item 31: metric on $\reals^{\omega}$
\begin{definition}\label{omega_metric}
    $\omegaMetric = \{w \mid \exists p\in\realsPower{\naturalsPlus}\times\realsPower{\naturalsPlus}.
        \exists x\in\realsPower{\naturalsPlus}. \exists y\in\realsPower{\naturalsPlus}. \exists r\in\reals.
        w = (p,r) \land p = (x,y) \land
        \realSupremum{r}{\omegaMetricSet{x}{y}}\}$.
\end{definition}

% from §20 Item 32: $\omega$-metric is a metric
\begin{theorem}\label{omega_metric_is_metric}
    Then $\omegaMetric$ is a metric on $\realsPower{\naturalsPlus}$.
\end{theorem}
\begin{proof}
    Omitted. % do not touch
\end{proof}

% from §20 Item 33: $\omega$-metric induces the product topology
\begin{theorem}\label{omega_metric_product_topology}
    Let $X = \{(i,\reals) \mid i\in\naturalsPlus\}$.
    Let $T = \{(i,\standardTopologyR) \mid i\in\naturalsPlus\}$.
    Let $T_0 = \productTopologyIndexed{T}{X}$.
    Then $\metricTopology{\omegaMetric}{\realsPower{\naturalsPlus}} = T_0$.
\end{theorem}
\begin{proof}
    Omitted.  % do not touch
\end{proof}
 
\section{§21 The Metric Topology (continued)}

% from §21 helper lemma 1: positive reciprocal of naturalsplus
\begin{lemma}\label{positive_reciprocal_naturalsplus}
    Suppose $n\in\naturalsPlus$.
    Then $0 < 1 / n$.
\end{lemma}
\begin{proof}
    Follows by \cref{real_zero_lt_one,real_add_ident,real_naturals_subset_reals,real_one_in_naturals,subseteq,real_one_leq_naturalsplus,real_lt_subst_right,real_lt_trans,real_inv_pos,real_pos_nonzero,real_mul_inverse,real_div,real_mul_ident}.
\end{proof}

% from §21 helper lemma 2: metric ball monotonic in radius
\begin{lemma}\label{metric_ball_mono}
    Assume $d$ is a metric on $X$.
    Suppose $x\in X$.
    Suppose $r_1\in\reals$.
    Suppose $r_2\in\reals$.
    Suppose $r_1 < r_2$.
    Then $\metricBall{d}{X}{x}{r_1} \subseteq \metricBall{d}{X}{x}{r_2}$.
\end{lemma}
\begin{proof}
    Follows by \cref{metric_ball,metric_on,times_tuple_intro,funs_type_apply,real_lt_trans,subseteq}.
\end{proof}

% from §21 Item 1: epsilon-delta characterization of continuity in metric spaces
\begin{theorem}\label{metric_continuity_eps_delta}
    Suppose $d_X$ is a metric on $X$.
    Suppose $d_Y$ is a metric on $Y$.
    Let $T_X = \metricTopology{d_X}{X}$.
    Let $T_Y = \metricTopology{d_Y}{Y}$.
    Suppose $f$ is a function from $X$ to $Y$.
    Then $f$ is continuous from $X$ to $Y$ with respect to $T_X$ and $T_Y$ iff
        for all $x\in X$ we have
        for all $\epsilon\in\reals$ such that $0 < \epsilon$
        there exists $\delta\in\reals$ such that $0 < \delta$ and
        for all $y\in X$ if $d_X((x,y)) < \delta$ then $d_Y((f(x),f(y))) < \epsilon$.
\end{theorem}
\begin{proof}
    Show that if $f$ is continuous from $X$ to $Y$ with respect to $T_X$ and $T_Y$ then
        for all $x\in X$ we have
        for all $\epsilon\in\reals$ such that $0 < \epsilon$
        there exists $\delta\in\reals$ such that $0 < \delta$ and
        for all $y\in X$ if $d_X((x,y)) < \delta$ then $d_Y((f(x),f(y))) < \epsilon$.
    \begin{subproof}
        Assume $f$ is continuous from $X$ to $Y$ with respect to $T_X$ and $T_Y$.
        Fix $x\in X$.
        Fix $\epsilon\in\reals$.
        Assume $0 < \epsilon$.
        Let $V = \metricBall{d_Y}{Y}{f(x)}{\epsilon}$.
        $f(x)\in Y$ by \cref{continuous,funs_type_apply}.
        Therefore $V$ is open in $Y$ with respect to $T_Y$ by
            \cref{continuous,funs_type_apply,metric_ball,subseteq,powerset_intro,metric_basis,metric_basis_is_basis,basis_elem_open_in_generated,metric_topology,basis_for_topology,basis_topology_is_topology,open_in}.
        Hence $x\in\preimg{f}{V}$ by
            \cref{metric_zero_characterization,metric_on,real_lt_subst_left,metric_ball,funs,funs_is_function,function_apply_iff,preimg_iff}.
        Take $\delta\in\reals$ such that $0 < \delta$ and
            $\metricBall{d_X}{X}{x}{\delta} \subseteq \preimg{f}{V}$
            by \cref{continuous,preimg_subseteq_dom,funs,metric_open_iff}.
        For all $y\in X$ if $d_X((x,y)) < \delta$ then $d_Y((f(x),f(y))) < \epsilon$ by
            \cref{metric_ball,subseteq,preimg_iff,funs_is_function,function_apply_iff}.
        Therefore there exists $\delta\in\reals$ such that $0 < \delta$ and
            for all $y\in X$ if $d_X((x,y)) < \delta$ then $d_Y((f(x),f(y))) < \epsilon$.
    \end{subproof}
    Show that if
        for all $x\in X$ we have
        for all $\epsilon\in\reals$ such that $0 < \epsilon$
        there exists $\delta\in\reals$ such that $0 < \delta$ and
        for all $y\in X$ if $d_X((x,y)) < \delta$ then $d_Y((f(x),f(y))) < \epsilon$
        then $f$ is continuous from $X$ to $Y$ with respect to $T_X$ and $T_Y$.
    \begin{subproof}
        Assume for all $x\in X$ we have
            for all $\epsilon\in\reals$ such that $0 < \epsilon$
            there exists $\delta\in\reals$ such that $0 < \delta$ and
            for all $y\in X$ if $d_X((x,y)) < \delta$ then $d_Y((f(x),f(y))) < \epsilon$.
        Show that for all $V$ such that $V$ is open in $Y$ with respect to $T_Y$ we have
            $\preimg{f}{V}$ is open in $X$ with respect to $T_X$.
        \begin{subproof}
            Fix $V$ such that $V$ is open in $Y$ with respect to $T_Y$.
            Hence $V\subseteq Y$ by
                \cref{open_in,basis_topology_is_topology,metric_topology,metric_basis_is_basis,basis_for_topology,topology_on,subseteq,pow_iff}.
            Show that for all $x\in\preimg{f}{V}$ there exists $\delta\in\reals$ such that
                $0 < \delta$ and $\metricBall{d_X}{X}{x}{\delta} \subseteq \preimg{f}{V}$.
            \begin{subproof}
                Fix $x\in\preimg{f}{V}$.
                Take $z\in V$ such that $(x,z)\in f$ by \cref{preimg_iff}.
                Take $\epsilon\in\reals$ such that $0 < \epsilon$ and
                    $\metricBall{d_Y}{Y}{f(x)}{\epsilon} \subseteq V$
                    by \cref{preimg_iff,funs_is_function,preimg_subseteq_dom,elem_subseteq,function_apply_iff,subseteq,metric_open_iff}.
                Take $\delta\in\reals$ such that $0 < \delta$ and
                    for all $y\in X$ if $d_X((x,y)) < \delta$ then $d_Y((f(x),f(y))) < \epsilon$.
                Hence $\metricBall{d_X}{X}{x}{\delta} \subseteq \preimg{f}{V}$ by
                    \cref{funs_type_apply,metric_ball,subseteq,funs_is_function,funs,function_apply_iff,preimg_iff}.
            \end{subproof}
            Therefore $\preimg{f}{V}$ is open in $X$ with respect to $T_X$ by
                \cref{preimg_subseteq_dom,funs,metric_open_iff}.
        \end{subproof}
        Therefore $f$ is continuous from $X$ to $Y$ with respect to $T_X$ and $T_Y$ by
            \cref{basis_topology_is_topology,metric_topology,metric_basis_is_basis,basis_for_topology,continuous}.
    \end{subproof}
\end{proof}

% from §21 Item 2: sequence implies closure
\begin{lemma}\label{sequence_implies_closure}
    Assume $T$ is a topology on $X$.
    Suppose $A\subseteq X$.
    Suppose $s$ is a sequence in $A$.
    Suppose $x\in X$.
    Suppose $s$ converges to $x$ in $X$ with respect to $T$.
    Then $x\in\closure{A}{X}{T}$.
\end{lemma}
\begin{proof}
    For all $U$ such that $U$ is open in $X$ with respect to $T$ and $x\in U$ we have $U$ intersects $A$ by \cref{converges_to,neighborhood,real_naturals_closed_succ,real_naturals_subset_reals,subseteq,real_add_ident,real_one_in_naturals,real_zero_lt_one,real_lt_add,real_add_comm,real_add_closed,real_lt_subst_left,real_lt_imp_leq,sequence_in,funs_type_apply,inter_intro,emptyset,intersects}.
    Hence $x\in\closure{A}{X}{T}$ by \cref{closure_iff_intersects_open}.
\end{proof}

% from §21 Item 3: closure implies sequence in a metrizable space
\begin{lemma}\label{closure_implies_sequence_metrizable}
    Assume $T$ is a topology on $X$.
    Suppose $A\subseteq X$.
    Suppose $x\in X$.
    Suppose $X$ is metrizable with respect to $T$.
    Suppose $x\in\closure{A}{X}{T}$.
    Then there exists $s$ such that $s$ is a sequence in $A$ and $s$ converges to $x$ in $X$ with respect to $T$.
\end{lemma}
\begin{proof}
    Take $d$ such that $d$ is a metric on $X$ and $T = \metricTopology{d}{X}$ by \cref{metrizable}.
    Let $b = \{(k,\metricBall{d}{X}{x}{1 / k}) \mid k\in\naturalsPlus\}$.
    For all $k,U_1,U_2$ such that $(k,U_1)\in b$ and $(k,U_2)\in b$ we have $U_1 = U_2$ by \cref{pair_eq_iff}.
    Hence $b$ is right-unique by \cref{rightunique}.
    Therefore $b$ is a function.
    Hence $\dom{b} = \naturalsPlus$ by \cref{pair_eq_iff,dom_intro,dom_iff,subseteq,subseteq_antisymmetric}.
    Therefore $b$ is a function on $\naturalsPlus$.

    Let $F = \{(n,\img{\restrl{b}{\initialSegment{n}}}{\initialSegment{n}}) \mid n\in\naturalsPlus\}$.
    Let $B = \{(n,\inters{\apply{F}{n}}) \mid n\in\naturalsPlus\}$.

    For all $n,U_1,U_2$ such that $(n,U_1)\in F$ and $(n,U_2)\in F$ we have $U_1 = U_2$
        by \cref{pair_eq_iff}.
    Hence $F$ is right-unique by \cref{rightunique}.
    Therefore $F$ is a function.
    Hence $\dom{F} = \naturalsPlus$ by \cref{pair_eq_iff,dom_intro,dom_iff,subseteq,subseteq_antisymmetric}.
    Therefore $F$ is a function on $\naturalsPlus$.

    For all $n\in\naturalsPlus$ we have $F(n) = \img{\restrl{b}{\initialSegment{n}}}{\initialSegment{n}}$
        by \cref{pair_eq_iff,function_apply_intro}.

    For all $n,U_1,U_2$ such that $(n,U_1)\in B$ and $(n,U_2)\in B$ we have $U_1 = U_2$
        by \cref{pair_eq_iff}.
    Hence $B$ is right-unique by \cref{rightunique}.
    Therefore $B$ is a function.
    Hence $\dom{B} = \naturalsPlus$ by \cref{pair_eq_iff,dom_intro,dom_iff,subseteq,subseteq_antisymmetric}.
    Therefore $B$ is a function on $\naturalsPlus$.

    For all $n\in\naturalsPlus$ we have $B(n) = \inters{F(n)}$ by \cref{pair_eq_iff,function_apply_intro}.

    Show that for all $n\in\naturalsPlus$ we have $B(n)$ is a neighborhood of $x$ in $X$ with respect to $T$.
    \begin{subproof}
        Fix $n\in\naturalsPlus$.
        $\initialSegment{n}$ is finite by \cref{initial_segment_finite}.
        Hence $\initialSegment{n}\subseteq\dom{b}$
            by \cref{initial_segment_naturalsplus,pair_eq_iff,dom_intro,subseteq}.
        Let $g = \restrl{b}{\initialSegment{n}}$.
        $g$ is a function by \cref{restrl_of_function_is_function}.
        Hence $\dom{g} = \initialSegment{n}$ by
            \cref{restrl_dom,initial_segment_naturalsplus,pair_eq_iff,restrl_iff,dom_intro,subseteq,subseteq_antisymmetric}.
        Therefore $g$ is a function on $\initialSegment{n}$.
        Hence $F(n) = \img{g}{\initialSegment{n}}$.
        Therefore $F(n)$ is finite by \cref{finite_image_function}.
        Hence $1\in\initialSegment{n}$ by \cref{real_one_in_naturals,real_one_leq_naturalsplus,initial_segment_naturalsplus}.
        Let $U_1 = \metricBall{d}{X}{x}{1 / 1}$.
        $(1,U_1)\in b$ by \cref{initial_segment_naturalsplus,pair_eq_iff}.
        Hence $1\in\dom{g}\inter\initialSegment{n}$ by \cref{subseteq,restrl_iff,dom_intro,inter_intro}.
        Therefore $b(1)\in F(n)$ by \cref{restrl_of_function_apply,img_of_function_intro}.
        Hence $F(n)$ is inhabited.
        For all $U$ such that $U\in F(n)$ we have $U\in T$ by
            \cref{pair_eq_iff,function_apply_intro,img_iff,restrl_iff,initial_segment_naturalsplus,positive_reciprocal_naturalsplus,real_naturals_subset_reals,real_one_in_naturals,subseteq,real_add_ident,real_zero_lt_one,real_one_leq_naturalsplus,real_lt_subst_right,real_lt_trans,real_pos_nonzero,real_mul_inverse,real_div,real_mul_closed,metric_ball,powerset_intro,metric_basis,metric_basis_is_basis,basis_elem_open_in_generated,metric_topology,basis_for_topology}.
        Hence $F(n)\subseteq T$ by \cref{subseteq}.
        Therefore $B(n)\in T$ by \cref{finite_intersections_open_in_topology}.
        For all $U$ such that $U\in F(n)$ we have $x\in U$ by
            \cref{pair_eq_iff,function_apply_intro,img_iff,restrl_iff,initial_segment_naturalsplus,positive_reciprocal_naturalsplus,metric_zero_characterization,metric_on,real_lt_subst_left,metric_ball}.
        Hence $B(n)$ is a neighborhood of $x$ in $X$ with respect to $T$ by \cref{inters_iff_forall,neighborhood,open_in}.
    \end{subproof}

    For all $n\in\naturalsPlus$ we have $B(n)$ intersects $A$ by \cref{closure_iff_intersects_open,neighborhood}.

    Let $R = \{w\in\naturalsPlus\times X \mid \snd{w}\in A\inter \apply{B}{\fst{w}}\}$.
    Hence $R$ is a relation by \cref{times_elem_is_tuple,relation}.
    For all $n\in\naturalsPlus$ there exists $y$ such that $n\mathrel{R} y$ by
        \cref{intersects,inter_comm,nonempty_inhabited,inter_elim_left,inter_elim_right,subseteq,times_tuple_intro,fst_eq,snd_eq,inter_intro,pair_eq_iff}.
    Take $s$ such that $s$ is a function on $\naturalsPlus$ and
        for all $n\in\naturalsPlus$ we have $n\mathrel{R} s(n)$ by \cref{axiom_of_choice}.

    Therefore $s$ is a sequence in $A$ by \cref{funs_intro,sequence_in,subseteq,pair_eq_iff,fst_eq,snd_eq,inter}.

    Therefore $s$ is a sequence in $X$ by \cref{sequence_in,funs_weaken_codom}.
    Show that for all $U$ such that $U$ is a neighborhood of $x$ in $X$ with respect to $T$
        there exists $N\in\naturalsPlus$ such that for all $n\in\naturalsPlus$ if $N \leq n$ then $s(n)\in U$.
    \begin{subproof}
        Fix $U$ such that $U$ is a neighborhood of $x$ in $X$ with respect to $T$.
        Hence $U$ is open in $X$ with respect to $T$ and $x\in U$ by \cref{neighborhood}.
        Therefore $U\in T$ and $U\subseteq X$ by \cref{open_in,topology_on,subseteq,pow_iff}.
        Take $\epsilon\in\reals$ such that $0 < \epsilon$ and
            $\metricBall{d}{X}{x}{\epsilon} \subseteq U$ by \cref{metric_open_iff}.
        Take $N\in\naturalsPlus$ such that $0 < 1 / N$ and $1 / N < \epsilon$ by \cref{real_small_reciprocal}.
        Therefore $\metricBall{d}{X}{x}{1 / N} \subseteq U$ by
            \cref{real_naturals_subset_reals,real_one_in_naturals,subseteq,real_add_ident,real_zero_lt_one,real_one_leq_naturalsplus,real_lt_subst_right,real_lt_trans,real_pos_nonzero,real_mul_inverse,real_div,real_mul_closed,metric_ball_mono,subseteq_transitive}.
        Show that for all $n\in\naturalsPlus$ if $N \leq n$ then $s(n)\in U$.
        \begin{subproof}
            Fix $n\in\naturalsPlus$.
            Assume $N \leq n$.
            Therefore $F(N)\subseteq F(n)$ by
                \cref{pair_eq_iff,function_apply_intro,img_iff,restrl_iff,initial_segment_naturalsplus,real_naturals_subset_reals,subseteq,real_leq_trans}.
            For all $U_0$ such that $U_0\in F(N)$ we have $B(n)\subseteq U_0$ by
                \cref{subseteq,inters_subseteq_elem}.
            Hence $F(N)$ is inhabited.
            Therefore $B(n)\subseteq B(N)$ by \cref{subseteq_inters_iff}.
            $N\in\initialSegment{N}$ by \cref{initial_segment_naturalsplus}.
            For all $i$ such that $i\in\initialSegment{N}$ we have $i\in\dom{b}$ by
                \cref{initial_segment_naturalsplus,pair_eq_iff,dom_intro}.
            Hence $\initialSegment{N}\subseteq\dom{b}$ by \cref{subseteq}.
            Therefore $(N,b(N))\in b$ by \cref{subseteq,function_apply_elim}.
            Therefore $b(N)\in F(N)$ by \cref{restrl_iff,img_iff}.
            Therefore $B(N)\subseteq \metricBall{d}{X}{x}{1 / N}$ by
                \cref{inters_subseteq_elem,pair_eq_iff}.
            Therefore $B(n)\subseteq \metricBall{d}{X}{x}{1 / N}$ by \cref{subseteq_transitive}.
            Hence $s(n)\in U$ by \cref{pair_eq_iff,fst_eq,snd_eq,inter,subseteq}.
        \end{subproof}
        Therefore there exists $N\in\naturalsPlus$ such that for all $n\in\naturalsPlus$ if $N \leq n$ then $s(n)\in U$.
    \end{subproof}
    Hence $s$ converges to $x$ in $X$ with respect to $T$ by \cref{converges_to}.

    Therefore there exists $s$ such that $s$ is a sequence in $A$ and $s$ converges to $x$ in $X$ with respect to $T$.
\end{proof}

% from §21 Item 4: continuous implies sequential continuity
\begin{theorem}\label{continuous_implies_sequential}
    Assume $T$ is a topology on $X$.
    Assume $T_1$ is a topology on $Y$.
    Suppose $f$ is a function from $X$ to $Y$.
    Suppose $f$ is continuous from $X$ to $Y$ with respect to $T$ and $T_1$.
    Then for all $s$ such that $s$ is a sequence in $X$ we have
        for all $x\in X$ if $s$ converges to $x$ in $X$ with respect to $T$ then
        $f\circ s$ converges to $f(x)$ in $Y$ with respect to $T_1$.
\end{theorem}
\begin{proof}
    Fix $s$ such that $s$ is a sequence in $X$.
    Fix $x\in X$.
    Assume $s$ converges to $x$ in $X$ with respect to $T$.
    Hence $f\circ s$ is a sequence in $Y$ by \cref{continuous,funs_circ,sequence_in}.
    For all $V$ such that $V$ is a neighborhood of $f(x)$ in $Y$ with respect to $T_1$
        there exists $N\in\naturalsPlus$ such that for all $n\in\naturalsPlus$ if $N\leq n$ then
        $(f\circ s)(n)\in V$ by \cref{neighborhood,continuous,funs,funs_is_function,function_apply_iff,preimg_iff,converges_to,sequence_in,funs_elim,function_to_ran,subseteq,circ_apply}.
    Therefore $f\circ s$ converges to $f(x)$ in $Y$ with respect to $T_1$ by \cref{converges_to}.
\end{proof}
% from §21 Item 5: sequential continuity implies continuity in metrizable spaces
\begin{theorem}\label{sequential_implies_continuous_metrizable}
    Assume $T$ is a topology on $X$.
    Assume $T_1$ is a topology on $Y$.
    Suppose $f$ is a function from $X$ to $Y$.
    Suppose $X$ is metrizable with respect to $T$.
    Suppose for all $s$ such that $s$ is a sequence in $X$ we have
        for all $x\in X$ if $s$ converges to $x$ in $X$ with respect to $T$ then
        $f\circ s$ converges to $f(x)$ in $Y$ with respect to $T_1$.
    Then $f$ is continuous from $X$ to $Y$ with respect to $T$ and $T_1$.
\end{theorem}
\begin{proof}
    Show that for all $A$ such that $A\subseteq X$ we have
        $\img{f}{\closure{A}{X}{T}}\subseteq \closure{\img{f}{A}}{Y}{T_1}$.
    \begin{subproof}
        Fix $A$ such that $A\subseteq X$.
        Show that for all $y$ such that $y\in\img{f}{\closure{A}{X}{T}}$ we have
            $y\in\closure{\img{f}{A}}{Y}{T_1}$.
        \begin{subproof}
            Fix $y$ such that $y\in\img{f}{\closure{A}{X}{T}}$.
            Take $x\in\dom{f}\inter\closure{A}{X}{T}$ such that $y = f(x)$ by \cref{funs_is_function,img_of_function_elim}.
            Take $s$ such that $s$ is a sequence in $A$ and
                $s$ converges to $x$ in $X$ with respect to $T$
                by \cref{inter_elim_right,closure,closure_implies_sequence_metrizable}.
            Hence $x\in X$ and $s$ is a sequence in $X$ and
                $f\circ s$ converges to $f(x)$ in $Y$ with respect to $T_1$ by
                \cref{inter_elim_right,closure,sequence_in,funs_weaken_codom}.
            Hence $f\circ s\in\funs{\naturalsPlus}{Y}$ by \cref{funs_circ,sequence_in}.
            For all $n\in\dom{(f\circ s)}$ we have $(f\circ s)(n)\in\img{f}{A}$ by
                \cref{funs,sequence_in,funs_elim,funs_type_apply,function_to_ran,subseteq,subseteq_transitive,inter_intro,img_of_function_intro,circ_apply}.
            Hence $f\circ s$ is a function to $\img{f}{A}$.
            Hence $f\circ s$ is a sequence in $\img{f}{A}$ by \cref{funs,funs_intro,sequence_in}.
        Then $f(x)\in\closure{\img{f}{A}}{Y}{T_1}$ by
            \cref{img_subseteq_ran,funs_elim,function_to_ran,subseteq_transitive,funs_type_apply,sequence_implies_closure}.
        Therefore $y\in\closure{\img{f}{A}}{Y}{T_1}$.
        \end{subproof}
        Therefore $\img{f}{\closure{A}{X}{T}}\subseteq \closure{\img{f}{A}}{Y}{T_1}$ by \cref{subseteq}.
    \end{subproof}
    Hence for all $B$ such that $B$ is closed in $Y$ with respect to $T_1$
        we have $\preimg{f}{B}$ is closed in $X$ with respect to $T$ by
        \cref{closure_image_subset_implies_preimg_closed}.
    Therefore $f$ is continuous from $X$ to $Y$ with respect to $T$ and $T_1$ by
        \cref{preimg_closed_implies_continuous}.
\end{proof}

% from §21 Item 6: countable basis at a point
\begin{definition}\label{countable_basis_at_point}
    $b$ is a countable basis at $x$ in $X$ with respect to $T$ iff
        $b$ is a sequence in $\pow{X}$ and
        for all $n\in\naturalsPlus$ we have
            $b(n)$ is a neighborhood of $x$ in $X$ with respect to $T$ and
        for all $U$ such that $U$ is a neighborhood of $x$ in $X$ with respect to $T$
            there exists $n\in\naturalsPlus$ such that $b(n)\subseteq U$.
\end{definition}

% from §21 Item 7: first countability axiom
\begin{definition}\label{first_countability_axiom}
    $X$ satisfies the first countability axiom with respect to $T$ iff
        for all $x\in X$ there exists $b$ such that
            $b$ is a countable basis at $x$ in $X$ with respect to $T$.
\end{definition}

% from §21 helper definition 1: addition operation on reals
\begin{definition}\label{real_addition_operation}
    $\realAddition = \{w \mid \exists p\in\reals\times\reals. \exists z\in\reals.
        w = (p,z) \land z = \fst{p} + \snd{p}\}$.
\end{definition}

% from §21 helper definition 2: subtraction operation on reals
\begin{definition}\label{real_subtraction_operation}
    $\realSubtraction = \{w \mid \exists p\in\reals\times\reals. \exists z\in\reals.
        w = (p,z) \land z = \fst{p} - \snd{p}\}$.
\end{definition}

% from §21 helper definition 3: multiplication operation on reals
\begin{definition}\label{real_multiplication_operation}
    $\realMultiplication = \{w \mid \exists p\in\reals\times\reals. \exists z\in\reals.
        w = (p,z) \land z = \fst{p} \rmul \snd{p}\}$.
\end{definition}

% from §21 helper definition 4: quotient operation on reals
\begin{definition}\label{real_division_operation}
    $\realDivision = \{w \mid \exists p\in\reals\times(\reals\setminus\{0\}). \exists z\in\reals.
        w = (p,z) \land z = \fst{p} \rmul \inv{\snd{p}}\}$.
\end{definition}

% from §21 helper lemma 3: standard metric value
\begin{lemma}\label{standard_metric_value}
    Suppose $x\in\reals$.
    Suppose $y\in\reals$.
    Then $\apply{\standardMetricR}{(x,y)} = \abs{x - y}$.
\end{lemma}
\begin{proof}
    Therefore $\apply{\standardMetricR}{(x,y)} = \abs{x - y}$ by
        \cref{standard_metric_is_metric,metric_on,times_tuple_intro,fst_eq,snd_eq,real_add_inverse,real_sub,real_add_closed,real_abs_in_reals,standard_metric_r,funs_elim,function_apply_intro}.
\end{proof}

% from §21 Item 8: addition operation is continuous
\begin{lemma}\label{real_addition_continuous}
    Let $T = \productTopology{\standardTopologyR}{\standardTopologyR}{\reals}{\reals}$.
    Then $\realAddition$ is continuous from $\reals\times\reals$ to $\reals$ with respect to $T$ and $\standardTopologyR$.
\end{lemma}
\begin{proof}
    Let $f = \realAddition$.
    For all $p,z_1,z_2$ such that $(p,z_1)\in f$ and $(p,z_2)\in f$ we have $z_1 = z_2$
        by \cref{real_addition_operation,pair_eq_iff}.
    Hence $f$ is right-unique by \cref{rightunique}.
    Therefore $f$ is a relation by \cref{real_addition_operation,relation}.
    Hence $f$ is a function.
    For all $p\in\dom{f}$ we have $f(p)\in\reals$
        by \cref{dom_iff,real_addition_operation,pair_eq_iff,function_apply_intro}.
    Therefore $f$ is a function to $\reals$.
    Hence $\dom{f} = \reals\times\reals$ by
        \cref{dom_iff,real_addition_operation,pair_eq_iff,subseteq,times_elem_is_tuple,real_add_closed,fst_eq,snd_eq,dom_intro,subseteq_antisymmetric}.
    Therefore $f\in\funs{\reals\times\reals}{\reals}$ by \cref{funs_intro}.
    Show that for all $p\in\reals\times\reals$ we have
        $f$ is continuous at $p$ from $\reals\times\reals$ to $\reals$ with respect to $T$ and $\standardTopologyR$.
        \begin{subproof}
            Fix $p$ such that $p\in\reals\times\reals$.
            Take $x_0,y_0$ such that $p = (x_0,y_0)$ and $x_0\in\reals$ and $y_0\in\reals$ by \cref{times_elem_is_tuple}.
            Hence $f(p) = x_0 + y_0$ by
                \cref{fst_eq,snd_eq,real_add_closed,real_addition_operation,function_apply_intro}.
        Show that for all $V$ such that $V$ is a neighborhood of $f(p)$ in $\reals$ with respect to $\standardTopologyR$
            there exists $U$ such that $U$ is a neighborhood of $p$ in $\reals\times\reals$ with respect to $T$ and $\img{f}{U}\subseteq V$.
        \begin{subproof}
            Fix $V$ such that $V$ is a neighborhood of $f(p)$ in $\reals$ with respect to $\standardTopologyR$.
            Let $d = \standardMetricR$.
            $\metricTopology{d}{\reals} = \standardTopologyR$ by \cref{standard_metric_topology}.
            $\standardBasisR$ is a basis on $\reals$ and $\standardTopologyR$ is a topology on $\reals$
                by \cref{standard_basis_is_basis,basis_topology_is_topology,standard_topology_reals}.
            Hence $V\subseteq\reals$ by
                \cref{neighborhood,open_in,standard_basis_is_basis,basis_topology_is_topology,standard_topology_reals,topology_on,subseteq,pow_iff}.
            $d$ is a metric on $\reals$ by \cref{standard_metric_is_metric}.
            Take $\epsilon\in\reals$ such that $0 < \epsilon$ and
                $\metricBall{d}{\reals}{f(p)}{\epsilon} \subseteq V$ by \cref{neighborhood,standard_metric_topology,metric_open_iff}.
            Let $\delta = \epsilon / (1 + 1)$.
            $1\in\reals$ and $0\in\reals$ and $0 < 1$
                by \cref{real_mul_ident,real_add_ident,real_zero_lt_one}.
            $0 + 1 < 1 + 1$ and $0 + 1 = 1$ and $1 + 1\in\reals$
                by \cref{real_lt_add,real_add_ident,real_add_closed}.
            Thus $0 < 1 + 1$ by \cref{real_lt_subst_left,real_lt_trans}.
            Thus $(1 + 1)\neq 0$ by \cref{real_pos_nonzero}.
            Thus $\inv{(1 + 1)}\in\reals$ by \cref{real_mul_inverse}.
            $\delta = \epsilon \rmul \inv{(1 + 1)}$ by \cref{real_div}.
            $\delta\in\reals$ by \cref{real_mul_closed}.
            $0\rmul \inv{(1 + 1)} = 0$ by \cref{real_mul_zero}.
            $0 < \inv{(1 + 1)}$ by \cref{real_inv_pos}.
            Thus $0 < \delta$ by \cref{real_lt_mul_pos}.
            Let $U_1 = \metricBall{d}{\reals}{x_0}{\delta}$.
            Let $U_2 = \metricBall{d}{\reals}{y_0}{\delta}$.
            $x_0\in\reals$ and $\delta\in\reals$ and $0 < \delta$.
            Thus $U_1$ is open in $\reals$ with respect to $\standardTopologyR$
                by \cref{metric_ball,subseteq,pow_iff,metric_basis,metric_basis_is_basis,basis_elem_open_in_generated,metric_topology,basis_for_topology,open_in}.
            $y_0\in\reals$ and $\delta\in\reals$ and $0 < \delta$.
            Thus $U_2$ is open in $\reals$ with respect to $\standardTopologyR$
                by \cref{metric_ball,subseteq,pow_iff,metric_basis,metric_basis_is_basis,basis_elem_open_in_generated,metric_topology,basis_for_topology,open_in}.
            Let $U = U_1\times U_2$.
            Thus $U\in\productBasis{\standardTopologyR}{\standardTopologyR}{\reals}{\reals}$ by
                \cref{open_in,topology_on,subseteq,pow_iff,times_subseteq_left,times_subseteq_right,subseteq_transitive,product_basis}.
            Thus $U$ is open in $\reals\times\reals$ with respect to $T$ by
                \cref{product_basis_is_basis,basis_topology_is_topology,product_topology,basis_elem_open_in_generated,open_in}.
            Thus $p\in U$ by
                \cref{metric_zero_characterization,standard_metric_is_metric,metric_on,real_lt_subst_left,metric_ball,times_tuple_intro}.
            Show that for all $q$ such that $q\in U$ we have
                $f(q)\in\metricBall{d}{\reals}{f(p)}{\epsilon}$.
            \begin{subproof}
                    Fix $q$ such that $q\in U$.
                    Take $x,y$ such that $q = (x,y)$ and $x\in U_1$ and $y\in U_2$ by \cref{times_elem_is_tuple}.
                    $x\in\reals$ and $y\in\reals$ by \cref{metric_ball}.
                    Thus $\abs{x - x_0} < \delta$ and $\abs{y - y_0} < \delta$
                        by \cref{metric_ball,standard_metric_value,real_abs_sub_swap,real_lt_subst_left}.
                    $(q, x + y)\in f$ by \cref{real_add_closed,times_tuple_intro,fst_eq,snd_eq,real_addition_operation}.
                    Thus $f(q) = x + y$ by \cref{function_apply_intro}.
                    $x - x_0 = x + \neg{x_0}$ by \cref{real_sub}.
                    $y - y_0 = y + \neg{y_0}$ by \cref{real_sub}.
                    $(x + y) - (x_0 + y_0) = (x + y) + \neg{(x_0 + y_0)}$ by \cref{real_sub}.
                    $\neg{(x_0 + y_0)} = \neg{x_0} + \neg{y_0}$ by \cref{real_neg_addition}.
                    $(x + y) + (\neg{x_0} + \neg{y_0}) = (x + \neg{x_0}) + (y + \neg{y_0})$ by \cref{real_add_assoc,real_add_comm,real_add_inverse,real_add_closed}.
                    Thus $(x + y) - (x_0 + y_0) = (x - x_0) + (y - y_0)$.
                    $x - x_0\in\reals$ and $y - y_0\in\reals$ by
                        \cref{real_sub,real_add_inverse,real_add_closed}.
                    $x + y\in\reals$ by \cref{real_add_closed}.
                    $x_0 + y_0\in\reals$ by \cref{real_add_closed}.
                    Thus $(x + y) - (x_0 + y_0)\in\reals$ and $\abs{(x + y) - (x_0 + y_0)}\in\reals$ by
                        \cref{real_sub,real_add_closed,real_abs_in_reals}.
                    Thus $\abs{(x + y) - (x_0 + y_0)} \leq \abs{x - x_0} + \abs{y - y_0}$ by \cref{real_triangle_inequality}.
                    $\abs{x - x_0}\in\reals$ by \cref{real_abs_in_reals}.
                    $\abs{y - y_0}\in\reals$ by \cref{real_abs_in_reals}.
                    $\abs{x - x_0} + \abs{y - y_0}\in\reals$ by \cref{real_add_closed}.
                    $\delta + \delta\in\reals$ by \cref{real_add_closed}.
                    $\abs{x - x_0} + \abs{y - y_0} < \delta + \delta$ by \cref{real_lt_add_two}.
                    Thus $\abs{(x + y) - (x_0 + y_0)} < \delta + \delta$
                        by \cref{real_leq_split,real_lt_trans,real_lt_subst_left}.
                    $1\rmul \delta = \delta$ by \cref{real_mul_ident}.
                    $(1 + 1) \rmul \delta = (1 \rmul \delta) + (1 \rmul \delta)$ by \cref{real_right_distrib}.
                    Thus $(1 \rmul \delta) + (1 \rmul \delta) = \delta + (1 \rmul \delta)$ by \cref{real_add_eq_subst}.
                    Thus $(1 \rmul \delta) + (1 \rmul \delta) = \delta + \delta$ by \cref{real_add_eq_subst}.
                    Hence $(1 + 1) \rmul \delta = \delta + \delta$.
                    Thus $(1 + 1) \rmul \delta = (1 + 1) \rmul (\epsilon \rmul \inv{(1 + 1)})$.
                    $(1 + 1) \rmul (\epsilon \rmul \inv{(1 + 1)}) = \epsilon \rmul ((1 + 1) \rmul \inv{(1 + 1)})$ by \cref{real_mul_assoc,real_mul_comm}.
                    $(1 + 1) \rmul \inv{(1 + 1)} = 1$ by \cref{real_mul_inverse}.
                    Thus $\epsilon \rmul 1 = \epsilon$ by \cref{real_mul_comm,real_mul_ident}.
                    Then $\delta + \delta = \epsilon$.
                    Therefore $\abs{(x + y) - (x_0 + y_0)} < \epsilon$ by \cref{real_lt_subst_right}.
                    Hence $\abs{f(q) - f(p)} = \abs{(x + y) - (x_0 + y_0)}$.
                    Therefore $\abs{f(q) - f(p)} < \epsilon$ by \cref{real_lt_subst_left}.
                $f$ is a function to $\reals$ and $\dom{f} = \reals\times\reals$ by \cref{funs_elim}.
                Thus $f(p)\in\reals$ and $f(q)\in\reals$.
                Hence $d((f(p),f(q))) < \epsilon$ by \cref{standard_metric_value,real_abs_sub_swap,real_lt_subst_left}.
                Thus $f(q)\in\metricBall{d}{\reals}{f(p)}{\epsilon}$ by \cref{metric_ball}.
            \end{subproof}
            Thus $\img{f}{U}\subseteq V$ by
                \cref{img_iff,function_apply_intro,subseteq,subseteq_transitive}.
            Thus $U$ is a neighborhood of $p$ in $\reals\times\reals$ with respect to $T$ by \cref{neighborhood}.
        \end{subproof}
        Therefore $f$ is continuous at $p$ from $\reals\times\reals$ to $\reals$ with respect to $T$ and $\standardTopologyR$ by
            \cref{standard_basis_is_basis,basis_topology_is_topology,standard_topology_reals,product_basis_is_basis,product_topology,continuous_at}.
    \end{subproof}
    Hence $f$ is continuous from $\reals\times\reals$ to $\reals$ with respect to $T$ and $\standardTopologyR$ by
        \cref{standard_basis_is_basis,basis_topology_is_topology,standard_topology_reals,product_basis_is_basis,product_topology,continuous_at_implies_continuous}.
\end{proof}

% from §21 Item 9: subtraction operation is continuous
\begin{lemma}\label{real_subtraction_continuous}
    Let $T = \productTopology{\standardTopologyR}{\standardTopologyR}{\reals}{\reals}$.
    Then $\realSubtraction$ is continuous from $\reals\times\reals$ to $\reals$ with respect to $T$ and $\standardTopologyR$.
\end{lemma}
\begin{proof}
    Let $f = \realSubtraction$.
    For all $p,z_1,z_2$ such that $(p,z_1)\in f$ and $(p,z_2)\in f$ we have $z_1 = z_2$
        by \cref{real_subtraction_operation,pair_eq_iff}.
    Hence $f$ is right-unique by \cref{rightunique}.
    Therefore $f$ is a relation by \cref{real_subtraction_operation,relation}.
    Hence $f$ is a function.
    For all $p\in\dom{f}$ we have $f(p)\in\reals$
        by \cref{dom_iff,real_subtraction_operation,pair_eq_iff,function_apply_intro}.
    Therefore $f$ is a function to $\reals$.
    Hence $\dom{f} = \reals\times\reals$ by
        \cref{dom_iff,real_subtraction_operation,pair_eq_iff,subseteq,times_elem_is_tuple,real_sub,real_add_closed,real_add_inverse,fst_eq,snd_eq,dom_intro,subseteq_antisymmetric}.
    Therefore $f\in\funs{\reals\times\reals}{\reals}$ by \cref{funs_intro}.
    Show that for all $p\in\reals\times\reals$ we have
        $f$ is continuous at $p$ from $\reals\times\reals$ to $\reals$ with respect to $T$ and $\standardTopologyR$.
        \begin{subproof}
            Fix $p$ such that $p\in\reals\times\reals$.
            Take $x_0,y_0$ such that $p = (x_0,y_0)$ and $x_0\in\reals$ and $y_0\in\reals$ by \cref{times_elem_is_tuple}.
            Hence $f(p) = x_0 - y_0$ by
                \cref{fst_eq,snd_eq,real_sub,real_add_closed,real_add_inverse,real_subtraction_operation,function_apply_intro}.
        Show that for all $V$ such that $V$ is a neighborhood of $f(p)$ in $\reals$ with respect to $\standardTopologyR$
            there exists $U$ such that $U$ is a neighborhood of $p$ in $\reals\times\reals$ with respect to $T$ and $\img{f}{U}\subseteq V$.
        \begin{subproof}
            Fix $V$ such that $V$ is a neighborhood of $f(p)$ in $\reals$ with respect to $\standardTopologyR$.
            Let $d = \standardMetricR$.
            $\metricTopology{d}{\reals} = \standardTopologyR$ by \cref{standard_metric_topology}.
            $\standardBasisR$ is a basis on $\reals$ and $\standardTopologyR$ is a topology on $\reals$
                by \cref{standard_basis_is_basis,basis_topology_is_topology,standard_topology_reals}.
            Hence $V\subseteq\reals$ by
                \cref{neighborhood,open_in,standard_basis_is_basis,basis_topology_is_topology,standard_topology_reals,topology_on,subseteq,pow_iff}.
            $d$ is a metric on $\reals$ by \cref{standard_metric_is_metric}.
            Take $\epsilon\in\reals$ such that $0 < \epsilon$ and
                $\metricBall{d}{\reals}{f(p)}{\epsilon} \subseteq V$ by \cref{neighborhood,standard_metric_topology,metric_open_iff}.
            Let $\delta = \epsilon / (1 + 1)$.
            $1\in\reals$ and $0\in\reals$ and $0 < 1$
                by \cref{real_mul_ident,real_add_ident,real_zero_lt_one}.
            $0 + 1 < 1 + 1$ and $0 + 1 = 1$ and $1 + 1\in\reals$
                by \cref{real_lt_add,real_add_ident,real_add_closed}.
            Thus $0 < 1 + 1$ by \cref{real_lt_subst_left,real_lt_trans}.
            Thus $(1 + 1)\neq 0$ by \cref{real_pos_nonzero}.
            Thus $\inv{(1 + 1)}\in\reals$ by \cref{real_mul_inverse}.
            $\delta = \epsilon \rmul \inv{(1 + 1)}$ by \cref{real_div}.
            $0\rmul \inv{(1 + 1)} = 0$ by \cref{real_mul_zero}.
            $0 < \inv{(1 + 1)}$ by \cref{real_inv_pos}.
            Thus $0 < \delta$ by \cref{real_lt_mul_pos}.
            Let $U_1 = \metricBall{d}{\reals}{x_0}{\delta}$.
            Let $U_2 = \metricBall{d}{\reals}{y_0}{\delta}$.
            $x_0\in\reals$ and $\delta\in\reals$ and $0 < \delta$.
            Thus $U_1$ is open in $\reals$ with respect to $\standardTopologyR$
                by \cref{metric_ball,subseteq,pow_iff,metric_basis,metric_basis_is_basis,basis_elem_open_in_generated,metric_topology,basis_for_topology,open_in}.
            $y_0\in\reals$ and $\delta\in\reals$ and $0 < \delta$.
            Thus $U_2$ is open in $\reals$ with respect to $\standardTopologyR$
                by \cref{metric_ball,subseteq,pow_iff,metric_basis,metric_basis_is_basis,basis_elem_open_in_generated,metric_topology,basis_for_topology,open_in}.
            Let $U = U_1\times U_2$.
            Thus $U\in\productBasis{\standardTopologyR}{\standardTopologyR}{\reals}{\reals}$ by
                \cref{open_in,topology_on,subseteq,pow_iff,times_subseteq_left,times_subseteq_right,subseteq_transitive,product_basis}.
            Thus $U$ is open in $\reals\times\reals$ with respect to $T$ by
                \cref{product_basis_is_basis,basis_topology_is_topology,product_topology,basis_elem_open_in_generated,open_in}.
            Thus $p\in U$ by
                \cref{metric_zero_characterization,standard_metric_is_metric,metric_on,real_lt_subst_left,metric_ball,times_tuple_intro}.
            Show that for all $q$ such that $q\in U$ we have
                $f(q)\in\metricBall{d}{\reals}{f(p)}{\epsilon}$.
            \begin{subproof}
                    Fix $q$ such that $q\in U$.
                    Take $x,y$ such that $q = (x,y)$ and $x\in U_1$ and $y\in U_2$ by \cref{times_elem_is_tuple}.
                    $x\in\reals$ and $y\in\reals$ by \cref{metric_ball}.
                    Thus $\abs{x - x_0} < \delta$ and $\abs{y - y_0} < \delta$
                        by \cref{metric_ball,standard_metric_value,real_abs_sub_swap,real_lt_subst_left}.
                    $(q, x - y)\in f$ by \cref{real_sub,real_add_closed,real_add_inverse,times_tuple_intro,fst_eq,snd_eq,real_subtraction_operation}.
                    Thus $f(q) = x - y$ by \cref{function_apply_intro}.
                    $x - y = x + \neg{y}$ by \cref{real_sub}.
                    $x_0 - y_0 = x_0 + \neg{y_0}$ by \cref{real_sub}.
                    $(x - y) - (x_0 - y_0) = (x + \neg{y}) + \neg{(x_0 + \neg{y_0})}$ by \cref{real_sub}.
                    $\neg{y_0}\in\reals$ by \cref{real_add_inverse}.
                    $\neg{(x_0 + \neg{y_0})} = \neg{x_0} + \neg{(\neg{y_0})}$ by \cref{real_neg_addition}.
                    $\neg{(\neg{y_0})} = y_0$ by \cref{real_neg_neg}.
                    Thus $(x - y) - (x_0 - y_0) = (x + \neg{y}) + (\neg{x_0} + y_0)$.
                    $\neg{y}\in\reals$ by \cref{real_add_inverse}.
                    $\neg{x_0}\in\reals$ by \cref{real_add_inverse}.
                    $\neg{x_0} + y_0\in\reals$ by \cref{real_add_closed}.
                    $\neg{y} + \neg{x_0}\in\reals$ by \cref{real_add_closed}.
                    $\neg{y} + y_0\in\reals$ by \cref{real_add_closed}.
                    $(x + \neg{y}) + (\neg{x_0} + y_0) = x + (\neg{y} + (\neg{x_0} + y_0))$ by \cref{real_add_assoc}.
                    $\neg{y} + (\neg{x_0} + y_0) = (\neg{y} + \neg{x_0}) + y_0$ by \cref{real_add_assoc}.
                    $\neg{y} + \neg{x_0} = \neg{x_0} + \neg{y}$ by \cref{real_add_comm}.
                    Thus $\neg{y} + (\neg{x_0} + y_0) = (\neg{x_0} + \neg{y}) + y_0$.
                    $(\neg{x_0} + \neg{y}) + y_0 = \neg{x_0} + (\neg{y} + y_0)$ by \cref{real_add_assoc}.
                    Thus $(x - y) - (x_0 - y_0) = x + (\neg{x_0} + (\neg{y} + y_0))$ by \cref{real_add_assoc}.
                    $\neg{y} + y_0 = y_0 + \neg{y}$ by \cref{real_add_comm}.
                    Thus $(x - y) - (x_0 - y_0) = x + (\neg{x_0} + (y_0 + \neg{y}))$.
                    $y_0 + \neg{y} = y_0 - y$ by \cref{real_sub}.
                    $x + \neg{x_0} = x - x_0$ by \cref{real_sub}.
                    Thus $(x - y) - (x_0 - y_0) = (x - x_0) + (y_0 - y)$ by \cref{real_add_assoc}.
                    $x - x_0\in\reals$ and $y_0 - y\in\reals$ by \cref{real_sub,real_add_closed}.
                    $\abs{(x - x_0) + (y_0 - y)} \leq \abs{x - x_0} + \abs{y_0 - y}$ by \cref{real_triangle_inequality}.
                    $\abs{y_0 - y} = \abs{y - y_0}$ by \cref{real_abs_sub_swap}.
                    Thus $\abs{y_0 - y} < \delta$ by \cref{real_lt_subst_left}.
                    $\abs{x - x_0}\in\reals$ by \cref{real_abs_in_reals}.
                    $\abs{y_0 - y}\in\reals$ by \cref{real_abs_in_reals}.
                    $\abs{x - x_0} + \abs{y_0 - y} < \delta + \delta$ by \cref{real_lt_add_two}.
                    $\abs{x - x_0} + \abs{y_0 - y}\in\reals$ by \cref{real_add_closed}.
                    $\delta + \delta\in\reals$ by \cref{real_add_closed}.
                    $(x - x_0) + (y_0 - y)\in\reals$ by \cref{real_add_closed}.
                    $\abs{(x - x_0) + (y_0 - y)}\in\reals$ by \cref{real_abs_in_reals}.
                    Thus $\abs{(x - x_0) + (y_0 - y)} < \delta + \delta$
                        by \cref{real_leq_split,real_lt_trans,real_lt_subst_left}.
                    $1\rmul \delta = \delta$ by \cref{real_mul_ident}.
                    $(1 + 1) \rmul \delta = (1 \rmul \delta) + (1 \rmul \delta)$ by \cref{real_right_distrib}.
                    Thus $(1 \rmul \delta) + (1 \rmul \delta) = \delta + (1 \rmul \delta)$ by \cref{real_add_eq_subst}.
                    Thus $(1 \rmul \delta) + (1 \rmul \delta) = \delta + \delta$ by \cref{real_add_eq_subst}.
                    Hence $(1 + 1) \rmul \delta = \delta + \delta$.
                    Thus $(1 + 1) \rmul \delta = (1 + 1) \rmul (\epsilon \rmul \inv{(1 + 1)})$.
                    $(1 + 1) \rmul (\epsilon \rmul \inv{(1 + 1)}) = \epsilon \rmul ((1 + 1) \rmul \inv{(1 + 1)})$ by \cref{real_mul_assoc,real_mul_comm}.
                    $(1 + 1) \rmul \inv{(1 + 1)} = 1$ by \cref{real_mul_inverse}.
                    Thus $\epsilon \rmul 1 = \epsilon$ by \cref{real_mul_comm,real_mul_ident}.
                    Then $\delta + \delta = \epsilon$.
                    Therefore $\abs{(x - x_0) + (y_0 - y)} < \epsilon$ by \cref{real_lt_subst_right}.
                    Hence $\abs{f(q) - f(p)} = \abs{(x - y) - (x_0 - y_0)}$.
                    Therefore $\abs{f(q) - f(p)} < \epsilon$ by \cref{real_lt_subst_left}.
                $f$ is a function to $\reals$ and $\dom{f} = \reals\times\reals$ by \cref{funs_elim}.
                Thus $f(p)\in\reals$ and $f(q)\in\reals$.
                Hence $d((f(p),f(q))) < \epsilon$ by \cref{standard_metric_value,real_abs_sub_swap,real_lt_subst_left}.
                Thus $f(q)\in\metricBall{d}{\reals}{f(p)}{\epsilon}$ by \cref{metric_ball}.
            \end{subproof}
            Thus $\img{f}{U}\subseteq V$ by
                \cref{img_iff,function_apply_intro,subseteq,subseteq_transitive}.
            Thus $U$ is a neighborhood of $p$ in $\reals\times\reals$ with respect to $T$ by \cref{neighborhood}.
        \end{subproof}
        Therefore $f$ is continuous at $p$ from $\reals\times\reals$ to $\reals$ with respect to $T$ and $\standardTopologyR$ by
            \cref{standard_basis_is_basis,basis_topology_is_topology,standard_topology_reals,product_basis_is_basis,product_topology,continuous_at}.
    \end{subproof}
    Hence $f$ is continuous from $\reals\times\reals$ to $\reals$ with respect to $T$ and $\standardTopologyR$ by
        \cref{standard_basis_is_basis,basis_topology_is_topology,standard_topology_reals,product_basis_is_basis,product_topology,continuous_at_implies_continuous}.
\end{proof}

% from §21 helper lemma 4: absolute value bound by nearby point
\begin{lemma}\label{real_abs_sub_lt_bound}
    Suppose $x\in\reals$ and $y\in\reals$ and $\delta\in\reals$.
    Suppose $\abs{x - y} < \delta$.
    Then $\abs{x} < \abs{y} + \delta$.
\end{lemma}
\begin{proof}
    $x - y = x + \neg{y}$ and $\neg{y}\in\reals$ and $x - y\in\reals$ by
        \cref{real_sub,real_add_inverse,real_add_closed}.
    Then $(x - y) + y = x$ by \cref{real_add_assoc,real_add_inverse,real_add_comm,real_add_ident,real_add_eq_subst}.
    Hence $\abs{x} \leq \abs{x - y} + \abs{y}$ by
        \cref{real_sub,real_add_inverse,real_add_closed,real_triangle_inequality,real_lt_subst_left}.
    Therefore $\abs{x - y} + \abs{y}\in\reals$ and $\delta + \abs{y}\in\reals$ and
        $\abs{x - y} + \abs{y} < \delta + \abs{y}$ by \cref{real_abs_in_reals,real_add_closed,real_lt_add}.
    Hence $\abs{x} < \abs{y} + \delta$ by
        \cref{real_abs_in_reals,real_leq_split,real_lt_trans,real_lt_subst_left,real_add_comm,real_lt_subst_right}.
\end{proof}
% from §21 helper lemma 5: positive sum with nonnegative term
\begin{lemma}\label{real_pos_add_nonneg}
    Suppose $a\in\reals$ and $b\in\reals$.
    Suppose $a \geq 0$ and $0 < b$.
    Then $0 < a + b$.
\end{lemma}
\begin{proof}
    Therefore $0 < a + b$ by
        \cref{real_add_closed,real_add_ident,real_leq_split,real_add_eq_subst,real_lt_subst_right,real_lt_add,real_lt_subst_left,real_lt_trans}.
\end{proof}

% from §21 helper lemma 6: substitute equality in multiplication
\begin{lemma}\label{real_mul_eq_subst}
    Suppose $a\in\reals$ and $b\in\reals$ and $c\in\reals$.
    Suppose $a = b$.
    Then $a \rmul c = b \rmul c$.
\end{lemma}
\begin{proof}
    Trivial.
\end{proof}

% from §21 Item 10: multiplication operation is continuous
\begin{lemma}\label{real_multiplication_continuous}
    Let $T = \productTopology{\standardTopologyR}{\standardTopologyR}{\reals}{\reals}$.
    Then $\realMultiplication$ is continuous from $\reals\times\reals$ to $\reals$ with respect to $T$ and $\standardTopologyR$.
\end{lemma}
\begin{proof}
    Let $f = \realMultiplication$.
    For all $p,z_1,z_2$ such that $(p,z_1)\in f$ and $(p,z_2)\in f$ we have $z_1 = z_2$
        by \cref{real_multiplication_operation,pair_eq_iff}.
    Hence $f$ is a function by \cref{rightunique,real_multiplication_operation,relation}.
    For all $p\in\dom{f}$ we have $f(p)\in\reals$
        by \cref{dom_iff,real_multiplication_operation,pair_eq_iff,function_apply_intro}.
    Therefore $f$ is a function to $\reals$.
    Hence $\dom{f}\subseteq\reals\times\reals$ by \cref{dom_iff,real_multiplication_operation,pair_eq_iff,subseteq}.
    For all $p\in\reals\times\reals$ we have $p\in\dom{f}$
        by \cref{times_elem_is_tuple,real_mul_closed,fst_eq,snd_eq,real_multiplication_operation,dom_intro}.
    Thus $\reals\times\reals\subseteq\dom{f}$ by \cref{subseteq}.
    Hence $\dom{f} = \reals\times\reals$ by \cref{subseteq_antisymmetric}.
    Therefore $f\in\funs{\reals\times\reals}{\reals}$ by \cref{funs_intro}.
    Show that for all $p\in\reals\times\reals$ we have
        $f$ is continuous at $p$ from $\reals\times\reals$ to $\reals$ with respect to $T$ and $\standardTopologyR$.
    \begin{subproof}
        Fix $p$ such that $p\in\reals\times\reals$.
        Take $x_0,y_0$ such that $p = (x_0,y_0)$ and $x_0\in\reals$ and $y_0\in\reals$ by \cref{times_elem_is_tuple}.
        $\fst{p} = x_0$ and $\snd{p} = y_0$ by \cref{fst_eq,snd_eq}.
        $(p,x_0 \rmul y_0)\in f$ by \cref{real_mul_closed,real_multiplication_operation}.
        Hence $f(p) = x_0 \rmul y_0$ by \cref{function_apply_intro}.
        Show that for all $V$ such that $V$ is a neighborhood of $f(p)$ in $\reals$ with respect to $\standardTopologyR$
            there exists $U$ such that $U$ is a neighborhood of $p$ in $\reals\times\reals$ with respect to $T$ and $\img{f}{U}\subseteq V$.
        \begin{subproof}
            Fix $V$ such that $V$ is a neighborhood of $f(p)$ in $\reals$ with respect to $\standardTopologyR$.
            $V$ is open in $\reals$ with respect to $\standardTopologyR$ by \cref{neighborhood}.
            Let $d = \standardMetricR$.
            $\metricTopology{d}{\reals} = \standardTopologyR$ by \cref{standard_metric_topology}.
            Therefore $V\in\metricTopology{d}{\reals}$ by \cref{standard_metric_topology,open_in,subseteq}.
            $f(p)\in V$ by \cref{neighborhood}.
            $\standardBasisR$ is a basis on $\reals$ and $\standardTopologyR$ is a topology on $\reals$
                by \cref{standard_basis_is_basis,basis_topology_is_topology,standard_topology_reals}.
            Hence $V\subseteq\reals$ by
                \cref{open_in,neighborhood,standard_basis_is_basis,basis_topology_is_topology,standard_topology_reals,topology_on,subseteq,pow_iff}.
            $d$ is a metric on $\reals$ by \cref{standard_metric_is_metric}.
            Take $\epsilon\in\reals$ such that $0 < \epsilon$ and
                $\metricBall{d}{\reals}{f(p)}{\epsilon} \subseteq V$ by \cref{metric_open_iff}.
            Let $m = \abs{x_0} + \abs{y_0} + 1$.
            $\abs{x_0}\in\reals$ by \cref{real_abs_in_reals}.
            $\abs{y_0}\in\reals$ by \cref{real_abs_in_reals}.
            $\abs{x_0} + \abs{y_0}\in\reals$ by \cref{real_add_closed}.
            $1\in\reals$ by \cref{real_mul_ident}.
            Therefore $m\in\reals$ by \cref{real_add_closed}.
            $\abs{y_0} \geq 0$ by \cref{real_abs_nonnegative}.
            $0 < 1$ by \cref{real_zero_lt_one}.
            $\abs{y_0} + 1\in\reals$ by \cref{real_add_closed}.
            Hence $0 < \abs{y_0} + 1$ by \cref{real_pos_add_nonneg}.
            $\abs{x_0} \geq 0$ by \cref{real_abs_nonnegative}.
            Therefore $0 < \abs{x_0} + (\abs{y_0} + 1)$ by \cref{real_pos_add_nonneg}.
            Hence $0 < m$ by \cref{real_add_assoc,real_lt_subst_right}.
            Therefore $m \geq 0$ by \cref{real_lt_imp_leq}.
            Let $\delta = \epsilon / (m + \epsilon)$.
            $m + \epsilon\in\reals$ by \cref{real_add_closed}.
            $0 < m + \epsilon$ by \cref{real_pos_add_nonneg}.
            Thus $(m + \epsilon)\neq 0$ by \cref{real_pos_nonzero}.
            Thus $\inv{(m + \epsilon)}\in\reals$ by \cref{real_mul_inverse}.
            $0 < \inv{(m + \epsilon)}$ by \cref{real_inv_pos}.
            $0\in\reals$ and $0 \rmul \inv{(m + \epsilon)} = 0$ by \cref{real_add_ident,real_mul_zero}.
            Thus $0 \rmul \inv{(m + \epsilon)} < \epsilon \rmul \inv{(m + \epsilon)}$ by \cref{real_lt_mul_pos}.
            Thus $0 < \epsilon \rmul \inv{(m + \epsilon)}$ by \cref{real_lt_subst_left}.
            $\delta = \epsilon \rmul \inv{(m + \epsilon)}$ by \cref{real_div}.
            Thus $0 < \delta$ by \cref{real_lt_subst_left}.
            Thus $\epsilon < m + \epsilon$ by \cref{real_lt_add,real_add_ident,real_lt_subst_left}.
            $\epsilon \rmul \inv{(m + \epsilon)} < (m + \epsilon) \rmul \inv{(m + \epsilon)}$ by \cref{real_lt_mul_pos}.
            $(m + \epsilon) \rmul \inv{(m + \epsilon)} = 1$ by \cref{real_mul_inverse}.
            Thus $\epsilon \rmul \inv{(m + \epsilon)} < 1$ by \cref{real_lt_subst_right}.
            Thus $\delta < 1$ by \cref{real_lt_subst_left}.
            Let $U_1 = \metricBall{d}{\reals}{x_0}{\delta}$.
            Let $U_2 = \metricBall{d}{\reals}{y_0}{\delta}$.
            $x_0\in\reals$ and $y_0\in\reals$ and $\delta\in\reals$ and $0 < \delta$.
            Thus $U_1\in\metricBasis{d}{\reals}$ by \cref{metric_ball,subseteq,pow_iff,metric_basis}.
            Thus $U_1$ is open in $\reals$ with respect to $\standardTopologyR$
                by \cref{metric_basis_is_basis,basis_elem_open_in_generated,metric_topology,basis_for_topology,subseteq,open_in}.
            Thus $U_2\in\metricBasis{d}{\reals}$ by \cref{metric_ball,subseteq,pow_iff,metric_basis}.
            Thus $U_2$ is open in $\reals$ with respect to $\standardTopologyR$
                by \cref{metric_basis_is_basis,basis_elem_open_in_generated,metric_topology,basis_for_topology,subseteq,open_in}.
            Let $U = U_1\times U_2$.
            $U_1\in\standardTopologyR$ and $U_2\in\standardTopologyR$ by \cref{open_in}.
            Thus $U_1\subseteq\reals$ and $U_2\subseteq\reals$ by \cref{open_in,topology_on,subseteq,pow_iff}.
            Thus $U\in\pow{\reals\times\reals}$ by
                \cref{times_subseteq_left,times_subseteq_right,subseteq_transitive,pow_iff}.
            Thus $U\in\productBasis{\standardTopologyR}{\standardTopologyR}{\reals}{\reals}$ by \cref{product_basis}.
            $\productBasis{\standardTopologyR}{\standardTopologyR}{\reals}{\reals}$ is a basis on $\reals\times\reals$ and
                $T$ is a topology on $\reals\times\reals$ by
                \cref{product_basis_is_basis,basis_topology_is_topology,product_topology}.
            Thus $U$ is open in $\reals\times\reals$ with respect to $T$ by
                \cref{basis_elem_open_in_generated,product_topology,open_in}.
            $d((x_0,x_0)) = 0$ and $d((y_0,y_0)) = 0$ by \cref{metric_zero_characterization,standard_metric_is_metric,metric_on}.
            Thus $d((x_0,x_0)) < \delta$ and $d((y_0,y_0)) < \delta$ by \cref{real_lt_subst_left}.
            Thus $p\in U$ by \cref{metric_ball,times_tuple_intro}.
            Show that for all $q$ such that $q\in U$ we have
                $f(q)\in\metricBall{d}{\reals}{f(p)}{\epsilon}$.
            \begin{subproof}
                Fix $q$ such that $q\in U$.
                Take $x,y$ such that $q = (x,y)$ and $x\in U_1$ and $y\in U_2$ by \cref{times_elem_is_tuple}.
                $x\in\reals$ and $y\in\reals$ by \cref{metric_ball}.
                Thus $\abs{x - x_0} < \delta$ and $\abs{y - y_0} < \delta$
                    by \cref{metric_ball,standard_metric_value,real_abs_sub_swap,real_lt_subst_left}.
                $(q, x \rmul y)\in f$ by \cref{real_mul_closed,times_tuple_intro,fst_eq,snd_eq,real_multiplication_operation}.
                Thus $f(q) = x \rmul y$ by \cref{function_apply_intro}.
                Show that $\abs{f(q) - f(p)} < \epsilon$.
                \begin{subproof}
                        $x - x_0 = x + \neg{x_0}$ by \cref{real_sub}.
                        $y - y_0 = y + \neg{y_0}$ by \cref{real_sub}.
                        $\neg{x_0}\in\reals$ by \cref{real_add_inverse}.
                        $\neg{y_0}\in\reals$ by \cref{real_add_inverse}.
                        Thus $(x - x_0) \rmul y = (x + \neg{x_0}) \rmul y$.
                        Thus $x_0 \rmul (y - y_0) = x_0 \rmul (y + \neg{y_0})$.
                        $y + \neg{y_0}\in\reals$ by \cref{real_add_closed}.
                        $(x + \neg{x_0}) \rmul y = (x \rmul y) + (\neg{x_0} \rmul y)$ by \cref{real_right_distrib}.
                        $x_0 \rmul (y + \neg{y_0}) = (y + \neg{y_0}) \rmul x_0$ by \cref{real_mul_comm}.
                        $(y + \neg{y_0}) \rmul x_0 = (y \rmul x_0) + (\neg{y_0} \rmul x_0)$ by \cref{real_right_distrib}.
                        $y \rmul x_0 = x_0 \rmul y$ by \cref{real_mul_comm}.
                        $\neg{y_0} \rmul x_0 = x_0 \rmul \neg{y_0}$ by \cref{real_mul_comm}.
                        $x_0 \rmul \neg{y_0} = \neg{(x_0 \rmul y_0)}$ by \cref{real_mul_neg_right}.
                        Thus $(x - x_0) \rmul y + x_0 \rmul (y - y_0) = (x \rmul y) + (\neg{x_0} \rmul y) + (x_0 \rmul y) + \neg{(x_0 \rmul y_0)}$.
                        $(\neg{x_0} \rmul y) + (x_0 \rmul y) = (\neg{x_0} + x_0) \rmul y$ by \cref{real_right_distrib}.
                        $\neg{x_0} + x_0 = 0$ by \cref{real_add_inverse,real_add_comm}.
                        $0 \rmul y = 0$ by \cref{real_mul_zero}.
                        Thus $(\neg{x_0} \rmul y) + (x_0 \rmul y) = 0$.
                        $0\in\reals$ by \cref{real_add_ident}.
                        $x \rmul y\in\reals$ by \cref{real_mul_closed}.
                        $\neg{x_0} \rmul y\in\reals$ by \cref{real_mul_closed}.
                        $x_0 \rmul y\in\reals$ by \cref{real_mul_closed}.
                        $(\neg{x_0} \rmul y) + (x_0 \rmul y)\in\reals$ by \cref{real_add_closed}.
                        Thus $(x \rmul y) + (\neg{x_0} \rmul y) + (x_0 \rmul y) = x \rmul y$
                            by \cref{real_add_assoc,real_add_comm,real_add_eq_subst,real_add_ident}.
                        $(x \rmul y) + (\neg{x_0} \rmul y)\in\reals$ by \cref{real_add_closed}.
                        $(x \rmul y) + (\neg{x_0} \rmul y) + (x_0 \rmul y)\in\reals$ by \cref{real_add_closed}.
                        $\neg{(x_0 \rmul y_0)}\in\reals$ by \cref{real_add_inverse}.
                        Thus $(x \rmul y) + (\neg{x_0} \rmul y) + (x_0 \rmul y) + \neg{(x_0 \rmul y_0)} = (x \rmul y) + \neg{(x_0 \rmul y_0)}$ by \cref{real_add_eq_subst}.
                        Thus $(x - x_0) \rmul y + x_0 \rmul (y - y_0) = (x \rmul y) + \neg{(x_0 \rmul y_0)}$.
                        $(x \rmul y) + \neg{(x_0 \rmul y_0)} = (x \rmul y) - (x_0 \rmul y_0)$ by \cref{real_sub}.
                        Thus $(x \rmul y) - (x_0 \rmul y_0) = (x - x_0) \rmul y + x_0 \rmul (y - y_0)$.
                        $x - x_0\in\reals$ by \cref{real_sub,real_add_closed}.
                        $y - y_0\in\reals$ by \cref{real_sub,real_add_closed}.
                        $(x - x_0) \rmul y\in\reals$ by \cref{real_mul_closed}.
                        $x_0 \rmul (y - y_0)\in\reals$ by \cref{real_mul_closed}.
                        $\abs{(x - x_0) \rmul y + x_0 \rmul (y - y_0)} \leq \abs{(x - x_0) \rmul y} + \abs{x_0 \rmul (y - y_0)}$ by \cref{real_triangle_inequality}.
                        $\abs{(x - x_0) \rmul y} = \abs{x - x_0} \rmul \abs{y}$ by \cref{real_abs_mul}.
                        $\abs{x_0 \rmul (y - y_0)} = \abs{x_0} \rmul \abs{y - y_0}$ by \cref{real_abs_mul}.
                        Thus $\abs{(x - x_0) \rmul y + x_0 \rmul (y - y_0)} \leq \abs{x - x_0} \rmul \abs{y} + \abs{x_0} \rmul \abs{y - y_0}$.
                        $\delta\in\reals$ by \cref{real_mul_closed}.
                        $\abs{y} < \abs{y_0} + \delta$ by \cref{real_abs_sub_lt_bound}.
                        Thus $0 < \abs{y_0} + \delta$ by \cref{real_abs_nonnegative,real_pos_add_nonneg}.
                        Thus $0 < \abs{x - x_0}$ or $0 = \abs{x - x_0}$ by \cref{real_abs_nonnegative,real_leq_split}.
                        Show that $\abs{x - x_0} \rmul \abs{y} < \delta \rmul (\abs{y_0} + \delta)$.
                        \begin{subproof}
                            \begin{byCase}
                            \caseOf{$0 = \abs{x - x_0}$.}
                                Thus $\abs{x - x_0} = 0$ by assumption.
                                $\abs{x - x_0}\in\reals$ by \cref{real_abs_in_reals}.
                                $\abs{y}\in\reals$ by \cref{real_abs_in_reals}.
                                $0\in\reals$ by \cref{real_add_ident}.
                                $\abs{x - x_0} \rmul \abs{y} = 0 \rmul \abs{y}$ by \cref{real_mul_eq_subst}.
                                $0 \rmul \abs{y} = 0$ by \cref{real_mul_zero}.
                                Hence $\abs{x - x_0} \rmul \abs{y} = 0$.
                                $\delta$ is positive.
                                $\abs{y_0} + \delta$ is positive.
                                $\abs{y_0}\in\reals$ by \cref{real_abs_in_reals}.
                                $\abs{y_0} + \delta\in\reals$ by \cref{real_add_closed}.
                                Therefore $\delta \rmul (\abs{y_0} + \delta)$ is positive by \cref{real_mul_pos_pos}.
                                Then $0 < \delta \rmul (\abs{y_0} + \delta)$.
                                Hence $\abs{x - x_0} \rmul \abs{y} < \delta \rmul (\abs{y_0} + \delta)$ by \cref{real_lt_subst_left}.
                            \caseOf{$0 < \abs{x - x_0}$.}
                                $\abs{y}\in\reals$ by \cref{real_abs_in_reals}.
                                $\abs{y_0}\in\reals$ by \cref{real_abs_in_reals}.
                                $\abs{x - x_0}\in\reals$ by \cref{real_abs_in_reals}.
                                $\abs{y_0} + \delta\in\reals$ by \cref{real_add_closed}.
                                $\abs{x - x_0} \rmul \abs{y}\in\reals$ by \cref{real_mul_closed}.
                                $\abs{x - x_0} \rmul (\abs{y_0} + \delta)\in\reals$ by \cref{real_mul_closed}.
                                $\delta \rmul (\abs{y_0} + \delta)\in\reals$ by \cref{real_mul_closed}.
                                $\abs{y} \rmul \abs{x - x_0} < (\abs{y_0} + \delta) \rmul \abs{x - x_0}$ by \cref{real_lt_mul_pos}.
                                $\abs{y} \rmul \abs{x - x_0} = \abs{x - x_0} \rmul \abs{y}$ by \cref{real_mul_comm}.
                                $(\abs{y_0} + \delta) \rmul \abs{x - x_0} = \abs{x - x_0} \rmul (\abs{y_0} + \delta)$ by \cref{real_mul_comm}.
                                Hence $\abs{x - x_0} \rmul \abs{y} < \abs{x - x_0} \rmul (\abs{y_0} + \delta)$.
                                $\abs{x - x_0} \rmul (\abs{y_0} + \delta) < \delta \rmul (\abs{y_0} + \delta)$ by \cref{real_lt_mul_pos}.
                                Therefore $\abs{x - x_0} \rmul \abs{y} < \delta \rmul (\abs{y_0} + \delta)$ by \cref{real_lt_trans}.
                            \end{byCase}
                        \end{subproof}
                        Thus $0 < \abs{x_0}$ or $0 = \abs{x_0}$ by \cref{real_abs_nonnegative,real_leq_split}.
                        Show that $\abs{x - x_0} \rmul \abs{y} + \abs{x_0} \rmul \abs{y - y_0} < \delta \rmul m$.
                        \begin{subproof}
                            \begin{byCase}
                            \caseOf{$0 = \abs{x_0}$.}
                                Thus $\abs{x_0} = 0$ by assumption.
                                $\abs{x_0}\in\reals$ by \cref{real_abs_in_reals}.
                                $\abs{y - y_0}\in\reals$ by \cref{real_abs_in_reals}.
                                $0\in\reals$ by \cref{real_add_ident}.
                                $\abs{x_0} \rmul \abs{y - y_0} = 0 \rmul \abs{y - y_0}$ by \cref{real_mul_eq_subst}.
                                $0 \rmul \abs{y - y_0} = 0$ by \cref{real_mul_zero}.
                                Thus $\abs{x_0} \rmul \abs{y - y_0} = 0$.
                                $\abs{x - x_0}\in\reals$ by \cref{real_abs_in_reals}.
                                $\abs{y}\in\reals$ by \cref{real_abs_in_reals}.
                                $\abs{x - x_0} \rmul \abs{y}\in\reals$ by \cref{real_mul_closed}.
                                Thus $\abs{x - x_0} \rmul \abs{y} + \abs{x_0} \rmul \abs{y - y_0} = \abs{x - x_0} \rmul \abs{y} + 0$ by \cref{real_add_eq_subst}.
                                $\abs{x - x_0} \rmul \abs{y} + 0 = 0 + \abs{x - x_0} \rmul \abs{y}$ by \cref{real_add_comm}.
                                $0 + \abs{x - x_0} \rmul \abs{y} = \abs{x - x_0} \rmul \abs{y}$ by \cref{real_add_ident}.
                                Thus $\abs{x - x_0} \rmul \abs{y} + \abs{x_0} \rmul \abs{y - y_0} = \abs{x - x_0} \rmul \abs{y}$.
                                Thus $\abs{x - x_0} \rmul \abs{y} < \delta \rmul (\abs{y_0} + \delta)$.
                                $\delta + \abs{y_0} < 1 + \abs{y_0}$ by \cref{real_lt_add}.
                                Thus $\abs{y_0} + \delta < \abs{y_0} + 1$ by \cref{real_add_comm,real_add_assoc}.
                                $\abs{y_0} + \delta\in\reals$ by \cref{real_add_closed}.
                                $\abs{y_0} + 1\in\reals$ by \cref{real_add_closed}.
                                $\delta\in\reals$ by \cref{real_mul_closed}.
                                $\delta$ is positive.
                                $\delta \rmul (\abs{y_0} + \delta) = (\abs{y_0} + \delta) \rmul \delta$ by \cref{real_mul_comm}.
                                $\delta \rmul (\abs{y_0} + 1) = (\abs{y_0} + 1) \rmul \delta$ by \cref{real_mul_comm}.
                                $(\abs{y_0} + \delta) \rmul \delta < (\abs{y_0} + 1) \rmul \delta$ by \cref{real_lt_mul_pos}.
                                Thus $\delta \rmul (\abs{y_0} + \delta) < (\abs{y_0} + 1) \rmul \delta$ by \cref{real_lt_subst_left}.
                                Thus $\delta \rmul (\abs{y_0} + \delta) < \delta \rmul (\abs{y_0} + 1)$ by \cref{real_lt_subst_right}.
                                $\delta \rmul (\abs{y_0} + \delta)\in\reals$ by \cref{real_mul_closed}.
                                $\delta \rmul (\abs{y_0} + 1)\in\reals$ by \cref{real_mul_closed}.
                                Thus $\abs{x - x_0} \rmul \abs{y} < \delta \rmul (\abs{y_0} + 1)$ by \cref{real_lt_trans}.
                                Thus $\abs{x_0} = 0$.
                                Thus $\abs{x_0} + \abs{y_0} + 1 = \abs{y_0} + 1$ by \cref{real_add_ident,real_add_comm,real_add_assoc,real_add_eq_subst}.
                                $\delta \rmul (\abs{y_0} + 1) = (\abs{y_0} + 1) \rmul \delta$ by \cref{real_mul_comm}.
                                $(\abs{y_0} + 1) \rmul \delta = (\abs{x_0} + \abs{y_0} + 1) \rmul \delta$ by \cref{real_mul_eq_subst}.
                                Thus $\delta \rmul (\abs{y_0} + 1) = \delta \rmul (\abs{x_0} + \abs{y_0} + 1)$ by \cref{real_mul_comm}.
                                Thus $\abs{x - x_0} \rmul \abs{y} + \abs{x_0} \rmul \abs{y - y_0} < \delta \rmul (\abs{x_0} + \abs{y_0} + 1)$ by \cref{real_lt_subst_right}.
                                Thus $\abs{x - x_0} \rmul \abs{y} + \abs{x_0} \rmul \abs{y - y_0} < \delta \rmul m$ by \cref{real_lt_subst_right}.
                            \caseOf{$0 < \abs{x_0}$.}
                                $\abs{y - y_0}\in\reals$ by \cref{real_abs_in_reals}.
                                $\abs{x_0}\in\reals$ by \cref{real_abs_in_reals}.
                                $\delta\in\reals$ by \cref{real_mul_closed}.
                                $\abs{y - y_0} \rmul \abs{x_0} < \delta \rmul \abs{x_0}$ by \cref{real_lt_mul_pos}.
                                $\abs{y - y_0} \rmul \abs{x_0} = \abs{x_0} \rmul \abs{y - y_0}$ by \cref{real_mul_comm}.
                                $\delta \rmul \abs{x_0} = \abs{x_0} \rmul \delta$ by \cref{real_mul_comm}.
                                Thus $\abs{x_0} \rmul \abs{y - y_0} < \delta \rmul \abs{x_0}$ by \cref{real_lt_subst_left}.
                                Thus $\abs{x_0} \rmul \abs{y - y_0} < \abs{x_0} \rmul \delta$ by \cref{real_lt_subst_right}.
                                $\abs{x - x_0}\in\reals$ by \cref{real_abs_in_reals}.
                                $\abs{y}\in\reals$ by \cref{real_abs_in_reals}.
                                $\abs{x - x_0} \rmul \abs{y}\in\reals$ by \cref{real_mul_closed}.
                                $\abs{y_0}\in\reals$ by \cref{real_abs_in_reals}.
                                $\abs{y_0} + \delta\in\reals$ by \cref{real_add_closed}.
                                $\delta \rmul (\abs{y_0} + \delta)\in\reals$ by \cref{real_mul_closed}.
                                $\abs{y - y_0}\in\reals$ by \cref{real_abs_in_reals}.
                                $\abs{x_0}\in\reals$ by \cref{real_abs_in_reals}.
                                $\abs{x_0} \rmul \abs{y - y_0}\in\reals$ by \cref{real_mul_closed}.
                                $\abs{x_0} \rmul \delta\in\reals$ by \cref{real_mul_closed}.
                                Thus $\abs{x - x_0} \rmul \abs{y} + \abs{x_0} \rmul \abs{y - y_0} < \delta \rmul (\abs{y_0} + \delta) + \abs{x_0} \rmul \delta$ by \cref{real_lt_add_two}.
                                $(\abs{x_0} + \abs{y_0}) + \delta = \abs{x_0} + (\abs{y_0} + \delta)$ by \cref{real_add_assoc}.
                                $\delta \rmul (\abs{y_0} + \delta) = (\abs{y_0} + \delta) \rmul \delta$ by \cref{real_mul_comm}.
                                $(\abs{y_0} + \delta) \rmul \delta = (\abs{y_0} \rmul \delta) + (\delta \rmul \delta)$ by \cref{real_right_distrib}.
                                $\abs{x_0} \rmul \delta = \delta \rmul \abs{x_0}$ by \cref{real_mul_comm}.
                                $\abs{y_0} \rmul \delta = \delta \rmul \abs{y_0}$ by \cref{real_mul_comm}.
                                Thus $\delta \rmul (\abs{y_0} + \delta) + \abs{x_0} \rmul \delta = (\delta \rmul \abs{x_0}) + (\delta \rmul \abs{y_0}) + (\delta \rmul \delta)$
                                    by \cref{real_add_eq_subst,real_right_distrib,real_mul_comm,real_add_assoc,real_add_comm}.
                                Thus $(\delta \rmul \abs{x_0}) + (\delta \rmul \abs{y_0}) = \delta \rmul (\abs{x_0} + \abs{y_0})$
                                    by \cref{real_mul_comm,real_right_distrib,real_add_eq_subst}.
                                Thus $\delta \rmul (\abs{y_0} + \delta) + \abs{x_0} \rmul \delta = (\delta \rmul (\abs{x_0} + \abs{y_0})) + (\delta \rmul \delta)$ by \cref{real_add_assoc}.
                                Show that $(\delta \rmul (\abs{x_0} + \abs{y_0})) + (\delta \rmul \delta) = \delta \rmul ((\abs{x_0} + \abs{y_0}) + \delta)$.
                                \begin{subproof}
                                    $\abs{x_0}\in\reals$ by \cref{real_abs_in_reals}.
                                    $\abs{y_0}\in\reals$ by \cref{real_abs_in_reals}.
                                    $\abs{x_0} + \abs{y_0}\in\reals$ by \cref{real_add_closed}.
                                    $\delta\in\reals$ by \cref{real_mul_closed}.
                                    $(\abs{x_0} + \abs{y_0}) + \delta\in\reals$ by \cref{real_add_closed}.
                                    $\delta \rmul ((\abs{x_0} + \abs{y_0}) + \delta) = ((\abs{x_0} + \abs{y_0}) + \delta) \rmul \delta$ by \cref{real_mul_comm}.
                                    $((\abs{x_0} + \abs{y_0}) + \delta) \rmul \delta = (\abs{x_0} + \abs{y_0}) \rmul \delta + \delta \rmul \delta$ by \cref{real_right_distrib}.
                                    $(\abs{x_0} + \abs{y_0}) \rmul \delta = \delta \rmul (\abs{x_0} + \abs{y_0})$ by \cref{real_mul_comm}.
                                    Thus $\delta \rmul ((\abs{x_0} + \abs{y_0}) + \delta) = \delta \rmul (\abs{x_0} + \abs{y_0}) + \delta \rmul \delta$ by \cref{real_add_eq_subst}.
                                    Thus $(\delta \rmul (\abs{x_0} + \abs{y_0})) + (\delta \rmul \delta) = \delta \rmul ((\abs{x_0} + \abs{y_0}) + \delta)$.
                                \end{subproof}
                                Thus $\delta \rmul (\abs{y_0} + \delta) + \abs{x_0} \rmul \delta = \delta \rmul ((\abs{x_0} + \abs{y_0}) + \delta)$.
                                $\abs{x_0} + \abs{y_0} = \abs{y_0} + \abs{x_0}$ by \cref{real_add_comm}.
                                Thus $(\abs{x_0} + \abs{y_0}) + \delta = (\abs{y_0} + \abs{x_0}) + \delta$ by \cref{real_add_eq_subst}.
                                Thus $\delta \rmul ((\abs{x_0} + \abs{y_0}) + \delta) = \delta \rmul (\abs{x_0} + \abs{y_0} + \delta)$ by \cref{real_add_assoc}.
                                Thus $\abs{x - x_0} \rmul \abs{y} + \abs{x_0} \rmul \abs{y - y_0} < \delta \rmul (\abs{x_0} + \abs{y_0} + \delta)$ by \cref{real_lt_subst_right}.
                                $\delta + (\abs{x_0} + \abs{y_0}) < 1 + (\abs{x_0} + \abs{y_0})$ by \cref{real_lt_add}.
                                Thus $\abs{x_0} + \abs{y_0} + \delta < \abs{x_0} + \abs{y_0} + 1$ by \cref{real_add_comm,real_add_assoc}.
                                $\abs{x_0}\in\reals$ by \cref{real_abs_in_reals}.
                                $\abs{y_0}\in\reals$ by \cref{real_abs_in_reals}.
                                $\abs{x_0} + \abs{y_0}\in\reals$ by \cref{real_add_closed}.
                                $\abs{x_0} + \abs{y_0} + \delta\in\reals$ by \cref{real_add_closed}.
                                $1\in\reals$ by \cref{real_mul_ident}.
                                $\abs{x_0} + \abs{y_0} + 1\in\reals$ by \cref{real_add_closed}.
                                $(\abs{x_0} + \abs{y_0} + \delta) \rmul \delta < (\abs{x_0} + \abs{y_0} + 1) \rmul \delta$ by \cref{real_lt_mul_pos}.
                                $(\abs{x_0} + \abs{y_0} + \delta) \rmul \delta = \delta \rmul (\abs{x_0} + \abs{y_0} + \delta)$ by \cref{real_mul_comm}.
                                Thus $\delta \rmul (\abs{x_0} + \abs{y_0} + \delta) < (\abs{x_0} + \abs{y_0} + 1) \rmul \delta$ by \cref{real_lt_subst_left}.
                                $(\abs{x_0} + \abs{y_0} + 1) \rmul \delta = \delta \rmul (\abs{x_0} + \abs{y_0} + 1)$ by \cref{real_mul_comm}.
                                Thus $\delta \rmul (\abs{x_0} + \abs{y_0} + \delta) < \delta \rmul (\abs{x_0} + \abs{y_0} + 1)$ by \cref{real_lt_subst_right}.
                                $\abs{x - x_0} \rmul \abs{y}\in\reals$ by \cref{real_mul_closed}.
                                $\abs{x_0} \rmul \abs{y - y_0}\in\reals$ by \cref{real_mul_closed}.
                                $\abs{x - x_0} \rmul \abs{y} + \abs{x_0} \rmul \abs{y - y_0}\in\reals$ by \cref{real_add_closed}.
                                $\delta \rmul (\abs{x_0} + \abs{y_0} + \delta)\in\reals$ by \cref{real_mul_closed}.
                                $\delta \rmul (\abs{x_0} + \abs{y_0} + 1)\in\reals$ by \cref{real_mul_closed}.
                                Thus $\abs{x - x_0} \rmul \abs{y} + \abs{x_0} \rmul \abs{y - y_0} < \delta \rmul (\abs{x_0} + \abs{y_0} + 1)$ by \cref{real_lt_trans}.
                                Thus $\abs{x - x_0} \rmul \abs{y} + \abs{x_0} \rmul \abs{y - y_0} < \delta \rmul m$ by \cref{real_lt_subst_right}.
                            \end{byCase}
                        \end{subproof}
                        Therefore $m < \epsilon + m$ by \cref{real_lt_add,real_add_ident,real_lt_subst_left}.
                        Then $m < m + \epsilon$ by \cref{real_add_comm,real_lt_subst_right}.
                        $m \rmul \inv{(m + \epsilon)} < (m + \epsilon) \rmul \inv{(m + \epsilon)}$ by \cref{real_lt_mul_pos}.
                        $(m + \epsilon) \rmul \inv{(m + \epsilon)} = 1$ by \cref{real_mul_inverse}.
                        Thus $m \rmul \inv{(m + \epsilon)} < 1$ by \cref{real_lt_subst_right}.
                        $m \rmul \inv{(m + \epsilon)}\in\reals$ by \cref{real_mul_closed}.
                        $1\in\reals$ by \cref{real_mul_ident}.
                        $(m \rmul \inv{(m + \epsilon)}) \rmul \epsilon < 1 \rmul \epsilon$ by \cref{real_lt_mul_pos}.
                        $(m \rmul \inv{(m + \epsilon)}) \rmul \epsilon = \epsilon \rmul (m \rmul \inv{(m + \epsilon)})$ by \cref{real_mul_comm}.
                        Thus $\epsilon \rmul (m \rmul \inv{(m + \epsilon)}) < 1 \rmul \epsilon$ by \cref{real_lt_subst_left}.
                        $1 \rmul \epsilon = \epsilon \rmul 1$ by \cref{real_mul_comm}.
                        Thus $\epsilon \rmul (m \rmul \inv{(m + \epsilon)}) < \epsilon \rmul 1$ by \cref{real_lt_subst_right}.
                        Therefore $\epsilon \rmul 1 = \epsilon$ by \cref{real_mul_comm,real_mul_ident}.
                        $\delta = \epsilon \rmul \inv{(m + \epsilon)}$ by \cref{real_div}.
                        Hence $\delta \rmul m = \epsilon \rmul (m \rmul \inv{(m + \epsilon)})$ by \cref{real_div,real_mul_assoc,real_mul_comm}.
                        Therefore $\delta \rmul m < \epsilon$ by \cref{real_lt_subst_left}.
                        $x - x_0\in\reals$ and $y - y_0\in\reals$ by \cref{real_sub,real_add_closed}.
                        $\abs{x - x_0}\in\reals$ and $\abs{y}\in\reals$ and $\abs{y - y_0}\in\reals$ by \cref{real_abs_in_reals}.
                        $\abs{x - x_0} \rmul \abs{y}\in\reals$ by \cref{real_mul_closed}.
                        $\abs{x_0} \rmul \abs{y - y_0}\in\reals$ by \cref{real_mul_closed}.
                        $\abs{x - x_0} \rmul \abs{y} + \abs{x_0} \rmul \abs{y - y_0}\in\reals$ by \cref{real_add_closed}.
                        $\delta \rmul m\in\reals$ by \cref{real_mul_closed}.
                        Thus $\abs{x - x_0} \rmul \abs{y} + \abs{x_0} \rmul \abs{y - y_0} < \epsilon$ by \cref{real_lt_trans}.
                        $(x - x_0) \rmul y + x_0 \rmul (y - y_0)\in\reals$ and
                            $\abs{(x - x_0) \rmul y + x_0 \rmul (y - y_0)}\in\reals$ by
                            \cref{real_add_closed,real_abs_in_reals}.
                        Thus $\abs{(x - x_0) \rmul y + x_0 \rmul (y - y_0)} < \epsilon$
                            by \cref{real_leq_split,real_lt_trans,real_lt_subst_left}.
                        Hence $\abs{(x \rmul y) - (x_0 \rmul y_0)} < \epsilon$ by \cref{real_lt_subst_left}.
                        Therefore $\abs{f(q) - f(p)} < \epsilon$ by \cref{real_lt_subst_left}.
                \end{subproof}
                $f$ is a function to $\reals$ and $\dom{f} = \reals\times\reals$ by \cref{funs_elim}.
                Thus $f(p)\in\reals$ and $f(q)\in\reals$.
                Hence $d((f(p),f(q))) < \epsilon$ by \cref{standard_metric_value,real_abs_sub_swap,real_lt_subst_left}.
                Thus $f(q)\in\metricBall{d}{\reals}{f(p)}{\epsilon}$ by \cref{metric_ball}.
            \end{subproof}
            Thus $\img{f}{U}\subseteq \metricBall{d}{\reals}{f(p)}{\epsilon}$ by \cref{img_iff,function_apply_intro,subseteq}.
            Thus $\img{f}{U}\subseteq V$ by \cref{subseteq,subseteq_transitive}.
            Thus $U$ is a neighborhood of $p$ in $\reals\times\reals$ with respect to $T$ by \cref{neighborhood}.
        \end{subproof}
        $\standardBasisR$ is a basis on $\reals$ and $\standardTopologyR$ is a topology on $\reals$
            by \cref{standard_basis_is_basis,basis_topology_is_topology,standard_topology_reals}.
        $\productBasis{\standardTopologyR}{\standardTopologyR}{\reals}{\reals}$ is a basis on $\reals\times\reals$ and
            $T$ is a topology on $\reals\times\reals$ by
            \cref{product_basis_is_basis,basis_topology_is_topology,product_topology}.
        Hence $f$ is continuous at $p$ from $\reals\times\reals$ to $\reals$ with respect to $T$ and $\standardTopologyR$ by \cref{continuous_at}.
    \end{subproof}
    $\standardBasisR$ is a basis on $\reals$ and $\standardTopologyR$ is a topology on $\reals$
        by \cref{standard_basis_is_basis,basis_topology_is_topology,standard_topology_reals}.
    $\productBasis{\standardTopologyR}{\standardTopologyR}{\reals}{\reals}$ is a basis on $\reals\times\reals$ and
        $T$ is a topology on $\reals\times\reals$ by
        \cref{product_basis_is_basis,basis_topology_is_topology,product_topology}.
    Therefore $f$ is continuous from $\reals\times\reals$ to $\reals$ with respect to $T$ and $\standardTopologyR$ by \cref{continuous_at_implies_continuous}.
\end{proof}

% from §21 helper definition 5: reciprocal operation on reals
\begin{definition}\label{real_inverse_operation}
    $\realInverse = \{w \mid \exists y\in\reals\setminus\{0\}. \exists z\in\reals.
        w = (y,z) \land z = \inv{y}\}$.
\end{definition}
% from §21 helper lemma 7: reciprocal operation is continuous
\begin{lemma}\label{real_inverse_continuous}
    Let $T_1 = \standardTopologyR$.
    Let $T_2 = \subspaceTopology{T_1}{\reals\setminus\{0\}}$.
    Then $\realInverse$ is continuous from $\reals\setminus\{0\}$ to $\reals$
        with respect to $T_2$ and $T_1$.
\end{lemma}
\begin{proof}
    Let $f = \realInverse$.
    For all $p,z_1,z_2$ such that $(p,z_1)\in f$ and $(p,z_2)\in f$ we have $z_1 = z_2$
        by \cref{real_inverse_operation,pair_eq_iff}.
    Then $f$ is a function by \cref{rightunique,real_inverse_operation,relation}.
    For all $p\in\dom{f}$ we have $f(p)\in\reals$
        by \cref{dom_iff,real_inverse_operation,pair_eq_iff,function_apply_intro}.
    Hence $f$ is a function to $\reals$.
    Therefore $\dom{f}\subseteq\reals\setminus\{0\}$ by \cref{dom_iff,real_inverse_operation,pair_eq_iff,subseteq}.
    For all $p\in\reals\setminus\{0\}$ we have $p\in\dom{f}$
        by \cref{setminus_elim_left,setminus_elim_right,singleton_iff,real_mul_inverse,real_inverse_operation,dom_intro}.
    Hence $\reals\setminus\{0\}\subseteq\dom{f}$ by \cref{subseteq}.
    Then $f\in\funs{\reals\setminus\{0\}}{\reals}$ by \cref{subseteq_antisymmetric,funs_intro}.
    Show that for all $p\in\reals\setminus\{0\}$ we have
        $f$ is continuous at $p$ from $\reals\setminus\{0\}$ to $\reals$ with respect to $T_2$ and $T_1$.
        \begin{subproof}
            Fix $p$ such that $p\in\reals\setminus\{0\}$.
            $p\in\reals$ and $p\notin\{0\}$ by \cref{setminus_elim_left,setminus_elim_right}.
            Hence $p\neq 0$ and $\inv{p}\in\reals$ by \cref{singleton_iff,real_mul_inverse}.
        Then $f(p) = \inv{p}$ by \cref{real_inverse_operation,function_apply_intro}.
        Show that for all $V$ such that $V$ is a neighborhood of $f(p)$ in $\reals$ with respect to $T_1$
            there exists $U$ such that $U$ is a neighborhood of $p$ in $\reals\setminus\{0\}$ with respect to $T_2$
            and $\img{f}{U}\subseteq V$.
        \begin{subproof}
            Fix $V$ such that $V$ is a neighborhood of $f(p)$ in $\reals$ with respect to $T_1$.
            $V$ is open in $\reals$ with respect to $T_1$ by \cref{neighborhood}.
            Let $d = \standardMetricR$.
            Let $T_d = \metricTopology{d}{\reals}$.
            $T_d = \standardTopologyR$ by \cref{standard_metric_topology}.
            Hence $V\in T_d$ by \cref{open_in,subseteq}.
            $f(p)\in V$ by \cref{neighborhood}.
            $\standardBasisR$ is a basis on $\reals$ and $\standardTopologyR$ is a topology on $\reals$
                by \cref{standard_basis_is_basis,basis_topology_is_topology,standard_topology_reals}.
            Therefore $V\subseteq\reals$ by \cref{topology_on,subseteq,pow_iff}.
            $d$ is a metric on $\reals$ by \cref{standard_metric_is_metric}.
            Take $\epsilon\in\reals$ such that $0 < \epsilon$ and
                $\metricBall{d}{\reals}{f(p)}{\epsilon} \subseteq V$ by \cref{metric_open_iff}.
            $\abs{p}\in\reals$ by \cref{real_abs_in_reals}.
            Hence $\abs{p} > 0$ by \cref{real_abs_eq_zero,real_lt_trichotomy,real_abs_nonnegative}.
            Let $a = \abs{p} \rmul \inv{(1 + 1)}$.
            $1\in\reals$ and $0\in\reals$ and $0 < 1$ and $1 + 1\in\reals$
                by \cref{real_mul_ident,real_add_ident,real_zero_lt_one,real_add_closed}.
            $0 + 1 < 1 + 1$ and $0 + 1 = 1$ by \cref{real_lt_add,real_add_ident}.
            Therefore $1 < 1 + 1$ by \cref{real_lt_subst_left}.
            Hence $0 < 1 + 1$ by \cref{real_lt_trans}.
            Therefore $(1 + 1)\neq 0$ by \cref{real_pos_nonzero}.
            Then $\inv{(1 + 1)}\in\reals$ by \cref{real_mul_inverse}.
            $0 < \inv{(1 + 1)}$ by \cref{real_inv_pos}.
            $0 \rmul \inv{(1 + 1)} < \abs{p} \rmul \inv{(1 + 1)}$ by \cref{real_lt_mul_pos}.
            Hence $0 \rmul \inv{(1 + 1)} < a$ by \cref{real_lt_subst_right}.
            $0 \rmul \inv{(1 + 1)} = 0$ by \cref{real_mul_zero}.
            Therefore $0 < a$ by \cref{real_lt_subst_left}.
            Let $b = (\epsilon \rmul (\abs{p} \rmul \abs{p})) \rmul \inv{(1 + 1)}$.
            $\abs{p} \rmul \abs{p}\in\reals$ and $\epsilon \rmul (\abs{p} \rmul \abs{p})\in\reals$ and $(1 + 1)\in\reals$
                by \cref{real_mul_closed,real_add_closed}.
            Hence $\abs{p} \rmul \abs{p} > 0$ by \cref{real_mul_pos_pos_gt}.
            $0 \rmul (\abs{p} \rmul \abs{p}) < \epsilon \rmul (\abs{p} \rmul \abs{p})$ by \cref{real_lt_mul_pos}.
            $0 \rmul (\abs{p} \rmul \abs{p}) = 0$ by \cref{real_mul_zero}.
            Therefore $0 < \epsilon \rmul (\abs{p} \rmul \abs{p})$ by \cref{real_lt_subst_left}.
            $0 \rmul \inv{(1 + 1)} < (\epsilon \rmul (\abs{p} \rmul \abs{p})) \rmul \inv{(1 + 1)}$ by \cref{real_lt_mul_pos}.
            Hence $0 \rmul \inv{(1 + 1)} < b$ by \cref{real_lt_subst_right}.
            $0 \rmul \inv{(1 + 1)} = 0$ by \cref{real_mul_zero}.
            Therefore $0 < b$ by \cref{real_lt_subst_left}.
            $a\in\reals$ and $b\in\reals$ by \cref{real_mul_closed}.
            Let $E = \{a,b\}$.
            For all $x$ such that $x\in E$ we have $x\in\reals$.
            Hence $E\subseteq\reals$ by \cref{subseteq}.
            $E$ is finite by \cref{pair_is_finite}.
            $E\neq\emptyset$ by \cref{upair_intro_left,emptyset}.
            Take $\delta\in E$ such that $\delta$ is a minimum of $E$ by \cref{finite_real_minimum}.
            Therefore $\delta\in\reals$ by \cref{subseteq}.
            $a\in E$ and $b\in E$ by \cref{upair_intro_left,upair_intro_right}.
            $0 < \delta$ by \cref{upair_elim,real_lt_subst_right}.
            Hence $\delta \leq a$ and $\delta \leq b$ by \cref{real_minimum}.
            Let $U = \metricBall{d}{\reals}{p}{\delta}$.
            Show that for all $y$ such that $y\in U$ we have $y\in\reals\setminus\{0\}$.
            \begin{subproof}
                Fix $y$ such that $y\in U$.
                Hence $y\in\reals$ by \cref{metric_ball}.
                $d((p,y)) < \delta$ by \cref{metric_ball}.
                $d((p,y)) = \abs{p - y}$ by \cref{standard_metric_value}.
                Then $\abs{p - y} < \delta$.
                Show that $y\neq 0$.
                \begin{subproof}
                    $\neg{y}\in\reals$ by \cref{real_add_inverse}.
                    $p - y\in\reals$ by \cref{real_sub,real_add_closed}.
                    Hence $\abs{p - y}\in\reals$ by \cref{real_abs_in_reals}.
                    Therefore $\abs{p - y} < a$ by \cref{real_leq_split,real_lt_trans,real_lt_subst_right}.
                    $a + a = \abs{p}$ by
                        \cref{real_right_distrib,real_mul_ident,real_mul_eq_subst,real_mul_assoc,real_mul_comm,real_mul_inverse}.
                    $\abs{p} < \abs{y} + a$ by \cref{real_abs_sub_lt_bound}.
                    Show that it is wrong that $\abs{y} \leq a$.
                    \begin{subproof}
                        Suppose not.
                        Then $\abs{y} \leq a$.
                        $\abs{y}\in\reals$ by \cref{real_abs_in_reals}.
                        Hence $\abs{y} + a \leq a + a$ by \cref{real_leq_add}.
                        $\abs{y} + a\in\reals$ by \cref{real_add_closed}.
                        Then $\abs{y} + a < a + a$ or $\abs{y} + a = a + a$ by \cref{real_leq_split}.
                        \begin{byCase}
                        \caseOf{$\abs{y} + a < a + a$.}
                            Hence $\abs{y} + a < \abs{p}$ by \cref{real_lt_subst_right}.
                            Therefore $\abs{p} < \abs{p}$ by \cref{real_lt_trans}.
                            Contradiction by \cref{real_lt_irreflexive}.
                        \caseOf{$\abs{y} + a = a + a$.}
                            Hence $\abs{y} + a = \abs{p}$.
                            Therefore $\abs{p} < \abs{p}$ by \cref{real_lt_subst_right}.
                            Contradiction by \cref{real_lt_irreflexive}.
                        \end{byCase}
                    \end{subproof}
                    Hence $\abs{y} > a$.
                    Therefore $\abs{y}\neq 0$ by \cref{real_lt_trichotomy,real_abs_nonnegative}.
                    Show that $y\neq 0$.
                    \begin{subproof}
                        Suppose not.
                        Then $y = 0$.
                        Then $0 < y$ or $0 = y$.
                        Therefore $0 \leq y$.
                        Hence $\abs{y} = y$ by \cref{real_abs_nonneg}.
                        Then $\abs{y} = 0$.
                        Contradiction.
                    \end{subproof}
                \end{subproof}
                Therefore $y\in\reals\setminus\{0\}$ by \cref{singleton_iff,setminus_intro}.
            \end{subproof}
            Hence $U\subseteq\reals\setminus\{0\}$ by \cref{subseteq}.
            $p\in\reals$ and $\delta\in\reals$ and $0 < \delta$.
            For all $y$ such that $y\in U$ we have $y\in\reals$ by \cref{metric_ball}.
            Therefore $U\in\metricBasis{d}{\reals}$ by \cref{subseteq,powerset_intro,metric_basis}.
            Hence $U\in T_d$ by \cref{metric_basis_is_basis,basis_elem_open_in_generated,basis_for_topology,metric_topology}.
            $T_d\subseteq\standardTopologyR$ by \cref{subseteq}.
            Therefore $U\in\standardTopologyR$ by \cref{subseteq}.
            $\standardBasisR$ is a basis on $\reals$ and $\standardTopologyR$ is a topology on $\reals$
                by \cref{standard_basis_is_basis,basis_topology_is_topology,standard_topology_reals}.
            Hence $U$ is open in $\reals$ with respect to $\standardTopologyR$ by \cref{open_in}.
            $U\in T_1$ by \cref{open_in}.
            Therefore $U\subseteq(\reals\setminus\{0\})\inter U$ by \cref{subseteq,subseteq_inter_iff}.
            $(\reals\setminus\{0\})\inter U\subseteq U$ by \cref{inter_lower_right}.
            Hence $U = (\reals\setminus\{0\})\inter U$ by \cref{subseteq_antisymmetric}.
            Then $U\in T_2$ by \cref{subspace_topology}.
            $\standardBasisR$ is a basis on $\reals$ and $T_1$ is a topology on $\reals$
                by \cref{standard_basis_is_basis,basis_topology_is_topology,standard_topology_reals}.
            $\reals\setminus\{0\}\subseteq\reals$ and $T_2$ is a topology on $\reals\setminus\{0\}$
                by \cref{setminus_elim_left,subseteq,subspace_topology_is_topology}.
            Therefore $U$ is open in $\reals\setminus\{0\}$ with respect to $T_2$ by \cref{open_in}.
            $d((p,p)) = 0$ by \cref{metric_zero_characterization,standard_metric_is_metric,metric_on}.
            Hence $d((p,p)) < \delta$ by \cref{real_lt_subst_left}.
            Then $p\in U$ by \cref{metric_ball}.
            Therefore $U$ is a neighborhood of $p$ in $\reals\setminus\{0\}$ with respect to $T_2$ by \cref{neighborhood}.
            Show that for all $q$ such that $q\in U$ we have
                $f(q)\in\metricBall{d}{\reals}{f(p)}{\epsilon}$.
            \begin{subproof}
                Fix $q$ such that $q\in U$.
                $q\in\reals$ by \cref{metric_ball}.
                $d((p,q)) < \delta$ by \cref{metric_ball}.
                $d((p,q)) = \abs{p - q}$ by \cref{standard_metric_value}.
                Then $\abs{p - q} < \delta$.
                Hence $q\neq 0$ by \cref{subseteq,setminus_elim_right,metric_ball,singleton_iff}.
                Therefore $q\in\reals\setminus\{0\}$ by \cref{setminus_intro}.
                $\inv{q}\in\reals$ by \cref{real_mul_inverse}.
                Thus $f(q) = \inv{q}$ by \cref{real_inverse_operation,function_apply_intro}.
                    Show that $q \rmul p \neq 0$.
                    \begin{subproof}
                        Suppose not.
                        Then $q \rmul p = 0$.
                        Hence $q = 0$ or $p = 0$ by \cref{real_mul_zero_product}.
                        Contradiction.
                    \end{subproof}
                    Therefore $q \rmul p \neq 0$.
                    $q \rmul p\in\reals$ by \cref{real_mul_closed}.
                    $\inv{(q \rmul p)}\in\reals$ by \cref{real_mul_inverse}.
                    $q \rmul (p \rmul \inv{(q \rmul p)}) = 1$ by \cref{real_mul_assoc,real_mul_inverse}.
                    $p \rmul \inv{(q \rmul p)}\in\reals$ by \cref{real_mul_closed}.
                    Hence $\inv{q} = p \rmul \inv{(q \rmul p)}$ by \cref{real_mul_inverse_unique}.
                    $p \rmul (q \rmul \inv{(q \rmul p)}) = 1$ by \cref{real_mul_assoc,real_mul_comm,real_mul_inverse}.
                    $q \rmul \inv{(q \rmul p)}\in\reals$ by \cref{real_mul_closed}.
                    Therefore $\inv{p} = q \rmul \inv{(q \rmul p)}$ by \cref{real_mul_inverse_unique}.
                        $\inv{q} - \inv{p} = (p \rmul \inv{(q \rmul p)}) - (q \rmul \inv{(q \rmul p)})$ by \cref{real_lt_subst_left}.
                        $(p \rmul \inv{(q \rmul p)}) - (q \rmul \inv{(q \rmul p)}) =
                            (p \rmul \inv{(q \rmul p)}) + \neg{(q \rmul \inv{(q \rmul p)})}$ by \cref{real_sub}.
                        Show that $\neg{(q \rmul \inv{(q \rmul p)})} = (\neg{q}) \rmul \inv{(q \rmul p)}$.
                        \begin{subproof}
                            $\neg{(q \rmul \inv{(q \rmul p)})} = \neg{1} \rmul (q \rmul \inv{(q \rmul p)})$ by \cref{real_neg_as_mul}.
                            $q \rmul \inv{(q \rmul p)} = \inv{(q \rmul p)} \rmul q$ by \cref{real_mul_comm}.
                            Thus $\neg{1} \rmul (q \rmul \inv{(q \rmul p)}) = \neg{1} \rmul (\inv{(q \rmul p)} \rmul q)$ by \cref{real_mul_eq_subst}.
                            Thus $\neg{(q \rmul \inv{(q \rmul p)})} = \neg{1} \rmul (\inv{(q \rmul p)} \rmul q)$ by \cref{real_mul_eq_subst}.
                            $\neg{1} \rmul (\inv{(q \rmul p)} \rmul q) = \neg{(\inv{(q \rmul p)} \rmul q)}$ by \cref{real_neg_as_mul}.
                            $\inv{(q \rmul p)} \rmul \neg{q} = \neg{(\inv{(q \rmul p)} \rmul q)}$ by \cref{real_mul_neg_right}.
                            $\neg{q}\in\reals$ by \cref{real_add_inverse}.
                            $\inv{(q \rmul p)} \rmul \neg{q} = (\neg{q}) \rmul \inv{(q \rmul p)}$ by \cref{real_mul_comm}.
                            Thus $\neg{(q \rmul \inv{(q \rmul p)})} = (\neg{q}) \rmul \inv{(q \rmul p)}$.
                        \end{subproof}
                        Hence $(p \rmul \inv{(q \rmul p)}) + \neg{(q \rmul \inv{(q \rmul p)})} =
                            (p \rmul \inv{(q \rmul p)}) + ((\neg{q}) \rmul \inv{(q \rmul p)})$ by \cref{real_add_eq_subst}.
                        $\neg{q}\in\reals$ by \cref{real_add_inverse}.
                        $(p \rmul \inv{(q \rmul p)}) + ((\neg{q}) \rmul \inv{(q \rmul p)}) =
                            (p + \neg{q}) \rmul \inv{(q \rmul p)}$ by \cref{real_right_distrib}.
                        $p + \neg{q} = p - q$ by \cref{real_sub}.
                        Thus $(p \rmul \inv{(q \rmul p)}) + ((\neg{q}) \rmul \inv{(q \rmul p)}) =
                            (p - q) \rmul \inv{(q \rmul p)}$ by \cref{real_mul_eq_subst}.
                        Thus $(p \rmul \inv{(q \rmul p)}) - (q \rmul \inv{(q \rmul p)}) = (p - q) \rmul \inv{(q \rmul p)}$
                            by \cref{real_add_eq_subst}.
                        Thus $\abs{\inv{q} - \inv{p}} = \abs{(p - q) \rmul \inv{(q \rmul p)}}$ by \cref{real_lt_subst_left}.
                        $\neg{q}\in\reals$ by \cref{real_add_inverse}.
                        $p - q\in\reals$ by \cref{real_sub,real_add_closed}.
                        $\abs{(p - q) \rmul \inv{(q \rmul p)}} = \abs{p - q} \rmul \abs{\inv{(q \rmul p)}}$ by \cref{real_abs_mul}.
                        $\abs{p - q} = \abs{q - p}$ by \cref{real_abs_sub_swap}.
                        Show that $\abs{p - q} < \epsilon \rmul \abs{q \rmul p}$.
                        \begin{subproof}
                            $\neg{q}\in\reals$ by \cref{real_add_inverse}.
                            $p - q\in\reals$ by \cref{real_sub,real_add_closed}.
                            Thus $\abs{p - q}\in\reals$ by \cref{real_abs_in_reals}.
                            $a + a = \abs{p}$ by
                                \cref{real_right_distrib,real_mul_ident,real_mul_eq_subst,real_mul_assoc,real_mul_comm,real_mul_inverse}.
                            Show that it is wrong that $\abs{q} \leq a$.
                            \begin{subproof}
                                Suppose not.
                                Then $\abs{q} \leq a$.
                                $\abs{q}\in\reals$ by \cref{real_abs_in_reals}.
                                Thus $\abs{p - q} < a$ by \cref{real_leq_split,real_lt_trans,real_lt_subst_right}.
                                $\abs{p} < \abs{q} + a$ by \cref{real_abs_sub_lt_bound}.
                                Hence $\abs{q} + a \leq a + a$ by \cref{real_leq_add}.
                                $\abs{q} + a\in\reals$ by \cref{real_add_closed}.
                                Then $\abs{q} + a < a + a$ or $\abs{q} + a = a + a$ by \cref{real_leq_split}.
                                \begin{byCase}
                                \caseOf{$\abs{q} + a < a + a$.}
                                    Thus $\abs{q} + a < \abs{p}$ by \cref{real_lt_subst_right}.
                                    Thus $\abs{p} < \abs{p}$ by \cref{real_lt_trans}.
                                    Contradiction by \cref{real_lt_irreflexive}.
                                \caseOf{$\abs{q} + a = a + a$.}
                                    Thus $\abs{q} + a = \abs{p}$.
                                    Thus $\abs{p} < \abs{p}$ by \cref{real_lt_subst_right}.
                                    Contradiction by \cref{real_lt_irreflexive}.
                                \end{byCase}
                            \end{subproof}
                            Hence $\abs{q} > a$.
                            $\abs{p}\in\reals$ by \cref{real_abs_in_reals}.
                            Therefore $0 < \abs{p}$ by \cref{real_abs_eq_zero,real_abs_nonnegative,real_lt_trichotomy}.
                            $a\in\reals$ by \cref{real_mul_closed}.
                            $\abs{q}\in\reals$ by \cref{real_abs_in_reals}.
                            $a \rmul \abs{p} < \abs{q} \rmul \abs{p}$ by \cref{real_lt_mul_pos}.
                            $\abs{q \rmul p} = \abs{q} \rmul \abs{p}$ by \cref{real_abs_mul}.
                            Hence $\abs{q \rmul p} > a \rmul \abs{p}$ by \cref{real_lt_subst_left}.
                            Therefore $\abs{p - q} < b$ by \cref{real_leq_split,real_lt_trans,real_lt_subst_right}.
                            $a \rmul \abs{p} = (\abs{p} \rmul \abs{p}) \rmul \inv{(1 + 1)}$
                                by \cref{real_mul_eq_subst,real_mul_assoc,real_mul_comm}.
                            Thus $\epsilon \rmul (a \rmul \abs{p}) = (\epsilon \rmul (\abs{p} \rmul \abs{p})) \rmul \inv{(1 + 1)}$
                                by \cref{real_mul_assoc,real_mul_comm}.
                            Hence $b = \epsilon \rmul (a \rmul \abs{p})$.
                            $\abs{q \rmul p} = \abs{q} \rmul \abs{p}$ by \cref{real_abs_mul}.
                            Therefore $a \rmul \abs{p} < \abs{q \rmul p}$ by \cref{real_lt_subst_right}.
                            $a \rmul \abs{p}\in\reals$ by \cref{real_mul_closed}.
                            $\abs{q \rmul p}\in\reals$ by \cref{real_abs_in_reals}.
                            $(a \rmul \abs{p}) \rmul \epsilon < \abs{q \rmul p} \rmul \epsilon$ by \cref{real_lt_mul_pos}.
                            $\epsilon \rmul (a \rmul \abs{p}) = (a \rmul \abs{p}) \rmul \epsilon$ by \cref{real_mul_comm}.
                            $\epsilon \rmul \abs{q \rmul p} = \abs{q \rmul p} \rmul \epsilon$ by \cref{real_mul_comm}.
                            Hence $\epsilon \rmul (a \rmul \abs{p}) < \abs{q \rmul p} \rmul \epsilon$ by \cref{real_lt_subst_left}.
                            Therefore $\epsilon \rmul (a \rmul \abs{p}) < \epsilon \rmul \abs{q \rmul p}$ by \cref{real_lt_subst_right}.
                            Then $b < \epsilon \rmul \abs{q \rmul p}$ by \cref{real_lt_subst_left}.
                            $\epsilon \rmul \abs{q \rmul p}\in\reals$ by \cref{real_mul_closed}.
                            Hence $\abs{p - q} < \epsilon \rmul \abs{q \rmul p}$ by \cref{real_lt_trans}.
                        \end{subproof}
                        $\inv{(q \rmul p)} \rmul (q \rmul p) = 1$ by \cref{real_mul_inverse,real_mul_comm}.
                        Therefore $\inv{(q \rmul p)} \neq 0$ by \cref{real_mul_unit_left_nonzero}.
                        Hence $\abs{\inv{(q \rmul p)}}\neq 0$ by \cref{real_abs_eq_zero}.
                        $\abs{\inv{(q \rmul p)}} \geq 0$ by \cref{real_abs_nonnegative}.
                        Therefore $0 < \abs{\inv{(q \rmul p)}}$ by \cref{real_lt_trichotomy}.
                        $\abs{p - q}\in\reals$ by \cref{real_abs_in_reals}.
                        $\abs{\inv{(q \rmul p)}}\in\reals$ by \cref{real_abs_in_reals}.
                        $\epsilon\in\reals$.
                        $\abs{q \rmul p}\in\reals$ by \cref{real_abs_in_reals}.
                        $\epsilon \rmul \abs{q \rmul p}\in\reals$ by \cref{real_mul_closed}.
                        Thus $\abs{p - q} \rmul \abs{\inv{(q \rmul p)}} < \epsilon \rmul \abs{q \rmul p} \rmul \abs{\inv{(q \rmul p)}}$
                            by \cref{real_lt_mul_pos}.
                        $\abs{q \rmul p} \rmul \abs{\inv{(q \rmul p)}} = \abs{(q \rmul p) \rmul \inv{(q \rmul p)}}$ by \cref{real_abs_mul,real_mul_assoc}.
                        $(q \rmul p) \rmul \inv{(q \rmul p)} = 1$ by \cref{real_mul_inverse}.
                        $\abs{1} = 1$ by \cref{real_abs_nonneg,real_zero_lt_one,real_lt_imp_leq}.
                        Hence $\abs{q \rmul p} \rmul \abs{\inv{(q \rmul p)}} = 1$ by \cref{real_lt_subst_right}.
                        $\epsilon \rmul \abs{q \rmul p} \rmul \abs{\inv{(q \rmul p)}} =
                            \epsilon \rmul (\abs{q \rmul p} \rmul \abs{\inv{(q \rmul p)}})$ by \cref{real_mul_assoc}.
                        Therefore $\epsilon \rmul \abs{q \rmul p} \rmul \abs{\inv{(q \rmul p)}} = \epsilon \rmul 1$ by \cref{real_mul_eq_subst}.
                        $\epsilon \rmul 1 = 1 \rmul \epsilon$ by \cref{real_mul_comm}.
                        Then $\epsilon \rmul 1 = \epsilon$ by \cref{real_mul_ident}.
                        Hence $\epsilon \rmul \abs{q \rmul p} \rmul \abs{\inv{(q \rmul p)}} = \epsilon$ by \cref{real_mul_eq_subst}.
                        Therefore $\abs{p - q} \rmul \abs{\inv{(q \rmul p)}} < \epsilon$ by \cref{real_lt_subst_right}.
                    Hence $\abs{\inv{q} - \inv{p}} < \epsilon$ by \cref{real_lt_subst_left}.
                Therefore $\abs{f(q) - f(p)} < \epsilon$ by \cref{real_lt_subst_left,real_div}.
                $f(p)\in\reals$ and $f(q)\in\reals$ by \cref{funs_type_apply}.
                Hence $d((f(p),f(q))) < \epsilon$ by \cref{standard_metric_value,real_abs_sub_swap,real_lt_subst_left}.
                Therefore $f(q)\in\metricBall{d}{\reals}{f(p)}{\epsilon}$ by \cref{metric_ball}.
            \end{subproof}
            Hence $\img{f}{U}\subseteq \metricBall{d}{\reals}{f(p)}{\epsilon}$ by \cref{img_iff,function_apply_intro,subseteq}.
            $\metricBall{d}{\reals}{f(p)}{\epsilon}\subseteq V$ by \cref{subseteq}.
            Therefore $\img{f}{U}\subseteq V$ by \cref{subseteq_transitive}.
        \end{subproof}
        $\standardBasisR$ is a basis on $\reals$ and $T_1$ is a topology on $\reals$
            by \cref{standard_basis_is_basis,basis_topology_is_topology,standard_topology_reals}.
        $\reals\setminus\{0\}\subseteq\reals$ and $T_2$ is a topology on $\reals\setminus\{0\}$
            by \cref{setminus_elim_left,subseteq,subspace_topology_is_topology}.
        Hence $f$ is continuous at $p$ from $\reals\setminus\{0\}$ to $\reals$ with respect to $T_2$ and $T_1$ by \cref{continuous_at}.
    \end{subproof}
    $\standardBasisR$ is a basis on $\reals$ and $T_1$ is a topology on $\reals$
        by \cref{standard_basis_is_basis,basis_topology_is_topology,standard_topology_reals}.
    $\reals\setminus\{0\}\subseteq\reals$ and $T_2$ is a topology on $\reals\setminus\{0\}$
        by \cref{setminus_elim_left,subseteq,subspace_topology_is_topology}.
    Therefore $f$ is continuous from $\reals\setminus\{0\}$ to $\reals$ with respect to $T_2$ and $T_1$ by \cref{continuous_at_implies_continuous}.
\end{proof}

% from §21 Item 11: quotient operation is continuous
\begin{lemma}\label{real_division_continuous}
    Let $T_1 = \standardTopologyR$.
    Let $T_2 = \subspaceTopology{T_1}{\reals\setminus\{0\}}$.
    Let $T = \productTopology{T_1}{T_2}{\reals}{\reals\setminus\{0\}}$.
    Then $\realDivision$ is continuous from $\reals\times(\reals\setminus\{0\})$ to $\reals$
        with respect to $T$ and $\standardTopologyR$.
\end{lemma}
\begin{proof}
    Let $T_3 = \productTopology{T_1}{T_1}{\reals}{\reals}$.
    $\standardBasisR$ is a basis on $\reals$ and $T_1$ is a topology on $\reals$
        by \cref{standard_basis_is_basis,basis_topology_is_topology,standard_topology_reals}.
    $\reals\setminus\{0\}\subseteq\reals$ and $T_2$ is a topology on $\reals\setminus\{0\}$
        by \cref{setminus_elim_left,subseteq,subspace_topology_is_topology}.
    Let $f_1 = \piOne{\reals}{\reals\setminus\{0\}}$.
    Let $f_2 = \piTwo{\reals}{\reals\setminus\{0\}}$.
    $f_1$ is continuous from $\reals\times(\reals\setminus\{0\})$ to $\reals$ with respect to $T$ and $T_1$
        by \cref{projection_one_continuous}.
    $f_2$ is continuous from $\reals\times(\reals\setminus\{0\})$ to $\reals\setminus\{0\}$ with respect to $T$ and $T_2$
        by \cref{projection_two_continuous}.
    $\realInverse$ is continuous from $\reals\setminus\{0\}$ to $\reals$ with respect to $T_2$ and $T_1$
        by \cref{real_inverse_continuous}.
    Let $g = \realInverse\circ f_2$.
    Hence $g$ is continuous from $\reals\times(\reals\setminus\{0\})$ to $\reals$ with respect to $T$ and $T_1$
        by \cref{continuous_composite}.
    Let $k = \{(p,(f_1(p), g(p))) \mid p\in\reals\times(\reals\setminus\{0\})\}$.
    Therefore $k$ is continuous from $\reals\times(\reals\setminus\{0\})$ to $\reals\times\reals$
        with respect to $T$ and $T_3$ by \cref{continuous_map_into_product_backward}.
    $\realMultiplication$ is continuous from $\reals\times\reals$ to $\reals$ with respect to $T_3$ and $T_1$
        by \cref{real_multiplication_continuous}.
    Let $m = \realMultiplication\circ k$.
    Hence $m$ is continuous from $\reals\times(\reals\setminus\{0\})$ to $\reals$ with respect to $T$ and $T_1$
        by \cref{continuous_composite}.
        $m$ is a function and $\dom{m} = \reals\times(\reals\setminus\{0\})$ by \cref{continuous,funs_elim}.
            For all $p,z_1,z_2$ such that $(p,z_1)\in\realDivision$ and $(p,z_2)\in\realDivision$ we have $z_1 = z_2$.
            Thus $\realDivision$ is a function by \cref{rightunique,real_division_operation,relation}.
            For all $p\in\dom{\realDivision}$ we have $\apply{\realDivision}{p}\in\reals$
                by \cref{dom_iff,real_division_operation,pair_eq_iff,function_apply_intro}.
            Thus $\realDivision$ is a function to $\reals$.
            Thus $\dom{\realDivision}\subseteq\reals\times(\reals\setminus\{0\})$
                by \cref{dom_iff,real_division_operation,pair_eq_iff,subseteq}.
            For all $p\in\reals\times(\reals\setminus\{0\})$ we have $p\in\dom{\realDivision}$
                by \cref{times_elem_is_tuple,fst_eq,snd_eq,setminus_elim_left,setminus_elim_right,singleton_iff,real_mul_inverse,real_mul_closed,real_division_operation,dom_intro}.
            Thus $\reals\times(\reals\setminus\{0\})\subseteq\dom{\realDivision}$ by \cref{subseteq}.
            Thus $\dom{\realDivision} = \reals\times(\reals\setminus\{0\})$ by \cref{subseteq_antisymmetric}.
            Thus $\realDivision\in\funs{\reals\times(\reals\setminus\{0\})}{\reals}$ by \cref{funs_intro}.
        Show that for all $p\in\reals\times(\reals\setminus\{0\})$ we have $m(p) = \apply{\realDivision}{p}$.
        \begin{subproof}
            Fix $p\in\reals\times(\reals\setminus\{0\})$.
            $k\in\funs{\reals\times(\reals\setminus\{0\})}{\reals\times\reals}$ by \cref{continuous}.
            $(p,(f_1(p), g(p)))\in k$.
            $k$ is a function by \cref{continuous,funs_is_function}.
            Thus $k(p) = (f_1(p), g(p))$ by \cref{continuous,funs_is_function,function_apply_intro}.
            $\realMultiplication\in\funs{\reals\times\reals}{\reals}$ and
                $\dom{\realMultiplication} = \reals\times\reals$ by \cref{continuous,funs}.
            $\ran{k}\subseteq\reals\times\reals$ by \cref{funs_ran}.
            Thus $\realMultiplication$ is a function by \cref{funs_is_function}.
            Thus $p\in\dom{k}$ by \cref{funs}.
            Thus $m(p) = \apply{\realMultiplication}{k(p)}$ by \cref{circ_apply}.
            $f_1\in\funs{\reals\times(\reals\setminus\{0\})}{\reals}$ and
                $g\in\funs{\reals\times(\reals\setminus\{0\})}{\reals}$ by \cref{continuous}.
            $f_1(p)\in\reals$ and $g(p)\in\reals$ and $(f_1(p), g(p))\in\reals\times\reals$
                by \cref{funs_type_apply,times_tuple_intro}.
            $\fst{k(p)} = f_1(p)$ and $\snd{k(p)} = g(p)$ by \cref{fst_eq,snd_eq}.
            Therefore $\apply{\realMultiplication}{k(p)} = f_1(p) \rmul g(p)$
                by \cref{continuous,real_multiplication_operation,real_multiplication_continuous,funs_is_function,function_apply_intro,real_mul_closed}.
            Then $m(p) = f_1(p) \rmul g(p)$.
            Take $a,b$ such that $p = (a,b)$ and $a\in\reals$ and $b\in\reals\setminus\{0\}$ by \cref{times_elem_is_tuple}.
            Hence $f_1(p) = a$ by \cref{projection_first,continuous,funs_is_function,function_apply_intro}.
            Then $f_2(p) = b$ by \cref{projection_second,continuous,funs_is_function,function_apply_intro}.
            $f_2\in\funs{\reals\times(\reals\setminus\{0\})}{\reals\setminus\{0\}}$ and
                $\realInverse\in\funs{\reals\setminus\{0\}}{\reals}$ by \cref{continuous}.
            $f_2$ is right-unique and $f_2$ is a relation and
                $\realInverse$ is right-unique and $\realInverse$ is a relation.
            Hence $\ran{f_2}\subseteq\dom{\realInverse}$ by \cref{funs_ran,funs,subseteq}.
            Hence $p\in\dom{f_2}$ by \cref{funs}.
            Therefore $g(p) = \apply{\realInverse}{f_2(p)}$ by \cref{circ_apply}.
            $\fst{p} = a$ and $\snd{p} = b$ by \cref{fst_eq,snd_eq}.
            Then $f_1(p) = \fst{p}$ and $f_2(p) = \snd{p}$.
            Hence $g(p) = \apply{\realInverse}{\snd{p}}$.
            $\snd{p}\in\reals\setminus\{0\}$.
            Hence $\apply{\realInverse}{\snd{p}} = \inv{\snd{p}}$
                by \cref{continuous,setminus_elim_left,setminus_elim_right,singleton_iff,real_mul_inverse,real_inverse_operation,funs_is_function,function_apply_intro}.
            Then $g(p) = \inv{\snd{p}}$.
            Therefore $m(p) = \fst{p} \rmul \inv{\snd{p}}$.
            Hence $m(p) = \apply{\realDivision}{p}$
                by \cref{real_mul_closed,real_mul_inverse,real_division_operation,funs_is_function,function_apply_intro}.
        \end{subproof}
        $m$ is a function and $\realDivision$ is a function and
            $\dom{m} = \reals\times(\reals\setminus\{0\})$ and
            $\dom{\realDivision} = \reals\times(\reals\setminus\{0\})$ by \cref{funs}.
        For all $x\in\dom{m}$ we have $m(x) = \apply{\realDivision}{x}$.
        Therefore $m\subseteq\realDivision$ by \cref{subseteq,fun_subseteq}.
        For all $x\in\dom{\realDivision}$ we have $m(x) = \apply{\realDivision}{x}$.
        Hence $\realDivision\subseteq m$ by \cref{subseteq,fun_subseteq}.
        Therefore $m = \realDivision$ by \cref{subseteq_antisymmetric}.
    Therefore $\realDivision$ is continuous from $\reals\times(\reals\setminus\{0\})$ to $\reals$
        with respect to $T$ and $\standardTopologyR$ by \cref{continuous_composite}.
\end{proof}
% from §21 Item 12: sum of continuous functions
\begin{theorem}\label{real_function_addition_continuous}
    Suppose $f$ is continuous from $X$ to $\reals$ with respect to $T$ and $\standardTopologyR$.
    Suppose $g$ is continuous from $X$ to $\reals$ with respect to $T$ and $\standardTopologyR$.
    Let $h = \{(x, f(x) + g(x)) \mid x\in X\}$.
    Then $h$ is continuous from $X$ to $\reals$ with respect to $T$ and $\standardTopologyR$.
\end{theorem}
\begin{proof}
    Let $T_1 = \standardTopologyR$.
    Let $T_2 = \productTopology{T_1}{T_1}{\reals}{\reals}$.
    Let $k = \{(x,(f(x), g(x))) \mid x\in X\}$.
    Hence $k$ is continuous from $X$ to $\reals\times\reals$ with respect to $T$ and $T_2$
        by \cref{continuous_map_into_product_backward}.
    $\realAddition$ is continuous from $\reals\times\reals$ to $\reals$ with respect to $T_2$ and $T_1$
        by \cref{real_addition_continuous}.
    Let $m = \realAddition\circ k$.
    Therefore $m$ is continuous from $X$ to $\reals$ with respect to $T$ and $T_1$ by \cref{continuous_composite}.
        For all $x,z_1,z_2$ such that $(x,z_1)\in h$ and $(x,z_2)\in h$ we have $z_1 = z_2$ and
            for all $w$ such that $w\in h$ there exist $x,y$ such that $w = (x,y)$.
        Hence $h$ is a function by \cref{rightunique,relation}.
        Show that for all $x\in\dom{h}$ we have $h(x)\in\reals$.
        \begin{subproof}
            Fix $x\in\dom{h}$.
            Take $z$ such that $(x,z)\in h$ by \cref{dom_iff}.
            Hence $x\in X$ and $z = f(x) + g(x)$.
            Therefore $h(x)\in\reals$ by \cref{function_apply_intro,continuous,funs_type_apply,real_add_closed}.
        \end{subproof}
        Hence $h$ is a function to $\reals$.
        For all $x\in\dom{h}$ we have $x\in X$ and for all $x\in X$ we have $(x, f(x) + g(x))\in h$.
        Therefore $\dom{h} = X$ by \cref{dom_intro,subseteq,subseteq_antisymmetric}.
        Therefore $h\in\funs{X}{\reals}$ by \cref{funs_intro}.
    Show that for all $x\in X$ we have $m(x) = \apply{h}{x}$.
    \begin{subproof}
        Fix $x\in X$.
        $k\in\funs{X}{\reals\times\reals}$ by \cref{continuous}.
        $k$ is a function by \cref{continuous,funs_is_function}.
        Hence $k(x) = (f(x), g(x))$ by \cref{continuous,funs_is_function,function_apply_intro}.
        Therefore $f(x)\in\reals$ and $g(x)\in\reals$ and $(f(x), g(x))\in\reals\times\reals$
            by \cref{continuous,funs_type_apply,times_tuple_intro}.
        Hence $\fst{k(x)} = f(x)$ and $\snd{k(x)} = g(x)$ by \cref{fst_eq,snd_eq}.
        $\realAddition$ is a function by \cref{continuous,real_addition_continuous,funs_is_function}.
        Therefore $\apply{\realAddition}{k(x)} = f(x) + g(x)$
            by \cref{continuous,real_addition_operation,real_addition_continuous,funs_is_function,function_apply_intro,real_add_closed}.
        Hence $m(x) = \apply{\realAddition}{k(x)}$
            by \cref{funs_ran,continuous,real_addition_continuous,funs,subseteq,circ_apply}.
        Hence $m(x) = \apply{h}{x}$ by \cref{funs_is_function,function_apply_intro}.
    \end{subproof}
    $m$ is a function and $h$ is a function and $\dom{m} = X$ and $\dom{h} = X$
        by \cref{funs,continuous,funs_is_function}.
    For all $x\in\dom{m}$ we have $m(x) = \apply{h}{x}$.
    Therefore $m\subseteq h$ by \cref{subseteq,fun_subseteq}.
    For all $x\in\dom{h}$ we have $m(x) = \apply{h}{x}$.
    Hence $h\subseteq m$ by \cref{subseteq,fun_subseteq}.
    Therefore $m = h$ by \cref{subseteq_antisymmetric}.
    Therefore $h$ is continuous from $X$ to $\reals$ with respect to $T$ and $\standardTopologyR$ by \cref{continuous_composite}.
\end{proof}

% from §21 Item 13: difference of continuous functions
\begin{theorem}\label{real_function_subtraction_continuous}
    Suppose $f$ is continuous from $X$ to $\reals$ with respect to $T$ and $\standardTopologyR$.
    Suppose $g$ is continuous from $X$ to $\reals$ with respect to $T$ and $\standardTopologyR$.
    Let $h = \{(x, f(x) - g(x)) \mid x\in X\}$.
    Then $h$ is continuous from $X$ to $\reals$ with respect to $T$ and $\standardTopologyR$.
\end{theorem}
\begin{proof}
    Let $T_1 = \standardTopologyR$.
    Let $T_2 = \productTopology{T_1}{T_1}{\reals}{\reals}$.
    Let $k = \{(x,(f(x), g(x))) \mid x\in X\}$.
    Hence $k$ is continuous from $X$ to $\reals\times\reals$ with respect to $T$ and $T_2$
        by \cref{continuous_map_into_product_backward}.
    $\realSubtraction$ is continuous from $\reals\times\reals$ to $\reals$ with respect to $T_2$ and $T_1$
        by \cref{real_subtraction_continuous}.
    Let $m = \realSubtraction\circ k$.
    Therefore $m$ is continuous from $X$ to $\reals$ with respect to $T$ and $T_1$ by \cref{continuous_composite}.
        For all $x,z_1,z_2$ such that $(x,z_1)\in h$ and $(x,z_2)\in h$ we have $z_1 = z_2$ and
            for all $w$ such that $w\in h$ there exist $x,y$ such that $w = (x,y)$.
        Hence $h$ is a function by \cref{rightunique,relation}.
        Show that for all $x\in\dom{h}$ we have $h(x)\in\reals$.
        \begin{subproof}
            Fix $x\in\dom{h}$.
            Take $z$ such that $(x,z)\in h$ by \cref{dom_iff}.
            Hence $x\in X$ and $z = f(x) - g(x)$.
            Therefore $h(x)\in\reals$
                by \cref{function_apply_intro,continuous,funs_type_apply,real_add_inverse,real_sub,real_add_closed}.
        \end{subproof}
        Hence $h$ is a function to $\reals$.
        For all $x\in\dom{h}$ we have $x\in X$ and for all $x\in X$ we have $(x, f(x) - g(x))\in h$.
        Therefore $\dom{h} = X$ by \cref{dom_intro,subseteq,subseteq_antisymmetric}.
        Therefore $h\in\funs{X}{\reals}$ by \cref{funs_intro}.
    Show that for all $x\in X$ we have $m(x) = \apply{h}{x}$.
    \begin{subproof}
        Fix $x\in X$.
        $k\in\funs{X}{\reals\times\reals}$ by \cref{continuous}.
        $k$ is a function by \cref{continuous,funs_is_function}.
        Hence $k(x) = (f(x), g(x))$ by \cref{continuous,funs_is_function,function_apply_intro}.
        Therefore $f(x)\in\reals$ and $g(x)\in\reals$ and $(f(x), g(x))\in\reals\times\reals$
            by \cref{continuous,funs_type_apply,times_tuple_intro}.
        Hence $\fst{k(x)} = f(x)$ and $\snd{k(x)} = g(x)$ by \cref{fst_eq,snd_eq}.
        $\realSubtraction$ is a function by \cref{continuous,real_subtraction_continuous,funs_is_function}.
        Therefore $\apply{\realSubtraction}{k(x)} = f(x) - g(x)$
            by \cref{continuous,real_subtraction_operation,real_subtraction_continuous,funs_is_function,function_apply_intro,real_add_inverse,real_sub,real_add_closed}.
        Hence $m(x) = \apply{\realSubtraction}{k(x)}$
            by \cref{funs_ran,continuous,real_subtraction_continuous,funs,subseteq,circ_apply}.
        Hence $m(x) = \apply{h}{x}$ by \cref{funs_is_function,function_apply_intro}.
    \end{subproof}
    $m$ is a function and $h$ is a function and $\dom{m} = X$ and $\dom{h} = X$
        by \cref{funs,continuous,funs_is_function}.
    For all $x\in\dom{m}$ we have $m(x) = \apply{h}{x}$.
    Therefore $m\subseteq h$ by \cref{subseteq,fun_subseteq}.
    For all $x\in\dom{h}$ we have $m(x) = \apply{h}{x}$.
    Hence $h\subseteq m$ by \cref{subseteq,fun_subseteq}.
    Therefore $m = h$ by \cref{subseteq_antisymmetric}.
    Therefore $h$ is continuous from $X$ to $\reals$ with respect to $T$ and $\standardTopologyR$ by \cref{continuous_composite}.
\end{proof}

% from §21 Item 14: product of continuous functions
\begin{theorem}\label{real_function_multiplication_continuous}
    Suppose $f$ is continuous from $X$ to $\reals$ with respect to $T$ and $\standardTopologyR$.
    Suppose $g$ is continuous from $X$ to $\reals$ with respect to $T$ and $\standardTopologyR$.
    Let $h = \{(x, f(x) \rmul g(x)) \mid x\in X\}$.
    Then $h$ is continuous from $X$ to $\reals$ with respect to $T$ and $\standardTopologyR$.
\end{theorem}
\begin{proof}
    Let $T_1 = \standardTopologyR$.
    Let $T_2 = \productTopology{T_1}{T_1}{\reals}{\reals}$.
    Let $k = \{(x,(f(x), g(x))) \mid x\in X\}$.
    Hence $k$ is continuous from $X$ to $\reals\times\reals$ with respect to $T$ and $T_2$
        by \cref{continuous_map_into_product_backward}.
    $\realMultiplication$ is continuous from $\reals\times\reals$ to $\reals$ with respect to $T_2$ and $T_1$
        by \cref{real_multiplication_continuous}.
    Let $m = \realMultiplication\circ k$.
    Therefore $m$ is continuous from $X$ to $\reals$ with respect to $T$ and $T_1$ by \cref{continuous_composite}.
    For all $x,z_1,z_2$ such that $(x,z_1)\in h$ and $(x,z_2)\in h$ we have $z_1 = z_2$ and
        for all $w$ such that $w\in h$ there exist $x,y$ such that $w = (x,y)$.
    Hence $h$ is a function by \cref{rightunique,relation}.
    Show that for all $x\in\dom{h}$ we have $h(x)\in\reals$.
    \begin{subproof}
        Fix $x\in\dom{h}$.
        Take $z$ such that $(x,z)\in h$ by \cref{dom_iff}.
        Hence $x\in X$ and $z = f(x) \rmul g(x)$.
        Therefore $h(x)\in\reals$ by \cref{function_apply_intro,continuous,funs_type_apply,real_mul_closed}.
    \end{subproof}
    Hence $h$ is a function to $\reals$.
    For all $x\in\dom{h}$ we have $x\in X$ and for all $x\in X$ we have $(x, f(x) \rmul g(x))\in h$.
    Hence $h\in\funs{X}{\reals}$ by \cref{subseteq,dom_intro,subseteq_antisymmetric,funs_intro}.
    Show that for all $x\in X$ we have $m(x) = \apply{h}{x}$.
    \begin{subproof}
        Fix $x\in X$.
        $k\in\funs{X}{\reals\times\reals}$ and $k$ is a function by \cref{continuous,funs_is_function}.
        Hence $k(x) = (f(x), g(x))$ by \cref{continuous,function_apply_intro}.
        Therefore $f(x)\in\reals$ and $g(x)\in\reals$ and $(f(x), g(x))\in\reals\times\reals$
            by \cref{continuous,funs_type_apply,times_tuple_intro}.
        Hence $\fst{k(x)} = f(x)$ and $\snd{k(x)} = g(x)$ by \cref{fst_eq,snd_eq}.
        $\realMultiplication$ is a function by \cref{continuous,real_multiplication_continuous,funs_is_function}.
        Therefore $\apply{\realMultiplication}{k(x)} = f(x) \rmul g(x)$
            by \cref{continuous,real_multiplication_operation,real_multiplication_continuous,funs_is_function,function_apply_intro,real_mul_closed}.
        Hence $m(x) = \apply{\realMultiplication}{k(x)}$
            by \cref{funs_ran,continuous,real_multiplication_continuous,funs,subseteq,circ_apply}.
        Hence $m(x) = \apply{h}{x}$ by \cref{funs_is_function,function_apply_intro}.
    \end{subproof}
    $m$ is a function and $h$ is a function and $\dom{m} = X$ and $\dom{h} = X$
        by \cref{funs,continuous,funs_is_function}.
    For all $x\in\dom{m}$ we have $m(x) = \apply{h}{x}$.
    Therefore $m\subseteq h$ by \cref{subseteq,fun_subseteq}.
    For all $x\in\dom{h}$ we have $m(x) = \apply{h}{x}$.
    Hence $h\subseteq m$ by \cref{subseteq,fun_subseteq}.
    Therefore $m = h$ by \cref{subseteq_antisymmetric}.
    Therefore $h$ is continuous from $X$ to $\reals$ with respect to $T$ and $\standardTopologyR$ by \cref{continuous_composite}.
\end{proof}

% from §21 Item 15: quotient of continuous functions
\begin{theorem}\label{real_function_division_continuous}
    Suppose $f$ is continuous from $X$ to $\reals$ with respect to $T$ and $\standardTopologyR$.
    Suppose $g$ is continuous from $X$ to $\reals$ with respect to $T$ and $\standardTopologyR$.
    Suppose for all $x\in X$ we have $g(x)\neq 0$.
    Let $h = \{(x, f(x) \rmul \inv{(g(x))}) \mid x\in X\}$.
    Then $h$ is continuous from $X$ to $\reals$ with respect to $T$ and $\standardTopologyR$.
\end{theorem}
\begin{proof}
    Let $T_1 = \standardTopologyR$.
    Let $T_2 = \subspaceTopology{T_1}{\reals\setminus\{0\}}$.
    For all $y$ such that $y\in\img{g}{X}$ we have $y\in\reals\setminus\{0\}$
        by \cref{img_iff,continuous,funs_is_function,function_apply_intro,continuous,funs_type_apply,singleton_iff,setminus_intro}.
    Hence $\img{g}{X}\subseteq \reals\setminus\{0\}$ by \cref{subseteq}.
    $\standardBasisR$ is a basis on $\reals$ by \cref{standard_basis_is_basis}.
    Therefore $T_1$ is a topology on $\reals$ by \cref{basis_topology_is_topology,standard_topology_reals}.
    Hence $\reals\setminus\{0\}$ is a subspace of $\reals$ with respect to $T_1$
        by \cref{setminus_elim_left,subseteq,subspace_of}.
    Therefore $g$ is continuous from $X$ to $\reals\setminus\{0\}$ with respect to $T$ and $T_2$
        by \cref{continuous_restrict_range}.
    $\realInverse$ is continuous from $\reals\setminus\{0\}$ to $\reals$ with respect to $T_2$ and $T_1$
        by \cref{real_inverse_continuous}.
    Let $j = \realInverse\circ g$.
    Hence $j$ is continuous from $X$ to $\reals$ with respect to $T$ and $T_1$ by \cref{continuous_composite}.
    Let $k = \{(x,(f(x), j(x))) \mid x\in X\}$.
    Therefore $k$ is continuous from $X$ to $\reals\times\reals$ with respect to $T$ and
        $\productTopology{T_1}{T_1}{\reals}{\reals}$ by \cref{continuous_map_into_product_backward}.
    $\realMultiplication$ is continuous from $\reals\times\reals$ to $\reals$ with respect to
        $\productTopology{T_1}{T_1}{\reals}{\reals}$ and $T_1$ by \cref{real_multiplication_continuous}.
    Let $m = \realMultiplication\circ k$.
    Hence $m$ is continuous from $X$ to $\reals$ with respect to $T$ and $T_1$ by \cref{continuous_composite}.
    For all $x,z_1,z_2$ such that $(x,z_1)\in h$ and $(x,z_2)\in h$ we have $z_1 = z_2$ and
        for all $w$ such that $w\in h$ there exist $x,y$ such that $w = (x,y)$.
    Hence $h$ is a function by \cref{rightunique,relation}.
    Show that for all $x\in\dom{h}$ we have $h(x)\in\reals$.
    \begin{subproof}
        Fix $x\in\dom{h}$.
        Take $z$ such that $(x,z)\in h$ by \cref{dom_iff}.
        Take $x_0\in X$ such that $(x,z) = (x_0, f(x_0) \rmul \inv{(g(x_0))})$.
        Hence $x\in X$ and $z = f(x) \rmul \inv{(g(x))}$ by \cref{pair_eq_iff}.
        Therefore $h(x)\in\reals$ by \cref{function_apply_intro,continuous,funs_type_apply,real_mul_inverse,real_mul_closed}.
    \end{subproof}
    Hence $h$ is a function to $\reals$.
    For all $x\in\dom{h}$ we have $x\in X$.
    For all $x\in X$ we have $(x, f(x) \rmul \inv{(g(x))})\in h$.
    Hence $\dom{h} = X$ by \cref{subseteq,dom_intro,subseteq_antisymmetric}.
    Therefore $h\in\funs{X}{\reals}$ by \cref{funs_intro}.
    Show that for all $x\in X$ we have $m(x) = \apply{h}{x}$.
    \begin{subproof}
        Fix $x\in X$.
        $k\in\funs{X}{\reals\times\reals}$ and $k$ is a function by \cref{continuous,funs_is_function}.
        Hence $k(x) = (f(x), j(x))$ by \cref{continuous,function_apply_intro}.
        Therefore $f(x)\in\reals$ and $j(x)\in\reals$ and $(f(x), j(x))\in\reals\times\reals$
            by \cref{continuous,funs_type_apply,times_tuple_intro}.
        Hence $\fst{k(x)} = f(x)$ and $\snd{k(x)} = j(x)$ by \cref{fst_eq,snd_eq}.
        $\realMultiplication$ is a function by \cref{continuous,real_multiplication_continuous,funs_is_function}.
        Therefore $\apply{\realMultiplication}{k(x)} = f(x) \rmul j(x)$
            by \cref{continuous,real_multiplication_operation,real_multiplication_continuous,funs_is_function,function_apply_intro,real_mul_closed}.
        Hence $m(x) = \apply{\realMultiplication}{k(x)}$
            by \cref{funs_ran,continuous,real_multiplication_continuous,funs,subseteq,circ_apply}.
        $g\in\funs{X}{\reals\setminus\{0\}}$ by \cref{continuous}.
        Hence $g(x)\in\reals\setminus\{0\}$ by \cref{continuous,funs_type_apply}.
        $\realInverse\in\funs{\reals\setminus\{0\}}{\reals}$ and
            $\dom{\realInverse} = \reals\setminus\{0\}$ by \cref{continuous,real_inverse_continuous,funs}.
        Hence $\ran{g}\subseteq\dom{\realInverse}$ by \cref{funs_ran,subseteq}.
        $g$ is a function and $\realInverse$ is a function by \cref{funs_is_function}.
        Then $x\in\dom{g}$ by \cref{funs}.
        Therefore $j(x) = \apply{\realInverse}{g(x)}$ by \cref{circ_apply}.
        Hence $g(x)\neq 0$ and $g(x)\in\reals$ and $\inv{(g(x))}\in\reals$
            by \cref{singleton_iff,setminus_elim_left,setminus_elim_right,real_mul_inverse}.
        Therefore $\apply{\realInverse}{g(x)} = \inv{(g(x))}$
            by \cref{real_inverse_operation,funs_is_function,function_apply_intro}.
        Hence $m(x) = \apply{h}{x}$ by \cref{funs_is_function,function_apply_intro}.
    \end{subproof}
    $m$ is a function and $h$ is a function and $\dom{m} = X$ and $\dom{h} = X$
        by \cref{funs,continuous,funs_is_function}.
    For all $x\in\dom{m}$ we have $m(x) = \apply{h}{x}$.
    Therefore $m\subseteq h$ by \cref{subseteq,fun_subseteq}.
    For all $x\in\dom{h}$ we have $m(x) = \apply{h}{x}$.
    Hence $h\subseteq m$ by \cref{subseteq,fun_subseteq}.
    Therefore $m = h$ by \cref{subseteq_antisymmetric}.
    Therefore $h$ is continuous from $X$ to $\reals$ with respect to $T$ and $\standardTopologyR$ by \cref{continuous_composite}.
\end{proof}

% from §21 Item 16: uniform convergence
\begin{definition}\label{uniform_converges_to}
    $s$ converges uniformly to $f$ from $X$ to $Y$ with respect to $d$ iff
        $s$ is a sequence in $\funs{X}{Y}$ and
        $f\in\funs{X}{Y}$ and
        for all $\epsilon\in\reals$ such that $0 < \epsilon$
        there exists $N\in\naturalsPlus$ such that for all $n\in\naturalsPlus$ if $N\leq n$ then
            for all $x\in X$ we have $d((\apply{\apply{s}{n}}{x}, f(x))) < \epsilon$.
\end{definition}

% from §21 Item 17: uniform limit theorem
\begin{theorem}\label{uniform_limit_theorem}
    Suppose $T_X$ is a topology on $X$.
    Suppose $s$ is a sequence in $\funs{X}{Y}$.
    Suppose $f\in\funs{X}{Y}$.
    Suppose $d$ is a metric on $Y$.
    Let $T = \metricTopology{d}{Y}$.
    Suppose for all $n\in\naturalsPlus$ we have
        $s(n)$ is continuous from $X$ to $Y$ with respect to $T_X$ and $T$.
    Suppose $s$ converges uniformly to $f$ from $X$ to $Y$ with respect to $d$.
    Then $f$ is continuous from $X$ to $Y$ with respect to $T_X$ and $T$.
\end{theorem}
\begin{proof}
    $f$ is a function from $X$ to $Y$ by \cref{funs}.
    $d$ is a metric on $Y$.
    $\metricBasis{d}{Y}$ is a basis on $Y$ by \cref{metric_basis_is_basis}.
    Hence $T$ is a topology on $Y$ by \cref{basis_topology_is_topology,metric_topology,basis_for_topology}.
    Show that for all $V$ such that $V$ is open in $Y$ with respect to $T$ we have
        $\preimg{f}{V}$ is open in $X$ with respect to $T_X$.
    \begin{subproof}
        Fix $V$ such that $V$ is open in $Y$ with respect to $T$.
        Therefore $V\subseteq Y$ by \cref{open_in,topology_on,subseteq,pow_iff}.
        Show that for all $x\in\preimg{f}{V}$ there exists $U$ such that
            $U$ is open in $X$ with respect to $T_X$ and $x\in U$ and $U\subseteq\preimg{f}{V}$.
        \begin{subproof}
            Fix $x$ such that $x\in\preimg{f}{V}$.
            Take $y\in V$ such that $(x,y)\in f$ by \cref{preimg_iff}.
            Hence $f(x)\in V$ by \cref{funs_is_function,function_apply_iff}.
            Take $\epsilon\in\reals$ such that $0 < \epsilon$ and
                $\metricBall{d}{Y}{f(x)}{\epsilon} \subseteq V$ by \cref{metric_open_iff}.
            Let $\epsilon_1 = \epsilon / (1 + 1)$.
            $0 < 1$ and $1\in\reals$ by \cref{real_zero_lt_one,real_naturals_subset_reals,real_one_in_naturals,subseteq}.
            Therefore $1$ is positive.
            Hence $1 + 1$ is positive by \cref{real_pos_add}.
            Therefore $1 + 1\in\reals$ and $1 + 1 \neq 0$ by \cref{real_add_closed,real_pos_nonzero}.
            Hence $\inv{(1 + 1)}\in\reals$ by \cref{real_mul_inverse}.
            Hence $0 < \inv{(1 + 1)}$ by \cref{real_inv_pos}.
            Therefore $\epsilon \rmul \inv{(1 + 1)}$ is positive by \cref{real_mul_pos_pos}.
            $\epsilon / (1 + 1) = \epsilon \rmul \inv{(1 + 1)}$ by \cref{real_div}.
            Hence $0 < \epsilon_1$ by \cref{real_lt_subst_right}.
            Therefore $\epsilon_1\in\reals$ by \cref{real_div,real_mul_closed}.
            Hence $s\in\funs{\naturalsPlus}{\funs{X}{Y}}$ by \cref{sequence_in}.
            Take $N\in\naturalsPlus$ such that for all $n\in\naturalsPlus$ if $N\leq n$ then
                for all $z\in X$ we have $d((\apply{\apply{s}{n}}{z}, f(z))) < \epsilon_1$
                by \cref{uniform_converges_to}.
            Let $g = s(N)$.
            Therefore $g\in\funs{X}{Y}$ by \cref{funs_type_apply}.
            Hence $g$ is a function and $f$ is a function by \cref{funs_is_function}.
            $N\leq N$.
            Therefore for all $z\in X$ we have $d((g(z), f(z))) < \epsilon_1$.
            $N\in\naturalsPlus$.
            Hence $g$ is continuous from $X$ to $Y$ with respect to $T_X$ and $T$.
            Let $W = \metricBall{d}{Y}{f(x)}{\epsilon_1}$.
            Therefore $W\in\pow{Y}$ by \cref{subseteq,powerset_intro,metric_ball}.
            Hence $x\in X$ by \cref{preimg_subseteq_dom,elem_subseteq,funs}.
            Hence $f(x)\in Y$ by \cref{preimg_subseteq_dom,elem_subseteq,funs,funs_type_apply}.
            Therefore $W$ is open in $Y$ with respect to $T$
                by \cref{metric_basis,metric_basis_is_basis,basis_elem_open_in_generated,metric_topology,basis_for_topology,open_in}.
            Let $U = \preimg{g}{W}$.
            Hence $U$ is open in $X$ with respect to $T_X$ by \cref{continuous}.
            $g(x)\in Y$ and $f(x)\in Y$ by \cref{funs_type_apply}.
            Therefore $d((f(x), g(x))) = d((g(x), f(x)))$ by \cref{metric_on,metric_symmetric}.
            Hence $d((f(x), g(x))) < \epsilon_1$ by \cref{real_lt_subst_left}.
            Therefore $g(x)\in W$ by \cref{metric_ball}.
            Hence $x\in U$ by \cref{funs,function_apply_iff,preimg_iff}.
            Show that for all $y$ such that $y\in U$ we have $y\in\preimg{f}{V}$.
            \begin{subproof}
                Fix $y$ such that $y\in U$.
                Take $z\in W$ such that $(y,z)\in g$ by \cref{preimg_iff}.
                Hence $g(y)\in Y$ and $d((f(x), g(y))) < \epsilon_1$ by \cref{function_apply_iff,metric_ball}.
                Therefore $d((g(y), f(x))) < \epsilon_1$
                    by \cref{subseteq,metric_on,metric_symmetric,real_lt_subst_left}.
                Therefore $y\in X$ by \cref{preimg_subseteq_dom,elem_subseteq,funs}.
                Hence $d((g(y), f(y))) < \epsilon_1$.
                $f(y)\in Y$ by \cref{funs_type_apply}.
                Therefore $d((f(y), g(y))) < \epsilon_1$
                    by \cref{metric_on,metric_symmetric,real_lt_subst_left}.
                $d\in\funs{Y\times Y}{\reals}$ by \cref{metric_on}.
                $(f(y), g(y))\in Y\times Y$ and $(g(y), f(x))\in Y\times Y$ and $(f(y), f(x))\in Y\times Y$
                    by \cref{times_tuple_intro}.
                Thus $d((f(y), g(y)))\in\reals$ and $d((g(y), f(x)))\in\reals$ and $d((f(y), f(x)))\in\reals$
                    by \cref{funs_type_apply}.
                Therefore $d((f(y), g(y))) + d((g(y), f(x))) \geq d((f(y), f(x)))$
                    by \cref{metric_on,metric_triangle_inequality}.
                Hence $d((f(y), f(x))) \leq d((f(y), g(y))) + d((g(y), f(x)))$.
                Therefore $d((f(y), g(y))) + d((g(y), f(x))) < \epsilon_1 + \epsilon_1$
                    by \cref{real_add_closed,real_lt_add_two}.
                $1 \rmul \epsilon_1 = \epsilon_1$ by \cref{real_mul_ident}.
                $(1 + 1) \rmul \epsilon_1 = (1 \rmul \epsilon_1) + (1 \rmul \epsilon_1)$ by \cref{real_right_distrib}.
                Hence $(1 \rmul \epsilon_1) + (1 \rmul \epsilon_1) = \epsilon_1 + (1 \rmul \epsilon_1)$ by \cref{real_add_eq_subst}.
                Therefore $(1 \rmul \epsilon_1) + (1 \rmul \epsilon_1) = \epsilon_1 + \epsilon_1$ by \cref{real_add_eq_subst}.
                Then $(1 + 1) \rmul \epsilon_1 = \epsilon_1 + \epsilon_1$.
                $\epsilon_1 = \epsilon \rmul \inv{(1 + 1)}$ by \cref{real_div}.
                Hence $(1 + 1) \rmul \epsilon_1 = (1 + 1) \rmul (\epsilon \rmul \inv{(1 + 1)})$.
                $(1 + 1) \rmul (\epsilon \rmul \inv{(1 + 1)}) = \epsilon \rmul ((1 + 1) \rmul \inv{(1 + 1)})$ by \cref{real_mul_assoc,real_mul_comm}.
                $(1 + 1) \rmul \inv{(1 + 1)} = 1$ by \cref{real_mul_inverse}.
                Therefore $\epsilon \rmul 1 = \epsilon$ by \cref{real_mul_comm,real_mul_ident}.
                Hence $\epsilon_1 + \epsilon_1 = \epsilon$.
                Therefore $d((f(y), g(y))) + d((g(y), f(x))) < \epsilon$ by \cref{real_lt_subst_right}.
                $d((f(y), f(x))) < \epsilon$.
                Hence $d((f(x), f(y))) = d((f(y), f(x)))$ by \cref{metric_on,metric_symmetric}.
                Therefore $d((f(x), f(y))) < \epsilon$ by \cref{real_lt_subst_left}.
                Hence $f(y)\in\metricBall{d}{Y}{f(x)}{\epsilon}$ by \cref{metric_ball}.
                Therefore $f(y)\in V$ by \cref{subseteq}.
                Hence $y\in\preimg{f}{V}$ by \cref{funs,function_apply_iff,preimg_iff}.
            \end{subproof}
            Therefore $U\subseteq\preimg{f}{V}$ by \cref{subseteq}.
            Hence there exists $U$ such that
                $U$ is open in $X$ with respect to $T_X$ and $x\in U$ and $U\subseteq\preimg{f}{V}$.
        \end{subproof}
        Let $M = \{U\in T_X \mid U\subseteq\preimg{f}{V}\}$.
        Therefore $\unions{M}\in T_X$ by \cref{subseteq,topology_on}.
        For all $x$ such that $x\in\preimg{f}{V}$ we have $x\in\unions{M}$.
        Hence $\preimg{f}{V}\subseteq\unions{M}$ by \cref{subseteq}.
        For all $x$ such that $x\in\unions{M}$ we have $x\in\preimg{f}{V}$.
        Therefore $\unions{M}\subseteq\preimg{f}{V}$ by \cref{subseteq}.
        Hence $\preimg{f}{V} = \unions{M}$ by \cref{subseteq_antisymmetric}.
        Therefore $\preimg{f}{V}\in T_X$.
        Hence $\preimg{f}{V}$ is open in $X$ with respect to $T_X$ by \cref{open_in}.
    \end{subproof}
    Therefore $f$ is continuous from $X$ to $Y$ with respect to $T_X$ and $T$ by \cref{continuous}.
\end{proof}
\section{§22 The Quotient Topology}

% from §22 Item 1: quotient map
\begin{definition}\label{quotient_map}
    $p$ is a quotient map from $X$ to $Y$ with respect to $T_X$ and $T_Y$ iff
        $p$ is a surjection from $X$ to $Y$ and
        $p$ is continuous from $X$ to $Y$ with respect to $T_X$ and $T_Y$ and
        for all $U$ such that $U\subseteq Y$ if $\preimg{p}{U}$ is open in $X$ with respect to $T_X$ then
            $U$ is open in $Y$ with respect to $T_Y$.
\end{definition}

% from §22 Item 2: quotient map closed set characterization
\begin{lemma}\label{quotient_map_closed_iff}
    Suppose $p$ is a surjection from $X$ to $Y$.
    Suppose $T_X$ is a topology on $X$.
    Suppose $T_Y$ is a topology on $Y$.
    Then $p$ is a quotient map from $X$ to $Y$ with respect to $T_X$ and $T_Y$ iff
        for all $A$ such that $A\subseteq Y$ we have
            $A$ is closed in $Y$ with respect to $T_Y$ iff
            $\preimg{p}{A}$ is closed in $X$ with respect to $T_X$.
\end{lemma}
\begin{proof}
    Hence $p$ is a function from $X$ to $Y$ by \cref{surj,funs}.
    Show that if $p$ is a quotient map from $X$ to $Y$ with respect to $T_X$ and $T_Y$ then
        for all $A$ such that $A\subseteq Y$ we have
            $A$ is closed in $Y$ with respect to $T_Y$ iff
            $\preimg{p}{A}$ is closed in $X$ with respect to $T_X$.
    \begin{subproof}
        Assume $p$ is a quotient map from $X$ to $Y$ with respect to $T_X$ and $T_Y$.
        For all $A$ such that $A\subseteq Y$ we have
            $A$ is closed in $Y$ with respect to $T_Y$ iff
            $\preimg{p}{A}$ is closed in $X$ with respect to $T_X$ by
            \cref{closed_in,setminus_subseteq,quotient_map,continuous,preimg_setminus_function,preimg_subseteq_dom,funs,subseteq}.
    \end{subproof}
    Show that if for all $A$ such that $A\subseteq Y$ we have
            $A$ is closed in $Y$ with respect to $T_Y$ iff
            $\preimg{p}{A}$ is closed in $X$ with respect to $T_X$
        then $p$ is a quotient map from $X$ to $Y$ with respect to $T_X$ and $T_Y$.
    \begin{subproof}
        Assume that for all $A$ such that $A\subseteq Y$ we have
            $A$ is closed in $Y$ with respect to $T_Y$ iff
            $\preimg{p}{A}$ is closed in $X$ with respect to $T_X$.
        For all $B$ such that $B$ is closed in $Y$ with respect to $T_Y$ we have
            $\preimg{p}{B}$ is closed in $X$ with respect to $T_X$.
        Hence $p$ is continuous from $X$ to $Y$ with respect to $T_X$ and $T_Y$
            by \cref{preimg_closed_implies_continuous}.
        For all $U$ such that $U\subseteq Y$ we have
            $U$ is open in $Y$ with respect to $T_Y$ iff
            $\preimg{p}{U}$ is open in $X$ with respect to $T_X$ by
            \cref{continuous,preimg_subseteq_dom,funs,subseteq,setminus_subseteq,double_relative_complement,open_in,closed_in,preimg_setminus_function}.
        Therefore $p$ is a quotient map from $X$ to $Y$ with respect to $T_X$ and $T_Y$ by \cref{quotient_map}.
    \end{subproof}
\end{proof}

% from §22 Item 3: saturated subset
\begin{definition}\label{saturated}
    $C$ is saturated with respect to $p$ iff
        $C\subseteq\dom{p}$ and
        for all $y\in\ran{p}$ if $C$ intersects $\preimg{p}{\{y\}}$ then
            $\preimg{p}{\{y\}}\subseteq C$.
\end{definition}

% from §22 Item 4: saturated sets are complete inverse images
\begin{lemma}\label{saturated_iff_preimg}
    Suppose $p$ is a function.
    Then $C$ is saturated with respect to $p$ iff there exists $B$ such that
        $C = \preimg{p}{B}$.
\end{lemma}
\begin{proof}
    Show that if $C$ is saturated with respect to $p$ then there exists $B$ such that
        $C = \preimg{p}{B}$.
    \begin{subproof}
        Assume $C$ is saturated with respect to $p$.
        Let $B = \img{p}{C}$.
        For all $x$ such that $x\in C$ we have $x\in\preimg{p}{B}$
            by \cref{saturated,subseteq,img_of_function_intro,inter_intro,function_apply_elim,preimg_iff}.
        Hence $C\subseteq \preimg{p}{B}$ by \cref{subseteq}.
        For all $x$ such that $x\in\preimg{p}{B}$ we have $x\in C$ by
            \cref{preimg_iff,img_of_function_elim,inter_elim_right,inter_elim_left,function_apply_iff,singleton_iff,inter_intro,emptyset,intersects,ran_iff,saturated,subseteq}.
        Hence $\preimg{p}{B}\subseteq C$ by \cref{subseteq}.
        Therefore $C = \preimg{p}{B}$ by \cref{subseteq_antisymmetric}.
    \end{subproof}
    Show that if there exists $B$ such that $C = \preimg{p}{B}$ then
        $C$ is saturated with respect to $p$.
    \begin{subproof}
        Assume that there exists $B$ such that $C = \preimg{p}{B}$.
        Take $B$ such that $C = \preimg{p}{B}$.
        Hence $C\subseteq\dom{p}$ by \cref{preimg_subseteq_dom,subseteq}.
        Show that for all $y\in\ran{p}$ if $C$ intersects $\preimg{p}{\{y\}}$ then
            $\preimg{p}{\{y\}}\subseteq C$.
        \begin{subproof}
            Fix $y\in\ran{p}$.
            Assume $C$ intersects $\preimg{p}{\{y\}}$.
            Take $c$ such that $c\in C\inter \preimg{p}{\{y\}}$ by \cref{intersects,nonempty_inhabited}.
            Take $b\in B$ such that $c\mathrel{p} b$ by \cref{preimg_iff,subseteq,inter_elim_left}.
            Therefore $c\mathrel{p} y$ by \cref{preimg_iff,singleton_iff,inter_elim_right}.
            Hence $b = y$ by \cref{function_rightunique,rightunique}.
            Hence $\preimg{p}{\{y\}}\subseteq C$ by \cref{function_rightunique,rightunique,preimg_iff,singleton_iff,subseteq}.
        \end{subproof}
        Therefore $C$ is saturated with respect to $p$ by \cref{saturated}.
    \end{subproof}
\end{proof}

% from §22 helper lemma 1: image of a preimage for surjections
\begin{lemma}\label{img_preimg_surjection}
    Suppose $p$ is a surjection from $X$ to $Y$.
    Suppose $B\subseteq Y$.
    Then $\img{p}{\preimg{p}{B}} = B$.
\end{lemma}
\begin{proof}
    $p$ is a function and $\img{p}{\preimg{p}{B}}\subseteq B$
        by \cref{surj,funs_is_function,img_preimg_subseteq}.
    Hence $\img{p}{\preimg{p}{B}} = B$
        by \cref{subseteq,surj,funs_is_function,img_preimg_subseteq,funs,function_apply_iff,preimg_iff,img_iff,subseteq_antisymmetric}.
\end{proof}

% from §22 Item 5: quotient maps and saturated open sets
\begin{theorem}\label{quotient_map_saturated_open_iff}
    Suppose $p$ is a surjection from $X$ to $Y$.
    Suppose $T_X$ is a topology on $X$.
    Suppose $T_Y$ is a topology on $Y$.
    Then $p$ is a quotient map from $X$ to $Y$ with respect to $T_X$ and $T_Y$ iff
        $p$ is continuous from $X$ to $Y$ with respect to $T_X$ and $T_Y$ and
        for all $C$ such that
            $C$ is saturated with respect to $p$ and
            $C$ is open in $X$ with respect to $T_X$
        we have $\img{p}{C}$ is open in $Y$ with respect to $T_Y$.
\end{theorem}
\begin{proof}
    Show that if $p$ is a quotient map from $X$ to $Y$ with respect to $T_X$ and $T_Y$ then
        $p$ is continuous from $X$ to $Y$ with respect to $T_X$ and $T_Y$ and
        for all $C$ such that
            $C$ is saturated with respect to $p$ and
            $C$ is open in $X$ with respect to $T_X$
        we have $\img{p}{C}$ is open in $Y$ with respect to $T_Y$.
    \begin{subproof}
        Assume $p$ is a quotient map from $X$ to $Y$ with respect to $T_X$ and $T_Y$.
        Show that for all $C$ such that
            $C$ is saturated with respect to $p$ and
            $C$ is open in $X$ with respect to $T_X$
        we have $\img{p}{C}$ is open in $Y$ with respect to $T_Y$.
        \begin{subproof}
            Fix $C$ such that
                $C$ is saturated with respect to $p$ and
                $C$ is open in $X$ with respect to $T_X$.
            Hence $p\in\funs{X}{Y}$ and $p$ is a function by \cref{surj,funs_is_function}.
            Take $B$ such that $C = \preimg{p}{B}$ by \cref{saturated_iff_preimg}.
            Let $B_1 = B\inter Y$.
            Hence $\preimg{p}{B_1} = \preimg{p}{B}$ by
                \cref{inter_lower_left,preimg_subseteq,preimg_iff,ran_iff,surj,funs_ran,subseteq,inter_intro,subseteq_antisymmetric}.
            Hence $\img{p}{C}$ is open in $Y$ with respect to $T_Y$
                by \cref{subseteq,quotient_map,inter_lower_right,img_preimg_surjection,open_in}.
        \end{subproof}
    \end{subproof}
    Show that if
        $p$ is continuous from $X$ to $Y$ with respect to $T_X$ and $T_Y$ and
        for all $C$ such that
            $C$ is saturated with respect to $p$ and
            $C$ is open in $X$ with respect to $T_X$
        we have $\img{p}{C}$ is open in $Y$ with respect to $T_Y$
    then $p$ is a quotient map from $X$ to $Y$ with respect to $T_X$ and $T_Y$.
    \begin{subproof}
        Assume $p$ is continuous from $X$ to $Y$ with respect to $T_X$ and $T_Y$.
        Assume that for all $C$ such that
            $C$ is saturated with respect to $p$ and
            $C$ is open in $X$ with respect to $T_X$
        we have $\img{p}{C}$ is open in $Y$ with respect to $T_Y$.
        Show that for all $U$ such that $U\subseteq Y$ we have
            $U$ is open in $Y$ with respect to $T_Y$ iff
            $\preimg{p}{U}$ is open in $X$ with respect to $T_X$.
        \begin{subproof}
            Fix $U$ such that $U\subseteq Y$.
            If $U$ is open in $Y$ with respect to $T_Y$ then
                $\preimg{p}{U}$ is open in $X$ with respect to $T_X$ by \cref{continuous}.
            If $\preimg{p}{U}$ is open in $X$ with respect to $T_X$ then
                $\img{p}{\preimg{p}{U}}$ is open in $Y$ with respect to $T_Y$.
            Therefore if $\preimg{p}{U}$ is open in $X$ with respect to $T_X$ then
                $U$ is open in $Y$ with respect to $T_Y$ by \cref{img_preimg_surjection}.
            Hence $U$ is open in $Y$ with respect to $T_Y$ iff
                $\preimg{p}{U}$ is open in $X$ with respect to $T_X$.
        \end{subproof}
        Therefore $p$ is a quotient map from $X$ to $Y$ with respect to $T_X$ and $T_Y$ by \cref{quotient_map}.
    \end{subproof}
\end{proof}

% from §22 Item 6: open map
\begin{definition}\label{open_map}
    $f$ is an open map from $X$ to $Y$ with respect to $T$ and $T_1$ iff
        $f$ is a function from $X$ to $Y$ and
        for all $U$ such that $U$ is open in $X$ with respect to $T$ we have
            $\img{f}{U}$ is open in $Y$ with respect to $T_1$.
\end{definition}

% from §22 Item 7: closed map
\begin{definition}\label{closed_map}
    $f$ is a closed map from $X$ to $Y$ with respect to $T$ and $T_1$ iff
        $f$ is a function from $X$ to $Y$ and
        for all $A$ such that $A$ is closed in $X$ with respect to $T$ we have
            $\img{f}{A}$ is closed in $Y$ with respect to $T_1$.
\end{definition}

% from §22 Item 8: open or closed surjections are quotient maps
\begin{lemma}\label{open_or_closed_surjection_quotient}
    Suppose $p$ is a surjection from $X$ to $Y$.
    Suppose $T_X$ is a topology on $X$.
    Suppose $T_Y$ is a topology on $Y$.
    Suppose $p$ is continuous from $X$ to $Y$ with respect to $T_X$ and $T_Y$.
    Suppose $p$ is an open map from $X$ to $Y$ with respect to $T_X$ and $T_Y$
        or $p$ is a closed map from $X$ to $Y$ with respect to $T_X$ and $T_Y$.
    Then $p$ is a quotient map from $X$ to $Y$ with respect to $T_X$ and $T_Y$.
\end{lemma}
\begin{proof}
    \begin{byCase}
    \caseOf{$p$ is an open map from $X$ to $Y$ with respect to $T_X$ and $T_Y$.}
        For all $U$ such that $U\subseteq Y$ we have
            $U$ is open in $Y$ with respect to $T_Y$ iff
            $\preimg{p}{U}$ is open in $X$ with respect to $T_X$
            by \cref{continuous,open_map,img_preimg_surjection}.
        Hence $p$ is a quotient map from $X$ to $Y$ with respect to $T_X$ and $T_Y$ by \cref{quotient_map}.
    \caseOf{$p$ is a closed map from $X$ to $Y$ with respect to $T_X$ and $T_Y$.}
        For all $A$ such that $A\subseteq Y$ we have
            $A$ is closed in $Y$ with respect to $T_Y$ iff
            $\preimg{p}{A}$ is closed in $X$ with respect to $T_X$ by
            \cref{closed_in,continuous,preimg_setminus_function,preimg_subseteq_dom,funs,subseteq,closed_map,img_preimg_surjection}.
        Therefore $p$ is a quotient map from $X$ to $Y$ with respect to $T_X$ and $T_Y$
            by \cref{quotient_map_closed_iff}.
    \end{byCase}
\end{proof}

% from §22 Item 9: quotient topology
\begin{definition}\label{quotient_topology}
    $\quotientTopology{p}{T}{X}{A}
        = \{U\in\pow{A} \mid \text{$\preimg{p}{U}$ is open in $X$ with respect to $T$}\}$.
\end{definition}

% from §22 Item 10: quotient topology is a topology
\begin{lemma}\label{quotient_topology_is_topology}
    Assume $p$ is a function from $X$ to $A$.
    Assume $T$ is a topology on $X$.
    Then $\quotientTopology{p}{T}{X}{A}$ is a topology on $A$.
\end{lemma}
\begin{proof}
    Let $T_Q = \quotientTopology{p}{T}{X}{A}$.
    $p\in\funs{X}{A}$ and $\dom{p} = X$ and $p$ is a function by \cref{funs,funs_is_function}.
    $T_Q\subseteq\pow{A}$ by \cref{quotient_topology,subseteq}.
    $\preimg{p}{\emptyset}\subseteq\emptyset$ by \cref{preimg_iff,emptyset,subseteq}.
    Therefore $\emptyset\in T_Q$
        by \cref{emptyset_subseteq,subseteq_antisymmetric,topology_on,open_in,pow_iff,quotient_topology}.
    $\preimg{p}{A} = X$ by
        \cref{preimg_subseteq_dom,funs,subseteq,dom_iff,ran_iff,funs_ran,subseteq_antisymmetric,preimg_iff}.
    Therefore $A\in T_Q$ by \cref{topology_on,open_in,subseteq_refl,pow_iff,quotient_topology}.
    Show that for all $M$ such that $M\subseteq T_Q$ we have $\unions{M}\in T_Q$.
    \begin{subproof}
        Fix $M$ such that $M\subseteq T_Q$.
        Let $N = \{V\in T \mid \text{there exists $U\in M$ such that $V = \preimg{p}{U}$}\}$.
        Now $\unions{N}\in T$ by \cref{subseteq,topology_on}.
        Hence $\preimg{p}{\unions{M}}\subseteq \unions{N}$
            by \cref{preimg_iff,unions_iff,subseteq,quotient_topology,open_in}.
        Hence $\unions{N}\subseteq \preimg{p}{\unions{M}}$ by \cref{unions_iff,subseteq,preimg_iff}.
        Hence $\preimg{p}{\unions{M}} = \unions{N}$ by \cref{subseteq,subseteq_antisymmetric}.
        Therefore $\unions{M}\in T_Q$ by \cref{open_in,subseteq,quotient_topology,unions_subseteq_of_powerset_is_subseteq,pow_iff}.
    \end{subproof}
    For all $U,V$ such that $U\in T_Q$ and $V\in T_Q$ we have $U\inter V\in T_Q$
        by \cref{quotient_topology,open_in,topology_on,preimg_inter_function,pow_iff,inter_subseteq}.
    $T_Q$ is closed under arbitrary unions.
    $T_Q$ is closed under binary intersections.
    Therefore $\quotientTopology{p}{T}{X}{A}$ is a topology on $A$ by \cref{topology_on}.
\end{proof}

% from §22 Item 11: quotient topology induces a quotient map
\begin{lemma}\label{quotient_topology_quotient_map}
    Suppose $p$ is a surjection from $X$ to $A$.
    Suppose $T$ is a topology on $X$.
    Let $T_Q = \quotientTopology{p}{T}{X}{A}$.
    Then $p$ is a quotient map from $X$ to $A$ with respect to $T$ and $T_Q$.
\end{lemma}
\begin{proof}
    $p$ is a function from $X$ to $A$ and $T_Q$ is a topology on $A$
        by \cref{surj,funs,quotient_topology_is_topology}.
    For all $V$ such that $V$ is open in $A$ with respect to $T_Q$
        we have $\preimg{p}{V}$ is open in $X$ with respect to $T$ by \cref{open_in,quotient_topology}.
    Then $p$ is continuous from $X$ to $A$ with respect to $T$ and $T_Q$ by \cref{continuous}.
    For all $U$ such that $U\subseteq A$ if $\preimg{p}{U}$ is open in $X$ with respect to $T$
        then $U$ is open in $A$ with respect to $T_Q$ by \cref{pow_iff,quotient_topology,open_in}.
    Therefore $p$ is a quotient map from $X$ to $A$ with respect to $T$ and $T_Q$
        by \cref{quotient_map}.
\end{proof}
% from §22 Item 12: quotient topology is unique
\begin{lemma}\label{quotient_topology_unique}
    Suppose $p$ is a surjection from $X$ to $A$.
    Suppose $T$ is a topology on $X$.
    Suppose $T_A$ is a topology on $A$.
    Suppose $p$ is a quotient map from $X$ to $A$ with respect to $T$ and $T_A$.
    Then $T_A = \quotientTopology{p}{T}{X}{A}$.
\end{lemma}
\begin{proof}
    Let $T_Q = \quotientTopology{p}{T}{X}{A}$.
    Hence $T_A\subseteq T_Q$ by
        \cref{open_in,topology_on,subseteq,continuous,quotient_map,quotient_topology}.
    Hence $T_Q\subseteq T_A$ by \cref{quotient_topology,pow_iff,quotient_map,open_in,subseteq}.
    Therefore $T_A = T_Q$ by \cref{subseteq_antisymmetric}.
\end{proof}

% from §22 helper definition 1: partition map
\begin{definition}\label{partition_map}
    $\partitionMap{P}{X} = \{w\in X\times P \mid \fst{w}\in \snd{w}\}$.
\end{definition}

% from §22 Item 13: quotient space
\begin{definition}\label{quotient_space}
    $P$ is a quotient space of $X$ with respect to $T$ and $T_P$ iff
        $P$ is a partition of $X$ and
        $T$ is a topology on $X$ and
        $T_P = \quotientTopology{\partitionMap{P}{X}}{T}{X}{P}$.
\end{definition}

% from §22 helper lemma 2: restriction preimage for saturated subsets
\begin{lemma}\label{preimg_restrl_saturated}
    Suppose $p$ is a surjection from $X$ to $Y$.
    Suppose $A\subseteq X$.
    Suppose $A$ is saturated with respect to $p$.
    Let $q = \restrl{p}{A}$.
    Suppose $V\subseteq \img{p}{A}$.
    Then $\preimg{q}{V} = \preimg{p}{V}$.
\end{lemma}
\begin{proof}
    Hence $\preimg{q}{V}\subseteq \preimg{p}{V}$ by \cref{preimg_iff,restrl_iff,subseteq}.
    For all $x$ such that $x\in\preimg{p}{V}$ we have $x\in\preimg{q}{V}$ by
        \cref{preimg_iff,img_iff,subseteq,intersects,emptyset,inter_intro,ran_iff,saturated,singleton_iff,restrl_iff}.
    Therefore $\preimg{q}{V} = \preimg{p}{V}$ by \cref{subseteq,subseteq_antisymmetric}.
\end{proof}

% from §22 helper lemma 3: image of intersection with saturated set
\begin{lemma}\label{img_inter_saturated}
    Suppose $p$ is a function.
    Suppose $A\subseteq \dom{p}$.
    Suppose $A$ is saturated with respect to $p$.
    Suppose $U\subseteq \dom{p}$.
    Then $\img{p}{U\inter A} = \img{p}{U}\inter \img{p}{A}$.
\end{lemma}
\begin{proof}
    Hence $\img{p}{U\inter A}\subseteq \img{p}{U}\inter \img{p}{A}$ by \cref{img_iff,inter,inter_intro,subseteq}.
    For all $y$ such that $y\in\img{p}{U}\inter \img{p}{A}$ we have $y\in\img{p}{U\inter A}$ by
        \cref{inter,img_iff,preimg_iff,singleton_iff,inter_intro,emptyset,intersects,ran_iff,saturated,subseteq}.
    Therefore $\img{p}{U\inter A} = \img{p}{U}\inter \img{p}{A}$ by \cref{subseteq,subseteq_antisymmetric}.
\end{proof}

% from §22 helper lemma 4: preimage of image of a saturated set
\begin{lemma}\label{preimg_img_saturated}
    Suppose $p$ is a function.
    Suppose $A\subseteq \dom{p}$.
    Suppose $A$ is saturated with respect to $p$.
    Then $\preimg{p}{\img{p}{A}} = A$.
\end{lemma}
\begin{proof}
    Hence $A\subseteq \preimg{p}{\img{p}{A}}$ by
        \cref{subseteq,function_apply_elim,img_iff,preimg_iff}.
    For all $x$ such that $x\in\preimg{p}{\img{p}{A}}$ we have $x\in A$ by
        \cref{preimg_iff,img_iff,singleton_iff,inter_intro,emptyset,intersects,ran_iff,saturated,subseteq}.
    Therefore $\preimg{p}{\img{p}{A}} = A$ by \cref{subseteq,subseteq_antisymmetric}.
\end{proof}

% from §22 Item 14: restriction to a saturated subspace
\begin{theorem}\label{quotient_map_restriction_open_or_closed}
    Suppose $p$ is a quotient map from $X$ to $Y$ with respect to $T_X$ and $T_Y$.
    Suppose $A$ is a subspace of $X$ with respect to $T_X$.
    Suppose $A$ is saturated with respect to $p$.
    Let $q = \restrl{p}{A}$.
    Let $T_A = \subspaceTopology{T_X}{A}$.
    Let $T_B = \subspaceTopology{T_Y}{\img{p}{A}}$.
    Suppose $A$ is open in $X$ with respect to $T_X$
        or $A$ is closed in $X$ with respect to $T_X$.
    Then $q$ is a quotient map from $A$ to $\img{p}{A}$ with respect to $T_A$ and $T_B$.
\end{theorem}
    \begin{proof}
        $p\in\surj{X}{Y}$ and $p\in\funs{X}{Y}$ and $p$ is a function
            by \cref{quotient_map,surj,funs_is_function}.
    $p$ is continuous from $X$ to $Y$ with respect to $T_X$ and $T_Y$ by \cref{quotient_map}.
    Therefore $T_X$ is a topology on $X$ and $T_Y$ is a topology on $Y$ by \cref{continuous}.
    $A\subseteq X$ and $A$ is a subspace of $X$ with respect to $T_X$ by \cref{subspace_of}.
    Hence $\img{p}{A}\subseteq Y$ and $\img{p}{A}$ is a subspace of $Y$ with respect to $T_Y$
        by \cref{img_subseteq_ran,funs_ran,subseteq_transitive,subspace_of}.
    Therefore $T_A$ is a topology on $A$ and $T_B$ is a topology on $\img{p}{A}$
        by \cref{subspace_topology_is_topology,subspace_of}.
    Hence $q$ is a function by \cref{restrl_of_function_is_function}.
    $q$ is right-unique and $q$ is a relation.
        $A\subseteq \dom{p}$ by \cref{funs,subseteq_transitive,subspace_of}.
        $\dom{q} = A$ and $\ran{q} = \img{p}{A}$
            by \cref{dom_of_restrl_eq_inter,funs,inter_absorb_supseteq_left,subspace_of,restrl_ran}.
        Hence $\dom{q}\subseteq A$ and $\ran{q}\subseteq \img{p}{A}$ by \cref{subseteq_refl}.
        $q\in\rels{A}{\img{p}{A}}$ by \cref{rels_intro_dom_and_ran}.
        Then $q\in\funs{A}{\img{p}{A}}$ by \cref{funs}.
        Therefore $q\in\surj{A}{\img{p}{A}}$ by \cref{funs_surj_iff,restrl_ran}.
        Hence $q$ is continuous from $A$ to $Y$ with respect to $T_A$ and $T_Y$
            by \cref{continuous_restrict_domain}.
    Therefore $q$ is continuous from $A$ to $\img{p}{A}$ with respect to $T_A$ and $T_B$
        by \cref{restrl_img,inter_absorb_supseteq_left,subseteq,subseteq_refl,continuous_restrict_range}.
    \begin{byCase}
    \caseOf{$A$ is open in $X$ with respect to $T_X$.}
        For all $V$ such that $V\subseteq \img{p}{A}$ if $\preimg{q}{V}$ is open in $A$ with respect to $T_A$ then
            $V$ is open in $\img{p}{A}$ with respect to $T_B$
            by \cref{open_in_subspace_open_in_space,subspace_of,preimg_restrl_saturated,open_in,subseteq_transitive,quotient_map,subspace_topology,inter_absorb_supseteq_left}.
        Therefore $q$ is a quotient map from $A$ to $\img{p}{A}$ with respect to $T_A$ and $T_B$ by \cref{quotient_map}.
    \caseOf{$A$ is closed in $X$ with respect to $T_X$.}
        Show that for all $B$ such that $B\subseteq \img{p}{A}$ we have
            $B$ is closed in $\img{p}{A}$ with respect to $T_B$ iff
            $\preimg{q}{B}$ is closed in $A$ with respect to $T_A$.
        \begin{subproof}
            Fix $B$ such that $B\subseteq \img{p}{A}$.
            Show that if $B$ is closed in $\img{p}{A}$ with respect to $T_B$ then
                $\preimg{q}{B}$ is closed in $A$ with respect to $T_A$.
            \begin{subproof}
                Assume $B$ is closed in $\img{p}{A}$ with respect to $T_B$.
                Take $C$ such that $C$ is closed in $Y$ with respect to $T_Y$ and
                    $B = C\inter \img{p}{A}$ by \cref{closed_in_subspace_iff,subspace_of}.
                Hence $\preimg{q}{B} = \preimg{p}{C}\inter A$ and
                    $\preimg{p}{C}$ is closed in $X$ with respect to $T_X$
                    by \cref{closed_in,continuous,preimg_setminus_function,preimg_subseteq_dom,funs,subseteq,preimg_restrl_saturated,preimg_inter_function,subseteq_transitive,subspace_of,preimg_img_saturated}.
                Therefore $\preimg{q}{B}$ is closed in $A$ with respect to $T_A$ by \cref{closed_in_subspace_iff,subspace_of}.
            \end{subproof}
            Show that if $\preimg{q}{B}$ is closed in $A$ with respect to $T_A$ then
                $B$ is closed in $\img{p}{A}$ with respect to $T_B$.
            \begin{subproof}
                Assume $\preimg{q}{B}$ is closed in $A$ with respect to $T_A$.
                Take $C$ such that $C$ is closed in $X$ with respect to $T_X$ and
                    $\preimg{q}{B} = C\inter A$ by \cref{closed_in_subspace_iff,subspace_of}.
                $\preimg{q}{B} = \preimg{p}{B}$ by \cref{preimg_restrl_saturated}.
                Show that $\preimg{p}{B}$ is closed in $X$ with respect to $T_X$.
                \begin{subproof}
                    Let $M = \{X\setminus C, X\setminus A\}$.
                    For all $U$ such that $U\in M$ we have $U\in T_X$ by \cref{upair_elim,closed_in,open_in}.
                    $(X\setminus C)\union (X\setminus A)$ is open in $X$ with respect to $T_X$
                        by \cref{union_as_unions,subseteq,topology_on,open_in}.
                    Then $X\setminus (C\inter A)$ is open in $X$ with respect to $T_X$
                        by \cref{setminus_inter,open_in}.
                    Hence $C\inter A$ is closed in $X$ with respect to $T_X$
                        by \cref{closed_in,inter_lower_left,subseteq_transitive}.
                    Therefore $\preimg{p}{B}$ is closed in $X$ with respect to $T_X$ by \cref{subseteq}.
                \end{subproof}
                Then $B$ is closed in $Y$ with respect to $T_Y$
                    by \cref{subseteq_transitive,quotient_map_closed_iff,quotient_map}.
                Hence $B$ is closed in $\img{p}{A}$ with respect to $T_B$
                    by \cref{inter_absorb_supseteq_left,inter_comm,closed_in_subspace_iff,subspace_of}.
            \end{subproof}
            Therefore $B$ is closed in $\img{p}{A}$ with respect to $T_B$ iff
                $\preimg{q}{B}$ is closed in $A$ with respect to $T_A$.
        \end{subproof}
        Therefore $q$ is a quotient map from $A$ to $\img{p}{A}$ with respect to $T_A$ and $T_B$
            by \cref{quotient_map_closed_iff}.
    \end{byCase}
\end{proof}

% from §22 Item 15: restriction of a quotient map with open or closed map
\begin{theorem}\label{quotient_map_restriction_open_or_closed_map}
    Suppose $p$ is a quotient map from $X$ to $Y$ with respect to $T_X$ and $T_Y$.
    Suppose $A$ is a subspace of $X$ with respect to $T_X$.
    Suppose $A$ is saturated with respect to $p$.
    Let $q = \restrl{p}{A}$.
    Let $T_A = \subspaceTopology{T_X}{A}$.
    Let $T_B = \subspaceTopology{T_Y}{\img{p}{A}}$.
    Suppose $p$ is an open map from $X$ to $Y$ with respect to $T_X$ and $T_Y$
        or $p$ is a closed map from $X$ to $Y$ with respect to $T_X$ and $T_Y$.
    Then $q$ is a quotient map from $A$ to $\img{p}{A}$ with respect to $T_A$ and $T_B$.
\end{theorem}
    \begin{proof}
        $p\in\surj{X}{Y}$ and $p\in\funs{X}{Y}$ and $p$ is a function
            by \cref{quotient_map,surj,funs_is_function}.
    $p$ is continuous from $X$ to $Y$ with respect to $T_X$ and $T_Y$ by \cref{quotient_map}.
    Therefore $T_X$ is a topology on $X$ and $T_Y$ is a topology on $Y$ by \cref{continuous}.
    $A\subseteq X$ and $A$ is a subspace of $X$ with respect to $T_X$ by \cref{subspace_of}.
    Hence $\img{p}{A}\subseteq Y$ and $\img{p}{A}$ is a subspace of $Y$ with respect to $T_Y$
        by \cref{img_subseteq_ran,funs_ran,subseteq_transitive,subspace_of}.
    Therefore $T_A$ is a topology on $A$ and $T_B$ is a topology on $\img{p}{A}$
        by \cref{subspace_topology_is_topology,subspace_of}.
    Hence $q$ is a function by \cref{restrl_of_function_is_function}.
    $q$ is right-unique and $q$ is a relation.
        $A\subseteq \dom{p}$ by \cref{funs,subseteq_transitive,subspace_of}.
        $\dom{q} = A$ and $\ran{q} = \img{p}{A}$
            by \cref{dom_of_restrl_eq_inter,funs,inter_absorb_supseteq_left,subspace_of,restrl_ran}.
        Hence $\dom{q}\subseteq A$ and $\ran{q}\subseteq \img{p}{A}$ by \cref{subseteq_refl}.
        $q\in\rels{A}{\img{p}{A}}$ by \cref{rels_intro_dom_and_ran}.
        Then $q\in\funs{A}{\img{p}{A}}$ by \cref{funs}.
        Therefore $q\in\surj{A}{\img{p}{A}}$ by \cref{funs_surj_iff,restrl_ran}.
        Hence $q$ is continuous from $A$ to $Y$ with respect to $T_A$ and $T_Y$
            by \cref{continuous_restrict_domain}.
    Therefore $q$ is continuous from $A$ to $\img{p}{A}$ with respect to $T_A$ and $T_B$
        by \cref{restrl_img,inter_absorb_supseteq_left,subseteq,subseteq_refl,continuous_restrict_range}.
    \begin{byCase}
    \caseOf{$p$ is an open map from $X$ to $Y$ with respect to $T_X$ and $T_Y$.}
        Show that for all $U$ such that $U$ is open in $A$ with respect to $T_A$ we have
            $\img{q}{U}$ is open in $\img{p}{A}$ with respect to $T_B$.
        \begin{subproof}
            Fix $U$ such that $U$ is open in $A$ with respect to $T_A$.
            Then $U\in T_A$ and $U\subseteq A$ by \cref{open_in,topology_on,subseteq,pow_iff}.
            Take $W\in T_X$ such that $U = A\inter W$ by \cref{subspace_topology}.
            $\img{q}{U} = \img{p}{A\inter U}$ by \cref{funs,subseteq_transitive,subspace_of,restrl_img}.
            Then $\img{q}{U} = \img{p}{A\inter W}$ by \cref{inter_absorb_supseteq_left,subseteq}.
            Hence $\img{q}{U} = \img{p}{W}\inter \img{p}{A}$
                by \cref{inter_comm,topology_on,subseteq,pow_iff,funs,subseteq_transitive,img_inter_saturated}.
            Therefore $\img{q}{U}$ is open in $\img{p}{A}$ with respect to $T_B$
                by \cref{inter_comm,open_in,open_map,subspace_topology}.
        \end{subproof}
        Hence $q$ is an open map from $A$ to $\img{p}{A}$ with respect to $T_A$ and $T_B$
            by \cref{surj_to_fun,funs,open_map}.
        Therefore $q$ is a quotient map from $A$ to $\img{p}{A}$ with respect to $T_A$ and $T_B$
            by \cref{open_or_closed_surjection_quotient}.
    \caseOf{$p$ is a closed map from $X$ to $Y$ with respect to $T_X$ and $T_Y$.}
        For all $U$ such that $U$ is closed in $A$ with respect to $T_A$ we have
            $\img{q}{U}$ is closed in $\img{p}{A}$ with respect to $T_B$
            by \cref{closed_in_subspace_iff,subspace_of,closed_in,funs,subseteq_transitive,restrl_img,inter_absorb_supseteq_left,subseteq,img_inter_saturated,closed_map}.
        Hence $q$ is a closed map from $A$ to $\img{p}{A}$ with respect to $T_A$ and $T_B$
            by \cref{surj_to_fun,funs,closed_map}.
        Therefore $q$ is a quotient map from $A$ to $\img{p}{A}$ with respect to $T_A$ and $T_B$
            by \cref{open_or_closed_surjection_quotient}.
    \end{byCase}
\end{proof}

% from §22 Item 16: composite of quotient maps
\begin{lemma}\label{quotient_map_composite}
    Suppose $p$ is a quotient map from $X$ to $Y$ with respect to $T_X$ and $T_Y$.
    Suppose $q$ is a quotient map from $Y$ to $Z$ with respect to $T_Y$ and $T_Z$.
    Then $q\circ p$ is a quotient map from $X$ to $Z$ with respect to $T_X$ and $T_Z$.
\end{lemma}
\begin{proof}
    $q\circ p\in\surj{X}{Z}$ and $q\circ p$ is continuous from $X$ to $Z$
        with respect to $T_X$ and $T_Z$ by \cref{quotient_map,surjections_circ,continuous_composite}.
    Show that for all $U$ such that $U\subseteq Z$ if $\preimg{(q\circ p)}{U}$ is open in $X$ with respect to $T_X$
        then $U$ is open in $Z$ with respect to $T_Z$.
        \begin{subproof}
            Fix $U$ such that $U\subseteq Z$.
            Assume $\preimg{(q\circ p)}{U}$ is open in $X$ with respect to $T_X$.
            Then $\preimg{p}{\preimg{q}{U}}$ is open in $X$ with respect to $T_X$
                by \cref{preimg_iff,circ_iff,subseteq,subseteq_antisymmetric,open_in}.
            Therefore $U$ is open in $Z$ with respect to $T_Z$
                by \cref{preimg_subseteq_dom,quotient_map,surj_to_fun,funs,subseteq}.
        \end{subproof}
    Therefore $q\circ p$ is a quotient map from $X$ to $Z$ with respect to $T_X$ and $T_Z$ by \cref{quotient_map}.
\end{proof}

% from §22 helper definition 2: induced map
\begin{definition}\label{induced_map}
    $\inducedMap{p}{g}{X}{Y}{Z} =
        \{w\in Y\times Z \mid \snd{w}\in \img{g}{\preimg{p}{\{\fst{w}\}}}\}$.
\end{definition}

% from §22 Item 17: induced map exists
\begin{lemma}\label{quotient_map_induced_map}
    Suppose $p$ is a quotient map from $X$ to $Y$ with respect to $T_X$ and $T_Y$.
    Suppose $g$ is a function from $X$ to $Z$.
    Suppose for all $y\in Y$ there exists $z$ such that for all $x\in\preimg{p}{\{y\}}$ we have $g(x) = z$.
    Then there exists $f$ such that
        $f$ is a function from $Y$ to $Z$ and
        $f\circ p = g$.
\end{lemma}
\begin{proof}
    $p\in\surj{X}{Y}$ by \cref{quotient_map}.
    Hence $p\in\funs{X}{Y}$ and $p$ is a function by \cref{surj,funs_is_function}.
    Let $f = \inducedMap{p}{g}{X}{Y}{Z}$.
    Show that $f$ is a function from $Y$ to $Z$.
    \begin{subproof}
        Show that for all $y,z_1,z_2$ such that $y\mathrel{f} z_1$ and $y\mathrel{f} z_2$ we have $z_1 = z_2$.
        \begin{subproof}
            Fix $y$.
            Fix $z_1$.
            Fix $z_2$ such that $y\mathrel{f} z_1$ and $y\mathrel{f} z_2$.
            Thus $y\in Y$ and $z_1\in \img{g}{\preimg{p}{\{y\}}}$ and $z_2\in \img{g}{\preimg{p}{\{y\}}}$
                by \cref{induced_map,times_elem_is_tuple,pair_eq_iff,fst_eq,snd_eq}.
            Take $x_1\in \preimg{p}{\{y\}}$ such that $x_1\mathrel{g} z_1$ by \cref{img_iff}.
            Take $x_2\in \preimg{p}{\{y\}}$ such that $x_2\mathrel{g} z_2$ by \cref{img_iff}.
            Thus $g(x_1) = z_1$ and $g(x_2) = z_2$ by \cref{funs,funs_is_function,function_apply_intro}.
            Take $z$ such that for all $x\in\preimg{p}{\{y\}}$ we have $g(x) = z$.
            Thus $g(x_1) = z$ and $g(x_2) = z$.
            Therefore $z_1 = z_2$.
        \end{subproof}
        Therefore $f$ is right-unique.
        Show that $\dom{f} = Y$.
        \begin{subproof}
            Hence $\dom{f}\subseteq Y$ by
                \cref{dom,induced_map,times_elem_is_tuple,pair_eq_iff,subseteq}.
            Show that for all $y$ such that $y\in Y$ we have $y\in\dom{f}$.
            \begin{subproof}
                Fix $y$ such that $y\in Y$.
                Take $x\in X$ such that $p(x) = y$ by \cref{surj}.
                Then $x\mathrel{p} y$ by \cref{funs,function_apply_elim}.
                Hence $g$ is a function to $Z$ and $X = \dom{g}$ and $x\in\dom{g}$ and $g(x)\in Z$
                    by \cref{funs,funs_elim,subseteq}.
                Then $(y,g(x))\in Y\times Z$ by \cref{times_tuple_intro}.
                Hence $x\in\preimg{p}{\{y\}}$ by \cref{singleton_intro,preimg_iff}.
                Hence $(y,g(x))\in f$ by
                    \cref{function_apply_elim,img_iff,fst_eq,snd_eq,induced_map}.
                Therefore $y\in\dom{f}$ by \cref{dom}.
            \end{subproof}
            Therefore $\dom{f} = Y$ by \cref{subseteq,subseteq_antisymmetric}.
        \end{subproof}
        For all $w$ such that $w\in f$ we have $w\in Y\times Z$ by \cref{induced_map}.
        Therefore $f\in\funs{Y}{Z}$ and $f$ is a function from $Y$ to $Z$
            by \cref{subseteq,rels_intro,funs}.
    \end{subproof}
    Show that $f\circ p = g$.
    \begin{subproof}
        Show that for all $w$ such that $w\in g$ we have $w\in f\circ p$.
        \begin{subproof}
            Fix $w$ such that $w\in g$.
            Take $x,z$ such that $w = (x,z)$ and $x\in X$ and $z\in Z$
                by \cref{funs,rels_elim,subseteq,times_elem_is_tuple}.
            Thus $(x,z)\in g$.
            Then $z = g(x)$ by \cref{funs,funs_is_function,function_apply_intro}.
            Take $y$ such that $x\mathrel{p} y$ by \cref{funs,subseteq,dom_iff}.
            Hence $y\in Y$ by \cref{ran_iff,funs_ran,subseteq}.
            Then $x\in\preimg{p}{\{y\}}$ by \cref{singleton_intro,preimg_iff}.
            Hence $z\in \img{g}{\preimg{p}{\{y\}}}$ by \cref{img_iff}.
            Then $(y,z)\in f$ by \cref{times_tuple_intro,fst_eq,snd_eq,induced_map}.
            Therefore $w\in f\circ p$ by \cref{circ_iff}.
        \end{subproof}
        Hence $g\subseteq f\circ p$ by \cref{subseteq}.
        Show that for all $w$ such that $w\in f\circ p$ we have $w\in g$.
        \begin{subproof}
            Fix $w$ such that $w\in f\circ p$.
            Take $x,z$ such that $w = (x,z)$ by \cref{circ_is_relation,relation}.
            Thus $(x,z)\in f\circ p$.
            Take $y$ such that $x\mathrel{p} y$ and $y\mathrel{f} z$ by \cref{circ_iff}.
            Hence $z\in \img{g}{\preimg{p}{\{y\}}}$ by \cref{induced_map,fst_eq,snd_eq}.
            Then $g$ is a function by \cref{funs,funs_is_function}.
            Take $x_1\in \preimg{p}{\{y\}}$ such that $x_1\mathrel{g} z$ by \cref{img_iff}.
            Hence $g(x_1) = z$ by \cref{funs,funs_is_function,function_apply_intro}.
            Then $x\in\preimg{p}{\{y\}}$ by \cref{singleton_intro,preimg_iff}.
            Hence $y\in Y$ by \cref{funs_ran,ran_iff,subseteq}.
            Take $z_0$ such that for all $x\in\preimg{p}{\{y\}}$ we have $g(x) = z_0$.
            Then $g(x) = z_0$ and $g(x_1) = z_0$.
            Hence $g(x) = z$.
            Therefore $w\in g$ by \cref{preimg_subseteq_dom,funs,subseteq,function_apply_elim}.
        \end{subproof}
        Hence $f\circ p\subseteq g$ by \cref{subseteq}.
        Therefore $f\circ p = g$ by \cref{subseteq_antisymmetric}.
    \end{subproof}
\end{proof}

% from §22 Item 18: induced map continuity
\begin{lemma}\label{quotient_map_induced_map_continuous}
    Suppose $p$ is a quotient map from $X$ to $Y$ with respect to $T_X$ and $T_Y$.
    Suppose $g$ is a function from $X$ to $Z$.
    Suppose for all $y\in Y$ there exists $z$ such that for all $x\in\preimg{p}{\{y\}}$ we have $g(x) = z$.
    Suppose $f$ is a function from $Y$ to $Z$.
    Suppose $f\circ p = g$.
    Then $f$ is continuous from $Y$ to $Z$ with respect to $T_Y$ and $T_Z$
        iff $g$ is continuous from $X$ to $Z$ with respect to $T_X$ and $T_Z$.
\end{lemma}
\begin{proof}
    If $f$ is continuous from $Y$ to $Z$ with respect to $T_Y$ and $T_Z$ then
        $g$ is continuous from $X$ to $Z$ with respect to $T_X$ and $T_Z$
        by \cref{quotient_map,continuous_composite}.
    Show that if $g$ is continuous from $X$ to $Z$ with respect to $T_X$ and $T_Z$
        then $f$ is continuous from $Y$ to $Z$ with respect to $T_Y$ and $T_Z$.
    \begin{subproof}
        Assume $g$ is continuous from $X$ to $Z$ with respect to $T_X$ and $T_Z$.
        Then $T_Y$ is a topology on $Y$ and $T_Z$ is a topology on $Z$ by \cref{quotient_map,continuous}.
        Therefore for all $V$ such that $V$ is open in $Z$ with respect to $T_Z$
            we have $\preimg{f}{V}$ is open in $Y$ with respect to $T_Y$
            by \cref{continuous,preimg_iff,circ_iff,subseteq,subseteq_antisymmetric,preimg_subseteq_dom,funs,quotient_map}.
        Therefore $f$ is continuous from $Y$ to $Z$ with respect to $T_Y$ and $T_Z$ by \cref{continuous}.
    \end{subproof}
\end{proof}

% from §22 Item 19: induced map and quotient maps
\begin{lemma}\label{quotient_map_induced_map_quotient}
    Suppose $p$ is a quotient map from $X$ to $Y$ with respect to $T_X$ and $T_Y$.
    Suppose $g$ is a function from $X$ to $Z$.
    Suppose for all $y\in Y$ there exists $z$ such that for all $x\in\preimg{p}{\{y\}}$ we have $g(x) = z$.
    Suppose $f$ is a function from $Y$ to $Z$.
    Suppose $f\circ p = g$.
    Then $f$ is a quotient map from $Y$ to $Z$ with respect to $T_Y$ and $T_Z$
        iff $g$ is a quotient map from $X$ to $Z$ with respect to $T_X$ and $T_Z$.
\end{lemma}
\begin{proof}
    If $f$ is a quotient map from $Y$ to $Z$ with respect to $T_Y$ and $T_Z$ then
        $g$ is a quotient map from $X$ to $Z$ with respect to $T_X$ and $T_Z$
        by \cref{quotient_map_composite}.
    Show that if $g$ is a quotient map from $X$ to $Z$ with respect to $T_X$ and $T_Z$
        then $f$ is a quotient map from $Y$ to $Z$ with respect to $T_Y$ and $T_Z$.
    \begin{subproof}
        Assume $g$ is a quotient map from $X$ to $Z$ with respect to $T_X$ and $T_Z$.
        % For all $z\in Z$ there exists $y\in Y$ such that $f(y) = z$.
         Show that for all $z\in Z$ there exists $y\in Y$ such that $f(y) = z$. 
        \begin{subproof}
            Fix $z\in Z$.
            $g\in\surj{X}{Z}$ and $p\in\surj{X}{Y}$ by \cref{quotient_map}.
            Take $x\in X$ such that $g(x) = z$ by \cref{surj}.
            Let $y = p(x)$.
            Hence $p\in\funs{X}{Y}$ and $g\in\funs{X}{Z}$ and $f\in\funs{Y}{Z}$ and
                $p$ is a function and $g$ is a function and $f$ is a function by \cref{surj,funs_is_function}.
            Hence $y\in Y$ and $x\mathrel{p} y$ and $x\mathrel{g} z$
                by \cref{funs_type_apply,funs,function_apply_elim,function_apply_iff}.
            $x\mathrel{(f\circ p)} z$.
            Take $y_0$ such that $x\mathrel{p} y_0$ and $y_0\mathrel{f} z$ by \cref{circ_iff}.
            Therefore $f(y) = z$ by \cref{function_rightunique,function_apply_iff}.
        \end{subproof}

        Hence $f\in\surj{Y}{Z}$ by \cref{surj}.
        Therefore $f$ is continuous from $Y$ to $Z$ with respect to $T_Y$ and $T_Z$
            by \cref{quotient_map,quotient_map_induced_map_continuous}.
        Show that for all $V$ such that $V\subseteq Z$ if $\preimg{f}{V}$ is open in $Y$ with respect to $T_Y$
            then $V$ is open in $Z$ with respect to $T_Z$.
        \begin{subproof}
            Fix $V$ such that $V\subseteq Z$.
            Assume $\preimg{f}{V}$ is open in $Y$ with respect to $T_Y$.
            Hence $\preimg{g}{V}$ is open in $X$ with respect to $T_X$ by
                \cref{quotient_map,continuous,preimg_iff,circ_iff,subseteq,subseteq_antisymmetric}.
            Then $V$ is open in $Z$ with respect to $T_Z$ by \cref{quotient_map}.
        \end{subproof}
        Therefore $f$ is a quotient map from $Y$ to $Z$ with respect to $T_Y$ and $T_Z$ by \cref{quotient_map}.
    \end{subproof}
\end{proof}
% from §22 helper definition 3: fiber partition
\begin{definition}\label{fiber_partition}
    $\fiberPartition{g}{Z} = \{\preimg{g}{\{z\}} \mid z\in Z\}$.
\end{definition}

% from §22 helper lemma 5: fibers form a partition
\begin{lemma}\label{fiber_partition_is_partition}
    Suppose $g$ is a surjection from $X$ to $Z$.
    Let $X_s = \fiberPartition{g}{Z}$.
    Then $X_s$ is a partition of $X$.
\end{lemma}
\begin{proof}
    Hence $g\in\funs{X}{Z}$ and $g$ is a function by \cref{surj,funs_is_function}.
    For all $A$ such that $A\in X_s$ we have $A\neq\emptyset$
        by \cref{fiber_partition,surj,funs,function_apply_iff,singleton_intro,preimg_iff,emptyset}.
    Hence $\emptyset\notin X_s$ by \cref{emptyset}.
    For all $A,B$ such that $A\in X_s$ and $B\in X_s$ and $A\neq B$ we have $A$ is disjoint from $B$
        by \cref{fiber_partition,preimg_iff,singleton_elim,function_rightunique,disjoint}.
    Hence $X_s$ is a partition by \cref{partition}.
    For all $x$ such that $x\in X$ we have $x\in\unions{X_s}$ by
        \cref{funs_type_apply,fiber_partition,funs,function_apply_elim,preimg_iff,singleton_intro,unions_intro}.
    For all $x$ such that $x\in\unions{X_s}$ we have $x\in X$
        by \cref{unions_iff,fiber_partition,preimg_subseteq_dom,funs,subseteq}.
    Therefore $\unions{X_s} = X$ by \cref{subseteq,subseteq_antisymmetric}.
    Therefore $X_s$ is a partition of $X$.
\end{proof}

% from §22 helper lemma 6: partition map is a surjection
\begin{lemma}\label{partition_map_surjection}
    Suppose $P$ is a partition of $X$.
    Let $p = \partitionMap{P}{X}$.
    Then $p$ is a surjection from $X$ to $P$.
\end{lemma}
\begin{proof}
    $P$ is a partition and $\unions{P} = X$.
    Show that $p\in\funs{X}{P}$.
    \begin{subproof}
        $p\in\rels{X}{P}$ and $\dom{p}\subseteq X$ by \cref{partition_map,subseteq,rels_intro,rels_dom_subseteq}.
        For all $x$ such that $x\in X$ we have $x\in\dom{p}$
            by \cref{subseteq,unions_iff,times_tuple_intro,fst_eq,snd_eq,partition_map,dom}.
        Therefore $\dom{p} = X$ by \cref{subseteq,subseteq_antisymmetric}.
        For all $x,C,D$ such that $x\mathrel{p} C$ and $x\mathrel{p} D$ we have $C = D$
            by \cref{partition_map,times_tuple_elim,fst_eq,snd_eq,partition,disjoint}.
        Therefore $p\in\funs{X}{P}$ by \cref{rightunique,funs}.
    \end{subproof}
    For all $C\in P$ there exists $x\in X$ such that $p(x) = C$ by
        \cref{partition,nonempty_inhabited,unions_intro,subseteq,times_tuple_intro,fst_eq,snd_eq,partition_map,funs_is_function,function_apply_intro}.
    Therefore $p\in\surj{X}{P}$ by \cref{surj}.
\end{proof}

% from §22 Item 20: Corollary 22.3(a) induced bijective map
\begin{lemma}\label{corollary_quotient_space_induced_bijection}
    Suppose $g$ is a surjection from $X$ to $Z$.
    Suppose $g$ is continuous from $X$ to $Z$ with respect to $T_X$ and $T_Z$.
    Let $X_s = \fiberPartition{g}{Z}$.
    Let $p = \partitionMap{X_s}{X}$.
    Let $T_s = \quotientTopology{p}{T_X}{X}{X_s}$.
    Then there exists $f$ such that
        $f$ is a bijection from $X_s$ to $Z$ and
        $f$ is continuous from $X_s$ to $Z$ with respect to $T_s$ and $T_Z$ and
        $f\circ p = g$.
\end{lemma}
\begin{proof}
    $T_X$ is a topology on $X$ and $X_s$ is a partition of $X$
        by \cref{continuous,fiber_partition_is_partition}.
    Hence $p$ is a surjection from $X$ to $X_s$ by \cref{partition_map_surjection}.
    Therefore $p$ is a quotient map from $X$ to $X_s$ with respect to $T_X$ and $T_s$ by \cref{quotient_topology_quotient_map}.
    Hence $g\in\funs{X}{Z}$ and $g$ is a function by \cref{surj,funs_is_function}.
    For all $y\in X_s$ there exists $z$ such that for all $x\in\preimg{p}{\{y\}}$ we have $g(x) = z$
        by \cref{fiber_partition,preimg_iff,singleton_elim,partition_map,fst_eq,snd_eq,function_apply_intro}.
    Take $f$ such that $f$ is a function from $X_s$ to $Z$ and $f\circ p = g$ by \cref{quotient_map_induced_map}.
    Show that for all $z\in Z$ there exists $y\in X_s$ such that $f(y) = z$.
    \begin{subproof}
        Fix $z\in Z$.
        Take $x\in X$ such that $g(x) = z$ by \cref{surj}.
        $p\in\funs{X}{X_s}$ by \cref{surj}.
        Hence $p$ is a function and $f$ is a function and $p(x)\in X_s$ and
            $x\mathrel{p} p(x)$ and $x\mathrel{g} z$
            by \cref{funs_is_function,funs_type_apply,funs,function_apply_elim,function_apply_iff}.
        $x\mathrel{(f\circ p)} z$.
        Take $y$ such that $x\mathrel{p} y$ and $y\mathrel{f} z$ by \cref{circ_iff}.
        Therefore $f(p(x)) = z$ by \cref{function_rightunique,function_apply_iff}.
    \end{subproof}
    Hence $f\in\surj{X_s}{Z}$ by \cref{surj}.
    Show that for all $y_1,y_2,z$ such that $y_1\mathrel{f} z$ and $y_2\mathrel{f} z$ we have $y_1 = y_2$.
    \begin{subproof}
        Fix $y_1$.
        Fix $y_2$.
        Fix $z$.
        Assume $y_1\mathrel{f} z$ and $y_2\mathrel{f} z$.
        Hence $y_1\in X_s$ and $y_2\in X_s$
            by \cref{funs_is_function,function_apply_intro,function_apply_iff,funs}.
        Take $z_1\in Z$ such that $y_1 = \preimg{g}{\{z_1\}}$ by \cref{fiber_partition}.
        Take $z_2\in Z$ such that $y_2 = \preimg{g}{\{z_2\}}$ by \cref{fiber_partition}.
        Then $y_1\neq\emptyset$ and $y_2\neq\emptyset$ by \cref{partition}.
        Take $x_1$ such that $x_1\in y_1$ by \cref{nonempty_inhabited}.
        Take $x_2$ such that $x_2\in y_2$ by \cref{nonempty_inhabited}.
        Hence $x_1\in X$ and $x_2\in X$ by \cref{unions_intro,subseteq}.
        Then $x_1\mathrel{p} y_1$ and $x_2\mathrel{p} y_2$ by \cref{times_tuple_intro,fst_eq,snd_eq,partition_map}.
        Hence $x_1\mathrel{(f\circ p)} z$ and $x_2\mathrel{(f\circ p)} z$ by \cref{circ_iff}.
        Then $x_1\mathrel{g} z$ and $x_2\mathrel{g} z$.
        Hence $g(x_1) = z$ and $g(x_2) = z$ by \cref{function_apply_intro}.
        Then $x_1\in\preimg{g}{\{z_1\}}$ and $x_2\in\preimg{g}{\{z_2\}}$.
        Hence $g(x_1) = z_1$ and $g(x_2) = z_2$
            by \cref{preimg_iff,singleton_elim,function_apply_intro}.
        Hence $z_1 = z$ and $z_2 = z$.
        Therefore $y_1 = \preimg{g}{\{z\}}$ and $y_2 = \preimg{g}{\{z\}}$.
        Hence $y_1 = y_2$.
    \end{subproof}
    Hence $f$ is injective by \cref{injective}.
    Therefore $f\in\bijections{X_s}{Z}$ by \cref{bijections}.
    Therefore $f$ is continuous from $X_s$ to $Z$ with respect to $T_s$ and $T_Z$
        by \cref{quotient_map_induced_map_continuous}.
\end{proof}

% from §22 Item 21: Corollary 22.3(a) homeomorphism criterion
\begin{lemma}\label{corollary_quotient_space_homeomorphism}
    Suppose $g$ is a surjection from $X$ to $Z$.
    Suppose $g$ is continuous from $X$ to $Z$ with respect to $T_X$ and $T_Z$.
    Let $X_s = \fiberPartition{g}{Z}$.
    Let $p = \partitionMap{X_s}{X}$.
    Let $T_s = \quotientTopology{p}{T_X}{X}{X_s}$.
    Suppose $f$ is a function from $X_s$ to $Z$.
    Suppose $f\circ p = g$.
    Suppose $f$ is a bijection from $X_s$ to $Z$.
    Then $f$ is a homeomorphism between $X_s$ and $Z$ with respect to $T_s$ and $T_Z$
        iff $g$ is a quotient map from $X$ to $Z$ with respect to $T_X$ and $T_Z$.
\end{lemma}
\begin{proof}
    $T_X$ is a topology on $X$ and $X_s$ is a partition of $X$
        by \cref{continuous,fiber_partition_is_partition}.
    Hence $p$ is a surjection from $X$ to $X_s$ by \cref{partition_map_surjection}.
    Therefore $p$ is a quotient map from $X$ to $X_s$ with respect to $T_X$ and $T_s$ by \cref{quotient_topology_quotient_map}.
    Show that if $f$ is a homeomorphism between $X_s$ and $Z$ with respect to $T_s$ and $T_Z$
        then $g$ is a quotient map from $X$ to $Z$ with respect to $T_X$ and $T_Z$.
    \begin{subproof}
        Assume $f$ is a homeomorphism between $X_s$ and $Z$ with respect to $T_s$ and $T_Z$.
        Hence $f$ is continuous from $X_s$ to $Z$ with respect to $T_s$ and $T_Z$ by \cref{homeomorphism}.
        Therefore for all $U$ such that $U$ is open in $X_s$ with respect to $T_s$
            we have $\img{f}{U}$ is open in $Z$ with respect to $T_Z$
            by \cref{homeomorphism,continuous,funs_is_relation,funs,preim_eq_img_of_converse,converse_converse_eq}.
        Hence $f$ is an open map from $X_s$ to $Z$ with respect to $T_s$ and $T_Z$ by \cref{open_map}.
        Then $f$ is a surjection from $X_s$ to $Z$ by \cref{bijections}.
        $T_s$ is a topology on $X_s$ and $T_Z$ is a topology on $Z$
            by \cref{surj,quotient_topology_is_topology,continuous}.
        Therefore $f$ is a quotient map from $X_s$ to $Z$ with respect to $T_s$ and $T_Z$ by \cref{open_or_closed_surjection_quotient}.
        Therefore $g$ is a quotient map from $X$ to $Z$ with respect to $T_X$ and $T_Z$
            by \cref{quotient_map_composite}.
    \end{subproof}
    Show that if $g$ is a quotient map from $X$ to $Z$ with respect to $T_X$ and $T_Z$
        then $f$ is a homeomorphism between $X_s$ and $Z$ with respect to $T_s$ and $T_Z$.
    \begin{subproof}
        Assume $g$ is a quotient map from $X$ to $Z$ with respect to $T_X$ and $T_Z$.
        Hence $g\in\funs{X}{Z}$ and $g$ is a function by \cref{quotient_map,surj,funs_is_function}.
        For all $y\in X_s$ there exists $z$ such that for all $x\in\preimg{p}{\{y\}}$ we have $g(x) = z$
            by \cref{fiber_partition,preimg_iff,singleton_elim,partition_map,fst_eq,snd_eq,function_apply_intro}.
        Therefore $f$ is a quotient map from $X_s$ to $Z$ with respect to $T_s$ and $T_Z$
            by \cref{quotient_map_induced_map_quotient}.
        Hence $f$ is continuous from $X_s$ to $Z$ with respect to $T_s$ and $T_Z$ by \cref{quotient_map}.
        Show that $\converse{f}$ is continuous from $Z$ to $X_s$ with respect to $T_Z$ and $T_s$.
        \begin{subproof}
            Show that for all $U$ such that $U$ is open in $X_s$ with respect to $T_s$
                we have $\preimg{\converse{f}}{U}$ is open in $Z$ with respect to $T_Z$.
            \begin{subproof}
                Fix $U$ such that $U$ is open in $X_s$ with respect to $T_s$.
                For all $x$ such that $x\in\preimg{f}{\img{f}{U}}$ we have $x\in U$
                    by \cref{bijections,preimg_iff,img_iff,injective}.
                For all $x$ such that $x\in U$ we have $x\in\preimg{f}{\img{f}{U}}$ by
                    \cref{open_in,topology_on,subseteq,powerset_elim,funs,funs_is_function,function_apply_elim,img_iff,preimg_iff}.
                Therefore $\preimg{f}{\img{f}{U}} = U$ by \cref{subseteq,subseteq_antisymmetric}.
                Hence $\img{f}{U}$ is open in $Z$ with respect to $T_Z$
                    by \cref{funs,rels_ran_subseteq,img_subseteq_ran,subseteq_transitive,quotient_map}.
                Then $\preimg{\converse{f}}{U} = \img{f}{U}$
                    by \cref{funs_is_relation,funs,preim_eq_img_of_converse,converse_converse_eq}.
            \end{subproof}
            $\converse{f}\in\funs{Z}{X_s}$ by \cref{bijective_converse_are_funs}.
            $T_Z$ is a topology on $Z$ and $T_s$ is a topology on $X_s$ by \cref{continuous}.
            Therefore $\converse{f}$ is continuous from $Z$ to $X_s$ with respect to $T_Z$ and $T_s$ by \cref{continuous}.
        \end{subproof}
        Therefore $f$ is a homeomorphism between $X_s$ and $Z$ with respect to $T_s$ and $T_Z$ by \cref{homeomorphism}.
    \end{subproof}
\end{proof}

% from §22 Item 22: Corollary 22.3(b) Hausdorff
\begin{lemma}\label{corollary_quotient_space_hausdorff}
    Suppose $g$ is a surjection from $X$ to $Z$.
    Suppose $g$ is continuous from $X$ to $Z$ with respect to $T_X$ and $T_Z$.
    Suppose $Z$ is Hausdorff with respect to $T_Z$.
    Let $X_s = \fiberPartition{g}{Z}$.
    Let $p = \partitionMap{X_s}{X}$.
    Let $T_s = \quotientTopology{p}{T_X}{X}{X_s}$.
    Suppose $f$ is a bijection from $X_s$ to $Z$.
    Suppose $f$ is continuous from $X_s$ to $Z$ with respect to $T_s$ and $T_Z$.
    Then $X_s$ is Hausdorff with respect to $T_s$.
\end{lemma}
\begin{proof}
    $T_X$ is a topology on $X$ and $X_s$ is a partition of $X$
        by \cref{continuous,fiber_partition_is_partition}.
    Hence $p$ is a surjection from $X$ to $X_s$ by \cref{partition_map_surjection}.
    Therefore $T_s$ is a topology on $X_s$ by \cref{surj,quotient_topology_is_topology}.
    Show that for all $x_1,x_2$ such that $x_1\in X_s$ and $x_2\in X_s$ and $x_1\neq x_2$ there exist $U_1,U_2$ such that
        $U_1$ is a neighborhood of $x_1$ in $X_s$ with respect to $T_s$ and
        $U_2$ is a neighborhood of $x_2$ in $X_s$ with respect to $T_s$ and
        $U_1$ is disjoint from $U_2$.
    \begin{subproof}
        Fix $x_1$.
        Fix $x_2$.
        Assume $x_1\in X_s$ and $x_2\in X_s$ and $x_1\neq x_2$.
        $f$ is injective and $f\in\funs{X_s}{Z}$ by \cref{bijections,bijections_to_funs}.
        Hence $f$ is a function and $x_1\in\dom{f}$ and $x_2\in\dom{f}$
            by \cref{bijections_to_funs,funs_is_function,funs,subseteq}.
        Then $f(x_1) \neq f(x_2)$ and $f(x_1)\in Z$ and $f(x_2)\in Z$
            by \cref{injective_function,funs_type_apply}.
        Therefore there exist $V_1,V_2$ such that
            $V_1$ is a neighborhood of $f(x_1)$ in $Z$ with respect to $T_Z$ and
            $V_2$ is a neighborhood of $f(x_2)$ in $Z$ with respect to $T_Z$ and
            $V_1$ is disjoint from $V_2$ by \cref{hausdorff}.
        Hence $V_1$ and $V_2$ are open in $Z$ with respect to $T_Z$ and
            $f(x_1)\in V_1$ and $f(x_2)\in V_2$ by \cref{neighborhood}.
        Let $U_1 = \preimg{f}{V_1}$.
        Let $U_2 = \preimg{f}{V_2}$.
        Therefore $U_1$ and $U_2$ are open in $X_s$ with respect to $T_s$ and
            $x_1\in U_1$ and $x_2\in U_2$ by \cref{continuous,preimg_iff,function_apply_elim}.
        Then $U_1$ is a neighborhood of $x_1$ in $X_s$ with respect to $T_s$ and
            $U_2$ is a neighborhood of $x_2$ in $X_s$ with respect to $T_s$ by \cref{neighborhood}.
        Therefore $U_1$ is disjoint from $U_2$ by \cref{preimg_iff,function_apply_iff,disjoint}.
    \end{subproof}
    Therefore $X_s$ is Hausdorff with respect to $T_s$ by \cref{hausdorff}.
\end{proof}

\section{§23 Connected Spaces}

% from §23 Item 1: separation
\begin{definition}\label{separation}
    $X$ is separated by $U$ and $V$ with respect to $T$ iff
        $U$ is open in $X$ with respect to $T$ and
        $V$ is open in $X$ with respect to $T$ and
        $U\neq\emptyset$ and $V\neq\emptyset$ and
        $U$ is disjoint from $V$ and
        $U\union V = X$.
\end{definition}

% from §23 Item 2: connected
\begin{definition}\label{connected}
    $X$ is connected with respect to $T$ iff
        there exists no $U$ such that there exists $V$ such that
            $X$ is separated by $U$ and $V$ with respect to $T$.
\end{definition}
% from §23 Item 3: connected iff only clopen sets are empty or whole space
\begin{lemma}\label{connected_iff_clopen}
    Assume $T$ is a topology on $X$.
    Then $X$ is connected with respect to $T$ iff
        for all $A$ such that $A\subseteq X$ we have
            if $A$ is open in $X$ with respect to $T$ and $A$ is closed in $X$ with respect to $T$ then
                $A = \emptyset$ or $A = X$.
\end{lemma}
\begin{proof}
    Show that if $X$ is connected with respect to $T$ then
        for all $A$ such that $A\subseteq X$ we have
            if $A$ is open in $X$ with respect to $T$ and $A$ is closed in $X$ with respect to $T$ then
                $A = \emptyset$ or $A = X$.
    \begin{subproof}
        Assume $X$ is connected with respect to $T$.
        Fix $A$ such that $A\subseteq X$.
        Assume $A$ is open in $X$ with respect to $T$ and $A$ is closed in $X$ with respect to $T$.
        \begin{byCase}
        \caseOf{$A = \emptyset$.}
            Then $A = \emptyset$ or $A = X$.
        \caseOf{$A\neq \emptyset$.}
            Suppose not.
            Then $A\neq X$.
            Let $U = A$.
            Let $V = X\setminus A$.
            $U$ is open in $X$ with respect to $T$.
            $V$ is open in $X$ with respect to $T$ by \cref{closed_in}.
            $U\neq\emptyset$.
            Then $V\neq\emptyset$ by \cref{setminus_eq_emptyset_iff_subseteq,subseteq_antisymmetric,subseteq}.
            Hence $U$ is disjoint from $V$ by \cref{setminus_elim_right,disjoint}.
            $U\union V = X$ by \cref{setminus_partition}.
            Therefore there exists $U$ such that there exists $V$ such that
                $X$ is separated by $U$ and $V$ with respect to $T$ by \cref{separation}.
            Contradiction by \cref{connected}.
            Hence $A = X$.
            Therefore $A = \emptyset$ or $A = X$.
        \end{byCase}
    \end{subproof}
    Show that if
        for all $A$ such that $A\subseteq X$ we have
            if $A$ is open in $X$ with respect to $T$ and $A$ is closed in $X$ with respect to $T$ then
                $A = \emptyset$ or $A = X$
        then $X$ is connected with respect to $T$.
    \begin{subproof}
        Assume that for all $A$ such that $A\subseteq X$ we have
            if $A$ is open in $X$ with respect to $T$ and $A$ is closed in $X$ with respect to $T$ then
                $A = \emptyset$ or $A = X$.
        Show that there exists no $U$ such that there exists $V$ such that
            $X$ is separated by $U$ and $V$ with respect to $T$.
        \begin{subproof}
            Suppose not.
            Take $U$ such that there exists $V$ such that
                $X$ is separated by $U$ and $V$ with respect to $T$.
            Take $V$ such that $X$ is separated by $U$ and $V$ with respect to $T$.
            $U$ and $V$ are open in $X$ with respect to $T$ and $U\neq\emptyset$ and $V\neq\emptyset$
                by \cref{separation}.
            $U$ is disjoint from $V$ and $U\union V = X$ by \cref{separation}.
            Hence $V\subseteq X\setminus U$ by
                \cref{disjoint,union_intro_right,union_upper_left,subseteq,setminus_intro}.
            Hence $X\setminus U\subseteq V$ by
                \cref{setminus_elim_left,union_upper_left,subseteq,union_iff,setminus_elim_right}.
            Then $U$ is open in $X$ with respect to $T$ and
                $U$ is closed in $X$ with respect to $T$
                by \cref{subseteq_antisymmetric,separation,union_upper_left,subseteq,open_in,closed_in}.
            Take $y$ such that $y\in V$ by \cref{nonempty_inhabited}.
            Hence $y\in X$ and $y\notin U$ by \cref{union_intro_right,subseteq,disjoint}.
            Therefore $U\neq X$ by \cref{subseteq}.
            Hence $U = \emptyset$ or $U = X$.
            Contradiction.
        \end{subproof}
        Therefore $X$ is connected with respect to $T$ by \cref{connected}.
    \end{subproof}
\end{proof}

% from §23 Item 4: separation of a subspace
\begin{lemma}\label{separation_subspace_limit_points}
    Assume $Y$ is a subspace of $X$ with respect to $T$.
    Let $T_Y = \subspaceTopology{T}{Y}$.
    Suppose $A\subseteq Y$.
    Suppose $B\subseteq Y$.
    Then $Y$ is separated by $A$ and $B$ with respect to $T_Y$ iff
        $A\neq\emptyset$ and $B\neq\emptyset$ and
        $A$ is disjoint from $B$ and
        $A\union B = Y$ and
        $\limitPoints{A}{X}{T}$ is disjoint from $B$ and
        $\limitPoints{B}{X}{T}$ is disjoint from $A$.
\end{lemma}
\begin{proof}
    Show that if $Y$ is separated by $A$ and $B$ with respect to $T_Y$ then
        $A\neq\emptyset$ and $B\neq\emptyset$ and
        $A$ is disjoint from $B$ and
        $A\union B = Y$ and
        $\limitPoints{A}{X}{T}$ is disjoint from $B$ and
        $\limitPoints{B}{X}{T}$ is disjoint from $A$.
    \begin{subproof}
        Assume $Y$ is separated by $A$ and $B$ with respect to $T_Y$.
        $A$ and $B$ are open in $Y$ with respect to $T_Y$ and $A\neq\emptyset$ and $B\neq\emptyset$
            by \cref{separation}.
        $A$ is disjoint from $B$ and $A\union B = Y$ by \cref{separation}.
        $T$ is a topology on $X$ and $Y\subseteq X$ and $T_Y$ is a topology on $Y$
            by \cref{subspace_of,subspace_topology_is_topology}.
        Therefore $B = Y\setminus A$ by
            \cref{subseteq,disjoint,setminus_intro,setminus_elim_left,union_iff,setminus_elim_right,subseteq_antisymmetric}.
        Hence $Y\setminus A$ is open in $Y$ with respect to $T_Y$ and
            $A$ is closed in $Y$ with respect to $T_Y$ by \cref{open_in,closed_in}.
        Now $A = Y\setminus B$ by \cref{double_complement,subseteq_refl}.
        Hence $Y\setminus B$ is open in $Y$ with respect to $T_Y$ and
            $B$ is closed in $Y$ with respect to $T_Y$ by \cref{open_in,closed_in}.
        Therefore $\closure{A}{X}{T}\inter Y = A$ and $\closure{B}{X}{T}\inter Y = B$
            by \cref{closed_eq_closure,closure_in_subspace,subspace_of}.
        Therefore $\limitPoints{A}{X}{T}$ is disjoint from $B$ by
            \cref{subseteq_transitive,subseteq,closure_union_limit_points,union_intro_right,inter_intro,disjoint}.
        Therefore $\limitPoints{B}{X}{T}$ is disjoint from $A$ by
            \cref{subseteq_transitive,subseteq,closure_union_limit_points,union_intro_right,inter_intro,disjoint}.
    \end{subproof}
    Show that if
        $A\neq\emptyset$ and $B\neq\emptyset$ and
        $A$ is disjoint from $B$ and
        $A\union B = Y$ and
        $\limitPoints{A}{X}{T}$ is disjoint from $B$ and
        $\limitPoints{B}{X}{T}$ is disjoint from $A$
        then $Y$ is separated by $A$ and $B$ with respect to $T_Y$.
    \begin{subproof}
        Assume $A\neq\emptyset$ and $B\neq\emptyset$ and $A$ is disjoint from $B$ and $A\union B = Y$ and
            $\limitPoints{A}{X}{T}$ is disjoint from $B$ and $\limitPoints{B}{X}{T}$ is disjoint from $A$.
        $T$ is a topology on $X$ and $Y\subseteq X$ and $T_Y$ is a topology on $Y$
            by \cref{subspace_of,subspace_topology_is_topology}.
        Therefore $B = Y\setminus A$ by
            \cref{subseteq,disjoint,setminus_intro,setminus_elim_left,union_iff,setminus_elim_right,subseteq_antisymmetric}.
        Now $A = Y\setminus B$ by \cref{double_complement,subseteq_refl}.
        Show that for all $x$ such that $x\in\closure{A}{X}{T}\inter Y$ we have $x\in A$.
        \begin{subproof}
            Fix $x$ such that $x\in\closure{A}{X}{T}\inter Y$.
            Then $x\in\closure{A}{X}{T}$ and $x\in Y$ by \cref{inter}.
            Hence $x\in A$ or $x\in\limitPoints{A}{X}{T}$ by
                \cref{closure_union_limit_points,subseteq,union_iff}.
            Suppose not.
            Then $x\in\limitPoints{A}{X}{T}$.
            Therefore $x\in A$ by \cref{union_iff,disjoint}.
        \end{subproof}
        Therefore $\closure{A}{X}{T}\inter Y = A$ by
            \cref{subseteq,subseteq_transitive,subseteq_closure,subseteq_inter_iff,subseteq_antisymmetric}.
        Hence $\closure{A}{Y}{T_Y} = A$ by \cref{closure_in_subspace,subspace_of}.
        Therefore $Y\setminus A$ is open in $Y$ with respect to $T_Y$ and
            $A$ is closed in $Y$ with respect to $T_Y$ by \cref{closure_closed_in,closed_in,open_in}.
        Show that for all $x$ such that $x\in\closure{B}{X}{T}\inter Y$ we have $x\in B$.
        \begin{subproof}
            Fix $x$ such that $x\in\closure{B}{X}{T}\inter Y$.
            Then $x\in\closure{B}{X}{T}$ and $x\in Y$ by \cref{inter}.
            Hence $x\in B$ or $x\in\limitPoints{B}{X}{T}$ by
                \cref{closure_union_limit_points,subseteq,union_iff}.
            Suppose not.
            Then $x\in\limitPoints{B}{X}{T}$.
            Therefore $x\in B$ by \cref{union_iff,disjoint}.
        \end{subproof}
        Therefore $\closure{B}{X}{T}\inter Y = B$ by
            \cref{subseteq,subseteq_transitive,subseteq_closure,subseteq_inter_iff,subseteq_antisymmetric}.
        Hence $\closure{B}{Y}{T_Y} = B$ by \cref{closure_in_subspace,subspace_of}.
        Therefore $B$ is closed in $Y$ with respect to $T_Y$ and
            $Y\setminus B$ is open in $Y$ with respect to $T_Y$ by \cref{closure_closed_in,closed_in,open_in}.
        Hence $A$ and $B$ are open in $Y$ with respect to $T_Y$ by \cref{open_in}.
        Therefore $Y$ is separated by $A$ and $B$ with respect to $T_Y$ by \cref{separation}.
    \end{subproof}
\end{proof}

% from §23 Item 5: connected subspace lies in a separation set
\begin{lemma}\label{connected_subspace_in_separation}
    Assume $X$ is separated by $C$ and $D$ with respect to $T$.
    Assume $Y$ is a subspace of $X$ with respect to $T$.
    Let $T_Y = \subspaceTopology{T}{Y}$.
    Assume $Y$ is connected with respect to $T_Y$.
    Then $Y\subseteq C$ or $Y\subseteq D$.
\end{lemma}
\begin{proof}
    $C$ and $D$ are open in $X$ with respect to $T$ and $C$ is disjoint from $D$
        and $C\union D = X$ by \cref{separation}.
    $T$ is a topology on $X$ and $Y\subseteq X$ and $T_Y$ is a topology on $Y$
        by \cref{subspace_of,subspace_topology_is_topology}.
    Then $C\inter Y$ and $D\inter Y$ are open in $Y$ with respect to $T_Y$
        by \cref{open_in,inter_comm,subspace_topology}.
    Hence $C\inter Y$ is disjoint from $D\inter Y$ by \cref{inter,disjoint}.
    $(C\inter Y)\union(D\inter Y)\subseteq Y$ by \cref{union_iff,inter,subseteq}.
    For all $x$ such that $x\in Y$ we have $x\in (C\inter Y)\union(D\inter Y)$
        by \cref{subseteq,union_iff,inter_intro,union_intro_left,union_intro_right}.
    Hence $(C\inter Y)\union(D\inter Y) = Y$ by \cref{subseteq,subseteq_antisymmetric}.
    Show that $Y\subseteq C$ or $Y\subseteq D$.
    \begin{subproof}
        \begin{byCase}
        \caseOf{$C\inter Y = \emptyset$.}
            Therefore $Y\subseteq D$ by \cref{union_comm,union_emptyset,inter_lower_left,subseteq}.
        \caseOf{$C\inter Y\neq \emptyset$.}
            Then $D\inter Y = \emptyset$ by \cref{separation,connected}.
            Therefore $Y\subseteq C$ by \cref{union_emptyset,inter_lower_left,subseteq}.
        \end{byCase}
    \end{subproof}
\end{proof}

% from §23 helper lemma 1: subspace topology of a subspace
\begin{lemma}\label{subspace_subspace_topology}
    Assume $T$ is a topology on $X$.
    Suppose $A\subseteq Y$.
    Suppose $Y\subseteq X$.
    Let $T_Y = \subspaceTopology{T}{Y}$.
    Then $\subspaceTopology{T_Y}{A} = \subspaceTopology{T}{A}$.
\end{lemma}
\begin{proof}
    For all $U$ such that $U\in\subspaceTopology{T_Y}{A}$ we have $U\in\subspaceTopology{T}{A}$
        by \cref{subspace_topology,inter_assoc,inter_absorb_supseteq_right,subseteq}.
    Therefore $\subspaceTopology{T_Y}{A}\subseteq \subspaceTopology{T}{A}$ by \cref{subseteq}.
    For all $U$ such that $U\in\subspaceTopology{T}{A}$ we have $U\in\subspaceTopology{T_Y}{A}$
        by \cref{subspace_topology,inter_assoc,inter_absorb_supseteq_right,subseteq}.
    Hence $\subspaceTopology{T}{A}\subseteq \subspaceTopology{T_Y}{A}$ by \cref{subseteq}.
    Therefore $\subspaceTopology{T_Y}{A} = \subspaceTopology{T}{A}$ by \cref{subseteq_antisymmetric}.
\end{proof}

% from §23 Item 6: union of connected subspaces with a common point
\begin{theorem}\label{union_connected_common_point}
    Assume $M\neq\emptyset$.
    Assume $T$ is a topology on $X$.
    Assume for all $A$ such that $A\in M$ we have $A\subseteq X$.
    Assume for all $A$ such that $A\in M$ we have
        $A$ is connected with respect to $\subspaceTopology{T}{A}$.
    Assume $\inters{M}\neq\emptyset$.
    Then $\unions{M}$ is connected with respect to $\subspaceTopology{T}{\unions{M}}$.
\end{theorem}
\begin{proof}
    Let $Y = \unions{M}$.
    Let $T_Y = \subspaceTopology{T}{Y}$.
    For all $x$ such that $x\in Y$ we have $x\in X$.
    Hence $Y$ is a subspace of $X$ with respect to $T$ and $T_Y$ is a topology on $Y$
        by \cref{subseteq,subspace_of,subspace_topology_is_topology}.
    Show that there exists no $C$ such that there exists $D$ such that
        $Y$ is separated by $C$ and $D$ with respect to $T_Y$.
    \begin{subproof}
        Suppose not.
        Take $C$ such that there exists $D$ such that
            $Y$ is separated by $C$ and $D$ with respect to $T_Y$.
        Take $D$ such that $Y$ is separated by $C$ and $D$ with respect to $T_Y$.
        $C$ and $D$ are open in $Y$ with respect to $T_Y$ and $C\neq\emptyset$ and $D\neq\emptyset$
            and $C$ is disjoint from $D$ and $C\union D = Y$ by \cref{separation}.
        Take $p$ such that $p\in\inters{M}$ by \cref{nonempty_inhabited}.
        Then $p\in C$ or $p\in D$ by \cref{inters_iff_forall,unions_iff,union_iff}.
        \begin{byCase}
            \caseOf{$p\in C$.}
                Show that for all $A$ such that $A\in M$ we have $A\subseteq C$.
            \begin{subproof}
                Fix $A$ such that $A\in M$.
                Hence $A\subseteq C$ or $A\subseteq D$ by
                    \cref{subseteq,unions_iff,subspace_subspace_topology,subspace_of,connected_subspace_in_separation}.
                Therefore $A\subseteq C$ by \cref{inters_iff_forall,subseteq,disjoint}.
            \end{subproof}
            Therefore $Y\subseteq C$ by \cref{subseteq,unions_iff}.
            Hence $C = Y$ and $D = \emptyset$ by
                \cref{nonempty_inhabited,open_in,topology_on,subseteq,powerset_elim,subseteq_antisymmetric,disjoint}.
            Contradiction.
            \caseOf{$p\in D$.}
                Show that for all $A$ such that $A\in M$ we have $A\subseteq D$.
            \begin{subproof}
                Fix $A$ such that $A\in M$.
                Hence $A\subseteq C$ or $A\subseteq D$ by
                    \cref{subseteq,unions_iff,subspace_subspace_topology,subspace_of,connected_subspace_in_separation}.
                Therefore $A\subseteq D$ by \cref{inters_iff_forall,subseteq,disjoint}.
            \end{subproof}
            Therefore $Y\subseteq D$ by \cref{subseteq,unions_iff}.
            Hence $D = Y$ and $C = \emptyset$ by
                \cref{nonempty_inhabited,open_in,topology_on,subseteq,powerset_elim,subseteq_antisymmetric,disjoint}.
            Contradiction.
        \end{byCase}
    \end{subproof}
    Therefore $Y$ is connected with respect to $T_Y$ by \cref{connected}.
\end{proof}

% from §23 Item 7: connectedness preserved under adjoining limit points
\begin{theorem}\label{connected_superset_closure}
    Assume $T$ is a topology on $X$.
    Suppose $A\subseteq X$.
    Suppose $B\subseteq X$.
    Assume $A$ is connected with respect to $\subspaceTopology{T}{A}$.
    Suppose $A\subseteq B$.
    Suppose $B\subseteq \closure{A}{X}{T}$.
    Then $B$ is connected with respect to $\subspaceTopology{T}{B}$.
\end{theorem}
\begin{proof}
    Let $T_B = \subspaceTopology{T}{B}$.
    $T$ is a topology on $X$ and $A\subseteq X$ and $B\subseteq X$.
    Then $T_B$ is a topology on $B$ by \cref{subspace_topology_is_topology}.
    Show that there exists no $C$ such that there exists $D$ such that
        $B$ is separated by $C$ and $D$ with respect to $T_B$.
    \begin{subproof}
        Suppose not.
        Take $C$ such that there exists $D$ such that
            $B$ is separated by $C$ and $D$ with respect to $T_B$.
        Take $D$ such that $B$ is separated by $C$ and $D$ with respect to $T_B$.
        $C$ and $D$ are open in $B$ with respect to $T_B$ and $C\neq\emptyset$ and $D\neq\emptyset$
            and $C$ is disjoint from $D$ and $C\union D = B$ by \cref{separation}.
        Hence $C\subseteq B$ and $D\subseteq B$ and $A$ is a subspace of $B$ with respect to $T_B$
            by \cref{open_in,topology_on,subseteq,powerset_elim,subspace_of}.
        Therefore $A\subseteq C$ or $A\subseteq D$ by
            \cref{subspace_subspace_topology,subspace_of,connected_subspace_in_separation}.
        \begin{byCase}
        \caseOf{$A\subseteq C$.}
            Hence $\limitPoints{C}{X}{T}$ is disjoint from $D$ by \cref{subspace_of,separation_subspace_limit_points}.
            Therefore $\closure{C}{X}{T}$ is disjoint from $D$ and $C\subseteq X$ by
                \cref{subseteq_transitive,subseteq,closure_union_limit_points,union_iff,disjoint}.
            Therefore $D\subseteq \closure{C}{X}{T}$ by
                \cref{subseteq_transitive,subseteq,subseteq_closure,closure_closed_in,closure_subseteq_closed}.
            Then $D = \emptyset$ by \cref{nonempty_inhabited,subseteq,disjoint}.
            Contradiction by \cref{emptyset}.
        \caseOf{$A\subseteq D$.}
            Hence $\limitPoints{D}{X}{T}$ is disjoint from $C$ by \cref{subspace_of,separation_subspace_limit_points}.
            Therefore $\closure{D}{X}{T}$ is disjoint from $C$ and $D\subseteq X$ by
                \cref{subseteq_transitive,subseteq,closure_union_limit_points,union_iff,disjoint}.
            Therefore $C\subseteq \closure{D}{X}{T}$ by
                \cref{subseteq_transitive,subseteq,subseteq_closure,closure_closed_in,closure_subseteq_closed}.
            Then $C = \emptyset$ by \cref{nonempty_inhabited,subseteq,disjoint}.
            Contradiction by \cref{emptyset}.
        \end{byCase}
    \end{subproof}
    Therefore $B$ is connected with respect to $T_B$ by \cref{connected}.
\end{proof}
% from §23 Item 8: continuous image of connected space
\begin{theorem}\label{connected_image_continuous}
    Assume $T$ is a topology on $X$.
    Assume $T_1$ is a topology on $Y$.
    Suppose $f$ is continuous from $X$ to $Y$ with respect to $T$ and $T_1$.
    Assume $X$ is connected with respect to $T$.
    Let $Z = \img{f}{X}$.
    Let $T_Z = \subspaceTopology{T_1}{Z}$.
    Then $Z$ is connected with respect to $T_Z$.
\end{theorem}
\begin{proof}
    Show that there exists no $A$ such that there exists $B$ such that
        $Z$ is separated by $A$ and $B$ with respect to $T_Z$.
    \begin{subproof}
        Suppose not.
        Take $A$ such that there exists $B$ such that
            $Z$ is separated by $A$ and $B$ with respect to $T_Z$.
        Take $B$ such that $Z$ is separated by $A$ and $B$ with respect to $T_Z$.
        $A$ and $B$ are open in $Z$ with respect to $T_Z$ and $A\neq\emptyset$ and $B\neq\emptyset$
            and $A$ is disjoint from $B$ and $A\union B = Z$ by \cref{separation}.
        Hence $Z$ is a subspace of $Y$ with respect to $T_1$ and $T_Z$ is a topology on $Z$ and
            $A\subseteq Z$ and $B\subseteq Z$ by
            \cref{continuous,funs_ran,img_subseteq_ran,subseteq_transitive,subspace_of,subspace_topology_is_topology,open_in,topology_on,subseteq,powerset_elim}.
        Hence $f$ is continuous from $X$ to $Z$ with respect to $T$ and $T_Z$
            by \cref{subseteq_refl,continuous_restrict_range}.
        Then $f\in\funs{X}{Z}$ by \cref{continuous}.
        Hence $\preimg{f}{A}$ and $\preimg{f}{B}$ are open in $X$ with respect to $T$ by \cref{continuous}.
        Therefore $\preimg{f}{A}\neq\emptyset$ and $\preimg{f}{B}\neq\emptyset$
            by \cref{nonempty_inhabited,subseteq,img_iff,preimg_iff,emptyset}.
        Hence $\preimg{f}{A}$ is disjoint from $\preimg{f}{B}$ by
            \cref{preimg_iff,funs_is_function,function_apply_iff,disjoint}.
        For all $x$ such that $x\in X$ we have $x\in\preimg{f}{Z}$.
        Therefore $\preimg{f}{A}\union\preimg{f}{B} = X$
            by \cref{subseteq,preimg_subseteq_dom,funs,subseteq_antisymmetric,separation,preimg_union}.
        Therefore $X$ is separated by $\preimg{f}{A}$ and $\preimg{f}{B}$ with respect to $T$ by \cref{separation}.
        Contradiction by \cref{connected,separation}.
    \end{subproof}
    Therefore $Z$ is connected with respect to $T_Z$ by \cref{connected}.
\end{proof}

% from §23 helper lemma 2: identity is continuous
\begin{lemma}\label{identity_continuous}
    Assume $T$ is a topology on $X$.
    Let $j = \identity{X}$.
    Then $j$ is continuous from $X$ to $X$ with respect to $T$ and $T$.
\end{lemma}
\begin{proof}
    Therefore for all $U$ such that $U$ is open in $X$ with respect to $T$
        we have $\preimg{j}{U} = U$ and $\preimg{j}{U}$ is open in $X$ with respect to $T$
        by \cref{preimg_iff,id_iff,subseteq,open_in,topology_on,pow_iff,subseteq_antisymmetric}.
    Hence $j$ is continuous from $X$ to $X$ with respect to $T$ and $T$ by \cref{id_elem_funs,continuous}.
\end{proof}

% from §23 helper lemma 3: constant map is a function
\begin{lemma}\label{constant_map_function}
    Suppose $c\in B$.
    Let $f = \{(a,c) \mid a\in A\}$.
    Then $f$ is a function from $A$ to $B$.
\end{lemma}
\begin{proof}
    For all $a,b_1,b_2$ such that $(a,b_1)\in f$ and $(a,b_2)\in f$ we have $b_1 = b_2$.
    For all $w\in f$ there exist $x,y$ such that $w = (x,y)$.
    Therefore $f$ is a function by \cref{relation,rightunique}.
    For all $x\in\dom{f}$ we have $f(x)\in B$.
    Hence $f$ is a function to $B$.
    For all $x$ such that $x\in\dom{f}$ we have $x\in A$.
    For all $a$ such that $a\in A$ we have $a\in\dom{f}$.
    Therefore $\dom{f} = A$ by \cref{subseteq,subseteq_antisymmetric}.
    Hence $f$ is a function from $A$ to $B$ by \cref{funs_intro,funs}.
\end{proof}

% from §23 Item 9: finite products of connected spaces
\begin{theorem}\label{connected_product_finite}
    Assume $T_1$ is a topology on $X$.
    Assume $T_2$ is a topology on $Y$.
    Assume $X$ is connected with respect to $T_1$.
    Assume $Y$ is connected with respect to $T_2$.
    Let $T = \productTopology{T_1}{T_2}{X}{Y}$.
    Then $X\times Y$ is connected with respect to $T$.
\end{theorem}
\begin{proof}
    Then $T$ is a topology on $X\times Y$ by
        \cref{product_basis_is_basis,basis_topology_is_topology,product_topology}.
    \begin{byCase}
    \caseOf{$X = \emptyset$.}
        Hence $X\times Y = \emptyset$ by \cref{times_empty_left}.
        Show that there exists no $U$ such that there exists $V$ such that
            $X\times Y$ is separated by $U$ and $V$ with respect to $T$.
        \begin{subproof}
            Suppose not.
            Take $U$ such that there exists $V$ such that
                $X\times Y$ is separated by $U$ and $V$ with respect to $T$.
            Take $V$ such that $X\times Y$ is separated by $U$ and $V$ with respect to $T$.
            Contradiction by \cref{separation,open_in,topology_on,pow_iff,subseteq,subseteq_emptyset_iff}.
        \end{subproof}
        Therefore $X\times Y$ is connected with respect to $T$ by \cref{connected}.
    \caseOf{$X\neq\emptyset$.}
        \begin{byCase}
        \caseOf{$Y = \emptyset$.}
            Hence $X\times Y = \emptyset$ by \cref{times_empty_right}.
            Show that there exists no $U$ such that there exists $V$ such that
                $X\times Y$ is separated by $U$ and $V$ with respect to $T$.
            \begin{subproof}
                Suppose not.
                Take $U$ such that there exists $V$ such that
                    $X\times Y$ is separated by $U$ and $V$ with respect to $T$.
                Take $V$ such that $X\times Y$ is separated by $U$ and $V$ with respect to $T$.
                Contradiction by \cref{separation,open_in,topology_on,pow_iff,subseteq,subseteq_emptyset_iff}.
            \end{subproof}
            Therefore $X\times Y$ is connected with respect to $T$ by \cref{connected}.
        \caseOf{$Y\neq\emptyset$.}
            Take $a$ such that $a\in X$ by \cref{nonempty_inhabited}.
            Take $b$ such that $b\in Y$ by \cref{nonempty_inhabited}.
            Let $j_X = \identity{X}$.
            Then $j_X$ is continuous from $X$ to $X$ with respect to $T_1$ and $T_1$
                by \cref{identity_continuous}.
            Let $k = \{(x,b) \mid x\in X\}$.
            Hence $k$ is a function from $X$ to $Y$ by \cref{constant_map_function}.
            For all $x$ such that $x\in X$ we have $j_X(x) = x$ and $k(x) = b$.
            Therefore $k$ is continuous from $X$ to $Y$ with respect to $T_1$ and $T_2$
                by \cref{constant_continuous}.
            Let $f = \{(x,(j_X(x),k(x))) \mid x\in X\}$.
            Hence $f$ is continuous from $X$ to $X\times Y$ with respect to $T_1$ and $T$
                by \cref{continuous_map_into_product_backward}.
            Let $Z = \img{f}{X}$.
            Show that $Z = X\times\{b\}$.
            \begin{subproof}
                Show that for all $p$ such that $p\in Z$ we have $p\in X\times\{b\}$.
                \begin{subproof}
                    Fix $p$ such that $p\in Z$.
                    Take $x\in X$ such that $x\mathrel{f} p$ by \cref{img_iff}.
                    Take $x_0\in X$ such that $(x,p) = (x_0,(j_X(x_0),k(x_0)))$.
                    Then $p = (x,b)$ by \cref{pair_eq_iff}.
                    Hence $p\in X\times\{b\}$ by \cref{singleton_intro,times_tuple_intro}.
                \end{subproof}
                Show that for all $p$ such that $p\in X\times\{b\}$ we have $p\in Z$.
                \begin{subproof}
                    Fix $p$ such that $p\in X\times\{b\}$.
                    Take $x,y$ such that $x\in X$ and $y\in\{b\}$ and $p = (x,y)$
                        by \cref{times_elem_is_tuple}.
                    Then $p = (x,b)$ by \cref{singleton_iff,pair_eq_iff}.
                    $(x,b)\in k$.
                    Hence $(x,(j_X(x),k(x))) = (x,p)$ by \cref{function_apply_intro}.
                    Then $(x,p)\in f$.
                    Hence $x\mathrel{f} p$.
                    Therefore $p\in Z$ by \cref{img_iff}.
                \end{subproof}
                Therefore $Z = X\times\{b\}$ by \cref{subseteq,subseteq_antisymmetric}.
            \end{subproof}
            Let $T_{Xb} = \subspaceTopology{T}{X\times\{b\}}$.
            Therefore $X\times\{b\}$ is connected with respect to $T_{Xb}$ by \cref{connected_image_continuous}.
            Show that for all $x$ such that $x\in X$ we have
                $\{x\}\times Y$ is connected with respect to $\subspaceTopology{T}{\{x\}\times Y}$.
            \begin{subproof}
                Fix $x$ such that $x\in X$.
                Let $j_Y = \identity{Y}$.
                Then $j_Y$ is continuous from $Y$ to $Y$ with respect to $T_2$ and $T_2$
                    by \cref{identity_continuous}.
                Let $h = \{(y,x) \mid y\in Y\}$.
                Hence $h$ is a function from $Y$ to $X$ by \cref{constant_map_function}.
                For all $y$ such that $y\in Y$ we have $j_Y(y) = y$ and $h(y) = x$.
                Therefore $h$ is continuous from $Y$ to $X$ with respect to $T_2$ and $T_1$
                    by \cref{constant_continuous}.
                Let $g = \{(y,(h(y),j_Y(y))) \mid y\in Y\}$.
                Hence $g$ is continuous from $Y$ to $X\times Y$ with respect to $T_2$ and $T$
                    by \cref{continuous_map_into_product_backward}.
                Let $Z_x = \img{g}{Y}$.
                Show that $Z_x = \{x\}\times Y$.
                \begin{subproof}
                    Show that for all $p$ such that $p\in Z_x$ we have $p\in \{x\}\times Y$.
                    \begin{subproof}
                        Fix $p$ such that $p\in Z_x$.
                        Take $y\in Y$ such that $y\mathrel{g} p$ by \cref{img_iff}.
                        Take $y_0\in Y$ such that $(y,p) = (y_0,(h(y_0),j_Y(y_0)))$.
                        Then $p = (x,y)$ by \cref{pair_eq_iff}.
                        Hence $p\in \{x\}\times Y$ by \cref{singleton_intro,times_tuple_intro}.
                    \end{subproof}
                    Show that for all $p$ such that $p\in \{x\}\times Y$ we have $p\in Z_x$.
                    \begin{subproof}
                        Fix $p$ such that $p\in \{x\}\times Y$.
                        Take $y_0,y$ such that $y_0\in\{x\}$ and $y\in Y$ and $p = (y_0,y)$
                            by \cref{times_elem_is_tuple}.
                        Then $p = (x,y)$ by \cref{singleton_iff,pair_eq_iff}.
                        $(y,x)\in h$.
                        Hence $(y,(h(y),j_Y(y))) = (y,p)$ by \cref{function_apply_intro}.
                        Then $(y,p)\in g$.
                        Hence $y\mathrel{g} p$.
                        Therefore $p\in Z_x$ by \cref{img_iff}.
                    \end{subproof}
                    Therefore $Z_x = \{x\}\times Y$ by \cref{subseteq,subseteq_antisymmetric}.
                \end{subproof}
                Therefore $\{x\}\times Y$ is connected with respect to $\subspaceTopology{T}{\{x\}\times Y}$
                    by \cref{connected_image_continuous}.
            \end{subproof}
            Show that for all $x$ such that $x\in X$ we have
                $(X\times\{b\})\union(\{x\}\times Y)$ is connected with respect to
                $\subspaceTopology{T}{(X\times\{b\})\union(\{x\}\times Y)}$.
            \begin{subproof}
                Fix $x$ such that $x\in X$.
                Let $M = \{X\times\{b\},\{x\}\times Y\}$.
                $M\neq\emptyset$ by \cref{upair_intro_left,emptyset}.
                For all $A$ such that $A\in M$ we have $A\subseteq X\times Y$
                    by \cref{upair_iff,singleton_iff,subseteq,times_subseteq_right,times_subseteq_left}.
                For all $A$ such that $A\in M$ we have
                    $A$ is connected with respect to $\subspaceTopology{T}{A}$ by \cref{upair_iff}.
                Hence $\inters{M}\neq\emptyset$
                    by \cref{upair_intro_left,upair_iff,singleton_intro,times_tuple_intro,inters_iff_forall,emptyset}.
                Let $T_x = (X\times\{b\})\union(\{x\}\times Y)$.
                Therefore $T_x$ is connected with respect to $\subspaceTopology{T}{T_x}$
                    by \cref{union_as_unions,union_connected_common_point}.
            \end{subproof}
            Show that there exists no $C$ such that there exists $D$ such that
                $X\times Y$ is separated by $C$ and $D$ with respect to $T$.
            \begin{subproof}
                Suppose not.
                Take $C$ such that there exists $D$ such that
                    $X\times Y$ is separated by $C$ and $D$ with respect to $T$.
                Take $D$ such that $X\times Y$ is separated by $C$ and $D$ with respect to $T$.
                For all $x$ such that $x\in X$ we have
                    $(X\times\{b\})\union(\{x\}\times Y)\subseteq C$ or
                    $(X\times\{b\})\union(\{x\}\times Y)\subseteq D$ by
                    \cref{union_iff,singleton_iff,times_subseteq_right,times_subseteq_left,subseteq,subspace_of,connected_subspace_in_separation}.
                Then $(a,b)\in C$ or $(a,b)\in D$ by \cref{times_tuple_intro,separation,union_iff,subseteq}.
                \begin{byCase}
                \caseOf{$(a,b)\in C$.}
                    For all $x$ such that $x\in X$ we have
                        $(X\times\{b\})\union(\{x\}\times Y)\subseteq C$ by
                        \cref{singleton_intro,times_tuple_intro,union_intro_left,subseteq,separation,disjoint}.
                    For all $p$ such that $p\in X\times Y$ we have $p\in C$ by
                        \cref{times_elem_is_tuple,singleton_intro,times_tuple_intro,union_intro_right,subseteq}.
                    Hence $X\times Y\subseteq C$ by \cref{subseteq}.
                    Then $D = \emptyset$ by \cref{nonempty_inhabited,separation,open_in,topology_on,pow_iff,subseteq,disjoint}.
                    Contradiction by \cref{separation}.
                \caseOf{$(a,b)\in D$.}
                    For all $x$ such that $x\in X$ we have
                        $(X\times\{b\})\union(\{x\}\times Y)\subseteq D$ by
                        \cref{singleton_intro,times_tuple_intro,union_intro_left,subseteq,separation,disjoint}.
                    For all $p$ such that $p\in X\times Y$ we have $p\in D$ by
                        \cref{times_elem_is_tuple,singleton_intro,times_tuple_intro,union_intro_right,subseteq}.
                    Hence $X\times Y\subseteq D$ by \cref{subseteq}.
                    Then $C = \emptyset$ by \cref{nonempty_inhabited,separation,open_in,topology_on,pow_iff,subseteq,disjoint}.
                    Contradiction by \cref{separation}.
                \end{byCase}
            \end{subproof}
            Therefore $X\times Y$ is connected with respect to $T$ by \cref{connected}.
        \end{byCase}
    \end{byCase}
\end{proof}

\section{§24 Connected Subspaces of the Real Line}

% from §24 helper definition 1: upper bound in an ordered set
\begin{definition}\label{upper_bound_rel}
    $b$ is an upper bound of $A$ in $X$ with respect to $R$ iff
        $b\in X$ and $A\subseteq X$ and
        for all $a\in A$ we have $a\mathrel{R}b$ or $a = b$.
\end{definition}

% from §24 helper definition 2: least upper bound in an ordered set
\begin{definition}\label{least_upper_bound_rel}
    $b$ is a least upper bound of $A$ in $X$ with respect to $R$ iff
        $b$ is an upper bound of $A$ in $X$ with respect to $R$ and
        for all $c$ such that $c$ is an upper bound of $A$ in $X$ with respect to $R$
            we have $b\mathrel{R}c$ or $b = c$.
\end{definition}

% from §24 helper definition 3: least upper bound property
\begin{definition}\label{least_upper_bound_property}
    $X$ has the least upper bound property with respect to $R$ iff
        for all $A$ such that $A\subseteq X$ and $A$ is inhabited and
            there exists $b$ such that $b$ is an upper bound of $A$ in $X$ with respect to $R$
        there exists $c$ such that $c$ is a least upper bound of $A$ in $X$ with respect to $R$.
\end{definition}

% from §24 Item 1: linear continuum
\begin{definition}\label{linear_continuum}
    $X$ is a linear continuum with respect to $R$ iff
        $R$ is a simple order relation on $X$ and
        $X$ has more than one element and
        $X$ has the least upper bound property with respect to $R$ and
        for all $x,y\in X$ if $x\mathrel{R}y$ then
            there exists $z\in X$ such that $x\mathrel{R}z$ and $z\mathrel{R}y$.
\end{definition}

% from §24 helper lemma 1: open interval is convex
\begin{lemma}\label{intervalopenrel_convex}
    Assume $R$ is transitive.
    Suppose $a,b\in X$.
    Then $\intervalopenrel{R}{X}{a}{b}$ is convex in $X$ with respect to $R$.
\end{lemma}
\begin{proof}
    Follows by \cref{intervalopen_rel,transitive,subseteq,convex_in}.
\end{proof}

% from §24 helper lemma 2: open-closed interval is convex
\begin{lemma}\label{intervalopenclosedrel_convex}
    Assume $R$ is transitive.
    Suppose $a,b\in X$.
    Then $\intervalopenclosedrel{R}{X}{a}{b}$ is convex in $X$ with respect to $R$.
\end{lemma}
\begin{proof}
    Follows by \cref{intervalopenclosed_rel,transitive,subseteq,convex_in}.
\end{proof}

% from §24 helper lemma 3: closed-open interval is convex
\begin{lemma}\label{intervalclosedopenrel_convex}
    Assume $R$ is transitive.
    Suppose $a,b\in X$.
    Then $\intervalclosedopenrel{R}{X}{a}{b}$ is convex in $X$ with respect to $R$.
\end{lemma}
\begin{proof}
    Follows by \cref{intervalclosedopen_rel,transitive,subseteq,convex_in}.
\end{proof}

% from §24 helper lemma 4: closed interval is convex
\begin{lemma}\label{intervalclosedrel_convex}
    Assume $R$ is transitive.
    Suppose $a,b\in X$.
    Then $\intervalclosedrel{R}{X}{a}{b}$ is convex in $X$ with respect to $R$.
\end{lemma}
\begin{proof}
    Follows by \cref{intervalclosed_rel,transitive,subseteq,convex_in}.
\end{proof}

% from §24 helper lemma 5: open right ray is convex
\begin{lemma}\label{openrayright_convex}
    Assume $R$ is transitive.
    Suppose $a\in X$.
    Then $\openrayright{R}{X}{a}$ is convex in $X$ with respect to $R$.
\end{lemma}
\begin{proof}
    Follows by \cref{openray_right,transitive,subseteq,convex_in}.
\end{proof}

% from §24 helper lemma 6: open left ray is convex
\begin{lemma}\label{openrayleft_convex}
    Assume $R$ is transitive.
    Suppose $a\in X$.
    Then $\openrayleft{R}{X}{a}$ is convex in $X$ with respect to $R$.
\end{lemma}
\begin{proof}
    Follows by \cref{openray_left,transitive,subseteq,convex_in}.
\end{proof}

% from §24 helper lemma 7: closed right ray is convex
\begin{lemma}\label{closedrayright_convex}
    Assume $R$ is transitive.
    Suppose $a\in X$.
    Then $\closedrayright{R}{X}{a}$ is convex in $X$ with respect to $R$.
\end{lemma}
\begin{proof}
    Follows by \cref{closedray_right,transitive,subseteq,convex_in}.
\end{proof}

% from §24 helper lemma 8: closed left ray is convex
\begin{lemma}\label{closedrayleft_convex}
    Assume $R$ is transitive.
    Suppose $a\in X$.
    Then $\closedrayleft{R}{X}{a}$ is convex in $X$ with respect to $R$.
\end{lemma}
\begin{proof}
    Follows by \cref{closedray_left,transitive,subseteq,convex_in}.
\end{proof}
% from §24 helper lemma 9: convex separation contradiction
\begin{lemma}\label{convex_separation_contradiction}
    Assume $T$ is the order topology on $L$ with respect to $R$.
    Assume $L$ is a linear continuum with respect to $R$.
    Suppose $Y$ is convex in $L$ with respect to $R$.
    Let $T_Y = \subspaceTopology{T}{Y}$.
    Suppose $Y$ is separated by $A$ and $B$ with respect to $T_Y$.
    Suppose $a\in A$ and $b\in B$ and $a\mathrel{R}b$.
    Contradiction.
\end{lemma}
\begin{proof}
    $R$ is a simple order relation on $L$ by \cref{linear_continuum}.
    $R$ is transitive and $R$ is asymmetric and $R$ is connex on $L$ by \cref{simple_order_relation}.
    $L$ has more than one element and $L$ has the least upper bound property with respect to $R$
        by \cref{linear_continuum}.
    $Y\subseteq L$ by \cref{convex_in}.

    $T = \basisTopology{\orderBasis{R}{L}}{L}$ and $\orderBasis{R}{L}$ is a basis on $L$
        by \cref{order_topology,order_basis_is_basis}.
    Hence $T$ is a topology on $L$ by \cref{basis_topology_is_topology}.
    Therefore $T_Y$ is a topology on $Y$ by \cref{subspace_topology_is_topology,subspace_of,convex_in}.
    $A$ is open in $Y$ with respect to $T_Y$ and $B$ is open in $Y$ with respect to $T_Y$ by \cref{separation}.
    Then $a\in Y$ and $b\in Y$ and $a\in L$ and $b\in L$
        by \cref{open_in,separation,subspace_topology_is_topology,subspace_of,convex_in,topology_on,subseteq,powerset_elim}.

    Let $I = \intervalclosedrel{R}{L}{a}{b}$.
    For all $x$ such that $x\in I$ we have $x\in Y$ by \cref{intervalclosed_rel,convex_in}.
    Hence $I\subseteq Y$ by \cref{subseteq}.
    Then $I\subseteq L$ by \cref{convex_in,subseteq_transitive}.

    Let $A_0 = A\inter I$.
    Let $B_0 = B\inter I$.
    Let $T_I = \subspaceTopology{T}{I}$.

    There exists $U\in T$ such that $A = Y\inter U$ by \cref{subspace_topology,open_in}.
    Thus $A_0$ is open in $I$ with respect to $T_I$
        by \cref{subspace_topology,subspace_topology_is_topology,subspace_of,open_in,inter_assoc,inter_comm,inter_absorb_supseteq_right,subseteq}.

    There exists $V\in T$ such that $B = Y\inter V$ by \cref{subspace_topology,open_in}.
    Thus $B_0$ is open in $I$ with respect to $T_I$
        by \cref{subspace_topology,subspace_topology_is_topology,subspace_of,open_in,inter_assoc,inter_comm,inter_absorb_supseteq_right,subseteq}.

    $a\in I$ and $b\in I$ by \cref{intervalclosed_rel}.
    Then $a\in A_0$ and $b\in B_0$ by \cref{inter_intro}.
    Therefore $A_0\neq\emptyset$ and $B_0\neq\emptyset$ by \cref{emptyset}.

    Hence $A_0$ is disjoint from $B_0$ by \cref{inter,disjoint,separation}.

    $A_0\subseteq I$ and $B_0\subseteq I$ by \cref{inter_elim_right,subseteq}.
    For all $x$ such that $x\in I$ we have $x\in A_0\union B_0$
        by \cref{subseteq,separation,union_iff,inter_intro,union_intro_left,union_intro_right}.
    Therefore $A_0\union B_0 = I$ by \cref{union_subsets_is_subset,subseteq,subseteq_antisymmetric}.

    $I$ is separated by $A_0$ and $B_0$ with respect to $T_I$ by \cref{separation}.

    Then $A_0\subseteq L$ by \cref{inter_elim_right,subseteq,subseteq_transitive}.
    Therefore $A_0$ is inhabited by \cref{nonempty_inhabited}.

    For all $x\in A_0$ we have $x\mathrel{R}b$ or $x = b$
        by \cref{inter_elim_right,intervalclosed_rel}.
    Hence $b$ is an upper bound of $A_0$ in $L$ with respect to $R$ by \cref{upper_bound_rel}.
    Take $c$ such that $c$ is a least upper bound of $A_0$ in $L$ with respect to $R$ by
        \cref{upper_bound_rel,least_upper_bound_property,linear_continuum}.

    Hence $c\in I$ by \cref{least_upper_bound_rel,upper_bound_rel,intervalclosed_rel}.

    Then $c\in A_0$ or $c\in B_0$ by \cref{union_iff}.
    \begin{byCase}
    \caseOf{$c\in B_0$.}
        Suppose not.
        $c\notin A_0$ by \cref{disjoint}.
        Therefore $c\neq a$.
        Hence $a\mathrel{R}c$ by \cref{intervalclosed_rel}.

        Take $U\in T$ such that $B_0 = I\inter U$ by \cref{subspace_topology,open_in}.
        Then $c\in U$ by \cref{inter_elim_right}.
        Take $B_1\in\orderBasis{R}{L}$ such that $c\in B_1$ and $B_1\subseteq U$
            by \cref{topology_generated_by_basis}.
        Therefore $B_1\inter I \subseteq B_0$ by \cref{subseteq_inter_iff,subseteq}.

        Show that there exists $d\in I$ such that $d\mathrel{R}c$ and
            $\intervalopenclosedrel{R}{I}{d}{c} \subseteq B_0$.
        \begin{subproof}
            $B_1\in\pow{L}$ and
                $((\exists u, v\in L. u\mathrel{R}v \land B_1 = \intervalopenrel{R}{L}{u}{v}) \lor
                (\exists v\in L. \exists u_0\in L. (\forall x\in L. u_0\mathrel{R}x \lor u_0 = x) \land B_1 = \intervalclosedopenrel{R}{L}{u_0}{v}) \lor
                (\exists u\in L. \exists v_0\in L. (\forall x\in L. x\mathrel{R}v_0 \lor x = v_0) \land B_1 = \intervalopenclosedrel{R}{L}{u}{v_0}))$
                by \cref{order_basis}.
            \begin{byCase}
            \caseOf{$\exists u, v\in L. u\mathrel{R}v \land B_1 = \intervalopenrel{R}{L}{u}{v}$.}
                Take $u, v\in L$ such that $u\mathrel{R}v$ and $B_1 = \intervalopenrel{R}{L}{u}{v}$.
                Then $u\mathrel{R}c$ and $c\mathrel{R}v$ by \cref{intervalopen_rel}.
                Take $d\in L$ such that
                    $(d = u \lor d = a) \land (u\mathrel{R}d \lor u = d) \land (a\mathrel{R}d \lor a = d)$
                    by \cref{max_of_two_elements}.
                Hence $d\mathrel{R}c$ by \cref{max_of_two_elements,intervalopen_rel,transitive}.
                Therefore $d\in I$ by \cref{max_of_two_elements,intervalclosed_rel,transitive}.
                For all $x$ such that $x\in \intervalopenclosedrel{R}{I}{d}{c}$ we have $x\in B_0$
                    by \cref{intervalopenclosed_rel,transitive,intervalopen_rel,inter_intro,subseteq}.
                Hence $\intervalopenclosedrel{R}{I}{d}{c} \subseteq B_0$ by \cref{subseteq}.
                Therefore there exists $d\in I$ such that $d\mathrel{R}c$ and
                    $\intervalopenclosedrel{R}{I}{d}{c} \subseteq B_0$.
            \caseOf{$\exists v\in L. \exists u_0\in L. (\forall x\in L. u_0\mathrel{R}x \lor u_0 = x) \land B_1 = \intervalclosedopenrel{R}{L}{u_0}{v}$.}
                Take $u_0, v\in L$ such that
                    $(\forall x\in L. u_0\mathrel{R}x \lor u_0 = x)$ and $B_1 = \intervalclosedopenrel{R}{L}{u_0}{v}$.
                Then $c\mathrel{R}v$ by \cref{intervalclosedopen_rel}.
                Let $d = a$.
                $d\in I$ by \cref{intervalclosed_rel}.
                Hence $d\mathrel{R}c$ by \cref{intervalclosed_rel}.
                For all $x$ such that $x\in \intervalopenclosedrel{R}{I}{d}{c}$ we have $x\in B_0$
                    by \cref{intervalopenclosed_rel,subseteq,transitive,intervalclosedopen_rel,inter_intro}.
                Therefore $\intervalopenclosedrel{R}{I}{d}{c} \subseteq B_0$ by \cref{subseteq}.
                Hence there exists $d\in I$ such that $d\mathrel{R}c$ and
                    $\intervalopenclosedrel{R}{I}{d}{c} \subseteq B_0$.
            \caseOf{$\exists u\in L. \exists v_0\in L. (\forall x\in L. x\mathrel{R}v_0 \lor x = v_0) \land B_1 = \intervalopenclosedrel{R}{L}{u}{v_0}$.}
                Take $u, v_0\in L$ such that
                    $(\forall x\in L. x\mathrel{R}v_0 \lor x = v_0)$ and $B_1 = \intervalopenclosedrel{R}{L}{u}{v_0}$.
                Then $u\mathrel{R}c$ by \cref{intervalopenclosed_rel}.
                Take $d\in L$ such that
                    $(d = u \lor d = a) \land (u\mathrel{R}d \lor u = d) \land (a\mathrel{R}d \lor a = d)$
                    by \cref{max_of_two_elements}.
                $a\mathrel{R}d$ or $d = a$.
                $d\mathrel{R}c$ by \cref{intervalopenclosed_rel}.
                $c\mathrel{R}b$ or $c = b$ by \cref{intervalclosed_rel}.
                Therefore $d\mathrel{R}b$ or $d = b$ by \hyperref[transitive]{transitivity}.
                Hence $d\in I$ by \cref{intervalclosed_rel}.
                For all $x$ such that $x\in \intervalopenclosedrel{R}{I}{d}{c}$ we have $x\in B_0$
                    by \cref{intervalopenclosed_rel,transitive,inter_intro,subseteq}.
                Therefore $\intervalopenclosedrel{R}{I}{d}{c} \subseteq B_0$ by \cref{subseteq}.
                Hence there exists $d\in I$ such that $d\mathrel{R}c$ and
                    $\intervalopenclosedrel{R}{I}{d}{c} \subseteq B_0$.
            \end{byCase}
        \end{subproof}

        Take $d\in I$ such that $d\mathrel{R}c$ and $\intervalopenclosedrel{R}{I}{d}{c} \subseteq B_0$.
        Show that for all $x\in A_0$ we have $x\mathrel{R}d$ or $x = d$.
        \begin{subproof}
            Fix $x$ such that $x\in A_0$.
            Then $x\mathrel{R}c$ or $x = c$ by \cref{least_upper_bound_rel,upper_bound_rel}.
            Show that it is wrong that $d\mathrel{R}x$.
            \begin{subproof}
                Suppose not.
                Then $d\mathrel{R}x$.
                \begin{byCase}
                \caseOf{$x = c$.}
                    Then $c\in A_0$.
                    Contradiction by \cref{disjoint}.
                \caseOf{$x\mathrel{R}c$.}
                    Hence $x\in B_0$ by \cref{intervalopenclosed_rel,subseteq}.
                    Contradiction by \cref{disjoint}.
                \end{byCase}
            \end{subproof}
            \begin{byCase}
            \caseOf{$x = d$.}
                Hence $x\mathrel{R}d$ or $x = d$.
            \caseOf{$x \neq d$.}
                $x\mathrel{R}d$ or $d\mathrel{R}x$ by \cref{subseteq,connex}.
                Therefore $x\mathrel{R}d$.
                Hence $x\mathrel{R}d$ or $x = d$.
            \end{byCase}
        \end{subproof}
        Therefore $d$ is an upper bound of $A_0$ in $L$ with respect to $R$ by \cref{intervalclosed_rel,upper_bound_rel}.
        Hence $c\mathrel{R}d$ or $c = d$ by \cref{least_upper_bound_rel}.
        Contradiction by \cref{asymmetric}.
    \caseOf{$c\in A_0$.}
        Suppose not.
        $c\notin B_0$ by \cref{disjoint}.
        Therefore $c\neq b$.

        Take $U\in T$ such that $A_0 = I\inter U$ by \cref{subspace_topology,open_in}.
        Then $c\in U$ by \cref{inter_elim_right}.
        Take $B_1\in\orderBasis{R}{L}$ such that $c\in B_1$ and $B_1\subseteq U$
            by \cref{topology_generated_by_basis}.
        Therefore $B_1\inter I \subseteq A_0$ by \cref{subseteq_inter_iff,subseteq}.

        Show that there exists $e\in I$ such that $c\mathrel{R}e$ and
            $\intervalclosedopenrel{R}{I}{c}{e} \subseteq A_0$.
        \begin{subproof}
            $B_1\in\pow{L}$ and
                $((\exists u, v\in L. u\mathrel{R}v \land B_1 = \intervalopenrel{R}{L}{u}{v}) \lor
                (\exists v\in L. \exists u_0\in L. (\forall x\in L. u_0\mathrel{R}x \lor u_0 = x) \land B_1 = \intervalclosedopenrel{R}{L}{u_0}{v}) \lor
                (\exists u\in L. \exists v_0\in L. (\forall x\in L. x\mathrel{R}v_0 \lor x = v_0) \land B_1 = \intervalopenclosedrel{R}{L}{u}{v_0}))$
                by \cref{order_basis}.
            \begin{byCase}
            \caseOf{$\exists u, v\in L. u\mathrel{R}v \land B_1 = \intervalopenrel{R}{L}{u}{v}$.}
                Take $u, v\in L$ such that $u\mathrel{R}v$ and $B_1 = \intervalopenrel{R}{L}{u}{v}$.
                Then $u\mathrel{R}c$ and $c\mathrel{R}v$ by \cref{intervalopen_rel}.
                Take $e\in L$ such that
                    $(e = v \lor e = b) \land (e\mathrel{R}v \lor e = v) \land (e\mathrel{R}b \lor e = b)$
                    by \cref{min_of_two_elements}.
                Hence $c\mathrel{R}e$ by \cref{min_of_two_elements,intervalclosed_rel,transitive}.
                Therefore $e\in I$ by \cref{min_of_two_elements,intervalclosed_rel,transitive}.
                For all $x$ such that $x\in \intervalclosedopenrel{R}{I}{c}{e}$ we have $x\in A_0$
                    by \cref{intervalclosedopen_rel,transitive,intervalopen_rel,inter_intro,subseteq}.
                Hence $\intervalclosedopenrel{R}{I}{c}{e} \subseteq A_0$ by \cref{subseteq}.
                Therefore there exists $e\in I$ such that $c\mathrel{R}e$ and
                    $\intervalclosedopenrel{R}{I}{c}{e} \subseteq A_0$.
            \caseOf{$\exists v\in L. \exists u_0\in L. (\forall x\in L. u_0\mathrel{R}x \lor u_0 = x) \land B_1 = \intervalclosedopenrel{R}{L}{u_0}{v}$.}
                Take $u_0, v\in L$ such that
                    $(\forall x\in L. u_0\mathrel{R}x \lor u_0 = x)$ and $B_1 = \intervalclosedopenrel{R}{L}{u_0}{v}$.
                Then $c\mathrel{R}v$ by \cref{intervalclosedopen_rel}.
                Take $e\in L$ such that
                    $(e = v \lor e = b) \land (e\mathrel{R}v \lor e = v) \land (e\mathrel{R}b \lor e = b)$
                    by \cref{min_of_two_elements}.
                $e\mathrel{R}b$ or $e = b$.
                $c\mathrel{R}b$ by \cref{intervalclosed_rel}.
                Therefore $c\mathrel{R}e$.
                $a\mathrel{R}c$ or $a = c$ by \cref{intervalclosed_rel}.
                Hence $a\mathrel{R}e$ or $e = a$ by \hyperref[transitive]{transitivity}.
                Therefore $e\in I$ by \cref{intervalclosed_rel}.
                For all $x$ such that $x\in \intervalclosedopenrel{R}{I}{c}{e}$ we have $x\in A_0$
                    by \cref{intervalclosedopen_rel,transitive,inter_intro,subseteq}.
                Hence $\intervalclosedopenrel{R}{I}{c}{e} \subseteq A_0$ by \cref{subseteq}.
                Therefore there exists $e\in I$ such that $c\mathrel{R}e$ and
                    $\intervalclosedopenrel{R}{I}{c}{e} \subseteq A_0$.
            \caseOf{$\exists u\in L. \exists v_0\in L. (\forall x\in L. x\mathrel{R}v_0 \lor x = v_0) \land B_1 = \intervalopenclosedrel{R}{L}{u}{v_0}$.}
                Take $u, v_0\in L$ such that
                    $(\forall x\in L. x\mathrel{R}v_0 \lor x = v_0)$ and $B_1 = \intervalopenclosedrel{R}{L}{u}{v_0}$.
                Then $u\mathrel{R}c$ by \cref{intervalopenclosed_rel}.
                Take $e\in L$ such that
                    $(e = v_0 \lor e = b) \land (e\mathrel{R}v_0 \lor e = v_0) \land (e\mathrel{R}b \lor e = b)$
                    by \cref{min_of_two_elements}.
                Show that $c\mathrel{R}e$.
                \begin{subproof}
                    \begin{byCase}
                    \caseOf{$e = b$.}
                        $c\mathrel{R}b$ or $c = b$ by \cref{intervalclosed_rel}.
                        $c \neq b$ by \cref{disjoint}.
                        Therefore $c\mathrel{R}e$.
                    \caseOf{$e = v_0$.}
                        Then $c\mathrel{R}v_0$ or $c = v_0$.
                        Show that $c\mathrel{R}v_0$.
                        \begin{subproof}
                            Suppose not.
                            Then $c = v_0$.
                            $b\mathrel{R}v_0$ or $b = v_0$.
                            \begin{byCase}
                            \caseOf{$b = v_0$.} Contradiction by \cref{disjoint}.
                            \caseOf{$b\mathrel{R}v_0$.}
                                $c\mathrel{R}b$ or $c = b$ by \cref{intervalclosed_rel}.
                                \begin{byCase}
                                \caseOf{$c = b$.} Contradiction by \cref{disjoint}.
                                \caseOf{$c\mathrel{R}b$.}
                                    Then $b\mathrel{R}c$.
                                    Contradiction by \cref{asymmetric}.
                                \end{byCase}
                            \end{byCase}
                        \end{subproof}
                        Therefore $c\mathrel{R}e$.
                    \end{byCase}
                \end{subproof}
                $a\mathrel{R}c$ or $a = c$ by \cref{intervalclosed_rel}.
                Hence $a\mathrel{R}e$ or $e = a$ by \hyperref[transitive]{transitivity}.
                Therefore $e\in I$ by \cref{intervalclosed_rel}.
                For all $x$ such that $x\in \intervalclosedopenrel{R}{I}{c}{e}$ we have $x\in A_0$
                    by \cref{intervalclosedopen_rel,transitive,intervalopenclosed_rel,inter_intro,subseteq}.
                Hence $\intervalclosedopenrel{R}{I}{c}{e} \subseteq A_0$ by \cref{subseteq}.
                Therefore there exists $e\in I$ such that $c\mathrel{R}e$ and
                    $\intervalclosedopenrel{R}{I}{c}{e} \subseteq A_0$.
            \end{byCase}
        \end{subproof}

        Take $e\in I$ such that $c\mathrel{R}e$ and $\intervalclosedopenrel{R}{I}{c}{e} \subseteq A_0$.
        Take $z\in L$ such that $c\mathrel{R}z$ and $z\mathrel{R}e$ by \cref{subseteq,linear_continuum}.
        $a\mathrel{R}c$ or $a = c$ by \cref{intervalclosed_rel}.
        Hence $a\mathrel{R}z$ or $z = a$ by \hyperref[transitive]{transitivity}.
        $e\mathrel{R}b$ or $e = b$ by \cref{intervalclosed_rel}.
        Therefore $z\mathrel{R}b$ or $z = b$ by \hyperref[transitive]{transitivity}.
        Hence $z\in I$ by \cref{intervalclosed_rel}.
        Then $z\in A_0$ by \cref{intervalclosedopen_rel,subseteq}.
        Hence $z\mathrel{R}c$ or $z = c$ by \cref{least_upper_bound_rel,upper_bound_rel}.
        Contradiction by \cref{asymmetric}.
    \end{byCase}
\end{proof}

% from §24 helper lemma 10: convex subset of a linear continuum is connected
\begin{lemma}\label{convex_connected_linear_continuum}
    Assume $T$ is the order topology on $L$ with respect to $R$.
    Assume $L$ is a linear continuum with respect to $R$.
    Suppose $Y$ is convex in $L$ with respect to $R$.
    Let $T_Y = \subspaceTopology{T}{Y}$.
    Then $Y$ is connected with respect to $T_Y$.
\end{lemma}
\begin{proof}
    Show that there exists no $A$ such that there exists $B$ such that
        $Y$ is separated by $A$ and $B$ with respect to $T_Y$.
    \begin{subproof}
        Suppose not.
        Then there exists $A$ such that there exists $B$ such that
            $Y$ is separated by $A$ and $B$ with respect to $T_Y$.
        Take $A$ such that there exists $B$ such that
            $Y$ is separated by $A$ and $B$ with respect to $T_Y$.
        Take $B$ such that $Y$ is separated by $A$ and $B$ with respect to $T_Y$.
        Take $a$ such that $a\in A$ by \cref{nonempty_inhabited,separation}.
        Take $b$ such that $b\in B$ by \cref{nonempty_inhabited,separation}.
        Hence $a\mathrel{R}b$ or $b\mathrel{R}a$
            by \cref{separation,union_upper_left,union_upper_right,convex_in,subseteq_transitive,subseteq,disjoint,linear_continuum,simple_order_relation,connex}.
        Contradiction by \cref{convex_separation_contradiction,separation,disjoint_symmetric,union_comm}.
    \end{subproof}
    Therefore $Y$ is connected with respect to $T_Y$ by \cref{connected}.
\end{proof}

% from §24 Item 2: linear continuum connected, intervals and rays connected
\begin{theorem}\label{linear_continuum_connected}
    Assume $T$ is the order topology on $L$ with respect to $R$.
    Assume $L$ is a linear continuum with respect to $R$.
    Then $L$ is connected with respect to $T$ and
        for all $a,b\in L$ if $a\mathrel{R}b$ then
            $\intervalopenrel{R}{L}{a}{b}$ is connected with respect to $\subspaceTopology{T}{\intervalopenrel{R}{L}{a}{b}}$ and
            $\intervalopenclosedrel{R}{L}{a}{b}$ is connected with respect to $\subspaceTopology{T}{\intervalopenclosedrel{R}{L}{a}{b}}$ and
            $\intervalclosedopenrel{R}{L}{a}{b}$ is connected with respect to $\subspaceTopology{T}{\intervalclosedopenrel{R}{L}{a}{b}}$ and
            $\intervalclosedrel{R}{L}{a}{b}$ is connected with respect to $\subspaceTopology{T}{\intervalclosedrel{R}{L}{a}{b}}$ and
        for all $a\in L$ we have
            $\openrayright{R}{L}{a}$ is connected with respect to $\subspaceTopology{T}{\openrayright{R}{L}{a}}$ and
            $\openrayleft{R}{L}{a}$ is connected with respect to $\subspaceTopology{T}{\openrayleft{R}{L}{a}}$ and
            $\closedrayright{R}{L}{a}$ is connected with respect to $\subspaceTopology{T}{\closedrayright{R}{L}{a}}$ and
            $\closedrayleft{R}{L}{a}$ is connected with respect to $\subspaceTopology{T}{\closedrayleft{R}{L}{a}}$.
\end{theorem}
\begin{proof}
    $T$ is a topology on $L$ by \cref{order_topology,order_basis_is_basis,basis_topology_is_topology}.
    Let $T_L = \subspaceTopology{T}{L}$.
    Hence $L$ is connected with respect to $T_L$ by \cref{convex_in,subseteq,convex_connected_linear_continuum}.
    Then $T_L\subseteq T$
        by \cref{subspace_topology,topology_on,powerset_elim,inter_absorb_supseteq_left,subseteq}.
    For all $U$ such that $U\in T$ we have $U\in T_L$
        by \cref{topology_on,subseteq,powerset_elim,inter_absorb_supseteq_left,subspace_topology}.
    Therefore $T\subseteq T_L$ by \cref{subseteq}.
    Hence $L$ is connected with respect to $T$ by \cref{subseteq_antisymmetric,subseteq}.

    For all $a,b\in L$ if $a\mathrel{R}b$ then
        $\intervalopenrel{R}{L}{a}{b}$ is connected with respect to $\subspaceTopology{T}{\intervalopenrel{R}{L}{a}{b}}$ and
        $\intervalopenclosedrel{R}{L}{a}{b}$ is connected with respect to $\subspaceTopology{T}{\intervalopenclosedrel{R}{L}{a}{b}}$ and
        $\intervalclosedopenrel{R}{L}{a}{b}$ is connected with respect to $\subspaceTopology{T}{\intervalclosedopenrel{R}{L}{a}{b}}$ and
        $\intervalclosedrel{R}{L}{a}{b}$ is connected with respect to $\subspaceTopology{T}{\intervalclosedrel{R}{L}{a}{b}}$
        by \cref{linear_continuum,simple_order_relation,intervalopenrel_convex,intervalopenclosedrel_convex,intervalclosedopenrel_convex,intervalclosedrel_convex,convex_connected_linear_continuum}.

    For all $a\in L$ we have
        $\openrayright{R}{L}{a}$ is connected with respect to $\subspaceTopology{T}{\openrayright{R}{L}{a}}$ and
        $\openrayleft{R}{L}{a}$ is connected with respect to $\subspaceTopology{T}{\openrayleft{R}{L}{a}}$ and
        $\closedrayright{R}{L}{a}$ is connected with respect to $\subspaceTopology{T}{\closedrayright{R}{L}{a}}$ and
        $\closedrayleft{R}{L}{a}$ is connected with respect to $\subspaceTopology{T}{\closedrayleft{R}{L}{a}}$
        by \cref{linear_continuum,simple_order_relation,openrayright_convex,openrayleft_convex,closedrayright_convex,closedrayleft_convex,convex_connected_linear_continuum}.
\end{proof}
% from §24 helper definition 4: real order relation
\begin{definition}\label{real_order_relation}
    $\realOrderRel = \{w\in\reals\times\reals \mid \fst{w} < \snd{w}\}$.
\end{definition}

% from §24 helper definition 5: order topology on reals
\begin{definition}\label{real_order_topology}
    $\realOrderTopology = \basisTopology{\orderBasis{\realOrderRel}{\reals}}{\reals}$.
\end{definition}

% from §24 helper definition 6: between set in an ordered set
\begin{definition}\label{between_rel}
    $\betweenrel{R}{Y}{x}{y} = \intervalclosedrel{R}{Y}{x}{y} \union \intervalclosedrel{R}{Y}{y}{x}$.
\end{definition}

% from §24 helper lemma 11: real order relation intro
\begin{lemma}\label{real_order_relation_intro}
    Suppose $a\in\reals$ and $b\in\reals$ and $a < b$.
    Then $a\mathrel{\realOrderRel} b$.
\end{lemma}
\begin{proof}
    Follows by \cref{times_tuple_intro,fst_eq,snd_eq,real_lt_subst_left,real_lt_subst_right,real_order_relation}.
\end{proof}

% from §24 helper lemma 12: real order relation elim
\begin{lemma}\label{real_order_relation_elim}
    Suppose $a\mathrel{\realOrderRel} b$.
    Then $a\in\reals$ and $b\in\reals$ and $a < b$.
\end{lemma}
\begin{proof}
    Follows by \cref{real_order_relation,times_tuple_elim,fst_eq,snd_eq,real_lt_subst_left,real_lt_subst_right}.
\end{proof}

% from §24 helper lemma 13: real order relation is simple
\begin{lemma}\label{real_order_relation_simple}
    $\realOrderRel$ is a simple order relation on $\reals$.
\end{lemma}
\begin{proof}
    $\realOrderRel$ is transitive
        by \cref{real_order_relation_elim,real_lt_trans,real_order_relation_intro,transitive}.
    For all $a,b$ such that $a\mathrel{\realOrderRel} b$ we have it is wrong that $b\mathrel{\realOrderRel} a$
        by \cref{real_order_relation_elim,real_lt_trans,real_lt_irreflexive}.
    Hence $\realOrderRel$ is a simple order relation on $\reals$
        by \cref{real_order_relation,subseteq,asymmetric,real_lt_trichotomy,real_order_relation_intro,connex,simple_order_relation}.
\end{proof}

% from §24 helper lemma 14: inhabited implies nonempty
\begin{lemma}\label{inhabited_nonempty}
    Suppose $A$ is inhabited.
    Then $A\neq\emptyset$.
\end{lemma}
\begin{proof}
    Take $a$ such that $a\in A$.
    Then $A\neq\emptyset$ by \cref{emptyset}.
\end{proof}

% from §24 helper lemma 15: real line is a linear continuum
\begin{lemma}\label{reals_linear_continuum}
    $\reals$ is a linear continuum with respect to $\realOrderRel$.
\end{lemma}
\begin{proof}
    $\realOrderRel$ is a simple order relation on $\reals$ by \cref{real_order_relation_simple}.
    Hence $\reals$ has more than one element
        by \cref{real_add_ident,real_mul_ident,real_zero_lt_one,real_lt_irreflexive,more_than_one_element}.
    For all $A$ such that $A\subseteq\reals$ and $A$ is inhabited and
        there exists $u$ such that $u$ is an upper bound of $A$ in $\reals$ with respect to $\realOrderRel$
        there exists $c$ such that $c$ is a least upper bound of $A$ in $\reals$ with respect to $\realOrderRel$
        by \cref{upper_bound_rel,real_order_relation_elim,real_upper_bound,real_bounded_above,inhabited_nonempty,real_completeness_axiom,real_supremum,subseteq,real_order_relation_intro,least_upper_bound_rel}.
    Hence $\reals$ has the least upper bound property with respect to $\realOrderRel$ by \cref{least_upper_bound_property}.
    For all $a,b\in\reals$ if $a\mathrel{\realOrderRel} b$ then
        there exists $z\in\reals$ such that $a\mathrel{\realOrderRel} z$ and $z\mathrel{\realOrderRel} b$
        by \cref{real_order_relation_elim,real_rational_open_interval,intervalopen,real_order_relation_intro}.
    Therefore $\reals$ is a linear continuum with respect to $\realOrderRel$ by \cref{linear_continuum}.
\end{proof}

% from §24 Item 3: real line connected, intervals and rays connected
\begin{corollary}\label{reals_connected}
    $\reals$ is connected with respect to $\realOrderTopology$ and
        for all $a,b\in\reals$ if $a\mathrel{\realOrderRel} b$ then
            $\intervalopenrel{\realOrderRel}{\reals}{a}{b}$ is connected with respect to $\subspaceTopology{\realOrderTopology}{\intervalopenrel{\realOrderRel}{\reals}{a}{b}}$ and
            $\intervalopenclosedrel{\realOrderRel}{\reals}{a}{b}$ is connected with respect to $\subspaceTopology{\realOrderTopology}{\intervalopenclosedrel{\realOrderRel}{\reals}{a}{b}}$ and
            $\intervalclosedopenrel{\realOrderRel}{\reals}{a}{b}$ is connected with respect to $\subspaceTopology{\realOrderTopology}{\intervalclosedopenrel{\realOrderRel}{\reals}{a}{b}}$ and
            $\intervalclosedrel{\realOrderRel}{\reals}{a}{b}$ is connected with respect to $\subspaceTopology{\realOrderTopology}{\intervalclosedrel{\realOrderRel}{\reals}{a}{b}}$ and
        for all $a\in\reals$ we have
            $\openrayright{\realOrderRel}{\reals}{a}$ is connected with respect to $\subspaceTopology{\realOrderTopology}{\openrayright{\realOrderRel}{\reals}{a}}$ and
            $\openrayleft{\realOrderRel}{\reals}{a}$ is connected with respect to $\subspaceTopology{\realOrderTopology}{\openrayleft{\realOrderRel}{\reals}{a}}$ and
            $\closedrayright{\realOrderRel}{\reals}{a}$ is connected with respect to $\subspaceTopology{\realOrderTopology}{\closedrayright{\realOrderRel}{\reals}{a}}$ and
            $\closedrayleft{\realOrderRel}{\reals}{a}$ is connected with respect to $\subspaceTopology{\realOrderTopology}{\closedrayleft{\realOrderRel}{\reals}{a}}$.
\end{corollary}
\begin{proof}
    Follows by \cref{reals_linear_continuum,linear_continuum,real_order_topology,order_topology,linear_continuum_connected}.
\end{proof}

% from §24 Item 4: intermediate value theorem
\begin{theorem}\label{intermediate_value_theorem}
    Suppose $f$ is continuous from $X$ to $Y$ with respect to $T$ and $T_1$.
    Assume $X$ is connected with respect to $T$.
    Assume $R$ is a simple order relation on $Y$.
    Assume $Y$ has more than one element.
    Assume $T_1$ is the order topology on $Y$ with respect to $R$.
    Suppose $a\in X$.
    Suppose $b\in X$.
    Suppose $r\in Y$.
    Suppose $r\in\betweenrel{R}{Y}{f(a)}{f(b)}$.
    Then there exists $c\in X$ such that $f(c) = r$.
\end{theorem}
\begin{proof}
    $f\in\funs{X}{Y}$ and $f$ is a function by \cref{continuous,funs_elim}.
    $T$ is a topology on $X$ and $T_1$ is a topology on $Y$ by \cref{continuous}.
    Let $Z = \img{f}{X}$.
    Let $T_Z = \subspaceTopology{T_1}{Z}$.
    $Z$ is connected with respect to $T_Z$ and
        $Z$ is a subspace of $Y$ with respect to $T_1$ and
        $T_Z$ is a topology on $Z$
        by \cref{connected_image_continuous,funs_ran,img_subseteq_ran,subseteq_transitive,subspace_of,subspace_topology_is_topology}.
    Suppose not.
    Then there exists no $c\in X$ such that $f(c) = r$.
    Then $r\notin Z$.
    Hence $f(a)\in Z$ and $f(b)\in Z$ and $f(a)\in Y$ and $f(b)\in Y$
        by \cref{funs,function_apply_elim,img_elem_intro,funs_type_apply}.
    Let $A = Z\inter\openrayleft{R}{Y}{r}$.
    Let $B = Z\inter\openrayright{R}{Y}{r}$.
    Therefore $A$ is open in $Z$ with respect to $T_Z$ and $B$ is open in $Z$ with respect to $T_Z$
        by \cref{openrayleft_open_in_order_topology,openrayright_open_in_order_topology,open_in,subspace_topology}.
    $A\neq\emptyset$ and $B\neq\emptyset$
        by \cref{between_rel,union_iff,intervalclosed_rel,openray_left,openray_right,inter_intro,emptyset}.
    $A$ is disjoint from $B$ by \cref{inter,openray_left,openray_right,simple_order_relation,asymmetric,disjoint}.
    Hence $A\union B\subseteq Z$ by \cref{union_iff,inter,subseteq}.
    For all $y$ such that $y\in Z$ we have $y\in A\union B$
        by \cref{subspace_of,subseteq,simple_order_relation,connex,openray_left,openray_right,inter_intro,union_iff}.
    Hence $A\union B = Z$ by \cref{union_iff,inter,subseteq,subseteq_antisymmetric}.
    Therefore $Z$ is separated by $A$ and $B$ with respect to $T_Z$ by \cref{separation}.
    Contradiction by \cref{connected}.
\end{proof}

% from §24 Item 5: path in a space
\begin{definition}\label{path_in_space}
    $f$ is a path in $X$ from $x$ to $y$ with respect to $T$ iff
        there exist $a,b\in\reals$ such that
            $a\mathrel{\realOrderRel} b$ and
            $f$ is continuous from $\intervalclosedrel{\realOrderRel}{\reals}{a}{b}$ to $X$ with respect to
                $\subspaceTopology{\realOrderTopology}{\intervalclosedrel{\realOrderRel}{\reals}{a}{b}}$ and $T$ and
            $f(a) = x$ and $f(b) = y$.
\end{definition}

% from §24 Item 6: path connected
\begin{definition}\label{path_connected}
    $X$ is path connected with respect to $T$ iff
        for all $x,y\in X$ there exists $f$ such that
            $f$ is a path in $X$ from $x$ to $y$ with respect to $T$.
\end{definition}

% from §24 Item 7: path-connected implies connected
\begin{theorem}\label{path_connected_implies_connected}
    Assume $T$ is a topology on $X$.
    Assume $X$ is path connected with respect to $T$.
    Then $X$ is connected with respect to $T$.
\end{theorem}
\begin{proof}
    Suppose not.
    Then there exists $A$ such that there exists $B$ such that
        $X$ is separated by $A$ and $B$ with respect to $T$.
    Take $A$ such that there exists $B$ such that
        $X$ is separated by $A$ and $B$ with respect to $T$.
    Take $B$ such that $X$ is separated by $A$ and $B$ with respect to $T$.
    $A\neq\emptyset$ and $B\neq\emptyset$ by \cref{separation}.
    Take $x$ such that $x\in A$ by \cref{nonempty_inhabited}.
    Take $y$ such that $y\in B$ by \cref{nonempty_inhabited}.
    Take $f$ such that $f$ is a path in $X$ from $x$ to $y$ with respect to $T$
        by \cref{separation,open_in,topology_on,powerset_elim,subseteq,path_connected}.
    Take $a,b\in\reals$ such that
        $a\mathrel{\realOrderRel} b$ and
        $f$ is continuous from $\intervalclosedrel{\realOrderRel}{\reals}{a}{b}$ to $X$ with respect to
            $\subspaceTopology{\realOrderTopology}{\intervalclosedrel{\realOrderRel}{\reals}{a}{b}}$ and $T$ and
        $f(a) = x$ and $f(b) = y$
        by \cref{path_in_space}.
    Let $I = \intervalclosedrel{\realOrderRel}{\reals}{a}{b}$.
    Let $T_I = \subspaceTopology{\realOrderTopology}{I}$.
    $I$ is connected with respect to $T_I$ by \cref{continuous,reals_connected}.
    Let $Z = \img{f}{I}$.
    Let $T_Z = \subspaceTopology{T}{Z}$.
    Hence $Z\subseteq A$ or $Z\subseteq B$
        by \cref{connected_image_continuous,continuous,funs_elim,funs_ran,img_subseteq_ran,subseteq_transitive,subspace_of,connected_subspace_in_separation}.
    Therefore $x\in Z$ and $y\in Z$
        by \cref{intervalclosed_rel,continuous,funs_elim,function_apply_elim,img_elem_intro}.
    Contradiction by \cref{subseteq,separation,disjoint}.
\end{proof}

\section{§25 Components and Local Connectedness}

% from §25 Item 1: component relation
\begin{definition}\label{component_relation}
    $\componentRelation{T}{X} = \{w\in X\times X \mid
        \text{there exists $A$ such that
            $A$ is a subspace of $X$ with respect to $T$ and
            $\fst{w}\in A$ and $\snd{w}\in A$ and
            $A$ is connected with respect to $\subspaceTopology{T}{A}$}\}$.
\end{definition}

% from §25 Item 2: component
\begin{definition}\label{component}
    Assume $T$ is a topology on $X$.
    $C$ is a component of $X$ with respect to $T$ iff
        there exists $x\in X$ such that
            $C = \equivalenceClass{\componentRelation{T}{X}}{x}$.
\end{definition}

% from §25 helper lemma 1: component relation is a binary relation
\begin{lemma}\label{component_relation_binary}
    Assume $T$ is a topology on $X$.
    Then $\componentRelation{T}{X}$ is a binary relation on $X$.
\end{lemma}
\begin{proof}
    Follows by \cref{component_relation,subseteq}.
\end{proof}

% from §25 helper lemma 2: component relation is reflexive
\begin{lemma}\label{component_relation_reflexive}
    Assume $T$ is a topology on $X$.
    Then $\componentRelation{T}{X}$ is reflexive on $X$.
\end{lemma}
\begin{proof}
    Show that for all $x\in X$ we have $x\mathrel{\componentRelation{T}{X}}x$.
    \begin{subproof}
        Fix $x\in X$.
        Let $A = \{x\}$.
        Thus $A$ is a subspace of $X$ with respect to $T$ by \cref{singleton_subset_intro,subspace_of}.
        Show that $A$ is connected with respect to $\subspaceTopology{T}{A}$.
        \begin{subproof}
            Suppose not.
            Then there exists $U$ such that there exists $V$ such that
                $A$ is separated by $U$ and $V$ with respect to $\subspaceTopology{T}{A}$.
            Take $U$ such that there exists $V$ such that
                $A$ is separated by $U$ and $V$ with respect to $\subspaceTopology{T}{A}$.
            Take $V$ such that $A$ is separated by $U$ and $V$ with respect to $\subspaceTopology{T}{A}$.
            $U\neq\emptyset$ and $V\neq\emptyset$ by \cref{separation}.
            Take $u$ such that $u\in U$ by \cref{nonempty_inhabited}.
            Take $v$ such that $v\in V$ by \cref{nonempty_inhabited}.
            Then $u = x$ and $v = x$
                by \cref{separation,open_in,subspace_topology_is_topology,subspace_of,topology_on,powerset_elim,subseteq,singleton_iff}.
            Contradiction by \cref{disjoint,separation}.
        \end{subproof}
        Therefore $x\mathrel{\componentRelation{T}{X}}x$
            by \cref{singleton_intro,fst_eq,snd_eq,times_tuple_intro,component_relation}.
    \end{subproof}
    Hence $\componentRelation{T}{X}$ is reflexive on $X$ by \cref{reflexive_on}.
\end{proof}

% from §25 helper lemma 3: component relation is symmetric
\begin{lemma}\label{component_relation_symmetric}
    Assume $T$ is a topology on $X$.
    Then $\componentRelation{T}{X}$ is symmetric.
\end{lemma}
\begin{proof}
    Follows by \cref{component_relation,subspace_of,subseteq,times_tuple_intro,fst_eq,snd_eq,symmetric}.
\end{proof}

% from §25 helper lemma 4: component relation is transitive
\begin{lemma}\label{component_relation_transitive}
    Assume $T$ is a topology on $X$.
    Then $\componentRelation{T}{X}$ is transitive.
\end{lemma}
\begin{proof}
    Show that for all $x,y,z$ such that
        $x\mathrel{\componentRelation{T}{X}}y$ and $y\mathrel{\componentRelation{T}{X}}z$
        we have $x\mathrel{\componentRelation{T}{X}}z$.
    \begin{subproof}
        Fix $x$.
        Fix $y$.
        Fix $z$.
        Assume $x\mathrel{\componentRelation{T}{X}}y$ and $y\mathrel{\componentRelation{T}{X}}z$.
        Take $A$ such that
            $A$ is a subspace of $X$ with respect to $T$ and
            $\fst{(x,y)}\in A$ and $\snd{(x,y)}\in A$ and
            $A$ is connected with respect to $\subspaceTopology{T}{A}$
            by \cref{component_relation}.
        Then $x\in A$ and $y\in A$ by \cref{fst_eq,snd_eq}.
        Take $B$ such that
            $B$ is a subspace of $X$ with respect to $T$ and
            $\fst{(y,z)}\in B$ and $\snd{(y,z)}\in B$ and
            $B$ is connected with respect to $\subspaceTopology{T}{B}$
            by \cref{component_relation}.
        Then $y\in B$ and $z\in B$ by \cref{fst_eq,snd_eq}.
        Let $M = \{A,B\}$.
        Hence $M\neq\emptyset$ by \cref{upair_intro_left,emptyset}.
        For all $C$ such that $C\in M$ we have $C\subseteq X$ by \cref{upair_iff,subspace_of}.
        For all $C$ such that $C\in M$ we have
            $C$ is connected with respect to $\subspaceTopology{T}{C}$ by \cref{upair_iff}.
        For all $C$ such that $C\in M$ we have $y\in C$ by \cref{upair_iff}.
        Therefore $\inters{M}\neq\emptyset$ by \cref{upair_intro_left,nonempty_inhabited,inters_iff_forall,emptyset}.
        Hence $A\union B$ is connected with respect to $\subspaceTopology{T}{A\union B}$
            by \cref{union_connected_common_point,union_as_unions}.
        Then $A\union B$ is a subspace of $X$ with respect to $T$ by \cref{union_iff,subspace_of,subseteq}.
        Hence $x\in A\union B$ and $z\in A\union B$ by \cref{union_iff}.
        Hence $x\mathrel{\componentRelation{T}{X}}z$
            by \cref{subspace_of,subseteq,times_tuple_intro,fst_eq,snd_eq,component_relation}.
    \end{subproof}
    Therefore $\componentRelation{T}{X}$ is transitive by \cref{transitive}.
\end{proof}

% from §25 helper lemma 5: component relation is an equivalence
\begin{lemma}\label{component_relation_equivalence}
    Assume $T$ is a topology on $X$.
    Then $\componentRelation{T}{X}$ is an equivalence on $X$.
\end{lemma}
\begin{proof}
    Follows by \cref{component_relation_binary,component_relation_reflexive,component_relation_transitive,component_relation_symmetric}.
\end{proof}

% from §25 Item 3: components are connected
\begin{theorem}\label{component_connected}
    Assume $T$ is a topology on $X$.
    Suppose $C$ is a component of $X$ with respect to $T$.
    Then $C$ is connected with respect to $\subspaceTopology{T}{C}$.
\end{theorem}
    \begin{proof}
        Take $x\in X$ such that $C = \equivalenceClass{\componentRelation{T}{X}}{x}$ by \cref{component}.
        Let $M = \{A\in\pow{X} \mid \text{$x\in A$ and $A$ is connected with respect to $\subspaceTopology{T}{A}$}\}$.
        Take $A$ such that
            $A$ is a subspace of $X$ with respect to $T$ and
            $\fst{(x,x)}\in A$ and $\snd{(x,x)}\in A$ and
            $A$ is connected with respect to $\subspaceTopology{T}{A}$
            by \cref{component_relation_reflexive,reflexive_on,component_relation}.
        Therefore $M\neq\emptyset$ by \cref{fst_eq,snd_eq,subspace_of,powerset_intro,emptyset}.
    For all $A$ such that $A\in M$ we have $A\subseteq X$ by \cref{powerset_elim}.
    For all $A$ such that $A\in M$ we have
        $A$ is connected with respect to $\subspaceTopology{T}{A}$ and $x\in A$.
    Hence $\inters{M}\neq\emptyset$ by \cref{fst_eq,snd_eq,subspace_of,powerset_intro,nonempty_inhabited,inters_iff_forall,emptyset}.
    Let $U = \unions{M}$.
    $U$ is connected with respect to $\subspaceTopology{T}{U}$ by \cref{union_connected_common_point}.
    For all $y$ such that $y\in U$ we have $y\in C$
        by \cref{unions_iff,powerset_elim,subspace_of,subseteq,times_tuple_intro,fst_eq,snd_eq,component_relation,component_relation_symmetric,symmetric,downward_closure_iff}.
    Therefore $U\subseteq C$ by \cref{subseteq}.
    For all $y$ such that $y\in C$ we have $y\in U$
        by \cref{downward_closure_iff,component_relation,fst_eq,snd_eq,subspace_of,powerset_intro,unions_iff}.
    Therefore $U = C$ by \cref{subseteq,subseteq_antisymmetric}.
    Hence $C$ is connected with respect to $\subspaceTopology{T}{C}$.
\end{proof}
% from §25 Item 4: distinct components are disjoint
\begin{theorem}\label{components_disjoint}
    Assume $T$ is a topology on $X$.
    Suppose $C$ is a component of $X$ with respect to $T$.
    Suppose $D$ is a component of $X$ with respect to $T$.
    Suppose $C\neq D$.
    Then $C$ is disjoint from $D$.
\end{theorem}
\begin{proof}
    Suppose not.
    Then there exists $z$ such that $z\in C$ and $z\in D$.
    Take $x\in X$ such that $C = \equivalenceClass{\componentRelation{T}{X}}{x}$ by \cref{component}.
    Take $y\in X$ such that $D = \equivalenceClass{\componentRelation{T}{X}}{y}$ by \cref{component}.
    Therefore $C = D$
        by \cref{downward_closure_iff,component_relation_symmetric,component_relation_transitive,symmetric,transitive,component_relation_equivalence,equiv_iff_equivclasses_eq}.
    Contradiction.
\end{proof}

% from §25 Item 5: components cover the space
\begin{theorem}\label{components_cover}
    Assume $T$ is a topology on $X$.
    Then for all $x\in X$ there exists $C$ such that
        $C$ is a component of $X$ with respect to $T$ and $x\in C$.
\end{theorem}
\begin{proof}
    Follows by \cref{component_relation_equivalence,equivclasses_inhabited_equivalenceon,component}.
\end{proof}

% from §25 Item 6: connected subspaces intersect only one component
\begin{theorem}\label{connected_subspace_single_component}
    Assume $T$ is a topology on $X$.
    Suppose $A$ is a subspace of $X$ with respect to $T$.
    Suppose $A$ is connected with respect to $\subspaceTopology{T}{A}$.
    Suppose $A\neq\emptyset$.
    Then for all $C,D$ such that
        $C$ is a component of $X$ with respect to $T$ and
        $D$ is a component of $X$ with respect to $T$ and
        $A$ intersects $C$ and $A$ intersects $D$
        we have $C = D$.
\end{theorem}
\begin{proof}
    Fix $C$.
    Fix $D$ such that
        $C$ is a component of $X$ with respect to $T$ and
        $D$ is a component of $X$ with respect to $T$ and
        $A$ intersects $C$ and $A$ intersects $D$.
    Take $x$ such that $x\in A\inter C$ by \cref{intersects,nonempty_inhabited}.
    Take $y$ such that $y\in A\inter D$ by \cref{intersects,nonempty_inhabited}.
    Then $(x,y)\in \componentRelation{T}{X}$
        by \cref{inter,subspace_of,subseteq,times_tuple_intro,fst_eq,snd_eq,component_relation}.
    Take $c\in X$ such that $C = \equivalenceClass{\componentRelation{T}{X}}{c}$ by \cref{component}.
    Take $d\in X$ such that $D = \equivalenceClass{\componentRelation{T}{X}}{d}$ by \cref{component}.
    Therefore $C = D$
        by \cref{inter,downward_closure_iff,component_relation_symmetric,component_relation_transitive,symmetric,transitive,component_relation_equivalence,equiv_iff_equivclasses_eq}.
\end{proof}

% from §25 Item 7: path relation
\begin{definition}\label{path_relation}
    $\pathRelation{T}{X} = \{w\in X\times X \mid
        \text{there exists $f$ such that
            $f$ is a path in $X$ from $\fst{w}$ to $\snd{w}$ with respect to $T$}\}$.
\end{definition}

% from §25 Item 8: path component
\begin{definition}\label{path_component}
    Assume $T$ is a topology on $X$.
    $C$ is a path component of $X$ with respect to $T$ iff
        there exists $x\in X$ such that
            $C = \equivalenceClass{\pathRelation{T}{X}}{x}$.
\end{definition}

% from §25 helper lemma 6: real line has smaller elements
\begin{lemma}\label{reals_has_smaller}
    For all $a\in\reals$ there exists $b\in\reals$ such that $b < a$.
\end{lemma}
\begin{proof}
    Follows by \cref{real_add_ident,real_mul_ident,real_add_inverse,real_zero_lt_one,real_lt_neg,real_neg_zero,real_lt_subst_right,real_lt_add,real_add_comm,real_lt_subst_left,real_add_closed}.
\end{proof}

% from §25 helper lemma 7: real line has larger elements
\begin{lemma}\label{reals_has_larger}
    For all $a\in\reals$ there exists $b\in\reals$ such that $a < b$.
\end{lemma}
\begin{proof}
    Follows by \cref{real_add_ident,real_mul_ident,real_zero_lt_one,real_lt_add,real_add_comm,real_lt_subst_left,real_lt_subst_right,real_add_closed}.
\end{proof}

% from §25 helper lemma 8: order interval equals real open interval
\begin{lemma}\label{intervalopenrel_real_eq_intervalopen}
    Suppose $a\in\reals$ and $b\in\reals$.
    Then $\intervalopenrel{\realOrderRel}{\reals}{a}{b} = \intervalopen{a}{b}$.
\end{lemma}
\begin{proof}
    Follows by \cref{intervalopen_rel,real_order_relation_elim,intervalopen,real_order_relation_intro,subseteq,subseteq_antisymmetric}.
\end{proof}

% from §25 helper lemma 9: order basis on reals equals the standard basis
\begin{lemma}\label{order_basis_reals_standard}
    $\orderBasis{\realOrderRel}{\reals} = \standardBasisR$.
\end{lemma}
\begin{proof}
    Show that for all $B$ such that $B\in\orderBasis{\realOrderRel}{\reals}$ we have
        $B\in\standardBasisR$.
    \begin{subproof}
            Fix $B$ such that $B\in\orderBasis{\realOrderRel}{\reals}$.
            Then
                $(\exists a, b\in\reals. a\mathrel{\realOrderRel}b \land B = \intervalopenrel{\realOrderRel}{\reals}{a}{b}) \lor
                (\exists b\in\reals. \exists a_0\in\reals. (\forall x\in\reals. a_0\mathrel{\realOrderRel}x \lor a_0 = x) \land B = \intervalclosedopenrel{\realOrderRel}{\reals}{a_0}{b}) \lor
                (\exists a\in\reals. \exists b_0\in\reals. (\forall x\in\reals. x\mathrel{\realOrderRel}b_0 \lor x = b_0) \land B = \intervalopenclosedrel{\realOrderRel}{\reals}{a}{b_0})$
                by \cref{order_basis}.
            \begin{byCase}
            \caseOf{$\exists a, b\in\reals. a\mathrel{\realOrderRel}b \land B = \intervalopenrel{\realOrderRel}{\reals}{a}{b}$.}
                Take $a,b\in\reals$ such that $a\mathrel{\realOrderRel}b$ and
                    $B = \intervalopenrel{\realOrderRel}{\reals}{a}{b}$.
                Hence $B\in\standardBasisR$
                    by \cref{real_order_relation_elim,intervalopenrel_real_eq_intervalopen,intervalopen_subset_reals,subseteq,powerset_intro,standard_basis_reals}.
            \caseOf{$\exists b\in\reals. \exists a_0\in\reals. (\forall x\in\reals. a_0\mathrel{\realOrderRel}x \lor a_0 = x) \land B = \intervalclosedopenrel{\realOrderRel}{\reals}{a_0}{b}$.}
                Suppose not.
                Take $b,a_0\in\reals$ such that
                    $(\forall x\in\reals. a_0\mathrel{\realOrderRel}x \lor a_0 = x)$ and
                    $B = \intervalclosedopenrel{\realOrderRel}{\reals}{a_0}{b}$.
                Take $c\in\reals$ such that $c < a_0$ by \cref{reals_has_smaller}.
                Then $a_0\mathrel{\realOrderRel}c$ or $a_0 = c$.
                Contradiction by \cref{real_order_relation_elim,real_lt_trans,real_lt_irreflexive,real_lt_subst_left,real_lt_subst_right}.
            \caseOf{$\exists a\in\reals. \exists b_0\in\reals. (\forall x\in\reals. x\mathrel{\realOrderRel}b_0 \lor x = b_0) \land B = \intervalopenclosedrel{\realOrderRel}{\reals}{a}{b_0}$.}
                Suppose not.
                Take $a,b_0\in\reals$ such that
                    $(\forall x\in\reals. x\mathrel{\realOrderRel}b_0 \lor x = b_0)$ and
                    $B = \intervalopenclosedrel{\realOrderRel}{\reals}{a}{b_0}$.
                Take $c\in\reals$ such that $b_0 < c$ by \cref{reals_has_larger}.
                Then $c\mathrel{\realOrderRel}b_0$ or $c = b_0$.
                Contradiction by \cref{real_order_relation_elim,real_lt_trans,real_lt_irreflexive,real_lt_subst_left,real_lt_subst_right}.
            \end{byCase}
    \end{subproof}
    Hence $\orderBasis{\realOrderRel}{\reals} \subseteq \standardBasisR$ by \cref{subseteq}.
    For all $B$ such that $B\in\standardBasisR$ there exist $a,b\in\reals$ such that
        $a < b$ and $B = \intervalopen{a}{b}$ by \cref{standard_basis_reals}.
    Thus for all $B$ such that $B\in\standardBasisR$ we have $B\in\orderBasis{\realOrderRel}{\reals}$
        by \cref{real_order_relation_intro,intervalopenrel_real_eq_intervalopen,intervalopen_subset_reals,subseteq,powerset_intro,order_basis}.
    Therefore $\standardBasisR \subseteq \orderBasis{\realOrderRel}{\reals}$ by \cref{subseteq}.
    Hence $\orderBasis{\realOrderRel}{\reals} = \standardBasisR$ by \cref{subseteq_antisymmetric}.
\end{proof}

% from §25 helper lemma 10: real order topology equals standard topology
\begin{lemma}\label{real_order_topology_eq_standard}
    $\realOrderTopology = \standardTopologyR$.
\end{lemma}
\begin{proof}
    Follows by \cref{real_order_topology,standard_topology_reals,order_basis_reals_standard}.
\end{proof}

% from §25 helper lemma 11: closed right ray is complement of open left ray
\begin{lemma}\label{closedrayright_complement_openrayleft}
    Assume $R$ is a simple order relation on $X$.
    Suppose $a\in X$.
    Then $\closedrayright{R}{X}{a} = X\setminus \openrayleft{R}{X}{a}$.
\end{lemma}
\begin{proof}
    For all $x$ such that $x\in\closedrayright{R}{X}{a}$ we have
        $x\in X\setminus \openrayleft{R}{X}{a}$
        by \cref{closedray_right,openray_left,simple_order_relation,asymmetric,real_lt_subst_left,setminus_intro}.
    Therefore $\closedrayright{R}{X}{a} \subseteq X\setminus \openrayleft{R}{X}{a}$ by \cref{subseteq}.
    For all $x$ such that $x\in X\setminus \openrayleft{R}{X}{a}$ we have
        $x\in\closedrayright{R}{X}{a}$
        by \cref{setminus_elim_left,setminus_elim_right,openray_left,simple_order_relation,connex,closedray_right}.
    Hence $X\setminus \openrayleft{R}{X}{a} \subseteq \closedrayright{R}{X}{a}$ by \cref{subseteq}.
    Therefore $\closedrayright{R}{X}{a} = X\setminus \openrayleft{R}{X}{a}$ by \cref{subseteq_antisymmetric}.
\end{proof}

% from §25 helper lemma 12: closed left ray is complement of open right ray
\begin{lemma}\label{closedrayleft_complement_openrayright}
    Assume $R$ is a simple order relation on $X$.
    Suppose $a\in X$.
    Then $\closedrayleft{R}{X}{a} = X\setminus \openrayright{R}{X}{a}$.
\end{lemma}
\begin{proof}
    For all $x$ such that $x\in\closedrayleft{R}{X}{a}$ we have
        $x\in X\setminus \openrayright{R}{X}{a}$
        by \cref{closedray_left,openray_right,simple_order_relation,asymmetric,real_lt_subst_left,setminus_intro}.
    Therefore $\closedrayleft{R}{X}{a} \subseteq X\setminus \openrayright{R}{X}{a}$ by \cref{subseteq}.
    For all $x$ such that $x\in X\setminus \openrayright{R}{X}{a}$ we have
        $x\in\closedrayleft{R}{X}{a}$ by
        \cref{setminus_elim_left,setminus_elim_right,openray_right,simple_order_relation,connex,closedray_left}.
    Hence $X\setminus \openrayright{R}{X}{a} \subseteq \closedrayleft{R}{X}{a}$ by \cref{subseteq}.
    Therefore $\closedrayleft{R}{X}{a} = X\setminus \openrayright{R}{X}{a}$ by \cref{subseteq_antisymmetric}.
\end{proof}
% from §25 helper lemma 13: closed right ray is closed in the order topology
\begin{lemma}\label{closedrayright_closed_in_order_topology}
    Assume $R$ is a simple order relation on $X$.
    Assume $X$ has more than one element.
    Assume $T$ is the order topology on $X$ with respect to $R$.
    Suppose $a\in X$.
    Then $\closedrayright{R}{X}{a}$ is closed in $X$ with respect to $T$.
\end{lemma}
\begin{proof}
    Hence $X\setminus \closedrayright{R}{X}{a}$ is open in $X$ with respect to $T$
        by \cref{openrayleft_open_in_order_topology,closedray_right,openray_left,subseteq,closedrayright_complement_openrayleft,setminus_setminus,inter_absorb_supseteq_left,open_in}.
    Hence $\closedrayright{R}{X}{a}$ is closed in $X$ with respect to $T$
        by \cref{closedray_right,subseteq,closed_in}.
\end{proof}

% from §25 helper lemma 14: closed left ray is closed in the order topology
\begin{lemma}\label{closedrayleft_closed_in_order_topology}
    Assume $R$ is a simple order relation on $X$.
    Assume $X$ has more than one element.
    Assume $T$ is the order topology on $X$ with respect to $R$.
    Suppose $a\in X$.
    Then $\closedrayleft{R}{X}{a}$ is closed in $X$ with respect to $T$.
\end{lemma}
\begin{proof}
    Hence $X\setminus \closedrayleft{R}{X}{a}$ is open in $X$ with respect to $T$
        by \cref{openrayright_open_in_order_topology,closedray_left,openray_right,subseteq,closedrayleft_complement_openrayright,setminus_setminus,inter_absorb_supseteq_left,open_in}.
    Hence $\closedrayleft{R}{X}{a}$ is closed in $X$ with respect to $T$
        by \cref{closedray_left,subseteq,closed_in}.
\end{proof}

% from §25 helper lemma 15: closed interval is closed in the order topology
\begin{lemma}\label{intervalclosedrel_closed_in_order_topology}
    Assume $R$ is a simple order relation on $X$.
    Assume $X$ has more than one element.
    Assume $T$ is the order topology on $X$ with respect to $R$.
    Suppose $a\in X$.
    Suppose $b\in X$.
    Then $\intervalclosedrel{R}{X}{a}{b}$ is closed in $X$ with respect to $T$.
\end{lemma}
\begin{proof}
    Let $A = \closedrayright{R}{X}{a}$.
    Let $B = \closedrayleft{R}{X}{b}$.
    Then $T$ is a topology on $X$ by \cref{order_basis_is_basis,basis_topology_is_topology,order_topology}.
    Let $C = \{A,B\}$.
    For all $D$ such that $D\in C$ we have $D$ is closed in $X$ with respect to $T$
        by \cref{upair_iff,closedrayright_closed_in_order_topology,closedrayleft_closed_in_order_topology}.
    Therefore $\inters{C}$ is closed in $X$ with respect to $T$
        by \cref{upair_iff,closedray_right,closedray_left,subseteq,powerset_intro,closed_inters}.
    Then $\intervalclosedrel{R}{X}{a}{b} \subseteq \inters{C}$
        by \cref{intervalclosed_rel,closedray_right,closedray_left,upair_iff,inters_iff_forall,subseteq}.
    Hence $\inters{C} \subseteq \intervalclosedrel{R}{X}{a}{b}$
        by \cref{inters_iff_forall,upair_iff,closedray_right,closedray_left,intervalclosed_rel,subseteq}.
    Therefore $\intervalclosedrel{R}{X}{a}{b} = \inters{C}$ by \cref{subseteq_antisymmetric}.
    Hence $\intervalclosedrel{R}{X}{a}{b}$ is closed in $X$ with respect to $T$ by \cref{subseteq}.
\end{proof}

% from §25 helper lemma 16: reverse a path
\begin{lemma}\label{path_reverse}
    Assume $T$ is a topology on $X$.
    Suppose $f$ is a path in $X$ from $x$ to $y$ with respect to $T$.
    Then there exists $g$ such that $g$ is a path in $X$ from $y$ to $x$ with respect to $T$.
\end{lemma}
\begin{proof}
    Take $a,b\in\reals$ such that
        $a\mathrel{\realOrderRel} b$ and
        $f$ is continuous from $\intervalclosedrel{\realOrderRel}{\reals}{a}{b}$ to $X$ with respect to
            $\subspaceTopology{\realOrderTopology}{\intervalclosedrel{\realOrderRel}{\reals}{a}{b}}$ and $T$ and
        $f(a) = x$ and $f(b) = y$
        by \cref{path_in_space}.
    Let $I = \intervalclosedrel{\realOrderRel}{\reals}{a}{b}$.
    Let $T_I = \subspaceTopology{\realOrderTopology}{I}$.
    Then $\realOrderTopology$ is a topology on $\reals$ by
        \cref{standard_basis_is_basis,standard_topology_reals,basis_topology_is_topology,real_order_topology_eq_standard}.
    Hence $I$ is a subspace of $\reals$ with respect to $\realOrderTopology$ by \cref{intervalclosed_rel,subseteq,subspace_of}.
    Therefore $T_I$ is a topology on $I$ by \cref{subspace_topology_is_topology,subspace_of}.
    Let $c_0 = a + b$.
    $c_0\in\reals$ by \cref{real_add_closed}.
    Let $k = \{(t,c_0) \mid t\in I\}$.
    $k$ is a function from $I$ to $\reals$ by \cref{constant_map_function}.
    For all $t\in I$ we have $k(t) = c_0$.
    Hence $k$ is continuous from $I$ to $\reals$ with respect to $T_I$ and $\standardTopologyR$
        by \cref{constant_continuous,real_order_topology_eq_standard}.
    Let $j = \identity{I}$.
    $j$ is continuous from $I$ to $\reals$ with respect to $T_I$ and $\standardTopologyR$
        by \cref{inclusion_continuous,real_order_topology_eq_standard}.
    Let $r = \{(t, k(t) - j(t)) \mid t\in I\}$.
    $r$ is continuous from $I$ to $\reals$ with respect to $T_I$ and $\standardTopologyR$
        by \cref{real_function_subtraction_continuous}.
    Therefore $r$ is continuous from $I$ to $\reals$ with respect to $T_I$ and $\realOrderTopology$
        by \cref{real_order_topology_eq_standard}.
    Show that for all $t$ such that $t\in I$ we have $r(t)\in I$.
    \begin{subproof}
            Fix $t$ such that $t\in I$.
            Then $t\in\reals$ by \cref{intervalclosed_rel}.
            Hence $a\mathrel{\realOrderRel}t$ or $t = a$ by \cref{intervalclosed_rel}.
            Then $t\mathrel{\realOrderRel}b$ or $t = b$ by \cref{intervalclosed_rel}.
            Therefore $j(t) = t$ by \cref{id_elem_funs,id_iff,funs_is_function,function_apply_iff}.
            $(t, k(t) - j(t))\in r$.
            Then $r(t) = k(t) - j(t)$ by \cref{continuous,funs_is_function,function_apply_iff}.
            Hence $k(t) = c_0$.
            Then $r(t) = c_0 - t$.
            Therefore $r(t)\in\reals$ by \cref{real_add_inverse,real_add_closed,real_sub}.
            Show that $a\mathrel{\realOrderRel}r(t)$ or $r(t) = a$.
            \begin{subproof}
                \begin{byCase}
                \caseOf{$t = b$.}
                    $b + \neg{b} = 0$ by \cref{real_add_inverse}.
                    Then $c_0 - b = (a + b) + \neg{b}$ by \cref{real_sub,real_add_eq_subst}.
                    $a\in\reals$ and $b\in\reals$.
                    $\neg{b}\in\reals$ by \cref{real_add_inverse}.
                    $(a + b) + \neg{b} = a + (b + \neg{b})$ by \cref{real_add_assoc}.
                    Hence $r(t) = a + 0$ by \cref{real_add_eq_subst}.
                    Then $r(t) = a$ by \cref{real_add_ident,real_add_comm}.
                \caseOf{$t\mathrel{\realOrderRel}b$.}
                    Then $t < b$ by \cref{real_order_relation_elim}.
                    $b\in\reals$.
                    $\neg{t}\in\reals$ by \cref{real_add_inverse}.
                    Hence $t + \neg{t} < b + \neg{t}$ by \cref{real_lt_add}.
                    $t + \neg{t} = 0$ by \cref{real_add_inverse}.
                    Then $0 < b + \neg{t}$ by \cref{real_lt_subst_left}.
                    Hence $0 < b - t$ by \cref{real_sub,real_lt_subst_right}.
                    $0\in\reals$ by \cref{real_add_ident}.
                    $a\in\reals$.
                    Then $b - t\in\reals$ by \cref{real_sub,real_add_closed}.
                    Hence $0 + a < (b - t) + a$ by \cref{real_lt_add}.
                    $0 + a = a$ by \cref{real_add_ident}.
                    $(b - t) + a = a + (b - t)$ by \cref{real_add_comm}.
                    Therefore $a < a + (b - t)$ by \cref{real_lt_subst_left,real_lt_subst_right}.
                    $a + (b - t) = a + (b + \neg{t})$ by \cref{real_sub,real_add_eq_subst}.
                    $a + (b + \neg{t}) = (a + b) + \neg{t}$ by \cref{real_add_assoc}.
                    $(a + b) + \neg{t} = (a + b) - t$ by \cref{real_sub}.
                    Then $a + (b - t) = c_0 - t$.
                    Hence $a < c_0 - t$ by \cref{real_lt_subst_right}.
                    Therefore $a\mathrel{\realOrderRel}r(t)$ by \cref{real_order_relation_intro}.
                \end{byCase}
            \end{subproof}
            Show that $r(t)\mathrel{\realOrderRel}b$ or $r(t) = b$.
            \begin{subproof}
                \begin{byCase}
                \caseOf{$t = a$.}
                    Then $c_0 = b + a$ by \cref{real_add_comm,real_add_eq_subst}.
                    Hence $c_0 - a = (b + a) + \neg{a}$ by \cref{real_sub,real_add_eq_subst}.
                    $a\in\reals$ and $b\in\reals$.
                    $\neg{a}\in\reals$ by \cref{real_add_inverse}.
                    $(b + a) + \neg{a} = b + (a + \neg{a})$ by \cref{real_add_assoc}.
                    $a + \neg{a} = 0$ by \cref{real_add_inverse}.
                    Therefore $r(t) = b + 0$ by \cref{real_add_eq_subst}.
                    Then $r(t) = b$ by \cref{real_add_ident,real_add_comm}.
                \caseOf{$a\mathrel{\realOrderRel}t$.}
                    Then $a < t$ by \cref{real_order_relation_elim}.
                    $a\in\reals$ and $t\in\reals$ and $b\in\reals$.
                    $\neg{t}\in\reals$ by \cref{real_add_inverse}.
                    Hence $b - t\in\reals$ by \cref{real_sub,real_add_closed}.
                    Then $a + (b - t) < t + (b - t)$ by \cref{real_lt_add}.
                    $t + (b - t) = t + (b + \neg{t})$ by \cref{real_sub,real_add_eq_subst}.
                    $t + (b + \neg{t}) = (t + b) + \neg{t}$ by \cref{real_add_assoc}.
                    $(t + b) + \neg{t} = (b + t) + \neg{t}$ by \cref{real_add_comm}.
                    $(b + t) + \neg{t} = b + (t + \neg{t})$ by \cref{real_add_assoc}.
                    $t + \neg{t} = 0$ by \cref{real_add_inverse}.
                    Therefore $t + (b - t) = b + 0$ by \cref{real_add_eq_subst}.
                    Then $t + (b - t) = b$ by \cref{real_add_ident,real_add_comm}.
                    Hence $a + (b - t) < b$ by \cref{real_lt_subst_right}.
                    $a + (b - t) = a + (b + \neg{t})$ by \cref{real_sub,real_add_eq_subst}.
                    $a + (b + \neg{t}) = (a + b) + \neg{t}$ by \cref{real_add_assoc}.
                    $(a + b) + \neg{t} = (a + b) - t$ by \cref{real_sub}.
                    Then $a + (b - t) = c_0 - t$.
                    Therefore $c_0 - t < b$ by \cref{real_lt_subst_left}.
                    Hence $r(t)\mathrel{\realOrderRel}b$ by \cref{real_order_relation_intro}.
                \end{byCase}
            \end{subproof}
            Therefore $r(t)\in I$ by \cref{intervalclosed_rel}.
    \end{subproof}
    Hence $\img{r}{I}\subseteq I$
        by \cref{img_iff,continuous,funs_is_function,function_apply_iff,subseteq}.
    Therefore $r$ is continuous from $I$ to $I$ with respect to $T_I$ and $T_I$ by \cref{continuous_restrict_range}.
    Let $g = f\circ r$.
    Then $f$ is composable with $r$ by \cref{continuous,funs,img_dom,subseteq,funs_is_function}.
    Hence $g$ is continuous from $I$ to $X$ with respect to $T_I$ and $T$ by \cref{continuous_composite}.
    Show that $r(a) = b$ and $r(b) = a$.
    \begin{subproof}
        Then $a\in I$ and $b\in I$ by \cref{intervalclosed_rel}.
        $(a, k(a) - j(a))\in r$ and $(b, k(b) - j(b))\in r$.
        Hence $r(a) = k(a) - j(a)$ and $r(b) = k(b) - j(b)$
            by \cref{continuous,funs_is_function,function_apply_iff}.
        Then $k(a) = c_0$ and $k(b) = c_0$.
        Therefore $j(a) = a$ and $j(b) = b$
            by \cref{id_elem_funs,id_iff,funs_is_function,function_apply_iff}.
        Then $r(a) = c_0 - a$ and $r(b) = c_0 - b$.
        Hence $c_0 = b + a$ by \cref{real_add_comm,real_add_eq_subst}.
        Then $c_0 - a = (b + a) + \neg{a}$ by \cref{real_sub,real_add_eq_subst}.
        $\neg{a}\in\reals$ by \cref{real_add_inverse}.
        $(b + a) + \neg{a} = b + (a + \neg{a})$ by \cref{real_add_assoc}.
        $a + \neg{a} = 0$ by \cref{real_add_inverse}.
        Therefore $r(a) = b + 0$ by \cref{real_add_eq_subst}.
        Then $r(a) = b$ by \cref{real_add_ident,real_add_comm}.
        Hence $c_0 - b = (a + b) + \neg{b}$ by \cref{real_sub,real_add_eq_subst}.
        $\neg{b}\in\reals$ by \cref{real_add_inverse}.
        $(a + b) + \neg{b} = a + (b + \neg{b})$ by \cref{real_add_assoc}.
        $b + \neg{b} = 0$ by \cref{real_add_inverse}.
        Therefore $r(b) = a + 0$ by \cref{real_add_eq_subst}.
        Then $r(b) = a$ by \cref{real_add_ident,real_add_comm}.
    \end{subproof}
    Hence $g(a) = f(r(a))$ and $g(b) = f(r(b))$
        by \cref{intervalclosed_rel,continuous,funs,funs_is_function,subseteq,circ_apply}.
    Then $g(a) = y$ and $g(b) = x$.
    Therefore $g$ is a path in $X$ from $y$ to $x$ with respect to $T$ by \cref{path_in_space}.
    Hence there exists $g$ such that $g$ is a path in $X$ from $y$ to $x$ with respect to $T$.
\end{proof}

% from §25 helper lemma 17: concatenate paths
\begin{lemma}\label{path_concatenation}
    Assume $T$ is a topology on $X$.
    Suppose $f$ is a path in $X$ from $x$ to $y$ with respect to $T$.
    Suppose $g$ is a path in $X$ from $y$ to $z$ with respect to $T$.
    Then there exists $h$ such that $h$ is a path in $X$ from $x$ to $z$ with respect to $T$.
\end{lemma}
\begin{proof}
    Take $a,b\in\reals$ such that
        $a\mathrel{\realOrderRel} b$ and
        $f$ is continuous from $\intervalclosedrel{\realOrderRel}{\reals}{a}{b}$ to $X$ with respect to
            $\subspaceTopology{\realOrderTopology}{\intervalclosedrel{\realOrderRel}{\reals}{a}{b}}$ and $T$ and
        $f(a) = x$ and $f(b) = y$
        by \cref{path_in_space}.
    Let $I_1 = \intervalclosedrel{\realOrderRel}{\reals}{a}{b}$.
    Let $T_1 = \subspaceTopology{\realOrderTopology}{I_1}$.
    Take $c,d\in\reals$ such that
        $c\mathrel{\realOrderRel} d$ and
        $g$ is continuous from $\intervalclosedrel{\realOrderRel}{\reals}{c}{d}$ to $X$ with respect to
            $\subspaceTopology{\realOrderTopology}{\intervalclosedrel{\realOrderRel}{\reals}{c}{d}}$ and $T$ and
        $g(c) = y$ and $g(d) = z$
        by \cref{path_in_space}.
    Let $I_2 = \intervalclosedrel{\realOrderRel}{\reals}{c}{d}$.
    Let $T_2 = \subspaceTopology{\realOrderTopology}{I_2}$.
    Let $s = c - b$.
    $b\in\reals$ and $c\in\reals$.
    $\neg{b}\in\reals$ by \cref{real_add_inverse}.
    $s\in\reals$ by \cref{real_sub,real_add_closed}.
    Let $e = b + (d - c)$.
    $\neg{c}\in\reals$ by \cref{real_add_inverse}.
    Then $d - c\in\reals$ by \cref{real_sub,real_add_closed}.
    Hence $e\in\reals$ by \cref{real_add_closed}.
    Let $I_3 = \intervalclosedrel{\realOrderRel}{\reals}{b}{e}$.
    Let $I = \intervalclosedrel{\realOrderRel}{\reals}{a}{e}$.
    Let $T_I = \subspaceTopology{\realOrderTopology}{I}$.
    Let $T_3 = \subspaceTopology{\realOrderTopology}{I_3}$.
    Now $\realOrderTopology$ is a topology on $\reals$ by
        \cref{standard_basis_is_basis,standard_topology_reals,basis_topology_is_topology,real_order_topology_eq_standard}.
    Hence $I$ is a subspace of $\reals$ with respect to $\realOrderTopology$
        by \cref{intervalclosed_rel,subseteq,subspace_of}.
    Therefore $T_I$ is a topology on $I$ by \cref{subspace_topology_is_topology,subspace_of}.
    $I_1\subseteq\reals$ by \cref{intervalclosed_rel,subseteq}.
    $I_2\subseteq\reals$ by \cref{intervalclosed_rel,subseteq}.
    $I_3\subseteq\reals$ by \cref{intervalclosed_rel,subseteq}.
    Hence $I_1$ is a subspace of $\reals$ with respect to $\realOrderTopology$ and
        $I_2$ is a subspace of $\reals$ with respect to $\realOrderTopology$ and
        $I_3$ is a subspace of $\reals$ with respect to $\realOrderTopology$ by \cref{subspace_of}.
    Therefore $T_1$ is a topology on $I_1$ and $T_2$ is a topology on $I_2$ and $T_3$ is a topology on $I_3$
        by \cref{subspace_topology_is_topology,subspace_of}.
    Let $k = \{(t,s) \mid t\in I_3\}$.
    $k$ is a function from $I_3$ to $\reals$ by \cref{constant_map_function}.
    For all $t\in I_3$ we have $k(t) = s$.
    Hence $k$ is continuous from $I_3$ to $\reals$ with respect to $T_3$ and $\standardTopologyR$
        by \cref{constant_continuous,real_order_topology_eq_standard}.
    Let $j = \identity{I_3}$.
    $j$ is continuous from $I_3$ to $\reals$ with respect to $T_3$ and $\standardTopologyR$
        by \cref{inclusion_continuous,real_order_topology_eq_standard}.
    Let $u = \{(t, j(t) + k(t)) \mid t\in I_3\}$.
    $u$ is continuous from $I_3$ to $\reals$ with respect to $T_3$ and $\standardTopologyR$
        by \cref{real_function_addition_continuous}.
    Hence $u$ is continuous from $I_3$ to $\reals$ with respect to $T_3$ and $\realOrderTopology$
        by \cref{real_order_topology_eq_standard}.
    Show that for all $t$ such that $t\in I_3$ we have $u(t)\in I_2$.
    \begin{subproof}
            Fix $t$ such that $t\in I_3$.
            Thus $t\in\reals$ by \cref{intervalclosed_rel}.
            Thus $b\mathrel{\realOrderRel}t$ or $t = b$ by \cref{intervalclosed_rel}.
            Thus $t\mathrel{\realOrderRel}e$ or $t = e$ by \cref{intervalclosed_rel}.
            Thus $j(t) = t$ by \cref{id_elem_funs,id_iff,funs_is_function,function_apply_iff}.
            $(t, j(t) + k(t))\in u$.
            Thus $u(t) = j(t) + k(t)$ by \cref{continuous,funs_is_function,function_apply_iff}.
            Thus $u(t) = t + s$.
            $s\in\reals$.
            Thus $u(t)\in\reals$ by \cref{real_add_closed}.
            Show that $c\mathrel{\realOrderRel}u(t)$ or $u(t) = c$.
            \begin{subproof}
                \begin{byCase}
                \caseOf{$t = b$.}
                    $s = c - b$.
                    $c - b = c + \neg{b}$ by \cref{real_sub}.
                    Thus $b + s = b + (c + \neg{b})$ by \cref{real_add_eq_subst}.
                    $b + (c + \neg{b}) = (b + c) + \neg{b}$ by \cref{real_add_assoc}.
                    $(b + c) + \neg{b} = (c + b) + \neg{b}$ by \cref{real_add_comm}.
                    $(c + b) + \neg{b} = c + (b + \neg{b})$ by \cref{real_add_assoc}.
                    $b + \neg{b} = 0$ by \cref{real_add_inverse}.
                    Thus $b + s = c + 0$ by \cref{real_add_eq_subst}.
                    Hence $b + s = c$ by \cref{real_add_ident,real_add_comm}.
                    Therefore $u(t) = c$.
                \caseOf{$b\mathrel{\realOrderRel}t$.}
                    Thus $b < t$ by \cref{real_order_relation_elim}.
                    $b\in\reals$ and $t\in\reals$ and $s\in\reals$.
                    Thus $b + s < t + s$ by \cref{real_lt_add}.
                    Thus $b + s = b + (c + \neg{b})$ by \cref{real_sub,real_add_eq_subst}.
                    $b + (c + \neg{b}) = (b + c) + \neg{b}$ by \cref{real_add_assoc}.
                    $(b + c) + \neg{b} = (c + b) + \neg{b}$ by \cref{real_add_comm}.
                    $(c + b) + \neg{b} = c + (b + \neg{b})$ by \cref{real_add_assoc}.
                    $b + \neg{b} = 0$ by \cref{real_add_inverse}.
                    Thus $b + s = c + 0$ by \cref{real_add_eq_subst}.
                    Hence $b + s = c$ by \cref{real_add_ident,real_add_comm}.
                    Therefore $c < u(t)$ by \cref{real_lt_subst_left,real_lt_subst_right}.
                    Then $c\mathrel{\realOrderRel}u(t)$ by \cref{real_order_relation_intro}.
                \end{byCase}
            \end{subproof}
            Show that $u(t)\mathrel{\realOrderRel}d$ or $u(t) = d$.
            \begin{subproof}
                \begin{byCase}
                \caseOf{$t = e$.}
                    $s = c - b$.
                    $e = b + (d - c)$.
                    $d - c = d + \neg{c}$ by \cref{real_sub}.
                    $c - b = c + \neg{b}$ by \cref{real_sub}.
                    Thus $e + s = (b + (d + \neg{c})) + (c + \neg{b})$ by \cref{real_add_eq_subst}.
                    $(b + (d + \neg{c})) + (c + \neg{b}) = b + ((d + \neg{c}) + (c + \neg{b}))$ by \cref{real_add_assoc}.
                    $(d + \neg{c}) + (c + \neg{b}) = ((d + \neg{c}) + c) + \neg{b}$ by \cref{real_add_assoc}.
                    $(d + \neg{c}) + c = d + (\neg{c} + c)$ by \cref{real_add_assoc}.
                    $\neg{c} + c = 0$ by \cref{real_add_comm,real_add_inverse}.
                    Thus $d + (\neg{c} + c) = d + 0$ by \cref{real_add_eq_subst}.
                    Thus $(d + \neg{c}) + c = d + 0$ by \cref{real_add_eq_subst}.
                    Thus $((d + \neg{c}) + c) + \neg{b} = (d + 0) + \neg{b}$ by \cref{real_add_eq_subst}.
                    $d + 0 = d$ by \cref{real_add_ident,real_add_comm}.
                    Thus $(d + 0) + \neg{b} = d + \neg{b}$ by \cref{real_add_eq_subst}.
                    Thus $(d + \neg{c}) + (c + \neg{b}) = d + \neg{b}$.
                    Thus $e + s = b + (d + \neg{b})$ by \cref{real_add_eq_subst}.
                    $b + (d + \neg{b}) = (b + d) + \neg{b}$ by \cref{real_add_assoc}.
                    $(b + d) + \neg{b} = (d + b) + \neg{b}$ by \cref{real_add_comm}.
                    $(d + b) + \neg{b} = d + (b + \neg{b})$ by \cref{real_add_assoc}.
                    $b + \neg{b} = 0$ by \cref{real_add_inverse}.
                    Thus $e + s = d + 0$ by \cref{real_add_eq_subst}.
                    Hence $e + s = d$ by \cref{real_add_ident,real_add_comm}.
                    Therefore $u(t) = d$.
                \caseOf{$t\mathrel{\realOrderRel}e$.}
                    Thus $t < e$ by \cref{real_order_relation_elim}.
                    $t\in\reals$ and $e\in\reals$ and $s\in\reals$.
                    Thus $t + s < e + s$ by \cref{real_lt_add}.
                    $s = c - b$.
                    $e = b + (d - c)$.
                    $d - c = d + \neg{c}$ by \cref{real_sub}.
                    $c - b = c + \neg{b}$ by \cref{real_sub}.
                    Thus $e + s = (b + (d + \neg{c})) + (c + \neg{b})$ by \cref{real_add_eq_subst}.
                    $(b + (d + \neg{c})) + (c + \neg{b}) = b + ((d + \neg{c}) + (c + \neg{b}))$ by \cref{real_add_assoc}.
                    $(d + \neg{c}) + (c + \neg{b}) = ((d + \neg{c}) + c) + \neg{b}$ by \cref{real_add_assoc}.
                    $(d + \neg{c}) + c = d + (\neg{c} + c)$ by \cref{real_add_assoc}.
                    $\neg{c} + c = 0$ by \cref{real_add_comm,real_add_inverse}.
                    Thus $d + (\neg{c} + c) = d + 0$ by \cref{real_add_eq_subst}.
                    Thus $(d + \neg{c}) + c = d + 0$ by \cref{real_add_eq_subst}.
                    Thus $((d + \neg{c}) + c) + \neg{b} = (d + 0) + \neg{b}$ by \cref{real_add_eq_subst}.
                    $d + 0 = d$ by \cref{real_add_ident,real_add_comm}.
                    Thus $(d + 0) + \neg{b} = d + \neg{b}$ by \cref{real_add_eq_subst}.
                    Thus $(d + \neg{c}) + (c + \neg{b}) = d + \neg{b}$.
                    Thus $e + s = b + (d + \neg{b})$ by \cref{real_add_eq_subst}.
                    $b + (d + \neg{b}) = (b + d) + \neg{b}$ by \cref{real_add_assoc}.
                    $(b + d) + \neg{b} = (d + b) + \neg{b}$ by \cref{real_add_comm}.
                    $(d + b) + \neg{b} = d + (b + \neg{b})$ by \cref{real_add_assoc}.
                    $b + \neg{b} = 0$ by \cref{real_add_inverse}.
                    Thus $e + s = d + 0$ by \cref{real_add_eq_subst}.
                    Hence $e + s = d$ by \cref{real_add_ident,real_add_comm}.
                    Therefore $u(t) < d$ by \cref{real_lt_subst_right}.
                    Then $u(t)\mathrel{\realOrderRel}d$ by \cref{real_order_relation_intro}.
                \end{byCase}
            \end{subproof}
            Thus $u(t)\in I_2$ by \cref{intervalclosed_rel}.
    \end{subproof}
    Therefore $\img{u}{I_3}\subseteq I_2$
        by \cref{img_iff,continuous,funs_is_function,function_apply_iff,subseteq}.
    Hence $u$ is continuous from $I_3$ to $I_2$ with respect to $T_3$ and $T_2$
        by \cref{continuous_restrict_range}.
    Let $g_1 = g\circ u$.
    Then $g$ is composable with $u$ by \cref{continuous,funs,img_dom,subseteq,funs_is_function}.
    Hence $g_1$ is continuous from $I_3$ to $X$ with respect to $T_3$ and $T$ by \cref{continuous_composite}.
    Thus $c < d$ by \cref{real_order_relation_elim}.
    $c\in\reals$ and $d\in\reals$.
    Thus $c + \neg{c} < d + \neg{c}$ by \cref{real_lt_add}.
    $c + \neg{c} = 0$ by \cref{real_add_inverse}.
    Then $0 < d + \neg{c}$ by \cref{real_lt_subst_left}.
    $d + \neg{c} = d - c$ by \cref{real_sub}.
    Hence $0 < d - c$ by \cref{real_lt_subst_right}.
    $b\in\reals$ and $d - c\in\reals$.
    $0\in\reals$ by \cref{real_add_ident}.
    Therefore $0 + b < (d - c) + b$ by \cref{real_lt_add}.
    $0 + b = b$ by \cref{real_add_comm,real_add_ident}.
    Then $b < (d - c) + b$ by \cref{real_lt_subst_left}.
    $(d - c) + b = b + (d - c)$ by \cref{real_add_comm}.
    Hence $b < b + (d - c)$ by \cref{real_lt_subst_right}.
    Therefore $b < e$ by \cref{real_lt_subst_right}.
    Then $b\mathrel{\realOrderRel}e$ by \cref{real_order_relation_intro}.
    Thus $b\in I_3$ and $e\in I_3$ by \cref{intervalclosed_rel}.
    $(b, j(b) + k(b))\in u$ and $(e, j(e) + k(e))\in u$.
    Thus $u(b) = j(b) + k(b)$ and $u(e) = j(e) + k(e)$
        by \cref{continuous,funs_is_function,function_apply_iff}.
    Thus $j(b) = b$ and $j(e) = e$
        by \cref{id_elem_funs,id_iff,funs_is_function,function_apply_iff}.
    Thus $u(b) = b + s$ and $u(e) = e + s$.
    Thus $b + s = b + (c + \neg{b})$ by \cref{real_sub,real_add_eq_subst}.
    $b + (c + \neg{b}) = (b + c) + \neg{b}$ by \cref{real_add_assoc}.
    $(b + c) + \neg{b} = (c + b) + \neg{b}$ by \cref{real_add_comm}.
    $(c + b) + \neg{b} = c + (b + \neg{b})$ by \cref{real_add_assoc}.
    $b + \neg{b} = 0$ by \cref{real_add_inverse}.
    Thus $b + s = c + 0$ by \cref{real_add_eq_subst}.
    Hence $u(b) = c$ by \cref{real_add_ident,real_add_comm}.
    Thus $e + s = (b + (d + \neg{c})) + (c + \neg{b})$ by \cref{real_sub,real_add_eq_subst}.
    $d + \neg{c}\in\reals$ and $c + \neg{b}\in\reals$ by \cref{real_add_inverse,real_add_closed}.
    $(b + (d + \neg{c})) + (c + \neg{b}) = b + ((d + \neg{c}) + (c + \neg{b}))$ by \cref{real_add_assoc}.
    $(d + \neg{c}) + (c + \neg{b}) = (d + (\neg{c} + c)) + \neg{b}$ by \cref{real_add_assoc}.
    $\neg{c} + c = 0$ by \cref{real_add_comm,real_add_inverse}.
    Thus $(d + \neg{c}) + (c + \neg{b}) = (d + 0) + \neg{b}$ by \cref{real_add_eq_subst}.
    $d + 0 = d$ by \cref{real_add_ident,real_add_comm}.
    Thus $(d + 0) + \neg{b} = d + \neg{b}$ by \cref{real_add_eq_subst}.
    Thus $(d + \neg{c}) + (c + \neg{b}) = d + \neg{b}$ by \cref{real_add_eq_subst}.
    Thus $e + s = b + (d + \neg{b})$ by \cref{real_add_eq_subst}.
    $b + (d + \neg{b}) = (b + d) + \neg{b}$ by \cref{real_add_assoc}.
    $(b + d) + \neg{b} = (d + b) + \neg{b}$ by \cref{real_add_comm}.
    $(d + b) + \neg{b} = d + (b + \neg{b})$ by \cref{real_add_assoc}.
    $b + \neg{b} = 0$ by \cref{real_add_inverse}.
    Thus $e + s = d + 0$ by \cref{real_add_eq_subst}.
    Therefore $u(e) = d$ by \cref{real_add_ident,real_add_comm}.
    Then $g_1(b) = g(u(b))$ and $g_1(e) = g(u(e))$
        by \cref{intervalclosed_rel,continuous,funs,funs_is_function,subseteq,circ_apply}.
    Hence $g_1(b) = y$ and $g_1(e) = z$.
    Show that for all $t$ such that $t\in I_1\union I_3$ we have $t\in I$.
    \begin{subproof}
            Fix $t$ such that $t\in I_1\union I_3$.
            Thus $t\in I_1$ or $t\in I_3$ by \cref{union_iff}.
            \begin{byCase}
            \caseOf{$t\in I_1$.}
                Thus $a\mathrel{\realOrderRel}t$ or $t = a$ by \cref{intervalclosed_rel}.
                Thus $t\mathrel{\realOrderRel}b$ or $t = b$ by \cref{intervalclosed_rel}.
                \begin{byCase}
                \caseOf{$t = b$.}
                    Thus $t\mathrel{\realOrderRel}e$ by \cref{real_lt_subst_right}.
                    Thus $t\in I$ by \cref{intervalclosed_rel}.
                \caseOf{$t\mathrel{\realOrderRel}b$.}
                    Thus $t < b$ by \cref{real_order_relation_elim}.
                    Thus $b < e$ by \cref{real_order_relation_elim}.
                    Thus $t\in\reals$ by \cref{intervalclosed_rel}.
                    $b\in\reals$ and $e\in\reals$.
                    Thus $t < e$ by \cref{real_lt_trans}.
                    Thus $t\mathrel{\realOrderRel}e$ by \cref{real_order_relation_intro}.
                    Thus $t\in I$ by \cref{intervalclosed_rel}.
                \end{byCase}
            \caseOf{$t\in I_3$.}
                Thus $b\mathrel{\realOrderRel}t$ or $t = b$ by \cref{intervalclosed_rel}.
                Thus $t\mathrel{\realOrderRel}e$ or $t = e$ by \cref{intervalclosed_rel}.
                \begin{byCase}
                \caseOf{$t = b$.}
                    Thus $a\mathrel{\realOrderRel}t$ by \cref{real_lt_subst_right}.
                    Thus $t\in I$ by \cref{intervalclosed_rel}.
                \caseOf{$b\mathrel{\realOrderRel}t$.}
                    Thus $a < b$ by \cref{real_order_relation_elim}.
                    Thus $b < t$ by \cref{real_order_relation_elim}.
                    $a\in\reals$ and $b\in\reals$.
                    Thus $t\in\reals$ by \cref{intervalclosed_rel}.
                    Thus $a < t$ by \cref{real_lt_trans}.
                    Thus $a\mathrel{\realOrderRel}t$ by \cref{real_order_relation_intro}.
                    Thus $t\in I$ by \cref{intervalclosed_rel}.
                \end{byCase}
            \end{byCase}
    \end{subproof}
    Hence $I_1\union I_3\subseteq I$ by \cref{subseteq}.
    Show that for all $t$ such that $t\in I$ we have $t\in I_1\union I_3$.
    \begin{subproof}
            Fix $t$ such that $t\in I$.
            Thus $t\in\reals$ by \cref{intervalclosed_rel}.
            Thus $a\mathrel{\realOrderRel}t$ or $t = a$ by \cref{intervalclosed_rel}.
            Thus $t\mathrel{\realOrderRel}e$ or $t = e$ by \cref{intervalclosed_rel}.
            \begin{byCase}
            \caseOf{$t = b$.}
                Thus $t\in I_1$ by \cref{intervalclosed_rel}.
                Thus $t\in I_1\union I_3$ by \cref{union_iff}.
            \caseOf{$t\neq b$.}
                $b\in\reals$ and $t\in\reals$.
                $t\neq b$.
                Thus $t < b$ or $b < t$ by \cref{real_lt_trichotomy}.
                \begin{byCase}
                \caseOf{$t < b$.}
                    Thus $t\mathrel{\realOrderRel}b$ by \cref{real_order_relation_intro}.
                    Thus $t\in I_1$ by \cref{intervalclosed_rel}.
                    Thus $t\in I_1\union I_3$ by \cref{union_iff}.
                \caseOf{$b < t$.}
                    Thus $b\mathrel{\realOrderRel}t$ by \cref{real_order_relation_intro}.
                    Thus $t\in I_3$ by \cref{intervalclosed_rel}.
                    Thus $t\in I_1\union I_3$ by \cref{union_iff}.
                \end{byCase}
            \end{byCase}
    \end{subproof}
    Therefore $I\subseteq I_1\union I_3$ by \cref{subseteq}.
    Hence $I = I_1\union I_3$ by \cref{subseteq_antisymmetric}.
    Then $I_1\subseteq I$ and $I_3\subseteq I$ by \cref{union_upper_left,union_upper_right,subseteq}.
    $\realOrderRel$ is a simple order relation on $\reals$ and $\reals$ has more than one element
        by \cref{reals_linear_continuum,linear_continuum}.
    Hence $\realOrderTopology$ is the order topology on $\reals$ with respect to $\realOrderRel$
        by \cref{real_order_topology,order_topology}.
    $I$ is a subspace of $\reals$ with respect to $\realOrderTopology$ by \cref{subspace_of}.
    $I_1$ is closed in $\reals$ with respect to $\realOrderTopology$ by \cref{intervalclosedrel_closed_in_order_topology}.
    Therefore there exists $C$ such that $C$ is closed in $\reals$ with respect to $\realOrderTopology$ and $I_1 = C\inter I$.
    Hence $I_1$ is closed in $I$ with respect to $T_I$ by \cref{closed_in_subspace_iff,subspace_of}.
    $I_3$ is closed in $\reals$ with respect to $\realOrderTopology$ by \cref{intervalclosedrel_closed_in_order_topology}.
    Therefore there exists $C$ such that $C$ is closed in $\reals$ with respect to $\realOrderTopology$ and $I_3 = C\inter I$.
    Hence $I_3$ is closed in $I$ with respect to $T_I$ by \cref{closed_in_subspace_iff,subspace_of}.
    For all $t$ such that $t\in I_1\inter I_3$ we have $f(t) = g_1(t)$
        by \cref{inter,intervalclosed_rel,real_order_relation_elim,real_lt_asym,real_lt_subst_right}.
    Let $T_{I1} = \subspaceTopology{T_I}{I_1}$.
    Let $T_{I3} = \subspaceTopology{T_I}{I_3}$.
    Hence $T_{I1} = T_1$ by \cref{intervalclosed_rel,union_upper_left,subseteq,subspace_subspace_topology,subspace_of,subspace_topology}.
    Therefore $T_{I3} = T_3$ by \cref{intervalclosed_rel,union_upper_right,subseteq,subspace_subspace_topology,subspace_of,subspace_topology}.
    Hence $f$ is continuous from $I_1$ to $X$ with respect to $T_{I1}$ and $T$.
    $g_1$ is continuous from $I_3$ to $X$ with respect to $T_{I3}$ and $T$.
    Let $h = f\union g_1$.
    Therefore $h$ is continuous from $I$ to $X$ with respect to $T_I$ and $T$ by \cref{pasting_lemma}.
    $a\in I_1$ and $e\in I_3$ by \cref{intervalclosed_rel,real_order_relation_elim}.
    $(a, f(a))\in f$ and $(e, g_1(e))\in g_1$ by \cref{continuous,funs_elim,function_apply_iff}.
    Thus $h(a) = f(a)$ and $h(e) = g_1(e)$
        by \cref{continuous,funs_is_function,function_apply_iff,union_iff}.
    Hence $h(a) = x$ and $h(e) = z$.
    Then $a < b$ by \cref{real_order_relation_elim}.
    Hence $b < e$ by \cref{real_lt_subst_right}.
    $a\in\reals$ and $b\in\reals$ and $e\in\reals$.
    Therefore $a < e$ by \cref{real_lt_trans}.
    Then $a\mathrel{\realOrderRel}e$ by \cref{real_order_relation_intro}.
    Hence $h$ is continuous from $\intervalclosedrel{\realOrderRel}{\reals}{a}{e}$ to $X$ with respect to
        $\subspaceTopology{\realOrderTopology}{\intervalclosedrel{\realOrderRel}{\reals}{a}{e}}$ and $T$.
    Therefore $h$ is a path in $X$ from $x$ to $z$ with respect to $T$ by \cref{path_in_space}.
    Hence there exists $h$ such that $h$ is a path in $X$ from $x$ to $z$ with respect to $T$.
\end{proof}

% from §25 helper lemma 18: path relation is a binary relation
\begin{lemma}\label{path_relation_binary}
    Assume $T$ is a topology on $X$.
    Then $\pathRelation{T}{X}$ is a binary relation on $X$.
\end{lemma}
\begin{proof}
    Follows by \cref{path_relation,subseteq}.
\end{proof}

% from §25 helper lemma 19: path relation is reflexive
\begin{lemma}\label{path_relation_reflexive}
    Assume $T$ is a topology on $X$.
    Then $\pathRelation{T}{X}$ is reflexive on $X$.
\end{lemma}
\begin{proof}
    Show that for all $x\in X$ we have $x\mathrel{\pathRelation{T}{X}}x$.
    \begin{subproof}
        Fix $x\in X$.
        Thus $0\mathrel{\realOrderRel}1$ by
            \cref{real_add_ident,real_mul_ident,real_zero_lt_one,real_order_relation_intro}.
        Let $I = \intervalclosedrel{\realOrderRel}{\reals}{0}{1}$.
        Let $T_I = \subspaceTopology{\realOrderTopology}{I}$.
        For all $t$ such that $t\in I$ we have $t\in\reals$ by \cref{intervalclosed_rel}.
        Thus $I\subseteq\reals$ by \cref{subseteq}.
        Thus $\realOrderTopology$ is a topology on $\reals$ by
            \cref{standard_basis_is_basis,standard_topology_reals,basis_topology_is_topology,real_order_topology_eq_standard}.
        Thus $I$ is a subspace of $\reals$ with respect to $\realOrderTopology$ by \cref{subspace_of}.
        Thus $T_I$ is a topology on $I$ by \cref{subspace_topology_is_topology,subspace_of}.
        Let $f = \{(t,x) \mid t\in I\}$.
        $f$ is a function from $I$ to $X$ by \cref{constant_map_function}.
        For all $t\in I$ we have $f(t) = x$.
        Thus $f$ is continuous from $I$ to $X$ with respect to $T_I$ and $T$ by \cref{constant_continuous}.
        $0\in I$ and $1\in I$ by \cref{intervalclosed_rel,real_order_relation_elim}.
        Hence $f(0) = x$ and $f(1) = x$.
        Therefore $f$ is a path in $X$ from $x$ to $x$ with respect to $T$ by \cref{path_in_space}.
        Then $x\mathrel{\pathRelation{T}{X}}x$ by \cref{path_relation,path_in_space,times_tuple_intro,fst_eq,snd_eq}.
    \end{subproof}
    Thus $\pathRelation{T}{X}$ is reflexive on $X$ by \cref{reflexive_on}.
\end{proof}

% from §25 helper lemma 20: path relation is symmetric
\begin{lemma}\label{path_relation_symmetric}
    Assume $T$ is a topology on $X$.
    Then $\pathRelation{T}{X}$ is symmetric.
\end{lemma}
\begin{proof}
    Follows by \cref{path_relation,times_tuple_elim,fst_eq,snd_eq,path_reverse,times_tuple_intro,symmetric}.
\end{proof}

% from §25 helper lemma 21: path relation is transitive
\begin{lemma}\label{path_relation_transitive}
    Assume $T$ is a topology on $X$.
    Then $\pathRelation{T}{X}$ is transitive.
\end{lemma}
\begin{proof}
    Follows by \cref{path_relation,times_tuple_elim,fst_eq,snd_eq,path_concatenation,times_tuple_intro,transitive}.
\end{proof}
% from §25 helper lemma 22: path relation is an equivalence
\begin{lemma}\label{path_relation_equivalence}
    Assume $T$ is a topology on $X$.
    Then $\pathRelation{T}{X}$ is an equivalence on $X$.
\end{lemma}
\begin{proof}
    Follows by \cref{path_relation_binary,path_relation_reflexive,path_relation_transitive,path_relation_symmetric}.
\end{proof}

% from §25 helper lemma 23: union of path-connected subspaces with a common point
\begin{lemma}\label{union_path_connected_common_point}
    Assume $M\neq\emptyset$.
    Assume $T$ is a topology on $X$.
    Assume for all $A$ such that $A\in M$ we have $A\subseteq X$.
    Assume for all $A$ such that $A\in M$ we have
        $A$ is path connected with respect to $\subspaceTopology{T}{A}$.
    Assume $\inters{M}\neq\emptyset$.
    Then $\unions{M}$ is path connected with respect to $\subspaceTopology{T}{\unions{M}}$.
\end{lemma}
\begin{proof}
    Let $U = \unions{M}$.
    Let $T_U = \subspaceTopology{T}{U}$.
    Hence $U$ is a subspace of $X$ with respect to $T$ by \cref{unions_iff,subseteq,subspace_of}.
    Therefore $T_U$ is a topology on $U$ by \cref{subspace_topology_is_topology,subspace_of}.
    Show that for all $x,y\in U$ there exists $h$ such that
        $h$ is a path in $U$ from $x$ to $y$ with respect to $T_U$.
    \begin{subproof}
        Fix $x$.
        Fix $y$ such that $x\in U$ and $y\in U$.
        Take $A\in M$ such that $x\in A$ by \cref{unions_iff}.
        Take $B\in M$ such that $y\in B$ by \cref{unions_iff}.
        Take $p$ such that $p\in\inters{M}$ by \cref{nonempty_inhabited}.
        Then $p\in A$ and $p\in B$ by \cref{inters_iff_forall}.
        Let $T_A = \subspaceTopology{T}{A}$.
        Let $T_B = \subspaceTopology{T}{B}$.
        $A$ is path connected with respect to $T_A$ and $B$ is path connected with respect to $T_B$.
        Take $f$ such that $f$ is a path in $A$ from $x$ to $p$ with respect to $T_A$
            by \cref{path_connected}.
        Take $g$ such that $g$ is a path in $B$ from $p$ to $y$ with respect to $T_B$
            by \cref{path_connected}.
        Let $j_A = \identity{A}$.
        Let $j_B = \identity{B}$.
        Thus $j_A$ is continuous from $A$ to $U$ with respect to $T_A$ and $T_U$ and
            $j_B$ is continuous from $B$ to $U$ with respect to $T_B$ and $T_U$
            by \cref{unions_iff,subseteq,subspace_of,inclusion_continuous,subspace_subspace_topology}.
        Let $f_1 = j_A\circ f$.
        Let $g_1 = j_B\circ g$.
        Take $a,b\in\reals$ such that
                $a\mathrel{\realOrderRel} b$ and
                $f$ is continuous from $\intervalclosedrel{\realOrderRel}{\reals}{a}{b}$ to $A$ with respect to
                    $\subspaceTopology{\realOrderTopology}{\intervalclosedrel{\realOrderRel}{\reals}{a}{b}}$ and $T_A$ and
                $f(a) = x$ and $f(b) = p$
                by \cref{path_in_space}.
            Thus $j_A$ is composable with $f$ by \cref{continuous,funs,funs_ran,subseteq,funs_is_function}.
            Thus $f_1$ is continuous from $\intervalclosedrel{\realOrderRel}{\reals}{a}{b}$ to $U$ with respect to
                $\subspaceTopology{\realOrderTopology}{\intervalclosedrel{\realOrderRel}{\reals}{a}{b}}$ and $T_U$
                by \cref{continuous_composite}.
            Thus $f_1(a) = x$ and $f_1(b) = p$
                by \cref{intervalclosed_rel,continuous,funs,funs_is_function,funs_type_apply,subseteq,circ_apply,id_apply}.
        Hence $f_1$ is a path in $U$ from $x$ to $p$ with respect to $T_U$ by \cref{path_in_space}.
        Take $c,d\in\reals$ such that
                $c\mathrel{\realOrderRel} d$ and
                $g$ is continuous from $\intervalclosedrel{\realOrderRel}{\reals}{c}{d}$ to $B$ with respect to
                    $\subspaceTopology{\realOrderTopology}{\intervalclosedrel{\realOrderRel}{\reals}{c}{d}}$ and $T_B$ and
                $g(c) = p$ and $g(d) = y$
                by \cref{path_in_space}.
            Thus $j_B$ is composable with $g$ by \cref{continuous,funs,funs_ran,subseteq,funs_is_function}.
            Thus $g_1$ is continuous from $\intervalclosedrel{\realOrderRel}{\reals}{c}{d}$ to $U$ with respect to
                $\subspaceTopology{\realOrderTopology}{\intervalclosedrel{\realOrderRel}{\reals}{c}{d}}$ and $T_U$
                by \cref{continuous_composite}.
            Thus $g_1(c) = p$ and $g_1(d) = y$
                by \cref{intervalclosed_rel,continuous,funs,funs_is_function,funs_type_apply,subseteq,circ_apply,id_apply}.
        Therefore $g_1$ is a path in $U$ from $p$ to $y$ with respect to $T_U$ by \cref{path_in_space}.
        Hence there exists $h$ such that $h$ is a path in $U$ from $x$ to $y$ with respect to $T_U$
            by \cref{path_concatenation}.
    \end{subproof}
    Therefore $U$ is path connected with respect to $T_U$ by \cref{path_connected}.
\end{proof}


% from §25 helper lemma 24: initial segment of a path
\begin{lemma}\label{path_initial_segment}
    Assume $T$ is a topology on $X$.
    Suppose $f$ is a path in $X$ from $x$ to $y$ with respect to $T$.
    Suppose $t\in\dom{f}$.
    Then there exists $g$ such that $g$ is a path in $X$ from $x$ to $f(t)$ with respect to $T$.
\end{lemma}
\begin{proof}
    Take $a,b\in\reals$ such that
        $a\mathrel{\realOrderRel} b$ and
        $f$ is continuous from $\intervalclosedrel{\realOrderRel}{\reals}{a}{b}$ to $X$ with respect to
            $\subspaceTopology{\realOrderTopology}{\intervalclosedrel{\realOrderRel}{\reals}{a}{b}}$ and $T$ and
        $f(a) = x$ and $f(b) = y$
        by \cref{path_in_space}.
        Let $I = \intervalclosedrel{\realOrderRel}{\reals}{a}{b}$.
        Let $T_I = \subspaceTopology{\realOrderTopology}{I}$.
        Thus $\dom{f} = I$ by \cref{continuous,funs}.
        Then $t\in I$ by \cref{continuous,funs,subseteq}.
        Hence $f$ is a function by \cref{continuous,funs_is_function}.
    Therefore $x\in X$ by \cref{intervalclosed_rel,continuous,funs_type_apply}.
    \begin{byCase}
    \caseOf{$t = a$.}
        Then $f(t) = x$.
        Hence there exists $g$ such that $g$ is a path in $X$ from $x$ to $f(t)$ with respect to $T$
            by \cref{path_relation_reflexive,reflexive_on,path_relation,fst_eq,snd_eq}.
    \caseOf{$t\neq a$.}
        Then $a\mathrel{\realOrderRel} t$ by \cref{intervalclosed_rel}.
        Hence $t\in\reals$ by \cref{real_order_relation_elim}.
        Let $J = \intervalclosedrel{\realOrderRel}{\reals}{a}{t}$.
        $J\subseteq I$ by \cref{real_order_relation_simple,simple_order_relation,intervalclosedrel_convex,intervalclosed_rel,convex_in,subseteq}.
        Thus $J$ is a subspace of $I$ with respect to $T_I$ by \cref{continuous,subspace_of}.
        Let $T_J = \subspaceTopology{T_I}{J}$.
        Let $g = \restrl{f}{J}$.
        Hence $g$ is continuous from $J$ to $X$ with respect to $T_J$ and $T$
            by \cref{continuous_restrict_domain}.
        Now $\realOrderTopology$ is a topology on $\reals$ by
            \cref{standard_basis_is_basis,standard_topology_reals,basis_topology_is_topology,real_order_topology_eq_standard}.
        Hence $T_J = \subspaceTopology{\realOrderTopology}{J}$ by \cref{intervalclosed_rel,subseteq,subspace_subspace_topology}.
        Therefore $g$ is continuous from $J$ to $X$ with respect to $\subspaceTopology{\realOrderTopology}{J}$ and $T$.
        Then $g(a) = x$ and $g(t) = f(t)$ by \cref{intervalclosed_rel,funs_is_function,restrl_of_function_apply}.
        Hence there exists $g$ such that $g$ is a path in $X$ from $x$ to $f(t)$ with respect to $T$
            by \cref{path_in_space}.
    \end{byCase}
\end{proof}
% from §25 Item 9: path components are path connected
\begin{theorem}\label{path_component_path_connected}
    Assume $T$ is a topology on $X$.
    Suppose $C$ is a path component of $X$ with respect to $T$.
    Then $C$ is path connected with respect to $\subspaceTopology{T}{C}$.
\end{theorem}
\begin{proof}
    Take $x\in X$ such that $C = \equivalenceClass{\pathRelation{T}{X}}{x}$ by \cref{path_component}.
    Hence $C$ is a subspace of $X$ with respect to $T$
        by \cref{downward_closure_iff,path_relation,times_tuple_elim,subseteq,subspace_of}.
    Let $T_C = \subspaceTopology{T}{C}$.
    Therefore $T_C$ is a topology on $C$ by \cref{subspace_topology_is_topology,subspace_of}.
    Show that for all $y,z\in C$ there exists $h$ such that
        $h$ is a path in $C$ from $y$ to $z$ with respect to $T_C$.
    \begin{subproof}
        Fix $y$.
        Fix $z$ such that $y\in C$ and $z\in C$.
        Then $x\mathrel{\pathRelation{T}{X}}y$ and $x\mathrel{\pathRelation{T}{X}}z$
            by \cref{downward_closure_iff,path_relation_symmetric,symmetric}.
        Take $f$ such that $f$ is a path in $X$ from $x$ to $y$ with respect to $T$
            by \cref{path_relation,fst_eq,snd_eq}.
        Take $g$ such that $g$ is a path in $X$ from $x$ to $z$ with respect to $T$
            by \cref{path_relation,fst_eq,snd_eq}.
        Take $a,b\in\reals$ such that
            $a\mathrel{\realOrderRel} b$ and
            $f$ is continuous from $\intervalclosedrel{\realOrderRel}{\reals}{a}{b}$ to $X$ with respect to
                $\subspaceTopology{\realOrderTopology}{\intervalclosedrel{\realOrderRel}{\reals}{a}{b}}$ and $T$ and
            $f(a) = x$ and $f(b) = y$
            by \cref{path_in_space}.
        Let $I = \intervalclosedrel{\realOrderRel}{\reals}{a}{b}$.
        Let $T_I = \subspaceTopology{\realOrderTopology}{I}$.
        Take $c,d\in\reals$ such that
            $c\mathrel{\realOrderRel} d$ and
            $g$ is continuous from $\intervalclosedrel{\realOrderRel}{\reals}{c}{d}$ to $X$ with respect to
                $\subspaceTopology{\realOrderTopology}{\intervalclosedrel{\realOrderRel}{\reals}{c}{d}}$ and $T$ and
            $g(c) = x$ and $g(d) = z$
            by \cref{path_in_space}.
        Let $J = \intervalclosedrel{\realOrderRel}{\reals}{c}{d}$.
        Let $T_J = \subspaceTopology{\realOrderTopology}{J}$.
        For all $u$ such that $u\in\ran{f}$ we have $u\in C$
            by \cref{fun_ran_iff,function_apply_elim,ran_intro,continuous,funs_is_function,funs_ran,subseteq,path_initial_segment,times_tuple_intro,path_relation,fst_eq,snd_eq,path_relation_symmetric,symmetric,downward_closure_iff}.
        For all $u$ such that $u\in\ran{g}$ we have $u\in C$
            by \cref{fun_ran_iff,function_apply_elim,ran_intro,continuous,funs_is_function,funs_ran,subseteq,path_initial_segment,times_tuple_intro,path_relation,fst_eq,snd_eq,path_relation_symmetric,symmetric,downward_closure_iff}.
        Hence $f$ is continuous from $I$ to $C$ with respect to $T_I$ and $T_C$
            and $g$ is continuous from $J$ to $C$ with respect to $T_J$ and $T_C$
            by \cref{img_subseteq_ran,subseteq,continuous_restrict_range}.
        Therefore $f$ is a path in $C$ from $x$ to $y$ with respect to $T_C$
            and $g$ is a path in $C$ from $x$ to $z$ with respect to $T_C$
            by \cref{path_in_space}.
        Take $f_1$ such that $f_1$ is a path in $C$ from $y$ to $x$ with respect to $T_C$
            by \cref{path_reverse}.
        Hence there exists $h$ such that $h$ is a path in $C$ from $y$ to $z$ with respect to $T_C$
            by \cref{path_concatenation}.
    \end{subproof}
    Thus $C$ is path connected with respect to $T_C$ by \cref{path_connected}.
\end{proof}

% from §25 Item 10: distinct path components are disjoint
\begin{theorem}\label{path_components_disjoint}
    Assume $T$ is a topology on $X$.
    Suppose $C$ is a path component of $X$ with respect to $T$.
    Suppose $D$ is a path component of $X$ with respect to $T$.
    Suppose $C\neq D$.
    Then $C$ is disjoint from $D$.
\end{theorem}
\begin{proof}
    Take $x\in X$ such that $C = \equivalenceClass{\pathRelation{T}{X}}{x}$ by \cref{path_component}.
    Take $y\in X$ such that $D = \equivalenceClass{\pathRelation{T}{X}}{y}$ by \cref{path_component}.
    Hence $C$ is disjoint from $D$ by \cref{path_relation_equivalence,equivalenceon_equivclasses_diseq_implies_disjoint}.
\end{proof}

% from §25 Item 11: path components cover the space
\begin{theorem}\label{path_components_cover}
    Assume $T$ is a topology on $X$.
    Then for all $x\in X$ there exists $C$ such that
        $C$ is a path component of $X$ with respect to $T$ and $x\in C$.
\end{theorem}
\begin{proof}
    Fix $x\in X$.
    Follows by \cref{path_relation_equivalence,equivclasses_inhabited_equivalenceon,path_component}.
\end{proof}

% from §25 Item 12: path-connected subspaces intersect only one path component
\begin{theorem}\label{path_connected_subspace_single_component}
    Assume $T$ is a topology on $X$.
    Suppose $A$ is a subspace of $X$ with respect to $T$.
    Suppose $A$ is path connected with respect to $\subspaceTopology{T}{A}$.
    Suppose $A\neq\emptyset$.
    Then for all $C,D$ such that
        $C$ is a path component of $X$ with respect to $T$ and
        $D$ is a path component of $X$ with respect to $T$ and
        $A$ intersects $C$ and $A$ intersects $D$
        we have $C = D$.
\end{theorem}
\begin{proof}
    Fix $C$.
    Fix $D$ such that
        $C$ is a path component of $X$ with respect to $T$ and
        $D$ is a path component of $X$ with respect to $T$ and
        $A$ intersects $C$ and $A$ intersects $D$.
    Take $x$ such that $x\in A\inter C$ by \cref{intersects,nonempty_inhabited}.
    Take $y$ such that $y\in A\inter D$ by \cref{intersects,nonempty_inhabited}.
    Then $x\in A$ and $x\in C$ and $y\in A$ and $y\in D$ by \cref{inter}.
    Let $T_A = \subspaceTopology{T}{A}$.
    Take $f$ such that $f$ is a path in $A$ from $x$ to $y$ with respect to $T_A$
        by \cref{path_connected}.
    Take $a,b\in\reals$ such that
        $a\mathrel{\realOrderRel} b$ and
        $f$ is continuous from $\intervalclosedrel{\realOrderRel}{\reals}{a}{b}$ to $A$ with respect to
            $\subspaceTopology{\realOrderTopology}{\intervalclosedrel{\realOrderRel}{\reals}{a}{b}}$ and $T_A$ and
        $f(a) = x$ and $f(b) = y$
        by \cref{path_in_space}.
    Let $I = \intervalclosedrel{\realOrderRel}{\reals}{a}{b}$.
    Let $T_I = \subspaceTopology{\realOrderTopology}{I}$.
    Let $j = \identity{A}$.
    Therefore $j$ is continuous from $A$ to $X$ with respect to $T_A$ and $T$ by \cref{inclusion_continuous}.
    Let $f_1 = j\circ f$.
    Then $j$ is composable with $f$
        by \cref{continuous,funs,funs_is_function,funs_ran,subseteq}.
    Hence $f_1$ is continuous from $I$ to $X$ with respect to $T_I$ and $T$ by \cref{continuous_composite}.
    Therefore $f_1(a) = x$ and $f_1(b) = y$
        by \cref{intervalclosed_rel,continuous,funs,funs_is_function,funs_type_apply,subseteq,circ_apply,id_apply}.
    Then $f_1$ is a path in $X$ from $x$ to $y$ with respect to $T$ by \cref{path_in_space}.
    Hence $x\mathrel{\pathRelation{T}{X}}y$ by \cref{subspace_of,subseteq,times_tuple_intro,path_relation,fst_eq,snd_eq}.
    Take $c\in X$ such that $C = \equivalenceClass{\pathRelation{T}{X}}{c}$ by \cref{path_component}.
    Take $d\in X$ such that $D = \equivalenceClass{\pathRelation{T}{X}}{d}$ by \cref{path_component}.
    Therefore $x\in \downward{\pathRelation{T}{X}}{c}$ and $y\in \downward{\pathRelation{T}{X}}{d}$.
    Therefore $C = D$
        by \cref{downward_closure_iff,path_relation_symmetric,path_relation_transitive,symmetric,transitive,path_relation_equivalence,equiv_iff_equivclasses_eq}.
\end{proof}

% from §25 Item 13: locally connected at a point
\begin{definition}\label{locally_connected_at}
    $X$ is locally connected at $x$ with respect to $T$ iff
        for all $U$ such that $U$ is a neighborhood of $x$ in $X$ with respect to $T$
        there exists $V$ such that
            $V$ is a neighborhood of $x$ in $X$ with respect to $T$ and
            $V\subseteq U$ and
            $V$ is connected with respect to $\subspaceTopology{T}{V}$.
\end{definition}

% from §25 Item 14: locally connected
\begin{definition}\label{locally_connected}
    $X$ is locally connected with respect to $T$ iff
        for all $x\in X$ we have $X$ is locally connected at $x$ with respect to $T$.
\end{definition}

% from §25 Item 15: locally path connected at a point
\begin{definition}\label{locally_path_connected_at}
    $X$ is locally path connected at $x$ with respect to $T$ iff
        for all $U$ such that $U$ is a neighborhood of $x$ in $X$ with respect to $T$
        there exists $V$ such that
            $V$ is a neighborhood of $x$ in $X$ with respect to $T$ and
            $V\subseteq U$ and
            $V$ is path connected with respect to $\subspaceTopology{T}{V}$.
\end{definition}

% from §25 Item 16: locally path connected
\begin{definition}\label{locally_path_connected}
    $X$ is locally path connected with respect to $T$ iff
        for all $x\in X$ we have $X$ is locally path connected at $x$ with respect to $T$.
\end{definition}

% from §25 Item 17: locally connected iff components of open sets are open
\begin{theorem}\label{locally_connected_iff_components_open}
    Assume $T$ is a topology on $X$.
    Then $X$ is locally connected with respect to $T$ iff
        for all $U$ such that $U$ is open in $X$ with respect to $T$
        for all $C$ such that $C$ is a component of $U$ with respect to $\subspaceTopology{T}{U}$
        we have $C$ is open in $X$ with respect to $T$.
\end{theorem}
\begin{proof}
    Show that if $X$ is locally connected with respect to $T$ then
        for all $U$ such that $U$ is open in $X$ with respect to $T$
        for all $C$ such that $C$ is a component of $U$ with respect to $\subspaceTopology{T}{U}$
        we have $C$ is open in $X$ with respect to $T$.
    \begin{subproof}
        Assume $X$ is locally connected with respect to $T$.
        Fix $U$ such that $U$ is open in $X$ with respect to $T$.
        Thus $U$ is a subspace of $X$ with respect to $T$
            by \cref{topology_on,open_in,powerset_elim,subseteq,subspace_of}.
        Let $T_U = \subspaceTopology{T}{U}$.
        Thus $T_U$ is a topology on $U$ by \cref{subspace_topology_is_topology,subspace_of}.
        Fix $C$ such that $C$ is a component of $U$ with respect to $\subspaceTopology{T}{U}$.
        Thus $C$ is a component of $U$ with respect to $T_U$.
        Show that $C$ is open in $X$ with respect to $T$.
        \begin{subproof}
            Show that for all $x$ such that $x\in C$ we have
                there exists $V$ such that $V$ is open in $X$ with respect to $T$ and $x\in V$ and $V\subseteq C$.
            \begin{subproof}
                Fix $x$ such that $x\in C$.
                Take $x_0\in U$ such that $C = \equivalenceClass{\componentRelation{T_U}{U}}{x_0}$ by \cref{component}.
                Thus $C\subseteq U$ by \cref{downward_closure_iff,component_relation,times_tuple_elim,subseteq}.
                Thus $x\in U$ and $x\in X$ by \cref{subseteq,subspace_of}.
                Thus $U$ is a neighborhood of $x$ in $X$ with respect to $T$ by \cref{neighborhood}.
                Take $V$ such that
                    $V$ is a neighborhood of $x$ in $X$ with respect to $T$ and
                    $V\subseteq U$ and
                    $V$ is connected with respect to $\subspaceTopology{T}{V}$
                    by \cref{locally_connected,locally_connected_at}.
                Thus $V$ is open in $X$ with respect to $T$ and $x\in V$ by \cref{neighborhood}.
                Thus $V\subseteq C$
                    by \cref{subseteq,subspace_of,subspace_subspace_topology,components_cover,inter,emptyset,intersects,connected_subspace_single_component}.
                Thus there exists $V$ such that $V$ is open in $X$ with respect to $T$ and $x\in V$ and $V\subseteq C$.
            \end{subproof}
            Let $M = \{V\in T \mid V\subseteq C\}$.
            Thus for all $x$ such that $x\in C$ we have $x\in \unions{M}$ by \cref{subseteq,open_in,unions_iff}.
            Thus $C\subseteq \unions{M}$ by \cref{subseteq}.
            Thus $\unions{M}\subseteq C$ by \cref{unions_iff,subseteq}.
                Thus $C$ is open in $X$ with respect to $T$ by \cref{subseteq_antisymmetric,topology_on,subseteq,open_in}.
        \end{subproof}
    \end{subproof}
    Show that if for all $U$ such that $U$ is open in $X$ with respect to $T$
        for all $C$ such that $C$ is a component of $U$ with respect to $\subspaceTopology{T}{U}$
        we have $C$ is open in $X$ with respect to $T$ then
        $X$ is locally connected with respect to $T$.
    \begin{subproof}
        Assume for all $U$ such that $U$ is open in $X$ with respect to $T$
            for all $C$ such that $C$ is a component of $U$ with respect to $\subspaceTopology{T}{U}$
            we have $C$ is open in $X$ with respect to $T$.
        Show that for all $x\in X$ we have $X$ is locally connected at $x$ with respect to $T$.
        \begin{subproof}
            Fix $x\in X$.
            Show that for all $U$ such that $U$ is a neighborhood of $x$ in $X$ with respect to $T$
                there exists $V$ such that
                    $V$ is a neighborhood of $x$ in $X$ with respect to $T$ and
                    $V\subseteq U$ and
                    $V$ is connected with respect to $\subspaceTopology{T}{V}$.
            \begin{subproof}
                Fix $U$ such that $U$ is a neighborhood of $x$ in $X$ with respect to $T$.
                Thus $U$ is open in $X$ with respect to $T$ and $x\in U$ by \cref{neighborhood}.
                Thus $U$ is a subspace of $X$ with respect to $T$
                    by \cref{topology_on,open_in,powerset_elim,subseteq,subspace_of}.
                Let $T_U = \subspaceTopology{T}{U}$.
                Thus $T_U$ is a topology on $U$ by \cref{subspace_topology_is_topology,subspace_of}.
                Take $C$ such that
                    $C$ is a component of $U$ with respect to $T_U$ and $x\in C$
                    by \cref{components_cover}.
                Thus $C$ is open in $X$ with respect to $T$.
                Take $x_0\in U$ such that $C = \equivalenceClass{\componentRelation{T_U}{U}}{x_0}$ by \cref{component}.
                Thus $C\subseteq U$ by \cref{downward_closure_iff,component_relation,times_tuple_elim,subseteq}.
                Let $T_C = \subspaceTopology{T}{C}$.
                Thus $C$ is connected with respect to $T_C$
                    by \cref{component_connected,subspace_of,subspace_subspace_topology,subseteq}.
                Thus $C$ is a neighborhood of $x$ in $X$ with respect to $T$ by \cref{neighborhood,open_in}.
                Thus there exists $V$ such that
                    $V$ is a neighborhood of $x$ in $X$ with respect to $T$ and
                    $V\subseteq U$ and
                    $V$ is connected with respect to $\subspaceTopology{T}{V}$.
            \end{subproof}
            Thus $X$ is locally connected at $x$ with respect to $T$ by \cref{locally_connected_at}.
        \end{subproof}
        Thus $X$ is locally connected with respect to $T$ by \cref{locally_connected}.
    \end{subproof}
\end{proof}

% from §25 Item 18: locally path connected iff path components of open sets are open
\begin{theorem}\label{locally_path_connected_iff_components_open}
    Assume $T$ is a topology on $X$.
    Then $X$ is locally path connected with respect to $T$ iff
        for all $U$ such that $U$ is open in $X$ with respect to $T$
        for all $C$ such that $C$ is a path component of $U$ with respect to $\subspaceTopology{T}{U}$
        we have $C$ is open in $X$ with respect to $T$.
\end{theorem}
\begin{proof}
    Show that if $X$ is locally path connected with respect to $T$ then
        for all $U$ such that $U$ is open in $X$ with respect to $T$
        for all $C$ such that $C$ is a path component of $U$ with respect to $\subspaceTopology{T}{U}$
        we have $C$ is open in $X$ with respect to $T$.
    \begin{subproof}
        Assume $X$ is locally path connected with respect to $T$.
        Fix $U$ such that $U$ is open in $X$ with respect to $T$.
        Thus $U$ is a subspace of $X$ with respect to $T$
            by \cref{topology_on,open_in,powerset_elim,subseteq,subspace_of}.
        Let $T_U = \subspaceTopology{T}{U}$.
        Thus $T_U$ is a topology on $U$ by \cref{subspace_topology_is_topology,subspace_of}.
        Fix $C$ such that $C$ is a path component of $U$ with respect to $\subspaceTopology{T}{U}$.
        Thus $C$ is a path component of $U$ with respect to $T_U$.
        Show that for all $x$ such that $x\in C$ we have
            there exists $V$ such that $V$ is open in $X$ with respect to $T$ and $x\in V$ and $V\subseteq C$.
        \begin{subproof}
                Fix $x$ such that $x\in C$.
                Take $x_0\in U$ such that $C = \equivalenceClass{\pathRelation{T_U}{U}}{x_0}$ by \cref{path_component}.
                Thus $C\subseteq U$ by \cref{downward_closure_iff,path_relation,times_tuple_elim,subseteq}.
                Thus $x\in U$ and $x\in X$ by \cref{subseteq,subspace_of}.
                Thus $U$ is a neighborhood of $x$ in $X$ with respect to $T$ by \cref{neighborhood}.
                Take $V$ such that
                    $V$ is a neighborhood of $x$ in $X$ with respect to $T$ and
                    $V\subseteq U$ and
                    $V$ is path connected with respect to $\subspaceTopology{T}{V}$
                    by \cref{locally_path_connected,locally_path_connected_at}.
                Thus $V$ is open in $X$ with respect to $T$ and $x\in V$ by \cref{neighborhood}.
                Show that for all $y$ such that $y\in V$ we have $y\in C$.
                \begin{subproof}
                    Fix $y$ such that $y\in V$.
                    Thus $V$ is path connected with respect to $\subspaceTopology{T_U}{V}$
                        by \cref{subseteq,subspace_of,subspace_subspace_topology}.
                    Thus there exists $D$ such that
                        $D$ is a path component of $U$ with respect to $T_U$ and $y\in D$
                        by \cref{subseteq,path_components_cover}.
                    Thus $V$ intersects $D$ by \cref{inter,emptyset,intersects}.
                    Thus $V$ intersects $C$ by \cref{inter,emptyset,intersects}.
                    Therefore $y\in C$
                        by \cref{inter,emptyset,intersects,path_connected_subspace_single_component,subspace_of}.
                \end{subproof}
                Thus there exists $V$ such that $V$ is open in $X$ with respect to $T$ and $x\in V$ and $V\subseteq C$
                    by \cref{subseteq}.
        \end{subproof}
        Let $M = \{V\in T \mid V\subseteq C\}$.
        Thus for all $x$ such that $x\in C$ we have $x\in \unions{M}$ by \cref{subseteq,open_in,unions_iff}.
        Thus $C\subseteq \unions{M}$ by \cref{subseteq}.
        Thus $\unions{M}\subseteq C$ by \cref{unions_iff,subseteq}.
        Thus $C$ is open in $X$ with respect to $T$ by \cref{subseteq_antisymmetric,topology_on,subseteq,open_in}.
    \end{subproof}
    Show that if for all $U$ such that $U$ is open in $X$ with respect to $T$
        for all $C$ such that $C$ is a path component of $U$ with respect to $\subspaceTopology{T}{U}$
        we have $C$ is open in $X$ with respect to $T$ then
        $X$ is locally path connected with respect to $T$.
    \begin{subproof}
        Assume for all $U$ such that $U$ is open in $X$ with respect to $T$
            for all $C$ such that $C$ is a path component of $U$ with respect to $\subspaceTopology{T}{U}$
            we have $C$ is open in $X$ with respect to $T$.
        Show that for all $x\in X$ we have $X$ is locally path connected at $x$ with respect to $T$.
        \begin{subproof}
            Fix $x\in X$.
            Show that for all $U$ such that $U$ is a neighborhood of $x$ in $X$ with respect to $T$
                there exists $V$ such that
                    $V$ is a neighborhood of $x$ in $X$ with respect to $T$ and
                    $V\subseteq U$ and
                    $V$ is path connected with respect to $\subspaceTopology{T}{V}$.
            \begin{subproof}
                Fix $U$ such that $U$ is a neighborhood of $x$ in $X$ with respect to $T$.
                Thus $U$ is open in $X$ with respect to $T$ and $x\in U$ by \cref{neighborhood}.
                Thus $U$ is a subspace of $X$ with respect to $T$
                    by \cref{topology_on,open_in,powerset_elim,subseteq,subspace_of}.
                Let $T_U = \subspaceTopology{T}{U}$.
                Thus $T_U$ is a topology on $U$ by \cref{subspace_topology_is_topology,subspace_of}.
                Take $C$ such that
                    $C$ is a path component of $U$ with respect to $T_U$ and $x\in C$
                    by \cref{path_components_cover}.
                Thus $C$ is open in $X$ with respect to $T$.
                Take $x_0\in U$ such that $C = \equivalenceClass{\pathRelation{T_U}{U}}{x_0}$ by \cref{path_component}.
                Thus $C\subseteq U$ by \cref{downward_closure_iff,path_relation,times_tuple_elim,subseteq}.
                Let $T_C = \subspaceTopology{T}{C}$.
                Thus $C$ is path connected with respect to $T_C$
                    by \cref{path_component_path_connected,subspace_of,subspace_subspace_topology,subseteq}.
                Thus $C$ is a neighborhood of $x$ in $X$ with respect to $T$ by \cref{neighborhood,open_in}.
                Thus there exists $V$ such that
                    $V$ is a neighborhood of $x$ in $X$ with respect to $T$ and
                    $V\subseteq U$ and
                    $V$ is path connected with respect to $\subspaceTopology{T}{V}$.
            \end{subproof}
            Thus $X$ is locally path connected at $x$ with respect to $T$ by \cref{locally_path_connected_at}.
        \end{subproof}
        Thus $X$ is locally path connected with respect to $T$ by \cref{locally_path_connected}.
    \end{subproof}
\end{proof}

% from §25 Item 19: path components lie in components
\begin{theorem}\label{path_component_subset_component}
    Assume $T$ is a topology on $X$.
    Suppose $P$ is a path component of $X$ with respect to $T$.
    Then there exists $C$ such that $C$ is a component of $X$ with respect to $T$ and $P\subseteq C$.
\end{theorem}
\begin{proof}
    Take $x\in X$ such that $P = \equivalenceClass{\pathRelation{T}{X}}{x}$ by \cref{path_component}.
    For all $y$ such that $y\in P$ we have
        $(y,x)\in X\times X$ and there exists $f$ such that
        $f$ is a path in $X$ from $\fst{(y,x)}$ to $\snd{(y,x)}$ with respect to $T$
        by \cref{downward_closure_iff,path_relation}.
    Thus $P$ is a subspace of $X$ with respect to $T$ by \cref{times_tuple_elim,subseteq,subspace_of}.
    Let $T_P = \subspaceTopology{T}{P}$.
    Thus $T_P$ is a topology on $P$ by \cref{subspace_topology_is_topology,subspace_of}.
    Thus $P$ is connected with respect to $T_P$
        by \cref{path_component_path_connected,path_connected_implies_connected}.
    Take $C$ such that
        $C$ is a component of $X$ with respect to $T$ and $x\in C$
        by \cref{components_cover}.
    For all $y$ such that $y\in P$ we have $y\in C$
        by \cref{times_tuple_elim,components_cover,path_relation_equivalence,equivclasses_inhabited_equivalenceon,inter,emptyset,intersects,connected_subspace_single_component,subspace_of}.
    Thus there exists $C$ such that $C$ is a component of $X$ with respect to $T$ and $P\subseteq C$
        by \cref{subseteq}.
\end{proof}


% from §25 Item 20: components equal path components in locally path connected spaces
\begin{theorem}\label{components_equal_path_components_locally_path_connected}
    Assume $T$ is a topology on $X$.
    Suppose $X$ is locally path connected with respect to $T$.
    Then for all $C$ such that $C$ is a component of $X$ with respect to $T$
        there exists $P$ such that $P$ is a path component of $X$ with respect to $T$ and $C = P$.
\end{theorem}
\begin{proof}
    Fix $C$ such that $C$ is a component of $X$ with respect to $T$.
    Take $x\in X$ such that $C = \equivalenceClass{\componentRelation{T}{X}}{x}$ by \cref{component}.
    Hence $x\in C$ by \cref{component_relation_equivalence,equivclasses_inhabited_equivalenceon}.
    Therefore $C$ is a subspace of $X$ with respect to $T$
        by \cref{downward_closure_iff,component_relation,times_tuple_elim,subseteq,subspace_of}.
    Let $T_X = \subspaceTopology{T}{X}$.
    Hence $T_X$ is a topology on $X$ by \cref{subseteq_refl,subspace_of,subspace_topology_is_topology}.
    Therefore $T\subseteq T_X$
        by \cref{topology_on,powerset_elim,subseteq,inter_absorb_supseteq_left,subspace_topology}.
    Hence $T_X\subseteq T$
        by \cref{subspace_topology,topology_on,powerset_elim,subseteq,inter_absorb_supseteq_left}.
    Therefore $T_X = T$ by \cref{subseteq_antisymmetric}.
    Hence for all $D$ such that $D$ is a path component of $X$ with respect to $T$
        we have $D$ is open in $X$ with respect to $T$
        by \cref{topology_on,open_in,locally_path_connected_iff_components_open}.
    Take $P$ such that $P$ is a path component of $X$ with respect to $T$ and $x\in P$
        by \cref{path_components_cover}.
    Suppose not.
    Take $C_1$ such that
        $C_1$ is a component of $X$ with respect to $T$ and $P\subseteq C_1$
        by \cref{path_component_subset_component}.
    Then $C = C_1$ by \cref{components_disjoint,subseteq,disjoint}.
    Thus $P\subset C$ by \cref{subseteq,subset}.
    Take $y$ such that $y\in C\setminus P$
        by \cref{difference_with_proper_subset_is_inhabited,nonempty_inhabited}.
    Thus $y\in C$ and $y\notin P$ by \cref{setminus_elim_left,setminus_elim_right}.
    Let $M = \{D\in\pow{X} \mid
        \text{$D$ is a path component of $X$ with respect to $T$ and $D\subseteq C$ and $D\neq P$}\}$.
    Let $Q = \unions{M}$.
    Hence $Q\subseteq C\setminus P$
        by \cref{unions_iff,subseteq,path_components_disjoint,disjoint,setminus_intro}.
    Show that for all $z$ such that $z\in C\setminus P$ we have $z\in Q$.
    \begin{subproof}
        Fix $z$ such that $z\in C\setminus P$.
        Thus $z\in C$ and $z\notin P$ by \cref{setminus_elim_left,setminus_elim_right}.
        Take $D$ such that $D$ is a path component of $X$ with respect to $T$ and $z\in D$
            by \cref{subspace_of,subseteq,path_components_cover}.
        Take $C_1$ such that
            $C_1$ is a component of $X$ with respect to $T$ and $D\subseteq C_1$
            by \cref{path_component_subset_component}.
        Then $C = C_1$ by \cref{components_disjoint,subseteq,disjoint}.
        Thus $D\in M$ by \cref{subspace_of,subseteq_transitive,powerset_intro}.
        Thus $z\in Q$ by \cref{unions_iff}.
    \end{subproof}
    Therefore $C\setminus P\subseteq Q$ by \cref{subseteq}.
    Hence $Q = C\setminus P$ by \cref{subseteq_antisymmetric}.
    Let $T_C = \subspaceTopology{T}{C}$.
    Thus $T_C$ is a topology on $C$ by \cref{subspace_topology_is_topology,subspace_of}.
    Then $P$ is open in $C$ with respect to $T_C$
        by \cref{open_in,inter_absorb_supseteq_left,subspace_topology}.
    Therefore $Q$ is open in $C$ with respect to $T_C$
        by \cref{open_in,subseteq,topology_on,inter_absorb_supseteq_left,setminus_subseteq,subspace_topology}.
    Hence $x\in P$ and $y\in Q$.
    Therefore $P\neq\emptyset$ and $Q\neq\emptyset$ by \cref{emptyset}.
    Hence $P$ is disjoint from $Q$ by \cref{setminus_disjoint,inter,emptyset,disjoint}.
    Hence $P\union Q = C$ by \cref{setminus_partition,subseteq}.
    Therefore $C$ is separated by $P$ and $Q$ with respect to $T_C$ by \cref{separation}.
    Then $C$ is not connected with respect to $T_C$ by \cref{connected}.
    Contradiction by \cref{component_connected}.
    Thus $C = P$.
    Thus there exists $P$ such that $P$ is a path component of $X$ with respect to $T$ and $C = P$.
\end{proof}


\section{§26 Compact Spaces}

% from §26 Item 1: covering
\begin{definition}\label{covering}
    $A$ is a covering of $X$ iff $X\subseteq\unions{A}$.
\end{definition}

% from §26 Item 2: open covering
\begin{definition}\label{open_covering}
    $A$ is an open covering of $X$ with respect to $T$ iff
        $A$ is a covering of $X$ and
        for all $U$ such that $U\in A$ we have $U$ is open in $X$ with respect to $T$.
\end{definition}

% from §26 Item 3: compact space
\begin{definition}\label{compact}
    $X$ is compact with respect to $T$ iff
        for all $A$ such that $A$ is an open covering of $X$ with respect to $T$
        there exists $B$ such that
            $B\subseteq A$ and
            $B$ is finite and
            $B$ is a covering of $X$.
\end{definition}

% from §26 Item 4: compactness via coverings open in the ambient space
\begin{lemma}\label{compact_subspace_open_covering}
    Assume $Y$ is a subspace of $X$ with respect to $T$.
    Let $T_Y = \subspaceTopology{T}{Y}$.
    Then $Y$ is compact with respect to $T_Y$ iff
        for all $A$ such that
            $A$ is a covering of $Y$ and
            for all $U$ such that $U\in A$ we have $U$ is open in $X$ with respect to $T$
        there exists $B$ such that
            $B\subseteq A$ and
            $B$ is finite and
            $B$ is a covering of $Y$.
\end{lemma}
\begin{proof}
    Let $T_Y = \subspaceTopology{T}{Y}$.
    $T$ is a topology on $X$ and $Y\subseteq X$ by \cref{subspace_of}.
    Hence $T_Y$ is a topology on $Y$ by \cref{subspace_topology_is_topology}.
    Show that if $Y$ is compact with respect to $T_Y$ then
        for all $A$ such that
            $A$ is a covering of $Y$ and
            for all $U$ such that $U\in A$ we have $U$ is open in $X$ with respect to $T$
        there exists $B$ such that
            $B\subseteq A$ and
            $B$ is finite and
            $B$ is a covering of $Y$.
    \begin{subproof}
        Assume $Y$ is compact with respect to $T_Y$.
        Fix $A$ such that
            $A$ is a covering of $Y$ and
            for all $U$ such that $U\in A$ we have $U$ is open in $X$ with respect to $T$.
        Let $A_Y = \{V \mid \exists U\in A. V = Y\inter U\}$.
        Thus for all $V$ such that $V\in A_Y$ we have $V$ is open in $Y$ with respect to $T_Y$
            by \cref{open_in,subspace_topology}.
        $Y\subseteq\unions{A}$ by \cref{covering}.
        Thus for all $y$ such that $y\in Y$ we have $y\in\unions{A_Y}$.
        Therefore $A_Y$ is a covering of $Y$ by \cref{subseteq,covering}.
        Hence $A_Y$ is an open covering of $Y$ with respect to $T_Y$ by \cref{open_covering}.
        Take $B_Y$ such that
            $B_Y\subseteq A_Y$ and $B_Y$ is finite and $B_Y$ is a covering of $Y$
            by \cref{compact}.
        Let $R = \{w\in B_Y\times A \mid \exists V\in B_Y.\exists U\in A. w = (V,U) \land V = Y\inter U\}$.
        For all $w$ such that $w\in R$ there exist $V,U$ such that $w = (V,U)$.
        Then $R$ is a relation by \cref{relation}.
        For all $V\in B_Y$ there exists $U$ such that $V\mathrel{R} U$.
        Take $g$ such that $g$ is a function on $B_Y$ and for all $V\in B_Y$ we have $V\mathrel{R}\apply{g}{V}$ by \cref{choice_function}.
        Let $B = \img{g}{B_Y}$.
        Therefore $B\subseteq A$ by \cref{img_iff,function_apply_iff,pair_eq_iff,subseteq}.
        $B$ is finite by \cref{finite_image_function}.
        Thus for all $y$ such that $y\in Y$ we have $y\in\unions{B}$
            by \cref{covering,subseteq,function_apply_iff,pair_eq_iff,inter_intro,inter_elim_right,function_apply_elim,img_iff,unions_iff}.
        Hence $B$ is a covering of $Y$ by \cref{subseteq,covering}.
    \end{subproof}
    Show that if for all $A$ such that
            $A$ is a covering of $Y$ and
            for all $U$ such that $U\in A$ we have $U$ is open in $X$ with respect to $T$
        there exists $B$ such that
            $B\subseteq A$ and
            $B$ is finite and
            $B$ is a covering of $Y$
        then $Y$ is compact with respect to $T_Y$.
    \begin{subproof}
        Assume for all $A$ such that
                $A$ is a covering of $Y$ and
                for all $U$ such that $U\in A$ we have $U$ is open in $X$ with respect to $T$
            there exists $B$ such that
                $B\subseteq A$ and
                $B$ is finite and
                $B$ is a covering of $Y$.
        Show that for all $A$ such that $A$ is an open covering of $Y$ with respect to $T_Y$
            there exists $B$ such that $B\subseteq A$ and $B$ is finite and $B$ is a covering of $Y$.
        \begin{subproof}
            Fix $A$ such that $A$ is an open covering of $Y$ with respect to $T_Y$.
            Let $R = \{w\in A\times T \mid \exists V\in A.\exists U\in T. w = (V,U) \land V = Y\inter U\}$.
            For all $w$ such that $w\in R$ there exist $V,U$ such that $w = (V,U)$.
            Then $R$ is a relation by \cref{relation}.
            For all $V\in A$ there exists $U\in T$ such that $V = Y\inter U$
                by \cref{open_covering,open_in,subspace_topology}.
            Therefore for all $V\in A$ there exists $U$ such that $V\mathrel{R} U$.
            Take $g$ such that $g$ is a function on $A$ and for all $V\in A$ we have $V\mathrel{R}\apply{g}{V}$ by \cref{choice_function}.
            Let $C = \img{g}{A}$.
            For all $U$ such that $U\in C$ we have $U$ is open in $X$ with respect to $T$
                by \cref{img_iff,function_apply_iff,pair_eq_iff,open_in}.
            Thus for all $y$ such that $y\in Y$ we have $y\in\unions{C}$
                by \cref{open_covering,covering,subseteq,unions_iff,function_apply_iff,pair_eq_iff,inter_intro,inter_elim_right,function_apply_elim,img_iff}.
            Hence $C$ is a covering of $Y$ by \cref{subseteq,covering}.
            Take $B$ such that $B\subseteq C$ and $B$ is finite and $B$ is a covering of $Y$.
            Let $S = \{w\in B\times A \mid \exists U\in B.\exists V\in A. w = (U,V) \land \apply{g}{V} = U\}$.
            For all $w$ such that $w\in S$ there exist $U,V$ such that $w = (U,V)$.
            Then $S$ is a relation by \cref{relation}.
            Therefore for all $U\in B$ there exists $V\in A$ such that $\apply{g}{V} = U$
                by \cref{subseteq,img_iff,function_apply_iff}.
            Hence for all $U\in B$ there exists $V$ such that $U\mathrel{S} V$.
            Take $h$ such that $h$ is a function on $B$ and for all $U\in B$ we have $U\mathrel{S}\apply{h}{U}$ by \cref{choice_function}.
            Let $B_1 = \img{h}{B}$.
            Therefore $B_1\subseteq A$ by \cref{img_iff,function_apply_iff,pair_eq_iff,subseteq}.
            $B_1$ is finite by \cref{finite_image_function}.
            Thus for all $y$ such that $y\in Y$ we have $y\in\unions{B_1}$
                by \cref{covering,subseteq,unions_iff,function_apply_iff,pair_eq_iff,inter_intro,open_covering,open_in,subspace_topology,function_apply_elim,img_iff}.
            Therefore $B_1$ is a covering of $Y$ by \cref{subseteq,covering}.
            Thus there exists $B$ such that $B\subseteq A$ and $B$ is finite and $B$ is a covering of $Y$.
        \end{subproof}
        Hence $Y$ is compact with respect to $T_Y$ by \cref{compact}.
    \end{subproof}
\end{proof}
% from §26 Item 5: closed subspace of a compact space is compact
\begin{theorem}\label{closed_subspace_compact}
    Assume $T$ is a topology on $X$.
    Assume $X$ is compact with respect to $T$.
    Suppose $Y\subseteq X$.
    Suppose $Y$ is closed in $X$ with respect to $T$.
    Let $T_Y = \subspaceTopology{T}{Y}$.
    Then $Y$ is compact with respect to $T_Y$.
\end{theorem}
\begin{proof}
    $Y$ is a subspace of $X$ with respect to $T$ by \cref{subspace_of}.
    Show that for all $A$ such that
        $A$ is a covering of $Y$ and
        for all $U$ such that $U\in A$ we have $U$ is open in $X$ with respect to $T$
        there exists $B$ such that
            $B\subseteq A$ and
            $B$ is finite and
            $B$ is a covering of $Y$.
    \begin{subproof}
        Fix $A$ such that
            $A$ is a covering of $Y$ and
            for all $U$ such that $U\in A$ we have $U$ is open in $X$ with respect to $T$.
        Let $C = A\union \{X\setminus Y\}$.
        For all $U$ such that $U\in C$ we have $U$ is open in $X$ with respect to $T$
            by \cref{union_iff,singleton_iff,closed_in}.
        Show that for all $x$ such that $x\in X$ we have $x\in\unions{C}$.
        \begin{subproof}
            Fix $x$ such that $x\in X$.
            \begin{byCase}
            \caseOf{$x\in Y$.}
                Take $U\in A$ such that $x\in U$ by \cref{covering,subseteq,unions_iff}.
                Hence $x\in\unions{C}$ by \cref{union_iff,unions_iff}.
            \caseOf{$x\notin Y$.}
                Therefore $x\in\unions{C}$ by \cref{setminus_intro,singleton_iff,union_iff,unions_iff}.
            \end{byCase}
        \end{subproof}
        Hence $C$ is a covering of $X$ by \cref{subseteq,covering}.
        Therefore $C$ is an open covering of $X$ with respect to $T$ by \cref{open_covering}.
        Take $B$ such that $B\subseteq C$ and $B$ is finite and $B$ is a covering of $X$ by \cref{compact}.
        \begin{byCase}
        \caseOf{$X\setminus Y\in B$.}
            Let $B_1 = B\setminus \{X\setminus Y\}$.
            For all $y$ such that $y\in Y$ there exists $U\in B$ such that $y\in U$
                by \cref{covering,subseteq,unions_iff}.
            Thus for all $y$ such that $y\in Y$ there exists $U\in B_1$ such that $y\in U$
                by \cref{setminus_intro,singleton_iff,setminus_elim_right}.
            Hence for all $y$ such that $y\in Y$ we have $y\in\unions{B_1}$ by \cref{unions_iff}.
            Hence $B_1$ is a covering of $Y$ by \cref{subseteq,covering}.
            $B_1\subseteq B$ by \cref{setminus_subseteq}.
            Therefore $B_1$ is finite by \cref{setminus_subseteq,subseteq_finite_is_finite}.
            Hence $B_1\subseteq A$ by \cref{subseteq,setminus_elim_right,union_iff,subseteq}.
            Therefore there exists $B$ such that $B\subseteq A$ and $B$ is finite and $B$ is a covering of $Y$.
        \caseOf{$X\setminus Y\notin B$.}
            Therefore there exists $B$ such that $B\subseteq A$ and $B$ is finite and $B$ is a covering of $Y$
                by \cref{subseteq,union_iff,singleton_iff,covering}.
        \end{byCase}
    \end{subproof}
    Hence $Y$ is compact with respect to $T_Y$ by \cref{compact_subspace_open_covering}.
\end{proof}
% from §26 Item 6: compact subspace of Hausdorff is closed
\begin{theorem}\label{compact_subspace_closed}
    Assume $T$ is a topology on $X$.
    Assume $X$ is Hausdorff with respect to $T$.
    Suppose $Y\subseteq X$.
    Let $T_Y = \subspaceTopology{T}{Y}$.
    Suppose $Y$ is compact with respect to $T_Y$.
    Then $Y$ is closed in $X$ with respect to $T$.
\end{theorem}
\begin{proof}
    \begin{byCase}
    \caseOf{$Y = \emptyset$.}
        Thus $Y$ is closed in $X$ with respect to $T$ by \cref{closed_emptyset}.
    \caseOf{$Y \neq \emptyset$.}
        Let $M = \{U\in T \mid \text{$U$ is disjoint from $Y$}\}$.
        Show that for all $x_0$ such that $x_0\in X\setminus Y$ we have $x_0\in\unions{M}$.
        \begin{subproof}
                        Fix $x_0$ such that $x_0\in X\setminus Y$.
                        Then $x_0\in X$ and $x_0\notin Y$ by \cref{setminus_elim_left,setminus_elim_right}.
                        Let $A = \{V\in T \mid \text{there exists $y\in Y$ such that there exists $U$ such that
                                $U$ is a neighborhood of $x_0$ in $X$ with respect to $T$ and
                                $V$ is a neighborhood of $y$ in $X$ with respect to $T$ and
                                $U$ is disjoint from $V$}\}$.
                        For all $y$ such that $y\in Y$ there exists $V\in A$ such that $y\in V$
                            by \cref{subseteq,hausdorff,neighborhood,open_in}.
                        Thus for all $y$ such that $y\in Y$ we have $y\in\unions{A}$ by \cref{unions_iff}.
                        Therefore $Y\subseteq\unions{A}$ and $A$ is a covering of $Y$ by \cref{subseteq,covering}.
                        For all $U$ such that $U\in A$ we have $U$ is open in $X$ with respect to $T$
                            by \cref{subseteq,open_in}.
                        Take $B$ such that
                            $B\subseteq A$ and $B$ is finite and $B$ is a covering of $Y$
                            by \cref{subspace_of,compact_subspace_open_covering,subseteq}.
                        Let $R = \{w\in B\times T \mid \text{there exists $V\in B$ such that there exists $U\in T$ such that
                                $w = (V,U)$ and
                                $U$ is a neighborhood of $x_0$ in $X$ with respect to $T$ and
                                $U$ is disjoint from $V$}\}$.
                        For all $w$ such that $w\in R$ there exist $V,U$ such that $w = (V,U)$.
                        Then $R$ is a relation by \cref{relation}.
                        Show that for all $V\in B$ there exists $U$ such that $V\mathrel{R} U$.
                        \begin{subproof}
                            Fix $V$ such that $V\in B$.
                            Take $y,U$ such that $y\in Y$ and
                                $U$ is a neighborhood of $x_0$ in $X$ with respect to $T$ and
                                $V$ is a neighborhood of $y$ in $X$ with respect to $T$ and
                                $U$ is disjoint from $V$
                                by \cref{subseteq}.
                            Hence $U\in T$ by \cref{neighborhood,open_in}.
                            Hence $V\mathrel{R} U$.
                        \end{subproof}
                        Take $g$ such that $g$ is a function on $B$ and for all $V\in B$ we have $V\mathrel{R}\apply{g}{V}$ by \cref{choice_function}.
                        Let $G = \img{g}{B}$.
                        Therefore $G\subseteq T$ by \cref{img_iff,function_apply_iff,pair_eq_iff,subseteq}.
                        $G$ is finite by \cref{finite_image_function}.
                        Take $y$ such that $y\in Y$ by \cref{nonempty_inhabited}.
                        Take $V\in B$ such that $y\in V$ by \cref{covering,subseteq,unions_iff}.
                        Hence $G$ is inhabited by \cref{function_apply_elim,subseteq,img_iff}.
                        Let $U_0 = \inters{G}$.
                        Therefore $U_0\in T$ by \cref{hausdorff,finite_intersections_open_in_topology}.
                        For all $U$ such that $U\in G$ there exists $V\in B$ such that $(V,U)\in g$
                            by \cref{img_iff}.
                        Thus for all $U$ such that $U\in G$ we have $x_0\in U$
                            by \cref{function_apply_iff,pair_eq_iff,neighborhood}.
                        Therefore $x_0\in U_0$ by \cref{inters_intro}.
                        Suppose not.
                        Then there exists $z$ such that $z\in U_0$ and $z\in Y$.
                        Take $V\in B$ such that $z\in V$ by \cref{covering,subseteq,unions_iff}.
                        Hence $z\in \apply{g}{V}$ by \cref{function_apply_elim,subseteq,img_iff,inters_subseteq_elem}.
                        Take $V_1,U_1$ such that $V_1\in B$ and $U_1\in T$ and
                            $(V,\apply{g}{V}) = (V_1,U_1)$ and
                            $U_1$ is a neighborhood of $x_0$ in $X$ with respect to $T$ and
                            $U_1$ is disjoint from $V_1$
                            by \cref{subseteq,function_apply_iff}.
                        Hence $z\notin \apply{g}{V}$ by \cref{pair_eq_iff,disjoint}.
                        Contradiction.
                        Therefore $U_0$ is disjoint from $Y$.
                        Therefore $U_0\in M$.
                        Hence $x_0\in\unions{M}$ by \cref{unions_iff}.
        \end{subproof}
        Therefore $X\setminus Y\subseteq \unions{M}$ by \cref{subseteq}.
        For all $x$ such that $x\in\unions{M}$ we have $x\in X\setminus Y$
            by \cref{unions_iff,disjoint,subseteq,topology_on,powerset_elim,setminus_intro}.
        Hence $X\setminus Y = \unions{M}$ by \cref{subseteq,subseteq_antisymmetric}.
        Therefore $X\setminus Y$ is open in $X$ with respect to $T$ by \cref{subseteq,topology_on,open_in}.
        Hence $Y$ is closed in $X$ with respect to $T$ by \cref{closed_in}.
    \end{byCase}
\end{proof}

% from §26 Item 7: separation of a point and a compact set in a Hausdorff space
\begin{lemma}\label{compact_hausdorff_separation}
    Assume $T$ is a topology on $X$.
    Assume $X$ is Hausdorff with respect to $T$.
    Suppose $Y\subseteq X$.
    Let $T_Y = \subspaceTopology{T}{Y}$.
    Suppose $Y$ is compact with respect to $T_Y$.
    Suppose $x_0\in X\setminus Y$.
    Then there exist $U,V$ such that
        $U$ is open in $X$ with respect to $T$ and
        $V$ is open in $X$ with respect to $T$ and
        $x_0\in U$ and
        $Y\subseteq V$ and
        $U$ is disjoint from $V$.
\end{lemma}
\begin{proof}
    \begin{byCase}
    \caseOf{$Y = \emptyset$.}
        Let $U = X$.
        Let $V = \emptyset$.
        Hence there exist $U,V$ such that
            $U$ is open in $X$ with respect to $T$ and
            $V$ is open in $X$ with respect to $T$ and
            $x_0\in U$ and
            $Y\subseteq V$ and
            $U$ is disjoint from $V$
            by \cref{open_in,topology_on,setminus_elim_left,emptyset_subseteq,emptyset,disjoint}.
    \caseOf{$Y \neq \emptyset$.}
        Then $x_0\in X$ and $x_0\notin Y$ by \cref{setminus_elim_left,setminus_elim_right}.
        Let $A = \{V\in T \mid \text{there exists $y\in Y$ such that there exists $U$ such that
                $U$ is a neighborhood of $x_0$ in $X$ with respect to $T$ and
                $V$ is a neighborhood of $y$ in $X$ with respect to $T$ and
                $U$ is disjoint from $V$}\}$.
        For all $y$ such that $y\in Y$ there exists $V\in A$ such that $y\in V$
            by \cref{subseteq,hausdorff,neighborhood,open_in}.
        Thus for all $y$ such that $y\in Y$ we have $y\in\unions{A}$ by \cref{unions_iff}.
        Therefore $Y\subseteq\unions{A}$ and $A$ is a covering of $Y$ by \cref{subseteq,covering}.
        For all $U$ such that $U\in A$ we have $U$ is open in $X$ with respect to $T$
            by \cref{subseteq,open_in}.
        Take $B$ such that
            $B\subseteq A$ and $B$ is finite and $B$ is a covering of $Y$
            by \cref{subspace_of,compact_subspace_open_covering,subseteq}.
        Let $R = \{w\in B\times T \mid \text{there exists $V\in B$ such that there exists $U\in T$ such that
                $w = (V,U)$ and
                $U$ is a neighborhood of $x_0$ in $X$ with respect to $T$ and
                $U$ is disjoint from $V$}\}$.
        For all $w$ such that $w\in R$ there exist $V,U$ such that $w = (V,U)$.
        Then $R$ is a relation by \cref{relation}.
        Show that for all $V\in B$ there exists $U$ such that $V\mathrel{R} U$.
        \begin{subproof}
            Fix $V$ such that $V\in B$.
            Take $y,U$ such that $y\in Y$ and
                $U$ is a neighborhood of $x_0$ in $X$ with respect to $T$ and
                $V$ is a neighborhood of $y$ in $X$ with respect to $T$ and
                $U$ is disjoint from $V$
                by \cref{subseteq}.
            Therefore $U\in T$ by \cref{neighborhood,open_in}.
            Hence $V\mathrel{R} U$.
        \end{subproof}
        Take $g$ such that $g$ is a function on $B$ and for all $V\in B$ we have $V\mathrel{R}\apply{g}{V}$ by \cref{choice_function}.
        Let $G = \img{g}{B}$.
        Therefore $G\subseteq T$ by \cref{img_iff,function_apply_iff,pair_eq_iff,subseteq}.
        $G$ is finite by \cref{finite_image_function}.
        Take $y$ such that $y\in Y$ by \cref{nonempty_inhabited}.
        Take $V\in B$ such that $y\in V$ by \cref{covering,subseteq,unions_iff}.
        Hence $G$ is inhabited by \cref{function_apply_elim,subseteq,img_iff}.
        Let $U_0 = \inters{G}$.
        Therefore $U_0\in T$ by \cref{hausdorff,finite_intersections_open_in_topology}.
        For all $U$ such that $U\in G$ there exists $V\in B$ such that $(V,U)\in g$
            by \cref{img_iff}.
        Thus for all $U$ such that $U\in G$ we have $x_0\in U$
            by \cref{function_apply_iff,pair_eq_iff,neighborhood}.
        Therefore $x_0\in U_0$ by \cref{inters_intro}.
        Let $V_0 = \unions{B}$.
        Hence $V_0\in T$ and $V_0$ is open in $X$ with respect to $T$ by \cref{subseteq,topology_on,open_in}.
        $Y\subseteq V_0$ by \cref{covering}.
        Show that $U_0$ is disjoint from $V_0$.
        \begin{subproof}
            Suppose not.
            Then there exists $z$ such that $z\in U_0$ and $z\in V_0$.
            Take $V\in B$ such that $z\in V$ by \cref{unions_iff}.
            Therefore $z\in \apply{g}{V}$ by \cref{function_apply_elim,subseteq,img_iff,inters_subseteq_elem}.
            Take $V_1,U_1$ such that $V_1\in B$ and $U_1\in T$ and
                $(V,\apply{g}{V}) = (V_1,U_1)$ and
                $U_1$ is a neighborhood of $x_0$ in $X$ with respect to $T$ and
                $U_1$ is disjoint from $V_1$
                by \cref{subseteq,function_apply_iff}.
            Therefore $z\notin \apply{g}{V}$ by \cref{pair_eq_iff,disjoint}.
            Contradiction.
        \end{subproof}
        Therefore there exist $U,V$ such that
            $U$ is open in $X$ with respect to $T$ and
            $V$ is open in $X$ with respect to $T$ and
            $x_0\in U$ and
            $Y\subseteq V$ and
            $U$ is disjoint from $V$.
    \end{byCase}
\end{proof}

% from §26 Item 8: image of a compact space under a continuous map is compact
\begin{theorem}\label{continuous_image_compact}
    Assume $f$ is continuous from $X$ to $Y$ with respect to $T$ and $T_1$.
    Assume $X$ is compact with respect to $T$.
    Let $A = \img{f}{X}$.
    Let $T_A = \subspaceTopology{T_1}{A}$.
    Then $A$ is compact with respect to $T_A$.
\end{theorem}
\begin{proof}
    Show that for all $C$ such that
        $C$ is a covering of $A$ and
        for all $U$ such that $U\in C$ we have $U$ is open in $Y$ with respect to $T_1$
        there exists $B$ such that
            $B\subseteq C$ and
            $B$ is finite and
            $B$ is a covering of $A$.
    \begin{subproof}
        Fix $C$ such that
            $C$ is a covering of $A$ and
            for all $U$ such that $U\in C$ we have $U$ is open in $Y$ with respect to $T_1$.
        Let $D = \{W \mid \exists U\in C. W = \preimg{f}{U}\}$.
        For all $x$ such that $x\in X$ there exists $U\in C$ such that $f(x)\in U$
            by \cref{continuous,funs_is_function,funs,function_apply_elim,img_iff,covering,subseteq,unions_iff}.
        Thus for all $x$ such that $x\in X$ we have $x\in\unions{D}$
            by \cref{continuous,funs_is_function,funs,function_apply_elim,preimg_iff,unions_iff}.
        Therefore $X\subseteq\unions{D}$ and $D$ is a covering of $X$ by \cref{subseteq,covering}.
        For all $W$ such that $W\in D$ we have $W$ is open in $X$ with respect to $T$.
        Hence $D$ is an open covering of $X$ with respect to $T$ by \cref{open_covering}.
        Take $E$ such that
            $E\subseteq D$ and $E$ is finite and $E$ is a covering of $X$
            by \cref{compact}.
        Let $R = \{w\in E\times C \mid \exists W\in E.\exists U\in C. w = (W,U) \land W = \preimg{f}{U}\}$.
        For all $w$ such that $w\in R$ there exist $W,U$ such that $w = (W,U)$.
        Then $R$ is a relation by \cref{relation}.
        For all $W\in E$ there exists $U\in C$ such that $W = \preimg{f}{U}$ by \cref{subseteq}.
        Hence for all $W\in E$ there exists $U$ such that $W\mathrel{R} U$.
        Take $h$ such that $h$ is a function on $E$ and for all $W\in E$ we have $W\mathrel{R}\apply{h}{W}$ by \cref{choice_function}.
        Let $B = \img{h}{E}$.
        Therefore $B\subseteq C$ by \cref{img_iff,function_apply_intro,pair_eq_iff,subseteq}.
        $B$ is finite by \cref{finite_image_function}.
        Thus for all $a$ such that $a\in A$ we have $a\in\unions{B}$
            by \cref{img_iff,continuous,funs_is_function,function_apply_intro,covering,subseteq,unions_iff,function_apply_iff,pair_eq_iff,preimg_iff,function_apply_elim}.
        Hence $A\subseteq\unions{B}$ and $B$ is a covering of $A$ by \cref{subseteq,covering}.
    \end{subproof}
    Therefore $A\subseteq Y$ and $A$ is a subspace of $Y$ with respect to $T_1$
        by \cref{continuous,funs,rels_ran_subseteq,img_subseteq_ran,subseteq_transitive,subspace_of}.
    Hence $A$ is compact with respect to $T_A$ by \cref{compact_subspace_open_covering}.
\end{proof}

% from §26 Item 9: bijective continuous map from compact to Hausdorff is a homeomorphism
\begin{theorem}\label{compact_hausdorff_bijection_homeomorphism}
    Assume $f$ is a bijection from $X$ to $Y$.
    Assume $f$ is continuous from $X$ to $Y$ with respect to $T$ and $T_1$.
    Assume $X$ is compact with respect to $T$.
    Assume $Y$ is Hausdorff with respect to $T_1$.
    Then $f$ is a homeomorphism between $X$ and $Y$ with respect to $T$ and $T_1$.
\end{theorem}
\begin{proof}
    $T$ is a topology on $X$ and $T_1$ is a topology on $Y$ by \cref{continuous}.
    Show that for all $A$ such that $A$ is closed in $X$ with respect to $T$
        we have $\img{f}{A}$ is closed in $Y$ with respect to $T_1$.
    \begin{subproof}
        Fix $A$ such that $A$ is closed in $X$ with respect to $T$.
        Hence $A\subseteq X$ and $A$ is a subspace of $X$ with respect to $T$
            by \cref{closed_in,subspace_of}.
        Let $T_A = \subspaceTopology{T}{A}$.
        Therefore $A$ is compact with respect to $T_A$ by \cref{closed_subspace_compact}.
        Let $g = \restrl{f}{A}$.
        Hence $g$ is continuous from $A$ to $Y$ with respect to $T_A$ and $T_1$
            by \cref{continuous_restrict_domain}.
        Let $B = \img{g}{A}$.
        Therefore $B = \img{f}{A}$ by \cref{bijections_to_funs,subspace_of,subseteq,funs,restrl_img,inter_idempotent}.
        Let $T_B = \subspaceTopology{T_1}{B}$.
        Hence $B$ is compact with respect to $T_B$ by \cref{continuous_image_compact}.
        Therefore $B$ is closed in $Y$ with respect to $T_1$ by \cref{bijections_to_funs,funs,rels_ran_subseteq,img_subseteq_ran,subseteq_transitive,compact_subspace_closed}.
        Hence $\img{f}{A}$ is closed in $Y$ with respect to $T_1$.
    \end{subproof}
    Therefore $f$ is a closed map from $X$ to $Y$ with respect to $T$ and $T_1$ by \cref{bijections_to_funs,closed_map}.
    $\converse{f}\in\funs{Y}{X}$ by \cref{bijective_converse_are_funs}.
    For all $C$ such that $C$ is closed in $X$ with respect to $T$
        we have $\preimg{\converse{f}}{C}$ is closed in $Y$ with respect to $T_1$
        by \cref{closed_map,bijective_converse_are_funs,funs_is_relation,bijections_to_funs,preim_eq_img_of_converse,converse_converse_eq}.
    Hence $\converse{f}$ is continuous from $Y$ to $X$ with respect to $T_1$ and $T$
        by \cref{preimg_closed_implies_continuous}.
    Therefore $f$ is a homeomorphism between $X$ and $Y$ with respect to $T$ and $T_1$
        by \cref{homeomorphism}.
\end{proof}

% from §26 helper lemma 1: slice is compact
\begin{lemma}\label{slice_compact}
    Assume $T$ is a topology on $X$.
    Assume $T_1$ is a topology on $Y$.
    Assume $Y$ is compact with respect to $T_1$.
    Assume $x_0\in X$.
    Let $T_2 = \productTopology{T}{T_1}{X}{Y}$.
    Let $S = \{x_0\}\times Y$.
    Let $T_S = \subspaceTopology{T_2}{S}$.
    Then $S$ is compact with respect to $T_S$.
\end{lemma}
\begin{proof}
    Let $f_1 = \{(y,x_0) \mid y\in Y\}$.
    $f_1\in\funs{Y}{X}$ and for all $y$ such that $y\in Y$ we have $f_1(y) = x_0$
        by \cref{constant_map_function,funs_is_function,function_apply_intro}.
    Hence $f_1$ is continuous from $Y$ to $X$ with respect to $T_1$ and $T$
        by \cref{constant_continuous}.
    Let $f_2 = \identity{Y}$.
    Therefore $f_2$ is continuous from $Y$ to $Y$ with respect to $T_1$ and $T_1$
        by \cref{identity_continuous}.
    $f_2\in\funs{Y}{Y}$ and for all $y$ such that $y\in Y$ we have $f_2(y) = y$
        by \cref{id_iff,id_elem_funs,funs_is_function,function_apply_intro}.
    Let $k = \{(y,(f_1(y),f_2(y))) \mid y\in Y\}$.
    Hence $k$ is continuous from $Y$ to $X\times Y$ with respect to $T_1$ and $T_2$
        by \cref{continuous_map_into_product_backward}.
    Let $S_0 = \img{k}{Y}$.
    Therefore $S_0$ is compact with respect to $\subspaceTopology{T_2}{S_0}$ by \cref{continuous_image_compact}.
    For all $p$ such that $p\in S_0$ we have $p\in S$ by \cref{img_iff,pair_eq_iff,singleton_iff,times_tuple_intro}.
    For all $p$ such that $p\in S$ we have $p\in S_0$ by \cref{times_elem_is_tuple,singleton_iff,img_iff}.
    Therefore $S_0 = S$ by \cref{setext}.
    Hence $S$ is compact with respect to $T_S$.
\end{proof}

% from §26 helper lemma 2: tube lemma
\begin{lemma}\label{tube_lemma_helper}
    Assume $T$ is a topology on $X$.
    Assume $T_1$ is a topology on $Y$.
    Assume $Y$ is compact with respect to $T_1$.
    Let $T_2 = \productTopology{T}{T_1}{X}{Y}$.
    Assume $N$ is open in $X\times Y$ with respect to $T_2$.
    Assume $x_0\in X$.
    Suppose $\{x_0\}\times Y\subseteq N$.
    Then there exists $W$ such that
        $W$ is open in $X$ with respect to $T$ and
        $x_0\in W$ and
        $W\times Y\subseteq N$.
\end{lemma}
\begin{proof}
    Let $S = \{x_0\}\times Y$.
    Let $T_S = \subspaceTopology{T_2}{S}$.
    \begin{byCase}
    \caseOf{$Y = \emptyset$.}
        Let $W = X$.
        Hence there exists $W$ such that
            $W$ is open in $X$ with respect to $T$ and
            $x_0\in W$ and
            $W\times Y\subseteq N$
            by \cref{open_in,topology_on,subseteq_refl,times_empty_right,emptyset_subseteq}.
        \caseOf{$Y\neq\emptyset$.}
            Take $y_0$ such that $y_0\in Y$ by \cref{nonempty_inhabited}.
            Thus $x_0\in\{x_0\}$ by \cref{singleton_iff}.
            Thus $(x_0,y_0)\in S$ and $S$ is inhabited by \cref{times_tuple_intro}.
        Hence $S$ is compact with respect to $T_S$ by \cref{slice_compact}.
        Let $B = \productBasis{T}{T_1}{X}{Y}$.
        $B$ is a basis on $X\times Y$ by \cref{product_basis_is_basis}.
        Therefore $T_2$ is a topology on $X\times Y$ by \cref{basis_topology_is_topology,product_topology}.
        For all $p$ such that $p\in S$ we have $p\in X\times Y$
            by \cref{times_elem_is_tuple,singleton_iff,times_tuple_intro}.
        Thus $S$ is a subspace of $X\times Y$ with respect to $T_2$ by \cref{subseteq,subspace_of}.
        Let $R = \{w\in S\times B \mid \exists p\in S.\exists B_1\in B. w = (p,B_1) \land p\in B_1 \land B_1\subseteq N\}$.
        For all $w$ such that $w\in R$ there exist $p,B_1$ such that $w = (p,B_1)$.
        Hence $R$ is a relation and for all $p$ such that $p\in S$ we have $p\in N$ by \cref{relation,subseteq}.
        Therefore for all $p\in S$ there exists $B_1\in B$ such that $p\in B_1$ and $B_1\subseteq N$
            by \cref{open_in,product_topology,topology_generated_by_basis}.
        Therefore for all $p\in S$ there exists $B_1$ such that $p\mathrel{R} B_1$.
        Take $g$ such that $g$ is a function on $S$ and for all $p\in S$ we have $p\mathrel{R}\apply{g}{p}$
            by \cref{choice_function}.
        Let $C = \img{g}{S}$.
        For all $p$ such that $p\in S$ we have $p\in\unions{C}$ by \cref{function_apply_iff,pair_eq_iff,function_apply_elim,img_iff,unions_iff}.
        Hence $S\subseteq\unions{C}$ and $C$ is a covering of $S$ by \cref{subseteq,covering}.
        Hence for all $U$ such that $U\in C$ we have $U$ is open in $X\times Y$ with respect to $T_2$
            by \cref{img_iff,function_apply_iff,pair_eq_iff,basis_elem_open_in_generated,product_topology,open_in}.
        Take $F$ such that
            $F\subseteq C$ and $F$ is finite and $F$ is a covering of $S$
            by \cref{compact_subspace_open_covering}.
        Let $R_1 = \{w\in F\times T \mid \exists E\in F.\exists U\in T.\exists V\in T_1. w = (E,U) \land E = U\times V\}$.
        For all $w$ such that $w\in R_1$ there exist $E,U$ such that $w = (E,U)$.
        Hence $R_1$ is a relation by \cref{relation}.
        For all $E$ such that $E\in F$ there exists $p\in S$ such that $p\mathrel{R} E$ by \cref{subseteq,img_iff,function_apply_iff}.
        Therefore for all $E\in F$ we have $E\in B$.
        Therefore for all $E\in F$ there exists $U\in T$ such that there exists $V\in T_1$ such that $E = U\times V$
            by \cref{product_basis}.
        Hence for all $E\in F$ there exists $U$ such that $E\mathrel{R_1} U$.
        Take $h$ such that $h$ is a function on $F$ and for all $E\in F$ we have $E\mathrel{R_1}\apply{h}{E}$
            by \cref{choice_function}.
        Let $G = \img{h}{F}$.
        Thus $G$ is finite by \cref{finite_image_function}.
        Take $p_0$ such that $p_0\in S$ by \cref{nonempty_inhabited}.
        Take $E\in F$ such that $p_0\in E$ by \cref{covering,subseteq,unions_iff}.
        Therefore $G$ is inhabited by \cref{function_apply_elim,img_iff}.
        Therefore $G\subseteq T$ by \cref{img_iff,function_apply_iff,pair_eq_iff,subseteq}.
        Let $W = \inters{G}$.
        Hence $W$ is open in $X$ with respect to $T$ by \cref{finite_intersections_open_in_topology,open_in}.
        For all $U$ such that $U\in G$ there exists $E\in F$ such that $(E,U)\in h$
            by \cref{img_iff}.
        Thus for all $U$ such that $U\in G$ we have $x_0\in U$
            by \cref{function_apply_iff,subseteq,img_iff,pair_eq_iff,times_elem_is_tuple,singleton_iff,times_tuple_elim}.
        Hence $x_0\in W$ by \cref{inters_intro}.
        Show that for all $p$ such that $p\in W\times Y$ we have $p\in N$.
        \begin{subproof}
            Fix $p$ such that $p\in W\times Y$.
            Take $x,y$ such that $p = (x,y)$ and $x\in W$ and $y\in Y$
                by \cref{times_elem_is_tuple}.
            Therefore $(x_0,y)\in S$ by \cref{times_tuple_intro}.
            Take $E\in F$ such that $(x_0,y)\in E$ by \cref{covering,subseteq,unions_iff}.
            Take $E_1,U_1,V_1$ such that $E_1\in F$ and $U_1\in T$ and $V_1\in T_1$ and
                $(E,\apply{h}{E}) = (E_1,U_1)$ and $E_1 = U_1\times V_1$
                by \cref{function_apply_iff}.
            Hence $E = E_1$ and $\apply{h}{E} = U_1$ by \cref{pair_eq_iff}.
            Thus $E = U_1\times V_1$.
            Thus $(x_0,y)\in U_1\times V_1$.
            Thus $y\in V_1$ by \cref{times_tuple_elim}.
            Therefore $(E,\apply{h}{E})\in h$ and $U_1\in G$ by \cref{function_apply_elim,img_iff,pair_eq_iff}.
            Hence $W\subseteq U_1$ by \cref{inters_subseteq_elem}.
            Therefore $p\in U_1\times V_1$ by \cref{subseteq,times_tuple_intro}.
            Hence $p\in E$ and $E\in C$ by \cref{subseteq}.
            Take $p_0\in S$ such that $(p_0,E)\in g$ by \cref{img_iff}.
            Take $p_1,B_1$ such that $p_1\in S$ and $B_1\in B$ and
                $(p_0,E) = (p_1,B_1)$ and $p_1\in B_1$ and $B_1\subseteq N$
                by \cref{function_apply_iff}.
            Therefore $p\in N$ by \cref{pair_eq_iff,subseteq}.
        \end{subproof}
        Hence $W\times Y\subseteq N$ by \cref{subseteq}.
        Therefore there exists $W$ such that
            $W$ is open in $X$ with respect to $T$ and
            $x_0\in W$ and
            $W\times Y\subseteq N$.
    \end{byCase}
\end{proof}

% from §26 helper lemma 3: union of a finite family of finite sets is finite
\begin{lemma}\label{finite_union_of_finite_family}
    Suppose $F$ is finite.
    Suppose for all $A$ such that $A\in F$ we have $A$ is finite.
    Then $\unions{F}$ is finite.
\end{lemma}
\begin{proof}
    Take $k$ such that $k$ is a natural number and there exists a bijection from $k$ to $F$ by \cref{finite}.
    Take $f$ such that $f$ is a bijection from $k$ to $F$.
    Let $A = \{ n\in\naturals \mid \text{for all $G,g$ such that $g$ is a bijection from $n$ to $G$ and for all $B\in G$ we have $B$ is finite we have $\unions{G}$ is finite} \}$.
    For all $G,g$ such that $g$ is a bijection from $\emptyset$ to $G$ and
        for all $B\in G$ we have $B$ is finite we have $\unions{G} = \emptyset$ and $\unions{G}$ is finite
        by \cref{bijections,surj,emptyset,nonempty_inhabited,unions_emptyset,emptyset_is_finite}.
    Therefore $\emptyset\in A$.
        Show that for all $n\in A$ we have $\suc{n}\in A$.
        \begin{subproof}
            Fix $n$ such that $n\in A$.
            Show that for all $G,g$ such that $g$ is a bijection from $\suc{n}$ to $G$
                and for all $B\in G$ we have $B$ is finite we have $\unions{G}$ is finite.
            \begin{subproof}
                Fix $G$.
                Fix $g$ such that $g$ is a bijection from $\suc{n}$ to $G$ and for all $B\in G$ we have $B$ is finite.
                $g$ is a function and $\dom{g} = \suc{n}$ by \cref{bijection_is_function,bijections_dom}.
                Hence $n\in\dom{g}$ by \cref{suc_intro_self}.
                Let $a = g(n)$.
                Let $G_1 = \img{g}{n}$.
                Therefore $G_1\subseteq G$ by \cref{suc,subseteq,img_subseteq,img_dom,bijections,ran_of_surj}.
                $a\in G$ by \cref{function_apply_elim,ran_iff,bijections,ran_of_surj}.
                For all $z$ such that $z\in G$ there exists $x\in\suc{n}$ such that $g(x)=z$ by \cref{bijections,surj}.
                Thus for all $z$ such that $z\in G$ we have $z\in G_1\union\{a\}$
                    by \cref{suc,function_apply_elim,img_iff,union_intro_left,union_intro_right,singleton_intro}.
                Therefore $G\subseteq G_1\union\{a\}$ by \cref{subseteq}.
                Hence $G_1\union\{a\}\subseteq G$ by \cref{singleton_subset_intro,union_subsets_is_subset}.
                Therefore $G = G_1\union\{a\}$ by \cref{subseteq_antisymmetric}.
                Thus for all $B$ such that $B\in G_1$ we have $B\in G$ and $B$ is finite by \cref{subseteq}.
                Let $h = \restrl{g}{n}$.
                $h$ is a function by \cref{restrl_of_function_is_function}.
                Thus $n\subseteq\dom{g}$ by \cref{suc,bijections_dom,subseteq}.
                Hence $\dom{h} = n$ by \cref{dom_of_restrl_eq_inter,inter_absorb_supseteq_left}.
                Thus $h\in\funs{n}{G_1}$ by \cref{restrl_of_function_is_function,function_apply_elim,restrl_subseteq,elem_subseteq,img_iff,funs_intro}.
                For all $y\in G_1$ there exists $x\in n$ such that $h(x) = y$
                    by \cref{img_iff,function_apply_iff,restrl_of_function_apply}.
                Hence $h\in\surj{n}{G_1}$ by \cref{surj}.
                For all $x,y$ such that $x\in\dom{h}$ and $y\in\dom{h}$ and $h(x) = h(y)$ we have $x\in n$ and $y\in n$.
                Thus for all $x,y$ such that $x\in\dom{h}$ and $y\in\dom{h}$ and $h(x) = h(y)$
                    we have $h(x) = g(x)$ and $h(y) = g(y)$ by \cref{restrl_of_function_apply}.
                Thus for all $x,y$ such that $x\in\dom{h}$ and $y\in\dom{h}$ and $h(x) = h(y)$ we have $g(x) = g(y)$.
                Thus for all $x,y$ such that $x\in\dom{h}$ and $y\in\dom{h}$ and $h(x) = h(y)$
                    we have $x\in\dom{g}$ and $y\in\dom{g}$ by \cref{subseteq}.
                Thus for all $x,y$ such that $x\in\dom{h}$ and $y\in\dom{h}$ and $h(x) = h(y)$ we have $x = y$
                    by \cref{bijections,injective_function}.
                Therefore $h$ is injective by \cref{injective_function}.
                Hence $h$ is a bijection from $n$ to $G_1$ by \cref{bijections}.
                Thus $\unions{G_1}$ is finite and $a$ is finite.
                Therefore $\unions{G} = a\union\unions{G_1}$ by \cref{union_comm,union_eq_cons,unions_cons}.
                Hence $\unions{G}$ is finite by \cref{union_of_finite_is_finite}.
            \end{subproof}
        \end{subproof}
    Therefore $A$ is an inductive set.
    Hence $\naturals\subseteq A$ and $k\in\naturals$ and $k\in A$ by \cref{num_naturals_smallest_inductive_set,subseteq}.
    Therefore for all $G,g$ such that $g$ is a bijection from $k$ to $G$ and for all $B\in G$ we have $B$ is finite
        we have $\unions{G}$ is finite.
    Thus $\unions{F}$ is finite.
\end{proof}

% from §26 Item 10: product of finitely many compact spaces is compact
\begin{theorem}\label{product_compact}
    Assume $T$ is a topology on $X$.
    Assume $T_1$ is a topology on $Y$.
    Assume $X$ is compact with respect to $T$.
    Assume $Y$ is compact with respect to $T_1$.
    Let $T_2 = \productTopology{T}{T_1}{X}{Y}$.
    Then $X\times Y$ is compact with respect to $T_2$.
\end{theorem}
\begin{proof}
    Hence $T_2$ is a topology on $X\times Y$ by \cref{product_basis_is_basis,basis_topology_is_topology,product_topology}.
    Show that for all $A$ such that $A$ is an open covering of $X\times Y$ with respect to $T_2$
        there exists $B$ such that
            $B\subseteq A$ and
            $B$ is finite and
            $B$ is a covering of $X\times Y$.
    \begin{subproof}
        Fix $A$ such that $A$ is an open covering of $X\times Y$ with respect to $T_2$.
        Let $C = \{W\in T \mid \text{there exists $B\subseteq A$ such that $B$ is finite and $W\times Y\subseteq \unions{B}$}\}$.
        For all $W$ such that $W\in C$ we have $W\in T$ and $W$ is open in $X$ with respect to $T$ by \cref{open_in}.
        Show that for all $x$ such that $x\in X$ we have $x\in\unions{C}$.
        \begin{subproof}
            Fix $x$ such that $x\in X$.
            Let $S = \{x\}\times Y$.
            Let $T_S = \subspaceTopology{T_2}{S}$.
            Hence $S$ is compact with respect to $T_S$ by \cref{slice_compact}.
            For all $p$ such that $p\in S$ we have $p\in X\times Y$
                by \cref{times_elem_is_tuple,singleton_iff,times_tuple_intro}.
            Hence $S$ is a subspace of $X\times Y$ with respect to $T_2$ by \cref{subseteq,subspace_of}.
            Therefore $S\subseteq\unions{A}$ and $A$ is a covering of $S$ by \cref{open_covering,covering,subseteq}.
            For all $U$ such that $U\in A$ we have $U$ is open in $X\times Y$ with respect to $T_2$
                by \cref{open_covering}.
            Take $B$ such that
                $B\subseteq A$ and $B$ is finite and $B$ is a covering of $S$
                by \cref{compact_subspace_open_covering}.
            Let $N = \unions{B}$.
            Therefore $N$ is open in $X\times Y$ with respect to $T_2$ by \cref{subseteq,open_covering,open_in,topology_on}.
            Take $W$ such that
                $W$ is open in $X$ with respect to $T$ and
                $x\in W$ and
                $W\times Y\subseteq N$ by \cref{covering,tube_lemma_helper}.
            Hence $x\in\unions{C}$ by \cref{open_in,unions_iff}.
        \end{subproof}
        Therefore $C$ is an open covering of $X$ with respect to $T$ by \cref{subseteq,covering,open_covering}.
        Take $F$ such that
            $F\subseteq C$ and $F$ is finite and $F$ is a covering of $X$
            by \cref{compact}.
        Let $R = \{w\in F\times\pow{A} \mid \text{there exists $W\in F$ such that there exists $B\subseteq A$ such that $w = (W,B)$ and $B$ is finite and $W\times Y\subseteq \unions{B}$}\}$.
        For all $w$ such that $w\in R$ there exist $W,B$ such that $w = (W,B)$.
        Therefore $R$ is a relation by \cref{relation}.
        For all $W\in F$ there exists $B\subseteq A$ such that $B$ is finite and $W\times Y\subseteq \unions{B}$ by \cref{subseteq}.
        Hence for all $W\in F$ there exists $B$ such that $(W,B)\in R$.
        Therefore for all $W\in F$ there exists $B$ such that $W\mathrel{R} B$.
        Take $h$ such that $h$ is a function on $F$ and for all $W\in F$ we have $W\mathrel{R}\apply{h}{W}$
            by \cref{choice_function}.
        Let $H = \img{h}{F}$.
        For all $B$ such that $B\in H$ we have $B$ is finite by \cref{img_iff,function_apply_iff,pair_eq_iff}.
        Let $B = \unions{H}$.
        Hence $B$ is finite by \cref{finite_image_function,finite_union_of_finite_family}.
        Therefore $B\subseteq A$ by \cref{img_iff,function_apply_iff,pair_eq_iff,subseteq,powerset_intro,unions_subseteq_of_powerset_is_subseteq}.
        Show that for all $p$ such that $p\in X\times Y$ we have $p\in\unions{B}$.
        \begin{subproof}
            Fix $p$ such that $p\in X\times Y$.
            Take $x,y$ such that $p = (x,y)$ and $x\in X$ and $y\in Y$
                by \cref{times_elem_is_tuple}.
            Take $W\in F$ such that $x\in W$ by \cref{covering,subseteq,unions_iff}.
            Hence $p\in W\times Y$ by \cref{times_tuple_intro}.
            Take $W_1,B_1$ such that $W_1\in F$ and $B_1\subseteq A$ and
                $(W,\apply{h}{W}) = (W_1,B_1)$ and $B_1$ is finite and $W_1\times Y\subseteq \unions{B_1}$
                by \cref{function_apply_iff}.
            Take $U\in\apply{h}{W}$ such that $p\in U$ by \cref{pair_eq_iff,subseteq,unions_iff}.
            Hence $p\in\unions{B}$ by \cref{function_apply_elim,img_iff,unions_iff}.
        \end{subproof}
        Therefore $X\times Y\subseteq\unions{B}$ and $B$ is a covering of $X\times Y$ by \cref{subseteq,covering}.
    \end{subproof}
    Hence $X\times Y$ is compact with respect to $T_2$ by \cref{compact}.
\end{proof}

% from §26 Item 11: tube lemma
\begin{lemma}\label{tube_lemma}
    Assume $T$ is a topology on $X$.
    Assume $T_1$ is a topology on $Y$.
    Assume $Y$ is compact with respect to $T_1$.
    Let $T_2 = \productTopology{T}{T_1}{X}{Y}$.
    Assume $N$ is open in $X\times Y$ with respect to $T_2$.
    Assume $x_0\in X$.
    Suppose $\{x_0\}\times Y\subseteq N$.
    Then there exists $W$ such that
        $W$ is open in $X$ with respect to $T$ and
        $x_0\in W$ and
        $W\times Y\subseteq N$.
\end{lemma}
\begin{proof}
    Follows by \cref{tube_lemma_helper}.
\end{proof}
% from §26 Item 12: finite intersection property
\begin{definition}\label{finite_intersection_property}
    $C$ has the finite intersection property iff
        $C$ is inhabited and
        for all $F$ such that $F\subseteq C$ and $F$ is finite and inhabited
        we have $\inters{F}\neq\emptyset$.
\end{definition}
% from §26 Item 13: compactness via finite intersection property
\begin{theorem}\label{compact_iff_finite_intersection_property}
    Assume $T$ is a topology on $X$.
    Then $X$ is compact with respect to $T$ iff
        for all $C$ if
            $C$ has the finite intersection property and
            for all $A$ such that $A\in C$ we have $A$ is closed in $X$ with respect to $T$
        then $\inters{C}\neq\emptyset$.
\end{theorem}
\begin{proof}
    Show that if $X$ is compact with respect to $T$ then
        for all $C$ if
            $C$ has the finite intersection property and
            for all $A$ such that $A\in C$ we have $A$ is closed in $X$ with respect to $T$
        then $\inters{C}\neq\emptyset$.
    \begin{subproof}
        Assume $X$ is compact with respect to $T$.
        Fix $C$ such that $C$ has the finite intersection property and
            for all $A$ such that $A\in C$ we have $A$ is closed in $X$ with respect to $T$.
        Suppose not.
        Then $\inters{C} = \emptyset$ by \cref{emptyset}.
        Take $A_0$ such that $A_0\in C$ by \cref{finite_intersection_property}.
        Let $F_0 = \{A_0\}$.
        Hence $F_0$ is finite and $F_0\subseteq C$
            by \cref{pair_is_finite,upair_intro_left,singleton_subset_intro,subseteq_finite_is_finite}.
        $F_0$ is inhabited.
        Therefore $\inters{F_0}\neq\emptyset$ by \cref{finite_intersection_property}.
        Hence $A_0\neq\emptyset$ by \cref{inters_singleton,emptyset}.
        Let $A = \{U \mid \exists D\in C. U = X\setminus D\}$.
        Show that $A$ is an open covering of $X$ with respect to $T$.
        \begin{subproof}
            For all $U$ such that $U\in A$ we have $U$ is open in $X$ with respect to $T$ by \cref{closed_in}.
            Show that $A$ is a covering of $X$.
            \begin{subproof}
                Show that for all $x$ such that $x\in X$ we have $x\in\unions{A}$.
                \begin{subproof}
                    Fix $x$ such that $x\in X$.
                    Suppose not.
                    For all $D\in C$ we have $X\setminus D\in A$.
                    Thus for all $D$ such that $D\in C$ we have $x\in D$.
                    Thus $x\in\inters{C}$ and $\inters{C}\neq\emptyset$ by \cref{inters_intro,emptyset}.
                    Contradiction.
                \end{subproof}
                $X\subseteq\unions{A}$ and $A$ is a covering of $X$ by \cref{subseteq,covering}.
            \end{subproof}
            Hence $A$ is an open covering of $X$ with respect to $T$ by \cref{open_covering}.
        \end{subproof}
        Take $B$ such that $B\subseteq A$ and $B$ is finite and $B$ is a covering of $X$
            by \cref{compact}.
        Thus $B$ is inhabited by \cref{closed_in,subseteq,nonempty_inhabited,covering,unions_iff}.
        Let $R = \{w\in B\times C \mid \exists U\in B.\exists D\in C. w = (U,D) \land U = X\setminus D\}$.
        For all $w$ such that $w\in R$ there exist $U,D$ such that $w = (U,D)$ and
            $R$ is a relation by \cref{times_elem_is_tuple,relation}.
        For all $U\in B$ there exists $D\in C$ such that $U = X\setminus D$ by \cref{subseteq}.
        Thus for all $U\in B$ there exists $D$ such that $U\mathrel{R} D$.
        Take $h$ such that $h$ is a function on $B$ and for all $U\in B$ we have $U\mathrel{R}\apply{h}{U}$
            by \cref{choice_function}.
        Let $F = \img{h}{B}$.
        Now $F$ is inhabited and $F\subseteq C$
            by \cref{nonempty_inhabited,function_apply_elim,img_iff,function_apply_iff,pair_eq_iff,subseteq}.
        Thus for all $x$ such that $x\in X\setminus\unions{B}$ we have $x\in\inters{F}$
            by \cref{img_iff,function_apply_iff,setminus_elim_left,setminus_elim_right,pair_eq_iff,setminus_intro,unions_iff,inters_intro}.
        Thus $X\setminus\unions{B}\subseteq\inters{F}$ by \cref{subseteq}.
        Show that for all $x$ such that $x\in\inters{F}$ we have $x\in X\setminus\unions{B}$.
        \begin{subproof}
            Fix $x$ such that $x\in\inters{F}$.
            Thus $x\in X$ by \cref{nonempty_inhabited,inters_destr,subseteq,closed_in}.
            Suppose not.
            Then $x\in\unions{B}$.
            Take $U\in B$ such that $x\in U$ by \cref{unions_iff}.
            Thus $U\mathrel{R}\apply{h}{U}$.
            Take $U_1,D_1$ such that $U_1\in B$ and $D_1\in C$ and
                $(U,\apply{h}{U}) = (U_1,D_1)$ and $U_1 = X\setminus D_1$.
            Thus $U = U_1$ and $\apply{h}{U} = D_1$ by \cref{pair_eq_iff}.
            Thus $x\in D_1$ by \cref{img_iff,function_apply_elim,inters_destr}.
            Thus $x\in X\setminus D_1$ by \cref{pair_eq_iff}.
            Thus $x\notin D_1$ by \cref{setminus_elim_right}.
            Contradiction.
        \end{subproof}
        Therefore $\inters{F}\subseteq X\setminus\unions{B}$ by \cref{subseteq}.
        Hence $X\setminus\unions{B} = \inters{F}$ by \cref{subseteq_antisymmetric}.
        For all $U$ such that $U\in B$ we have $U\subseteq X$.
        Therefore $X\setminus\unions{B} = \emptyset$
            by \cref{covering,powerset_intro,subseteq,unions_subseteq_of_powerset_is_subseteq,subseteq_antisymmetric,setminus_self}.
        Thus $\inters{F} = \emptyset$ by \cref{subseteq_antisymmetric,subseteq_refl}.
        Finally $\inters{F}\neq\emptyset$ by \cref{finite_image_function,finite_intersection_property}.
        Contradiction.
    \end{subproof}
    Show that if
        for all $C$ if
            $C$ has the finite intersection property and
            for all $A$ such that $A\in C$ we have $A$ is closed in $X$ with respect to $T$
        then $\inters{C}\neq\emptyset$
        then $X$ is compact with respect to $T$.
    \begin{subproof}
        Assume for all $C$ if
            $C$ has the finite intersection property and
            for all $A$ such that $A\in C$ we have $A$ is closed in $X$ with respect to $T$
        then $\inters{C}\neq\emptyset$.
        Show that for all $A$ such that $A$ is an open covering of $X$ with respect to $T$
            there exists $B$ such that
                $B\subseteq A$ and
                $B$ is finite and
                $B$ is a covering of $X$.
        \begin{subproof}
            Fix $A$ such that $A$ is an open covering of $X$ with respect to $T$.
            \begin{byCase}
            \caseOf{$X = \emptyset$.}
                Let $B = \emptyset$.
                $B\subseteq A$ and $B$ is finite by \cref{emptyset_subseteq,emptyset_is_finite}.
                Thus $B$ is a covering of $X$ by \cref{unions_emptyset,subseteq_refl,covering}.
                Thus there exists $B$ such that
                    $B\subseteq A$ and $B$ is finite and $B$ is a covering of $X$.
            \caseOf{$X \neq \emptyset$.}
                Show that there exists $B$ such that
                    $B\subseteq A$ and $B$ is finite and $B$ is a covering of $X$.
                \begin{subproof}
                    Suppose not.
                    Then for all $B$ such that $B\subseteq A$ and $B$ is finite we have
                        $B$ is not a covering of $X$.
                    Let $C = \{D \mid \exists U\in A. D = X\setminus U\}$.
                    For all $D$ such that $D\in C$ we have $D$ is closed in $X$ with respect to $T$
                        by \cref{open_covering,open_in,topology_on,subseteq,pow_iff,double_relative_complement,setminus_subseteq,closed_in}.
                    Show that $C$ has the finite intersection property.
                    \begin{subproof}
                        Take $x_0,U_1$ such that $x_0\in X$ and $U_1\in A$ and $x_0\in U_1$
                            by \cref{open_covering,nonempty_inhabited,covering,subseteq,unions_iff}.
                        Thus $C$ is inhabited.
                        Show that for all $F$ such that $F\subseteq C$ and $F$ is finite and inhabited
                            we have $\inters{F}\neq\emptyset$.
                        \begin{subproof}
                            Fix $F$ such that $F\subseteq C$ and $F$ is finite and inhabited.
                            Suppose not.
                            Then $\inters{F} = \emptyset$ by \cref{emptyset}.
                            Let $R_0 = \{w\in F\times A \mid \exists D\in F.\exists U\in A. w = (D,U) \land D = X\setminus U\}$.
                            For all $w$ such that $w\in R_0$ there exist $D,U$ such that $w = (D,U)$ and
                                $R_0$ is a relation by \cref{times_elem_is_tuple,relation}.
                            For all $D\in F$ there exists $U\in A$ such that $D = X\setminus U$ by \cref{subseteq}.
                            Thus for all $D\in F$ there exists $U$ such that $D\mathrel{R_0} U$.
                            Take $h$ such that $h$ is a function on $F$ and for all $D\in F$ we have $D\mathrel{R_0}\apply{h}{D}$
                                by \cref{choice_function}.
                            Let $B = \img{h}{F}$.
                            Thus $B$ is finite by \cref{finite_image_function}.
                            Thus $B\subseteq A$ by \cref{img_iff,function_apply_iff,pair_eq_iff,subseteq}.
                            For all $x$ such that $x\in X\setminus\unions{B}$ we have $x\in X$ and $x\notin\unions{B}$
                                by \cref{setminus_elim_left,setminus_elim_right}.
                            Thus for all $x$ such that $x\in X\setminus\unions{B}$ we have $x\in\inters{F}$
                                by \cref{pair_eq_iff,setminus_intro,function_apply_iff,function_apply_elim,img_iff,unions_iff,inters_intro}.
                            Thus $X\setminus\unions{B}\subseteq\inters{F}$ by \cref{subseteq}.
                            Show that for all $x$ such that $x\in\inters{F}$ we have $x\in X\setminus\unions{B}$.
                            \begin{subproof}
                                Fix $x$ such that $x\in\inters{F}$.
                                Thus $x\in X$ by \cref{nonempty_inhabited,inters_destr,subseteq,closed_in}.
                                Suppose not.
                                Then $x\in\unions{B}$.
                                Take $U\in B$ such that $x\in U$ by \cref{unions_iff}.
                                Take $D\in F$ such that $(D,U)\in h$ by \cref{img_iff}.
                                Thus $D\mathrel{R_0} U$ by \cref{function_apply_iff}.
                                Take $D_1,U_1$ such that $D_1\in F$ and $U_1\in A$ and
                                    $(D,U) = (D_1,U_1)$ and $D_1 = X\setminus U_1$.
                                Thus $D = D_1$ and $U = U_1$ by \cref{pair_eq_iff}.
                                Thus $x\in X\setminus D$ by \cref{open_covering,open_in,topology_on,subseteq,pow_iff,double_relative_complement,pair_eq_iff}.
                                Thus $x\notin D$ by \cref{setminus_elim_right}.
                                Thus $x\in D$ by \cref{inters_destr}.
                                Contradiction.
                            \end{subproof}
                            Therefore $\inters{F}\subseteq X\setminus\unions{B}$ by \cref{subseteq}.
                            Hence $X\setminus\unions{B} = \inters{F}$ by \cref{subseteq_antisymmetric}.
                            Then $X\setminus\unions{B} = \emptyset$.
                            Therefore $B$ is a covering of $X$ by \cref{setminus_eq_emptyset_iff_subseteq,covering}.
                            Contradiction.
                        \end{subproof}
                        Hence $C$ has the finite intersection property by \cref{finite_intersection_property}.
                    \end{subproof}
                    Therefore $\inters{C}\neq\emptyset$.
                    Show that for all $x$ such that $x\in X\setminus\unions{A}$ we have $x\in\inters{C}$.
                    \begin{subproof}
                        Fix $x$ such that $x\in X\setminus\unions{A}$.
                        Thus $x\in X$ and $x\notin\unions{A}$ by \cref{setminus_elim_left,setminus_elim_right}.
                        Thus for all $D$ such that $D\in C$ we have $x\in D$.
                        Thus $x\in\inters{C}$ by \cref{finite_intersection_property,setminus_elim_right,inters_intro}.
                    \end{subproof}
                    Thus $X\setminus\unions{A}\subseteq\inters{C}$ by \cref{subseteq}.
                    Show that for all $x$ such that $x\in\inters{C}$ we have $x\in X\setminus\unions{A}$.
                    \begin{subproof}
                        Fix $x$ such that $x\in\inters{C}$.
                        Suppose not.
                        Then $x\in\unions{A}$.
                        Take $U\in A$ such that $x\in U$ by \cref{unions_iff}.
                        Thus $X\setminus U\in C$ and $x\in X\setminus U$ by \cref{inters_destr}.
                        Thus $x\notin U$ by \cref{setminus_elim_right}.
                        Contradiction.
                    \end{subproof}
                    Therefore $\inters{C}\subseteq X\setminus\unions{A}$ by \cref{subseteq}.
                    Hence $X\setminus\unions{A} = \inters{C}$ by \cref{subseteq_antisymmetric}.
                    Therefore $X\setminus\unions{A} = \emptyset$
                        by \cref{covering,powerset_intro,open_covering,open_in,topology_on,subseteq,pow_iff,unions_subseteq_of_powerset_is_subseteq,subseteq_antisymmetric,setminus_self}.
                    Thus $\inters{C} = \emptyset$ by \cref{subseteq_antisymmetric,subseteq_refl}.
                    Contradiction.
                \end{subproof}
            \end{byCase}
        \end{subproof}
        Therefore $X$ is compact with respect to $T$ by \cref{compact}.
    \end{subproof}
\end{proof}

\section{§27 Compact Subspaces of the Real Line}

% from §27 helper definition 1: immediate successor
\begin{definition}\label{immediate_successor}
    $y$ is an immediate successor of $x$ in $X$ with respect to $R$ iff
        $x\in X$ and $y\in X$ and $x\mathrel{R}y$ and
        $\intervalclosedrel{R}{X}{x}{y} = \{x,y\}$.
\end{definition}

% from §27 helper lemma 1: interval of an immediate successor
\begin{lemma}\label{immediate_successor_interval}
    Assume $R$ is asymmetric.
    Suppose $y$ is an immediate successor of $x$ in $X$ with respect to $R$.
    Let $I = \intervalclosedrel{R}{X}{x}{y}$.
    Then $I = \{x,y\}$.
\end{lemma}
\begin{proof}
    Follows by \cref{immediate_successor}.
\end{proof}

% from §27 helper lemma 2: non-immediate successor gives an intermediate point
\begin{lemma}\label{no_immediate_successor_between}
    Suppose $x\in X$.
    Suppose $y\in X$.
    Suppose $x\mathrel{R}y$.
    Suppose it is wrong that $y$ is an immediate successor of $x$ in $X$ with respect to $R$.
    Then there exists $z\in X$ such that $x\mathrel{R}z$ and $z\mathrel{R}y$.
\end{lemma}
\begin{proof}
    Let $I = \intervalclosedrel{R}{X}{x}{y}$.
    Suppose not.
    $\{x,y\}\subseteq I$ by \cref{upair_elim,intervalclosed_rel,subseteq}.
    Show that for all $z$ such that $z\in I$ we have $z\in\{x,y\}$.
    \begin{subproof}
        Fix $z$ such that $z\in I$.
        Thus $(x\mathrel{R}z \lor z = x)$ and $(z\mathrel{R}y \lor z = y)$
            by \cref{intervalclosed_rel}.
        \begin{byCase}
        \caseOf{$z = x$.}
            Thus $z\in\{x,y\}$ by \cref{upair_intro_left}.
        \caseOf{$z = y$.}
            Thus $z\in\{x,y\}$ by \cref{upair_intro_right}.
        \caseOf{$x\mathrel{R}z$ and $z\mathrel{R}y$.}
            Contradiction.
        \end{byCase}
    \end{subproof}
    Hence $y$ is an immediate successor of $x$ in $X$ with respect to $R$
        by \cref{subseteq,subseteq_antisymmetric,immediate_successor}.
    Contradiction.
\end{proof}

% from §27 Item 1: closed interval is compact in an ordered set with the least upper bound property
\begin{theorem}\label{ordered_closed_interval_compact}
    Assume $R$ is a simple order relation on $X$.
    Assume $X$ has the least upper bound property with respect to $R$.
    Assume $T$ is the order topology on $X$ with respect to $R$.
    Suppose $a\in X$.
    Suppose $b\in X$.
    Suppose $a\mathrel{R}b$ or $a = b$.
    Let $Y = \intervalclosedrel{R}{X}{a}{b}$.
    Let $T_Y = \subspaceTopology{T}{Y}$.
    Then $Y$ is compact with respect to $T_Y$.
\end{theorem}
\begin{proof}
    $R$ is transitive and $R$ is asymmetric and $R$ is connex on $X$ by \cref{simple_order_relation}.
    Show that for all $A$ such that $A$ is an open covering of $Y$ with respect to $T_Y$
        there exists $B$ such that $B\subseteq A$ and $B$ is finite and $B$ is a covering of $Y$.
    \begin{subproof}
        Fix $A$ such that $A$ is an open covering of $Y$ with respect to $T_Y$.
        \begin{byCase}
        \caseOf{$a = b$.}
            Take $U\in A$ such that $a\in U$
                by \cref{open_covering,covering,intervalclosed_rel,subseteq,unions_iff}.
            Let $B = \{U\}$.
            Hence $B\subseteq A$ and $B$ is finite
                by \cref{singleton_subset_intro,pair_is_finite,upair_intro_left,subseteq_finite_is_finite}.
            Thus for all $x$ such that $x\in Y$ we have $x\in\unions{B}$
                by \cref{intervalclosed_rel,asymmetric,singleton_intro,unions_iff}.
            Therefore $B$ is a covering of $Y$ by \cref{subseteq,covering}.
            Hence there exists $B$ such that $B\subseteq A$ and $B$ is finite and $B$ is a covering of $Y$.
        \caseOf{$a\mathrel{R}b$.}
            $a\in Y$ and $b\in Y$ by \cref{intervalclosed_rel}.
            Hence $a\neq b$ by \cref{asymmetric}.
            Therefore $X$ has more than one element and $Y$ has more than one element by \cref{more_than_one_element}.
            $Y$ is convex in $X$ with respect to $R$ by \cref{intervalclosedrel_convex}.
            For all $y$ such that $y\in Y$ we have $y\in X$ by \cref{intervalclosed_rel}.
            Hence $Y\subseteq X$ by \cref{subseteq}.
            Then $a$ is a smallest element of $Y$ with respect to $R$ by \cref{intervalclosed_rel,smallest_element}.
            Finally $b$ is a largest element of $Y$ with respect to $R$ by \cref{intervalclosed_rel,largest_element}.
            $T = \basisTopology{\orderBasis{R}{X}}{X}$ by \cref{order_topology}.
            Therefore $T$ is a topology on $X$ by \cref{order_basis_is_basis,basis_topology_is_topology}.
            Hence $T_Y$ is a topology on $Y$ by \cref{subspace_topology_is_topology,subspace_of}.

            Show that for all $x$ such that $x\in Y$ and $x\neq b$ there exists $y\in Y$ such that
                $x\mathrel{R}y$ and
                there exists $B$ such that $B\subseteq A$ and $B$ is finite and
                    $B$ is a covering of $\intervalclosedrel{R}{X}{x}{y}$.
            \begin{subproof}
                Fix $x$ such that $x\in Y$ and $x\neq b$.
                Take $U\in A$ such that $x\in U$ by \cref{open_covering,covering,subseteq,unions_iff}.
                Take $V\in T$ such that $U = Y\inter V$ by \cref{open_covering,open_in,subspace_topology}.
                Thus $x\in V$ by \cref{inter_elim_right}.
                Take $B_1\in\orderBasis{R}{X}$ such that $x\in B_1$ and $B_1\subseteq V$
                    by \cref{topology_generated_by_basis}.
                Show that there exists $c\in Y$ such that $x\mathrel{R}c$ and
                    $\intervalclosedopenrel{R}{X}{x}{c} \subseteq V$.
                \begin{subproof}
                    $B_1\in\pow{X}$ and
                        $((\exists u, v\in X. u\mathrel{R}v \land B_1 = \intervalopenrel{R}{X}{u}{v}) \lor
                        (\exists v\in X. \exists a_0\in X. (\forall z\in X. a_0\mathrel{R}z \lor a_0 = z) \land
                            B_1 = \intervalclosedopenrel{R}{X}{a_0}{v}) \lor
                        (\exists u\in X. \exists b_0\in X. (\forall z\in X. z\mathrel{R}b_0 \lor z = b_0) \land
                            B_1 = \intervalopenclosedrel{R}{X}{u}{b_0}))$
                        by \cref{order_basis}.
                    \begin{byCase}
                    \caseOf{$\exists u, v\in X. u\mathrel{R}v \land B_1 = \intervalopenrel{R}{X}{u}{v}$.}
                        Take $u, v\in X$ such that $u\mathrel{R}v$ and $B_1 = \intervalopenrel{R}{X}{u}{v}$.
                        Thus $u\mathrel{R}x$ and $x\mathrel{R}v$ by \cref{intervalopen_rel}.
                        Show that there exists $c\in Y$ such that $x\mathrel{R}c$ and
                            $\intervalclosedopenrel{R}{X}{x}{c} \subseteq B_1$.
                        \begin{subproof}
                            \begin{byCase}
                            \caseOf{$v\in Y$.}
                                Let $c = v$.
                                Thus $\intervalclosedopenrel{R}{X}{x}{c} \subseteq B_1$
                                    by \cref{intervalclosedopen_rel,transitive,intervalopen_rel,subseteq}.
                            \caseOf{it is wrong that $v\in Y$.}
                                Thus $a\mathrel{R}v$ by \cref{intervalclosed_rel,transitive}.
                                Thus $b\mathrel{R}v$ by \cref{connex,intervalclosed_rel}.
                                Thus $\intervalclosedopenrel{R}{X}{x}{b} \subseteq B_1$
                                    by \cref{intervalclosedopen_rel,transitive,intervalopen_rel,subseteq}.
                                Thus there exists $c\in Y$ such that $x\mathrel{R}c$ and
                                    $\intervalclosedopenrel{R}{X}{x}{c} \subseteq B_1$ by \cref{intervalclosed_rel}.
                            \end{byCase}
                        \end{subproof}
                    \caseOf{$\exists v\in X. \exists a_0\in X. (\forall z\in X. a_0\mathrel{R}z \lor a_0 = z) \land
                        B_1 = \intervalclosedopenrel{R}{X}{a_0}{v}$.}
                        Take $a_0, v\in X$ such that
                            $(\forall z\in X. a_0\mathrel{R}z \lor a_0 = z)$ and
                            $B_1 = \intervalclosedopenrel{R}{X}{a_0}{v}$.
                        Thus $x\in X$ and $(a_0\mathrel{R}x \lor x = a_0)$ and $x\mathrel{R}v$
                            by \cref{intervalclosedopen_rel}.
                        Show that there exists $c\in Y$ such that $x\mathrel{R}c$ and
                            $\intervalclosedopenrel{R}{X}{x}{c} \subseteq B_1$.
                        \begin{subproof}
                            \begin{byCase}
                            \caseOf{$v\in Y$.}
                                Let $c = v$.
                                Thus $\intervalclosedopenrel{R}{X}{x}{c} \subseteq B_1$
                                    by \cref{intervalclosedopen_rel,smallest_element,intervalclosed_rel,subseteq}.
                            \caseOf{it is wrong that $v\in Y$.}
                                Thus $a\mathrel{R}v$ by \cref{intervalclosed_rel,transitive}.
                                Thus $b\mathrel{R}v$ by \cref{connex,intervalclosed_rel}.
                                Thus $\intervalclosedopenrel{R}{X}{x}{b} \subseteq B_1$
                                    by \cref{intervalclosedopen_rel,smallest_element,transitive,subseteq}.
                                Thus there exists $c\in Y$ such that $x\mathrel{R}c$ and
                                    $\intervalclosedopenrel{R}{X}{x}{c} \subseteq B_1$ by \cref{intervalclosed_rel}.
                            \end{byCase}
                        \end{subproof}
                    \caseOf{$\exists u\in X. \exists b_0\in X. (\forall z\in X. z\mathrel{R}b_0 \lor z = b_0) \land
                        B_1 = \intervalopenclosedrel{R}{X}{u}{b_0}$.}
                        Take $u, b_0\in X$ such that
                            $(\forall z\in X. z\mathrel{R}b_0 \lor z = b_0)$ and
                            $B_1 = \intervalopenclosedrel{R}{X}{u}{b_0}$.
                        Thus $u\mathrel{R}x$ and $(x\mathrel{R}b_0 \lor x = b_0)$ by \cref{intervalopenclosed_rel}.
                        Thus $\intervalclosedopenrel{R}{X}{x}{b} \subseteq B_1$
                            by \cref{intervalclosedopen_rel,transitive,intervalopenclosed_rel,subseteq}.
                        Thus there exists $c\in Y$ such that $x\mathrel{R}c$ and
                            $\intervalclosedopenrel{R}{X}{x}{c} \subseteq B_1$ by \cref{intervalclosed_rel}.
                    \end{byCase}
                \end{subproof}
                Take $c\in Y$ such that $x\mathrel{R}c$ and $\intervalclosedopenrel{R}{X}{x}{c} \subseteq V$.
                \begin{byCase}
                \caseOf{there exists $y\in X$ such that $y$ is an immediate successor of $x$ in $X$ with respect to $R$.}
                    Take $y\in X$ such that $y$ is an immediate successor of $x$ in $X$ with respect to $R$.
                    Show that $y\in Y$.
                    \begin{subproof}
                        Show that $y\mathrel{R}b$ or $y = b$.
                        \begin{subproof}
                            Thus $y\mathrel{R}b$ or $y = b$ or $b\mathrel{R}y$ by \cref{connex}.
                            \begin{byCase}
                            \caseOf{$b\mathrel{R}y$.}
                                Then $b\in\intervalclosedrel{R}{X}{x}{y}$ by \cref{intervalclosed_rel}.
                                Thus $b\in\{x,y\}$.
                                Thus $b = x$ or $b = y$ by \cref{upair_elim}.
                                Thus $y = b$.
                            \caseOf{$y\mathrel{R}b$ or $y = b$.}
                            \end{byCase}
                        \end{subproof}
                        Thus $y\in Y$ by \cref{immediate_successor,intervalclosed_rel,transitive}.
                    \end{subproof}
                    Take $U_1\in A$ such that $y\in U_1$
                        by \cref{open_covering,covering,subseteq,unions_iff}.
                    Let $B = \{U,U_1\}$.
                    Hence $B\subseteq A$ and $B$ is finite by \cref{upair_elim,subseteq,pair_is_finite}.
                    Show that for all $z$ such that $z\in\intervalclosedrel{R}{X}{x}{y}$ we have $z\in\unions{B}$.
                    \begin{subproof}
                        Fix $z$ such that $z\in\intervalclosedrel{R}{X}{x}{y}$.
                        Thus $z = x$ or $z = y$ by \cref{immediate_successor_interval,upair_elim}.
                        \begin{byCase}
                        \caseOf{$z = x$.}
                            Thus $z\in U$.
                            Thus $z\in\unions{B}$ by \cref{upair_intro_left,unions_iff}.
                        \caseOf{$z = y$.}
                            Thus $z\in U_1$.
                            Thus $z\in\unions{B}$ by \cref{upair_intro_right,unions_iff}.
                        \end{byCase}
                    \end{subproof}
                    Therefore $B$ is a covering of $\intervalclosedrel{R}{X}{x}{y}$ by \cref{subseteq,covering}.
                \caseOf{it is wrong that there exists $y\in X$ such that $y$ is an immediate successor of $x$ in $X$ with respect to $R$.}
                    Therefore it is wrong that $c$ is an immediate successor of $x$ in $X$ with respect to $R$
                        by \cref{subseteq}.
                    Take $y\in X$ such that $x\mathrel{R}y$ and $y\mathrel{R}c$ by \cref{no_immediate_successor_between}.
                    Hence $y\in Y$ by \cref{convex_in}.
                    Let $B = \{U\}$.
                    Thus $B\subseteq A$ and $B$ is finite
                        by \cref{singleton_subset_intro,pair_is_finite,upair_intro_left,subseteq_finite_is_finite}.
                    For all $z$ such that $z\in\intervalclosedrel{R}{X}{x}{y}$ we have
                        $z\in X$ and $(x\mathrel{R}z \lor z = x)$ and $(z\mathrel{R}y \lor z = y)$
                        by \cref{intervalclosed_rel}.
                    For all $z$ such that $z\in\intervalclosedrel{R}{X}{x}{y}$ we have $z\mathrel{R}c$
                        by \hyperref[transitive]{transitivity}.
                    Thus for all $z$ such that $z\in\intervalclosedrel{R}{X}{x}{y}$ we have $z\in\unions{B}$
                        by \cref{intervalclosedopen_rel,subseteq,convex_in,inter_intro,singleton_intro,unions_iff}.
                    Therefore $B$ is a covering of $\intervalclosedrel{R}{X}{x}{y}$ by \cref{subseteq,covering}.
                \end{byCase}
            \end{subproof}

            Let $C = \{y\in Y \mid \text{$a\mathrel{R}y$ and there exists $B\subseteq A$ such that
                $B$ is finite and $B$ is a covering of $\intervalclosedrel{R}{X}{a}{y}$}\}$.
            Now $a\in Y$.
            Take $y\in Y$ such that $a\mathrel{R}y$ and
                there exists $B$ such that $B\subseteq A$ and $B$ is finite and
                    $B$ is a covering of $\intervalclosedrel{R}{X}{a}{y}$.
            Hence $y\in C$.
            Therefore $C$ is inhabited.
            Hence $C\subseteq X$ by \cref{intervalclosed_rel,subseteq}.
            For all $y\in C$ we have $y\mathrel{R}b$ or $y = b$ by \cref{intervalclosed_rel}.
            Therefore $b$ is an upper bound of $C$ in $X$ with respect to $R$ by \cref{upper_bound_rel}.
            Take $c$ such that $c$ is a least upper bound of $C$ in $X$ with respect to $R$
                by \cref{least_upper_bound_property}.
            Show that $a\mathrel{R}c$.
            \begin{subproof}
                Take $y\in C$ such that $a\mathrel{R}y$ by \cref{nonempty_inhabited}.
                Suppose not.
                \begin{byCase}
                \caseOf{$c = a$.}
                    Thus $y\mathrel{R}c$ or $y = c$ by \cref{least_upper_bound_rel,upper_bound_rel}.
                    Thus $y\mathrel{R}c$ by \cref{least_upper_bound_rel,upper_bound_rel,asymmetric}.
                    Thus $y\mathrel{R}a$.
                    Contradiction by \cref{asymmetric}.
                \caseOf{$c\neq a$.}
                    Thus $c\mathrel{R}a$ or $a\mathrel{R}c$ by \cref{least_upper_bound_rel,upper_bound_rel,connex}.
                    Thus $c\mathrel{R}a$ by \cref{least_upper_bound_rel,upper_bound_rel,connex}.
                    Thus $c\mathrel{R}y$ by \hyperref[transitive]{transitivity}.
                    Thus $y\mathrel{R}c$ or $y = c$ by \cref{least_upper_bound_rel,upper_bound_rel}.
                    Contradiction by \cref{asymmetric}.
                \end{byCase}
            \end{subproof}
            Thus $c\in X$ and $c\in Y$ by \cref{least_upper_bound_rel,upper_bound_rel,intervalclosed_rel}.

            Show that $c\in C$.
            \begin{subproof}
                Take $U\in A$ such that $c\in U$ by \cref{open_covering,covering,subseteq,unions_iff}.
                Take $V\in T$ such that $U = Y\inter V$ by \cref{open_covering,open_in,subspace_topology}.
                Thus $c\in V$ by \cref{inter_elim_right}.
                Take $B_1\in\orderBasis{R}{X}$ such that $c\in B_1$ and $B_1\subseteq V$
                    by \cref{topology_generated_by_basis}.
                Show that there exists $d\in Y$ such that $d\mathrel{R}c$ and
                    $\intervalopenclosedrel{R}{X}{d}{c} \subseteq V$.
                \begin{subproof}
                    $B_1\in\pow{X}$ and
                        $((\exists u, v\in X. u\mathrel{R}v \land B_1 = \intervalopenrel{R}{X}{u}{v}) \lor
                        (\exists v\in X. \exists a_0\in X. (\forall z\in X. a_0\mathrel{R}z \lor a_0 = z) \land
                            B_1 = \intervalclosedopenrel{R}{X}{a_0}{v}) \lor
                        (\exists u\in X. \exists b_0\in X. (\forall z\in X. z\mathrel{R}b_0 \lor z = b_0) \land
                            B_1 = \intervalopenclosedrel{R}{X}{u}{b_0}))$
                        by \cref{order_basis}.
                    \begin{byCase}
                    \caseOf{$\exists u, v\in X. u\mathrel{R}v \land B_1 = \intervalopenrel{R}{X}{u}{v}$.}
                        Take $u, v\in X$ such that $u\mathrel{R}v$ and $B_1 = \intervalopenrel{R}{X}{u}{v}$.
                        Thus $u\mathrel{R}c$ and $c\mathrel{R}v$ by \cref{intervalopen_rel}.
                        Show that there exists $d\in Y$ such that $d\mathrel{R}c$ and
                            $\intervalopenclosedrel{R}{X}{d}{c} \subseteq B_1$.
                        \begin{subproof}
                            \begin{byCase}
                            \caseOf{$u\in Y$.}
                                Let $d = u$.
                                Thus $\intervalopenclosedrel{R}{X}{d}{c} \subseteq B_1$
                                    by \cref{intervalopenclosed_rel,transitive,intervalopen_rel,subseteq}.
                            \caseOf{it is wrong that $u\in Y$.}
                                Thus $u\mathrel{R}a$ by \cref{connex,least_upper_bound_rel,transitive,intervalclosed_rel}.
                                Let $d = a$.
                                Thus $\intervalopenclosedrel{R}{X}{d}{c} \subseteq B_1$
                                    by \cref{intervalopenclosed_rel,transitive,intervalopen_rel,subseteq}.
                            \end{byCase}
                        \end{subproof}
                    \caseOf{$\exists v\in X. \exists a_0\in X. (\forall z\in X. a_0\mathrel{R}z \lor a_0 = z) \land
                        B_1 = \intervalclosedopenrel{R}{X}{a_0}{v}$.}
                        Take $a_0, v\in X$ such that
                            $(\forall z\in X. a_0\mathrel{R}z \lor a_0 = z)$ and
                            $B_1 = \intervalclosedopenrel{R}{X}{a_0}{v}$.
                        Thus $c\in X$ and $(a_0\mathrel{R}c \lor c = a_0)$ and $c\mathrel{R}v$
                            by \cref{intervalclosedopen_rel}.
                        Let $d = a$.
                        Thus $\intervalopenclosedrel{R}{X}{d}{c} \subseteq B_1$
                            by \cref{intervalopenclosed_rel,smallest_element,transitive,intervalclosedopen_rel,subseteq}.
                    \caseOf{$\exists u\in X. \exists b_0\in X. (\forall z\in X. z\mathrel{R}b_0 \lor z = b_0) \land
                        B_1 = \intervalopenclosedrel{R}{X}{u}{b_0}$.}
                        Take $u, b_0\in X$ such that
                            $(\forall z\in X. z\mathrel{R}b_0 \lor z = b_0)$ and
                            $B_1 = \intervalopenclosedrel{R}{X}{u}{b_0}$.
                        Thus $u\mathrel{R}c$ and $(c\mathrel{R}b_0 \lor c = b_0)$ by \cref{intervalopenclosed_rel}.
                        Show that there exists $d\in Y$ such that $d\mathrel{R}c$ and
                            $\intervalopenclosedrel{R}{X}{d}{c} \subseteq B_1$.
                        \begin{subproof}
                            \begin{byCase}
                            \caseOf{$u\in Y$.}
                                Let $d = u$.
                                Thus $\intervalopenclosedrel{R}{X}{d}{c} \subseteq B_1$
                                    by \cref{intervalopenclosed_rel,transitive,intervalopenclosed_rel,subseteq}.
                            \caseOf{it is wrong that $u\in Y$.}
                                Thus $u\mathrel{R}a$ by \cref{connex,least_upper_bound_rel,transitive,intervalclosed_rel}.
                                Let $d = a$.
                                Thus $\intervalopenclosedrel{R}{X}{d}{c} \subseteq B_1$
                                    by \cref{intervalopenclosed_rel,transitive,intervalopenclosed_rel,subseteq}.
                            \end{byCase}
                        \end{subproof}
                    \end{byCase}
                \end{subproof}
                Take $d\in Y$ such that $d\mathrel{R}c$ and $\intervalopenclosedrel{R}{X}{d}{c} \subseteq V$.
                Suppose not.
                Show that there exists $z\in C$ such that $d\mathrel{R}z$.
                \begin{subproof}
                    Suppose not.
                    For all $y\in C$ we have $y\mathrel{R}d$ or $y = d$ by \cref{subseteq,connex}.
                    Thus $c\mathrel{R}d$ or $c = d$ by \cref{subseteq,upper_bound_rel,least_upper_bound_rel}.
                    Contradiction by \cref{asymmetric}.
                \end{subproof}
                Take $z\in C$ such that $d\mathrel{R}z$.
                Thus $z\mathrel{R}c$ by \cref{least_upper_bound_rel,upper_bound_rel}.
                Take $B_0$ such that $B_0\subseteq A$ and $B_0$ is finite and
                    $B_0$ is a covering of $\intervalclosedrel{R}{X}{a}{z}$.
                Let $B = B_0\union\{U\}$.
                Hence $B$ is finite by \cref{pair_is_finite,upair_intro_left,singleton_subset_intro,subseteq_finite_is_finite,union_of_finite_is_finite}.
                Show that for all $w$ such that $w\in\intervalclosedrel{R}{X}{a}{c}$ we have $w\in\unions{B}$.
                \begin{subproof}
                    Fix $w$ such that $w\in\intervalclosedrel{R}{X}{a}{c}$.
                    Thus $w\in X$ and $(a\mathrel{R}w \lor w = a)$ and $(w\mathrel{R}c \lor w = c)$
                        by \cref{intervalclosed_rel}.
                    Thus $w\mathrel{R}z$ or $w = z$ or $z\mathrel{R}w$ by \cref{connex}.
                    \begin{byCase}
                    \caseOf{$w\mathrel{R}z$ or $w = z$.}
                        Thus $w\in\unions{B}$ by \cref{subseteq,intervalclosed_rel,covering,unions_iff,union_intro_left}.
                    \caseOf{$z\mathrel{R}w$.}
                        Thus $w\in V$ by \cref{transitive,intervalopenclosed_rel,subseteq}.
                        Thus $w\in Y$ by \cref{convex_in}.
                        Thus $w\in Y\inter V$ by \cref{inter_intro}.
                        Thus $w\in U$.
                        Thus $w\in\unions{B}$ by \cref{singleton_intro,union_intro_right,unions_iff}.
                    \end{byCase}
                \end{subproof}
                Therefore $B$ is a covering of $\intervalclosedrel{R}{X}{a}{c}$ by \cref{subseteq,covering}.
                Hence there exists $B$ such that
                    $B\subseteq A$ and $B$ is finite and
                    $B$ is a covering of $\intervalclosedrel{R}{X}{a}{c}$.
                Therefore $c\in C$.
                Contradiction.
            \end{subproof}

            Show that $c = b$.
            \begin{subproof}
                Suppose not.
                Thus $c\mathrel{R}b$ by \cref{least_upper_bound_rel}.
                Take $y,B_1$ such that $y\in Y$ and $c\mathrel{R}y$ and $B_1\subseteq A$ and
                    $B_1$ is finite and $B_1$ is a covering of $\intervalclosedrel{R}{X}{c}{y}$.
                Thus $a\mathrel{R}y$ by \hyperref[transitive]{transitivity}.
                Take $B_0$ such that $B_0\subseteq A$ and $B_0$ is finite and
                    $B_0$ is a covering of $\intervalclosedrel{R}{X}{a}{c}$.
                Let $B = B_0\union B_1$.
                Hence $B$ is finite by \cref{union_of_finite_is_finite}.
                Show that for all $w$ such that $w\in\intervalclosedrel{R}{X}{a}{y}$ we have $w\in\unions{B}$.
                \begin{subproof}
                    Fix $w$ such that $w\in\intervalclosedrel{R}{X}{a}{y}$.
                    Thus $w\in X$ and $(a\mathrel{R}w \lor w = a)$ and $(w\mathrel{R}y \lor w = y)$
                        by \cref{intervalclosed_rel}.
                    Thus $w\mathrel{R}c$ or $w = c$ or $c\mathrel{R}w$ by \cref{connex}.
                \begin{byCase}
                \caseOf{$w\mathrel{R}c$ or $w = c$.}
                    Thus $w\in\unions{B}$ by \cref{intervalclosed_rel,covering,subseteq,unions_iff,union_intro_left}.
                \caseOf{$c\mathrel{R}w$.}
                    Thus $w\in\unions{B}$ by \cref{intervalclosed_rel,covering,subseteq,unions_iff,union_intro_right}.
                \end{byCase}
                \end{subproof}
                Therefore $B$ is a covering of $\intervalclosedrel{R}{X}{a}{y}$ by \cref{subseteq,covering}.
                Hence $y\in C$.
                Therefore $y\mathrel{R}c$ or $y = c$ by \cref{least_upper_bound_rel,upper_bound_rel}.
                Contradiction by \cref{asymmetric}.
            \end{subproof}
            Take $B$ such that $B\subseteq A$ and $B$ is finite and
                $B$ is a covering of $\intervalclosedrel{R}{X}{a}{b}$.
        \end{byCase}
    \end{subproof}
    Therefore $Y$ is compact with respect to $T_Y$ by \cref{compact}.
\end{proof}
% from §27 Item 2: closed interval in reals is compact
\begin{corollary}\label{real_closed_interval_compact}
    Suppose $a\in\reals$.
    Suppose $b\in\reals$.
    Suppose $a < b$ or $a = b$.
    Let $R = \realOrderRel$.
    Let $T = \realOrderTopology$.
    Let $Y = \intervalclosedrel{R}{\reals}{a}{b}$.
    Let $T_Y = \subspaceTopology{T}{Y}$.
    Then $Y$ is compact with respect to $T_Y$.
\end{corollary}
\begin{proof}
    Let $X = \reals$.
    $R$ is a simple order relation on $X$ by \cref{real_order_relation_simple}.
    Hence $X$ has the least upper bound property with respect to $R$ and $X$ has more than one element
        by \cref{reals_linear_continuum,linear_continuum}.
    Therefore $T$ is the order topology on $X$ with respect to $R$ by \cref{real_order_topology,order_topology}.
    Hence $Y$ is compact with respect to $T_Y$ by \cref{real_order_relation_intro,ordered_closed_interval_compact}.
\end{proof}

% from §27 helper lemma 3: four-term real sum bound
\begin{lemma}\label{real_four_sum_bound}
    Suppose $r_1,r_2,r_3,r_4,M_0\in\reals$.
    Suppose $r_1 \leq 1$ and $r_4 \leq 1$ and $r_2 \leq M_0$ and $r_3 \leq M_0$.
    Then $(r_1 + r_2) + (r_3 + r_4) \leq (M_0 + 1) + (M_0 + 1)$.
\end{lemma}
\begin{proof}
    $1\in\reals$ and $r_2 + 1\in\reals$ and $r_3 + 1\in\reals$ and $1 + r_2\in\reals$ and $1 + r_3\in\reals$ and
        $M_0 + 1\in\reals$ by \cref{real_naturals_subset_reals,real_one_in_naturals,subseteq,real_add_closed}.
    $r_1 + r_2\in\reals$ and $r_3 + r_4\in\reals$ and $r_4 + r_3\in\reals$ and
        $(r_1 + r_2) + (r_3 + r_4)\in\reals$ by \cref{real_add_closed}.
    $(M_0 + 1) + (r_3 + r_4)\in\reals$ and $(r_3 + r_4) + (M_0 + 1)\in\reals$ and
        $(M_0 + 1) + (M_0 + 1)\in\reals$ by \cref{real_add_closed}.
    Hence $r_1 + r_2 \leq M_0 + 1$ and $r_3 + r_4 \leq M_0 + 1$
        by \cref{real_leq_add,real_add_comm,real_leq_trans}.
    Therefore $(r_1 + r_2) + (r_3 + r_4) \leq (M_0 + 1) + (M_0 + 1)$
        by \cref{real_leq_add,real_add_comm,real_leq_trans}.
\end{proof}

% from §27 helper lemma 4: compact standard closed interval
\begin{lemma}\label{real_intervalclosed_compact_standard}
    Suppose $a\in\reals$.
    Suppose $b\in\reals$.
    Suppose $a \leq b$.
    Let $I = \intervalclosed{a}{b}$.
    Then $I$ is compact with respect to $\subspaceTopology{\standardTopologyR}{I}$.
\end{lemma}
\begin{proof}
    Let $R = \realOrderRel$.
    Hence $\intervalclosedrel{R}{\reals}{a}{b} = I$
        by \cref{intervalclosed_rel,real_order_relation_elim,intervalclosed,subseteq,real_order_relation_intro,subseteq_antisymmetric}.
    Hence $I$ is compact with respect to $\subspaceTopology{\standardTopologyR}{I}$
        by \cref{real_closed_interval_compact,real_order_topology_eq_standard}.
\end{proof}

% from §27 helper lemma 5: indexed projection map is continuous
\begin{lemma}\label{projection_map_indexed_continuous}
    Assume $X$ is a function.
    Assume $T$ is a function.
    Assume $\dom{T} = \dom{X}$.
    Assume $\dom{X}$ is inhabited.
    Assume that for all $\alpha\in\dom{X}$ we have $\apply{T}{\alpha}$ is a topology on $\apply{X}{\alpha}$.
    Let $T_0 = \productTopologyIndexed{T}{X}$.
    Then for all $\alpha\in\dom{X}$ we have $\piIndex{X}{\alpha}$ is continuous from $\productFamily{X}$ to $\apply{X}{\alpha}$
        with respect to $T_0$ and $\apply{T}{\alpha}$.
\end{lemma}
\begin{proof}
    Hence $T_0$ is a topology on $\productFamily{X}$
        by \cref{product_subbasis_is_subbasis,product_subbasis_finite_intersections_basis,basis_topology_is_topology,basis_for_topology}.
    Let $j = \identity{\productFamily{X}}$.
    $j\in\funs{\productFamily{X}}{\productFamily{X}}$ by \cref{id_elem_funs}.
    $j$ is continuous from $\productFamily{X}$ to $\productFamily{X}$ with respect to $T_0$ and $T_0$
        by \cref{identity_continuous}.
    Therefore for all $\alpha\in\dom{X}$ we have $\piIndex{X}{\alpha}\circ j$ is continuous from $\productFamily{X}$
        to $\apply{X}{\alpha}$ with respect to $T_0$ and $\apply{T}{\alpha}$
        by \cref{continuous_map_into_indexed_product}.
    Fix $\alpha$ such that $\alpha\in\dom{X}$.
    For all $x,z$ we have $x\mathrel{\piIndex{X}{\alpha}\circ j} z$ iff $x\mathrel{\piIndex{X}{\alpha}} z$
        by \cref{circ_iff,id_iff,projection_map_indexed,pair_eq_iff}.
    Hence $\piIndex{X}{\alpha}\circ j = \piIndex{X}{\alpha}$
        by \cref{circ_is_relation,projection_map_indexed_function,id_elem_funs,funs_is_relation,relext}.
    Therefore $\piIndex{X}{\alpha}$ is continuous from $\productFamily{X}$ to $\apply{X}{\alpha}$
        with respect to $T_0$ and $\apply{T}{\alpha}$.
\end{proof}

% from §27 helper definition 2: indexed product compact index
\begin{definition}\label{indexed_product_compact_index}
    A natural number $m$ is indexed-product-compact iff
        for all $Y_1,T_1$ such that $Y_1$ is a function and $T_1$ is a function and
            $\dom{Y_1} = \initialSegment{m}$ and $\dom{T_1} = \initialSegment{m}$,
            if for all $i\in\initialSegment{m}$ we have $\apply{T_1}{i}$ is a topology on $\apply{Y_1}{i}$ then
            if for all $i\in\initialSegment{m}$ we have $\apply{Y_1}{i}$ is compact with respect to $\apply{T_1}{i}$ then
            $\productFamily{Y_1}$ is compact with respect to $\productTopologyIndexed{T_1}{Y_1}$.
\end{definition}

% from §27 helper definition 3: finite indexed product compactness property
\begin{definition}\label{finite_indexed_product_compact_property}
    $\finiteIndexedProductCompactProp{m}$ iff
        for all $Y_1,T_1$ such that
            $Y_1$ is a function and $T_1$ is a function and
            $\dom{Y_1} = \initialSegment{m}$ and $\dom{T_1} = \initialSegment{m}$ and
            for all $i\in\dom{Y_1}$ we have $\apply{T_1}{i}$ is a topology on $\apply{Y_1}{i}$ and
                $\apply{Y_1}{i}$ is compact with respect to $\apply{T_1}{i}$,
        we have $\productFamily{Y_1}$ is compact with respect to $\productTopologyIndexed{T_1}{Y_1}$.
\end{definition}

% from §27 helper lemma 7: continuity preserved under equality
\begin{lemma}\label{continuous_eq_subst}
    Let $f,g$ be functions.
    Suppose $f = g$.
    Suppose $g$ is continuous from $X$ to $Y$ with respect to $T$ and $T_1$.
    Then $f$ is continuous from $X$ to $Y$ with respect to $T$ and $T_1$.
\end{lemma}
\begin{proof}
    Trivial.
\end{proof}

% from §27 helper lemma 8: continuity respects equal codomain/topology
\begin{lemma}\label{continuous_eq_codomain}
    Suppose $Y = Y_2$.
    Suppose $T_1 = T_2$.
    Suppose $f$ is continuous from $X$ to $Y$ with respect to $T$ and $T_1$.
    Then $f$ is continuous from $X$ to $Y_2$ with respect to $T$ and $T_2$.
\end{lemma}
\begin{proof}
    Trivial.
\end{proof}

\begin{lemma}\label{finite_indexed_product_compact_property_succ}
    Assume $k\in\naturalsPlus$.
    Assume $\finiteIndexedProductCompactProp{k}$.
    Then $\finiteIndexedProductCompactProp{k + 1}$.
\end{lemma}
\begin{proof}
    Show that for all $Y_1,T_1$ such that $Y_1$ is a function and $T_1$ is a function and
        $\dom{Y_1} = \initialSegment{k + 1}$ and $\dom{T_1} = \initialSegment{k + 1}$ and
        for all $i\in\dom{Y_1}$ we have $\apply{T_1}{i}$ is a topology on $\apply{Y_1}{i}$ and
            $\apply{Y_1}{i}$ is compact with respect to $\apply{T_1}{i}$,
        we have $\productFamily{Y_1}$ is compact with respect to $\productTopologyIndexed{T_1}{Y_1}$.
    \begin{subproof}
        Fix $Y_1$.
        Fix $T_1$.
        Assume $Y_1$ is a function.
        Assume $T_1$ is a function.
        Assume $\dom{Y_1} = \initialSegment{k + 1}$.
        Assume $\dom{T_1} = \initialSegment{k + 1}$.
        Assume that for all $i\in\dom{Y_1}$ we have $\apply{T_1}{i}$ is a topology on $\apply{Y_1}{i}$ and
            $\apply{Y_1}{i}$ is compact with respect to $\apply{T_1}{i}$.
        Let $J_1 = \initialSegment{k}$.
        Let $J = \initialSegment{k + 1}$.
        For all $i$ such that $i\in J_1$ we have $i\in J$
            by \cref{initial_segment_naturalsplus,real_naturals_subset_reals,subseteq,real_naturals_closed_succ,real_add_ident,real_one_in_naturals,real_zero_lt_one,real_lt_add,real_add_comm,real_lt_subst_left,real_lt_imp_leq,real_leq_trans}.
        Hence $J_1\subseteq J$ by \cref{subseteq}.
        $k\in\reals$ by \cref{real_naturals_subset_reals,subseteq}.
        $k + 1\in\naturalsPlus$ and $k + 1\in\reals$
            by \cref{real_naturals_closed_succ,real_naturals_subset_reals,subseteq}.
        Thus $k < k + 1$
            by \cref{real_zero_lt_one,real_add_ident,real_naturals_subset_reals,real_one_in_naturals,subseteq,real_lt_add,real_add_comm,real_lt_subst_left}.
        Show that $k \neq k + 1$.
        \begin{subproof}
            Suppose not.
            Thus $k = k + 1$.
            Thus $k < k$ by \cref{real_lt_subst_left}.
            Contradiction by \cref{real_lt_irreflexive}.
        \end{subproof}
        Let $Y_0 = \{(i,\apply{Y_1}{i}) \mid i\in J_1\}$.
        Let $T_0 = \{(i,\apply{T_1}{i}) \mid i\in J_1\}$.
        $\dom{Y_0} = J_1$ by \cref{dom_intro,dom_iff,pair_eq_iff,subseteq,subseteq_antisymmetric}.
        Then $\dom{T_0} = J_1$ by \cref{dom_intro,dom_iff,pair_eq_iff,subseteq,subseteq_antisymmetric}.
        $Y_0$ is a relation by \cref{relation}.
        For all $a,b_1,b_2$ such that $(a,b_1)\in Y_0$ and $(a,b_2)\in Y_0$ we have $b_1 = b_2$
            by \cref{pair_eq_iff}.
        Then $Y_0$ is a function by \cref{rightunique}.
        $T_0$ is a relation by \cref{relation}.
        For all $a,b_1,b_2$ such that $(a,b_1)\in T_0$ and $(a,b_2)\in T_0$ we have $b_1 = b_2$
            by \cref{pair_eq_iff}.
        Then $T_0$ is a function by \cref{rightunique}.
        For all $i$ such that $i\in J_1$ we have $\apply{T_0}{i}$ is a topology on $\apply{Y_0}{i}$
            by \cref{subseteq,function_apply_intro}.
        For all $i$ such that $i\in J_1$ we have $\apply{Y_0}{i}$ is compact with respect to $\apply{T_0}{i}$
            by \cref{subseteq,function_apply_intro}.
        Therefore $\productFamily{Y_0}$ is compact with respect to $\productTopologyIndexed{T_0}{Y_0}$
            by \cref{finite_indexed_product_compact_property}.
        Let $X_0 = \productFamily{Y_0}$.
        Let $T_X = \productTopologyIndexed{T_0}{Y_0}$.
        Let $X_1 = \apply{Y_1}{k + 1}$.
        Let $T_{1p} = \apply{T_1}{k + 1}$.
        Then $k + 1\in J$ by \cref{real_naturals_closed_succ,initial_segment_naturalsplus}.
        $T_{1p}$ is a topology on $X_1$ and $X_1$ is compact with respect to $T_{1p}$ by \cref{subseteq}.
        Now $J_1$ is inhabited by \cref{real_one_leq_naturalsplus,real_one_in_naturals,initial_segment_naturalsplus}.
        Hence $T_X$ is a topology on $X_0$
            by \cref{product_subbasis_finite_intersections_basis,basis_for_topology,basis_topology_is_topology}.
        Let $T_P = \productTopology{T_X}{T_{1p}}{X_0}{X_1}$.
        Therefore $X_0\times X_1$ is compact with respect to $T_P$ by \cref{product_compact}.
        Let $g = \{(p,\fst{p}\union\{(k + 1,\snd{p})\}) \mid p\in X_0\times X_1\}$.
        $g$ is a relation by \cref{relation}.
        For all $p,u,v$ such that $(p,u)\in g$ and $(p,v)\in g$ we have $u = v$ by \cref{pair_eq_iff}.
        Hence $g$ is a function by \cref{rightunique}.
        Thus $\dom{g} = X_0\times X_1$ by \cref{dom_intro,dom_iff,pair_eq_iff,subseteq,subseteq_antisymmetric}.
            Show that for all $p\in\dom{g}$ we have $g(p)\in\productFamily{Y_1}$.
            \begin{subproof}
                Fix $p$ such that $p\in\dom{g}$.
                Take $v$ such that $(p,v)\in g$ by \cref{dom_iff}.
                Take $p_0\in X_0\times X_1$ such that $(p,v) = (p_0,\fst{p_0}\union\{(k + 1,\snd{p_0})\})$.
                Thus $p = p_0$ and $v = \fst{p}\union\{(k + 1,\snd{p})\}$ by \cref{pair_eq_iff}.
                Let $x = \fst{p}$.
                    Let $y = \snd{p}$.
                    Then $x\in X_0$ and $y\in X_1$ by \cref{subseteq,times_proj_elim}.
                    Hence $x\in\funs{\dom{Y_0}}{\unions{\ran{Y_0}}}$ and $x$ is a function and $\dom{x} = \dom{Y_0}$
                        by \cref{subseteq,indexed_product,funs_elim}.
                    Then $x\in\rels{\dom{Y_0}}{\unions{\ran{Y_0}}}$ and $x$ is a relation
                        by \cref{funs_subseteq_rels,subseteq,rels_is_relation}.
                    Show that $v\in\funs{\dom{Y_1}}{\unions{\ran{Y_1}}}$.
                    \begin{subproof}
                            Show that for all $w$ such that $w\in v$ there exist $a,b$ such that $w = (a,b)$.
                            \begin{subproof}
                                Fix $w$ such that $w\in v$.
                                Then $w\in x$ or $w = (k + 1,y)$ by \cref{union_iff,singleton_iff}.
                                \begin{byCase}
                                \caseOf{$w\in x$.}
                                    Take $a,b$ such that $w = (a,b)$ by \cref{relation}.
                                    Hence there exist $a,b$ such that $w = (a,b)$.
                                \caseOf{$w = (k + 1,y)$.}
                                    Hence there exist $a,b$ such that $w = (a,b)$.
                                \end{byCase}
                            \end{subproof}
                            Hence $v$ is a relation by \cref{relation}.
                            Show that for all $a,b_1,b_2$ such that $(a,b_1)\in v$ and $(a,b_2)\in v$ we have $b_1 = b_2$.
                            \begin{subproof}
                                Fix $a$.
                                Fix $b_1$.
                                Fix $b_2$ such that $(a,b_1)\in v$ and $(a,b_2)\in v$.
                                \begin{byCase}
                                \caseOf{$a = k + 1$.}
                                    Thus $(a,b_1)\in x$ or $(a,b_1) = (k + 1,y)$ by \cref{union_iff,singleton_iff}.
                                    Show that $(a,b_1)\notin x$.
                                    \begin{subproof}
                                        Suppose not.
                                        Thus $(a,b_1)\in x$.
                                        Thus $a\in\dom{x}$ by \cref{dom_intro}.
                                        Thus $\dom{x} = \dom{Y_0}$ by \cref{funs_elim}.
                                        Thus $a\in J_1$ by \cref{subseteq}.
                                        Thus $a\in\naturalsPlus$ and $a \leq k$ by \cref{initial_segment_naturalsplus}.
                                        Then $k + 1 \leq k$.
                                        Hence $k + 1 < k$ or $k + 1 = k$ by \cref{real_leq_split}.
                                        Contradiction by \cref{real_lt_asym,real_lt_subst_left,real_lt_irreflexive}.
                                    \end{subproof}
                                    Hence $(a,b_1) = (k + 1,y)$.
                                    Then $(a,b_2)\in x$ or $(a,b_2) = (k + 1,y)$ by \cref{union_iff,singleton_iff}.
                                    Show that $(a,b_2)\notin x$.
                                    \begin{subproof}
                                        Suppose not.
                                        Thus $(a,b_2)\in x$.
                                        Thus $a\in\dom{x}$ by \cref{dom_intro}.
                                        Thus $\dom{x} = \dom{Y_0}$ by \cref{funs_elim}.
                                        Thus $a\in J_1$ by \cref{subseteq}.
                                        Thus $a\in\naturalsPlus$ and $a \leq k$ by \cref{initial_segment_naturalsplus}.
                                        Then $k + 1 \leq k$.
                                        Hence $k + 1 < k$ or $k + 1 = k$ by \cref{real_leq_split}.
                                        Contradiction by \cref{real_lt_asym,real_lt_subst_left,real_lt_irreflexive}.
                                    \end{subproof}
                                    Hence $(a,b_2) = (k + 1,y)$.
                                    Then $b_1 = y$ and $b_2 = y$ by \cref{pair_eq_iff}.
                                    Hence $b_1 = b_2$.
                                \caseOf{$a \neq k + 1$.}
                                    Then $(a,b_1)\in x$ by \cref{union_iff,singleton_iff,pair_eq_iff}.
                                    Then $(a,b_2)\in x$ by \cref{union_iff,singleton_iff,pair_eq_iff}.
                                    Thus $b_1 = b_2$ by \cref{rightunique,relation}.
                                \end{byCase}
                            \end{subproof}
                            Then $v$ is a function by \cref{rightunique}.
                            For all $i$ such that $i\in\dom{Y_1}$ we have $i\in\dom{v}$
                                by \cref{subseteq,union_iff,singleton_iff,dom_intro,real_naturals_leq_succ_elim,initial_segment_naturalsplus,indexed_product,funs_elim,dom_iff}.
                            Hence $\dom{Y_1}\subseteq\dom{v}$ by \cref{subseteq}.
                            Show that for all $i$ such that $i\in\dom{v}$ we have $i\in\dom{Y_1}$.
                            \begin{subproof}
                                Fix $i$ such that $i\in\dom{v}$.
                                Take $b$ such that $(i,b)\in v$ by \cref{dom_iff}.
                                Thus $(i,b)\in x$ or $(i,b) = (k + 1,y)$ by \cref{union_iff,singleton_iff}.
                                \begin{byCase}
                                \caseOf{$(i,b) = (k + 1,y)$.}
                                    Thus $i = k + 1$ by \cref{pair_eq_iff}.
                                    Thus $i\in J$.
                                    Thus $i\in\dom{Y_1}$ by \cref{subseteq}.
                                \caseOf{$(i,b)\in x$.}
                                    Thus $i\in\dom{x}$ by \cref{dom_intro}.
                                    Thus $i\in J_1$ by \cref{indexed_product,funs_elim}.
                                    Thus $i\in J$ by \cref{subseteq,initial_segment_naturalsplus}.
                                    Thus $i\in\dom{Y_1}$ by \cref{subseteq}.
                                \end{byCase}
                            \end{subproof}
                            Therefore $\dom{v}\subseteq\dom{Y_1}$ by \cref{subseteq}.
                            Hence $\dom{v} = \dom{Y_1}$ by \cref{subseteq_antisymmetric}.
                        Show that for all $i\in\dom{Y_1}$ we have $v(i)\in\unions{\ran{Y_1}}$.
                        \begin{subproof}
                            Fix $i$ such that $i\in\dom{Y_1}$.
                            Thus $i\in J$ by \cref{subseteq}.
                            \begin{byCase}
                            \caseOf{$i = k + 1$.}
                                Then $(k + 1,y)\in v$ by \cref{union_iff,singleton_iff}.
                                Hence $v(k + 1) = y$ by \cref{function_apply_intro}.
                                Thus $(k + 1,X_1)\in Y_1$ by \cref{function_apply_elim,subseteq}.
                                Thus $X_1\in\ran{Y_1}$ by \cref{ran_intro}.
                                Thus $y\in\unions{\ran{Y_1}}$ by \cref{unions_iff}.
                                Therefore $v(i)\in\unions{\ran{Y_1}}$.
                            \caseOf{$i \neq k + 1$.}
                                Thus $i\in J_1$ by \cref{real_naturals_leq_succ_elim,initial_segment_naturalsplus}.
                                Thus $i\in\dom{x}$ by \cref{subseteq,indexed_product,funs_elim}.
                                Thus $x(i)\in\apply{Y_0}{i}$ by \cref{indexed_product}.
                                Thus $(i,\apply{Y_1}{i})\in Y_0$.
                                Then $x(i)\in\apply{Y_1}{i}$ by \cref{function_apply_intro}.
                                Thus $(i,\apply{Y_1}{i})\in Y_1$ by \cref{function_apply_elim}.
                                Thus $\apply{Y_1}{i}\in\ran{Y_1}$ by \cref{ran_intro}.
                                Thus $x(i)\in\unions{\ran{Y_1}}$ by \cref{unions_iff}.
                                Thus $(i,x(i))\in v$ by \cref{function_apply_elim,union_iff}.
                                Thus $v(i) = x(i)$ by \cref{function_apply_intro}.
                                Hence $v(i)\in\unions{\ran{Y_1}}$.
                            \end{byCase}
                        \end{subproof}
                        Therefore $v\in\funs{\dom{Y_1}}{\unions{\ran{Y_1}}}$ by \cref{funs_intro}.
                    \end{subproof}
                    Then $v$ is right-unique and $v$ is a relation by \cref{funs,rels_is_relation}.
                    Show that for all $i\in\dom{Y_1}$ we have $v(i)\in\apply{Y_1}{i}$.
                    \begin{subproof}
                        Fix $i$ such that $i\in\dom{Y_1}$.
                        Thus $i\in J$ by \cref{subseteq}.
                        \begin{byCase}
                        \caseOf{$i = k + 1$.}
                            Thus $(k + 1,y)\in v$ by \cref{union_iff,singleton_iff}.
                            Thus $v(k + 1) = y$ by \cref{function_apply_intro}.
                            Thus $y\in X_1$ by \cref{subseteq}.
                            Hence $v(i)\in\apply{Y_1}{i}$.
                        \caseOf{$i \neq k + 1$.}
                            Thus $i\in J_1$ by \cref{real_naturals_leq_succ_elim,initial_segment_naturalsplus}.
                            Thus $x(i)\in\apply{Y_0}{i}$ by \cref{indexed_product}.
                            Then $x(i)\in\apply{Y_1}{i}$ by \cref{function_apply_intro}.
                            Thus $(i,x(i))\in v$ by \cref{function_apply_elim,union_iff}.
                            Thus $v(i) = x(i)$ by \cref{function_apply_intro}.
                            Then $v(i)\in\apply{Y_1}{i}$.
                        \end{byCase}
                    \end{subproof}
                    Therefore $v\in\productFamily{Y_1}$ by \cref{indexed_product}.
                Hence $g(p)\in\productFamily{Y_1}$ by \cref{function_apply_intro}.
            \end{subproof}
            Therefore $g\in\funs{X_0\times X_1}{\productFamily{Y_1}}$ by \cref{funs_intro}.
        Then $g$ is right-unique and $g$ is a relation by \cref{funs,rels_is_relation}.
        Show that for all $i\in J$ we have $\piIndex{Y_1}{i}\circ g$ is continuous from $X_0\times X_1$
            to $\apply{Y_1}{i}$ with respect to $T_P$ and $\apply{T_1}{i}$.
        \begin{subproof}
            Fix $i$ such that $i\in J$.
            \begin{byCase}
            \caseOf{$i = k + 1$.}
                    Let $h = \piTwo{X_0}{X_1}$.
                    $h$ is continuous from $X_0\times X_1$ to $X_1$ with respect to $T_P$ and $T_{1p}$
                        by \cref{projection_two_continuous}.
                    Show that for all $p\in X_0\times X_1$ we have $\apply{\piIndex{Y_1}{i}\circ g}{p} = \apply{h}{p}$.
                    \begin{subproof}
                        Fix $p$ such that $p\in X_0\times X_1$.
                        Thus $g$ is a function to $\productFamily{Y_1}$ and $g$ is a function
                            by \cref{funs_elim,funs_is_function}.
                        Thus $g(p)\in\productFamily{Y_1}$
                            by \cref{funs,function_apply_elim,ran_intro,function_to_ran,subseteq}.
                        Thus $(p,\fst{p}\union\{(k + 1,\snd{p})\})\in g$.
                        Thus $g(p) = \fst{p}\union\{(k + 1,\snd{p})\}$ by \cref{function_apply_intro}.
                        Thus $(i,\snd{p})\in g(p)$ by \cref{union_iff,singleton_iff}.
                        Thus $g(p)\in\funs{\dom{Y_1}}{\unions{\ran{Y_1}}}$
                            and $g(p)$ is a function by \cref{indexed_product,funs_is_function}.
                        Thus $\apply{g(p)}{i} = \snd{p}$ by \cref{function_apply_intro}.
                        Thus $(g(p),\snd{p})\in\piIndex{Y_1}{i}$ by \cref{pair_eq_iff,projection_map_indexed}.
                        Thus $(p,\snd{p})\in\piIndex{Y_1}{i}\circ g$ by \cref{circ_iff}.
                        Thus $\piIndex{Y_1}{i}\in\funs{\productFamily{Y_1}}{\apply{Y_1}{i}}$
                            and $\piIndex{Y_1}{i}$ is a function by \cref{projection_map_indexed_function,funs_is_function}.
                        Thus $\piIndex{Y_1}{i}\circ g$ is a function by \cref{function_circ}.
                        Thus $\apply{\piIndex{Y_1}{i}\circ g}{p} = \snd{p}$ by \cref{function_apply_intro}.
                        Thus $h\in\funs{X_0\times X_1}{X_1}$ and $h$ is a function by \cref{continuous,funs_is_function}.
                        Take $x,y$ such that $p = (x,y)$ by \cref{times_elem_is_tuple}.
                        Thus $(x,y)\in X_0\times X_1$.
                        Thus $x\in X_0$ and $y\in X_1$ by \cref{times_tuple_elim}.
                        Thus $\snd{p} = y$ by \cref{snd_eq}.
                        Thus $(p,\snd{p})\in h$ by \cref{pair_eq_iff,projection_second}.
                        Thus $\apply{h}{p} = \snd{p}$ by \cref{function_apply_intro}.
                        Thus $\apply{\piIndex{Y_1}{i}\circ g}{p} = \apply{h}{p}$.
                    \end{subproof}
                    $\piIndex{Y_1}{i}\circ g$ is a function and $\dom{\piIndex{Y_1}{i}\circ g} = X_0\times X_1$
                        by \cref{projection_map_indexed_function,funs_is_function,function_circ,dom_iff,circ_iff,dom_intro,funs,subseteq,function_apply_elim,ran_intro,funs_elim,function_to_ran,projection_map_indexed,subseteq_antisymmetric}.
                    $h$ is a function and $\dom{h} = X_0\times X_1$ by \cref{continuous,funs_is_function,funs_elim}.
                    Hence $\piIndex{Y_1}{i}\circ g = h$ by \cref{funs_elim,subseteq,fun_subseteq,subseteq_antisymmetric}.
                    Thus $\piIndex{Y_1}{i}\circ g$ is continuous from $X_0\times X_1$ to $\apply{Y_1}{i}$ with respect to $T_P$ and $\apply{T_1}{i}$
                        by \cref{continuous_eq_subst}.
                \caseOf{$i \neq k + 1$.}
                    Then $i\in J_1$ by \cref{real_naturals_leq_succ_elim,initial_segment_naturalsplus}.
                    Let $h_0 = \piOne{X_0}{X_1}$.
                    $h_0$ is continuous from $X_0\times X_1$ to $X_0$ with respect to $T_P$ and $T_X$
                        by \cref{projection_one_continuous}.
                    $\piIndex{Y_0}{i}$ is continuous from $X_0$ to $\apply{Y_0}{i}$ with respect to $T_X$ and $\apply{T_0}{i}$
                        by \cref{projection_map_indexed_continuous,subseteq,subseteq_antisymmetric}.
                    Let $h = \piIndex{Y_0}{i}\circ h_0$.
                    $h$ is continuous from $X_0\times X_1$ to $\apply{Y_0}{i}$ with respect to $T_P$ and $\apply{T_0}{i}$
                        by \cref{continuous_composite}.
                    Show that for all $p\in X_0\times X_1$ we have $\apply{\piIndex{Y_1}{i}\circ g}{p} = \apply{h}{p}$.
                    \begin{subproof}
                            Fix $p$ such that $p\in X_0\times X_1$.
                            Thus $(p,\fst{p}\union\{(k + 1,\snd{p})\})\in g$.
                            Thus $g(p) = \fst{p}\union\{(k + 1,\snd{p})\}$ by \cref{function_apply_intro}.
                            Let $x = \fst{p}$.
                            Take $a,b$ such that $p = (a,b)$ and $a\in X_0$ and $b\in X_1$ by \cref{times_elem_is_tuple}.
                            Thus $\fst{p} = a$ by \cref{fst_eq}.
                            Thus $x = a$.
                            Thus $x\in X_0$.
                            Thus $(p,x)\in h_0$ by \cref{pair_eq_iff,projection_first}.
                            Thus $(x,\apply{x}{i})\in\piIndex{Y_0}{i}$ by \cref{projection_map_indexed}.
                            Thus $(p,\apply{x}{i})\in\piIndex{Y_0}{i}\circ h_0$ by \cref{circ_iff}.
                            Thus $h\in\funs{X_0\times X_1}{\apply{Y_0}{i}}$ and $h$ is a function
                                by \cref{continuous,funs_is_function}.
                            Thus $\apply{h}{p} = \apply{x}{i}$ by \cref{function_apply_intro}.
                            Thus $g$ is a function to $\productFamily{Y_1}$ and $g$ is a function
                                by \cref{funs_elim,funs_is_function}.
                            Thus $g(p)\in\productFamily{Y_1}$
                                by \cref{funs,function_apply_elim,ran_intro,function_to_ran,subseteq}.
                            Thus $g(p)\in\funs{\dom{Y_1}}{\unions{\ran{Y_1}}}$
                                and $g(p)$ is a function by \cref{indexed_product,funs_is_function}.
                            Thus $x\in\funs{\dom{Y_0}}{\unions{\ran{Y_0}}}$ and $x$ is a function
                                by \cref{indexed_product,funs_is_function}.
                            Thus $i\in\dom{x}$ by \cref{funs_elim,subseteq}.
                            Thus $(i,\apply{x}{i})\in x$ by \cref{function_apply_elim}.
                            Thus $(i,\apply{x}{i})\in g(p)$ by \cref{union_iff}.
                            Thus $\apply{g(p)}{i} = \apply{x}{i}$ by \cref{function_apply_intro}.
                            Thus $(g(p),\apply{x}{i})\in\piIndex{Y_1}{i}$ by \cref{projection_map_indexed}.
                            Thus $(p,\apply{x}{i})\in\piIndex{Y_1}{i}\circ g$ by \cref{circ_iff}.
                            Thus $\piIndex{Y_1}{i}\in\funs{\productFamily{Y_1}}{\apply{Y_1}{i}}$ and
                                $\piIndex{Y_1}{i}$ is a function by \cref{projection_map_indexed_function,funs_is_function}.
                            Thus $\piIndex{Y_1}{i}\circ g$ is a function by \cref{function_circ}.
                            Thus $\apply{\piIndex{Y_1}{i}\circ g}{p} = \apply{x}{i}$ by \cref{function_apply_intro}.
                            Thus $\apply{\piIndex{Y_1}{i}\circ g}{p} = \apply{h}{p}$.
                    \end{subproof}
                    $\piIndex{Y_1}{i}\circ g$ is a function and $\dom{\piIndex{Y_1}{i}\circ g} = X_0\times X_1$
                        by \cref{projection_map_indexed_function,funs_is_function,function_circ,dom_iff,circ_iff,dom_intro,funs,subseteq,function_apply_elim,ran_intro,funs_elim,function_to_ran,projection_map_indexed,subseteq_antisymmetric}.
                    $h$ is a function and $\dom{h} = X_0\times X_1$ by \cref{continuous,funs_is_function,funs_elim}.
                    Hence $\piIndex{Y_1}{i}\circ g = h$ by \cref{funs_elim,subseteq,fun_subseteq,subseteq_antisymmetric}.
                    Hence $h$ is continuous from $X_0\times X_1$ to $\apply{Y_1}{i}$ with respect to $T_P$ and $\apply{T_1}{i}$
                        by \cref{function_apply_intro,continuous_eq_codomain}.
                    Then $h\in\funs{X_0\times X_1}{\apply{Y_1}{i}}$ and $h$ is a function
                        by \cref{continuous,funs_is_function}.
                    Thus $\piIndex{Y_1}{i}\in\funs{\productFamily{Y_1}}{\apply{Y_1}{i}}$ and
                        $\piIndex{Y_1}{i}$ is a function by \cref{projection_map_indexed_function,funs_is_function}.
                    Hence $\piIndex{Y_1}{i}\circ g$ is a function by \cref{function_circ}.
                    Therefore $\piIndex{Y_1}{i}\circ g$ is continuous from $X_0\times X_1$ to $\apply{Y_1}{i}$ with respect to $T_P$ and $\apply{T_1}{i}$
                        by \cref{continuous_eq_subst}.
            \end{byCase}
        \end{subproof}
        $J$ is inhabited.
        $T_P$ is a topology on $X_0\times X_1$.
        Hence $g$ is continuous from $X_0\times X_1$ to $\productFamily{Y_1}$ with respect to $T_P$ and $\productTopologyIndexed{T_1}{Y_1}$
            by \cref{continuous_map_into_indexed_product}.
        Show that for all $z\in\productFamily{Y_1}$ there exists $p\in X_0\times X_1$ such that $g(p) = z$.
        \begin{subproof}
                Fix $z$ such that $z\in\productFamily{Y_1}$.
                Thus $z$ is a function and $\dom{z} = \dom{Y_1}$ by \cref{indexed_product,funs_elim}.
                Let $x = \{(i,z(i)) \mid i\in J_1\}$.
                $x$ is a relation by \cref{relation}.
                        For all $a,b_1,b_2$ such that $(a,b_1)\in x$ and $(a,b_2)\in x$ we have $b_1 = b_2$
                            by \cref{pair_eq_iff}.
                        Thus $x$ is a function by \cref{rightunique}.
                        Thus $\dom{x} = \dom{Y_0}$ by \cref{dom_intro,dom_iff,pair_eq_iff,subseteq,subseteq_antisymmetric}.
                        For all $i$ such that $i\in\dom{Y_0}$ we have $x(i)\in\unions{\ran{Y_0}}$
                            by \cref{subseteq,function_apply_intro,indexed_product,ran_intro,unions_iff}.
                        Thus $x\in\funs{\dom{Y_0}}{\unions{\ran{Y_0}}}$ and $x$ is a function
                            by \cref{funs_intro,funs_is_function}.
                        Show that for all $i\in\dom{Y_0}$ we have $x(i)\in\apply{Y_0}{i}$.
                        \begin{subproof}
                        Fix $i$ such that $i\in\dom{Y_0}$.
                        Thus $i\in J_1$ by \cref{subseteq}.
                        Thus $(i,z(i))\in x$.
                        Thus $x(i) = z(i)$ by \cref{function_apply_intro}.
                        Thus $(i,\apply{Y_1}{i})\in Y_0$.
                        Thus $\apply{Y_0}{i} = \apply{Y_1}{i}$ by \cref{function_apply_intro}.
                        Thus $z(i)\in\apply{Y_1}{i}$ by \cref{indexed_product}.
                        Thus $x(i)\in\apply{Y_0}{i}$.
                    \end{subproof}
                Thus $x\in\productFamily{Y_0}$ by \cref{indexed_product}.
                Hence $x\in X_0$ by \cref{subseteq}.
                Let $y = z(k + 1)$.
                Then $y\in X_1$ by \cref{indexed_product,subseteq}.
                Let $p = (x,y)$.
                Hence $p\in X_0\times X_1$ by \cref{times_tuple_intro}.
                Show that for all $i\in\dom{z}$ we have $\apply{g(p)}{i} = z(i)$.
                    \begin{subproof}
                        Fix $i$ such that $i\in\dom{z}$.
                        Thus $i\in J$ by \cref{subseteq}.
                        Thus $p\in\dom{g}$ by \cref{funs}.
                        Hence $g(p)\in\productFamily{Y_1}$
                            by \cref{funs_elim,funs,function_apply_elim,ran_intro,function_to_ran,subseteq}.
                        Thus $g(p)\in\funs{\dom{Y_1}}{\unions{\ran{Y_1}}}$
                            and $g(p)$ is a function by \cref{indexed_product,funs_is_function}.
                        Thus $(p,\fst{p}\union\{(k + 1,\snd{p})\})\in g$.
                        Thus $g(p) = \fst{p}\union\{(k + 1,\snd{p})\}$ by \cref{function_apply_intro}.
                        Thus $g(p) = x\union\{(k + 1,y)\}$ by \cref{fst_eq,snd_eq}.
                        \begin{byCase}
                        \caseOf{$i = k + 1$.}
                            Thus $(k + 1,y)\in\{(k + 1,y)\}$ by \cref{singleton_intro}.
                            Thus $(k + 1,y)\in g(p)$ by \cref{union_iff}.
                            Thus $\apply{g(p)}{k + 1} = y$ by \cref{function_apply_intro}.
                            Thus $\apply{g(p)}{i} = z(i)$.
                        \caseOf{$i \neq k + 1$.}
                            Thus $i\in J_1$ by \cref{real_naturals_leq_succ_elim,initial_segment_naturalsplus}.
                            Thus $(i,z(i))\in x$.
                            Thus $(i,z(i))\in g(p)$ by \cref{union_iff}.
                            Thus $\apply{g(p)}{i} = z(i)$ by \cref{function_apply_intro}.
                        \end{byCase}
                    \end{subproof}
                Thus $g$ is a function to $\productFamily{Y_1}$ by \cref{funs_elim}.
                Thus $p\in\dom{g}$ by \cref{funs}.
                Hence $g(p)\in\productFamily{Y_1}$
                    by \cref{funs_elim,funs,function_apply_elim,ran_intro,function_to_ran,subseteq}.
                Thus $g(p)$ is a function and $\dom{g(p)} = \dom{Y_1}$ by \cref{indexed_product,funs_elim}.
                Therefore $g(p) = z$ by \cref{subseteq,fun_subseteq,subseteq_antisymmetric}.
                Hence there exists $p\in X_0\times X_1$ such that $g(p) = z$.
            \end{subproof}
        Therefore $g\in\surj{X_0\times X_1}{\productFamily{Y_1}}$ by \cref{surj}.
        Hence $\img{g}{X_0\times X_1} = \productFamily{Y_1}$ by \cref{surj,funs_elim,img_dom,funs_surj_iff}.
        Let $A = \img{g}{X_0\times X_1}$.
        Let $T_Q = \productTopologyIndexed{T_1}{Y_1}$.
        Let $T_A = \subspaceTopology{T_Q}{A}$.
        Then $T_A = \subspaceTopology{T_Q}{\productFamily{Y_1}}$.
        $T_Q$ is a topology on $\productFamily{Y_1}$
            by \cref{product_subbasis_finite_intersections_basis,basis_for_topology,basis_topology_is_topology}.
        For all $U$ such that $U\in T_A$ we have $U\in T_Q$
            by \cref{subspace_topology,topology_on,subseteq,powerset_elim,inter_absorb_supseteq_left}.
        For all $U$ such that $U\in T_Q$ we have $U\in T_A$
            by \cref{topology_on,subseteq,powerset_elim,inter_absorb_supseteq_left,subspace_topology}.
        Therefore $T_A = T_Q$ by \cref{subseteq,subseteq_antisymmetric}.
        Hence $\productFamily{Y_1}$ is compact with respect to $T_Q$ by \cref{continuous_image_compact}.
    \end{subproof}
    Hence $\finiteIndexedProductCompactProp{k + 1}$ by \cref{finite_indexed_product_compact_property}.
\end{proof}

% from §27 helper lemma 6: compact indexed product on an initial segment
\begin{lemma}\label{compact_indexed_product_initial_segment}
    Suppose $n\in\naturalsPlus$.
    Let $J = \initialSegment{n}$.
    Assume $Y$ is a function.
    Assume $T_Y$ is a function.
    Assume $\dom{Y} = J$.
    Assume $\dom{T_Y} = J$.
    Assume that for all $i\in J$ we have $\apply{T_Y}{i}$ is a topology on $\apply{Y}{i}$.
    Assume that for all $i\in J$ we have $\apply{Y}{i}$ is compact with respect to $\apply{T_Y}{i}$.
    Then $\productFamily{Y}$ is compact with respect to $\productTopologyIndexed{T_Y}{Y}$.
\end{lemma}
\begin{proof}
    Let $A = \{m\in\naturalsPlus \mid \finiteIndexedProductCompactProp{m}\}$.
    $A\subseteq\naturalsPlus$ by \cref{subseteq}.
    Show that $1\in A$.
    \begin{subproof}
        It suffices to show that $1\in\naturalsPlus$ and $\finiteIndexedProductCompactProp{1}$.
        Show that for all $Y_1,T_1$ such that $Y_1$ is a function and $T_1$ is a function and
            $\dom{Y_1} = \initialSegment{1}$ and $\dom{T_1} = \initialSegment{1}$ and
            for all $i\in\dom{Y_1}$ we have $\apply{T_1}{i}$ is a topology on $\apply{Y_1}{i}$ and
                $\apply{Y_1}{i}$ is compact with respect to $\apply{T_1}{i}$,
            we have $\productFamily{Y_1}$ is compact with respect to $\productTopologyIndexed{T_1}{Y_1}$.
        \begin{subproof}
            Fix $Y_1$.
            Fix $T_1$.
            Assume $Y_1$ is a function.
            Assume $T_1$ is a function.
            Assume $\dom{Y_1} = \initialSegment{1}$.
            Assume $\dom{T_1} = \initialSegment{1}$.
            Assume that for all $i\in\dom{Y_1}$ we have $\apply{T_1}{i}$ is a topology on $\apply{Y_1}{i}$ and
                $\apply{Y_1}{i}$ is compact with respect to $\apply{T_1}{i}$.
            Let $J_1 = \initialSegment{1}$.
            Then $1\in J_1$ by \cref{real_one_in_naturals,initial_segment_naturalsplus}.
            For all $i\in J_1$ we have $i < 1$ or $i = 1$ by \cref{initial_segment_naturalsplus}.
            For all $i\in J_1$ we have $1\leq i$ by \cref{initial_segment_naturalsplus,real_one_leq_naturalsplus}.
            Thus for all $i\in J_1$ we have $i = 1$.
            Therefore $J_1 = \{1,1\}$ by \cref{subseteq_antisymmetric,subseteq,upair_iff}.
            Let $X = \apply{Y_1}{1}$.
            Let $T_X = \apply{T_1}{1}$.
            $T_X$ is a topology on $X$ and $X$ is compact with respect to $T_X$ by \cref{subseteq}.
            Let $g = \{(x,\{(1,x)\}) \mid x\in X\}$.
            $g$ is a relation by \cref{relation}.
            For all $x,u,v$ such that $(x,u)\in g$ and $(x,v)\in g$ we have $u = v$ by \cref{pair_eq_iff}.
            Hence $g$ is a function by \cref{rightunique}.
            Hence $\dom{g} = X$ by \cref{dom_intro,dom_iff,pair_eq_iff,subseteq,subseteq_antisymmetric}.
            Show that for all $x\in\dom{g}$ we have $g(x)\in\productFamily{Y_1}$.
            \begin{subproof}
                Fix $x$ such that $x\in\dom{g}$.
                Take $y$ such that $(x,y)\in g$ by \cref{dom_iff}.
                Take $x_0\in X$ such that $(x,y) = (x_0,\{(1,x_0)\})$.
                Thus $y = \{(1,x_0)\}$ by \cref{pair_eq_iff}.
                Thus $y$ is a relation by \cref{singleton_iff,relation}.
                For all $a,b_1,b_2$ such that $(a,b_1)\in y$ and $(a,b_2)\in y$ we have $b_1 = b_2$
                    by \cref{singleton_iff,pair_eq_iff}.
                Then $y$ is a function by \cref{rightunique}.
                Therefore $\dom{y} = \dom{Y_1}$
                    by \cref{subseteq,singleton_iff,dom_intro,dom_iff,pair_eq_iff,subseteq_antisymmetric}.
                For all $i$ such that $i\in\dom{Y_1}$ we have $y(i)\in\unions{\ran{Y_1}}$
                    by \cref{subseteq,singleton_iff,function_apply_intro,function_apply_elim,ran_intro,unions_iff}.
                Hence $y\in\funs{\dom{Y_1}}{\unions{\ran{Y_1}}}$ by \cref{funs_intro}.
                For all $i$ such that $i\in\dom{Y_1}$ we have $y(i)\in\apply{Y_1}{i}$
                    by \cref{subseteq,singleton_iff,funs_elim,function_apply_intro}.
                Therefore $y\in\productFamily{Y_1}$ by \cref{indexed_product}.
                Hence $g(x)\in\productFamily{Y_1}$ by \cref{function_apply_intro}.
            \end{subproof}
            Therefore $g\in\funs{X}{\productFamily{Y_1}}$ by \cref{funs_intro}.
            Let $T_P = \productTopologyIndexed{T_1}{Y_1}$.
        $J_1$ is inhabited.
        Show that for all $i\in J_1$ we have $\piIndex{Y_1}{i}\circ g$ is continuous from $X$ to $\apply{Y_1}{i}$ with respect to $T_X$ and $\apply{T_1}{i}$.
        \begin{subproof}
            Fix $i$ such that $i\in J_1$.
            Then $i = 1$.
            Let $j = \identity{X}$.
            $j$ is continuous from $X$ to $X$ with respect to $T_X$ and $T_X$ by \cref{identity_continuous}.
                    Show that for all $x\in X$ we have $\apply{\piIndex{Y_1}{i}\circ g}{x} = \apply{j}{x}$.
                    \begin{subproof}
                        Fix $x$ such that $x\in X$.
                        Thus $g$ is a function to $\productFamily{Y_1}$ and $g$ is a function and $\dom{g} = X$
                            by \cref{funs_elim,funs_is_function}.
                        Thus $g(x)\in\productFamily{Y_1}$ by \cref{subseteq,funs,function_apply_elim,ran_intro,function_to_ran}.
                        Thus $(x,\{(1,x)\})\in g$.
                        Thus $g(x) = \{(1,x)\}$ by \cref{function_apply_intro}.
                        Thus $g(x)\in\funs{\dom{Y_1}}{\unions{\ran{Y_1}}}$ and $g(x)$ is a function
                            by \cref{indexed_product,funs_elim}.
                        Thus $(1,x)\in g(x)$ by \cref{singleton_iff}.
                        Hence $\apply{g(x)}{i} = x$ by \cref{function_apply_intro}.
                        Thus $(g(x),x)\in\piIndex{Y_1}{i}$ by \cref{pair_eq_iff,projection_map_indexed}.
                        Thus $(x,x)\in\piIndex{Y_1}{i}\circ g$ by \cref{circ_iff}.
                        Thus $\piIndex{Y_1}{i}\in\funs{\productFamily{Y_1}}{\apply{Y_1}{i}}$
                            and $\piIndex{Y_1}{i}$ is a function by \cref{projection_map_indexed_function,funs_is_function}.
                        Thus $\piIndex{Y_1}{i}\circ g$ is a function by \cref{function_circ}.
                        Thus $\apply{\piIndex{Y_1}{i}\circ g}{x} = x$ by \cref{function_apply_intro}.
                        Therefore $\apply{\piIndex{Y_1}{i}\circ g}{x} = \apply{j}{x}$ by \cref{id_apply}.
                    \end{subproof}
            $\piIndex{Y_1}{i}\circ g$ is a function and $\dom{\piIndex{Y_1}{i}\circ g} = X$
                by \cref{funs_elim,projection_map_indexed_function,funs_is_function,function_circ,dom_iff,circ_iff,dom_intro,subseteq,function_apply_elim,ran_intro,function_to_ran,projection_map_indexed,subseteq_antisymmetric}.
            $j$ is a function and $\dom{j} = X$ by \cref{id_is_function,id_dom}.
            Therefore $\piIndex{Y_1}{i}\circ g = j$ by \cref{subseteq,fun_subseteq,subseteq_antisymmetric}.
            Hence $\piIndex{Y_1}{i}\circ g$ is continuous from $X$ to $\apply{Y_1}{i}$ with respect to $T_X$ and $\apply{T_1}{i}$.
        \end{subproof}
        Therefore $g$ is continuous from $X$ to $\productFamily{Y_1}$ with respect to $T_X$ and $T_P$
            by \cref{continuous_map_into_indexed_product}.
        Show that for all $y\in\productFamily{Y_1}$ there exists $x\in X$ such that $g(x) = y$.
        \begin{subproof}
            Fix $y$ such that $y\in\productFamily{Y_1}$.
            Then $y$ is a function and $\dom{y} = \dom{Y_1}$ by \cref{indexed_product,funs_elim}.
            Hence $1\in\dom{y}$ by \cref{subseteq,real_one_in_naturals,initial_segment_naturalsplus}.
            Take $x\in X$ such that $x = y(1)$.
            Thus $(x,\{(1,x)\})\in g$.
            For all $i$ such that $i\in\dom{y}$ we have $\apply{g(x)}{i} = y(i)$
                by \cref{subseteq,funs_elim,function_apply_intro,singleton_iff,function_apply_elim,ran_intro,function_to_ran,indexed_product}.
            Thus $g(x)$ is a function and $\dom{g(x)} = \dom{Y_1}$
                by \cref{subseteq,funs,funs_is_function,function_apply_elim,ran_intro,function_to_ran,indexed_product,funs_elim}.
            Therefore $g(x) = y$ by \cref{subseteq,fun_subseteq,subseteq_antisymmetric}.
            Hence there exists $x\in X$ such that $g(x) = y$.
        \end{subproof}
        Therefore $g\in\surj{X}{\productFamily{Y_1}}$ by \cref{surj}.
        Hence $\img{g}{X} = \productFamily{Y_1}$ by \cref{surj,funs_elim,img_dom,funs_surj_iff}.
        Let $A = \img{g}{X}$.
        Let $T_A = \subspaceTopology{T_P}{A}$.
        Then $T_A = \subspaceTopology{T_P}{\productFamily{Y_1}}$.
        $T_P$ is a topology on $\productFamily{Y_1}$
            by \cref{product_subbasis_finite_intersections_basis,basis_for_topology,basis_topology_is_topology}.
        For all $U$ such that $U\in T_A$ we have $U\in T_P$
            by \cref{subspace_topology,topology_on,subseteq,powerset_elim,inter_absorb_supseteq_left}.
        For all $U$ such that $U\in T_P$ we have $U\in T_A$
            by \cref{topology_on,subseteq,powerset_elim,inter_absorb_supseteq_left,subspace_topology}.
        Therefore $T_A = T_P$ by \cref{subseteq,subseteq_antisymmetric}.
        Hence $\productFamily{Y_1}$ is compact with respect to $T_P$ by \cref{continuous_image_compact}.
    \end{subproof}
        Hence $\finiteIndexedProductCompactProp{1}$ by \cref{finite_indexed_product_compact_property}.
        Hence $1\in A$ by \cref{real_one_in_naturals}.
    \end{subproof}
    For all $k$ such that $k\in A$ we have $k\in\naturalsPlus$ and $\finiteIndexedProductCompactProp{k}$.
    Thus for all $k\in A$ we have $k + 1\in A$
        by \cref{real_naturals_closed_succ,finite_indexed_product_compact_property_succ}.
    Hence $n\in A$ by \cref{real_naturals_induction}.
    Therefore $\productFamily{Y}$ is compact with respect to $\productTopologyIndexed{T_Y}{Y}$
        by \cref{subseteq,finite_indexed_product_compact_property}.
\end{proof}
% from §27 helper lemma 9: square and euclidean metric topologies agree
\begin{lemma}\label{square_metric_topology_equals_euclidean}
    Suppose $n\in\naturalsPlus$.
    Let $J = \initialSegment{n}$.
    Let $X = \realsPower{J}$.
    Let $T_E = \metricTopology{\euclideanMetric{n}}{X}$.
    Let $T_S = \metricTopology{\squareMetric{n}}{X}$.
    Then $T_E = T_S$.
\end{lemma}
\begin{proof}
    Let $X_0 = X$.
    Let $d = \euclideanMetric{n}$.
    Let $d_1 = \squareMetric{n}$.
    $d$ is a metric on $X_0$ and $d_1$ is a metric on $X_0$ by \cref{euclidean_metric_is_metric,square_metric_is_metric}.
    Show that for all $x\in X_0$ we have
        for all $\epsilon\in\reals$ such that $0 < \epsilon$ there exists $\delta\in\reals$
            such that $0 < \delta$ and
            $\metricBall{d}{X_0}{x}{\delta} \subseteq \metricBall{d_1}{X_0}{x}{\epsilon}$.
    \begin{subproof}
        Fix $x$ such that $x\in X_0$.
        Fix $\epsilon$ such that $\epsilon\in\reals$ and $0 < \epsilon$.
        Let $\delta = \epsilon$. Then $0 < \delta$.
        For all $y$ such that $y\in\metricBall{d}{X_0}{x}{\delta}$ we have $y\in\metricBall{d_1}{X_0}{x}{\epsilon}$
            by \cref{metric_ball,euclidean_metric_value_exists,square_metric_value_exists,euclidean_square_metric_bounds,euclidean_metric_function,square_metric_function,real_leq_split,real_lt_subst_left,real_lt_trans,function_apply_intro}.
        Therefore $\metricBall{d}{X_0}{x}{\delta} \subseteq \metricBall{d_1}{X_0}{x}{\epsilon}$ by \cref{subseteq}.
    \end{subproof}
    Hence $T_E$ is finer than $T_S$ on $X_0$ by \cref{metric_topology_finer_iff}.
    Show that for all $x\in X_0$ we have
        for all $\epsilon\in\reals$ such that $0 < \epsilon$ there exists $\delta\in\reals$
            such that $0 < \delta$ and
            $\metricBall{d_1}{X_0}{x}{\delta} \subseteq \metricBall{d}{X_0}{x}{\epsilon}$.
    \begin{subproof}
        Fix $x$ such that $x\in X_0$.
        Fix $\epsilon$ such that $\epsilon\in\reals$ and $0 < \epsilon$.
        $0 < 1$ and $0\in\reals$ and $1\in\reals$ and $n\in\reals$
            by \cref{real_zero_lt_one,real_add_ident,real_naturals_subset_reals,real_one_in_naturals,subseteq}.
        Hence $0 < n$ by \cref{real_one_leq_naturalsplus,real_lt_subst_right,real_lt_trans}.
        Then $\sqrt{n}\in\reals$ and $\sqrt{n} \geq 0$ and $\sqrt{n} \rmul \sqrt{n} = n$ by \cref{positive_square_root}.
        Hence $\sqrt{n}$ is positive.
        Let $\delta = \epsilon \rmul \inv{\sqrt{n}}$.
        $\sqrt{n} \neq 0$ and $\inv{\sqrt{n}}\in\reals$ and $\inv{\sqrt{n}} > 0$
            by \cref{real_pos_nonzero,real_mul_inverse,real_inv_pos}.
        Therefore $\delta\in\reals$ and $\delta > 0$ by \cref{real_mul_closed,real_mul_pos_pos_gt}.
        Show that for all $y$ such that $y\in\metricBall{d_1}{X_0}{x}{\delta}$ we have
            $y\in\metricBall{d}{X_0}{x}{\epsilon}$.
        \begin{subproof}
            Fix $y$ such that $y\in\metricBall{d_1}{X_0}{x}{\delta}$.
            Thus $y\in X_0$ and $d_1((x,y)) < \delta$ by \cref{metric_ball}.
            Take $m\in\reals$ such that $((x,y),m)\in d_1$ by \cref{square_metric_value_exists}.
            Take $r\in\reals$ such that $((x,y),r)\in d$ by \cref{euclidean_metric_value_exists}.
            Hence $m < \delta$ by \cref{square_metric_function,function_apply_intro,real_lt_subst_left}.
            Then $r \leq \sqrt{n} \rmul m$ by \cref{euclidean_square_metric_bounds}.
            $m \rmul \sqrt{n} < \delta \rmul \sqrt{n}$ by \cref{real_lt_mul_pos}.
            Hence $\sqrt{n} \rmul m < \sqrt{n} \rmul \delta$ by \cref{real_mul_comm}.
            Then $\sqrt{n} \rmul \delta = \epsilon$ by
                \cref{real_mul_assoc,real_mul_comm,real_mul_inverse,real_mul_ident}.
            Hence $\sqrt{n} \rmul m < \epsilon$ by \cref{real_lt_subst_right}.
            Therefore $d((x,y)) < \epsilon$ by
                \cref{euclidean_metric_function,real_mul_closed,real_leq_split,real_lt_subst_left,real_lt_trans,function_apply_intro}.
            Hence $y\in\metricBall{d}{X_0}{x}{\epsilon}$ by \cref{metric_ball}.
        \end{subproof}
        Therefore $\metricBall{d_1}{X_0}{x}{\delta} \subseteq \metricBall{d}{X_0}{x}{\epsilon}$ by \cref{subseteq}.
    \end{subproof}
    Hence $T_S$ is finer than $T_E$ on $X_0$ by \cref{metric_topology_finer_iff}.
    Hence $T_E = T_S$ by \cref{finer_topology,subseteq_antisymmetric}.
\end{proof}

% from §27 helper lemma 10: indexed real/standard families are functions with domain J
\begin{lemma}\label{reals_standard_indexed_families_data}
    Suppose $n\in\naturalsPlus$.
    Let $J = \initialSegment{n}$.
    Let $X_1 = \{(i,\reals) \mid i\in J\}$.
    Let $T_1 = \{(i,\standardTopologyR) \mid i\in J\}$.
    Then $X_1$ is a function and $T_1$ is a function and $\dom{X_1} = J$ and $\dom{T_1} = J$.
\end{lemma}
\begin{proof}
    For all $i,y_1,y_2$ such that $(i,y_1)\in X_1$ and $(i,y_2)\in X_1$ we have $y_1 = y_2$.
    For all $i,y_1,y_2$ such that $(i,y_1)\in T_1$ and $(i,y_2)\in T_1$ we have $y_1 = y_2$.
    Hence $X_1$ is a function and $T_1$ is a function and $\dom{X_1} = J$ and $\dom{T_1} = J$
        by \cref{relation,rightunique,dom_intro,dom_iff,pair_eq_iff,subseteq,subseteq_antisymmetric}.
\end{proof}

% from §27 Item 3: compact subsets of $\reals^n$
\begin{theorem}\label{compact_reals_power_iff_closed_bounded}
    Suppose $n\in\naturalsPlus$.
    Let $J = \initialSegment{n}$.
    Let $X = \realsPower{J}$.
    Let $d = \euclideanMetric{n}$.
    Let $\rho = \squareMetric{n}$.
    Let $T = \metricTopology{\rho}{X}$.
    Suppose $A\subseteq X$.
    Let $T_A = \subspaceTopology{T}{A}$.
    Then $A$ is compact with respect to $T_A$ iff
        $A$ is closed in $X$ with respect to $T$ and
        $A$ is bounded in $X$ with respect to $d$ or $A$ is bounded in $X$ with respect to $\rho$.
\end{theorem}
\begin{proof}
    Let $X_0 = X$.
    Let $d_1 = \rho$.
    Show that if $A$ is compact with respect to $T_A$ then
        $A$ is closed in $X$ with respect to $T$ and
        $A$ is bounded in $X$ with respect to $d$ or $A$ is bounded in $X$ with respect to $\rho$.
    \begin{subproof}
        Assume $A$ is compact with respect to $T_A$.
        $d_1$ is a metric on $X$ by \cref{square_metric_is_metric}.
        $T$ is a topology on $X$ by \cref{basis_topology_is_topology,metric_topology,metric_basis_is_basis,basis_for_topology}.
        Let $T_S = \metricTopology{d_1}{X_0}$.
        Let $X_1 = \{(i,\reals) \mid i\in J\}$.
        Let $T_1 = \{(i,\standardTopologyR) \mid i\in J\}$.
        Let $T_0 = \productTopologyIndexed{T_1}{X_1}$.
        $d$ is a metric on $X_0$ and $d_1$ is a metric on $X_0$ by \cref{euclidean_metric_is_metric,square_metric_is_metric}.
        Hence $J$ is inhabited by \cref{real_one_in_naturals,real_one_leq_naturalsplus,initial_segment_naturalsplus}.
        $X_1$ is a function and $T_1$ is a function and $\dom{X_1} = J$ and $\dom{T_1} = J$
            by \cref{reals_standard_indexed_families_data}.
        For all $\alpha\in\dom{X_1}$ we have $\apply{X_1}{\alpha} = \reals$ and $\apply{T_1}{\alpha} = \standardTopologyR$
            by \cref{subseteq,function_apply_intro}.
        $\realOrderTopology = \standardTopologyR$ and $\realOrderTopology = \basisTopology{\orderBasis{\realOrderRel}{\reals}}{\reals}$
            by \cref{real_order_topology_eq_standard,real_order_topology}.
        $\realOrderRel$ is a simple order relation on $\reals$ by \cref{real_order_relation_simple}.
        $\reals$ has more than one element by \cref{linear_continuum,reals_linear_continuum}.
        Hence $\realOrderTopology$ is the order topology on $\reals$ with respect to $\realOrderRel$ by \cref{order_topology}.
        Therefore for all $\alpha\in\dom{X_1}$ we have $\apply{X_1}{\alpha}$ is Hausdorff with respect to $\apply{T_1}{\alpha}$
            by \cref{order_topology_hausdorff,real_order_topology_eq_standard}.
        Hence $X_0$ is Hausdorff with respect to $T_0$ by \cref{product_hausdorff_product,product_family_reals_power}.
        Therefore $T_S = T_0$ by \cref{square_metric_topology_equals_euclidean,euclidean_metric_product_topology}.
        Then $X$ is Hausdorff with respect to $T$.
        Therefore $A$ is closed in $X$ with respect to $T$ by \cref{compact_subspace_closed}.
        Show that $A$ is bounded in $X$ with respect to $\rho$.
        \begin{subproof}
            \begin{byCase}
            \caseOf{$A = \emptyset$.}
                Let $M = 0$.
                For all $a_1,a_2\in A$ we have $\rho((a_1,a_2)) \leq M$
                    by \cref{subseteq_emptyset_iff,elem_subseteq,emptyset}.
                Thus $A$ is bounded in $X$ with respect to $\rho$ by \cref{real_add_ident,metric_bounded}.
            \caseOf{$A \neq \emptyset$.}
                $A$ is inhabited by \cref{nonempty_inhabited}.
                Let $C = \{\metricBall{\rho}{X}{a}{1} \mid a\in A\}$.
                For all $U$ such that $U\in C$ we have $U$ is open in $X$ with respect to $T$
                    by \cref{real_naturals_subset_reals,real_one_in_naturals,subseteq,real_zero_lt_one,metric_ball,powerset_intro,metric_basis,metric_basis_is_basis,basis_elem_open_in_generated,metric_topology,basis_for_topology,open_in}.
                For all $x$ such that $x\in A$ we have $x\in\unions{C}$
                    by \cref{subseteq,square_metric_is_metric,metric_zero_characterization,metric_on,real_zero_lt_one,real_lt_subst_left,metric_ball,unions_iff}.
                Hence $C$ is a covering of $A$ by \cref{subseteq,covering}.
                $A$ is a subspace of $X$ with respect to $T$ by \cref{subspace_of}.
                Thus there exists $B$ such that
                    $B\subseteq C$ and
                    $B$ is finite and
                    $B$ is a covering of $A$
                    by \cref{compact_subspace_open_covering}.
                Let $R = \{w\in B\times A \mid \exists U\in B.\exists a\in A. w = (U,a) \land U = \metricBall{\rho}{X}{a}{1}\}$.
                For all $w$ such that $w\in R$ there exist $x,y$ such that $w = (x,y)$.
                Thus $R$ is a relation by \cref{relation}.
                For all $U\in B$ there exists $a$ such that $U\mathrel{R} a$.
                Take $g$ such that $g$ is a function on $B$ and for all $U\in B$ we have $U\mathrel{R}\apply{g}{U}$
                    by \cref{choice_function}.
                Let $F = \img{g}{B}$.
                $F$ is finite by \cref{finite_image_function}.
                For all $f$ such that $f\in F$ we have $f\in A$
                    by \cref{img_iff,function_apply_intro,subseteq,times_tuple_elim}.
                Hence $F\subseteq A$ by \cref{subseteq}.
                Take $a_0$ such that $a_0\in A$.
                Thus $a_0\in\unions{B}$ by \cref{covering,subseteq}.
                Take $U_0\in B$ such that $a_0\in U_0$ by \cref{unions_iff}.
                Take $a_0$ such that $a_0\in F$ by \cref{subseteq,function_apply_elim,img_iff}.
                Let $h = \{(f,\rho((a_0,f))) \mid f\in F\}$.
                Show that $h$ is a function on $F$.
                \begin{subproof}
                    For all $w$ such that $w\in h$ there exist $x,y$ such that $w = (x,y)$.
                    Thus $h$ is a relation by \cref{relation}.
                    Show that for all $x,y_1,y_2$ such that $(x,y_1)\in h$ and $(x,y_2)\in h$ we have $y_1 = y_2$.
                    \begin{subproof}
                        Fix $x$.
                        Fix $y_1$.
                        Fix $y_2$ such that $(x,y_1)\in h$ and $(x,y_2)\in h$.
                        Take $f_1\in F$ such that $(x,y_1) = (f_1,\rho((a_0,f_1)))$.
                        Take $f_2\in F$ such that $(x,y_2) = (f_2,\rho((a_0,f_2)))$.
                        Thus $\rho((a_0,f_1)) = \rho((a_0,f_2))$ by \cref{pair_eq_iff}.
                        Thus $y_1 = y_2$ by \cref{pair_eq_iff}.
                    \end{subproof}
                    Hence $h$ is right-unique by \cref{rightunique}.
                    For all $f$ such that $f\in F$ we have $f\in\dom{h}$.
                    For all $f$ such that $f\in\dom{h}$ we have $f\in F$ by \cref{dom_iff,pair_eq_iff}.
                    Therefore $\dom{h} = F$ by \cref{subseteq,subseteq_antisymmetric}.
                    Hence $h$ is a function on $F$.
                \end{subproof}
                Let $D = \img{h}{F}$.
                $D$ is finite by \cref{finite_image_function}.
                For all $r$ such that $r\in D$ we have $r\in\reals$
                    by \cref{img_iff,subseteq,times_tuple_intro,metric_on,funs_type_apply,square_metric_is_metric,pair_eq_iff}.
                Hence $D\subseteq\reals$ by \cref{subseteq}.
                Take $f_0$ such that $f_0\in F$.
                Thus $(f_0,\rho((a_0,f_0)))\in h$.
                Therefore $D\neq\emptyset$ by \cref{img_iff,nonempty_inhabited,inhabited_nonempty}.
                Take $M_0\in D$ such that $M_0$ is a maximum of $D$ by \cref{finite_real_maximum}.
                Let $M = (M_0 + 1) + (M_0 + 1)$.
                Show that for all $x,y\in A$ we have $\rho((x,y)) \leq M$.
                \begin{subproof}
                    Fix $x$.
                    Fix $y$ such that $x\in A$ and $y\in A$.
                    Take $U_x\in B$ such that $x\in U_x$ by \cref{covering,subseteq,unions_iff}.
                    Take $U_y\in B$ such that $y\in U_y$ by \cref{covering,subseteq,unions_iff}.
                    Let $f_x = g(U_x)$.
                    Let $f_y = g(U_y)$.
                    $U_x\mathrel{R} f_x$ by \cref{function_apply_elim}.
                    $U_y\mathrel{R} f_y$ by \cref{function_apply_elim}.
                    Thus $(U_x,f_x)\in R$.
                    Take $U_0, a_0$ such that
                        $U_0\in B$ and $a_0\in A$ and
                        $(U_x,f_x) = (U_0,a_0)$ and $U_0 = \metricBall{\rho}{X}{a_0}{1}$.
                    Thus $U_x = \metricBall{\rho}{X}{a_0}{1}$ and $f_x = a_0$ by \cref{pair_eq_iff}.
                    Thus $U_x = \metricBall{\rho}{X}{f_x}{1}$.
                    Thus $(U_y,f_y)\in R$.
                    Take $U_1, a_1$ such that
                        $U_1\in B$ and $a_1\in A$ and
                        $(U_y,f_y) = (U_1,a_1)$ and $U_1 = \metricBall{\rho}{X}{a_1}{1}$.
                    Thus $U_y = \metricBall{\rho}{X}{a_1}{1}$ and $f_y = a_1$ by \cref{pair_eq_iff}.
                    Thus $U_y = \metricBall{\rho}{X}{f_y}{1}$.
                    $x\in U_x$.
                    Hence $x\in \metricBall{\rho}{X}{f_x}{1}$.
                    Thus $x\in X$ and $\rho((f_x,x)) < 1$ by \cref{metric_ball}.
                    Thus $f_x\in A$ and $f_x\in X$ by \cref{subseteq,times_tuple_elim}.
                    Thus $\rho((x,f_x)) = \rho((f_x,x))$ by \cref{metric_symmetric,metric_on,square_metric_is_metric}.
                    Therefore $\rho((x,f_x)) < 1$ by \cref{real_lt_subst_left}.
                    $y\in U_y$.
                    Hence $y\in \metricBall{\rho}{X}{f_y}{1}$.
                    Thus $y\in X$ and $\rho((f_y,y)) < 1$ by \cref{metric_ball}.
                    Thus $f_y\in A$ and $f_y\in X$ by \cref{subseteq,times_tuple_elim}.
                    Thus $\rho((y,f_y)) = \rho((f_y,y))$ by \cref{metric_symmetric,metric_on,square_metric_is_metric}.
                    Therefore $\rho((y,f_y)) < 1$ by \cref{real_lt_subst_left}.
                    Then $f_x\in F$ and $f_y\in F$ by \cref{subseteq,function_apply_elim,img_iff}.
                    Hence $h(f_x)\in D$ and $h(f_y)\in D$ by \cref{subseteq,function_apply_elim,img_iff}.
                    Therefore $h(f_x) \leq M_0$ and $h(f_y) \leq M_0$ by \cref{real_maximum}.
                    Then $\rho((a_0,f_x)) \leq M_0$ and $\rho((a_0,f_y)) \leq M_0$ by \cref{function_apply_intro}.
                    Show that $\rho((x,y)) \leq \rho((x,f_x)) + \rho((f_x,a_0)) + \rho((a_0,f_y)) + \rho((f_y,y))$.
                    \begin{subproof}
                        $x\in X$ and $y\in X$ and $f_x\in X$ and $f_y\in X$ and $a_0\in X$ by \cref{subseteq}.
                        Thus $\rho((x,y)) \leq \rho((x,f_x)) + \rho((f_x,y))$
                            by \cref{square_metric_is_metric,metric_on,metric_triangle_inequality}.
                        Thus $\rho((f_x,y)) \leq \rho((f_x,a_0)) + \rho((a_0,y))$
                            by \cref{square_metric_is_metric,metric_on,metric_triangle_inequality}.
                        Thus $\rho((a_0,y)) \leq \rho((a_0,f_y)) + \rho((f_y,y))$
                            by \cref{square_metric_is_metric,metric_on,metric_triangle_inequality}.
                        Thus $x\in X$ and $f_x\in A$ and $f_x\in X$ and $f_y\in A$ by \cref{elem_subseteq}.
                        Thus $f_y\in X$ and $y\in X$ and $a_0\in A$ and $a_0\in X$ by \cref{elem_subseteq}.
                        $(x,f_x)\in X\times X$ by \cref{times_tuple_intro}.
                        $(f_x,y)\in X\times X$ by \cref{times_tuple_intro}.
                        $(f_x,a_0)\in X\times X$ by \cref{times_tuple_intro}.
                        $(a_0,y)\in X\times X$ by \cref{times_tuple_intro}.
                        $(a_0,f_y)\in X\times X$ by \cref{times_tuple_intro}.
                        $(f_y,y)\in X\times X$ by \cref{times_tuple_intro}.
                        $(x,y)\in X\times X$ by \cref{times_tuple_intro}.
                        Thus $\rho((x,y))\in\reals$ by \cref{metric_on,funs_type_apply}.
                        Thus $\rho((x,f_x))\in\reals$ by \cref{metric_on,funs_type_apply}.
                        Thus $\rho((f_x,y))\in\reals$ by \cref{metric_on,funs_type_apply}.
                        Thus $\rho((f_x,a_0))\in\reals$ by \cref{metric_on,funs_type_apply}.
                        Thus $\rho((a_0,y))\in\reals$ by \cref{metric_on,funs_type_apply}.
                        Thus $\rho((a_0,f_y))\in\reals$ by \cref{metric_on,funs_type_apply}.
                        Thus $\rho((f_y,y))\in\reals$ by \cref{metric_on,funs_type_apply}.
                        Let $r_1 = \rho((x,f_x))$.
                        Let $r_2 = \rho((f_x,y))$.
                        Let $r_3 = \rho((f_x,a_0))$.
                        Let $r_4 = \rho((a_0,y))$.
                        Thus $r_1\in\reals$ and $r_2\in\reals$ and $r_3\in\reals$ and $r_4\in\reals$.
                        Thus $r_3 + r_4\in\reals$ by \cref{real_add_closed}.
                        Thus $r_1 + r_2\in\reals$ by \cref{real_add_closed}.
                        Thus $r_1 + (r_3 + r_4)\in\reals$ by \cref{real_add_closed}.
                        $r_2 \leq r_3 + r_4$.
                        Hence $r_2 + r_1 \leq (r_3 + r_4) + r_1$ by \cref{real_leq_add}.
                        Then $r_1 + r_2 = r_2 + r_1$ by \cref{real_add_comm}.
                        $(r_3 + r_4) + r_1 = r_1 + (r_3 + r_4)$ by \cref{real_add_comm}.
                        Therefore $r_1 + r_2 \leq (r_3 + r_4) + r_1$ by \cref{real_leq_trans}.
                        Hence $r_1 + r_2 \leq r_1 + (r_3 + r_4)$ by \cref{real_leq_trans}.
                        Then $\rho((x,y)) \leq r_1 + (r_3 + r_4)$ by \cref{real_leq_trans}.
                        Let $u = \rho((a_0,f_y)) + \rho((f_y,y))$.
                        Thus $u\in\reals$ by \cref{real_add_closed}.
                        Thus $r_1 + r_3\in\reals$ by \cref{real_add_closed}.
                        Thus $(r_1 + r_3) + u\in\reals$ by \cref{real_add_closed}.
                        $r_4 \leq u$.
                        Therefore $r_4 + (r_1 + r_3) \leq u + (r_1 + r_3)$ by \cref{real_leq_add}.
                        Thus $(r_1 + r_3) + r_4 = r_4 + (r_1 + r_3)$ by \cref{real_add_comm}.
                        Thus $u + (r_1 + r_3) = (r_1 + r_3) + u$ by \cref{real_add_comm}.
                        Hence $(r_1 + r_3) + r_4 \leq u + (r_1 + r_3)$ by \cref{real_leq_trans}.
                        Then $(r_1 + r_3) + r_4 \leq (r_1 + r_3) + u$ by \cref{real_leq_trans}.
                        $r_1 + (r_3 + r_4) = (r_1 + r_3) + r_4$ by \cref{real_add_assoc}.
                        Therefore $r_1 + (r_3 + r_4) \leq (r_1 + r_3) + u$ by \cref{real_leq_trans}.
                        Hence $\rho((x,y)) \leq (r_1 + r_3) + u$ by \cref{real_leq_trans}.
                        Therefore $\rho((x,y)) \leq \rho((x,f_x)) + \rho((f_x,a_0)) + \rho((a_0,f_y)) + \rho((f_y,y))$
                            by \cref{real_add_assoc,real_add_comm}.
                    \end{subproof}
                    Show that $\rho((x,y)) \leq M$.
                    \begin{subproof}
                        Let $s_1 = \rho((x,f_x))$.
                        Let $s_2 = \rho((f_x,a_0))$.
                        Let $s_3 = \rho((a_0,f_y))$.
                        Let $s_4 = \rho((f_y,y))$.
                        Thus $x\in X$ and $y\in X$ and $f_x\in A$ and $f_x\in X$ and $f_y\in A$ and $f_y\in X$ and $a_0\in A$ and $a_0\in X$
                            by \cref{elem_subseteq}.
                        $(x,y)\in X\times X$ by \cref{times_tuple_intro}.
                        $(x,f_x)\in X\times X$ by \cref{times_tuple_intro}.
                        $(f_x,a_0)\in X\times X$ by \cref{times_tuple_intro}.
                        $(a_0,f_y)\in X\times X$ by \cref{times_tuple_intro}.
                        $(f_y,y)\in X\times X$ by \cref{times_tuple_intro}.
                        Thus $s_1\in\reals$ and $s_2\in\reals$ and $s_3\in\reals$ and $s_4\in\reals$ by \cref{metric_on,funs_type_apply}.
                        $s_1 + s_2\in\reals$ and $s_3 + s_4\in\reals$ by \cref{real_add_closed}.
                        $1\in\reals$ by \cref{real_naturals_subset_reals,real_one_in_naturals,subseteq}.
                        Hence $s_2 = \rho((a_0,f_x))$ by \cref{metric_symmetric,metric_on,square_metric_is_metric}.
                        Therefore $s_1 \leq 1$ by \cref{real_lt_imp_leq}.
                        Then $s_4 \leq 1$ by \cref{metric_symmetric,metric_on,square_metric_is_metric,real_lt_subst_left,real_lt_imp_leq}.
                        Hence $s_2 \leq M_0$ and $s_3 \leq M_0$.
                        Thus $M_0\in\reals$ by \cref{real_maximum,subseteq}.
                        $M_0 + 1\in\reals$ by \cref{real_add_closed}.
                        $(s_1 + s_2) + (s_3 + s_4)\in\reals$ by \cref{real_add_closed}.
                        $(M_0 + 1) + (M_0 + 1)\in\reals$ by \cref{real_add_closed}.
                        Therefore $\rho((x,y))\in\reals$ by \cref{metric_on,funs_type_apply}.
                        $\rho((x,y)) \leq s_1 + s_2 + s_3 + s_4$.
                        Hence $s_1 + s_2 + s_3 + s_4 = (s_1 + s_2) + (s_3 + s_4)$ by \cref{real_add_assoc}.
                        Therefore $\rho((x,y)) \leq (s_1 + s_2) + (s_3 + s_4)$ by \cref{real_leq_trans}.
                        Then $(s_1 + s_2) + (s_3 + s_4) \leq (M_0 + 1) + (M_0 + 1)$ by \cref{real_four_sum_bound}.
                        Hence $\rho((x,y)) \leq (M_0 + 1) + (M_0 + 1)$ by \cref{real_leq_trans}.
                        Therefore $\rho((x,y)) \leq M$.
                    \end{subproof}
                \end{subproof}
                Hence $M\in\reals$ by \cref{real_maximum,subseteq,real_naturals_subset_reals,real_one_in_naturals,real_add_closed}.
                Therefore $A$ is bounded in $X$ with respect to $\rho$ by \cref{metric_bounded}.
            \end{byCase}
        \end{subproof}
        Hence $A$ is bounded in $X$ with respect to $d$ or $A$ is bounded in $X$ with respect to $\rho$.
    \end{subproof}
    Show that if
        $A$ is closed in $X$ with respect to $T$ and
        $A$ is bounded in $X$ with respect to $d$ or $A$ is bounded in $X$ with respect to $\rho$
        then $A$ is compact with respect to $T_A$.
    \begin{subproof}
        Assume $A$ is closed in $X$ with respect to $T$ and
            $A$ is bounded in $X$ with respect to $d$ or $A$ is bounded in $X$ with respect to $\rho$.
        \begin{byCase}
            \caseOf{$A = \emptyset$.}
                For all $C$ such that $C$ is an open covering of $A$ with respect to $T_A$ there exists $B$ such that
                    $B\subseteq C$ and $B$ is finite and $B$ is a covering of $A$
                    by \cref{emptyset_is_finite,emptyset_subseteq,covering}.
            Therefore $A$ is compact with respect to $T_A$ by \cref{compact}.
        \caseOf{$A \neq \emptyset$.}
            $A$ is inhabited by \cref{nonempty_inhabited}.
            Show that $A$ is bounded in $X$ with respect to $\rho$.
            \begin{subproof}
                \begin{byCase}
                \caseOf{$A$ is bounded in $X$ with respect to $\rho$.}
                    Thus $A$ is bounded in $X$ with respect to $\rho$.
                \caseOf{$A$ is bounded in $X$ with respect to $d$.}
                    Take $M\in\reals$ such that for all $a_1,a_2\in A$ we have $d((a_1,a_2)) \leq M$
                        by \cref{metric_bounded}.
                    For all $a_1,a_2\in A$ we have $\rho((a_1,a_2)) \leq M$
                        by \cref{elem_subseteq,euclidean_metric_value_exists,square_metric_value_exists,times_tuple_intro,euclidean_square_metric_bounds,euclidean_metric_function,square_metric_function,function_apply_intro,real_leq_trans}.
                    Hence $A$ is bounded in $X$ with respect to $\rho$ by \cref{square_metric_is_metric,metric_bounded}.
                \end{byCase}
            \end{subproof}
            Take $M\in\reals$ such that for all $a_1,a_2\in A$ we have $\rho((a_1,a_2)) \leq M$
                by \cref{metric_bounded}.
            Take $x_0$ such that $x_0\in A$.
            Let $Y = \{(i,\intervalclosed{x_0(i) - M}{x_0(i) + M}) \mid i\in J\}$.
            For all $w$ such that $w\in Y$ there exist $x,y$ such that $w = (x,y)$.
            For all $i,u,v$ such that $(i,u)\in Y$ and $(i,v)\in Y$ we have $u = v$.
            Hence $Y$ is a function by \cref{relation,rightunique}.
            Then $\dom{Y} = J$ by \cref{dom_intro,dom_iff,pair_eq_iff,subseteq,subseteq_antisymmetric}.
            Show that for all $x$ such that $x\in A$ we have $x\in \productFamily{Y}$.
            \begin{subproof}
                Fix $x$ such that $x\in A$.
                Show that for all $i\in J$ we have $x(i)\in \intervalclosed{x_0(i) - M}{x_0(i) + M}$.
                \begin{subproof}
                    Fix $i$ such that $i\in J$.
                    Thus $x_0\in\realsPower{J}$ and $x\in\realsPower{J}$ by \cref{elem_subseteq}.
                    Thus $x_0(i)\in\reals$ and $x(i)\in\reals$ by \cref{reals_power,tuple_power,funs_type_apply}.
                    $\neg{x(i)}\in\reals$ by \cref{real_add_inverse}.
                    $x_0(i) - x(i) = x_0(i) + \neg{x(i)}$ by \cref{real_sub}.
                    Thus $x_0(i) - x(i)\in\reals$ by \cref{real_add_closed}.
                    Thus $\abs{x_0(i) - x(i)}\in\reals$ by \cref{real_abs_in_reals}.
                    Take $m\in\reals$ such that $((x_0,x),m)\in \rho$ by \cref{square_metric_value_exists}.
                    Thus $m\in\coordDiffSet{n}{x_0}{x}$ and for all $t\in\coordDiffSet{n}{x_0}{x}$ we have $t \leq m$
                        by \cref{square_metric,pair_eq_iff}.
                    Thus $\abs{x_0(i) - x(i)} \leq m$ by \cref{coord_diff_set}.
                    Thus $m \leq M$ by \cref{square_metric_function,function_apply_intro,real_leq_trans}.
                    Thus $\abs{x_0(i) - x(i)} \leq M$ by \cref{real_leq_trans}.
                    Thus $\abs{x(i) - x_0(i)} \leq M$ by \cref{real_abs_sub_swap,real_leq_trans}.
                    Thus $M$ is nonnegative by \cref{real_abs_nonnegative,real_add_ident,real_leq_trans}.
                    Thus $x_0(i) - M \leq x(i)$ by \cref{real_abs_sub_leq_imp_left_bound}.
                    Thus $x(i) \leq x_0(i) + M$ by \cref{real_abs_sub_leq_imp_right_bound}.
                    Thus $x(i)\in\intervalclosed{x_0(i) - M}{x_0(i) + M}$ by \cref{intervalclosed}.
                \end{subproof}
                Thus $x$ is a function to $\reals$ and $\dom{x} = \dom{Y}$ by \cref{subseteq,reals_power,tuple_power,funs_elim}.
                For all $i\in\dom{Y}$ we have $x(i)\in\unions{\ran{Y}}$ by \cref{subseteq,ran_intro,unions_iff}.
                Hence $x\in\funs{\dom{Y}}{\unions{\ran{Y}}}$ by \cref{funs_intro}.
                For all $i\in\dom{Y}$ we have $x(i)\in\apply{Y}{i}$ by \cref{dom_iff,pair_eq_iff,function_apply_intro}.
                Hence $x\in\productFamily{Y}$ by \cref{indexed_product}.
            \end{subproof}
            Hence $A\subseteq \productFamily{Y}$ by \cref{subseteq}.
            Let $T_Y = \{(i,\subspaceTopology{\standardTopologyR}{\intervalclosed{x_0(i) - M}{x_0(i) + M}}) \mid i\in J\}$.
            For all $w$ such that $w\in T_Y$ there exist $x,y$ such that $w = (x,y)$.
            For all $i,u,v$ such that $(i,u)\in T_Y$ and $(i,v)\in T_Y$ we have $u = v$.
            Therefore $T_Y$ is a function by \cref{relation,rightunique}.
            Hence $\dom{T_Y} = J$ by \cref{dom_intro,dom_iff,pair_eq_iff,subseteq,subseteq_antisymmetric}.
            % product of finitely many compact intervals
            Show that for all $i\in J$ we have $\intervalclosed{x_0(i) - M}{x_0(i) + M}$ is compact
                with respect to $\subspaceTopology{\standardTopologyR}{\intervalclosed{x_0(i) - M}{x_0(i) + M}}$.
            \begin{subproof}
                Fix $i$ such that $i\in J$.
                Thus $x_0(i)\in\reals$ by \cref{elem_subseteq,reals_power,tuple_power,funs_type_apply}.
                $0\in\reals$ by \cref{real_add_ident}.
                $\rho$ is a metric on $X$ by \cref{square_metric_is_metric}.
                Hence $\rho$ is nonnegative on $X$ by \cref{metric_on}.
                Thus $x_0\in X$ by \cref{elem_subseteq}.
                $(x_0,x_0)\in X\times X$ by \cref{times_tuple_intro}.
                Thus $\rho((x_0,x_0))\in\reals$ by \cref{metric_on,funs_type_apply}.
                Thus $\rho((x_0,x_0)) \geq 0$ by \cref{metric_nonnegative}.
                Thus $\rho((x_0,x_0)) \leq M$.
                Therefore $0 \leq M$ by \cref{real_leq_trans}.
                $\neg{M}\in\reals$ and $\neg{0}\in\reals$ by \cref{real_add_inverse}.
                $0 + \neg{0} = 0$ and $0 + \neg{0} = \neg{0}$ by \cref{real_add_inverse,real_add_ident}.
                Hence $\neg{0} = 0$.
                Therefore $\neg{M} \leq \neg{0}$ by \cref{real_leq_neg}.
                Then $\neg{M} \leq 0$ by \cref{real_leq_trans}.
                Hence $\neg{M} \leq M$ by \cref{real_leq_trans}.
                Therefore $\neg{M} + x_0(i) \leq M + x_0(i)$ by \cref{real_leq_add}.
                $x_0(i) + \neg{M} = \neg{M} + x_0(i)$ and $x_0(i) + M = M + x_0(i)$ by \cref{real_add_comm}.
                Then $x_0(i) + \neg{M} \leq x_0(i) + M$ by \cref{real_leq_trans}.
                Thus $x_0(i) + M\in\reals$ by \cref{real_add_closed}.
                $x_0(i) - M = x_0(i) + \neg{M}$ by \cref{real_sub}.
                Thus $x_0(i) - M\in\reals$ by \cref{real_add_closed}.
                Thus $x_0(i) - M \leq x_0(i) + M$ by \cref{real_sub}.
                Hence $\intervalclosed{x_0(i) - M}{x_0(i) + M}$ is compact
                    with respect to $\subspaceTopology{\standardTopologyR}{\intervalclosed{x_0(i) - M}{x_0(i) + M}}$
                    by \cref{real_intervalclosed_compact_standard}.
            \end{subproof}
            For all $i\in J$ we have $\apply{T_Y}{i}$ is a topology on $\apply{Y}{i}$
                by \cref{subseteq,function_apply_intro,standard_basis_is_basis,basis_topology_is_topology,standard_topology_reals,intervalclosed,subspace_topology_is_topology}.
            For all $i\in J$ we have $\apply{Y}{i}$ is compact with respect to $\apply{T_Y}{i}$
                by \cref{subseteq,function_apply_intro}.
            Therefore $\productFamily{Y}$ is compact with respect to $\productTopologyIndexed{T_Y}{Y}$ by \cref{compact_indexed_product_initial_segment}.
            Let $T_{P0} = \productTopologyIndexed{T_Y}{Y}$.
            Let $T_{Y0} = \subspaceTopology{T}{\productFamily{Y}}$.
            Let $X_1 = \{(i,\reals) \mid i\in J\}$.
            Let $T_1 = \{(i,\standardTopologyR) \mid i\in J\}$.
            Let $T_0 = \productTopologyIndexed{T_1}{X_1}$.
            Hence $T = T_0$ by \cref{euclidean_metric_product_topology,square_metric_topology_equals_euclidean}.
            $X_1$ is a function and $T_1$ is a function and $\dom{X_1} = J$ and $\dom{T_1} = J$
                by \cref{reals_standard_indexed_families_data}.
            Therefore $\dom{X_1}$ is inhabited by \cref{real_one_in_naturals,real_one_leq_naturalsplus,initial_segment_naturalsplus,subseteq}.
            For all $i\in\dom{X_1}$ we have $\apply{Y}{i}\subseteq\apply{X_1}{i}$
                by \cref{subseteq,function_apply_intro,intervalclosed,subseteq_transitive}.
            For all $i\in\dom{X_1}$ we have $\apply{T_1}{i}$ is a topology on $\apply{X_1}{i}$
                by \cref{subseteq,function_apply_intro,standard_basis_is_basis,basis_topology_is_topology,standard_topology_reals}.
            Hence $\productFamily{Y}$ is a subspace of $\productFamily{X_1}$ with respect to $T$ by \cref{product_subspace_product}.
            Let $j = \identity{\productFamily{Y}}$.
            Therefore $\dom{Y}$ is inhabited by \cref{real_one_in_naturals,real_one_leq_naturalsplus,initial_segment_naturalsplus,subseteq}.
            For all $i\in\dom{Y}$ we have $\apply{T_Y}{i}$ is a topology on $\apply{Y}{i}$
                by \cref{subseteq,function_apply_intro,standard_basis_is_basis,basis_topology_is_topology,standard_topology_reals,intervalclosed,subspace_topology_is_topology}.
            Therefore $T_{P0}$ is a topology on $\productFamily{Y}$
                by \cref{product_subbasis_finite_intersections_basis,basis_for_topology,basis_topology_is_topology}.
            Then $\productFamily{X_1} = X$ by \cref{product_family_reals_power}.
            Hence $\productFamily{Y}\subseteq X$ by \cref{subspace_of,subseteq_transitive}.
            $j\in\funs{\productFamily{Y}}{X}$ by \cref{id_elem_funs,funs_weaken_codom}.
            Show that for all $i\in\dom{X_1}$ we have $\piIndex{X_1}{i}\circ j$ is continuous from $\productFamily{Y}$
                to $\apply{X_1}{i}$ with respect to $T_{P0}$ and $\apply{T_1}{i}$.
            \begin{subproof}
                Fix $i$ such that $i\in\dom{X_1}$.
                Let $g = \piIndex{Y}{i}$.
                Let $k = \identity{\apply{Y}{i}}$.
                $k$ is continuous from $\apply{Y}{i}$ to $\apply{X_1}{i}$ with respect to $\apply{T_Y}{i}$ and $\apply{T_1}{i}$
                    by \cref{subseteq,subspace_of,function_apply_intro,inclusion_continuous}.
                $g$ is continuous from $\productFamily{Y}$ to $\apply{Y}{i}$ with respect to $T_{P0}$ and $\apply{T_Y}{i}$
                    by \cref{projection_map_indexed_continuous}.
                Therefore $k\circ g$ is continuous from $\productFamily{Y}$ to $\apply{X_1}{i}$ with respect to $T_{P0}$ and $\apply{T_1}{i}$
                    by \cref{continuous_composite}.
                Show that $k\circ g = \piIndex{X_1}{i}\circ j$.
                \begin{subproof}
                    Show that for all $y,z$ if $y\mathrel{k\circ g} z$ then $y\mathrel{\piIndex{X_1}{i}\circ j} z$.
                    \begin{subproof}
                        Fix $y$.
                        Fix $z$ such that $y\mathrel{k\circ g} z$.
                        Take $x$ such that $y\mathrel{g} x$ and $x\mathrel{k} z$ by \cref{circ_iff}.
                        Thus $x = z$ and $x\in\apply{Y}{i}$ by \cref{id_iff}.
                        Take $y_0\in\productFamily{Y}$ such that $(y,x) = (y_0,\apply{y_0}{i})$ by \cref{projection_map_indexed}.
                        Hence $y\in\productFamily{Y}$ and $z = y(i)$ by \cref{pair_eq_iff}.
                        Therefore $y\in\productFamily{X_1}$ by \cref{subspace_of,subseteq}.
                        Then $(y,\apply{y}{i})\in\piIndex{X_1}{i}$ by \cref{projection_map_indexed}.
                        Hence $(y,y)\in j$ by \cref{id_iff}.
                        Therefore $(y,z)\in\piIndex{X_1}{i}\circ j$ by \cref{circ_iff}.
                    \end{subproof}
                    Show that for all $y,z$ if $y\mathrel{\piIndex{X_1}{i}\circ j} z$ then $y\mathrel{k\circ g} z$.
                    \begin{subproof}
                        Fix $y$.
                        Fix $z$ such that $y\mathrel{\piIndex{X_1}{i}\circ j} z$.
                        Take $x$ such that $y\mathrel{j} x$ and $x\mathrel{\piIndex{X_1}{i}} z$ by \cref{circ_iff}.
                        Thus $y = x$ and $y\in\productFamily{Y}$ by \cref{id_iff}.
                        Take $x_1\in\productFamily{X_1}$ such that $(x,z) = (x_1,\apply{x_1}{i})$ by \cref{projection_map_indexed}.
                        Hence $z = y(i)$ by \cref{pair_eq_iff}.
                        Therefore $(y,z)\in g$ by \cref{projection_map_indexed}.
                        Hence $(z,z)\in k$ by \cref{subseteq,indexed_product,id_iff}.
                        Therefore $y\mathrel{k\circ g} z$ by \cref{circ_iff}.
                    \end{subproof}
                    Follows by set extensionality.
                \end{subproof}
                Hence $\piIndex{X_1}{i}\circ j$ is continuous from $\productFamily{Y}$ to $\apply{X_1}{i}$ with respect to $T_{P0}$ and $\apply{T_1}{i}$
                    by \cref{continuous_eq_subst}.
            \end{subproof}
            Therefore $j$ is continuous from $\productFamily{Y}$ to $\productFamily{X_1}$ with respect to $T_{P0}$ and $T_0$
                by \cref{continuous_map_into_indexed_product}.
            Then $j$ is continuous from $\productFamily{Y}$ to $X$ with respect to $T_{P0}$ and $T$ by \cref{continuous_eq_codomain}.
            Show that for all $U$ such that $U\in T_{Y0}$ we have $U\in T_{P0}$.
            \begin{subproof}
                Fix $U$ such that $U\in T_{Y0}$.
                Take $V\in T$ such that $U = \productFamily{Y}\inter V$ by \cref{subspace_topology}.
                Let $j = \identity{\productFamily{Y}}$.
                Thus $V$ is open in $X$ with respect to $T$ by \cref{continuous,open_in}.
                Hence $\preimg{j}{V}$ is open in $\productFamily{Y}$ with respect to $T_{P0}$ by \cref{continuous}.
                For all $y$ such that $y\in\preimg{j}{V}$ we have $y\in \productFamily{Y}\inter V$
                    by \cref{preimg_iff,id_iff,inter_intro}.
                For all $y$ such that $y\in \productFamily{Y}\inter V$ we have $y\in\preimg{j}{V}$
                    by \cref{inter,id_iff,preimg_iff}.
                Therefore $\preimg{j}{V} = \productFamily{Y}\inter V$ by \cref{subseteq,subseteq_antisymmetric}.
                Then $U\in T_{P0}$ by \cref{open_in}.
            \end{subproof}
            Hence $T_{Y0}\subseteq T_{P0}$ by \cref{subseteq}.
            Show that for all $C$ such that $C$ is an open covering of $\productFamily{Y}$ with respect to $T_{Y0}$ there exists $B$ such that
                $B\subseteq C$ and $B$ is finite and $B$ is a covering of $\productFamily{Y}$.
            \begin{subproof}
                Fix $C$ such that $C$ is an open covering of $\productFamily{Y}$ with respect to $T_{Y0}$.
                Hence $\dom{Y}$ is inhabited by \cref{real_one_in_naturals,real_one_leq_naturalsplus,initial_segment_naturalsplus,subseteq}.
                For all $i\in\dom{Y}$ we have $\apply{T_Y}{i}$ is a topology on $\apply{Y}{i}$
                    by \cref{subseteq,function_apply_intro,standard_basis_is_basis,basis_topology_is_topology,standard_topology_reals,intervalclosed,subspace_topology_is_topology}.
                Hence $T_{P0}$ is a topology on $\productFamily{Y}$
                    by \cref{product_subbasis_finite_intersections_basis,basis_for_topology,basis_topology_is_topology}.
                For all $U$ such that $U\in C$ we have $U$ is open in $\productFamily{Y}$ with respect to $T_{P0}$
                    by \cref{open_covering,open_in,subseteq,continuous}.
                Hence $C$ is an open covering of $\productFamily{Y}$ with respect to $T_{P0}$ by \cref{open_covering}.
                Take $B$ such that $B\subseteq C$ and $B$ is finite and $B$ is a covering of $\productFamily{Y}$
                    by \cref{compact}.
            \end{subproof}
            Therefore $\productFamily{Y}$ is compact with respect to $T_{Y0}$ by \cref{compact}.
            Hence $T$ is a topology on $X$ by \cref{square_metric_is_metric,metric_basis_is_basis,basis_topology_is_topology,metric_topology,basis_for_topology}.
            Hence $\productFamily{Y}\subseteq X$
                by \cref{indexed_product,funs_elim,subseteq,function_apply_intro,intervalclosed,funs_intro,tuple_power,reals_power}.
            Hence $\productFamily{Y}$ is a subspace of $X$ with respect to $T$ by \cref{subspace_of}.
            Then $A = A\inter\productFamily{Y}$ by \cref{subseteq,inter_absorb_supseteq_right}.
            Hence $A$ is closed in $\productFamily{Y}$ with respect to $T_{Y0}$ by \cref{closed_in_subspace_iff,subspace_of}.
            Therefore $T$ is a topology on $X$ and $\productFamily{Y}\subseteq X$ by \cref{subspace_of}.
            Then $T_{Y0}$ is a topology on $\productFamily{Y}$ by \cref{subspace_topology_is_topology}.
            Hence $A$ is compact with respect to $T_A$ by \cref{closed_subspace_compact,subspace_subspace_topology}.
        \end{byCase}
    \end{subproof}
\end{proof}
% from §27 helper lemma 10: finite inhabited subset has a largest element
\begin{lemma}\label{finite_inhabited_largest_element}
    Assume $R$ is a simple order relation on $X$.
    Suppose $F\subseteq X$.
    Suppose $F$ is finite.
    Suppose $F$ is inhabited.
    Then there exists $b$ such that $b$ is a largest element of $F$ with respect to $R$.
\end{lemma}
\begin{proof}
    Take $k$ such that $k$ is a natural number and there exists a bijection from $k$ to $F$ by \cref{finite}.
    Take $f$ such that $f$ is a bijection from $k$ to $F$.
    Let $A = \{ n\in\naturals \mid \text{for all $G,g$ such that $g$ is a bijection from $n$ to $G$ and $G\subseteq X$ and $G$ is inhabited
        there exists $b$ such that $b$ is a largest element of $G$ with respect to $R$}\}$.
    For all $G,g$ such that $g$ is a bijection from $\emptyset$ to $G$ and $G\subseteq X$ and $G$ is inhabited
        there exists $b$ such that $b$ is a largest element of $G$ with respect to $R$
        by \cref{bijections,surj,emptyset}.
    Hence $\emptyset\in A$.
    Show that for all $n\in A$ we have $\suc{n}\in A$.
    \begin{subproof}
        Fix $n\in A$.
        Show that for all $G,g$ such that $g$ is a bijection from $\suc{n}$ to $G$ and $G\subseteq X$ and $G$ is inhabited
            there exists $b$ such that $b$ is a largest element of $G$ with respect to $R$.
        \begin{subproof}
            Fix $G$.
            Fix $g$ such that $g$ is a bijection from $\suc{n}$ to $G$ and $G\subseteq X$ and $G$ is inhabited.
            $g$ is a function and $\dom{g} = \suc{n}$ by \cref{bijection_is_function,bijections_dom}.
            Thus $n\in\dom{g}$ by \cref{suc_intro_self}.
            Let $a = g(n)$.
            Let $G_1 = \img{g}{n}$.
            Thus $G_1\subseteq G$ and $G_1\subseteq X$
                by \cref{suc,subseteq,img_subseteq,img_dom,bijections,ran_of_surj,subseteq_transitive}.
            Thus $a\in G$ by \cref{function_apply_elim,ran_iff,bijections,ran_of_surj}.
            Show that for all $z$ such that $z\in G$ we have $z\in G_1\union\{a\}$.
            \begin{subproof}
                Fix $z\in G$.
                Thus there exists $x\in\suc{n}$ such that $g(x)=z$ by \cref{bijections,surj}.
                $x\in n$ or $x = n$ by \cref{suc}.
                \begin{byCase}
                \caseOf{$x\in n$.}
                    $(x,z)\in g$ by \cref{function_apply_elim}.
                    Thus $z\in G_1$ by \cref{img_iff}.
                    Thus $z\in G_1\union\{a\}$ by \cref{union_intro_left}.
                \caseOf{$x = n$.}
                    Thus $z = a$.
                    Thus $z\in G_1\union\{a\}$ by \cref{union_intro_right,singleton_intro}.
                \end{byCase}
            \end{subproof}
            Hence $G\subseteq G_1\union\{a\}$ by \cref{subseteq}.
            Therefore $G_1\union\{a\}\subseteq G$ by \cref{singleton_subset_intro,union_subsets_is_subset}.
            Then $G = G_1\union\{a\}$ by \cref{subseteq_antisymmetric}.
            Let $h = \restrl{g}{n}$.
            $h$ is a function and $\dom{h} = \dom{g}\inter n$ by \cref{restrl_of_function_is_function,dom_of_restrl_eq_inter}.
            For all $x$ such that $x\in n$ we have $x\in\dom{g}$ by \cref{suc,bijections_dom}.
            Hence $n\subseteq\dom{g}$ by \cref{subseteq}.
            Therefore $\dom{h} = n$ by \cref{inter_absorb_supseteq_left}.
            For all $x\in\dom{h}$ we have $h(x)\in G_1$
                by \cref{function_apply_elim,restrl_subseteq,elem_subseteq,img_iff}.
            Hence $h\in\funs{n}{G_1}$ by \cref{funs_intro}.
            For all $y\in G_1$ there exists $x\in n$ such that $h(x) = y$
                by \cref{img_iff,function_apply_iff,restrl_of_function_apply}.
            Therefore $h\in\surj{n}{G_1}$ by \cref{surj}.
            Hence $h$ is injective by \cref{restrl_of_function_apply,bijections,injective_function}.
            Hence $h$ is a bijection from $n$ to $G_1$ by \cref{bijections}.
            \begin{byCase}
            \caseOf{$n = \emptyset$.}
                Thus $a$ is a largest element of $G$ with respect to $R$
                    by \cref{img_emptyset,union_emptyset,union_comm,singleton_intro,singleton_iff,largest_element}.
            \caseOf{$n\neq \emptyset$.}
                Hence $G_1$ is inhabited by \cref{nonempty_inhabited,suc,bijections_dom,function_apply_elim,img_iff}.
                Take $b$ such that $b$ is a largest element of $G_1$ with respect to $R$.
                Therefore $b\mathrel{R}a$ or $a\mathrel{R}b$ or $b = a$
                    by \cref{largest_element,elem_subseteq,simple_order_relation,connex}.
                \begin{byCase}
                \caseOf{$b\mathrel{R}a$ or $b = a$.}
                    Show that for all $y\in G$ we have $y\mathrel{R}a$ or $y = a$.
                    \begin{subproof}
                        Fix $y\in G$.
                        Thus $y\in G_1$ or $y = a$ by \cref{union_iff,singleton_iff}.
                        \begin{byCase}
                        \caseOf{$y\in G_1$.}
                            Thus $y\mathrel{R}b$ or $y = b$ by \cref{largest_element}.
                            \begin{byCase}
                            \caseOf{$b = a$.}
                                Then $y\mathrel{R}a$ or $y = a$.
                            \caseOf{$b\mathrel{R}a$.}
                                \begin{byCase}
                                \caseOf{$y\mathrel{R}b$.}
                                    Therefore $y\mathrel{R}a$ by \cref{simple_order_relation,transitive}.
                                \caseOf{$y = b$.}
                                    Then $y\mathrel{R}a$.
                                \end{byCase}
                            \end{byCase}
                        \caseOf{$y = a$.}
                            Then $y = a$.
                        \end{byCase}
                    \end{subproof}
                    Hence $a$ is a largest element of $G$ with respect to $R$
                        by \cref{union_intro_right,singleton_intro,largest_element}.
                \caseOf{$a\mathrel{R}b$.}
                    Show that for all $y\in G$ we have $y\mathrel{R}b$ or $y = b$.
                    \begin{subproof}
                        Fix $y\in G$.
                        Thus $y\in G_1$ or $y = a$ by \cref{union_iff,singleton_iff}.
                        \begin{byCase}
                        \caseOf{$y\in G_1$.}
                            Thus $y\mathrel{R}b$ or $y = b$ by \cref{largest_element}.
                        \caseOf{$y = a$.}
                            Thus $y\mathrel{R}b$.
                        \end{byCase}
                    \end{subproof}
                    Therefore $b$ is a largest element of $G$ with respect to $R$ by \cref{elem_subseteq,largest_element}.
                \end{byCase}
            \end{byCase}
        \end{subproof}
        Hence $\suc{n}\in A$.
    \end{subproof}
    Hence $A$ is an inductive set.
    Therefore $k\in A$ by \cref{num_naturals_smallest_inductive_set,subseteq}.
    Then for all $G,g$ such that $g$ is a bijection from $k$ to $G$ and $G\subseteq X$ and $G$ is inhabited
        there exists $b$ such that $b$ is a largest element of $G$ with respect to $R$.
    Hence there exists $b$ such that $b$ is a largest element of $F$ with respect to $R$.
\end{proof}

% from §27 helper lemma 11: finite inhabited subset has a smallest element
\begin{lemma}\label{finite_inhabited_smallest_element}
    Assume $R$ is a simple order relation on $X$.
    Suppose $F\subseteq X$.
    Suppose $F$ is finite.
    Suppose $F$ is inhabited.
    Then there exists $a$ such that $a$ is a smallest element of $F$ with respect to $R$.
\end{lemma}
\begin{proof}
    Then there exists $a$ such that $a$ is a smallest element of $F$ with respect to $R$
        by \cref{simple_order_relation,converse_type,converse_iff,transitive,asymmetric,connex,finite_inhabited_largest_element,largest_element,smallest_element}.
\end{proof}

% from §27 Item 4: extreme value theorem
\begin{theorem}\label{extreme_value_theorem}
    Assume $T_X$ is a topology on $X$.
    Assume $T_Y$ is the order topology on $Y$ with respect to $R$.
    Suppose $f$ is continuous from $X$ to $Y$ with respect to $T_X$ and $T_Y$.
    Suppose $X$ is compact with respect to $T_X$.
    Suppose $X\neq\emptyset$.
    Then there exist $c,d\in X$ such that
        $\apply{f}{c}$ is a smallest element of $\img{f}{X}$ with respect to $R$ and
        $\apply{f}{d}$ is a largest element of $\img{f}{X}$ with respect to $R$.
\end{theorem}
\begin{proof}
    Let $A = \img{f}{X}$.
    Let $T_A = \subspaceTopology{T_Y}{A}$.
    Hence $A$ is compact with respect to $T_A$ by \cref{continuous_image_compact}.
    $T_Y = \basisTopology{\orderBasis{R}{Y}}{Y}$ and $R$ is a simple order relation on $Y$
        and $Y$ has more than one element by \cref{order_topology}.
    Therefore $T_Y$ is a topology on $Y$ by \cref{order_topology,order_basis_is_basis,basis_topology_is_topology}.
    Then $f$ is a function to $Y$ and $\dom{f} = X$ by \cref{continuous,funs_elim}.
    Hence $A\subseteq Y$ by \cref{img_dom,function_to_ran}.
    Therefore $A$ is a subspace of $Y$ with respect to $T_Y$ by \cref{subspace_of}.
    Then $A$ is inhabited by \cref{nonempty_inhabited,continuous,funs_elim,function_apply_elim,img_iff}.
    Show that there exists $M$ such that $M$ is a largest element of $A$ with respect to $R$.
    \begin{subproof}
        Suppose not.
        Let $C = \{U \mid \exists a\in A. U = \openrayleft{R}{Y}{a}\}$.
        Show that for all $y$ such that $y\in A$ we have $y\in\unions{C}$.
        \begin{subproof}
            Fix $y\in A$.
            Take $a\in A$ such that it is wrong that $a\mathrel{R}y$ or $a = y$ by \cref{largest_element}.
            Then $a\neq y$ and it is wrong that $a\mathrel{R}y$.
            Thus $y\in Y$ and $a\in Y$ by \cref{subseteq}.
            Therefore $y\mathrel{R}a$ or $a\mathrel{R}y$ by \cref{order_topology,simple_order_relation,connex}.
            \begin{byCase}
            \caseOf{$y\mathrel{R}a$.}
                Thus $y\in\unions{C}$ by \cref{openray_left,unions_iff}.
            \caseOf{$a\mathrel{R}y$.}
                Contradiction.
            \end{byCase}
        \end{subproof}
        Hence $C$ is a covering of $A$ by \cref{subseteq,covering}.
        Take $B$ such that $B\subseteq C$ and $B$ is finite and $B$ is a covering of $A$
            by \cref{subseteq,openrayleft_open_in_order_topology,compact_subspace_open_covering,compact}.
        Thus $B$ is inhabited by \cref{nonempty_inhabited,covering,subseteq,unions_iff}.
        Let $R_B = \{w\in B\times A \mid \exists U\in B.\exists a\in A. w = (U,a) \land U = \openrayleft{R}{Y}{a}\}$.
        For all $w$ such that $w\in R_B$ there exist $U,a$ such that $w = (U,a)$.
        Therefore $R_B$ is a relation by \cref{relation}.
        For all $U\in B$ there exists $a$ such that $U\mathrel{R_B}a$.
        Take $g$ such that $g$ is a function on $B$ and for all $U\in B$ we have $U\mathrel{R_B} \apply{g}{U}$
            by \cref{axiom_of_choice}.
        Let $E = \img{g}{B}$.
        Therefore $E\subseteq A$ by \cref{img_iff,function_apply_intro,relation,pair_eq_iff,subseteq}.
        Then $E$ is finite by \cref{finite_image_function}.
        Hence $E$ is inhabited by \cref{nonempty_inhabited,function_apply_elim,img_iff}.
        Therefore $E\subseteq Y$ by \cref{subseteq_transitive}.
        Take $a_0$ such that $a_0$ is a largest element of $E$ with respect to $R$
            by \cref{finite_inhabited_largest_element}.
        Show that for all $U\in B$ we have $a_0\notin U$.
        \begin{subproof}
            Fix $U\in B$.
            Hence $U\mathrel{R_B} g(U)$.
            Thus $(U,\apply{g}{U})\in R_B$ by \cref{relation}.
            Take $U_1, a$ such that $U_1\in B$ and $a\in A$ and $(U,\apply{g}{U}) = (U_1,a)$ and
                $U_1 = \openrayleft{R}{Y}{a}$.
            Thus $U = \openrayleft{R}{Y}{a}$ and $g(U) = a$ by \cref{pair_eq_iff}.
            Thus $a\in E$ by \cref{function_apply_elim,img_iff}.
            Then $a\mathrel{R}a_0$ or $a = a_0$ by \cref{largest_element}.
            Suppose not.
            Then $a_0\in U$.
            Hence $a_0\mathrel{R}a$ by \cref{openray_left}.
            \begin{byCase}
            \caseOf{$a\mathrel{R}a_0$.}
                Contradiction by \cref{simple_order_relation,order_topology,asymmetric}.
            \caseOf{$a = a_0$.}
                Thus $a_0\mathrel{R}a_0$.
                Contradiction by \cref{simple_order_relation,order_topology,asymmetric}.
            \end{byCase}
        \end{subproof}
        Hence $a_0\notin\unions{B}$ by \cref{unions_iff}.
        Therefore $a_0\in\unions{B}$ by \cref{largest_element,elem_subseteq,covering,subseteq}.
        Contradiction.
    \end{subproof}
    Show that there exists $m$ such that $m$ is a smallest element of $A$ with respect to $R$.
    \begin{subproof}
        Suppose not.
        Let $C = \{U \mid \exists a\in A. U = \openrayright{R}{Y}{a}\}$.
        Show that for all $y$ such that $y\in A$ we have $y\in\unions{C}$.
        \begin{subproof}
            Fix $y\in A$.
            Take $a\in A$ such that it is wrong that $y\mathrel{R}a$ or $y = a$ by \cref{smallest_element}.
            Then $a\neq y$ and it is wrong that $y\mathrel{R}a$.
            Thus $y\in Y$ and $a\in Y$ by \cref{subseteq}.
            Therefore $y\mathrel{R}a$ or $a\mathrel{R}y$ by \cref{order_topology,simple_order_relation,connex}.
            \begin{byCase}
            \caseOf{$a\mathrel{R}y$.}
                Thus $y\in\unions{C}$ by \cref{openray_right,unions_iff}.
            \caseOf{$y\mathrel{R}a$.}
                Contradiction.
            \end{byCase}
        \end{subproof}
        Hence $C$ is a covering of $A$ by \cref{subseteq,covering}.
        Take $B$ such that $B\subseteq C$ and $B$ is finite and $B$ is a covering of $A$
            by \cref{subseteq,openrayright_open_in_order_topology,compact_subspace_open_covering,compact}.
        Thus $B$ is inhabited by \cref{nonempty_inhabited,covering,subseteq,unions_iff}.
        Let $R_B = \{w\in B\times A \mid \exists U\in B.\exists a\in A. w = (U,a) \land U = \openrayright{R}{Y}{a}\}$.
        For all $w$ such that $w\in R_B$ there exist $U,a$ such that $w = (U,a)$.
        Therefore $R_B$ is a relation by \cref{relation}.
        For all $U\in B$ there exists $a$ such that $U\mathrel{R_B}a$.
        Take $g$ such that $g$ is a function on $B$ and for all $U\in B$ we have $U\mathrel{R_B} \apply{g}{U}$
            by \cref{axiom_of_choice}.
        Let $E = \img{g}{B}$.
        Therefore $E\subseteq A$ by \cref{img_iff,function_apply_intro,relation,pair_eq_iff,subseteq}.
        Then $E$ is finite by \cref{finite_image_function}.
        Hence $E$ is inhabited by \cref{nonempty_inhabited,function_apply_elim,img_iff}.
        Therefore $E\subseteq Y$ by \cref{subseteq_transitive}.
        Take $a_0$ such that $a_0$ is a smallest element of $E$ with respect to $R$
            by \cref{finite_inhabited_smallest_element}.
        Show that for all $U\in B$ we have $a_0\notin U$.
        \begin{subproof}
            Fix $U\in B$.
            Hence $U\mathrel{R_B} g(U)$.
            Thus $(U,\apply{g}{U})\in R_B$ by \cref{relation}.
            Take $U_1, a$ such that $U_1\in B$ and $a\in A$ and $(U,\apply{g}{U}) = (U_1,a)$ and
                $U_1 = \openrayright{R}{Y}{a}$.
            Thus $U = \openrayright{R}{Y}{a}$ and $g(U) = a$ by \cref{pair_eq_iff}.
            Thus $a\in E$ by \cref{function_apply_elim,img_iff}.
            Then $a_0\mathrel{R}a$ or $a_0 = a$ by \cref{smallest_element}.
            Suppose not.
            Then $a_0\in U$.
            Hence $a\mathrel{R}a_0$ by \cref{openray_right}.
            \begin{byCase}
            \caseOf{$a_0\mathrel{R}a$.}
                Contradiction by \cref{simple_order_relation,order_topology,asymmetric}.
            \caseOf{$a_0 = a$.}
                Thus $a_0\mathrel{R}a_0$.
                Contradiction by \cref{simple_order_relation,order_topology,asymmetric}.
            \end{byCase}
        \end{subproof}
        Hence $a_0\notin\unions{B}$ by \cref{unions_iff}.
        Therefore $a_0\in\unions{B}$ by \cref{smallest_element,elem_subseteq,covering,subseteq}.
        Contradiction.
    \end{subproof}
    Take $m$ such that $m$ is a smallest element of $A$ with respect to $R$.
    Take $c\in X$ such that $f(c) = m$ by \cref{smallest_element,img_iff,function_apply_intro}.
    Take $M$ such that $M$ is a largest element of $A$ with respect to $R$.
    Take $d\in X$ such that $f(d) = M$ by \cref{largest_element,img_iff,function_apply_intro}.
    Therefore $\apply{f}{c}$ is a smallest element of $A$ with respect to $R$ and
        $\apply{f}{d}$ is a largest element of $A$ with respect to $R$ by \cref{smallest_element,largest_element}.
\end{proof}

% from §27 helper definition 4: distance set from a point to a set
\begin{definition}\label{point_distance_set}
    $\pointDistanceSet{d}{A}{x} = \{d((x,a)) \mid a\in A\}$.
\end{definition}

% from §27 Item 5: distance from a point to a set
\begin{definition}\label{point_set_distance}
    $r$ is the distance from $x$ to $A$ in $X$ with respect to $d$ iff
        $x\in X$ and $A\subseteq X$ and $r\in\reals$ and
        $r$ is an infimum of $\pointDistanceSet{d}{A}{x}$.
\end{definition}

% from §27 helper lemma 3: distance-to-set map is continuous
\begin{lemma}\label{point_set_distance_continuous}
    Suppose $d$ is a metric on $X$.
    Suppose $A\subseteq X$.
    Suppose $A\neq\emptyset$.
    Let $T = \metricTopology{d}{X}$.
    Then there exists $f$ such that
        $f$ is a function from $X$ to $\reals$ and
        for all $x\in X$ we have $f(x)$ is the distance from $x$ to $A$ in $X$ with respect to $d$ and
        $f$ is continuous from $X$ to $\reals$ with respect to $T$ and $\standardTopologyR$.
\end{lemma}
\begin{proof}
    Let $R = \{w\in X\times\reals \mid \text{there exists $x\in X$ such that there exists $r\in\reals$ such that
        $w = (x,r)$ and $r$ is the distance from $x$ to $A$ in $X$ with respect to $d$}\}$.
    For all $w$ such that $w\in R$ there exist $x,r$ such that $w = (x,r)$.
    Hence $R$ is a relation by \cref{relation}.
    Show that for all $x\in X$ there exists $r$ such that $x\mathrel{R} r$.
    \begin{subproof}
        Fix $x\in X$.
        Let $S = \pointDistanceSet{d}{A}{x}$.
        Thus $S\subseteq\reals$ by \cref{point_distance_set,subseteq,times_tuple_intro,metric_on,funs_type_apply}.
        Therefore $S\neq\emptyset$ by \cref{nonempty_inhabited,point_distance_set,emptyset}.
        Thus $0$ is a lower bound of $S$ and $S$ is bounded below
            by \cref{point_distance_set,subseteq,metric_on,metric_nonnegative,real_add_ident,real_lower_bound,real_bounded_below}.
        Take $r$ such that $r$ is an infimum of $S$ by \cref{real_infimum_exists}.
        Thus $r$ is the distance from $x$ to $A$ in $X$ with respect to $d$
            by \cref{real_infimum,real_lower_bound,point_set_distance}.
        Hence $(x,r)\in R$.
        Therefore $x\mathrel{R} r$ by \cref{relation}.
    \end{subproof}
    Take $f$ such that $f$ is a function on $X$ and for all $x\in X$ we have $x\mathrel{R} f(x)$
        by \cref{axiom_of_choice}.
    For all $x\in X$ we have $f(x)$ is the distance from $x$ to $A$ in $X$ with respect to $d$ and $f(x)\in\reals$
        by \cref{relation,pair_eq_iff,point_set_distance}.
    Hence $f$ is a function from $X$ to $\reals$ by \cref{funs_intro,funs}.
    Let $T_Y = \metricTopology{\standardMetricR}{\reals}$.
    Show that for all $x\in X$ we have
        for all $\epsilon\in\reals$ such that $0 < \epsilon$
        there exists $\delta\in\reals$ such that $0 < \delta$ and
        for all $y\in X$ if $d((x,y)) < \delta$ then
            $\apply{\standardMetricR}{(f(x),f(y))} < \epsilon$.
    \begin{subproof}
        Fix $x\in X$.
        Fix $\epsilon$ such that $\epsilon\in\reals$ and $0 < \epsilon$.
        Let $\delta = \epsilon$.
        Then $0 < \delta$.
        Show that for all $y\in X$ if $d((x,y)) < \delta$ then
            $\apply{\standardMetricR}{(f(x),f(y))} < \epsilon$.
        \begin{subproof}
            Fix $y\in X$.
            Assume $d((x,y)) < \delta$.
            Thus $d((x,y)) < \epsilon$ by \cref{real_lt_subst_right}.
            Let $S_x = \pointDistanceSet{d}{A}{x}$.
            Let $S_y = \pointDistanceSet{d}{A}{y}$.
                $f(x)$ is an infimum of $S_x$ and $f(y)$ is an infimum of $S_y$ by \cref{point_set_distance}.
                Therefore $f(x)$ is a lower bound of $S_x$ and $f(y)$ is a lower bound of $S_y$ by \cref{real_infimum}.
                Show that $f(x) \geq f(y) - d((x,y))$.
                \begin{subproof}
                    Let $b = f(y) - d((x,y))$.
                    $f(y)\in\reals$ by \cref{point_set_distance}.
                    $(x,y)\in X\times X$ by \cref{times_tuple_intro}.
                    $d\in\funs{X\times X}{\reals}$ by \cref{metric_on}.
                    Thus $d((x,y))\in\reals$ by \cref{funs_type_apply}.
                    $b = f(y) + \neg{(d((x,y)))}$ by \cref{real_sub}.
                    $\neg{(d((x,y)))}\in\reals$ by \cref{real_add_inverse}.
                    Thus $b\in\reals$ by \cref{real_add_closed}.
                    Show that for all $s\in S_x$ we have $b \leq s$.
                    \begin{subproof}
                        Fix $s\in S_x$.
                        Take $a\in A$ such that $s = d((x,a))$ by \cref{point_distance_set}.
                        Thus $a\in X$ by \cref{subseteq}.
                        Thus $d((y,x)) = d((x,y))$ by \cref{metric_on,metric_symmetric}.
                        Thus $d((y,x)) + d((x,a)) \geq d((y,a))$ by \cref{metric_on,metric_triangle_inequality}.
                        Hence $d((y,a)) \leq d((y,x)) + d((x,a))$.
                        Therefore $d((y,a)) \leq d((x,y)) + d((x,a))$.
                        $(y,a)\in X\times X$ and $(x,y)\in X\times X$ and $(x,a)\in X\times X$
                            by \cref{times_tuple_intro}.
                        Thus $d((y,a))\in\reals$ and $d((x,y))\in\reals$ and $d((x,a))\in\reals$
                            by \cref{funs_type_apply}.
                        Thus $d((x,y)) + d((x,a))\in\reals$ by \cref{real_add_closed}.
                        Thus $f(y) \leq d((y,a))$ by \cref{point_set_distance,real_lower_bound,point_distance_set}.
                        Thus $f(y) \leq d((x,y)) + d((x,a))$ by \cref{real_leq_trans}.
                        Let $u = d((x,y))$.
                        Let $v = d((x,a))$.
                        Then $f(y) \leq u + v$ by \cref{real_leq_trans}.
                        Hence $f(y) + \neg{u} \leq (u + v) + \neg{u}$ by \cref{real_leq_add}.
                        Therefore $u\in\reals$ and $v\in\reals$.
                        $u + \neg{u} = 0$ by \cref{real_add_inverse}.
                        $(u + v) + \neg{u} = v + (u + \neg{u})$ by \cref{real_add_assoc,real_add_comm}.
                        Thus $(u + v) + \neg{u} = v + 0$ by \cref{real_add_inverse}.
                        $0\in\reals$ by \cref{real_add_ident}.
                        $v + 0 = 0 + v$ by \cref{real_add_comm}.
                        $0 + v = v$ by \cref{real_add_ident}.
                        Then $(u + v) + \neg{u} = v$ by \cref{real_add_comm,real_add_ident}.
                        Hence $f(y) + \neg{u} \leq v$ by \cref{real_leq_trans}.
                        Therefore $f(y) - u \leq v$ by \cref{real_sub}.
                        Then $f(y) - d((x,y)) \leq d((x,a))$.
                        Hence $b \leq s$.
                    \end{subproof}
                    Therefore $b$ is a lower bound of $S_x$ by \cref{real_lower_bound,subseteq}.
                    Then $f(x) \geq b$ by \cref{real_infimum}.
                \end{subproof}
                Show that $f(y) \geq f(x) - d((x,y))$.
                \begin{subproof}
                    Let $b = f(x) - d((x,y))$.
                    $f(x)\in\reals$ and $(x,y)\in X\times X$ and $d\in\funs{X\times X}{\reals}$
                        by \cref{point_set_distance,times_tuple_intro,metric_on}.
                    Thus $d((x,y))\in\reals$ by \cref{funs_type_apply}.
                    $b = f(x) + \neg{(d((x,y)))}$ by \cref{real_sub}.
                    $\neg{(d((x,y)))}\in\reals$ by \cref{real_add_inverse}.
                    Thus $b\in\reals$ by \cref{real_add_closed}.
                    Show that for all $s\in S_y$ we have $b \leq s$.
                    \begin{subproof}
                        Fix $s\in S_y$.
                        Take $a\in A$ such that $s = d((y,a))$ by \cref{point_distance_set}.
                        Thus $a\in X$ by \cref{subseteq}.
                        Thus $d((x,y)) = d((y,x))$ by \cref{metric_on,metric_symmetric}.
                        Thus $d((x,y)) + d((y,a)) \geq d((x,a))$ by \cref{metric_on,metric_triangle_inequality}.
                        Hence $d((x,a)) \leq d((x,y)) + d((y,a))$.
                        $(x,a)\in X\times X$ and $(x,y)\in X\times X$ and $(y,a)\in X\times X$
                            by \cref{times_tuple_intro}.
                        Thus $d((x,a))\in\reals$ and $d((x,y))\in\reals$ and $d((y,a))\in\reals$
                            by \cref{funs_type_apply}.
                        Thus $d((x,y)) + d((y,a))\in\reals$ by \cref{real_add_closed}.
                        Thus $f(x) \leq d((x,a))$ by \cref{point_set_distance,real_lower_bound,point_distance_set}.
                        Thus $f(x) \leq d((x,y)) + d((y,a))$ by \cref{real_leq_trans}.
                        Let $u = d((x,y))$.
                        Let $v = d((y,a))$.
                        Then $f(x) \leq u + v$ by \cref{real_leq_trans}.
                        Hence $f(x) + \neg{u} \leq (u + v) + \neg{u}$ by \cref{real_leq_add}.
                        Therefore $u\in\reals$ and $v\in\reals$.
                        $u + \neg{u} = 0$ by \cref{real_add_inverse}.
                        $(u + v) + \neg{u} = v + (u + \neg{u})$ by \cref{real_add_assoc,real_add_comm}.
                        Thus $(u + v) + \neg{u} = v + 0$ by \cref{real_add_inverse}.
                        $0\in\reals$ by \cref{real_add_ident}.
                        $v + 0 = 0 + v$ by \cref{real_add_comm}.
                        $0 + v = v$ by \cref{real_add_ident}.
                        Then $(u + v) + \neg{u} = v$ by \cref{real_add_comm,real_add_ident}.
                        Hence $f(x) + \neg{u} \leq v$ by \cref{real_leq_trans}.
                        Therefore $f(x) - u \leq v$ by \cref{real_sub}.
                        Then $f(x) - d((x,y)) \leq d((y,a))$.
                        Hence $b \leq s$.
                    \end{subproof}
                    Therefore $b$ is a lower bound of $S_y$ by \cref{real_lower_bound,subseteq}.
                    Then $f(y) \geq b$ by \cref{real_infimum}.
                \end{subproof}
                Thus $f(x) - d((x,y)) \leq f(y)$.
                $f(x)\in\reals$ and $f(y)\in\reals$ by \cref{point_set_distance}.
                $(x,y)\in X\times X$ and $d\in\funs{X\times X}{\reals}$
                    by \cref{times_tuple_intro,metric_on}.
                Thus $d((x,y))\in\reals$ by \cref{funs_type_apply}.
                $\neg{(d((x,y)))}\in\reals$ by \cref{real_add_inverse}.
                Hence $f(x) + \neg{(d((x,y)))}\in\reals$ by \cref{real_add_closed}.
                Therefore $f(x) - d((x,y))\in\reals$ by \cref{real_sub}.
                Then $(f(x) - d((x,y))) + d((x,y)) \leq f(y) + d((x,y))$ by \cref{real_leq_add}.
                Let $s_0 = (f(x) - d((x,y))) + d((x,y))$.
                $f(x) - d((x,y)) = f(x) + \neg{(d((x,y)))}$ by \cref{real_sub}.
                Thus $s_0 = f(x) + (\neg{(d((x,y)))} + d((x,y)))$ by \cref{real_add_assoc}.
                $\neg{(d((x,y)))} + d((x,y)) = 0$ by \cref{real_add_inverse,real_add_comm}.
                Thus $s_0 = f(x) + 0$ by \cref{real_add_assoc}.
                $f(x) + 0 = 0 + f(x)$ by \cref{point_set_distance,real_add_ident,real_add_comm}.
                $0 + f(x) = f(x)$ by \cref{point_set_distance,real_add_ident}.
                Hence $s_0 = f(x)$ by \cref{real_add_comm,real_add_ident}.
                Then $f(x) \leq s_0$.
                Therefore $f(x) \leq f(y) + d((x,y))$ by \cref{real_leq_trans}.
                Then $f(y) - d((x,y)) \leq f(x)$.
                $f(y)\in\reals$ and $f(x)\in\reals$ by \cref{point_set_distance}.
                $(x,y)\in X\times X$ and $d\in\funs{X\times X}{\reals}$
                    by \cref{times_tuple_intro,metric_on}.
                Thus $d((x,y))\in\reals$ by \cref{funs_type_apply}.
                $\neg{(d((x,y)))}\in\reals$ by \cref{real_add_inverse}.
                Hence $f(y) + \neg{(d((x,y)))}\in\reals$ by \cref{real_add_closed}.
                Therefore $f(y) - d((x,y))\in\reals$ by \cref{real_sub}.
                Thus $(f(y) - d((x,y))) + d((x,y)) \leq f(x) + d((x,y))$ by \cref{real_leq_add}.
                Let $s_1 = (f(y) - d((x,y))) + d((x,y))$.
                $f(y) - d((x,y)) = f(y) + \neg{(d((x,y)))}$ by \cref{real_sub}.
                Thus $s_1 = f(y) + (\neg{(d((x,y)))} + d((x,y)))$ by \cref{real_add_assoc}.
                $\neg{(d((x,y)))} + d((x,y)) = 0$ by \cref{real_add_inverse,real_add_comm}.
                Thus $s_1 = f(y) + 0$ by \cref{real_add_assoc}.
                $f(y) + 0 = 0 + f(y)$ by \cref{point_set_distance,real_add_ident,real_add_comm}.
                $0 + f(y) = f(y)$ by \cref{point_set_distance,real_add_ident}.
                Hence $s_1 = f(y)$ by \cref{real_add_comm,real_add_ident}.
                Then $f(y) \leq s_1$.
                Therefore $f(y) \leq f(x) + d((x,y))$ by \cref{real_leq_trans}.
                $f(x)\in\reals$ and $f(y)\in\reals$ by \cref{point_set_distance}.
                $(x,y)\in X\times X$ and $d\in\funs{X\times X}{\reals}$
                    by \cref{times_tuple_intro,metric_on}.
                Thus $d((x,y))\in\reals$ by \cref{funs_type_apply}.
                $\neg{(f(y))}\in\reals$ by \cref{real_add_inverse}.
                Thus $f(x) + \neg{(f(y))}\in\reals$ by \cref{real_add_closed}.
                $f(x) - f(y)\in\reals$ by \cref{real_sub}.
                $\neg{(f(x) - f(y))}\in\reals$ by \cref{real_add_inverse}.
                Thus $f(y) + d((x,y))\in\reals$ by \cref{real_add_closed}.
                Thus $f(x) + \neg{(f(y))} \leq (f(y) + d((x,y))) + \neg{(f(y))}$ by \cref{real_leq_add}.
                $(f(y) + d((x,y))) + \neg{(f(y))} = f(y) + (d((x,y)) + \neg{(f(y))})$ by \cref{real_add_assoc}.
                $d((x,y)) + \neg{(f(y))} = \neg{(f(y))} + d((x,y))$ by \cref{real_add_comm}.
                $f(y) + (\neg{(f(y))} + d((x,y))) = (f(y) + \neg{(f(y))}) + d((x,y))$ by \cref{real_add_assoc}.
                $f(y) + \neg{(f(y))} = 0$ by \cref{real_add_inverse}.
                $0\in\reals$ by \cref{real_add_ident}.
                Therefore $(f(y) + d((x,y))) + \neg{(f(y))} = d((x,y))$ by \cref{real_add_ident,real_add_assoc,real_add_comm}.
                Then $f(x) + \neg{(f(y))} \leq d((x,y))$ by \cref{real_leq_trans}.
                Hence $f(x) - f(y) \leq d((x,y))$ by \cref{real_sub}.
                Thus $f(x) + d((x,y))\in\reals$ and $\neg{(f(x))}\in\reals$
                    by \cref{real_add_closed,real_add_inverse}.
                Thus $f(y) + \neg{(f(x))} \leq (f(x) + d((x,y))) + \neg{(f(x))}$ by \cref{real_leq_add}.
                $(f(x) + d((x,y))) + \neg{(f(x))} = f(x) + (d((x,y)) + \neg{(f(x))})$ by \cref{real_add_assoc}.
                $d((x,y)) + \neg{(f(x))} = \neg{(f(x))} + d((x,y))$ by \cref{real_add_comm}.
                $f(x) + (\neg{(f(x))} + d((x,y))) = (f(x) + \neg{(f(x))}) + d((x,y))$ by \cref{real_add_assoc}.
                $f(x) + \neg{(f(x))} = 0$ by \cref{real_add_inverse}.
                $0\in\reals$ by \cref{real_add_ident}.
                Therefore $(f(x) + d((x,y))) + \neg{(f(x))} = d((x,y))$ by \cref{real_add_ident,real_add_assoc,real_add_comm}.
                Then $f(y) + \neg{(f(x))} \leq d((x,y))$ by \cref{real_leq_trans}.
                Hence $f(y) - f(x) \leq d((x,y))$ by \cref{real_sub}.
                Thus $f(y) - f(x) = \neg{(f(x) - f(y))}$ by \cref{real_sub_swap_neg}.
                Then $\neg{(f(x) - f(y))} \leq d((x,y))$.
                Therefore $\abs{(f(x) - f(y))} \leq d((x,y))$ by \cref{real_abs_leq_if_pm_leq}.
                $f(x)\in\reals$ and $f(y)\in\reals$ by \cref{point_set_distance}.
                $\neg{(f(y))}\in\reals$ by \cref{real_add_inverse}.
                Thus $f(x) + \neg{(f(y))}\in\reals$ by \cref{real_add_closed}.
                $f(x) - f(y)\in\reals$ by \cref{real_sub}.
                $(f(x),f(y))\in\reals\times\reals$ by \cref{times_tuple_intro}.
                $\abs{(f(x) - f(y))}\in\reals$ by \cref{real_abs_in_reals}.
                Let $p = (f(x),f(y))$.
                $\fst{p} = f(x)$ and $\snd{p} = f(y)$ by \cref{fst_eq,snd_eq}.
                Then $\abs{(\fst{p} - \snd{p})} = \abs{(f(x) - f(y))}$.
                Hence $(p,\abs{(f(x) - f(y))})\in\standardMetricR$ by \cref{standard_metric_r}.
                Therefore $\apply{\standardMetricR}{(f(x),f(y))} = \abs{(f(x) - f(y))}$
                    by \cref{standard_metric_is_metric,metric_on,funs_elim,function_apply_intro}.
                Then $\apply{\standardMetricR}{(f(x),f(y))} \leq d((x,y))$.
                $(x,y)\in X\times X$ by \cref{times_tuple_intro}.
                $d\in\funs{X\times X}{\reals}$ by \cref{metric_on}.
                Thus $d((x,y))\in\reals$ and $\apply{\standardMetricR}{(f(x),f(y))}\in\reals$
                    by \cref{funs_type_apply}.
                Hence $\apply{\standardMetricR}{(f(x),f(y))} < d((x,y))$ or
                    $\apply{\standardMetricR}{(f(x),f(y))} = d((x,y))$ by \cref{real_leq_split}.
                \begin{byCase}
                \caseOf{$\apply{\standardMetricR}{(f(x),f(y))} < d((x,y))$.}
                    Therefore $\apply{\standardMetricR}{(f(x),f(y))} < \epsilon$ by \cref{real_lt_trans}.
                \caseOf{$\apply{\standardMetricR}{(f(x),f(y))} = d((x,y))$.}
                    Then $\apply{\standardMetricR}{(f(x),f(y))} < \epsilon$ by \cref{real_lt_subst_left}.
                \end{byCase}
        \end{subproof}
    \end{subproof}
    Hence $f$ is continuous from $X$ to $\reals$ with respect to $T$ and $T_Y$
        by \cref{standard_metric_is_metric,metric_continuity_eps_delta}.
    Therefore $f$ is continuous from $X$ to $\reals$ with respect to $T$ and $\standardTopologyR$
        by \cref{standard_metric_topology,continuous_eq_codomain}.
\end{proof}

from §27 Item 6: Lebesgue number lemma
\begin{lemma}\label{lebesgue_number_lemma}
    Suppose $d$ is a metric on $X$.
    Let $T = \metricTopology{d}{X}$.
    Suppose $A$ is an open covering of $X$ with respect to $T$.
    Suppose $X\neq\emptyset$.
    Suppose $X$ is compact with respect to $T$.
    Then there exists $\delta\in\reals$ such that $0 < \delta$ and
        for all $B$ such that $B\subseteq X$ and
            for all $x,y\in B$ we have $d((x,y)) < \delta$
        there exists $U\in A$ such that $B\subseteq U$.
\end{lemma}
\begin{proof}
    Omitted. % do not touch
\end{proof}

% from §27 Item 7: Lebesgue number
\begin{definition}\label{lebesgue_number}
    $\delta$ is a Lebesgue number for $A$ in $X$ with respect to $d$ iff
        $d$ is a metric on $X$ and
        $A$ is an open covering of $X$ with respect to $\metricTopology{d}{X}$ and
        $X\neq\emptyset$ and
        $\delta\in\reals$ and $0 < \delta$ and
        for all $B$ such that $B\subseteq X$ and
            for all $x,y\in B$ we have $d((x,y)) < \delta$
        there exists $U\in A$ such that $B\subseteq U$.
\end{definition}

% from §27 Item 8: uniformly continuous function
\begin{definition}\label{uniformly_continuous}
    $f$ is uniformly continuous from $X$ to $Y$ with respect to $d_X$ and $d_Y$ iff
        $f\in\funs{X}{Y}$ and
        $d_X$ is a metric on $X$ and
        $d_Y$ is a metric on $Y$ and
        for all $\epsilon\in\reals$ such that $0 < \epsilon$
            there exists $\delta\in\reals$ such that $0 < \delta$ and
                for all $x_0,x_1\in X$ such that $d_X((x_0,x_1)) < \delta$
                    we have $d_Y((\apply{f}{x_0},\apply{f}{x_1})) < \epsilon$.
\end{definition}

% from §27 Item 9: uniform continuity theorem
\begin{theorem}\label{uniform_continuity_theorem}
    Suppose $d_X$ is a metric on $X$.
    Suppose $d_Y$ is a metric on $Y$.
    Let $T_X = \metricTopology{d_X}{X}$.
    Let $T_Y = \metricTopology{d_Y}{Y}$.
    Suppose $f$ is continuous from $X$ to $Y$ with respect to $T_X$ and $T_Y$.
    Suppose $X$ is compact with respect to $T_X$.
    Then $f$ is uniformly continuous from $X$ to $Y$ with respect to $d_X$ and $d_Y$.
\end{theorem}
\begin{proof}
    $f\in\funs{X}{Y}$ by \cref{continuous,funs}.
    Show that for all $\epsilon\in\reals$ such that $0 < \epsilon$
        there exists $\delta\in\reals$ such that $0 < \delta$ and
            for all $x_0,x_1\in X$ such that $d_X((x_0,x_1)) < \delta$
                we have $d_Y((\apply{f}{x_0},\apply{f}{x_1})) < \epsilon$.
    \begin{subproof}
        Fix $\epsilon\in\reals$.
        Assume $0 < \epsilon$.
        Let $\epsilon_0 = \epsilon / (1 + 1)$.
        $1\in\reals$ by \cref{real_mul_ident}.
        $0\in\reals$ by \cref{real_add_ident}.
        $0 < 1$ by \cref{real_zero_lt_one}.
        Hence $1 < 1 + 1$ by \cref{real_lt_add,real_add_ident}.
        Therefore $(1 + 1)\in\reals$ and $0 < 1 + 1$ by \cref{real_add_closed,real_lt_trans}.
        Then $(1 + 1)\neq 0$ by \cref{real_pos_nonzero}.
        Hence $\inv{(1 + 1)}\in\reals$ and $\inv{(1 + 1)} > 0$ by \cref{real_mul_inverse,real_inv_pos}.
        $\epsilon_0 = \epsilon \rmul \inv{(1 + 1)}$ and $\epsilon_0\in\reals$ by \cref{real_div,real_mul_closed}.
        Hence $0 < \epsilon_0$ by \cref{real_div,real_mul_zero,real_lt_mul_pos,real_lt_subst_left}.
        \begin{byCase}
        \caseOf{$X = \emptyset$.}
            Let $\delta = \epsilon_0$.
            $\delta\in\reals$ and $0 < \delta$.
            For all $x_0,x_1$ such that $x_0\in X$ and $x_1\in X$ and $d_X((x_0,x_1)) < \delta$
                we have $d_Y((\apply{f}{x_0},\apply{f}{x_1})) < \epsilon$.
            Hence there exists $\delta\in\reals$ such that $0 < \delta$ and
                for all $x_0,x_1\in X$ such that $d_X((x_0,x_1)) < \delta$
                    we have $d_Y((\apply{f}{x_0},\apply{f}{x_1})) < \epsilon$.
        \caseOf{$X \neq \emptyset$.}
            Let $R = \{w\in X\times\reals \mid \text{$0 < \snd{w}$ and
                for all $y\in X$ if $d_X((\fst{w},y)) < \snd{w}$ then
                $d_Y((\apply{f}{\fst{w}},\apply{f}{y})) < \epsilon_0$}\}$.
            Hence $R$ is a relation by \cref{times_elem_is_tuple,relation}.
            For all $x\in X$ there exists $\delta\in\reals$ such that $x\mathrel{R}\delta$
                by \cref{metric_continuity_eps_delta,times_tuple_intro,fst_eq,snd_eq,relation}.
            Take $g$ such that $g$ is a function on $X$ and for all $x\in X$ we have $x\mathrel{R}\apply{g}{x}$
                by \cref{choice_function}.
            Let $A = \{U\in\pow{X} \mid \exists x\in X. U = \metricBall{d_X}{X}{x}{\apply{g}{x}}\}$.

            Show that for all $x\in X$ we have $x\in\unions{A}$.
            \begin{subproof}
                Fix $x\in X$.
                $x\mathrel{R}\apply{g}{x}$.
                Hence $0 < \apply{g}{x}$.
                Therefore $x\in\unions{A}$ by \cref{metric_zero_characterization,metric_on,real_lt_subst_left,metric_ball,subseteq,powerset_intro,unions_iff}.
            \end{subproof}
            Show that for all $U$ such that $U\in A$ we have $U$ is open in $X$ with respect to $T_X$.
            \begin{subproof}
                Fix $U$.
                Assume $U\in A$.
                Take $x\in X$ such that $U = \metricBall{d_X}{X}{x}{\apply{g}{x}}$.
                $x\mathrel{R}\apply{g}{x}$.
                Hence $(x,\apply{g}{x})\in R$.
                Therefore $(x,\apply{g}{x})\in X\times\reals$.
                Then $\apply{g}{x}\in\reals$ by \cref{times_tuple_elim}.
                Thus $0 < \apply{g}{x}$.
                Hence $U\in\metricBasis{d_X}{X}$ by \cref{metric_ball,subseteq,powerset_intro,metric_basis}.
                Therefore $U\in T_X$ by \cref{metric_basis_is_basis,basis_elem_open_in_generated,basis_for_topology,metric_topology}.
                Hence $U$ is open in $X$ with respect to $T_X$
                    by \cref{metric_basis_is_basis,basis_topology_is_topology,metric_topology,open_in}.
            \end{subproof}
            Therefore $A$ is an open covering of $X$ with respect to $T_X$ by \cref{subseteq,covering,open_covering}.

            Take $\delta\in\reals$ such that $0 < \delta$ and
                for all $B$ such that $B\subseteq X$ and
                    for all $x,y\in B$ we have $d_X((x,y)) < \delta$
                there exists $U\in A$ such that $B\subseteq U$
                by \cref{lebesgue_number_lemma}.

            Show that for all $x_0,x_1$ such that $x_0\in X$ and $x_1\in X$ and $d_X((x_0,x_1)) < \delta$
                we have $d_Y((\apply{f}{x_0},\apply{f}{x_1})) < \epsilon$.
            \begin{subproof}
                Fix $x_0$.
                Fix $x_1$.
                Assume $x_0\in X$ and $x_1\in X$ and $d_X((x_0,x_1)) < \delta$.
                Let $B = \{x_0,x_1\}$.
                $B\subseteq X$ by \cref{upair_iff,subseteq}.
                Show that for all $x\in B$ we have for all $y\in B$ we have $d_X((x,y)) < \delta$.
                \begin{subproof}
                    Fix $x$.
                    Assume $x\in B$.
                    Fix $y$.
                    Assume $y\in B$.
                    $x = x_0$ or $x = x_1$ by \cref{upair_iff}.
                    Then $y = x_0$ or $y = x_1$ by \cref{upair_iff}.
                    \begin{byCase}
                    \caseOf{$x = x_0$.}
                        \begin{byCase}
                        \caseOf{$y = x_0$.}
                            Thus $d_X((x,y)) < \delta$ by \cref{metric_zero_characterization,metric_on,real_lt_subst_left}.
                        \caseOf{$y \neq x_0$.}
                            Then $y = x_1$.
                            $d_X((x,y)) < \delta$.
                        \end{byCase}
                    \caseOf{$x \neq x_0$.}
                        Then $x = x_1$.
                        \begin{byCase}
                        \caseOf{$y = x_0$.}
                            $d_X((x_1,x_0)) = d_X((x_0,x_1))$ by \cref{metric_symmetric,metric_on}.
                            $d_X((x,y)) < \delta$.
                        \caseOf{$y \neq x_0$.}
                            Then $y = x_1$.
                            Therefore $d_X((x,y)) < \delta$ by \cref{metric_zero_characterization,metric_on,real_lt_subst_left}.
                        \end{byCase}
                    \end{byCase}
                \end{subproof}
                Take $U\in A$ such that $B\subseteq U$.
                Take $x\in X$ such that $U = \metricBall{d_X}{X}{x}{\apply{g}{x}}$.
                Thus $d_X((x,x_0)) < \apply{g}{x}$ and
                    $d_X((x,x_1)) < \apply{g}{x}$ by \cref{upair_intro_left,upair_intro_right,subseteq,metric_ball}.
                $x\mathrel{R}\apply{g}{x}$.
                Hence $(x,\apply{g}{x})\in R$.
                Therefore $(x,\apply{g}{x})\in X\times\reals$.
                For all $y\in X$ if $d_X((x,y)) < \apply{g}{x}$ then
                    $d_Y((\apply{f}{x},\apply{f}{y})) < \epsilon_0$
                    by \cref{fst_eq,snd_eq}.
                Therefore $d_Y((\apply{f}{x},\apply{f}{x_0})) < \epsilon_0$ and
                    $d_Y((\apply{f}{x},\apply{f}{x_1})) < \epsilon_0$.
                $\apply{f}{x_0}\in Y$ and $\apply{f}{x}\in Y$ and $\apply{f}{x_1}\in Y$ by \cref{funs_type_apply}.
                Hence $d_Y((\apply{f}{x_0},\apply{f}{x})) < \epsilon_0$ by \cref{metric_symmetric,metric_on,real_lt_subst_left}.
                Let $s = d_Y((\apply{f}{x_0},\apply{f}{x})) + d_Y((\apply{f}{x},\apply{f}{x_1}))$.
                Then $d_Y((\apply{f}{x_0},\apply{f}{x_1})) \leq s$ by \cref{metric_on,metric_triangle_inequality}.
                $d_Y\in\funs{Y\times Y}{\reals}$ by \cref{metric_on}.
                $(\apply{f}{x_0},\apply{f}{x})\in Y\times Y$ and
                    $(\apply{f}{x},\apply{f}{x_1})\in Y\times Y$ by \cref{times_tuple_intro}.
                Hence $d_Y((\apply{f}{x_0},\apply{f}{x}))\in\reals$ and
                    $d_Y((\apply{f}{x},\apply{f}{x_1}))\in\reals$ by \cref{funs_type_apply}.
                Therefore $s < \epsilon_0 + \epsilon_0$ by \cref{real_lt_add_two}.
                Hence $d_Y((\apply{f}{x_0},\apply{f}{x_1})) < \epsilon_0 + \epsilon_0$
                    by \cref{times_tuple_intro,funs_type_apply,real_add_closed,real_leq_split,real_lt_trans,real_lt_subst_left}.

                $\epsilon_0 = \epsilon \rmul \inv{(1 + 1)}$ by \cref{real_div}.
                Hence $\epsilon_0 + \epsilon_0 =
                    (\epsilon \rmul \inv{(1 + 1)}) + (\epsilon \rmul \inv{(1 + 1)})$.
                $(\epsilon + \epsilon) \rmul \inv{(1 + 1)}
                    = (\epsilon \rmul \inv{(1 + 1)}) + (\epsilon \rmul \inv{(1 + 1)})$
                    by \cref{real_right_distrib}.
                Therefore $\epsilon_0 + \epsilon_0 = (\epsilon + \epsilon) \rmul \inv{(1 + 1)}$.
                $\epsilon \rmul (1 + 1) = (\epsilon \rmul 1) + (\epsilon \rmul 1)$ by \cref{real_distrib}.
                Hence $\epsilon \rmul 1 = \epsilon$ by \cref{real_mul_ident,real_mul_comm}.
                Therefore $\epsilon \rmul (1 + 1) = \epsilon + \epsilon$.
                Then $(\epsilon + \epsilon) \rmul \inv{(1 + 1)} = (\epsilon \rmul (1 + 1)) \rmul \inv{(1 + 1)}$.
                $(\epsilon \rmul (1 + 1)) \rmul \inv{(1 + 1)}
                    = \epsilon \rmul ((1 + 1) \rmul \inv{(1 + 1)})$ by \cref{real_mul_assoc}.
                $(1 + 1) \rmul \inv{(1 + 1)} = 1$ by \cref{real_mul_inverse}.
                Hence $\epsilon_0 + \epsilon_0 = \epsilon \rmul 1$.
                Therefore $\epsilon_0 + \epsilon_0 = \epsilon$ by \cref{real_mul_ident}.
                Therefore $d_Y((\apply{f}{x_0},\apply{f}{x_1})) < \epsilon$ by \cref{real_lt_subst_right}.
            \end{subproof}
            Hence there exists $\delta\in\reals$ such that $0 < \delta$ and
                for all $x_0,x_1\in X$ such that $d_X((x_0,x_1)) < \delta$
                    we have $d_Y((\apply{f}{x_0},\apply{f}{x_1})) < \epsilon$.
        \end{byCase}
    \end{subproof}
    Therefore $f$ is uniformly continuous from $X$ to $Y$ with respect to $d_X$ and $d_Y$ by \cref{uniformly_continuous}.
\end{proof}
% from §27 Item 10: isolated point
\begin{definition}\label{isolated_point}
    $x$ is an isolated point of $X$ with respect to $T$ iff
        $x\in X$ and $\{x\}$ is open in $X$ with respect to $T$.
\end{definition}

% from §27 Item 11: compact Hausdorff without isolated points is uncountable
\begin{theorem}\label{compact_hausdorff_no_isolated_uncountable}
    Assume $T$ is a topology on $X$.
    Assume $X$ is compact with respect to $T$.
    Assume $X$ is Hausdorff with respect to $T$.
    Assume $X\neq\emptyset$.
    Suppose for all $x\in X$ we have $x$ is not an isolated point of $X$ with respect to $T$.
    Then $X$ is not countable.
\end{theorem}
\begin{proof}
    Omitted. % do not touch
\end{proof}

% from §27 Item 12: closed interval in reals is uncountable
\begin{corollary}\label{real_closed_interval_uncountable}
    Suppose $a\in\reals$.
    Suppose $b\in\reals$.
    Suppose $a < b$.
    Let $R = \realOrderRel$.
    Let $T = \realOrderTopology$.
    Let $Y = \intervalclosedrel{R}{\reals}{a}{b}$.
    Let $T_Y = \subspaceTopology{T}{Y}$.
    Then $Y$ is not countable.
\end{corollary}
\begin{proof}
    $R$ is a simple order relation on $\reals$ and
        $\reals$ is a linear continuum with respect to $\realOrderRel$
        by \cref{real_order_relation_simple,reals_linear_continuum}.
    Hence $\reals$ has more than one element by \cref{linear_continuum}.
    Therefore $T$ is the order topology on $\reals$ with respect to $R$ and $T$ is a topology on $\reals$
        by \cref{real_order_topology,order_topology,order_basis_is_basis,basis_topology_is_topology}.

    $a\mathrel{R}b$ by \cref{real_order_relation_intro}.
    Hence $a\neq b$ by \cref{real_lt_subst_right,real_lt_irreflexive}.
    Therefore $a\in Y$ and $b\in Y$ by \cref{intervalclosed_rel}.
    Hence $Y$ has more than one element by \cref{more_than_one_element}.
    Therefore $Y\neq\emptyset$ by \cref{emptyset}.

    $Y\subseteq\reals$ by \cref{intervalclosed_rel,subseteq}.

    Hence $\reals$ is Hausdorff with respect to $T$ and $Y$ is Hausdorff with respect to $T_Y$
        by \cref{order_topology_hausdorff,subspace_hausdorff}.
    $T_Y$ is a topology on $Y$ by \cref{subspace_topology_is_topology}.

    Therefore $Y$ is compact with respect to $T_Y$ and $Y$ is connected with respect to $T_Y$
        by \cref{real_closed_interval_compact,linear_continuum_connected}.

    Show that for all $x\in Y$ we have $x$ is not an isolated point of $Y$ with respect to $T_Y$.
    \begin{subproof}
        Fix $x\in Y$.
        Suppose not.
        Hence $\{x\}$ is open in $Y$ with respect to $T_Y$ by \cref{isolated_point}.
        Take $y$ such that $y\neq x$ by \cref{more_than_one_element_has_other}.
        Then $\{x\}\subseteq Y$ by \cref{singleton_subset_intro}.
        Therefore $\{x\}$ is closed in $Y$ with respect to $T_Y$
            by \cref{singleton_subset_intro,pair_is_finite,upair_intro_left,subseteq_finite_is_finite,finite_point_set_closed_hausdorff}.
        Hence $\{x\} = \emptyset$ or $\{x\} = Y$ by \cref{connected_iff_clopen}.
        Therefore $\{x\} = Y$ by \cref{singleton_iff,emptyset}.
        Then $y = x$ by \cref{subseteq,singleton_iff}.
        Contradiction.
    \end{subproof}

    Therefore $Y$ is not countable by \cref{compact_hausdorff_no_isolated_uncountable}.
\end{proof}


\section{§28 Limit Point Compactness}

% from §28 Item 1: limit point compact
\begin{definition}\label{limit_point_compact}
    $X$ is limit point compact with respect to $T$ iff
        $T$ is a topology on $X$ and
        for all $A$ such that $A\subseteq X$ and $A$ is infinite
            there exists $x\in X$ such that $x$ is a limit point of $A$ in $X$ with respect to $T$.
\end{definition}

% from §28 Item 2: compact implies limit point compact
\begin{theorem}\label{compact_implies_limit_point_compact}
    Assume $T$ is a topology on $X$.
    Suppose $X$ is compact with respect to $T$.
    Then $X$ is limit point compact with respect to $T$.
\end{theorem}
\begin{proof}
    Show that for all $A$ such that $A\subseteq X$ and $A$ is infinite
        there exists $x\in X$ such that $x$ is a limit point of $A$ in $X$ with respect to $T$.
    \begin{subproof}
        Fix $A$.
        Assume $A\subseteq X$ and $A$ is infinite.
        Suppose not.
        Then for all $x\in X$ we have $x$ is not a limit point of $A$ in $X$ with respect to $T$.

        Let $C = \{U\in\pow{X} \mid \text{there exists $a\in A$ such that
            $U$ is a neighborhood of $a$ in $X$ with respect to $T$ and
            $U\inter (A\setminus\{a\}) = \emptyset$}\}$.

        For all $U$ such that $U\in C$ we have $U$ is open in $X$ with respect to $T$
            by \cref{neighborhood}.

        For all $a$ such that $a\in A$ we have $a\in\unions{C}$
            by \cref{subseteq,limit_point,intersects,neighborhood,open_in,topology_on,unions_iff}.
        Therefore $A\subseteq\unions{C}$ by \cref{subseteq}.

        For all $x$ such that $x\in\limitPoints{A}{X}{T}$ we have $x\in\emptyset$ by \cref{limit_points}.
        Hence $X\setminus A$ is open in $X$ with respect to $T$
            by \cref{subseteq,emptyset_subseteq,subseteq_transitive,closed_iff_contains_limit_points,closed_in}.

        Let $D = C\union\{X\setminus A\}$.

        For all $U$ such that $U\in D$ we have $U$ is open in $X$ with respect to $T$
            by \cref{union_iff,singleton_iff,closed_in}.
        For all $x$ such that $x\in X$ we have $x\in\unions{D}$
            by \cref{subseteq,union_iff,unions_iff,setminus_intro,singleton_intro}.
        Hence $D$ is an open covering of $X$ with respect to $T$ by \cref{subseteq,covering,open_covering}.

        Take $B$ such that $B\subseteq D$ and $B$ is finite and $B$ is a covering of $X$ by \cref{compact}.
        Let $B_A = \{U\in B \mid U\neq X\setminus A\}$.
        $B_A\subseteq B$ and $B_A$ is finite by \cref{subseteq,subseteq_finite_is_finite}.

        Show that for all $a$ such that $a\in A$ we have $a\in\unions{B_A}$.
        \begin{subproof}
            Fix $a$.
            Assume $a\in A$.
            Take $U\in B$ such that $a\in U$ by \cref{subseteq,covering,unions_iff}.
            $U\neq X\setminus A$.
            Hence $a\in\unions{B_A}$ by \cref{unions_iff}.
        \end{subproof}
        Therefore $A\subseteq\unions{B_A}$ by \cref{subseteq}.

        Let $h = \{(U, U\inter A) \mid U\in B_A\}$.
        For all $U,V_1,V_2$ such that $(U,V_1)\in h$ and $(U,V_2)\in h$ we have $V_1 = V_2$
            by \cref{pair_eq_iff}.
        Hence $h$ is right-unique by \cref{rightunique}.
        Therefore $h$ is a function.
        For all $U$ such that $U\in B_A$ we have $U\in\dom{h}$ by \cref{pair_eq_iff,dom_intro}.
        Hence $B_A\subseteq\dom{h}$ by \cref{subseteq}.
        For all $U$ such that $U\in\dom{h}$ we have $U\in B_A$ by \cref{dom_iff,pair_eq_iff}.
        Therefore $\dom{h}\subseteq B_A$ by \cref{subseteq}.
        Hence $\dom{h} = B_A$ by \cref{subseteq_antisymmetric}.
        Therefore $h$ is a function on $B_A$.

        Let $F = \img{h}{B_A}$.

        Show that for all $V$ such that $V\in F$ we have $V$ is finite.
        \begin{subproof}
            Fix $V$.
            Assume $V\in F$.
            Take $U\in B_A$ such that $(U,V)\in h$ by \cref{img_iff}.
            Take $U_0$ such that $(U,V) = (U_0,U_0\inter A)$ by \cref{pair_eq_iff}.
            Take $a$ such that
                $U$ is a neighborhood of $a$ in $X$ with respect to $T$ and
                $U\inter (A\setminus\{a\}) = \emptyset$
                by \cref{subseteq,union_iff,singleton_iff}.
            Show that for all $x$ such that $x\in V$ we have $x\in\{a\}$.
            \begin{subproof}
                Fix $x$.
                Assume $x\in V$.
                $x\in U$ and $x\in A$ by \cref{pair_eq_iff,inter}.
                \begin{byCase}
                \caseOf{$x = a$.}
                    Then $x\in\{a\}$ by \cref{singleton_intro}.
                \caseOf{$x\neq a$.}
                    Hence $x\notin\{a\}$ and $x\in A\setminus\{a\}$ by \cref{singleton_iff,setminus_intro}.
                    Therefore $x\in U\inter (A\setminus\{a\})$ by \cref{inter_intro}.
                    Contradiction by \cref{emptyset}.
                \end{byCase}
            \end{subproof}
            Hence $V$ is finite
                by \cref{subseteq,pair_is_finite,upair_intro_left,singleton_subset_intro,subseteq_finite_is_finite}.
        \end{subproof}

        For all $a$ such that $a\in A$ we have $a\in\unions{F}$
            by \cref{subseteq,unions_iff,pair_eq_iff,img_iff,inter_intro}.
        Therefore $A\subseteq\unions{F}$ by \cref{subseteq}.
        For all $x$ such that $x\in\unions{F}$ we have $x\in A$
            by \cref{unions_iff,img_iff,pair_eq_iff,inter}.
        Therefore $\unions{F}\subseteq A$ by \cref{subseteq}.
        Hence $\unions{F} = A$ by \cref{subseteq_antisymmetric}.

        Therefore $A$ is finite by \cref{finite_image_function,finite_union_of_finite_family,subseteq_antisymmetric}.
        Contradiction.
    \end{subproof}
    Hence $X$ is limit point compact with respect to $T$ by \cref{limit_point_compact}.
\end{proof}

% from §28 helper lemma 1: naturalsPlus is infinite
\begin{lemma}\label{naturalsplus_infinite}
    $\naturalsPlus$ is infinite.
\end{lemma}
\begin{proof}
    Suppose not.
    Then $\naturalsPlus$ is finite.
    Take $m\in\naturalsPlus$ such that for all $a\in\naturalsPlus$ we have $a \leq m$
        by \cref{real_one_in_naturals,emptyset,subseteq_refl,finite_naturals_maximum}.
    Hence $m + 1 \leq m$ and $m\in\reals$ and $m + 1\in\reals$
        by \cref{real_naturals_closed_succ,real_naturals_subset_reals,subseteq}.
    Therefore $m < m + 1$
        by \cref{real_add_ident,real_one_in_naturals,real_naturals_subset_reals,subseteq,real_zero_lt_one,real_lt_add,real_add_comm}.
    Contradiction by \cref{real_leq_antisym,real_lt_subst_right,real_lt_irreflexive}.
\end{proof}

% from §28 Item 3: limit point compact does not imply compact
\begin{theorem}\label{limit_point_compact_not_compact}
    There exists $X,T$ such that
        $X$ is limit point compact with respect to $T$ and
        $X$ is not compact with respect to $T$.
\end{theorem}
\begin{proof}
    Let $Y = \{0,1\}$.
    Let $X = \naturalsPlus\times Y$.
    Let $T_0 = \discrete{\naturalsPlus}$.
    Let $T_1 = \{\emptyset, Y\}$.
    Let $T = \productTopology{T_0}{T_1}{\naturalsPlus}{Y}$.

    $T_0 = \pow{\naturalsPlus}$ and $T_0\subseteq\pow{\naturalsPlus}$
        by \cref{discrete_topology,subseteq_refl}.
    $\emptyset\in T_0$ and $\naturalsPlus\in T_0$ by \cref{discrete_topology,powerset_bottom,powerset_top}.
    Hence $T_0$ is closed under arbitrary unions by \cref{subseteq_transitive,powerset_closed_unions,discrete_topology}.
    Therefore $T_0$ is closed under binary intersections by \cref{discrete_topology,subseteq,inter_powerset}.
    Hence $T_0$ is a topology on $\naturalsPlus$ by \cref{topology_on}.

    Therefore $T_1\subseteq\pow{Y}$ by \cref{upair_iff,emptyset_subseteq,subseteq_refl,powerset_intro,subseteq}.
    $\emptyset\in T_1$ and $Y\in T_1$ by \cref{upair_intro_left,upair_intro_right}.
    Show that $T_1$ is closed under arbitrary unions.
    \begin{subproof}
        Fix $M$ such that $M\subseteq T_1$.
        \begin{byCase}
        \caseOf{$Y\in M$.}
            Hence $\unions{M}\subseteq Y$
                by \cref{subseteq_transitive,powerset_closed_unions,pow_iff}.
            Therefore $Y\subseteq\unions{M}$ by \cref{unions_iff,subseteq}.
            Hence $\unions{M}\in T_1$ by \cref{subseteq_antisymmetric,upair_intro_right}.
        \caseOf{$Y\notin M$.}
            For all $U$ such that $U\in M$ we have $U\subseteq\emptyset$
                by \cref{subseteq,upair_iff,subseteq_refl}.
            Therefore $\unions{M}\in T_1$ by \cref{unions_family,subseteq_emptyset_iff,upair_intro_left}.
        \end{byCase}
    \end{subproof}
    Hence $T_1$ is closed under binary intersections
        by \cref{upair_iff,inter_comm,inter_emptyset,inter_idempotent,upair_intro_left,upair_intro_right}.
    Therefore $T_1$ is a topology on $Y$ by \cref{topology_on}.

    $\productBasis{T_0}{T_1}{\naturalsPlus}{Y}$ is a basis on $X$ by \cref{product_basis_is_basis}.
    Hence $T$ is a topology on $X$ by \cref{basis_topology_is_topology,product_topology}.

    Show that for all $A$ such that $A\subseteq X$ and $A$ is infinite
        there exists $x\in X$ such that $x$ is a limit point of $A$ in $X$ with respect to $T$.
        \begin{subproof}
            Fix $A$ such that $A\subseteq X$ and $A$ is infinite.
            Take $a$ such that $a\in A$ by \cref{emptyset_is_finite,nonempty_inhabited}.
            Take $n,y$ such that $a = (n,y)$ and $n\in\naturalsPlus$ and $y\in Y$ by \cref{subseteq,times_elem_is_tuple}.

        Show that there exists $y'$ such that $y'\in Y$ and $y'\neq y$.
        \begin{subproof}
            $y = 0$ or $y = 1$ by \cref{upair_iff}.
            \begin{byCase}
            \caseOf{$y = 0$.}
                Let $y' = 1$.
                Hence there exists $y'$ such that $y'\in Y$ and $y'\neq y$.
            \caseOf{$y = 1$.}
                Let $y' = 0$.
                Therefore there exists $y'$ such that $y'\in Y$ and $y'\neq y$.
            \end{byCase}
        \end{subproof}
        Take $y'$ such that $y'\in Y$ and $y'\neq y$.
        Let $x = (n,y')$.
        $x\in X$ by \cref{times_tuple_intro}.

        Show that for all $U$ such that $U$ is a neighborhood of $x$ in $X$ with respect to $T$
            we have $U$ intersects $A\setminus\{x\}$.
        \begin{subproof}
            Fix $U$ such that $U$ is a neighborhood of $x$ in $X$ with respect to $T$.
            $x\in U$ and $U\in\basisTopology{\productBasis{T_0}{T_1}{\naturalsPlus}{Y}}{X}$
                by \cref{neighborhood,open_in,product_topology}.
            Take $B\in\productBasis{T_0}{T_1}{\naturalsPlus}{Y}$ such that $x\in B$ and $B\subseteq U$
                by \cref{topology_generated_by_basis}.
            Take $U_0,V_0$ such that $V_0\in T_1$ and $B = U_0\times V_0$
                by \cref{product_basis}.
            Then $n\in U_0$ and $y'\in V_0$ by \cref{pair_eq_iff,times_tuple_elim}.
            Hence $V_0 = Y$ by \cref{upair_iff,emptyset}.
            Therefore $B = U_0\times Y$.
            Then $(n,y)\in B$ by \cref{times_tuple_intro}.
            Hence $a\in U$ by \cref{subseteq}.
            Therefore $a\in A\setminus\{x\}$ by \cref{pair_eq_iff,setminus_intro,singleton_iff}.
            Hence $U$ intersects $A\setminus\{x\}$ by \cref{inter_intro,emptyset,intersects}.
        \end{subproof}
        Therefore $x$ is a limit point of $A$ in $X$ with respect to $T$ by \cref{limit_point}.
    \end{subproof}
    Hence $X$ is limit point compact with respect to $T$ by \cref{limit_point_compact}.

    Suppose not.
    Then $X$ is compact with respect to $T$.
    Let $\pi = \piOne{\naturalsPlus}{Y}$.
    $\pi$ is continuous from $X$ to $\naturalsPlus$ with respect to $T$ and $T_0$
        by \cref{projection_one_continuous}.
    Hence $\img{\pi}{X}\subseteq\naturalsPlus$
        by \cref{img_iff,times_elem_is_tuple,projection_first,pair_eq_iff,subseteq}.
    Therefore $\naturalsPlus\subseteq\img{\pi}{X}$
        by \cref{upair_intro_left,times_tuple_intro,projection_first,img_iff,subseteq}.
    Hence $\img{\pi}{X} = \naturalsPlus$ by \cref{subseteq_antisymmetric}.
    Let $T_A = \subspaceTopology{T_0}{\naturalsPlus}$.
    Therefore $\naturalsPlus$ is compact with respect to $T_A$ by \cref{continuous_image_compact}.
    Hence $T_A = T_0$
        by \cref{subspace_topology,discrete_topology,pow_iff,inter_absorb_supseteq_left,subseteq,subseteq_antisymmetric}.
    Therefore $\naturalsPlus$ is compact with respect to $T_0$.

    Let $C = \{\cons{n}{\emptyset} \mid n\in\naturalsPlus\}$.
    Therefore $C$ is an open covering of $\naturalsPlus$ with respect to $T_0$
        by \cref{emptyset_subseteq,cons_subseteq_intro,powerset_intro,discrete_topology,open_in,cons_left,unions_iff,subseteq,covering,open_covering}.
    Take $B$ such that
        $B\subseteq C$ and
        $B$ is finite and
        $B$ is a covering of $\naturalsPlus$ by \cref{compact}.
    Therefore $\naturalsPlus$ is finite
        by \cref{subseteq,pair_is_finite,upair_intro_left,emptyset_subseteq,cons_subseteq_intro,subseteq_finite_is_finite,covering,finite_union_of_finite_family}.
    Contradiction by \cref{naturalsplus_infinite}.
    Therefore there exists $X,T$ such that
        $X$ is limit point compact with respect to $T$ and
        $X$ is not compact with respect to $T$.
\end{proof}

% from §28 helper definition 1: strictly increasing on a set
\begin{definition}\label{strictly_increasing_on_set}
    $u$ is strictly increasing on $A$ iff
        $u\in\funs{A}{A}$ and
        for all $m$ such that $m\in A$ we have
            for all $n$ such that $n\in A$ and $m < n$ we have $u(m) < u(n)$.
\end{definition}

% from §28 Item 4: subsequence
\begin{definition}\label{subsequence}
    $t$ is a subsequence of $s$ iff
        there exists $u$ such that
            $u$ is strictly increasing on $\naturalsPlus$ and
            for all $n\in\naturalsPlus$ we have $t(n) = \apply{s}{\apply{u}{n}}$.
\end{definition}

% from §28 Item 5: sequentially compact
\begin{definition}\label{sequentially_compact}
    $X$ is sequentially compact with respect to $T$ iff
        $T$ is a topology on $X$ and
        for all $s$ such that $s$ is a sequence in $X$
            there exist $t,x$ such that
                $t$ is a sequence in $X$ and
                $t$ is a subsequence of $s$ and
                $x\in X$ and
                $t$ converges to $x$ in $X$ with respect to $T$.
\end{definition}

% from §28 helper lemma 2: metric spaces are Hausdorff
\begin{lemma}\label{metric_space_hausdorff}
    Suppose $d$ is a metric on $X$.
    Let $T = \metricTopology{d}{X}$.
    Then $X$ is Hausdorff with respect to $T$.
\end{lemma}
\begin{proof}
    $\metricBasis{d}{X}$ is a basis on $X$ by \cref{metric_basis_is_basis}.
    Hence $T$ is a topology on $X$ by \cref{basis_topology_is_topology,metric_topology}.
    Show that for all $x_1,x_2$ such that $x_1\in X$ and $x_2\in X$ and $x_1\neq x_2$
        there exist $U_1,U_2$ such that
            $U_1$ is a neighborhood of $x_1$ in $X$ with respect to $T$ and
            $U_2$ is a neighborhood of $x_2$ in $X$ with respect to $T$ and
            $U_1$ is disjoint from $U_2$.
    \begin{subproof}
        Fix $x_1$.
        Fix $x_2$ such that $x_1\in X$ and $x_2\in X$ and $x_1\neq x_2$.
        Let $r = d((x_1,x_2))$.
        $r\in\reals$ by \cref{metric_on,times_tuple_intro,funs_type_apply}.
        Hence $0 < r$ by \cref{metric_on,metric_nonnegative,metric_zero_characterization,real_leq_split}.
        Let $m_2 = 1 + 1$.
        $m_2\in\reals$ and $0 < m_2$
            by \cref{real_one_in_naturals,real_naturals_subset_reals,subseteq,real_add_closed,real_add_ident,real_zero_lt_one,real_lt_add,real_lt_trans}.
        Hence $m_2 \neq 0$.
        Let $\epsilon = r / m_2$.
        $\epsilon\in\reals$ and $0 < \epsilon$
            by \cref{real_div,real_mul_inverse,real_mul_closed,real_inv_pos,real_add_ident,real_lt_mul_pos,real_mul_zero}.
        Let $U_1 = \metricBall{d}{X}{x_1}{\epsilon}$.
        Let $U_2 = \metricBall{d}{X}{x_2}{\epsilon}$.
        Therefore $U_1$ is a neighborhood of $x_1$ in $X$ with respect to $T$ and
            $U_2$ is a neighborhood of $x_2$ in $X$ with respect to $T$
            by \cref{metric_ball,subseteq,powerset_intro,metric_basis,metric_basis_is_basis,basis_elem_open_in_generated,basis_for_topology,basis_topology_is_topology,metric_topology,open_in,metric_on,metric_zero_characterization,neighborhood}.
        Show that for all $y$ such that $y\in U_1$ we have $y\notin U_2$.
        \begin{subproof}
            Fix $y$ such that $y\in U_1$.
            Suppose not.
            $y\in X$ and $d((x_1,y)) < \epsilon$ and $d((x_2,y)) < \epsilon$ and
                $d((x_1,y))\in\reals$ and $d((x_2,y))\in\reals$
                by \cref{metric_ball,metric_on,times_tuple_intro,funs_type_apply}.
            Hence $r \leq d((x_1,y)) + d((x_2,y))$ by \cref{metric_on,metric_triangle_inequality,metric_symmetric}.
            Then $d((x_1,y)) + \epsilon\in\reals$ and
                $d((x_1,y)) + d((x_2,y))\in\reals$ and
                $\epsilon + \epsilon\in\reals$ by \cref{real_add_closed}.
            Therefore $d((x_1,y)) + d((x_2,y)) < \epsilon + \epsilon$
                by \cref{real_lt_add,real_add_comm,real_lt_trans}.
            Hence $r < \epsilon + \epsilon$ by \cref{real_leq_split,real_lt_trans,real_lt_subst_left}.
            Hence $r < r$
                by \cref{real_div,real_right_distrib,real_mul_ident,real_mul_inverse,real_mul_assoc,real_mul_comm,real_lt_subst_right}.
            Contradiction by \cref{real_lt_irreflexive}.
        \end{subproof}
        Therefore $U_1$ is disjoint from $U_2$ by \cref{disjoint}.
    \end{subproof}
    Therefore $X$ is Hausdorff with respect to $T$ by \cref{hausdorff}.
\end{proof}


% from §28 helper definition 1: iteration sequence predicate
\begin{definition}\label{iteration_sequence}
    $u$ is an iteration sequence for $f$ on $n$ starting at $a_0$ iff
        $u$ is a function and
        $\dom{u} = \initialSegment{n}$ and
        $u(1) = a_0$ and
        for all $k\in\naturalsPlus$ such that $k + 1 \leq n$ we have $u(k + 1) = \apply{f}{u(k)}$.
\end{definition}

% from §28 helper definition 2: iteration set
\begin{definition}\label{iteration_set}
    $\iterationSet{f}{a_0} = \{ n\in\naturalsPlus \mid \text{there exists $u$ such that
        $u$ is an iteration sequence for $f$ on $n$ starting at $a_0$} \}$.
\end{definition}

% from §28 helper lemma 3: iteration set introduction
\begin{lemma}\label{iteration_set_intro}
    Suppose $n\in\naturalsPlus$.
    Suppose there exists $u$ such that
        $u$ is an iteration sequence for $f$ on $n$ starting at $a_0$.
    Then $n\in\iterationSet{f}{a_0}$.
\end{lemma}
\begin{proof}
    Hence $n\in\iterationSet{f}{a_0}$ by \cref{iteration_set}.
\end{proof}

% from §28 helper lemma 4: iteration set elimination
\begin{lemma}\label{iteration_set_elim}
    Suppose $n\in\iterationSet{f}{a_0}$.
    Then $n\in\naturalsPlus$ and there exists $u$ such that
        $u$ is an iteration sequence for $f$ on $n$ starting at $a_0$.
\end{lemma}
\begin{proof}
    Hence $n\in\naturalsPlus$ and there exists $u$ such that
        $u$ is an iteration sequence for $f$ on $n$ starting at $a_0$
        by \cref{iteration_set}.
\end{proof}

% from §28 helper lemma 5: iteration sequence introduction
\begin{lemma}\label{iteration_sequence_intro}
    Suppose $u$ is a function and $\dom{u} = \initialSegment{n}$ and $u(1) = a_0$.
    Suppose for all $k\in\naturalsPlus$ such that $k + 1 \leq n$ we have $u(k + 1) = \apply{f}{u(k)}$.
    Then $u$ is an iteration sequence for $f$ on $n$ starting at $a_0$.
\end{lemma}
\begin{proof}
    Hence $u$ is an iteration sequence for $f$ on $n$ starting at $a_0$ by \cref{iteration_sequence}.
\end{proof}

% from §28 helper lemma 5a: naturalsPlus successor is larger
\begin{lemma}\label{real_naturals_lt_succ_local}
    Suppose $n\in\naturalsPlus$.
    Then $n < n + 1$.
\end{lemma}
\begin{proof}
    Hence $n < n + 1$
        by \cref{real_naturals_subset_reals,real_one_in_naturals,subseteq,real_add_ident,real_zero_lt_one,real_lt_add,real_add_comm,real_lt_subst_left}.
\end{proof}

% from §28 helper lemma 5b: naturalsPlus successor is not smaller
\begin{lemma}\label{real_naturals_leq_succ_local}
    Suppose $n\in\naturalsPlus$.
    Then $n \leq n + 1$.
\end{lemma}
\begin{proof}
    Hence $n \leq n + 1$ by \cref{real_naturals_lt_succ_local}.
\end{proof}

% from §28 helper lemma 5c: predecessor bound elimination
\begin{lemma}\label{real_naturals_succ_bound_elim_local}
    Suppose $k\in\naturalsPlus$ and $n\in\naturalsPlus$.
    Suppose $k + 1 \leq n$.
    Then $k \leq n$.
\end{lemma}
\begin{proof}
    Hence $k \leq n$
        by \cref{real_naturals_leq_succ_local,real_naturals_subset_reals,real_naturals_closed_succ,subseteq,real_leq_trans}.
\end{proof}

% from §28 helper lemma 6: iteration sequence uniqueness
\begin{lemma}\label{iteration_sequence_unique}
    Suppose $n\in\naturalsPlus$.
    Suppose $u$ and $v$ are iteration sequences for $f$ on $n$ starting at $a_0$.
    Then $u = v$.
\end{lemma}
\begin{proof}
    Hence $u$ is a function and $v$ is a function and
        $\dom{u} = \initialSegment{n}$ and $\dom{v} = \initialSegment{n}$ by \cref{iteration_sequence}.
    Let $A = \{k\in\naturalsPlus \mid \text{if $k\leq n$ then $u(k) = v(k)$}\}$.
    Then $1\in A$ by \cref{real_one_in_naturals,iteration_sequence}.
    Show that for all $k\in A$ we have $k + 1\in A$.
    \begin{subproof}
        Fix $k$ such that $k\in A$.
        If $k + 1\leq n$ then $u(k + 1) = v(k + 1)$
            by \cref{real_naturals_succ_bound_elim_local,iteration_sequence}.
        Therefore $k + 1\in A$.
    \end{subproof}
    Therefore for all $k\in\initialSegment{n}$ we have $u(k) = v(k)$
        by \cref{subseteq,real_naturals_induction,initial_segment_naturalsplus}.
    For all $w$ such that $w\in u$ we have $w\in v$ by \cref{function_member_elim,subseteq,function_apply_elim}.
    For all $w$ such that $w\in v$ we have $w\in u$ by \cref{function_member_elim,subseteq,function_apply_elim}.
    Hence $u = v$ by \cref{subseteq,subseteq_antisymmetric}.
\end{proof}
% from §28 helper lemma 7: infinite subset of naturalsPlus has increasing sequence
\begin{lemma}\label{infinite_naturalsplus_has_strictly_increasing_sequence}
    Suppose $A\subseteq\naturalsPlus$.
    Suppose $A$ is infinite.
    Then there exists $u$ such that
        $u$ is strictly increasing on $\naturalsPlus$ and
        for all $n\in\naturalsPlus$ we have $u(n)\in A$.
\end{lemma}
\begin{proof}
    Show that $A\neq\emptyset$.
    \begin{subproof}
        Suppose not.
        Then $A$ is finite by \cref{emptyset_is_finite}.
        Contradiction.
    \end{subproof}
    Take $a_0\in A$ such that for all $a\in A$ we have $a_0 \leq a$
        by \cref{real_naturals_least_element}.
    Show that for all $a\in A$ there exists $b\in A$ such that $a < b$.
    \begin{subproof}
        Fix $a$ such that $a\in A$.
        Thus $a\in\naturalsPlus$ by \cref{subseteq}.
        Suppose not.
        Then for all $b\in A$ we have $b \leq a$.
        Thus $A\subseteq\initialSegment{a}$ by \cref{subseteq,initial_segment_naturalsplus}.
        Thus $A$ is finite by \cref{initial_segment_finite,subseteq_finite_is_finite}.
        Contradiction.
    \end{subproof}
    Let $R = \{w\in A\times A \mid \fst{w} < \snd{w}\}$.
    For all $x\in A$ there exists $y$ such that $x\mathrel{R} y$.
    Take $f$ such that $f$ is a function on $A$ and for all $x\in A$ we have
        $x\mathrel{R}\apply{f}{x}$ by \cref{times_elem_is_tuple,relation,choice_function}.
    For all $x\in A$ we have $\apply{f}{x}\in A$ and $x < \apply{f}{x}$.
    Let $S = \iterationSet{f}{a_0}$.
    Show that for all $n\in\naturalsPlus$ we have $n\in S$.
    \begin{subproof}
        $S\subseteq\naturalsPlus$ by \cref{iteration_set_elim,subseteq}.
        $1\in\naturalsPlus$ by \cref{real_one_in_naturals}.
        Show that $1\in S$.
        \begin{subproof}
            Let $u = \{(1,a_0)\}$.
            For all $x,y,y'$ such that $(x,y)\in u$ and $(x,y')\in u$ we have $y = y'$
                by \cref{singleton_iff,pair_eq_iff}.
            Hence $u$ is a function by \cref{singleton_iff,relation,rightunique}.
            For all $x$ such that $x\in\dom{u}$ we have $x\in\initialSegment{1}$
                by \cref{dom,singleton_iff,pair_eq_iff,real_one_in_naturals,initial_segment_naturalsplus}.
            For all $x$ such that $x\in\initialSegment{1}$ we have $x\in\dom{u}$.
            Hence $\dom{u} = \initialSegment{1}$ by \cref{subseteq,subseteq_antisymmetric}.
            Thus $u(1) = a_0$ by \cref{singleton_intro,function_apply_intro}.
            Show that for all $k\in\naturalsPlus$ such that $k + 1 \leq 1$ we have $u(k + 1) = \apply{f}{u(k)}$.
            \begin{subproof}
                Fix $k$ such that $k\in\naturalsPlus$.
                Assume $k + 1 \leq 1$.
                Suppose not.
                Then $k < k + 1$ by \cref{real_naturals_lt_succ_local}.
                Show that $1 < k + 1$.
                \begin{subproof}
                    $1\in\reals$ and $k\in\reals$ and $k + 1\in\reals$
                        by \cref{real_naturals_subset_reals,real_one_in_naturals,real_naturals_closed_succ,subseteq}.
                    $1 < k$ or $1 = k$ by \cref{real_one_leq_naturalsplus}.
                    \begin{byCase}
                    \caseOf{$1 = k$.}
                        Thus $1 < k + 1$ by \cref{real_lt_subst_left}.
                    \caseOf{$1 < k$.}
                        Thus $1 < k + 1$ by \cref{real_lt_trans}.
                    \end{byCase}
                \end{subproof}
                $k + 1 < 1$ or $k + 1 = 1$ by \cref{real_leq_split}.
                Show that $1 < 1$.
                \begin{subproof}
                    $1\in\reals$ and $k + 1\in\reals$
                        by \cref{real_naturals_subset_reals,real_one_in_naturals,real_naturals_closed_succ,subseteq}.
                    \begin{byCase}
                    \caseOf{$k + 1 < 1$.}
                        Thus $1 < 1$ by \cref{real_lt_trans}.
                    \caseOf{$k + 1 = 1$.}
                        Thus $1 < 1$ by \cref{real_lt_subst_right}.
                    \end{byCase}
                \end{subproof}
                Contradiction by \cref{real_naturals_subset_reals,real_one_in_naturals,subseteq,real_lt_irreflexive}.
            \end{subproof}
            Hence $1\in S$ by \cref{iteration_sequence,iteration_set_intro}.
        \end{subproof}
        Show that for all $n\in S$ we have $n + 1\in S$.
        \begin{subproof}
            Fix $n$ such that $n\in S$.
            Thus $n\in\naturalsPlus$ by \cref{iteration_set_elim}.
            Take $u$ such that
                $u$ is an iteration sequence for $f$ on $n$ starting at $a_0$.
            Thus for all $k\in\naturalsPlus$ such that $k + 1 \leq n$ we have $u(k + 1) = \apply{f}{u(k)}$
                by \cref{iteration_sequence}.
            Thus $n + 1\in\naturalsPlus$ by \cref{real_naturals_closed_succ}.
            Let $v = u\union\{(n + 1,\apply{f}{u(n)})\}$.
            Thus $v$ is a relation by \cref{singleton_iff,relation,iteration_sequence,union_relations_is_relation}.
            Show that for all $x,y,y'$ such that $(x,y)\in v$ and $(x,y')\in v$ we have $y = y'$.
            \begin{subproof}
                Fix $x$.
                Fix $y$.
                Fix $y'$.
                Assume $(x,y)\in v$ and $(x,y')\in v$.
                \begin{byCase}
                \caseOf{$x = n + 1$.}
                    $(n + 1,y')\in v$.
                    Show that $(n + 1,y)\in\{(n + 1,\apply{f}{u(n)})\}$.
                    \begin{subproof}
                        Suppose not.
                        Thus $(n + 1,y)\in u$ by \cref{union_iff}.
                        Thus $n + 1\in\dom{u}$ by \cref{dom}.
                        Thus $\dom{u} = \initialSegment{n}$ by \cref{iteration_sequence}.
                        Thus $n + 1\in\initialSegment{n}$.
                        Thus $n + 1 \leq n$ by \cref{initial_segment_naturalsplus}.
                        Thus $n < n + 1$ by \cref{real_naturals_lt_succ_local}.
                        $n\in\reals$ and $n + 1\in\reals$
                            by \cref{real_naturals_subset_reals,real_naturals_closed_succ,subseteq}.
                        $n + 1 < n$ or $n + 1 = n$ by \cref{real_leq_split}.
                        \begin{byCase}
                        \caseOf{$n + 1 < n$.}
                            Thus $n < n$ by \cref{real_lt_trans}.
                            Contradiction by \cref{real_lt_irreflexive}.
                        \caseOf{$n + 1 = n$.}
                            Thus $n < n$ by \cref{real_lt_subst_right}.
                            Contradiction by \cref{real_lt_irreflexive}.
                        \end{byCase}
                    \end{subproof}
                    Thus $(n + 1,y) = (n + 1,\apply{f}{u(n)})$ by \cref{singleton_iff}.
                    Thus $y = \apply{f}{u(n)}$ by \cref{pair_eq_iff}.
                    Show that $(n + 1,y')\in\{(n + 1,\apply{f}{u(n)})\}$.
                    \begin{subproof}
                        Suppose not.
                        Thus $(n + 1,y')\in u$ by \cref{union_iff}.
                        Thus $n + 1\in\dom{u}$ by \cref{dom}.
                        Thus $\dom{u} = \initialSegment{n}$ by \cref{iteration_sequence}.
                        Thus $n + 1\in\initialSegment{n}$.
                        Thus $n + 1 \leq n$ by \cref{initial_segment_naturalsplus}.
                        Thus $n < n + 1$ by \cref{real_naturals_lt_succ_local}.
                        $n\in\reals$ and $n + 1\in\reals$
                            by \cref{real_naturals_subset_reals,real_naturals_closed_succ,subseteq}.
                        $n + 1 < n$ or $n + 1 = n$ by \cref{real_leq_split}.
                        \begin{byCase}
                        \caseOf{$n + 1 < n$.}
                            Thus $n < n$ by \cref{real_lt_trans}.
                            Contradiction by \cref{real_lt_irreflexive}.
                        \caseOf{$n + 1 = n$.}
                            Thus $n < n$ by \cref{real_lt_subst_right}.
                            Contradiction by \cref{real_lt_irreflexive}.
                        \end{byCase}
                    \end{subproof}
                    Thus $(n + 1,y') = (n + 1,\apply{f}{u(n)})$ by \cref{singleton_iff}.
                    Thus $y' = \apply{f}{u(n)}$ by \cref{pair_eq_iff}.
                    Thus $y = y'$.
                \caseOf{$x \neq n + 1$.}
                    Thus $(x,y)\in u$ by \cref{union_iff,singleton_iff,pair_eq_iff}.
                    Thus $(x,y')\in u$ by \cref{union_iff,singleton_iff,pair_eq_iff}.
                    Thus $y = y'$ by \cref{iteration_sequence,function_rightunique}.
                \end{byCase}
            \end{subproof}
            Hence $v$ is a function by \cref{rightunique}.
            Show that for all $x$ such that $x\in\dom{v}$ we have $x\in\initialSegment{n + 1}$.
            \begin{subproof}
                Fix $x$ such that $x\in\dom{v}$.
                Take $y$ such that $(x,y)\in v$ by \cref{dom}.
                Thus $(x,y)\in u$ or $(x,y)\in\{(n + 1,\apply{f}{u(n)})\}$ by \cref{union_iff}.
                \begin{byCase}
                \caseOf{$(x,y)\in u$.}
                    Thus $x\in\initialSegment{n}$ by \cref{dom,iteration_sequence}.
                    Thus $x\in\initialSegment{n + 1}$
                        by \cref{initial_segment_naturalsplus,real_naturals_leq_succ_local,real_leq_trans,real_naturals_subset_reals,real_naturals_closed_succ,subseteq}.
                \caseOf{$(x,y)\in\{(n + 1,\apply{f}{u(n)})\}$.}
                    Thus $(x,y) = (n + 1,\apply{f}{u(n)})$ by \cref{singleton_iff}.
                    Thus $x = n + 1$ by \cref{pair_eq_iff}.
                    Thus $x\in\initialSegment{n + 1}$ by \cref{initial_segment_naturalsplus}.
                \end{byCase}
            \end{subproof}
            Show that for all $x$ such that $x\in\initialSegment{n + 1}$ we have $x\in\dom{v}$.
            \begin{subproof}
                Fix $x$ such that $x\in\initialSegment{n + 1}$.
                Thus $x\in\naturalsPlus$ and $x \leq n + 1$ by \cref{initial_segment_naturalsplus}.
                \begin{byCase}
                \caseOf{$x = n + 1$.}
                    Thus $x\in\dom{v}$ by \cref{cons_left,union_intro_right,dom}.
                \caseOf{$x \neq n + 1$.}
                    Thus $x \leq n$ by \cref{real_naturals_leq_succ_elim}.
                    Thus $x\in\dom{u}$ by \cref{initial_segment_naturalsplus,iteration_sequence}.
                    Thus $x\in\dom{v}$ by \cref{dom,union_iff}.
                \end{byCase}
            \end{subproof}
            Hence $\dom{v} = \initialSegment{n + 1}$ by \cref{subseteq,subseteq_antisymmetric}.
            Thus $v(1) = a_0$
                by \cref{real_one_in_naturals,real_one_leq_naturalsplus,initial_segment_naturalsplus,iteration_sequence,function_apply_elim,union_intro_left,function_apply_intro}.
            Show that for all $k\in\naturalsPlus$ such that $k + 1 \leq n + 1$ we have $v(k + 1) = \apply{f}{v(k)}$.
            \begin{subproof}
                Fix $k$ such that $k\in\naturalsPlus$.
                Assume $k + 1 \leq n + 1$.
                \begin{byCase}
                \caseOf{$k + 1 = n + 1$.}
                    Thus $k = n$.
                    Thus $v(k + 1) = \apply{f}{u(n)}$
                        by \cref{cons_left,union_intro_right,function_apply_intro}.
                    Thus $v(k + 1) = \apply{f}{u(k)}$.
                    Thus $k \leq n$.
                    Thus $v(k) = u(k)$
                        by \cref{initial_segment_naturalsplus,iteration_sequence,function_apply_elim,union_intro_left,function_apply_intro}.
                    Thus $v(k + 1) = \apply{f}{v(k)}$.
                \caseOf{$k + 1 \neq n + 1$.}
                    Thus $k + 1\in\naturalsPlus$ by \cref{real_naturals_closed_succ}.
                    Thus $k + 1 \leq n$ by \cref{real_naturals_leq_succ_elim}.
                    Thus $k + 1\in\initialSegment{n}$ by \cref{initial_segment_naturalsplus}.
                    Thus $(k + 1,u(k + 1))\in v$
                        by \cref{iteration_sequence,function_apply_elim,union_intro_left}.
                    Thus $k + 1\in\dom{v}$ by \cref{dom}.
                    Thus $(k + 1,v(k + 1))\in v$ by \cref{function_apply_elim}.
                    Thus $u(k + 1) = v(k + 1)$ by \cref{function_rightunique}.
                    Thus $v(k + 1) = u(k + 1)$.
                    Thus $k\in\naturalsPlus$ and $k + 1 \leq n$.
                    Thus $u(k + 1) = \apply{f}{u(k)}$.
                    Thus $v(k + 1) = \apply{f}{u(k)}$.
                    Thus $k \leq n$ by \cref{real_naturals_succ_bound_elim_local}.
                    Thus $(k,u(k))\in v$
                        by \cref{initial_segment_naturalsplus,iteration_sequence,function_apply_elim,union_intro_left}.
                    Thus $k\in\dom{v}$ by \cref{dom}.
                    Thus $(k,v(k))\in v$ by \cref{function_apply_elim}.
                    Thus $u(k) = v(k)$ by \cref{function_rightunique}.
                    Thus $v(k + 1) = \apply{f}{v(k)}$.
                \end{byCase}
            \end{subproof}
            Hence $v$ is an iteration sequence for $f$ on $n + 1$ starting at $a_0$ by \cref{iteration_sequence_intro}.
            Therefore $n + 1\in S$ by \cref{iteration_set_intro}.
        \end{subproof}
        Hence for all $n\in\naturalsPlus$ we have $n\in S$ by \cref{real_naturals_induction}.
    \end{subproof}
    Let $F = \{u \mid \text{there exists $n\in\naturalsPlus$ such that
        $u$ is an iteration sequence for $f$ on $n$ starting at $a_0$} \}$.
    For all $u\in F$ we have $u$ is a function by \cref{iteration_sequence}.
    Show that for all $u,v$ such that $u\in F$ and $v\in F$ we have $u\subseteq v$ or $v\subseteq u$.
    \begin{subproof}
        Fix $u$.
        Fix $v$.
        Assume $u\in F$ and $v\in F$.
        Take $m\in\naturalsPlus$ such that
            $u$ is an iteration sequence for $f$ on $m$ starting at $a_0$.
        Take $n\in\naturalsPlus$ such that
            $v$ is an iteration sequence for $f$ on $n$ starting at $a_0$.
        Thus $u$ is a function and $\dom{u} = \initialSegment{m}$ and $u(1) = a_0$
            by \cref{iteration_sequence}.
        Thus for all $k\in\naturalsPlus$ such that $k + 1 \leq m$ we have $u(k + 1) = \apply{f}{u(k)}$
            by \cref{iteration_sequence}.
        Thus $v$ is a function and $\dom{v} = \initialSegment{n}$ and $v(1) = a_0$
            by \cref{iteration_sequence}.
        Thus for all $k\in\naturalsPlus$ such that $k + 1 \leq n$ we have $v(k + 1) = \apply{f}{v(k)}$
            by \cref{iteration_sequence}.
        $m = n$ or $m \neq n$.
        \begin{byCase}
        \caseOf{$m = n$.}
            Then $u\subseteq v$ or $v\subseteq u$ by \cref{iteration_sequence_unique,subseteq_refl}.
        \caseOf{$m \neq n$.}
            $m\in\reals$ and $n\in\reals$ by \cref{real_naturals_subset_reals,subseteq}.
            Thus $m < n$ or $n < m$ by \cref{real_lt_trichotomy}.
            \begin{byCase}
        \caseOf{$m < n$.}
            Let $A = \{k\in\naturalsPlus \mid \text{if $k\leq m$ then $u(k) = v(k)$}\}$.
            Show that for all $k\in\naturalsPlus$ we have $k\in A$.
            \begin{subproof}
                $1\in\naturalsPlus$ by \cref{real_one_in_naturals}.
                If $1\leq m$ then $u(1) = v(1)$ by \cref{iteration_sequence}.
                Thus $1\in A$.
                Show that for all $k\in A$ we have $k + 1\in A$.
                \begin{subproof}
                    Fix $k$ such that $k\in A$.
                    If $k + 1\leq m$ then $u(k + 1) = v(k + 1)$
                        by \cref{real_naturals_succ_bound_elim_local,real_leq_trans,real_lt_imp_leq,real_add_closed,real_naturals_subset_reals,subseteq,iteration_sequence}.
                    Thus $k + 1\in A$.
                \end{subproof}
                Therefore for all $k\in\initialSegment{m}$ we have $u(k) = v(k)$
                    by \cref{subseteq,real_naturals_induction,initial_segment_naturalsplus}.
            \end{subproof}
            For all $w$ such that $w\in u$ we have $w\in v$
                by \cref{iteration_sequence,function_member_elim,initial_segment_naturalsplus,real_naturals_subset_reals,subseteq,real_leq_trans,real_lt_imp_leq,function_apply_elim}.
            Hence $u\subseteq v$ by \cref{subseteq}.
        \caseOf{$n < m$.}
            Let $A = \{k\in\naturalsPlus \mid \text{if $k\leq n$ then $v(k) = u(k)$}\}$.
            Show that for all $k\in\naturalsPlus$ we have $k\in A$.
            \begin{subproof}
                $1\in\naturalsPlus$ by \cref{real_one_in_naturals}.
                If $1\leq n$ then $v(1) = u(1)$ by \cref{iteration_sequence}.
                Thus $1\in A$.
                Show that for all $k\in A$ we have $k + 1\in A$.
                \begin{subproof}
                    Fix $k$ such that $k\in A$.
                    If $k + 1\leq n$ then $v(k + 1) = u(k + 1)$
                        by \cref{real_naturals_succ_bound_elim_local,real_leq_trans,real_add_closed,real_naturals_subset_reals,subseteq,iteration_sequence}.
                    Thus $k + 1\in A$.
                \end{subproof}
                Therefore for all $k\in\initialSegment{n}$ we have $v(k) = u(k)$
                    by \cref{subseteq,real_naturals_induction,initial_segment_naturalsplus}.
            \end{subproof}
            For all $w$ such that $w\in v$ we have $w\in u$
                by \cref{iteration_sequence,function_member_elim,initial_segment_naturalsplus,real_naturals_subset_reals,subseteq,real_leq_trans,real_lt_imp_leq,function_apply_elim}.
            Hence $v\subseteq u$ by \cref{subseteq}.
        \end{byCase}
        \end{byCase}
    \end{subproof}
    Let $u = \unions{F}$.
    $F$ is a family of functions.
    Hence $u$ is a function by \cref{unions_of_compatible_family_of_function_is_function}.
    For all $x$ such that $x\in\dom{u}$ we have $x\in\naturalsPlus$
        by \cref{dom,unions_iff,iteration_sequence,initial_segment_naturalsplus}.
    For all $n\in\naturalsPlus$ we have $n\in\dom{u}$
        by \cref{subseteq,iteration_set_elim,iteration_sequence,initial_segment_naturalsplus,dom,unions_iff}.
    Therefore $\dom{u} = \naturalsPlus$ by \cref{subseteq,subseteq_antisymmetric}.
    Show that for all $n\in\naturalsPlus$ we have $u(n)\in A$.
    \begin{subproof}
        Fix $n\in\naturalsPlus$.
        Take $v\in F$ such that $(n,u(n))\in v$ by \cref{function_apply_elim,unions_iff}.
        Take $n_0\in\naturalsPlus$ such that
            $v$ is an iteration sequence for $f$ on $n_0$ starting at $a_0$.
        Thus $n \leq n_0$ by \cref{dom,iteration_sequence,initial_segment_naturalsplus}.
        Thus $v(1) = a_0$ and for all $k\in\naturalsPlus$ such that $k + 1 \leq n_0$ we have $v(k + 1) = \apply{f}{v(k)}$
            by \cref{iteration_sequence}.
        Let $B = \{k\in\naturalsPlus \mid \text{if $k\leq n$ then $v(k)\in A$}\}$.
        Show that for all $k\in\naturalsPlus$ we have $k\in B$.
        \begin{subproof}
            $1\in\naturalsPlus$ by \cref{real_one_in_naturals}.
            If $1\leq n$ then $v(1)\in A$ by \cref{iteration_sequence}.
            Thus $1\in B$.
            Show that for all $k\in B$ we have $k + 1\in B$.
            \begin{subproof}
                Fix $k$ such that $k\in B$.
                If $k + 1\leq n$ then $v(k + 1)\in A$
                    by \cref{real_naturals_succ_bound_elim_local,real_leq_trans,real_naturals_closed_succ,real_naturals_subset_reals,subseteq,iteration_sequence}.
                Thus $k + 1\in B$.
            \end{subproof}
            Thus $v(n)\in A$ by \cref{subseteq,real_naturals_induction,initial_segment_naturalsplus}.
        \end{subproof}
        Thus $u(n)\in A$ by \cref{function_apply_intro}.
    \end{subproof}
    Show that $u$ is strictly increasing on $\naturalsPlus$.
    \begin{subproof}
        Show that for all $n\in\naturalsPlus$ we have $u(n) < u(n + 1)$.
        \begin{subproof}
            Fix $n\in\naturalsPlus$.
            Thus $(n + 1,u(n + 1))\in\unions{F}$ by \cref{real_naturals_closed_succ,function_apply_elim}.
            Take $v\in F$ such that $(n + 1,u(n + 1))\in v$ by \cref{unions_iff}.
            Thus $n + 1\in\dom{v}$ by \cref{dom}.
            Thus $v(n + 1) = u(n + 1)$ by \cref{function_apply_intro}.
            Take $n_0\in\naturalsPlus$ such that
                $v$ is an iteration sequence for $f$ on $n_0$ starting at $a_0$.
            Thus $n + 1 \leq n_0$ by \cref{iteration_sequence,initial_segment_naturalsplus}.
            Thus $v(n + 1) = \apply{f}{v(n)}$ by \cref{iteration_sequence}.
            Thus $u(n + 1) = \apply{f}{v(n)}$.
            Thus $n \leq n_0$ by \cref{real_naturals_succ_bound_elim_local}.
            Thus $(n,v(n))\in u$
                by \cref{initial_segment_naturalsplus,iteration_sequence,function_apply_elim,unions_iff}.
            Thus $u(n) = v(n)$ by \cref{function_apply_intro}.
            Thus $u(n + 1) = \apply{f}{u(n)}$.
            Thus $u(n) < u(n + 1)$ by \cref{function_apply_intro}.
        \end{subproof}
        Show that for all $m,n\in\naturalsPlus$ such that $m < n$ we have $u(m) < u(n)$.
        \begin{subproof}
            Let $T = \{n\in\naturalsPlus \mid \text{for all $m\in\naturalsPlus$ such that $m < n$ we have $u(m) < u(n)$}\}$.
            Show that for all $n\in\naturalsPlus$ we have $n\in T$.
            \begin{subproof}
                $1\in\naturalsPlus$ by \cref{real_one_in_naturals}.
                For all $m\in\naturalsPlus$ such that $m < 1$ we have $u(m) < u(1)$
                    by \cref{real_naturals_subset_reals,real_one_in_naturals,subseteq,real_one_leq_naturalsplus,real_lt_subst_right,real_lt_trans,real_lt_irreflexive}.
                Thus $1\in T$.
                Show that for all $n\in T$ we have $n + 1\in T$.
                \begin{subproof}
                    Fix $n$ such that $n\in T$.
                    Show that for all $m\in\naturalsPlus$ such that $m < n + 1$ we have $u(m) < u(n + 1)$.
                    \begin{subproof}
                        Fix $m$.
                        Assume $m\in\naturalsPlus$ and $m < n + 1$.
                        Suppose not.
                        Thus $u(m)\in A$ and $u(n)\in A$ and $u(n + 1)\in A$.
                        Thus $u(m)\in\reals$ and $u(n)\in\reals$ and $u(n + 1)\in\reals$
                            by \cref{subseteq,real_naturals_subset_reals}.
                        Thus $m < n$ or $m = n$ or $m = n + 1$ by \cref{real_succ_discrete}.
                        \begin{byCase}
                        \caseOf{$m = n$.}
                            Thus $u(m) < u(n + 1)$ by \cref{real_lt_subst_left}.
                            Contradiction.
                        \caseOf{$m < n$.}
                            Thus $u(m) < u(n + 1)$ by \cref{subseteq,real_lt_trans}.
                            Contradiction.
                        \caseOf{$m = n + 1$.}
                            Contradiction by \cref{real_lt_subst_left,real_naturals_subset_reals,subseteq,real_lt_irreflexive}.
                        \end{byCase}
                    \end{subproof}
                    Thus $n + 1\in T$.
                \end{subproof}
                Thus for all $m,n\in\naturalsPlus$ such that $m < n$ we have $u(m) < u(n)$
                    by \cref{subseteq,real_naturals_induction}.
            \end{subproof}
        \end{subproof}
        For all $n\in\dom{u}$ we have $u(n)\in\naturalsPlus$.
        Thus $u\in\funs{\naturalsPlus}{\naturalsPlus}$ by \cref{funs_intro}.
        Hence $u$ is strictly increasing on $\naturalsPlus$ by \cref{strictly_increasing_on_set}.
    \end{subproof}
    Therefore there exists $u$ such that
        $u$ is strictly increasing on $\naturalsPlus$ and
        for all $n\in\naturalsPlus$ we have $u(n)\in A$.
\end{proof}


% from §28 helper lemma 8: constant subsequence in a finite range
\begin{lemma}\label{finite_sequence_constant_subsequence}
    Suppose $A\subseteq X$.
    Suppose $A$ is finite.
    Suppose $s$ is a sequence in $A$.
    Then there exist $t,a$ such that
        $a\in A$ and
        $t$ is a sequence in $A$ and
        $t$ is a subsequence of $s$ and
        for all $n\in\naturalsPlus$ we have $t(n) = a$.
\end{lemma}
\begin{proof}
    $s\in\funs{\naturalsPlus}{A}$ by \cref{sequence_in}.
    Then $\dom{s} = \naturalsPlus$ and $s$ is right-unique and $s$ is a function by \cref{funs,funs_is_function}.
    Let $g = \{ (a,\preimg{s}{\{a\}}) \mid a\in A \}$.
    Let $F = \img{g}{A}$.
    Show that $g$ is a function on $A$.
    \begin{subproof}
        For all $w$ such that $w\in g$ there exist $a,B$ such that $w = (a,B)$.
        Show that for all $a,B_1,B_2$ such that $(a,B_1)\in g$ and $(a,B_2)\in g$ we have $B_1 = B_2$.
        \begin{subproof}
            Fix $a$.
            Fix $B_1$.
            Fix $B_2$ such that $(a,B_1)\in g$ and $(a,B_2)\in g$.
            Take $a_1\in A$ such that $(a,B_1) = (a_1,\preimg{s}{\{a_1\}})$.
            Take $a_2\in A$ such that $(a,B_2) = (a_2,\preimg{s}{\{a_2\}})$.
            Thus $a = a_1$ and $B_1 = \preimg{s}{\{a_1\}}$ by \cref{pair_eq_iff}.
            Thus $a = a_2$ and $B_2 = \preimg{s}{\{a_2\}}$ by \cref{pair_eq_iff}.
            Thus $B_1 = \preimg{s}{\{a\}}$ and $B_2 = \preimg{s}{\{a\}}$.
            Thus $B_1 = B_2$.
        \end{subproof}
        For all $a$ such that $a\in\dom{g}$ we have $a\in A$.
        For all $a$ such that $a\in A$ we have $a\in\dom{g}$.
        Therefore $\dom{g} = A$ by \cref{subseteq,subseteq_antisymmetric}.
        Hence $g$ is a function on $A$ by \cref{relation,rightunique}.
    \end{subproof}
    Hence $F$ is finite by \cref{finite_image_function}.
    For all $n$ such that $n\in\naturalsPlus$ we have $n\in\unions{F}$
        by \cref{funs_type_apply,img_iff,singleton_intro,preimg_iff,function_apply_iff,unions_iff}.
    For all $n$ such that $n\in\unions{F}$ we have $n\in\naturalsPlus$
        by \cref{unions_iff,img_iff,function_apply_iff,preimg_subseteq_dom,subseteq}.
    Hence $\unions{F} = \naturalsPlus$ by \cref{subseteq,subseteq_antisymmetric}.
    Show that there exists $a\in A$ such that $\preimg{s}{\{a\}}$ is infinite.
    \begin{subproof}
        Suppose not.
        Then for all $a\in A$ we have $\preimg{s}{\{a\}}$ is finite.
        For all $B\in F$ we have $B$ is finite by \cref{img_iff,function_apply_iff}.
        Hence $\unions{F}$ is finite by \cref{finite_union_of_finite_family}.
        Then $\naturalsPlus$ is finite.
        Contradiction by \cref{naturalsplus_infinite}.
    \end{subproof}
    Take $a\in A$ such that $\preimg{s}{\{a\}}$ is infinite.
    Let $I = \preimg{s}{\{a\}}$.
    $I\subseteq\naturalsPlus$ by \cref{preimg_subseteq_dom,subseteq}.
    Take $u$ such that
        $u$ is strictly increasing on $\naturalsPlus$ and
        for all $n\in\naturalsPlus$ we have $u(n)\in I$
        by \cref{infinite_naturalsplus_has_strictly_increasing_sequence}.
    Let $t = s\circ u$.
    Show that $t\in\funs{\naturalsPlus}{A}$.
    \begin{subproof}
        Hence $t$ is a function by \cref{strictly_increasing_on_set,funs,function_circ}.
        For all $n$ such that $n\in\dom{t}$ we have $n\in\naturalsPlus$
            by \cref{dom,circ_elem_elim,strictly_increasing_on_set,funs}.
        For all $n$ such that $n\in\naturalsPlus$ we have $n\in\dom{t}$
            by \cref{strictly_increasing_on_set,funs,subseteq,funs_is_function,function_apply_elim,circ_elem_intro,dom}.
        Hence $\dom{t} = \naturalsPlus$ by \cref{subseteq,subseteq_antisymmetric}.
        For all $n\in\dom{t}$ we have $t(n)\in A$
            by \cref{dom,strictly_increasing_on_set,funs_is_function,function_apply_iff,function_apply_elim,circ_elem_elim,subseteq,funs_type_apply}.
        Hence $t\in\funs{\naturalsPlus}{A}$ by \cref{funs_intro}.
    \end{subproof}
    Hence $t$ is a sequence in $A$ by \cref{sequence_in}.
    For all $n\in\naturalsPlus$ we have $t(n) = \apply{s}{\apply{u}{n}}$
        by \cref{sequence_in,funs,funs_is_function,function_apply_elim,circ_elem_elim,strictly_increasing_on_set,function_apply_iff}.
    Therefore $t$ is a subsequence of $s$ by \cref{subsequence}.
    For all $n\in\naturalsPlus$ we have $t(n) = a$
        by \cref{strictly_increasing_on_set,funs,preimg_iff,singleton_iff,function_apply_iff}.
    Therefore there exist $t,a$ such that
        $a\in A$ and
        $t$ is a sequence in $A$ and
        $t$ is a subsequence of $s$ and
        for all $n\in\naturalsPlus$ we have $t(n) = a$.
\end{proof}

% from §28 helper lemma: reciprocal reverses order on naturalsPlus
\begin{lemma}\label{naturalsplus_reciprocal_antitone}
    Suppose $m,n\in\naturalsPlus$.
    Suppose $m \leq n$.
    Then $1 / n \leq 1 / m$.
\end{lemma}
\begin{proof}
    $m < n$ or $m = n$ by \cref{real_leq_split}.
    $m\in\reals$ and $n\in\reals$ and $0\in\reals$ and $1\in\reals$
        by \cref{real_naturals_subset_reals,subseteq,real_add_ident,real_mul_ident}.

    Hence $0 < m$ by \cref{real_add_ident,real_zero_lt_one,real_one_leq_naturalsplus,real_lt_subst_right,real_lt_trans}.
    Hence $0 < n$ by \cref{real_add_ident,real_zero_lt_one,real_one_leq_naturalsplus,real_lt_subst_right,real_lt_trans}.
    Suppose not.
    Then $1 / m < 1 / n$.
    Hence $0 < m \rmul n$ by \cref{real_lt_mul_pos,real_mul_zero,real_lt_subst_left}.
    $m\neq 0$ and $n\neq 0$ by \cref{real_pos_nonzero}.
    Then $1 / m\in\reals$ and $1 / n\in\reals$ and $m \rmul n\in\reals$
        by \cref{real_mul_inverse,real_div,real_mul_closed}.
    Hence $(1 / m) \rmul (m \rmul n) < (1 / n) \rmul (m \rmul n)$ by \cref{real_lt_mul_pos}.
    Therefore $(m \rmul n) \rmul (1 / m) < (m \rmul n) \rmul (1 / n)$
        by \cref{real_mul_comm,real_lt_subst_left,real_lt_subst_right}.
    Then $(m \rmul n) \rmul (1 / m) = n$
        by \cref{real_div,real_mul_assoc,real_mul_comm,real_mul_closed,real_mul_inverse,real_mul_ident}.
    Then $(m \rmul n) \rmul (1 / n) = m$
        by \cref{real_div,real_mul_assoc,real_mul_closed,real_mul_comm,real_mul_inverse,real_mul_ident}.
    Therefore $n < m$.
    Contradiction by \cref{real_lt_trans,real_lt_subst_left,real_lt_irreflexive}.
\end{proof}
% from §28 helper lemma 9: sequential compactness implies Lebesgue number lemma
\begin{lemma}\label{sequentially_compact_lebesgue}
    Suppose $d$ is a metric on $X$.
    Let $T = \metricTopology{d}{X}$.
    Suppose $A$ is an open covering of $X$ with respect to $T$.
    Suppose $X\neq\emptyset$.
    Suppose $X$ is sequentially compact with respect to $T$.
    Then there exists $\delta\in\reals$ such that $0 < \delta$ and
        for all $B$ such that $B\subseteq X$ and
            for all $x,y\in B$ we have $d((x,y)) < \delta$
        there exists $U\in A$ such that $B\subseteq U$.
\end{lemma}
\begin{proof}
    Suppose not.
    Then for all $\delta\in\reals$ such that $0 < \delta$ there exists $B$ such that
        $B\subseteq X$ and
        for all $U\in A$ we have $B\not\subseteq U$ and
        for all $x,y\in B$ we have $d((x,y)) < \delta$.
    Take $x_0$ such that $x_0\in X$ by \cref{nonempty_inhabited}.
    Hence $x_0\in\unions{A}$ by \cref{open_covering,covering,subseteq}.
    Take $U_0\in A$ such that $x_0\in U_0$ by \cref{unions_iff}.

    Let $R = \{w\in\naturalsPlus\times X \mid
        \text{for all $U\in A$ we have $\metricBall{d}{X}{\snd{w}}{1 / (\fst{w})}\not\subseteq U$}\}$.
    Hence $R$ is a relation by \cref{times_elem_is_tuple,relation}.

    Show that for all $n\in\naturalsPlus$ there exists $x$ such that $n\mathrel{R} x$.
    \begin{subproof}
        Fix $n\in\naturalsPlus$.
        Let $\delta = 1 / n$.
        $0 < \delta$ by \cref{positive_reciprocal_naturalsplus}.
        Thus $n\in\reals$ and $0 < n$
            by \cref{real_naturals_subset_reals,subseteq,real_mul_ident,real_add_ident,real_zero_lt_one,real_one_leq_naturalsplus,real_lt_subst_right,real_lt_trans}.
        Thus $\inv{n}\in\reals$ and $\delta\in\reals$
            by \cref{real_pos_nonzero,real_mul_inverse,real_div,real_mul_closed}.
        Take $B$ such that
            $B\subseteq X$ and
            for all $U\in A$ we have $B\not\subseteq U$ and
            for all $x,y\in B$ we have $d((x,y)) < \delta$.
        Then $B\neq\emptyset$ by \cref{emptyset_subseteq}.
        Take $x$ such that $x\in B$.
        For all $U\in A$ we have $\metricBall{d}{X}{x}{\delta}\not\subseteq U$
            by \cref{subseteq,metric_ball,subseteq_transitive}.
        Thus $x\in X$ by \cref{subseteq}.
        Hence $n\mathrel{R} x$.
    \end{subproof}

    Take $s$ such that $s$ is a function on $\naturalsPlus$ and
        for all $n\in\naturalsPlus$ we have $n\mathrel{R} \apply{s}{n}$ by \cref{choice_function}.
    For all $n\in\dom{s}$ we have $s(n)\in X$ by \cref{times_tuple_elim}.
    Hence $s\in\funs{\naturalsPlus}{X}$ by \cref{funs_intro}.
    Hence $s$ is a sequence in $X$ by \cref{sequence_in}.

    Take $t,x$ such that
        $t$ is a sequence in $X$ and
        $t$ is a subsequence of $s$ and
        $x\in X$ and
        $t$ converges to $x$ in $X$ with respect to $T$ by \cref{sequentially_compact}.
    Take $u$ such that
        $u$ is strictly increasing on $\naturalsPlus$ and
        for all $n\in\naturalsPlus$ we have $t(n) = \apply{s}{\apply{u}{n}}$ by \cref{subsequence}.

    Show that for all $n\in\naturalsPlus$ we have $n \leq u(n)$.
    \begin{subproof}
        Let $P = \{n\in\naturalsPlus \mid n \leq u(n)\}$.
        Show that for all $n\in\naturalsPlus$ we have $n\in P$.
        \begin{subproof}
            For all $n$ such that $n\in P$ we have $n\in\naturalsPlus$.
            Thus $P\subseteq\naturalsPlus$ by \cref{subseteq}.
            Thus $1\in P$ by \cref{real_one_in_naturals,strictly_increasing_on_set,funs_type_apply,real_one_leq_naturalsplus}.
            Show that for all $n\in P$ we have $n + 1\in P$.
            \begin{subproof}
                Fix $n$ such that $n\in P$.
                Thus $n \leq u(n)$.
                Thus $n\in\naturalsPlus$.
                Thus $n + 1\in\naturalsPlus$ and $n < n + 1$
                    by \cref{real_naturals_closed_succ,real_naturals_lt_succ_local}.
                $n\in\reals$ by \cref{real_naturals_subset_reals,subseteq}.
                Thus $u(n) < u(n + 1)$ by \cref{strictly_increasing_on_set}.
                Show that it is wrong that $u(n + 1) < n + 1$.
                \begin{subproof}
                    Suppose not.
                    Then $u(n + 1) < n + 1$.
                    Thus $u(n + 1)\in\naturalsPlus$ by \cref{strictly_increasing_on_set,funs_type_apply}.
                    Thus $u(n + 1) < n$ or $u(n + 1) = n$ or $u(n + 1) = n + 1$
                        by \cref{real_succ_discrete}.
                    \begin{byCase}
                    \caseOf{$u(n + 1) < n$.}
                        Thus $u(n)\in\reals$ and $u(n + 1)\in\reals$
                            by \cref{strictly_increasing_on_set,funs_type_apply,real_naturals_subset_reals,subseteq}.
                        $n < u(n)$ or $n = u(n)$ by \cref{real_leq_split}.
                        \begin{byCase}
                        \caseOf{$n < u(n)$.}
                            Contradiction by \cref{real_lt_trans,real_lt_irreflexive}.
                        \caseOf{$n = u(n)$.}
                            Contradiction by \cref{real_lt_subst_right,real_lt_trans,real_lt_irreflexive}.
                        \end{byCase}
                    \caseOf{$u(n + 1) = n$.}
                        Thus $u(n)\in\reals$
                            by \cref{strictly_increasing_on_set,funs_type_apply,real_naturals_subset_reals,subseteq}.
                        Thus $u(n) < n$ by \cref{real_lt_subst_right}.
                        $n < u(n)$ or $n = u(n)$ by \cref{real_leq_split}.
                        \begin{byCase}
                        \caseOf{$n < u(n)$.}
                            Contradiction by \cref{real_lt_trans,real_lt_irreflexive}.
                        \caseOf{$n = u(n)$.}
                            Contradiction by \cref{real_lt_subst_left,real_lt_irreflexive}.
                        \end{byCase}
                    \caseOf{$u(n + 1) = n + 1$.}
                        Contradiction by \cref{real_lt_subst_left,real_naturals_closed_succ,real_naturals_subset_reals,subseteq,real_lt_irreflexive}.
                    \end{byCase}
                \end{subproof}
                Thus $u(n + 1)\in\reals$ and $n + 1\in\reals$
                    by \cref{strictly_increasing_on_set,funs_type_apply,real_naturals_closed_succ,real_naturals_subset_reals,subseteq}.
                Thus $n + 1 < u(n + 1)$ or $n + 1 = u(n + 1)$ or $u(n + 1) < n + 1$
                    by \cref{real_lt_trichotomy}.
                Thus $n + 1 < u(n + 1)$ or $n + 1 = u(n + 1)$.
                Thus $n + 1 \leq u(n + 1)$.
                Thus $n + 1\in P$.
            \end{subproof}
            Thus for all $n\in\naturalsPlus$ we have $n\in P$ by \cref{real_naturals_induction}.
        \end{subproof}
        Thus for all $n\in\naturalsPlus$ we have $n \leq u(n)$ by \cref{subseteq}.
    \end{subproof}

    Hence $x\in\unions{A}$ by \cref{open_covering,covering,subseteq}.
    Take $U\in A$ such that $x\in U$ by \cref{unions_iff}.
    Therefore $U$ is open in $X$ with respect to $T$ and $U\subseteq X$
        by \cref{metric_topology,basis_topology_is_topology,metric_basis_is_basis,open_covering,open_in,topology_on,subseteq,powerset_elim}.
    Take $\delta\in\reals$ such that $0 < \delta$ and $\metricBall{d}{X}{x}{\delta}\subseteq U$
        by \cref{metric_open_iff}.
    Let $m_2 = 1 + 1$.
    Let $\epsilon = \delta / m_2$.
    $1\in\reals$ and $0\in\reals$ and $0 < 1$ by \cref{real_mul_ident,real_add_ident,real_zero_lt_one}.
    Then $m_2\in\reals$ and $0 < m_2$
        by \cref{real_add_closed,real_lt_add,real_add_ident,real_lt_subst_left,real_lt_trans}.
    $m_2\neq 0$ by \cref{real_pos_nonzero}.
    Then $\inv{m_2}\in\reals$ and $0 < \inv{m_2}$ by \cref{real_mul_inverse,real_inv_pos}.
    $\epsilon = \delta \rmul \inv{m_2}$ by \cref{real_div}.
    Then $\epsilon\in\reals$ and $0 < \epsilon$
        by \cref{real_mul_closed,real_add_ident,real_lt_mul_pos,real_mul_zero,real_lt_subst_left}.
    Take $N\in\naturalsPlus$ such that $0 < 1 / N$ and $1 / N < \epsilon$ by \cref{real_small_reciprocal}.

    Show that $\metricBall{d}{X}{x}{\epsilon}$ is open in $X$ with respect to $T$.
    \begin{subproof}
            Thus $\metricBall{d}{X}{x}{\epsilon}\subseteq X$ by \cref{metric_ball,subseteq}.
            Show that for all $y\in\metricBall{d}{X}{x}{\epsilon}$ there exists $\eta\in\reals$ such that
                $0 < \eta$ and $\metricBall{d}{X}{y}{\eta}\subseteq\metricBall{d}{X}{x}{\epsilon}$.
            \begin{subproof}
                Fix $y$ such that $y\in\metricBall{d}{X}{x}{\epsilon}$.
                Thus $y\in X$ and $d((x,y)) < \epsilon$ by \cref{metric_ball}.
                Thus $d((x,y))\in\reals$ by \cref{metric_on,times_tuple_intro,funs_type_apply}.
                Let $\eta = \epsilon - d((x,y))$.
                $0 < \eta$ by \cref{real_lt_imp_pos_diff}.
                Show that for all $z$ such that $z\in\metricBall{d}{X}{y}{\eta}$ we have $z\in\metricBall{d}{X}{x}{\epsilon}$.
                \begin{subproof}
                        Fix $z$ such that $z\in\metricBall{d}{X}{y}{\eta}$.
                        Thus $z\in X$ and $d((y,z)) < \eta$ by \cref{metric_ball}.
                        Thus $d((y,z))\in\reals$ by \cref{metric_on,times_tuple_intro,funs_type_apply}.
                        Thus $\eta\in\reals$ by \cref{real_add_inverse,real_sub,real_add_closed}.
                        Thus $d((x,z)) \leq d((x,y)) + d((y,z))$ by \cref{metric_on,metric_triangle_inequality}.
                        Thus $d((x,z))\in\reals$ by \cref{metric_on,times_tuple_intro,funs_type_apply}.
                        Thus $\eta + d((x,y))\in\reals$ and $d((x,y)) + d((y,z))\in\reals$
                            by \cref{real_add_closed}.
                        Thus $d((x,y)) + d((y,z)) < \eta + d((x,y))$
                            by \cref{real_lt_add,real_add_comm,real_lt_subst_left}.
                        $d((x,z)) < d((x,y)) + d((y,z))$ or
                            $d((x,z)) = d((x,y)) + d((y,z))$ by \cref{real_leq_split}.
                        Thus $d((x,z)) < \eta + d((x,y))$
                            by \cref{real_lt_trans,metric_on,times_tuple_intro,funs_type_apply,real_add_closed,real_lt_subst_left}.
                        Thus $z\in\metricBall{d}{X}{x}{\epsilon}$
                            by \cref{real_sub,real_add_assoc,real_add_comm,real_add_inverse,real_add_ident,metric_ball,real_lt_subst_right}.
                \end{subproof}
                Thus $\metricBall{d}{X}{y}{\eta}\subseteq\metricBall{d}{X}{x}{\epsilon}$ by \cref{subseteq}.
            \end{subproof}
        Hence $\metricBall{d}{X}{x}{\epsilon}$ is open in $X$ with respect to $T$ by \cref{metric_open_iff}.
    \end{subproof}
    Therefore $\metricBall{d}{X}{x}{\epsilon}$ is a neighborhood of $x$ in $X$ with respect to $T$
        by \cref{metric_zero_characterization,metric_on,metric_ball,real_lt_subst_left,neighborhood}.

    Take $N_0\in\naturalsPlus$ such that for all $n\in\naturalsPlus$ if $N_0 \leq n$ then
        $t(n)\in\metricBall{d}{X}{x}{\epsilon}$ by \cref{converges_to}.
    Show that there exists $n\in\naturalsPlus$ such that $N_0 \leq n$ and $N \leq n$.
    \begin{subproof}
        \begin{byCase}
        \caseOf{$N_0 = N$.}
            Take $n$ such that $n = N$.
            Thus $N_0 \leq n$ and $N \leq n$.
        \caseOf{$N_0 \neq N$.}
            Thus $N_0 < N$ or $N < N_0$ by \cref{real_naturals_subset_reals,subseteq,real_lt_trichotomy}.
            \begin{byCase}
            \caseOf{$N_0 < N$.}
                Take $n$ such that $n = N$.
                Thus $N_0 \leq n$ and $N \leq n$ by \cref{real_lt_imp_leq}.
            \caseOf{$N < N_0$.}
                Take $n$ such that $n = N_0$.
                Thus $N_0 \leq n$ and $N \leq n$ by \cref{real_lt_imp_leq}.
            \end{byCase}
        \end{byCase}
    \end{subproof}
    Take $n\in\naturalsPlus$ such that $N_0 \leq n$ and $N \leq n$.

    Hence $t(n)\in\metricBall{d}{X}{x}{\epsilon}$.
    Then $d((x,t(n))) < \epsilon$ by \cref{metric_ball}.
    $u(n)\in\naturalsPlus$ by \cref{strictly_increasing_on_set,funs_type_apply}.
    Hence $n \leq u(n)$.
    Therefore $N \leq u(n)$
        by \cref{real_naturals_subset_reals,subseteq,strictly_increasing_on_set,funs_type_apply,real_leq_trans}.
    Hence $1 / u(n) \leq 1 / N$ by \cref{naturalsplus_reciprocal_antitone}.
    $u(n)\in\reals$ and $N\in\reals$ by \cref{real_naturals_subset_reals,subseteq}.
    $0 < u(n)$ and $0 < N$
        by \cref{real_zero_lt_one,real_one_in_naturals,real_one_leq_naturalsplus,real_lt_trans,strictly_increasing_on_set,funs_type_apply}.
    Thus $1 / u(n)\in\reals$ and $1 / N\in\reals$
        by \cref{real_pos_nonzero,real_div,real_mul_closed,real_mul_inverse}.
    Thus $1 / u(n) < \epsilon$ by \cref{real_lt_trans,real_lt_subst_left}.

    Show that for all $z$ such that $z\in\metricBall{d}{X}{t(n)}{1 / u(n)}$ we have $z\in U$.
    \begin{subproof}
            Fix $z$ such that $z\in\metricBall{d}{X}{t(n)}{1 / u(n)}$.
            Thus $z\in X$ and $d((t(n),z)) < 1 / u(n)$ by \cref{metric_ball}.
            Thus $t(n)\in X$ by \cref{sequence_in,funs_type_apply}.
            Thus $d((x,z)) \leq d((x,t(n))) + d((t(n),z))$ by \cref{metric_on,metric_triangle_inequality}.
            Show that $d((x,z)) < \epsilon + \epsilon$.
            \begin{subproof}
                Thus $d((t(n),z))\in\reals$ by \cref{metric_on,times_tuple_intro,funs_type_apply}.
                $1 / u(n)\in\reals$.
                $\epsilon\in\reals$.
                Thus $d((t(n),z)) < \epsilon$ by \cref{real_lt_trans}.
                Thus $d((x,t(n)))\in\reals$ by \cref{metric_on,times_tuple_intro,funs_type_apply}.
                Thus $d((x,t(n))) + d((t(n),z)) < d((x,t(n))) + \epsilon$
                    by \cref{real_lt_add,real_add_comm,real_lt_subst_left,real_lt_subst_right}.
                Thus $d((x,t(n))) + \epsilon < \epsilon + \epsilon$ by \cref{real_lt_add}.
                Thus $d((x,t(n))) + d((t(n),z))\in\reals$ and $d((x,t(n))) + \epsilon\in\reals$
                    and $\epsilon + \epsilon\in\reals$ by \cref{real_add_closed}.
                Thus $d((x,t(n))) + d((t(n),z)) < \epsilon + \epsilon$ by \cref{real_lt_trans}.
                Thus $d((x,z))\in\reals$ by \cref{metric_on,times_tuple_intro,funs_type_apply}.
                $d((x,z)) < d((x,t(n))) + d((t(n),z))$ or
                    $d((x,z)) = d((x,t(n))) + d((t(n),z))$ by \cref{real_leq_split}.
                Thus $d((x,z)) < \epsilon + \epsilon$ by \cref{real_lt_trans,real_lt_subst_left}.
            \end{subproof}
            Thus $d((x,z)) < \delta$
                by \cref{real_div,real_mul_assoc,real_mul_inverse,real_mul_comm,real_mul_ident,real_right_distrib,real_lt_subst_right}.
            Thus $z\in U$ by \cref{metric_ball,subseteq}.
    \end{subproof}
    Hence $\metricBall{d}{X}{t(n)}{1 / u(n)}\subseteq U$ by \cref{subseteq}.

    Then $(u(n),t(n))\in R$ by \cref{pair_eq_iff}.
    Therefore for all $V\in A$ we have $\metricBall{d}{X}{t(n)}{1 / u(n)}\not\subseteq V$.
    Then $\metricBall{d}{X}{t(n)}{1 / u(n)}\not\subseteq U$.
    Contradiction.
\end{proof}

% from §28 helper lemma 11: iteration sequence on naturalsPlus
\begin{lemma}\label{iteration_sequence_naturalsplus}
    Suppose $f\in\funs{A}{A}$.
    Suppose $a_0\in A$.
    Then there exists $u$ such that
        $u$ is a sequence in $A$ and
        $u(1) = a_0$ and
        for all $n\in\naturalsPlus$ we have $u(n + 1) = \apply{f}{u(n)}$.
\end{lemma}
\begin{proof}
    Let $S = \iterationSet{f}{a_0}$.
    Show that for all $n\in\naturalsPlus$ we have $n\in S$.
    \begin{subproof}
        $S\subseteq\naturalsPlus$ by \cref{iteration_set_elim,subseteq}.
        $1\in\naturalsPlus$ by \cref{real_one_in_naturals}.
        Show that $1\in S$.
        \begin{subproof}
            Let $u = \{(1,a_0)\}$.
            Thus $u$ is a relation by \cref{singleton_iff,relation}.
            For all $x,y,y'$ such that $(x,y)\in u$ and $(x,y')\in u$ we have $y = y'$
                by \cref{singleton_elim,pair_eq_iff}.
            Therefore $u$ is a function by \cref{rightunique}.
            For all $x$ such that $x\in\dom{u}$ we have $x\in\initialSegment{1}$
                by \cref{dom,singleton_iff,pair_eq_iff,real_one_in_naturals,initial_segment_naturalsplus}.
            Thus $\dom{u}\subseteq\initialSegment{1}$ by \cref{subseteq}.
            Show that for all $x$ such that $x\in\initialSegment{1}$ we have $x\in\dom{u}$.
            \begin{subproof}
                Fix $x$ such that $x\in\initialSegment{1}$.
                Thus $x\in\naturalsPlus$ and $x \leq 1$ by \cref{initial_segment_naturalsplus}.
                $1 < x$ or $1 = x$ by \cref{real_one_leq_naturalsplus}.
                Thus $x = 1$.
                Thus $x\in\dom{u}$ by \cref{singleton_intro,dom}.
            \end{subproof}
            Hence $\dom{u} = \initialSegment{1}$ by \cref{subseteq,subseteq_antisymmetric}.
            Then $u(1) = a_0$ by \cref{singleton_intro,function_apply_intro}.
            Show that for all $k\in\naturalsPlus$ such that $k + 1 \leq 1$ we have $u(k + 1) = \apply{f}{u(k)}$.
            \begin{subproof}
                Fix $k\in\naturalsPlus$.
                Assume $k + 1 \leq 1$.
                Suppose not.
                Then $k < k + 1$ by \cref{real_naturals_lt_succ_local}.
                Show that $1 < k + 1$.
                \begin{subproof}
                    $1\in\reals$ and $k\in\reals$ and $k + 1\in\reals$
                        by \cref{real_naturals_subset_reals,real_one_in_naturals,real_naturals_closed_succ,subseteq}.
                    $1 < k$ or $1 = k$ by \cref{real_one_leq_naturalsplus}.
                    \begin{byCase}
                    \caseOf{$1 = k$.}
                        Thus $1 < k + 1$ by \cref{real_lt_subst_left}.
                    \caseOf{$1 < k$.}
                        Thus $1 < k + 1$ by \cref{real_lt_trans}.
                    \end{byCase}
                \end{subproof}
                $k + 1 < 1$ or $k + 1 = 1$ by \cref{real_leq_split}.
                Show that $1 < 1$.
                \begin{subproof}
                    $1\in\reals$ and $k + 1\in\reals$
                        by \cref{real_naturals_subset_reals,real_one_in_naturals,real_naturals_closed_succ,subseteq}.
                    \begin{byCase}
                    \caseOf{$k + 1 < 1$.}
                        Thus $1 < 1$ by \cref{real_lt_trans}.
                    \caseOf{$k + 1 = 1$.}
                        Thus $1 < 1$ by \cref{real_lt_subst_right}.
                    \end{byCase}
                \end{subproof}
                Contradiction by \cref{real_naturals_subset_reals,real_one_in_naturals,subseteq,real_lt_irreflexive}.
            \end{subproof}
            Therefore $1\in S$ by \cref{iteration_sequence,iteration_set_intro}.
        \end{subproof}
        Show that for all $n\in S$ we have $n + 1\in S$.
        \begin{subproof}
            Fix $n$ such that $n\in S$.
            Thus $n\in\naturalsPlus$ by \cref{iteration_set_elim}.
            Take $u$ such that
                $u$ is an iteration sequence for $f$ on $n$ starting at $a_0$
                by \cref{iteration_set_elim}.
            Thus $n + 1\in\naturalsPlus$ by \cref{real_naturals_closed_succ}.
            Let $v = u\union\{(n + 1,\apply{f}{u(n)})\}$.
            Thus $v$ is a relation by \cref{singleton_iff,relation,iteration_sequence,union_relations_is_relation}.
            Show that $v$ is right-unique.
            \begin{subproof}
                Show that for all $x,y,y'$ such that $(x,y)\in v$ and $(x,y')\in v$ we have $y = y'$.
                \begin{subproof}
                    Fix $x$.
                    Fix $y$.
                    Fix $y'$.
                    Assume $(x,y)\in v$ and $(x,y')\in v$.
                    \begin{byCase}
                    \caseOf{$x = n + 1$.}
                        $(n + 1,y)\in v$.
                        $(n + 1,y')\in v$.
                        Show that $(n + 1,y)\in\{(n + 1,\apply{f}{u(n)})\}$.
                        \begin{subproof}
                            Suppose not.
                            Thus $(n + 1,y)\in u$ by \cref{union_iff}.
                            Thus $n + 1\in\dom{u}$ by \cref{dom}.
                            Thus $\dom{u} = \initialSegment{n}$ by \cref{iteration_sequence}.
                            Thus $n + 1\in\initialSegment{n}$.
                            Thus $n + 1 \leq n$ by \cref{initial_segment_naturalsplus}.
                            Thus $n < n + 1$ by \cref{real_naturals_lt_succ_local}.
                            $n\in\reals$ and $n + 1\in\reals$
                                by \cref{real_naturals_subset_reals,real_naturals_closed_succ,subseteq}.
                            $n + 1 < n$ or $n + 1 = n$ by \cref{real_leq_split}.
                            \begin{byCase}
                            \caseOf{$n + 1 < n$.}
                                Thus $n < n$ by \cref{real_lt_trans}.
                                Contradiction by \cref{real_lt_irreflexive}.
                            \caseOf{$n + 1 = n$.}
                                Thus $n < n$ by \cref{real_lt_subst_right}.
                                Contradiction by \cref{real_lt_irreflexive}.
                            \end{byCase}
                        \end{subproof}
                        Thus $(n + 1,y) = (n + 1,\apply{f}{u(n)})$ by \cref{singleton_iff}.
                        Thus $y = \apply{f}{u(n)}$ by \cref{pair_eq_iff}.
                        Show that $(n + 1,y')\in\{(n + 1,\apply{f}{u(n)})\}$.
                        \begin{subproof}
                            Suppose not.
                            Thus $(n + 1,y')\in u$ by \cref{union_iff}.
                            Thus $n + 1\in\dom{u}$ by \cref{dom}.
                            Thus $\dom{u} = \initialSegment{n}$ by \cref{iteration_sequence}.
                            Thus $n + 1\in\initialSegment{n}$.
                            Thus $n + 1 \leq n$ by \cref{initial_segment_naturalsplus}.
                            Thus $n < n + 1$ by \cref{real_naturals_lt_succ_local}.
                            $n\in\reals$ and $n + 1\in\reals$
                                by \cref{real_naturals_subset_reals,real_naturals_closed_succ,subseteq}.
                            $n + 1 < n$ or $n + 1 = n$ by \cref{real_leq_split}.
                            \begin{byCase}
                            \caseOf{$n + 1 < n$.}
                                Thus $n < n$ by \cref{real_lt_trans}.
                                Contradiction by \cref{real_lt_irreflexive}.
                            \caseOf{$n + 1 = n$.}
                                Thus $n < n$ by \cref{real_lt_subst_right}.
                                Contradiction by \cref{real_lt_irreflexive}.
                            \end{byCase}
                        \end{subproof}
                        Thus $(n + 1,y') = (n + 1,\apply{f}{u(n)})$ by \cref{singleton_iff}.
                        Thus $y' = \apply{f}{u(n)}$ by \cref{pair_eq_iff}.
                        Thus $y = y'$.
                    \caseOf{$x \neq n + 1$.}
                        Show that $(x,y)\in u$.
                        \begin{subproof}
                            Suppose not.
                            Thus $x = n + 1$ by \cref{union_iff,singleton_iff,pair_eq_iff}.
                            Contradiction.
                        \end{subproof}
                        Show that $(x,y')\in u$.
                        \begin{subproof}
                            Suppose not.
                            Thus $x = n + 1$ by \cref{union_iff,singleton_iff,pair_eq_iff}.
                            Contradiction.
                        \end{subproof}
                        Thus $y = y'$ by \cref{iteration_sequence,function_rightunique}.
                    \end{byCase}
                \end{subproof}
            \end{subproof}
            Therefore $v$ is a function by \cref{rightunique}.
            Show that for all $x$ such that $x\in\dom{v}$ we have $x\in\initialSegment{n + 1}$.
            \begin{subproof}
                Fix $x$ such that $x\in\dom{v}$.
                Take $y$ such that $(x,y)\in v$ by \cref{dom}.
                Thus $(x,y)\in u$ or $(x,y)\in\{(n + 1,\apply{f}{u(n)})\}$ by \cref{union_iff}.
                \begin{byCase}
                \caseOf{$(x,y)\in u$.}
                    Thus $x\in\initialSegment{n + 1}$
                        by \cref{dom,iteration_sequence,initial_segment_naturalsplus,real_naturals_subset_reals,real_naturals_closed_succ,subseteq,real_leq_trans,real_naturals_lt_succ_local}.
                \caseOf{$(x,y)\in\{(n + 1,\apply{f}{u(n)})\}$.}
                    Thus $x\in\initialSegment{n + 1}$
                        by \cref{singleton_iff,pair_eq_iff,initial_segment_naturalsplus}.
                \end{byCase}
            \end{subproof}
            Thus $\dom{v}\subseteq\initialSegment{n + 1}$ by \cref{subseteq}.
            Show that for all $x$ such that $x\in\initialSegment{n + 1}$ we have $x\in\dom{v}$.
            \begin{subproof}
                Fix $x$ such that $x\in\initialSegment{n + 1}$.
                Thus $x\in\naturalsPlus$ and $x \leq n + 1$ by \cref{initial_segment_naturalsplus}.
                \begin{byCase}
                \caseOf{$x = n + 1$.}
                    Thus $x\in\dom{v}$ by \cref{cons_left,union_intro_right,dom}.
                \caseOf{$x \neq n + 1$.}
                    Thus $x \leq n$ by \cref{real_naturals_leq_succ_elim}.
                    Thus $x\in\dom{u}$ by \cref{initial_segment_naturalsplus,iteration_sequence}.
                    Thus $x\in\dom{v}$ by \cref{dom,union_iff}.
                \end{byCase}
            \end{subproof}
            Hence $\dom{v} = \initialSegment{n + 1}$ by \cref{subseteq,subseteq_antisymmetric}.
            Thus $v(1) = a_0$
                by \cref{real_one_in_naturals,real_one_leq_naturalsplus,initial_segment_naturalsplus,iteration_sequence,function_apply_elim,union_intro_left,function_apply_intro}.
            Show that for all $k\in\naturalsPlus$ such that $k + 1 \leq n + 1$ we have $v(k + 1) = \apply{f}{v(k)}$.
            \begin{subproof}
                Fix $k\in\naturalsPlus$.
                Assume $k + 1 \leq n + 1$.
                \begin{byCase}
                \caseOf{$k + 1 = n + 1$.}
                    Thus $k = n$.
                    Thus $v(k + 1) = \apply{f}{u(n)}$
                        by \cref{cons_left,union_intro_right,function_apply_intro}.
                    Thus $v(k + 1) = \apply{f}{u(k)}$.
                    Thus $k \leq n$.
                    Thus $v(k) = u(k)$
                        by \cref{initial_segment_naturalsplus,iteration_sequence,function_apply_elim,union_intro_left,function_apply_intro}.
                    Thus $v(k + 1) = \apply{f}{v(k)}$.
                \caseOf{$k + 1 \neq n + 1$.}
                    Thus $k + 1\in\naturalsPlus$ by \cref{real_naturals_closed_succ}.
                    Thus $k + 1 \leq n$ by \cref{real_naturals_leq_succ_elim}.
                    Thus $(k + 1,u(k + 1))\in u$
                        by \cref{initial_segment_naturalsplus,iteration_sequence,function_apply_elim}.
                    Thus $v(k + 1) = u(k + 1)$
                        by \cref{union_intro_left,dom,function_apply_elim,function_rightunique}.
                    Thus $u(k + 1) = \apply{f}{u(k)}$.
                    Thus $v(k + 1) = \apply{f}{u(k)}$.
                    Thus $(k,u(k))\in v$
                        by \cref{real_naturals_succ_bound_elim_local,initial_segment_naturalsplus,iteration_sequence,function_apply_elim,union_intro_left}.
                    Thus $u(k) = v(k)$ by \cref{dom,function_apply_elim,function_rightunique}.
                    Thus $v(k + 1) = \apply{f}{v(k)}$.
                \end{byCase}
            \end{subproof}
            Hence $v$ is an iteration sequence for $f$ on $n + 1$ starting at $a_0$ by \cref{iteration_sequence_intro}.
            Therefore $n + 1\in S$ by \cref{iteration_set_intro}.
        \end{subproof}
        Thus for all $n\in\naturalsPlus$ we have $n\in S$ by \cref{real_naturals_induction}.
    \end{subproof}
    Let $F = \{u \mid \text{there exists $n\in\naturalsPlus$ such that
        $u$ is an iteration sequence for $f$ on $n$ starting at $a_0$} \}$.
    For all $u\in F$ we have $u$ is a function by \cref{iteration_sequence}.
    Show that for all $u,v$ such that $u\in F$ and $v\in F$ we have $u\subseteq v$ or $v\subseteq u$.
    \begin{subproof}
        Fix $u$.
        Fix $v$.
        Assume $u\in F$ and $v\in F$.
        Take $m\in\naturalsPlus$ such that
            $u$ is an iteration sequence for $f$ on $m$ starting at $a_0$.
        Take $n\in\naturalsPlus$ such that
            $v$ is an iteration sequence for $f$ on $n$ starting at $a_0$.
        Thus $u$ is a function and $\dom{u} = \initialSegment{m}$ and $u(1) = a_0$
            by \cref{iteration_sequence}.
        Thus for all $k\in\naturalsPlus$ such that $k + 1 \leq m$ we have $u(k + 1) = \apply{f}{u(k)}$
            by \cref{iteration_sequence}.
        Thus $v$ is a function and $\dom{v} = \initialSegment{n}$ and $v(1) = a_0$
            by \cref{iteration_sequence}.
        Thus for all $k\in\naturalsPlus$ such that $k + 1 \leq n$ we have $v(k + 1) = \apply{f}{v(k)}$
            by \cref{iteration_sequence}.
        $m = n$ or $m \neq n$.
        \begin{byCase}
        \caseOf{$m = n$.}
            Thus $u\subseteq v$ or $v\subseteq u$ by \cref{iteration_sequence_unique,subseteq_refl}.
        \caseOf{$m \neq n$.}
            Thus $m < n$ or $n < m$ by \cref{real_naturals_subset_reals,subseteq,real_lt_trichotomy}.
            \begin{byCase}
        \caseOf{$m < n$.}
            Let $A_1 = \{k\in\naturalsPlus \mid \text{if $k\leq m$ then $u(k) = v(k)$}\}$.
            Show that for all $k\in\naturalsPlus$ we have $k\in A_1$.
            \begin{subproof}
                $A_1\subseteq\naturalsPlus$.
                Thus $1\in A_1$ by \cref{real_one_in_naturals,iteration_sequence}.
                Show that for all $k\in A_1$ we have $k + 1\in A_1$.
                \begin{subproof}
                    Fix $k$ such that $k\in A_1$.
                    If $k + 1\leq m$ then $u(k + 1) = v(k + 1)$
                        by \cref{real_naturals_succ_bound_elim_local,real_lt_imp_leq,real_add_closed,real_naturals_subset_reals,subseteq,real_leq_trans,iteration_sequence}.
                    Thus $k + 1\in A_1$.
                \end{subproof}
                Thus for all $k\in\naturalsPlus$ we have $k\in A_1$ by \cref{real_naturals_induction}.
            \end{subproof}
            For all $k\in\initialSegment{m}$ we have $u(k) = v(k)$ by \cref{initial_segment_naturalsplus}.
            Show that for all $w$ such that $w\in u$ we have $w\in v$.
            \begin{subproof}
                Fix $w$ such that $w\in u$.
                Take $k\in\dom{u}$ such that $w = (k,u(k))$ by \cref{iteration_sequence,function_member_elim}.
                Thus $k\in\initialSegment{m}$ by \cref{iteration_sequence}.
                Thus $w\in v$
                    by \cref{initial_segment_naturalsplus,real_naturals_subset_reals,subseteq,real_leq_trans,real_lt_imp_leq,iteration_sequence,function_apply_elim}.
            \end{subproof}
            Thus $u\subseteq v$ by \cref{subseteq}.
        \caseOf{$n < m$.}
            Let $A_2 = \{k\in\naturalsPlus \mid \text{if $k\leq n$ then $v(k) = u(k)$}\}$.
            Show that for all $k\in\naturalsPlus$ we have $k\in A_2$.
            \begin{subproof}
                $A_2\subseteq\naturalsPlus$.
                Thus $1\in A_2$ by \cref{real_one_in_naturals,iteration_sequence}.
                Show that for all $k\in A_2$ we have $k + 1\in A_2$.
                \begin{subproof}
                    Fix $k$ such that $k\in A_2$.
                    If $k + 1\leq n$ then $v(k + 1) = u(k + 1)$
                        by \cref{real_naturals_succ_bound_elim_local,real_lt_imp_leq,real_add_closed,real_naturals_subset_reals,subseteq,real_leq_trans,iteration_sequence}.
                    Thus $k + 1\in A_2$.
                \end{subproof}
                Thus for all $k\in\naturalsPlus$ we have $k\in A_2$ by \cref{real_naturals_induction}.
            \end{subproof}
            For all $k\in\initialSegment{n}$ we have $v(k) = u(k)$ by \cref{initial_segment_naturalsplus}.
            Show that for all $w$ such that $w\in v$ we have $w\in u$.
            \begin{subproof}
                Fix $w$ such that $w\in v$.
                Take $k\in\dom{v}$ such that $w = (k,v(k))$ by \cref{iteration_sequence,function_member_elim}.
                Thus $k\in\initialSegment{n}$ by \cref{iteration_sequence}.
                Thus $w\in u$
                    by \cref{initial_segment_naturalsplus,real_naturals_subset_reals,subseteq,real_leq_trans,real_lt_imp_leq,iteration_sequence,function_apply_elim}.
            \end{subproof}
            Thus $v\subseteq u$ by \cref{subseteq}.
        \end{byCase}
        \end{byCase}
    \end{subproof}
    Let $u = \unions{F}$.
    Hence $u$ is a function by \cref{iteration_sequence,unions_of_compatible_family_of_function_is_function}.
    For all $x$ such that $x\in\dom{u}$ we have $x\in\naturalsPlus$
        by \cref{dom,unions_iff,iteration_sequence,initial_segment_naturalsplus}.
    For all $n\in\naturalsPlus$ we have $n\in\dom{u}$
        by \cref{subseteq,iteration_set_elim,iteration_sequence,initial_segment_naturalsplus,dom,unions_iff}.
    Therefore $\dom{u} = \naturalsPlus$ by \cref{subseteq,subseteq_antisymmetric}.
    Show that for all $n\in\naturalsPlus$ we have $u(n)\in A$.
    \begin{subproof}
        Fix $n\in\naturalsPlus$.
        Take $v\in F$ such that $(n,u(n))\in v$
            by \cref{function_apply_elim,unions_iff}.
        Take $n_0\in\naturalsPlus$ such that
            $v$ is an iteration sequence for $f$ on $n_0$ starting at $a_0$.
        Thus $n \leq n_0$ by \cref{dom,iteration_sequence,initial_segment_naturalsplus}.
        Let $B = \{k\in\naturalsPlus \mid \text{if $k\leq n$ then $v(k)\in A$}\}$.
        Show that for all $k\in\naturalsPlus$ we have $k\in B$.
        \begin{subproof}
            $B\subseteq\naturalsPlus$.
            Thus $1\in B$ by \cref{real_one_in_naturals,iteration_sequence}.
            Show that for all $k\in B$ we have $k + 1\in B$.
            \begin{subproof}
                Fix $k$ such that $k\in B$.
                If $k + 1\leq n$ then $v(k + 1)\in A$
                    by \cref{real_naturals_succ_bound_elim_local,funs_type_apply,real_naturals_closed_succ,real_naturals_subset_reals,subseteq,real_leq_trans,iteration_sequence}.
                Thus $k + 1\in B$.
            \end{subproof}
            Thus for all $k\in\naturalsPlus$ we have $k\in B$ by \cref{real_naturals_induction}.
        \end{subproof}
        Thus $v(n)\in A$ by \cref{initial_segment_naturalsplus}.
        Thus $u(n)\in A$ by \cref{function_apply_intro}.
    \end{subproof}
    Take $v\in F$ such that $(1,u(1))\in v$
        by \cref{real_one_in_naturals,function_apply_elim,unions_iff}.
    Take $n_0\in\naturalsPlus$ such that
        $v$ is an iteration sequence for $f$ on $n_0$ starting at $a_0$.
    Thus $u(1) = a_0$ by \cref{iteration_sequence,dom,function_apply_intro}.
    Show that for all $n\in\naturalsPlus$ we have $u(n + 1) = \apply{f}{u(n)}$.
    \begin{subproof}
        Fix $n\in\naturalsPlus$.
        Take $v\in F$ such that $(n + 1,u(n + 1))\in v$
            by \cref{real_naturals_closed_succ,function_apply_elim,unions_iff}.
        Take $n_0\in\naturalsPlus$ such that
            $v$ is an iteration sequence for $f$ on $n_0$ starting at $a_0$.
        Thus $n + 1\in\dom{v}$ by \cref{dom}.
        Thus $v(n + 1) = u(n + 1)$ by \cref{function_apply_intro}.
        Thus $n + 1 \leq n_0$ by \cref{iteration_sequence,initial_segment_naturalsplus}.
        Thus $u(n + 1) = \apply{f}{v(n)}$ by \cref{function_apply_intro,iteration_sequence}.
        Thus $u(n) = v(n)$
            by \cref{real_naturals_succ_bound_elim_local,initial_segment_naturalsplus,iteration_sequence,function_apply_elim,unions_iff,function_apply_intro}.
        Thus $u(n + 1) = \apply{f}{u(n)}$.
    \end{subproof}
    Hence $u\in\funs{\naturalsPlus}{A}$ and $u$ is a sequence in $A$ by \cref{funs_intro,sequence_in}.
    Therefore there exists $u$ such that
        $u$ is a sequence in $A$ and
        $u(1) = a_0$ and
        for all $n\in\naturalsPlus$ we have $u(n + 1) = \apply{f}{u(n)}$.
\end{proof}
% from §28 helper lemma 10: sequential compactness implies finite epsilon-ball cover
\begin{lemma}\label{sequentially_compact_finite_ball_cover}
    Suppose $d$ is a metric on $X$.
    Let $T = \metricTopology{d}{X}$.
    Suppose $X$ is sequentially compact with respect to $T$.
    Suppose $\epsilon\in\reals$.
    Suppose $0 < \epsilon$.
    Then there exists $B$ such that
        $B$ is finite and
        $B$ is a covering of $X$ and
        for all $U\in B$ there exists $x\in X$ such that $U = \metricBall{d}{X}{x}{\epsilon}$.
\end{lemma}
\begin{proof}
    \begin{byCase}
    \caseOf{$X = \emptyset$.}
        Let $B = \emptyset$.
        $B$ is finite by \cref{emptyset_is_finite}.
        For all $U\in B$ we have $\exists x\in X. U = \metricBall{d}{X}{x}{\epsilon}$.
        Hence $B$ is a covering of $X$ by \cref{emptyset_subseteq,unions_emptyset,covering}.
        Therefore there exists $B$ such that
            $B$ is finite and
            $B$ is a covering of $X$ and
            for all $U\in B$ there exists $x\in X$ such that $U = \metricBall{d}{X}{x}{\epsilon}$.
    \caseOf{$X \neq \emptyset$.}
        Suppose not.
        Then for all $B$ such that
            $B$ is finite and
            for all $U\in B$ we have $\exists x\in X. U = \metricBall{d}{X}{x}{\epsilon}$
        we have $B$ is not a covering of $X$.
        Let $F = \{B\in\pow{X} \mid \text{$B$ is finite}\}$.
        Show that for all $B\in F$ there exists $x\in X$ such that
            for all $y\in B$ we have $x\notin\metricBall{d}{X}{y}{\epsilon}$.
        \begin{subproof}
            Fix $B$ such that $B\in F$.
            Thus $B$ is finite and $B\subseteq X$ by \cref{powerset_elim}.
            Let $g = \{(x,\metricBall{d}{X}{x}{\epsilon}) \mid x\in B\}$.
            Show that $g$ is a function on $B$.
            \begin{subproof}
                Show that $g$ is a relation.
                \begin{subproof}
                    For all $w$ such that $w\in g$ there exist $x,y$ such that $w = (x,y)$.
                    Thus $g$ is a relation by \cref{relation}.
                \end{subproof}
                Show that $g$ is right-unique.
                \begin{subproof}
                    For all $x,U_1,U_2$ such that $(x,U_1)\in g$ and $(x,U_2)\in g$ we have $U_1 = U_2$.
                    Thus $g$ is right-unique by \cref{rightunique}.
                \end{subproof}
                Show that $\dom{g} = B$.
                \begin{subproof}
                    Show that for all $x$ such that $x\in\dom{g}$ we have $x\in B$.
                    \begin{subproof}
                        Fix $x$ such that $x\in\dom{g}$.
                        Take $U$ such that $(x,U)\in g$ by \cref{dom}.
                        Take $x_1\in B$ such that $(x,U) = (x_1,\metricBall{d}{X}{x_1}{\epsilon})$.
                        Thus $x\in B$ by \cref{pair_eq_iff}.
                    \end{subproof}
                    Thus $\dom{g}\subseteq B$ by \cref{subseteq}.
                    For all $x$ such that $x\in B$ we have $x\in\dom{g}$ by \cref{dom}.
                    Hence $\dom{g} = B$ by \cref{subseteq,subseteq_antisymmetric}.
                \end{subproof}
                Thus $g$ is a function on $B$.
            \end{subproof}
            Let $G = \img{g}{B}$.
            Thus $G$ is finite by \cref{finite_image_function}.
            Show that for all $U\in G$ we have $\exists x\in X. U = \metricBall{d}{X}{x}{\epsilon}$.
            \begin{subproof}
                Fix $U$ such that $U\in G$.
                Take $x\in B$ such that $(x,U)\in g$ by \cref{img_iff}.
                Take $x_1\in B$ such that $(x,U) = (x_1,\metricBall{d}{X}{x_1}{\epsilon})$.
                Thus there exists $x\in X$ such that $U = \metricBall{d}{X}{x}{\epsilon}$
                    by \cref{pair_eq_iff,subseteq}.
            \end{subproof}
            Hence $X\not\subseteq\unions{G}$ by \cref{covering}.
            For all $U\in G$ we have $U\subseteq X$.
            Thus $\unions{G}\subseteq X$ by \cref{unions_family}.
            Show that $\unions{G}\neq X$.
            \begin{subproof}
                Suppose not.
                Then $X\subseteq\unions{G}$ by \cref{subseteq_refl}.
                Contradiction.
            \end{subproof}
            Thus $\unions{G}\subset X$ by \cref{subset}.
            Take $x\in X$ such that $x\notin\unions{G}$ by \cref{subset_witness}.
            Show that for all $y\in B$ we have $x\notin\metricBall{d}{X}{y}{\epsilon}$.
            \begin{subproof}
                Fix $y$ such that $y\in B$.
                $(y,\metricBall{d}{X}{y}{\epsilon})\in g$.
                Thus $\metricBall{d}{X}{y}{\epsilon}\in G$ by \cref{function_apply_iff,img_iff}.
                Suppose not.
                Then $x\in\metricBall{d}{X}{y}{\epsilon}$.
                Thus $x\in\unions{G}$ by \cref{unions_iff}.
                Contradiction.
            \end{subproof}
            Thus there exists $x\in X$ such that for all $y\in B$ we have $x\notin\metricBall{d}{X}{y}{\epsilon}$.
        \end{subproof}
        Let $R = \{ w\in F\times X \mid
            \text{for all $y\in\fst{w}$ we have $\snd{w}\notin\metricBall{d}{X}{y}{\epsilon}$} \}$.
        $R$ is a relation by \cref{relation,subseteq,times_elem_is_tuple}.
        Show that for all $B\in F$ there exists $x$ such that $B\mathrel{R} x$.
        \begin{subproof}
            Fix $B\in F$.
            Take $x\in X$ such that for all $y\in B$ we have $x\notin\metricBall{d}{X}{y}{\epsilon}$.
            Hence there exists $x$ such that $B\mathrel{R} x$ by \cref{times_tuple_intro,fst_eq,snd_eq}.
        \end{subproof}
        Take $c$ such that $c$ is a function on $F$ and for all $B\in F$ we have $B\mathrel{R} \apply{c}{B}$ by \cref{choice_function}.
        Show that for all $B\in F$ we have $\apply{c}{B}\in X$ and for all $y\in B$ we have $\apply{c}{B}\notin\metricBall{d}{X}{y}{\epsilon}$.
        \begin{subproof}
            Fix $B\in F$.
            Thus $B\mathrel{R}\apply{c}{B}$.
            Thus $\apply{c}{B}\in X$ and for all $y\in B$ we have $\apply{c}{B}\notin\metricBall{d}{X}{y}{\epsilon}$
                by \cref{subseteq,times_tuple_elim,fst_eq,snd_eq}.
        \end{subproof}
        $c\in\funs{F}{X}$ by \cref{funs_intro,subseteq}.
        Let $S = \{ w\in F\times F \mid
            \text{$\snd{w} = \fst{w}\union\{\apply{c}{\fst{w}},\apply{c}{\fst{w}}\}$} \}$.
        $S$ is a relation by \cref{relation,subseteq,times_elem_is_tuple}.
        Show that for all $B\in F$ there exists $D$ such that $B\mathrel{S} D$.
        \begin{subproof}
            Fix $B\in F$.
            Let $D = B\union\{\apply{c}{B},\apply{c}{B}\}$.
            Thus $B$ is finite and $B\subseteq X$ by \cref{powerset_elim}.
            Hence $D\in F$
                by \cref{upair_elim,subseteq,pair_is_finite,union_of_finite_is_finite,union_subsets_is_subset,powerset_intro}.
            Thus there exists $D$ such that $B\mathrel{S} D$ by \cref{times_tuple_intro,fst_eq,snd_eq}.
        \end{subproof}
        Take $g$ such that $g$ is a function on $F$ and for all $B\in F$ we have $B\mathrel{S}\apply{g}{B}$
            by \cref{choice_function}.
        Show that for all $B\in F$ we have $g(B)\in F$ and $g(B) = B\union\{\apply{c}{B},\apply{c}{B}\}$.
        \begin{subproof}
            Fix $B\in F$.
            Thus $B\mathrel{S}\apply{g}{B}$.
            Thus $\apply{g}{B}\in F$ and $g(B) = B\union\{\apply{c}{B},\apply{c}{B}\}$
                by \cref{times_tuple_elim,fst_eq,snd_eq}.
        \end{subproof}
        Hence $g\in\funs{F}{F}$ by \cref{subseteq,funs_intro}.
        Take $x_0$ such that $x_0\in X$ by \cref{nonempty_inhabited}.
        Let $a_0 = \{x_0,x_0\}$.
        Thus $a_0\in F$ by \cref{pair_is_finite,upair_elim,subseteq,powerset_intro}.
        Take $u$ such that
            $u$ is a sequence in $F$ and
            $u(1) = a_0$ and
            for all $n\in\naturalsPlus$ we have $u(n + 1) = \apply{g}{u(n)}$ by \cref{iteration_sequence_naturalsplus}.
        Let $s = c\circ u$.
        $u\in\funs{\naturalsPlus}{F}$ and $c\in\funs{F}{X}$ by \cref{sequence_in,funs}.
        Hence $s$ is a sequence in $X$ by \cref{funs_circ,sequence_in}.
        Then $\ran{u}\subseteq\dom{c}$ by \cref{funs_ran,funs,subseteq}.
        Now $u$ is right-unique and $u$ is a relation and $c$ is right-unique and $c$ is a relation
            by \cref{funs_is_function}.
        Show that for all $n\in\naturalsPlus$ we have
            for all $y\in u(n)$ we have $s(n)\notin\metricBall{d}{X}{y}{\epsilon}$.
        \begin{subproof}
            Fix $n\in\naturalsPlus$.
            Fix $y$ such that $y\in u(n)$.
            Thus $u(n)\in F$ by \cref{funs_type_apply,sequence_in}.
            Thus $u(n)\mathrel{R} c(u(n))$.
            Thus $n\in\dom{u}$ by \cref{funs,subseteq}.
            Thus for all $z\in u(n)$ we have $c(u(n))\notin\metricBall{d}{X}{z}{\epsilon}$ and $s(n) = c(u(n))$
                by \cref{circ_apply}.
            Thus $s(n)\notin\metricBall{d}{X}{y}{\epsilon}$.
        \end{subproof}
        For all $n\in\naturalsPlus$ we have $u(n + 1) = u(n)\union\{c(u(n)),c(u(n))\}$
            by \cref{iteration_sequence_naturalsplus,sequence_in,funs_type_apply,funs_is_function,function_apply_intro}.
        Thus for all $n\in\naturalsPlus$ we have $u(n)\subseteq u(n + 1)$ by \cref{union_upper_left}.
        For all $n\in\naturalsPlus$ we have $s(n)\in u(n + 1)$
            by \cref{iteration_sequence_naturalsplus,sequence_in,funs_type_apply,funs_is_function,function_apply_intro,funs,subseteq,circ_apply,upair_intro_left,union_iff}.
        Show that for all $n\in\naturalsPlus$ we have
            for all $m\in\naturalsPlus$ such that $m < n$ we have $s(m)\in u(n)$.
        \begin{subproof}
            Let $P = \{n\in\naturalsPlus \mid \text{for all $m\in\naturalsPlus$ such that $m < n$ we have $s(m)\in u(n)$}\}$.
            Show that for all $n\in\naturalsPlus$ we have $n\in P$.
            \begin{subproof}
                $P\subseteq\naturalsPlus$ by \cref{subseteq}.
                $1\in\naturalsPlus$ by \cref{real_one_in_naturals}.
                Show that for all $m\in\naturalsPlus$ such that $m < 1$ we have $s(m)\in u(1)$.
                \begin{subproof}
                    Fix $m$ such that $m\in\naturalsPlus$ and $m < 1$.
                    Suppose not.
                    $1 < m$ or $1 = m$ by \cref{real_one_leq_naturalsplus}.
                    Contradiction
                        by \cref{real_naturals_subset_reals,subseteq,real_one_in_naturals,real_lt_subst_right,real_lt_trans,real_lt_irreflexive}.
                \end{subproof}
                Thus $1\in P$.
                Show that for all $n\in P$ we have $n + 1\in P$.
                \begin{subproof}
                    Fix $n$ such that $n\in P$.
                    Show that for all $m\in\naturalsPlus$ such that $m < n + 1$ we have $s(m)\in u(n + 1)$.
                    \begin{subproof}
                        Fix $m$ such that $m\in\naturalsPlus$ and $m < n + 1$.
                        Thus $m \neq n + 1$.
                        Thus $m \leq n$ by \cref{real_naturals_leq_succ_elim}.
                        $m < n$ or $m = n$ by \cref{real_leq_split}.
                        \begin{byCase}
                        \caseOf{$m < n$.}
                            Thus $s(m)\in u(n)$.
                            Thus $s(m)\in u(n + 1)$ by \cref{union_upper_left,subseteq}.
                        \caseOf{$m = n$.}
                            Thus $s(n)\in u(n + 1)$.
                        \end{byCase}
                    \end{subproof}
                    Thus $n + 1\in P$.
                \end{subproof}
                Thus for all $n\in\naturalsPlus$ we have $n\in P$ by \cref{real_naturals_induction}.
            \end{subproof}
            Thus for all $n\in\naturalsPlus$ we have
                for all $m\in\naturalsPlus$ such that $m < n$ we have $s(m)\in u(n)$ by \cref{subseteq}.
        \end{subproof}
        Show that for all $m\in\naturalsPlus$ we have
            for all $n\in\naturalsPlus$ such that $m < n$ we have
                it is wrong that $d((s(m),s(n))) < \epsilon$.
        \begin{subproof}
            Fix $m\in\naturalsPlus$.
            Fix $n\in\naturalsPlus$.
            Assume $m < n$.
            Thus $s(m)\in u(n)$.
            Thus $s(n)\notin\metricBall{d}{X}{s(m)}{\epsilon}$.
            Suppose not.
            Then $s(n)\in\metricBall{d}{X}{s(m)}{\epsilon}$
                by \cref{sequence_in,funs_type_apply,metric_ball}.
            Contradiction.
        \end{subproof}
        Take $t,x$ such that
            $t$ is a sequence in $X$ and
            $t$ is a subsequence of $s$ and
            $x\in X$ and
            $t$ converges to $x$ in $X$ with respect to $T$ by \cref{sequentially_compact}.
        Take $v$ such that
            $v$ is strictly increasing on $\naturalsPlus$ and
            for all $n\in\naturalsPlus$ we have $t(n) = \apply{s}{\apply{v}{n}}$ by \cref{subsequence}.
        Let $m_2 = 1 + 1$.
        Let $\delta = \epsilon / m_2$.
        $1\in\reals$ and $m_2\in\reals$ and $0\in\reals$ and $0 < 1$
            by \cref{real_naturals_subset_reals,real_one_in_naturals,subseteq,real_add_closed,real_add_ident,real_zero_lt_one}.
        Then $0 < m_2$ by \cref{real_lt_add,real_add_ident,real_lt_subst_right,real_lt_trans}.
        $m_2\neq 0$ by \cref{real_pos_nonzero}.
        Thus $\inv{m_2}\in\reals$ and $0 < \inv{m_2}$ by \cref{real_mul_inverse,real_inv_pos}.
        $\delta = \epsilon \rmul \inv{m_2}$ by \cref{real_div}.
        Thus $0 < \delta$ by \cref{real_lt_mul_pos,real_mul_zero,real_lt_subst_right}.
        Show that $\metricBall{d}{X}{x}{\delta}$ is a neighborhood of $x$ in $X$ with respect to $T$.
        \begin{subproof}
                Show that $\metricBall{d}{X}{x}{\delta}$ is open in $X$ with respect to $T$.
                \begin{subproof}
                    Thus $\metricBall{d}{X}{x}{\delta}\in\pow{X}$ by \cref{metric_ball,subseteq,powerset_intro}.
                    $x\in X$ and $\delta\in\reals$ and $0 < \delta$.
                    Thus $\metricBall{d}{X}{x}{\delta}\in\metricBasis{d}{X}$ by \cref{metric_basis}.
                    Hence $\metricBall{d}{X}{x}{\delta}\in T$
                        by \cref{metric_basis_is_basis,basis_elem_open_in_generated,metric_topology}.
                Therefore $\metricBall{d}{X}{x}{\delta}$ is open in $X$ with respect to $T$
                    by \cref{open_in,basis_topology_is_topology,metric_topology,metric_basis_is_basis,basis_for_topology}.
            \end{subproof}
            Thus $\metricBall{d}{X}{x}{\delta}$ is a neighborhood of $x$ in $X$ with respect to $T$
                by \cref{metric_ball,metric_zero_characterization,real_lt_subst_left,metric_on,neighborhood}.
        \end{subproof}
        Take $N\in\naturalsPlus$ such that for all $n\in\naturalsPlus$ if $N\leq n$ then $t(n)\in\metricBall{d}{X}{x}{\delta}$ by \cref{converges_to}.
        $N + 1\in\naturalsPlus$ by \cref{real_naturals_closed_succ}.
        Then $N\leq N + 1$.
        Hence $t(N)\in\metricBall{d}{X}{x}{\delta}$ and $t(N + 1)\in\metricBall{d}{X}{x}{\delta}$ by \cref{converges_to}.
        Therefore $d((x,t(N))) < \delta$ and $d((x,t(N + 1))) < \delta$ by \cref{metric_ball}.
        $t(N)\in X$ and $t(N + 1)\in X$ by \cref{sequence_in,funs_type_apply}.
        Hence $d((t(N),x)) < \delta$ by \cref{metric_on,metric_symmetric,real_lt_subst_left}.
        $d$ satisfies the triangle inequality on $X$ by \cref{metric_on}.
        Thus $d((t(N),t(N + 1))) \leq d((t(N),x)) + d((x,t(N + 1)))$ by \cref{metric_triangle_inequality}.
        Show that $d((t(N),t(N + 1))) < \epsilon$.
        \begin{subproof}
            Show that $d((t(N),x)) + d((x,t(N + 1))) < \epsilon$.
            \begin{subproof}
                $(t(N),x)\in X\times X$ and $(x,t(N + 1))\in X\times X$ and
                    $d\in\funs{X\times X}{\reals}$ by \cref{times_tuple_intro,metric_on}.
                Thus $d((t(N),x))\in\reals$ and $d((x,t(N + 1)))\in\reals$ by \cref{funs_type_apply}.
                $\delta\in\reals$.
                Thus $d((x,t(N + 1))) + d((t(N),x)) < \delta + d((t(N),x))$ by \cref{real_lt_add}.
                $d((x,t(N + 1))) + d((t(N),x)) = d((t(N),x)) + d((x,t(N + 1)))$ and
                    $\delta + d((t(N),x)) = d((t(N),x)) + \delta$ by \cref{real_add_comm}.
                Thus $d((t(N),x)) + d((x,t(N + 1))) < d((t(N),x)) + \delta$
                    by \cref{real_lt_subst_left,real_lt_subst_right}.
                Thus $d((t(N),x)) + \delta < \delta + \delta$ by \cref{real_lt_add}.
                Thus $d((t(N),x)) + d((x,t(N + 1))) < \delta + \delta$
                    by \cref{real_add_closed,real_lt_trans}.
                Therefore $d((t(N),x)) + d((x,t(N + 1))) < \epsilon$
                    by \cref{real_mul_assoc,real_mul_inverse,real_mul_comm,real_mul_ident,real_right_distrib,real_lt_subst_right}.
            \end{subproof}
            Thus $d((t(N),t(N + 1))) < \epsilon$
                by \cref{real_leq_split,times_tuple_intro,metric_on,funs_type_apply,real_add_closed,real_lt_trans,real_lt_subst_left}.
        \end{subproof}
        Thus $v(N)\in\naturalsPlus$ and $v(N + 1)\in\naturalsPlus$ by \cref{strictly_increasing_on_set,funs_type_apply}.
        Thus $N < N + 1$ by \cref{real_naturals_lt_succ_local}.
        Thus $v(N) < v(N + 1)$ by \cref{strictly_increasing_on_set}.
        Thus $t(N) = \apply{s}{\apply{v}{N}}$ and $t(N + 1) = \apply{s}{\apply{v}{(N + 1)}}$.
        Thus it is wrong that $d((\apply{s}{\apply{v}{N}},\apply{s}{\apply{v}{(N + 1)}})) < \epsilon$.
        Thus it is wrong that $d((t(N),t(N + 1))) < \epsilon$
            by \cref{pair_eq_iff,real_lt_subst_left}.
        Contradiction.
    \end{byCase}
\end{proof}

% from §28 Item 6: compactness, limit point compactness, and sequential compactness
\begin{theorem}\label{metrizable_compact_iff_limit_point_compact_iff_sequentially_compact}
    Assume $T$ is a topology on $X$.
    Suppose $X$ is metrizable with respect to $T$.
    Then $X$ is compact with respect to $T$ iff
        $X$ is limit point compact with respect to $T$ and
        $X$ is sequentially compact with respect to $T$.
\end{theorem}
\begin{proof}
    Take $d$ such that $d$ is a metric on $X$ and $T = \metricTopology{d}{X}$.

    If $X$ is compact with respect to $T$ then $X$ is limit point compact with respect to $T$
        by \cref{compact_implies_limit_point_compact}.

    Show that if $X$ is sequentially compact with respect to $T$ then $X$ is compact with respect to $T$.
    \begin{subproof}
        Assume $X$ is sequentially compact with respect to $T$.
        Show that for all $A$ if $A$ is an open covering of $X$ with respect to $T$ then
            there exists $B$ such that $B\subseteq A$ and $B$ is finite and $B$ is a covering of $X$.
        \begin{subproof}
            Fix $A$ such that $A$ is an open covering of $X$ with respect to $T$.
            \begin{byCase}
            \caseOf{$X = \emptyset$.}
                Let $B = \emptyset$.
                Therefore there exists $B$ such that $B\subseteq A$ and $B$ is finite and $B$ is a covering of $X$
                    by \cref{emptyset_subseteq,emptyset_is_finite,unions_emptyset,covering}.
            \caseOf{$X\neq\emptyset$.}
                Take $\delta\in\reals$ such that $0 < \delta$ and
                    for all $B$ such that $B\subseteq X$ and
                        for all $x,y\in B$ we have $d((x,y)) < \delta$
                    there exists $U\in A$ such that $B\subseteq U$
                    by \cref{sequentially_compact_lebesgue}.
                Let $m_2 = 1 + 1$.
                Let $\epsilon = \delta / m_2$.
                $1\in\reals$ and $0\in\reals$ and $0 < 1$
                    by \cref{real_mul_ident,real_add_ident,real_zero_lt_one}.
                Hence $m_2\in\reals$ by \cref{real_add_closed}.
                Therefore $0 < m_2$ by \cref{real_lt_add,real_add_ident,real_lt_subst_left,real_lt_trans}.
                $m_2\neq 0$ by \cref{real_pos_nonzero}.
                Thus $\inv{m_2}\in\reals$ by \cref{real_mul_inverse}.
                Thus $0 < \inv{m_2}$ by \cref{real_inv_pos}.
                $\epsilon = \delta \rmul \inv{m_2}$ by \cref{real_div}.
                Thus $\epsilon\in\reals$ by \cref{real_mul_closed}.
                Thus $0 < \epsilon$ by \cref{real_add_ident,real_lt_mul_pos,real_mul_zero,real_lt_subst_left}.
                Take $B$ such that
                    $B$ is finite and
                    $B$ is a covering of $X$ and
                    for all $U\in B$ there exists $x\in X$ such that $U = \metricBall{d}{X}{x}{\epsilon}$
                    by \cref{sequentially_compact_finite_ball_cover}.

                Show that for all $V\in B$ there exists $U\in A$ such that $V\subseteq U$.
                \begin{subproof}
                    Fix $V$ such that $V\in B$.
                    Take $x_0\in X$ such that $V = \metricBall{d}{X}{x_0}{\epsilon}$.
                    Show that for all $x,y\in V$ we have $d((x,y)) < \delta$.
                    \begin{subproof}
                        Fix $x$.
                        Fix $y$ such that $x\in V$ and $y\in V$.
                        Thus $x\in X$ and $y\in X$ by \cref{metric_ball,subseteq}.
                        Thus $d((x_0,x)) < \epsilon$ and $d((x_0,y)) < \epsilon$ by \cref{metric_ball}.
                        Thus $d((x,y)) \leq d((x,x_0)) + d((x_0,y))$ by \cref{metric_on,metric_triangle_inequality}.
                        Thus $d((x,x_0)) < \epsilon$ by \cref{metric_on,metric_symmetric,real_lt_subst_left}.
                        Show that $d((x,y)) < \epsilon + \epsilon$.
                        \begin{subproof}
                            $(x,x_0)\in X\times X$ and $(x_0,y)\in X\times X$ and $(x,y)\in X\times X$
                                by \cref{times_tuple_intro}.
                            Thus $d((x,x_0))\in\reals$ and $d((x_0,y))\in\reals$ and $d((x,y))\in\reals$
                                by \cref{metric_on,funs_type_apply}.
                            Thus $d((x,x_0)) + d((x_0,y)) < \epsilon + d((x_0,y))$ by \cref{real_lt_add}.
                            $d((x,x_0)) + d((x_0,y)) = d((x_0,y)) + d((x,x_0))$ and
                                $\epsilon + d((x_0,y)) = d((x_0,y)) + \epsilon$ by \cref{real_add_comm}.
                            Thus $d((x,x_0)) + d((x_0,y)) < d((x_0,y)) + \epsilon$
                                by \cref{real_lt_subst_left,real_lt_subst_right}.
                            Thus $d((x_0,y)) + \epsilon < \epsilon + \epsilon$ by \cref{metric_ball,real_lt_add}.
                            Thus $d((x,x_0)) + d((x_0,y))\in\reals$ and $d((x_0,y)) + \epsilon\in\reals$ and
                                $\epsilon + \epsilon\in\reals$ by \cref{real_add_closed}.
                            Thus $d((x,x_0)) + d((x_0,y)) < \epsilon + \epsilon$ by \cref{real_lt_trans}.
                            Thus $d((x,y)) < \epsilon + \epsilon$
                                by \cref{real_leq_split,real_lt_trans,real_lt_subst_left}.
                        \end{subproof}
                        Thus $d((x,y)) < \delta$
                            by \cref{real_div,real_mul_assoc,real_mul_inverse,real_mul_comm,real_mul_ident,real_right_distrib,real_lt_subst_right}.
                    \end{subproof}
                    Thus there exists $U\in A$ such that $V\subseteq U$.
                \end{subproof}

                Let $R = \{w\in B\times A \mid \fst{w}\subseteq \snd{w}\}$.
                For all $V\in B$ there exists $U$ such that $V\mathrel{R} U$.
                Take $h$ such that $h$ is a function on $B$ and for all $V\in B$ we have $V\mathrel{R}\apply{h}{V}$
                    by \cref{subseteq,times_elem_is_tuple,relation,choice_function}.
                Let $C = \img{h}{B}$.
                Hence $C$ is finite by \cref{finite_image_function}.
                Therefore $C\subseteq A$ by \cref{img_iff,function_apply_intro,subseteq,times_tuple_elim}.
                Show that for all $x$ such that $x\in X$ we have $x\in\unions{C}$.
                \begin{subproof}
                    Fix $x$ such that $x\in X$.
                    Thus $x\in\unions{B}$ by \cref{covering,subseteq}.
                    Take $V\in B$ such that $x\in V$ by \cref{unions_iff}.
                    Thus $V\mathrel{R}\apply{h}{V}$.
                    Thus $(V,\apply{h}{V})\in R$.
                    Thus $\fst{(V,\apply{h}{V})}\subseteq \snd{(V,\apply{h}{V})}$.
                    Thus $x\in\apply{h}{V}$ and $\apply{h}{V}\in C$
                        by \cref{fst_eq,snd_eq,subseteq,function_apply_elim,img_iff}.
                    Thus $x\in\unions{C}$ by \cref{unions_iff}.
                \end{subproof}
                Hence $C$ is a covering of $X$ by \cref{subseteq,covering}.
                Therefore there exists $B$ such that $B\subseteq A$ and $B$ is finite and $B$ is a covering of $X$.
            \end{byCase}
        \end{subproof}
        Hence $X$ is compact with respect to $T$ by \cref{compact}.
    \end{subproof}

    Show that if $X$ is limit point compact with respect to $T$ then $X$ is sequentially compact with respect to $T$.
    \begin{subproof}
        Assume $X$ is limit point compact with respect to $T$.
        Show that for all $s$ such that $s$ is a sequence in $X$
            there exist $t,x$ such that
                $t$ is a sequence in $X$ and
                $t$ is a subsequence of $s$ and
                $x\in X$ and
                $t$ converges to $x$ in $X$ with respect to $T$.
        \begin{subproof}
            Fix $s$ such that $s$ is a sequence in $X$.
            $s\in\funs{\naturalsPlus}{X}$ by \cref{sequence_in}.
            Let $A = \ran{s}$.
            Thus $A\subseteq X$ by \cref{funs_ran}.
            \begin{byCase}
            \caseOf{$A$ is finite.}
                Show that $s\in\funs{\naturalsPlus}{A}$.
                \begin{subproof}
                    Thus for all $n\in\dom{s}$ we have $s(n)\in A$
                        by \cref{funs_is_function,function_apply_elim,ran_iff}.
                    Thus $s$ is a function to $A$.
                \end{subproof}
                Thus $s$ is a sequence in $A$ by \cref{sequence_in}.
                Take $t,a$ such that
                    $a\in A$ and
                    $t$ is a sequence in $A$ and
                    $t$ is a subsequence of $s$ and
                    for all $n\in\naturalsPlus$ we have $t(n) = a$
                    by \cref{finite_sequence_constant_subsequence}.
                Thus $a\in X$ by \cref{subseteq}.
                Show that $t\in\funs{\naturalsPlus}{X}$.
                \begin{subproof}
                    Thus for all $n\in\dom{t}$ we have $t(n)\in X$
                        by \cref{sequence_in,funs,funs_type_apply,subseteq}.
                    Thus $t$ is a function to $X$.
                \end{subproof}
                Thus $t$ is a sequence in $X$ by \cref{sequence_in}.
                Hence $t$ converges to $a$ in $X$ with respect to $T$
                    by \cref{neighborhood,real_one_in_naturals,subseteq,converges_to}.
                Therefore there exist $t,a$ such that
                    $t$ is a sequence in $X$ and
                    $t$ is a subsequence of $s$ and
                    $a\in X$ and
                    $t$ converges to $a$ in $X$ with respect to $T$.
            \caseOf{$A$ is infinite.}
                Take $x\in X$ such that $x$ is a limit point of $A$ in $X$ with respect to $T$
                    by \cref{limit_point_compact}.

                Hence $X$ satisfies the T-one axiom with respect to $T$
                    by \cref{metrizable,metric_space_hausdorff,finite_point_set_closed_hausdorff,t1_axiom}.

                For all $U$ such that $U$ is a neighborhood of $x$ in $X$ with respect to $T$
                    we have $U\inter A$ is infinite by \cref{limit_point_infinite_neighborhood}.

                Show that for all $n\in\naturalsPlus$ we have
                    $\metricBall{d}{X}{x}{1 / n}$ is a neighborhood of $x$ in $X$ with respect to $T$.
                \begin{subproof}
                    Fix $n\in\naturalsPlus$.
                    Thus $\metricBall{d}{X}{x}{1 / n}$ is a neighborhood of $x$ in $X$ with respect to $T$
                        by \cref{positive_reciprocal_naturalsplus,real_naturals_subset_reals,subseteq,real_one_in_naturals,real_add_ident,real_zero_lt_one,real_one_leq_naturalsplus,real_lt_subst_right,real_lt_trans,real_pos_nonzero,real_mul_inverse,real_div,real_mul_closed,metric_ball,powerset_intro,metric_basis,metric_basis_is_basis,basis_elem_open_in_generated,metric_topology,basis_for_topology,basis_topology_is_topology,open_in,metric_zero_characterization,metric_on,real_lt_subst_left,neighborhood}.
                \end{subproof}

                Let $D = \{w\in\naturalsPlus\times\naturalsPlus \mid
                    \apply{s}{\snd{w}}\in\metricBall{d}{X}{x}{1 / (\fst{w})}\}$.

                Show that there exists $a_0\in D$ such that $\fst{a_0} = 1$.
                \begin{subproof}
                    Let $U_1 = \metricBall{d}{X}{x}{1 / 1}$.
                    Hence $U_1\inter A$ is infinite by \cref{real_one_in_naturals,subseteq}.
                    Take $y$ such that $y\in U_1\inter A$ by \cref{emptyset_is_finite,nonempty_inhabited}.
                    Take $k$ such that $(k,y)\in s$ by \cref{inter_elim_right,ran_iff}.
                    Thus $k\in\naturalsPlus$ and $s(k) = y$
                        by \cref{dom,funs,subseteq,sequence_in,function_apply_iff,funs_is_function}.
                    Let $a_0 = (1,k)$.
                    Therefore there exists $a_0\in D$ such that $\fst{a_0} = 1$
                        by \cref{real_one_in_naturals,times_tuple_intro,inter_elim_left,fst_eq,snd_eq}.
                \end{subproof}

                Let $R = \{w\in D\times D \mid
                    \exists n,k,n_1,k_1\in\naturalsPlus.
                        w = ((n,k),(n_1,k_1)) \land n_1 = n + 1 \land k < k_1\}$.

                Show that for all $p\in D$ there exists $q$ such that $p\mathrel{R} q$.
                \begin{subproof}
                    Fix $p\in D$.
                    Take $n,k$ such that $p = (n,k)$ by \cref{times_elem_is_tuple}.
                    Thus $n\in\naturalsPlus$ and $k\in\naturalsPlus$ by \cref{times_tuple_elim}.
                    Let $V = \metricBall{d}{X}{x}{1 / (n + 1)}$.
                    Thus $V$ is a neighborhood of $x$ in $X$ with respect to $T$ by \cref{real_naturals_closed_succ,subseteq}.
                    Thus $V\inter A$ is infinite.
                    Let $J = \{m\in\naturalsPlus \mid \apply{s}{m}\in V\}$.
                    Show that $J$ is infinite.
                    \begin{subproof}
                        Suppose not.
                        Then $J$ is finite.
                        Show that $\img{s}{J} = V\inter A$.
                        \begin{subproof}
                            Show that $\img{s}{J}\subseteq V\inter A$.
                            \begin{subproof}
                                Show that for all $y$ such that $y\in\img{s}{J}$ we have $y\in V\inter A$.
                                \begin{subproof}
                                    Fix $y$ such that $y\in\img{s}{J}$.
                                    Take $m\in J$ such that $(m,y)\in s$ by \cref{img_iff}.
                                    Thus $y\in V\inter A$ by \cref{function_apply_intro,funs_is_function,subseteq,ran_iff,inter_intro}.
                                \end{subproof}
                                Thus $\img{s}{J}\subseteq V\inter A$ by \cref{subseteq}.
                            \end{subproof}
                            Show that $V\inter A\subseteq \img{s}{J}$.
                            \begin{subproof}
                                Show that for all $y$ such that $y\in V\inter A$ we have $y\in\img{s}{J}$.
                                \begin{subproof}
                                    Fix $y$ such that $y\in V\inter A$.
                                    Take $m$ such that $(m,y)\in s$ by \cref{inter_elim_right,ran_iff}.
                                    Thus $m\in J$
                                        by \cref{dom,funs,subseteq,sequence_in,function_apply_intro,funs_is_function,inter_elim_left}.
                                    Thus $y\in\img{s}{J}$ by \cref{img_iff}.
                                \end{subproof}
                                Thus $V\inter A\subseteq \img{s}{J}$ by \cref{subseteq}.
                            \end{subproof}
                            Thus $\img{s}{J} = V\inter A$ by \cref{subseteq_antisymmetric}.
                        \end{subproof}
                        Let $g = \restrl{s}{J}$.
                        Thus $g$ is a function on $J$
                            by \cref{sequence_in,funs_is_function,restrl_of_function_is_function,dom_of_restrl_eq_inter,funs_elim,subseteq,inter_absorb_supseteq_left}.
                        Thus $\img{g}{J} = V\inter A$
                            by \cref{sequence_in,funs_elim,subseteq,restrl_img,subseteq_refl,inter_absorb_supseteq_right}.
                        Thus $\img{g}{J}$ is finite by \cref{finite_image_function}.
                        Thus $V\inter A$ is finite.
                        Contradiction.
                    \end{subproof}
                    Show that there exists $k_1\in J$ such that $k < k_1$.
                    \begin{subproof}
                        Suppose not.
                        Then for all $m\in J$ we have $m \leq k$.
                        Thus $J$ is finite
                            by \cref{initial_segment_naturalsplus,subseteq,initial_segment_finite,subseteq_finite_is_finite}.
                        Contradiction.
                    \end{subproof}
                    Take $k_1\in J$ such that $k < k_1$.
                    Let $q = (n + 1,k_1)$.
                    Thus $q\in\naturalsPlus\times\naturalsPlus$ by \cref{real_naturals_closed_succ,subseteq,times_tuple_intro}.
                    Thus $\apply{s}{\snd{q}}\in\metricBall{d}{X}{x}{1 / (n + 1)}$ by \cref{snd_eq}.
                    Thus $q\in D$.
                    Thus $(p,q) = ((n,k),(n + 1,k_1))$ by \cref{real_naturals_closed_succ,pair_eq_iff}.
                    Thus there exists $q$ such that $p\mathrel{R} q$.
                \end{subproof}

                Take $f$ such that $f$ is a function on $D$ and for all $p\in D$ we have
                    $p\mathrel{R}\apply{f}{p}$ by \cref{subseteq,times_elem_is_tuple,relation,choice_function}.
                Hence $f\in\funs{D}{D}$ by \cref{funs_intro,subseteq,times_tuple_elim}.
                Take $a_0\in D$ such that $\fst{a_0} = 1$.
                Take $v$ such that
                    $v$ is a sequence in $D$ and
                    $v(1) = a_0$ and
                    for all $n\in\naturalsPlus$ we have $v(n + 1) = \apply{f}{v(n)}$
                    by \cref{iteration_sequence_naturalsplus}.

                Let $g = \{(p,\snd{p}) \mid p\in D\}$.
                For all $p,y_1,y_2$ such that $(p,y_1)\in g$ and $(p,y_2)\in g$ we have $y_1 = y_2$.
                Then $g$ is a function by \cref{rightunique,relation}.
                Hence $\dom{g} = D$ by \cref{dom_iff,pair_eq_iff,subseteq,subseteq_antisymmetric}.
                Show that for all $p\in\dom{g}$ we have $g(p)\in\naturalsPlus$.
                \begin{subproof}
                    Fix $p\in\dom{g}$.
                    Take $p_1\in D$ such that $(p,g(p)) = (p_1,\snd{p_1})$ by \cref{pair_eq_iff,function_apply_intro}.
                    Thus $g(p)\in\naturalsPlus$ by \cref{pair_eq_iff,subseteq,times_elem_is_tuple,snd_eq}.
                \end{subproof}
                Hence $g\in\funs{D}{\naturalsPlus}$ by \cref{funs_intro}.

                Let $u = g\circ v$.
                Thus $u$ is a sequence in $\naturalsPlus$ by \cref{sequence_in,funs_circ}.

                Show that for all $n\in\naturalsPlus$ we have $u(n) = \snd{v(n)}$.
                \begin{subproof}
                    Fix $n\in\naturalsPlus$.
                    Thus $(n,v(n))\in v$ by \cref{funs,sequence_in,funs_is_function,function_apply_elim}.
                    Thus $(v(n),g(v(n)))\in g$
                        by \cref{sequence_in,funs,funs_type_apply,funs_is_function,function_apply_elim}.
                    Thus $(n,g(v(n)))\in u$ by \cref{circ_iff}.
                    Thus $u(n) = g(v(n))$ by \cref{sequence_in,funs_is_function,function_apply_intro}.
                    Thus $(v(n),\snd{v(n)})\in g$.
                    Thus $g(v(n)) = \snd{v(n)}$ by \cref{function_apply_intro}.
                    Thus $u(n) = \snd{v(n)}$.
                \end{subproof}

                Let $P = \{n\in\naturalsPlus \mid \fst{v(n)} = n\}$.
                Show that for all $n\in\naturalsPlus$ we have $n\in P$.
                \begin{subproof}
                    $P\subseteq\naturalsPlus$ by \cref{subseteq}.
                    Thus $1\in P$ by \cref{real_one_in_naturals,fst_eq,pair_eq_pair_of_proj}.
                    Show that for all $n\in P$ we have $n + 1\in P$.
                    \begin{subproof}
                        Fix $n$ such that $n\in P$.
                        Thus $\fst{v(n)} = n$.
                        Thus $v(n + 1) = \apply{f}{v(n)}$ by \cref{real_naturals_closed_succ,iteration_sequence_naturalsplus}.
                        Thus $v(n)\mathrel{R} v(n + 1)$.
                        Take $n_0,k_0,n_1,k_1\in\naturalsPlus$ such that
                            $(v(n),v(n + 1)) = ((n_0,k_0),(n_1,k_1))$ and
                            $n_1 = n_0 + 1$ and $k_0 < k_1$.
                        Thus $\fst{v(n + 1)} = n + 1$
                            by \cref{pair_eq_iff,fst_eq,real_add_comm,real_add_ident}.
                        Thus $n + 1\in P$.
                    \end{subproof}
                    Thus for all $n\in\naturalsPlus$ we have $n\in P$ by \cref{real_naturals_induction}.
                \end{subproof}
                For all $n\in\naturalsPlus$ we have $\apply{s}{u(n)}\in\metricBall{d}{X}{x}{1 / n}$
                    by \cref{funs_type_apply,sequence_in,snd_eq,subseteq}.

                Show that for all $n\in\naturalsPlus$ we have $u(n) < u(n + 1)$.
                \begin{subproof}
                    Fix $n\in\naturalsPlus$.
                    Thus $v(n + 1) = \apply{f}{v(n)}$ by \cref{iteration_sequence_naturalsplus}.
                    Thus $v(n)\mathrel{R} v(n + 1)$.
                    Take $n_0,k_0,n_1,k_1\in\naturalsPlus$ such that
                        $(v(n),v(n + 1)) = ((n_0,k_0),(n_1,k_1))$ and
                        $n_1 = n_0 + 1$ and $k_0 < k_1$.
                    Thus $u(n) < u(n + 1)$
                        by \cref{pair_eq_iff,snd_eq,real_naturals_closed_succ,real_lt_subst_left,real_lt_subst_right}.
                \end{subproof}

                Show that for all $m,n\in\naturalsPlus$ such that $m < n$ we have $u(m) < u(n)$.
                \begin{subproof}
                    Let $T_0 = \{n\in\naturalsPlus \mid \text{for all $m\in\naturalsPlus$ such that $m < n$ we have $u(m) < u(n)$}\}$.
                    Show that for all $n\in\naturalsPlus$ we have $n\in T_0$.
                    \begin{subproof}
                        $T_0\subseteq\naturalsPlus$ by \cref{subseteq}.
                        Show that $1\in T_0$.
                        \begin{subproof}
                            Show that for all $m\in\naturalsPlus$ such that $m < 1$ we have $u(m) < u(1)$.
                            \begin{subproof}
                                Fix $m$ such that $m\in\naturalsPlus$ and $m < 1$.
                                Suppose not.
                                $1 < m$ or $1 = m$ by \cref{real_one_leq_naturalsplus}.
                                Contradiction
                                    by \cref{real_naturals_subset_reals,subseteq,real_one_in_naturals,real_lt_subst_right,real_lt_trans,real_lt_irreflexive}.
                            \end{subproof}
                            Thus $1\in T_0$ by \cref{real_one_in_naturals}.
                        \end{subproof}
                        Show that for all $n\in T_0$ we have $n + 1\in T_0$.
                        \begin{subproof}
                            Fix $n$ such that $n\in T_0$.
                            Show that for all $m\in\naturalsPlus$ such that $m < n + 1$ we have $u(m) < u(n + 1)$.
                            \begin{subproof}
                                Fix $m$ such that $m\in\naturalsPlus$ and $m < n + 1$.
                                Suppose not.
                                Thus $m < n$ or $m = n$ or $m = n + 1$ by \cref{real_succ_discrete}.
                                \begin{byCase}
                                \caseOf{$m = n$.}
                                    Contradiction by \cref{real_lt_subst_left}.
                                \caseOf{$m < n$.}
                                    Thus $u(m) < u(n)$ by \cref{subseteq}.
                                    Thus $u(n) < u(n + 1)$.
                                    Thus $u(m) < u(n + 1)$
                                        by \cref{sequence_in,real_naturals_closed_succ,funs_type_apply,real_naturals_subset_reals,subseteq,real_lt_trans}.
                                    Contradiction.
                                \caseOf{$m = n + 1$.}
                                    Contradiction
                                        by \cref{real_lt_subst_left,real_naturals_subset_reals,subseteq,real_lt_irreflexive}.
                                \end{byCase}
                            \end{subproof}
                            Thus $n + 1\in T_0$.
                        \end{subproof}
                        Thus for all $n\in\naturalsPlus$ we have $n\in T_0$ by \cref{real_naturals_induction}.
                    \end{subproof}
                    Thus for all $m,n\in\naturalsPlus$ such that $m < n$ we have $u(m) < u(n)$ by \cref{subseteq}.
                \end{subproof}
                Hence $u$ is strictly increasing on $\naturalsPlus$ by \cref{sequence_in,strictly_increasing_on_set}.

                Let $t = s\circ u$.
                Thus $t$ is a sequence in $X$ by \cref{sequence_in,funs_circ}.

                Show that for all $n\in\naturalsPlus$ we have $t(n) = \apply{s}{\apply{u}{n}}$.
                \begin{subproof}
                    Fix $n\in\naturalsPlus$.
                    Thus $(n,u(n))\in u$ by \cref{funs,sequence_in,funs_is_function,function_apply_elim}.
                    Thus $(u(n),s(u(n)))\in s$ by \cref{sequence_in,funs,funs_type_apply,funs_is_function,function_apply_elim}.
                    Thus $(n,s(u(n)))\in t$ by \cref{circ_iff}.
                    Thus $t(n) = s(u(n))$ by \cref{sequence_in,funs_is_function,function_apply_intro}.
                    Thus $t(n) = \apply{s}{\apply{u}{n}}$.
                \end{subproof}
                Hence $t$ is a subsequence of $s$ by \cref{subsequence}.

                Show that for all $U$ if $U$ is a neighborhood of $x$ in $X$ with respect to $T$ then
                    there exists $N\in\naturalsPlus$ such that for all $n\in\naturalsPlus$ if $N\leq n$ then $t(n)\in U$.
                \begin{subproof}
                    Fix $U$ such that $U$ is a neighborhood of $x$ in $X$ with respect to $T$.
                    Thus $U$ is open in $X$ with respect to $T$ and $x\in U$ by \cref{neighborhood}.
                    Take $\epsilon\in\reals$ such that $0 < \epsilon$ and $\metricBall{d}{X}{x}{\epsilon}\subseteq U$
                        by \cref{open_in,topology_on,subseteq,powerset_elim,metric_open_iff}.
                    Take $N\in\naturalsPlus$ such that $0 < 1 / N$ and $1 / N < \epsilon$ by \cref{real_small_reciprocal}.

                        Show that for all $n\in\naturalsPlus$ if $N\leq n$ then $t(n)\in U$.
                        \begin{subproof}
                            Fix $n\in\naturalsPlus$.
                            Assume $N\leq n$.
                            Thus $t(n) = \apply{s}{\apply{u}{n}}$.
                            Thus $\metricBall{d}{X}{x}{1 / n}\subseteq U$
                                by \cref{naturalsplus_reciprocal_antitone,real_leq_split,real_naturals_subset_reals,real_one_in_naturals,subseteq,real_add_ident,real_zero_lt_one,real_one_leq_naturalsplus,real_lt_subst_right,real_lt_trans,real_pos_nonzero,real_mul_inverse,real_div,real_mul_closed,metric_ball_mono,subseteq_transitive}.
                            Thus $t(n)\in U$ by \cref{funs_type_apply,sequence_in,snd_eq,subseteq}.
                        \end{subproof}
                    Thus there exists $N\in\naturalsPlus$ such that for all $n\in\naturalsPlus$ if $N\leq n$ then $t(n)\in U$.
                \end{subproof}
                Hence $t$ converges to $x$ in $X$ with respect to $T$ by \cref{converges_to}.
                Therefore there exist $t,x$ such that
                    $t$ is a sequence in $X$ and
                    $t$ is a subsequence of $s$ and
                    $x\in X$ and
                    $t$ converges to $x$ in $X$ with respect to $T$.
            \end{byCase}
        \end{subproof}
        Hence $X$ is sequentially compact with respect to $T$ by \cref{sequentially_compact}.
    \end{subproof}

    If $X$ is compact with respect to $T$ then
        $X$ is limit point compact with respect to $T$ and
        $X$ is sequentially compact with respect to $T$.

    If $X$ is limit point compact with respect to $T$ and $X$ is sequentially compact with respect to $T$ then $X$ is compact with respect to $T$.
\end{proof}
    
\section{§29 Local Compactness}

% from §29 helper definition 1: compact subspace
\begin{definition}\label{compact_subspace}
    $C$ is a compact subspace of $X$ with respect to $T$ iff
        $C\subseteq X$ and
        $C$ is compact with respect to $\subspaceTopology{T}{C}$.
\end{definition}

% from §29 Item 1: locally compact at a point
\begin{definition}\label{locally_compact_at_point}
    $X$ is locally compact at $x$ with respect to $T$ iff
        $x\in X$ and
        there exists $C$ such that
            $C$ is a compact subspace of $X$ with respect to $T$ and
            there exists $U$ such that
                $U$ is a neighborhood of $x$ in $X$ with respect to $T$ and
                $U\subseteq C$.
\end{definition}

% from §29 Item 2: locally compact
\begin{definition}\label{locally_compact}
    $X$ is locally compact with respect to $T$ iff
        for all $x\in X$ we have $X$ is locally compact at $x$ with respect to $T$.
\end{definition}

% from §29 Item 3: compact implies locally compact
\begin{lemma}\label{compact_implies_locally_compact}
    Assume $T$ is a topology on $X$.
    Suppose $X$ is compact with respect to $T$.
    Then $X$ is locally compact with respect to $T$.
\end{lemma}
\begin{proof}
    Let $T_X = \subspaceTopology{T}{X}$.
    For all $U$ such that $U\in T$ we have $U\in T_X$
        by \cref{topology_on,powerset_elim,subseteq,inter_absorb_supseteq_left,subspace_topology}.
    Hence $T\subseteq T_X$ by \cref{subseteq}.
    For all $U$ such that $U\in T_X$ we have $U\in T$
        by \cref{subspace_topology,topology_on,powerset_elim,subseteq,inter_absorb_supseteq_left}.
    Therefore $T_X = T$ by \cref{subseteq,subseteq_antisymmetric}.
    Hence $X$ is a compact subspace of $X$ with respect to $T$ by \cref{subseteq_refl,compact_subspace}.
    For all $x\in X$ we have $X$ is locally compact at $x$ with respect to $T$
        by \cref{subseteq_refl,topology_on,open_in,neighborhood,locally_compact_at_point}.
    Hence $X$ is locally compact with respect to $T$ by \cref{locally_compact}.
\end{proof}

% from §29 helper lemma 1: set not member of itself
\begin{lemma}\label{set_notin_self}
    For all $x$ we have $x\notin x$.
\end{lemma}
\begin{proof}
    Fix $x$.
    Suppose not.
    Take $a\in\{x\}$ such that for all $y\in a$ we have $y\notin\{x\}$ by \cref{singleton_intro,emptyset,foundation}.
    Hence $x\notin\{x\}$ by \cref{singleton_elim}.
    Contradiction by \cref{singleton_intro}.
\end{proof}

% from §29 helper lemma 2: empty set is compact
\begin{lemma}\label{emptyset_compact}
    Assume $T$ is a topology on $X$.
    Then $\emptyset$ is compact with respect to $T$.
\end{lemma}
\begin{proof}
    Trivial.
\end{proof}

% from §29 helper lemma 3: intersection of compact subspaces is compact
\begin{lemma}\label{compact_subspace_inters}
    Assume $T$ is a topology on $X$.
    Assume $X$ is Hausdorff with respect to $T$.
    Suppose $C\subseteq\pow{X}$.
    Suppose $C$ is inhabited.
    Suppose for all $K$ such that $K\in C$ we have $K$ is a compact subspace of $X$ with respect to $T$.
    Then $\inters{C}$ is a compact subspace of $X$ with respect to $T$.
\end{lemma}
\begin{proof}
    Take $K_0$ such that $K_0\in C$.
    For all $K$ such that $K\in C$ we have $K$ is closed in $X$ with respect to $T$
        by \cref{compact_subspace,compact_subspace_closed}.
    Hence $\inters{C}$ is closed in $X$ with respect to $T$ by \cref{closed_inters}.
    $\inters{C}\subseteq K_0$ by \cref{inters,subseteq}.
    Let $T_{K0} = \subspaceTopology{T}{K_0}$.
    $K_0$ is compact with respect to $T_{K0}$ and $K_0\subseteq X$ by \cref{compact_subspace}.
    $\inters{C}\inter K_0 = \inters{C}$ by \cref{closed_in_subspace_iff,inter_absorb_supseteq_right,subseteq}.
    Hence $\inters{C}$ is compact with respect to $\subspaceTopology{T_{K0}}{\inters{C}}$
        by \cref{compact_subspace,subspace_of,subspace_topology_is_topology,closed_in_subspace_iff,closed_subspace_compact}.
    Hence $\inters{C}$ is compact with respect to $\subspaceTopology{T}{\inters{C}}$
        by \cref{closed_in_subspace_iff,closed_subspace_compact,subspace_subspace_topology,subseteq}.
    Finally $\inters{C}$ is a compact subspace of $X$ with respect to $T$ by \cref{subseteq_transitive,compact_subspace}.
\end{proof}

% from §29 helper lemma 4: union of compact subspaces is compact
\begin{lemma}\label{compact_subspace_union}
    Assume $T$ is a topology on $X$.
    Suppose $A$ is a compact subspace of $X$ with respect to $T$.
    Suppose $B$ is a compact subspace of $X$ with respect to $T$.
    Let $C = A\union B$.
    Then $C$ is a compact subspace of $X$ with respect to $T$.
\end{lemma}
\begin{proof}
    Hence $C\subseteq X$ by \cref{compact_subspace,union_subsets_is_subset}.
    Let $T_C = \subspaceTopology{T}{C}$.
    Show that for all $U$ such that $U$ is an open covering of $C$ with respect to $T_C$
        there exists $D$ such that
            $D\subseteq U$ and
            $D$ is finite and
            $D$ is a covering of $C$.
    \begin{subproof}
        Fix $U$ such that $U$ is an open covering of $C$ with respect to $T_C$.
        Let $V = \{W\in T \mid \text{there exists $U_0\in U$ such that $U_0 = C\inter W$}\}$.
        For all $a$ such that $a\in A$ we have $a\in\unions{V}$
            by \cref{union_intro_left,covering,elem_subseteq,unions_iff,open_covering,open_in,subspace_topology,inter_lower_right}.
        Hence $V$ is a covering of $A$ by \cref{subseteq,covering}.
        For all $V_0$ such that $V_0\in V$ we have $V_0$ is open in $X$ with respect to $T$ by \cref{open_in}.
        Let $T_A = \subspaceTopology{T}{A}$.
        $A$ is a subspace of $X$ with respect to $T$ and $A$ is compact with respect to $T_A$
            by \cref{subspace_of,compact_subspace}.
        Take $V_A$ such that
            $V_A\subseteq V$ and
            $V_A$ is finite and
            $V_A$ is a covering of $A$
            by \cref{compact_subspace_open_covering}.

        For all $b$ such that $b\in B$ we have $b\in\unions{V}$
            by \cref{union_intro_right,covering,elem_subseteq,unions_iff,open_covering,open_in,subspace_topology,inter_lower_right}.
        Thus $V$ is a covering of $B$ by \cref{subseteq,covering}.
        Let $T_B = \subspaceTopology{T}{B}$.
        $B$ is a subspace of $X$ with respect to $T$ and $B$ is compact with respect to $T_B$
            by \cref{subspace_of,compact_subspace}.
        Take $V_B$ such that
            $V_B\subseteq V$ and
            $V_B$ is finite and
            $V_B$ is a covering of $B$
            by \cref{compact_subspace_open_covering}.

        Let $V_0 = V_A\union V_B$.
        $V_0$ is finite by \cref{union_of_finite_is_finite}.
        For all $x$ such that $x\in C$ we have $x\in\unions{V_0}$
            by \cref{union_iff,covering,elem_subseteq,unions_iff,union_intro_left,union_intro_right}.
        Thus $V_0$ is a covering of $C$ by \cref{subseteq,covering}.

        Let $g = \{(v, C\inter v) \mid v\in V_0\}$.
        Show that $g$ is a function from $V_0$ to $\pow{C}$.
        \begin{subproof}
            For all $w$ such that $w\in g$ there exist $x,y$ such that $w = (x,y)$.
            For all $x,y,z$ such that $(x,y)\in g$ and $(x,z)\in g$ we have $y = z$
                by \cref{pair_eq_iff}.
            Then $g$ is a function by \cref{rightunique,relation}.
            Hence $\dom{g} = V_0$ by \cref{dom_iff,pair_eq_iff,subseteq,subseteq_antisymmetric}.
            For all $v$ such that $v\in\dom{g}$ we have $g(v)\in\pow{C}$
                by \cref{function_apply_iff,inter_lower_left,powerset_intro}.
            Therefore $g\in\funs{V_0}{\pow{C}}$ by \cref{function_apply_iff,inter_lower_left,powerset_intro,funs_intro}.
        \end{subproof}
        Let $D = \{w \mid \exists v\in V_0. w = C\inter v\}$.
        Therefore $D$ is finite by \cref{funs_is_function,funs,finite_image_function,img_iff,subseteq,subseteq_finite_is_finite}.
        $D\subseteq U$ by \cref{union_subsets_is_subset,elem_subseteq,subseteq}.
        Hence $D$ is a covering of $C$ by \cref{covering,elem_subseteq,unions_iff,inter_intro,subseteq}.
        Therefore there exists $D$ such that
            $D\subseteq U$ and
            $D$ is finite and
            $D$ is a covering of $C$.
    \end{subproof}
    Finally $C$ is a compact subspace of $X$ with respect to $T$ by \cref{compact,compact_subspace}.
\end{proof}

% from §29 helper lemma 5: compact subspace minus open set is compact
\begin{lemma}\label{compact_subspace_minus_open}
    Assume $T$ is a topology on $X$.
    Suppose $C$ is a compact subspace of $X$ with respect to $T$.
    Suppose $U$ is open in $X$ with respect to $T$.
    Let $D = C\setminus U$.
    Then $D$ is a compact subspace of $X$ with respect to $T$.
\end{lemma}
\begin{proof}
    Hence $D\subseteq X$ by \cref{compact_subspace,setminus_subseteq,subseteq_transitive}.
    Let $T_C = \subspaceTopology{T}{C}$.
    Let $T_D = \subspaceTopology{T}{D}$.
    Therefore $X\setminus U$ is closed in $X$ with respect to $T$
        by \cref{open_in,topology_on,subseteq,pow_iff,setminus_subseteq,setminus_setminus,inter_absorb_supseteq_left,closed_in}.
    Hence $D$ is closed in $C$ with respect to $T_C$
        by \cref{open_in,topology_on,subseteq,pow_iff,setminus_eq_inter_complement,inter_comm,closed_in_subspace_iff,subspace_of,compact_subspace}.
    Hence $D$ is compact with respect to $T_D$
        by \cref{compact_subspace,setminus_subseteq,subspace_topology_is_topology,closed_subspace_compact,subspace_subspace_topology}.
    Therefore $D$ is a compact subspace of $X$ with respect to $T$ by \cref{compact_subspace}.
\end{proof}

% from §29 helper definition 2: open sets from X
\begin{definition}\label{one_point_open_sets}
    $\onePointOpenSets{T}{X}{p}
        = \{U\in\pow{X\union\{p\}} \mid \text{$U\subseteq X$ and $U\in T$}\}$.
\end{definition}

% from §29 helper definition 3: complements of compact sets
\begin{definition}\label{one_point_compact_complements}
    $\onePointCompactComplements{T}{X}{p}
        = \{U\in\pow{X\union\{p\}} \mid \text{there exists $C$ such that
            $C$ is a compact subspace of $X$ with respect to $T$ and
            $U = (X\union\{p\})\setminus C$}\}$.
\end{definition}

% from §29 helper definition 4: one-point compactification topology
\begin{definition}\label{one_point_compactification_topology}
    $\onePointTopology{T}{X}{p}
        = \onePointOpenSets{T}{X}{p}\union \onePointCompactComplements{T}{X}{p}$.
\end{definition}

% from §29 helper lemma 6: one-point topology is contained in powerset
\begin{lemma}\label{one_point_topology_subseteq_pow}
    Assume $p\notin X$.
    Let $Y = X\union\{p\}$.
    Let $S = \onePointTopology{T}{X}{p}$.
    Then $S\subseteq\pow{Y}$.
\end{lemma}
\begin{proof}
    Hence $S\subseteq\pow{Y}$
        by \cref{one_point_compactification_topology,union_iff,one_point_open_sets,one_point_compact_complements,subseteq}.
\end{proof}

% from §29 helper lemma 7: open set from X is open in one-point topology
\begin{lemma}\label{one_point_topology_intro_left}
    Assume $p\notin X$.
    Let $Y = X\union\{p\}$.
    Let $S = \onePointTopology{T}{X}{p}$.
    Suppose $U\subseteq X$.
    Suppose $U\in T$.
    Then $U\in S$.
\end{lemma}
\begin{proof}
    Then $U\in S$ by \cref{one_point_compactification_topology,one_point_open_sets,union_upper_left,subseteq_transitive,pow_iff,union_intro_left}.
\end{proof}

% from §29 helper lemma 8: complement of compact set is open in one-point topology
\begin{lemma}\label{one_point_topology_intro_right}
    Assume $p\notin X$.
    Let $Y = X\union\{p\}$.
    Let $S = \onePointTopology{T}{X}{p}$.
    Suppose there exists $C$ such that
        $C$ is a compact subspace of $X$ with respect to $T$ and
        $U = Y\setminus C$.
    Then $U\in S$.
\end{lemma}
\begin{proof}
    Take $C$ such that
        $C$ is a compact subspace of $X$ with respect to $T$ and
        $U = Y\setminus C$.
    Then $U\in S$ by \cref{one_point_compactification_topology,one_point_compact_complements,setminus_subseteq,pow_iff,union_intro_right}.
\end{proof}

% from §29 helper lemma 9: open set without added point
\begin{lemma}\label{one_point_topology_without_p}
    Assume $p\notin X$.
    Let $Y = X\union\{p\}$.
    Let $S = \onePointTopology{T}{X}{p}$.
    Suppose $U\in S$.
    Suppose $p\notin U$.
    Then $U\subseteq X$ and $U\in T$.
\end{lemma}
\begin{proof}
    Show that $U\in \onePointOpenSets{T}{X}{p}$.
    \begin{subproof}
        Suppose not.
        Then $U\in \onePointCompactComplements{T}{X}{p}$ by \cref{one_point_compactification_topology,union_iff}.
        Take $C$ such that
            $C$ is a compact subspace of $X$ with respect to $T$ and
            $U = Y\setminus C$
            by \cref{one_point_compact_complements}.
        Therefore $p\in U$ by \cref{compact_subspace,subseteq,singleton_intro,union_intro_right,setminus_intro}.
        Contradiction.
    \end{subproof}
    Hence $U\subseteq X$ and $U\in T$ by \cref{one_point_open_sets}.
\end{proof}

% from §29 helper lemma 10: open set containing added point
\begin{lemma}\label{one_point_topology_with_p}
    Assume $p\notin X$.
    Let $Y = X\union\{p\}$.
    Let $S = \onePointTopology{T}{X}{p}$.
    Suppose $U\in S$.
    Suppose $p\in U$.
    Then there exists $C$ such that $C$ is a compact subspace of $X$ with respect to $T$
        and $U = (X\union\{p\})\setminus C$.
\end{lemma}
\begin{proof}
    Show that $U\in \onePointCompactComplements{T}{X}{p}$.
    \begin{subproof}
        Suppose not.
        Then $U\in \onePointOpenSets{T}{X}{p}$ by \cref{one_point_compactification_topology,union_iff}.
        Therefore $p\in X$ by \cref{one_point_open_sets,subseteq}.
        Contradiction.
    \end{subproof}
    Hence there exists $C$ such that $C$ is a compact subspace of $X$ with respect to $T$
        and $U = Y\setminus C$
        by \cref{one_point_compact_complements}.
\end{proof}

% from §29 helper lemma 11: one-point topology is a topology
\begin{lemma}\label{one_point_topology_is_topology}
    Assume $T$ is a topology on $X$.
    Assume $X$ is Hausdorff with respect to $T$.
    Assume $p\notin X$.
    Let $Y = X\union\{p\}$.
    Let $S = \onePointTopology{T}{X}{p}$.
    Then $S$ is a topology on $Y$.
\end{lemma}
\begin{proof}
    $\emptyset\in S$ by \cref{emptyset_subseteq,topology_on,one_point_topology_intro_left}.
    Let $T_0 = \subspaceTopology{T}{\emptyset}$.
    Hence $\emptyset$ is a compact subspace of $X$ with respect to $T$
        by \cref{emptyset_subseteq,subspace_topology_is_topology,emptyset_compact,compact_subspace}.
    Therefore $Y\in S$ by \cref{setminus_emptyset,one_point_topology_intro_right}.
    Show that for all $A$ such that $A\subseteq S$ we have $\unions{A}\in S$.
    \begin{subproof}
        Fix $A$ such that $A\subseteq S$.
        \begin{byCase}
        \caseOf{$p\notin\unions{A}$.}
            Therefore $A\subseteq T$ by \cref{unions_iff,subseteq,one_point_topology_without_p}.
            $\unions{A}\subseteq X$
                by \cref{topology_on,subseteq_transitive,unions_subseteq_of_powerset_is_subseteq}.
            Hence $\unions{A}\in S$ by \cref{topology_on,one_point_topology_intro_left}.
        \caseOf{$p\in\unions{A}$.}
            Take $U_0\in A$ such that $p\in U_0$ by \cref{unions_iff}.
            Take $C_0$ such that
                $C_0$ is a compact subspace of $X$ with respect to $T$ and
                $U_0 = (X\union\{p\})\setminus C_0$
                by \cref{subseteq,one_point_topology_with_p}.
            Let $A_1 = \{U\in A \mid p\notin U\}$.
            Hence $A_1\subseteq T$ by \cref{subseteq,one_point_topology_without_p}.
            Let $U = \unions{A_1}$.
            Then $U\subseteq X$ and $U\in T$ by \cref{topology_on,subseteq,pow_iff}.
            Let $A_2 = \{U\in A \mid p\in U\}$.
            Let $F = \{Y\setminus U \mid U\in A_2\}$.
            $U_0\in A_2$.
            Hence $Y\setminus U_0\in F$.
            Therefore $F$ is inhabited.
            For all $C$ such that $C\in F$ we have $C$ is a compact subspace of $X$ with respect to $T$
                by \cref{subseteq,one_point_topology_with_p,compact_subspace,union_upper_left,subseteq_transitive,setminus_setminus,inter_absorb_supseteq_left}.
            Hence $F\subseteq\pow{X}$ by \cref{subseteq,compact_subspace,pow_iff}.
            Let $C = \inters{F}$.
            $C$ is a compact subspace of $X$ with respect to $T$ by \cref{compact_subspace_inters}.
            Show that for all $x$ such that $x\in\unions{A_2}$ we have
                $x\in (X\union\{p\})\setminus C$.
            \begin{subproof}
                Fix $x$ such that $x\in\unions{A_2}$.
                Take $U_1\in A_2$ such that $x\in U_1$ by \cref{unions_iff}.
                Let $C_1 = Y\setminus U_1$.
                Then $C_1\in F$.
                Hence $U_1\subseteq Y$ by \cref{one_point_topology_subseteq_pow,subseteq,pow_iff}.
                Then $Y\setminus C_1 = Y\inter U_1$ by \cref{setminus_setminus}.
                Hence $Y\inter U_1 = U_1$ by \cref{inter_absorb_supseteq_left}.
                Therefore $x\notin C_1$ by \cref{setminus_elim_right}.
                Hence $x\notin C$ by \cref{inters_iff_forall}.
                Therefore $x\in (X\union\{p\})\setminus C$ by \cref{setminus_elim_left,setminus_intro}.
            \end{subproof}
            Hence $\unions{A_2}\subseteq (X\union\{p\})\setminus C$ by \cref{subseteq}.
            Show that for all $x$ such that $x\in (X\union\{p\})\setminus C$ we have
                $x\in\unions{A_2}$.
            \begin{subproof}
                Fix $x$ such that $x\in (X\union\{p\})\setminus C$.
                Thus $x\notin C$ by \cref{setminus_elim_right}.
                Suppose not.
                Show that for all $C_1$ such that $C_1\in F$ we have $x\in C_1$.
                \begin{subproof}
                    Fix $C_1$ such that $C_1\in F$.
                    Take $U_1\in A_2$ such that $C_1 = Y\setminus U_1$.
                    Suppose not.
                    Therefore $x\in U_1$
                        by \cref{subseteq,one_point_topology_subseteq_pow,pow_iff,setminus_elim_left,setminus_intro,setminus_setminus,inter_absorb_supseteq_left}.
                    Hence $x\in\unions{A_2}$ by \cref{unions_iff}.
                    Contradiction.
                \end{subproof}
                Therefore $x\in C$ by \cref{inters_iff_forall}.
                Contradiction.
                Hence $x\in\unions{A_2}$.
            \end{subproof}
            Therefore $(X\union\{p\})\setminus C\subseteq \unions{A_2}$ by \cref{subseteq}.
            Hence $\unions{A_2} = (X\union\{p\})\setminus C$ by \cref{subseteq_antisymmetric}.
            For all $x$ such that $x\in\unions{A}$ we have $x\in\unions{A_1}\union\unions{A_2}$.
            Therefore $\unions{A}\subseteq \unions{A_1}\union\unions{A_2}$ by \cref{subseteq}.
            Therefore $\unions{A_1}\union\unions{A_2}\subseteq \unions{A}$
                by \cref{union_iff,unions_iff,subseteq}.
            Hence $\unions{A} = \unions{A_1}\union\unions{A_2}$ by \cref{subseteq_antisymmetric}.
            Let $D = C\setminus U$.
            $D$ is a compact subspace of $X$ with respect to $T$
                by \cref{open_in,compact_subspace_minus_open}.
            Show that for all $x$ such that $x\in\unions{A}$ we have
                $x\in (X\union\{p\})\setminus D$.
            \begin{subproof}
                Fix $x$ such that $x\in\unions{A}$.
                Then $x\in\unions{A_1}\union\unions{A_2}$.
                Therefore $x\in\unions{A_1}$ or $x\in\unions{A_2}$ by \cref{union_iff}.
                \begin{byCase}
                \caseOf{$x\in\unions{A_1}$.}
                    Then $x\in U$.
                    Show that $x\notin D$.
                    \begin{subproof}
                        Suppose not.
                        Thus $x\in C$ and $x\notin U$ by \cref{setminus_elim_left,setminus_elim_right}.
                        Contradiction.
                    \end{subproof}
                    Hence $x\in (X\union\{p\})\setminus D$
                        by \cref{subseteq,union_intro_left,setminus_intro}.
                \caseOf{$x\in\unions{A_2}$.}
                    Then $x\in (X\union\{p\})\setminus C$.
                    Hence $x\notin C$ by \cref{setminus_elim_right}.
                    Therefore $x\in (X\union\{p\})\setminus D$
                        by \cref{setminus_elim_left,setminus_subseteq,subseteq,setminus_intro}.
                \end{byCase}
            \end{subproof}
            Hence $\unions{A}\subseteq (X\union\{p\})\setminus D$ by \cref{subseteq}.
            Show that for all $x$ such that $x\in (X\union\{p\})\setminus D$ we have
                $x\in\unions{A}$.
            \begin{subproof}
                Fix $x$ such that $x\in (X\union\{p\})\setminus D$.
                Thus $x\notin D$ by \cref{setminus_elim_right}.
                \begin{byCase}
                \caseOf{$x\in U$.}
                    Take $U_1\in A_1$ such that $x\in U_1$ by \cref{unions_iff}.
                    Hence $U_1\in A$.
                    Therefore $x\in\unions{A}$ by \cref{unions_iff}.
                \caseOf{$x\notin U$.}
                    Show that $x\in (X\union\{p\})\setminus C$.
                    \begin{subproof}
                        Suppose not.
                        Then $x\in X\union\{p\}$ by \cref{setminus_elim_left}.
                        Show that $x\in C$.
                        \begin{subproof}
                            Suppose not.
                            Hence $x\in (X\union\{p\})\setminus C$ by \cref{setminus_intro}.
                            Contradiction.
                        \end{subproof}
                        Therefore $x\in C$ and $x\notin U$.
                        Hence $x\in D$ by \cref{setminus_intro}.
                        Contradiction.
                    \end{subproof}
                    Hence $x\in\unions{A_2}$.
                    Take $U_1\in A_2$ such that $x\in U_1$ by \cref{unions_iff}.
                    Then $U_1\in A$.
                    Hence $x\in\unions{A}$ by \cref{unions_iff}.
                \end{byCase}
            \end{subproof}
            Therefore $(X\union\{p\})\setminus D\subseteq \unions{A}$ by \cref{subseteq}.
            Hence $\unions{A} = (X\union\{p\})\setminus D$ by \cref{subseteq_antisymmetric}.
            Therefore there exists $C$ such that $C$ is a compact subspace of $X$ with respect to $T$
                and $\unions{A} = Y\setminus C$.
            Hence $\unions{A}\in S$ by \cref{one_point_topology_intro_right}.
        \end{byCase}
    \end{subproof}
    Show that for all $U,V$ such that $U\in S$ and $V\in S$ we have $U\inter V\in S$.
    \begin{subproof}
        Fix $U$.
        Fix $V$ such that $U\in S$ and $V\in S$.
        \begin{byCase}
        \caseOf{$p\notin U$ and $p\notin V$.}
            Hence $U\subseteq X$ and $U\in T$ and $V\subseteq X$ and $V\in T$
                by \cref{one_point_topology_without_p}.
            Therefore $U\inter V\in S$
                by \cref{inter_lower_left,subseteq_transitive,topology_on,one_point_topology_intro_left}.
        \caseOf{$p\in U$ and $p\in V$.}
            Take $C_1$ such that
                $C_1$ is a compact subspace of $X$ with respect to $T$ and
                $U = (X\union\{p\})\setminus C_1$
                by \cref{one_point_topology_with_p}.
            Take $C_2$ such that
                $C_2$ is a compact subspace of $X$ with respect to $T$ and
                $V = (X\union\{p\})\setminus C_2$
                by \cref{one_point_topology_with_p}.
            Let $C_3 = C_1\union C_2$.
            Then $U\inter V\in S$ by \cref{compact_subspace_union,setminus_union,one_point_topology_intro_right}.
        \caseOf{$p\in U$ and $p\notin V$.}
            Take $C_1$ such that
                $C_1$ is a compact subspace of $X$ with respect to $T$ and
                $U = (X\union\{p\})\setminus C_1$
                by \cref{one_point_topology_with_p}.
            Hence $V\subseteq X$ and $V\in T$ by \cref{one_point_topology_without_p}.
            $C_1\subseteq X$ by \cref{compact_subspace}.
            $C_1$ is closed in $X$ with respect to $T$ by \cref{compact_subspace,compact_subspace_closed}.
            Therefore $X\setminus C_1\in T$ by \cref{closed_in,setminus_subseteq,open_in}.
            Hence $V\setminus C_1\in T$ by \cref{setminus_eq_inter_complement,topology_on}.
            $V\subseteq X\union\{p\}$ and $C_1\subseteq X\union\{p\}$
                by \cref{subseteq_transitive,union_upper_left}.
            Hence $V\setminus C_1 = V\inter U$ by \cref{setminus_eq_inter_complement}.
            Therefore $U\inter V = V\setminus C_1$ by \cref{inter_comm}.
            Then $U\inter V\subseteq X$ by \cref{inter_lower_left,subseteq_transitive}.
            Hence $U\inter V\in T$.
            Therefore $U\inter V\in S$ by \cref{one_point_topology_intro_left}.
        \caseOf{$p\notin U$ and $p\in V$.}
            Take $C_1$ such that
                $C_1$ is a compact subspace of $X$ with respect to $T$ and
                $V = (X\union\{p\})\setminus C_1$
                by \cref{one_point_topology_with_p}.
            Hence $U\subseteq X$ and $U\in T$ by \cref{one_point_topology_without_p}.
            $C_1\subseteq X$ by \cref{compact_subspace}.
            $C_1$ is closed in $X$ with respect to $T$ by \cref{compact_subspace,compact_subspace_closed}.
            Therefore $X\setminus C_1\in T$ by \cref{closed_in,setminus_subseteq,open_in}.
            Hence $U\setminus C_1\in T$ by \cref{setminus_eq_inter_complement,topology_on}.
            $U\subseteq X\union\{p\}$ and $C_1\subseteq X\union\{p\}$
                by \cref{subseteq_transitive,union_upper_left}.
            Hence $U\setminus C_1 = U\inter V$ by \cref{setminus_eq_inter_complement}.
            Therefore $U\inter V = U\setminus C_1$.
            Then $U\inter V\subseteq X$ by \cref{inter_lower_left,subseteq_transitive}.
            Hence $U\inter V\in T$.
            Therefore $U\inter V\in S$ by \cref{one_point_topology_intro_left}.
        \end{byCase}
    \end{subproof}
    Finally $S$ is a topology on $Y$ by \cref{one_point_topology_subseteq_pow,topology_on}.
\end{proof}

% from §29 helper lemma 12: subspace topology from one-point topology
\begin{lemma}\label{one_point_topology_subspace}
    Assume $T$ is a topology on $X$.
    Assume $X$ is Hausdorff with respect to $T$.
    Assume $p\notin X$.
    Let $Y = X\union\{p\}$.
    Let $S = \onePointTopology{T}{X}{p}$.
    Let $T_X = \subspaceTopology{S}{X}$.
    Then $T_X = T$.
\end{lemma}
\begin{proof}
    Hence $T\subseteq T_X$
        by \cref{topology_on,pow_iff,subseteq,one_point_topology_intro_left,inter_absorb_supseteq_left,subspace_topology}.
    Show that for all $U$ such that $U\in T_X$ we have $U\in T$.
    \begin{subproof}
        Fix $U$ such that $U\in T_X$.
        Take $V\in S$ such that $U = X\inter V$ by \cref{subspace_topology}.
        \begin{byCase}
        \caseOf{$p\notin V$.}
            Then $U\in T$ by \cref{one_point_topology_without_p,inter_absorb_supseteq_left}.
        \caseOf{$p\in V$.}
            Take $C$ such that
                $C$ is a compact subspace of $X$ with respect to $T$ and
                $V = (X\union\{p\})\setminus C$
                by \cref{one_point_topology_with_p}.
            $C\subseteq X$ by \cref{compact_subspace}.
            $C$ is closed in $X$ with respect to $T$ by \cref{compact_subspace,compact_subspace_closed}.
            Therefore $X\setminus C\in T$ by \cref{closed_in,setminus_subseteq,open_in}.
            $X\subseteq Y$ by \cref{union_upper_left}.
            Hence $C\subseteq Y$ by \cref{subseteq_transitive}.
            Then $U = X\inter (Y\setminus C)$.
            Therefore $U = X\setminus C$ by \cref{setminus_eq_inter_complement}.
            Hence $U\in T$.
        \end{byCase}
    \end{subproof}
    Therefore $T_X\subseteq T$ by \cref{subseteq}.
    Hence $T_X = T$ by \cref{subseteq_antisymmetric}.
\end{proof}

% from §29 helper lemma 13: one-point topology is compact
\begin{lemma}\label{one_point_topology_compact}
    Assume $T$ is a topology on $X$.
    Assume $X$ is Hausdorff with respect to $T$.
    Assume $p\notin X$.
    Let $Y = X\union\{p\}$.
    Let $S = \onePointTopology{T}{X}{p}$.
    Then $Y$ is compact with respect to $S$.
\end{lemma}
\begin{proof}
    Show that for all $A$ such that $A$ is an open covering of $Y$ with respect to $S$
        there exists $B$ such that $B\subseteq A$ and $B$ is finite and $B$ is a covering of $Y$.
    \begin{subproof}
        Fix $A$ such that $A$ is an open covering of $Y$ with respect to $S$.
        Thus $A$ is a covering of $Y$ and
            for all $U$ such that $U\in A$ we have $U$ is open in $Y$ with respect to $S$
            by \cref{open_covering}.
        Hence $p\in\unions{A}$ by \cref{singleton_intro,union_intro_right,covering,subseteq}.
        Take $U_0\in A$ such that $p\in U_0$ by \cref{unions_iff}.
        Take $C_0$ such that
            $C_0$ is a compact subspace of $X$ with respect to $T$ and
            $U_0 = (X\union\{p\})\setminus C_0$
            by \cref{open_covering,open_in,one_point_topology_with_p}.
        Let $A_0 = A\setminus\{U_0\}$.
        Therefore $A_0$ is a covering of $C_0$
            by \cref{setminus_elim_right,compact_subspace,subseteq,union_upper_left,covering,unions_iff,singleton_iff,setminus_intro}.
        $T = \subspaceTopology{S}{X}$ by \cref{one_point_topology_subspace}.
        Let $T_C = \subspaceTopology{T}{C_0}$.
        Let $T_D = \subspaceTopology{S}{C_0}$.
        $S$ is a topology on $Y$ by \cref{one_point_topology_is_topology}.
        Hence $T_D = T_C$ by \cref{compact_subspace,union_upper_left,subspace_subspace_topology}.
        $C_0$ is compact with respect to $T_C$ by \cref{compact_subspace}.
        Therefore $C_0$ is compact with respect to $T_D$.
        Hence $C_0$ is a subspace of $Y$ with respect to $S$
            by \cref{compact_subspace,union_upper_left,subseteq_transitive,subspace_of}.
        For all $U$ such that $U\in A_0$ we have $U$ is open in $Y$ with respect to $S$
            by \cref{setminus_elim_left,open_covering}.
        Take $B$ such that
            $B\subseteq A_0$ and
            $B$ is finite and
            $B$ is a covering of $C_0$
            by \cref{compact_subspace_open_covering}.
        Let $B_0 = B\union\{U_0\}$.
        $B_0$ is finite by \cref{pair_is_finite,upair_intro_left,singleton_subset_intro,subseteq_finite_is_finite,union_of_finite_is_finite}.
        Therefore $B_0\subseteq A$ by \cref{setminus_subseteq,subseteq_transitive,singleton_subset_intro,union_subsets_is_subset}.
        Hence $B_0$ is a covering of $Y$
            by \cref{singleton_intro,union_intro_right,unions_iff,setminus_intro,compact_subspace,subseteq_transitive,setminus_setminus,inter_absorb_supseteq_left,subseteq,covering,union_intro_left}.
        Therefore there exists $B_0$ such that
            $B_0\subseteq A$ and
            $B_0$ is finite and
            $B_0$ is a covering of $Y$.
    \end{subproof}
    Finally $Y$ is compact with respect to $S$ by \cref{compact}.
\end{proof}

% from §29 helper lemma 14: one-point topology is Hausdorff
\begin{lemma}\label{one_point_topology_hausdorff}
    Assume $T$ is a topology on $X$.
    Assume $X$ is Hausdorff with respect to $T$.
    Assume $X$ is locally compact with respect to $T$.
    Assume $p\notin X$.
    Let $Y = X\union\{p\}$.
    Let $S = \onePointTopology{T}{X}{p}$.
    Then $Y$ is Hausdorff with respect to $S$.
\end{lemma}
\begin{proof}
    $S$ is a topology on $Y$ by \cref{one_point_topology_is_topology}.
    Show that for all $x,y$ such that $x\in Y$ and $y\in Y$ and $x\neq y$ there exist $U,V$ such that
        $U$ is a neighborhood of $x$ in $Y$ with respect to $S$ and
        $V$ is a neighborhood of $y$ in $Y$ with respect to $S$ and
        $U$ is disjoint from $V$.
    \begin{subproof}
        Fix $x$.
        Fix $y$ such that $x\in Y$ and $y\in Y$ and $x\neq y$.
        \begin{byCase}
        \caseOf{$x\in X$ and $y\in X$.}
            Take $U,V$ such that
                $U$ is a neighborhood of $x$ in $X$ with respect to $T$ and
                $V$ is a neighborhood of $y$ in $X$ with respect to $T$ and
                $U$ is disjoint from $V$
                by \cref{hausdorff}.
            Then $U$ is a neighborhood of $x$ in $Y$ with respect to $S$ and
                $V$ is a neighborhood of $y$ in $Y$ with respect to $S$
                by \cref{neighborhood,open_in,topology_on,subseteq,pow_iff,one_point_topology_intro_left}.
            Therefore there exist $U,V$ such that
                $U$ is a neighborhood of $x$ in $Y$ with respect to $S$ and
                $V$ is a neighborhood of $y$ in $Y$ with respect to $S$ and
                $U$ is disjoint from $V$.
        \caseOf{$x\in X$ and $y\notin X$.}
            Take $C$ such that
                $C$ is a compact subspace of $X$ with respect to $T$ and
                there exists $U$ such that
                    $U$ is a neighborhood of $x$ in $X$ with respect to $T$ and
                    $U\subseteq C$
                by \cref{locally_compact,locally_compact_at_point}.
            Take $U$ such that
                $U$ is a neighborhood of $x$ in $X$ with respect to $T$ and
                $U\subseteq C$.
            Let $V = (X\union\{p\})\setminus C$.
            Therefore $V\in S$ by \cref{one_point_topology_intro_right}.
            Hence $y\in V$ by \cref{compact_subspace,subseteq,union_iff,singleton_iff,setminus_intro}.
            Then $U$ is disjoint from $V$ by \cref{subseteq,setminus_elim_right,disjoint}.
            Then $U$ is a neighborhood of $x$ in $Y$ with respect to $S$ and
                $V$ is a neighborhood of $y$ in $Y$ with respect to $S$
                by \cref{neighborhood,open_in,topology_on,subseteq,pow_iff,one_point_topology_intro_left,compact_subspace,setminus_intro}.
            Therefore there exist $U,V$ such that
                $U$ is a neighborhood of $x$ in $Y$ with respect to $S$ and
                $V$ is a neighborhood of $y$ in $Y$ with respect to $S$ and
                $U$ is disjoint from $V$.
        \caseOf{$x\notin X$ and $y\in X$.}
            Take $C$ such that
                $C$ is a compact subspace of $X$ with respect to $T$ and
                there exists $U$ such that
                    $U$ is a neighborhood of $y$ in $X$ with respect to $T$ and
                    $U\subseteq C$
                by \cref{locally_compact,locally_compact_at_point}.
            Take $U$ such that
                $U$ is a neighborhood of $y$ in $X$ with respect to $T$ and
                $U\subseteq C$.
            Let $V = (X\union\{p\})\setminus C$.
            Therefore $V\in S$ by \cref{one_point_topology_intro_right}.
            Hence $x\in V$ by \cref{compact_subspace,subseteq,union_iff,singleton_iff,setminus_intro}.
            Then $U$ is disjoint from $V$ by \cref{subseteq,setminus_elim_right,disjoint}.
            Then $U$ is a neighborhood of $y$ in $Y$ with respect to $S$ and
                $V$ is a neighborhood of $x$ in $Y$ with respect to $S$
                by \cref{neighborhood,open_in,topology_on,subseteq,pow_iff,one_point_topology_intro_left,compact_subspace,setminus_intro}.
            Therefore $V$ is disjoint from $U$ by \cref{disjoint_symmetric}.
            Hence there exist $U,V$ such that
                $U$ is a neighborhood of $x$ in $Y$ with respect to $S$ and
                $V$ is a neighborhood of $y$ in $Y$ with respect to $S$ and
                $U$ is disjoint from $V$.
        \end{byCase}
    \end{subproof}
    Finally $Y$ is Hausdorff with respect to $S$ by \cref{hausdorff}.
\end{proof}
% from §29 Item 4: one-point compactification existence (forward)
\begin{theorem}\label{one_point_compactification_existence_forward}
    Assume $T$ is a topology on $X$.
    If $X$ is locally compact with respect to $T$ and $X$ is Hausdorff with respect to $T$, then
        there exist $Y,S,p$ such that
            $S$ is a topology on $Y$ and
            $X\subseteq Y$ and
            $T = \subspaceTopology{S}{X}$ and
            $Y\setminus X = \{p\}$ and
            $Y$ is compact with respect to $S$ and
            $Y$ is Hausdorff with respect to $S$.
\end{theorem}
\begin{proof}
    Assume $X$ is locally compact with respect to $T$ and $X$ is Hausdorff with respect to $T$.
    Let $p = X$.
    $p\notin X$ by \cref{set_notin_self}.
    Let $Y = X\union\{p\}$.
    Let $S = \onePointTopology{T}{X}{p}$.
    $S$ is a topology on $Y$ by \cref{one_point_topology_is_topology}.
    $X\subseteq Y$ by \cref{union_upper_left}.
    $T = \subspaceTopology{S}{X}$ by \cref{one_point_topology_subspace}.
    Hence $Y\setminus X\subseteq \{p\}$
        by \cref{setminus_elim_left,setminus_elim_right,union_iff,subseteq}.
    Therefore $\{p\}\subseteq Y\setminus X$
        by \cref{singleton_iff,union_intro_right,setminus_intro,subseteq}.
    Hence $Y\setminus X = \{p\}$ by \cref{subseteq_antisymmetric}.
    $Y$ is compact with respect to $S$ and $Y$ is Hausdorff with respect to $S$
        by \cref{one_point_topology_compact,one_point_topology_hausdorff}.
    Therefore there exist $Y,S,p$ such that
        $S$ is a topology on $Y$ and
        $X\subseteq Y$ and
        $T = \subspaceTopology{S}{X}$ and
        $Y\setminus X = \{p\}$ and
        $Y$ is compact with respect to $S$ and
        $Y$ is Hausdorff with respect to $S$.
\end{proof}

% from §29 Item 4: one-point compactification existence (reverse)
\begin{theorem}\label{one_point_compactification_existence_reverse}
    Assume $T$ is a topology on $X$.
    If there exist $Y,S,p$ such that
            $S$ is a topology on $Y$ and
            $X\subseteq Y$ and
            $T = \subspaceTopology{S}{X}$ and
            $Y\setminus X = \{p\}$ and
            $Y$ is compact with respect to $S$ and
            $Y$ is Hausdorff with respect to $S$
        then $X$ is locally compact with respect to $T$ and $X$ is Hausdorff with respect to $T$.
\end{theorem}
\begin{proof}
    Assume there exist $Y,S,p$ such that
        $S$ is a topology on $Y$ and
        $X\subseteq Y$ and
        $T = \subspaceTopology{S}{X}$ and
        $Y\setminus X = \{p\}$ and
        $Y$ is compact with respect to $S$ and
        $Y$ is Hausdorff with respect to $S$.
    Take $Y,S,p$ such that
        $S$ is a topology on $Y$ and
        $X\subseteq Y$ and
        $T = \subspaceTopology{S}{X}$ and
        $Y\setminus X = \{p\}$ and
        $Y$ is compact with respect to $S$ and
        $Y$ is Hausdorff with respect to $S$.
    Then $Y = X\union\{p\}$ by \cref{setminus_partition}.
    Hence $X$ is Hausdorff with respect to $T$ by \cref{subspace_of,subspace_hausdorff}.
    Show that for all $x\in X$ we have $X$ is locally compact at $x$ with respect to $T$.
    \begin{subproof}
            Fix $x$ such that $x\in X$.
            Then $p\in Y$ and $p\notin X$ by \cref{singleton_intro,setminus_elim_left,setminus_elim_right}.
            $x\in Y$ by \cref{subseteq}.
            Therefore $x\neq p$.
            Take $U,V$ such that
                $U$ is a neighborhood of $x$ in $Y$ with respect to $S$ and
                $V$ is a neighborhood of $p$ in $Y$ with respect to $S$ and
                $U$ is disjoint from $V$
                by \cref{hausdorff}.
            Let $C = Y\setminus V$.
            $C\subseteq Y$ and $V\subseteq Y$ and $V$ is open in $Y$ with respect to $S$
                by \cref{setminus_subseteq,neighborhood,open_in,topology_on,subseteq,pow_iff}.
            Hence $Y\setminus C = V$ by \cref{setminus_setminus,inter_absorb_supseteq_left}.
            Therefore $C$ is closed in $Y$ with respect to $S$ and
                $C$ is compact with respect to $\subspaceTopology{S}{C}$
                by \cref{neighborhood,closed_in,closed_subspace_compact}.
            Hence $p\notin C$ by \cref{neighborhood,setminus_elim_right}.
            Hence $C\subseteq X$ by \cref{setminus_elim_left,union_iff,singleton_iff,setminus_elim_right,subseteq}.
            Hence $C$ is a compact subspace of $X$ with respect to $T$
                by \cref{subspace_subspace_topology,compact_subspace}.
            Then $U\inter X$ is a neighborhood of $x$ in $X$ with respect to $T$
                by \cref{neighborhood,inter_comm,open_in,topology_on,subseteq,pow_iff,subspace_topology,inter_intro}.
            For all $y$ such that $y\in U\inter X$ we have $y\in C$.
            Therefore $U\inter X\subseteq C$ by \cref{subseteq}.
            Hence $X$ is locally compact at $x$ with respect to $T$ by \cref{locally_compact_at_point}.
    \end{subproof}
    Therefore $X$ is locally compact with respect to $T$ by \cref{locally_compact}.
    Hence $X$ is locally compact with respect to $T$ and $X$ is Hausdorff with respect to $T$.
\end{proof}

 
% from §29 Item 5: one-point compactification uniqueness
\begin{theorem}\label{one_point_compactification_uniqueness}
    Assume $T$ is a topology on $X$.
    Assume $S$ is a topology on $Y$.
    Assume $S_1$ is a topology on $Y_1$.
    Suppose $X$ is a subspace of $Y$ with respect to $S$.
    Suppose $X$ is a subspace of $Y_1$ with respect to $S_1$.
    Suppose $T = \subspaceTopology{S}{X}$.
    Suppose $T = \subspaceTopology{S_1}{X}$.
    Suppose $Y\setminus X = \{p\}$.
    Suppose $Y_1\setminus X = \{q\}$.
    Suppose $Y$ is compact with respect to $S$ and $Y$ is Hausdorff with respect to $S$.
    Suppose $Y_1$ is compact with respect to $S_1$ and $Y_1$ is Hausdorff with respect to $S_1$.
    Then there exists $h$ such that
        $h$ is a homeomorphism between $Y$ and $Y_1$ with respect to $S$ and $S_1$ and
        $\restrl{h}{X} = \identity{X}$.
\end{theorem}
\begin{proof}
    Hence $X\subseteq Y$ and $X\subseteq Y_1$ by \cref{subspace_of}.
    Therefore $Y = X\union\{p\}$ and $Y_1 = X\union\{q\}$ by \cref{setminus_partition}.
    Then $p\in Y$ and $q\in Y_1$ and $p\notin X$ and $q\notin X$
        by \cref{singleton_intro,setminus_elim_left,setminus_elim_right}.

    Let $f = \identity{X}$.
    Let $g = \{(p,q)\}$.
    Let $h = f\union g$.

    $f$ is a function by \cref{id_is_function}.
    Then $g$ is a function by \cref{singleton_iff,relation,pair_eq_iff,rightunique}.

    $\dom{f} = X$ by \cref{id_dom}.
    $g = \cons{(p,q)}{\emptyset}$.
    Hence $\dom{g} = \{p\}$ by \cref{dom_cons,dom_emptyset}.
    For all $x$ such that $x\in\dom{f}\inter\dom{g}$ we have $f(x) = g(x)$.
    Hence $h$ is a function by \cref{union_compatible_functions}.

    Therefore $\dom{h} = Y$ by \cref{dom_union}.

    For all $x$ such that $x\in\dom{h}$ we have $h(x)\in Y_1$
        by \cref{subseteq,union_iff,id_elem_intro,union_intro_left,function_apply_intro,singleton_iff,singleton_intro,union_intro_right}.
    Therefore $h\in\funs{Y}{Y_1}$ by \cref{funs_intro}.

    For all $y$ such that $y\in Y_1$ there exists $x$ such that $x\in Y$ and $h(x) = y$
        by \cref{union_iff,union_intro_left,id_elem_intro,union_intro_left,function_apply_intro,singleton_iff,union_intro_right,singleton_intro}.
    Hence $h\in\surj{Y}{Y_1}$ by \cref{surj}.

    Show that for all $x_1,x_2$ such that $x_1\in\dom{h}$ and $x_2\in\dom{h}$ and $h(x_1) = h(x_2)$
        we have $x_1 = x_2$.
    \begin{subproof}
        Fix $x_1$.
        Fix $x_2$ such that $x_1\in\dom{h}$ and $x_2\in\dom{h}$ and $h(x_1) = h(x_2)$.
        Then $x_1\in Y$ and $x_2\in Y$.
        Hence $x_1\in X$ or $x_1\in\{p\}$ by \cref{union_iff}.
        Thus $x_2\in X$ or $x_2\in\{p\}$ by \cref{union_iff}.
        \begin{byCase}
        \caseOf{$x_1\in X$ and $x_2\in X$.}
            Therefore $x_1 = x_2$
                by \cref{id_elem_intro,union_intro_left,function_apply_intro}.
        \caseOf{$x_1\in X$ and $x_2\in\{p\}$.}
            Suppose not.
            Hence $h(x_1) = x_1$ and $h(x_2) = q$
                by \cref{id_elem_intro,union_intro_left,function_apply_intro,singleton_iff,singleton_intro,union_intro_right}.
            Hence $x_1 = q$.
            Thus $x_1\notin X$.
            Contradiction.
            Therefore $x_1 = x_2$.
        \caseOf{$x_1\in\{p\}$ and $x_2\in X$.}
            Suppose not.
            Hence $h(x_1) = q$ and $h(x_2) = x_2$
                by \cref{singleton_iff,singleton_intro,union_intro_right,function_apply_intro,id_elem_intro,union_intro_left}.
            Hence $x_2 = q$.
            Thus $x_2\notin X$.
            Contradiction.
            Therefore $x_1 = x_2$.
        \caseOf{$x_1\in\{p\}$ and $x_2\in\{p\}$.}
            Then $x_1 = x_2$ by \cref{singleton_iff}.
        \end{byCase}
    \end{subproof}
    Therefore $h$ is injective by \cref{injective_function}.
    Hence $h\in\bijections{Y}{Y_1}$ by \cref{bijections}.

    Show that for all $U$ such that $U$ is open in $Y_1$ with respect to $S_1$
        we have $\preimg{h}{U}$ is open in $Y$ with respect to $S$.
    \begin{subproof}
        Fix $U$ such that $U$ is open in $Y_1$ with respect to $S_1$.
        \begin{byCase}
        \caseOf{$q\notin U$.}
            Hence $U\subseteq X$
                by \cref{open_in,topology_on,subseteq,pow_iff,union_iff,singleton_iff,subseteq}.
            Show that for all $y$ such that $y\in\preimg{h}{U}$ we have $y\in U$.
            \begin{subproof}
                Fix $y$ such that $y\in\preimg{h}{U}$.
                Take $b\in U$ such that $y\mathrel{h} b$ by \cref{preimg_iff}.
                Show that $y\notin\{p\}$.
                \begin{subproof}
                    Suppose not.
                    $(p,q)\in h$ by \cref{singleton_intro,union_intro_right}.
                    Then $h(y) = q$ and $h(y) = b$ by \cref{singleton_iff,function_apply_intro}.
                    Therefore $q\in U$.
                    Contradiction.
                \end{subproof}
                Therefore $y\in X$ by \cref{preimg_subseteq_dom,subseteq,union_iff,singleton_iff}.
                $(y,y)\in h$ by \cref{id_elem_intro,union_intro_left}.
                Hence $h(y) = y$ and $y = b$ by \cref{function_apply_intro}.
                Therefore $y\in U$.
            \end{subproof}
            Hence $\preimg{h}{U}\subseteq U$ by \cref{subseteq}.
            Therefore $U\subseteq\preimg{h}{U}$
                by \cref{subseteq,id_elem_intro,union_intro_left,function_apply_intro,preimg_iff}.
            Hence $\preimg{h}{U} = U$ by \cref{subseteq_antisymmetric}.
            Therefore $U$ is open in $X$ with respect to $T$
                by \cref{open_in,inter_absorb_supseteq_left,subspace_topology}.
            Hence $\{p\}$ is closed in $Y$ with respect to $S$
                by \cref{singleton_subset_intro,pair_is_finite,upair_intro_left,subseteq_finite_is_finite,finite_point_set_closed_hausdorff}.
            Therefore $Y\setminus\{p\}$ is open in $Y$ with respect to $S$ by \cref{closed_in}.
            Hence $X\subseteq Y\setminus\{p\}$
                by \cref{union_intro_left,singleton_iff,setminus_intro,subseteq}.
            Therefore $Y\setminus\{p\}\subseteq X$
                by \cref{setminus_elim_left,union_iff,singleton_iff,setminus_elim_right,subseteq}.
            Hence $X = Y\setminus\{p\}$ by \cref{subseteq_antisymmetric}.
            Hence $U$ is open in $Y$ with respect to $S$ by \cref{open_in_subspace_open_in_space,subspace_of}.
            Therefore $\preimg{h}{U}$ is open in $Y$ with respect to $S$.
        \caseOf{$q\in U$.}
            Let $C = Y_1\setminus U$.
            $C\subseteq Y_1$ and $U\in S_1$ by \cref{setminus_subseteq,open_in}.
            Hence $U\subseteq Y_1$ by \cref{topology_on,subseteq,pow_iff}.
            Therefore $Y_1\setminus C = U$ by \cref{setminus_setminus,inter_absorb_supseteq_left}.
            Hence $C$ is closed in $Y_1$ with respect to $S_1$ by \cref{open_in,closed_in}.
            Therefore $C\subseteq X$
                by \cref{setminus_elim_left,union_iff,singleton_iff,setminus_elim_right,subseteq}.
            Hence $C$ is compact with respect to $\subspaceTopology{S_1}{C}$ by \cref{closed_subspace_compact}.
            Therefore $C$ is a compact subspace of $X$ with respect to $T$
                by \cref{subspace_subspace_topology,compact_subspace}.
            Then $C$ is compact with respect to $\subspaceTopology{S}{C}$ by \cref{subspace_subspace_topology}.
            Hence $C$ is closed in $Y$ with respect to $S$
                by \cref{subseteq_transitive,compact_subspace,compact_subspace_closed}.
            Therefore $Y\setminus C$ is open in $Y$ with respect to $S$ by \cref{closed_in}.
            Show that for all $y$ such that $y\in\preimg{h}{U}$ we have $y\in Y\setminus C$.
            \begin{subproof}
                Fix $y$ such that $y\in\preimg{h}{U}$.
                Hence $y\in Y$ by \cref{preimg_subseteq_dom,subseteq}.
                Take $b\in U$ such that $y\mathrel{h} b$ by \cref{preimg_iff}.
                Show that $y\notin C$.
                \begin{subproof}
                    Suppose not.
                    Then $y\in X$ by \cref{subseteq}.
                    $(y,y)\in h$ by \cref{id_elem_intro,union_intro_left}.
                    Hence $h(y) = y$ and $y = b$ by \cref{function_apply_intro}.
                    Therefore $y\in U$.
                    Contradiction by \cref{setminus_elim_right}.
                \end{subproof}
                Hence $y\in Y\setminus C$ by \cref{setminus_intro}.
            \end{subproof}
            Therefore $\preimg{h}{U}\subseteq Y\setminus C$ by \cref{subseteq}.
            Hence $Y\setminus C\subseteq\preimg{h}{U}$
                by \cref{setminus_elim_left,setminus_elim_right,union_iff,union_intro_left,setminus_intro,id_elem_intro,function_apply_intro,preimg_iff,singleton_intro,union_intro_right,singleton_iff,subseteq}.
            Therefore $\preimg{h}{U} = Y\setminus C$ by \cref{subseteq_antisymmetric}.
            Hence $\preimg{h}{U}$ is open in $Y$ with respect to $S$ by \cref{subseteq}.
        \end{byCase}
    \end{subproof}
    Therefore $h$ is continuous from $Y$ to $Y_1$ with respect to $S$ and $S_1$ by \cref{continuous}.

    Hence $h$ is a homeomorphism between $Y$ and $Y_1$ with respect to $S$ and $S_1$
        by \cref{bijections,compact_hausdorff_bijection_homeomorphism}.

    Then $\restrl{h}{X}\subseteq \identity{X}$
        by \cref{restrl,id_elem_intro,union_intro_left,function_apply_intro,subseteq}.
    Therefore $\identity{X}\subseteq\restrl{h}{X}$
        by \cref{id_elem_inspect,id_elem_intro,union_intro_left,restrl,subseteq}.
    Hence $\restrl{h}{X} = \identity{X}$ by \cref{subseteq_antisymmetric}.
    Finally there exists $h$ such that
        $h$ is a homeomorphism between $Y$ and $Y_1$ with respect to $S$ and $S_1$ and
        $\restrl{h}{X} = \identity{X}$.
\end{proof}


% from §29 Item 6: compactification
\begin{definition}\label{compactification}
    $Y$ is a compactification of $X$ with respect to $S$ and $T$ iff
        $S$ is a topology on $Y$ and
        $T$ is a topology on $X$ and
        $X$ is a subspace of $Y$ with respect to $S$ and
        $T = \subspaceTopology{S}{X}$ and
        $X\neq Y$ and
        $\closure{X}{Y}{S} = Y$ and
        $Y$ is compact with respect to $S$ and
        $Y$ is Hausdorff with respect to $S$.
\end{definition}

% from §29 Item 7: one-point compactification
\begin{definition}\label{one_point_compactification}
    $Y$ is a one-point compactification of $X$ with respect to $S$ and $T$ iff
        $Y$ is a compactification of $X$ with respect to $S$ and $T$ and
        there exists $p$ such that $Y\setminus X = \{p\}$.
\end{definition}

% from §29 Item 8: one-point compactification characterization
\begin{theorem}\label{one_point_compactification_characterization}
    Assume $T$ is a topology on $X$.
    Then there exist $Y,S$ such that $Y$ is a one-point compactification of $X$ with respect to $S$ and $T$
        iff $X$ is locally compact with respect to $T$ and $X$ is Hausdorff with respect to $T$
            and $X$ is not compact with respect to $T$.
\end{theorem}
\begin{proof}
    Show that if there exist $Y,S$ such that $Y$ is a one-point compactification of $X$ with respect to $S$ and $T$ then
        $X$ is locally compact with respect to $T$ and $X$ is Hausdorff with respect to $T$ and
            $X$ is not compact with respect to $T$.
    \begin{subproof}
        Assume there exist $Y,S$ such that $Y$ is a one-point compactification of $X$ with respect to $S$ and $T$.
        Take $Y,S$ such that $Y$ is a one-point compactification of $X$ with respect to $S$ and $T$.
        Take $p$ such that $Y\setminus X = \{p\}$ and
            $Y$ is a compactification of $X$ with respect to $S$ and $T$
            by \cref{one_point_compactification}.
        $T = \subspaceTopology{S}{X}$ by \cref{compactification}.
        $Y$ is compact with respect to $S$ and $Y$ is Hausdorff with respect to $S$ by \cref{compactification}.
        $X\subseteq Y$ and $S$ is a topology on $Y$ by \cref{subspace_of,compactification}.
        Hence $X$ is locally compact with respect to $T$ and $X$ is Hausdorff with respect to $T$
            by \cref{one_point_compactification_existence_reverse}.

        Show that $X$ is not compact with respect to $T$.
        \begin{subproof}
            Suppose not.
            Then $X$ is compact with respect to $T$.
            Therefore $X$ is closed in $Y$ with respect to $S$
                by \cref{compact_subspace,compactification,compact_subspace_closed}.
            Hence $\closure{X}{Y}{S} = X$
                by \cref{subseteq_refl,closure_subseteq_closed,subseteq_closure,subseteq_antisymmetric}.
            Therefore $X = Y$ by \cref{compactification}.
            Contradiction by \cref{compactification}.
        \end{subproof}
        Hence $X$ is locally compact with respect to $T$ and $X$ is Hausdorff with respect to $T$
            and $X$ is not compact with respect to $T$.
    \end{subproof}

    Show that if $X$ is locally compact with respect to $T$ and $X$ is Hausdorff with respect to $T$
        and $X$ is not compact with respect to $T$
        then there exist $Y,S$ such that $Y$ is a one-point compactification of $X$ with respect to $S$ and $T$.
    \begin{subproof}
        Assume $X$ is locally compact with respect to $T$ and $X$ is Hausdorff with respect to $T$
            and $X$ is not compact with respect to $T$.
        Take $Y,S,p$ such that
            $S$ is a topology on $Y$ and
            $X\subseteq Y$ and
            $T = \subspaceTopology{S}{X}$ and
            $Y\setminus X = \{p\}$ and
            $Y$ is compact with respect to $S$ and
            $Y$ is Hausdorff with respect to $S$
            by \cref{one_point_compactification_existence_forward}.
        Then $Y = X\union\{p\}$ by \cref{setminus_partition}.

        Show that $p\in \closure{X}{Y}{S}$.
        \begin{subproof}
            Suppose not.
            Then $p\notin \closure{X}{Y}{S}$.
            For all $y$ such that $y\in\closure{X}{Y}{S}$ we have $y\in X$.
            Therefore $\closure{X}{Y}{S}\subseteq X$ by \cref{subseteq}.
            Therefore $X$ is closed in $Y$ with respect to $S$
                by \cref{subseteq_closure,subseteq_antisymmetric,closure_closed_in}.
            Hence $X$ is compact with respect to $T$ by \cref{closed_subspace_compact}.
            Contradiction.
        \end{subproof}
        Hence $\closure{X}{Y}{S} = Y$
            by \cref{closure_closed_in,subseteq_closure,singleton_subset_intro,union_subsets_is_subset,closed_in,subseteq_antisymmetric}.

        Then $X\neq Y$ by \cref{singleton_intro,setminus_elim_right}.
        Hence $X$ is a subspace of $Y$ with respect to $S$ by \cref{subspace_of}.
        Therefore $Y$ is a compactification of $X$ with respect to $S$ and $T$ by \cref{compactification}.
        Hence there exist $Y,S$ such that $Y$ is a one-point compactification of $X$ with respect to $S$ and $T$
            by \cref{one_point_compactification}.
    \end{subproof}
\end{proof}

% from §29 Item 9: local compactness via neighborhoods
\begin{theorem}\label{locally_compact_neighborhood_characterization}
    Assume $T$ is a topology on $X$.
    Assume $X$ is Hausdorff with respect to $T$.
    Then $X$ is locally compact with respect to $T$ iff
        for all $x\in X$ we have
            for all $U$ such that $U$ is a neighborhood of $x$ in $X$ with respect to $T$ we have
                there exists $V$ such that
                    $V$ is a neighborhood of $x$ in $X$ with respect to $T$ and
                    $\closure{V}{X}{T}$ is a compact subspace of $X$ with respect to $T$ and
                    $\closure{V}{X}{T}\subseteq U$.
\end{theorem}
\begin{proof}
    Show that if for all $x\in X$ we have
        for all $U$ such that $U$ is a neighborhood of $x$ in $X$ with respect to $T$ we have
            there exists $V$ such that
                $V$ is a neighborhood of $x$ in $X$ with respect to $T$ and
                $\closure{V}{X}{T}$ is a compact subspace of $X$ with respect to $T$ and
                $\closure{V}{X}{T}\subseteq U$
        then $X$ is locally compact with respect to $T$.
    \begin{subproof}
        Assume for all $x\in X$ we have
            for all $U$ such that $U$ is a neighborhood of $x$ in $X$ with respect to $T$ we have
                there exists $V$ such that
                    $V$ is a neighborhood of $x$ in $X$ with respect to $T$ and
                    $\closure{V}{X}{T}$ is a compact subspace of $X$ with respect to $T$ and
                    $\closure{V}{X}{T}\subseteq U$.
        Let $U = X$.
        For all $x$ such that $x\in X$ we have $U$ is a neighborhood of $x$ in $X$ with respect to $T$
            by \cref{topology_on,open_in,neighborhood}.
        Thus for all $x$ such that $x\in X$ there exists $V$ such that
            $V$ is a neighborhood of $x$ in $X$ with respect to $T$ and
            $\closure{V}{X}{T}$ is a compact subspace of $X$ with respect to $T$ and
            $\closure{V}{X}{T}\subseteq U$.
        Thus for all $x\in X$ we have $X$ is locally compact at $x$ with respect to $T$
            by \cref{neighborhood,open_in,topology_on,subseteq,pow_iff,subseteq_closure,locally_compact_at_point}.
        Therefore $X$ is locally compact with respect to $T$ by \cref{locally_compact}.
    \end{subproof}

    Show that if $X$ is locally compact with respect to $T$ then
        for all $x\in X$ we have
            for all $U$ such that $U$ is a neighborhood of $x$ in $X$ with respect to $T$ we have
                there exists $V$ such that
                    $V$ is a neighborhood of $x$ in $X$ with respect to $T$ and
                    $\closure{V}{X}{T}$ is a compact subspace of $X$ with respect to $T$ and
                    $\closure{V}{X}{T}\subseteq U$.
    \begin{subproof}
        Assume $X$ is locally compact with respect to $T$.
        Let $p = X$.
        $p\notin X$ by \cref{set_notin_self}.
        Let $Y = X\union\{p\}$.
        Let $S = \onePointTopology{T}{X}{p}$.
        $S$ is a topology on $Y$ and $Y$ is compact with respect to $S$ and $Y$ is Hausdorff with respect to $S$
            by \cref{one_point_topology_is_topology,one_point_topology_compact,one_point_topology_hausdorff}.
        $T = \subspaceTopology{S}{X}$ by \cref{one_point_topology_subspace}.
        $X\subseteq Y$ by \cref{union_upper_left}.
        Hence $X$ is a subspace of $Y$ with respect to $S$ by \cref{subspace_of}.

        Show that for all $x\in X$ we have
            for all $U$ such that $U$ is a neighborhood of $x$ in $X$ with respect to $T$ we have
                there exists $V$ such that
                    $V$ is a neighborhood of $x$ in $X$ with respect to $T$ and
                    $\closure{V}{X}{T}$ is a compact subspace of $X$ with respect to $T$ and
                    $\closure{V}{X}{T}\subseteq U$.
        \begin{subproof}
            Fix $x$ such that $x\in X$.
            Fix $U$ such that $U$ is a neighborhood of $x$ in $X$ with respect to $T$.
            $U$ is open in $X$ with respect to $T$ by \cref{neighborhood}.
            Thus $U\subseteq X$ by \cref{open_in,topology_on,subseteq,pow_iff}.

            Hence $X$ is open in $Y$ with respect to $S$
                by \cref{topology_on,subseteq_refl,one_point_topology_intro_left,open_in}.
            Therefore $U$ is open in $Y$ with respect to $S$ by
                \cref{open_in_subspace_open_in_space,subspace_of}.

            Let $C = Y\setminus U$.
            $C\subseteq Y$ by \cref{setminus_subseteq}.
            Then $U\in S$ and $U\subseteq Y$ by \cref{open_in,topology_on,subseteq,pow_iff}.
            Hence $Y\setminus C = U$ by \cref{setminus_setminus,inter_absorb_supseteq_left}.
            Therefore $C$ is closed in $Y$ with respect to $S$ and
                $C$ is compact with respect to $\subspaceTopology{S}{C}$
                by \cref{setminus_subseteq,open_in,closed_in,closed_subspace_compact}.

            $x\in U$ by \cref{neighborhood}.
            Hence $x\notin C$ by \cref{setminus_elim_right}.
            Take $V,W$ such that
                $V$ is open in $Y$ with respect to $S$ and
                $W$ is open in $Y$ with respect to $S$ and
                $x\in V$ and
                $C\subseteq W$ and
                $V$ is disjoint from $W$
                by \cref{compact_hausdorff_separation}.

            Then $p\notin U$ by \cref{subseteq}.
            Therefore $p\in W$ by \cref{singleton_intro,union_intro_right,setminus_intro,subseteq}.
            Hence $p\notin V$ by \cref{disjoint}.

            For all $y$ such that $y\in V$ we have $y\in X$.
            Therefore $V\subseteq X$ by \cref{subseteq}.

            Hence $V$ is a neighborhood of $x$ in $X$ with respect to $T$
                by \cref{inter_absorb_supseteq_left,subspace_topology,open_in,neighborhood}.
            Then $V\subseteq Y$ by \cref{subseteq_transitive}.

            Let $B = \closure{V}{Y}{S}$.
            Then $B$ is closed in $Y$ with respect to $S$ and
                $B$ is compact with respect to $\subspaceTopology{S}{B}$
                by \cref{closure_closed_in,closed_in,closed_subspace_compact}.

            Hence $Y\setminus W$ is closed in $Y$ with respect to $S$ by \cref{setminus_subseteq,open_in,topology_on,subseteq,pow_iff,setminus_setminus,inter_absorb_supseteq_left,closed_in}.
            Therefore $V\subseteq Y\setminus W$ by \cref{disjoint,subseteq,setminus_intro,subseteq}.
            Hence $B\subseteq Y\setminus W$ by \cref{closure_subseteq_closed}.

            For all $y$ such that $y\in Y\setminus W$ we have $y\in Y\setminus C$.
            Therefore $Y\setminus W\subseteq Y\setminus C$ by \cref{subseteq}.
            Therefore $Y\setminus W\subseteq U$ by \cref{subseteq}.
            Hence $B\subseteq U$ and $B\subseteq X$ by \cref{subseteq_transitive}.

            Therefore $B$ is a compact subspace of $X$ with respect to $T$
                by \cref{subspace_subspace_topology,compact_subspace}.

            Then $\closure{V}{X}{T} = B$ by
                \cref{closure_in_subspace,subspace_of,inter_comm,inter_absorb_supseteq_left}.
            Hence $\closure{V}{X}{T}$ is a compact subspace of $X$ with respect to $T$.
            Therefore $\closure{V}{X}{T}\subseteq U$ by \cref{subseteq}.
        \end{subproof}
    \end{subproof}
\end{proof}

% from §29 Item 10: locally compact subspaces
\begin{corollary}\label{locally_compact_subspace}
    Assume $T$ is a topology on $X$.
    Assume $X$ is locally compact with respect to $T$.
    Assume $X$ is Hausdorff with respect to $T$.
    Suppose $A$ is a subspace of $X$ with respect to $T$.
    If $A$ is closed in $X$ with respect to $T$ or $A$ is open in $X$ with respect to $T$ then
        $A$ is locally compact with respect to $\subspaceTopology{T}{A}$.
\end{corollary}
\begin{proof}
    Let $T_A = \subspaceTopology{T}{A}$.
    Assume $A$ is closed in $X$ with respect to $T$ or $A$ is open in $X$ with respect to $T$.
    \begin{byCase}
    \caseOf{$A$ is closed in $X$ with respect to $T$.}
        Show that for all $x\in A$ we have $A$ is locally compact at $x$ with respect to $T_A$.
        \begin{subproof}
            Fix $x$ such that $x\in A$.
            $A\subseteq X$ by \cref{subspace_of}.
            Hence $x\in X$ by \cref{subseteq}.
            $X$ is locally compact at $x$ with respect to $T$ by \cref{locally_compact}.
            Take $C$ such that
                $C$ is a compact subspace of $X$ with respect to $T$ and
                there exists $U$ such that
                    $U$ is a neighborhood of $x$ in $X$ with respect to $T$ and
                    $U\subseteq C$
                by \cref{locally_compact_at_point}.
            Take $U$ such that
                $U$ is a neighborhood of $x$ in $X$ with respect to $T$ and
                $U\subseteq C$.
            Let $D = C\inter A$.
            $D\subseteq A$ by \cref{inter_lower_right}.
            Let $T_C = \subspaceTopology{T}{C}$.
            $C$ is compact with respect to $T_C$ by \cref{compact_subspace}.
            $C$ is a subspace of $X$ with respect to $T$ by \cref{subspace_of,compact_subspace}.
            Then $D$ is closed in $C$ with respect to $T_C$ by \cref{inter_comm,closed_in_subspace_iff}.
            $D\subseteq C$ by \cref{inter_lower_left}.
            Hence $T_C$ is a topology on $C$ by \cref{subspace_topology_is_topology,subspace_of}.
            $C\subseteq X$ by \cref{compact_subspace}.
            Therefore $D$ is compact with respect to $\subspaceTopology{T_C}{D}$ by \cref{closed_subspace_compact}.
            Let $T_D = \subspaceTopology{T}{D}$.
            Hence $\subspaceTopology{T_C}{D} = T_D$ by \cref{subspace_subspace_topology}.
            Let $T_{AD} = \subspaceTopology{T_A}{D}$.
            Therefore $T_D = T_{AD}$ by \cref{subspace_subspace_topology}.
            Hence $D$ is a compact subspace of $A$ with respect to $T_A$
                by \cref{subspace_subspace_topology,compact_subspace}.

            Then $U\inter A$ is a neighborhood of $x$ in $A$ with respect to $T_A$
                by \cref{inter_comm,neighborhood,open_in,subspace_topology,subspace_topology_is_topology,subspace_of,inter_intro}.
            Therefore $U\inter A\subseteq C\inter A$
                by \cref{inter_lower_left,subseteq_transitive,inter_lower_right,subseteq_inter_iff}.
            Hence $A$ is locally compact at $x$ with respect to $T_A$ by \cref{locally_compact_at_point}.
        \end{subproof}
        Therefore $A$ is locally compact with respect to $T_A$ by \cref{locally_compact}.
    \caseOf{$A$ is open in $X$ with respect to $T$.}
        Show that for all $x\in A$ we have $A$ is locally compact at $x$ with respect to $T_A$.
        \begin{subproof}
            Fix $x$ such that $x\in A$.
            $A\subseteq X$ by \cref{subspace_of}.
            Hence $x\in X$ by \cref{subseteq}.
            $A$ is a neighborhood of $x$ in $X$ with respect to $T$ by \cref{neighborhood}.
            Take $V$ such that
                $V$ is a neighborhood of $x$ in $X$ with respect to $T$ and
                $\closure{V}{X}{T}$ is a compact subspace of $X$ with respect to $T$ and
                $\closure{V}{X}{T}\subseteq A$
                by \cref{locally_compact_neighborhood_characterization}.

            $V$ is open in $X$ with respect to $T$ by \cref{neighborhood}.
            Then $V\subseteq X$ by \cref{open_in,topology_on,subseteq,pow_iff}.
            $V\subseteq \closure{V}{X}{T}$ by \cref{subseteq_closure}.
            Hence $V\subseteq A$ by \cref{subseteq_transitive}.
            Therefore $V$ is open in $A$ with respect to $T_A$ and
                $V$ is a neighborhood of $x$ in $A$ with respect to $T_A$
                by \cref{inter_absorb_supseteq_left,neighborhood,open_in,subspace_topology,subspace_topology_is_topology,subspace_of}.

            Let $D = \closure{V}{X}{T}$.
            $D\subseteq A$ by \cref{subseteq}.
            Let $T_D = \subspaceTopology{T}{D}$.
            Let $T_{AD} = \subspaceTopology{T_A}{D}$.
            Hence $T_D = T_{AD}$ by \cref{subspace_subspace_topology}.
            Therefore $D$ is compact with respect to $T_{AD}$ and
                $D$ is a compact subspace of $A$ with respect to $T_A$ by \cref{compact_subspace}.
            Hence $A$ is locally compact at $x$ with respect to $T_A$ by \cref{locally_compact_at_point}.
        \end{subproof}
        Therefore $A$ is locally compact with respect to $T_A$ by \cref{locally_compact}.
    \end{byCase}
\end{proof}

% from §29 helper lemma 15: homeomorphism preserves local compactness
\begin{lemma}\label{homeomorphism_preserves_local_compactness}
    Assume $T$ is a topology on $X$.
    Assume $S$ is a topology on $Y$.
    Suppose $f$ is a homeomorphism between $X$ and $Y$ with respect to $T$ and $S$.
    If $Y$ is locally compact with respect to $S$ then $X$ is locally compact with respect to $T$.
\end{lemma}
\begin{proof}
    Assume $Y$ is locally compact with respect to $S$.
    Let $g = \converse{f}$.
    $f\in\funs{X}{Y}$ and $f$ is a function and $\dom{f} = X$
        by \cref{homeomorphism,bijections,bijections_to_funs,funs_is_function,funs}.
    Hence $g$ is a bijection from $Y$ to $X$ and $g\in\bijections{Y}{X}$ and $\dom{g} = Y$ and $g$ is a function
        by \cref{homeomorphism,bijection_converse_is_bijection,bijections,bijections_dom,bijection_is_function}.
    $g$ is continuous from $Y$ to $X$ with respect to $S$ and $T$ by \cref{homeomorphism}.

    Show that for all $x\in X$ we have $X$ is locally compact at $x$ with respect to $T$.
        \begin{subproof}
        Fix $x$ such that $x\in X$.
        Let $y = f(x)$.
        $Y$ is locally compact at $y$ with respect to $S$ by \cref{funs_type_apply,locally_compact}.
        Take $C$ such that
            $C$ is a compact subspace of $Y$ with respect to $S$ and
            there exists $U$ such that
                $U$ is a neighborhood of $y$ in $Y$ with respect to $S$ and
                $U\subseteq C$
            by \cref{locally_compact_at_point}.
        Take $U$ such that
            $U$ is a neighborhood of $y$ in $Y$ with respect to $S$ and
            $U\subseteq C$.

        Let $g_C = \restrl{g}{C}$.
        Then $C\subseteq \dom{g}$ and $C$ is a subspace of $Y$ with respect to $S$
            by \cref{compact_subspace,bijections_dom,subspace_of}.
        Let $T_C = \subspaceTopology{S}{C}$.
        $g_C$ is continuous from $C$ to $X$ with respect to $T_C$ and $T$ by \cref{continuous_restrict_domain}.
        $C$ is compact with respect to $T_C$ by \cref{compact_subspace}.
        Let $A = \img{g_C}{C}$.
        $A$ is compact with respect to $\subspaceTopology{T}{A}$ by \cref{continuous_image_compact}.
        $g_C\in\funs{C}{X}$ by \cref{continuous}.
        Hence $A\subseteq X$
            by \cref{continuous,funs_subseteq_rels,subseteq,rels_ran_subseteq,img_subseteq_ran,subseteq_transitive}.
        Therefore $A$ is a compact subspace of $X$ with respect to $T$ by \cref{compact_subspace}.

        Let $V = \img{g}{U}$.
        $V = \preimg{f}{U}$ by \cref{preim_eq_img_of_converse}.
        $U$ is open in $Y$ with respect to $S$ and $y\in U$ by \cref{neighborhood}.
        $f$ is continuous from $X$ to $Y$ with respect to $T$ and $S$ by \cref{homeomorphism}.
        Hence $V$ is a neighborhood of $x$ in $X$ with respect to $T$
            by \cref{continuous,function_apply_elim,preimg_iff,neighborhood}.

        Therefore $V\subseteq A$
            by \cref{subseteq,img_subseteq,continuous_restrict_domain,funs,img_dom,restrl_ran}.
        Hence $X$ is locally compact at $x$ with respect to $T$ by \cref{locally_compact_at_point}.
    \end{subproof}
    Therefore $X$ is locally compact with respect to $T$ by \cref{locally_compact}.
\end{proof}

% from §29 Item 11: locally compact Hausdorff embeds as open subspace
\begin{corollary}\label{locally_compact_open_subspace_compact_hausdorff_forward}
    Assume $T$ is a topology on $X$.
    If $X$ is locally compact with respect to $T$ and $X$ is Hausdorff with respect to $T$ then
        there exist $Y,S,A,f$ such that
            $S$ is a topology on $Y$ and
            $Y$ is compact with respect to $S$ and
            $Y$ is Hausdorff with respect to $S$ and
            $A$ is open in $Y$ with respect to $S$ and
            $f$ is a homeomorphism between $X$ and $A$ with respect to $T$ and $\subspaceTopology{S}{A}$.
\end{corollary}
\begin{proof}
    Assume $X$ is locally compact with respect to $T$ and $X$ is Hausdorff with respect to $T$.
    Let $p = X$.
    $p\notin X$ by \cref{set_notin_self}.
    Let $Y = X\union\{p\}$.
    Let $S = \onePointTopology{T}{X}{p}$.
    $S$ is a topology on $Y$ and $Y$ is compact with respect to $S$ and $Y$ is Hausdorff with respect to $S$
        by \cref{one_point_topology_is_topology,one_point_topology_compact,one_point_topology_hausdorff}.
    Let $A = X$.
    Hence $A$ is open in $Y$ with respect to $S$ and $T = \subspaceTopology{S}{A}$
        by \cref{topology_on,subseteq,one_point_topology_intro_left,open_in,one_point_topology_subspace}.
    Let $f = \identity{A}$.
    $f$ is a bijection from $X$ to $A$ by \cref{id_elem_bijections}.
    $f$ is continuous from $X$ to $A$ with respect to $T$ and $\subspaceTopology{S}{A}$ by
        \cref{identity_continuous}.
    Therefore $\converse{f}$ is a function from $A$ to $X$ by \cref{bijective_converse_are_funs}.
    Hence $A$ is a subspace of $Y$ with respect to $S$ and $\subspaceTopology{S}{A}$ is a topology on $A$
        by \cref{open_in,topology_on,subseteq,pow_iff,subspace_of,subspace_topology_is_topology}.
    For all $U$ such that $U$ is open in $X$ with respect to $T$ we have
        $\preimg{\converse{f}}{U}$ is open in $A$ with respect to $\subspaceTopology{S}{A}$
        by \cref{open_in,topology_on,subseteq,pow_iff,id_is_relation,preim_eq_img_of_converse,converse_converse_eq,id_img,inter_absorb_supseteq_left}.
    Therefore $\converse{f}$ is continuous from $A$ to $X$ with respect to $\subspaceTopology{S}{A}$ and $T$
        by \cref{continuous}.
    Hence $f$ is a homeomorphism between $X$ and $A$ with respect to $T$ and $\subspaceTopology{S}{A}$
        by \cref{homeomorphism}.
    Therefore there exist $Y,S,A,f$ such that
        $S$ is a topology on $Y$ and
        $Y$ is compact with respect to $S$ and
        $Y$ is Hausdorff with respect to $S$ and
        $A$ is open in $Y$ with respect to $S$ and
        $f$ is a homeomorphism between $X$ and $A$ with respect to $T$ and $\subspaceTopology{S}{A}$.
\end{proof}

% from §29 Item 12: open subspace of compact Hausdorff is locally compact Hausdorff
\begin{corollary}\label{open_subspace_compact_hausdorff_locally_compact_hausdorff}
    Assume $T$ is a topology on $X$.
    If there exist $Y,S,A,f$ such that
        $S$ is a topology on $Y$ and
        $Y$ is compact with respect to $S$ and
        $Y$ is Hausdorff with respect to $S$ and
        $A$ is open in $Y$ with respect to $S$ and
        $f$ is a homeomorphism between $X$ and $A$ with respect to $T$ and $\subspaceTopology{S}{A}$
    then $X$ is locally compact with respect to $T$ and $X$ is Hausdorff with respect to $T$.
\end{corollary}
\begin{proof}
    Assume there exist $Y,S,A,f$ such that
        $S$ is a topology on $Y$ and
        $Y$ is compact with respect to $S$ and
        $Y$ is Hausdorff with respect to $S$ and
        $A$ is open in $Y$ with respect to $S$ and
        $f$ is a homeomorphism between $X$ and $A$ with respect to $T$ and $\subspaceTopology{S}{A}$.
    Take $Y,S,A,f$ such that
        $S$ is a topology on $Y$ and
        $Y$ is compact with respect to $S$ and
        $Y$ is Hausdorff with respect to $S$ and
        $A$ is open in $Y$ with respect to $S$ and
    $f$ is a homeomorphism between $X$ and $A$ with respect to $T$ and $\subspaceTopology{S}{A}$.

    $Y$ is locally compact with respect to $S$ by \cref{compact_implies_locally_compact}.
    Hence $A\subseteq Y$ by \cref{open_in,topology_on,subseteq,pow_iff}.
    Let $T_A = \subspaceTopology{S}{A}$.
    Therefore $A$ is a subspace of $Y$ with respect to $S$ and $T_A$ is a topology on $A$
        by \cref{subspace_of,subspace_topology_is_topology}.
    $A$ is Hausdorff with respect to $T_A$ by \cref{subspace_hausdorff}.
    $A$ is locally compact with respect to $T_A$ by \cref{locally_compact_subspace}.
    Hence $X$ is locally compact with respect to $T$ by \cref{homeomorphism_preserves_local_compactness}.

    Show that for all $x_1,x_2$ such that $x_1\in X$ and $x_2\in X$ and $x_1\neq x_2$ there exist $U_1,U_2$ such that
        $U_1$ is a neighborhood of $x_1$ in $X$ with respect to $T$ and
        $U_2$ is a neighborhood of $x_2$ in $X$ with respect to $T$ and
        $U_1$ is disjoint from $U_2$.
    \begin{subproof}
            Fix $x_1$.
            Fix $x_2$ such that $x_1\in X$ and $x_2\in X$ and $x_1\neq x_2$.
            $f\in\funs{X}{A}$ and $f$ is a function and $f$ is injective
                by \cref{bijections_to_funs,homeomorphism,bijection_is_function,bijections}.
            Let $y_1 = f(x_1)$.
            Let $y_2 = f(x_2)$.
            $y_1\in A$ and $y_2\in A$ by \cref{funs_type_apply,homeomorphism}.
            Hence $y_1\neq y_2$ by \cref{funs,subseteq,injective_function}.
            Take $V_1,V_2$ such that
                $V_1$ is a neighborhood of $y_1$ in $A$ with respect to $T_A$ and
                $V_2$ is a neighborhood of $y_2$ in $A$ with respect to $T_A$ and
                $V_1$ is disjoint from $V_2$
                by \cref{hausdorff}.
            Let $U_1 = \preimg{f}{V_1}$.
            Let $U_2 = \preimg{f}{V_2}$.
            Hence $U_1$ is a neighborhood of $x_1$ in $X$ with respect to $T$ and
                $U_2$ is a neighborhood of $x_2$ in $X$ with respect to $T$
                by \cref{neighborhood,continuous,homeomorphism,funs,subseteq,function_apply_elim,preimg_iff}.
            Therefore $U_1$ is disjoint from $U_2$
                by \cref{preimg_iff,funs_is_function,function_apply_iff,disjoint}.
            Hence there exist $U_1,U_2$ such that
                $U_1$ is a neighborhood of $x_1$ in $X$ with respect to $T$ and
                $U_2$ is a neighborhood of $x_2$ in $X$ with respect to $T$ and
                $U_1$ is disjoint from $U_2$.
    \end{subproof}
    Therefore $X$ is Hausdorff with respect to $T$ by \cref{hausdorff}.
    Hence $X$ is locally compact with respect to $T$ and $X$ is Hausdorff with respect to $T$.
\end{proof}
\section{§30 The Countability Axioms}

% from §30 Item 1: sequence implies closure
\begin{lemma}\label{sequence_implies_closure_first_countable}
    Assume $T$ is a topology on $X$.
    Suppose $A\subseteq X$.
    Suppose $s$ is a sequence in $A$.
    Suppose $x\in X$.
    Suppose $s$ converges to $x$ in $X$ with respect to $T$.
    Then $x\in\closure{A}{X}{T}$.
\end{lemma}
\begin{proof}
    Follows by \cref{sequence_implies_closure}.
\end{proof}

% from §30 helper lemma 1: nested countable basis at a point
\begin{lemma}\label{nested_countable_basis}
    Suppose $b$ is a countable basis at $x$ in $X$ with respect to $T$.
    Then there exists $c$ such that
        $c$ is a countable basis at $x$ in $X$ with respect to $T$ and
        for all $m,n\in\naturalsPlus$ if $m\leq n$ then $c(n)\subseteq c(m)$.
\end{lemma}
\begin{proof}
    $1\in\naturalsPlus$ by \cref{real_one_in_naturals}.
    $b$ is a sequence in $\pow{X}$ by \cref{countable_basis_at_point}. $b(1)$ is a neighborhood of $x$ in $X$ with respect to $T$ by \cref{countable_basis_at_point}.
    Then $T$ is a topology on $X$ and $b(1)$ is open in $X$ with respect to $T$ by \cref{neighborhood,open_in}.
    Let $F = \{(n,\img{\restrl{b}{\initialSegment{n}}}{\initialSegment{n}}) \mid n\in\naturalsPlus\}$.
    Let $C = \{(n,\inters{\apply{F}{n}}) \mid n\in\naturalsPlus\}$.

    For all $n,U_1,U_2$ such that $(n,U_1)\in F$ and $(n,U_2)\in F$ we have $U_1 = U_2$
        by \cref{pair_eq_iff}.
    Hence $F$ is right-unique by \cref{rightunique}.
    Therefore $F$ is a function.
    For all $n$ such that $n\in\naturalsPlus$ we have $n\in\dom{F}$
        by \cref{pair_eq_iff,dom_intro}.
    For all $n$ such that $n\in\dom{F}$ we have $n\in\naturalsPlus$
        by \cref{dom_iff,pair_eq_iff}.
    Therefore $F$ is a function on $\naturalsPlus$ by \cref{subseteq,subseteq_antisymmetric}.

    For all $n\in\naturalsPlus$ we have
        $F(n) = \img{\restrl{b}{\initialSegment{n}}}{\initialSegment{n}}$
        by \cref{pair_eq_iff,function_apply_intro}.

    For all $n,U_1,U_2$ such that $(n,U_1)\in C$ and $(n,U_2)\in C$ we have $U_1 = U_2$
        by \cref{pair_eq_iff}.
    Hence $C$ is right-unique by \cref{rightunique}.
    Therefore $C$ is a function.
    For all $n$ such that $n\in\naturalsPlus$ we have $n\in\dom{C}$
        by \cref{pair_eq_iff,dom_intro}.
    For all $n$ such that $n\in\dom{C}$ we have $n\in\naturalsPlus$
        by \cref{dom_iff,pair_eq_iff}.
    Therefore $C$ is a function on $\naturalsPlus$ by \cref{subseteq,subseteq_antisymmetric}.

    For all $n\in\naturalsPlus$ we have $C(n) = \inters{F(n)}$
        by \cref{pair_eq_iff,function_apply_intro}.

    Show that for all $n\in\naturalsPlus$ we have $C(n)$ is a neighborhood of $x$ in $X$ with respect to $T$.
    \begin{subproof}
        Fix $n\in\naturalsPlus$.
        $\initialSegment{n}$ is finite by \cref{initial_segment_finite}.
        $b\in\funs{\naturalsPlus}{\pow{X}}$ and $\dom{b} = \naturalsPlus$
            by \cref{sequence_in,funs}.
        Thus $\initialSegment{n}\subseteq\dom{b}$ by \cref{initial_segment_naturalsplus,subseteq}.
        Let $g = \restrl{b}{\initialSegment{n}}$.
        $g$ is a function by \cref{sequence_in,funs_is_function,restrl_of_function_is_function}.
        Thus $\dom{g}\subseteq\initialSegment{n}$ by \cref{restrl_dom}.
        For all $i$ such that $i\in\initialSegment{n}$ we have $i\in\dom{g}$
            by \cref{initial_segment_naturalsplus,subseteq,dom_iff,restrl_iff,dom_intro}.
        Thus $\initialSegment{n}\subseteq\dom{g}$ by \cref{subseteq}.
        Thus $\dom{g} = \initialSegment{n}$ by \cref{subseteq_antisymmetric}.
        Thus $g$ is a function on $\initialSegment{n}$.
        Thus $F(n) = \img{g}{\initialSegment{n}}$.
        Thus $F(n)$ is finite by \cref{finite_image_function}.
        Thus $1\in\initialSegment{n}$ by \cref{real_one_in_naturals,real_one_leq_naturalsplus,initial_segment_naturalsplus}.
        Take $U_1$ such that $1\mathrel{b} U_1$ by \cref{dom_iff,funs}.
        Thus $F(n)$ is inhabited by \cref{img_iff,restrl_iff,nonempty_inhabited}.
        For all $U$ such that $U\in F(n)$ we have $U\in T$
            by \cref{img_iff,restrl_iff,sequence_in,funs,initial_segment_naturalsplus,subseteq,funs_is_function,function_apply_iff,countable_basis_at_point,neighborhood,open_in}.
        Therefore $C(n)\in T$ by \cref{subseteq,finite_intersections_open_in_topology}.
        For all $U$ such that $U\in F(n)$ we have $x\in U$
            by \cref{img_iff,restrl_iff,sequence_in,funs,initial_segment_naturalsplus,subseteq,funs_is_function,function_apply_iff,countable_basis_at_point,neighborhood}.
        Hence $x\in C(n)$ by \cref{inters_iff_forall}.
        Therefore $C(n)$ is a neighborhood of $x$ in $X$ with respect to $T$ by \cref{neighborhood,open_in}.
    \end{subproof}

    For all $n\in\dom{C}$ we have $C(n)\in\pow{X}$
        by \cref{subseteq,neighborhood,open_in,topology_on}.
    Hence $C$ is a sequence in $\pow{X}$ by \cref{funs_intro,sequence_in}.
    Thus for all $U$ such that $U$ is a neighborhood of $x$ in $X$ with respect to $T$
        there exists $n\in\naturalsPlus$ such that $C(n)\subseteq U$
        by \cref{countable_basis_at_point,initial_segment_naturalsplus,sequence_in,funs,subseteq,funs_is_function,function_apply_iff,restrl_iff,img_iff,inters_subseteq_elem,subseteq_transitive}.
    Therefore $C$ is a countable basis at $x$ in $X$ with respect to $T$ by \cref{countable_basis_at_point}.

    Show that for all $m,n$ such that $m\in\naturalsPlus$ and $n\in\naturalsPlus$ and $m\leq n$ we have $C(n)\subseteq C(m)$.
        \begin{subproof}
            Fix $m$.
            Fix $n$ such that $m\in\naturalsPlus$ and $n\in\naturalsPlus$ and $m\leq n$.
            Show that for all $z$ such that $z\in C(n)$ we have $z\in C(m)$.
            \begin{subproof}
                Fix $z$ such that $z\in C(n)$.
                Then for all $U$ such that $U\in F(n)$ we have $z\in U$ by \cref{inters_iff_forall}.
                For all $i$ such that $i\in\initialSegment{m}$ we have $i\in\initialSegment{n}$
                    by \cref{initial_segment_naturalsplus,real_naturals_subset_reals,subseteq,real_leq_trans}.
                Hence $\initialSegment{m}\subseteq\initialSegment{n}$ by \cref{subseteq}.
                For all $y$ such that $y\in\img{\restrl{b}{\initialSegment{m}}}{\initialSegment{m}}$ we have
                    $y\in\img{\restrl{b}{\initialSegment{n}}}{\initialSegment{n}}$
                    by \cref{img_iff,restrl_iff,subseteq}.
                Therefore $\img{\restrl{b}{\initialSegment{m}}}{\initialSegment{m}}\subseteq
                    \img{\restrl{b}{\initialSegment{n}}}{\initialSegment{n}}$ by \cref{subseteq}.
                Hence for all $U$ such that $U\in F(m)$ we have $U\in F(n)$ by \cref{subseteq}.
                Therefore for all $U$ such that $U\in F(m)$ we have $z\in U$.
                Hence $1\in\initialSegment{m}$ by \cref{real_one_in_naturals,real_one_leq_naturalsplus,initial_segment_naturalsplus}.
                Thus $(1,b(1))\in\restrl{b}{\initialSegment{m}}$ by \cref{sequence_in,funs,subseteq,funs_is_function,function_apply_iff,restrl_iff}.
                Therefore $F(m)$ is inhabited by \cref{img_iff,nonempty_inhabited}.
                Hence $z\in C(m)$ by \cref{inters_iff_forall}.
            \end{subproof}
        Therefore $C(n)\subseteq C(m)$ by \cref{subseteq}.
    \end{subproof}

    Finally there exists $c$ such that
        $c$ is a countable basis at $x$ in $X$ with respect to $T$ and
        for all $m,n\in\naturalsPlus$ if $m\leq n$ then $c(n)\subseteq c(m)$.
\end{proof}

% from §30 Item 2: closure implies sequence in a first-countable space
\begin{lemma}\label{closure_implies_sequence_first_countable}
    Assume $T$ is a topology on $X$.
    Suppose $A\subseteq X$.
    Suppose $x\in X$.
    Suppose $X$ satisfies the first countability axiom with respect to $T$.
    Suppose $x\in\closure{A}{X}{T}$.
    Then there exists $s$ such that $s$ is a sequence in $A$ and $s$ converges to $x$ in $X$ with respect to $T$.
\end{lemma}
\begin{proof}
    Take $b$ such that $b$ is a countable basis at $x$ in $X$ with respect to $T$ by \cref{first_countability_axiom}.
    Take $c$ such that
        $c$ is a countable basis at $x$ in $X$ with respect to $T$ and
        for all $m,n\in\naturalsPlus$ if $m\leq n$ then $c(n)\subseteq c(m)$ by \cref{nested_countable_basis}.

    Let $R = \{w\in\naturalsPlus\times X \mid \snd{w}\in A\inter \apply{c}{\fst{w}}\}$.
    Hence $R$ is a relation by \cref{times_elem_is_tuple,relation}.
    Thus for all $n\in\naturalsPlus$ there exists $y$ such that $n\mathrel{R} y$
        by \cref{countable_basis_at_point,closure_iff_intersects_open,neighborhood,intersects,nonempty_inhabited,inter_elim_right,inter_elim_left,subseteq,times_tuple_intro,fst_eq,snd_eq,inter_intro,pair_eq_iff}.
    Take $s$ such that $s$ is a function on $\naturalsPlus$ and
        for all $n\in\naturalsPlus$ we have $n\mathrel{R} s(n)$ by \cref{axiom_of_choice}.

    Then for all $n$ such that $n\in\dom{s}$ we have $s(n)\in A$
        by \cref{choice_function,subseteq,pair_eq_iff,fst_eq,snd_eq,inter}.
    Hence $s$ is a sequence in $A$ by \cref{funs_intro,sequence_in}.

    Therefore $s$ is a sequence in $X$ by \cref{sequence_in,funs_weaken_codom}.
    Show that for all $U$ such that $U$ is a neighborhood of $x$ in $X$ with respect to $T$
        there exists $N\in\naturalsPlus$ such that for all $n\in\naturalsPlus$ if $N\leq n$ then $s(n)\in U$.
    \begin{subproof}
        Fix $U$ such that $U$ is a neighborhood of $x$ in $X$ with respect to $T$.
        $c$ is a countable basis at $x$ in $X$ with respect to $T$.
        Hence there exists $N\in\naturalsPlus$ such that $c(N)\subseteq U$.
        Take $N\in\naturalsPlus$ such that $c(N)\subseteq U$.
        For all $n\in\naturalsPlus$ if $N\leq n$ then $s(n)\in U$
            by \cref{nested_countable_basis,subseteq_transitive,pair_eq_iff,fst_eq,snd_eq,inter,subseteq}.
        Hence there exists $N\in\naturalsPlus$ such that for all $n\in\naturalsPlus$ if $N\leq n$ then $s(n)\in U$.
    \end{subproof}
    Therefore $s$ converges to $x$ in $X$ with respect to $T$ by \cref{converges_to}.

    Finally there exists $s$ such that $s$ is a sequence in $A$ and $s$ converges to $x$ in $X$ with respect to $T$.
\end{proof}

% from §30 Item 3: continuous implies sequential continuity
\begin{theorem}\label{continuous_implies_sequential_first_countable}
    Assume $T$ is a topology on $X$.
    Assume $T_1$ is a topology on $Y$.
    Suppose $f$ is a function from $X$ to $Y$.
    Suppose $f$ is continuous from $X$ to $Y$ with respect to $T$ and $T_1$.
    Then for all $s$ such that $s$ is a sequence in $X$ we have
        for all $x\in X$ if $s$ converges to $x$ in $X$ with respect to $T$ then
        $f\circ s$ converges to $f(x)$ in $Y$ with respect to $T_1$.
\end{theorem}
\begin{proof}
    Follows by \cref{continuous_implies_sequential}.
\end{proof}

% from §30 Item 4: sequential continuity implies continuity in a first-countable space
\begin{theorem}\label{sequential_implies_continuous_first_countable}
    Assume $T$ is a topology on $X$.
    Assume $T_1$ is a topology on $Y$.
    Suppose $f$ is a function from $X$ to $Y$.
    Suppose $X$ satisfies the first countability axiom with respect to $T$.
    Suppose for all $s$ such that $s$ is a sequence in $X$ we have
        for all $x\in X$ if $s$ converges to $x$ in $X$ with respect to $T$ then
        $f\circ s$ converges to $f(x)$ in $Y$ with respect to $T_1$.
    Then $f$ is continuous from $X$ to $Y$ with respect to $T$ and $T_1$.
\end{theorem}
\begin{proof}
    Show that for all $A$ such that $A\subseteq X$ we have
        $\img{f}{\closure{A}{X}{T}}\subseteq \closure{\img{f}{A}}{Y}{T_1}$.
    \begin{subproof}
        Fix $A$ such that $A\subseteq X$.
        Show that for all $y$ such that $y\in\img{f}{\closure{A}{X}{T}}$ we have
            $y\in\closure{\img{f}{A}}{Y}{T_1}$.
        \begin{subproof}
            Fix $y$ such that $y\in\img{f}{\closure{A}{X}{T}}$.
            Then $f\in\funs{X}{Y}$ and $f$ is a function by \cref{funs,funs_is_function}.
            Take $x\in\dom{f}\inter\closure{A}{X}{T}$ such that $y = f(x)$ by \cref{img_of_function_elim}.
            Hence $x\in\closure{A}{X}{T}$ by \cref{inter_elim_right}.
            Therefore $x\in X$ by \cref{closure}.
            Take $s$ such that $s$ is a sequence in $A$ and $s$ converges to $x$ in $X$ with respect to $T$ by
                \cref{closure_implies_sequence_first_countable}.
            Hence $s$ is a sequence in $X$ by \cref{sequence_in,funs_weaken_codom}.
            Therefore $f\circ s$ converges to $f(x)$ in $Y$ with respect to $T_1$ by assumption.
            Hence $f\circ s$ converges to $y$ in $Y$ with respect to $T_1$.
            Then $f\circ s\in\funs{\naturalsPlus}{Y}$ by \cref{funs_circ,sequence_in}.
            For all $n\in\dom{(f\circ s)}$ we have $(f\circ s)(n)\in\img{f}{A}$
                by \cref{funs,sequence_in,subseteq,inter_intro,img_of_function_intro,funs_elim,function_to_ran,subseteq_transitive,circ_apply}.
            Hence $f\circ s$ is a function to $\img{f}{A}$.
            $\dom{(f\circ s)} = \naturalsPlus$ by \cref{funs}.
            Therefore $f\circ s$ is a sequence in $\img{f}{A}$ by \cref{funs_intro,sequence_in}.
            Hence $\img{f}{A}\subseteq Y$ and $f(x)\in Y$ by
                \cref{img_subseteq_ran,funs_elim,function_to_ran,subseteq_transitive,funs_type_apply}.
            Therefore $f(x)\in\closure{\img{f}{A}}{Y}{T_1}$ by \cref{sequence_implies_closure}.
            Hence $y\in\closure{\img{f}{A}}{Y}{T_1}$.
        \end{subproof}
        Therefore $\img{f}{\closure{A}{X}{T}}\subseteq \closure{\img{f}{A}}{Y}{T_1}$ by \cref{subseteq}.
    \end{subproof}
    Hence for all $B$ such that $B$ is closed in $Y$ with respect to $T_1$
        we have $\preimg{f}{B}$ is closed in $X$ with respect to $T$ by
        \cref{closure_image_subset_implies_preimg_closed}.
    Therefore $f$ is continuous from $X$ to $Y$ with respect to $T$ and $T_1$ by
        \cref{preimg_closed_implies_continuous}.
\end{proof}
% from §30 Item 5: second countability axiom
\begin{definition}\label{second_countability_axiom}
    $X$ satisfies the second countability axiom with respect to $T$ iff
        there exists $b$ such that
            $b$ is a sequence in $\pow{X}$ and
            for all $n\in\naturalsPlus$ we have $b(n)$ is open in $X$ with respect to $T$ and
            for all $x\in X$ we have
                for all $U$ such that $U$ is open in $X$ with respect to $T$ and $x\in U$
                    there exists $n\in\naturalsPlus$ such that $x\in b(n)$ and $b(n)\subseteq U$.
\end{definition}

% from §30 Item 6: second countability implies first countability
\begin{lemma}\label{second_countability_implies_first}
    Assume $T$ is a topology on $X$.
    Suppose $X$ satisfies the second countability axiom with respect to $T$.
    Then $X$ satisfies the first countability axiom with respect to $T$.
\end{lemma}
\begin{proof}
    Take $b$ such that
        $b$ is a sequence in $\pow{X}$ and
        for all $n\in\naturalsPlus$ we have $b(n)$ is open in $X$ with respect to $T$ and
        for all $x\in X$ we have
            for all $U$ such that $U$ is open in $X$ with respect to $T$ and $x\in U$
                there exists $n\in\naturalsPlus$ such that $x\in b(n)$ and $b(n)\subseteq U$
        by \cref{second_countability_axiom}.
    Show that for all $x\in X$ there exists $c$ such that
        $c$ is a countable basis at $x$ in $X$ with respect to $T$.
    \begin{subproof}
        Fix $x\in X$.
        Let $R = \{(n,y) \mid n\in\naturalsPlus \mid
            (x\in b(n)\land y = b(n))\lor(x\notin b(n)\land y = X)\}$.
        Hence $R$ is a relation by \cref{pair_eq_iff,relation}.
        Show that for all $n\in\naturalsPlus$ there exists $y$ such that $n\mathrel{R} y$.
        \begin{subproof}
            Fix $n\in\naturalsPlus$.
            \begin{byCase}
                \caseOf{$x\in b(n)$.}
                    Let $y = b(n)$.
                    Then $(n,y)\in R$ and $n\mathrel{R} y$ by \cref{pair_eq_iff}.
                \caseOf{$x\notin b(n)$.}
                    Let $y = X$.
                    Then $(n,y)\in R$ and $n\mathrel{R} y$ by \cref{pair_eq_iff}.
            \end{byCase}
        \end{subproof}
        Take $c$ such that $c$ is a function on $\naturalsPlus$ and for all $n\in\naturalsPlus$ we have
            $n\mathrel{R} c(n)$ by \cref{axiom_of_choice}.
        Show that $c$ is a countable basis at $x$ in $X$ with respect to $T$.
        \begin{subproof}
            Show that for all $n\in\dom{c}$ we have $c(n)\in\pow{X}$.
            \begin{subproof}
                Fix $n\in\dom{c}$.
                Then $(x\in b(n)\land c(n) = b(n))\lor(x\notin b(n)\land c(n) = X)$
                    by \cref{subseteq,choice_function,pair_eq_iff}.
                \begin{byCase}
                    \caseOf{$x\in b(n)\land c(n) = b(n)$.}
                        Hence $c(n)\in\pow{X}$ by \cref{sequence_in,funs_type_apply}.
                    \caseOf{$x\notin b(n)\land c(n) = X$.}
                        Therefore $c(n)\in\pow{X}$ by \cref{subseteq_refl,powerset_intro}.
                \end{byCase}
            \end{subproof}
            Hence $c$ is a sequence in $\pow{X}$ by \cref{funs_intro,sequence_in}.
            Show that for all $n\in\naturalsPlus$ we have $c(n)$ is a neighborhood of $x$ in $X$ with respect to $T$.
            \begin{subproof}
                Fix $n\in\naturalsPlus$.
                Then $(x\in b(n)\land c(n) = b(n))\lor(x\notin b(n)\land c(n) = X)$
                    by \cref{choice_function,pair_eq_iff}.
                \begin{byCase}
                    \caseOf{$x\in b(n)\land c(n) = b(n)$.}
                        Hence $c(n)$ is a neighborhood of $x$ in $X$ with respect to $T$ by \cref{neighborhood}.
                    \caseOf{$x\notin b(n)\land c(n) = X$.}
                        Therefore $c(n)$ is a neighborhood of $x$ in $X$ with respect to $T$ by \cref{topology_on,open_in,neighborhood}.
                \end{byCase}
            \end{subproof}
            For all $U$ such that $U$ is a neighborhood of $x$ in $X$ with respect to $T$
                we have $U$ is open in $X$ with respect to $T$ and $x\in U$ by \cref{neighborhood}.
            Thus for all $U$ such that $U$ is a neighborhood of $x$ in $X$ with respect to $T$
                there exists $n\in\naturalsPlus$ such that $x\in b(n)$ and $b(n)\subseteq U$.
            Thus for all $U$ such that $U$ is a neighborhood of $x$ in $X$ with respect to $T$
                there exists $n\in\naturalsPlus$ such that $c(n)\subseteq U$
                by \cref{choice_function,pair_eq_iff,subseteq}.
            Hence $c$ is a countable basis at $x$ in $X$ with respect to $T$ by \cref{countable_basis_at_point}.
        \end{subproof}
    \end{subproof}
    Therefore $X$ satisfies the first countability axiom with respect to $T$ by \cref{first_countability_axiom}.
\end{proof}
% from §30 Item 7: subspace of a first-countable space is first-countable
\begin{lemma}\label{subspace_first_countable}
    Assume $T$ is a topology on $X$.
    Suppose $A\subseteq X$.
    Suppose $X$ satisfies the first countability axiom with respect to $T$.
    Then $A$ satisfies the first countability axiom with respect to $\subspaceTopology{T}{A}$.
\end{lemma}
\begin{proof}
    Let $T_1 = \subspaceTopology{T}{A}$.
    $T_1$ is a topology on $A$ by \cref{subspace_topology_is_topology}.
    Show that for all $x\in A$ there exists $c$ such that
        $c$ is a countable basis at $x$ in $A$ with respect to $T_1$.
    \begin{subproof}
        Fix $x\in A$.
        Hence $x\in X$ by \cref{subseteq}.
        Take $b$ such that $b$ is a countable basis at $x$ in $X$ with respect to $T$
            by \cref{first_countability_axiom}.
        Let $c = \{(n, A\inter b(n)) \mid n\in\naturalsPlus\}$.
        Show that $c$ is a countable basis at $x$ in $A$ with respect to $T_1$.
        \begin{subproof}
            For all $n,U_1,U_2$ such that $(n,U_1)\in c$ and $(n,U_2)\in c$ we have $U_1 = U_2$
                by \cref{pair_eq_iff}.
            Hence $c$ is right-unique by \cref{rightunique}.
            Therefore $c$ is a function.
            For all $n\in\dom{c}$ we have $c(n)\in\pow{A}$
                by \cref{dom_iff,pair_eq_iff,inter_elim_left,subseteq,powerset_intro,function_apply_iff}.
            For all $n$ such that $n\in\naturalsPlus$ we have $n\in\dom{c}$
                by \cref{pair_eq_iff,dom_intro}.
            For all $n$ such that $n\in\dom{c}$ we have $n\in\naturalsPlus$ by \cref{dom_iff,pair_eq_iff}.
            Hence $\dom{c} = \naturalsPlus$ by \cref{subseteq,subseteq_antisymmetric}.
            Therefore $c$ is a sequence in $\pow{A}$ by \cref{funs_intro,sequence_in}.
            For all $n\in\naturalsPlus$ we have $c(n)$ is a neighborhood of $x$ in $A$ with respect to $T_1$
                by \cref{countable_basis_at_point,neighborhood,open_in,subspace_topology,inter_intro,sequence_in,funs_elim,subseteq,pair_eq_iff,function_apply_iff}.
            Show that for all $U$ such that $U$ is a neighborhood of $x$ in $A$ with respect to $T_1$
                there exists $n\in\naturalsPlus$ such that $c(n)\subseteq U$.
            \begin{subproof}
                Fix $U$ such that $U$ is a neighborhood of $x$ in $A$ with respect to $T_1$.
                Then $U\in T_1$ and $x\in U$ by \cref{neighborhood,open_in}.
                Take $V\in T$ such that $U = A\inter V$ by \cref{subspace_topology}.
                Therefore $V$ is a neighborhood of $x$ in $X$ with respect to $T$ by \cref{inter_elim_right,open_in,neighborhood}.
                $1\in\naturalsPlus$ by \cref{real_one_in_naturals}.
                Now for all $U$ such that $U$ is a neighborhood of $x$ in $X$ with respect to $T$
                    there exists $n\in\naturalsPlus$ such that $b(n)\subseteq U$
                    by \cref{countable_basis_at_point}.
                Hence there exists $n\in\naturalsPlus$ such that $b(n)\subseteq V$.
                Take $n\in\naturalsPlus$ such that $b(n)\subseteq V$.
                For all $z$ such that $z\in A\inter b(n)$ we have $z\in A\inter V$
                    by \cref{inter_elim_left,inter_elim_right,subseteq,inter_intro}.
                Therefore $A\inter b(n)\subseteq A\inter V$ by \cref{subseteq}.
                Hence $c(n)\subseteq A\inter V$ by \cref{sequence_in,funs_elim,subseteq,pair_eq_iff,function_apply_iff}.
                Therefore $c(n)\subseteq U$ by \cref{subseteq}.
            \end{subproof}
            Hence $c$ is a countable basis at $x$ in $A$ with respect to $T_1$ by \cref{countable_basis_at_point}.
        \end{subproof}
    \end{subproof}
    Therefore $A$ satisfies the first countability axiom with respect to $T_1$ by \cref{first_countability_axiom}.
\end{proof}
% from §30 helper lemma 2: basis elements are neighborhoods
\begin{lemma}\label{countable_basis_elem_neighborhood}
    Assume $b$ is a countable basis at $x$ in $X$ with respect to $T$.
    Assume $n\in\naturalsPlus$.
    Then $b(n)$ is a neighborhood of $x$ in $X$ with respect to $T$.
\end{lemma}
\begin{proof}
    Thus $b(n)$ is a neighborhood of $x$ in $X$ with respect to $T$ by \cref{countable_basis_at_point}.
\end{proof}
% from §30 Item 8: countable product of first-countable spaces is first-countable
\begin{lemma}\label{countable_product_first_countable}
    Assume $X$ is a function.
    Assume $T$ is a function.
    Assume $\dom{X} = \naturalsPlus$.
    Assume $\dom{T} = \naturalsPlus$.
    Assume that for all $n\in\naturalsPlus$ we have $\apply{T}{n}$ is a topology on $\apply{X}{n}$.
    Suppose for all $n\in\naturalsPlus$ we have $\apply{X}{n}$ satisfies the first countability axiom with respect to $\apply{T}{n}$.
    Then $\productFamily{X}$ satisfies the first countability axiom with respect to $\productTopologyIndexed{T}{X}$.
\end{lemma}
\begin{proof}
    Let $T_0 = \productTopologyIndexed{T}{X}$.
    Show that for all $x\in\productFamily{X}$ there exists $d$ such that
        $d$ is a countable basis at $x$ in $\productFamily{X}$ with respect to $T_0$.
    \begin{subproof}
        Fix $x\in\productFamily{X}$.
        Let $R = \{(n,b) \mid n\in\naturalsPlus \mid
            \text{$b$ is a countable basis at $\apply{x}{n}$ in $\apply{X}{n}$ with respect to $\apply{T}{n}$}\}$.
        Hence $R$ is a relation by \cref{pair_eq_iff,relation}.
        Show that for all $n\in\naturalsPlus$ there exists $b$ such that $n\mathrel{R} b$.
        \begin{subproof}
            Fix $n\in\naturalsPlus$.
            Thus $\apply{x}{n}\in\apply{X}{n}$ by \cref{indexed_product,funs_elim,subseteq}.
            $\apply{X}{n}$ satisfies the first countability axiom with respect to $\apply{T}{n}$.
            Take $b$ such that $b$ is a countable basis at $\apply{x}{n}$ in $\apply{X}{n}$ with respect to $\apply{T}{n}$.
            Therefore $(n,b)\in R$ and $n\mathrel{R} b$ by \cref{pair_eq_iff}.
        \end{subproof}
        Take $b$ such that $b$ is a function on $\naturalsPlus$ and for all $n\in\naturalsPlus$ we have
            $n\mathrel{R}\apply{b}{n}$ by \cref{choice_function}.
        Let $R_1 = \{(n,c) \mid n\in\naturalsPlus \mid
            \text{$c$ is a countable basis at $\apply{x}{n}$ in $\apply{X}{n}$ with respect to $\apply{T}{n}$ and
            for all $m,k\in\naturalsPlus$ if $m\leq k$ then $c(k)\subseteq c(m)$}\}$.
        Hence $R_1$ is a relation by \cref{pair_eq_iff,relation}.
        Show that for all $n\in\naturalsPlus$ there exists $c$ such that $n\mathrel{R_1} c$.
        \begin{subproof}
            Fix $n\in\naturalsPlus$.
            Take $c$ such that
                $c$ is a countable basis at $\apply{x}{n}$ in $\apply{X}{n}$ with respect to $\apply{T}{n}$ and
                for all $m,k\in\naturalsPlus$ if $m\leq k$ then $c(k)\subseteq c(m)$
                by \cref{nested_countable_basis,choice_function,pair_eq_iff}.
            Therefore $(n,c)\in R_1$ and $n\mathrel{R_1} c$ by \cref{pair_eq_iff}.
        \end{subproof}
        Take $c$ such that $c$ is a function on $\naturalsPlus$ and for all $n\in\naturalsPlus$ we have
            $n\mathrel{R_1}c(n)$ by \cref{choice_function}.
        For all $n\in\naturalsPlus$ we have $c(n)$ is a countable basis at
            $\apply{x}{n}$ in $\apply{X}{n}$ with respect to $\apply{T}{n}$
            by \cref{choice_function,pair_eq_iff}.
        For all $n\in\naturalsPlus$ we have for all $m,k\in\naturalsPlus$ if $m\leq k$ then
            $\apply{c(n)}{k}\subseteq \apply{c(n)}{m}$ by \cref{choice_function,pair_eq_iff}.

        Let $R_d = \{(n,g) \mid n\in\naturalsPlus \mid
            \text{$g$ is a function on $\initialSegment{n}$ and for all $i\in\initialSegment{n}$ we have
            $g(i) = \preimg{\piIndex{X}{i}}{\apply{c(i)}{n}}$}\}$.
        Hence $R_d$ is a relation by \cref{pair_eq_iff,relation}.
        Show that for all $n\in\naturalsPlus$ there exists $g$ such that $n\mathrel{R_d} g$.
        \begin{subproof}
            Fix $n\in\naturalsPlus$.
            Let $g = \{(i, \preimg{\piIndex{X}{i}}{\apply{c(i)}{n}}) \mid i\in\initialSegment{n}\}$.
            $g$ is a function by \cref{rightunique,relation,pair_eq_iff}.
            For all $i$ such that $i\in\initialSegment{n}$ we have $i\in\dom{g}$
                by \cref{pair_eq_iff,dom_intro}.
            Thus $\initialSegment{n}\subseteq\dom{g}$ by \cref{subseteq}.
            For all $i$ such that $i\in\dom{g}$ we have $i\in\initialSegment{n}$ by \cref{dom_iff,pair_eq_iff}.
            Thus $\dom{g}\subseteq\initialSegment{n}$ by \cref{subseteq}.
            Therefore $\dom{g} = \initialSegment{n}$ by \cref{subseteq_antisymmetric}.
            Thus $g$ is a function on $\initialSegment{n}$.
            For all $i\in\initialSegment{n}$ we have $\apply{g}{i} = \preimg{\piIndex{X}{i}}{\apply{c(i)}{n}}$
                by \cref{pair_eq_iff,dom_intro,function_apply_iff}.
            Thus $(n,g)\in R_d$ and $n\mathrel{R_d} g$ by \cref{pair_eq_iff}.
        \end{subproof}
        Take $h$ such that $h$ is a function on $\naturalsPlus$ and for all $n\in\naturalsPlus$ we have
            $n\mathrel{R_d}\apply{h}{n}$ by \cref{choice_function}.
        Let $d = \{(n, \inters{\img{\apply{h}{n}}{\initialSegment{n}}}) \mid n\in\naturalsPlus\}$.

        Show that $d$ is a countable basis at $x$ in $\productFamily{X}$ with respect to $T_0$.
        \begin{subproof}
            $d$ is a function by \cref{rightunique,relation,pair_eq_iff}.
            Show that for all $n\in\dom{d}$ we have $d(n)\in\pow{\productFamily{X}}$.
            \begin{subproof}
                        Fix $n\in\dom{d}$.
                        Take $U$ such that $(n,U)\in d$ by \cref{dom_iff}.
                        Thus $U = \inters{\img{\apply{h}{n}}{\initialSegment{n}}}$ by \cref{pair_eq_iff}.
                        Thus $n\in\naturalsPlus$ by \cref{pair_eq_iff}.
                        Let $F = \img{\apply{h}{n}}{\initialSegment{n}}$.
                        Show that for all $A$ such that $A\in F$ we have $A\in\productSubbasis{T}{X}$.
                        \begin{subproof}
                            Fix $A$ such that $A\in F$.
                            Take $i$ such that $i\in\initialSegment{n}$ and $(i,A)\in\apply{h}{n}$ by \cref{img_iff}.
                            Thus $A = \preimg{\piIndex{X}{i}}{\apply{c(i)}{n}}$
                                by \cref{choice_function,pair_eq_iff,subseteq,function_apply_iff}.
                            Thus $i\in\dom{X}$ by \cref{initial_segment_naturalsplus,subseteq}.
                            $c(i)$ is a countable basis at $\apply{x}{i}$ in $\apply{X}{i}$ with respect to $\apply{T}{i}$.
                            Thus $\apply{c(i)}{n}\in\apply{T}{i}$ by \cref{countable_basis_at_point,neighborhood,open_in}.
                            Thus $A\in\pow{\productFamily{X}}$ by
                                \cref{preimg_iff,projection_map_indexed,pair_eq_iff,subseteq,powerset_intro}.
                            Thus $A\in\productSubbasisAt{T}{X}{i}$ by \cref{product_subbasis_indexed}.
                            Thus $A\in\productSubbasis{T}{X}$ by \cref{product_subbasis_union_indexed}.
                        \end{subproof}
                        Hence $F\subseteq\productSubbasis{T}{X}$ by \cref{subseteq}.
                        Thus $F$ is finite by \cref{initial_segment_finite,choice_function,pair_eq_iff,finite_image_function}.
                        $n\in\initialSegment{n}$ by \cref{initial_segment_naturalsplus}.
                        Thus $\apply{h}{n}$ is a function and $\dom{\apply{h}{n}} = \initialSegment{n}$ by \cref{pair_eq_iff}.
                        Thus $n\in\dom{\apply{h}{n}}$ by \cref{subseteq}.
                        Thus for all $i\in\initialSegment{n}$ we have
                            $\apply{\apply{h}{n}}{i} = \preimg{\piIndex{X}{i}}{\apply{c(i)}{n}}$ by \cref{pair_eq_iff}.
                        Thus $(n, \preimg{\piIndex{X}{n}}{\apply{c(n)}{n}})\in\apply{h}{n}$ by \cref{pair_eq_iff,function_apply_iff}.
                        Therefore $F$ is inhabited by \cref{img_iff,nonempty_inhabited}.
                        Take $A$ such that $A\in F$.
                        Thus $A\in\productSubbasis{T}{X}$ by \cref{subseteq}.
                        For all $z$ such that $z\in A$ we have $z\in\unions{\productSubbasis{T}{X}}$ by \cref{unions_intro}.
                        Thus $A\subseteq\unions{\productSubbasis{T}{X}}$ by \cref{subseteq}.
                        Thus $U\subseteq A$ by \cref{inters_subseteq_elem}.
                        Thus $U\subseteq\unions{\productSubbasis{T}{X}}$ by \cref{subseteq_transitive}.
                        $\unions{\productSubbasis{T}{X}} = \productFamily{X}$ by \cref{product_subbasis_is_subbasis,subbasis_on}.
                        Thus $U\subseteq\productFamily{X}$ by \cref{subseteq}.
                        Thus $U\in\pow{\productFamily{X}}$ by \cref{powerset_intro}.
                        Thus $d(n)\in\pow{\productFamily{X}}$ by \cref{function_apply_iff}.
            \end{subproof}
            Thus $\dom{d} = \naturalsPlus$ by \cref{subseteq,pair_eq_iff,dom_intro,dom_iff,subseteq_antisymmetric}.
            Thus $d$ is a sequence in $\pow{\productFamily{X}}$ by \cref{funs_intro,sequence_in}.

            Show that for all $n\in\naturalsPlus$ we have $d(n)$ is a neighborhood of $x$ in $\productFamily{X}$ with respect to $T_0$.
            \begin{subproof}
                Fix $n\in\naturalsPlus$.
                Let $F = \img{\apply{h}{n}}{\initialSegment{n}}$.
                Show that for all $A$ such that $A\in F$ we have $x\in A$.
                \begin{subproof}
                    Fix $A$ such that $A\in F$.
                    Take $i$ such that $i\in\initialSegment{n}$ and $(i,A)\in\apply{h}{n}$ by \cref{img_iff}.
                    Thus $A = \preimg{\piIndex{X}{i}}{\apply{c(i)}{n}}$
                        by \cref{choice_function,pair_eq_iff,subseteq,function_apply_iff}.
                    $c(i)$ is a countable basis at $\apply{x}{i}$ in $\apply{X}{i}$ with respect to $\apply{T}{i}$.
                    Thus $\apply{c(i)}{n}$ is a neighborhood of $\apply{x}{i}$ in $\apply{X}{i}$ with respect to $\apply{T}{i}$
                        by \cref{countable_basis_elem_neighborhood}.
                    Thus $\apply{x}{i}\in\apply{c(i)}{n}$ by \cref{neighborhood}.
                    Thus $x\in A$ by \cref{preimg_iff,projection_map_indexed}.
                \end{subproof}
                Show that $F$ is inhabited.
                \begin{subproof}
                    $n\in\initialSegment{n}$ by \cref{initial_segment_naturalsplus}.
                    Thus $n\mathrel{R_d}\apply{h}{n}$ by \cref{choice_function}.
                    Thus $\apply{h}{n}$ is a function and $\dom{\apply{h}{n}} = \initialSegment{n}$ by \cref{pair_eq_iff}.
                    Thus $n\in\dom{\apply{h}{n}}$ by \cref{subseteq}.
                    Thus for all $j\in\initialSegment{n}$ we have
                        $\apply{\apply{h}{n}}{j} = \preimg{\piIndex{X}{j}}{\apply{c(j)}{n}}$ by \cref{pair_eq_iff}.
                    Thus $(n, \preimg{\piIndex{X}{n}}{\apply{c(n)}{n}})\in\apply{h}{n}$ by \cref{pair_eq_iff,function_apply_iff}.
                    Hence $\preimg{\piIndex{X}{n}}{\apply{c(n)}{n}}\in F$ by \cref{img_iff}.
                \end{subproof}
                Thus $x\in\inters{F}$ by \cref{inters_intro}.
                Show that for all $A$ such that $A\in F$ we have $A\in\productSubbasis{T}{X}$.
                \begin{subproof}
                    Fix $A$ such that $A\in F$.
                    Take $i$ such that $i\in\initialSegment{n}$ and $(i,A)\in\apply{h}{n}$ by \cref{img_iff}.
                    Thus $A = \preimg{\piIndex{X}{i}}{\apply{c(i)}{n}}$
                        by \cref{choice_function,pair_eq_iff,subseteq,function_apply_iff}.
                    Thus $i\in\dom{X}$ by \cref{initial_segment_naturalsplus,subseteq}.
                    $c(i)$ is a countable basis at $\apply{x}{i}$ in $\apply{X}{i}$ with respect to $\apply{T}{i}$.
                    Thus $\apply{c(i)}{n}\in\apply{T}{i}$ by \cref{countable_basis_elem_neighborhood,neighborhood,open_in}.
                    Thus $A\in\pow{\productFamily{X}}$ by
                        \cref{preimg_iff,projection_map_indexed,pair_eq_iff,subseteq,powerset_intro}.
                    Thus $A\in\productSubbasisAt{T}{X}{i}$ by \cref{product_subbasis_indexed}.
                    Thus $A\in\productSubbasis{T}{X}$ by \cref{product_subbasis_union_indexed}.
                \end{subproof}
                Hence $F\subseteq\productSubbasis{T}{X}$ by \cref{subseteq}.
                Thus $F$ is finite by \cref{initial_segment_finite,choice_function,pair_eq_iff,finite_image_function}.
                $n\in\initialSegment{n}$ by \cref{initial_segment_naturalsplus}.
                Thus $\apply{h}{n}$ is a function and $\dom{\apply{h}{n}} = \initialSegment{n}$ by \cref{pair_eq_iff}.
                Thus $n\in\dom{\apply{h}{n}}$ by \cref{subseteq}.
                Thus for all $j\in\initialSegment{n}$ we have
                    $\apply{\apply{h}{n}}{j} = \preimg{\piIndex{X}{j}}{\apply{c(j)}{n}}$ by \cref{pair_eq_iff}.
                Thus $(n, \preimg{\piIndex{X}{n}}{\apply{c(n)}{n}})\in\apply{h}{n}$ by \cref{pair_eq_iff,function_apply_iff}.
                Therefore $F$ is inhabited by \cref{img_iff,nonempty_inhabited}.
                Take $A$ such that $A\in F$.
                Thus $A\in\productSubbasis{T}{X}$ by \cref{subseteq}.
                For all $z$ such that $z\in A$ we have $z\in\unions{\productSubbasis{T}{X}}$ by \cref{unions_intro}.
                Thus $A\subseteq\unions{\productSubbasis{T}{X}}$ by \cref{subseteq}.
                Thus $\inters{F}\subseteq A$ by \cref{inters_subseteq_elem}.
                Thus $\inters{F}\subseteq\unions{\productSubbasis{T}{X}}$ by \cref{subseteq_transitive}.
                Thus $\inters{F}\in\pow{\unions{\productSubbasis{T}{X}}}$ by \cref{powerset_intro}.
                Thus $\inters{F}\in\finiteIntersections{\productSubbasis{T}{X}}$ by \cref{finite_intersections}.
                $\finiteIntersections{\productSubbasis{T}{X}}$ is a basis for $T_0$ on $\productFamily{X}$ by
                    \cref{product_subbasis_finite_intersections_basis,product_topology_indexed,basis_for_topology,topology_generated_by_subbasis}.
                Thus $T_0 = \basisTopology{\finiteIntersections{\productSubbasis{T}{X}}}{\productFamily{X}}$ by \cref{basis_for_topology}.
                Thus $\inters{F}\in\basisTopology{\finiteIntersections{\productSubbasis{T}{X}}}{\productFamily{X}}$
                    by \cref{basis_elem_open_in_generated,basis_for_topology}.
                Thus $\inters{F}\in T_0$ by \cref{basis_for_topology}.
                $\basisTopology{\finiteIntersections{\productSubbasis{T}{X}}}{\productFamily{X}}$ is a topology on $\productFamily{X}$
                    by \cref{basis_topology_is_topology,basis_for_topology}.
                Thus $T_0$ is a topology on $\productFamily{X}$ by \cref{basis_for_topology}.
                Thus $\inters{F}$ is open in $\productFamily{X}$ with respect to $T_0$ by \cref{open_in}.
                Thus $\inters{F}$ is a neighborhood of $x$ in $\productFamily{X}$ with respect to $T_0$ by \cref{neighborhood}.
                Thus $n\in\dom{d}$ by \cref{sequence_in,funs_elim,subseteq}.
                $(n, \inters{F})\in d$ by \cref{pair_eq_iff}.
                Thus $d(n) = \inters{F}$ by \cref{function_apply_iff}.
                Thus $d(n)$ is a neighborhood of $x$ in $\productFamily{X}$ with respect to $T_0$.
            \end{subproof}

            Show that for all $U$ such that $U$ is a neighborhood of $x$ in $\productFamily{X}$ with respect to $T_0$
                there exists $n\in\naturalsPlus$ such that $d(n)\subseteq U$.
            \begin{subproof}
                Fix $U$ such that $U$ is a neighborhood of $x$ in $\productFamily{X}$ with respect to $T_0$.
                Thus $U$ is open in $\productFamily{X}$ with respect to $T_0$ and $x\in U$ by \cref{neighborhood}.
                $X$ is a function.
                $T$ is a function.
                $\dom{T} = \dom{X}$.
                $1\in\naturalsPlus$ by \cref{real_one_in_naturals}.
                $\dom{X} = \naturalsPlus$.
                Thus $1\in\dom{X}$ by \cref{subseteq}.
                Thus $\dom{X}$ is inhabited.
                Thus for all $n\in\dom{X}$ we have $\apply{T}{n}$ is a topology on $\apply{X}{n}$.
                Thus $\finiteIntersections{\productSubbasis{T}{X}}$ is a basis for $\productTopologyIndexed{T}{X}$ on $\productFamily{X}$
                    by \cref{product_subbasis_finite_intersections_basis}.
                $T_0 = \productTopologyIndexed{T}{X}$.
                Thus $\finiteIntersections{\productSubbasis{T}{X}}$ is a basis for $T_0$ on $\productFamily{X}$.
                Thus $U\in T_0$ by \cref{open_in}.
                Thus $T_0 = \basisTopology{\finiteIntersections{\productSubbasis{T}{X}}}{\productFamily{X}}$
                    by \cref{basis_for_topology}.
                Thus $U\in\basisTopology{\finiteIntersections{\productSubbasis{T}{X}}}{\productFamily{X}}$.
                Take $B\in\finiteIntersections{\productSubbasis{T}{X}}$ such that $x\in B$ and $B\subseteq U$
                    by \cref{topology_generated_by_basis}.
                Take $F\subseteq\productSubbasis{T}{X}$ such that $F$ is finite and inhabited and $B = \inters{F}$ by \cref{finite_intersections}.

                Let $R_2 = \{(A,i) \mid A\in F \mid
                    \text{$i\in\dom{X}$ and there exists $V\in\apply{T}{i}$ such that
                    $A = \preimg{\piIndex{X}{i}}{V}$}\}$.
                Thus $R_2$ is a relation by \cref{pair_eq_iff,relation}.
                Show that for all $A\in F$ there exists $i$ such that $A\mathrel{R_2} i$.
                \begin{subproof}
                    Fix $A\in F$.
                    Take $i$ such that $i\in\dom{X}$ and $A\in\productSubbasisAt{T}{X}{i}$
                        by \cref{subseteq,product_subbasis_union_indexed}.
                    Take $V\in\apply{T}{i}$ such that $A = \preimg{\piIndex{X}{i}}{V}$ by \cref{product_subbasis_indexed}.
                    Thus $(A,i)\in R_2$ and $A\mathrel{R_2} i$ by \cref{pair_eq_iff}.
                \end{subproof}
                Take $p$ such that $p$ is a function on $F$ and for all $A\in F$ we have $A\mathrel{R_2}\apply{p}{A}$
                    by \cref{choice_function}.
                Let $J = \img{p}{F}$.
                Thus $J$ is finite by \cref{choice_function,finite_image_function}.
                Take $A$ such that $A\in F$.
                Thus $J$ is inhabited by \cref{choice_function,subseteq,function_apply_iff,img_iff,nonempty_inhabited}.

                Let $R_3 = \{(A,k) \mid A\in F \mid
                    \text{there exists $V\in\apply{T}{\apply{p}{A}}$ such that
                    $A = \preimg{\piIndex{X}{\apply{p}{A}}}{V}$ and
                    $k\in\naturalsPlus$ and $\apply{c(\apply{p}{A})}{k}\subseteq V$}\}$.
                Thus $R_3$ is a relation by \cref{pair_eq_iff,relation}.
                Show that for all $A\in F$ there exists $k$ such that $A\mathrel{R_3} k$.
                \begin{subproof}
                    Fix $A\in F$.
                    Thus $\apply{p}{A}\in\dom{X}$ and there exists $V\in\apply{T}{\apply{p}{A}}$ such that
                        $A = \preimg{\piIndex{X}{\apply{p}{A}}}{V}$ by \cref{choice_function,pair_eq_iff}.
                    Take $V\in\apply{T}{\apply{p}{A}}$ such that $A = \preimg{\piIndex{X}{\apply{p}{A}}}{V}$.
                    Thus $B\subseteq A$ by \cref{inters_subseteq_elem}.
                    Thus $x\in A$ and $x\in \preimg{\piIndex{X}{\apply{p}{A}}}{V}$ by \cref{subseteq}.
                    Take $y\in V$ such that $(x,y)\in\piIndex{X}{\apply{p}{A}}$ by \cref{preimg_iff}.
                    Take $z\in\productFamily{X}$ such that $(z,\apply{z}{\apply{p}{A}}) = (x,y)$ by \cref{projection_map_indexed}.
                    Thus $\apply{x}{\apply{p}{A}}\in V$ by \cref{pair_eq_iff}.
                    $\dom{X} = \naturalsPlus$.
                    Thus $\apply{p}{A}\in\naturalsPlus$ by \cref{subseteq}.
                    Thus $\apply{T}{\apply{p}{A}}$ is a topology on $\apply{X}{\apply{p}{A}}$.
                    Thus $V$ is a neighborhood of $\apply{x}{\apply{p}{A}}$ in $\apply{X}{\apply{p}{A}}$ with respect to $\apply{T}{\apply{p}{A}}$ by \cref{open_in,neighborhood}.
                    $c(\apply{p}{A})$ is a countable basis at $\apply{x}{\apply{p}{A}}$ in $\apply{X}{\apply{p}{A}}$ with respect to $\apply{T}{\apply{p}{A}}$.
                    Take $k\in\naturalsPlus$ such that $\apply{c(\apply{p}{A})}{k}\subseteq V$ by \cref{countable_basis_at_point}.
                    Thus $(A,k)\in R_3$ and $A\mathrel{R_3} k$ by \cref{pair_eq_iff}.
                \end{subproof}
                Take $q$ such that $q$ is a function on $F$ and for all $A\in F$ we have $A\mathrel{R_3}\apply{q}{A}$
                    by \cref{choice_function}.
                Let $K = \img{q}{F}$.
                Thus $K$ is finite by \cref{choice_function,finite_image_function}.
                Take $A$ such that $A\in F$.
                Thus $K$ is inhabited by \cref{choice_function,subseteq,function_apply_iff,img_iff,nonempty_inhabited}.

                For all $j$ such that $j\in J$ we have $j\in\naturalsPlus$
                    by \cref{img_iff,choice_function,subseteq,function_apply_iff,pair_eq_iff}.
                Thus $J\subseteq\naturalsPlus$ by \cref{subseteq}.

                For all $k$ such that $k\in K$ we have $k\in\naturalsPlus$
                    by \cref{img_iff,choice_function,subseteq,function_apply_iff,pair_eq_iff}.
                Thus $K\subseteq\naturalsPlus$ by \cref{subseteq}.

                Let $G = J\union K$.
                Thus $G\subseteq\naturalsPlus$ by \cref{union_subsets_is_subset}.
                $\naturalsPlus\subseteq\reals$ by \cref{real_naturals_subset_reals}.
                Thus $G\subseteq\reals$ by \cref{subseteq_transitive}.
                $G$ is finite by \cref{union_of_finite_is_finite}.
                $G$ is inhabited.
                Take $m$ such that $m$ is a largest element of $G$ with respect to $\realOrderRel$
                    by \cref{finite_inhabited_largest_element,real_order_relation_simple}.
                Thus $m\in G$ by \cref{largest_element}.
                Thus $m\in\naturalsPlus$ by \cref{subseteq}.

                Show that for all $A\in F$ we have $d(m)\subseteq A$.
                \begin{subproof}
                    Fix $A\in F$.
                    Thus $\apply{p}{A}\in J$ by \cref{choice_function,subseteq,function_apply_iff,img_iff}.
                    Thus $\apply{p}{A}\in G$ by \cref{union_intro_left}.
                    Thus $\apply{p}{A}\leq m$ by \cref{largest_element,real_order_relation_elim}.
                    Thus $\apply{p}{A}\in\dom{X}$ by \cref{pair_eq_iff}.
                    $\dom{X} = \naturalsPlus$.
                    Thus $\apply{p}{A}\in\naturalsPlus$.
                    Thus $\apply{p}{A}\in\initialSegment{m}$ by \cref{initial_segment_naturalsplus}.

                    Thus $\apply{q}{A}\in K$ by \cref{choice_function,subseteq,function_apply_iff,img_iff}.
                    Thus $\apply{q}{A}\in G$ by \cref{union_intro_right}.
                    Thus $\apply{q}{A}\leq m$ by \cref{largest_element,real_order_relation_elim}.

                    Let $i = \apply{p}{A}$.
                    Thus $i\in\initialSegment{m}$.
                    Thus $A\mathrel{R_2} i$ by \cref{choice_function}.
                    Thus $i\in\dom{X}$ by \cref{pair_eq_iff}.
                    Let $k = \apply{q}{A}$.
                    Thus $A\mathrel{R_3} k$ by \cref{choice_function}.
                    Thus there exists $V\in\apply{T}{i}$ such that
                        $A = \preimg{\piIndex{X}{i}}{V}$ and $k\in\naturalsPlus$ and $\apply{c(i)}{k}\subseteq V$
                        by \cref{pair_eq_iff}.
                    Take $V\in\apply{T}{i}$ such that
                        $A = \preimg{\piIndex{X}{i}}{V}$ and $k\in\naturalsPlus$ and $\apply{c(i)}{k}\subseteq V$.

                    Thus $k\leq m$.
                    Thus $i\mathrel{R_1}\apply{c}{i}$ by \cref{choice_function}.
                    Thus for all $m_1,k_1\in\naturalsPlus$ if $m_1\leq k_1$ then
                        $\apply{c(i)}{k_1}\subseteq \apply{c(i)}{m_1}$ by \cref{pair_eq_iff}.
                    Thus $\apply{c(i)}{m}\subseteq \apply{c(i)}{k}$.
                    Thus $\apply{c(i)}{m}\subseteq V$ by \cref{subseteq}.

                    Let $F_m = \img{\apply{h}{m}}{\initialSegment{m}}$.
                    Thus $m\mathrel{R_d}\apply{h}{m}$ by \cref{choice_function}.
                    Thus $\apply{h}{m}$ is right-unique and $\apply{h}{m}$ is a relation and
                        $\initialSegment{m} = \dom{\apply{h}{m}}$ and
                        for all $j\in\initialSegment{m}$ we have
                        $\apply{\apply{h}{m}}{j} = \preimg{\piIndex{X}{j}}{\apply{c(j)}{m}}$ by \cref{pair_eq_iff}.
                    Thus $\apply{h}{m}$ is a function.
                    Thus $i\in\dom{\apply{h}{m}}$ by \cref{subseteq}.
                    Thus $\preimg{\piIndex{X}{i}}{\apply{c(i)}{m}}\in F_m$ by \cref{pair_eq_iff,function_apply_iff,img_iff}.
                    Thus $m\in\dom{d}$ by \cref{sequence_in,funs_elim,subseteq}.
                    $(m, \inters{F_m})\in d$ by \cref{pair_eq_iff}.
                    Thus $d(m) = \inters{F_m}$ by \cref{function_apply_iff}.
                    Thus $d(m)\subseteq \preimg{\piIndex{X}{i}}{\apply{c(i)}{m}}$ by \cref{inters_subseteq_elem}.
                    Thus $d(m)\subseteq A$ by \cref{subseteq_transitive,preimg_subseteq,subseteq}.
                \end{subproof}
                Therefore $d(m)\subseteq U$ by \cref{subseteq_inters_iff,subseteq,subseteq_transitive}.
            \end{subproof}
            Hence $d$ is a countable basis at $x$ in $\productFamily{X}$ with respect to $T_0$ by \cref{countable_basis_at_point}.
        \end{subproof}
    \end{subproof}
    Therefore $\productFamily{X}$ satisfies the first countability axiom with respect to $T_0$ by \cref{first_countability_axiom}.
\end{proof}
% from §30 Item 9: subspace of a second-countable space is second-countable
\begin{lemma}\label{subspace_second_countable}
    Assume $T$ is a topology on $X$.
    Suppose $A\subseteq X$.
    Suppose $X$ satisfies the second countability axiom with respect to $T$.
    Then $A$ satisfies the second countability axiom with respect to $\subspaceTopology{T}{A}$.
\end{lemma}
\begin{proof}
    Let $T_1 = \subspaceTopology{T}{A}$.
    $T_1$ is a topology on $A$ by \cref{subspace_topology_is_topology}.
    Take $b$ such that
        $b$ is a sequence in $\pow{X}$ and
        for all $n\in\naturalsPlus$ we have $b(n)$ is open in $X$ with respect to $T$ and
        for all $x\in X$ we have
            for all $U$ such that $U$ is open in $X$ with respect to $T$ and $x\in U$
                there exists $n\in\naturalsPlus$ such that $x\in b(n)$ and $b(n)\subseteq U$
        by \cref{second_countability_axiom}.
    Let $c = \{(n, A\inter b(n)) \mid n\in\naturalsPlus\}$.
    For all $n,U_1,U_2$ such that $(n,U_1)\in c$ and $(n,U_2)\in c$ we have $U_1 = U_2$ by \cref{pair_eq_iff}.
    Therefore $c$ is a function by \cref{rightunique,relation,pair_eq_iff}.
    Hence $c\in\funs{\naturalsPlus}{\pow{A}}$
        by \cref{funs_intro,dom_iff,pair_eq_iff,inter_elim_left,subseteq,powerset_intro,function_apply_iff,dom_intro,subseteq_antisymmetric}.
    Hence $c$ is a sequence in $\pow{A}$ by \cref{sequence_in}.
    Thus for all $n\in\naturalsPlus$ we have $c(n)$ is open in $A$ with respect to $T_1$
        by \cref{sequence_in,funs_is_function,funs_elim,subseteq,pair_eq_iff,function_apply_iff,open_in,subspace_topology}.
    Show that for all $x\in A$ we have
        for all $U$ such that $U$ is open in $A$ with respect to $T_1$ and $x\in U$
            there exists $n\in\naturalsPlus$ such that $x\in c(n)$ and $c(n)\subseteq U$.
    \begin{subproof}
        Fix $x\in A$.
        Fix $U$ such that $U$ is open in $A$ with respect to $T_1$ and $x\in U$.
        Take $V\in T$ such that $U = A\inter V$ by \cref{open_in,subspace_topology}.
        Thus $V$ is open in $X$ with respect to $T$ and $x\in V$ by \cref{inter_elim_right,open_in}.
        Take $n\in\naturalsPlus$ such that $x\in b(n)$ and $b(n)\subseteq V$.
        Thus $x\in c(n)$ and $c(n)\subseteq U$ by
            \cref{sequence_in,funs_is_function,funs_elim,subseteq,pair_eq_iff,function_apply_iff,inter_elim_left,inter_elim_right,inter_intro}.
    \end{subproof}
    Therefore $A$ satisfies the second countability axiom with respect to $T_1$ by \cref{second_countability_axiom}.
\end{proof}

% from §30 helper lemma 3: naturalsPlus closed under addition
\begin{lemma}\label{naturalsplus_add_closed}
    Suppose $a\in\naturalsPlus$ and $b\in\naturalsPlus$.
    Then $a + b\in\naturalsPlus$.
\end{lemma}
\begin{proof}
    Let $A = \{b\in\naturalsPlus \mid a + b\in\naturalsPlus\}$.
    $1\in A$ by \cref{real_one_in_naturals,real_naturals_closed_succ}.
    For all $n\in A$ we have $n + 1\in A$
        by \cref{real_naturals_closed_succ,real_naturals_subset_reals,subseteq,real_one_in_naturals,real_add_assoc}.
    Hence $a + b\in\naturalsPlus$ by \cref{subseteq,real_naturals_induction}.
\end{proof}

% from §30 helper lemma 4: naturalsPlus closed under multiplication
\begin{lemma}\label{naturalsplus_mul_closed}
    Suppose $a\in\naturalsPlus$ and $b\in\naturalsPlus$.
    Then $a \rmul b\in\naturalsPlus$.
\end{lemma}
\begin{proof}
    Let $A = \{b\in\naturalsPlus \mid a \rmul b\in\naturalsPlus\}$.
    $1\in A$ by \cref{real_one_in_naturals,real_naturals_subset_reals,subseteq,real_mul_comm,real_mul_ident}.
    For all $n\in A$ we have $n + 1\in A$
        by \cref{naturalsplus_add_closed,real_naturals_subset_reals,subseteq,real_one_in_naturals,real_distrib,real_mul_comm,real_mul_ident}.
    Hence $a \rmul b\in\naturalsPlus$ by \cref{subseteq,real_naturals_induction}.
\end{proof}

% from §30 helper lemma 5: successor product below next square
\begin{lemma}\label{naturalsplus_mul_succ_lt_square}
    Suppose $s\in\naturalsPlus$.
    Then $s \rmul (s + 1) < (s + 1) \rmul (s + 1)$.
\end{lemma}
\begin{proof}
    $s\in\reals$ and $1\in\reals$ and $0\in\reals$ and $0 < 1$ by \cref{real_naturals_subset_reals,real_one_in_naturals,subseteq,real_add_ident,real_zero_lt_one}.
    Hence $s < s + 1$ and $s + 1\in\reals$
        by \cref{real_naturals_closed_succ,real_naturals_subset_reals,subseteq,real_lt_add,real_add_ident,real_add_comm}.
    Therefore $0 < s + 1$ by \cref{real_one_leq_naturalsplus,real_lt_trans,real_lt_subst_right,real_zero_lt_one}.
    Hence $s \rmul (s + 1) < (s + 1) \rmul (s + 1)$ by \cref{real_lt_mul_pos}.
\end{proof}

% from §30 helper lemma 6: bound for square against successor product
\begin{lemma}\label{naturalsplus_mul_succ_lt_square_bound}
    Suppose $s\in\naturalsPlus$ and $t\in\naturalsPlus$ and $s < t$.
    Then $s \rmul (s + 1) < t \rmul t$.
\end{lemma}
\begin{proof}
    $t = 1$ or there exists $m\in\naturalsPlus$ such that $t = m + 1$ by \cref{real_naturals_predecessor}.
    \begin{byCase}
    \caseOf{$t = 1$.}
        Show that $s \rmul (s + 1) < t \rmul t$.
        \begin{subproof}
            Suppose not.
            $1 < s$ or $1 = s$ by \cref{real_one_leq_naturalsplus}.
            Hence $s < 1$.
            $1\in\reals$ and $s\in\reals$ by \cref{real_naturals_subset_reals,real_one_in_naturals,subseteq}.
            Thus $1 < 1$ by \cref{real_lt_trans,real_lt_subst_left}.
            Contradiction by \cref{real_lt_irreflexive}.
        \end{subproof}
    \caseOf{there exists $m\in\naturalsPlus$ such that $t = m + 1$.}
        Take $m\in\naturalsPlus$ such that $t = m + 1$.
        Thus $s < m + 1$.
        Thus $s \leq m + 1$.
        Thus $s \neq m + 1$.
        Thus $s \leq m$ by \cref{real_naturals_leq_succ_elim}.
        $s < m$ or $s = m$.
        \begin{byCase}
        \caseOf{$s = m$.}
            Hence $s \rmul (s + 1) < t \rmul t$ by \cref{naturalsplus_mul_succ_lt_square}.
        \caseOf{$s < m$.}
            $s\in\reals$ and $t\in\reals$ by \cref{real_naturals_subset_reals,subseteq}.
            $s + 1\in\naturalsPlus$ and $s + 1\in\reals$ by \cref{real_naturals_closed_succ,real_naturals_subset_reals,subseteq}.
            $1\in\reals$ and $0\in\reals$ and $0 < 1$ by \cref{real_naturals_subset_reals,real_one_in_naturals,subseteq,real_add_ident,real_zero_lt_one}.
            $1 < s + 1$ or $1 = s + 1$ by \cref{real_one_leq_naturalsplus}.
            Hence $0 < s + 1$ by \cref{real_lt_trans,real_lt_subst_right}.
            Thus $s \rmul (s + 1) < t \rmul (s + 1)$ by \cref{real_lt_mul_pos}.
            $s + 1 < m + 1$ by \cref{real_naturals_subset_reals,real_one_in_naturals,subseteq,real_add_ident,real_zero_lt_one,real_lt_add}.
            Then $s + 1 < t$.
            Therefore $0 < t$ by \cref{real_one_leq_naturalsplus,real_lt_trans,real_lt_subst_right,real_zero_lt_one}.
            Thus $(s + 1) \rmul t < t \rmul t$ by \cref{real_lt_mul_pos}.
            Thus $t \rmul (s + 1) < t \rmul t$ by \cref{real_mul_comm}.
            $s \rmul (s + 1)\in\reals$ and $t \rmul (s + 1)\in\reals$ and $t \rmul t\in\reals$ by \cref{real_mul_closed}.
            Hence $s \rmul (s + 1) < t \rmul t$ by \cref{real_lt_trans}.
        \end{byCase}
    \end{byCase}
\end{proof}
% from §30 helper lemma 7: injective pairing of naturalsPlus
\begin{lemma}\label{naturalsplus_pairing_injection}
    There exists $f\in\inj{\naturalsPlus\times\naturalsPlus}{\naturalsPlus}$.
\end{lemma}
    \begin{proof}
        Let $g = \{((i,j), (i + j) \rmul (i + j) + i) \mid i\in\naturalsPlus \mid j\in\naturalsPlus\}$.
        For all $p,a,b$ such that $(p,a)\in g$ and $(p,b)\in g$ we have $a = b$ by \cref{pair_eq_iff}.
        Thus $g$ is a function by \cref{rightunique,relation,pair_eq_iff}.
            For all $p\in\dom{g}$ we have $g(p)\in\naturalsPlus$ by \cref{dom_iff,pair_eq_iff,naturalsplus_add_closed,naturalsplus_mul_closed,function_apply_iff}.
                Show that for all $p$ such that $p\in\naturalsPlus\times\naturalsPlus$ we have $p\in\dom{g}$.
                \begin{subproof}
                    Fix $p$ such that $p\in\naturalsPlus\times\naturalsPlus$.
                    Take $i,j$ such that $i\in\naturalsPlus$ and $j\in\naturalsPlus$ and $p = (i,j)$ by \cref{times}.
                    Thus $p\in\dom{g}$ by \cref{pair_eq_iff,dom_intro}.
                \end{subproof}
                Thus $\dom{g} = \naturalsPlus\times\naturalsPlus$ by \cref{subseteq,dom_iff,pair_eq_iff,times_tuple_intro,dom_intro,subseteq_antisymmetric}.
        Hence $g\in\funs{\naturalsPlus\times\naturalsPlus}{\naturalsPlus}$ by \cref{funs_intro}.
        Show that for all $p_1,p_2$ such that $p_1\in\naturalsPlus\times\naturalsPlus$ and $p_2\in\naturalsPlus\times\naturalsPlus$ and $g(p_1) = g(p_2)$ we have $p_1 = p_2$.
        \begin{subproof}
            Fix $p_1$.
            Fix $p_2$ such that $p_1\in\naturalsPlus\times\naturalsPlus$ and $p_2\in\naturalsPlus\times\naturalsPlus$ and $g(p_1) = g(p_2)$.
            Thus $p_1\in\dom{g}$ and $p_2\in\dom{g}$ by \cref{funs_elim,subseteq}.
            Thus $(p_1,g(p_1))\in g$ and $(p_2,g(p_2))\in g$ by \cref{function_apply_iff}.
            Take $i_1,j_1\in\naturalsPlus$ such that
                $p_1 = (i_1,j_1)$ and $g(p_1) = (i_1 + j_1) \rmul (i_1 + j_1) + i_1$ by \cref{pair_eq_iff}.
            Take $i_2,j_2\in\naturalsPlus$ such that
                $p_2 = (i_2,j_2)$ and $g(p_2) = (i_2 + j_2) \rmul (i_2 + j_2) + i_2$ by \cref{pair_eq_iff}.
            Thus $(i_1 + j_1) \rmul (i_1 + j_1) + i_1 = (i_2 + j_2) \rmul (i_2 + j_2) + i_2$.
            Let $s = i_1 + j_1$.
            Let $t = i_2 + j_2$.
            Hence $s \rmul s + i_1 = t \rmul t + i_2$.
            \begin{byCase}
            \caseOf{$s = t$.}
            $i_1\in\reals$ and $j_1\in\reals$ and $i_2\in\reals$ and $j_2\in\reals$ and $s\in\reals$ and $t\in\reals$ by \cref{naturalsplus_add_closed,real_naturals_subset_reals,subseteq}.
            $t \rmul t\in\reals$ by \cref{real_mul_closed}.
            Thus $i_1 + t \rmul t = i_2 + t \rmul t$ by \cref{real_add_comm}.
            Hence $i_1 = i_2$ by \cref{real_add_cancel}.
            Thus $i_1 + j_1 = i_2 + j_2$.
            Thus $j_1 = j_2$ by \cref{real_add_comm,real_add_cancel}.
            Thus $p_1 = p_2$ by \cref{pair_eq_iff}.
            \caseOf{$s < t$.}
                Suppose not.
                $j_1\in\naturalsPlus$ and $j_1\in\reals$ by \cref{subseteq,real_naturals_subset_reals}.
                $1\in\reals$ and $0\in\reals$ and $0 < 1$ by \cref{real_naturals_subset_reals,real_one_in_naturals,subseteq,real_add_ident,real_zero_lt_one}.
                $0 + j_1 < 1 + j_1$ by \cref{real_lt_add}.
                $0 + j_1 = j_1$ by \cref{real_add_ident}.
                $1 + j_1 = j_1 + 1$ by \cref{real_add_comm}.
                Hence $0 < j_1$.
                $i_1\in\reals$ by \cref{real_naturals_subset_reals,subseteq}.
                Thus $s\in\naturalsPlus$ and $s\in\reals$ by \cref{naturalsplus_add_closed,real_naturals_subset_reals,subseteq}.
                Thus $i_1 < s$ by \cref{real_lt_add,real_add_ident,real_add_comm}.
                $s \rmul s\in\reals$ by \cref{real_mul_closed}.
                Thus $s \rmul s + i_1 < s \rmul s + s$ by \cref{real_lt_add,real_add_comm}.
                Thus $s \rmul s + i_1 < s \rmul (s + 1)$ by \cref{real_lt_subst_right,real_distrib,real_mul_comm,real_mul_ident}.
                $s \rmul s + i_1\in\reals$ and $s + 1\in\reals$ and $s \rmul (s + 1)\in\reals$ by \cref{real_add_closed,real_mul_closed}.
                Thus $t\in\naturalsPlus$ and $t\in\reals$ and $t \rmul t\in\reals$
                    by \cref{naturalsplus_add_closed,real_naturals_subset_reals,subseteq,real_mul_closed}.
                Thus $s \rmul (s + 1) < t \rmul t$ by \cref{naturalsplus_mul_succ_lt_square_bound}.
                Thus $s \rmul s + i_1 < t \rmul t$ by \cref{real_lt_trans}.
                $i_2\in\reals$ by \cref{real_naturals_subset_reals,subseteq}.
                Thus $0 < i_2$
                    by \cref{real_one_leq_naturalsplus,real_one_in_naturals,real_naturals_subset_reals,subseteq,real_lt_subst_right,real_lt_trans,real_zero_lt_one}.
                Thus $t \rmul t < t \rmul t + i_2$ by \cref{real_add_ident,real_lt_add,real_lt_subst_left,real_add_comm}.
                Thus $s \rmul s + i_1 < t \rmul t + i_2$ by \cref{real_lt_trans}.
                Thus it is wrong that $s \rmul s + i_1 = t \rmul t + i_2$ by \cref{real_lt_irreflexive}.
                Contradiction.
            \caseOf{$t < s$.}
                Suppose not.
                $j_2\in\reals$ and $1\in\reals$ and $0\in\reals$ and $0 < 1$ by \cref{subseteq,real_naturals_subset_reals,real_one_in_naturals,real_add_ident,real_zero_lt_one}.
                $0 + j_2 < 1 + j_2$ by \cref{real_lt_add}.
                $0 + j_2 = j_2$ by \cref{real_add_ident}.
                $1 + j_2 = j_2 + 1$ by \cref{real_add_comm}.
                Hence $0 < j_2$.
                $i_2\in\reals$ by \cref{real_naturals_subset_reals,subseteq}.
                Thus $i_2 < t$ by \cref{real_lt_add,real_add_ident,real_add_comm}.
                Thus $t\in\naturalsPlus$ and $t\in\reals$ and $t \rmul t\in\reals$
                    by \cref{naturalsplus_add_closed,real_naturals_subset_reals,subseteq,real_mul_closed}.
                Thus $t \rmul t + i_2 < t \rmul t + t$ by \cref{real_lt_add,real_add_comm}.
                Thus $t \rmul t + i_2 < t \rmul (t + 1)$ by \cref{real_lt_subst_right,real_distrib,real_mul_comm,real_mul_ident}.
                $t \rmul t + i_2\in\reals$ and $t + 1\in\reals$ and $t \rmul (t + 1)\in\reals$ by \cref{real_add_closed,real_mul_closed}.
                Thus $s\in\naturalsPlus$ and $s\in\reals$ by \cref{naturalsplus_add_closed,real_naturals_subset_reals,subseteq}.
                $s \rmul s\in\reals$ by \cref{real_mul_closed}.
                Thus $t \rmul (t + 1) < s \rmul s$ by \cref{naturalsplus_mul_succ_lt_square_bound}.
                Thus $t \rmul t + i_2 < s \rmul s$ by \cref{real_lt_trans}.
                $i_1\in\reals$ by \cref{subseteq,real_naturals_subset_reals}.
                Thus $0 < i_1$
                    by \cref{real_one_leq_naturalsplus,real_one_in_naturals,real_naturals_subset_reals,subseteq,real_lt_subst_right,real_lt_trans,real_zero_lt_one}.
                Thus $s \rmul s < s \rmul s + i_1$ by \cref{real_add_ident,real_lt_add,real_lt_subst_left,real_add_comm}.
                Thus $t \rmul t + i_2 < s \rmul s + i_1$ by \cref{real_lt_trans}.
                Thus it is wrong that $s \rmul s + i_1 = t \rmul t + i_2$ by \cref{real_lt_irreflexive}.
                Contradiction.
            \end{byCase}
        \end{subproof}
        Hence there exists $f\in\inj{\naturalsPlus\times\naturalsPlus}{\naturalsPlus}$ by \cref{inj}.
    \end{proof}

% from §30 helper lemma 8: function equality by pointwise equality
\begin{lemma}\label{functions_eq_iff}
    Suppose $f\in\funs{X}{Y}$ and $g\in\funs{X}{Y}$.
    Suppose for all $x\in X$ we have $f(x) = g(x)$.
    Then $f = g$.
\end{lemma}
\begin{proof}
    Follows by \cref{funs_elim,fun_subseteq,subseteq_refl,subseteq_antisymmetric}.
\end{proof}

% from §30 helper lemma 8a: restriction of a sequence to a shorter initial segment
\begin{lemma}\label{restrl_initialsegment_funs}
    Suppose $n\in\naturalsPlus$ and $u\in\funs{\initialSegment{n + 1}}{\naturalsPlus}$.
    Then $\restrl{u}{\initialSegment{n}}\in\funs{\initialSegment{n}}{\naturalsPlus}$.
\end{lemma}
\begin{proof}
    Thus $u$ is a function and $\dom{u} = \initialSegment{n + 1}$ by \cref{funs_elim}.
    $\restrl{u}{\initialSegment{n}}$ is a function by \cref{restrl_of_function_is_function}.
    Thus $\dom{\restrl{u}{\initialSegment{n}}} = \initialSegment{n}$
        by \cref{initial_segment_subset_succ,subseteq,elem_dom_and_restr_implies_elem_of_restr,restrl_dom,subseteq_antisymmetric}.
    Hence $\restrl{u}{\initialSegment{n}}\in\funs{\initialSegment{n}}{\naturalsPlus}$
        by \cref{elem_dom_of_restrl_implies_elem_dom_and_restr,funs_elim,initial_segment_subset_succ,subseteq,restrl_of_function_apply,funs_intro}.
\end{proof}

% from §30 helper lemma 9: injective coding of finite sequences of naturalsPlus
\begin{lemma}\label{naturalsplus_finite_sequences_inj}
    Suppose $n\in\naturalsPlus$.
    Then there exists $f\in\inj{\funs{\initialSegment{n}}{\naturalsPlus}}{\naturalsPlus}$.
\end{lemma}
\begin{proof}
    Let $A = \{n\in\naturalsPlus \mid \exists f\in\inj{\funs{\initialSegment{n}}{\naturalsPlus}}{\naturalsPlus}. f = f\}$.
    $1\in\naturalsPlus$ by \cref{real_one_in_naturals}.
    Let $f_1 = \{(u, \apply{u}{1}) \mid u\in\funs{\initialSegment{1}}{\naturalsPlus}\}$.
    $f_1$ is a function by \cref{rightunique,relation,pair_eq_iff}.
    Hence $f_1\in\funs{\funs{\initialSegment{1}}{\naturalsPlus}}{\naturalsPlus}$
        by \cref{funs_intro,subseteq,dom_iff,pair_eq_iff,dom_intro,subseteq_antisymmetric,initial_segment_naturalsplus,real_one_in_naturals,funs_type_apply,function_apply_iff}.
    Show that for all $u_1,u_2$ such that $u_1\in\funs{\initialSegment{1}}{\naturalsPlus}$ and $u_2\in\funs{\initialSegment{1}}{\naturalsPlus}$ and $f_1(u_1) = f_1(u_2)$ we have $u_1 = u_2$.
    \begin{subproof}
                Fix $u_1$.
                Fix $u_2$ such that $u_1\in\funs{\initialSegment{1}}{\naturalsPlus}$ and $u_2\in\funs{\initialSegment{1}}{\naturalsPlus}$ and $f_1(u_1) = f_1(u_2)$.
                For all $k\in\initialSegment{1}$ we have $u_1(k) = u_2(k)$
                    by \cref{pair_eq_iff,funs_elim,function_apply_iff,initial_segment_naturalsplus,real_naturals_subset_reals,real_one_in_naturals,subseteq,real_one_leq_naturalsplus,real_leq_antisym}.
                Thus $u_1 = u_2$ by \cref{functions_eq_iff}.
            \end{subproof}
    Therefore $f_1\in\inj{\funs{\initialSegment{1}}{\naturalsPlus}}{\naturalsPlus}$ by \cref{inj}.
        Hence $1\in A$.
        Show that for all $n\in A$ we have $n + 1\in A$.
        \begin{subproof}
            Fix $n\in A$.
            Take $f\in\inj{\funs{\initialSegment{n}}{\naturalsPlus}}{\naturalsPlus}$ such that $f = f$.
            Take $g\in\inj{\naturalsPlus\times\naturalsPlus}{\naturalsPlus}$ such that $g = g$ by \cref{naturalsplus_pairing_injection}.
            Let $F = \{(u, \apply{g}{(\apply{f}{\restrl{u}{\initialSegment{n}}}, \apply{u}{n + 1})}) \mid u\in\funs{\initialSegment{n + 1}}{\naturalsPlus}\}$.
            $F$ is a function by \cref{rightunique,relation,pair_eq_iff}.
            Show that for all $u\in\dom{F}$ we have $F(u)\in\naturalsPlus$.
            \begin{subproof}
                Fix $u\in\dom{F}$.
                Take $a$ such that $(u,a)\in F$ by \cref{dom_iff}.
                Take $u_0\in\funs{\initialSegment{n + 1}}{\naturalsPlus}$ such that
                    $(u,a) = (u_0,\apply{g}{(\apply{f}{\restrl{u_0}{\initialSegment{n}}}, \apply{u_0}{n + 1})})$ by \cref{pair_eq_iff}.
                Thus $\restrl{u}{\initialSegment{n}}\in\funs{\initialSegment{n}}{\naturalsPlus}$
                    by \cref{pair_eq_iff,funs_elim,funs_intro,restrl_initialsegment_funs}.
                Thus $F(u)\in\naturalsPlus$ by \cref{pair_eq_iff,inj,funs_type_apply,function_apply_iff,real_naturals_closed_succ,initial_segment_naturalsplus,times_tuple_intro}.
            \end{subproof}
            Hence $F\in\funs{\funs{\initialSegment{n + 1}}{\naturalsPlus}}{\naturalsPlus}$
                by \cref{funs_intro,subseteq,dom_iff,pair_eq_iff,dom_intro,subseteq_antisymmetric}.
            Show that for all $u_1,u_2$ such that $u_1\in\funs{\initialSegment{n + 1}}{\naturalsPlus}$ and $u_2\in\funs{\initialSegment{n + 1}}{\naturalsPlus}$ and $F(u_1) = F(u_2)$ we have $u_1 = u_2$.
            \begin{subproof}
                    Fix $u_1$.
                    Fix $u_2$ such that $u_1\in\funs{\initialSegment{n + 1}}{\naturalsPlus}$ and $u_2\in\funs{\initialSegment{n + 1}}{\naturalsPlus}$ and $F(u_1) = F(u_2)$.
                    Thus $\restrl{u_1}{\initialSegment{n}}\in\funs{\initialSegment{n}}{\naturalsPlus}$ and
                        $\restrl{u_2}{\initialSegment{n}}\in\funs{\initialSegment{n}}{\naturalsPlus}$
                        by \cref{funs_elim,restrl_initialsegment_funs}.
                    $(u_1,\apply{g}{(\apply{f}{\restrl{u_1}{\initialSegment{n}}}, \apply{u_1}{n + 1})})\in F$ and
                        $(u_2,\apply{g}{(\apply{f}{\restrl{u_2}{\initialSegment{n}}}, \apply{u_2}{n + 1})})\in F$ by \cref{pair_eq_iff}.
                    Thus $\apply{g}{(\apply{f}{\restrl{u_1}{\initialSegment{n}}}, \apply{u_1}{n + 1})} = \apply{g}{(\apply{f}{\restrl{u_2}{\initialSegment{n}}}, \apply{u_2}{n + 1})}$ by \cref{funs_elim,function_apply_iff}.
                    Thus $\apply{f}{\restrl{u_1}{\initialSegment{n}}}\in\naturalsPlus$ and
                        $\apply{f}{\restrl{u_2}{\initialSegment{n}}}\in\naturalsPlus$ by \cref{inj,funs_type_apply}.
                    $n + 1\in\naturalsPlus$ and $n + 1\in\initialSegment{n + 1}$ by \cref{real_naturals_closed_succ,initial_segment_naturalsplus}.
                    Thus $\apply{u_1}{n + 1}\in\naturalsPlus$ and $\apply{u_2}{n + 1}\in\naturalsPlus$ by \cref{funs_type_apply}.
                    Thus $(\apply{f}{\restrl{u_1}{\initialSegment{n}}}, \apply{u_1}{n + 1})\in\naturalsPlus\times\naturalsPlus$ and
                        $(\apply{f}{\restrl{u_2}{\initialSegment{n}}}, \apply{u_2}{n + 1})\in\naturalsPlus\times\naturalsPlus$ by \cref{times_tuple_intro}.
                    Thus $(\apply{f}{\restrl{u_1}{\initialSegment{n}}}, \apply{u_1}{n + 1}) = (\apply{f}{\restrl{u_2}{\initialSegment{n}}}, \apply{u_2}{n + 1})$ by \cref{inj}.
                    Thus $\apply{f}{\restrl{u_1}{\initialSegment{n}}} = \apply{f}{\restrl{u_2}{\initialSegment{n}}}$ and
                        $\apply{u_1}{n + 1} = \apply{u_2}{n + 1}$ by \cref{pair_eq_iff}.
                    Thus $\restrl{u_1}{\initialSegment{n}} = \restrl{u_2}{\initialSegment{n}}$ by \cref{inj}.
                    $u_1$ is a function and $\dom{u_1} = \initialSegment{n + 1}$ and
                        $u_2$ is a function and $\dom{u_2} = \initialSegment{n + 1}$ by \cref{funs_elim}.
                    For all $k\in\initialSegment{n + 1}$ we have $u_1(k) = u_2(k)$
                        by \cref{real_naturals_leq_succ_elim,initial_segment_naturalsplus,initial_segment_subset_succ,funs_elim,subseteq,restrl_of_function_apply}.
                    Hence $u_1 = u_2$ by \cref{functions_eq_iff}.
                \end{subproof}
            Therefore $F\in\inj{\funs{\initialSegment{n + 1}}{\naturalsPlus}}{\naturalsPlus}$ by \cref{inj}.
            Hence $n + 1\in A$.
        \end{subproof}
    Hence there exists $f\in\inj{\funs{\initialSegment{n}}{\naturalsPlus}}{\naturalsPlus}$ by \cref{subseteq,real_naturals_induction,pair_eq_iff}.
\end{proof}

% from §30 helper lemma 10: injective coding of all finite sequences of naturalsPlus
\begin{lemma}\label{naturalsplus_finite_sequences_all_inj}
    Let $F = \{ \funs{\initialSegment{k}}{\naturalsPlus} \mid k\in\naturalsPlus \}$.
    Let $S = \unions{F}$.
    Then there exists $h\in\inj{S}{\naturalsPlus}$.
\end{lemma}
\begin{proof}
    Let $F = \{ \funs{\initialSegment{k}}{\naturalsPlus} \mid k\in\naturalsPlus \}$.
    Let $S = \unions{F}$.
    Let $R = \{(k,f) \mid k\in\naturalsPlus \mid f\in\inj{\funs{\initialSegment{k}}{\naturalsPlus}}{\naturalsPlus}\}$.
    Show that for all $k\in\naturalsPlus$ there exists $f$ such that $k\mathrel{R} f$.
    \begin{subproof}
        Fix $k\in\naturalsPlus$.
        Take $f\in\inj{\funs{\initialSegment{k}}{\naturalsPlus}}{\naturalsPlus}$ such that $f = f$ by \cref{naturalsplus_finite_sequences_inj}.
                Hence $k\mathrel{R} f$ by \cref{pair_eq_iff}.
    \end{subproof}
    $R$ is a relation by \cref{relation,pair_eq_iff}.
    Take $F$ such that $F$ is a function on $\naturalsPlus$ and for all $k\in\naturalsPlus$ we have $k\mathrel{R} \apply{F}{k}$
        by \cref{choice_function}.
    Take $g\in\inj{\naturalsPlus\times\naturalsPlus}{\naturalsPlus}$ such that $g = g$ by \cref{naturalsplus_pairing_injection}.
    Let $h = \{(u, \apply{g}{(k, \apply{\apply{F}{k}}{u})}) \mid k\in\naturalsPlus \mid u\in\funs{\initialSegment{k}}{\naturalsPlus}\}$.
    Show that $h\in\inj{S}{\naturalsPlus}$.
    \begin{subproof}
            Show that for all $u,a,b$ such that $(u,a)\in h$ and $(u,b)\in h$ we have $a = b$.
            \begin{subproof}
                Fix $u$.
                Fix $a$.
                Fix $b$ such that $(u,a)\in h$ and $(u,b)\in h$.
                Take $k_1\in\naturalsPlus$ such that $u\in\funs{\initialSegment{k_1}}{\naturalsPlus}$ and
                    $a = \apply{g}{(k_1, \apply{\apply{F}{k_1}}{u})}$ by \cref{pair_eq_iff}.
                Take $k_2\in\naturalsPlus$ such that $u\in\funs{\initialSegment{k_2}}{\naturalsPlus}$ and
                    $b = \apply{g}{(k_2, \apply{\apply{F}{k_2}}{u})}$ by \cref{pair_eq_iff}.
                Thus $\initialSegment{k_1} = \initialSegment{k_2}$ by \cref{funs_elim,subseteq_antisymmetric,subseteq}.
                Thus $k_1 \leq k_1$.
                Thus $k_1\in\initialSegment{k_1}$ by \cref{initial_segment_naturalsplus}.
                Thus $k_2 \leq k_2$.
                Thus $k_2\in\initialSegment{k_2}$ by \cref{initial_segment_naturalsplus}.
                Thus $k_1\in\initialSegment{k_2}$ and $k_2\in\initialSegment{k_1}$ by \cref{subseteq}.
                Thus $k_1 \leq k_2$ and $k_2 \leq k_1$ by \cref{initial_segment_naturalsplus}.
                $k_1\in\reals$ and $k_2\in\reals$ by \cref{real_naturals_subset_reals,subseteq}.
                Thus $k_1 = k_2$ by \cref{real_leq_antisym}.
                Therefore $a = b$.
            \end{subproof}
            Hence $h$ is a function by \cref{rightunique,relation,pair_eq_iff}.
            For all $u\in\dom{h}$ we have $h(u)\in\naturalsPlus$
                by \cref{dom_iff,pair_eq_iff,inj,funs_type_apply,function_apply_iff,times_tuple_intro}.
            For all $u$ such that $u\in\dom{h}$ we have $u\in S$
                by \cref{dom_iff,pair_eq_iff,unions_iff}.
            Show that for all $u$ such that $u\in S$ we have $u\in\dom{h}$.
            \begin{subproof}
                Fix $u$ such that $u\in S$.
                Take $A$ such that $A\in F$ and $u\in A$ by \cref{unions_iff}.
                Take $k\in\naturalsPlus$ such that $A = \funs{\initialSegment{k}}{\naturalsPlus}$ by \cref{pair_eq_iff}.
                Thus $u\in\dom{h}$ by \cref{pair_eq_iff,dom_intro}.
            \end{subproof}
            Hence $h\in\funs{S}{\naturalsPlus}$
                by \cref{funs_intro,dom_iff,pair_eq_iff,inj,funs_type_apply,function_apply_iff,times_tuple_intro,unions_iff,subseteq,dom_intro,subseteq_antisymmetric}.
        Show that for all $u_1,u_2$ such that $u_1\in S$ and $u_2\in S$ and $h(u_1) = h(u_2)$ we have $u_1 = u_2$.
        \begin{subproof}
            Fix $u_1$.
            Fix $u_2$ such that $u_1\in S$ and $u_2\in S$ and $h(u_1) = h(u_2)$.
            Thus $(u_1,\apply{h}{u_1})\in h$ and $(u_2,\apply{h}{u_2})\in h$ by \cref{funs_elim,subseteq,function_apply_iff}.
            Take $k_1\in\naturalsPlus$ such that $u_1\in\funs{\initialSegment{k_1}}{\naturalsPlus}$ and
                $\apply{h}{u_1} = \apply{g}{(k_1, \apply{\apply{F}{k_1}}{u_1})}$ by \cref{pair_eq_iff}.
            Take $k_2\in\naturalsPlus$ such that $u_2\in\funs{\initialSegment{k_2}}{\naturalsPlus}$ and
                $\apply{h}{u_2} = \apply{g}{(k_2, \apply{\apply{F}{k_2}}{u_2})}$ by \cref{pair_eq_iff}.
            Thus $\apply{g}{(k_1, \apply{\apply{F}{k_1}}{u_1})} = \apply{g}{(k_2, \apply{\apply{F}{k_2}}{u_2})}$ by \cref{funs_elim,function_apply_iff}.
            Thus $\apply{F}{k_1}\in\inj{\funs{\initialSegment{k_1}}{\naturalsPlus}}{\naturalsPlus}$ and
                $\apply{F}{k_2}\in\inj{\funs{\initialSegment{k_2}}{\naturalsPlus}}{\naturalsPlus}$ by \cref{pair_eq_iff}.
            Thus $\apply{\apply{F}{k_1}}{u_1}\in\naturalsPlus$ and $\apply{\apply{F}{k_2}}{u_2}\in\naturalsPlus$ by \cref{inj,funs_type_apply}.
            Thus $(k_1, \apply{\apply{F}{k_1}}{u_1})\in\naturalsPlus\times\naturalsPlus$ and
                $(k_2, \apply{\apply{F}{k_2}}{u_2})\in\naturalsPlus\times\naturalsPlus$ by \cref{times_tuple_intro}.
            Thus $(k_1, \apply{\apply{F}{k_1}}{u_1}) = (k_2, \apply{\apply{F}{k_2}}{u_2})$ by \cref{inj}.
            Thus $u_1 = u_2$ by \cref{pair_eq_iff,inj}.
        \end{subproof}
        Therefore $h\in\inj{S}{\naturalsPlus}$ by \cref{inj}.
    \end{subproof}
    Hence there exists $h\in\inj{S}{\naturalsPlus}$ by \cref{inj}.
\end{proof}

% from §30 helper lemma 11: enumerate finite sequences of naturalsPlus
\begin{lemma}\label{naturalsplus_finite_sequence_enumeration}
    Let $F = \{ \funs{\initialSegment{k}}{\naturalsPlus} \mid k\in\naturalsPlus \}$.
    Let $S = \unions{F}$.
    There exists $q$ such that
        $q$ is a function on $\naturalsPlus$ and
        for all $n\in\naturalsPlus$ we have $q(n)\in S$ and
        for all $u$ such that $u\in S$
            there exists $n\in\naturalsPlus$ such that $q(n) = u$.
\end{lemma}
\begin{proof}
    Let $F = \{ \funs{\initialSegment{k}}{\naturalsPlus} \mid k\in\naturalsPlus \}$.
    Let $S = \unions{F}$.
    Take $h\in\inj{S}{\naturalsPlus}$ such that $h = h$ by \cref{naturalsplus_finite_sequences_all_inj}.
    Let $u_0 = \{(1,1)\}$.
    $1\in\naturalsPlus$ by \cref{real_one_in_naturals}.
        $u_0$ is a function by \cref{singleton_elim,pair_eq_iff,rightunique,relation}.
        Thus for all $x$ such that $x\in\dom{u_0}$ we have $x\in\initialSegment{1}$
            by \cref{dom,singleton_elim,pair_eq_iff,real_one_in_naturals,initial_segment_naturalsplus}.
        For all $x$ such that $x\in\initialSegment{1}$ we have $x\in\dom{u_0}$
            by \cref{initial_segment_naturalsplus,real_naturals_subset_reals,real_one_in_naturals,subseteq,real_one_leq_naturalsplus,real_leq_antisym,pair_eq_iff,singleton_intro,dom}.
        Therefore $\dom{u_0} = \initialSegment{1}$ by \cref{subseteq,subseteq_antisymmetric}.
    For all $x\in\dom{u_0}$ we have $u_0(x)\in\naturalsPlus$
        by \cref{dom,singleton_elim,pair_eq_iff,function_apply_intro,real_one_in_naturals}.
    Hence $u_0\in\funs{\initialSegment{1}}{\naturalsPlus}$ by \cref{funs_intro}.
    Therefore $u_0\in S$ by \cref{unions_iff}.
    Let $R_a = \{(n,u) \mid n\in\ran{h} \mid u\in S \land h(u) = n\}$.
    Let $R_b = (\naturalsPlus\setminus\ran{h})\times\{u_0\}$.
    Let $R_1 = R_a\union R_b$.
    $R_a$ is a relation by \cref{relation,pair_eq_iff}.
    $R_b$ is a relation by \cref{relation,times_elem_is_tuple}.
    Hence $R_1$ is a relation by \cref{union_relations_is_relation}.
    Show that for all $n\in\naturalsPlus$ there exists $u$ such that $n\mathrel{R_1} u$.
    \begin{subproof}
        Fix $n\in\naturalsPlus$.
        \begin{byCase}
        \caseOf{$n\in\ran{h}$.}
        Take $u$ such that $(u,n)\in h$ by \cref{ran_iff}.
        Therefore $u\in S$ and $h(u) = n$ by \cref{dom,inj,funs_elim,function_apply_intro}.
        Hence $n\mathrel{R_1} u$ by \cref{pair_eq_iff,union_intro_left}.
        \caseOf{$n\notin\ran{h}$.}
        Hence $n\mathrel{R_1} u_0$ by \cref{setminus_intro,singleton_intro,times_tuple_intro,union_intro_right}.
        \end{byCase}
    \end{subproof}
    Take $q$ such that $q$ is a function on $\naturalsPlus$ and for all $n\in\naturalsPlus$ we have $n\mathrel{R_1} \apply{q}{n}$
        by \cref{choice_function}.
    Show that for all $n\in\naturalsPlus$ we have $q(n)\in S$.
    \begin{subproof}
        Fix $n\in\naturalsPlus$.
        \begin{byCase}
        \caseOf{$n\in\ran{h}$.}
        Hence $\apply{q}{n}\in S$ by \cref{choice_function,union_iff,pair_eq_iff,times_tuple_elim,singleton_elim}.
        \caseOf{$n\notin\ran{h}$.}
        Therefore $\apply{q}{n}\in S$ by \cref{choice_function,union_iff,pair_eq_iff,times_tuple_elim,singleton_elim}.
        \end{byCase}
    \end{subproof}
    Show that for all $u\in S$ there exists $n\in\naturalsPlus$ such that $q(n) = u$.
    \begin{subproof}
        Fix $u\in S$.
        Let $n = h(u)$.
        Thus $n\in\naturalsPlus$ and $n\in\ran{h}$
            by \cref{inj,funs_elim,funs_type_apply,function_apply_elim,ran_intro}.
        Thus $(n,\apply{q}{n})\in R_a$ or $(n,\apply{q}{n})\in R_b$ by \cref{choice_function,union_iff}.
        Hence $\apply{q}{n} = u$ by \cref{pair_eq_iff,inj,times_tuple_elim,setminus_elim_right}.
    \end{subproof}
    Therefore there exists $q$ such that
        $q$ is a function on $\naturalsPlus$ and
        for all $n\in\naturalsPlus$ we have $q(n)\in S$ and
        for all $u\in S$ there exists $n\in\naturalsPlus$ such that $q(n) = u$.
\end{proof}
% from §30 Item 10: countable product of second-countable spaces is second-countable
\begin{lemma}\label{countable_product_second_countable}
    Assume $X$ is a function.
    Assume $T$ is a function.
    Assume $\dom{X} = \naturalsPlus$.
    Assume $\dom{T} = \naturalsPlus$.
    Assume that for all $n\in\naturalsPlus$ we have $\apply{T}{n}$ is a topology on $\apply{X}{n}$.
    Suppose for all $n\in\naturalsPlus$ we have $\apply{X}{n}$ satisfies the second countability axiom with respect to $\apply{T}{n}$.
    Then $\productFamily{X}$ satisfies the second countability axiom with respect to $\productTopologyIndexed{T}{X}$.
\end{lemma}
\begin{proof}
    Let $T_0 = \productTopologyIndexed{T}{X}$.
    Let $R = \{(n,b) \mid n\in\naturalsPlus \mid
        \text{$b$ is a sequence in $\pow{\apply{X}{n}}$ and
        for all $k\in\naturalsPlus$ we have $b(k)$ is open in $\apply{X}{n}$ with respect to $\apply{T}{n}$ and
        for all $x\in\apply{X}{n}$ we have
            for all $U$ such that $U$ is open in $\apply{X}{n}$ with respect to $\apply{T}{n}$ and $x\in U$
                there exists $k\in\naturalsPlus$ such that $x\in b(k)$ and $b(k)\subseteq U$}\}$.
    $R$ is a relation by \cref{pair_eq_iff,relation}.
    Show that for all $n\in\naturalsPlus$ there exists $b$ such that $n\mathrel{R} b$.
    \begin{subproof}
        Fix $n\in\naturalsPlus$.
        Take $b$ such that
            $b$ is a sequence in $\pow{\apply{X}{n}}$ and
            for all $k\in\naturalsPlus$ we have $b(k)$ is open in $\apply{X}{n}$ with respect to $\apply{T}{n}$ and
            for all $x\in\apply{X}{n}$ we have
                for all $U$ such that $U$ is open in $\apply{X}{n}$ with respect to $\apply{T}{n}$ and $x\in U$
                    there exists $k\in\naturalsPlus$ such that $x\in b(k)$ and $b(k)\subseteq U$
            by \cref{second_countability_axiom}.
        Hence $n\mathrel{R} b$ by \cref{pair_eq_iff}.
    \end{subproof}
    Take $b$ such that $b$ is a function on $\naturalsPlus$ and for all $n\in\naturalsPlus$ we have
        $n\mathrel{R}\apply{b}{n}$ by \cref{choice_function}.

    Let $F = \{ \funs{\initialSegment{k}}{\naturalsPlus} \mid k\in\naturalsPlus \}$.
    Let $S = \unions{F}$.

    Take $q$ such that
        $q$ is a function on $\naturalsPlus$ and
        for all $n\in\naturalsPlus$ we have $q(n)\in S$ and
        for all $u$ such that $u\in S$
            there exists $n\in\naturalsPlus$ such that $q(n) = u$
        by \cref{naturalsplus_finite_sequence_enumeration}.

    Let $R_d = \{(n,g) \mid n\in\naturalsPlus \mid \text{$g$ is a function on $\dom{\apply{q}{n}}$ and
        for all $i\in\dom{\apply{q}{n}}$ we have
            $(i, \preimg{\piIndex{X}{i}}{\apply{\apply{b}{i}}{\apply{\apply{q}{n}}{i}}})\in g$}\}$.
    $R_d$ is a relation by \cref{pair_eq_iff,relation}.
    Show that for all $n\in\naturalsPlus$ there exists $g$ such that $n\mathrel{R_d} g$.
    \begin{subproof}
        Fix $n\in\naturalsPlus$.
        Take $m\in\naturalsPlus$ such that $\apply{q}{n}\in\funs{\initialSegment{m}}{\naturalsPlus}$
            by \cref{naturalsplus_finite_sequence_enumeration,unions_iff}.
        Then $\apply{q}{n}$ is a function and $\dom{\apply{q}{n}} = \initialSegment{m}$ by \cref{funs_elim}.
        Let $g = \{(i, \preimg{\piIndex{X}{i}}{\apply{\apply{b}{i}}{\apply{\apply{q}{n}}{i}}}) \mid i\in\dom{\apply{q}{n}}\}$.
            For all $i,A_1,A_2$ such that $(i,A_1)\in g$ and $(i,A_2)\in g$ we have $A_1 = A_2$ by \cref{pair_eq_iff}.
            Hence $g$ is a function on $\dom{\apply{q}{n}}$
                by \cref{rightunique,relation,pair_eq_iff,subseteq,dom_intro,dom_iff,subseteq_antisymmetric}.
        For all $i\in\dom{\apply{q}{n}}$ we have
            $\apply{g}{i} = \preimg{\piIndex{X}{i}}{\apply{\apply{b}{i}}{\apply{\apply{q}{n}}{i}}}$
            by \cref{pair_eq_iff,dom_intro,function_apply_iff}.
        Therefore $n\mathrel{R_d} g$ by \cref{pair_eq_iff}.
    \end{subproof}
    Take $h$ such that $h$ is a function on $\naturalsPlus$ and for all $n\in\naturalsPlus$ we have
        $n\mathrel{R_d}\apply{h}{n}$ by \cref{choice_function}.

    Let $d = \{(n, \inters{\img{\apply{h}{n}}{\dom{\apply{q}{n}}}}) \mid n\in\naturalsPlus\}$.

    Hence $1\in\dom{X}$ by \cref{real_one_in_naturals,subseteq}.
    Then there exists $a$ such that $a\in\dom{X}$.
    For all $n\in\dom{X}$ we have $\apply{T}{n}$ is a topology on $\apply{X}{n}$.
    Hence $\productSubbasis{T}{X}$ is a subbasis on $\productFamily{X}$ by \cref{product_subbasis_is_subbasis}.

    For all $n,U_1,U_2$ such that $(n,U_1)\in d$ and $(n,U_2)\in d$ we have $U_1 = U_2$ by \cref{pair_eq_iff}.
    Thus $d$ is a function by \cref{rightunique,relation,pair_eq_iff}.
    Show that for all $n\in\naturalsPlus$ we have $d(n)\in\pow{\productFamily{X}}$.
    \begin{subproof}
        Fix $n\in\naturalsPlus$.
        Let $U = \inters{\img{\apply{h}{n}}{\dom{\apply{q}{n}}}}$.
        Then $d(n) = U$ by \cref{pair_eq_iff,dom_intro,function_apply_iff}.
        Let $F = \img{\apply{h}{n}}{\dom{\apply{q}{n}}}$.
        Show that for all $A$ such that $A\in F$ we have $A\in\productSubbasis{T}{X}$.
        \begin{subproof}
            Fix $A$ such that $A\in F$.
            Take $i$ such that $i\in\dom{\apply{q}{n}}$ and $(i,A)\in\apply{h}{n}$ by \cref{img_iff}.
            Now $\apply{h}{n}$ is a function on $\dom{\apply{q}{n}}$ and
                for all $j\in\dom{\apply{q}{n}}$ we have
                    $(j, \preimg{\piIndex{X}{j}}{\apply{\apply{b}{j}}{\apply{\apply{q}{n}}{j}}})\in\apply{h}{n}$
                by \cref{choice_function,pair_eq_iff}.
            Thus $i\in\dom{\apply{h}{n}}$ and
                $\apply{\apply{h}{n}}{i} = \preimg{\piIndex{X}{i}}{\apply{\apply{b}{i}}{\apply{\apply{q}{n}}{i}}}$
                by \cref{dom_intro,function_apply_iff}.
            Hence $A = \preimg{\piIndex{X}{i}}{\apply{\apply{b}{i}}{\apply{\apply{q}{n}}{i}}}$.
            Take $m\in\naturalsPlus$ such that $\apply{q}{n}\in\funs{\initialSegment{m}}{\naturalsPlus}$
                by \cref{naturalsplus_finite_sequence_enumeration,unions_iff}.
            Then $\apply{q}{n}$ is a function and $\dom{\apply{q}{n}} = \initialSegment{m}$ by \cref{funs_elim}.
            Thus $i\in\dom{X}$ by \cref{funs_elim,subseteq,initial_segment_naturalsplus}.
            Now $\apply{b}{i}$ is a sequence in $\pow{\apply{X}{i}}$ and
                for all $k\in\naturalsPlus$ we have $\apply{\apply{b}{i}}{k}$ is open in $\apply{X}{i}$ with respect to $\apply{T}{i}$
                by \cref{choice_function,pair_eq_iff}.
            Therefore $\apply{\apply{b}{i}}{\apply{\apply{q}{n}}{i}}\in\apply{T}{i}$
                by \cref{funs_elim,function_apply_elim,ran_iff,function_to_ran,subseteq,open_in}.
            Hence $A\in\pow{\productFamily{X}}$ by \cref{projection_map_indexed_function,funs_elim,preimg_subseteq_dom,powerset_intro}.
            Therefore $A\in\productSubbasisAt{T}{X}{i}$ by \cref{product_subbasis_indexed}.
            Hence $A\in\productSubbasis{T}{X}$ by \cref{product_subbasis_union_indexed}.
        \end{subproof}
        Hence $F\subseteq\productSubbasis{T}{X}$ by \cref{subseteq}.
        Hence $F$ is finite
            by \cref{naturalsplus_finite_sequence_enumeration,unions_iff,funs_elim,initial_segment_finite,choice_function,pair_eq_iff,finite_image_function}.
        Take $m\in\naturalsPlus$ such that $\apply{q}{n}\in\funs{\initialSegment{m}}{\naturalsPlus}$
            by \cref{naturalsplus_finite_sequence_enumeration,unions_iff}.
        Then $\apply{q}{n}$ is a function and $\dom{\apply{q}{n}} = \initialSegment{m}$ by \cref{funs_elim}.
        Hence $\apply{\apply{h}{n}}{m}\in F$
            by \cref{funs_elim,subseteq,initial_segment_naturalsplus,choice_function,pair_eq_iff,function_apply_iff,img_iff}.
        For all $x$ such that $x\in U$ we have $x\in\unions{F}$ by \cref{inters}.
        Hence $U\subseteq\unions{F}$ by \cref{subseteq}.
        For all $x$ such that $x\in\unions{F}$ we have $x\in\unions{\productSubbasis{T}{X}}$.
        Therefore $\unions{F}\subseteq\unions{\productSubbasis{T}{X}}$ by \cref{subseteq}.
        Hence $U\subseteq\unions{\productSubbasis{T}{X}}$ by \cref{subseteq_transitive}.
        Therefore $U\in\pow{\unions{\productSubbasis{T}{X}}}$ by \cref{powerset_intro}.
        $\unions{\productSubbasis{T}{X}} = \productFamily{X}$ by \cref{subbasis_on}.
        Hence $U\in\pow{\productFamily{X}}$ by \cref{subseteq}.
        Therefore $d(n)\in\pow{\productFamily{X}}$.
    \end{subproof}
    Therefore $d\in\funs{\naturalsPlus}{\pow{\productFamily{X}}}$
        by \cref{funs_intro,subseteq,pair_eq_iff,dom_intro,dom_iff,subseteq_antisymmetric}.
    Hence $d$ is a sequence in $\pow{\productFamily{X}}$ by \cref{sequence_in}.

    Show that for all $n\in\naturalsPlus$ we have $d(n)$ is open in $\productFamily{X}$ with respect to $T_0$.
    \begin{subproof}
        Fix $n\in\naturalsPlus$.
        Let $F = \img{\apply{h}{n}}{\dom{\apply{q}{n}}}$.
        Show that for all $A$ such that $A\in F$ we have $A\in\productSubbasis{T}{X}$.
        \begin{subproof}
            Fix $A$ such that $A\in F$.
            Take $i$ such that $i\in\dom{\apply{q}{n}}$ and $(i,A)\in\apply{h}{n}$ by \cref{img_iff}.
            Now $\apply{h}{n}$ is a function on $\dom{\apply{q}{n}}$ and
                for all $j\in\dom{\apply{q}{n}}$ we have
                    $(j,\preimg{\piIndex{X}{j}}{\apply{\apply{b}{j}}{\apply{\apply{q}{n}}{j}}})\in\apply{h}{n}$
                by \cref{choice_function,pair_eq_iff}.
            Thus $i\in\dom{\apply{h}{n}}$ and
                $\apply{\apply{h}{n}}{i} = \preimg{\piIndex{X}{i}}{\apply{\apply{b}{i}}{\apply{\apply{q}{n}}{i}}}$
                by \cref{dom_intro,function_apply_iff}.
            Hence $A = \preimg{\piIndex{X}{i}}{\apply{\apply{b}{i}}{\apply{\apply{q}{n}}{i}}}$.
            Take $m\in\naturalsPlus$ such that $\apply{q}{n}\in\funs{\initialSegment{m}}{\naturalsPlus}$
                by \cref{naturalsplus_finite_sequence_enumeration,unions_iff}.
            Then $\apply{q}{n}$ is a function to $\naturalsPlus$ and $\dom{\apply{q}{n}} = \initialSegment{m}$ by \cref{funs_elim}.
            Thus $i\in\dom{X}$ by \cref{funs_elim,subseteq,initial_segment_naturalsplus}.
            Now $\apply{b}{i}$ is a sequence in $\pow{\apply{X}{i}}$ and
                for all $k\in\naturalsPlus$ we have $\apply{\apply{b}{i}}{k}$ is open in $\apply{X}{i}$ with respect to $\apply{T}{i}$
                by \cref{choice_function,pair_eq_iff}.
            Therefore $\apply{\apply{b}{i}}{\apply{\apply{q}{n}}{i}}\in\apply{T}{i}$
                by \cref{function_to_ran,function_apply_elim,ran_iff,subseteq,open_in}.
            Hence $A\in\pow{\productFamily{X}}$ by \cref{projection_map_indexed_function,funs_elim,preimg_subseteq_dom,powerset_intro}.
            Therefore $A\in\productSubbasisAt{T}{X}{i}$ by \cref{product_subbasis_indexed}.
            Hence $A\in\productSubbasis{T}{X}$ by \cref{product_subbasis_union_indexed}.
        \end{subproof}
        Hence $F\subseteq\productSubbasis{T}{X}$ by \cref{subseteq}.
        Hence $F$ is finite
            by \cref{naturalsplus_finite_sequence_enumeration,unions_iff,funs_elim,initial_segment_finite,choice_function,pair_eq_iff,finite_image_function}.
        Take $m\in\naturalsPlus$ such that $\apply{q}{n}\in\funs{\initialSegment{m}}{\naturalsPlus}$
            by \cref{naturalsplus_finite_sequence_enumeration,unions_iff}.
        Then $\apply{q}{n}$ is a function and $\dom{\apply{q}{n}} = \initialSegment{m}$ by \cref{funs_elim}.
        Hence $\apply{\apply{h}{n}}{m}\in F$
            by \cref{funs_elim,subseteq,initial_segment_naturalsplus,choice_function,pair_eq_iff,function_apply_iff,img_iff}.
        For all $x$ such that $x\in\inters{F}$ we have $x\in\unions{\productSubbasis{T}{X}}$.
        Hence $\inters{F}\subseteq\unions{\productSubbasis{T}{X}}$ by \cref{subseteq}.
        Hence $\inters{F}\in\finiteIntersections{\productSubbasis{T}{X}}$ by \cref{powerset_intro,finite_intersections}.
        Therefore $\finiteIntersections{\productSubbasis{T}{X}}$ is a basis for $T_0$ on $\productFamily{X}$
            by \cref{product_subbasis_finite_intersections_basis}.
        Hence $T_0 = \basisTopology{\finiteIntersections{\productSubbasis{T}{X}}}{\productFamily{X}}$
            by \cref{product_topology_indexed,topology_generated_by_subbasis,basis_for_topology}.
        Therefore $\inters{F}\in\basisTopology{\finiteIntersections{\productSubbasis{T}{X}}}{\productFamily{X}}$
            by \cref{basis_elem_open_in_generated,basis_for_topology}.
        Hence $\inters{F}\in T_0$ by \cref{basis_for_topology}.
        Therefore $T_0$ is a topology on $\productFamily{X}$ by \cref{basis_for_topology,basis_topology_is_topology}.
        Hence $d(n)$ is open in $\productFamily{X}$ with respect to $T_0$
            by \cref{basis_for_topology,basis_topology_is_topology,sequence_in,funs_is_function,funs_elim,subseteq,pair_eq_iff,function_apply_iff,open_in}.
    \end{subproof}

    Show that for all $x\in\productFamily{X}$ we have
        for all $U$ such that $U$ is open in $\productFamily{X}$ with respect to $T_0$ and $x\in U$
            there exists $n\in\naturalsPlus$ such that $x\in d(n)$ and $d(n)\subseteq U$.
    \begin{subproof}
        Fix $x\in\productFamily{X}$.
        Fix $U$ such that $U$ is open in $\productFamily{X}$ with respect to $T_0$ and $x\in U$.
        $U\in T_0$ by \cref{open_in}.
        Hence $1\in\dom{X}$ by \cref{real_one_in_naturals,subseteq}.
        Then there exists $a$ such that $a\in\dom{X}$.
        For all $n\in\dom{X}$ we have $\apply{T}{n}$ is a topology on $\apply{X}{n}$.
        Hence $\finiteIntersections{\productSubbasis{T}{X}}$ is a basis for $T_0$ on $\productFamily{X}$
            by \cref{product_subbasis_finite_intersections_basis}.
        Therefore $T_0 = \basisTopology{\finiteIntersections{\productSubbasis{T}{X}}}{\productFamily{X}}$
            by \cref{product_topology_indexed,topology_generated_by_subbasis,basis_for_topology}.
        Hence $U\in\basisTopology{\finiteIntersections{\productSubbasis{T}{X}}}{\productFamily{X}}$.
        Take $B\in\finiteIntersections{\productSubbasis{T}{X}}$ such that $x\in B$ and $B\subseteq U$
            by \cref{topology_generated_by_basis}.
        Take $F\subseteq\productSubbasis{T}{X}$ such that $F$ is finite and inhabited and $B = \inters{F}$ by \cref{finite_intersections}.

        Let $R_2 = \{(A,i) \mid A\in F \mid
            \text{$i\in\dom{X}$ and there exists $V\in\apply{T}{i}$ such that $A = \preimg{\piIndex{X}{i}}{V}$}\}$.
        $R_2$ is a relation by \cref{pair_eq_iff,relation}.
        Show that for all $A\in F$ there exists $i$ such that $A\mathrel{R_2} i$.
        \begin{subproof}
            Fix $A\in F$.
            Take $i$ such that $i\in\dom{X}$ and $A\in\productSubbasisAt{T}{X}{i}$
                by \cref{subseteq,product_subbasis_union_indexed}.
            Take $V\in\apply{T}{i}$ such that $A = \preimg{\piIndex{X}{i}}{V}$ by \cref{product_subbasis_indexed}.
            Hence $A\mathrel{R_2} i$ by \cref{pair_eq_iff}.
        \end{subproof}
        Take $p$ such that $p$ is a function on $F$ and for all $A\in F$ we have $A\mathrel{R_2}\apply{p}{A}$
            by \cref{choice_function}.
        Let $J = \img{p}{F}$.
        Hence $J$ is finite and inhabited by \cref{finite_image_function,choice_function,subseteq,function_apply_iff,img_iff}.
        For all $j$ such that $j\in J$ we have $j\in\naturalsPlus$
            by \cref{img_iff,choice_function,subseteq,function_apply_iff,pair_eq_iff}.
        Hence $J\subseteq\naturalsPlus$ by \cref{subseteq}.

        Let $R_3 = \{(A,V) \mid A\in F \mid
            \text{$V\in\apply{T}{\apply{p}{A}}$ and $A = \preimg{\piIndex{X}{\apply{p}{A}}}{V}$}\}$.
        $R_3$ is a relation by \cref{pair_eq_iff,relation}.
        Show that for all $A\in F$ there exists $V$ such that $A\mathrel{R_3} V$.
        \begin{subproof}
            Fix $A\in F$.
            Take $V\in\apply{T}{\apply{p}{A}}$ such that $A = \preimg{\piIndex{X}{\apply{p}{A}}}{V}$
                by \cref{choice_function,pair_eq_iff}.
            Hence $A\mathrel{R_3} V$ by \cref{pair_eq_iff}.
        \end{subproof}
        Take $r$ such that $r$ is a function on $F$ and for all $A\in F$ we have $A\mathrel{R_3}\apply{r}{A}$
            by \cref{choice_function}.

        Let $R_4 = \{(i,V) \mid i\in J \mid
            V = \inters{\img{r}{\preimg{p}{\{i\}}}}\}$.
        $R_4$ is a relation by \cref{pair_eq_iff,relation}.
        Show that for all $i\in J$ there exists $V$ such that $i\mathrel{R_4} V$.
        \begin{subproof}
            Fix $i\in J$.
            Let $H = \img{r}{\preimg{p}{\{i\}}}$.
            Let $V = \inters{H}$.
            Hence $i\mathrel{R_4} V$ by \cref{pair_eq_iff}.
        \end{subproof}
        Take $v$ such that $v$ is a function on $J$ and for all $i\in J$ we have $i\mathrel{R_4}\apply{v}{i}$
            by \cref{choice_function}.

        Show that for all $i\in J$ we have $\apply{v}{i}\in\apply{T}{i}$.
        \begin{subproof}
            Fix $i\in J$.
            Let $H = \img{r}{\preimg{p}{\{i\}}}$.
            Then $\apply{v}{i} = \inters{H}$ by \cref{choice_function,pair_eq_iff}.
            Show that for all $V$ such that $V\in H$ we have $V\in\apply{T}{i}$.
            \begin{subproof}
                Fix $V$ such that $V\in H$.
                Take $A$ such that $A\in\preimg{p}{\{i\}}$ and $(A,V)\in r$ by \cref{img_iff}.
                Take $y\in\{i\}$ such that $A\mathrel{p} y$ by \cref{preimg_iff}.
                Thus $(A,y)\in p$.
                Hence $A\in F$ and $\apply{p}{A} = i$ by \cref{choice_function,dom_intro,subseteq,singleton_elim,function_apply_iff}.
                Therefore $V\in\apply{T}{i}$ by \cref{choice_function,pair_eq_iff,dom_intro,subseteq,singleton_elim,function_apply_iff}.
            \end{subproof}
            Hence $H\subseteq\apply{T}{i}$ by \cref{subseteq}.
            Therefore $\img{r}{F}$ is finite by \cref{finite_image_function}.
            Show that for all $V$ such that $V\in H$ we have $V\in\img{r}{F}$.
            \begin{subproof}
                Fix $V$ such that $V\in H$.
                Take $A$ such that $A\in\preimg{p}{\{i\}}$ and $(A,V)\in r$ by \cref{img_iff}.
                Take $y\in\{i\}$ such that $A\mathrel{p} y$ by \cref{preimg_iff}.
                Thus $(A,y)\in p$.
                Hence $A\in F$ and $V\in\img{r}{F}$ by \cref{choice_function,dom_intro,subseteq,img_iff}.
            \end{subproof}
            Therefore $H$ is finite by \cref{subseteq,subseteq_finite_is_finite,finite_image_function}.
            Hence $H$ is inhabited
                by \cref{img_iff,singleton_intro,preimg_iff,choice_function,subseteq,function_apply_iff}.
            Therefore $\apply{v}{i}\in\apply{T}{i}$ by \cref{subseteq,finite_intersections_open_in_topology}.
        \end{subproof}

        Show that for all $i\in J$ we have $\apply{x}{i}\in\apply{v}{i}$.
        \begin{subproof}
            Fix $i\in J$.
            Let $H = \img{r}{\preimg{p}{\{i\}}}$.
            Thus $\apply{v}{i} = \inters{H}$ by \cref{choice_function,pair_eq_iff}.
            Show that for all $V\in H$ we have $\apply{x}{i}\in V$.
            \begin{subproof}
                Fix $V\in H$.
                Take $A$ such that $A\in\preimg{p}{\{i\}}$ and $(A,V)\in r$ by \cref{img_iff}.
                Take $y\in\{i\}$ such that $A\mathrel{p} y$ by \cref{preimg_iff}.
                Thus $(A,y)\in p$.
                Thus $A\in F$ and $\apply{p}{A} = i$ by \cref{choice_function,dom_intro,subseteq,singleton_elim,function_apply_iff}.
                Thus $\apply{r}{A} = V$ by \cref{choice_function,subseteq,function_apply_iff}.
                Thus $A = \preimg{\piIndex{X}{i}}{V}$ by \cref{choice_function,pair_eq_iff}.
                Thus $x\in A$ by \cref{inters_subseteq_elem,subseteq}.
                Take $y\in V$ such that $(x,y)\in\piIndex{X}{i}$ by \cref{preimg_iff}.
                Take $z\in\productFamily{X}$ such that $(z,\apply{z}{i}) = (x,y)$ by \cref{projection_map_indexed}.
                Thus $z = x$ and $\apply{z}{i} = y$ by \cref{pair_eq_iff}.
                Thus $\apply{x}{i} = y$.
                Thus $\apply{x}{i}\in V$.
            \end{subproof}
            Hence $H$ is inhabited
                by \cref{img_iff,singleton_intro,preimg_iff,choice_function,subseteq,function_apply_iff}.
            Thus $\apply{x}{i}\in\apply{v}{i}$ by \cref{choice_function,pair_eq_iff,inters_intro}.
        \end{subproof}

        Let $R_5 = \{(i,k) \mid i\in J \mid
            \text{$k\in\naturalsPlus$ and $\apply{x}{i}\in\apply{\apply{b}{i}}{k}$ and $\apply{\apply{b}{i}}{k}\subseteq\apply{v}{i}$}\}$.
        Thus $R_5$ is a relation by \cref{pair_eq_iff,relation}.
        Show that for all $i\in J$ there exists $k$ such that $i\mathrel{R_5} k$.
        \begin{subproof}
            Fix $i\in J$.
            Thus $\apply{v}{i}$ is open in $\apply{X}{i}$ with respect to $\apply{T}{i}$ by \cref{open_in}.
            $\apply{x}{i}\in\apply{v}{i}$.
            Thus $\apply{b}{i}$ is a sequence in $\pow{\apply{X}{i}}$ and
                for all $k\in\naturalsPlus$ we have $\apply{\apply{b}{i}}{k}$ is open in $\apply{X}{i}$ with respect to $\apply{T}{i}$ and
                for all $y\in\apply{X}{i}$ we have
                    for all $U$ such that $U$ is open in $\apply{X}{i}$ with respect to $\apply{T}{i}$ and $y\in U$
                        there exists $k\in\naturalsPlus$ such that $y\in\apply{\apply{b}{i}}{k}$ and $\apply{\apply{b}{i}}{k}\subseteq U$
                by \cref{choice_function,pair_eq_iff}.
            Take $k\in\naturalsPlus$ such that
                $\apply{x}{i}\in\apply{\apply{b}{i}}{k}$ and $\apply{\apply{b}{i}}{k}\subseteq\apply{v}{i}$.
            Thus $i\mathrel{R_5} k$ by \cref{pair_eq_iff}.
        \end{subproof}
        Take $s$ such that $s$ is a function on $J$ and for all $i\in J$ we have $i\mathrel{R_5}\apply{s}{i}$
            by \cref{choice_function}.

        Thus $J\subseteq\reals$ by \cref{real_naturals_subset_reals,subseteq_transitive}.
        Take $m$ such that $m$ is a largest element of $J$ with respect to $\realOrderRel$
            by \cref{finite_inhabited_largest_element,real_order_relation_simple}.
        Thus $m\in\naturalsPlus$ by \cref{largest_element,subseteq}.

        Let $R_6 = \{(i,k) \mid i\in\initialSegment{m} \mid
            \text{$k\in\naturalsPlus$ and
            $((i\in J \land \apply{x}{i}\in\apply{\apply{b}{i}}{k} \land \apply{\apply{b}{i}}{k}\subseteq\apply{v}{i})
            \lor (i\notin J \land \apply{x}{i}\in\apply{\apply{b}{i}}{k}))$}\}$.
        Thus $R_6$ is a relation by \cref{pair_eq_iff,relation}.
        Show that for all $i\in\initialSegment{m}$ there exists $k$ such that $i\mathrel{R_6} k$.
        \begin{subproof}
            Fix $i\in\initialSegment{m}$.
            \begin{byCase}
            \caseOf{$i\in J$.}
                Thus $i\mathrel{R_6}\apply{s}{i}$ by \cref{choice_function,pair_eq_iff}.
            \caseOf{$i\notin J$.}
                Thus $i\in\dom{X}$ by \cref{subseteq}.
                Thus $\apply{X}{i}$ is open in $\apply{X}{i}$ with respect to $\apply{T}{i}$ by \cref{topology_on,open_in}.
                Thus $\apply{x}{i}\in\apply{X}{i}$ by \cref{indexed_product}.
                Thus $\apply{b}{i}$ is a sequence in $\pow{\apply{X}{i}}$ and
                    for all $k\in\naturalsPlus$ we have $\apply{\apply{b}{i}}{k}$ is open in $\apply{X}{i}$ with respect to $\apply{T}{i}$ and
                    for all $y\in\apply{X}{i}$ we have
                        for all $U$ such that $U$ is open in $\apply{X}{i}$ with respect to $\apply{T}{i}$ and $y\in U$
                            there exists $k\in\naturalsPlus$ such that $y\in\apply{\apply{b}{i}}{k}$ and $\apply{\apply{b}{i}}{k}\subseteq U$
                    by \cref{choice_function,pair_eq_iff}.
                Take $k\in\naturalsPlus$ such that $\apply{x}{i}\in\apply{\apply{b}{i}}{k}$ and $\apply{\apply{b}{i}}{k}\subseteq\apply{X}{i}$.
                Thus $i\mathrel{R_6} k$ by \cref{pair_eq_iff}.
            \end{byCase}
        \end{subproof}
        Take $u$ such that $u$ is a function on $\initialSegment{m}$ and for all $i\in\initialSegment{m}$ we have
            $i\mathrel{R_6}\apply{u}{i}$ by \cref{choice_function}.
        Thus $u\in\funs{\initialSegment{m}}{\naturalsPlus}$ by \cref{funs_intro,subseteq,choice_function,pair_eq_iff}.
        Thus $u\in S$ by \cref{unions_intro}.
        Take $n\in\naturalsPlus$ such that $\apply{q}{n} = u$ by \cref{naturalsplus_finite_sequence_enumeration}.

        Let $F_n = \img{\apply{h}{n}}{\dom{\apply{q}{n}}}$.
        Thus $\apply{q}{n} = u$.
        Thus $\dom{\apply{q}{n}} = \initialSegment{m}$ by \cref{funs_elim}.
        For all $i\in\initialSegment{m}$ we have
            $x\in\preimg{\piIndex{X}{i}}{\apply{\apply{b}{i}}{\apply{u}{i}}}$
            by \cref{choice_function,pair_eq_iff,preimg_iff,projection_map_indexed}.
        Show that for all $A\in F_n$ we have $x\in A$.
        \begin{subproof}
            Fix $A\in F_n$.
            Take $i$ such that $i\in\dom{\apply{q}{n}}$ and $(i,A)\in\apply{h}{n}$ by \cref{img_iff}.
            Thus $i\in\initialSegment{m}$ by \cref{subseteq}.
            Thus $\apply{h}{n}$ is a function on $\dom{\apply{q}{n}}$ and
                for all $j\in\dom{\apply{q}{n}}$ we have
                    $(j,\preimg{\piIndex{X}{j}}{\apply{\apply{b}{j}}{\apply{\apply{q}{n}}{j}}})\in\apply{h}{n}$
                by \cref{choice_function,pair_eq_iff}.
            Thus $i\in\dom{\apply{h}{n}}$ and
                $\apply{\apply{h}{n}}{i} = \preimg{\piIndex{X}{i}}{\apply{\apply{b}{i}}{\apply{\apply{q}{n}}{i}}}$
                by \cref{dom_intro,function_apply_iff}.
            Thus $A = \preimg{\piIndex{X}{i}}{\apply{\apply{b}{i}}{\apply{\apply{q}{n}}{i}}}$.
            Thus $\apply{\apply{q}{n}}{i} = \apply{u}{i}$.
            Thus $x\in A$ by \cref{subseteq}.
        \end{subproof}
        Thus $m\in\dom{\apply{q}{n}}$ by \cref{subseteq,initial_segment_naturalsplus}.
        Hence $F_n$ is inhabited
            by \cref{choice_function,pair_eq_iff,subseteq,function_apply_iff,img_iff}.
        Thus $x\in d(n)$ by \cref{choice_function,pair_eq_iff,subseteq,function_apply_iff,img_iff,inters_intro,sequence_in,funs_is_function,funs_elim}.

        Show that for all $A\in F$ we have $d(n)\subseteq A$.
        \begin{subproof}
            Fix $A\in F$.
            Thus $(A,\apply{p}{A})\in p$ and $\apply{p}{A}\in J$
                by \cref{choice_function,subseteq,function_apply_elim,img_iff}.
            Let $i = \apply{p}{A}$.
            Thus $i\in\naturalsPlus$ by \cref{subseteq}.
            Thus $i\leq m$ by \cref{largest_element,real_order_relation_elim}.
            Thus $i\in\initialSegment{m}$ by \cref{initial_segment_naturalsplus}.
            Thus $\apply{u}{i}\in\naturalsPlus$ and
                $\apply{x}{i}\in\apply{\apply{b}{i}}{\apply{u}{i}}$ and
                $\apply{\apply{b}{i}}{\apply{u}{i}}\subseteq\apply{v}{i}$ by \cref{choice_function,pair_eq_iff}.
            Thus $A = \preimg{\piIndex{X}{i}}{\apply{r}{A}}$ by \cref{choice_function,pair_eq_iff}.
            Let $H = \img{r}{\preimg{p}{\{i\}}}$.
            Thus $\apply{v}{i} = \inters{H}$ by \cref{choice_function,pair_eq_iff}.
            Thus $A\in\preimg{p}{\{i\}}$ by \cref{singleton_intro,preimg_iff}.
            Thus $\apply{r}{A}\in H$ by \cref{choice_function,subseteq,function_apply_iff,img_iff}.
            Thus $\apply{v}{i}\subseteq\apply{r}{A}$ by \cref{inters_subseteq_elem}.
            Thus $\apply{\apply{b}{i}}{\apply{u}{i}}\subseteq\apply{r}{A}$ by \cref{subseteq_transitive}.
            Thus $\preimg{\piIndex{X}{i}}{\apply{\apply{b}{i}}{\apply{u}{i}}}\subseteq A$ by \cref{preimg_subseteq,subseteq}.
            Thus $\apply{h}{n}$ is a function on $\dom{\apply{q}{n}}$ and
                for all $j\in\dom{\apply{q}{n}}$ we have
                    $(j,\preimg{\piIndex{X}{j}}{\apply{\apply{b}{j}}{\apply{\apply{q}{n}}{j}}})\in\apply{h}{n}$
                by \cref{choice_function,pair_eq_iff}.
            Thus $\apply{h}{n}$ is a function and $\dom{\apply{h}{n}} = \dom{\apply{q}{n}}$.
            Thus for all $j\in\dom{\apply{q}{n}}$ we have
                $\apply{\apply{h}{n}}{j} = \preimg{\piIndex{X}{j}}{\apply{\apply{b}{j}}{\apply{\apply{q}{n}}{j}}}$
                by \cref{dom_intro,function_apply_iff}.
            Thus $i\in\dom{\apply{h}{n}}$ and
                $\apply{\apply{h}{n}}{i} = \preimg{\piIndex{X}{i}}{\apply{\apply{b}{i}}{\apply{\apply{q}{n}}{i}}}$
                by \cref{subseteq}.
            Thus $\apply{\apply{q}{n}}{i} = \apply{u}{i}$.
            Thus $\preimg{\piIndex{X}{i}}{\apply{\apply{b}{i}}{\apply{u}{i}}}\in F_n$ by \cref{img_iff}.
            Thus $d(n) = \inters{F_n}$ by \cref{sequence_in,funs_is_function,funs_elim,subseteq,pair_eq_iff,function_apply_iff}.
            Thus $d(n)\subseteq \preimg{\piIndex{X}{i}}{\apply{\apply{b}{i}}{\apply{u}{i}}}$ by \cref{inters_subseteq_elem}.
            Thus $d(n)\subseteq A$ by \cref{subseteq_transitive}.
        \end{subproof}
        Thus $d(n)\subseteq U$ by \cref{subseteq_inters_iff,subseteq,subseteq_transitive}.
    \end{subproof}
    Thus $\productFamily{X}$ satisfies the second countability axiom with respect to $T_0$ by \cref{second_countability_axiom}.
\end{proof}

% from §30 Item 11: dense subset
\begin{definition}\label{dense_in}
    $A$ is dense in $X$ with respect to $T$ iff $\closure{A}{X}{T} = X$.
\end{definition}

% from §30 Item 12: countable subcover in a second-countable space
\begin{lemma}\label{second_countable_open_covering_countable_subcover}
    Assume $T$ is a topology on $X$.
    Suppose $X$ satisfies the second countability axiom with respect to $T$.
    Then for all $C$ such that $C$ is an open covering of $X$ with respect to $T$
        there exists $J$ such that
            $J\subseteq\naturalsPlus$ and
            there exists $f$ such that
                $f$ is a function on $J$ and
                for all $j\in J$ we have $\apply{f}{j}\in C$ and
                $X\subseteq\unions{\img{f}{J}}$.
\end{lemma}
\begin{proof}
    Take $b$ such that
        $b$ is a sequence in $\pow{X}$ and
        for all $n\in\naturalsPlus$ we have $b(n)$ is open in $X$ with respect to $T$ and
        for all $x\in X$ we have
            for all $U$ such that $U$ is open in $X$ with respect to $T$ and $x\in U$
                there exists $n\in\naturalsPlus$ such that $x\in b(n)$ and $b(n)\subseteq U$
        by \cref{second_countability_axiom}.
    Fix $C$ such that $C$ is an open covering of $X$ with respect to $T$.
    Let $J = \{n\in\naturalsPlus \mid \exists A\in C. b(n)\subseteq A\}$.
    Let $R = \{(n,A) \mid n\in J \mid A\in C \land b(n)\subseteq A\}$.
    $R$ is a relation by \cref{pair_eq_iff,relation}.
    Show that for all $n\in J$ there exists $A$ such that $n\mathrel{R} A$.
    \begin{subproof}
        Fix $n\in J$.
        Take $A\in C$ such that $b(n)\subseteq A$.
        Hence $n\mathrel{R} A$ by \cref{pair_eq_iff}.
    \end{subproof}
    Take $f$ such that $f$ is a function on $J$ and for all $n\in J$ we have $n\mathrel{R}\apply{f}{n}$
        by \cref{choice_function}.
    For all $j\in J$ we have $\apply{f}{j}\in C$ by \cref{choice_function,pair_eq_iff}.
    Show that for all $x$ such that $x\in X$ we have $x\in\unions{\img{f}{J}}$.
    \begin{subproof}
        Fix $x$ such that $x\in X$.
        Take $A_0$ such that $A_0\in C$ and $x\in A_0$ by \cref{open_covering,covering,subseteq,unions_iff}.
        $A_0$ is open in $X$ with respect to $T$ by \cref{open_covering}.
        Take $n\in\naturalsPlus$ such that $x\in b(n)$ and $b(n)\subseteq A_0$.
        Now $\apply{f}{n}\in C$ and $b(n)\subseteq\apply{f}{n}$ by \cref{choice_function,pair_eq_iff}.
        Hence $x\in\unions{\img{f}{J}}$ by \cref{choice_function,pair_eq_iff,subseteq,function_apply_iff,img_iff,unions_intro}.
    \end{subproof}
    Therefore $X\subseteq\unions{\img{f}{J}}$ by \cref{subseteq}.
    Hence there exists $J$ such that
        $J\subseteq\naturalsPlus$ and
        there exists $f$ such that
            $f$ is a function on $J$ and
            for all $j\in J$ we have $\apply{f}{j}\in C$ and
            $X\subseteq\unions{\img{f}{J}}$.
\end{proof}

% from §30 Item 13: countable dense subset in a second-countable space
\begin{lemma}\label{second_countable_countable_dense_subset}
    Assume $T$ is a topology on $X$.
    Suppose $X$ satisfies the second countability axiom with respect to $T$.
    Then there exists $J$ such that
        $J\subseteq\naturalsPlus$ and
        there exists $f$ such that
            $f$ is a function on $J$ and
            for all $j\in J$ we have $\apply{f}{j}\in X$ and
            $\img{f}{J}$ is dense in $X$ with respect to $T$.
\end{lemma}
\begin{proof}
    Take $b$ such that
        $b$ is a sequence in $\pow{X}$ and
        for all $n\in\naturalsPlus$ we have $b(n)$ is open in $X$ with respect to $T$ and
        for all $x\in X$ we have
            for all $U$ such that $U$ is open in $X$ with respect to $T$ and $x\in U$
                there exists $n\in\naturalsPlus$ such that $x\in b(n)$ and $b(n)\subseteq U$
        by \cref{second_countability_axiom}.
    Let $J = \{n\in\naturalsPlus \mid b(n)\neq\emptyset\}$.
    Let $R = \{(n,x) \mid n\in J \mid x\in b(n)\}$.
    $R$ is a relation by \cref{pair_eq_iff,relation}.
    Show that for all $n\in J$ there exists $x$ such that $n\mathrel{R} x$.
        \begin{subproof}
            Fix $n\in J$.
            Take $x$ such that $x\in b(n)$ by \cref{pair_eq_iff,nonempty_inhabited}.
            Hence $n\mathrel{R} x$ by \cref{pair_eq_iff}.
        \end{subproof}
    Take $f$ such that $f$ is a function on $J$ and for all $n\in J$ we have $n\mathrel{R}\apply{f}{n}$
        by \cref{choice_function}.
    $\img{f}{J}\subseteq X$
        by \cref{img_iff,choice_function,pair_eq_iff,sequence_in,funs_elim,subseteq,function_apply_iff,function_to_ran,ran_iff,powerset_elim}.
        Show that for all $x\in X$ we have $x\in\closure{\img{f}{J}}{X}{T}$.
        \begin{subproof}
            Fix $x\in X$.
            Show that for all $U$ such that $U$ is open in $X$ with respect to $T$ and $x\in U$
                we have $U$ intersects $\img{f}{J}$.
            \begin{subproof}
                Fix $U$ such that $U$ is open in $X$ with respect to $T$ and $x\in U$.
                Take $n\in\naturalsPlus$ such that $x\in b(n)$ and $b(n)\subseteq U$.
                Hence $U$ intersects $\img{f}{J}$
                    by \cref{choice_function,pair_eq_iff,subseteq,function_apply_iff,img_iff,inter_intro,emptyset,intersects}.
            \end{subproof}
            Therefore $x\in\closure{\img{f}{J}}{X}{T}$ by \cref{closure_iff_intersects_open}.
        \end{subproof}
        Therefore $\img{f}{J}$ is dense in $X$ with respect to $T$
            by \cref{subseteq,closed_whole_space,closure_subseteq_closed,subseteq_antisymmetric,dense_in}.
    Hence there exists $J$ such that
        $J\subseteq\naturalsPlus$ and
        there exists $f$ such that
            $f$ is a function on $J$ and
            for all $j\in J$ we have $\apply{f}{j}\in X$ and
            $\img{f}{J}$ is dense in $X$ with respect to $T$.
\end{proof}

% from §30 Item 14: Lindelöf space
\begin{definition}\label{lindelof_space}
    $X$ is Lindelöf with respect to $T$ iff
        for all $C$ such that $C$ is an open covering of $X$ with respect to $T$
            there exists $J$ such that
                $J\subseteq\naturalsPlus$ and
                there exists $f$ such that
                    $f$ is a function on $J$ and
                    for all $j\in J$ we have $\apply{f}{j}\in C$ and
                    $X\subseteq\unions{\img{f}{J}}$.
\end{definition}

% from §30 Item 15: separable space
\begin{definition}\label{separable_space}
    $X$ is separable with respect to $T$ iff
        there exists $J$ such that
            $J\subseteq\naturalsPlus$ and
            there exists $f$ such that
                $f$ is a function on $J$ and
                for all $j\in J$ we have $\apply{f}{j}\in X$ and
                $\img{f}{J}$ is dense in $X$ with respect to $T$.
\end{definition}


\section{§31 The Separation Axioms}

% from §31 helper definition 1: closed singletons
\begin{definition}\label{one_point_closed}
    $X$ has closed singletons with respect to $T$ iff
        for all $x$ such that $x\in X$ we have $\{x\}$ is closed in $X$ with respect to $T$.
\end{definition}

% from §31 Item 1: regular space
\begin{definition}\label{regular_space}
    $X$ is regular with respect to $T$ iff
        $T$ is a topology on $X$ and
        $X$ has closed singletons with respect to $T$ and
        for all $x,B$ such that $x\in X$ and $B$ is closed in $X$ with respect to $T$ and $x\notin B$
            there exist $U,V$ such that
                $U$ is open in $X$ with respect to $T$ and
                $V$ is open in $X$ with respect to $T$ and
                $x\in U$ and
                $B\subseteq V$ and
                $U$ is disjoint from $V$.
\end{definition}

% from §31 Item 2: normal space
\begin{definition}\label{normal_space}
    $X$ is normal with respect to $T$ iff
        $T$ is a topology on $X$ and
        $X$ has closed singletons with respect to $T$ and
        for all $A,B$ such that $A$ is closed in $X$ with respect to $T$ and $B$ is closed in $X$ with respect to $T$
            and $A$ is disjoint from $B$
            there exist $U,V$ such that
                $U$ is open in $X$ with respect to $T$ and
                $V$ is open in $X$ with respect to $T$ and
                $A\subseteq U$ and
                $B\subseteq V$ and
                $U$ is disjoint from $V$.
\end{definition}

% from §31 helper lemma 1: regular implies Hausdorff
\begin{lemma}\label{regular_implies_hausdorff}
    Assume $X$ is regular with respect to $T$.
    Then $X$ is Hausdorff with respect to $T$.
\end{lemma}
\begin{proof}
    For all $x_1,x_2$ such that $x_1\in X$ and $x_2\in X$ and $x_1\neq x_2$
        there exist $U_1,U_2$ such that
            $U_1$ is a neighborhood of $x_1$ in $X$ with respect to $T$ and
            $U_2$ is a neighborhood of $x_2$ in $X$ with respect to $T$ and
            $U_1$ is disjoint from $U_2$
        by \cref{regular_space,one_point_closed,singleton_elim,singleton_intro,subseteq,neighborhood}.
    Hence $X$ is Hausdorff with respect to $T$ by \cref{regular_space,hausdorff}.
\end{proof}

% from §31 helper lemma 2: normal implies regular
\begin{lemma}\label{normal_implies_regular}
    Assume $X$ is normal with respect to $T$.
    Then $X$ is regular with respect to $T$.
\end{lemma}
\begin{proof}
    For all $x,B$ such that $x\in X$ and $B$ is closed in $X$ with respect to $T$ and $x\notin B$
        there exist $U,V$ such that
            $U$ is open in $X$ with respect to $T$ and
            $V$ is open in $X$ with respect to $T$ and
            $x\in U$ and
            $B\subseteq V$ and
            $U$ is disjoint from $V$
        by \cref{normal_space,one_point_closed,singleton_elim,disjoint,singleton_intro,subseteq}.
    Hence $X$ is regular with respect to $T$ by \cref{normal_space,regular_space}.
\end{proof}

% from §31 Item 3: regular neighborhood closure property
\begin{lemma}\label{regular_closure_neighborhood}
    Assume $X$ is regular with respect to $T$.
    Then for all $x,U$ such that $x\in X$ and $U$ is a neighborhood of $x$ in $X$ with respect to $T$
        there exists $V$ such that
            $V$ is a neighborhood of $x$ in $X$ with respect to $T$ and
            $\closure{V}{X}{T}\subseteq U$.
\end{lemma}
\begin{proof}
    $T$ is a topology on $X$ by \cref{regular_space}.
    Fix $x$.
    Fix $U$ such that $x\in X$ and $U$ is a neighborhood of $x$ in $X$ with respect to $T$.
    Let $B = X\setminus U$.
    $U$ is open in $X$ with respect to $T$ and $U\subseteq X$
        by \cref{neighborhood,open_in,topology_on,subseteq,pow_iff}.
    Hence $B$ is closed in $X$ with respect to $T$ by \cref{setminus_subseteq,double_relative_complement,subseteq,open_in,closed_in}.
    Take $U_1,V_1$ such that
        $U_1$ is open in $X$ with respect to $T$ and
        $V_1$ is open in $X$ with respect to $T$ and
        $x\in U_1$ and
        $B\subseteq V_1$ and
        $U_1$ is disjoint from $V_1$
        by \cref{setminus_elim_right,neighborhood,regular_space}.
    Let $V = U_1$.
    $V$ is a neighborhood of $x$ in $X$ with respect to $T$ by \cref{open_in,neighborhood}.
    $V$ is open in $X$ with respect to $T$ and $V\subseteq X$
        by \cref{open_in,topology_on,subseteq,pow_iff}.
    Then $V\subseteq X\setminus V_1$ by \cref{disjoint,subseteq,setminus_intro}.
    Therefore $X\setminus V_1$ is closed in $X$ with respect to $T$
        by \cref{open_in,topology_on,subseteq,pow_iff,setminus_subseteq,double_relative_complement,closed_in}.
    Hence $\closure{V}{X}{T}\subseteq X\setminus V_1$ by \cref{closure_subseteq_closed}.
    Therefore $X\setminus V_1\subseteq X\setminus B$ by \cref{subseteq_implies_setminus_supseteq,subseteq}.
    Finally $X\setminus V_1\subseteq U$ and $\closure{V}{X}{T}\subseteq U$
        by \cref{subseteq_transitive,double_relative_complement,subseteq}.
\end{proof}

% from §31 Item 4: regularity from neighborhood closure property
\begin{lemma}\label{regular_from_closure_neighborhood}
    Suppose $T$ is a topology on $X$.
    Suppose for all $x$ such that $x\in X$ we have $\{x\}$ is closed in $X$ with respect to $T$.
    Suppose for all $x,U$ such that $x\in X$ and $U$ is a neighborhood of $x$ in $X$ with respect to $T$
        there exists $V$ such that
            $V$ is a neighborhood of $x$ in $X$ with respect to $T$ and
            $\closure{V}{X}{T}\subseteq U$.
    Then $X$ is regular with respect to $T$.
\end{lemma}
\begin{proof}
    Show that for all $x,B$ such that $x\in X$ and $B$ is closed in $X$ with respect to $T$ and $x\notin B$
        there exist $U,V$ such that
            $U$ is open in $X$ with respect to $T$ and
            $V$ is open in $X$ with respect to $T$ and
            $x\in U$ and
            $B\subseteq V$ and
            $U$ is disjoint from $V$.
    \begin{subproof}
        Fix $x$.
        Fix $B$ such that $x\in X$ and $B$ is closed in $X$ with respect to $T$ and $x\notin B$.
        Let $U = X\setminus B$.
        $U$ is a neighborhood of $x$ in $X$ with respect to $T$ by \cref{closed_in,setminus_intro,neighborhood}.
        Take $V$ such that
            $V$ is a neighborhood of $x$ in $X$ with respect to $T$ and
            $\closure{V}{X}{T}\subseteq U$.
        Let $W = X\setminus \closure{V}{X}{T}$.
        $V$ is open in $X$ with respect to $T$ and $V\subseteq X$
            by \cref{neighborhood,open_in,topology_on,subseteq,pow_iff}.
        Hence $W$ is open in $X$ with respect to $T$ and $B\subseteq X$ by \cref{closure_closed_in,closed_in}.
        Therefore $B\subseteq W$ by \cref{subseteq_implies_setminus_supseteq,double_relative_complement,subseteq}.
        Now $V\subseteq \closure{V}{X}{T}$ by \cref{subseteq_closure}.
        Hence $V$ is disjoint from $W$ by \cref{subseteq,setminus_elim_right,disjoint}.
        Hence there exist $U,V$ such that
            $U$ is open in $X$ with respect to $T$ and
            $V$ is open in $X$ with respect to $T$ and
            $x\in U$ and
            $B\subseteq V$ and
            $U$ is disjoint from $V$.
    \end{subproof}
    Therefore $X$ is regular with respect to $T$ by \cref{one_point_closed,regular_space}.
\end{proof}

% from §31 Item 5: normal neighborhood closure property
\begin{lemma}\label{normal_closure_neighborhood}
    Assume $X$ is normal with respect to $T$.
    Then for all $A,U$ such that $A$ is closed in $X$ with respect to $T$ and $U$ is open in $X$ with respect to $T$
        and $A\subseteq U$
        there exists $V$ such that
            $V$ is open in $X$ with respect to $T$ and
            $A\subseteq V$ and
            $\closure{V}{X}{T}\subseteq U$.
\end{lemma}
\begin{proof}
    Fix $A$.
    Fix $U$ such that $A$ is closed in $X$ with respect to $T$ and $U$ is open in $X$ with respect to $T$ and $A\subseteq U$.
    Let $B = X\setminus U$.
    $U\in T$ and $U\subseteq X$ by \cref{open_in,normal_space,topology_on,subseteq,pow_iff}.
    Hence $B$ is closed in $X$ with respect to $T$ by \cref{setminus_subseteq,double_relative_complement,subseteq,open_in,closed_in}.
    Then $A$ is disjoint from $B$ by \cref{disjoint,subseteq,setminus_elim_right}.
    Take $V,W$ such that
        $V$ is open in $X$ with respect to $T$ and
        $W$ is open in $X$ with respect to $T$ and
        $A\subseteq V$ and
        $B\subseteq W$ and
        $V$ is disjoint from $W$
        by \cref{normal_space}.

    Thus $V\subseteq X$ by \cref{open_in,normal_space,topology_on,subseteq,pow_iff}.
    Then $V\subseteq X\setminus W$ by \cref{disjoint,subseteq,setminus_intro}.
    Therefore $X\setminus W$ is closed in $X$ with respect to $T$
        by \cref{open_in,normal_space,topology_on,subseteq,pow_iff,setminus_subseteq,double_relative_complement,closed_in}.
    Hence $\closure{V}{X}{T}\subseteq X\setminus W$ by \cref{normal_space,closure_subseteq_closed}.
    Therefore $X\setminus W\subseteq X\setminus B$ by \cref{subseteq,subseteq_implies_setminus_supseteq}.
    Finally $\closure{V}{X}{T}\subseteq U$ by \cref{subseteq_transitive,double_relative_complement,subseteq}.
    Hence there exists $V$ such that
        $V$ is open in $X$ with respect to $T$ and
        $A\subseteq V$ and
        $\closure{V}{X}{T}\subseteq U$.
\end{proof}
% from §31 Item 6: normality from neighborhood closure property
\begin{lemma}\label{normal_from_closure_neighborhood}
    Suppose $T$ is a topology on $X$.
    Suppose for all $x$ such that $x\in X$ we have $\{x\}$ is closed in $X$ with respect to $T$.
    Suppose for all $A,U$ such that $A$ is closed in $X$ with respect to $T$ and $U$ is open in $X$ with respect to $T$
        and $A\subseteq U$
        there exists $V$ such that
            $V$ is open in $X$ with respect to $T$ and
            $A\subseteq V$ and
            $\closure{V}{X}{T}\subseteq U$.
    Then $X$ is normal with respect to $T$.
\end{lemma}
\begin{proof}
    Show that for all $A,B$ such that $A$ is closed in $X$ with respect to $T$ and $B$ is closed in $X$ with respect to $T$
        and $A$ is disjoint from $B$
        there exist $U,V$ such that
            $U$ is open in $X$ with respect to $T$ and
            $V$ is open in $X$ with respect to $T$ and
            $A\subseteq U$ and
            $B\subseteq V$ and
            $U$ is disjoint from $V$.
    \begin{subproof}
        Fix $A$.
        Fix $B$ such that $A$ is closed in $X$ with respect to $T$ and $B$ is closed in $X$ with respect to $T$
            and $A$ is disjoint from $B$.
        Let $U = X\setminus B$.
        Hence $A\subseteq U$ by \cref{disjoint,closed_in,subseteq,setminus_intro}.
        $U$ is open in $X$ with respect to $T$ by \cref{closed_in}.
        Take $V$ such that
            $V$ is open in $X$ with respect to $T$ and
            $A\subseteq V$ and
            $\closure{V}{X}{T}\subseteq U$.
        Let $W = X\setminus \closure{V}{X}{T}$.
        $V\subseteq X$ by \cref{open_in,topology_on,subseteq,pow_iff}.
        $\closure{V}{X}{T}$ is closed in $X$ with respect to $T$
            by \cref{open_in,topology_on,subseteq,pow_iff,closure_closed_in}.
        Then $W$ is open in $X$ with respect to $T$ and $B\subseteq X$ by \cref{closed_in}.
        Therefore $B\subseteq W$
            by \cref{subseteq,double_relative_complement,subseteq_implies_setminus_supseteq}.
        Now $V\subseteq \closure{V}{X}{T}$ by \cref{subseteq_closure}.
        Hence $V$ is disjoint from $W$ by \cref{subseteq,setminus_elim_right,disjoint}.
        Hence there exist $U,V$ such that
            $U$ is open in $X$ with respect to $T$ and
            $V$ is open in $X$ with respect to $T$ and
            $A\subseteq U$ and
            $B\subseteq V$ and
            $U$ is disjoint from $V$.
    \end{subproof}
    Therefore $X$ is normal with respect to $T$ by \cref{one_point_closed,normal_space}.
\end{proof}

% from §31 Item 7: subspace of a Hausdorff space is Hausdorff
\begin{lemma}\label{subspace_hausdorff_31}
    Assume $X$ is Hausdorff with respect to $T$.
    Assume $Y$ is a subspace of $X$ with respect to $T$.
    Let $T_Y = \subspaceTopology{T}{Y}$.
    Then $Y$ is Hausdorff with respect to $T_Y$.
\end{lemma}
\begin{proof}
    Follows by \cref{subspace_of,subspace_hausdorff}.
\end{proof}

% from §31 Item 8: product of Hausdorff spaces is Hausdorff
\begin{lemma}\label{product_hausdorff_31}
    Assume $X$ is a function.
    Assume $T$ is a function.
    Assume $\dom{T} = \dom{X}$.
    Assume $\dom{X}$ is inhabited.
    Suppose for all $\alpha\in\dom{X}$ we have $\apply{X}{\alpha}$ is Hausdorff with respect to $\apply{T}{\alpha}$.
    Then $\productFamily{X}$ is Hausdorff with respect to $\productTopologyIndexed{T}{X}$.
\end{lemma}
\begin{proof}
    Follows by \cref{product_hausdorff_product}.
\end{proof}

% from §31 Item 9: subspace of a regular space is regular
\begin{lemma}\label{subspace_regular}
    Assume $X$ is regular with respect to $T$.
    Assume $Y$ is a subspace of $X$ with respect to $T$.
    Let $T_Y = \subspaceTopology{T}{Y}$.
    Then $Y$ is regular with respect to $T_Y$.
\end{lemma}
    \begin{proof}
    $T$ is a topology on $X$ and $Y\subseteq X$ by \cref{regular_space,subspace_of}.
    Hence $T_Y$ is a topology on $Y$ and $Y$ has closed singletons with respect to $T_Y$
        by \cref{subspace_topology_is_topology,regular_space,elem_subseteq,one_point_closed,singleton_subset_intro,inter_absorb_supseteq_right,closed_in_subspace_iff,subspace_of}.
    Show that for all $x,B$ such that $x\in Y$ and $B$ is closed in $Y$ with respect to $T_Y$ and $x\notin B$
        there exist $U,V$ such that
            $U$ is open in $Y$ with respect to $T_Y$ and
            $V$ is open in $Y$ with respect to $T_Y$ and
            $x\in U$ and
            $B\subseteq V$ and
            $U$ is disjoint from $V$.
    \begin{subproof}
        Fix $x$.
        Fix $B$ such that $x\in Y$ and $B$ is closed in $Y$ with respect to $T_Y$ and $x\notin B$.
        Take $C$ such that $C$ is closed in $X$ with respect to $T$ and $B = C\inter Y$
            by \cref{closed_in_subspace_iff,subspace_of}.
        Take $U,V$ such that
            $U$ is open in $X$ with respect to $T$ and
            $V$ is open in $X$ with respect to $T$ and
            $x\in U$ and
            $C\subseteq V$ and
            $U$ is disjoint from $V$
            by \cref{inter_intro,elem_subseteq,regular_space}.
        Let $U_1 = U\inter Y$.
        Let $V_1 = V\inter Y$.
        Hence $U_1$ is open in $Y$ with respect to $T_Y$ and $V_1$ is open in $Y$ with respect to $T_Y$
            and $x\in U_1$ by \cref{open_in,inter_comm,subspace_topology,inter_intro}.
        Therefore $B\subseteq V_1$ and $U_1$ is disjoint from $V_1$
            by \cref{inter_lower_left,subseteq_transitive,inter_lower_right,subseteq_inter_iff,inter,disjoint}.
    \end{subproof}
    Hence $Y$ is regular with respect to $T_Y$ by \cref{regular_space}.
\end{proof}

% from §31 Item 10: product of regular spaces is regular
\begin{lemma}\label{product_regular}
    Assume $X$ is a function.
    Assume $T$ is a function.
    Assume $\dom{T} = \dom{X}$.
    Assume $\dom{X}$ is inhabited.
    Suppose for all $\alpha\in\dom{X}$ we have $\apply{X}{\alpha}$ is regular with respect to $\apply{T}{\alpha}$.
    Then $\productFamily{X}$ is regular with respect to $\productTopologyIndexed{T}{X}$.
\end{lemma}
\begin{proof}
    Let $T_0 = \productTopologyIndexed{T}{X}$.
    Let $S = \productSubbasis{T}{X}$.
    For all $\alpha\in\dom{X}$ we have $\apply{T}{\alpha}$ is a topology on $\apply{X}{\alpha}$ by \cref{regular_space}.
    Hence for all $\alpha\in\dom{X}$ we have $\apply{X}{\alpha}$ is Hausdorff with respect to $\apply{T}{\alpha}$
        by \cref{regular_implies_hausdorff}.
    Therefore $\productFamily{X}$ is Hausdorff with respect to $T_0$ by \cref{product_hausdorff_31}.
    $T_0$ is a topology on $\productFamily{X}$ by \cref{hausdorff}.
    For all $x$ such that $x\in\productFamily{X}$ we have $\{x\}$ is closed in $\productFamily{X}$ with respect to $T_0$
        by \cref{singleton_subset_intro,pair_is_finite,upair_intro_left,subseteq_finite_is_finite,finite_point_set_closed_hausdorff}.
    Show that for all $x,U$ such that $x\in\productFamily{X}$ and
        $U$ is a neighborhood of $x$ in $\productFamily{X}$ with respect to $T_0$
        there exists $V$ such that
            $V$ is a neighborhood of $x$ in $\productFamily{X}$ with respect to $T_0$ and
            $\closure{V}{\productFamily{X}}{T_0}\subseteq U$.
    \begin{subproof}
        Fix $x$.
        Fix $U$.
        Assume $x\in\productFamily{X}$ and
            $U$ is a neighborhood of $x$ in $\productFamily{X}$ with respect to $T_0$.
        Then $U\in T_0$ and $x\in U$ by \cref{neighborhood,open_in}.
        $\finiteIntersections{S}$ is a basis for $T_0$ on $\productFamily{X}$ by \cref{product_subbasis_finite_intersections_basis}.
        Hence $T_0 = \basisTopology{\finiteIntersections{S}}{\productFamily{X}}$
            by \cref{product_topology_indexed,topology_generated_by_subbasis,basis_for_topology}.
        Therefore $U\in\basisTopology{\finiteIntersections{S}}{\productFamily{X}}$ by \cref{basis_for_topology}.
        Take $B_1\in\finiteIntersections{S}$ such that $x\in B_1$ and $B_1\subseteq U$ by \cref{topology_generated_by_basis}.
        Take $F\subseteq S$ such that $F$ is finite and inhabited and $B_1 = \inters{F}$ by \cref{finite_intersections}.
        Let $R = \{w\in F\times\pow{\productFamily{X}} \mid \exists \beta\in\dom{X}. \exists U_0\in\apply{T}{\beta}. \exists V_0\in\apply{T}{\beta}.
            \fst{w} = \preimg{\piIndex{X}{\beta}}{U_0} \land
            \apply{x}{\beta}\in V_0 \land
            \closure{V_0}{\apply{X}{\beta}}{\apply{T}{\beta}}\subseteq U_0 \land
            \snd{w} = \preimg{\piIndex{X}{\beta}}{V_0}\}$.
        $R$ is a relation by \cref{subseteq,times_elem_is_tuple,relation}.
        Show that for all $A\in F$ there exists $B$ such that $A\mathrel{R} B$.
        \begin{subproof}
            Fix $A$ such that $A\in F$.
            Take $\beta\in\dom{X}$ such that $A\in\productSubbasisAt{T}{X}{\beta}$ by \cref{subseteq,product_subbasis_union_indexed}.
            Take $U_0\in\apply{T}{\beta}$ such that $A = \preimg{\piIndex{X}{\beta}}{U_0}$ by \cref{product_subbasis_indexed}.
            Hence $x\in A$ by \cref{inters_subseteq_elem,subseteq}.
            Take $y\in U_0$ such that $x\mathrel{\piIndex{X}{\beta}} y$ by \cref{preimg_iff}.
            Take $x_0\in\productFamily{X}$ such that $(x,y) = (x_0,\apply{x_0}{\beta})$ by \cref{projection_map_indexed}.
            Then $\apply{x}{\beta}\in U_0$ by \cref{pair_eq_iff}.
            Take $V_0$ such that
                $V_0$ is a neighborhood of $\apply{x}{\beta}$ in $\apply{X}{\beta}$ with respect to $\apply{T}{\beta}$ and
                $\closure{V_0}{\apply{X}{\beta}}{\apply{T}{\beta}}\subseteq U_0$
                by \cref{open_in,neighborhood,indexed_product,regular_space,regular_closure_neighborhood}.
            Hence $V_0\in\apply{T}{\beta}$ and $\apply{x}{\beta}\in V_0$ by \cref{neighborhood,open_in}.
            Let $B = \preimg{\piIndex{X}{\beta}}{V_0}$.
            Let $w = (A,B)$.
            Therefore $w\in F\times\pow{\productFamily{X}}$
                by \cref{projection_map_indexed_function,funs_is_relation,preimg_subseteq_dom,funs,subseteq_transitive,powerset_intro,times_tuple_intro}.
            Thus $\fst{w} = \preimg{\piIndex{X}{\beta}}{U_0}$ and $\snd{w} = \preimg{\piIndex{X}{\beta}}{V_0}$
                by \cref{fst_eq,snd_eq}.
            Hence $w\in R$.
            Therefore $A\mathrel{R}B$ by \cref{pair_eq_iff}.
        \end{subproof}
        Take $g$ such that $g$ is a function on $F$ and for all $A\in F$ we have $A\mathrel{R}\apply{g}{A}$
            by \cref{choice_function}.
        Let $F_1 = \img{g}{F}$.
        Hence $F_1$ is finite and inhabited by \cref{finite_image_function,function_apply_elim,img_iff}.
        Show that for all $B$ such that $B\in F_1$ we have $B\in T_0$.
        \begin{subproof}
            Fix $B$ such that $B\in F_1$.
            Take $A$ such that $A\in F$ and $A\mathrel{R} B$ by \cref{img_iff,function_apply_iff,choice_function}.
            Take $\beta\in\dom{X}$ such that there exists $U_0\in\apply{T}{\beta}$ such that there exists $V_0\in\apply{T}{\beta}$ such that
                $A = \preimg{\piIndex{X}{\beta}}{U_0}$ and
                $\apply{x}{\beta}\in V_0$ and
                $\closure{V_0}{\apply{X}{\beta}}{\apply{T}{\beta}}\subseteq U_0$ and
                $B = \preimg{\piIndex{X}{\beta}}{V_0}$.
            Take $U_0\in\apply{T}{\beta}$ such that there exists $V_0\in\apply{T}{\beta}$ such that
                $A = \preimg{\piIndex{X}{\beta}}{U_0}$ and
                $\apply{x}{\beta}\in V_0$ and
                $\closure{V_0}{\apply{X}{\beta}}{\apply{T}{\beta}}\subseteq U_0$ and
                $B = \preimg{\piIndex{X}{\beta}}{V_0}$.
            Take $V_0\in\apply{T}{\beta}$ such that
                $A = \preimg{\piIndex{X}{\beta}}{U_0}$ and
                $\apply{x}{\beta}\in V_0$ and
                $\closure{V_0}{\apply{X}{\beta}}{\apply{T}{\beta}}\subseteq U_0$ and
                $B = \preimg{\piIndex{X}{\beta}}{V_0}$.
            Therefore $B\in T_0$ by \cref{projection_map_indexed_continuous,continuous,open_in}.
        \end{subproof}
        Hence $F_1\subseteq T_0$ by \cref{subseteq}.
        Let $V = \inters{F_1}$.
        Therefore $V\in T_0$ and $V$ is open in $\productFamily{X}$ with respect to $T_0$
            by \cref{finite_intersections_open_in_topology,open_in}.
        Show that for all $B$ such that $B\in F_1$ we have $x\in B$.
        \begin{subproof}
            Fix $B$ such that $B\in F_1$.
            Take $A$ such that $A\in F$ and $A\mathrel{R} B$ by \cref{img_iff,function_apply_iff,choice_function}.
            Take $\beta\in\dom{X}$ such that there exists $U_0$ such that there exists $V_0$ such that
                $U_0\in\apply{T}{\beta}$ and
                $V_0\in\apply{T}{\beta}$ and
                $A = \preimg{\piIndex{X}{\beta}}{U_0}$ and
                $\apply{x}{\beta}\in V_0$ and
                $\closure{V_0}{\apply{X}{\beta}}{\apply{T}{\beta}}\subseteq U_0$ and
                $B = \preimg{\piIndex{X}{\beta}}{V_0}$.
            Therefore $x\in B$ by \cref{projection_map_indexed,preimg_iff}.
        \end{subproof}
        Hence $x\in V$ by \cref{inters_iff_forall}.
        Therefore $V$ is a neighborhood of $x$ in $\productFamily{X}$ with respect to $T_0$ by \cref{neighborhood}.
        Show that for all $A\in F$ we have $\closure{V}{\productFamily{X}}{T_0}\subseteq A$.
        \begin{subproof}
            Fix $A$ such that $A\in F$.
            Let $B = \apply{g}{A}$.
            Hence $A\mathrel{R} B$ by \cref{choice_function}.
            Take $\beta\in\dom{X}$ such that there exists $U_0\in\apply{T}{\beta}$ such that there exists $V_0\in\apply{T}{\beta}$ such that
                $A = \preimg{\piIndex{X}{\beta}}{U_0}$ and
                $\apply{x}{\beta}\in V_0$ and
                $\closure{V_0}{\apply{X}{\beta}}{\apply{T}{\beta}}\subseteq U_0$ and
                $B = \preimg{\piIndex{X}{\beta}}{V_0}$.
            Take $U_0\in\apply{T}{\beta}$ such that there exists $V_0\in\apply{T}{\beta}$ such that
                $A = \preimg{\piIndex{X}{\beta}}{U_0}$ and
                $\apply{x}{\beta}\in V_0$ and
                $\closure{V_0}{\apply{X}{\beta}}{\apply{T}{\beta}}\subseteq U_0$ and
                $B = \preimg{\piIndex{X}{\beta}}{V_0}$.
            Take $V_0\in\apply{T}{\beta}$ such that
                $A = \preimg{\piIndex{X}{\beta}}{U_0}$ and
                $\apply{x}{\beta}\in V_0$ and
                $\closure{V_0}{\apply{X}{\beta}}{\apply{T}{\beta}}\subseteq U_0$ and
                $B = \preimg{\piIndex{X}{\beta}}{V_0}$.
            Hence $B\in F_1$ and $V\subseteq B$ by \cref{img_iff,function_apply_elim,inters_subseteq_elem}.
            Now $V_0\subseteq\apply{X}{\beta}$ by \cref{regular_space,topology_on,subseteq,powerset_elim}.
            $\closure{V_0}{\apply{X}{\beta}}{\apply{T}{\beta}}$ is closed in $\apply{X}{\beta}$ with respect to $\apply{T}{\beta}$
                by \cref{regular_space,topology_on,subseteq,powerset_elim,closure_closed_in}.
            Let $D = \apply{X}{\beta}\setminus \closure{V_0}{\apply{X}{\beta}}{\apply{T}{\beta}}$.
            Therefore $\preimg{\piIndex{X}{\beta}}{D}$ is open in $\productFamily{X}$ with respect to $T_0$
                by \cref{closed_in,projection_map_indexed_continuous,continuous}.
            Then $\closure{V_0}{\apply{X}{\beta}}{\apply{T}{\beta}}\subseteq \apply{X}{\beta}$
                by \cref{closed_whole_space,closure_subseteq_closed}.
            Therefore $\preimg{\piIndex{X}{\beta}}{D} = \productFamily{X}\setminus \preimg{\piIndex{X}{\beta}}{\closure{V_0}{\apply{X}{\beta}}{\apply{T}{\beta}}}$
                by \cref{projection_map_indexed_function,preimg_setminus_function}.
            Hence $\productFamily{X}\setminus \preimg{\piIndex{X}{\beta}}{\closure{V_0}{\apply{X}{\beta}}{\apply{T}{\beta}}}$
                is open in $\productFamily{X}$ with respect to $T_0$ by \cref{open_in}.
            Therefore $\preimg{\piIndex{X}{\beta}}{\closure{V_0}{\apply{X}{\beta}}{\apply{T}{\beta}}}$
                is closed in $\productFamily{X}$ with respect to $T_0$
                by \cref{projection_map_indexed_function,funs_is_relation,preimg_subseteq_dom,funs,subseteq_transitive,closed_in}.
            Then $B\subseteq \preimg{\piIndex{X}{\beta}}{\closure{V_0}{\apply{X}{\beta}}{\apply{T}{\beta}}}$
                by \cref{subseteq_closure,preimg_subseteq}.
            Hence $V\subseteq \preimg{\piIndex{X}{\beta}}{\closure{V_0}{\apply{X}{\beta}}{\apply{T}{\beta}}}$
                by \cref{subseteq_transitive}.
            Therefore $\closure{V}{\productFamily{X}}{T_0}\subseteq
                \preimg{\piIndex{X}{\beta}}{\closure{V_0}{\apply{X}{\beta}}{\apply{T}{\beta}}}$
                by \cref{hausdorff,topology_on,subseteq,powerset_elim,closure_subseteq_closed}.
            Then $\closure{V}{\productFamily{X}}{T_0}\subseteq \preimg{\piIndex{X}{\beta}}{U_0}$
                by \cref{preimg_subseteq,subseteq_transitive}.
            Hence $\closure{V}{\productFamily{X}}{T_0}\subseteq A$ by \cref{subseteq}.
        \end{subproof}
        Therefore $\closure{V}{\productFamily{X}}{T_0}\subseteq U$ by \cref{subseteq_inters_iff,subseteq,subseteq_transitive}.
    \end{subproof}
    Finally $\productFamily{X}$ is regular with respect to $T_0$ by \cref{regular_from_closure_neighborhood}.
\end{proof}

\section{§32 Normal Spaces}

% from §32 Item 1: regular space with countable basis is normal
\begin{theorem}\label{regular_second_countable_normal}
    Assume $X$ is regular with respect to $T$.
    Suppose $X$ satisfies the second countability axiom with respect to $T$.
    Then $X$ is normal with respect to $T$.
\end{theorem}
\begin{proof}
    $T$ is a topology on $X$ by \cref{regular_space}.
    Take $b$ such that
        $b$ is a sequence in $\pow{X}$ and
        for all $n\in\naturalsPlus$ we have $b(n)$ is open in $X$ with respect to $T$ and
        for all $x\in X$ we have
            for all $U$ such that $U$ is open in $X$ with respect to $T$ and $x\in U$
                there exists $n\in\naturalsPlus$ such that $x\in b(n)$ and $b(n)\subseteq U$
        by \cref{second_countability_axiom}.
    Show that for all $A,B$ such that $A$ is closed in $X$ with respect to $T$ and $B$ is closed in $X$ with respect to $T$
        and $A$ is disjoint from $B$
        there exist $U,V$ such that
            $U$ is open in $X$ with respect to $T$ and
            $V$ is open in $X$ with respect to $T$ and
            $A\subseteq U$ and
            $B\subseteq V$ and
            $U$ is disjoint from $V$.
    \begin{subproof}
        Fix $A$.
        Fix $B$.
        Assume $A$ is closed in $X$ with respect to $T$ and $B$ is closed in $X$ with respect to $T$
            and $A$ is disjoint from $B$.
        Let $P = \{n\in\naturalsPlus \mid b(n)\inter A \neq \emptyset \land \closure{b(n)}{X}{T}\inter B = \emptyset\}$.
        Let $Q = \{n\in\naturalsPlus \mid b(n)\inter B \neq \emptyset \land \closure{b(n)}{X}{T}\inter A = \emptyset\}$.
        Let $U_1 = \{w \mid \exists n\in P. w = (n,b(n))\}$. Let $U_2 = \{w \mid \exists n\in\naturalsPlus\setminus P. w = (n,\emptyset)\}$.
        Let $U = U_1 \union U_2$.
        Let $V_1 = \{w \mid \exists n\in Q. w = (n,b(n))\}$. Let $V_2 = \{w \mid \exists n\in\naturalsPlus\setminus Q. w = (n,\emptyset)\}$.
        Let $V = V_1 \union V_2$.

        Show that $U$ is a function.
        \begin{subproof}
                For all $w$ such that $w\in U$ there exist $x,y$ such that $w = (x,y)$.
                Hence $U$ is a relation by \cref{relation}.
                Show that for all $n,y_1,y_2$ such that $(n,y_1)\in U$ and $(n,y_2)\in U$ we have $y_1 = y_2$.
                \begin{subproof}
                    Fix $n$.
                    Fix $y_1$.
                    Fix $y_2$ such that $(n,y_1)\in U$ and $(n,y_2)\in U$.
                    \begin{byCase}
                    \caseOf{$n\in P$.}
                        Then $(n,y_1)\in U_1$ by \cref{union_iff,pair_eq_iff,setminus_elim_right}.
                        Take $n_1\in P$ such that $(n,y_1) = (n_1,b(n_1))$.
                        Then $(n,y_2)\in U_1$ by \cref{union_iff,pair_eq_iff,setminus_elim_right}.
                        Take $n_2\in P$ such that $(n,y_2) = (n_2,b(n_2))$.
                        Hence $y_1 = y_2$ by \cref{pair_eq_iff}.
                    \caseOf{$n\notin P$.}
                        Then $(n,y_1)\in U_2$ by \cref{union_iff,pair_eq_iff}.
                        Take $n_1\in\naturalsPlus\setminus P$ such that $(n,y_1) = (n_1,\emptyset)$.
                        Then $(n,y_2)\in U_2$ by \cref{union_iff,pair_eq_iff}.
                        Take $n_2\in\naturalsPlus\setminus P$ such that $(n,y_2) = (n_2,\emptyset)$.
                        Hence $y_1 = y_2$ by \cref{pair_eq_iff}.
                    \end{byCase}
                \end{subproof}
                Therefore $U$ is right-unique by \cref{rightunique}.
        \end{subproof}
        Hence $\dom{U}\subseteq\naturalsPlus$
            by \cref{dom_iff,union_iff,pair_eq_iff,setminus_elim_left,subseteq}.
        Therefore $\naturalsPlus\subseteq\dom{U}$
            by \cref{union_intro_left,setminus_intro,union_intro_right,dom_intro,subseteq}.
        Hence $\dom{U} = \naturalsPlus$ by \cref{subseteq_antisymmetric}.
        Therefore $U$ is a function on $\naturalsPlus$.

        Show that for all $n\in\dom{U}$ we have $U(n)\in\pow{X}$.
        \begin{subproof}
                Fix $n\in\dom{U}$.
                $\dom{U} = \naturalsPlus$.
                Hence $n\in\naturalsPlus$ by \cref{subseteq}.
                \begin{byCase}
                \caseOf{$n\in P$.}
                    $(n,U(n))\in U$ by \cref{function_apply_elim}.
                    Then $(n,U(n))\in U_1$ by \cref{union_iff,pair_eq_iff,setminus_elim_right}.
                    Take $n_0\in P$ such that $(n,U(n)) = (n_0,b(n_0))$.
                    Then $U(n)\in\pow{X}$ by \cref{pair_eq_iff,sequence_in,funs_type_apply}.
                \caseOf{$n\notin P$.}
                    $(n,U(n))\in U$ by \cref{function_apply_elim}.
                    Then $(n,U(n))\in U_2$ by \cref{union_iff,pair_eq_iff}.
                    Take $n_0\in\naturalsPlus\setminus P$ such that $(n,U(n)) = (n_0,\emptyset)$.
                    Then $U(n)\in\pow{X}$ by \cref{pair_eq_iff,emptyset_subseteq,powerset_intro}.
                \end{byCase}
        \end{subproof}
        Hence $U$ is a function to $\pow{X}$.
        Therefore $U\in\funs{\naturalsPlus}{\pow{X}}$ by \cref{funs_intro}.

        Show that $V$ is a function.
        \begin{subproof}
                For all $w$ such that $w\in V$ there exist $x,y$ such that $w = (x,y)$.
                Hence $V$ is a relation by \cref{relation}.
                Show that for all $n,y_1,y_2$ such that $(n,y_1)\in V$ and $(n,y_2)\in V$ we have $y_1 = y_2$.
                \begin{subproof}
                    Fix $n$.
                    Fix $y_1$.
                    Fix $y_2$ such that $(n,y_1)\in V$ and $(n,y_2)\in V$.
                    \begin{byCase}
                    \caseOf{$n\in Q$.}
                        Then $(n,y_1)\in V_1$ by \cref{union_iff,pair_eq_iff,setminus_elim_right}.
                        Take $n_1\in Q$ such that $(n,y_1) = (n_1,b(n_1))$.
                        Then $(n,y_2)\in V_1$ by \cref{union_iff,pair_eq_iff,setminus_elim_right}.
                        Take $n_2\in Q$ such that $(n,y_2) = (n_2,b(n_2))$.
                        Hence $y_1 = y_2$ by \cref{pair_eq_iff}.
                    \caseOf{$n\notin Q$.}
                        Then $(n,y_1)\in V_2$ by \cref{union_iff,pair_eq_iff}.
                        Take $n_1\in\naturalsPlus\setminus Q$ such that $(n,y_1) = (n_1,\emptyset)$.
                        Then $(n,y_2)\in V_2$ by \cref{union_iff,pair_eq_iff}.
                        Take $n_2\in\naturalsPlus\setminus Q$ such that $(n,y_2) = (n_2,\emptyset)$.
                        Hence $y_1 = y_2$ by \cref{pair_eq_iff}.
                    \end{byCase}
                \end{subproof}
                Therefore $V$ is right-unique by \cref{rightunique}.
        \end{subproof}
        For all $n\in\dom{V}$ we have $n\in\naturalsPlus$
            by \cref{dom_iff,union_iff,pair_eq_iff,setminus_elim_left}.
        Hence $\dom{V}\subseteq\naturalsPlus$ by \cref{subseteq}.
        Therefore $\naturalsPlus\subseteq\dom{V}$
            by \cref{union_intro_left,setminus_intro,union_intro_right,dom_intro,subseteq}.
        Hence $\dom{V} = \naturalsPlus$ by \cref{subseteq_antisymmetric}.
        Therefore $V$ is a function on $\naturalsPlus$.

        Show that for all $n\in\dom{V}$ we have $V(n)\in\pow{X}$.
        \begin{subproof}
                Fix $n\in\dom{V}$.
                $\dom{V} = \naturalsPlus$.
                Hence $n\in\naturalsPlus$ by \cref{subseteq}.
                \begin{byCase}
                \caseOf{$n\in Q$.}
                    $(n,V(n))\in V$ by \cref{function_apply_elim}.
                    Then $(n,V(n))\in V_1$ by \cref{union_iff,pair_eq_iff,setminus_elim_right}.
                    Take $n_0\in Q$ such that $(n,V(n)) = (n_0,b(n_0))$.
                    Then $V(n)\in\pow{X}$ by \cref{pair_eq_iff,sequence_in,funs_type_apply}.
                \caseOf{$n\notin Q$.}
                    $(n,V(n))\in V$ by \cref{function_apply_elim}.
                    Then $(n,V(n))\in V_2$ by \cref{union_iff,pair_eq_iff}.
                    Take $n_0\in\naturalsPlus\setminus Q$ such that $(n,V(n)) = (n_0,\emptyset)$.
                    Then $V(n)\in\pow{X}$ by \cref{pair_eq_iff,emptyset_subseteq,powerset_intro}.
                \end{byCase}
        \end{subproof}
        Hence $V$ is a function to $\pow{X}$.
        Therefore $V\in\funs{\naturalsPlus}{\pow{X}}$ by \cref{funs_intro}.

        Show that for all $n\in\naturalsPlus$ we have $U(n)$ is open in $X$ with respect to $T$ and
            $\closure{U(n)}{X}{T}$ is disjoint from $B$.
        \begin{subproof}
            Fix $n\in\naturalsPlus$.
            \begin{byCase}
            \caseOf{$n\in P$.}
                $b(n)$ is open in $X$ with respect to $T$ by \cref{second_countability_axiom}.
                Then $U(n) = b(n)$ by \cref{union_intro_left,function_apply_elim,function_rightunique}.
                Hence $U(n)$ is open in $X$ with respect to $T$.
                Then $\closure{U(n)}{X}{T}$ is disjoint from $B$ by \cref{inter,emptyset,disjoint}.
            \caseOf{$n\notin P$.}
                Then $U(n) = \emptyset$ by \cref{setminus_intro,union_intro_right,function_apply_elim,function_rightunique}.
                Hence $U(n)$ is open in $X$ with respect to $T$ by \cref{rightunique,topology_on,open_in}.
                Hence $\closure{U(n)}{X}{T} = \emptyset$
                    by \cref{closed_emptyset,emptyset_subseteq,subseteq_emptyset_iff,closure_subseteq_closed}.
                Hence $\closure{U(n)}{X}{T}$ is disjoint from $B$ by \cref{emptyset,disjoint}.
            \end{byCase}
        \end{subproof}

        Show that for all $n\in\naturalsPlus$ we have $V(n)$ is open in $X$ with respect to $T$ and
            $\closure{V(n)}{X}{T}$ is disjoint from $A$.
        \begin{subproof}
            Fix $n\in\naturalsPlus$.
            \begin{byCase}
            \caseOf{$n\in Q$.}
                $b(n)$ is open in $X$ with respect to $T$ by \cref{second_countability_axiom}.
                Then $V(n) = b(n)$ by \cref{union_intro_left,function_apply_elim,function_rightunique}.
                Hence $V(n)$ is open in $X$ with respect to $T$.
                Then $\closure{V(n)}{X}{T}$ is disjoint from $A$ by \cref{inter,emptyset,disjoint}.
            \caseOf{$n\notin Q$.}
                Then $V(n) = \emptyset$ by \cref{setminus_intro,union_intro_right,function_apply_elim,function_rightunique}.
                Hence $V(n)$ is open in $X$ with respect to $T$ by \cref{rightunique,topology_on,open_in}.
                Hence $\closure{V(n)}{X}{T} = \emptyset$
                    by \cref{closed_emptyset,emptyset_subseteq,subseteq_emptyset_iff,closure_subseteq_closed}.
                Hence $\closure{V(n)}{X}{T}$ is disjoint from $A$ by \cref{emptyset,disjoint}.
            \end{byCase}
        \end{subproof}

        Show that for all $x$ such that $x\in A$ we have $x\in\unions{\img{U}{\naturalsPlus}}$.
        \begin{subproof}
                Fix $x$ such that $x\in A$.
                Take $U_0,V_0$ such that
                    $U_0$ is open in $X$ with respect to $T$ and
                    $V_0$ is open in $X$ with respect to $T$ and
                    $x\in U_0$ and
                    $B\subseteq V_0$ and
                    $U_0$ is disjoint from $V_0$
                    by \cref{closed_in,subseteq,disjoint,regular_space}.
                Thus $x\in X$ by \cref{closed_in,subseteq}.
                Take $n\in\naturalsPlus$ such that $x\in b(n)$ and $b(n)\subseteq U_0$.
                Thus $b(n)\inter A\neq\emptyset$ by \cref{inter_intro,emptyset}.
                Thus $X\setminus V_0$ is closed in $X$ with respect to $T$
                    by \cref{open_in,topology_on,subseteq,powerset_elim,setminus_subseteq,double_relative_complement,closed_in}.
                Thus $U_0\subseteq X\setminus V_0$ by \cref{disjoint,subseteq,open_in,topology_on,powerset_elim,setminus_intro,subseteq}.
                Thus $\closure{b(n)}{X}{T}\subseteq X\setminus V_0$
                    by \cref{subseteq_transitive,sequence_in,funs_type_apply,powerset_elim,closure_subseteq_closed}.
                Thus $\closure{b(n)}{X}{T}$ is disjoint from $B$ by \cref{subseteq,setminus_elim_right,disjoint}.
                Show that $\closure{b(n)}{X}{T}\inter B = \emptyset$.
                \begin{subproof}
                    Suppose not.
                    Take $y$ such that $y\in\closure{b(n)}{X}{T}\inter B$ by \cref{nonempty_inhabited}.
                    Contradiction by \cref{inter,disjoint}.
                \end{subproof}
                Thus $n\in P$.
                Thus $(n,b(n))\in U$ by \cref{union_intro_left}.
                $(n,U(n))\in U$ by \cref{function_apply_elim}.
                Thus $x\in U(n)$ by \cref{function_rightunique}.
                Thus $x\in\unions{\img{U}{\naturalsPlus}}$ by \cref{unions_intro,img_iff}.
        \end{subproof}
        Thus $A\subseteq\unions{\img{U}{\naturalsPlus}}$ by \cref{subseteq}.

        Show that for all $x$ such that $x\in B$ we have $x\in\unions{\img{V}{\naturalsPlus}}$.
        \begin{subproof}
                Fix $x$ such that $x\in B$.
                Take $U_0,V_0$ such that
                    $U_0$ is open in $X$ with respect to $T$ and
                    $V_0$ is open in $X$ with respect to $T$ and
                    $x\in U_0$ and
                    $A\subseteq V_0$ and
                    $U_0$ is disjoint from $V_0$
                    by \cref{closed_in,subseteq,disjoint,regular_space}.
                Thus $x\in X$ by \cref{closed_in,subseteq}.
                Take $n\in\naturalsPlus$ such that $x\in b(n)$ and $b(n)\subseteq U_0$.
                Thus $b(n)\inter B\neq\emptyset$ by \cref{inter_intro,emptyset}.
                Thus $X\setminus V_0$ is closed in $X$ with respect to $T$
                    by \cref{open_in,topology_on,subseteq,powerset_elim,setminus_subseteq,double_relative_complement,closed_in}.
                Thus $U_0\subseteq X\setminus V_0$ by \cref{disjoint,subseteq,open_in,topology_on,powerset_elim,setminus_intro,subseteq}.
                Thus $\closure{b(n)}{X}{T}\subseteq X\setminus V_0$
                    by \cref{subseteq_transitive,sequence_in,funs_type_apply,powerset_elim,closure_subseteq_closed}.
                Thus $\closure{b(n)}{X}{T}$ is disjoint from $A$ by \cref{subseteq,setminus_elim_right,disjoint}.
                Show that $\closure{b(n)}{X}{T}\inter A = \emptyset$.
                \begin{subproof}
                    Suppose not.
                    Take $y$ such that $y\in\closure{b(n)}{X}{T}\inter A$ by \cref{nonempty_inhabited}.
                    Contradiction by \cref{inter,disjoint}.
                \end{subproof}
                Thus $n\in Q$.
                Thus $(n,b(n))\in V$ by \cref{union_intro_left}.
                $(n,V(n))\in V$ by \cref{function_apply_elim}.
                Thus $x\in V(n)$ by \cref{function_rightunique}.
                Thus $x\in\unions{\img{V}{\naturalsPlus}}$ by \cref{unions_intro,img_iff}.
        \end{subproof}
        Thus $B\subseteq\unions{\img{V}{\naturalsPlus}}$ by \cref{subseteq}.

        Let $c = \{w \mid \exists n\in\naturalsPlus. w = (n,\closure{V(n)}{X}{T})\}$. Let $d = \{w \mid \exists n\in\naturalsPlus. w = (n,\closure{U(n)}{X}{T})\}$.
        For all $w$ such that $w\in c$ there exist $x,y$ such that $w = (x,y)$.
        For all $n,y_1,y_2$ such that $(n,y_1)\in c$ and $(n,y_2)\in c$ we have $y_1 = y_2$.
        Thus $c$ is a function on $\naturalsPlus$
            by \cref{relation,rightunique,pair_eq_iff,dom_iff,subseteq,dom_intro,subseteq_antisymmetric}.

        For all $w$ such that $w\in d$ there exist $x,y$ such that $w = (x,y)$.
        For all $n,y_1,y_2$ such that $(n,y_1)\in d$ and $(n,y_2)\in d$ we have $y_1 = y_2$.
        Thus $d$ is a function on $\naturalsPlus$
            by \cref{relation,rightunique,pair_eq_iff,dom_iff,subseteq,dom_intro,subseteq_antisymmetric}.

        Let $F = \{W \mid \exists n\in\naturalsPlus. W = U(n)\setminus \unions{\img{\restrl{c}{\initialSegment{n}}}{\initialSegment{n}}}\}$. Let $G = \{W \mid \exists n\in\naturalsPlus. W = V(n)\setminus \unions{\img{\restrl{d}{\initialSegment{n}}}{\initialSegment{n}}}\}$.
        Show that for all $W$ such that $W\in F$ we have $W\in T$.
        \begin{subproof}
            Fix $W$ such that $W\in F$.
                Take $n\in\naturalsPlus$ such that $W = U(n)\setminus \unions{\img{\restrl{c}{\initialSegment{n}}}{\initialSegment{n}}}$.
                $\restrl{c}{\initialSegment{n}}$ is a function by \cref{restrl_of_function_is_function}.
                $\initialSegment{n}$ is finite by \cref{initial_segment_finite}.
                Thus $\dom{\restrl{c}{\initialSegment{n}}} = \initialSegment{n}$
                    by \cref{dom_of_restrl_eq_inter,initial_segment_naturalsplus,subseteq,inter_absorb_supseteq_left}.
                Thus $\img{\restrl{c}{\initialSegment{n}}}{\initialSegment{n}}$ is finite by \cref{finite_image_function}.
                For all $C$ such that $C\in \img{\restrl{c}{\initialSegment{n}}}{\initialSegment{n}}$ we have
                    $C$ is closed in $X$ with respect to $T$ by
                    \cref{img_iff,pair_eq_iff,restrl_iff,funs_type_apply,powerset_elim,closure_closed_in}.
                Thus $\img{\restrl{c}{\initialSegment{n}}}{\initialSegment{n}}\subseteq\pow{X}$
                    by \cref{closed_in,powerset_intro,subseteq}.
                Thus $\unions{\img{\restrl{c}{\initialSegment{n}}}{\initialSegment{n}}}$ is closed in $X$ with respect to $T$
                    by \cref{closed_finite_unions}.
                Thus $U(n)\in T$ and $U(n)\subseteq X$ by \cref{open_in,topology_on,subseteq,powerset_elim}.
                Thus $W\in T$
                    by \cref{closed_in,setminus_eq_inter_complement,open_in,topology_on}.
        \end{subproof}
        Thus $F\subseteq T$ by \cref{subseteq}.

        Show that for all $W$ such that $W\in G$ we have $W\in T$.
        \begin{subproof}
            Fix $W$ such that $W\in G$.
                Take $n\in\naturalsPlus$ such that $W = V(n)\setminus \unions{\img{\restrl{d}{\initialSegment{n}}}{\initialSegment{n}}}$.
                $\restrl{d}{\initialSegment{n}}$ is a function by \cref{restrl_of_function_is_function}.
                $\initialSegment{n}$ is finite by \cref{initial_segment_finite}.
                $d$ is a function on $\naturalsPlus$.
                Thus $\dom{\restrl{d}{\initialSegment{n}}} = \initialSegment{n}$
                    by \cref{dom_of_restrl_eq_inter,initial_segment_naturalsplus,subseteq,inter_absorb_supseteq_left}.
                Thus $\img{\restrl{d}{\initialSegment{n}}}{\initialSegment{n}}$ is finite by \cref{finite_image_function}.
                For all $C$ such that $C\in \img{\restrl{d}{\initialSegment{n}}}{\initialSegment{n}}$ we have
                    $C$ is closed in $X$ with respect to $T$ by
                    \cref{img_iff,pair_eq_iff,restrl_iff,funs_type_apply,powerset_elim,closure_closed_in}.
                Thus $\img{\restrl{d}{\initialSegment{n}}}{\initialSegment{n}}\subseteq\pow{X}$
                    by \cref{closed_in,powerset_intro,subseteq}.
                Thus $\unions{\img{\restrl{d}{\initialSegment{n}}}{\initialSegment{n}}}$ is closed in $X$ with respect to $T$
                    by \cref{closed_finite_unions}.
                Thus $V(n)\in T$ and $V(n)\subseteq X$ by \cref{open_in,topology_on,subseteq,powerset_elim}.
                Thus $W\in T$
                    by \cref{closed_in,setminus_eq_inter_complement,open_in,topology_on}.
        \end{subproof}
        Thus $G\subseteq T$ by \cref{subseteq}.

        Let $U_1 = \unions{F}$. Let $V_1 = \unions{G}$.
        Thus $U_1$ is open in $X$ with respect to $T$ and $V_1$ is open in $X$ with respect to $T$
            by \cref{topology_on,subseteq,open_in}.

        Show that for all $x$ such that $x\in A$ we have $x\in U_1$.
        \begin{subproof}
            Fix $x$ such that $x\in A$.
                Take $W$ such that $W\in \img{U}{\naturalsPlus}$ and $x\in W$ by \cref{subseteq,unions_iff}.
                Take $n\in\naturalsPlus$ such that $n\mathrel{U} W$ by \cref{img_iff}.
                $(n,U(n))\in U$ by \cref{function_apply_elim}.
                Thus $x\in U(n)$ by \cref{img_iff,pair_eq_iff,function_rightunique}.
                Show that $x\notin \unions{\img{\restrl{c}{\initialSegment{n}}}{\initialSegment{n}}}$.
                \begin{subproof}
                    Suppose not.
                    Then there exists $C$ such that $C\in \img{\restrl{c}{\initialSegment{n}}}{\initialSegment{n}}$ and $x\in C$ by \cref{unions_iff}.
                    Take $k\in\initialSegment{n}$ such that $k\mathrel{\restrl{c}{\initialSegment{n}}} C$ by \cref{img_iff}.
                    Thus $(k,C)\in c$ by \cref{pair_eq_iff,restrl_iff}.
                    Take $k_0\in\naturalsPlus$ such that $(k,C) = (k_0,\closure{V(k_0)}{X}{T})$.
                    Contradiction by \cref{pair_eq_iff,disjoint}.
                \end{subproof}
                Thus $x\in U(n)\setminus \unions{\img{\restrl{c}{\initialSegment{n}}}{\initialSegment{n}}}$ by \cref{setminus_intro}.
                Thus $U(n)\setminus \unions{\img{\restrl{c}{\initialSegment{n}}}{\initialSegment{n}}}\in F$.
                Thus $x\in U_1$ by \cref{unions_intro}.
        \end{subproof}
        Thus $A\subseteq U_1$ by \cref{subseteq}.

        Show that for all $x$ such that $x\in B$ we have $x\in V_1$.
        \begin{subproof}
            Fix $x$ such that $x\in B$.
                Take $W$ such that $W\in \img{V}{\naturalsPlus}$ and $x\in W$ by \cref{subseteq,unions_iff}.
                Take $n\in\naturalsPlus$ such that $n\mathrel{V} W$ by \cref{img_iff}.
                $(n,V(n))\in V$ by \cref{function_apply_elim}.
                Thus $x\in V(n)$ by \cref{img_iff,pair_eq_iff,function_rightunique}.
                Show that $x\notin \unions{\img{\restrl{d}{\initialSegment{n}}}{\initialSegment{n}}}$.
                \begin{subproof}
                    Suppose not.
                    Then there exists $C$ such that $C\in \img{\restrl{d}{\initialSegment{n}}}{\initialSegment{n}}$ and $x\in C$ by \cref{unions_iff}.
                    Take $k\in\initialSegment{n}$ such that $k\mathrel{\restrl{d}{\initialSegment{n}}} C$ by \cref{img_iff}.
                    Thus $(k,C)\in d$ by \cref{pair_eq_iff,restrl_iff}.
                    Take $k_0\in\naturalsPlus$ such that $(k,C) = (k_0,\closure{U(k_0)}{X}{T})$.
                    Contradiction by \cref{pair_eq_iff,disjoint}.
                \end{subproof}
                Thus $x\in V(n)\setminus \unions{\img{\restrl{d}{\initialSegment{n}}}{\initialSegment{n}}}$ by \cref{setminus_intro}.
                Thus $V(n)\setminus \unions{\img{\restrl{d}{\initialSegment{n}}}{\initialSegment{n}}}\in G$.
                Thus $x\in V_1$ by \cref{unions_intro}.
        \end{subproof}
        Thus $B\subseteq V_1$ by \cref{subseteq}.

        Show that $U_1$ is disjoint from $V_1$.
        \begin{subproof}
            Suppose not.
            Then there exists $x$ such that $x\in U_1$ and $x\in V_1$ by \cref{disjoint}.
            Take $W_1$ such that $W_1\in F$ and $x\in W_1$ by \cref{unions_iff}.
            Take $W_2$ such that $W_2\in G$ and $x\in W_2$ by \cref{unions_iff}.
            Take $j\in\naturalsPlus$ such that $W_1 = U(j)\setminus \unions{\img{\restrl{c}{\initialSegment{j}}}{\initialSegment{j}}}$.
            Take $k\in\naturalsPlus$ such that $W_2 = V(k)\setminus \unions{\img{\restrl{d}{\initialSegment{k}}}{\initialSegment{k}}}$.
            $x\in U(j)$ and $x\in V(k)$ by \cref{setminus_elim_left}.
            $x\notin \unions{\img{\restrl{c}{\initialSegment{j}}}{\initialSegment{j}}}$ and
                $x\notin \unions{\img{\restrl{d}{\initialSegment{k}}}{\initialSegment{k}}}$ by \cref{setminus_elim_right}.
            \begin{byCase}
            \caseOf{$j\leq k$.}
                $j\in\initialSegment{k}$ by \cref{initial_segment_naturalsplus}.
                Take $C$ such that $C = \closure{U(j)}{X}{T}$.
                $j\mathrel{\restrl{d}{\initialSegment{k}}} C$ by \cref{restrl_iff}.
                Thus $C\in \img{\restrl{d}{\initialSegment{k}}}{\initialSegment{k}}$ by \cref{img_iff}.
                Thus $x\in C$ by \cref{funs_type_apply,powerset_elim,subseteq_closure,elem_subseteq}.
                Thus $x\in \unions{\img{\restrl{d}{\initialSegment{k}}}{\initialSegment{k}}}$ by \cref{unions_intro}.
                Contradiction.
            \caseOf{$k\leq j$.}
                $k\in\initialSegment{j}$ by \cref{initial_segment_naturalsplus}.
                Take $C$ such that $C = \closure{V(k)}{X}{T}$.
                $k\mathrel{\restrl{c}{\initialSegment{j}}} C$ by \cref{restrl_iff}.
                Thus $C\in \img{\restrl{c}{\initialSegment{j}}}{\initialSegment{j}}$ by \cref{img_iff}.
                Thus $x\in C$ by \cref{funs_type_apply,powerset_elim,subseteq_closure,elem_subseteq}.
                Thus $x\in \unions{\img{\restrl{c}{\initialSegment{j}}}{\initialSegment{j}}}$ by \cref{unions_intro}.
                Contradiction.
            \end{byCase}
        \end{subproof}

        Thus there exist $U_1,V_1$ such that
            $U_1$ is open in $X$ with respect to $T$ and
            $V_1$ is open in $X$ with respect to $T$ and
            $A\subseteq U_1$ and
            $B\subseteq V_1$ and
            $U_1$ is disjoint from $V_1$.
    \end{subproof}
    Hence $X$ has closed singletons with respect to $T$ by \cref{regular_space}.
    Thus $X$ is normal with respect to $T$ by \cref{normal_space}.
\end{proof}
% from §32 helper definition 1: metric ball family avoiding a set
\begin{definition}\label{metric_ball_family}
    $\metricBallFamily{d}{X}{A}{B} = \{W\in\pow{X} \mid \exists a\in A. \exists \epsilon\in\reals.
        0 < \epsilon \land W = \metricBall{d}{X}{a}{\epsilon / (1 + 1)} \land
        \metricBall{d}{X}{a}{\epsilon}\inter B = \emptyset\}$.
\end{definition}

\begin{lemma}\label{subseteq_setminus_disjoint}
    Suppose $C\subseteq X\setminus B$.
    Then $C\inter B = \emptyset$.
\end{lemma}
\begin{proof}
    Suppose not.
    Take $x$ such that $x\in C\inter B$ by \cref{nonempty_inhabited}.
    Hence $x\in B$ and $x\notin B$ by \cref{inter,subseteq,setminus_elim_right}.
    Contradiction.
\end{proof}

\begin{lemma}\label{real_half_plus_half}
    Suppose $\epsilon\in\reals$.
    Then $\epsilon / (1 + 1) + \epsilon / (1 + 1) = \epsilon$.
\end{lemma}
\begin{proof}
    $1 + 1\in\reals$ and $0 < 1 + 1$
        by \cref{real_naturals_subset_reals,real_one_in_naturals,subseteq,real_add_ident,real_add_closed,real_zero_lt_one,real_lt_add,real_lt_trans}.
    Therefore $1 + 1\neq 0$ and $\inv{(1 + 1)}\in\reals$ by \cref{real_add_closed,real_pos_nonzero,real_mul_inverse}.
    Then $\epsilon / (1 + 1) + \epsilon / (1 + 1) =
        (\epsilon + \epsilon) \rmul \inv{(1 + 1)}$ by \cref{real_div,real_right_distrib}.
    Hence $\epsilon + \epsilon = \epsilon \rmul (1 + 1)$.
    Then $(\epsilon + \epsilon) \rmul \inv{(1 + 1)} =
        \epsilon \rmul ((1 + 1) \rmul \inv{(1 + 1)})$ by \cref{real_mul_assoc}.
    Hence $(\epsilon + \epsilon) \rmul \inv{(1 + 1)} = \epsilon$ by \cref{real_mul_inverse,real_mul_ident,real_mul_comm}.
    Therefore $\epsilon / (1 + 1) + \epsilon / (1 + 1) = \epsilon$.
\end{proof}

% from §32 Item 2: metrizable space is normal
\begin{theorem}\label{metrizable_normal}
    Suppose $X$ is metrizable with respect to $T$.
    Then $X$ is normal with respect to $T$.
\end{theorem}
\begin{proof}
    Take $d$ such that $d$ is a metric on $X$ and $T = \metricTopology{d}{X}$ by \cref{metrizable}.
    $1 + 1\in\reals$ and $0 < 1 + 1$
        by \cref{real_naturals_subset_reals,real_one_in_naturals,subseteq,real_add_ident,real_add_closed,real_zero_lt_one,real_lt_add,real_lt_trans}.
    Therefore $1 + 1\neq 0$ and $\inv{(1 + 1)}\in\reals$ by \cref{real_add_closed,real_pos_nonzero,real_mul_inverse}.
    Show that for all $A,B$ such that $A$ is closed in $X$ with respect to $T$ and $B$ is closed in $X$ with respect to $T$
        and $A$ is disjoint from $B$
        there exist $U,V$ such that
            $U$ is open in $X$ with respect to $T$ and
            $V$ is open in $X$ with respect to $T$ and
            $A\subseteq U$ and
            $B\subseteq V$ and
            $U$ is disjoint from $V$.
    \begin{subproof}
        Fix $A$.
        Fix $B$.
        Assume $A$ is closed in $X$ with respect to $T$ and $B$ is closed in $X$ with respect to $T$
            and $A$ is disjoint from $B$.
        Let $F = \metricBallFamily{d}{X}{A}{B}$. Let $G = \metricBallFamily{d}{X}{B}{A}$.
        Let $U = \unions{F}$. Let $V = \unions{G}$.
        $T$ is a topology on $X$ by \cref{metrizable}.

        Show that for all $W$ such that $W\in F$ we have $W\in T$.
        \begin{subproof}
            Fix $W$ such that $W\in F$.
            Take $a,\epsilon$ such that
                $a\in A$ and $\epsilon\in\reals$ and $0 < \epsilon$ and
                $W = \metricBall{d}{X}{a}{\epsilon / (1 + 1)}$ and
                $\metricBall{d}{X}{a}{\epsilon}\inter B = \emptyset$
                by \cref{metric_ball_family}.
            Then $W\in\pow{X}$ and $a\in X$ by \cref{metric_ball_family,closed_in,subseteq}.
            $\epsilon / (1 + 1)\in\reals$ by \cref{real_div,real_mul_closed}.
            $0 < \epsilon / (1 + 1)$ by \cref{real_add_ident,real_lt_mul_pos,real_inv_pos,real_mul_zero,real_lt_subst_left,real_div}.
            Hence $W\in\metricBasis{d}{X}$ by \cref{metric_basis}.
            Thus $W\in T$ by \cref{metric_on,metric_basis_is_basis,basis_elem_open_in_generated,metric_topology,basis_for_topology}.
        \end{subproof}
        Therefore $U$ is open in $X$ with respect to $T$ by \cref{topology_on,subseteq,open_in}.

        Show that for all $W$ such that $W\in G$ we have $W\in T$.
        \begin{subproof}
            Fix $W$ such that $W\in G$.
            Take $b,\epsilon$ such that
                $b\in B$ and $\epsilon\in\reals$ and $0 < \epsilon$ and
                $W = \metricBall{d}{X}{b}{\epsilon / (1 + 1)}$ and
                $\metricBall{d}{X}{b}{\epsilon}\inter A = \emptyset$
                by \cref{metric_ball_family}.
            Then $W\in\pow{X}$ and $b\in X$ by \cref{metric_ball_family,closed_in,subseteq}.
            $\epsilon / (1 + 1)\in\reals$ by \cref{real_div,real_mul_closed}.
            $0 < \epsilon / (1 + 1)$ by \cref{real_add_ident,real_lt_mul_pos,real_inv_pos,real_mul_zero,real_lt_subst_left,real_div}.
            Hence $W\in\metricBasis{d}{X}$ by \cref{metric_basis}.
            Thus $W\in T$ by \cref{metric_on,metric_basis_is_basis,basis_elem_open_in_generated,metric_topology,basis_for_topology}.
        \end{subproof}
        Therefore $V$ is open in $X$ with respect to $T$ by \cref{topology_on,subseteq,open_in}.

        For all $a$ such that $a\in A$ there exists $\epsilon\in\reals$ such that
            $0 < \epsilon$ and $\metricBall{d}{X}{a}{\epsilon}\subseteq X\setminus B$
            by \cref{closed_in,subseteq,setminus_intro,disjoint,setminus_subseteq,metric_open_iff}.
        Show that for all $a$ such that $a\in A$ we have $a\in U$.
        \begin{subproof}
            Fix $a$ such that $a\in A$.
            Take $\epsilon$ such that $\epsilon\in\reals$ and $0 < \epsilon$ and
                $\metricBall{d}{X}{a}{\epsilon}\subseteq X\setminus B$.
            Then $a\in X$ by \cref{closed_in,subseteq}.
            $d((a,a)) = 0$ by \cref{metric_on,metric_zero_characterization}.
            $0 < \epsilon / (1 + 1)$ by \cref{real_add_ident,real_lt_mul_pos,real_inv_pos,real_mul_zero,real_lt_subst_left,real_div}.
            Hence $a\in \metricBall{d}{X}{a}{\epsilon / (1 + 1)}$ by \cref{metric_ball,real_lt_subst_left}.
            Let $W = \metricBall{d}{X}{a}{\epsilon / (1 + 1)}$.
            Hence $W\in F$ by \cref{metric_ball,subseteq,powerset_intro,subseteq_setminus_disjoint,metric_ball_family}.
            Thus $a\in U$ by \cref{unions_intro}.
        \end{subproof}
        Therefore $A\subseteq U$ by \cref{subseteq}.

        For all $b$ such that $b\in B$ there exists $\epsilon\in\reals$ such that
            $0 < \epsilon$ and $\metricBall{d}{X}{b}{\epsilon}\subseteq X\setminus A$
            by \cref{closed_in,subseteq,setminus_intro,disjoint,setminus_subseteq,metric_open_iff}.
        Show that for all $b$ such that $b\in B$ we have $b\in V$.
        \begin{subproof}
            Fix $b$ such that $b\in B$.
            Take $\epsilon$ such that $\epsilon\in\reals$ and $0 < \epsilon$ and
                $\metricBall{d}{X}{b}{\epsilon}\subseteq X\setminus A$.
            Then $b\in X$ by \cref{closed_in,subseteq}.
            $d((b,b)) = 0$ by \cref{metric_on,metric_zero_characterization}.
            $0 < \epsilon / (1 + 1)$ by \cref{real_add_ident,real_lt_mul_pos,real_inv_pos,real_mul_zero,real_lt_subst_left,real_div}.
            Hence $b\in \metricBall{d}{X}{b}{\epsilon / (1 + 1)}$ by \cref{metric_ball,real_lt_subst_left}.
            Let $W = \metricBall{d}{X}{b}{\epsilon / (1 + 1)}$.
            Hence $W\in G$ by \cref{metric_ball,subseteq,powerset_intro,subseteq_setminus_disjoint,metric_ball_family}.
            Thus $b\in V$ by \cref{unions_intro}.
        \end{subproof}
        Therefore $B\subseteq V$ by \cref{subseteq}.

        Show that $U$ is disjoint from $V$.
        \begin{subproof}
            Suppose not.
            Take $z$ such that $z\in U$ and $z\in V$ by \cref{disjoint}.
            Take $W_1$ such that $W_1\in F$ and $z\in W_1$ by \cref{unions_iff}.
            Take $W_2$ such that $W_2\in G$ and $z\in W_2$ by \cref{unions_iff}.
            Take $a,\epsilon_a$ such that
                $a\in A$ and $\epsilon_a\in\reals$ and $0 < \epsilon_a$ and
                $W_1 = \metricBall{d}{X}{a}{\epsilon_a / (1 + 1)}$ and
                $\metricBall{d}{X}{a}{\epsilon_a}\inter B = \emptyset$
                by \cref{metric_ball_family}.
            Take $b,\epsilon_b$ such that
                $b\in B$ and $\epsilon_b\in\reals$ and $0 < \epsilon_b$ and
                $W_2 = \metricBall{d}{X}{b}{\epsilon_b / (1 + 1)}$ and
                $\metricBall{d}{X}{b}{\epsilon_b}\inter A = \emptyset$
                by \cref{metric_ball_family}.
            Thus $d((a,z)) < \epsilon_a / (1 + 1)$ and $d((b,z)) < \epsilon_b / (1 + 1)$ by \cref{metric_ball}.
            $d$ satisfies the triangle inequality on $X$ and $d$ is symmetric on $X$ and $z\in X$
                by \cref{metric_on,metric_symmetric,metric_ball}.
            Then $a\in X$ and $b\in X$ by \cref{closed_in,subseteq}.
            Therefore $d((z,b)) = d((b,z))$ by \cref{metric_symmetric}.
            $(a,z)\in X\times X$ and $(z,b)\in X\times X$ by \cref{times_tuple_intro}.
            Thus $d((a,z))\in\reals$ and $d((z,b))\in\reals$ by \cref{funs_type_apply,metric_on}.
            $\epsilon_a / (1 + 1)\in\reals$ and $\epsilon_b / (1 + 1)\in\reals$ by \cref{real_div,real_mul_closed}.
            Hence $d((a,b))\in\reals$ by \cref{times_tuple_intro,funs_type_apply,metric_on}.
            Then $d((a,z)) + d((z,b))\in\reals$ by \cref{real_add_closed}.
            Hence $\epsilon_a / (1 + 1) + \epsilon_b / (1 + 1)\in\reals$ by \cref{real_add_closed}.
            Therefore $d((a,b)) < d((a,z)) + d((z,b))$ or $d((a,b)) = d((a,z)) + d((z,b))$
                by \cref{metric_triangle_inequality}.
            Hence $d((a,z)) + d((z,b)) < \epsilon_a / (1 + 1) + \epsilon_b / (1 + 1)$ by \cref{real_lt_subst_left,real_lt_add_two}.
            Therefore $d((a,b)) < \epsilon_a / (1 + 1) + \epsilon_b / (1 + 1)$
                by \cref{real_lt_trans,real_lt_subst_left}.
            \begin{byCase}
            \caseOf{$\epsilon_a < \epsilon_b$.}
                $\epsilon_a / (1 + 1)\in\reals$ and $\epsilon_b / (1 + 1)\in\reals$ by \cref{real_div,real_mul_closed}.
                $\epsilon_a / (1 + 1) < \epsilon_b / (1 + 1)$ by \cref{real_lt_mul_pos,real_inv_pos,real_div}.
                Hence $\epsilon_a / (1 + 1) + \epsilon_b / (1 + 1) < \epsilon_b / (1 + 1) + \epsilon_b / (1 + 1)$ by \cref{real_lt_add}.
                Then $\epsilon_b / (1 + 1) + \epsilon_b / (1 + 1) = \epsilon_b$ by \cref{real_half_plus_half}.
                Therefore $d((a,b)) < \epsilon_b / (1 + 1) + \epsilon_b / (1 + 1)$ by \cref{real_lt_trans}.
                Hence $d((a,b)) < \epsilon_b$ by \cref{real_lt_subst_right}.
                Thus $a\in\metricBall{d}{X}{b}{\epsilon_b}$
                    by \cref{metric_symmetric,metric_on,real_lt_subst_left,metric_ball}.
                Contradiction by \cref{inter_intro,subseteq_emptyset_iff,subseteq,emptyset}.
            \caseOf{$\epsilon_a = \epsilon_b$.}
                Then $\epsilon_b = \epsilon_a$.
                Hence $\epsilon_b / (1 + 1) = \epsilon_a / (1 + 1)$.
                Therefore $\epsilon_a / (1 + 1) + \epsilon_b / (1 + 1) = \epsilon_a / (1 + 1) + \epsilon_a / (1 + 1)$ by \cref{real_add_eq_subst}.
                Then $\epsilon_a / (1 + 1) + \epsilon_a / (1 + 1) = \epsilon_a$ by \cref{real_half_plus_half}.
                Therefore $d((a,b)) < \epsilon_a / (1 + 1) + \epsilon_a / (1 + 1)$ by \cref{real_lt_subst_right}.
                Hence $d((a,b)) < \epsilon_a$ by \cref{real_lt_subst_right}.
                Thus $b\in\metricBall{d}{X}{a}{\epsilon_a}$ by \cref{metric_ball}.
                Contradiction by \cref{inter_intro,subseteq_emptyset_iff,subseteq,emptyset}.
            \caseOf{$\epsilon_b < \epsilon_a$.}
                $\epsilon_a / (1 + 1)\in\reals$ and $\epsilon_b / (1 + 1)\in\reals$ by \cref{real_div,real_mul_closed}.
                $\epsilon_b / (1 + 1) < \epsilon_a / (1 + 1)$ by \cref{real_lt_mul_pos,real_inv_pos,real_div}.
                Hence $\epsilon_b / (1 + 1) + \epsilon_a / (1 + 1) < \epsilon_a / (1 + 1) + \epsilon_a / (1 + 1)$ by \cref{real_lt_add}.
                Therefore $\epsilon_a / (1 + 1) + \epsilon_b / (1 + 1) < \epsilon_a / (1 + 1) + \epsilon_a / (1 + 1)$ by \cref{real_add_comm,real_lt_subst_left}.
                Then $\epsilon_a / (1 + 1) + \epsilon_a / (1 + 1) = \epsilon_a$ by \cref{real_half_plus_half}.
                Therefore $d((a,b)) < \epsilon_a / (1 + 1) + \epsilon_a / (1 + 1)$ by \cref{real_lt_trans}.
                Hence $d((a,b)) < \epsilon_a$ by \cref{real_lt_subst_right}.
                Thus $b\in\metricBall{d}{X}{a}{\epsilon_a}$ by \cref{metric_ball}.
                Contradiction by \cref{inter_intro,subseteq_emptyset_iff,subseteq,emptyset}.
            \end{byCase}
        \end{subproof}
    \end{subproof}
    $T$ is a topology on $X$ and $X$ is Hausdorff with respect to $T$ by \cref{metrizable,metric_space_hausdorff}.
    Hence $X$ has closed singletons with respect to $T$
        by \cref{singleton_subset_intro,pair_is_finite,singleton_iff,upair_intro_left,subseteq,subseteq_finite_is_finite,finite_point_set_closed_hausdorff,one_point_closed}.
    Therefore $X$ is normal with respect to $T$ by \cref{normal_space}.
\end{proof}

% from §32 Item 3: compact Hausdorff space is normal
\begin{theorem}\label{compact_hausdorff_normal}
    Suppose $X$ is compact with respect to $T$ and $X$ is Hausdorff with respect to $T$.
    Then $X$ is normal with respect to $T$.
\end{theorem}
\begin{proof}
    $T$ is a topology on $X$ by \cref{hausdorff}.
    Show that for all $A,B$ such that $A$ is closed in $X$ with respect to $T$ and $B$ is closed in $X$ with respect to $T$
        and $A$ is disjoint from $B$
        there exist $U,V$ such that
            $U$ is open in $X$ with respect to $T$ and
            $V$ is open in $X$ with respect to $T$ and
            $A\subseteq U$ and
            $B\subseteq V$ and
            $U$ is disjoint from $V$.
    \begin{subproof}
        Fix $A$.
        Fix $B$.
        Assume $A$ is closed in $X$ with respect to $T$ and $B$ is closed in $X$ with respect to $T$
            and $A$ is disjoint from $B$.
        \begin{byCase}
        \caseOf{$A = \emptyset$.}
            Let $U = \emptyset$. Let $V = X$.
            $U$ is open in $X$ with respect to $T$ and $V$ is open in $X$ with respect to $T$ by \cref{topology_on,open_in}.
            $A\subseteq U$ and $B\subseteq V$ by \cref{emptyset_subseteq,closed_in,subseteq}.
            $U$ is disjoint from $V$ by \cref{disjoint,emptyset}.
        \caseOf{$A\neq\emptyset$.}
            $A\subseteq X$ and $B\subseteq X$ by \cref{closed_in}.
            Let $T_A = \subspaceTopology{T}{A}$. Let $T_B = \subspaceTopology{T}{B}$.
            Hence $A$ is compact with respect to $T_A$ and $B$ is compact with respect to $T_B$ by \cref{closed_subspace_compact}.
            Let $F = \{U\in T \mid \text{there exists $a\in A$ such that there exists $V$ such that
                    $V$ is open in $X$ with respect to $T$ and
                    $a\in U$ and
                    $B\subseteq V$ and
                    $U$ is disjoint from $V$}\}$.
            Therefore $F$ is a covering of $A$
                by \cref{subseteq,disjoint,setminus_intro,compact_hausdorff_separation,open_in,unions_intro,covering}.
            Take $G$ such that $G\subseteq F$ and $G$ is finite and $G$ is a covering of $A$
                by \cref{open_in,subspace_of,compact_subspace_open_covering,subseteq}.
            $G$ is inhabited by \cref{nonempty_inhabited,covering,subseteq,unions_iff}.
            Let $R = \{w\in G\times T \mid \text{there exists $U\in G$ such that there exists $V\in T$ such that
                    $w = (U,V)$ and
                    $V$ is open in $X$ with respect to $T$ and
                    $B\subseteq V$ and
                    $U$ is disjoint from $V$}\}$.
            For all $U\in G$ there exists $V$ such that $U\mathrel{R} V$.
            Take $g$ such that $g$ is a function on $G$ and for all $U\in G$ we have $U\mathrel{R}\apply{g}{U}$
                by \cref{relation,choice_function}.
            Let $H = \img{g}{G}$. Hence $H\subseteq T$ by \cref{img_iff,function_apply_iff,pair_eq_iff,subseteq}.
            $H$ is finite and $H$ is inhabited by
                \cref{finite_image_function,nonempty_inhabited,inhabited_nonempty,function_apply_elim,subseteq,img_iff}.
            Let $U_0 = \unions{G}$. Let $V_0 = \inters{H}$.
            $U_0$ is open in $X$ with respect to $T$ and $V_0$ is open in $X$ with respect to $T$
                by \cref{topology_on,subseteq,finite_intersections_open_in_topology,open_in}.
            $A\subseteq U_0$ by \cref{covering}.
            Therefore $B\subseteq V_0$ by \cref{img_iff,function_apply_iff,pair_eq_iff,inters_greatest}.
            Show that $U_0$ is disjoint from $V_0$.
            \begin{subproof}
                Suppose not.
                Take $z$ such that $z\in U_0$ and $z\in V_0$ by \cref{disjoint}.
                Take $U\in G$ such that $z\in U$ by \cref{unions_iff}.
                Then $z\in \apply{g}{U}$ by \cref{function_apply_elim,subseteq,img_iff,inters_subseteq_elem}.
                Take $U_1,V_1$ such that $U_1\in G$ and $V_1\in T$ and
                    $(U,\apply{g}{U}) = (U_1,V_1)$ and
                    $V_1$ is open in $X$ with respect to $T$ and
                    $B\subseteq V_1$ and
                    $U_1$ is disjoint from $V_1$
                    by \cref{choice_function}.
                Contradiction by \cref{pair_eq_iff,disjoint}.
            \end{subproof}
        \end{byCase}
    \end{subproof}
    Hence $X$ has closed singletons with respect to $T$
        by \cref{singleton_subset_intro,pair_is_finite,singleton_iff,upair_intro_left,subseteq,subseteq_finite_is_finite,finite_point_set_closed_hausdorff,one_point_closed}.
    Therefore $X$ is normal with respect to $T$ by \cref{normal_space}.
\end{proof}

% from §32 helper definition 2: well order relation
\begin{definition}\label{well_order_relation}
    $R$ is a well order relation on $X$ iff
        $R$ is a simple order relation on $X$ and
        for all $A$ such that $A\subseteq X$ and $A\neq\emptyset$
            there exists $a$ such that $a$ is a smallest element of $A$ with respect to $R$.
\end{definition}

% from §32 Item 4: well-ordered set is normal in the order topology
\begin{theorem}\label{well_ordered_normal}
    Assume $R$ is a well order relation on $X$ and $T$ is the order topology on $X$ with respect to $R$.
    Then $X$ is normal with respect to $T$.
\end{theorem}
\begin{proof}
    Omitted. % do not touch
\end{proof}

\section{§33 The Urysohn Lemma}

% from §33 helper definition 1: unit rational set
\begin{definition}\label{unit_rationals}
    $\unitRationals =
        \{\fst{q} \rmul \inv{\snd{q}} \mid q\in\naturalsPlus\times\naturalsPlus\}$.
\end{definition}

% from §33 helper definition 2: unit rational set with zero
\begin{definition}\label{unit_rationals_zero}
    $\unitRationalsZero = \unitRationals \union \{0\}$.
\end{definition}

% from §33 helper definition 3: Urysohn admissible functions
\begin{definition}\label{urysohn_admissible}
    $\urysohnAdmissible{s}{X}{T}{A}{B}{n} =
        \{f\in\funs{\{0,1\}\union \img{s}{\initialSegment{n}}}{\pow{X}} \mid
            \text{for all $p\in\dom{f}$ we have $f(p)\in T$ and
            $f(1) = X\setminus B$ and
            $A\subseteq f(0)$ and
            $\closure{f(0)}{X}{T}\subseteq f(1)$ and
            for all $p\in\dom{f}$ for all $q\in\dom{f}$ if $p < q$ then
                $\closure{f(p)}{X}{T}\subseteq f(q)$ and
            for all $p\in\dom{f}$ if $1 < p$ then $f(p) = X$}\}$.
\end{definition}

% from §33 helper definition 4: Urysohn good pair predicate
\begin{definition}\label{urysohn_good_pair}
    $\urysohnGoodPair{s}{X}{T}{A}{B}{k}{g}$ iff
        $k\in\naturalsPlus$ and $g\in\urysohnAdmissible{s}{X}{T}{A}{B}{k}$.
\end{definition}

% from §33 helper definition 5: Urysohn good pair set
\begin{definition}\label{urysohn_good_pairs}
    $\urysohnGoodPairs{s}{X}{T}{A}{B} =
        \{(k,g) \mid k\in\naturalsPlus,
            g\in\rels{\{0,1\}\union \img{s}{\naturalsPlus}}{\pow{X}} \mid
            \text{$\urysohnGoodPair{s}{X}{T}{A}{B}{k}{g}$}\}$.
\end{definition}

% from §33 helper definition 6: Urysohn level set
\begin{definition}\label{urysohn_level_set}
    $\urysohnLevelSet{U}{x} = \{q\in\unitRationalsZero \mid x\in U(q)\}$.
\end{definition}

% from §33 helper lemma 1: unit rationals are dense in nonnegative intervals
\begin{lemma}\label{unit_rationals_dense_nonneg}
    Suppose $a\in\reals$ and $b\in\reals$ and $0 \leq a$ and $a < b$.
    Then there exists $p$ such that $p\in\unitRationals$ and $p\in\intervalopen{a}{b}$.
\end{lemma}
\begin{proof}
    Take $n, N$ such that $n\in\naturalsPlus$ and $N\in\naturalsPlus$ and $n / N\in\intervalopen{a}{b}$
        by \cref{real_rationals_nonneg_interval}.
    Let $p = n / N$. Then $p\in\unitRationals$ and $p\in\intervalopen{a}{b}$
        by \cref{times_tuple_intro,real_div,fst_eq,snd_eq,unit_rationals}.
\end{proof}

% from §33 helper lemma 2: sequence enumerating unit rationals
\begin{lemma}\label{unit_rationals_sequence}
    There exists $s$ such that
        $s$ is a sequence in $\unitRationals$ and
        for all $p$ such that $p\in\unitRationals$ there exists $n\in\naturalsPlus$ such that $s(n) = p$.
\end{lemma}
\begin{proof}
    Take $g$ such that $g\in\inj{\naturalsPlus\times\naturalsPlus}{\naturalsPlus}$
        by \cref{naturalsplus_pairing_injection}.
    Let $p_0 = (1,1)$. Let $R = \converse{g} \union ((\naturalsPlus\setminus\ran{g})\times\{p_0\})$.
    Show that for all $x,y_1,y_2$ such that $(x,y_1)\in R$ and $(x,y_2)\in R$ we have $y_1 = y_2$.
    \begin{subproof}
        Fix $x$.
        Fix $y_1$.
        Fix $y_2$ such that $(x,y_1)\in R$ and $(x,y_2)\in R$.
        \begin{byCase}
        \caseOf{$x\in\ran{g}$.}
            Hence $(x,y_1)\in\converse{g}$ and $(x,y_2)\in\converse{g}$
                by \cref{union_iff,times_tuple_elim,setminus_elim_right}.
            Then $y_1\mathrel{g} x$ and $y_2\mathrel{g} x$ by \cref{converse_iff}.
            Hence $g(y_1) = x$ and $g(y_2) = x$ by \cref{inj,funs_is_function,function_apply_intro}.
            Hence $y_1 = y_2$ by \cref{dom_iff,funs,subseteq,inj}.
        \caseOf{$x\notin\ran{g}$.}
            Hence $(x,y_1)\in(\naturalsPlus\setminus\ran{g})\times\{p_0\}$ and
                $(x,y_2)\in(\naturalsPlus\setminus\ran{g})\times\{p_0\}$ by \cref{union_iff,converse_iff,ran_iff}.
            Then $y_1 = y_2$ by \cref{times_tuple_elim,singleton_iff}.
        \end{byCase}
    \end{subproof}
    Therefore $R$ is a function by \cref{union_iff,converse_relation,times_elem_is_tuple,relation,rightunique}.
    For all $x\in\naturalsPlus$ we have $x\in\dom{R}$
        by \cref{ran_iff,converse_iff,union_iff,setminus_intro,singleton_intro,times_tuple_intro,dom_iff}.
    For all $x\in\dom{R}$ we have $x\in\naturalsPlus$
        by \cref{dom_iff,union_iff,converse_iff,inj,funs,funs_is_function,funs_type_apply,function_apply_intro,times_tuple_elim,setminus_elim_left}.
    Hence $\dom{R} = \naturalsPlus$ by \cref{subseteq,subseteq_antisymmetric}.
    Let $h = \{(n, \fst{R(n)} \rmul \inv{\snd{R(n)}}) \mid n\in\naturalsPlus\}$.
    For all $n,y_1,y_2$ such that $(n,y_1)\in h$ and $(n,y_2)\in h$ we have $y_1 = y_2$
        by \cref{pair_eq_iff}.
    Therefore $h$ is a function by \cref{pair_eq_iff,relation,rightunique}.
    For all $n\in\dom{h}$ we have $h(n)\in\unitRationals$ by
        \cref{dom_iff,pair_eq_iff,subseteq,function_apply_elim,union_iff,converse_iff,inj,funs,times_tuple_elim,real_one_in_naturals,times_tuple_intro,singleton_subset_intro,fst_type,snd_type,unit_rationals,function_apply_intro}.
    For all $n\in\naturalsPlus$ we have $n\in\dom{h}$ by \cref{pair_eq_iff,dom_iff}.
    For all $n\in\dom{h}$ we have $n\in\naturalsPlus$ by \cref{dom_iff,pair_eq_iff}.
    Therefore $\dom{h} = \naturalsPlus$ by \cref{subseteq,subseteq_antisymmetric}.
    Hence $h$ is a sequence in $\unitRationals$ by \cref{funs_intro,sequence_in}.
    Thus for all $p$ such that $p\in\unitRationals$ there exists $n\in\naturalsPlus$ such that $h(n) = p$ by
        \cref{unit_rationals,inj,funs_is_function,funs,subseteq,funs_type_apply,function_apply_elim,converse_iff,union_iff,dom_iff,function_apply_intro,pair_eq_iff,funs_is_relation}.
    Therefore there exists $s$ such that
        $s$ is a sequence in $\unitRationals$ and
        for all $p$ such that $p\in\unitRationals$ there exists $n\in\naturalsPlus$ such that $s(n) = p$.
\end{proof}


% from §33 helper lemma 3: good pair decomposition
\begin{lemma}\label{urysohn_good_pair_decomp}
    Let $G = \urysohnGoodPairs{s}{X}{T}{A}{B}$.
    Suppose $z\in G$.
    Then there exist $n,f$ such that
        $n\in\naturalsPlus$ and
        $f\in\urysohnAdmissible{s}{X}{T}{A}{B}{n}$ and
        $z = (n,f)$.
\end{lemma}
\begin{proof}
    Take $n,f$ such that
        $n\in\naturalsPlus$ and
        $\urysohnGoodPair{s}{X}{T}{A}{B}{n}{f}$ and
        $z = (n,f)$.
    Then $f\in\urysohnAdmissible{s}{X}{T}{A}{B}{n}$ by \cref{urysohn_good_pair}.
\end{proof}

% from §33 helper lemma 4: good pair introduction
\begin{lemma}\label{urysohn_good_pair_intro}
    Let $G = \urysohnGoodPairs{s}{X}{T}{A}{B}$.
    Suppose $n\in\naturalsPlus$ and $f\in\urysohnAdmissible{s}{X}{T}{A}{B}{n}$.
    Then $(n,f)\in G$.
\end{lemma}
\begin{proof}
    Then $f\in\rels{\{0,1\}\union \img{s}{\naturalsPlus}}{\pow{X}}$ by
        \cref{urysohn_admissible,funs_subseteq_rels,elem_subseteq,rels,powerset_elim,initial_segment_naturalsplus,subseteq,img_subseteq,union_upper_left,union_upper_right,subseteq_transitive,union_subsets_is_subset,times_subseteq_left,rels_intro}.
    Hence $(n,f)\in G$.
\end{proof}

% from §33 helper lemma 6: continuity at a point via open neighborhoods
\begin{lemma}\label{continuous_at_open}
    Assume $f$ is a function from $X$ to $Y$.
    Assume $T$ is a topology on $X$.
    Assume $T_1$ is a topology on $Y$.
    Assume $x\in X$.
    Suppose for all $V$ such that $V$ is open in $Y$ with respect to $T_1$ and $f(x)\in V$ there exists $U$ such that
        $U$ is open in $X$ with respect to $T$ and $x\in U$ and $\img{f}{U}\subseteq V$.
    Then $f$ is continuous at $x$ from $X$ to $Y$ with respect to $T$ and $T_1$.
\end{lemma}
\begin{proof}
    Follows by \cref{neighborhood,continuous_at}.
\end{proof}

% from §33 helper lemma 7: continuity to a subspace via range inclusion
\begin{lemma}\label{continuous_subspace_range}
    Assume $f\in\funs{X}{\reals}$.
    Assume $T$ is a topology on $X$.
    Assume $I_0\subseteq\reals$.
    Let $T_0 = \subspaceTopology{\standardTopologyR}{I_0}$.
    Assume $f$ is continuous from $X$ to $\reals$ with respect to $T$ and $\standardTopologyR$.
    Assume $\img{f}{X}\subseteq I_0$.
    Then $f$ is continuous from $X$ to $I_0$ with respect to $T$ and $T_0$.
\end{lemma}
\begin{proof}
    Hence $T_0$ is a topology on $I_0$
        by \cref{standard_basis_is_basis,basis_topology_is_topology,standard_topology_reals,subspace_of,subspace_topology_is_topology}.
    Show that for all $B$ such that $B$ is open in $I_0$ with respect to $T_0$
        we have $\preimg{f}{B}$ is open in $X$ with respect to $T$.
    \begin{subproof}
        Fix $B$ such that $B$ is open in $I_0$ with respect to $T_0$.
        Take $U\in\standardTopologyR$ such that $B = I_0\inter U$ by \cref{subspace_topology,open_in}.
        Therefore $\preimg{f}{B} = \preimg{f}{U}$ by
            \cref{preimg_iff,inter_elim_right,subseteq,preimg_subseteq_dom,funs,funs_is_function,function_rightunique,funs_is_relation,function_apply_iff,img_of_function_intro,inter_intro,subseteq_antisymmetric}.
        Hence $\preimg{f}{B}$ is open in $X$ with respect to $T$ by \cref{open_in,continuous}.
    \end{subproof}
    Therefore $f$ is continuous from $X$ to $I_0$ with respect to $T$ and $T_0$
        by \cref{funs_is_function,funs_elim,img_dom,funs,subseteq,function_apply_elim,ran_iff,funs_intro,continuous}.
\end{proof}

% from §33 helper lemma 8: initial good pair when $r_1 = 1$
\begin{lemma}\label{urysohn_u1_eq1}
    Assume $U_1 = X\setminus B$.
    Assume $U_1$ is open in $X$ with respect to $T$.
    Assume $U_0$ is open in $X$ with respect to $T$.
    Assume $A\subseteq U_0$.
    Assume $\closure{U_0}{X}{T}\subseteq U_1$.
    Assume $r_1 = 1$.
    Assume $\img{s}{\initialSegment{1}} = \{r_1\}$.
    Let $G = \urysohnGoodPairs{s}{X}{T}{A}{B}$.
    Then there exists $u_1$ such that $(1,u_1)\in G$.
\end{lemma}
\begin{proof}
    Let $u_1 = \{(0,U_0),(1,U_1)\}$.
    Show that for all $a,b,b'$ such that $(a,b)\in u_1$ and $(a,b')\in u_1$ we have $b = b'$.
    \begin{subproof}
        Fix $a$.
        Fix $b$.
        Fix $b'$ such that $(a,b)\in u_1$ and $(a,b')\in u_1$.
        Thus $(a,b) = (0,U_0)$ or $(a,b) = (1,U_1)$ by \cref{upair_elim}.
        \begin{byCase}
        \caseOf{$(a,b) = (0,U_0)$.}
            Thus $a = 0$ by \cref{pair_eq_iff}.
            Show that $b = b'$.
            \begin{subproof}
                Suppose not.
                Thus $(a,b') = (0,U_0)$ or $(a,b') = (1,U_1)$ by \cref{upair_elim}.
                \begin{byCase}
                \caseOf{$(a,b') = (0,U_0)$.}
                    Contradiction by \cref{pair_eq_iff}.
                \caseOf{$(a,b') = (1,U_1)$.}
                    Contradiction by \cref{pair_eq_iff,real_mul_ident}.
                \end{byCase}
            \end{subproof}
        \caseOf{$(a,b) = (1,U_1)$.}
            Thus $a = 1$ by \cref{pair_eq_iff}.
            Show that $b = b'$.
            \begin{subproof}
                Suppose not.
                Thus $(a,b') = (0,U_0)$ or $(a,b') = (1,U_1)$ by \cref{upair_elim}.
                \begin{byCase}
                \caseOf{$(a,b') = (1,U_1)$.}
                    Contradiction by \cref{pair_eq_iff}.
                \caseOf{$(a,b') = (0,U_0)$.}
                    Contradiction by \cref{pair_eq_iff,real_mul_ident}.
                \end{byCase}
            \end{subproof}
        \end{byCase}
    \end{subproof}
    Therefore $u_1$ is a function by \cref{upair_elim,relation,rightunique}.
    Thus $\dom{u_1}\subseteq \{0,1\}$
        by \cref{dom_iff,upair_elim,pair_eq_iff,upair_intro_left,upair_intro_right,subseteq}.
    Thus $\{0,1\}\subseteq \dom{u_1}$
        by \cref{upair_elim,upair_intro_left,upair_intro_right,dom_intro,subseteq}.
    Therefore $\dom{u_1} = \{0,1\}$ by \cref{subseteq_antisymmetric}.
    Then $\img{s}{\initialSegment{1}} = \{1\}$.
    Hence $\{0,1\}\union \img{s}{\initialSegment{1}}\subseteq \{0,1\}$ and
        $\{0,1\}\subseteq \{0,1\}\union \img{s}{\initialSegment{1}}$
        by \cref{upair_intro_right,subseteq_refl,singleton_subset_intro,subseteq_transitive,union_subsets_is_subset,union_intro_left,subseteq}.
    Therefore $\{0,1\}\union \img{s}{\initialSegment{1}} = \{0,1\}$ by \cref{subseteq_antisymmetric}.
    Hence $\dom{u_1} = \{0,1\}\union \img{s}{\initialSegment{1}}$ by \cref{subseteq_antisymmetric}.
    For all $p\in\dom{u_1}$ we have $u_1(p)\in\pow{X}$
        by \cref{subseteq,upair_elim,upair_intro_left,upair_intro_right,function_apply_intro,open_in,topology_on,elem_subseteq}.
    Therefore $u_1\in\funs{\{0,1\}\union \img{s}{\initialSegment{1}}}{\pow{X}}$ by \cref{funs_intro}.
    Show that for all $p\in\dom{u_1}$ we have
                $u_1(p)\in T$ and
                $u_1(1) = X\setminus B$ and
                $A\subseteq u_1(0)$ and
                $\closure{u_1(0)}{X}{T}\subseteq u_1(1)$ and
                for all $p\in\dom{u_1}$ for all $q\in\dom{u_1}$ if $p < q$ then
                    $\closure{u_1(p)}{X}{T}\subseteq u_1(q)$ and
                for all $p\in\dom{u_1}$ if $1 < p$ then $u_1(p) = X$.
    \begin{subproof}
        Fix $p$ such that $p\in\dom{u_1}$.
        Now $u_1$ is a function and $u_1$ is right-unique and $u_1$ is a relation
            by \cref{funs_is_function,function_rightunique,funs_is_relation}.
        Thus $u_1(p)\in T$ by \cref{subseteq,upair_elim,upair_intro_left,upair_intro_right,function_apply_intro,open_in}.
        Then $u_1(1) = X\setminus B$ and $A\subseteq u_1(0)$ and
            $\closure{u_1(0)}{X}{T}\subseteq u_1(1)$
            by \cref{upair_intro_right,upair_intro_left,function_apply_intro,subseteq}.
        Show that for all $p\in\dom{u_1}$ for all $q\in\dom{u_1}$ if $p < q$ then
            $\closure{u_1(p)}{X}{T}\subseteq u_1(q)$.
        \begin{subproof}
            Fix $p\in\dom{u_1}$.
            Fix $q\in\dom{u_1}$.
            Assume $p < q$.
            Now $p\in\{0,1\}$ and $q\in\{0,1\}$.
            Show that $p = 0$.
            \begin{subproof}
                Suppose not.
                Thus $p = 1$ by \cref{upair_elim}.
                Thus $q = 0$ or $q = 1$ by \cref{upair_elim}.
                \begin{byCase}
                \caseOf{$q = 0$.}
                    $0\in\reals$ and $1\in\reals$ and $0 < 1$
                        by \cref{real_add_ident,real_naturals_subset_reals,real_one_in_naturals,subseteq,real_zero_lt_one}.
                    Thus $1 < q$ by \cref{real_lt_subst_left}.
                    Thus $1 < 0$ by \cref{real_lt_subst_right}.
                    Thus $1 < 1$ by \cref{real_lt_trans}.
                    Contradiction by \cref{real_lt_irreflexive}.
                \caseOf{$q = 1$.}
                    Thus $1 < q$ by \cref{real_lt_subst_left}.
                    Thus $1 < 1$ by \cref{real_lt_subst_right}.
                    Contradiction by \cref{real_naturals_subset_reals,real_one_in_naturals,subseteq,real_lt_irreflexive}.
                \end{byCase}
            \end{subproof}
            Suppose not.
            Thus $q = 0$ by \cref{upair_elim}.
            Then $p = 0$.
            Hence $p < 0$ by \cref{real_lt_subst_right}.
            Therefore $0 < 0$ by \cref{real_lt_subst_left}.
            Contradiction by \cref{real_add_ident,real_lt_irreflexive}.
            Hence $q = 1$.
            Then $u_1(p) = U_0$ and $u_1(q) = U_1$
                by \cref{pair_eq_iff,upair_intro_left,upair_intro_right,function_apply_intro,funs_is_function,function_rightunique,funs_is_relation}.
            Hence $\closure{u_1(p)}{X}{T}\subseteq u_1(q)$ by \cref{subseteq}.
        \end{subproof}
        Show that for all $p\in\dom{u_1}$ if $1 < p$ then $u_1(p) = X$.
        \begin{subproof}
            Fix $p$ such that $p\in\dom{u_1}$.
            Assume $1 < p$.
            Suppose not.
            Then $p = 0$ or $p = 1$ by \cref{subseteq,upair_elim}.
            \begin{byCase}
            \caseOf{$p = 0$.}
                $0\in\reals$ and $1\in\reals$ and $0 < 1$
                    by \cref{real_add_ident,real_naturals_subset_reals,real_one_in_naturals,subseteq,real_zero_lt_one}.
                Thus $1 < 0$ by \cref{real_lt_subst_right}.
                Thus $1 < 1$ by \cref{real_lt_trans}.
                Contradiction by \cref{real_lt_irreflexive}.
            \caseOf{$p = 1$.}
                Thus $1 < 1$ by \cref{real_lt_subst_right}.
                Contradiction by \cref{real_naturals_subset_reals,real_one_in_naturals,subseteq,real_lt_irreflexive}.
            \end{byCase}
        \end{subproof}
    \end{subproof}
    Hence $u_1\in\urysohnAdmissible{s}{X}{T}{A}{B}{1}$ by \cref{urysohn_admissible}.
    Therefore there exists $u_1$ such that $(1,u_1)\in G$
        by \cref{urysohn_good_pair_intro,real_one_in_naturals}.
\end{proof}
% from §33 helper lemma 9: initial good pair when $1 < r_1$
\begin{lemma}\label{urysohn_u1_gt1}
    Assume $U_1 = X\setminus B$.
    Assume $U_1$ is open in $X$ with respect to $T$.
    Assume $U_0$ is open in $X$ with respect to $T$.
    Assume $A\subseteq U_0$.
    Assume $\closure{U_0}{X}{T}\subseteq U_1$.
    Assume $r_1\in\reals$.
    Assume $1 < r_1$.
    Assume $\img{s}{\initialSegment{1}} = \{r_1\}$.
    Let $G = \urysohnGoodPairs{s}{X}{T}{A}{B}$.
    Then there exists $u_1$ such that $(1,u_1)\in G$.
\end{lemma}
\begin{proof}
    Let $U_r = X$.
    Let $u_1 = \{(0,U_0),(1,U_1)\}\union\{(r_1,U_r)\}$.
    Show that $(1,u_1)\in G$.
    \begin{subproof}
        Show that $u_1\in\urysohnAdmissible{s}{X}{T}{A}{B}{1}$.
        \begin{subproof}
            Show that $u_1\in\funs{\{0,1\}\union \img{s}{\initialSegment{1}}}{\pow{X}}$.
            \begin{subproof}
                Thus $r_1 \neq 0$ by \cref{real_add_ident,real_mul_ident,real_zero_lt_one,real_lt_trans,real_pos_nonzero}.
                Show that for all $a,b,b'$ such that $(a,b)\in u_1$ and $(a,b')\in u_1$ we have $b = b'$.
                \begin{subproof}
                        Fix $a,b$.
                        Fix $b'$ such that $(a,b)\in u_1$ and $(a,b')\in u_1$.
                        Thus $(a,b) = (0,U_0)$ or $(a,b) = (1,U_1)$ or $(a,b) = (r_1,U_r)$
                            by \cref{union_iff,upair_elim,singleton_elim}.
                        \begin{byCase}
                        \caseOf{$(a,b) = (0,U_0)$.}
                            Thus $a = 0$ by \cref{pair_eq_iff}.
                            Show that $b = b'$.
                            \begin{subproof}
                                Suppose not.
                                Thus $(a,b')\in\{(0,U_0),(1,U_1)\}$ or $(a,b')\in\{(r_1,U_r)\}$ by \cref{union_iff}.
                                \begin{byCase}
                                \caseOf{$(a,b')\in\{(0,U_0),(1,U_1)\}$.}
                                    Thus $(a,b') = (0,U_0)$ or $(a,b') = (1,U_1)$ by \cref{upair_elim}.
                                    \begin{byCase}
                                    \caseOf{$(a,b') = (0,U_0)$.}
                                        Contradiction by \cref{pair_eq_iff}.
                                    \caseOf{$(a,b') = (1,U_1)$.}
                                        Contradiction by \cref{pair_eq_iff,real_mul_ident}.
                                    \end{byCase}
                                \caseOf{$(a,b')\in\{(r_1,U_r)\}$.}
                                    Thus $(a,b') = (r_1,U_r)$ by \cref{singleton_elim}.
                                    Thus $a = r_1$ by \cref{pair_eq_iff}.
                                    Contradiction.
                                \end{byCase}
                            \end{subproof}
                        \caseOf{$(a,b) = (1,U_1)$.}
                            Thus $a = 1$ by \cref{pair_eq_iff}.
                            Show that $b = b'$.
                            \begin{subproof}
                                Suppose not.
                                Thus $(a,b')\in\{(0,U_0),(1,U_1)\}$ or $(a,b')\in\{(r_1,U_r)\}$ by \cref{union_iff}.
                                \begin{byCase}
                                \caseOf{$(a,b')\in\{(0,U_0),(1,U_1)\}$.}
                                    Thus $(a,b') = (0,U_0)$ or $(a,b') = (1,U_1)$ by \cref{upair_elim}.
                                    \begin{byCase}
                                    \caseOf{$(a,b') = (1,U_1)$.}
                                        Contradiction by \cref{pair_eq_iff}.
                                    \caseOf{$(a,b') = (0,U_0)$.}
                                        Contradiction by \cref{pair_eq_iff,real_mul_ident}.
                                    \end{byCase}
                                \caseOf{$(a,b')\in\{(r_1,U_r)\}$.}
                                    Thus $(a,b') = (r_1,U_r)$ by \cref{singleton_elim}.
                                    Thus $a = r_1$ by \cref{pair_eq_iff}.
                                    Contradiction.
                                \end{byCase}
                            \end{subproof}
                        \caseOf{$(a,b) = (r_1,U_r)$.}
                            Thus $a = r_1$ by \cref{pair_eq_iff}.
                            Show that $b = b'$.
                            \begin{subproof}
                                Suppose not.
                                Thus $(a,b')\in\{(0,U_0),(1,U_1)\}$ or $(a,b')\in\{(r_1,U_r)\}$ by \cref{union_iff}.
                                \begin{byCase}
                                \caseOf{$(a,b')\in\{(r_1,U_r)\}$.}
                                    Thus $(a,b') = (r_1,U_r)$ by \cref{singleton_elim}.
                                    Contradiction by \cref{pair_eq_iff}.
                                \caseOf{$(a,b')\in\{(0,U_0),(1,U_1)\}$.}
                                    Thus $(a,b') = (0,U_0)$ or $(a,b') = (1,U_1)$ by \cref{upair_elim}.
                                    \begin{byCase}
                                    \caseOf{$(a,b') = (0,U_0)$.}
                                        Thus $a = 0$ by \cref{pair_eq_iff}.
                                        Contradiction.
                                    \caseOf{$(a,b') = (1,U_1)$.}
                                        Thus $a = 1$ by \cref{pair_eq_iff}.
                                        Contradiction.
                                    \end{byCase}
                                \end{byCase}
                            \end{subproof}
                        \end{byCase}
                \end{subproof}
                Therefore $u_1$ is a function by \cref{union_iff,upair_elim,singleton_elim,relation,rightunique}.
                    Thus $\dom{u_1}\subseteq \{0,1\}\union\{r_1\}$
                        by \cref{dom_iff,union_iff,upair_elim,singleton_elim,pair_eq_iff,upair_intro_left,upair_intro_right,singleton_intro,union_intro_left,union_intro_right,subseteq}.
                    Thus $\{0,1\}\union\{r_1\}\subseteq \dom{u_1}$
                        by \cref{union_iff,upair_elim,singleton_elim,upair_intro_left,upair_intro_right,singleton_intro,union_intro_left,union_intro_right,dom_intro,subseteq}.
                    Therefore $\dom{u_1} = \{0,1\}\union\{r_1\}$ by \cref{subseteq_antisymmetric}.
                    Hence $\{0,1\}\union \img{s}{\initialSegment{1}}\subseteq \{0,1\}\union\{r_1\}$ and
                        $\{0,1\}\union\{r_1\}\subseteq \{0,1\}\union \img{s}{\initialSegment{1}}$
                        by \cref{upair_intro_right,subseteq_refl,singleton_subset_intro,subseteq_transitive,union_subsets_is_subset,union_intro_left,union_intro_right,subseteq}.
                    Therefore $\{0,1\}\union \img{s}{\initialSegment{1}} = \{0,1\}\union\{r_1\}$ by \cref{subseteq_antisymmetric}.
                    Hence $\dom{u_1} = \{0,1\}\union \img{s}{\initialSegment{1}}$ by \cref{subseteq_antisymmetric}.
                For all $p\in\dom{u_1}$ we have $u_1(p)\in\pow{X}$
                    by \cref{subseteq,union_iff,upair_elim,upair_intro_left,upair_intro_right,singleton_elim,singleton_intro,union_intro_left,union_intro_right,function_apply_intro,open_in,topology_on,elem_subseteq}.
                Therefore $u_1\in\funs{\{0,1\}\union \img{s}{\initialSegment{1}}}{\pow{X}}$ by \cref{funs_intro}.
            \end{subproof}
            Show that for all $p\in\dom{u_1}$ we have
                $u_1(p)\in T$ and
                $u_1(1) = X\setminus B$ and
                $A\subseteq u_1(0)$ and
                $\closure{u_1(0)}{X}{T}\subseteq u_1(1)$ and
                for all $p\in\dom{u_1}$ for all $q\in\dom{u_1}$ if $p < q$ then
                    $\closure{u_1(p)}{X}{T}\subseteq u_1(q)$ and
                for all $p\in\dom{u_1}$ if $1 < p$ then $u_1(p) = X$.
            \begin{subproof}
                Fix $p$ such that $p\in\dom{u_1}$.
                Now $u_1$ is a function and $u_1$ is right-unique and $u_1$ is a relation
                    by \cref{funs_is_function,function_rightunique,funs_is_relation}.
                Show that $u_1(p)\in T$.
                \begin{subproof}
                    Thus $p\in\{0,1\}\union\{r_1\}$.
                    Thus $p\in\{0,1\}$ or $p\in\{r_1\}$ by \cref{union_iff}.
                    \begin{byCase}
                    \caseOf{$p\in\{0,1\}$.}
                        Thus $p = 0$ or $p = 1$ by \cref{upair_elim}.
                        \begin{byCase}
                        \caseOf{$p = 0$.}
                            Thus $u_1(p)\in T$ by \cref{upair_intro_left,union_intro_left,function_apply_intro,open_in}.
                        \caseOf{$p = 1$.}
                            Thus $u_1(p)\in T$ by \cref{upair_intro_right,union_intro_left,function_apply_intro,open_in}.
                        \end{byCase}
                    \caseOf{$p\in\{r_1\}$.}
                        Thus $u_1(p)\in T$ by \cref{singleton_elim,singleton_intro,union_intro_right,function_apply_intro,open_in,topology_on}.
                    \end{byCase}
                \end{subproof}
                Then $u_1(1) = X\setminus B$ and $A\subseteq u_1(0)$ and
                    $\closure{u_1(0)}{X}{T}\subseteq u_1(1)$
                    by \cref{upair_intro_right,upair_intro_left,union_intro_left,function_apply_intro,subseteq}.
                Show that for all $p\in\dom{u_1}$ for all $q\in\dom{u_1}$ if $p < q$ then
                    $\closure{u_1(p)}{X}{T}\subseteq u_1(q)$.
                \begin{subproof}
                    Fix $p\in\dom{u_1}$.
                    Fix $q\in\dom{u_1}$.
                    Assume $p < q$.
                    Thus $p\in\{0,1\}\union\{r_1\}$ and $q\in\{0,1\}\union\{r_1\}$.
                    Show that $q \neq 0$.
                    \begin{subproof}
                        Suppose not.
                        Thus $q = 0$.
                        Thus $p < 0$ by \cref{real_lt_subst_right}.
                        Thus $p\in\{0,1\}$ or $p\in\{r_1\}$ by \cref{union_iff}.
                        \begin{byCase}
                        \caseOf{$p\in\{0,1\}$.}
                            Thus $p = 0$ or $p = 1$ by \cref{upair_elim}.
                            \begin{byCase}
                            \caseOf{$p = 0$.}
                                Contradiction by \cref{real_lt_subst_left,real_add_ident,real_lt_irreflexive}.
                            \caseOf{$p = 1$.}
                                Thus $1 < 0$ by \cref{real_lt_subst_left}.
                                $0\in\reals$ and $1\in\reals$ and $0 < 1$ by \cref{real_add_ident,real_mul_ident,real_zero_lt_one}.
                                Thus $1 < 1$ by \cref{real_lt_trans}.
                                Contradiction by \cref{real_lt_irreflexive}.
                            \end{byCase}
                        \caseOf{$p\in\{r_1\}$.}
                            Thus $r_1 < 0$ by \cref{singleton_elim,real_lt_subst_left}.
                            $0\in\reals$ and $1\in\reals$ and $0 < 1$ by \cref{real_add_ident,real_mul_ident,real_zero_lt_one}.
                            Thus $1 < 0$ by \cref{real_lt_trans}.
                            Thus $1 < 1$ by \cref{real_lt_trans}.
                            Contradiction by \cref{real_lt_irreflexive}.
                        \end{byCase}
                    \end{subproof}
                    Hence $q\in\{1\}\union\{r_1\}$.
                    Then $q\in\{1\}$ or $q\in\{r_1\}$ by \cref{union_iff}.
                    \begin{byCase}
                    \caseOf{$q\in\{r_1\}$.}
                        Thus $u_1(q) = X$ by \cref{singleton_elim,singleton_intro,union_intro_right,function_apply_intro}.
                        Therefore $\closure{u_1(p)}{X}{T}\subseteq u_1(q)$
                            by \cref{open_in,topology_on,setminus_self,subseteq_refl,closed_in,elem_subseteq,pow_iff,closure_subseteq_closed,subseteq}.
                    \caseOf{$q\in\{1\}$.}
                        Thus $q = 1$ by \cref{singleton_elim}.
                        Show that $p = 0$.
                        \begin{subproof}
                            Suppose not.
                            Thus $p\in\{0,1\}$ or $p\in\{r_1\}$ by \cref{union_iff}.
                            \begin{byCase}
                            \caseOf{$p\in\{0,1\}$.}
                                Thus $p = 0$ or $p = 1$ by \cref{upair_elim}.
                                \begin{byCase}
                                \caseOf{$p = 1$.}
                                    Contradiction by \cref{real_lt_subst_left,real_mul_ident,real_lt_irreflexive}.
                                \caseOf{$p = 0$.}
                                    Contradiction.
                                \end{byCase}
                            \caseOf{$p\in\{r_1\}$.}
                                Thus $r_1 < 1$ by \cref{singleton_elim,real_lt_subst_left}.
                                Contradiction.
                            \end{byCase}
                        \end{subproof}
                        Then $p = 0$.
                        Hence $\closure{u_1(p)}{X}{T}\subseteq u_1(q)$
                            by \cref{singleton_elim,pair_eq_iff,union_intro_left,upair_intro_left,upair_intro_right,function_apply_intro,subseteq}.
                    \end{byCase}
                \end{subproof}
                Show that for all $p\in\dom{u_1}$ if $1 < p$ then $u_1(p) = X$.
                \begin{subproof}
                    Fix $p\in\dom{u_1}$.
                    Assume $1 < p$.
                    Suppose not.
                    Thus $p\in\{0,1\}\union\{r_1\}$.
                    Thus $p\in\{0,1\}$ or $p\in\{r_1\}$ by \cref{union_iff}.
                    \begin{byCase}
                    \caseOf{$p\in\{0,1\}$.}
                        Thus $p = 0$ or $p = 1$ by \cref{upair_elim}.
                        \begin{byCase}
                        \caseOf{$p = 0$.}
                            Thus $1 < 0$ by \cref{real_lt_subst_left}.
                            $0\in\reals$ and $1\in\reals$ and $0 < 1$ by \cref{real_add_ident,real_mul_ident,real_zero_lt_one}.
                            Thus $1 < 1$ by \cref{real_lt_trans}.
                            Contradiction by \cref{real_lt_irreflexive}.
                        \caseOf{$p = 1$.}
                            Contradiction by \cref{real_lt_subst_left,real_mul_ident,real_lt_irreflexive}.
                        \end{byCase}
                    \caseOf{$p\in\{r_1\}$.}
                        Thus $1 < r_1$ by \cref{singleton_elim,real_lt_subst_left}.
                        Contradiction.
                    \end{byCase}
                \end{subproof}
            \end{subproof}
            Hence $u_1\in\urysohnAdmissible{s}{X}{T}{A}{B}{1}$ by \cref{urysohn_admissible}.
        \end{subproof}
        Therefore $(1,u_1)\in G$ by \cref{urysohn_good_pair_intro,real_one_in_naturals}.
    \end{subproof}
\end{proof}

% from §33 helper lemma 12: unit rational is positive
\begin{lemma}\label{unit_rational_pos}
    Assume $r_1\in\unitRationals$.
    Then $0 < r_1$.
\end{lemma}
\begin{proof}
    Take $q$ such that
        $q\in\naturalsPlus\times\naturalsPlus$ and
        $r_1 = \fst{q} \rmul \inv{\snd{q}}$
        by \cref{unit_rationals}.
    Take $m,n$ such that $q = (m,n)$ and $m\in\naturalsPlus$ and $n\in\naturalsPlus$
        by \cref{times}.
    Therefore $r_1 > 0$
        by \cref{real_naturals_subset_reals,subseteq,real_add_ident,real_zero_lt_one,real_mul_ident,real_one_leq_naturalsplus,real_lt_subst_right,real_lt_trans,fst_eq,snd_eq,real_pos_nonzero,real_mul_inverse,real_inv_pos,real_mul_pos_pos}.
\end{proof}

% from §33 helper lemma 10: admissible good pair for $r_1 < 1$ given $U_r$
\begin{lemma}\label{urysohn_u1_lt1_admissible}
    Assume $U_1 = X\setminus B$.
    Assume $U_1$ is open in $X$ with respect to $T$.
    Assume $U_0$ is open in $X$ with respect to $T$.
    Assume $A\subseteq U_0$.
    Assume $\closure{U_0}{X}{T}\subseteq U_1$.
    Assume $r_1\in\reals$.
    Assume $r_1\in\unitRationals$.
    Assume $r_1 < 1$.
    Assume $\img{s}{\initialSegment{1}} = \{r_1\}$.
    Assume $U_r$ is open in $X$ with respect to $T$.
    Assume $\closure{U_0}{X}{T}\subseteq U_r$.
    Assume $\closure{U_r}{X}{T}\subseteq U_1$.
    Let $G = \urysohnGoodPairs{s}{X}{T}{A}{B}$.
    Then there exists $u_1$ such that $(1,u_1)\in G$.
\end{lemma}
\begin{proof}
    Let $u_1 = \{(0,U_0),(1,U_1)\}\union\{(r_1,U_r)\}$.
    Show that $(1,u_1)\in G$.
    \begin{subproof}
        Show that $u_1\in\urysohnAdmissible{s}{X}{T}{A}{B}{1}$.
        \begin{subproof}
            Show that $u_1\in\funs{\{0,1\}\union \img{s}{\initialSegment{1}}}{\pow{X}}$.
            \begin{subproof}
                Thus $r_1 \neq 0$ by \cref{unit_rational_pos,real_pos_nonzero}.
                Show that for all $a,b,b'$ such that $(a,b)\in u_1$ and $(a,b')\in u_1$ we have $b = b'$.
                \begin{subproof}
                        Fix $a,b$.
                        Fix $b'$ such that $(a,b)\in u_1$ and $(a,b')\in u_1$.
                        Thus $(a,b) = (0,U_0)$ or $(a,b) = (1,U_1)$ or $(a,b) = (r_1,U_r)$
                            by \cref{union_iff,upair_elim,singleton_elim}.
                        \begin{byCase}
                        \caseOf{$(a,b) = (0,U_0)$.}
                            Thus $a = 0$ by \cref{pair_eq_iff}.
                            Show that $b = b'$.
                            \begin{subproof}
                                Suppose not.
                                Thus $(a,b')\in\{(0,U_0),(1,U_1)\}$ or $(a,b')\in\{(r_1,U_r)\}$ by \cref{union_iff}.
                                \begin{byCase}
                                \caseOf{$(a,b')\in\{(0,U_0),(1,U_1)\}$.}
                                    Thus $(a,b') = (0,U_0)$ or $(a,b') = (1,U_1)$ by \cref{upair_elim}.
                                    \begin{byCase}
                                    \caseOf{$(a,b') = (0,U_0)$.}
                                        Thus $b = b'$ by \cref{pair_eq_iff}.
                                        Contradiction.
                                    \caseOf{$(a,b') = (1,U_1)$.}
                                        Thus $a = 1$ by \cref{pair_eq_iff}.
                                        Thus $1\neq 0$ by \cref{real_mul_ident}.
                                        Contradiction.
                                    \end{byCase}
                                \caseOf{$(a,b')\in\{(r_1,U_r)\}$.}
                                    Thus $(a,b') = (r_1,U_r)$ by \cref{singleton_elim}.
                                    Thus $a = r_1$ by \cref{pair_eq_iff}.
                                    Thus $r_1 \neq 0$.
                                    Contradiction.
                                \end{byCase}
                            \end{subproof}
                        \caseOf{$(a,b) = (1,U_1)$.}
                            Thus $a = 1$ by \cref{pair_eq_iff}.
                            Show that $b = b'$.
                            \begin{subproof}
                                Suppose not.
                                Thus $(a,b')\in\{(0,U_0),(1,U_1)\}$ or $(a,b')\in\{(r_1,U_r)\}$ by \cref{union_iff}.
                                \begin{byCase}
                                \caseOf{$(a,b')\in\{(0,U_0),(1,U_1)\}$.}
                                    Thus $(a,b') = (0,U_0)$ or $(a,b') = (1,U_1)$ by \cref{upair_elim}.
                                    \begin{byCase}
                                    \caseOf{$(a,b') = (1,U_1)$.}
                                        Thus $b = b'$ by \cref{pair_eq_iff}.
                                        Contradiction.
                                    \caseOf{$(a,b') = (0,U_0)$.}
                                        Thus $a = 0$ by \cref{pair_eq_iff}.
                                        Thus $1\neq 0$ by \cref{real_mul_ident}.
                                        Contradiction.
                                    \end{byCase}
                                \caseOf{$(a,b')\in\{(r_1,U_r)\}$.}
                                    Thus $(a,b') = (r_1,U_r)$ by \cref{singleton_elim}.
                                    Thus $a = r_1$ by \cref{pair_eq_iff}.
                                    Thus $r_1 \neq 1$.
                                    Contradiction.
                                \end{byCase}
                            \end{subproof}
                        \caseOf{$(a,b) = (r_1,U_r)$.}
                            Thus $a = r_1$ by \cref{pair_eq_iff}.
                            Show that $b = b'$.
                            \begin{subproof}
                                Suppose not.
                                Thus $(a,b')\in\{(0,U_0),(1,U_1)\}$ or $(a,b')\in\{(r_1,U_r)\}$ by \cref{union_iff}.
                                \begin{byCase}
                                \caseOf{$(a,b')\in\{(r_1,U_r)\}$.}
                                    Thus $(a,b') = (r_1,U_r)$ by \cref{singleton_elim}.
                                    Thus $b = b'$ by \cref{pair_eq_iff}.
                                    Contradiction.
                                \caseOf{$(a,b')\in\{(0,U_0),(1,U_1)\}$.}
                                    Thus $(a,b') = (0,U_0)$ or $(a,b') = (1,U_1)$ by \cref{upair_elim}.
                                    \begin{byCase}
                                    \caseOf{$(a,b') = (0,U_0)$.}
                                        Thus $a = 0$ by \cref{pair_eq_iff}.
                                        Thus $r_1 \neq 0$.
                                        Contradiction.
                                    \caseOf{$(a,b') = (1,U_1)$.}
                                        Thus $a = 1$ by \cref{pair_eq_iff}.
                                        Thus $r_1 \neq 1$.
                                        Contradiction.
                                    \end{byCase}
                                \end{byCase}
                            \end{subproof}
                        \end{byCase}
                \end{subproof}
                Hence $u_1$ is right-unique by \cref{rightunique}.
                Therefore $u_1$ is a function by \cref{union_iff,upair_elim,singleton_elim,relation,rightunique}.
                        For all $x\in\dom{u_1}$ we have $x\in\{0,1\}\union\{r_1\}$
                            by \cref{dom_iff,union_iff,upair_elim,singleton_elim,pair_eq_iff,upair_intro_left,upair_intro_right,singleton_intro,union_intro_left,union_intro_right}.
                        Thus $\dom{u_1}\subseteq \{0,1\}\union\{r_1\}$ by \cref{subseteq}.
                        For all $x\in\{0,1\}\union\{r_1\}$ we have $x\in\dom{u_1}$
                            by \cref{union_iff,upair_elim,singleton_elim,upair_intro_left,upair_intro_right,singleton_intro,union_intro_left,union_intro_right,dom_intro}.
                        Thus $\{0,1\}\union\{r_1\}\subseteq \dom{u_1}$ by \cref{subseteq}.
                        Therefore $\dom{u_1} = \{0,1\}\union\{r_1\}$ by \cref{subseteq_antisymmetric}.
                        Hence $\{0,1\}\union \img{s}{\initialSegment{1}}\subseteq \{0,1\}\union\{r_1\}$ and
                            $\{0,1\}\union\{r_1\}\subseteq \{0,1\}\union \img{s}{\initialSegment{1}}$
                            by \cref{upair_intro_right,subseteq_refl,singleton_subset_intro,subseteq_transitive,union_subsets_is_subset,union_intro_left,union_intro_right,subseteq}.
                        Therefore $\{0,1\}\union \img{s}{\initialSegment{1}} = \{0,1\}\union\{r_1\}$ by \cref{subseteq_antisymmetric}.
                    Hence $\dom{u_1} = \{0,1\}\union \img{s}{\initialSegment{1}}$ by \cref{subseteq_antisymmetric}.
                Thus for all $p\in\dom{u_1}$ we have $u_1(p)\in\pow{X}$
                    by \cref{subseteq,union_iff,upair_elim,upair_intro_left,upair_intro_right,singleton_elim,singleton_intro,union_intro_left,union_intro_right,function_apply_intro,open_in,topology_on,elem_subseteq}.
                Therefore $u_1\in\funs{\{0,1\}\union \img{s}{\initialSegment{1}}}{\pow{X}}$ by \cref{funs_intro}.
            \end{subproof}
            Show that for all $p\in\dom{u_1}$ we have
                $u_1(p)\in T$ and
                $u_1(1) = X\setminus B$ and
                $A\subseteq u_1(0)$ and
                $\closure{u_1(0)}{X}{T}\subseteq u_1(1)$ and
                for all $p\in\dom{u_1}$ for all $q\in\dom{u_1}$ if $p < q$ then
                    $\closure{u_1(p)}{X}{T}\subseteq u_1(q)$ and
                for all $p\in\dom{u_1}$ if $1 < p$ then $u_1(p) = X$.
            \begin{subproof}
                Fix $p$ such that $p\in\dom{u_1}$.
                Now $u_1$ is a function and $u_1$ is right-unique and $u_1$ is a relation
                    by \cref{funs_is_function,function_rightunique,funs_is_relation}.
                Show that $u_1(p)\in T$.
                \begin{subproof}
                    Thus $p\in\{0,1\}\union\{r_1\}$.
                    Thus $p\in\{0,1\}$ or $p\in\{r_1\}$ by \cref{union_iff}.
                    \begin{byCase}
                    \caseOf{$p\in\{0,1\}$.}
                        Thus $p = 0$ or $p = 1$ by \cref{upair_elim}.
                        \begin{byCase}
                        \caseOf{$p = 0$.}
                            Thus $u_1(p)\in T$ by \cref{upair_intro_left,union_intro_left,function_apply_intro,open_in}.
                        \caseOf{$p = 1$.}
                            Thus $u_1(p)\in T$ by \cref{upair_intro_right,union_intro_left,function_apply_intro,open_in}.
                        \end{byCase}
                    \caseOf{$p\in\{r_1\}$.}
                        Thus $u_1(p)\in T$ by \cref{singleton_elim,singleton_intro,union_intro_right,function_apply_intro,open_in}.
                    \end{byCase}
                \end{subproof}
                Then $u_1(1) = X\setminus B$ and $A\subseteq u_1(0)$ and
                    $\closure{u_1(0)}{X}{T}\subseteq u_1(1)$
                    by \cref{upair_intro_right,upair_intro_left,union_intro_left,function_apply_intro,subseteq}.
                Show that for all $p\in\dom{u_1}$ for all $q\in\dom{u_1}$ if $p < q$ then
                    $\closure{u_1(p)}{X}{T}\subseteq u_1(q)$.
                \begin{subproof}
                    Fix $p\in\dom{u_1}$.
                    Fix $q\in\dom{u_1}$.
                    Assume $p < q$.
                    Thus $p\in\{0,1\}\union\{r_1\}$ and $q\in\{0,1\}\union\{r_1\}$.
                    Show that $q \neq 0$.
                    \begin{subproof}
                        Suppose not.
                        Thus $q = 0$.
                        Thus $p < 0$ by \cref{real_lt_subst_right}.
                        Thus $p\in\{0,1\}$ or $p\in\{r_1\}$ by \cref{union_iff}.
                        \begin{byCase}
                        \caseOf{$p\in\{0,1\}$.}
                            Thus $p = 0$ or $p = 1$ by \cref{upair_elim}.
                            \begin{byCase}
                            \caseOf{$p = 0$.}
                                Thus $0 < 0$ and $0\in\reals$ by \cref{real_lt_subst_left,real_add_ident}.
                                Contradiction by \cref{real_lt_irreflexive}.
                            \caseOf{$p = 1$.}
                                Thus $1 < 0$ by \cref{real_lt_subst_left}.
                                $0\in\reals$ and $1\in\reals$ and $0 < 1$ by \cref{real_add_ident,real_mul_ident,real_zero_lt_one}.
                                Thus $1 < 1$ by \cref{real_lt_trans}.
                                Contradiction by \cref{real_lt_irreflexive}.
                            \end{byCase}
                        \caseOf{$p\in\{r_1\}$.}
                            Thus $r_1 < 0$ by \cref{singleton_elim,real_lt_subst_left}.
                            $0\in\reals$ and $0 < r_1$ by \cref{real_add_ident,unit_rational_pos}.
                            Thus $0 < 0$ by \cref{real_lt_trans}.
                            Contradiction by \cref{real_lt_irreflexive}.
                        \end{byCase}
                    \end{subproof}
                    Hence $q\in\{1\}\union\{r_1\}$.
                    Then $q\in\{1\}$ or $q\in\{r_1\}$ by \cref{union_iff}.
                    \begin{byCase}
                    \caseOf{$q\in\{r_1\}$.}
                        Thus $q = r_1$ by \cref{singleton_elim}.
                        Show that $p = 0$.
                        \begin{subproof}
                            Suppose not.
                            Thus $p\in\{0,1\}$ or $p\in\{r_1\}$ by \cref{union_iff}.
                            \begin{byCase}
                            \caseOf{$p\in\{0,1\}$.}
                                Thus $p = 0$ or $p = 1$ by \cref{upair_elim}.
                                \begin{byCase}
                                \caseOf{$p = 1$.}
                                    Thus $1 < r_1$ by \cref{real_lt_subst_left}.
                                    Contradiction.
                                \caseOf{$p = 0$.}
                                    Contradiction.
                                \end{byCase}
                            \caseOf{$p\in\{r_1\}$.}
                                Thus $r_1 < r_1$ by \cref{singleton_elim,real_lt_subst_left}.
                                Contradiction by \cref{real_lt_irreflexive}.
                            \end{byCase}
                        \end{subproof}
                        Then $p = 0$.
                        Hence $\closure{u_1(p)}{X}{T}\subseteq u_1(q)$
                            by \cref{singleton_elim,pair_eq_iff,union_intro_left,upair_intro_left,union_intro_right,singleton_intro,function_apply_intro,subseteq}.
                    \caseOf{$q\in\{1\}$.}
                        Thus $q = 1$ by \cref{singleton_elim}.
                        Thus $p\in\{0,1\}$ or $p\in\{r_1\}$ by \cref{union_iff}.
                        \begin{byCase}
                        \caseOf{$p\in\{0,1\}$.}
                            Suppose not.
                            Thus $p = 1$ by \cref{upair_elim}.
                            Thus $1 < 1$ and $1\in\reals$ by \cref{real_lt_subst_left,real_mul_ident}.
                            Contradiction by \cref{real_lt_irreflexive}.
                            Then $p = 0$.
                            Hence $\closure{u_1(p)}{X}{T}\subseteq u_1(q)$
                                by \cref{pair_eq_iff,union_intro_left,upair_intro_left,upair_intro_right,function_apply_intro,subseteq}.
                        \caseOf{$p\in\{r_1\}$.}
                            Thus $\closure{u_1(p)}{X}{T}\subseteq u_1(q)$
                                by \cref{singleton_elim,pair_eq_iff,union_intro_right,singleton_intro,union_intro_left,upair_intro_right,function_apply_intro,subseteq}.
                        \end{byCase}
                    \end{byCase}
                \end{subproof}
                Show that for all $p\in\dom{u_1}$ if $1 < p$ then $u_1(p) = X$.
                \begin{subproof}
                    Fix $p\in\dom{u_1}$.
                    Assume $1 < p$.
                    Suppose not.
                    Thus $p\in\{0,1\}\union\{r_1\}$.
                    Thus $p\in\{0,1\}$ or $p\in\{r_1\}$ by \cref{union_iff}.
                    \begin{byCase}
                    \caseOf{$p\in\{0,1\}$.}
                        Thus $p = 0$ or $p = 1$ by \cref{upair_elim}.
                        \begin{byCase}
                        \caseOf{$p = 0$.}
                            Thus $1 < 0$ by \cref{real_lt_subst_left}.
                            $0\in\reals$ and $1\in\reals$ and $0 < 1$ by \cref{real_add_ident,real_mul_ident,real_zero_lt_one}.
                            Thus $1 < 1$ by \cref{real_lt_trans}.
                            Contradiction by \cref{real_lt_irreflexive}.
                        \caseOf{$p = 1$.}
                            Thus $1 < 1$ and $1\in\reals$ by \cref{real_lt_subst_left,real_mul_ident}.
                            Contradiction by \cref{real_lt_irreflexive}.
                        \end{byCase}
                    \caseOf{$p\in\{r_1\}$.}
                        Thus $1 < r_1$ by \cref{singleton_elim,real_lt_subst_left}.
                        Contradiction.
                    \end{byCase}
                \end{subproof}
            \end{subproof}
            Hence $u_1\in\urysohnAdmissible{s}{X}{T}{A}{B}{1}$ by \cref{urysohn_admissible}.
        \end{subproof}
        Therefore $(1,u_1)\in G$ by \cref{urysohn_good_pair_intro,real_one_in_naturals}.
    \end{subproof}
\end{proof}

% from §33 helper lemma 11: initial good pair when $r_1 < 1$
\begin{lemma}\label{urysohn_u1_lt1}
    Assume $X$ is normal with respect to $T$.
    Assume $U_1 = X\setminus B$.
    Assume $U_1$ is open in $X$ with respect to $T$.
    Assume $U_0$ is open in $X$ with respect to $T$.
    Assume $A\subseteq U_0$.
    Assume $\closure{U_0}{X}{T}\subseteq U_1$.
    Assume $r_1\in\reals$.
    Assume $r_1\in\unitRationals$.
    Assume $r_1 < 1$.
    Assume $\img{s}{\initialSegment{1}} = \{r_1\}$.
    Let $G = \urysohnGoodPairs{s}{X}{T}{A}{B}$.
    Then there exists $u_1$ such that $(1,u_1)\in G$.
\end{lemma}
\begin{proof}
    Now $\closure{U_0}{X}{T}$ is closed in $X$ with respect to $T$ by \cref{open_in,topology_on,elem_subseteq,pow_iff,closure_closed_in}.
    Take $U_r$ such that
        $U_r$ is open in $X$ with respect to $T$ and
        $\closure{U_0}{X}{T}\subseteq U_r$ and
        $\closure{U_r}{X}{T}\subseteq U_1$
        by \cref{normal_closure_neighborhood}.
    Therefore there exists $u_1$ such that $(1,u_1)\in G$ by \cref{urysohn_u1_lt1_admissible}.
\end{proof}

% from §33 helper lemma 13: admissible chain is linearly ordered by inclusion
\begin{lemma}\label{urysohn_unit_interval_chain}
    Assume $s$ is a sequence in $\unitRationals$.
    Let $G = \urysohnGoodPairs{s}{X}{T}{A}{B}$.
    Let $R = \{(z,w) \mid z\in G, w\in G \mid
        \fst{w} = \fst{z} + 1 \land \snd{z}\subseteq\snd{w}\}$.
    Assume $c$ is a function on $G$.
    Assume for all $z\in G$ we have $z\mathrel{R} \apply{c}{z}$.
    Assume $u$ is a sequence in $G$.
    Assume for all $n\in\naturalsPlus$ we have $u(n + 1) = \apply{c}{u(n)}$.
    Let $F = \{\snd{u(n)} \mid n\in\naturalsPlus\}$.
    Then for all $f,g$ such that $f\in F$ and $g\in F$ we have $f\subseteq g$ or $g\subseteq f$.
\end{lemma}
\begin{proof}
    For all $n\in\naturalsPlus$ we have $\snd{u(n)}\subseteq \snd{u(n + 1)}$
        by \cref{sequence_in,funs_type_apply,pair_eq_iff}.
    Let $A_0 = \{m\in\naturalsPlus \mid \text{for all $n\in\naturalsPlus$ if $n \leq m$ then
        $\snd{u(n)}\subseteq \snd{u(m)}$}\}$.
    $A_0\subseteq\naturalsPlus$ by \cref{subseteq}.
    Hence $1\in A_0$
        by \cref{real_one_in_naturals,real_naturals_subset_reals,subseteq,real_one_leq_naturalsplus,real_leq_antisym,subseteq_refl}.
    Show that for all $k\in A_0$ we have $k + 1\in A_0$.
    \begin{subproof}
        Fix $k$ such that $k\in A_0$.
        Thus $k\in\naturalsPlus$ by \cref{subseteq}.
        Show that for all $n\in\naturalsPlus$ if $n \leq k + 1$ then $\snd{u(n)}\subseteq \snd{u(k + 1)}$.
        \begin{subproof}
            Fix $n\in\naturalsPlus$.
            Assume $n \leq k + 1$.
            \begin{byCase}
            \caseOf{$n = k + 1$.}
                Then $\snd{u(n)}\subseteq \snd{u(k + 1)}$ by \cref{subseteq_refl}.
            \caseOf{$n \neq k + 1$.}
                Therefore $\snd{u(n)}\subseteq \snd{u(k + 1)}$
                    by \cref{real_naturals_leq_succ_elim,subseteq_transitive}.
            \end{byCase}
        \end{subproof}
        Therefore $k + 1\in A_0$.
    \end{subproof}
    Hence for all $m\in\naturalsPlus$ we have $m\in A_0$ by \cref{real_naturals_induction}.
    Fix $f$.
    Fix $g$ such that $f\in F$ and $g\in F$.
    Take $n$ such that $n\in\naturalsPlus$ and $f = \snd{u(n)}$.
    Take $m$ such that $m\in\naturalsPlus$ and $g = \snd{u(m)}$.
    Thus $n < m$ or $n = m$ or $m < n$ by \cref{real_naturals_subset_reals,subseteq,real_lt_trichotomy}.
    \begin{byCase}
    \caseOf{$n = m$.}
        Then $f\subseteq g$ by \cref{subseteq_refl}.
    \caseOf{$n < m$.}
        Hence $f\subseteq g$ by \cref{real_lt_imp_leq}.
    \caseOf{$m < n$.}
        Hence $g\subseteq f$ by \cref{real_lt_imp_leq}.
    \end{byCase}
\end{proof}

% from §33 helper lemma 14: unions of the admissible chain are a function
\begin{lemma}\label{urysohn_unit_interval_union_function}
    Assume $s$ is a sequence in $\unitRationals$.
    Let $G = \urysohnGoodPairs{s}{X}{T}{A}{B}$.
    Let $R = \{(z,w) \mid z\in G, w\in G \mid
        \fst{w} = \fst{z} + 1 \land \snd{z}\subseteq\snd{w}\}$.
    Assume $c$ is a function on $G$.
    Assume for all $z\in G$ we have $z\mathrel{R} \apply{c}{z}$.
    Assume $u$ is a sequence in $G$.
    Assume for all $n\in\naturalsPlus$ we have $u(n + 1) = \apply{c}{u(n)}$.
    Assume for all $n\in\naturalsPlus$ we have $\fst{u(n)} = n$.
    Let $F = \{\snd{u(n)} \mid n\in\naturalsPlus\}$.
    Let $U = \unions{F}$.
    Then $U$ is a function.
\end{lemma}
\begin{proof}
    Hence $U$ is a function
        by \cref{sequence_in,funs_type_apply,urysohn_good_pair_decomp,snd_eq,urysohn_admissible,funs_is_function,urysohn_unit_interval_chain,unions_of_compatible_family_of_function_is_function}.
\end{proof}
% from §33 helper lemma 15: Urysohn function is a function to reals
\begin{lemma}\label{urysohn_unit_interval_f_funs}
    Assume $s$ is a sequence in $\unitRationals$.
    Assume for all $p$ such that $p\in\unitRationals$ there exists $n\in\naturalsPlus$ such that $s(n) = p$.
    Let $G = \urysohnGoodPairs{s}{X}{T}{A}{B}$.
    Assume $u$ is a sequence in $G$.
    Assume for all $n\in\naturalsPlus$ we have $\fst{u(n)} = n$.
    Let $F = \{\snd{u(n)} \mid n\in\naturalsPlus\}$.
    Let $U = \unions{F}$.
    Assume $U$ is a function.
    Let $f = \{(x,y) \mid x\in X, y\in\reals \mid
        \text{$y$ is an infimum of $\urysohnLevelSet{U}{x}$}\}$.
    Then $f$ is a function from $X$ to $\reals$.
\end{lemma}
\begin{proof}
        Show that for all $x,y_1,y_2$ such that $(x,y_1)\in f$ and $(x,y_2)\in f$ we have $y_1 = y_2$.
        \begin{subproof}
            Fix $x$.
            Let $Q_x = \urysohnLevelSet{U}{x}$.
            Fix $y_1$.
            Fix $y_2$.
            Assume $(x,y_1)\in f$ and $(x,y_2)\in f$.
            Then $y_1$ and $y_2$ are infimums of $Q_x$.
            Hence $y_1 = y_2$ by \cref{real_infimum_unique}.
        \end{subproof}
        Therefore $f$ is right-unique by \cref{rightunique}.
        Hence $f$ is a function by \cref{relation,rightunique}.
            Show that for all $x\in X$ we have $x\in\dom{f}$.
            \begin{subproof}
                Fix $x\in X$.
                Let $Q_x = \urysohnLevelSet{U}{x}$.
                Show that $Q_x\neq \emptyset$.
                \begin{subproof}
                    $0\in\reals$ and $1\in\reals$ and $0 < 1$
                        by \cref{real_add_ident,real_mul_ident,real_zero_lt_one}.
                    Hence $0 \leq 1$ by \cref{real_lt_imp_leq}.
                    Let $b = 1 + 1$.
                    Then $b\in\reals$ and $1 < b$
                        by \cref{real_add_closed,real_lt_add,real_add_ident,real_lt_subst_left}.
                    Take $p$ such that $p\in\unitRationals$ and $p\in\intervalopen{1}{b}$
                        by \cref{unit_rationals_dense_nonneg}.
                    Thus $1 < p$ by \cref{intervalopen}.
                    Take $k\in\naturalsPlus$ such that $s(k) = p$ by assumption.
                    Thus $u(k)\in G$ by \cref{sequence_in,funs_type_apply}.
                    Take $n,t$ such that
                        $t\in\urysohnAdmissible{s}{X}{T}{A}{B}{n}$ and
                        $u(k) = (n,t)$
                        by \cref{urysohn_good_pair_decomp}.
                    Thus $k = n$ and $\snd{u(k)} = t$ by \cref{fst_eq,snd_eq}.
                    Thus $t$ is a function and $\dom{t} = \{0,1\}\union \img{s}{\initialSegment{k}}$
                        by \cref{urysohn_admissible,funs_is_function,funs_elim}.
                    Thus $p\in\dom{t}$ by \cref{sequence_in,funs_is_function,funs_elim,initial_segment_naturalsplus,subseteq,function_apply_elim,img_iff,union_intro_right}.
                    $1\in\dom{t}$ by \cref{upair_intro_right,union_intro_left,subseteq}.
                    Thus for all $q\in\dom{t}$ if $1 < q$ then $t(q) = X$
                        by \cref{urysohn_admissible}.
                    Thus $t(p) = X$.
                    Thus $(p,t(p))\in U$ by \cref{function_apply_elim,unions_intro}.
                    Thus $U(p) = t(p)$ by \cref{function_apply_intro}.
                    Thus $x\in U(p)$.
                    Thus $p\in Q_x$ by \cref{unit_rationals_zero,union_intro_left,urysohn_level_set}.
                    Therefore $Q_x\neq \emptyset$.
                \end{subproof}
                Show that for all $p\in Q_x$ we have $p\in\reals$.
                \begin{subproof}
                    Fix $p\in Q_x$.
                    Thus $p\in\unitRationals$ or $p\in\{0\}$ by \cref{urysohn_level_set,unit_rationals_zero,union_iff}.
                    \begin{byCase}
                    \caseOf{$p\in\unitRationals$.}
                        Take $q$ such that $q\in\naturalsPlus\times\naturalsPlus$ and
                            $p = \fst{q} \rmul \inv{\snd{q}}$
                            by \cref{unit_rationals}.
                        Thus $p\in\reals$
                            by \cref{fst_type,snd_type,real_naturals_subset_reals,subseteq,real_add_ident,real_mul_ident,real_zero_lt_one,real_one_leq_naturalsplus,real_lt_subst_right,real_lt_trans,real_pos_nonzero,real_mul_inverse,real_mul_closed}.
                    \caseOf{$p\in\{0\}$.}
                        Thus $p\in\reals$ by \cref{singleton_elim,real_add_ident}.
                    \end{byCase}
                \end{subproof}
                Hence $Q_x\subseteq\reals$ by \cref{subseteq}.
                For all $p\in Q_x$ we have $0 \leq p$
                    by \cref{urysohn_level_set,unit_rationals_zero,union_iff,singleton_elim,real_add_ident,unit_rational_pos,real_lt_imp_leq}.
                Hence $0$ is a lower bound of $Q_x$
                    by \cref{real_add_ident,subseteq,real_lower_bound}.
                Hence $Q_x$ is bounded below.
                Take $p$ such that $p$ is an infimum of $Q_x$ by \cref{real_infimum_exists}.
                Thus $x\in\dom{f}$ by \cref{real_infimum,real_lower_bound,dom_intro}.
            \end{subproof}
            Therefore $X\subseteq\dom{f}$ by \cref{subseteq}.
        Hence $\dom{f}\subseteq X$ by \cref{dom_iff,pair_eq_iff,subseteq}.
        Therefore $\dom{f} = X$ by \cref{subseteq_antisymmetric}.
        Show that for all $x\in\dom{f}$ we have $f(x)\in\reals$.
        \begin{subproof}
            Fix $x\in\dom{f}$.
            Let $Q_x = \urysohnLevelSet{U}{x}$.
            Take $p$ such that $(x,p)\in f$ by \cref{dom_iff}.
            Thus $p$ is an infimum of $Q_x$.
            Thus $f(x) = p$ and $f(x)\in\reals$
                by \cref{real_infimum,real_lower_bound,function_apply_intro}.
        \end{subproof}
    Hence $f\in\funs{X}{\reals}$ by \cref{funs_intro}.
\end{proof}
% from §33 helper lemma 16: existence of a first good pair
\begin{lemma}\label{urysohn_unit_interval_u1}
    Assume $X$ is normal with respect to $T$.
    Assume $A$ is closed in $X$ with respect to $T$.
    Assume $B$ is closed in $X$ with respect to $T$.
    Assume $A$ is disjoint from $B$.
    Assume $s$ is a sequence in $\unitRationals$.
    Let $G = \urysohnGoodPairs{s}{X}{T}{A}{B}$.
    Then there exists $u_1$ such that $(1,u_1)\in G$.
\end{lemma}
\begin{proof}
    Let $U_1 = X\setminus B$.
    $U_1$ is open in $X$ with respect to $T$ by \cref{closed_in}.
    Hence $A\subseteq U_1$ by \cref{disjoint,closed_in,subseteq,setminus_intro}.
    Take $U_0$ such that
        $U_0$ is open in $X$ with respect to $T$ and
        $A\subseteq U_0$ and
        $\closure{U_0}{X}{T}\subseteq U_1$
        by \cref{normal_closure_neighborhood}.
    Let $r_1 = s(1)$.
    $s\in\funs{\naturalsPlus}{\unitRationals}$ and $1\in\naturalsPlus$
        by \cref{sequence_in,real_one_in_naturals}.
    Therefore $r_1\in\unitRationals$ by \cref{funs_type_apply}.
    Hence $\img{s}{\initialSegment{1}}\subseteq \{r_1\}$
        by \cref{img_iff,initial_segment_naturalsplus,real_one_leq_naturalsplus,real_naturals_subset_reals,subseteq,real_leq_antisym,funs_is_function,function_apply_intro,singleton_intro}.
    Hence $\{r_1\}\subseteq \img{s}{\initialSegment{1}}$
        by \cref{sequence_in,funs_is_function,funs_elim,initial_segment_naturalsplus,subseteq,function_apply_elim,img_iff,real_one_in_naturals,singleton_subset_intro}.
    Therefore $\img{s}{\initialSegment{1}} = \{r_1\}$ by \cref{subseteq_antisymmetric}.
    \begin{byCase}
    \caseOf{$r_1 = 1$.}
        Thus there exists $u_1$ such that $(1,u_1)\in G$ by \cref{urysohn_u1_eq1}.
    \caseOf{$r_1 \neq 1$.}
        Take $q$ such that
            $q\in\naturalsPlus\times\naturalsPlus$ and
            $r_1 = \fst{q} \rmul \inv{\snd{q}}$
            by \cref{unit_rationals}.
        Thus $r_1\in\reals$
            by \cref{fst_type,snd_type,real_naturals_subset_reals,subseteq,real_add_ident,real_mul_ident,real_zero_lt_one,real_one_leq_naturalsplus,real_lt_subst_right,real_lt_trans,real_pos_nonzero,real_mul_inverse,real_mul_closed}.
        Then $r_1 < 1$ or $1 < r_1$ by \cref{real_mul_ident,real_lt_trichotomy}.
        \begin{byCase}
        \caseOf{$1 < r_1$.}
            Thus there exists $u_1$ such that $(1,u_1)\in G$ by \cref{urysohn_u1_gt1}.
        \caseOf{$r_1 < 1$.}
            Thus there exists $u_1$ such that $(1,u_1)\in G$ by \cref{urysohn_u1_lt1}.
        \end{byCase}
    \end{byCase}
\end{proof}

% from §33 helper lemma 17: indices of the iteration sequence
\begin{lemma}\label{urysohn_unit_interval_fst_u}
    Assume $s$ is a sequence in $\unitRationals$.
    Let $G = \urysohnGoodPairs{s}{X}{T}{A}{B}$.
    Let $R = \{(z,w) \mid z\in G, w\in G \mid
        \fst{w} = \fst{z} + 1 \land \snd{z}\subseteq\snd{w}\}$.
    Assume $c$ is a function on $G$.
    Assume for all $z\in G$ we have $z\mathrel{R} \apply{c}{z}$.
    Assume $u$ is a sequence in $G$.
    Assume $u(1) = a_0$.
    Assume $a_0 = (1,u_1)$.
    Assume for all $n\in\naturalsPlus$ we have $u(n + 1) = \apply{c}{u(n)}$.
    Then for all $n\in\naturalsPlus$ we have $\fst{u(n)} = n$.
\end{lemma}
\begin{proof}
    Let $A_0 = \{n\in\naturalsPlus \mid \fst{u(n)} = n\}$.
    $A_0\subseteq\naturalsPlus$ by \cref{subseteq}.
    Hence $1\in A_0$ by \cref{real_one_in_naturals,fst_eq}.
    Show that for all $k\in A_0$ we have $k + 1\in A_0$.
    \begin{subproof}
        Fix $k\in A_0$.
        Hence $(u(k),u(k + 1))\in R$ by \cref{subseteq,sequence_in,funs_type_apply}.
        Therefore $k + 1\in A_0$ by \cref{pair_eq_iff,real_naturals_closed_succ}.
    \end{subproof}
    Therefore for all $n\in\naturalsPlus$ we have $n\in A_0$ by \cref{real_naturals_induction}.
    Hence for all $n\in\naturalsPlus$ we have $\fst{u(n)} = n$.
\end{proof}

% from §33 helper lemma 18: continuous subspace range
\begin{lemma}\label{urysohn_unit_interval_continuous_subspace}
    Assume $f\in\funs{X}{\reals}$.
    Assume $T$ is a topology on $X$.
    Assume $I_0\subseteq\reals$.
    Let $T_0 = \subspaceTopology{\standardTopologyR}{I_0}$.
    Assume $f$ is continuous from $X$ to $\reals$ with respect to $T$ and $\standardTopologyR$.
    Assume $\img{f}{X}\subseteq I_0$.
    Then $f$ is continuous from $X$ to $I_0$ with respect to $T$ and $T_0$.
\end{lemma}
\begin{proof}
    Hence $f$ is continuous from $X$ to $I_0$ with respect to $T$ and $T_0$
        by \cref{continuous_subspace_range}.
\end{proof}

% from §33 helper lemma 20a: unit rationals are real (local helper)
\begin{lemma}\label{unit_rational_real_aux}
    Assume $r\in\unitRationals$.
    Then $r\in\reals$.
\end{lemma}
\begin{proof}
    Take $q$ such that
        $q\in\naturalsPlus\times\naturalsPlus$ and
        $r = \fst{q} \rmul \inv{\snd{q}}$
        by \cref{unit_rationals}.
    Therefore $r\in\reals$
        by \cref{fst_type,snd_type,real_naturals_subset_reals,subseteq,real_add_ident,real_mul_ident,real_zero_lt_one,real_one_leq_naturalsplus,real_lt_subst_right,real_lt_trans,real_pos_nonzero,real_mul_inverse,real_mul_closed}.
\end{proof}

% from §33 helper lemma 20b: unit rationals zero are real
\begin{lemma}\label{unit_rational_zero_real}
    Assume $r\in\unitRationalsZero$.
    Then $r\in\reals$.
\end{lemma}
\begin{proof}
    Hence $r\in\reals$
        by \cref{unit_rationals_zero,union_iff,unit_rational_real_aux,singleton_elim,real_add_ident}.
\end{proof}

% from §33 helper lemma 18a: Urysohn level values are real
\begin{lemma}\label{urysohn_level_value_real}
    Assume $r\in\urysohnLevelSet{U}{x}$.
    Then $r\in\reals$.
\end{lemma}
\begin{proof}
    Hence $r\in\reals$ by \cref{urysohn_level_set,unit_rational_zero_real}.
\end{proof}

% from §33 helper lemma 20d: sequence values occur in the initial image
\begin{lemma}\label{sequence_value_in_initial_image}
    Assume $s$ is a sequence in $Y$.
    Assume $k\in\naturalsPlus$.
    Then $s(k)\in\img{s}{\initialSegment{k}}$.
\end{lemma}
\begin{proof}
    $s(k)\in\img{s}{\initialSegment{k}}$ by \cref{sequence_in,funs_is_function,funs_elim,initial_segment_naturalsplus,subseteq,function_apply_elim,img_iff}.
\end{proof}

% from §33 helper lemma 20da: rational values of U come from an admissible stage
\begin{lemma}\label{urysohn_union_rational_value_aux}
    Assume $s$ is a sequence in $\unitRationals$.
    Assume for all $p$ such that $p\in\unitRationals$ there exists $n\in\naturalsPlus$ such that $s(n) = p$.
    Let $G = \urysohnGoodPairs{s}{X}{T}{A}{B}$.
    Assume $u$ is a sequence in $G$.
    Assume for all $n\in\naturalsPlus$ we have $\fst{u(n)} = n$.
    Let $F = \{\snd{u(n)} \mid n\in\naturalsPlus\}$.
    Let $U = \unions{F}$.
    Assume $U$ is a function.
    Assume $q\in\unitRationals$.
    Then there exist $n,t$ such that
        $t\in\urysohnAdmissible{s}{X}{T}{A}{B}{n}$ and
        $t$ is a function and
        $q\in\dom{t}$ and
        $t\in F$ and
        $U(q) = t(q)$.
\end{lemma}
\begin{proof}
    Take $k$ such that $k\in\naturalsPlus$ and $s(k) = q$
        by \cref{unit_rationals_sequence}.
    Then $u(k)\in G$ by \cref{sequence_in,funs_type_apply}.
    Take $n,t$ such that
        $t\in\urysohnAdmissible{s}{X}{T}{A}{B}{n}$ and
        $u(k) = (n,t)$
        by \cref{urysohn_good_pair_decomp}.
    Hence $\fst{u(k)} = n$ and $\snd{u(k)} = t$ by \cref{fst_eq,snd_eq}.
    Then $n\in\naturalsPlus$ and $q = s(n)$ and $t\in F$ by \cref{subseteq}.
    Hence $t$ is a function by \cref{urysohn_admissible,funs_is_function}.
    Then $q\in\dom{t}$ by \cref{urysohn_admissible,sequence_value_in_initial_image,funs_elim,union_intro_right}.
    Therefore $U(q) = t(q)$ by \cref{function_apply_elim,unions_intro,function_apply_intro}.
\end{proof}

% from §33 helper lemma 18b: rational U-values are open
\begin{lemma}\label{urysohn_union_rational_open}
    Assume $X$ is normal with respect to $T$.
    Assume $A$ is closed in $X$ with respect to $T$.
    Assume $B$ is closed in $X$ with respect to $T$.
    Assume $A$ is disjoint from $B$.
    Assume $s$ is a sequence in $\unitRationals$.
    Assume for all $p$ such that $p\in\unitRationals$ there exists $n\in\naturalsPlus$ such that $s(n) = p$.
    Let $G = \urysohnGoodPairs{s}{X}{T}{A}{B}$.
    Assume $u$ is a sequence in $G$.
    Assume for all $n\in\naturalsPlus$ we have $\fst{u(n)} = n$.
    Let $F = \{\snd{u(n)} \mid n\in\naturalsPlus\}$.
    Let $U = \unions{F}$.
    Assume $U$ is a function.
    Assume $q\in\unitRationals$.
    Then $U(q)$ is open in $X$ with respect to $T$.
\end{lemma}
\begin{proof}
    Take $n,t$ such that
        $t\in\urysohnAdmissible{s}{X}{T}{A}{B}{n}$ and
        $q\in\dom{t}$ and
        $U(q) = t(q)$
        by \cref{urysohn_union_rational_value_aux}.
    Therefore $U(q)$ is open in $X$ with respect to $T$
        by \cref{urysohn_admissible,subseteq,function_apply_elim,unions_intro,function_apply_intro,normal_space,open_in}.
\end{proof}

% from §33 helper lemma 18c: zero-or-rational U-values lie in X
\begin{lemma}\label{urysohn_union_zero_subset_X_aux}
    Assume $s$ is a sequence in $\unitRationals$.
    Assume for all $p$ such that $p\in\unitRationals$ there exists $n\in\naturalsPlus$ such that $s(n) = p$.
    Let $G = \urysohnGoodPairs{s}{X}{T}{A}{B}$.
    Assume $u$ is a sequence in $G$.
    Assume for all $n\in\naturalsPlus$ we have $\fst{u(n)} = n$.
    Let $F = \{\snd{u(n)} \mid n\in\naturalsPlus\}$.
    Let $U = \unions{F}$.
    Assume $U$ is a function.
    Assume $r\in\unitRationalsZero$.
    Then $U(r)\subseteq X$.
\end{lemma}
\begin{proof}
    Then $r\in\unitRationals$ or $r\in\{0\}$ by \cref{unit_rationals_zero,union_iff}.
    \begin{byCase}
    \caseOf{$r\in\unitRationals$.}
        Take $n,t$ such that
            $t\in\urysohnAdmissible{s}{X}{T}{A}{B}{n}$ and
            $r\in\dom{t}$ and
            $U(r) = t(r)$
            by \cref{urysohn_union_rational_value_aux}.
        Thus $t\in\funs{\{0,1\}\union \img{s}{\initialSegment{n}}}{\pow{X}}$ and
            $\dom{t} = \{0,1\}\union \img{s}{\initialSegment{n}}$
            by \cref{urysohn_admissible,funs_elim}.
        Thus $r\in\{0,1\}\union \img{s}{\initialSegment{n}}$ by \cref{subseteq}.
        Hence $t(r)\subseteq X$ by \cref{funs_type_apply,powerset_elim}.
        Therefore $U(r)\subseteq X$ by \cref{subseteq}.
    \caseOf{$r\in\{0\}$.}
        Then $r = 0$ by \cref{singleton_elim}.
        $0\in\reals$ and $1\in\reals$ and $0 < 1$
            by \cref{real_add_ident,real_mul_ident,real_zero_lt_one}.
        Take $q$ such that $q\in\unitRationals$ and $q\in\intervalopen{0}{1}$
            by \cref{unit_rationals_dense_nonneg}.
        Take $k$ such that $k\in\naturalsPlus$ and $s(k) = q$.
        Then $u(k)\in G$ by \cref{sequence_in,funs_type_apply}.
        Take $n,t$ such that
            $t\in\urysohnAdmissible{s}{X}{T}{A}{B}{n}$ and
            $u(k) = (n,t)$
            by \cref{urysohn_good_pair_decomp}.
        Hence $\fst{u(k)} = n$ and $\snd{u(k)} = t$ by \cref{fst_eq,snd_eq}.
        Then $k = n$.
        Hence $t\in\funs{\{0,1\}\union \img{s}{\initialSegment{k}}}{\pow{X}}$
            by \cref{urysohn_admissible}.
        Then $t$ is a function and $\dom{t} = \{0,1\}\union \img{s}{\initialSegment{k}}$
            by \cref{funs_is_function,funs_elim}.
        Then $0\in\dom{t}$ by \cref{upair_intro_left,union_intro_left,subseteq}.
        Therefore $U(r)\subseteq X$
            by \cref{function_apply_elim,unions_intro,function_apply_intro,funs_type_apply,powerset_elim}.
    \end{byCase}
\end{proof}

% from §33 helper lemma 18d: closures of lower levels lie in higher rational levels
\begin{lemma}\label{urysohn_union_closure_subset}
    Assume $A$ is closed in $X$ with respect to $T$.
    Assume $B$ is closed in $X$ with respect to $T$.
    Assume $A$ is disjoint from $B$.
    Assume $s$ is a sequence in $\unitRationals$.
    Assume for all $p$ such that $p\in\unitRationals$ there exists $n\in\naturalsPlus$ such that $s(n) = p$.
    Let $G = \urysohnGoodPairs{s}{X}{T}{A}{B}$.
    Let $R = \{(z,w) \mid z\in G, w\in G \mid
        \fst{w} = \fst{z} + 1 \land \snd{z}\subseteq\snd{w}\}$.
    Assume $c$ is a function on $G$.
    Assume for all $z\in G$ we have $z\mathrel{R} \apply{c}{z}$.
    Assume $u$ is a sequence in $G$.
    Assume for all $n\in\naturalsPlus$ we have $u(n + 1) = \apply{c}{u(n)}$.
    Assume for all $n\in\naturalsPlus$ we have $\fst{u(n)} = n$.
    Let $F = \{\snd{u(n)} \mid n\in\naturalsPlus\}$.
    Let $U = \unions{F}$.
    Assume $U$ is a function.
    Assume $p\in\unitRationalsZero$.
    Assume $q\in\unitRationals$.
    Assume $p < q$.
    Then $\closure{U(p)}{X}{T}\subseteq U(q)$.
\end{lemma}
\begin{proof}
    Take $n_q,t_q$ such that
        $t_q\in\urysohnAdmissible{s}{X}{T}{A}{B}{n_q}$ and
        $t_q$ is a function and
        $q\in\dom{t_q}$ and
        $t_q\in F$ and
        $U(q) = t_q(q)$
        by \cref{urysohn_union_rational_value_aux}.
    Hence $(q,t_q(q))\in t_q$ by \cref{function_apply_elim}.
    Then $p\in\unitRationals$ or $p\in\{0\}$ by \cref{unit_rationals_zero,union_iff}.
    \begin{byCase}
    \caseOf{$p\in\unitRationals$.}
        Take $n_p,t_p$ such that
            $t_p\in\urysohnAdmissible{s}{X}{T}{A}{B}{n_p}$ and
            $t_p$ is a function and
            $p\in\dom{t_p}$ and
            $t_p\in F$ and
            $U(p) = t_p(p)$
            by \cref{urysohn_union_rational_value_aux}.
        Thus $(p,t_p(p))\in t_p$ by \cref{function_apply_elim}.
        Thus $t_p\subseteq t_q$ or $t_q\subseteq t_p$
            by \cref{urysohn_unit_interval_chain}.
        \begin{byCase}
        \caseOf{$t_p\subseteq t_q$.}
            Then $(p,t_p(p))\in t_q$ by \cref{subseteq}.
            Then $\closure{t_q(p)}{X}{T}\subseteq t_q(q)$ by \cref{dom_intro,urysohn_admissible}.
            Hence $U(p) = t_q(p)$ by \cref{function_apply_elim,function_apply_intro}.
            Therefore $\closure{U(p)}{X}{T}\subseteq U(q)$.
        \caseOf{$t_q\subseteq t_p$.}
            Then $(q,t_q(q))\in t_p$ by \cref{subseteq}.
            Then $\closure{t_p(p)}{X}{T}\subseteq t_p(q)$ by \cref{dom_intro,urysohn_admissible}.
            Hence $U(q) = t_p(q)$ by \cref{function_apply_elim,function_apply_intro}.
            Therefore $\closure{U(p)}{X}{T}\subseteq U(q)$.
        \end{byCase}
    \caseOf{$p\in\{0\}$.}
        Then $p = 0$ by \cref{singleton_elim}.
        Then $0\in\dom{t_q}$ by \cref{urysohn_admissible,funs_elim,upair_intro_left,union_intro_left}.
        Hence $\closure{t_q(0)}{X}{T}\subseteq t_q(q)$ by \cref{urysohn_admissible}.
        Hence $(0,t_q(0))\in U$ by \cref{function_apply_elim,unions_intro}.
        Therefore $U(p) = t_q(0)$ by \cref{function_apply_intro}.
        Hence $\closure{U(p)}{X}{T}\subseteq U(q)$ by \cref{subseteq}.
    \end{byCase}
\end{proof}

% from §33 helper lemma 20c: value of the Urysohn infimum function
\begin{lemma}\label{urysohn_infimum_value}
    Let $f = \{(x,y) \mid x\in X, y\in\reals \mid
        \text{$y$ is an infimum of $\urysohnLevelSet{U}{x}$}\}$.
    Assume $f\in\funs{X}{\reals}$.
    Assume $x\in X$.
    Then $f(x)$ is an infimum of $\urysohnLevelSet{U}{x}$.
\end{lemma}
\begin{proof}
    Take $x_0,y_0$ such that
        $(x,f(x)) = (x_0,y_0)$ and
        $y_0$ is an infimum of $\urysohnLevelSet{U}{x_0}$
        by \cref{funs_is_function,funs_elim,subseteq,function_apply_elim}.
    Hence $f(x)$ is an infimum of $\urysohnLevelSet{U}{x}$ by \cref{pair_eq_iff}.
\end{proof}

% from §33 helper lemma 19a: infimum below a rational level implies membership
\begin{lemma}\label{urysohn_infimum_lt_rational_in_union}
    Assume $X$ is normal with respect to $T$.
    Assume $A$ is closed in $X$ with respect to $T$.
    Assume $B$ is closed in $X$ with respect to $T$.
    Assume $A$ is disjoint from $B$.
    Assume $s$ is a sequence in $\unitRationals$.
    Assume for all $p$ such that $p\in\unitRationals$ there exists $n\in\naturalsPlus$ such that $s(n) = p$.
    Let $G = \urysohnGoodPairs{s}{X}{T}{A}{B}$.
    Let $R = \{(z,w) \mid z\in G, w\in G \mid
        \fst{w} = \fst{z} + 1 \land \snd{z}\subseteq\snd{w}\}$.
    Assume $c$ is a function on $G$.
    Assume for all $z\in G$ we have $z\mathrel{R} \apply{c}{z}$.
    Assume $u$ is a sequence in $G$.
    Assume $u(1) = a_0$.
    Assume $a_0 = (1,u_1)$.
    Assume for all $n\in\naturalsPlus$ we have $u(n + 1) = \apply{c}{u(n)}$.
    Assume for all $n\in\naturalsPlus$ we have $\fst{u(n)} = n$.
    Let $F = \{\snd{u(n)} \mid n\in\naturalsPlus\}$.
    Let $U = \unions{F}$.
    Assume $U$ is a function.
    Let $Q_x = \urysohnLevelSet{U}{x}$.
    Assume $m$ is an infimum of $Q_x$.
    Assume $q\in\unitRationals$.
    Assume $q\in\reals$.
    Assume $m < q$.
    Then $x\in U(q)$.
\end{lemma}
\begin{proof}
    Hence $m$ is a lower bound of $Q_x$ by \cref{real_infimum}.
    Then $m \leq q$ by \cref{real_lower_bound}.
    Hence $m < q$ by \cref{real_lt_subst_left}.
    Therefore it is wrong that $q$ is a lower bound of $Q_x$.
    Therefore there exists $r$ such that $r\in Q_x$ and $r < q$.
    Take $r$ such that $r\in Q_x$ and $r < q$.
    Then $x\in U(r)$ and $r\in\unitRationalsZero$ by \cref{urysohn_level_set}.
    Hence $\closure{U(r)}{X}{T}\subseteq U(q)$ by \cref{urysohn_union_closure_subset}.
    Therefore $x\in U(q)$ by \cref{urysohn_union_zero_subset_X_aux,subseteq_closure,normal_space,subseteq}.
\end{proof}


% from §33 helper lemma 19: continuity at zero
\begin{lemma}\label{urysohn_unit_interval_continuous_zero}
    Assume $X$ is normal with respect to $T$ and $A$ is closed in $X$ with respect to $T$ and
        $B$ is closed in $X$ with respect to $T$ and $A$ is disjoint from $B$.
    Assume $s$ is a sequence in $\unitRationals$ and for all $p$ such that $p\in\unitRationals$
        there exists $n\in\naturalsPlus$ such that $s(n) = p$.
    Let $G = \urysohnGoodPairs{s}{X}{T}{A}{B}$.
    Let $R = \{(z,w) \mid z\in G, w\in G \mid
        \fst{w} = \fst{z} + 1 \land \snd{z}\subseteq\snd{w}\}$.
    Assume $c$ is a function on $G$ and for all $z\in G$ we have $z\mathrel{R} \apply{c}{z}$.
    Assume $u$ is a sequence in $G$ and $u(1) = a_0$ and $a_0 = (1,u_1)$ and
        for all $n\in\naturalsPlus$ we have $u(n + 1) = \apply{c}{u(n)}$ and $\fst{u(n)} = n$.
    Let $F = \{\snd{u(n)} \mid n\in\naturalsPlus\}$.
    Let $U = \unions{F}$.
    Assume $U$ is a function.
    Let $I_0 = \intervalclosed{0}{1}$.
    Let $f = \{(x,y) \mid x\in X, y\in\reals \mid
        \text{$y$ is an infimum of $\urysohnLevelSet{U}{x}$}\}$.
    Assume $f\in\funs{X}{\reals}$.
    Assume $\img{f}{X}\subseteq I_0$.
    Assume $x\in X$.
    Assume $V$ is open in $\reals$ with respect to $\standardTopologyR$.
    Assume $a\in\reals$ and $b\in\reals$ and $a < b$.
    Assume $f(x)\in\intervalopen{a}{b}$.
    Assume $\intervalopen{a}{b}\subseteq V$.
    Assume $f(x) = 0$.
    Then there exists $U$ such that
        $U$ is open in $X$ with respect to $T$ and $x\in U$ and $\img{f}{U}\subseteq V$.
\end{lemma}
\begin{proof}
    Let $Q_x = \urysohnLevelSet{U}{x}$.
    $f(x) < b$ by \cref{intervalopen}.
    Hence $0 < b$ by \cref{real_lt_subst_left}.
    $0\in\reals$ by \cref{real_add_ident}.
    Take $q$ such that $q\in\unitRationals$ and $q\in\intervalopen{0}{b}$
        by \cref{unit_rationals_dense_nonneg}.
    Then $q\in\reals$ and $0 < q$ and $q < b$ by \cref{intervalopen}.
    Therefore $q\in\unitRationalsZero$ by \cref{unit_rationals_zero,union_intro_left}.
    Let $U = U(q)$.
    $U$ is open in $X$ with respect to $T$
        by \cref{urysohn_union_rational_open}.
    Then $f(x)$ is an infimum of $Q_x$ by \cref{urysohn_infimum_value}.
    Hence $x\in U(q)$ by \cref{urysohn_infimum_lt_rational_in_union,real_lt_subst_right}.
    Then $x\in U$.
    Show that for all $y\in U$ we have $f(y)\in V$.
    \begin{subproof}
            Fix $y\in U$.
            Let $Q_y = \urysohnLevelSet{U}{y}$.
            Then $y\in U(q)$.
            Then $q\in Q_y$ by \cref{urysohn_level_set}.
            Then $y\in X$ and $f(y)$ is an infimum of $Q_y$ by
                \cref{open_in,topology_on,subseteq,powerset_elim,urysohn_infimum_value}.
            Hence $f(y) \leq q$ by \cref{real_infimum,real_lower_bound}.
            Then $f(y)\in\reals$ and $0 \leq f(y)$
                by \cref{funs_is_function,funs_elim,function_apply_elim,subseteq,img_iff,intervalclosed}.
            $a < f(x)$ by \cref{intervalopen}.
            Hence $a < 0$ by \cref{real_lt_subst_right}.
            Hence $a < f(y)$ by \cref{real_lt_subst_right,real_leq_split,real_lt_trans}.
            Then $f(y) < b$ by \cref{real_lt_trans}.
            Hence $f(y)\in\intervalopen{a}{b}$ by \cref{intervalopen}.
            Therefore $f(y)\in V$ by \cref{subseteq}.
    \end{subproof}
    Hence $\img{f}{U}\subseteq V$
        by \cref{img_iff,funs_is_function,function_apply_intro,subseteq}.
    Therefore there exists $U$ such that
        $U$ is open in $X$ with respect to $T$ and $x\in U$ and $\img{f}{U}\subseteq V$.
\end{proof}

% from §33 helper lemma 20: positive case produces open neighborhood
\begin{lemma}\label{urysohn_unit_interval_continuous_pos_open}
    Assume $X$ is normal with respect to $T$.
    Assume $A$ is closed in $X$ with respect to $T$.
    Assume $B$ is closed in $X$ with respect to $T$.
    Assume $A$ is disjoint from $B$.
    Assume $s$ is a sequence in $\unitRationals$.
    Assume for all $p$ such that $p\in\unitRationals$ there exists $n\in\naturalsPlus$ such that $s(n) = p$.
    Let $G = \urysohnGoodPairs{s}{X}{T}{A}{B}$.
    Let $R = \{(z,w) \mid z\in G, w\in G \mid
        \fst{w} = \fst{z} + 1 \land \snd{z}\subseteq\snd{w}\}$.
    Assume $c$ is a function on $G$.
    Assume for all $z\in G$ we have $z\mathrel{R} \apply{c}{z}$.
    Assume $u$ is a sequence in $G$.
    Assume $u(1) = a_0$.
    Assume $a_0 = (1,u_1)$.
    Assume for all $n\in\naturalsPlus$ we have $u(n + 1) = \apply{c}{u(n)}$.
    Assume for all $n\in\naturalsPlus$ we have $\fst{u(n)} = n$.
    Let $F = \{\snd{u(n)} \mid n\in\naturalsPlus\}$.
    Let $U = \unions{F}$.
    Assume $U$ is a function.
    Let $I_0 = \intervalclosed{0}{1}$.
    Let $f = \{(x,y) \mid x\in X, y\in\reals \mid
        \text{$y$ is an infimum of $\urysohnLevelSet{U}{x}$}\}$.
    Assume $f\in\funs{X}{\reals}$.
    Assume $\img{f}{X}\subseteq I_0$.
    Assume $x\in X$.
    Assume $V$ is open in $\reals$ with respect to $\standardTopologyR$.
    Assume $a\in\reals$ and $b\in\reals$ and $a < b$.
    Assume $f(x)\in\intervalopen{a}{b}$.
    Assume $\intervalopen{a}{b}\subseteq V$.
    Assume $0 < f(x)$.
    Then there exist $p,q,r_0,W$ such that
        $p\in\unitRationals$ and
        $p\in\intervalopen{a}{f(x)}$ and
        $p\in\intervalopen{0}{f(x)}$ and
        $q\in\unitRationals$ and
        $q\in\intervalopen{f(x)}{b}$ and
        $r_0\in\unitRationals$ and
        $r_0\in\intervalopen{p}{f(x)}$ and
        $p\in\unitRationalsZero$ and
        $q\in\unitRationalsZero$ and
        $r_0\in\unitRationalsZero$ and
        $W = U(q)\setminus \closure{U(p)}{X}{T}$ and
        $W$ is open in $X$ with respect to $T$.
\end{lemma}
\begin{proof}
    Show that there exists $p$ such that
        $p\in\unitRationals$ and
        $p\in\intervalopen{a}{f(x)}$ and
        $p\in\intervalopen{0}{f(x)}$.
    \begin{subproof}
        $0\in\reals$ by \cref{real_add_ident}.
        Then $f(x)\in\reals$ by \cref{funs_type_apply}.
        Hence $a < 0$ or $a = 0$ or $0 < a$ by \cref{real_lt_trichotomy}.
        \begin{byCase}
        \caseOf{$a < 0$.}
            Take $p_0$ such that $p_0\in\unitRationals$ and $p_0\in\intervalopen{0}{f(x)}$
                by \cref{unit_rationals_dense_nonneg}.
            Then $p_0\in\reals$ and $0 < p_0$ and $p_0 < f(x)$ by \cref{intervalopen}.
            Hence $p_0\in\intervalopen{a}{f(x)}$ by \cref{intervalopen,real_lt_trans}.
            Therefore there exists $p$ such that
                $p\in\unitRationals$ and
                $p\in\intervalopen{a}{f(x)}$ and
                $p\in\intervalopen{0}{f(x)}$.
        \caseOf{$a = 0$.}
            Take $p_0$ such that $p_0\in\unitRationals$ and $p_0\in\intervalopen{0}{f(x)}$
                by \cref{unit_rationals_dense_nonneg}.
            Then $p_0\in\reals$ and $0 < p_0$ and $p_0 < f(x)$ by \cref{intervalopen}.
            Hence $p_0\in\intervalopen{a}{f(x)}$ by \cref{intervalopen,real_lt_subst_left}.
            Therefore there exists $p$ such that
                $p\in\unitRationals$ and
                $p\in\intervalopen{a}{f(x)}$ and
                $p\in\intervalopen{0}{f(x)}$.
        \caseOf{$0 < a$.}
            Then $a < f(x)$ by \cref{intervalopen}.
            Take $p_0$ such that $p_0\in\unitRationals$ and $p_0\in\intervalopen{a}{f(x)}$
                by \cref{unit_rationals_dense_nonneg}.
            Then $p_0\in\reals$ and $a < p_0$ and $p_0 < f(x)$ by \cref{intervalopen}.
            Hence $p_0\in\intervalopen{0}{f(x)}$ by \cref{intervalopen,real_lt_trans}.
            Therefore there exists $p$ such that
                $p\in\unitRationals$ and
                $p\in\intervalopen{a}{f(x)}$ and
                $p\in\intervalopen{0}{f(x)}$.
        \end{byCase}
    \end{subproof}
    Take $p$ such that
        $p\in\unitRationals$ and
        $p\in\intervalopen{a}{f(x)}$ and
        $p\in\intervalopen{0}{f(x)}$.
    Then $f(x)\in\reals$ and $f(x) < b$ by \cref{funs_type_apply,intervalopen}.
    Take $q$ such that $q\in\unitRationals$ and $q\in\intervalopen{f(x)}{b}$
        by \cref{unit_rationals_dense_nonneg}.
    Then $p\in\reals$ and $a < p$ and $0 < p$ and $p < f(x)$ by \cref{intervalopen}.
    Take $r_0$ such that $r_0\in\unitRationals$ and $r_0\in\intervalopen{p}{f(x)}$
        by \cref{unit_rationals_dense_nonneg}.
    Then $p < r_0$ and $r_0 < f(x)$ by \cref{intervalopen}.
    Therefore $p\in\unitRationalsZero$ and $q\in\unitRationalsZero$ and $r_0\in\unitRationalsZero$
        by \cref{unit_rationals_zero,union_intro_left}.
    Let $W = U(q)\setminus \closure{U(p)}{X}{T}$.
    Hence $U(q)\subseteq X$ and $U(q)\in T$
        by \cref{urysohn_union_rational_open,normal_space,open_in,topology_on,subseteq,powerset_elim}.
    $U(p)\subseteq X$ by \cref{urysohn_union_zero_subset_X_aux}.
    $T$ is a topology on $X$ by \cref{normal_space}.
    Hence $\closure{U(p)}{X}{T}$ is closed in $X$ with respect to $T$
        by \cref{closure_closed_in}.
    Then $X\setminus \closure{U(p)}{X}{T}\in T$ by \cref{closed_in,open_in}.
    Then $U(q)\inter (X\setminus \closure{U(p)}{X}{T})\in T$ by \cref{topology_on}.
    Hence $\closure{U(p)}{X}{T}\subseteq X$ by \cref{closed_in}.
    Then $W = U(q)\inter (X\setminus \closure{U(p)}{X}{T})$
        by \cref{setminus_eq_inter_complement}.
    Hence $W\in T$.
    Thus $W$ is open in $X$ with respect to $T$ by \cref{open_in}.
    Therefore there exist $p,q,r_0,W$ such that
        $p\in\unitRationals$ and
        $p\in\intervalopen{a}{f(x)}$ and
        $p\in\intervalopen{0}{f(x)}$ and
        $q\in\unitRationals$ and
        $q\in\intervalopen{f(x)}{b}$ and
        $r_0\in\unitRationals$ and
        $r_0\in\intervalopen{p}{f(x)}$ and
        $p\in\unitRationalsZero$ and
        $q\in\unitRationalsZero$ and
        $r_0\in\unitRationalsZero$ and
        $W = U(q)\setminus \closure{U(p)}{X}{T}$ and
        $W$ is open in $X$ with respect to $T$.
\end{proof}

% from §33 helper lemma 20e: rational U-values lie in X
\begin{lemma}\label{urysohn_union_rational_subset_X}
    Assume $s$ is a sequence in $\unitRationals$.
    Assume for all $p$ such that $p\in\unitRationals$ there exists $n\in\naturalsPlus$ such that $s(n) = p$.
    Let $G = \urysohnGoodPairs{s}{X}{T}{A}{B}$.
    Assume $u$ is a sequence in $G$.
    Assume for all $n\in\naturalsPlus$ we have $\fst{u(n)} = n$.
    Let $F = \{\snd{u(n)} \mid n\in\naturalsPlus\}$.
    Let $U = \unions{F}$.
    Assume $U$ is a function.
    Assume $q\in\unitRationals$.
    Then $U(q)\subseteq X$.
\end{lemma}
\begin{proof}
    Follows by \cref{unit_rationals_zero,union_intro_left,urysohn_union_zero_subset_X_aux}.
\end{proof}
% from §33 helper lemma 21: positive case point membership
\begin{lemma}\label{urysohn_unit_interval_continuous_pos_x_in_W}
    Assume $X$ is normal with respect to $T$.
    Assume $A$ is closed in $X$ with respect to $T$.
    Assume $B$ is closed in $X$ with respect to $T$.
    Assume $A$ is disjoint from $B$.
    Assume $s$ is a sequence in $\unitRationals$.
    Assume for all $p$ such that $p\in\unitRationals$ there exists $n\in\naturalsPlus$ such that $s(n) = p$.
    Let $G = \urysohnGoodPairs{s}{X}{T}{A}{B}$.
    Let $R = \{(z,w) \mid z\in G, w\in G \mid
        \fst{w} = \fst{z} + 1 \land \snd{z}\subseteq\snd{w}\}$.
    Assume $c$ is a function on $G$.
    Assume for all $z\in G$ we have $z\mathrel{R} \apply{c}{z}$.
    Assume $u$ is a sequence in $G$.
    Assume $u(1) = a_0$.
    Assume $a_0 = (1,u_1)$.
    Assume for all $n\in\naturalsPlus$ we have $u(n + 1) = \apply{c}{u(n)}$.
    Assume for all $n\in\naturalsPlus$ we have $\fst{u(n)} = n$.
    Let $F = \{\snd{u(n)} \mid n\in\naturalsPlus\}$.
    Let $U = \unions{F}$.
    Assume $U$ is a function.
    Let $I_0 = \intervalclosed{0}{1}$.
    Let $f = \{(x,y) \mid x\in X, y\in\reals \mid
        \text{$y$ is an infimum of $\urysohnLevelSet{U}{x}$}\}$.
    Assume $f\in\funs{X}{\reals}$ and $\img{f}{X}\subseteq I_0$.
    Assume $x\in X$ and $V$ is open in $\reals$ with respect to $\standardTopologyR$ and
        $a\in\reals$ and $b\in\reals$ and $a < b$ and
        $f(x)\in\intervalopen{a}{b}$ and $\intervalopen{a}{b}\subseteq V$ and $0 < f(x)$.
    Assume $p\in\unitRationals$ and $p\in\intervalopen{a}{f(x)}$ and $p\in\intervalopen{0}{f(x)}$ and
        $q\in\unitRationals$ and $q\in\intervalopen{f(x)}{b}$ and
        $r_0\in\unitRationals$ and $r_0\in\intervalopen{p}{f(x)}$ and
        $p\in\unitRationalsZero$ and $q\in\unitRationalsZero$ and $r_0\in\unitRationalsZero$.
    Let $W = U(q)\setminus \closure{U(p)}{X}{T}$.
    Then $x\in W$.
\end{lemma}
\begin{proof}
    Let $Q_x = \urysohnLevelSet{U}{x}$.
    Then $f(x)$ is an infimum of $Q_x$ by \cref{urysohn_infimum_value}.
    Hence $f(x)$ is a lower bound of $Q_x$ by \cref{real_infimum}.
    Then $f(x) < q$ by \cref{intervalopen}.
    Hence $x\in U(q)$ by \cref{urysohn_infimum_lt_rational_in_union,intervalopen}.
    Then $p < r_0$ by \cref{intervalopen}.
    Hence $\closure{U(p)}{X}{T}\subseteq U(r_0)$ by \cref{urysohn_union_closure_subset}.
    Show that $x\notin U(r_0)$.
    \begin{subproof}
        Suppose not.
        Then $x\in U(r_0)$.
        Hence $r_0\in Q_x$ by \cref{urysohn_level_set}.
        Then $f(x) \leq r_0$ by \cref{real_lower_bound}.
        Hence $f(x)\in\reals$ and $r_0\in\reals$ by \cref{funs_type_apply,intervalopen}.
        Then $f(x) < r_0$ or $f(x) = r_0$ by \cref{real_leq_split}.
        Hence $r_0 < f(x)$ by \cref{intervalopen}.
        Contradiction by \cref{real_lt_trans,real_lt_irreflexive,real_lt_subst_left}.
    \end{subproof}
    Therefore $x\notin\closure{U(p)}{X}{T}$ by \cref{not_in_subseteq}.
    Hence $x\in W$ by \cref{setminus_intro}.
\end{proof}
% from §33 helper lemma 22: positive case image control
\begin{lemma}\label{urysohn_unit_interval_continuous_pos_img}
    Assume $X$ is normal with respect to $T$ and $A$ is closed in $X$ with respect to $T$ and
        $B$ is closed in $X$ with respect to $T$ and $A$ is disjoint from $B$.
    Assume $s$ is a sequence in $\unitRationals$ and for all $p$ such that $p\in\unitRationals$
        there exists $n\in\naturalsPlus$ such that $s(n) = p$.
    Let $G = \urysohnGoodPairs{s}{X}{T}{A}{B}$.
    Let $R = \{(z,w) \mid z\in G, w\in G \mid
        \fst{w} = \fst{z} + 1 \land \snd{z}\subseteq\snd{w}\}$.
    Assume $c$ is a function on $G$ and for all $z\in G$ we have $z\mathrel{R} \apply{c}{z}$.
    Assume $u$ is a sequence in $G$ and $u(1) = a_0$ and $a_0 = (1,u_1)$ and
        for all $n\in\naturalsPlus$ we have $u(n + 1) = \apply{c}{u(n)}$ and $\fst{u(n)} = n$.
    Let $F = \{\snd{u(n)} \mid n\in\naturalsPlus\}$.
    Let $U = \unions{F}$.
    Assume $U$ is a function.
    Let $I_0 = \intervalclosed{0}{1}$.
    Let $f = \{(x,y) \mid x\in X, y\in\reals \mid
        \text{$y$ is an infimum of $\urysohnLevelSet{U}{x}$}\}$.
    Assume $f\in\funs{X}{\reals}$ and $\img{f}{X}\subseteq I_0$.
    Assume $x\in X$ and $V$ is open in $\reals$ with respect to $\standardTopologyR$ and
        $a\in\reals$ and $b\in\reals$ and $a < b$ and
        $f(x)\in\intervalopen{a}{b}$ and $\intervalopen{a}{b}\subseteq V$ and $0 < f(x)$.
    Assume $p\in\unitRationals$ and $p\in\intervalopen{a}{f(x)}$ and $p\in\intervalopen{0}{f(x)}$ and
        $q\in\unitRationals$ and $q\in\intervalopen{f(x)}{b}$ and
        $r_0\in\unitRationals$ and $r_0\in\intervalopen{p}{f(x)}$ and
        $p\in\unitRationalsZero$ and $q\in\unitRationalsZero$ and $r_0\in\unitRationalsZero$.
    Let $W = U(q)\setminus \closure{U(p)}{X}{T}$.
    Assume $W\subseteq X$.
    Then $\img{f}{W}\subseteq V$.
\end{lemma}
\begin{proof}
    Show that for all $y\in W$ we have $f(y)\in V$.
    \begin{subproof}
        Fix $y$ such that $y\in W$.
        Then $y\in X$ and $f(y)\in\reals$ by \cref{subseteq,funs_type_apply}.
        Let $Q_y = \urysohnLevelSet{U}{y}$.
        Then $y\notin\closure{U(p)}{X}{T}$ by \cref{setminus_elim_right}.
        Hence $q\in Q_y$ by \cref{setminus_elim_left,urysohn_level_set}.
        Therefore $f(y) \leq q$ by \cref{urysohn_infimum_value,real_infimum,real_lower_bound}.
        Show that $p \leq f(y)$.
        \begin{subproof}
            Show that $p$ is a lower bound of $Q_y$.
            \begin{subproof}
                Hence $p\in\reals$ and $Q_y\subseteq\reals$ by \cref{intervalopen,urysohn_level_value_real,subseteq}.
                Show that for all $r\in Q_y$ we have $p \leq r$.
                \begin{subproof}
                    Fix $r$ such that $r\in Q_y$.
                    Then $r\in\reals$ by \cref{subseteq}.
                    Hence $y\in U(r)$ and $r\in\unitRationalsZero$ by \cref{urysohn_level_set}.
                    Show that $p \neq r$.
                    \begin{subproof}
                        Suppose not.
                        Then $p = r$.
                        Hence $y\in U(p)$.
                        Therefore $y\in\closure{U(p)}{X}{T}$
                            by \cref{urysohn_union_zero_subset_X_aux,subseteq_closure,normal_space,subseteq}.
                        Contradiction.
                    \end{subproof}
                    Hence $p < r$ or $r < p$ by \cref{real_lt_trichotomy}.
                    \begin{byCase}
                    \caseOf{$p < r$.}
                        Therefore $p \leq r$ by \cref{real_lt_imp_leq}.
                    \caseOf{$r < p$.}
                        Show that $p \leq r$.
                        \begin{subproof}
                            Suppose not.
                            Then $\closure{U(r)}{X}{T}\subseteq U(p)$ by \cref{urysohn_union_closure_subset}.
                            Hence $y\in\closure{U(r)}{X}{T}$
                                by \cref{urysohn_union_zero_subset_X_aux,subseteq_closure,normal_space,subseteq}.
                            Hence $y\in U(p)$ by \cref{subseteq}.
                            Therefore $y\in\closure{U(p)}{X}{T}$
                                by \cref{urysohn_union_zero_subset_X_aux,subseteq_closure,normal_space,subseteq}.
                            Contradiction.
                        \end{subproof}
                    \end{byCase}
                \end{subproof}
                Therefore $p$ is a lower bound of $Q_y$ by \cref{real_lower_bound}.
            \end{subproof}
            Hence $p \leq f(y)$ by \cref{urysohn_infimum_value,real_infimum}.
        \end{subproof}
        Then $a < f(y)$ by \cref{intervalopen,real_leq_split,real_lt_subst_right,real_lt_trans}.
        Hence $f(y) < b$ by \cref{intervalopen,real_leq_split,real_lt_subst_left,real_lt_trans}.
        Therefore $f(y)\in V$ by \cref{intervalopen,subseteq}.
    \end{subproof}
    Hence $\img{f}{W}\subseteq V$ by \cref{img_iff,funs_is_function,function_apply_intro,subseteq}.
\end{proof}

% from §33 helper lemma 23: continuity at a point
\begin{lemma}\label{urysohn_unit_interval_continuous_at_point}
    Assume $X$ is normal with respect to $T$.
    Assume $A$ is closed in $X$ with respect to $T$.
    Assume $B$ is closed in $X$ with respect to $T$.
    Assume $A$ is disjoint from $B$.
    Assume $s$ is a sequence in $\unitRationals$.
    Assume for all $p$ such that $p\in\unitRationals$ there exists $n\in\naturalsPlus$ such that $s(n) = p$.
    Let $G = \urysohnGoodPairs{s}{X}{T}{A}{B}$.
    Let $R = \{(z,w) \mid z\in G, w\in G \mid
        \fst{w} = \fst{z} + 1 \land \snd{z}\subseteq\snd{w}\}$.
    Assume $c$ is a function on $G$.
    Assume for all $z\in G$ we have $z\mathrel{R} \apply{c}{z}$.
    Assume $u$ is a sequence in $G$.
    Assume $u(1) = a_0$.
    Assume $a_0 = (1,u_1)$.
    Assume for all $n\in\naturalsPlus$ we have $u(n + 1) = \apply{c}{u(n)}$.
    Assume for all $n\in\naturalsPlus$ we have $\fst{u(n)} = n$.
    Let $F = \{\snd{u(n)} \mid n\in\naturalsPlus\}$.
    Let $U = \unions{F}$.
    Assume $U$ is a function.
    Let $I_0 = \intervalclosed{0}{1}$.
    Let $f = \{(x,y) \mid x\in X, y\in\reals \mid
        \text{$y$ is an infimum of $\urysohnLevelSet{U}{x}$}\}$.
    Assume $f\in\funs{X}{\reals}$.
    Assume $\img{f}{X}\subseteq I_0$.
    Assume $x\in X$.
    Assume $V$ is open in $\reals$ with respect to $\standardTopologyR$.
    Assume $f(x)\in V$.
    Then there exists $U_0$ such that
        $U_0$ is open in $X$ with respect to $T$ and $x\in U_0$ and $\img{f}{U_0}\subseteq V$.
\end{lemma}
\begin{proof}
    $V\in\basisTopology{\standardBasisR}{\reals}$ by \cref{open_in,standard_topology_reals}.
    Take $B_1$ such that $B_1\in\standardBasisR$ and $f(x)\in B_1$ and $B_1\subseteq V$
        by \cref{topology_generated_by_basis}.
    Take $a,b\in\reals$ such that $a < b$ and $B_1 = \intervalopen{a}{b}$
        by \cref{standard_basis_reals}.
    Then $a < f(x)$ and $f(x) < b$ by \cref{intervalopen}.
    Hence $0 < f(x)$ or $0 = f(x)$
        by \cref{funs_elim,funs_is_function,subseteq,function_apply_elim,img_iff,funs_type_apply,intervalclosed,real_add_ident,real_leq_split}.
    \begin{byCase}
    \caseOf{$0 = f(x)$.}
        Therefore there exists $U_0$ such that
            $U_0$ is open in $X$ with respect to $T$ and $x\in U_0$ and $\img{f}{U_0}\subseteq V$
            by \cref{urysohn_unit_interval_continuous_zero}.
    \caseOf{$0 < f(x)$.}
        Take $p,q,r_0,W$ such that
            $p\in\unitRationals$ and
            $p\in\intervalopen{a}{f(x)}$ and
            $p\in\intervalopen{0}{f(x)}$ and
            $q\in\unitRationals$ and
            $q\in\intervalopen{f(x)}{b}$ and
            $r_0\in\unitRationals$ and
            $r_0\in\intervalopen{p}{f(x)}$ and
            $p\in\unitRationalsZero$ and
            $q\in\unitRationalsZero$ and
            $r_0\in\unitRationalsZero$ and
            $W = U(q)\setminus \closure{U(p)}{X}{T}$ and
            $W$ is open in $X$ with respect to $T$
            by \cref{urysohn_unit_interval_continuous_pos_open}.
        Then $x\in W$ by \cref{urysohn_unit_interval_continuous_pos_x_in_W}.
        Hence $\img{f}{W}\subseteq V$
            by \cref{urysohn_unit_interval_continuous_pos_img,open_in,normal_space,topology_on,subseteq,powerset_elim}.
        Therefore there exists $U_0$ such that
            $U_0$ is open in $X$ with respect to $T$ and $x\in U_0$ and $\img{f}{U_0}\subseteq V$.
    \end{byCase}
\end{proof}

% from §33 helper lemma 24: Urysohn function is continuous to reals
\begin{lemma}\label{urysohn_unit_interval_continuous_reals}
    Assume $X$ is normal with respect to $T$.
    Assume $A$ is closed in $X$ with respect to $T$.
    Assume $B$ is closed in $X$ with respect to $T$.
    Assume $A$ is disjoint from $B$.
    Assume $s$ is a sequence in $\unitRationals$.
    Assume for all $p$ such that $p\in\unitRationals$ there exists $n\in\naturalsPlus$ such that $s(n) = p$.
    Let $G = \urysohnGoodPairs{s}{X}{T}{A}{B}$.
    Let $R = \{(z,w) \mid z\in G, w\in G \mid
        \fst{w} = \fst{z} + 1 \land \snd{z}\subseteq\snd{w}\}$.
    Assume $c$ is a function on $G$.
    Assume for all $z\in G$ we have $z\mathrel{R} \apply{c}{z}$.
    Assume $u$ is a sequence in $G$.
    Assume $u(1) = a_0$.
    Assume $a_0 = (1,u_1)$.
    Assume for all $n\in\naturalsPlus$ we have $u(n + 1) = \apply{c}{u(n)}$.
    Assume for all $n\in\naturalsPlus$ we have $\fst{u(n)} = n$.
    Let $F = \{\snd{u(n)} \mid n\in\naturalsPlus\}$.
    Let $U = \unions{F}$.
    Assume $U$ is a function.
    Let $I_0 = \intervalclosed{0}{1}$.
    Let $f = \{(x,y) \mid x\in X, y\in\reals \mid
        \text{$y$ is an infimum of $\urysohnLevelSet{U}{x}$}\}$.
    Assume $f\in\funs{X}{\reals}$.
    Assume $\img{f}{X}\subseteq I_0$.
    Then $f$ is continuous from $X$ to $\reals$ with respect to $T$ and $\standardTopologyR$.
\end{lemma}
\begin{proof}
    For all $x$ such that $x\in X$ we have
        $f$ is continuous at $x$ from $X$ to $\reals$ with respect to $T$ and $\standardTopologyR$
        by \cref{urysohn_unit_interval_continuous_at_point,standard_basis_is_basis,basis_topology_is_topology,standard_topology_reals,normal_space,continuous_at_open}.
    Hence $f$ is continuous from $X$ to $\reals$ with respect to $T$ and $\standardTopologyR$
        by \cref{standard_basis_is_basis,basis_topology_is_topology,standard_topology_reals,normal_space,continuous_at_implies_continuous}.
\end{proof}

% from §33 helper lemma 34: admissible values above $1$
\begin{lemma}\label{urysohn_admissible_above_one}
    Assume $f\in\urysohnAdmissible{s}{X}{T}{A}{B}{n}$.
    Assume $p\in\dom{f}$.
    Assume $1 < p$.
    Then $f(p) = X$.
\end{lemma}
\begin{proof}
    Then $1\in\dom{f}$ by \cref{urysohn_admissible,funs_elim,upair_intro_right,union_intro_left}.
    Then for all $q\in\dom{f}$ if $1 < q$ then $f(q) = X$ by \cref{urysohn_admissible}.
    Therefore $f(p) = X$.
\end{proof}

% from §33 helper lemma 35: unit rationals are real
\begin{lemma}\label{unit_rational_real}
    Assume $r\in\unitRationals$.
    Then $r\in\reals$.
\end{lemma}
\begin{proof}
    Follows by \cref{unit_rational_real_aux}.
\end{proof}

% from §33 helper lemma 25: extension when $r\in\dom{f}$
\begin{lemma}\label{urysohn_unit_interval_extend_in_dom}
    Assume $X$ is normal with respect to $T$.
    Assume $A$ is closed in $X$ with respect to $T$.
    Assume $B$ is closed in $X$ with respect to $T$.
    Assume $A$ is disjoint from $B$.
    Assume $s$ is a sequence in $\unitRationals$.
    Let $G = \urysohnGoodPairs{s}{X}{T}{A}{B}$.
    Let $R = \{(z,w) \mid z\in G, w\in G \mid
        \fst{w} = \fst{z} + 1 \land \snd{z}\subseteq\snd{w}\}$.
    Suppose $n\in\naturalsPlus$ and $f\in\urysohnAdmissible{s}{X}{T}{A}{B}{n}$ and $z = (n,f)$.
    Let $r = s(n + 1)$.
    Assume $r\in\dom{f}$.
    Then there exists $w$ such that $z\mathrel{R} w$.
\end{lemma}
\begin{proof}
    $f\in\funs{\{0,1\}\union \img{s}{\initialSegment{n}}}{\pow{X}}$ and
        for all $p\in\dom{f}$ we have $f(p)\in T$ and
        $f(1) = X\setminus B$ and
        $A\subseteq f(0)$ and
        $\closure{f(0)}{X}{T}\subseteq f(1)$ and
        for all $p\in\dom{f}$ for all $q\in\dom{f}$ if $p < q$ then
            $\closure{f(p)}{X}{T}\subseteq f(q)$ and
        for all $p\in\dom{f}$ if $1 < p$ then $f(p) = X$
        by \cref{urysohn_admissible}.
    Show that $\dom{f} = \{0,1\}\union \img{s}{\initialSegment{n + 1}}$.
    \begin{subproof}
        Then $\{0,1\}\union \img{s}{\initialSegment{n}}\subseteq \{0,1\}\union \img{s}{\initialSegment{n + 1}}$
            by \cref{initial_segment_subset_succ,img_subseteq,union_intro_left,union_intro_right,subseteq,union_subsets_is_subset}.
        Show that for all $y\in \img{s}{\initialSegment{n + 1}}$ we have $y\in\dom{f}$.
        \begin{subproof}
            Fix $y$ such that $y\in \img{s}{\initialSegment{n + 1}}$.
            Take $k$ such that $k\in\initialSegment{n + 1}$ and $(k,y)\in s$ by \cref{img_iff}.
            Then $s(k) = y$ by \cref{sequence_in,funs_is_function,function_apply_intro}.
            \begin{byCase}
            \caseOf{$k\in\initialSegment{n}$.}
                Hence $y\in\dom{f}$ by \cref{img_iff,union_intro_right,funs_elim}.
            \caseOf{$k\notin\initialSegment{n}$.}
                $k\in\naturalsPlus$ and $k\leq n + 1$ by \cref{initial_segment_naturalsplus}.
                Then $k = n + 1$.
                Therefore $y = r$ and $y\in\dom{f}$ by \cref{subseteq}.
            \end{byCase}
        \end{subproof}
        Hence $\img{s}{\initialSegment{n + 1}}\subseteq \dom{f}$ by \cref{subseteq}.
        Hence $\{0,1\}\union \img{s}{\initialSegment{n + 1}}\subseteq \dom{f}$
            by \cref{funs_elim,union_intro_left,subseteq,union_subsets_is_subset}.
        For all $y$ such that $y\in\dom{f}$ we have $y\in \{0,1\}\union \img{s}{\initialSegment{n + 1}}$
            by \cref{funs_elim,subseteq}.
        Then $\dom{f}\subseteq \{0,1\}\union \img{s}{\initialSegment{n + 1}}$ by \cref{subseteq}.
        Therefore $\dom{f} = \{0,1\}\union \img{s}{\initialSegment{n + 1}}$ by \cref{subseteq_antisymmetric}.
    \end{subproof}
    Let $w = (n + 1,f)$.
    Then $n + 1\in\naturalsPlus$ by \cref{naturalsplus_add_closed,real_one_in_naturals}.
    Then $f\in\urysohnAdmissible{s}{X}{T}{A}{B}{n + 1}$ by \cref{funs_elim,funs_intro,urysohn_admissible}.
    Hence $z\in G$ and $w\in G$ and $\fst{w} = \fst{z} + 1$ and $\snd{z}\subseteq \snd{w}$
        by \cref{urysohn_good_pair_intro,fst_eq,snd_eq,subseteq_refl}.
    Therefore $z\mathrel{R} w$.
\end{proof}

% from §33 helper lemma 26: extension for $1 < r$ (setup)
\begin{lemma}\label{urysohn_unit_interval_extend_gt1_setup}
    Assume $X$ is normal with respect to $T$.
    Assume $A$ is closed in $X$ with respect to $T$.
    Assume $B$ is closed in $X$ with respect to $T$.
    Assume $A$ is disjoint from $B$.
    Assume $s$ is a sequence in $\unitRationals$.
    Suppose $n\in\naturalsPlus$ and $f\in\urysohnAdmissible{s}{X}{T}{A}{B}{n}$ and $z = (n,f)$.
    Let $r = s(n + 1)$.
    Assume $r\notin\dom{f}$.
    Let $U_r = X$.
    Let $h = \{(r,U_r)\}$.
    Let $g = f\union h$.
    Then $g$ is a function.
\end{lemma}
\begin{proof}
    For all $x$ such that $x\in\dom{f}\inter\dom{h}$ we have $f(x) = h(x)$.
    For all $a,b,b'$ such that $(a,b)\in\{(r,U_r)\}$ and $(a,b')\in\{(r,U_r)\}$ we have
        $b = b'$ by \cref{singleton_elim,pair_eq_iff}.
    For all $w\in\{(r,U_r)\}$ there exist $x, y$ such that $w = (x,y)$ by \cref{singleton_elim}.
    Therefore $h$ is a relation and $h$ is right-unique by \cref{relation,rightunique}.
    Hence $g$ is a function by \cref{urysohn_admissible,funs_is_function,union_compatible_functions}.
\end{proof}

% from §33 helper lemma 27: extension for $1 < r$
\begin{lemma}\label{urysohn_unit_interval_extend_gt1}
    Assume $X$ is normal with respect to $T$.
    Assume $A$ is closed in $X$ with respect to $T$.
    Assume $B$ is closed in $X$ with respect to $T$.
    Assume $A$ is disjoint from $B$.
    Assume $s$ is a sequence in $\unitRationals$.
    Let $G = \urysohnGoodPairs{s}{X}{T}{A}{B}$.
    Let $R = \{(z,w) \mid z\in G, w\in G \mid
        \fst{w} = \fst{z} + 1 \land \snd{z}\subseteq\snd{w}\}$.
    Suppose $n\in\naturalsPlus$ and $f\in\urysohnAdmissible{s}{X}{T}{A}{B}{n}$ and $z = (n,f)$.
    Let $r = s(n + 1)$.
    Assume $r\notin\dom{f}$.
    Assume $1 < r$.
    Then there exists $w$ such that $z\mathrel{R} w$.
\end{lemma}
\begin{proof}
    Then $f\in\funs{\{0,1\}\union \img{s}{\initialSegment{n}}}{\pow{X}}$ and
        for all $p\in\dom{f}$ we have $f(p)\in T$ and
        $f(1) = X\setminus B$ and
        $A\subseteq f(0)$ and
        $\closure{f(0)}{X}{T}\subseteq f(1)$ and
        for all $p\in\dom{f}$ for all $q\in\dom{f}$ if $p < q$ then
            $\closure{f(p)}{X}{T}\subseteq f(q)$ and
        for all $p\in\dom{f}$ if $1 < p$ then $f(p) = X$
        by \cref{urysohn_admissible}.
    Let $U_r = X$.
    Let $h = \{(r,U_r)\}$.
    Let $g = f\union h$.
    Hence $g$ is a function by \cref{urysohn_unit_interval_extend_gt1_setup}.
    $\dom{g} = \dom{f}\union\dom{h}$ by \cref{dom_union}.
    For all $x\in\dom{h}$ we have $x\in\{r\}$ by \cref{dom_iff,singleton_elim,pair_eq_iff,singleton_intro}.
    For all $x\in\{r\}$ we have $x\in\dom{h}$
        by \cref{singleton_elim,singleton_intro,pair_eq_iff,dom_intro}.
    Hence $\dom{h} = \{r\}$ by \cref{subseteq,subseteq_antisymmetric}.
    Therefore $\dom{g} = \dom{f}\union\{r\}$.
    $\dom{f} = \{0,1\}\union \img{s}{\initialSegment{n}}$ by \cref{funs_elim}.
    Show that for all $y\in\dom{g}$ we have $y\in \{0,1\}\union \img{s}{\initialSegment{n + 1}}$.
        \begin{subproof}
            Fix $y$ such that $y\in\dom{g}$.
            Then $y\in\dom{f}$ or $y\in\{r\}$ by \cref{union_iff}.
            \begin{byCase}
            \caseOf{$y\in\dom{f}$.}
                Then $y\in \{0,1\}$ or $y\in \img{s}{\initialSegment{n}}$ by \cref{funs_elim,union_iff}.
                \begin{byCase}
                \caseOf{$y\in \{0,1\}$.}
                    Thus $y\in \{0,1\}\union \img{s}{\initialSegment{n + 1}}$ by \cref{union_intro_left}.
                \caseOf{$y\in \img{s}{\initialSegment{n}}$.}
                    Take $k$ such that $k\in\initialSegment{n}$ and $(k,y)\in s$ by \cref{img_iff}.
                    Thus $y\in \{0,1\}\union \img{s}{\initialSegment{n + 1}}$
                        by \cref{initial_segment_subset_succ,subseteq,img_iff,union_intro_right}.
                \end{byCase}
            \caseOf{$y\in\{r\}$.}
                Then $y\in \{0,1\}\union \img{s}{\initialSegment{n + 1}}$
                    by \cref{singleton_elim,naturalsplus_add_closed,real_one_in_naturals,sequence_in,sequence_value_in_initial_image,union_intro_right}.
            \end{byCase}
    \end{subproof}
    Hence $\dom{g}\subseteq \{0,1\}\union \img{s}{\initialSegment{n + 1}}$ by \cref{subseteq}.
    Show that for all $y\in \{0,1\}\union \img{s}{\initialSegment{n + 1}}$ we have $y\in\dom{g}$.
    \begin{subproof}
            Fix $y$ such that $y\in \{0,1\}\union \img{s}{\initialSegment{n + 1}}$.
            Then $y\in \{0,1\}$ or $y\in \img{s}{\initialSegment{n + 1}}$ by \cref{union_iff}.
            \begin{byCase}
            \caseOf{$y\in \{0,1\}$.}
                Then $y\in\dom{g}$ by \cref{union_intro_left,funs_elim}.
            \caseOf{$y\in \img{s}{\initialSegment{n + 1}}$.}
                Take $k$ such that $k\in\initialSegment{n + 1}$ and $(k,y)\in s$ by \cref{img_iff}.
                \begin{byCase}
                \caseOf{$k\in\initialSegment{n}$.}
                    Thus $y\in\dom{g}$ by \cref{img_iff,union_intro_right,funs_elim,union_intro_left}.
                \caseOf{$k\notin\initialSegment{n}$.}
                    $k\in\naturalsPlus$ and $k\leq n + 1$ by \cref{initial_segment_naturalsplus}.
                    Hence $k = n + 1$.
                    Thus $y = r$.
                    Thus $y\in\dom{g}$ by \cref{singleton_intro,union_intro_right}.
                \end{byCase}
            \end{byCase}
    \end{subproof}
    Thus $\{0,1\}\union \img{s}{\initialSegment{n + 1}}\subseteq \dom{g}$ by \cref{subseteq}.
    Therefore $\dom{g} = \{0,1\}\union \img{s}{\initialSegment{n + 1}}$ by \cref{subseteq_antisymmetric}.
    Show that for all $p\in\dom{g}$ we have $g(p)\in\pow{X}$.
    \begin{subproof}
        Fix $p$ such that $p\in\dom{g}$.
        Then $p\in\dom{f}$ or $p\in\{r\}$ by \cref{union_iff}.
        \begin{byCase}
        \caseOf{$p\in\dom{f}$.}
            Thus $g(p)\in\pow{X}$
                by \cref{urysohn_admissible,funs_type_apply,funs_is_function,function_apply_elim,union_intro_left,function_apply_intro}.
        \caseOf{$p\in\{r\}$.}
            Thus $g(p)\in\pow{X}$
                by \cref{singleton_elim,singleton_intro,union_intro_right,function_apply_intro,normal_space,topology_on,subseteq}.
        \end{byCase}
    \end{subproof}
    Let $w = (n + 1,g)$.
    Then $n + 1\in\naturalsPlus$ by \cref{naturalsplus_add_closed,real_one_in_naturals}.
    Therefore $g\in\funs{\{0,1\}\union \img{s}{\initialSegment{n + 1}}}{\pow{X}}$
        by \cref{urysohn_unit_interval_extend_gt1_setup,funs_intro}.
    Show that for all $p\in\dom{g}$ we have
        $g(p)\in T$ and
        $g(1) = X\setminus B$ and
        $A\subseteq g(0)$ and
        $\closure{g(0)}{X}{T}\subseteq g(1)$ and
        for all $p\in\dom{g}$ for all $q\in\dom{g}$ if $p < q$ then
            $\closure{g(p)}{X}{T}\subseteq g(q)$ and
        for all $p\in\dom{g}$ if $1 < p$ then $g(p) = X$.
    \begin{subproof}
        Fix $p$ such that $p\in\dom{g}$.
        Show that $g(p)\in T$.
        \begin{subproof}
            Then $p\in\dom{f}$ or $p\in\dom{h}$ by \cref{dom_union,union_iff}.
            \begin{byCase}
            \caseOf{$p\in\dom{f}$.}
                Thus $g(p)\in T$
                    by \cref{urysohn_admissible,funs_is_function,function_apply_elim,union_intro_left,function_apply_intro}.
            \caseOf{$p\in\dom{h}$.}
                Take $y$ such that $(p,y)\in h$ by \cref{dom_iff}.
                Thus $(p,y) = (r,U_r)$ by \cref{singleton_elim}.
                Thus $p = r$ by \cref{pair_eq_iff}.
                Thus $g(p) = U_r$ by \cref{union_intro_right,singleton_intro,function_apply_intro}.
                Thus $g(p)\in T$ by \cref{normal_space,topology_on}.
            \end{byCase}
        \end{subproof}
        Hence $T$ is a topology on $X$ and $g(p)$ is open in $X$ with respect to $T$ by \cref{normal_space,open_in}.
        Thus $1\in\dom{f}$ by \cref{urysohn_admissible,funs_elim,upair_intro_right,union_intro_left}.
        Thus $(1,f(1))\in f$ and $(1,f(1))\in g$
            by \cref{urysohn_admissible,funs_is_function,function_apply_elim,union_intro_left}.
        Therefore $g(1) = X\setminus B$ by \cref{function_apply_intro}.
        Thus $(0,f(0))\in f$ and $(0,f(0))\in g$
            by \cref{urysohn_admissible,funs_is_function,funs_elim,function_apply_elim,upair_intro_left,union_intro_left}.
        Thus $g(0) = f(0)$ by \cref{function_apply_intro}.
        Hence $A\subseteq g(0)$ and $\closure{g(0)}{X}{T}\subseteq g(1)$ by \cref{subseteq}.
        Show that for all $p\in\dom{g}$ for all $q\in\dom{g}$ if $p < q$ then
            $\closure{g(p)}{X}{T}\subseteq g(q)$.
        \begin{subproof}
            Fix $p$ such that $p\in\dom{g}$.
            Fix $q$ such that $q\in\dom{g}$.
            Assume $p < q$.
            Then $p\in\dom{f}$ or $p\in\dom{h}$ by \cref{dom_union,union_iff}.
            Then $q\in\dom{f}$ or $q\in\dom{h}$ by \cref{dom_union,union_iff}.
            \begin{byCase}
            \caseOf{$q\in\dom{h}$.}
                Take $y$ such that $(q,y)\in h$ by \cref{dom_iff}.
                Thus $(q,y) = (r,U_r)$ by \cref{singleton_elim}.
                Thus $q = r$ by \cref{pair_eq_iff}.
                Then $g(q) = X$ by \cref{union_intro_right,singleton_intro,function_apply_intro}.
                Thus $\closure{g(p)}{X}{T}\subseteq X$
                    by \cref{closed_whole_space,funs_type_apply,powerset_elim,closure_subseteq_closed}.
                Therefore $\closure{g(p)}{X}{T}\subseteq g(q)$ by \cref{subseteq}.
            \caseOf{$q\in\dom{f}$.}
                \begin{byCase}
                \caseOf{$p\in\dom{f}$.}
                    Thus $g(p) = f(p)$ and $g(q) = f(q)$
                        by \cref{funs_is_function,function_apply_elim,union_intro_left,function_apply_intro}.
                    Hence $\closure{g(p)}{X}{T}\subseteq g(q)$ by \cref{urysohn_admissible,subseteq}.
                \caseOf{$p\in\dom{h}$.}
                    Take $y$ such that $(p,y)\in h$ by \cref{dom_iff}.
                    Thus $p = r$ by \cref{singleton_elim,pair_eq_iff}.
                    Then $r < q$ by \cref{real_lt_subst_left}.
                    Thus $r\in\unitRationals$ and $r\in\reals$
                        by \cref{sequence_in,naturalsplus_add_closed,real_one_in_naturals,funs_type_apply,unit_rational_real}.
                    Thus $q\in \{0,1\}\union \img{s}{\initialSegment{n}}$ by \cref{urysohn_admissible,funs_elim}.
                    Thus $q\in\reals$ by
                        \cref{union_iff,upair_elim,real_add_ident,real_mul_ident,img_iff,sequence_in,initial_segment_naturalsplus,funs_is_function,function_apply_intro,funs_type_apply,unit_rational_real}.
                    $1\in\reals$ and $1 < q$ by \cref{real_mul_ident,real_lt_trans}.
                    Hence $f(q) = X$ by \cref{urysohn_admissible_above_one}.
                    Thus $g(q) = f(q)$
                        by \cref{funs_is_function,function_apply_elim,union_intro_left,function_apply_intro}.
                    Then $g(p) = X$ by \cref{union_intro_right,singleton_intro,function_apply_intro}.
                    Thus $\closure{g(p)}{X}{T}\subseteq X$
                        by \cref{closed_whole_space,funs_type_apply,powerset_elim,closure_subseteq_closed}.
                    Therefore $\closure{g(p)}{X}{T}\subseteq g(q)$ by \cref{subseteq}.
                \end{byCase}
            \end{byCase}
        \end{subproof}
        Show that for all $p\in\dom{g}$ if $1 < p$ then $g(p) = X$.
        \begin{subproof}
            Fix $p$ such that $p\in\dom{g}$.
            Assume $1 < p$.
            Then $p\in\dom{f}$ or $p\in\{r\}$ by \cref{union_iff}.
            \begin{byCase}
            \caseOf{$p\in\{r\}$.}
                Hence $g(p) = X$ by \cref{singleton_elim,union_intro_right,singleton_intro,function_apply_intro}.
            \caseOf{$p\in\dom{f}$.}
                Thus $g(p) = X$
                    by \cref{funs_is_function,function_apply_elim,union_intro_left,function_apply_intro,urysohn_admissible_above_one}.
            \end{byCase}
        \end{subproof}
    \end{subproof}
    Therefore $g\in\urysohnAdmissible{s}{X}{T}{A}{B}{n + 1}$ by \cref{urysohn_admissible}.
    Thus $z\in G$ and $w\in G$ and $\fst{w} = \fst{z} + 1$ and $\snd{z}\subseteq \snd{w}$
        by \cref{urysohn_good_pair_intro,fst_eq,snd_eq,union_upper_left}.
    Hence $(z,w)\in R$.
    Therefore $z\mathrel{R} w$.
\end{proof}
% from §33 helper lemma 28: extension for $r < 1$ (setup)
\begin{lemma}\label{urysohn_unit_interval_extend_lt1_setup}
    Assume $X$ is normal with respect to $T$.
    Assume $A$ is closed in $X$ with respect to $T$.
    Assume $B$ is closed in $X$ with respect to $T$.
    Assume $A$ is disjoint from $B$.
    Assume $s$ is a sequence in $\unitRationals$.
    Let $G = \urysohnGoodPairs{s}{X}{T}{A}{B}$.
    Let $R = \{(z,w) \mid z\in G, w\in G \mid
        \fst{w} = \fst{z} + 1 \land \snd{z}\subseteq\snd{w}\}$.
    Suppose $n\in\naturalsPlus$ and $f\in\urysohnAdmissible{s}{X}{T}{A}{B}{n}$ and $z = (n,f)$.
    Let $r = s(n + 1)$.
    Assume $r\notin\dom{f}$.
    Assume $r < 1$.
    Let $S = \{p\in\dom{f} \mid p < r\}$.
    Let $T_1 = \{q\in\dom{f} \mid r < q\}$.
    Then there exist $p_0,q_0,U_r,h,g$ such that
        $p_0$ is a largest element of $S$ with respect to $\realOrderRel$ and
        $q_0$ is a smallest element of $T_1$ with respect to $\realOrderRel$ and
        $U_r$ is open in $X$ with respect to $T$ and
        $\closure{f(p_0)}{X}{T}\subseteq U_r$ and
        $\closure{U_r}{X}{T}\subseteq f(q_0)$ and
        $h = \{(r,U_r)\}$ and
        $g = f\union h$ and
        $g$ is a function.
\end{lemma}
\begin{proof}
    Then $f\in\funs{\{0,1\}\union \img{s}{\initialSegment{n}}}{\pow{X}}$ and
        for all $p\in\dom{f}$ we have $f(p)\in T$ and
        $f(1) = X\setminus B$ and
        $A\subseteq f(0)$ and
        $\closure{f(0)}{X}{T}\subseteq f(1)$ and
        for all $p\in\dom{f}$ for all $q\in\dom{f}$ if $p < q$ then
            $\closure{f(p)}{X}{T}\subseteq f(q)$ and
        for all $p\in\dom{f}$ if $1 < p$ then $f(p) = X$
        by \cref{urysohn_admissible}.
    Therefore $0\in S$ and $1\in T_1$
        by \cref{urysohn_admissible,funs_elim,upair_intro_left,upair_intro_right,union_intro_left,sequence_in,naturalsplus_add_closed,real_one_in_naturals,funs_type_apply,unit_rational_pos}.
    Let $D = \{0,1\}\union \img{s}{\initialSegment{n}}$.
    $\dom{f} = D$ by \cref{funs_elim}.
    $\{0,1\}$ is finite and $\initialSegment{n}$ is finite by \cref{pair_is_finite,initial_segment_finite}.
    $s\in\funs{\naturalsPlus}{\unitRationals}$ and $s$ is a function by \cref{sequence_in,funs_is_function}.
    Let $g_0 = \restrl{s}{\initialSegment{n}}$.
    $g_0$ is a function and $\dom{s} = \naturalsPlus$ by \cref{restrl_of_function_is_function,funs_elim}.
    $\dom{g_0} = \dom{s}\inter \initialSegment{n}$ by \cref{dom_of_restrl_eq_inter}.
    Hence $\dom{g_0}\subseteq\initialSegment{n}$ by \cref{restrl_dom}.
    Thus $\initialSegment{n}\subseteq\dom{g_0}$
        by \cref{initial_segment_naturalsplus,funs_elim,inter_intro,subseteq}.
    Therefore $\dom{g_0} = \initialSegment{n}$ by \cref{subseteq_antisymmetric}.
    Hence $g_0$ is a function on $\initialSegment{n}$.
    Thus $\img{g_0}{\initialSegment{n}}$ is finite by \cref{finite_image_function}.
    Therefore $\img{s}{\initialSegment{n}}$ is finite
        by \cref{funs_elim,initial_segment_naturalsplus,subseteq,restrl_img,inter_idempotent}.
    Thus $D$ is finite and $\dom{f}$ is finite by \cref{union_of_finite_is_finite}.
    Thus $S\subseteq\dom{f}$ and $T_1\subseteq\dom{f}$ by \cref{subseteq}.
    Hence $S$ is finite and $T_1$ is finite by \cref{subseteq_finite_is_finite}.
    Thus $\{0,1\}\subseteq\reals$
        by \cref{upair_elim,real_add_ident,real_naturals_subset_reals,real_one_in_naturals,subseteq}.
    Thus $\img{s}{\initialSegment{n}}\subseteq\reals$
        by \cref{img_iff,initial_segment_naturalsplus,sequence_in,funs_is_function,function_apply_intro,funs_type_apply,unit_rational_real,subseteq}.
    Therefore $\dom{f}\subseteq\reals$ by \cref{union_subsets_is_subset}.
    Thus $S\subseteq\reals$ and $T_1\subseteq\reals$ by \cref{subseteq_transitive}.
    $\realOrderRel$ is a simple order relation on $\reals$ by \cref{real_order_relation_simple}.
    Take $p_0$ such that $p_0$ is a largest element of $S$ with respect to $\realOrderRel$
        by \cref{finite_inhabited_largest_element}.
    Take $q_0$ such that $q_0$ is a smallest element of $T_1$ with respect to $\realOrderRel$
        by \cref{finite_inhabited_smallest_element}.
    Then $p_0\in\dom{f}$ and $p_0 < r$ and $q_0\in\dom{f}$ and $r < q_0$
        by \cref{largest_element,smallest_element,subseteq}.
    Hence $p_0 < q_0$
        by \cref{subseteq,sequence_in,naturalsplus_add_closed,real_one_in_naturals,funs_type_apply,unit_rational_real,real_lt_trans}.
    Thus $\closure{f(p_0)}{X}{T}\subseteq f(q_0)$ by \cref{urysohn_admissible}.
    Thus $\closure{f(p_0)}{X}{T}$ is closed in $X$ with respect to $T$ and
        $f(q_0)$ is open in $X$ with respect to $T$
        by \cref{sequence_in,naturalsplus_add_closed,real_one_in_naturals,funs_type_apply,unit_rational_real,subseteq,real_lt_trans,normal_space,urysohn_admissible,topology_on,elem_subseteq,powerset_elim,closure_closed_in,open_in}.
    Take $U_r$ such that
        $U_r$ is open in $X$ with respect to $T$ and
        $\closure{f(p_0)}{X}{T}\subseteq U_r$ and
        $\closure{U_r}{X}{T}\subseteq f(q_0)$
        by \cref{normal_closure_neighborhood}.
    Let $h = \{(r,U_r)\}$.
    Let $g = f\union h$.
    For all $x$ such that $x\in\dom{f}\inter\dom{h}$ we have $f(x) = h(x)$.
    For all $a,b,b'$ such that $(a,b)\in\{(r,U_r)\}$ and $(a,b')\in\{(r,U_r)\}$ we have
        $b = b'$ by \cref{singleton_elim,pair_eq_iff}.
    For all $w\in\{(r,U_r)\}$ there exist $x, y$ such that $w = (x,y)$ by \cref{singleton_elim}.
    Therefore $h$ is a relation and $h$ is right-unique by \cref{relation,rightunique}.
    Hence $g$ is a function by \cref{urysohn_admissible,funs_is_function,union_compatible_functions}.
    Finally there exist $p_0,q_0,U_r,h,g$ such that
        $p_0$ is a largest element of $S$ with respect to $\realOrderRel$ and
        $q_0$ is a smallest element of $T_1$ with respect to $\realOrderRel$ and
        $U_r$ is open in $X$ with respect to $T$ and
        $\closure{f(p_0)}{X}{T}\subseteq U_r$ and
        $\closure{U_r}{X}{T}\subseteq f(q_0)$ and
        $h = \{(r,U_r)\}$ and
        $g = f\union h$ and
        $g$ is a function.
\end{proof}
% from §33 helper lemma 29: extension for $r < 1$
\begin{lemma}\label{urysohn_unit_interval_extend_lt1}
    Assume $X$ is normal with respect to $T$.
    Assume $A$ is closed in $X$ with respect to $T$.
    Assume $B$ is closed in $X$ with respect to $T$.
    Assume $A$ is disjoint from $B$.
    Assume $s$ is a sequence in $\unitRationals$.
    Let $G = \urysohnGoodPairs{s}{X}{T}{A}{B}$.
    Let $R = \{(z,w) \mid z\in G, w\in G \mid
        \fst{w} = \fst{z} + 1 \land \snd{z}\subseteq\snd{w}\}$.
    Suppose $n\in\naturalsPlus$ and $f\in\urysohnAdmissible{s}{X}{T}{A}{B}{n}$ and $z = (n,f)$.
    Let $r = s(n + 1)$.
    Assume $r\notin\dom{f}$.
    Assume $r < 1$.
    Let $S = \{p\in\dom{f} \mid p < r\}$.
    Let $T_1 = \{q\in\dom{f} \mid r < q\}$.
    Then there exists $w$ such that $z\mathrel{R} w$.
\end{lemma}
\begin{proof}
    Then $f\in\funs{\{0,1\}\union \img{s}{\initialSegment{n}}}{\pow{X}}$ and
        for all $p\in\dom{f}$ we have $f(p)\in T$ and
        $f(1) = X\setminus B$ and
        $A\subseteq f(0)$ and
        $\closure{f(0)}{X}{T}\subseteq f(1)$ and
        for all $p\in\dom{f}$ for all $q\in\dom{f}$ if $p < q$ then
            $\closure{f(p)}{X}{T}\subseteq f(q)$ and
        for all $p\in\dom{f}$ if $1 < p$ then $f(p) = X$
        by \cref{urysohn_admissible}.
    Take $p_0,q_0,U_r,h,g$ such that
        $p_0$ is a largest element of $S$ with respect to $\realOrderRel$ and
        $q_0$ is a smallest element of $T_1$ with respect to $\realOrderRel$ and
        $U_r$ is open in $X$ with respect to $T$ and
        $\closure{f(p_0)}{X}{T}\subseteq U_r$ and
        $\closure{U_r}{X}{T}\subseteq f(q_0)$ and
        $h = \{(r,U_r)\}$ and
        $g = f\union h$ and
        $g$ is a function
        by \cref{urysohn_unit_interval_extend_lt1_setup}.
    $\dom{g} = \dom{f}\union\dom{h}$ by \cref{dom_union}.
    For all $x\in\dom{h}$ we have $x\in\{r\}$ by \cref{dom_iff,singleton_elim,pair_eq_iff,singleton_intro}.
    For all $x\in\{r\}$ we have $x\in\dom{h}$
        by \cref{singleton_elim,singleton_intro,pair_eq_iff,dom_intro}.
    Hence $\dom{g} = \dom{f}\union\{r\}$ by \cref{subseteq,subseteq_antisymmetric}.
        Show that for all $y\in\dom{g}$ we have $y\in \{0,1\}\union \img{s}{\initialSegment{n + 1}}$.
        \begin{subproof}
                Fix $y$ such that $y\in\dom{g}$.
                Thus $y\in\dom{f}\union\{r\}$.
                Then $y\in\dom{f}$ or $y\in\{r\}$ by \cref{union_iff}.
                \begin{byCase}
                \caseOf{$y\in\dom{f}$.}
                    Then $y\in \{0,1\}$ or $y\in \img{s}{\initialSegment{n}}$ by \cref{funs_elim,union_iff}.
                    \begin{byCase}
                    \caseOf{$y\in \{0,1\}$.}
                        Thus $y\in \{0,1\}\union \img{s}{\initialSegment{n + 1}}$ by \cref{union_intro_left}.
                    \caseOf{$y\in \img{s}{\initialSegment{n}}$.}
                        Take $k$ such that $k\in\initialSegment{n}$ and $(k,y)\in s$ by \cref{img_iff}.
                        Thus $y\in \{0,1\}\union \img{s}{\initialSegment{n + 1}}$
                            by \cref{initial_segment_subset_succ,subseteq,img_iff,union_intro_right}.
                    \end{byCase}
                \caseOf{$y\in\{r\}$.}
                    Then $y\in \{0,1\}\union \img{s}{\initialSegment{n + 1}}$
                        by \cref{singleton_elim,naturalsplus_add_closed,real_one_in_naturals,sequence_in,sequence_value_in_initial_image,union_intro_right}.
                \end{byCase}
        \end{subproof}
        Hence $\dom{g}\subseteq \{0,1\}\union \img{s}{\initialSegment{n + 1}}$ by \cref{subseteq}.
        Show that for all $y\in \{0,1\}\union \img{s}{\initialSegment{n + 1}}$ we have $y\in\dom{g}$.
        \begin{subproof}
                Fix $y$ such that $y\in \{0,1\}\union \img{s}{\initialSegment{n + 1}}$.
                Then $y\in \{0,1\}$ or $y\in \img{s}{\initialSegment{n + 1}}$ by \cref{union_iff}.
                \begin{byCase}
                \caseOf{$y\in \{0,1\}$.}
                    Then $y\in\dom{g}$ by \cref{union_intro_left,funs_elim}.
                \caseOf{$y\in \img{s}{\initialSegment{n + 1}}$.}
                    Take $k$ such that $k\in\initialSegment{n + 1}$ and $(k,y)\in s$ by \cref{img_iff}.
                    \begin{byCase}
                    \caseOf{$k\in\initialSegment{n}$.}
                        Thus $y\in\dom{g}$ by \cref{img_iff,union_intro_right,funs_elim,union_intro_left}.
                    \caseOf{$k\notin\initialSegment{n}$.}
                        $k\in\naturalsPlus$ and $k\leq n + 1$ by \cref{initial_segment_naturalsplus}.
                        Hence $k = n + 1$.
                        Thus $y = r$.
                        Thus $y\in\dom{g}$ by \cref{singleton_intro,union_intro_right}.
                    \end{byCase}
                \end{byCase}
        \end{subproof}
    Therefore $\{0,1\}\union \img{s}{\initialSegment{n + 1}}\subseteq \dom{g}$ by \cref{subseteq}.
    Hence $\dom{g} = \{0,1\}\union \img{s}{\initialSegment{n + 1}}$ by \cref{subseteq_antisymmetric}.
    Show that for all $p\in\dom{g}$ we have $g(p)$ is open in $X$ with respect to $T$.
    \begin{subproof}
        Fix $p$ such that $p\in\dom{g}$.
        Thus $p\in\dom{f}\union\{r\}$.
        Then $p\in\dom{f}$ or $p\in\{r\}$ by \cref{union_iff}.
        \begin{byCase}
        \caseOf{$p\in\dom{f}$.}
            Hence $g(p)$ is open in $X$ with respect to $T$
                by \cref{urysohn_admissible,open_in,topology_on,normal_space,funs_is_function,function_apply_elim,union_intro_left,function_apply_intro}.
        \caseOf{$p\in\{r\}$.}
            Thus $g(p)$ is open in $X$ with respect to $T$
                by \cref{singleton_elim,singleton_intro,union_intro_right,pair_eq_iff,function_apply_intro}.
        \end{byCase}
    \end{subproof}
    Show that for all $p,q$ such that $p\in\dom{g}$ and $q\in\dom{g}$ and $p < q$ we have
        $\closure{g(p)}{X}{T}\subseteq g(q)$.
    \begin{subproof}
        Fix $p$.
        Fix $q$.
        Assume $p\in\dom{g}$ and $q\in\dom{g}$ and $p < q$.
        Thus $p\in\dom{f}\union\{r\}$ and $q\in\dom{f}\union\{r\}$.
        Then $p\in\dom{f}$ or $p\in\{r\}$ by \cref{union_iff}.
        Then $q\in\dom{f}$ or $q\in\{r\}$ by \cref{union_iff}.
        \begin{byCase}
        \caseOf{$p\in\dom{f}$.}
            \begin{byCase}
            \caseOf{$q\in\dom{f}$.}
                Thus $g(p) = f(p)$ and $g(q) = f(q)$
                    by \cref{funs_is_function,function_apply_elim,union_intro_left,function_apply_intro}.
                Thus $\closure{g(p)}{X}{T}\subseteq g(q)$
                    by \cref{urysohn_admissible,funs_is_function,function_apply_elim,union_intro_left,function_apply_intro}.
            \caseOf{$q\in\{r\}$.}
                Thus $q = r$ by \cref{singleton_elim}.
                Then $p < r$.
                Thus $g(p) = f(p)$
                    by \cref{funs_is_function,function_apply_elim,union_intro_left,function_apply_intro}.
                Show that $\closure{f(p)}{X}{T}\subseteq U_r$.
                \begin{subproof}
                    Then $p\in S$.
                    Thus $p\mathrel{\realOrderRel} p_0$ or $p = p_0$ by \cref{largest_element}.
                    \begin{byCase}
                    \caseOf{$p\mathrel{\realOrderRel} p_0$.}
                        Thus $p < p_0$ by \cref{real_order_relation_elim}.
                        Thus $p_0\in\dom{f}$ by \cref{largest_element,subseteq}.
                        Thus $\closure{f(p)}{X}{T}\subseteq f(p_0)$ by \cref{urysohn_admissible}.
                        Thus $\closure{f(p)}{X}{T}\subseteq \closure{f(p_0)}{X}{T}$
                            by \cref{largest_element,subseteq,urysohn_admissible,normal_space,topology_on,elem_subseteq,powerset_elim,subseteq_closure,subseteq_transitive}.
                        Therefore $\closure{f(p)}{X}{T}\subseteq U_r$ by \cref{subseteq_transitive}.
                    \caseOf{$p = p_0$.}
                        Thus $\closure{f(p)}{X}{T}\subseteq U_r$.
                    \end{byCase}
                \end{subproof}
                Thus $g(q) = U_r$ by \cref{singleton_elim,singleton_intro,union_intro_right,function_apply_intro}.
                Hence $\closure{g(p)}{X}{T}\subseteq g(q)$.
            \end{byCase}
        \caseOf{$p\in\{r\}$.}
            \begin{byCase}
            \caseOf{$q\in\dom{f}$.}
                Thus $p = r$ by \cref{singleton_elim}.
                Thus $r < q$.
                Thus $g(p) = U_r$ by \cref{singleton_elim,singleton_intro,union_intro_right,function_apply_intro}.
                Thus $g(q) = f(q)$
                    by \cref{funs_is_function,function_apply_elim,union_intro_left,function_apply_intro}.
                Show that $\closure{U_r}{X}{T}\subseteq f(q)$.
                \begin{subproof}
                    Then $q\in T_1$.
                    Thus $q_0\mathrel{\realOrderRel} q$ or $q_0 = q$ by \cref{smallest_element}.
                    \begin{byCase}
                    \caseOf{$q_0\mathrel{\realOrderRel} q$.}
                        Thus $q_0 < q$ by \cref{real_order_relation_elim}.
                        Thus $q_0\in\dom{f}$ by \cref{smallest_element,subseteq}.
                        Thus $\closure{f(q_0)}{X}{T}\subseteq f(q)$ by \cref{urysohn_admissible}.
                        Thus $\closure{U_r}{X}{T}\subseteq \closure{f(q_0)}{X}{T}$
                            by \cref{smallest_element,subseteq,urysohn_admissible,normal_space,topology_on,elem_subseteq,powerset_elim,subseteq_closure,subseteq_transitive}.
                        Therefore $\closure{U_r}{X}{T}\subseteq f(q)$ by \cref{subseteq_transitive}.
                    \caseOf{$q_0 = q$.}
                        Thus $\closure{U_r}{X}{T}\subseteq f(q)$.
                    \end{byCase}
                \end{subproof}
                Hence $\closure{g(p)}{X}{T}\subseteq g(q)$.
            \caseOf{$q\in\{r\}$.}
                Suppose not.
                Thus $p = r$ and $q = r$ by \cref{singleton_elim}.
                Thus $r\in\reals$
                    by \cref{sequence_in,naturalsplus_add_closed,real_one_in_naturals,funs_type_apply,unit_rational_real}.
                Thus $p\in\reals$ and $p < p$ by \cref{real_lt_subst_left,real_lt_subst_right}.
                Contradiction by \cref{real_lt_irreflexive}.
            \end{byCase}
        \end{byCase}
    \end{subproof}
    Show that for all $p\in\dom{g}$ if $1 < p$ then $g(p) = X$.
    \begin{subproof}
        Fix $p$ such that $p\in\dom{g}$.
        Assume $1 < p$.
        Thus $p\in\dom{f}\union\{r\}$.
        Then $p\in\dom{f}$ or $p\in\{r\}$ by \cref{union_iff}.
        \begin{byCase}
        \caseOf{$p\in\dom{f}$.}
            Hence $g(p) = X$
                by \cref{urysohn_admissible_above_one,funs_is_function,function_apply_elim,union_intro_left,function_apply_intro}.
        \caseOf{$p\in\{r\}$.}
            Suppose not.
            Thus $p < 1$ by \cref{singleton_elim,real_lt_subst_left}.
            Thus $s\in\funs{\naturalsPlus}{\unitRationals}$ and $n + 1\in\naturalsPlus$
                by \cref{sequence_in,naturalsplus_add_closed,real_one_in_naturals}.
            Thus $p\in\reals$ by \cref{singleton_elim,funs_type_apply,unit_rational_real}.
            $1\in\reals$ by \cref{real_naturals_subset_reals,real_one_in_naturals,subseteq}.
            Thus $p < p$ by \cref{real_lt_trans}.
            Contradiction by \cref{real_lt_irreflexive}.
        \end{byCase}
    \end{subproof}
    For all $p\in\dom{g}$ we have $g(p)\in\pow{X}$ by \cref{open_in,normal_space,topology_on,elem_subseteq}.
    Therefore $g\in\funs{\{0,1\}\union \img{s}{\initialSegment{n + 1}}}{\pow{X}}$
        by \cref{urysohn_unit_interval_extend_lt1_setup,funs_intro}.
    Thus $0\in\dom{f}$ and $1\in\dom{f}$
        by \cref{funs_elim,upair_intro_left,upair_intro_right,union_intro_left}.
    Hence $g(1) = X\setminus B$ and $A\subseteq g(0)$ and
        $\closure{g(0)}{X}{T}\subseteq g(1)$
        by \cref{urysohn_admissible,funs_is_function,function_apply_elim,union_intro_left,function_apply_intro}.
    For all $p\in\dom{g}$ we have $g(p)\in T$ by \cref{open_in}.
    Therefore $g\in\urysohnAdmissible{s}{X}{T}{A}{B}{n + 1}$ by \cref{urysohn_admissible}.
    Let $w = (n + 1,g)$.
    Hence $w\in G$ by \cref{naturalsplus_add_closed,real_one_in_naturals,urysohn_good_pair_intro}.
    Now $z\in G$.
    Thus $\fst{w} = \fst{z} + 1$ by \cref{fst_eq}.
    Thus $\snd{z}\subseteq \snd{w}$ by \cref{snd_eq,union_upper_left}.
    Therefore $(z,w)\in R$.
    Thus $z\mathrel{R} w$.
\end{proof}

% from §33 helper lemma 30: extension relation on good pairs
\begin{lemma}\label{urysohn_unit_interval_extend}
    Assume $X$ is normal with respect to $T$.
    Assume $A$ is closed in $X$ with respect to $T$.
    Assume $B$ is closed in $X$ with respect to $T$.
    Assume $A$ is disjoint from $B$.
    Assume $s$ is a sequence in $\unitRationals$.
    Assume for all $p$ such that $p\in\unitRationals$ there exists $n\in\naturalsPlus$ such that $s(n) = p$.
    Let $G = \urysohnGoodPairs{s}{X}{T}{A}{B}$.
    Let $R = \{(z,w) \mid z\in G, w\in G \mid
        \fst{w} = \fst{z} + 1 \land \snd{z}\subseteq\snd{w}\}$.
    Then for all $z\in G$ there exists $w$ such that $z\mathrel{R} w$.
\end{lemma}
\begin{proof}
    Fix $z$ such that $z\in G$.
    Take $n,f$ such that
        $n\in\naturalsPlus$ and
        $f\in\urysohnAdmissible{s}{X}{T}{A}{B}{n}$ and
        $z = (n,f)$
        by \cref{urysohn_good_pair_decomp}.
    Let $r = s(n + 1)$.
    Hence $r\in\unitRationals$ and $r\in\reals$ and $0 < r$
        by \cref{sequence_in,funs_type_apply,naturalsplus_add_closed,real_one_in_naturals,unit_rational_real,unit_rational_pos}.
    \begin{byCase}
    \caseOf{$r\in\dom{f}$.}
        Therefore there exists $w$ such that $z\mathrel{R} w$ by \cref{urysohn_unit_interval_extend_in_dom}.
    \caseOf{$r\notin\dom{f}$.}
        Then $1\in\dom{f}$ by \cref{urysohn_admissible,funs_elim,upair_intro_right,union_intro_left}.
        Thus $r \neq 1$ by \cref{real_mul_ident}.
        Hence $r < 1$ or $1 < r$
            by \cref{real_naturals_subset_reals,real_one_in_naturals,subseteq,real_lt_trichotomy}.
        \begin{byCase}
        \caseOf{$1 < r$.}
            Therefore there exists $w$ such that $z\mathrel{R} w$ by \cref{urysohn_unit_interval_extend_gt1}.
        \caseOf{$r < 1$.}
            Let $S = \{p\in\dom{f} \mid p < r\}$.
            Let $T_1 = \{q\in\dom{f} \mid r < q\}$.
            Therefore there exists $w$ such that $z\mathrel{R} w$ by \cref{urysohn_unit_interval_extend_lt1}.
        \end{byCase}
    \end{byCase}
\end{proof}

% from §33 helper lemma 31: Urysohn function vanishes on $A$
\begin{lemma}\label{urysohn_unit_interval_value_A}
    Assume $X$ is normal with respect to $T$.
    Assume $A$ is closed in $X$ with respect to $T$.
    Assume $B$ is closed in $X$ with respect to $T$.
    Assume $A$ is disjoint from $B$.
    Assume $s$ is a sequence in $\unitRationals$.
    Assume for all $p$ such that $p\in\unitRationals$ there exists $n\in\naturalsPlus$ such that $s(n) = p$.
    Let $G = \urysohnGoodPairs{s}{X}{T}{A}{B}$.
    Let $R = \{(z,w) \mid z\in G, w\in G \mid
        \fst{w} = \fst{z} + 1 \land \snd{z}\subseteq\snd{w}\}$.
    Assume $c$ is a function on $G$.
    Assume for all $z\in G$ we have $z\mathrel{R} \apply{c}{z}$.
    Assume $u$ is a sequence in $G$.
    Assume $u(1) = a_0$.
    Assume $a_0 = (1,u_1)$.
    Assume for all $n\in\naturalsPlus$ we have $u(n + 1) = \apply{c}{u(n)}$.
    Assume for all $n\in\naturalsPlus$ we have $\fst{u(n)} = n$.
    Let $F = \{\snd{u(n)} \mid n\in\naturalsPlus\}$.
    Let $U = \unions{F}$.
    Assume $U$ is a function.
    Let $I_0 = \intervalclosed{0}{1}$.
    Let $f = \{(x,y) \mid x\in X, y\in\reals \mid
        \text{$y$ is an infimum of $\urysohnLevelSet{U}{x}$}\}$.
    Assume $f\in\funs{X}{\reals}$.
    Assume $\img{f}{X}\subseteq I_0$.
    Then for all $x\in A$ we have $f(x) = 0$.
\end{lemma}
\begin{proof}
    Fix $x\in A$.
    Let $Q_x = \urysohnLevelSet{U}{x}$.
    $1\in\naturalsPlus$ and $u$ is a sequence in $G$ and $u\in\funs{\naturalsPlus}{G}$
        by \cref{real_one_in_naturals,sequence_in}.
    Take $n,t$ such that
        $n\in\naturalsPlus$ and
        $t\in\urysohnAdmissible{s}{X}{T}{A}{B}{n}$ and
        $u(1) = (n,t)$
        by \cref{funs_type_apply,urysohn_good_pair_decomp}.
    Hence $t$ is a function and $0\in\dom{t}$
        by \cref{urysohn_admissible,funs_is_function,funs_elim,upair_intro_left,union_intro_left,subseteq}.
    Thus $x\in U(0)$ by \cref{snd_eq,function_apply_elim,unions_intro,function_apply_intro,urysohn_admissible,subseteq}.
    Hence $0\in Q_x$ by \cref{unit_rationals_zero,singleton_intro,union_intro_right,urysohn_level_set}.
    Thus $x\in X$ and $x\in\dom{f}$ and $(x,f(x))\in f$
        by \cref{funs_is_function,funs_elim,closed_in,subseteq,function_apply_elim}.
    Hence $f(x)$ is an infimum of $Q_x$ by \cref{urysohn_infimum_value}.
    Thus $f(x) \leq 0$ by \cref{real_infimum,real_lower_bound}.
    Therefore $0 \leq f(x)$ by \cref{img_iff,subseteq,intervalclosed}.
    Hence $f(x) = 0$ by \cref{real_add_ident,funs_type_apply,real_leq_antisym}.
\end{proof}

% from §33 helper lemma 32: Urysohn function equals $1$ on $B$
\begin{lemma}\label{urysohn_unit_interval_value_B}
    Assume $X$ is normal with respect to $T$.
    Assume $A$ is closed in $X$ with respect to $T$.
    Assume $B$ is closed in $X$ with respect to $T$.
    Assume $A$ is disjoint from $B$.
    Assume $s$ is a sequence in $\unitRationals$.
    Assume for all $p$ such that $p\in\unitRationals$ there exists $n\in\naturalsPlus$ such that $s(n) = p$.
    Let $G = \urysohnGoodPairs{s}{X}{T}{A}{B}$.
    Let $R = \{(z,w) \mid z\in G, w\in G \mid
        \fst{w} = \fst{z} + 1 \land \snd{z}\subseteq\snd{w}\}$.
    Assume $c$ is a function on $G$.
    Assume for all $z\in G$ we have $z\mathrel{R} \apply{c}{z}$.
    Assume $u$ is a sequence in $G$.
    Assume $u(1) = a_0$.
    Assume $a_0 = (1,u_1)$.
    Assume for all $n\in\naturalsPlus$ we have $u(n + 1) = \apply{c}{u(n)}$.
    Assume for all $n\in\naturalsPlus$ we have $\fst{u(n)} = n$.
    Let $F = \{\snd{u(n)} \mid n\in\naturalsPlus\}$.
    Let $U = \unions{F}$.
    Assume $U$ is a function.
    Let $I_0 = \intervalclosed{0}{1}$.
    Let $f = \{(x,y) \mid x\in X, y\in\reals \mid
        \text{$y$ is an infimum of $\urysohnLevelSet{U}{x}$}\}$.
    Assume $f\in\funs{X}{\reals}$.
    Assume $\img{f}{X}\subseteq I_0$.
    Then for all $x\in B$ we have $f(x) = 1$.
\end{lemma}
\begin{proof}
    Fix $x\in B$.
    Let $Q_x = \urysohnLevelSet{U}{x}$.
    $1\in\naturalsPlus$ and $u$ is a sequence in $G$ and $u\in\funs{\naturalsPlus}{G}$
        by \cref{real_one_in_naturals,sequence_in}.
    Take $n,t$ such that
        $n\in\naturalsPlus$ and
        $t\in\urysohnAdmissible{s}{X}{T}{A}{B}{n}$ and
        $u(1) = (n,t)$
        by \cref{funs_type_apply,urysohn_good_pair_decomp}.
    Hence $t$ is a function and $1\in\dom{t}$
        by \cref{urysohn_admissible,funs_is_function,funs_elim,upair_intro_right,union_intro_left,subseteq}.
    Thus $x\notin U(1)$
        by \cref{snd_eq,function_apply_elim,unions_intro,function_apply_intro,urysohn_admissible,subseteq_refl,subseteq,setminus_elim_right}.
    Hence $x\in X$ and $x\in\dom{f}$ and $(x,f(x))\in f$
        by \cref{funs_is_function,funs_elim,closed_in,subseteq,function_apply_elim}.
    Thus $f(x)$ is an infimum of $Q_x$ by \cref{urysohn_infimum_value}.
    Therefore $f(x)$ is a lower bound of $Q_x$ by \cref{real_infimum}.
    Show that for all $q\in Q_x$ we have $1 \leq q$.
    \begin{subproof}
            Fix $q$ such that $q\in Q_x$.
            Then $q\in\unitRationalsZero$ and $x\in U(q)$ by \cref{urysohn_level_set}.
            Suppose not.
            Hence $q < 1$.
            Take $y$ such that $y\in\img{U}{\{q\}}$ and $x\in y$ by \cref{apply,unions_iff}.
            Take $a$ such that $a\in\{q\}$ and $(a,y)\in U$ by \cref{img_iff}.
            Therefore $q\in\dom{U}$ by \cref{singleton_elim,dom_intro}.
            Thus $(q,U(q))\in U$ by \cref{function_apply_elim}.
            Thus there exists $t_q$ such that $t_q\in F$ and $(q,U(q))\in t_q$ by \cref{unions_iff}.
            Thus $q\in\dom{t_q}$ by \cref{dom_intro}.
            Hence there exists $n_q$ such that $n_q\in\naturalsPlus$ and $t_q = \snd{u(n_q)}$.
            Then $u(n_q)\in G$ by \cref{funs_type_apply}.
            Take $m_q,v_q$ such that
                $m_q\in\naturalsPlus$ and
                $v_q\in\urysohnAdmissible{s}{X}{T}{A}{B}{m_q}$ and
                $u(n_q) = (m_q,v_q)$
                by \cref{urysohn_good_pair_decomp}.
            Hence $t_q\in\urysohnAdmissible{s}{X}{T}{A}{B}{m_q}$ by \cref{snd_eq}.
            Then $\snd{u(1)} = u_1$ by \cref{snd_eq}.
            Hence $u_1\in F$.
            $1\in\naturalsPlus$ and $u(1)\in G$ by \cref{real_one_in_naturals,funs_type_apply}.
            Take $n_1,t_1$ such that
                $n_1\in\naturalsPlus$ and
                $t_1\in\urysohnAdmissible{s}{X}{T}{A}{B}{n_1}$ and
                $u(1) = (n_1,t_1)$
                by \cref{urysohn_good_pair_decomp}.
            Therefore $u_1\in\urysohnAdmissible{s}{X}{T}{A}{B}{n_1}$ by \cref{snd_eq}.
            Hence $t_q\subseteq u_1$ or $u_1\subseteq t_q$ by \cref{urysohn_unit_interval_chain}.
            \begin{byCase}
            \caseOf{$t_q\subseteq u_1$.}
                Thus $(q,U(q))\in u_1$ by \cref{subseteq}.
                Thus $q\in\dom{u_1}$ by \cref{dom_intro}.
                Thus $u_1$ is a function and $1\in\dom{u_1}$
                    by \cref{urysohn_admissible,funs_is_function,funs_elim,upair_intro_right,union_intro_left}.
                Hence $U(q) = u_1(q)$ by \cref{function_apply_intro}.
                Thus $\closure{u_1(q)}{X}{T}\subseteq u_1(1)$ and
                    $u_1(q)\subseteq \closure{u_1(q)}{X}{T}$
                    by \cref{urysohn_admissible,normal_space,topology_on,elem_subseteq,powerset_elim,subseteq_closure}.
                Therefore $U(q)\subseteq u_1(1)$ by \cref{subseteq_transitive}.
                Thus $(1,u_1(1))\in u_1$ by \cref{function_apply_elim}.
                Thus $(1,u_1(1))\in U$ by \cref{unions_intro}.
                Thus $U(1) = u_1(1)$ by \cref{function_apply_intro}.
                Hence $x\in U(1)$ by \cref{subseteq}.
                Contradiction.
            \caseOf{$u_1\subseteq t_q$.}
                Thus $u_1$ is a function and $1\in\dom{u_1}$
                    by \cref{urysohn_admissible,funs_is_function,funs_elim,upair_intro_right,union_intro_left}.
                Then $(1,u_1(1))\in u_1$ by \cref{function_apply_elim}.
                Thus $(1,u_1(1))\in t_q$ by \cref{subseteq}.
                Thus $1\in\dom{t_q}$ by \cref{dom_intro}.
                Thus $\closure{t_q(q)}{X}{T}\subseteq t_q(1)$ and
                    $t_q(q)\subseteq \closure{t_q(q)}{X}{T}$
                    by \cref{urysohn_admissible,normal_space,topology_on,elem_subseteq,powerset_elim,subseteq_closure}.
                Thus $t_q$ is a function by \cref{urysohn_admissible,funs_is_function}.
                Hence $U(q) = t_q(q)$ by \cref{function_apply_intro}.
                Thus $U(q)\subseteq t_q(1)$ by \cref{subseteq_transitive}.
                Thus $(1,t_q(1))\in t_q$ by \cref{function_apply_elim}.
                Thus $(1,t_q(1))\in U$ by \cref{unions_intro}.
                Thus $U(1) = t_q(1)$ by \cref{function_apply_intro}.
                Therefore $x\in U(1)$ by \cref{subseteq}.
                Contradiction.
            \end{byCase}
    \end{subproof}
    Therefore $1$ is a lower bound of $Q_x$
        by \cref{real_mul_ident,urysohn_level_set,unit_rational_zero_real,subseteq,real_lower_bound}.
    Hence $1 \leq f(x)$ by \cref{real_infimum}.
    Therefore $f(x) \leq 1$ by \cref{img_iff,subseteq,intervalclosed}.
    Hence $f(x) = 1$ by \cref{real_mul_ident,funs_type_apply,real_leq_antisym}.
\end{proof}

% from §33 helper lemma 33: Urysohn lemma for the unit interval
\begin{lemma}\label{urysohn_unit_interval}
    Assume $X$ is normal with respect to $T$.
    Assume $A$ is closed in $X$ with respect to $T$.
    Assume $B$ is closed in $X$ with respect to $T$.
    Assume $A$ is disjoint from $B$.
    Let $I_0 = \intervalclosed{0}{1}$.
    Let $T_0 = \subspaceTopology{\standardTopologyR}{I_0}$.
    Then there exists $f$ such that
        $f$ is continuous from $X$ to $I_0$ with respect to $T$ and $T_0$ and
        $\img{f}{A}\subseteq \{0\}$ and
        $\img{f}{B}\subseteq \{1\}$.
\end{lemma}
\begin{proof}
    Take $s$ such that
        $s$ is a sequence in $\unitRationals$ and
        for all $p$ such that $p\in\unitRationals$ there exists $n\in\naturalsPlus$ such that $s(n) = p$
        by \cref{unit_rationals_sequence}.
    Let $G = \urysohnGoodPairs{s}{X}{T}{A}{B}$.

    Let $R = \{(z,w) \mid z\in G, w\in G \mid
        \fst{w} = \fst{z} + 1 \land \snd{z}\subseteq\snd{w}\}$.

    Hence for all $z\in G$ there exists $w$ such that $z\mathrel{R} w$
        by \cref{urysohn_unit_interval_extend}.

    Therefore $R$ is a relation by \cref{relation}.

    Take $c$ such that $c$ is a function on $G$ and for all $z\in G$ we have $z\mathrel{R} \apply{c}{z}$
        by \cref{choice_function}.

    Take $u_1$ such that $(1,u_1)\in G$ by \cref{urysohn_unit_interval_u1}.
    Let $a_0 = (1,u_1)$.

    For all $x\in\dom{c}$ we have $c(x)\in G$.
    Hence $c$ is a function and $c$ is a function to $G$.
    Then $G = \dom{c}$.
    Thus $c\in\funs{G}{G}$ by \cref{funs_intro}.
    Take $u$ such that
        $u$ is a sequence in $G$ and
        $u(1) = a_0$ and
        for all $n\in\naturalsPlus$ we have $u(n + 1) = \apply{c}{u(n)}$
        by \cref{iteration_sequence_naturalsplus}.
    Hence for all $n\in\naturalsPlus$ we have $\fst{u(n)} = n$
        by \cref{urysohn_unit_interval_fst_u}.

    Let $F = \{\snd{u(n)} \mid n\in\naturalsPlus\}$.
    Let $U = \unions{F}$.
    Therefore $U$ is a function by \cref{urysohn_unit_interval_union_function}.

    Let $f = \{(x,y) \mid x\in X, y\in\reals \mid
        \text{$y$ is an infimum of $\urysohnLevelSet{U}{x}$}\}$.
    Hence $f$ is a function from $X$ to $\reals$ by \cref{urysohn_unit_interval_f_funs}.

    Show that for all $x\in X$ we have $f(x)\in I_0$.
    \begin{subproof}
        Fix $x\in X$.
        Let $Q_x = \urysohnLevelSet{U}{x}$.
        Then $f(x)$ is an infimum of $Q_x$ by \cref{urysohn_infimum_value}.
        $0\in\reals$ by \cref{real_add_ident}.
        Show that for all $p\in Q_x$ we have $p\in\reals$ and $0 \leq p$.
        \begin{subproof}
            Fix $p$ such that $p\in Q_x$.
            Then $p\in\unitRationalsZero$ by \cref{urysohn_level_set}.
            Thus $p\in\unitRationals$ or $p\in\{0\}$ by \cref{unit_rationals_zero,union_iff}.
            \begin{byCase}
            \caseOf{$p\in\unitRationals$.}
                Take $q$ such that $q\in\naturalsPlus\times\naturalsPlus$ and
                    $p = \fst{q} \rmul \inv{\snd{q}}$
                    by \cref{unit_rationals}.
                Thus $\fst{q}\in\naturalsPlus$ and $\snd{q}\in\naturalsPlus$ by \cref{fst_type,snd_type}.
                Thus $\fst{q}\in\reals$ and $\snd{q}\in\reals$
                    by \cref{real_naturals_subset_reals,subseteq}.
                $0\in\reals$ and $1\in\reals$ and $0 < 1$
                    by \cref{real_add_ident,real_mul_ident,real_zero_lt_one}.
                Thus $0 < \fst{q}$
                    by \cref{real_one_leq_naturalsplus,real_lt_trans,real_lt_subst_right}.
                Thus $0 < \snd{q}$
                    by \cref{real_one_leq_naturalsplus,real_lt_trans,real_lt_subst_right}.
                Therefore $\snd{q}\neq 0$ by \cref{real_pos_nonzero}.
                Thus $\inv{\snd{q}}\in\reals$ by \cref{real_mul_inverse}.
                Thus $\inv{\snd{q}}$ is positive by \cref{real_inv_pos}.
                Thus $0 < \inv{\snd{q}}$.
                Hence $0 < p$ by \cref{real_mul_pos_pos_gt}.
                Thus $p\in\reals$ by \cref{real_mul_closed}.
                Thus $0 \leq p$.
            \caseOf{$p\in\{0\}$.}
                Thus $p = 0$ by \cref{singleton_elim}.
                Thus $p\in\reals$ by \cref{real_add_ident}.
                Thus $0 \leq p$.
            \end{byCase}
        \end{subproof}
        Thus $Q_x\subseteq\reals$ by \cref{subseteq}.
        Hence $0$ is a lower bound of $Q_x$ by \cref{real_lower_bound}.
        Therefore $0 < f(x)$ or $0 = f(x)$ by \cref{real_infimum}.
        Hence $0 \leq f(x)$.
        Show that $f(x) \leq 1$.
        \begin{subproof}
            Suppose not.
            Then $1 < f(x)$.
            $1\in\reals$ and $0 < 1$ by \cref{real_mul_ident,real_zero_lt_one}.
            Thus $0 \leq 1$ by \cref{real_lt_imp_leq}.
            Thus $f(x)\in\reals$ by \cref{funs_type_apply}.
            Take $p$ such that $p\in\unitRationals$ and $p\in\intervalopen{1}{f(x)}$
                by \cref{unit_rationals_dense_nonneg}.
            Thus $p\in\reals$ by \cref{intervalopen}.
            Therefore $p > 1$.
            Thus $p < f(x)$ by \cref{intervalopen}.
            Thus $p\in\unitRationalsZero$ by \cref{unit_rationals_zero,union_iff}.
            Take $k$ such that $k\in\naturalsPlus$ and $s(k) = p$ by \cref{unit_rationals_sequence}.
            Thus $u(k)\in G$ by \cref{sequence_in,funs_type_apply}.
            Take $n,t$ such that
                $n\in\naturalsPlus$ and
                $t\in\urysohnAdmissible{s}{X}{T}{A}{B}{n}$ and
                $u(k) = (n,t)$
                by \cref{urysohn_good_pair_decomp}.
            Hence $\fst{u(k)} = n$ by \cref{fst_eq}.
            Then $n = k$.
            Thus $\snd{u(k)} = t$ by \cref{snd_eq}.
            Thus $t\in\funs{\{0,1\}\union \img{s}{\initialSegment{k}}}{\pow{X}}$
                by \cref{urysohn_admissible}.
            Therefore $p\in\dom{t}$ by \cref{funs_elim,sequence_value_in_initial_image,union_intro_right}.
            Hence $t(p) = X$ by \cref{urysohn_admissible_above_one}.
            Thus $t$ is a function and $U(p) = t(p)$
                by \cref{snd_eq,funs_is_function,function_apply_elim,unions_intro,function_apply_intro}.
            Therefore $U(p) = X$.
            Hence $x\in U(p)$.
            Therefore $p\in Q_x$ by \cref{urysohn_level_set}.
            Hence $f(x)$ is a lower bound of $Q_x$ by \cref{real_infimum}.
            Thus $f(x) < p$ or $f(x) = p$ by \cref{real_lower_bound}.
            \begin{byCase}
            \caseOf{$f(x) < p$.}
                Thus $f(x) < f(x)$ by \cref{real_lt_trans}.
                Contradiction by \cref{real_lt_irreflexive}.
            \caseOf{$f(x) = p$.}
                Thus $f(x) < f(x)$ by \cref{real_lt_subst_left}.
                Contradiction by \cref{real_lt_irreflexive}.
            \end{byCase}
        \end{subproof}
        Therefore $f(x)\in I_0$ by \cref{funs_type_apply,intervalclosed}.
    \end{subproof}
    For all $x\in\dom{f}$ we have $f(x)\in I_0$ by \cref{funs_elim}.
    Thus $f$ is right-unique and $f$ is a relation and $\dom{f} = X$
        by \cref{function_rightunique,funs_is_relation,funs_elim}.
    Therefore $f\in\funs{X}{I_0}$ by \cref{funs_intro}.
    Thus $\img{f}{X}\subseteq I_0$ by \cref{funs_ran,img_subseteq_ran,subseteq_transitive}.

    Now for all $x\in A$ we have $f(x) = 0$
        by \cref{urysohn_unit_interval_value_A}.

    Therefore for all $x\in B$ we have $f(x) = 1$
        by \cref{urysohn_unit_interval_value_B}.

    Hence $\img{f}{A}\subseteq \{0\}$
        by \cref{img_iff,funs_is_function,function_apply_iff,singleton_intro,subseteq}.
    Therefore $\img{f}{B}\subseteq \{1\}$
        by \cref{img_iff,funs_is_function,function_apply_iff,singleton_intro,subseteq}.

    $I_0\subseteq\reals$ by \cref{intervalclosed,subseteq}.
    Hence $f$ is continuous from $X$ to $\reals$ with respect to $T$ and $\standardTopologyR$
        by \cref{urysohn_unit_interval_continuous_reals}.
    $T$ is a topology on $X$ by \cref{normal_space}.
    Then $T_0 = \subspaceTopology{\standardTopologyR}{I_0}$.
    Therefore $f$ is continuous from $X$ to $I_0$ with respect to $T$ and $T_0$
        by \cref{urysohn_unit_interval_f_funs,urysohn_unit_interval_continuous_subspace}.
    Hence there exists $f$ such that
        $f$ is continuous from $X$ to $I_0$ with respect to $T$ and $T_0$ and
        $\img{f}{A}\subseteq \{0\}$ and
        $\img{f}{B}\subseteq \{1\}$.
\end{proof}

% from §33 Item 1: Urysohn lemma
\begin{theorem}\label{urysohn_lemma}
    Assume $X$ is normal with respect to $T$.
    Assume $A$ is closed in $X$ with respect to $T$.
    Assume $B$ is closed in $X$ with respect to $T$.
    Assume $A$ is disjoint from $B$.
    Assume $a\in\reals$ and $b\in\reals$ and $a < b$.
    Let $I = \intervalclosed{a}{b}$.
    Let $T_I = \subspaceTopology{\standardTopologyR}{I}$.
    Then there exists $f$ such that
        $f$ is continuous from $X$ to $I$ with respect to $T$ and $T_I$ and
        for all $x\in A$ we have $f(x) = a$ and
        for all $x\in B$ we have $f(x) = b$.
\end{theorem}
\begin{proof}
    Let $I_0 = \intervalclosed{0}{1}$.
    Let $T_0 = \subspaceTopology{\standardTopologyR}{I_0}$.
    Take $g$ such that
        $g$ is continuous from $X$ to $I_0$ with respect to $T$ and $T_0$ and
        $\img{g}{A}\subseteq \{0\}$ and
        $\img{g}{B}\subseteq \{1\}$
        by \cref{urysohn_unit_interval}.
    Let $c = b - a$.
    Let $j = \identity{I_0}$.
    Let $g_1 = j\circ g$.
    $\standardBasisR$ is a basis on $\reals$ by \cref{standard_basis_is_basis}.
    Hence $\basisTopology{\standardBasisR}{\reals}$ is a topology on $\reals$
        by \cref{basis_topology_is_topology}.
    Therefore $\standardTopologyR$ is a topology on $\reals$ by \cref{standard_topology_reals}.
    Hence $I_0\subseteq\reals$ by \cref{intervalclosed,subseteq}.
    Therefore $I_0$ is a subspace of $\reals$ with respect to $\standardTopologyR$ by \cref{subspace_of}.
    $g_1$ is continuous from $X$ to $\reals$ with respect to $T$ and $\standardTopologyR$
        by \cref{continuous_expand_range}.
    $g\in\funs{X}{I_0}$ and $g$ is a function by \cref{continuous,funs_is_function}.
    Hence $g$ is right-unique and $g$ is a relation and $\dom{g} = X$
        by \cref{function_rightunique,funs_is_relation,funs_elim}.
    Thus $g\in\rels{X}{I_0}$ by \cref{funs_subseteq_rels,subseteq}.
    Thus $\ran{g}\subseteq I_0$ by \cref{rels_ran_subseteq}.
    Therefore $\dom{j} = I_0$ by \cref{id_dom}.
    Hence $\ran{g}\subseteq\dom{j}$.
    $j$ is right-unique and $j$ is a relation by \cref{id_rightunique,id_is_relation}.
    For all $x\in X$ we have $g_1(x) = g(x)$ by \cref{circ_apply,funs_elim,id_apply}.
    Let $k = \{(x,a) \mid x\in X\}$.
    Let $h = \{(x,c) \mid x\in X\}$.
    $T$ is a topology on $X$ by \cref{normal_space}.
    Therefore $c\in\reals$ by \cref{real_add_inverse,real_add_closed,real_sub}.
    Hence $a + c = b$ by
        \cref{real_sub,real_add_assoc,real_add_comm,real_add_inverse,real_add_ident}.
    Therefore $I\subseteq\reals$ by \cref{intervalclosed,subseteq}.
    Hence $I$ is a subspace of $\reals$ with respect to $\standardTopologyR$ by \cref{subspace_of}.
    $T_I$ is a topology on $I$ by \cref{subspace_topology_is_topology}.
    $T_0$ is a topology on $I_0$ by \cref{intervalclosed,subseteq,subspace_topology_is_topology}.
    $k$ is a function from $X$ to $\reals$ and $h$ is a function from $X$ to $\reals$
        by \cref{constant_map_function}.
    Hence for all $x\in X$ we have $k(x) = a$ and $h(x) = c$
        by \cref{constant_map_function,funs_is_function,function_apply_intro}.
    $k$ is continuous from $X$ to $\reals$ with respect to $T$ and $\standardTopologyR$ and
        $h$ is continuous from $X$ to $\reals$ with respect to $T$ and $\standardTopologyR$
        by \cref{constant_continuous}.
    Let $m = \{(x, h(x) \rmul g_1(x)) \mid x\in X\}$.
    Let $f_0 = \{(x, k(x) + m(x)) \mid x\in X\}$.
    $m$ is continuous from $X$ to $\reals$ with respect to $T$ and $\standardTopologyR$
        by \cref{real_function_multiplication_continuous}.
    $f_0$ is continuous from $X$ to $\reals$ with respect to $T$ and $\standardTopologyR$
        by \cref{real_function_addition_continuous}.
    $m\in\funs{X}{\reals}$ and $m$ is a function and
        $f_0\in\funs{X}{\reals}$ and $f_0$ is a function
        by \cref{continuous,funs_is_function}.

    Show that for all $x\in X$ we have $f_0(x)\in I$.
        \begin{subproof}
            Fix $x\in X$.
            Then $g_1(x)\in I_0$ by \cref{circ_apply,funs_elim,id_apply}.
            Thus $g_1(x)\in\reals$ by \cref{intervalclosed}.
            Therefore $0 \leq g_1(x)$ and $g_1(x) \leq 1$ by \cref{intervalclosed}.
            Thus $(x, h(x) \rmul g_1(x))\in m$.
            Hence $m(x) = h(x) \rmul g_1(x)$ by \cref{function_apply_iff}.
            Thus $h(x) = c$.
            Therefore $m(x) = c \rmul g_1(x)$.
            Show that $0 \leq c \rmul g_1(x)$.
            \begin{subproof}
                \begin{byCase}
                \caseOf{$g_1(x) = 0$.}
                    Hence $c \rmul g_1(x) = 0$ by \cref{real_mul_zero,real_mul_comm}.
                    Thus $0 \leq c \rmul g_1(x)$.
                \caseOf{$g_1(x) \neq 0$.}
                    Thus $0 < g_1(x)$ or $g_1(x) < 0$ by \cref{real_lt_trichotomy}.
                    Then $0 < g_1(x)$.
                    Thus $0 < c$ by \cref{real_lt_imp_pos_diff}.
                    Thus $0 < c \rmul g_1(x)$ by \cref{real_mul_pos_pos_gt}.
                    Therefore $0 \leq c \rmul g_1(x)$.
                \end{byCase}
            \end{subproof}
            Hence $0 \leq m(x)$.
            Show that $c \rmul g_1(x) \leq c$.
            \begin{subproof}
                \begin{byCase}
                \caseOf{$g_1(x) = 1$.}
                    Hence $c \rmul g_1(x) = c$ by \cref{real_mul_comm,real_mul_ident}.
                    Thus $c \rmul g_1(x) \leq c$.
                \caseOf{$g_1(x) \neq 1$.}
                    Then $g_1(x) < 1$ by \cref{real_lt_trichotomy}.
                    Thus $0 < c$ by \cref{real_lt_imp_pos_diff}.
                    $1\in\reals$ by \cref{real_naturals_subset_reals,real_one_in_naturals,subseteq}.
                    Thus $g_1(x) \rmul c < 1 \rmul c$ by \cref{real_lt_mul_pos}.
                    Thus $c \rmul g_1(x) = g_1(x) \rmul c$ by \cref{real_mul_comm}.
                    Thus $c \rmul 1 = 1 \rmul c$ by \cref{real_mul_comm}.
                    Therefore $c \rmul g_1(x) < c \rmul 1$.
                    Thus $c \rmul g_1(x) < c$ by \cref{real_mul_ident}.
                    Thus $c \rmul g_1(x) \leq c$.
                \end{byCase}
            \end{subproof}
            Hence $m(x) \leq c$.
            Thus $m(x)\in\reals$ by \cref{real_mul_closed}.
            Thus $(x, k(x) + m(x))\in f_0$.
            Therefore $f_0(x) = a + m(x)$ by \cref{function_apply_iff,real_add_eq_subst}.
            Then $a + m(x)\in\reals$ and $a \leq a + m(x)$
                by \cref{real_add_closed,real_add_ident,real_leq_add,real_add_comm}.
            Thus $a + m(x) \leq b$ by \cref{real_leq_add,real_add_comm}.
            Hence $a + m(x)\in I$ by \cref{intervalclosed}.
            Hence $f_0(x)\in I$ by \cref{function_apply_intro}.
    \end{subproof}
    Therefore $\img{f_0}{X}\subseteq I$
        by \cref{subseteq,img_iff,continuous,funs_is_function,function_apply_iff}.

    $f_0$ is continuous from $X$ to $I$ with respect to $T$ and $T_I$
        by \cref{continuous_restrict_range}.

    Show that for all $x\in A$ we have $f_0(x) = a$.
    \begin{subproof}
        Fix $x\in A$.
        Thus $x\in X$ by \cref{closed_in,subseteq}.
        Thus $g(x)\in\img{g}{A}$ by \cref{function_apply_elim,img_iff}.
        Hence $g(x) = 0$ by \cref{subseteq,singleton_elim}.
        Thus $g(x)\in I_0$ by \cref{funs_elim}.
        Then $g_1(x) = g(x)$.
        Thus $g_1(x)\in\reals$ by \cref{intervalclosed}.
        Therefore $g_1(x) = 0$.
        Thus $(x, h(x) \rmul g_1(x))\in m$.
        Thus $m(x) = h(x) \rmul g_1(x)$ by \cref{function_apply_iff}.
        Hence $m(x) = 0$ by \cref{real_mul_comm,real_mul_zero}.
        Thus $(x, k(x) + m(x))\in f_0$.
        Therefore $f_0(x) = a$ by \cref{function_apply_iff,real_add_eq_subst,real_add_comm,real_add_ident}.
    \end{subproof}
    Show that for all $x\in B$ we have $f_0(x) = b$.
    \begin{subproof}
        Fix $x\in B$.
        Thus $x\in X$ by \cref{closed_in,subseteq}.
        Thus $g(x)\in\img{g}{B}$ by \cref{function_apply_elim,img_iff}.
        Hence $g(x) = 1$ by \cref{subseteq,singleton_elim}.
        Thus $g(x)\in I_0$ by \cref{funs_elim}.
        Then $g_1(x) = g(x)$ by \cref{circ_apply,funs_elim,id_apply}.
        Thus $g_1(x)\in\reals$ by \cref{intervalclosed}.
        Therefore $g_1(x) = 1$.
        Thus $(x, h(x) \rmul g_1(x))\in m$.
        Thus $m(x) = h(x) \rmul g_1(x)$ by \cref{function_apply_iff}.
        Hence $m(x) = c$ by \cref{real_mul_comm,real_mul_ident}.
        Thus $(x, k(x) + m(x))\in f_0$.
        Thus $f_0(x) = a + c$ by \cref{function_apply_iff,real_add_eq_subst}.
        Therefore $f_0(x) = b$.
    \end{subproof}
    Finally there exists $f$ such that
        $f$ is continuous from $X$ to $I$ with respect to $T$ and $T_I$ and
        for all $x\in A$ we have $f(x) = a$ and
        for all $x\in B$ we have $f(x) = b$.
\end{proof}
